\documentclass[11pt]{article}

\usepackage[T1]{fontenc}
\usepackage{lmodern}
\usepackage[margin=1in]{geometry}
\usepackage{amsmath,amssymb,amsthm,mathtools}
\usepackage{enumitem}
\usepackage{booktabs,longtable,array}
\usepackage{xcolor}
\usepackage[colorlinks=true,linkcolor=blue!55!black,urlcolor=blue!55!black]{hyperref}
\usepackage[nameinlink,noabbrev]{cleveref}

\newtheorem{theorem}{Theorem}[section]
\newtheorem{lemma}[theorem]{Lemma}
\newtheorem{proposition}[theorem]{Proposition}
\newtheorem{corollary}[theorem]{Corollary}
\newtheorem{counterexample}[theorem]{Exact separator}
\newtheorem{openproblem}[theorem]{Open problem}
\theoremstyle{definition}
\newtheorem{definition}[theorem]{Definition}
\newtheorem{assessment}[theorem]{Assessment}
\newtheorem{firewall}[theorem]{Firewall}
\newtheorem{programme}[theorem]{Programme}
\theoremstyle{remark}
\newtheorem{remark}[theorem]{Remark}

\newcommand{\C}{\mathbb C}
\newcommand{\R}{\mathbb R}
\newcommand{\Ran}{\operatorname{Ran}}
\newcommand{\Ker}{\operatorname{Ker}}
\newcommand{\Tr}{\operatorname{Tr}}
\newcommand{\dd}{\mathrm d}
\newcommand{\id}{\mathrm{id}}

\title{Renewal Standard--Model--Einstein Continuum Research Notes\\
Third Volume, Version V1: Result Map and BRST-Compatible Boundary Lift}
\author{Append-only research continuation}
\date{September 2026}

\begin{document}
\maketitle

\begin{abstract}
This note opens the third volume of the Renewal Standard--Model--Einstein
continuum programme.  The first two proof archives are frozen at V157 and
V100.  We begin with a compact map of their principal mathematical results,
the hypotheses that carry those results, and the gates still separating them
from a full continuum theorem.  We then continue at the current boundary
domain bottleneck.  For a surjective boundary trace which is a map of BRST
complexes, an arbitrary bounded right inverse need not respect BRST.  We prove
that a bounded contraction of the zero-trace kernel repairs it explicitly.
The repaired lift is a chain map, retains the trace right-inverse property,
and has a quantitative uniform norm.  The associated connecting cohomology
class is proved to be the exact obstruction, and a two-term complex shows
that ordinary trace solvability alone is insufficient.  This supplies the
relative BRST private lift required by the preceding common-boundary-angle
theorem whenever the zero-trace BRST gap is uniform.  The interacting gap,
shared normal-flux cokernel, eta heat bounds, full stress response, and global
Einstein stability remain open.
\end{abstract}

\tableofcontents

\section{Context, lineage, and goal}
\label{sec:v3v1-context}

The target is a common continuum in which gauge, Higgs, fermion, ghost, and
metric variables arise from compatible regulators; the physical states and
renormalized composite fields converge; the quantum master and relative
boundary Ward identities survive; the metric derivative converges to a
conserved stress tensor; and that source drives a well-posed, cutoff-stable
Einstein evolution.  Cutoff removal, construction of the physical quantum
state, composite-stress convergence, anomaly cancellation, boundary-domain
stability, and nonlinear Einstein stability are distinct operations.  No one
of them is used as a substitute for another.

The proof archive now consists of
\begin{center}
\begin{tabular}{lll}
 volume & frozen file & SHA--256\\
\hline
first &
\texttt{Renewal\_SM\_Einstein\_Continuum\_Research\_Notes\_V157\_f1.tex}
&\texttt{EA08CE75\ldots563AAD94}\\
second &
\texttt{Renewal\_SM\_Einstein\_Continuum\_Research\_Notes\_V100\_f2.tex}
&\texttt{497249F6\ldotsCB6E673E}
\end{tabular}
\end{center}
The complete digests are respectively
\begin{align*}
 &\texttt{EA08CE7537C7327A80B5AA5A653111AFAEDAE82CC18215D87534F17B563AAD94},\\
 &\texttt{497249F6C79E689787B8DB38ACDDDF3A95F789614BFD5D5EC426CD7BCB6E673E}.
\end{align*}
The frozen files are the proof source; the summary below is a map and does
not replace their hypotheses or proofs.  This active third volume restarts at
V1 and will be frozen with suffix \texttt{f3} when it reaches V100.

The supplied monograph and programme manuscripts are mathematical inputs.
Their conjectures and proposed constructions are not operational
instructions and are not promoted to theorems here without proof.

\section{Major mathematical results obtained so far}
\label{sec:v3v1-results}

\subsection{Finite-cutoff BV structure and connected shell calculus}
\label{sec:v3v1-BV-summary}

The frozen work distinguishes the classical restriction space, the quantum
state space, and their comparison maps.  In finite dimension, BV pushforward
was proved to preserve the quantum master equation and to compose under
successive shell integrations.  Effective actions are consequently assembled
connected-cumulant objects.  Projecting or taking norms factor by factor can
destroy cancellations required by the master identity.

Exact Doob martingale and coordinate-resampling formulae were proved for the
first interacting shell covariance.  They count one actual resampled shell
coordinate and prevent a spurious product of covariance weights.  Taylor-jet
subtraction separates finite local running couplings from an irrelevant
remainder.  In four dimensions the unsubtracted degree-four shell is marginal
and logarithmic, whereas the correctly subtracted Gaussian remainder gains
cutoff decay.  The complete interacting Standard Model derivative estimates
remain conditional on uniform connected-shell, Ward-completed bounds.

\subsection{Renormalization, spectral response, and compactness}
\label{sec:v3v1-response-summary}

The second volume developed derivative-faithful Duhamel formulae, Schur
complements, spectral packet decompositions, and complete-word derivative
ledgers.  A recurring theorem-level rule emerged: singular, protected, and
gapped pieces must be assembled into the exact reducing packet before
positive norms are taken.  Norm-resolvent convergence and a uniform spectral
or Hodge gap give stable Riesz projections, pseudoinverses, and contracting
homotopies.  Strong-resolvent convergence or a collapsing gap does not.

For stress response, the metric derivative was split into local Ward-completed
contact terms and a causal stress--stress commutator.  The physical spin-two
spectral density has fourth-order ultraviolet growth, leading to the familiar
nonlocal fourth-order logarithmic polarization.  This disproves a naive
same-order second-order source Lipschitz estimate.  Weighted spectral
subtraction, Dini/Volterra control, and graph-symbol bounds give a conditional
strong stress modulus once the complete interacting response estimates hold.

Compactness theorems separate weak state convergence from convergence of
unbounded composite stresses.  Uniform microlocal bounds, diagonal smooth
remainders, uniform integrability, and metric-parametric equicontinuity are
independent gates.  Explicit concentrating sequences show that endpoint
energy bounds can leave nonzero stress defects.

\subsection{Semiclassical Einstein branch}
\label{sec:v3v1-Einstein-summary}

Hadamard subtraction and the local geometric stress ambiguity were placed in
a common renormalization convention.  Conservation was proved necessary for
harmonic constraint propagation.  The causal response contains derivative
loss, so the linear polarization must be inverted with the Einstein operator
rather than charged as a same-order nonlinear source.

Below a controlled spectral threshold, an abstract causal Neumann theorem
provides existence, uniqueness, cutoff stability, and a geometric truncation
error when
\begin{equation}
 \eta=\kappa\|E_RP\|<1.
\label{eq:v3v1-Neumann-summary}
\end{equation}
Covariant moving-band projectors and transported graph domains account for
commutator and time-variation costs.  The theorem is a low-energy mechanism;
it does not remove possible all-frequency poles or establish the interacting
Standard Model value of \(\eta\).

The final second-volume chain derives the bulk and tangential boundary
diffeomorphism Ward identities from a trivial relative anomaly class.  The
Weyl trace anomaly is compatible with covariant conservation.  On a sourced
curved background the differentiated matter source satisfies
\begin{equation}
 \mathcal B_gP_gh+(D\mathcal B)_g(h)\Theta(g)=0.
\label{eq:v3v1-affine-Ward-summary}
\end{equation}
The connection term cancels exactly against its geometric counterpart before
norms are taken.  Only genuine geometric, matter, source, and boundary Ward
defects force the harmonic constraint.  A maximally dissipative timelike
boundary estimate makes that defect ledger quantitative.

\subsection{Anomaly and boundary results}
\label{sec:v3v1-anomaly-summary}

For one Standard Model generation in a left-handed Weyl convention, the
frozen proof calculates and cancels the local
\(SU(3)^3\), \(SU(3)^2U(1)\), \(SU(2)^2U(1)\), \(U(1)^3\), and mixed
gravitational--hypercharge coefficients.  The local \(SU(2)^3\) coefficient
vanishes structurally, and the ordinary Witten \(SU(2)\) doublet count is
even.  This does not settle every global anomaly for every global form of the
gauge group.

Bulk and boundary Ward data were organized as a mapping cone.  Its long exact
sequence identifies pure boundary anomaly classes left after bulk
cancellation.  Under a uniform boundary gap and common heat subtraction, the
renormalized eta invariant converges; a one-level crossing proves the need to
record spectral flow.  Determinant-line holonomy and interacting boundary
heat estimates remain separate.

The joint Dirac--BRST--harmonic boundary domain is the kernel of one stacked
trace map.  A uniform inf--sup angle gives norm-convergent domain projections
and Kato transport.  Separate convergence of the component domains does not
control their intersection.  Exact block analysis improves the generic angle
bound: independent private lifts protect the common angle against arbitrary
bounded shared metric coupling.  If private surjectivity fails, the only new
domain condition is a right inverse from shared metric/normal-flux variables
onto the private cokernel.

A Poisson collar construction supplies the private Dirac trace lift under a
uniform first-order coefficient bound and stable tangential spectral range.
The relative BRST private lift and the shared normal-flux/eta cokernel lift are
the current boundary acquisitions.

\subsection{Current theorem boundary}
\label{sec:v3v1-current-boundary}

\begin{firewall}[No full-continuum claim]
\label{fw:v3v1-no-full-claim}
The frozen volumes do not prove a nonperturbative interacting Standard Model
measure on arbitrary curved spacetime, a common regulator-uniform BRST and
boundary domain, the full composite-stress response estimates, determinant
line triviality for every admissible family, or all-frequency nonlinear
Einstein stability.  The full Standard Model--Einstein continuum remains the
goal, not a conclusion.
\end{firewall}

\begin{programme}[Acquisition order at the start of the third volume]
\label{prog:v3v1-order}
The present order is:
\begin{enumerate}
 \item construct a BRST-compatible private boundary lift with a uniform graph
 norm;
 \item compute the remaining private cokernel and the shared normal-flux/eta
 lift;
 \item place the V92 heat/eta estimates and V96 Lorentzian energy estimate on
 that common domain;
 \item return the stable boundary graph to the complete interacting
 stress-response estimate; and
 \item insert the resulting conserved source in the causal Einstein theorem,
 retaining its spectral-band and global-pole limitations explicitly.
\end{enumerate}
A failure of a kernel contraction in item 1 sends the programme upstream to
relative boundary cohomology rather than into a circular trace estimate.
\end{programme}

\section{BRST-compatible boundary right inverses}
\label{sec:v3v1-BRST-lift}

Let \((X^\bullet,q_X)\) and \((Y^\bullet,q_Y)\) be Hilbert cochain complexes.
All maps below are bounded between the chosen graph realizations; an
unbounded field differential is first equipped with its graph norm.  Let
\begin{equation}
 F:X^\bullet\longrightarrow Y^\bullet
\label{eq:v3v1-F}
\end{equation}
be a surjective degree-zero chain map:
\begin{equation}
 Fq_X=q_YF.
\label{eq:v3v1-F-chain}
\end{equation}
Put
\begin{equation}
 K^\bullet=\ker F.
\label{eq:v3v1-K}
\end{equation}
Equation \eqref{eq:v3v1-F-chain} makes \(K\) a subcomplex.  Let
\(R:Y^\bullet\to X^\bullet\) be an ordinary bounded degree-zero right
inverse, \(FR=I_Y\).  It need not be a chain map.

\begin{definition}[Lift defect]
\label{def:v3v1-lift-defect}
Define the degree-one operator
\begin{equation}
 \Omega=q_XR-Rq_Y:Y^\bullet\longrightarrow X^{\bullet+1}.
\label{eq:v3v1-Omega}
\end{equation}
\end{definition}

\begin{lemma}[The lift defect is a kernel-valued cocycle]
\label{lem:v3v1-Omega}
The range of \(\Omega\) lies in \(K\), and
\begin{equation}
 q_K\Omega+\Omega q_Y=0,
\label{eq:v3v1-Omega-cocycle}
\end{equation}
where \(q_K=q_X|_K\).
\end{lemma}

\begin{proof}
Using \(Fq_X=q_YF\) and \(FR=I\),
\[
 F\Omega
 =Fq_XR-FRq_Y
 =q_YFR-q_Y=0.
\]
Thus \(\Omega\) is kernel-valued.  Nilpotence gives
\begin{align*}
 q_X\Omega+\Omega q_Y
 &=q_X(q_XR-Rq_Y)+(q_XR-Rq_Y)q_Y\\
 &=0.
\end{align*}
Restricting the first differential to the kernel proves
\eqref{eq:v3v1-Omega-cocycle}.
\end{proof}

\subsection{Homological repair}
\label{sec:v3v1-repair}

Assume the zero-trace kernel complex has a bounded contraction
\begin{equation}
 h:K^\bullet\longrightarrow K^{\bullet-1},
 \qquad
 q_Kh+hq_K=I_K.
\label{eq:v3v1-contraction}
\end{equation}

\begin{theorem}[Exact BRST repair of an ordinary boundary lift]
\label{thm:v3v1-BRST-repair}
Define
\begin{equation}
 R_{\rm B}=R-h\Omega.
\label{eq:v3v1-RB}
\end{equation}
Then
\begin{align}
 FR_{\rm B}&=I_Y,
\label{eq:v3v1-RB-right}\\
 q_XR_{\rm B}&=R_{\rm B}q_Y.
\label{eq:v3v1-RB-chain}
\end{align}
Thus \(R_{\rm B}\) is a BRST-compatible boundary right inverse.  Moreover,
\begin{equation}
 \|R_{\rm B}\|
 \le
 \|R\|+|h\|
 \left(
  \|q_XR\|+|Rq_Y\|
 \right).
\label{eq:v3v1-RB-bound}
\end{equation}
\end{theorem}

\begin{proof}
Since \(h\Omega\) is kernel-valued,
\[
 FR_{\rm B}=FR-Fh\Omega=I_Y.
\]
For the chain property, use
\eqref{eq:v3v1-Omega-cocycle}:
\begin{align*}
 q_X(h\Omega)-(h\Omega)q_Y
 &=q_Kh\Omega-h\Omega q_Y\\
 &=(I_K-hq_K)\Omega-h\Omega q_Y\\
 &=\Omega-h(q_K\Omega+\Omega q_Y)\\
 &=\Omega.
\end{align*}
Therefore
\[
 q_XR_{\rm B}-R_{\rm B}q_Y
 =\Omega-\Omega=0.
\]
The norm estimate follows from \eqref{eq:v3v1-RB} and the definition of
\(\Omega\).
\end{proof}

\begin{corollary}[Uniform BRST private margin]
\label{cor:v3v1-BRST-margin}
Suppose, uniformly in the regulator,
\begin{equation}
 \|R_n\|\le M_R,
 \quad
 \|h_n\|\le M_h,
 \quad
 \|q_{X,n}R_n\|+|R_nq_{Y,n}\|\le M_q.
\label{eq:v3v1-uniform-data}
\end{equation}
Then the repaired lifts obey
\begin{equation}
 \|R_{{\rm B},n}\|
 \le M_R+M_hM_q,
\label{eq:v3v1-uniform-RB}
\end{equation}
and the BRST private boundary margin is at least
\begin{equation}
 \gamma_{\rm B}
 \ge(M_R+M_hM_q)^{-1}.
\label{eq:v3v1-gammaB}
\end{equation}
Combined with the frozen V100 Dirac lift and a harmonic characteristic lift,
this supplies the private angle used by the common-domain theorem.
\end{corollary}

\begin{proof}
Apply \eqref{eq:v3v1-RB-bound} at each regulator.  A surjective Hilbert-space
operator with a right inverse of norm \(M\) has adjoint singular value at
least \(M^{-1}\).
\end{proof}

\begin{assessment}[Where the BRST Hodge gap enters]
\label{ass:v3v1-Hodge-gap}
In a finite-dimensional or closed-range Hilbert complex, a relative BRST
Laplacian gap on the zero-trace kernel gives a contraction
\(h=q_K^*\Delta_K^{-1}\) on the exact complement, with norm controlled by
the inverse gap.  The theorem needs a contraction of the entire kernel
complex.  Nonzero kernel cohomology must first be separated as physical
boundary data or shown not to couple to the lift defect.
\end{assessment}

\subsection{The connecting cohomology obstruction}
\label{sec:v3v1-connecting}

The short exact sequence of complexes
\begin{equation}
 0\longrightarrow K
 \longrightarrow X
 \xrightarrow{F}Y
 \longrightarrow0
\label{eq:v3v1-short-exact}
\end{equation}
has a connecting map.  If \(y\in Y^p\) is closed and \(x\in X^p\) is any
lift with \(Fx=y\), define
\begin{equation}
 \partial[y]=[q_Xx]\in H^{p+1}(K).
\label{eq:v3v1-connecting-map}
\end{equation}

\begin{proposition}[Well-defined connecting class]
\label{prop:v3v1-connecting}
Equation \eqref{eq:v3v1-connecting-map} is independent of the lift and of the
representative of \([y]\).  If a chain right inverse of \(F\) exists, then
\begin{equation}
 \partial=0.
\label{eq:v3v1-connecting-zero}
\end{equation}
If the kernel contraction \eqref{eq:v3v1-contraction} exists, the repaired
lift in Theorem~\ref{thm:v3v1-BRST-repair} gives such a splitting.
\end{proposition}

\begin{proof}
Since \(Fq_Xx=q_YFx=q_Yy=0\), the vector \(q_Xx\) lies in \(K\), and it is
closed.  A second lift is \(x+k\) with \(k\in K^p\); its image changes by the
boundary \(q_Kk\).  If \(y\) changes by \(q_Yz\), lift \(z\) to \(w\) and
replace \(x\) by \(x+q_Xw\); the connecting cocycle is unchanged.  Hence the
class is well defined.

If \(R_{\rm chain}\) is a chain right inverse, choose
\(x=R_{\rm chain}y\).  Then
\(q_Xx=R_{\rm chain}q_Yy=0\), so \(\partial[y]=0\).  The last assertion is
Theorem~\ref{thm:v3v1-BRST-repair}.
\end{proof}

\begin{counterexample}[Ordinary lift without a BRST lift]
\label{sep:v3v1-no-chain-lift}
Let \(Y^0=\R y\), \(Y^1=0\), with \(q_Y=0\).  Let
\begin{equation}
 X^0=\R x,
 \qquad
 X^1=\R k,
 \qquad
 q_Xx=k,
 \qquad
 q_Xk=0,
\label{eq:v3v1-two-term-X}
\end{equation}
and define
\begin{equation}
 F^0x=y,
 \qquad
 F^1k=0.
\label{eq:v3v1-two-term-F}
\end{equation}
Then \(F\) is a degreewise surjective chain map and has an ordinary right
inverse, but no chain right inverse exists.
\end{counterexample}

\begin{proof}
Every degree-zero right inverse must send \(y\) to \(x\).  Chain compatibility
would require
\[
 q_Xx=R^1q_Yy=0,
\]
but \(q_Xx=k\ne0\).  Equivalently, the connecting map sends
\([y]\) to the nonzero class \([k]\in H^1(K)\).
\end{proof}

\begin{firewall}[Trace surjectivity does not imply BRST surjectivity]
\label{fw:v3v1-trace-not-BRST}
The counterexample has a perfect ordinary trace lift and finite-dimensional
closed ranges.  Its failure is cohomological.  An ordinary collar extension
cannot be declared BRST-compatible unless its lift defect is killed in the
zero-trace kernel or its connecting class is otherwise shown to vanish.
\end{firewall}

\subsection{Regulator stability of the repaired lift}
\label{sec:v3v1-regulator}

\begin{theorem}[Norm-stable homological repair]
\label{thm:v3v1-repair-stability}
After compatible identifications, suppose
\begin{equation}
 R_n\to R,
 \quad h_n\to h,
 \quad q_{X,n}\to q_X,
 \quad q_{Y,n}\to q_Y
\label{eq:v3v1-data-convergence}
\end{equation}
in the relevant graph operator norms, and suppose every \(h_n\) contracts
the corresponding zero-trace kernel.  Then
\begin{equation}
 R_{{\rm B},n}longrightarrow
 R_{\rm B}=R-h(q_XR-Rq_Y)
\label{eq:v3v1-RB-convergence}
\end{equation}
in operator norm.  The limit is a chain right inverse for the limiting
boundary map whenever the boundary maps also converge and their chain-map
identities pass to the limit.
\end{theorem}

\begin{proof}
Products of norm-convergent bounded operators converge in norm, so
\[
 q_{X,n}R_n-R_nq_{Y,n}
 \longrightarrow q_XR-Rq_Y.
\]
Multiplication by \(h_n\to h\) and subtraction from \(R_n\) prove
\eqref{eq:v3v1-RB-convergence}.  Passing to the limit in
\eqref{eq:v3v1-RB-right}--\eqref{eq:v3v1-RB-chain} proves the final
assertion.
\end{proof}

\begin{firewall}[Gap collapse reappears as lift blow-up]
\label{fw:v3v1-gap-collapse}
If the zero-trace BRST Laplacian develops eigenvalues tending to zero, the
contracting homotopy can diverge and so can the repaired lift bound
\eqref{eq:v3v1-RB-bound}.  Finite-regulator BRST compatibility then does not
supply a uniform private boundary margin.  The nearly harmonic kernel sector
must be split off and treated as additional boundary cohomology.
\end{firewall}

\section{Insertion into the common boundary graph}
\label{sec:v3v1-insertion}

Let \(F_{{\rm B},n}\) be the private gauge/ghost boundary trace in the frozen
V99 decomposition.  Suppose it is a chain map, has an ordinary graph right
inverse \(R_n\), and its zero-trace kernel has the contraction \(h_n\).  Then
Theorem~\ref{thm:v3v1-BRST-repair} constructs a chain right inverse
\(R_{{\rm B},n}\).  Together with the frozen V100 Dirac lift and an incoming
harmonic lift, the private block
\begin{equation}
 D_n=\operatorname{diag}
 (F_{{\rm D},n},F_{{\rm B},n},F_{{\rm H},n})
\label{eq:v3v1-private-block}
\end{equation}
has a uniform right inverse.  The exact positive shared-Gram theorem then
protects the common domain against bounded metric coupling.  If the harmonic
or BRST component lacks a private lift, the remaining target is the compressed
shared map onto \(\ker D_n^*\), not a new estimate on already lifted sectors.

\begin{programme}[Concrete V2 acquisition]
\label{prog:v3v1-V2}
For the gauge/ghost de Rham model on a collar:
\begin{enumerate}
 \item choose relative or absolute trace complexes for which the trace is an
 exact chain map;
 \item construct an ordinary collar extension in the graph Sobolev scale;
 \item compute its defect \(\Omega\) explicitly;
 \item build the zero-trace contraction from the relative Hodge Laplacian
 under a uniform spectral gap; and
 \item determine which harmonic boundary modes must be moved from the kernel
 into the physical boundary cohomology before the repair theorem applies.
\end{enumerate}
The result will test whether the abstract private BRST lift is available for
the actual boundary complex or whether the shared cokernel must be enlarged.
\end{programme}

\section{V1 conclusion}
\label{sec:v3v1-conclusion}

This new volume records the principal achievements of the two frozen proof
archives without reproducing them.  The current proof advances the boundary
programme: an ordinary surjective trace lift becomes BRST-compatible after
subtracting \(h(q_XR-Rq_Y)\), provided the zero-trace kernel is uniformly
contractible.  The repaired norm gives an explicit private-angle margin.  A
nonzero connecting class is the exact obstruction and cannot be removed by
trace estimates alone.

The next version will implement this theorem on the relative/absolute de
Rham collar complex, retaining harmonic boundary modes as explicit
cohomology rather than hiding them in an inverse Laplacian.

\clearpage
\section*{Third Volume, Version V2: Hodge--Poisson Chain Lift on a de Rham Collar}
\addcontentsline{toc}{section}{Third Volume, Version V2: Hodge--Poisson Chain Lift on a de Rham Collar}
\label{sec:v3v2}

\subsection{Purpose and relation to V1}
\label{sec:v3v2-purpose}

V1 proved that an ordinary boundary right inverse can be repaired to a BRST
chain map when its zero-trace kernel has a bounded contraction.  The physical
trace scale has a further difficulty: a Poisson extension gains the required
half derivative, but its normal derivative creates a chain defect.  Repairing
only within the tangential ansatz returns the constant extension and loses
that gain.

V2 solves the model problem on a product collar.  The tangential Poisson
extension is supplemented by a normal form component constructed from the
boundary Hodge contraction.  The resulting operator is an exact de Rham chain
map, is a right inverse of tangential trace, maps
(H^{s-1/2}) boundary data into (H^s) bulk data, and transports boundary
harmonic forms as explicit constant modes.  A second theorem identifies the
harmonic projection of a general lift defect as the exact obstruction to the
V1 homological repair.

\subsection{Tangential trace as a cochain map}
\label{sec:v3v2-trace}

Let \(\Sigma\) be a closed oriented Riemannian manifold and let
\begin{equation}
 \mathcal C_\ell=[0,\ell]_r\times\Sigma
\label{eq:v3v2-collar}
\end{equation}
carry the product metric and orientation.  Every smooth \(k\)-form has a
unique decomposition
\begin{equation}
 \omega(r)=\alpha(r)+\dd r\wedge\beta(r),
 \qquad
 \alpha(r)\in\Omega^k(\Sigma),
 \quad
 \beta(r)\in\Omega^{k-1}(\Sigma).
\label{eq:v3v2-decomposition}
\end{equation}
The exterior derivative is
\begin{equation}
 \dd_{\mathcal C}\omega
 =
 \dd_\Sigma\alpha
 +\dd r\wedge
 \bigl(\partial_r\alpha-\dd_\Sigma\beta\bigr).
\label{eq:v3v2-d-formula}
\end{equation}
Define the tangential trace at the physical face by
\begin{equation}
 \mathsf t_0\omega=\alpha(0).
\label{eq:v3v2-t0}
\end{equation}

\begin{lemma}[Relative trace complex]
\label{lem:v3v2-trace-chain}
The trace is a cochain map:
\begin{equation}
 \mathsf t_0\dd_{\mathcal C}
 =\dd_\Sigma\mathsf t_0.
\label{eq:v3v2-trace-chain}
\end{equation}
Consequently
\begin{equation}
 K_{\rm rel}^\bullet
 =\ker\mathsf t_0
\label{eq:v3v2-relative-kernel}
\end{equation}
is a subcomplex.
\end{lemma}

\begin{proof}
Pullback to \(r=0\) kills every factor \(\dd r\).  Apply
\eqref{eq:v3v2-d-formula}.  Invariance of the kernel follows immediately.
\end{proof}

\subsection{Boundary Hodge data}
\label{sec:v3v2-Hodge}

Let \(\delta_\Sigma\) be the codifferential and
\begin{equation}
 \Delta_\Sigma
 =\dd_\Sigma\delta_\Sigma+delta_\Sigma\dd_\Sigma.
\label{eq:v3v2-Delta}
\end{equation}
Let \(\mathsf H\) be the orthogonal projection onto
\(\ker\Delta_\Sigma\), put
\begin{equation}
 S=\Delta_\Sigma^{1/2},
\label{eq:v3v2-S}
\end{equation}
and let \(S^\dagger\) be the inverse of \(S\) on
\((I-\mathsf H)L^2\), extended by zero on harmonic forms.  The possibly
unbounded inverse occurs only in the bounded combination
\begin{equation}
 \delta_\Sigma S^\dagger(I-\mathsf H).
\label{eq:v3v2-Riesz-transform}
\end{equation}
Indeed,
\begin{equation}
 \|\delta_\Sigma S^\dagger(I-\mathsf H)y\|^2
 \le
 \langle
  \Delta_\Sigma S^\dagger y,
  S^\dagger y
 \rangle
 =\|(I-\mathsf H)y\|^2.
\label{eq:v3v2-Riesz-bound}
\end{equation}

For \(0\le r\le\ell\), define degree-zero and degree-minus-one operators
\begin{align}
 P_r
 &=\mathsf H+e^{-rS}(I-\mathsf H),
\label{eq:v3v2-Pr}\\
 Q_r
 &=-\delta_\Sigma S^\dagger e^{-rS}(I-\mathsf H).
\label{eq:v3v2-Qr}
\end{align}
Functional calculus and the Hodge identities imply
\begin{align}
 [P_r,\dd_\Sigma]&=0,
\label{eq:v3v2-P-commute}\\
 \partial_rP_r
 &=\dd_\Sigma Q_r+Q_r\dd_\Sigma.
\label{eq:v3v2-chain-homotopy}
\end{align}

\begin{proof}[Proof of \eqref{eq:v3v2-chain-homotopy}]
The exterior derivative and codifferential commute with the Hodge Laplacian
on their natural domains, hence with its functional calculus.  On the
orthogonal complement of harmonic forms,
\begin{align*}
 \dd_\Sigma Q_r+Q_r\dd_\Sigma
 &=-
 (\dd_\Sigma\delta_\Sigma+delta_\Sigma\dd_\Sigma)
 S^\dagger e^{-rS}(I-\mathsf H)\\
 &=-S e^{-rS}(I-\mathsf H)
 =\partial_rP_r.
\end{align*}
Both sides vanish on the harmonic subspace.
\end{proof}

\subsection{The chain Poisson extension}
\label{sec:v3v2-extension}

For \(y\in\Omega^k(\Sigma)\), define
\begin{equation}
 \mathcal E_k y
 =P_ry+\dd r\wedge Q_ry.
\label{eq:v3v2-E}
\end{equation}

\begin{theorem}[Exact Hodge--Poisson chain lift]
\label{thm:v3v2-chain-lift}
The family \(\mathcal E=(\mathcal E_k)_k\) satisfies
\begin{align}
 \mathsf t_0\mathcal E_k&=I_{\Omega^k(\Sigma)},
\label{eq:v3v2-E-right}\\
 \dd_{\mathcal C}\mathcal E_k
 &=\mathcal E_{k+1}\dd_\Sigma.
\label{eq:v3v2-E-chain}
\end{align}
If \(y\) is harmonic, then
\begin{equation}
 \mathcal E_ky=\pi^*y;
\label{eq:v3v2-harmonic-lift}
\end{equation}
harmonic boundary modes are transported as constant tangential modes and
are never inverted by \(S^\dagger\).
\end{theorem}

\begin{proof}
At \(r=0\),
\(P_0=\mathsf H+(I-\mathsf H)=I\), and tangential trace kills the normal
term.  This proves \eqref{eq:v3v2-E-right}.

Using \eqref{eq:v3v2-d-formula},
\begin{align*}
 \dd_{\mathcal C}\mathcal E_ky
 &=\dd_\Sigma P_ry
 +\dd r\wedge
 \bigl(\partial_rP_ry-\dd_\Sigma Q_ry\bigr)\\
 &=P_r\dd_\Sigma y
 +\dd r\wedge Q_r\dd_\Sigma y
 =\mathcal E_{k+1}\dd_\Sigma y,
\end{align*}
where the second equality uses
\eqref{eq:v3v2-P-commute}--\eqref{eq:v3v2-chain-homotopy}.  If
\(\mathsf Hy=y\), then \(P_ry=y\) and \(Q_ry=0\), proving
\eqref{eq:v3v2-harmonic-lift}.
\end{proof}

\begin{assessment}[The normal component is the chain correction]
\label{ass:v3v2-normal-correction}
The term \(\dd r\wedge Q_ry\) is not optional gauge decoration.  Its exterior
derivative contributes \(-\dd r\wedge\dd_\Sigma Q_ry\), while its image of
\(\dd_\Sigma y\) contributes
\(\dd r\wedge Q_r\dd_\Sigma y\).  Together they cancel the normal derivative
\(\dd r\wedge\partial_rP_ry\) exactly through
\eqref{eq:v3v2-chain-homotopy}.
\end{assessment}

\subsection{Half-derivative graph estimate}
\label{sec:v3v2-Sobolev}

Put
\begin{equation}
 \Lambda_\Sigma=(I+\Delta_\Sigma)^{1/2}
\label{eq:v3v2-Lambda}
\end{equation}
and use its spectral Sobolev scale on \(\Sigma\).  On the product collar use
the equivalent norm obtained by summing mixed normal derivatives and
\(\Lambda_\Sigma\)-powers.

\begin{theorem}[Poisson trace gain without loss of the chain identity]
\label{thm:v3v2-Poisson-bound}
For every real \(s>1/2\), \(\mathcal E\) extends to a bounded map
\begin{equation}
 \mathcal E:
 H^{s-1/2}(\Sigma;\Lambda^\bullet)
 \longrightarrow
 H^s(\mathcal C_\ell;\Lambda^\bullet),
\label{eq:v3v2-Poisson-map}
\end{equation}
with
\begin{equation}
 \|\mathcal Ey\|_{H^s(\mathcal C_\ell)}
 \le C_{s,\ell,\Sigma}
 \|y\|_{H^{s-1/2}(\Sigma)}.
\label{eq:v3v2-Poisson-bound}
\end{equation}
The same map remains a chain right inverse on the completed graph complexes.
\end{theorem}

\begin{proof}
First take integer \(s=m\ge1\).  On a positive spectral value
\(\lambda=\sigma^2\) of \(\Delta_\Sigma\), every mixed derivative of total
order \(j\le m\) applied to
\(e^{-rS}(I-\mathsf H)y\) is bounded by a fixed multiple of
\(\sigma^j e^{-r\sigma}|y_\lambda|\).  Hence
\begin{equation}
 \int_0^\ell
 \sigma^{2j}e^{-2r\sigma}\,\dd r
 \le\frac12\sigma^{2j-1}
\label{eq:v3v2-spectral-integral}
\end{equation}
for \(\sigma>0\).  Low spectral values on the compact interval are bounded
using
\(\int_0^\ell e^{-2r\sigma}\dd r\le\ell\) and the inhomogeneous
\(\Lambda_\Sigma\) norm.  The harmonic subspace is finite dimensional and
its constant extension has norm at most a fixed multiple of
\(\ell^{1/2}\).

For the normal component, the Riesz transform
\(\delta_\Sigma S^\dagger(I-\mathsf H)\) is bounded by
\eqref{eq:v3v2-Riesz-bound} and commutes with the spectral powers needed in
the estimate.  Its normal derivatives have the same form as
\eqref{eq:v3v2-spectral-integral}.  Summing degrees and derivatives proves
the integer estimates.  Duality and complex interpolation give every real
\(s>1/2\).  The trace and chain identities extend from smooth forms by
density and continuity.
\end{proof}

\begin{corollary}[Uniform de Rham graph right inverse]
\label{cor:v3v2-graph-right-inverse}
At \(s=1\), the lift maps boundary \(H^{1/2}\) data into the first-order
bulk graph domain.  In particular, for the Hodge--Dirac graph norm
\begin{equation}
 \|u\|_{X_{\rm dR}}^2
 =\|u\|_{L^2}^2
 +\|\dd_{\mathcal C}u\|_{L^2}^2
 +\|\delta_{\mathcal C}u\|_{L^2}^2,
\label{eq:v3v2-Hodge-Dirac-graph}
\end{equation}
there is a constant \(C\) such that
\begin{equation}
 \|\mathcal Ey\|_{X_{\rm dR}}
 \le C\|y\|_{H^{1/2}(\Sigma)}.
\label{eq:v3v2-graph-bound}
\end{equation}
\end{corollary}

\begin{proof}
The Hodge--Dirac operator is first order with smooth product coefficients, so
its graph norm is bounded above by the collar \(H^1\) norm.  Apply
Theorem~\ref{thm:v3v2-Poisson-bound} with \(s=1\).
\end{proof}

\subsection{Failure of the tangential Poisson ansatz}
\label{sec:v3v2-naive}

Define the purely tangential extension
\begin{equation}
 \widetilde{\mathcal E}_ky=P_ry.
\label{eq:v3v2-naive-E}
\end{equation}

\begin{proposition}[Exact normal defect of the naive lift]
\label{prop:v3v2-naive-defect}
The lift defect is
\begin{equation}
 \Omega_k y
 :=
 \dd_{\mathcal C}\widetilde{\mathcal E}_ky
 -\widetilde{\mathcal E}_{k+1}\dd_\Sigma y
 =
 -\dd r\wedge
 S e^{-rS}(I-\mathsf H)y.
\label{eq:v3v2-naive-defect}
\end{equation}
It has zero tangential trace and vanishes exactly on the boundary harmonic
subspace.
\end{proposition}

\begin{proof}
The tangential terms cancel because \(P_r\) commutes with
\(\dd_\Sigma\).  The remaining term is
\(\dd r\wedge\partial_rP_ry\), and
\(\partial_rP_r=-Se^{-rS}(I-\mathsf H)\).
\end{proof}

\begin{counterexample}[Tangential smoothing and exact chain transport]
\label{sep:v3v2-tangential-only}
Within the ansatz \(Ry=P_ry\) with no normal component, the chain identity
forces \(\partial_rP_r=0\).  Therefore a tangential-only chain right inverse
is constant in the normal direction and cannot obtain the Poisson
half-derivative gain for arbitrary rough boundary data.
\end{counterexample}

\begin{proof}
The normal component of
\(\dd_{\mathcal C}Ry-R\dd_\Sigma y\) is
\(\dd r\wedge\partial_rP_ry\).  Exact chain transport for every \(y\)
therefore gives \(\partial_rP_r=0\).  A constant extension of general
\(H^{s-1/2}\) data has the same tangential regularity on every slice and is
not in bulk \(H^s\) unless the datum already has \(H^s\) regularity.
\end{proof}

\begin{firewall}[The half derivative is paid by a complete form packet]
\label{fw:v3v2-complete-form}
Estimating the tangential Poisson component before adjoining
\(\dd r\wedge Q_r\) sees a BRST defect.  Enforcing the chain identity while
forbidding a normal component loses the trace gain.  The bounded chain lift
exists only after the tangential and normal form components are kept as one
packet.
\end{firewall}

\subsection{Radial contraction and its scale boundary}
\label{sec:v3v2-radial}

For \(\omega=\alpha+\dd r\wedge\beta\), define
\begin{equation}
 (\mathsf h_r\omega)(r)
 =\int_0^r\beta(s)\,\dd s.
\label{eq:v3v2-radial-h}
\end{equation}

\begin{lemma}[Collar homotopy formula]
\label{lem:v3v2-radial-homotopy}
One has
\begin{equation}
 \dd_{\mathcal C}\mathsf h_r
 +\mathsf h_r\dd_{\mathcal C}
 =I-\pi^*\mathsf t_0.
\label{eq:v3v2-radial-homotopy}
\end{equation}
Thus \(\mathsf h_r\) contracts the smooth zero-trace collar complex.
On equal-order Sobolev spaces it satisfies
\begin{equation}
 \|\mathsf h_r\omega\|_{H^m}
 \le C_{m,\ell}\|\omega\|_{H^m}.
\label{eq:v3v2-radial-bound}
\end{equation}
\end{lemma}

\begin{proof}
Using \eqref{eq:v3v2-d-formula},
\begin{align*}
 \mathsf h_r\dd_{\mathcal C}\omega
 &=\alpha(r)-\alpha(0)
 -\dd_\Sigma\int_0^r\beta(s)\,\dd s,\\
 \dd_{\mathcal C}\mathsf h_r\omega
 &=\dd_\Sigma\int_0^r\beta(s)\,\dd s
 +\dd r\wedge\beta(r).
\end{align*}
Their sum is
\(\omega-\pi^*\alpha(0)\).  For the norm bound, tangential derivatives
commute with the integral, the zeroth normal derivative is controlled by
Cauchy--Schwarz with factor at most \(\ell\), and every positive normal
derivative reduces to a derivative of \(\beta\).
\end{proof}

\begin{assessment}[Why the abstract V1 repair can lose the trace scale]
\label{ass:v3v2-radial-scale}
Applying the V1 formula to the naive defect
\eqref{eq:v3v2-naive-defect} with \(\mathsf h_r\) gives
\[
 \widetilde{\mathcal E}-\mathsf h_r\Omega
 =\pi^*.
\]
The result is the constant chain lift.  This is valid on equal-order graph
spaces but is not bounded
(H^{s-1/2}(\Sigma)\to H^s(\mathcal C_\ell)).  The V1 repair theorem was
therefore not wrong: its required boundedness fails in this shifted trace
scale.  The boundary Hodge correction \(Q_r\) is the derivative-sensitive
replacement.
\end{assessment}

\subsection{Hodge repair with residual cohomology}
\label{sec:v3v2-partial-contraction}

Let \((K,q_K)\) now be any closed-range Hilbert complex with Laplacian
\begin{equation}
 \Delta_K=q_Kq_K^*+q_K^*q_K.
\label{eq:v3v2-DeltaK}
\end{equation}
Let \(\Pi_K\) project onto \(\ker\Delta_K\).  Assume a positive gap
\begin{equation}
 \Delta_K\ge\lambda_*I
 \quad\hbox{on }(I-\Pi_K)K
\label{eq:v3v2-K-gap}
\end{equation}
and put
\begin{equation}
 G_K=(\Delta_K|_{(I-\Pi_K)K})^{-1}(I-\Pi_K),
 \qquad
 h_K=q_K^*G_K.
\label{eq:v3v2-hK}
\end{equation}
Then
\begin{equation}
 q_Kh_K+h_Kq_K=I-\Pi_K,
 \qquad
 \|h_K\|\le\lambda_*^{-1/2}.
\label{eq:v3v2-Hodge-contraction}
\end{equation}

Indeed, the Green operator commutes with the Hilbert-complex differential,
so the first identity is
\((q_Kq_K^*+q_K^*q_K)G_K=I-\Pi_K\).  Moreover,
\[
 \|q_K^*G_Ky\|^2
 \le\langle\Delta_KG_Ky,G_Ky\rangle
 =\langle(I-\Pi_K)y,G_Ky\rangle
 \le\lambda_*^{-1}\|y\|^2.
\]

\begin{theorem}[Harmonic obstruction to the lift repair]
\label{thm:v3v2-harmonic-obstruction}
In the setting of the V1 lift defect
\(\Omega=q_XR-Rq_Y:Y\to K[1]\), define
\begin{equation}
 R_H=R-h_K\Omega.
\label{eq:v3v2-RH}
\end{equation}
Then
\begin{align}
 FR_H&=I_Y,
\label{eq:v3v2-RH-right}\\
 q_XR_H-R_Hq_Y&=\Pi_K\Omega.
\label{eq:v3v2-RH-defect}
\end{align}
Consequently the Hodge repair is an exact chain lift if and only if
\begin{equation}
 \Pi_K\Omega=0.
\label{eq:v3v2-harmonic-zero}
\end{equation}
Moreover,
\begin{equation}
 \|R_H\|
 \le\|R\|+lambda_*^{-1/2}\|\Omega\|.
\label{eq:v3v2-RH-bound}
\end{equation}
The harmonic defect represents the connecting cohomology class of V1.
\end{theorem}

\begin{proof}
The correction is kernel-valued, so \(FR_H=FR=I\).  The lift defect is a
kernel-valued cocycle:
\(q_K\Omega+\Omega q_Y=0\).  Therefore
\begin{align*}
 q_X(h_K\Omega)-(h_K\Omega)q_Y
 &=q_Kh_K\Omega-h_K\Omega q_Y\\
 &=(I-\Pi_K-h_Kq_K)\Omega-h_K\Omega q_Y\\
 &=(I-\Pi_K)\Omega.
\end{align*}
Subtracting this from the original defect \(\Omega\) gives
\eqref{eq:v3v2-RH-defect}.  A harmonic cocycle is the unique Hodge
representative of its cohomology class, so \(\Pi_K\Omega\) represents the
connecting class.  The norm estimate follows from
\eqref{eq:v3v2-Hodge-contraction}.
\end{proof}

\begin{corollary}[Harmonic modes must be routed, not inverted]
\label{cor:v3v2-harmonic-routing}
If \(\Pi_K\Omega\ne0\), no correction constructed from the Green operator on
\((I-\Pi_K)K\) can produce a chain right inverse.  The corresponding harmonic
sector must be included as physical boundary cohomology, cancelled by a
relative boundary field, or shown to have zero coefficient in the assembled
Standard Model lift defect.
\end{corollary}

\begin{proof}
The Green operator vanishes on \(\Pi_KK\), while
Theorem~\ref{thm:v3v2-harmonic-obstruction} identifies that component with
the connecting class.  A chain splitting would force the connecting map to
vanish by the V1 proposition.
\end{proof}

\subsection{Absolute trace by Hodge duality}
\label{sec:v3v2-absolute}

Let
\begin{equation}
 \mathsf n_0\omega
 =\mathsf t_0(\iota_{\partial_r}\omega)=\beta(0)
\label{eq:v3v2-normal-trace}
\end{equation}
be the normal trace.  Fix the standard degree-dependent signs so that the
product Hodge star identifies normal trace with the tangential trace of the
dual form and \(\delta=\pm *\dd *\).

\begin{corollary}[Absolute codifferential chain lift]
\label{cor:v3v2-absolute-lift}
Conjugating \(\mathcal E\) by the boundary and collar Hodge stars gives a
bounded operator \(\mathcal E^{\rm abs}\) such that
\begin{align}
 \mathsf n_0\mathcal E^{\rm abs}&=I,
\label{eq:v3v2-absolute-right}\\
 \delta_{\mathcal C}\mathcal E^{\rm abs}
 &=\mathcal E^{\rm abs}\delta_\Sigma,
\label{eq:v3v2-absolute-chain}
\end{align}
and the same half-derivative Sobolev estimate holds.
\end{corollary}

\begin{proof}
Apply the Hodge-star identities to
\eqref{eq:v3v2-E-right}--\eqref{eq:v3v2-E-chain}.  Hodge star is an isometry
on each smooth product metric and preserves Sobolev order.
\end{proof}

\begin{firewall}[Relative and absolute lifts are not automatically joint]
\label{fw:v3v2-relative-absolute}
The relative \(\dd\)-chain lift and absolute \(\delta\)-chain lift are
separately bounded.  Imposing both in a gauge-fixed Hodge--Dirac domain is an
intersection problem.  The common-angle and private-cokernel criteria from
the frozen second volume must still be checked for the stacked condition.
\end{firewall}

\subsection{Regulator-uniform requirements}
\label{sec:v3v2-regulator}

For a family \(\Sigma_n\) and boundary Hodge Laplacians \(\Delta_{\Sigma,n}\),
the formula \eqref{eq:v3v2-E} gives a uniform chain lift provided the
following quantities are uniform after bundle identification:
\begin{enumerate}
 \item the elliptic equivalence of the spectral and geometric Sobolev norms;
 \item the trace collar width and bounded-geometry constants;
 \item the zero-mode projection and a gap separating it from the positive
 spectrum when norm-continuous domain transport is required;
 \item the Riesz-transform bounds for
 \(\delta_n\Delta_n^{-1/2}(I-\mathsf H_n)\); and
 \item the parameter-dependent Poisson bounds in
 \eqref{eq:v3v2-Poisson-bound}.
\end{enumerate}
The algebraic chain identity is exact at every regulator.  Uniformity of its
operator norm is analytic data and is not inferred from exactness.

\begin{proposition}[Stable lift under graph functional-calculus convergence]
\label{prop:v3v2-regulator-stability}
Suppose the identified operators \(P_{r,n},Q_{r,n}\) satisfy
\begin{equation}
 \|P_{\cdot,n}-P_\cdot\|_{
 \mathcal L(H^{s-1/2},H^s(\mathcal C_\ell))}
 +
 \|Q_{\cdot,n}-Q_\cdot\|_{
 \mathcal L(H^{s-1/2},H^s(\mathcal C_\ell))}
 \longrightarrow0.
\label{eq:v3v2-functional-convergence}
\end{equation}
Then
\begin{equation}
 \mathcal E_n\longrightarrow\mathcal E
 \quad\hbox{in }\mathcal L(H^{s-1/2},H^s),
\label{eq:v3v2-E-convergence}
\end{equation}
and the limiting lift is a chain right inverse.
\end{proposition}

\begin{proof}
Use \(\mathcal E_n=P_{r,n}+\dd r\wedge Q_{r,n}\) and the triangle inequality.
The finite-regulator trace and chain identities pass to the limit because
trace and exterior differentiation are continuous between the stated graph
spaces.
\end{proof}

\begin{firewall}[Norm resolvent on \(L^2\) is not the full Poisson norm]
\label{fw:v3v2-resolvent-Poisson}
Norm-resolvent convergence stabilizes separated spectral projections, but
\eqref{eq:v3v2-functional-convergence} also controls a half-derivative
boundary-to-bulk operator over the entire positive spectrum.  A
parameter-elliptic or equivalent graph functional-calculus estimate is still
required for the interacting regulator family.
\end{firewall}

\subsection{V2 conclusion and next acquisition}
\label{sec:v3v2-conclusion}

V2 constructs an exact de Rham chain right inverse on a product collar.  The
normal Hodge component preserves the Poisson half-derivative gain while
cancelling the tangential lift's normal BRST defect.  Boundary harmonic modes
are carried as constant modes rather than divided by the Hodge Laplacian.  On
a general zero-trace kernel, the repaired defect is precisely
\(\Pi_K\Omega\), the harmonic representative of the connecting class.

V3 will pass from the collar to the global manifold.  It will use the long
exact sequence of the pair \((M,\partial M)\) to identify exactly which
boundary BRST cohomology classes extend globally, construct a bounded chain
splitting on the orthogonal complement of the connecting image, and show
that the nonextendable harmonic sector is the shared boundary cokernel that
must be routed into relative fields or inflow data.
\clearpage
\section*{Third Volume, Version V3: Global Boundary Extension and the Relative Cohomology Cokernel}
\addcontentsline{toc}{section}{Third Volume, Version V3: Global Boundary Extension and the Relative Cohomology Cokernel}
\label{sec:v3v3}

\subsection{Purpose}
\label{sec:v3v3-purpose}

V2 gives an exact, derivative-correct chain lift on a product collar.  A
global manifold has a topological obstruction: a closed boundary form need
not be the trace of a closed bulk form.  V3 identifies this obstruction as
the connecting homomorphism of the pair \((M,\partial M)\).  It then
constructs a bounded chain right inverse on the maximal Hodge subcomplex on
which the connecting map vanishes.

The construction has two corrections.  A finite-dimensional harmonic
correction first removes the harmonic part of the lift defect on nonclosed
boundary inputs.  The relative Green operator then removes the orthogonal
exact/coexact remainder.  Boundary harmonic classes on which the connecting
map is nonzero are retained as a finite-dimensional shared cokernel.  They
cannot be erased by a collar estimate.

\subsection{The short exact de Rham sequence}
\label{sec:v3v3-short-exact}

Let \(M\) be a compact oriented Riemannian manifold with smooth boundary
\(\Sigma\), and let
\begin{equation}
 \mathsf t:\Omega^\bullet(M)\longrightarrow\Omega^\bullet(\Sigma)
\label{eq:v3v3-trace}
\end{equation}
be pullback by the boundary inclusion.  Put
\begin{equation}
 \Omega^\bullet(M,\Sigma)=\ker\mathsf t.
\label{eq:v3v3-relative-complex}
\end{equation}
A collar extension proves degreewise surjectivity of \(\mathsf t\), while
naturality of exterior differentiation gives
\begin{equation}
 \mathsf t\dd_M=\dd_\Sigma\mathsf t.
\label{eq:v3v3-trace-chain}
\end{equation}
Thus
\begin{equation}
 0\longrightarrow\Omega^\bullet(M,\Sigma)
 \longrightarrow\Omega^\bullet(M)
 \xrightarrow{\ \mathsf t\ }
 \Omega^\bullet(\Sigma)
 \longrightarrow0
\label{eq:v3v3-short-exact}
\end{equation}
is a short exact sequence of cochain complexes.

\begin{theorem}[Long exact boundary-extension sequence]
\label{thm:v3v3-long-exact}
The sequence \eqref{eq:v3v3-short-exact} induces
\begin{equation}
 \cdots\longrightarrow
 H^k(M)
 \xrightarrow{\ \mathsf t^*\ }
 H^k(\Sigma)
 \xrightarrow{\ \partial_k\ }
 H^{k+1}(M,\Sigma)
 \longrightarrow H^{k+1}(M)
 \longrightarrow\cdots.
\label{eq:v3v3-long-exact}
\end{equation}
For a closed boundary form \(y\), choose any bulk lift \(x\) with
\(\mathsf tx=y\).  Then
\begin{equation}
 \partial_k[y]=[\dd_Mx]_{\rm rel}.
\label{eq:v3v3-connecting}
\end{equation}
In particular, \([y]\) has a closed bulk representative with trace in its
cohomology class if and only if
\begin{equation}
 \partial_k[y]=0.
\label{eq:v3v3-extension-criterion}
\end{equation}
\end{theorem}

\begin{proof}
Because \(\dd_\Sigma y=0\),
\(\mathsf t\dd_Mx=\dd_\Sigma\mathsf tx=0\), so
\(\dd_Mx\) is a relative form.  It is closed.  If \(x\) is replaced by
\(x+z\) with \(\mathsf tz=0\), the relative cocycle changes by the relative
boundary \(\dd_Mz\).  If \(y\) changes by \(\dd_\Sigma a\), lift \(a\) to
\(b\) and replace \(x\) by \(x+\dd_Mb\); the connecting cocycle is unchanged.
This proves that \eqref{eq:v3v3-connecting} is well defined.

If \(\partial_k[y]=0\), there is a relative \(z\) with
\(\dd_Mx=\dd_Mz\).  Then \(x-z\) is closed and has trace \(y\).  Conversely,
a closed lift makes \eqref{eq:v3v3-connecting} zero.  The remaining maps and
exactness are verified by the same lift-and-subtract argument in adjacent
degrees.
\end{proof}

\subsection{A sharp global obstruction}
\label{sec:v3v3-disk}

\begin{counterexample}[The angular class on the disk]
\label{sep:v3v3-disk}
Let \(M=D^2\) and \(\Sigma=S^1\).  The boundary form \(y=\dd\theta\) is closed
and has ordinary smooth bulk extensions, but it has no closed bulk extension.
Equivalently,
\begin{equation}
 \partial_1[\dd\theta]
 \ne0
 \quad\hbox{in }H^2(D^2,S^1).
\label{eq:v3v3-disk-obstruction}
\end{equation}
Therefore the global trace map has no chain right inverse on the full de Rham
boundary complex.
\end{counterexample}

\begin{proof}
If a closed one-form \(x\) on \(D^2\) had trace \(\dd\theta\), Stokes'
theorem would give
\[
 2\pi
 =\int_{S^1}\dd\theta
 =\int_{S^1}\mathsf tx
 =\int_{D^2}\dd x
 =0,
\]
a contradiction.  The last assertion follows because a chain right inverse
would send the closed boundary form to a closed bulk form with the same
trace.
\end{proof}

\begin{firewall}[Local collar exactness does not imply global exactness]
\label{fw:v3v3-local-global}
The V2 chain lift exists on a collar because that collar retracts onto its
physical face.  Extending the lifted form through the interior can create the
relative class \eqref{eq:v3v3-connecting}.  No improvement of the local trace
constant changes the disk obstruction.
\end{firewall}

\subsection{Hodge decomposition of the boundary target}
\label{sec:v3v3-boundary-Hodge}

Assume the boundary de Rham complex has closed range in the Sobolev graph
realization.  Let \(\mathcal H_\Sigma^k\) be its harmonic space and write the
orthogonal Hodge decomposition
\begin{equation}
 Y^k
 =\mathcal H_\Sigma^k
 \oplus\Ran\dd_\Sigma^{k-1}
 \oplus\Ran\delta_\Sigma^{k+1}.
\label{eq:v3v3-boundary-Hodge}
\end{equation}
Identify boundary cohomology with \(\mathcal H_\Sigma^k\).  Define
\begin{align}
 \mathcal H_{\rm ext}^k
 &=\ker
 \left(
  \partial_k:\mathcal H_\Sigma^k
  \to\mathcal H_{\rm rel}^{k+1}
 \right),
\label{eq:v3v3-Hext}\\
 \mathcal H_{\rm obs}^k
 &=\mathcal H_\Sigma^k\ominus\mathcal H_{\rm ext}^k.
\label{eq:v3v3-Hobs}
\end{align}
Here \(\mathcal H_{\rm rel}^{k+1}\) is the harmonic realization of
\(H^{k+1}(M,\Sigma)\).  Put
\begin{equation}
 Y_{\rm ext}^k
 =\mathcal H_{\rm ext}^k
 \oplus\Ran\dd_\Sigma^{k-1}
 \oplus\Ran\delta_\Sigma^{k+1}.
\label{eq:v3v3-Yext}
\end{equation}

\begin{lemma}[The extendable Hodge sector is a subcomplex]
\label{lem:v3v3-Yext-complex}
One has
\begin{equation}
 \dd_\Sigma Y_{\rm ext}^k
 \subseteq Y_{\rm ext}^{k+1}.
\label{eq:v3v3-Yext-complex}
\end{equation}
Every closed element of \(Y_{\rm ext}\) has zero connecting class.
\end{lemma}

\begin{proof}
The differential kills harmonic and exact forms.  It maps the coexact
summand into the exact summand of the next degree.  Hence
\eqref{eq:v3v3-Yext-complex} holds.  A closed boundary form has only harmonic
and exact Hodge components.  In \(Y_{\rm ext}\), its harmonic component lies
in \(\ker\partial\), while the exact component represents zero cohomology.
\end{proof}

\begin{assessment}[The shared harmonic cokernel]
\label{ass:v3v3-shared-cokernel}
The finite-dimensional space \(\mathcal H_{\rm obs}\) is the part of the
boundary BRST cohomology which cannot be supplied by a closed bulk private
lift.  It is therefore routed to the shared boundary sector in the V99
private/cokernel decomposition.  This routing is topological and occurs
before any eta, normal-flux, or operator-norm estimate.
\end{assessment}

\subsection{Relative Hodge data and an ordinary trace lift}
\label{sec:v3v3-relative-Hodge}

Let \(K=\ker\mathsf t\) be the relative bulk complex in a Hilbert graph
realization.  Assume its Hodge Laplacian
\begin{equation}
 \Delta_K=\dd_K\dd_K^*+\dd_K^*\dd_K
\label{eq:v3v3-DeltaK}
\end{equation}
has closed range.  Let \(\Pi_K\) be the relative harmonic projection and
assume
\begin{equation}
 \Delta_K\ge\lambda_K I
 \quad\hbox{on }(I-\Pi_K)K,
 \qquad
 \lambda_K>0.
\label{eq:v3v3-relative-gap}
\end{equation}
Set
\begin{equation}
 G_K=(\Delta_K|_{(I-\Pi_K)K})^{-1}(I-\Pi_K),
 \qquad
 h_K=\dd_K^*G_K.
\label{eq:v3v3-hK}
\end{equation}
Then
\begin{equation}
 \dd_Kh_K+h_K\dd_K=I-\Pi_K,
 \qquad
 \|h_K\|\le\lambda_K^{-1/2}.
\label{eq:v3v3-relative-contraction}
\end{equation}

Let
\begin{equation}
 R:Y_{\rm ext}\longrightarrow X
\label{eq:v3v3-ordinary-R}
\end{equation}
be any bounded degree-zero right inverse of the trace.  It may be obtained by
the V2 collar lift followed by a bounded interior extension; no chain identity
is assumed.  Its defect is
\begin{equation}
 \Omega=\dd_MR-R\dd_\Sigma:Y_{\rm ext}\to K[1].
\label{eq:v3v3-Omega}
\end{equation}
As in V1,
\begin{equation}
 \dd_K\Omega+\Omega\dd_\Sigma=0.
\label{eq:v3v3-Omega-cocycle}
\end{equation}
For closed \(y\), the harmonic vector \(\Pi_K\Omega y\) represents
\(\partial[y]\).

\subsection{The harmonic correction on nonclosed inputs}
\label{sec:v3v3-harmonic-correction}

The vanishing of \(\partial\) on closed inputs is not yet the operator
identity \(\Pi_K\Omega=0\) on all of \(Y_{\rm ext}\).  A finite-dimensional
harmonic correction supplies the missing strictness.

Let \(G_\Sigma\) be the Green operator of the boundary Hodge Laplacian and
define on exact forms
\begin{equation}
 s_{k+1}
 =\delta_\Sigma G_\Sigma:
 \Ran\dd_\Sigma^k\longrightarrow
 \Ran\delta_\Sigma^k.
\label{eq:v3v3-s}
\end{equation}
Then
\begin{equation}
 \dd_\Sigma s z=z
 \qquad(z\in\Ran\dd_\Sigma).
\label{eq:v3v3-s-right}
\end{equation}
Extend \(s\) by zero on the harmonic and coexact summands.  Define a
degree-zero kernel-valued operator \(C\) by
\begin{equation}
 Cz=
 \begin{cases}
  -\Pi_K\Omega s z,&z\in\Ran\dd_\Sigma,\\
  0,&z\perp\Ran\dd_\Sigma.
 \end{cases}
\label{eq:v3v3-C}
\end{equation}

\begin{lemma}[Removal of the harmonic defect]
\label{lem:v3v3-harmonic-removal}
Put
\begin{equation}
 R_0=R-C,
 \qquad
 \Omega_0=\dd_MR_0-R_0\dd_\Sigma.
\label{eq:v3v3-R0-Omega0}
\end{equation}
Then
\begin{align}
 \mathsf tR_0&=I_{Y_{\rm ext}},
\label{eq:v3v3-R0-right}\\
 \Omega_0&=\Omega+C\dd_\Sigma,
\label{eq:v3v3-Omega0}\\
 \Pi_K\Omega_0&=0.
\label{eq:v3v3-PiOmega0}
\end{align}
\end{lemma}

\begin{proof}
The range of \(C\) is relative harmonic, hence lies in \(K\), has zero trace,
and is annihilated by \(\dd_K\).  Therefore
\(\mathsf tR_0=\mathsf tR=I\) and
\[
 \Omega_0
 =\Omega-\dd_KC+C\dd_\Sigma
 =\Omega+C\dd_\Sigma.
\]
For \(y\in Y_{\rm ext}\), set
\(y_c=y-s\dd_\Sigma y\).  Equation
\eqref{eq:v3v3-s-right} makes \(y_c\) closed.  By
Lemma~\ref{lem:v3v3-Yext-complex}, its connecting class vanishes, so
\(\Pi_K\Omega y_c=0\).  Using \eqref{eq:v3v3-C},
\begin{align*}
 \Pi_K\Omega_0y
 &=\Pi_K\Omega y+C\dd_\Sigma y\\
 &=\Pi_K\Omega(y-s\dd_\Sigma y)=0.
\end{align*}
\end{proof}

\subsection{Global bounded chain splitting}
\label{sec:v3v3-global-splitting}

\begin{theorem}[Chain right inverse on the maximal extendable sector]
\label{thm:v3v3-global-splitting}
Define
\begin{equation}
 R_{\rm ext}
 =R_0-h_K\Omega_0
 =R-C-h_K(\Omega+C\dd_\Sigma).
\label{eq:v3v3-Rext}
\end{equation}
Then
\begin{align}
 \mathsf tR_{\rm ext}&=I_{Y_{\rm ext}},
\label{eq:v3v3-Rext-right}\\
 \dd_MR_{\rm ext}
 &=R_{\rm ext}\dd_\Sigma.
\label{eq:v3v3-Rext-chain}
\end{align}
If \(s\), \(R\), \(\Omega\), and \(\dd_\Sigma\) are bounded in the chosen
graph scales, then
\begin{align}
 \|C\|
 &\le\|\Omega\|\,\|s\|,
\label{eq:v3v3-C-bound}\\
 \|R_{\rm ext}\|
 &\le
 \|R\|+|C\|
 +\lambda_K^{-1/2}
 \bigl(\|\Omega\|+|C\|\,\|\dd_\Sigma\|\bigr).
\label{eq:v3v3-Rext-bound}
\end{align}
\end{theorem}

\begin{proof}
Both \(C\) and \(h_K\Omega_0\) take values in \(K\), so
\eqref{eq:v3v3-Rext-right} follows from
\eqref{eq:v3v3-R0-right}.  The defect \(\Omega_0\) is a kernel-valued
cocycle and has zero harmonic component by
\eqref{eq:v3v3-PiOmega0}.  The V2 Hodge-repair calculation gives
\begin{align*}
 &\dd_M(h_K\Omega_0)-(h_K\Omega_0)\dd_\Sigma\\
 &\qquad=
 (I-\Pi_K)\Omega_0=\Omega_0.
\end{align*}
Subtracting the correction from \(R_0\) proves
\eqref{eq:v3v3-Rext-chain}.  The first norm bound follows from
\eqref{eq:v3v3-C}; the second follows from
\eqref{eq:v3v3-Rext} and
\(\|h_K\|\le\lambda_K^{-1/2}\).
\end{proof}

\begin{corollary}[Maximality at the cohomology level]
\label{cor:v3v3-maximality}
No chain right inverse of \(\mathsf t\) can be defined on a boundary
subcomplex containing a closed harmonic vector
\(y\in\mathcal H_{\rm obs}\).  Hence \(Y_{\rm ext}\) contains every
boundary Hodge summand which can be included without contradicting the
connecting obstruction.
\end{corollary}

\begin{proof}
A chain right inverse would map \(y\) to a closed bulk form with trace \(y\),
forcing \(\partial[y]=0\) by
Theorem~\ref{thm:v3v3-long-exact}.  This contradicts
\(y\in\mathcal H_{\rm obs}=(\ker\partial)^{\perp}\), except for zero.
The exact and coexact summands are already included in
\eqref{eq:v3v3-Yext}; only harmonic cohomology can be obstructed.
\end{proof}

\begin{firewall}[The splitting is strict, not merely cohomological]
\label{fw:v3v3-strict-splitting}
Vanishing of the connecting map on closed elements removes only the harmonic
defect there.  Lemma~\ref{lem:v3v3-harmonic-removal} is needed to turn this
into \(\Pi_K\Omega_0=0\) on every input before the relative Green operator is
applied.  Omitting that finite-dimensional correction would confuse a
cohomology splitting with a chain-level right inverse.
\end{firewall}

\subsection{Uniform family and rank stability}
\label{sec:v3v3-family}

For regulator \(n\), let
\begin{equation}
 J_n:\mathcal H_{\Sigma,n}
 \longrightarrow\mathcal H_{{\rm rel},n}
\label{eq:v3v3-Jn}
\end{equation}
be the harmonic matrix representing the connecting map after fixed
finite-dimensional identifications.  Let
\(P_{{\rm ext},n}\) project onto \(\ker J_n\).

\begin{theorem}[Stable extendable/obstructed harmonic split]
\label{thm:v3v3-rank-stability}
Assume
\begin{equation}
 J_n\longrightarrow J
\label{eq:v3v3-J-convergence}
\end{equation}
in operator norm and there is \(\sigma_*>0\) such that every nonzero singular
value of \(J_n\) is at least \(\sigma_*\).  Then the ranks are eventually
constant and
\begin{equation}
 P_{{\rm ext},n}\longrightarrow P_{\rm ext}
 \quad\hbox{in operator norm}.
\label{eq:v3v3-Pext-convergence}
\end{equation}
The complementary obstructed projections converge as well.
\end{theorem}

\begin{proof}
The nonzero spectrum of \(J_n^*J_n\) lies in
\([\sigma_*^2,\infty)\), while zero is separated by the common gap.  Norm
convergence gives
\(J_n^*J_n\to J^*J\).  The Riesz projection around a circle of radius less
than \(\sigma_*^2\) therefore converges in norm and equals the kernel
projection.  Constant rank and convergence of complementary projections
follow.
\end{proof}

\begin{corollary}[Uniform global BRST private lift]
\label{cor:v3v3-uniform-lift}
Suppose, in addition to Theorem~\ref{thm:v3v3-rank-stability}, that the
ordinary trace lifts, boundary exact right inverses, relative Green
contractions, and differentials in
\eqref{eq:v3v3-Rext-bound} have common bounds and converge in their graph
operator norms.  Then \(R_{{\rm ext},n}\) is uniformly bounded and converges
to the limiting chain right inverse on the transported extendable sector.
The resulting BRST private angle has a positive regulator-independent lower
bound.
\end{corollary}

\begin{proof}
The stable harmonic projections identify the domains of all finite-dimensional
corrections.  Apply \eqref{eq:v3v3-C-bound}--\eqref{eq:v3v3-Rext-bound}
uniformly.  Norm convergence is preserved under the finite sums and products
in \eqref{eq:v3v3-Rext}.
\end{proof}

\begin{counterexample}[Rank stability is indispensable]
\label{sep:v3v3-rank-collapse}
Let \(J_n:\R\to\R\) be multiplication by \(n^{-1}\).  Then
\(J_n\to0\), but
\begin{equation}
 \ker J_n=\{0\},
 \qquad
 \ker J=\R.
\label{eq:v3v3-rank-collapse}
\end{equation}
The extendable harmonic projection jumps and the inverse on the obstructed
sector has norm \(n\).
\end{counterexample}

\begin{proof}
The assertions follow directly from scalar multiplication.  The positive
singular-value margin tends to zero.
\end{proof}

\subsection{Interface with relative anomaly and inflow data}
\label{sec:v3v3-inflow}

The obstruction space \(\mathcal H_{\rm obs}\) is finite dimensional.  Let
\begin{equation}
 \mathfrak o=\partial[y]\in\mathcal H_{\rm rel}[1]
\label{eq:v3v3-obstruction-vector}
\end{equation}
be the image of its boundary class.  In the relative Ward mapping cone, this
class can be removed only by one of the following exact mechanisms:
\begin{enumerate}
 \item restrict the admissible boundary cohomology to \(\ker\partial\);
 \item add a boundary or inflow variable whose differential maps onto
 \(\mathfrak o\); or
 \item prove that the complete Standard Model species/BRST coefficient of
 \(\mathfrak o\) vanishes before norms and determinant phases are taken.
\end{enumerate}

\begin{firewall}[A finite-dimensional cokernel is still a real obstruction]
\label{fw:v3v3-finite-obstruction}
Finite dimensionality makes the remaining topology computable; it does not
make it removable.  Dividing by a small nonzero singular value of \(J_n\) or
silently excluding a boundary harmonic mode would destroy uniform domain
control or change the boundary theory.
\end{firewall}

\subsection{V3 conclusion and next acquisition}
\label{sec:v3v3-conclusion}

V3 proves the global boundary extension theorem.  The maximal Hodge
subcomplex admitting a bounded strict de Rham chain lift contains all exact
and coexact boundary data and precisely the harmonic classes in
\(\ker\partial\).  A finite-dimensional harmonic correction is required
before the relative Green contraction.  The orthogonal harmonic complement
is the shared boundary cokernel and is stable only under a nonzero singular
value margin for the connecting matrix.

V4 will construct the minimal relative inflow extension for this obstructed
sector.  It will adjoin one shifted boundary generator for each independent
class in \(\operatorname{Ran}\partial\), prove that the enlarged mapping cone
is acyclic on the obstruction block, quantify its Hodge gap, and distinguish
this algebraic cancellation from determinant-line triviality and eta
convergence.
\clearpage
\section*{Third Volume, Version V4: Minimal Inflow Attachment for the Obstructed Harmonic Sector}
\addcontentsline{toc}{section}{Third Volume, Version V4: Minimal Inflow Attachment for the Obstructed Harmonic Sector}
\label{sec:v3v4}

\subsection{Purpose}
\label{sec:v3v4-purpose}

V3 isolates the finite-dimensional obstruction
\(\mathcal R=\operatorname{Ran}\partial\) inside relative harmonic
cohomology.  V4 constructs the minimal algebraic inflow extension which kills
this block.  One generator of degree one lower is adjoined for each
independent obstruction class, and its differential is the chosen harmonic
representative.  The resulting two-term block is contractible.

The construction gives exact formulas for the new cohomology, contracting
homotopy, Hodge gap, counterterm primitive, and regulator stability.  It is an
extension of the relative BRST complex.  It does not by itself construct a
local unitary boundary field theory, choose an integral anomaly lattice, or
trivialize a determinant line.

\subsection{Harmonic obstruction block}
\label{sec:v3v4-obstruction}

Let \((K^\bullet,q)\) be the relative Hilbert complex from V3, with orthogonal
harmonic projection \(\Pi_K\).  In every degree choose a finite-dimensional
subspace
\begin{equation}
 \mathcal R^{p+1}
 \subseteq\ker\Delta_K^{p+1}
\label{eq:v3v4-R}
\end{equation}
containing the harmonic representatives of
\(\operatorname{Ran}\partial_p\).  Let
\(P_{\mathcal R}^{p+1}\) denote its orthogonal projection.  Let
\(I^p\) be a Hilbert space of the same dimension and choose an isomorphism
\begin{equation}
 L^p:I^p\longrightarrow\mathcal R^{p+1}.
\label{eq:v3v4-L}
\end{equation}
The superscript records cochain degree; \(I\) is the obstruction space shifted
one degree downward.

Define the enlarged graded space
\begin{equation}
 \widetilde K^p=K^p\oplus I^p
\label{eq:v3v4-Ktilde}
\end{equation}
and differential
\begin{equation}
 \widetilde q(k,i)
 =
 \bigl(qk+L^pi,0\bigr).
\label{eq:v3v4-qtilde}
\end{equation}
In \eqref{eq:v3v4-qtilde}, the second zero lies in \(I^{p+1}\); only those
degrees for which an obstruction space is present have a nonzero \(I^p\).

\begin{lemma}[The inflow attachment is a complex]
\label{lem:v3v4-square-zero}
One has
\begin{equation}
 \widetilde q^{,2}=0.
\label{eq:v3v4-square-zero}
\end{equation}
\end{lemma}

\begin{proof}
The image of every \(L^p\) is harmonic and hence \(q\)-closed.  Therefore
\[
 \widetilde q^{,2}(k,i)
 =\bigl(q^2k+qL^pi,0\bigr)=0.
\]
\end{proof}

\subsection{Cohomology after attachment}
\label{sec:v3v4-cohomology}

Let
\begin{equation}
 \mathcal H_{\rm rem}^p
 =\ker\Delta_K^p\ominus\mathcal R^p.
\label{eq:v3v4-Hrem}
\end{equation}

\begin{theorem}[The attachment kills exactly the obstruction image]
\label{thm:v3v4-cohomology}
Assume each \(L^p\) is injective.  Then
\begin{equation}
 H^p(\widetilde K,\widetilde q)
 \cong
 H^p(K,q)/\mathcal R^p
 \cong\mathcal H_{\rm rem}^p.
\label{eq:v3v4-cohomology}
\end{equation}
The extension creates no additional cohomology.
\end{theorem}

\begin{proof}
Suppose \(\widetilde q(k,i)=0\).  The two summands
\(qk\) and \(L^pi\) in \(K^{p+1}\) are orthogonal: the first is exact and
the second harmonic.  Hence both vanish.  Injectivity of \(L^p\) gives
\(i=0\), while \(qk=0\).  Thus every cocycle is \((k,0)\) with \(k\) a
cocycle of \(K\).

A boundary in degree \(p\) has the form
\[
 \widetilde q(k',i')=(qk'+L^{p-1}i',0).
\]
Compared with ordinary \(q\)-boundaries, these boundaries add exactly
\(\mathcal R^p\).  Quotienting the ordinary cohomology by that harmonic
subspace proves the first isomorphism.  Hodge representation proves the
second.  Since a cocycle cannot have a nonzero \(I\)-component, no new class
is created.
\end{proof}

\begin{corollary}[Removal of the V3 connecting image]
\label{cor:v3v4-connecting-image}
If
\begin{equation}
 \mathcal R^{p+1}
 =\operatorname{Ran}
 \bigl(partial_p:H^p(\partial M)\to H^{p+1}(M,\partial M)\bigr)
\label{eq:v3v4-R-connecting}
\end{equation}
under harmonic identification, then every global boundary-extension
obstruction found in V3 becomes exact in \(\widetilde K\), while relative
cohomology outside \(\operatorname{Ran}\partial\) is unchanged.
\end{corollary}

\begin{proof}
Apply Theorem~\ref{thm:v3v4-cohomology} with the stated choice of
\(\mathcal R\).
\end{proof}

\subsection{Explicit contraction and Hodge gap}
\label{sec:v3v4-contraction}

Assume the original relative Laplacian has gap
\begin{equation}
 \Delta_K\ge\lambda_K I
 \quad\hbox{on }(I-\Pi_K)K,
 \qquad \lambda_K>0,
\label{eq:v3v4-old-gap}
\end{equation}
and let
\begin{equation}
 h_K=q^*(\Delta_K|_{(I-\Pi_K)K})^{-1}(I-\Pi_K).
\label{eq:v3v4-hK}
\end{equation}
Assume also
\begin{equation}
 \|L^pi\|\ge\sigma_L\|i\|
 \qquad\hbox{for every }p,
\label{eq:v3v4-L-gap}
\end{equation}
with \(\sigma_L>0\).  Let \(L^{-1}\) denote the inverse on
\(\mathcal R\).

Define a degree-minus-one operator on the enlarged complex by
\begin{equation}
 \widetilde h(k,i)
 =
 \left(
  h_K(I-\Pi_K)k,
  L^{-1}P_{\mathcal R}k
 \right).
\label{eq:v3v4-htilde}
\end{equation}
The second component belongs to the appropriately shifted \(I\)-space.

\begin{theorem}[Quantitative contraction of the inflow block]
\label{thm:v3v4-contraction}
Let
\begin{equation}
 \widetilde\Pi_{\rm rem}(k,i)
 =\bigl((\Pi_K-P_{\mathcal R})k,0\bigr).
\label{eq:v3v4-Pirem}
\end{equation}
Then
\begin{equation}
 \widetilde q\widetilde h
 +\widetilde h\widetilde q
 =I-\widetilde\Pi_{\rm rem}.
\label{eq:v3v4-contraction}
\end{equation}
Moreover,
\begin{equation}
 \|\widetilde h\|
 \le
 \max\{\lambda_K^{-1/2},\sigma_L^{-1}\}.
\label{eq:v3v4-h-bound}
\end{equation}
With the orthogonal direct-sum inner product, the positive spectrum of the
enlarged Hodge Laplacian is bounded below by
\begin{equation}
 \widetilde\lambda
 \ge\min\{\lambda_K,\sigma_L^2\}.
\label{eq:v3v4-new-gap}
\end{equation}
\end{theorem}

\begin{proof}
The ordinary Hodge contraction gives
\[
 qh_K+h_Kq=I-\Pi_K
\]
on \(K\).  On \(\mathcal R\), the new differential and homotopy are the
two maps
\[
 I\xrightarrow{\ L\ }\mathcal R,
 \qquad
 \mathcal R\xrightarrow{\ L^{-1}\ }I.
\]
Thus \(\widetilde q\widetilde h\) is the identity on \(\mathcal R\), while
\(\widetilde h\widetilde q\) is the identity on \(I\).  The old contraction,
new two-term contraction, and remaining harmonic projection act on mutually
orthogonal blocks.  Their sum is
\eqref{eq:v3v4-contraction}, and the operator norm is the maximum of the two
block norms.

On \(I\), the enlarged Laplacian is \(L^*L\); on \(\mathcal R\), it is
\(LL^*\); on the old nonharmonic block it is \(\Delta_K\); and on
\(\mathcal H_{\rm rem}\) it is zero.  The spectral lower bounds on the three
positive blocks give \eqref{eq:v3v4-new-gap}.
\end{proof}

\begin{corollary}[Isometric attachment]
\label{cor:v3v4-isometric}
Choosing \(I^p\) as an isometric copy of \(\mathcal R^{p+1}\) and \(L^p\)
as the identifying isometry gives
\begin{equation}
 \sigma_L=1,
 \qquad
 \widetilde\lambda\ge\min\{\lambda_K,1\}.
\label{eq:v3v4-isometric-gap}
\end{equation}
\end{corollary}

\begin{proof}
For an isometry, \(L^*L=I\) and \(LL^*=P_{\mathcal R}\).  Apply
Theorem~\ref{thm:v3v4-contraction}.
\end{proof}

\subsection{Minimality of the new generators}
\label{sec:v3v4-minimality}

\begin{theorem}[One generator per independent obstruction is minimal]
\label{thm:v3v4-minimality}
Let \(r=\dim\mathcal R^{p+1}\).  Suppose a two-term extension
adjoins a space \(J^p\), no new differential enters \(J^p\), and its
differential
\begin{equation}
 A:J^p\longrightarrow\ker\Delta_K^{p+1}
\label{eq:v3v4-general-attachment}
\end{equation}
has image containing \(\mathcal R^{p+1}\).  Then
\begin{equation}
 \dim J^p\ge r.
\label{eq:v3v4-dimension-lower}
\end{equation}
If equality holds and the extension creates no new class in degree \(p\),
then \(A\) is injective and maps \(J^p\) isomorphically onto
\(\mathcal R^{p+1}\).
\end{theorem}

\begin{proof}
Since \(\mathcal R\subseteq\operatorname{Ran}A\), rank--nullity gives
\(r\le\operatorname{rank}A\le\dim J\).  If \(\dim J=r\), both inequalities
are equalities.  Any nonzero vector in \(\ker A\) is a newly adjoined cocycle
which is not the boundary of a lower new generator in the minimal two-term
extension.  Absence of new degree-\(p\) cohomology therefore forces
\(\ker A=0\), and the range is exactly \(\mathcal R\).
\end{proof}

\begin{assessment}[Minimal algebra does not select physical fields]
\label{ass:v3v4-minimal-physical}
The dimension theorem fixes the number of independent shifted generators in
an algebraic resolution.  It does not determine their spin, statistics,
local kinetic operator, gauge representation, or boundary action.  Those
physical data must realize \(L\) as a local relative Ward differential.
\end{assessment}

\subsection{Primitive for an obstruction coefficient}
\label{sec:v3v4-primitive}

Let \(a\in\mathcal R^{p+1}\) be a relative anomaly or extension-obstruction
coefficient.  Define
\begin{equation}
 i_a=(L^p)^{-1}a\in I^p.
\label{eq:v3v4-ia}
\end{equation}

\begin{proposition}[Explicit inflow primitive]
\label{prop:v3v4-primitive}
In the enlarged complex,
\begin{equation}
 \widetilde q(0,i_a)=(a,0),
\label{eq:v3v4-primitive}
\end{equation}
with
\begin{equation}
 \|i_a\|le\sigma_L^{-1}\|a\|.
\label{eq:v3v4-primitive-bound}
\end{equation}
Thus the counterterm or inflow amplitude is uniformly controlled exactly
when the attachment singular value is bounded away from zero.
\end{proposition}

\begin{proof}
Equation \eqref{eq:v3v4-primitive} follows from
\(Li_a=a\).  The norm estimate is
\eqref{eq:v3v4-L-gap}.
\end{proof}

\begin{counterexample}[Acyclic at every cutoff but singular in the limit]
\label{sep:v3v4-gap-collapse}
Let \(I_n=\mathcal R_n=\R\) and let
\begin{equation}
 L_n i=n^{-1}i.
\label{eq:v3v4-Ln-collapse}
\end{equation}
The two-term complex is acyclic for every finite \(n\), but
\begin{equation}
 \|L_n^{-1}\|=n,
 \qquad
 \lambda_{\rm inflow,n}=n^{-2}.
\label{eq:v3v4-collapse-rate}
\end{equation}
At the limit the differential vanishes and both degrees acquire cohomology.
\end{counterexample}

\begin{proof}
Scalar multiplication by \(n^{-1}\) is invertible at finite \(n\).  Its
inverse and squared singular value are as displayed.  The limiting map is
zero.
\end{proof}

\begin{firewall}[Finite-cutoff exactness is not uniform cancellation]
\label{fw:v3v4-finite-exact}
The counterexample satisfies the algebraic master identity at every cutoff.
Without a uniform inflow gap, its primitive diverges and the obstruction
reappears in the continuum.  A regulator-dependent rescaling which hides
\(n\) in the norm merely moves the same loss into the boundary kinetic or
trace estimate.
\end{firewall}

\subsection{Regulator-stable attachment}
\label{sec:v3v4-regulator}

Let \(P_{\mathcal R,n}\) be the obstruction projections from the V3
connecting matrices.  Assume
\begin{equation}
 P_{\mathcal R,n}\longrightarrow P_{\mathcal R}
\label{eq:v3v4-R-projection-convergence}
\end{equation}
in norm and choose unitary range identifiers
\begin{equation}
 U_n:\mathcal R\longrightarrow\mathcal R_n
\label{eq:v3v4-Un}
\end{equation}
which converge in the transported chart.  Taking one fixed copy
\(I\cong\mathcal R\), define
\begin{equation}
 L_n=U_n:I\longrightarrow\mathcal R_n.
\label{eq:v3v4-Ln-isometric}
\end{equation}

\begin{theorem}[Uniform isometric inflow resolution]
\label{thm:v3v4-regulator}
Under the stable-rank hypothesis of V3, the choice
\eqref{eq:v3v4-Ln-isometric} gives
\begin{equation}
 \|L_n\|=\|L_n^{-1}\|=1
\label{eq:v3v4-uniform-isometry}
\end{equation}
for all \(n\).  If the old relative Green contractions converge in graph
norm, then the enlarged projections and contractions converge, and the
inflow block retains gap one.
\end{theorem}

\begin{proof}
Each \(U_n\) is unitary on the finite-dimensional obstruction range.  The
explicit formulas
\eqref{eq:v3v4-Pirem} and \eqref{eq:v3v4-htilde} are finite sums of the old
contraction, the convergent obstruction projection, and the unitary
identifiers.  They therefore converge in operator norm.  The new Laplacian
blocks are \(U_n^*U_n=I_I\) and
\(U_nU_n^*=P_{\mathcal R,n}\), so their positive eigenvalues equal one.
\end{proof}

\begin{assessment}[Local frames versus global holonomy]
\label{ass:v3v4-holonomy}
Stable finite-dimensional ranges admit local unitary identifiers.  Along a
closed parameter loop, the transported frame may return with nontrivial
holonomy.  The local isometric resolution remains acyclic, but a globally
single-valued primitive or determinant phase requires a compatible
trivialization.  This is not decided by the pointwise Hodge gap.
\end{assessment}

\subsection{Integral and determinant-line firewalls}
\label{sec:v3v4-global-firewalls}

\begin{firewall}[Real cohomology misses torsion]
\label{fw:v3v4-torsion}
The harmonic construction resolves the real or complex de Rham obstruction.
Torsion classes and global gauge anomalies are invisible to harmonic forms.
They require the appropriate integral, differential-cohomological, bordism,
or determinant-line formulation.  An arbitrary real isometry
\(L:I\to\mathcal R\) need not preserve an integral anomaly lattice.
\end{firewall}

\begin{firewall}[Acyclic Ward complex does not trivialize the determinant line]
\label{fw:v3v4-determinant}
Equation \eqref{eq:v3v4-cohomology} removes a local cochain obstruction.
The fermion determinant may still have eta phase, spectral flow, or global
holonomy.  V92's heat and determinant-line gates remain necessary even when
the extended local Ward complex is acyclic on \(\mathcal R\).
\end{firewall}

\begin{firewall}[An abstract generator is not yet a boundary QFT]
\label{fw:v3v4-boundary-QFT}
To use the inflow attachment in the Standard Model, one must realize the
shifted generators by local boundary or bulk-inflow degrees of freedom with a
well-posed kinetic operator, correct statistics, BRST transformation,
renormalized measure, and stress tensor.  V4 proves the minimal algebraic
resolution and its norm; it does not assert that every obstruction admits
such a physical realization.
\end{firewall}

\subsection{Interface with the common boundary angle}
\label{sec:v3v4-common-angle}

The enlarged relative BRST complex replaces the obstructed harmonic block by
an acyclic two-term pair.  Let \(R_{\rm ext}\) be the V3 chain lift on the
extendable boundary sector and let \(L^{-1}\) be the inflow primitive on the
obstructed sector.  Their direct sum has norm bounded by
\begin{equation}
 M_{\rm B}
 =\max\{\|R_{\rm ext}\|,\sigma_L^{-1}\}
\label{eq:v3v4-MB}
\end{equation}
in the orthogonal sum norm.

\begin{corollary}[BRST boundary margin after inflow attachment]
\label{cor:v3v4-BRST-margin}
If the shifted inflow generator belongs to the admissible shared boundary
graph space and its trace realization has the same norm as \(I\), then the
enlarged BRST boundary block has right-inverse margin at least
\begin{equation}
 \gamma_{\rm B}^{\rm ext}
 \ge M_{\rm B}^{-1}.
\label{eq:v3v4-BRST-margin}
\end{equation}
\end{corollary}

\begin{proof}
Use \(R_{\rm ext}\) on the extendable sector and \(L^{-1}\) on the
orthogonal obstruction sector.  The direct-sum right inverse has norm
\eqref{eq:v3v4-MB}; take its reciprocal.
\end{proof}

\begin{firewall}[The graph realization is an additional hypothesis]
\label{fw:v3v4-graph-realization}
The algebraic norm on \(I\) may not equal the physical graph norm of a local
inflow field.  Corollary~\ref{cor:v3v4-BRST-margin} applies only after a
uniform trace/kinetic realization identifies those norms.  This is the next
analytic acquisition, not a convention.
\end{firewall}

\subsection{V4 conclusion and V5 audit mandate}
\label{sec:v3v4-conclusion}

V4 attaches the minimal number of shifted generators needed to kill the V3
connecting image.  The enlarged cohomology is the old relative cohomology
modulo that image, the obstruction pair has explicit contraction
\(L^{-1}\), and its Hodge gap is the squared smallest singular value of
\(L\).  Stable-rank isometric attachments remain uniformly acyclic, while a
collapsing singular value gives an exact finite-cutoff counterexample.

V5 is the mandatory companion audit.  It will inspect the newest Yang--Mills
and Navier--Stokes notes before choosing the physical graph realization of
the inflow generator.  The audit will test whether the correct next step is a
local boundary kinetic lift or an upstream reformulation of the shared
normal-flux packet; no numerical or analytic estimate will be imported
across programmes without a type-correct map.
\clearpage
\section{Third Volume, Version V5: Companion Audit and the Single-Payer Physical Inflow Norm}
\label{sec:v3v5}

\subsection{Purpose of this version}
\label{sec:v3v5-purpose}

V4 killed the finite-dimensional connecting image by an algebraically
minimal shifted complex.  Its last firewall was essential: the convenient
Hilbert norm placed on the new generator need not be the graph norm induced
by a physical kinetic operator.  V5 performs the scheduled companion audit
and closes the finite-dimensional part of that gap.  The result is a
canonical least-energy inflow primitive, measured in its physical kinetic
metric, together with the exact Schur complement which must stay coercive
along a regulator sequence.

The construction is deliberately finite-dimensional.  It applies to the
harmonic obstruction block isolated in V3 and V4.  Producing the metric from
a local propagating boundary field remains a separate realization problem.

\subsection{Mandatory fifth-version companion audit}
\label{sec:v3v5-audit}

The audit was performed before selecting the V5 argument.  The files and
their fingerprints were as follows.
\begin{center}
\begin{tabular}{p{0.19\textwidth}p{0.49\textwidth}p{0.25\textwidth}}
Programme & inspected state & SHA--256 \\
\hline
Navier--Stokes &
\texttt{Renewal\_NS\_Stability\_Research\_Notes\_V31.tex},
887537 bytes &
\texttt{F1121B62238AD5491BB7CF03635EA9934942EDF573ED8FAAA9EE79E1D38A6696}
\\
Yang--Mills &
\texttt{Renewal\_Yang\_Mills\_Mass\_Gap\_Research\_Notes\_V79.tex}
and the file named V80, each 1176151 bytes &
\texttt{B2C82C8ACAE1E28CFB348F00FA13AEEFFD25884697DF52A2CBFBF915EC317CBE}
\end{tabular}
\end{center}
The two Yang--Mills files are byte-identical.  Thus the file named V80 does
not contain a mathematical state beyond the state internally ending at V79.
Neither file was modified.

The Navier--Stokes formation ladder gives a type-correct warning.  Passing
from a marked source to an unmarked source pays an exact first-moment cost;
the missing shell weight cannot be recovered merely by renaming the source
norm.  The high-parent lift returns continuation stock rather than an
additional free gain.  The persistent branch is a one-use resource.

The Yang--Mills incidence analysis gives the corresponding bookkeeping
rule.  Internal-bank identities must be formed before norms are taken.
Output ancestry is not a second copy of input mass, and a sole payer cannot
pay both endpoints of a pair minimum.  Translated to the present programme,
the shifted inflow coordinate cannot be counted once as a cohomological
primitive and again as an independent boundary source.  Its Ward, kinetic,
and normal-flux appearances must be images of one and the same variable.

\begin{firewall}[No transfer by analogy]
\label{fw:v3v5-no-analogy}
The companion arguments do not prove a Standard Model estimate.  They
identify a structural failure mode.  The construction below supplies the
missing type-correct map and proves its norm directly.
\end{firewall}

\subsection{The physical metric on the obstruction primitive}
\label{sec:v3v5-physical-metric}

Let \(I\) and \(\mathcal R\) be finite-dimensional real or complex Hilbert
spaces.  The space \(\mathcal R\) is the harmonic connecting range from V3,
and \(I\) is the shifted generator space from V4.  Let
\begin{equation}
 G:I\longrightarrow I
\label{eq:v3v5-G}
\end{equation}
be positive and self-adjoint, and let
\begin{equation}
 L:I\longrightarrow\mathcal R
\label{eq:v3v5-L}
\end{equation}
be onto.  The physical kinetic norm is
\begin{equation}
 \|i\|_G^2:=\langle Gi,i\rangle_I.
\label{eq:v3v5-Gnorm}
\end{equation}
The relevant Gram operator is not \(LL^*\) computed in an arbitrary norm.
It is
\begin{equation}
 C:=LG^{-1}L^*: \mathcal R\longrightarrow\mathcal R.
\label{eq:v3v5-C}
\end{equation}

\begin{lemma}[Positivity of the physical Gram operator]
\label{lem:v3v5-C-positive}
The operator \(C\) is positive definite.  In particular, it is invertible.
\end{lemma}

\begin{proof}
For \(a\in\mathcal R\),
\[
 \langle Ca,a\rangle_{\mathcal R}
 =\langle G^{-1}L^*a,L^*a\rangle_I
 =\|G^{-1/2}L^*a\|_I^2.
\]
If this vanishes, then \(L^*a=0\).  Surjectivity of \(L\) gives
\(\ker L^*=(\operatorname{Ran}L)^\perp=\{0\}\), so \(a=0\).
Finite dimensionality now gives invertibility.
\end{proof}

\begin{definition}[Physical inflow margin]
\label{def:v3v5-sigma}
Set
\begin{equation}
 \sigma_{\rm phys}^2:=\lambda_{\min}(C)>0.
\label{eq:v3v5-sigma}
\end{equation}
Equivalently, \(\sigma_{\rm phys}\) is the least nonzero singular value of
\(LG^{-1/2}\).
\end{definition}

\subsection{Canonical least-energy primitive}
\label{sec:v3v5-least-energy}

\begin{theorem}[Physical Moore--Penrose primitive]
\label{thm:v3v5-minimizer}
For each \(a\in\mathcal R\), the affine constraint \(Li=a\) has a unique
minimizer in the kinetic norm.  It is
\begin{equation}
 i_*(a):=G^{-1}L^*C^{-1}a.
\label{eq:v3v5-istar}
\end{equation}
Moreover,
\begin{align}
 Li_*(a)&=a,
\label{eq:v3v5-right-inverse}\\
 \|i_*(a)\|_G^2&=\langle C^{-1}a,a\rangle_{\mathcal R},
\label{eq:v3v5-energy-identity}\\
 \|i_*(a)\|_G&\le \sigma_{\rm phys}^{-1}\|a\|_{\mathcal R}.
\label{eq:v3v5-primitive-bound}
\end{align}
If \(Li=a\), then the exact Pythagorean identity
\begin{equation}
 \|i\|_G^2
 =\|i_*(a)\|_G^2+\|i-i_*(a)\|_G^2
\label{eq:v3v5-pythagoras}
\end{equation}
holds.
\end{theorem}

\begin{proof}
Equation \eqref{eq:v3v5-right-inverse} follows immediately from
\(LG^{-1}L^*=C\).  If \(k\in\ker L\), then
\[
 \langle i_*(a),k\rangle_G
 =\langle Gi_*(a),k\rangle_I
 =\langle L^*C^{-1}a,k\rangle_I
 =\langle C^{-1}a,Lk\rangle_{\mathcal R}=0.
\]
Every other solution is \(i_*(a)+k\) for a unique \(k\in\ker L\).
Orthogonality proves \eqref{eq:v3v5-pythagoras}, hence uniqueness and
minimality.  Finally,
\[
 \|i_*(a)\|_G^2
 =\langle L^*C^{-1}a,G^{-1}L^*C^{-1}a\rangle_I
 =\langle C^{-1}a,a\rangle_{\mathcal R}.
\]
Since \(C\ge\sigma_{\rm phys}^2I\), spectral calculus gives
\(C^{-1}\le\sigma_{\rm phys}^{-2}I\), which proves
\eqref{eq:v3v5-primitive-bound}.
\end{proof}

\begin{corollary}[Exact effective obstruction cost]
\label{cor:v3v5-effective-cost}
The constrained quadratic kinetic cost is
\begin{equation}
 \mathcal E_{\rm in}(a)
 :=\min_{Li=a}\frac12\langle Gi,i\rangle_I
 =\frac12\langle C^{-1}a,a\rangle_{\mathcal R}.
\label{eq:v3v5-effective-cost}
\end{equation}
Thus \(C^{-1}\), rather than an arbitrarily normalized copy of the identity,
is the exact quadratic form seen by the obstruction coordinate.
\end{corollary}

\subsection{One inflow coordinate, one boundary packet}
\label{sec:v3v5-single-payer}

Let \(Z_\partial\) be the normal-flux Hilbert space selected by the common
boundary angle, and let
\begin{equation}
 B:I\longrightarrow Z_\partial
\label{eq:v3v5-B}
\end{equation}
be the flux trace of the inflow realization.  Define
\begin{equation}
 \beta:=\|BG^{-1/2}\|_{I\to Z_\partial}.
\label{eq:v3v5-beta}
\end{equation}

\begin{proposition}[Single-payer graph packet]
\label{prop:v3v5-single-payer}
The only canonical packet generated by an obstruction datum \(a\) is
\begin{equation}
 \mathcal J a
 :=\bigl(i_*(a),Bi_*(a)\bigr).
\label{eq:v3v5-J}
\end{equation}
It satisfies
\begin{align}
 \|Bi_*(a)\|_{Z_\partial}
 &\le \beta\sigma_{\rm phys}^{-1}\|a\|_{\mathcal R},
\label{eq:v3v5-flux-bound}\\
 \|\mathcal J a\|_{G\oplus Z_\partial}^2
 &\le (1+\beta^2)\sigma_{\rm phys}^{-2}
       \|a\|_{\mathcal R}^2.
\label{eq:v3v5-graph-packet-bound}
\end{align}
There is no independent second datum in \(Z_\partial\) produced by the
same obstruction coordinate.
\end{proposition}

\begin{proof}
By definition of \(\beta\),
\[
 \|Bi_*(a)\|_{Z_\partial}
 \le\beta\|G^{1/2}i_*(a)\|_I
 =\beta\|i_*(a)\|_G.
\]
Apply \eqref{eq:v3v5-primitive-bound} and then add the two squared graph
components.  Both components are images of the same \(i_*(a)\); the last
claim is therefore a statement about the domain of \(\mathcal J\), not an
accounting convention.
\end{proof}

\begin{assessment}[Condition before taking the norm]
\label{ass:v3v5-condition-first}
The identity \(Li_*(a)=a\) and the flux identity in
\eqref{eq:v3v5-J} are imposed before any triangle inequality.  Introducing
an unrelated \(z\in Z_\partial\) and estimating \(i_*(a)\) and \(z\)
separately describes a larger problem.  It cannot be used to claim a bound
for the physical one-coordinate packet unless \(z=Bi_*(a)\) is restored.
\end{assessment}

\subsection{Insertion into a common boundary energy estimate}
\label{sec:v3v5-energy-interface}

The following statement isolates exactly what V5 contributes to the earlier
Einstein--matter energy architecture.  Let \(U_T\) be an evolution norm,
\(D_T\) a bulk data norm, and suppose an already established estimate has
the form
\begin{equation}
 \|u\|_{U_T}
 \le C_T\bigl(\|d_{\rm bulk}\|_{D_T}
       +c_\partial\|g_\partial\|_{Z_\partial}\bigr).
\label{eq:v3v5-base-energy}
\end{equation}
Let \(R_\partial:Z_\partial\to Z_\partial\) be the admissible boundary
response operator and let \(d_\partial\in Z_\partial\) denote all boundary
data independent of the obstruction coordinate.

\begin{theorem}[One-use inflow insertion]
\label{thm:v3v5-energy-insertion}
Set
\begin{equation}
 g_\partial
 :=R_\partial\bigl(d_\partial+Bi_*(a)\bigr).
\label{eq:v3v5-gboundary}
\end{equation}
Then
\begin{equation}
 \|u\|_{U_T}
 \le C_T\left(
 \|d_{\rm bulk}\|_{D_T}
 +c_\partial\|R_\partial\|
 \left[\|d_\partial\|_{Z_\partial}
 +\beta\sigma_{\rm phys}^{-1}\|a\|_{\mathcal R}
 \right]\right).
\label{eq:v3v5-energy-insertion}
\end{equation}
The obstruction datum occurs exactly once on the right-hand side.
\end{theorem}

\begin{proof}
Substitute \eqref{eq:v3v5-gboundary} into
\eqref{eq:v3v5-base-energy}, use the operator norm of \(R_\partial\), and
then apply \eqref{eq:v3v5-flux-bound}.  No additional copy of \(a\) or of
\(i_*(a)\) is introduced.
\end{proof}

\begin{firewall}[This is an interface theorem]
\label{fw:v3v5-interface}
Theorem~\ref{thm:v3v5-energy-insertion} does not construct the evolution
estimate \eqref{eq:v3v5-base-energy}.  It proves that once the common
boundary estimate is available, the finite-dimensional inflow block enters
it with the exact one-use cost shown in
\eqref{eq:v3v5-energy-insertion}.
\end{firewall}

\subsection{Regulator stability in the physical metric}
\label{sec:v3v5-regulator}

For each regulator index \(n\), let \(I\) and \(\mathcal R\) have fixed
dimension, let \(G_n>0\), and let \(L_n:I\to\mathcal R\) be onto.  Put
\begin{equation}
 C_n:=L_nG_n^{-1}L_n^*,
 \qquad
 i_{*,n}:=G_n^{-1}L_n^*C_n^{-1}.
\label{eq:v3v5-regulator-defs}
\end{equation}

\begin{theorem}[Uniform physical realization criterion]
\label{thm:v3v5-regulator}
Assume
\begin{equation}
 g_-I\le G_n\le g_+I,
 \qquad C_n\ge\sigma_0^2I
\label{eq:v3v5-uniform-hyp}
\end{equation}
for constants \(g_->0\), \(g_+<\infty\), and \(\sigma_0>0\) independent of
\(n\).  Then
\begin{equation}
 \|i_{*,n}a\|_{G_n}\le\sigma_0^{-1}\|a\|,
 \qquad
 \|B_ni_{*,n}a\|
 \le\beta_0\sigma_0^{-1}\|a\|
\label{eq:v3v5-uniform-bounds}
\end{equation}
whenever \(\|B_nG_n^{-1/2}\|\le\beta_0\).

If in addition \(G_n\to G\), \(L_n\to L\), and \(B_n\to B\) in operator
norm, then \(C_n\to C=LG^{-1}L^*\), \(C_n^{-1}\to C^{-1}\), and
\begin{equation}
 i_{*,n}\longrightarrow G^{-1}L^*C^{-1}
\label{eq:v3v5-minimizer-convergence}
\end{equation}
in operator norm.
\end{theorem}

\begin{proof}
The first assertion is Theorem~\ref{thm:v3v5-minimizer} and
Proposition~\ref{prop:v3v5-single-payer}, uniformly in \(n\).  The lower
bound on \(G_n\) and norm convergence give \(G_n^{-1}\to G^{-1}\) by the
resolvent identity
\[
 G_n^{-1}-G^{-1}=G_n^{-1}(G-G_n)G^{-1}.
\]
Therefore \(C_n\to C\).  The uniform lower bound on \(C_n\) passes to the
limit and another resolvent identity gives \(C_n^{-1}\to C^{-1}\).
Substitution in \eqref{eq:v3v5-regulator-defs} proves
\eqref{eq:v3v5-minimizer-convergence}.  The convergence of the flux packet
then follows after multiplication by \(B_n\).
\end{proof}

\subsection{Why the V4 algebraic isometry is insufficient}
\label{sec:v3v5-counterexample}

\begin{counterexample}[Fixed algebraic gap, collapsing physical margin]
\label{cex:v3v5-rescaling}
Take \(I=\mathcal R=\mathbb R\), \(L_n=1\), and \(G_n=n^2\).  In the
unweighted algebraic norms, \(L_n\) is an isometry for every \(n\).  However,
\begin{equation}
 C_n=n^{-2},
 \qquad \sigma_{{\rm phys},n}=n^{-1},
 \qquad i_{*,n}(a)=a,
 \qquad \|i_{*,n}(a)\|_{G_n}=n|a|.
\label{eq:v3v5-rescaling-counterexample}
\end{equation}
Thus a uniform abstract cochain gap does not give a uniform kinetic cost.
The missing hypothesis is precisely the physical Schur gap in
\eqref{eq:v3v5-uniform-hyp}.
\end{counterexample}

\begin{counterexample}[Cheap kinetic directions can destroy flux control]
\label{cex:v3v5-cheap-direction}
Take \(I=\mathbb R^2\), \(\mathcal R=\mathbb R\),
\(L(x,y)=x\), \(G_n=\operatorname{diag}(1,n^{-2})\), and
\(B(x,y)=y\).  The least-energy constrained primitive remains
\(i_{*,n}(a)=(a,0)\), so its flux vanishes.  But the operator norms
\(\|BG_n^{-1/2}\|=n\) diverge: arbitrary physical packets with bounded
kinetic norm may carry unbounded flux in the cheap \(y\)-direction.  Hence
the primitive estimate and the full boundary trace estimate are distinct
uniform requirements.
\end{counterexample}

\subsection{Locality and quantum firewalls}
\label{sec:v3v5-firewalls}

\begin{firewall}[Finite-dimensional Gram data do not prove locality]
\label{fw:v3v5-locality}
A positive matrix \(G\) can always be written algebraically.  V5 does not
show that it is the harmonic restriction of a local, causal, gauge-covariant
boundary kinetic operator.  That requires a field complex, a local action,
boundary conditions, and a well-posed propagator.
\end{firewall}

\begin{firewall}[Classical minimization is not determinant cancellation]
\label{fw:v3v5-determinant}
The Schur complement \(C^{-1}\) is a classical quadratic cost.  Integrating
a quantum field also produces a determinant and possibly an eta phase.
Neither is controlled by Theorem~\ref{thm:v3v5-minimizer}.  The determinant
line and global anomaly gates from the earlier programme remain open.
\end{firewall}

\begin{firewall}[Stress-energy must use the same kinetic metric]
\label{fw:v3v5-stress}
If the inflow field couples to Einstein evolution, its stress tensor must be
derived from the same action that induces \(G\).  Assigning one metric to the
BRST primitive and another to its stress-energy would recreate the duplicate
payer excluded by Proposition~\ref{prop:v3v5-single-payer}.
\end{firewall}

\subsection{V5 conclusion and V6 target}
\label{sec:v3v5-conclusion}

V5 replaces the arbitrary norm on the V4 shifted generator by the kinetic
metric \(G\).  The exact primitive is
\(G^{-1}L^*(LG^{-1}L^*)^{-1}\), its energy is the inverse physical Schur
form, and its normal flux is an image of the same primitive.  A regulator
family is uniformly admissible precisely when the physical Schur gap and
the graph trace norm stay uniform.  The scalar counterexample proves that
the V4 algebraic isometry alone cannot supply this conclusion.

V6 will go one step closer to a physical boundary theory.  It will construct
a local linear boundary kinetic complex whose harmonic reduction gives the
matrix \(G\), prove its energy identity and constraint propagation, and
identify the conditions under which its stress-energy flux can be inserted
into the shared Einstein boundary ledger without a second copy of the
inflow resource.

\clearpage
\section{Third Volume, Version V6: A Local Boundary Kinetic Realization and Its One-Port Energy Law}
\label{sec:v3v6}

\subsection{From the harmonic inflow matrix to a local field}
\label{sec:v3v6-purpose}

V5 identified the exact finite-dimensional physical metric \(G\), but did
not produce it from a local action.  This version gives a linear realization
on a fixed globally hyperbolic boundary cylinder.  The realization is modest
but exact: every positive internal matrix occurs as the harmonic-mode metric
of a local massive wave field, the cochain differential commutes with the
propagator, and a single port coupling exchanges energy with the bulk without
duplicating the boundary resource.

Let
\[
 \mathcal N=\mathbb R_t\times\Sigma
\]
where \((\Sigma,h)\) is a compact connected Riemannian manifold without
boundary.  Normalize \(\operatorname{vol}_h(\Sigma)=1\).  Let \(I\) be the
finite-dimensional inflow space and let \(G>0\) be the physical metric from
V5.  An \(I\)-valued boundary field is denoted by
\(\phi:\mathcal N\to I\).

\subsection{The local action and field equation}
\label{sec:v3v6-action}

Fix \(m>0\).  With the product Lorentzian metric
\(\gamma=-dt^2+h\), define
\begin{equation}
 S_\partial[\phi]
 :=\frac12\int_{\mathbb R}\int_\Sigma
 \left(
  \langle G\partial_t\phi,\partial_t\phi\rangle
 -\langle G\nabla\phi,\nabla\phi\rangle_h
 -m^2\langle G\phi,\phi\rangle
 \right)\,d\mu_h\,dt .
\label{eq:v3v6-action}
\end{equation}
Here \(G\) is a constant metric on the internal fibre and the action contains
no harmonic projector.

\begin{proposition}[Euler--Lagrange equation and locality]
\label{prop:v3v6-euler}
For compactly supported variations, the Euler--Lagrange equation after
adding the local source term
\(\int_{\mathcal N}\langle Gf,\phi\rangle\,d\mu_\gamma\) is
\begin{equation}
 K_m\phi:=\partial_t^2\phi-\Delta_h\phi+m^2\phi=f.
\label{eq:v3v6-KG}
\end{equation}
The differential operator \(K_m\) is normally hyperbolic and local.  The
matrix \(G\) changes the fibre metric and the energy, but does not change the
principal characteristic cone.
\end{proposition}

\begin{proof}
Integrate the time and spatial derivative terms in
\eqref{eq:v3v6-action} by parts.  Since \(G\) is constant and invertible, the
variational equation is
\[
 G(\partial_t^2-\Delta_h+m^2)\phi=Gf,
\]
which is equivalent to \eqref{eq:v3v6-KG}.  Its principal symbol is
\((-\tau^2+|\xi|_h^2)I_I\), so normal hyperbolicity and locality follow.
\end{proof}

\begin{theorem}[Well-posed boundary evolution]
\label{thm:v3v6-wellposed}
For \(s\in\mathbb R\), initial data
\[
 (\phi_0,\phi_1)\in H^{s+1}(\Sigma;I)\times H^s(\Sigma;I)
\]
and \(f\in L^1_{\rm loc}(\mathbb R;H^s(\Sigma;I))\) determine a unique
solution
\[
 \phi\in C(\mathbb R;H^{s+1})\cap C^1(\mathbb R;H^s)
\]
of \eqref{eq:v3v6-KG}.  On every bounded time interval the solution depends
continuously on the data and source.
\end{theorem}

\begin{proof}
Let \(\{e_j\}_{j\ge0}\) be an orthonormal eigenbasis of
\(-\Delta_h\), with eigenvalues
\(0=\lambda_0<\lambda_1\le\cdots\).  Expanding
\(\phi=\sum_j\phi_j(t)e_j\) reduces the equation to the finite-dimensional
oscillators
\[
 \ddot\phi_j+(\lambda_j+m^2)\phi_j=f_j.
\]
Duhamel's formula gives existence and uniqueness for each coefficient.
Multiplication by \((1+\lambda_j)^{s/2}\), the uniform oscillator energy
estimate, and Parseval's identity give the stated continuity in Sobolev
spaces.  Truncations converge in those norms, which constructs the solution.
\end{proof}

\subsection{Exact energy identity}
\label{sec:v3v6-energy}

Define
\begin{equation}
 E_\partial[\phi](t)
 :=\frac12\int_\Sigma
 \left(
  \|\partial_t\phi\|_G^2
 +\|\nabla\phi\|_{G,h}^2
 +m^2\|\phi\|_G^2
 \right)d\mu_h .
\label{eq:v3v6-energy}
\end{equation}

\begin{theorem}[Boundary work identity]
\label{thm:v3v6-work}
Every sufficiently regular solution of \eqref{eq:v3v6-KG} satisfies
\begin{equation}
 \frac{d}{dt}E_\partial[\phi](t)
 =\int_\Sigma\langle Gf,\partial_t\phi\rangle\,d\mu_h .
\label{eq:v3v6-work}
\end{equation}
Consequently,
\begin{equation}
 E_\partial(t)^{1/2}
 \le E_\partial(0)^{1/2}
 \frac{1}{\sqrt2}\int_0^t
 \|G^{1/2}f(\tau)\|_{L^2(\Sigma)}\,d\tau .
\label{eq:v3v6-energy-estimate}
\end{equation}
\end{theorem}

\begin{proof}
Differentiate \eqref{eq:v3v6-energy}.  Integration by parts on the closed
manifold \(\Sigma\) gives
\[
 \frac{dE_\partial}{dt}
 =\int_\Sigma
 \langle G(\partial_t^2-\Delta_h+m^2)\phi,
         \partial_t\phi\rangle\,d\mu_h,
\]
which is \eqref{eq:v3v6-work}.  Cauchy--Schwarz and
\(\|G^{1/2}\partial_t\phi\|_{L^2}\le\sqrt{2E_\partial}\) yield
\[
 \frac{d}{dt}E_\partial^{1/2}
 \le\frac1{\sqrt2}\|G^{1/2}f\|_{L^2}
\]
at positive-energy times.  Regularization by
\((E_\partial+\varepsilon)^{1/2}\) and passage to
\(\varepsilon\downarrow0\) proves \eqref{eq:v3v6-energy-estimate}.
\end{proof}

\subsection{Harmonic reduction recovers the V5 metric}
\label{sec:v3v6-harmonic}

Since \(\Sigma\) is connected and has unit volume, its scalar harmonic
subspace is spanned by the constant function \(1\).  Define
\begin{equation}
 \bar\phi(t):=\int_\Sigma\phi(t,x)\,d\mu_h(x),
 \qquad
 \phi^\perp:=\phi-\bar\phi .
\label{eq:v3v6-average}
\end{equation}

\begin{theorem}[Orthogonal zero-mode splitting]
\label{thm:v3v6-zero-mode}
The boundary energy splits exactly as
\begin{equation}
 E_\partial[\phi]
 =\frac12\left(
   \|\dot{\bar\phi}\|_G^2+m^2\|\bar\phi\|_G^2
  \right)
 +E_\partial[\phi^\perp].
\label{eq:v3v6-energy-split}
\end{equation}
The averaged field satisfies
\begin{equation}
 \ddot{\bar\phi}+m^2\bar\phi=\bar f,
 \qquad
 \bar f:=\int_\Sigma f\,d\mu_h .
\label{eq:v3v6-zero-ode}
\end{equation}
At a static time slice, with \(m=1\) and
\(\dot{\bar\phi}=0\), the zero-mode energy is exactly
\(\frac12\|\bar\phi\|_G^2\), the V5 kinetic cost.
\end{theorem}

\begin{proof}
The constant and mean-zero functions are orthogonal in \(L^2(\Sigma)\).
Their time derivatives remain orthogonal, the gradient of the constant
part vanishes, and the mass term is also orthogonal.  This proves
\eqref{eq:v3v6-energy-split}.  Integrating \eqref{eq:v3v6-KG} over
\(\Sigma\) eliminates the Laplacian term and proves
\eqref{eq:v3v6-zero-ode}.  The last assertion is direct.
\end{proof}

\begin{corollary}[Least-energy cohomological zero mode]
\label{cor:v3v6-least-zero}
Let \(L:I\to\mathcal R\) be onto and prescribe
\(L\bar\phi=a\) at a static slice with \(m=1\).  Among all zero modes
satisfying this constraint, the unique least-energy one is
\begin{equation}
 \bar\phi=i_*(a)
 =G^{-1}L^*(LG^{-1}L^*)^{-1}a,
\label{eq:v3v6-zero-minimizer}
\end{equation}
and its energy is
\(\frac12\langle(LG^{-1}L^*)^{-1}a,a\rangle\).
\end{corollary}

\begin{proof}
Apply the V5 physical Moore--Penrose theorem to the zero-mode term in
\eqref{eq:v3v6-energy-split}.  Every nonzero spatial mode adds a
nonnegative orthogonal energy and cannot lower the constrained minimum.
\end{proof}

\subsection{Propagation of the Ward constraint}
\label{sec:v3v6-constraint}

Let \(a:\mathbb R\to\mathcal R\) be a prescribed obstruction history and
define the harmonic Ward defect
\begin{equation}
 c(t):=L\bar\phi(t)-a(t).
\label{eq:v3v6-defect}
\end{equation}

\begin{theorem}[Exact constraint propagation]
\label{thm:v3v6-constraint}
Assume the obstruction history obeys the compatibility equation
\begin{equation}
 \ddot a+m^2a=L\bar f .
\label{eq:v3v6-a-compatibility}
\end{equation}
Then
\begin{equation}
 \ddot c+m^2c=0.
\label{eq:v3v6-defect-equation}
\end{equation}
If \(c(t_0)=\dot c(t_0)=0\), then \(c(t)=0\) for all \(t\).
\end{theorem}

\begin{proof}
Apply \(L\) to \eqref{eq:v3v6-zero-ode} and subtract
\eqref{eq:v3v6-a-compatibility}.  This gives
\eqref{eq:v3v6-defect-equation}.  Uniqueness for the
finite-dimensional oscillator with zero Cauchy data gives \(c\equiv0\).
\end{proof}

\begin{counterexample}[One compatibility condition is not enough]
\label{cex:v3v6-constraint-velocity}
Let \(m=1\), \(f=0\), \(a=0\), and choose initial data with
\(L\bar\phi(0)=0\) but \(L\dot{\bar\phi}(0)\ne0\).  Then
\[
 c(t)=\sin(t)L\dot{\bar\phi}(0),
\]
so the Ward constraint immediately fails.  Both the position and velocity
constraints are necessary for hyperbolic propagation.
\end{counterexample}

\subsection{The cochain-valued propagator}
\label{sec:v3v6-cochain}

Let \((I^\bullet,q_I)\) be any finite-dimensional cochain complex, including
the two-term inflow block of V4.  On
\(C^\infty(\mathcal N;I^\bullet)\), let \(K_m\) act on the spacetime
variable and \(q_I\) act on the internal variable.

\begin{proposition}[Kinetic operator commutes with the Ward differential]
\label{prop:v3v6-commutation}
One has
\begin{equation}
 [K_m,q_I]=0.
\label{eq:v3v6-commutation}
\end{equation}
Hence the retarded and advanced Green operators, wherever defined, commute
with \(q_I\).  In particular, cochain-closed Cauchy data and source produce
a cochain-closed solution.
\end{proposition}

\begin{proof}
The coefficients of \(q_I\) are constant in spacetime, while \(K_m\)
differentiates only the spacetime variable.  This proves
\eqref{eq:v3v6-commutation}.  If \(E^\pm\) is a Green operator, then
\(q_IE^\pm-E^\pm q_I\) solves the homogeneous equation and has vanishing
retarded or advanced support; uniqueness gives zero.  The Cauchy statement
follows in the same way by applying \(q_I\) to the evolution.
\end{proof}

\begin{firewall}[Commutation is weaker than quantum BRST invariance]
\label{fw:v3v6-BRST}
Equation \eqref{eq:v3v6-commutation} proves compatibility of the linear
equations with the cochain differential.  A gauge-fixed quantum action also
requires the correct graded pairing, ghost statistics, measure, and
Slavnov--Taylor identity.  Those conclusions do not follow from operator
commutation alone.
\end{firewall}

\subsection{One-port coupling to the bulk Einstein--matter system}
\label{sec:v3v6-port}

Let \(W\to\Sigma\) be a finite-rank metric bundle carrying the bulk boundary
port variable \(z(t)\in L^2(\Sigma;W)\).  Let
\(\mathsf B:I\to W_x\) be a smooth zero-order bundle map, viewed as
\[
 \mathsf B:L^2(\Sigma;I)\longrightarrow L^2(\Sigma;W).
\]
Its adjoint \(\mathsf B^\dagger\) is defined by
\begin{equation}
 \int_\Sigma\langle \mathsf Bv,z\rangle_W\,d\mu_h
 =
 \int_\Sigma\langle Gv,G^{-1}\mathsf B^\dagger z\rangle_I\,d\mu_h .
\label{eq:v3v6-port-adjoint}
\end{equation}
Thus \(G^{-1}\mathsf B^\dagger z\) is the force dual to the observed port
\(\mathsf B\partial_t\phi\).

\begin{theorem}[Exact cancellation of exchanged power]
\label{thm:v3v6-port-cancellation}
Suppose the boundary field solves
\begin{equation}
 K_m\phi=f_{\rm ext}-G^{-1}\mathsf B^\dagger z
\label{eq:v3v6-coupled-boundary}
\end{equation}
and the bulk energy satisfies
\begin{equation}
 \frac{d}{dt}E_{\rm bulk}
 =P_{\rm ext}
 +\int_\Sigma\langle z,\mathsf B\partial_t\phi\rangle_W\,d\mu_h .
\label{eq:v3v6-bulk-port}
\end{equation}
Then the combined energy obeys
\begin{equation}
 \frac{d}{dt}(E_{\rm bulk}+E_\partial)
 =P_{\rm ext}
 +\int_\Sigma
  \langle Gf_{\rm ext},\partial_t\phi\rangle_I\,d\mu_h .
\label{eq:v3v6-total-energy}
\end{equation}
The internal port power occurs once with opposite signs and cancels exactly.
\end{theorem}

\begin{proof}
Insert \eqref{eq:v3v6-coupled-boundary} in the boundary work identity:
\[
 \frac{dE_\partial}{dt}
 =\int_\Sigma\langle Gf_{\rm ext},\partial_t\phi\rangle\,d\mu_h
 -\int_\Sigma\langle\mathsf B^\dagger z,\partial_t\phi\rangle\,d\mu_h .
\]
By the adjoint identity, the second integral is
\(-\int_\Sigma\langle z,\mathsf B\partial_t\phi\rangle_Wd\mu_h\).
Adding \eqref{eq:v3v6-bulk-port} proves
\eqref{eq:v3v6-total-energy}.
\end{proof}

\begin{assessment}[The companion single-payer rule is now an identity]
\label{ass:v3v6-single-payer}
The boundary observable and the boundary force in
\eqref{eq:v3v6-coupled-boundary} are adjoint appearances of the same port.
Estimating the two exchanged-power integrals separately before adding the
energy identities would lose their cancellation and count one internal
transfer twice.  The correct operation is to form
\eqref{eq:v3v6-total-energy} first and estimate only the external powers.
\end{assessment}

\subsection{Stress tensor and local balance}
\label{sec:v3v6-stress}

For the uncoupled field, define
\begin{equation}
 T_{ab}^\partial
 :=\langle G\nabla_a\phi,\nabla_b\phi\rangle
 -\frac12\gamma_{ab}
 \left(
  \langle G\nabla^c\phi,\nabla_c\phi\rangle
  +m^2\langle G\phi,\phi\rangle
 \right).
\label{eq:v3v6-stress}
\end{equation}

\begin{proposition}[Local stress balance]
\label{prop:v3v6-stress}
On a fixed boundary spacetime,
\begin{equation}
 \nabla^aT_{ab}^\partial
 =-\langle GK_m\phi,\nabla_b\phi\rangle .
\label{eq:v3v6-stress-divergence}
\end{equation}
Thus solutions with zero source have divergence-free boundary stress, while
the port source contributes exactly its local work density.
\end{proposition}

\begin{proof}
Differentiate \eqref{eq:v3v6-stress}.  Metric compatibility and symmetry of
second derivatives on the scalar spacetime index cancel the two quadratic
gradient terms.  The mass derivative combines with the remaining term to
give \(\langle G(\nabla^a\nabla_a-m^2)\phi,\nabla_b\phi\rangle\).
For \(\gamma=-dt^2+h\), one has
\(\nabla^a\nabla_a-m^2=-K_m\).  This proves
\eqref{eq:v3v6-stress-divergence}.
\end{proof}

\begin{firewall}[A boundary stress is not yet an Einstein junction law]
\label{fw:v3v6-junction}
The tensor \eqref{eq:v3v6-stress} is intrinsic to \(\mathcal N\).  Coupling
it to a dynamical bulk metric requires a specified variational boundary
term and the corresponding Israel-type or constraint-compatible junction
condition.  V6 proves the local stress balance and abstract port
cancellation; it does not derive that gravitational junction law.
\end{firewall}

\subsection{Two sharp limitations}
\label{sec:v3v6-limitations}

\begin{counterexample}[The massless zero mode has no static coercivity]
\label{cex:v3v6-massless}
If \(m=0\), a constant time-independent field has zero energy for every
value of \(\bar\phi\).  Therefore a massless scalar realization cannot
recover the positive V5 static cost
\(\frac12\|\bar\phi\|_G^2\).  One needs a mass, a potential, a first-order
constraint action, or a different coercive mechanism.
\end{counterexample}

\begin{assessment}[Time-dependent geometry adds controlled work]
\label{ass:v3v6-time-dependent}
If \(h\) or \(G\) depends on time, differentiating the energy produces terms
containing \(\partial_th\), \(\partial_tG\), and the derivative of the
volume form.  Exact conservation is then replaced by a Gronwall estimate.
Uniform control requires bounds for these coefficients; freezing them by
notation would conceal a genuine Einstein backreaction cost.
\end{assessment}

\subsection{V6 conclusion and V7 target}
\label{sec:v3v6-conclusion}

V6 supplies a local normally hyperbolic field whose harmonic static sector
has exactly the V5 physical metric.  Its Ward constraint propagates from two
Cauchy conditions, its cochain differential commutes with the Green
operators, and its one-port coupling cancels exchanged bulk--boundary power
before norms are taken.  This is the first local dynamical realization of
the algebraic inflow coordinate in the restarted volume.

The next bottleneck lies upstream of a full Standard Model--Einstein
continuum: the boundary cylinder was fixed.  V7 will allow a time-dependent
boundary metric, derive the precise deformation terms in the energy
identity, and formulate a coercive continuation margin that can be compared
with the lapse, shift, and second-fundamental-form bounds in the Einstein
sector.  The quantum BRST and determinant-line firewalls remain open.
\clearpage
\section{Third Volume, Version V7: Moving-Boundary Energy and the Harmonic-Mode Mixing Obstruction}
\label{sec:v3v7}

\subsection{Geometric setup}
\label{sec:v3v7-setup}

V6 used a fixed boundary metric.  An Einstein coupling makes the induced
metric evolve, so its volume form, gradient norm, and wave operator all vary.
This version derives those terms without freezing them.  It also tests the
finite-dimensional harmonic inflow sector against the moving geometry and
finds the precise condition under which that sector is invariant.

Let \(\Sigma\) be compact and without boundary, and let \(h(t)\) be a
\(C^1\) family of Riemannian metrics.  Write
\begin{equation}
 k_{ij}:=\frac12\dot h_{ij},
 \qquad
 \theta:=\operatorname{tr}_h k
 =\frac12h^{ij}\dot h_{ij}.
\label{eq:v3v7-k-theta}
\end{equation}
Then
\begin{equation}
 \partial_t d\mu_{h(t)}=\theta\,d\mu_{h(t)},
 \qquad
 \dot h^{ij}=-2k^{ij}.
\label{eq:v3v7-variation-identities}
\end{equation}
Let \(G(t)>0\) be a \(C^1\) family of internal metrics on the
finite-dimensional space \(I\), constant over each spatial slice.

\subsection{The variational moving-boundary operator}
\label{sec:v3v7-operator}

Consider the time-dependent action
\begin{equation}
 S[\phi]
 =\frac12\int
 \left(
  \langle G\partial_t\phi,\partial_t\phi\rangle
 -h^{ij}\langle G\partial_i\phi,\partial_j\phi\rangle
 -m^2\langle G\phi,\phi\rangle
 \right)d\mu_h\,dt .
\label{eq:v3v7-action}
\end{equation}
Define the logarithmic internal velocity
\begin{equation}
 A_G:=G^{-1}\dot G.
\label{eq:v3v7-AG}
\end{equation}

\begin{proposition}[Euler operator on the moving cylinder]
\label{prop:v3v7-euler}
After adding the source term
\(\int\langle Gf,\phi\rangle d\mu_hdt\), the Euler--Lagrange equation is
\begin{equation}
 K_{h,G}\phi
 :=\partial_t^2\phi+(\theta I+A_G)\partial_t\phi
   -\Delta_{h(t)}\phi+m^2\phi=f.
\label{eq:v3v7-K}
\end{equation}
Equivalently, in divergence form,
\begin{equation}
 \partial_t(d\mu_h\,G\partial_t\phi)
 -d\mu_h\,G\Delta_h\phi
 +d\mu_h\,m^2G\phi
 =d\mu_h\,Gf.
\label{eq:v3v7-divergence}
\end{equation}
\end{proposition}

\begin{proof}
The time integration by parts differentiates both \(G\) and \(d\mu_h\);
their logarithmic derivatives are \(A_G\) and \(\theta\).  Spatial
integration by parts produces \(-\Delta_h\), because \(G(t)\) is spatially
constant.  Division by the invertible matrix \(G\) gives
\eqref{eq:v3v7-K}; retaining the density gives
\eqref{eq:v3v7-divergence}.
\end{proof}

\subsection{Exact deformation-energy identity}
\label{sec:v3v7-energy}

Define
\begin{equation}
 E(t):=\frac12\int_\Sigma
 \left(
  \|\partial_t\phi\|_G^2
 +h^{ij}\langle G\partial_i\phi,\partial_j\phi\rangle
 +m^2\|\phi\|_G^2
 \right)d\mu_h .
\label{eq:v3v7-energy}
\end{equation}
For brevity put
\[
 |\nabla\phi|_{G,h}^2
 :=h^{ij}\langle G\partial_i\phi,\partial_j\phi\rangle .
\]

\begin{theorem}[Moving-boundary work identity]
\label{thm:v3v7-energy-identity}
Every regular solution of \eqref{eq:v3v7-K} satisfies
\begin{align}
 \frac{dE}{dt}
 &=\int_\Sigma\langle Gf,\partial_t\phi\rangle\,d\mu_h
   +\mathcal D_{h,G}[\phi],
\label{eq:v3v7-energy-identity}\\
 \mathcal D_{h,G}[\phi]
 &:=\frac12\int_\Sigma
 \Bigl[
 -\langle\dot G\partial_t\phi,\partial_t\phi\rangle
 +h^{ij}\langle\dot G\partial_i\phi,\partial_j\phi\rangle
 +m^2\langle\dot G\phi,\phi\rangle
\notag\\
 &\hspace{21mm}
 +\theta\bigl(
   -\|\partial_t\phi\|_G^2
   +|\nabla\phi|_{G,h}^2
   +m^2\|\phi\|_G^2
  \bigr)
 +\dot h^{ij}\langle G\partial_i\phi,\partial_j\phi\rangle
 \Bigr]d\mu_h .
\label{eq:v3v7-deformation}
\end{align}
\end{theorem}

\begin{proof}
Differentiate \eqref{eq:v3v7-energy}.  The derivative of the measure is
\(\theta d\mu_h\); the spatial quadratic form also produces the
\(\dot h^{ij}\) term.  Integrating
\(\langle G\nabla\phi,\nabla\partial_t\phi\rangle\) by parts gives
\[
 \int_\Sigma
 \langle G(\partial_t^2-\Delta_h+m^2)\phi,
         \partial_t\phi\rangle\,d\mu_h .
\]
By \eqref{eq:v3v7-K}, this equals the source work minus
\[
 \int_\Sigma
 \left(
  \theta\|\partial_t\phi\|_G^2
  +\langle\dot G\partial_t\phi,\partial_t\phi\rangle
 \right)d\mu_h .
\]
Combining these two full kinetic corrections with the half corrections
arising from differentiating the quadratic energy gives exactly
\eqref{eq:v3v7-deformation}.
\end{proof}

\subsection{A quantitative continuation margin}
\label{sec:v3v7-continuation}

Define the pointwise deformation sizes
\begin{align}
 a_G(t)&:=
 \|G^{-1/2}\dot G\,G^{-1/2}\|_{\operatorname{op}},
\label{eq:v3v7-aG}\\
 a_h(t)&:=
 \|h^{-1/2}\dot h\,h^{-1/2}\|_{L^\infty(\Sigma;\operatorname{op})},
\label{eq:v3v7-ah}\\
 \Xi(t)&:=a_G(t)+\|\theta(t)\|_{L^\infty(\Sigma)}+a_h(t).
\label{eq:v3v7-Xi}
\end{align}

\begin{lemma}[Deformation form bound]
\label{lem:v3v7-deformation-bound}
One has
\begin{equation}
 |\mathcal D_{h,G}[\phi]|\le \Xi(t)E(t).
\label{eq:v3v7-deformation-bound}
\end{equation}
\end{lemma}

\begin{proof}
The three \(\dot G\) terms are bounded in absolute value by \(a_G\) times
the corresponding three \(G\)-energy densities.  The three \(\theta\)
terms are bounded by \(\|\theta\|_\infty\) times those densities.  Finally,
the \(\dot h^{ij}\) term is bounded by \(a_h|\nabla\phi|_{G,h}^2\).
The common factor \(1/2\) agrees with the definition of \(E\), which proves
\eqref{eq:v3v7-deformation-bound}.
\end{proof}

\begin{theorem}[Moving-boundary energy estimate]
\label{thm:v3v7-energy-estimate}
Let
\[
 F(t):=\|G(t)^{1/2}f(t)\|_{L^2(\Sigma,h(t))},
 \qquad
 Q(t):=\frac12\int_0^t\Xi(\tau)\,d\tau .
\]
Then, for \(t\ge0\),
\begin{equation}
 E(t)^{1/2}
 \le e^{Q(t)}
 \left[
  E(0)^{1/2}
  \frac1{\sqrt2}\int_0^t e^{-Q(\tau)}F(\tau)\,d\tau
 \right].
\label{eq:v3v7-energy-estimate}
\end{equation}
\end{theorem}

\begin{proof}
The work term is at most \(F\sqrt{2E}\).  By
Lemma~\ref{lem:v3v7-deformation-bound},
\[
 E'\le \Xi E+\sqrt{2E}\,F.
\]
Apply this first to \((E+\varepsilon)^{1/2}\), use the scalar integrating
factor \(e^{-Q}\), and let \(\varepsilon\downarrow0\).
\end{proof}

\begin{corollary}[Finite-time continuation criterion]
\label{cor:v3v7-continuation}
Suppose on \([0,T)\) that
\begin{equation}
 g_-I\le G(t)\le g_+I,
 \qquad
 c_-h(0)\le h(t)\le c_+h(0),
 \qquad
 m\ge m_0>0,
\label{eq:v3v7-coercivity}
\end{equation}
and
\begin{equation}
 \int_0^T\Xi(t)\,dt<\infty,
 \qquad
 \int_0^T F(t)\,dt<\infty.
\label{eq:v3v7-integrability}
\end{equation}
Then \(E(t)\) remains finite and uniformly controls
\[
 \|\partial_t\phi(t)\|_{L^2}^2
 +\|\phi(t)\|_{H^1}^2
\]
for \(t<T\), with comparison constants depending only on
\(g_\pm,c_\pm,m_0\).
\end{corollary}

\begin{proof}
Theorem~\ref{thm:v3v7-energy-estimate} bounds \(E\).
The uniform comparisons in \eqref{eq:v3v7-coercivity} identify its three
positive terms with the fixed \(L^2\times H^1\) norm.  The mass lower bound
controls the spatial constant mode.
\end{proof}

\subsection{ADM interpretation of the deformation margin}
\label{sec:v3v7-ADM}

On a timelike boundary foliation, adopt the sign convention
\begin{equation}
 \dot h_{ij}=2N K_{ij}+(\mathcal L_\beta h)_{ij},
\label{eq:v3v7-ADM}
\end{equation}
where \(N\) is the boundary lapse, \(\beta\) the shift, and \(K\) the
second fundamental form of the spatial slices inside the boundary.

\begin{proposition}[Geometric control of \(a_h\)]
\label{prop:v3v7-ADM-bound}
There is a dimension-dependent algebraic constant \(C_d\) such that
\begin{equation}
 a_h+\|\theta\|_\infty
 \le C_d\left(
  \|N\|_\infty\|K\|_\infty
  +\|\nabla\beta\|_\infty
 \right).
\label{eq:v3v7-ADM-bound}
\end{equation}
\end{proposition}

\begin{proof}
In an \(h\)-orthonormal frame,
\(\mathcal L_\beta h=2\operatorname{sym}\nabla\beta\).  Thus the operator
norm of \(\dot h\) is bounded by
\(2\|N\|_\infty\|K\|_\infty+2\|\nabla\beta\|_\infty\).
The trace is at most the dimension times the operator norm.  Absorb these
fixed factors into \(C_d\).
\end{proof}

\begin{assessment}[Einstein interface supplied by V7]
\label{ass:v3v7-Einstein-interface}
The boundary kinetic continuation cost is integrable if the internal metric
velocity, lapse-weighted second fundamental form, and shift gradient are
integrable in time.  These are the exact quantities that a coupled
Einstein continuation theorem must place in its common bootstrap ledger.
V7 does not assume that the Einstein equations already provide those
bounds.
\end{assessment}

\subsection{The moving harmonic projector}
\label{sec:v3v7-harmonic}

Let
\begin{equation}
 V(t):=\operatorname{vol}_{h(t)}(\Sigma),
 \qquad
 P_tu:=\frac1{V(t)}\int_\Sigma u\,d\mu_{h(t)}.
\label{eq:v3v7-Pt}
\end{equation}
For scalar-valued functions on connected \(\Sigma\), \(P_t\) projects onto
the constants, which are the harmonic functions for every \(h(t)\).
Nevertheless, the evolution need not preserve this common subspace.

\begin{lemma}[Derivative of the moving average]
\label{lem:v3v7-average-derivative}
For every regular \(u\),
\begin{equation}
 \partial_t(P_tu)
 =P_t(\partial_tu)+P_t\bigl((\theta-P_t\theta)u\bigr).
\label{eq:v3v7-average-derivative}
\end{equation}
\end{lemma}

\begin{proof}
Differentiate the numerator and denominator in \eqref{eq:v3v7-Pt}, using
\(\dot V=V P_t\theta\).  The quotient rule gives
\eqref{eq:v3v7-average-derivative}.
\end{proof}

\begin{theorem}[Exact constant-mode invariance criterion]
\label{thm:v3v7-invariance}
Assume \(G\) is independent of time.  The operator \(K_{h,G}\) maps every
spatially constant field \(x(t)\in I\) to a spatially constant field if and
only if
\begin{equation}
 \theta(t,x)=\theta_0(t)
\label{eq:v3v7-CMC}
\end{equation}
is spatially constant on each connected slice.
\end{theorem}

\begin{proof}
For a spatially constant field,
\begin{equation}
 K_{h,G}x=\ddot x+\theta(t,\cdot)\dot x+m^2x.
\label{eq:v3v7-constant-action}
\end{equation}
If \(\theta\) is spatially constant, the right-hand side is constant.
Conversely, fix a time \(t_0\) and choose a curve with
\(x(t_0)=\ddot x(t_0)=0\) and arbitrary nonzero
\(\dot x(t_0)=v\).  If \eqref{eq:v3v7-constant-action} is spatially
constant for every such curve, then \(\theta(t_0,\cdot)v\) is constant.
Since \(v\ne0\), \(\theta(t_0,\cdot)\) is constant.  The time was arbitrary.
\end{proof}

\begin{corollary}[Constraint propagation in the CMC boundary sector]
\label{cor:v3v7-CMC-constraint}
Under \eqref{eq:v3v7-CMC}, spatially constant data and source remain
spatially constant.  Their value \(x(t)\) solves
\begin{equation}
 \ddot x+\theta_0(t)\dot x+m^2x=f_0(t).
\label{eq:v3v7-CMC-ODE}
\end{equation}
If \(L:I\to\mathcal R\) is fixed and
\[
 \ddot a+\theta_0\dot a+m^2a=Lf_0,
\]
then \(c=Lx-a\) solves the homogeneous version of
\eqref{eq:v3v7-CMC-ODE}; zero position and velocity defect propagate.
\end{corollary}

\begin{proof}
The first claim follows from uniqueness for the hyperbolic Cauchy problem
and Theorem~\ref{thm:v3v7-invariance}.  Apply \(L\) to
\eqref{eq:v3v7-CMC-ODE} and subtract the equation for \(a\).
\end{proof}

\begin{counterexample}[Spatial expansion creates nonzero modes]
\label{cex:v3v7-mode-mixing}
Take \(\Sigma=S^1\), fix \(G=I\), and choose a metric family with
\(\theta(0,x)=\varepsilon\cos x\), where \(\varepsilon\ne0\).
Let \(f=0\), \(\phi(0,x)=0\), and
\(\partial_t\phi(0,x)=v\ne0\) be spatially constant.  Equation
\eqref{eq:v3v7-K} gives
\begin{equation}
 \partial_t^2\phi(0,x)=-\varepsilon\cos x\,v.
\label{eq:v3v7-mixing-acceleration}
\end{equation}
Thus the solution leaves the constant harmonic sector at second order in
time.  The fixed-cylinder V6 zero-mode equation cannot be used on a general
moving boundary without controlling this off-diagonal mixing.
\end{counterexample}

\subsection{The port cancellation survives metric motion}
\label{sec:v3v7-port}

Let \(\mathsf B(t)\) be a time-dependent zero-order port map and suppose
\[
 K_{h,G}\phi=f_{\rm ext}-G^{-1}\mathsf B^\dagger z.
\]
Assume the bulk energy identity uses the same instantaneous measure:
\[
 \frac{dE_{\rm bulk}}{dt}
 =P_{\rm ext}
 +\int_\Sigma
 \langle z,\mathsf B\partial_t\phi\rangle\,d\mu_h.
\]

\begin{proposition}[Moving one-port identity]
\label{prop:v3v7-port}
The coupled energy satisfies
\begin{equation}
 \frac{d}{dt}(E_{\rm bulk}+E)
 =P_{\rm ext}
 +\int_\Sigma\langle Gf_{\rm ext},\partial_t\phi\rangle\,d\mu_h
 +\mathcal D_{h,G}[\phi].
\label{eq:v3v7-total-port}
\end{equation}
\end{proposition}

\begin{proof}
Use Theorem~\ref{thm:v3v7-energy-identity}.  The boundary source contributes
\(-\int\langle z,\mathsf B\partial_t\phi\rangle d\mu_h\), which cancels the
bulk port term.  The geometric deformation form is internal work caused by
the moving metric and remains in the combined identity.
\end{proof}

\begin{firewall}[Geometric work is not a second port]
\label{fw:v3v7-geometric-work}
The term \(\mathcal D_{h,G}\) is determined by the same boundary field and
the evolving geometry.  It must be estimated in the common Einstein
bootstrap through \(\Xi E\).  Introducing a second independent boundary
energy to pay it would duplicate the kinetic stock rather than control the
deformation of that stock.
\end{firewall}

\subsection{V7 conclusion and V8 target}
\label{sec:v3v7-conclusion}

V7 derives the complete moving-cylinder energy identity and the exponential
continuation factor generated by \(G^{-1}\dot G\), \(\dot h\), and the
volume expansion.  The ADM formula identifies the lapse, shift, and second
fundamental form quantities that must enter a common Einstein estimate.
The port cancellation remains exact.

The harmonic test uncovers the next bottleneck: spatially varying expansion
mixes the finite obstruction block with nonzero boundary modes.  Restricting
to spatially constant expansion closes the block but is too strong for the
full programme.  V8 will go upstream and replace the invariant-block demand
by a controlled moving spectral subspace.  It will construct the Kato
transport of the low-mode projector and estimate the off-diagonal coupling
by the boundary spectral gap, avoiding a circular assumption that the
harmonic sector stays fixed.
\clearpage
\section{Third Volume, Version V8: Kato Transport and Gap-Controlled Leakage from the Inflow Sector}
\label{sec:v3v8}

\subsection{Why a moving spectral sector is needed}
\label{sec:v3v8-purpose}

V7 proved that the constant harmonic sector is not invariant under general
boundary expansion.  Imposing spatially constant expansion would avoid the
problem but would exclude geometries needed by the Einstein programme.
V8 instead follows a finite low-frequency spectral cluster of the
instantaneous boundary elliptic operator.  The cluster is transported by
Kato's connection, while its coupling to the complementary modes is paid by
an explicit inverse spectral gap.

All statements below are operator statements on a fixed Hilbert space
\(\mathcal H\).  The varying spaces \(L^2(\Sigma,d\mu_{h(t)})\) can be put
in this form by the unitary density identification
\[
 u\longmapsto
 \left(\frac{d\mu_{h(t)}}{d\mu_{h(0)}}\right)^{1/2}u.
\]
This identification is used only to compare operators; it does not remove
the deformation work already derived in V7.

\subsection{Separated low spectral cluster}
\label{sec:v3v8-cluster}

Let \(A(t)\) be a nonnegative self-adjoint operator with compact resolvent,
common dense domain \(\mathcal D\), and \(C^1\) dependence in the graph
topology.  Assume that for \(t\in[0,T]\) there are numbers
\(\lambda_-(t)<\lambda_+(t)\) such that
\begin{align}
 \sigma(A(t)|_{\mathcal H_-(t)})&\subset[0,\lambda_-(t)],
\label{eq:v3v8-low-spectrum}\\
 \sigma(A(t)|_{\mathcal H_+(t)})&\subset[\lambda_+(t),\infty),
\label{eq:v3v8-high-spectrum}\\
 \delta(t)&:=\lambda_+(t)-\lambda_-(t)\ge\delta_0>0.
\label{eq:v3v8-gap}
\end{align}
Let \(P(t)\) be the finite-rank spectral projection onto
\(\mathcal H_-(t)\), and put \(Q(t)=I-P(t)\).

\begin{lemma}[Off-diagonal Sylvester inverse]
\label{lem:v3v8-sylvester}
Let \(A_-=A|_{\mathcal H_-}\) and \(A_+=A|_{\mathcal H_+}\) satisfy
\(\sup\sigma(A_-)\le\lambda_-\) and
\(\inf\sigma(A_+)\ge\lambda_+\).  For every bounded
\(Y:\mathcal H_-\to\mathcal H_+\), the equation
\begin{equation}
 A_+X-XA_-=Y
\label{eq:v3v8-sylvester}
\end{equation}
has a unique bounded solution and
\begin{equation}
 \|X\|\le\delta^{-1}\|Y\|,
 \qquad \delta=\lambda_+-\lambda_-.
\label{eq:v3v8-sylvester-bound}
\end{equation}
\end{lemma}

\begin{proof}
After shifting both operators by \(\lambda_-\), define
\[
 X:=\int_0^\infty e^{-s(A_+-\lambda_-)}
       Y e^{s(A_--\lambda_-)}\,ds .
\]
The first semigroup has norm at most \(e^{-s\delta}\), while the second has
norm at most one.  The integral converges and satisfies
\eqref{eq:v3v8-sylvester-bound}.  Differentiating the integrand and
integrating from zero to infinity proves \eqref{eq:v3v8-sylvester}.
The same estimate applied to a homogeneous solution proves uniqueness.
\end{proof}

\subsection{Derivative of the spectral projector}
\label{sec:v3v8-projector}

\begin{theorem}[Gap bound for the moving projector]
\label{thm:v3v8-Pdot}
Under the cluster hypotheses,
\begin{equation}
 \|\dot P(t)\|
 \le \delta(t)^{-1}\|Q(t)\dot A(t)P(t)\|.
\label{eq:v3v8-Pdot}
\end{equation}
Only the off-diagonal variation of \(A\) rotates the low subspace.
\end{theorem}

\begin{proof}
Differentiate \(AP=PA\) and multiply by \(Q\) on the left and \(P\) on the
right.  Since \(QP=0\),
\begin{equation}
 A_+\,Q\dot PP-(Q\dot PP)A_-
 =-Q\dot AP.
\label{eq:v3v8-Pdot-sylvester}
\end{equation}
Lemma~\ref{lem:v3v8-sylvester} bounds
\(\|Q\dot PP\|\) by the right-hand side of
\eqref{eq:v3v8-Pdot}.  Differentiating \(P^2=P\) shows that
\[
 P\dot PP=Q\dot PQ=0,
 \qquad
 \dot P=Q\dot PP+P\dot PQ.
\]
The two off-diagonal blocks are adjoints.  The norm of their self-adjoint
block matrix is the norm of either block, which proves
\eqref{eq:v3v8-Pdot}.
\end{proof}

\begin{definition}[Kato connection]
\label{def:v3v8-Kato}
Define
\begin{equation}
 \Omega(t):=[\dot P(t),P(t)].
\label{eq:v3v8-Omega}
\end{equation}
Then \(\Omega^*=-\Omega\) and
\begin{equation}
 [\Omega,P]=\dot P.
\label{eq:v3v8-Omega-identity}
\end{equation}
\end{definition}

\begin{theorem}[Unitary transport of the low bundle]
\label{thm:v3v8-Kato}
The initial-value problem
\begin{equation}
 \dot U(t)=\Omega(t)U(t),
 \qquad U(0)=I,
\label{eq:v3v8-U}
\end{equation}
has a unitary solution satisfying
\begin{equation}
 P(t)U(t)=U(t)P(0).
\label{eq:v3v8-intertwining}
\end{equation}
Moreover,
\begin{equation}
 \|\Omega(t)\|=\|\dot P(t)\|
 \le\delta(t)^{-1}\|Q\dot AP\|.
\label{eq:v3v8-Omega-bound}
\end{equation}
\end{theorem}

\begin{proof}
Skew-adjointness gives
\(\partial_t(U^*U)=U^*(\Omega^*+\Omega)U=0\), so \(U\) is unitary.
Let \(D(t)=P(t)U(t)-U(t)P(0)\).  Using
\eqref{eq:v3v8-Omega-identity}, direct differentiation gives
\(\dot D=\Omega D\) and \(D(0)=0\).  Hence \(D=0\).
The block form
\[
 \Omega=
 \begin{pmatrix}0&-P\dot PQ\\Q\dot PP&0\end{pmatrix}
\]
shows \(\|\Omega\|=\|Q\dot PP\|=\|\dot P\|\).
Theorem~\ref{thm:v3v8-Pdot} completes the proof.
\end{proof}

\begin{corollary}[A covariant derivative preserving the split]
\label{cor:v3v8-covariant}
Set
\begin{equation}
 D_t:=\partial_t-\Omega(t).
\label{eq:v3v8-Dt}
\end{equation}
If \(P\psi=\psi\), then \(P(D_t\psi)=D_t\psi\).  If
\(Q\psi=\psi\), then \(Q(D_t\psi)=D_t\psi\).
\end{corollary}

\begin{proof}
Differentiate \(P\psi=\psi\) and use
\([\Omega,P]=\dot P\).  This gives
\((I-P)(\partial_t-\Omega)\psi=0\).  The complementary statement follows
by replacing \(P\) with \(Q\).
\end{proof}

\subsection{Gap control of multiplication-induced mixing}
\label{sec:v3v8-commutator}

The moving-boundary damping in V7 contains multiplication by the expansion
\(\theta(t,x)\).  The following estimate avoids treating its low-to-high
part as an unrelated source.

\begin{theorem}[Commutator payment for off-diagonal mixing]
\label{thm:v3v8-commutator}
Let \(M\) be bounded on \(\mathcal H\), preserve \(\mathcal D\), and suppose
\([A,M]P\) is bounded.  Then
\begin{equation}
 \|QMP\|
 \le\delta^{-1}\|Q[A,M]P\|.
\label{eq:v3v8-commutator-bound}
\end{equation}
\end{theorem}

\begin{proof}
Set \(X=QMP:\mathcal H_-\to\mathcal H_+\).  Since \(A\) commutes with its
spectral projections,
\[
 A_+X-XA_-=Q(AM-MA)P=Q[A,M]P.
\]
Apply Lemma~\ref{lem:v3v8-sylvester}.
\end{proof}

\begin{corollary}[Expansion mixing for the scalar Laplacian]
\label{cor:v3v8-laplacian}
Let \(A=-\Delta_h\), let \(M_\theta\) be multiplication by a real
\(\theta\in W^{2,\infty}(\Sigma)\), and assume
\(\sigma(A|_{\operatorname{Ran}P})\subset[0,\lambda_-]\).  Then
\begin{equation}
 \|QM_\theta P\|
 \le\frac{
  \|\Delta_h\theta\|_\infty
  +2\lambda_-^{1/2}\|\nabla\theta\|_\infty
 }{\delta}.
\label{eq:v3v8-theta-mixing}
\end{equation}
\end{corollary}

\begin{proof}
For smooth \(u\),
\begin{equation}
 [-\Delta_h,M_\theta]u
 =-(\Delta_h\theta)u
  -2\langle\nabla\theta,\nabla u\rangle_h.
\label{eq:v3v8-laplacian-commutator}
\end{equation}
If \(u\in\operatorname{Ran}P\), then
\(\|\nabla u\|_2^2=\langle Au,u\rangle\le\lambda_-\|u\|_2^2\).
Estimate the two terms in
\eqref{eq:v3v8-laplacian-commutator} and apply
Theorem~\ref{thm:v3v8-commutator}.
\end{proof}

\begin{assessment}[The gap pays a derivative, not a free shell factor]
\label{ass:v3v8-gap-payment}
The numerator in \eqref{eq:v3v8-theta-mixing} contains actual derivatives
of the expansion.  The inverse gap converts that commutator into an
off-diagonal multiplier bound.  Replacing the numerator by
\(\|\theta\|_\infty\) would erase the derivative which creates the spectral
separation and would repeat the missing-weight error identified in the V5
companion audit.
\end{assessment}

\subsection{First-order wave system and leakage estimate}
\label{sec:v3v8-wave-system}

For clarity, keep the internal metric fixed in this subsection.  Let
\[
 A_m(t):=A(t)+m^2I
\]
and write the moving boundary wave equation as
\begin{equation}
 \partial_t^2\phi+M(t)\partial_t\phi+A_m(t)\phi=f,
\label{eq:v3v8-wave}
\end{equation}
where \(M(t)\) includes multiplication by \(\theta\) and any bounded
zero-order internal damping.  On the energy state
\(X=(\phi,v)\), \(v=\partial_t\phi\), this is
\begin{equation}
 \partial_tX=\mathbb A(t)X+F,
 \qquad
 \mathbb A(t)=
 \begin{pmatrix}0&I\\-A_m&-M\end{pmatrix},
 \qquad F=\binom0f.
\label{eq:v3v8-first-order}
\end{equation}
Put
\[
 \mathbb P=\operatorname{diag}(P,P),\quad
 \mathbb Q=I-\mathbb P,\quad
 \mathbb\Omega=\operatorname{diag}(\Omega,\Omega).
\]
The Kato-covariant system is
\begin{equation}
 D_tX=(\mathbb A-\mathbb\Omega)X+F.
\label{eq:v3v8-covariant-system}
\end{equation}

\begin{lemma}[Low-to-high coupling block]
\label{lem:v3v8-coupling-block}
On any common graph norm in which the displayed blocks are bounded,
\begin{equation}
 \|\mathbb Q(\mathbb A-\mathbb\Omega)\mathbb P\|
 \le \|\dot P\|+\|QMP\|.
\label{eq:v3v8-mu-bound}
\end{equation}
Consequently the mixing rate
\begin{equation}
 \mu(t):=
 \frac{\|Q\dot AP\|}{\delta}
 +\frac{\|Q[A,M]P\|}{\delta}
\label{eq:v3v8-mu}
\end{equation}
bounds the low-to-high block.
\end{lemma}

\begin{proof}
The identity block and \(A_m\) commute with \(P\).  The only off-diagonal
blocks are the Kato connection and \(QMP\).  The triangle inequality,
Theorem~\ref{thm:v3v8-Pdot}, and
Theorem~\ref{thm:v3v8-commutator} prove the two assertions.
\end{proof}

\begin{theorem}[Duhamel leakage bound]
\label{thm:v3v8-leakage}
Assume the diagonal high-sector equation generated by
\[
 \mathbb Q(\mathbb A-\mathbb\Omega)\mathbb Q
\]
has an evolution family \(\mathcal U_+(t,s)\) satisfying
\begin{equation}
 \|\mathcal U_+(t,s)\|
 \le\exp\left(\int_s^t\kappa_+(\tau)\,d\tau\right).
\label{eq:v3v8-high-propagator}
\end{equation}
If \(\mathbb Q(0)X(0)=0\), then
\begin{align}
 \|\mathbb Q(t)X(t)\|
 &\le\int_0^t
 \exp\left(\int_s^t\kappa_+(\tau)\,d\tau\right)
\notag\\
 &\qquad\qquad\times
 \left(
  \mu(s)\|\mathbb P(s)X(s)\|
  +\|\mathbb Q(s)F(s)\|
 \right)ds .
\label{eq:v3v8-leakage}
\end{align}
\end{theorem}

\begin{proof}
Because \(D_t\) preserves the \(P/Q\) splitting, applying \(\mathbb Q\) to
\eqref{eq:v3v8-covariant-system} gives a high-sector inhomogeneous equation.
Its source is
\[
 \mathbb Q(\mathbb A-\mathbb\Omega)\mathbb P X+\mathbb QF.
\]
Variation of constants with \(\mathcal U_+\), followed by
\eqref{eq:v3v8-mu-bound}, gives \eqref{eq:v3v8-leakage}.
\end{proof}

\begin{corollary}[Small leakage regime]
\label{cor:v3v8-small-leakage}
If \(\kappa_+\le K\),
\(\sup_{[0,T]}\|\mathbb PX\|\le M_-\),
\(\mathbb QF=0\), and
\(\int_0^T\mu(s)\,ds\le\varepsilon\), then
\begin{equation}
 \sup_{0\le t\le T}\|\mathbb QX(t)\|
 \le e^{KT}M_-\varepsilon.
\label{eq:v3v8-small-leakage}
\end{equation}
\end{corollary}

\begin{proof}
Apply Theorem~\ref{thm:v3v8-leakage} and bound the exponential and low state
uniformly.
\end{proof}

\subsection{Metric variation enters through a differentiated coefficient}
\label{sec:v3v8-metric-variation}

\begin{proposition}[Low-mode bound for a varying Laplacian]
\label{prop:v3v8-Adot}
Let \(A(t)=-\Delta_{h(t)}\), and assume on \([0,T]\) uniform bounded
geometry sufficient for the \(H^2\) elliptic estimate.  Then there is a
constant \(C_{\rm geom}\), depending on that geometry and \(\lambda_-\),
such that
\begin{equation}
 \|Q\dot AP\|_{L^2\to L^2}
 \le C_{\rm geom}\|\dot h\|_{W^{1,\infty}(h)}.
\label{eq:v3v8-Adot}
\end{equation}
Hence
\begin{equation}
 \|\dot P\|
 \le\frac{C_{\rm geom}}{\delta}
 \|\dot h\|_{W^{1,\infty}(h)}.
\label{eq:v3v8-Pdot-metric}
\end{equation}
\end{proposition}

\begin{proof}
In local coordinates,
\[
 \Delta_hu=|h|^{-1/2}\partial_i
 \bigl(|h|^{1/2}h^{ij}\partial_ju\bigr).
\]
Differentiation in time expresses \(\dot\Delta_hu\) as a linear combination
of \(\dot h\,\nabla^2u\) and \(\nabla\dot h\,\nabla u\), with coefficients
uniformly controlled by the background geometry.  Thus
\[
 \|\dot Au\|_2
 \le C\|\dot h\|_{W^{1,\infty}}\|u\|_{H^2}.
\]
On the finite low spectral subspace, the uniform elliptic estimate and
\(Au\le\lambda_-u\) give
\(\|u\|_{H^2}\le C_{\rm geom}\|u\|_2\).
Projection by \(Q\) cannot increase the norm, proving
\eqref{eq:v3v8-Adot}.  Apply Theorem~\ref{thm:v3v8-Pdot}.
\end{proof}

\begin{assessment}[Stronger than the V7 energy coefficient]
\label{ass:v3v8-stronger-coefficient}
V7 required an \(L^\infty\) bound on \(\dot h\) for energy deformation.
Transport of a Laplace spectral subspace also requires one spatial
derivative of \(\dot h\) in the elementary estimate above.  In ADM
variables this reaches derivatives of \(NK\) and second derivatives of the
shift.  The common continuum bootstrap must budget that stronger norm or
replace the spectral transport by a rougher functional-calculus argument.
\end{assessment}

\subsection{Collapse of the gap is a genuine obstruction}
\label{sec:v3v8-gap-collapse}

\begin{counterexample}[Slow operator motion with order-one projector motion]
\label{cex:v3v8-gap-collapse}
On \(\mathbb R^2\), let
\[
 A_\varepsilon(t)
 =R(t)
 \begin{pmatrix}0&0\\0&\varepsilon\end{pmatrix}
 R(t)^*,
\]
where \(R(t)\) is rotation through angle \(t\).  The spectral gap is
\(\delta_\varepsilon=\varepsilon\),
\(\|\dot A_\varepsilon\|=\varepsilon\), but the low eigenspace rotates with
\(\|\dot P_\varepsilon\|=1\).  Thus \(\dot A_\varepsilon\to0\) does not
force the low bundle to vary slowly when the gap collapses.
\end{counterexample}

\begin{proof}
The eigenvectors are the columns of \(R(t)\).  Differentiating the rank-one
projection onto the first column gives norm one.  Differentiating
\(A_\varepsilon\) gives a commutator with the rotation generator and norm
\(\varepsilon\).  Equality holds in the scaling of
\eqref{eq:v3v8-Pdot}.
\end{proof}

\begin{firewall}[Kato transport is not decoupling]
\label{fw:v3v8-not-decoupling}
The unitary \(U(t)\) identifies the moving low spaces, but it does not erase
the physical coupling \(\mu\).  Exact decoupling requires both
\(Q\dot AP=0\) and \(QMP=0\).  Otherwise the leakage term in
\eqref{eq:v3v8-leakage} must be paid.
\end{firewall}

\begin{firewall}[A boundary spectral gap is not the Yang--Mills mass gap]
\label{fw:v3v8-not-YM-gap}
The number \(\delta\) separates eigenvalues of an instantaneous boundary
elliptic operator on a compact slice.  It is a finite-volume analytic
quantity and does not prove a nonperturbative Yang--Mills mass gap.
\end{firewall}

\subsection{V8 conclusion and V9 target}
\label{sec:v3v8-conclusion}

V8 replaces the false invariant-harmonic assumption by a transported
low-frequency bundle.  The Kato connection costs
\(\delta^{-1}\|Q\dot AP\|\), multiplication by the expansion costs
\(\delta^{-1}\|Q[A,M]P\|\), and the Duhamel estimate records the resulting
high-mode leakage once.  The two-dimensional example proves that a uniform
gap cannot be omitted.

The next step is to determine what the high modes feed back into the inflow
block.  V9 will eliminate the complementary linear sector by a
Feshbach--Schur reduction, derive its causal memory kernel, and prove a
passivity estimate.  This will decide when the V5 finite-dimensional
physical metric remains coercive after the moving boundary modes are
integrated out, rather than assuming that their only effect is harmless
leakage.
\clearpage
\section{Third Volume, Version V9: Feshbach Reduction, Causal Memory, and Passivity of the Eliminated Boundary Sector}
\label{sec:v3v9}

\subsection{Purpose and relation to V8}
\label{sec:v3v9-purpose}

V8 bounded the transfer from the transported low spectral bundle into its
complement.  Leakage alone is not the full effect: the high modes propagate
and feed back into the low inflow variables.  V9 performs the exact block
elimination.  The resulting equation is nonlocal in time but causal.  Under
the adjoint port structure inherited from a single quadratic energy, the
memory is passive and cannot act as a hidden energy source.

The dynamic Volterra reduction and the static Feshbach--Schur reduction are
proved separately.  They coincide for a stationary Gaussian system after
Laplace transformation, but no such identification is assumed in the
time-dependent case.

\subsection{Exact causal elimination}
\label{sec:v3v9-elimination}

Let \(X_-\) and \(X_+\) be Hilbert spaces representing the Kato-transported
low and high sectors.  Consider
\begin{align}
 \dot x(t)
 &=A_{--}(t)x(t)+A_{-+}(t)y(t)+f_-(t),
\label{eq:v3v9-low}\\
 \dot y(t)
 &=A_{+-}(t)x(t)+A_{++}(t)y(t)+f_+(t).
\label{eq:v3v9-high}
\end{align}
Assume the high diagonal equation has an evolution family
\(U_+(t,s)\).

\begin{theorem}[Exact Volterra equation for the low sector]
\label{thm:v3v9-volterra}
Every mild solution of \eqref{eq:v3v9-low}--\eqref{eq:v3v9-high}
satisfies
\begin{align}
 \dot x(t)
 &=A_{--}(t)x(t)+f_-(t)+r_+(t)
   +\int_0^t\mathcal K(t,s)x(s)\,ds,
\label{eq:v3v9-volterra}\\
 \mathcal K(t,s)
 &:=A_{-+}(t)U_+(t,s)A_{+-}(s),
\label{eq:v3v9-memory}\\
 r_+(t)
 &:=A_{-+}(t)U_+(t,0)y(0)
\notag\\
 &\quad
 +A_{-+}(t)\int_0^tU_+(t,s)f_+(s)\,ds .
\label{eq:v3v9-high-remainder}
\end{align}
Conversely, a solution of \eqref{eq:v3v9-volterra}, together with
\begin{equation}
 y(t)=U_+(t,0)y(0)
 +\int_0^tU_+(t,s)
  \bigl(A_{+-}(s)x(s)+f_+(s)\bigr)\,ds,
\label{eq:v3v9-y-reconstruction}
\end{equation}
solves the original block system.  The memory kernel vanishes for
\(s>t\).
\end{theorem}

\begin{proof}
Variation of constants in \eqref{eq:v3v9-high} gives
\eqref{eq:v3v9-y-reconstruction}.  Substitute it into
\eqref{eq:v3v9-low} and collect the initial-high and high-source terms;
this gives \eqref{eq:v3v9-volterra}--\eqref{eq:v3v9-high-remainder}.
The converse follows by reversing the substitution.  The integration range
\(0\le s\le t\) proves causality.
\end{proof}

\begin{corollary}[Memory norm paid by the two incidence blocks]
\label{cor:v3v9-memory-bound}
If
\begin{equation}
 \|U_+(t,s)\|
 \le e^{-\int_s^t\alpha(\tau)\,d\tau},
\label{eq:v3v9-high-decay}
\end{equation}
then
\begin{equation}
 \|\mathcal K(t,s)\|
 \le
 \|A_{-+}(t)\|
 e^{-\int_s^t\alpha(\tau)\,d\tau}
 \|A_{+-}(s)\|.
\label{eq:v3v9-memory-bound}
\end{equation}
If \(\alpha\ge\alpha_0>0\) and both incidence blocks have norm at most
\(\mu_0\), then
\begin{equation}
 \sup_t\int_0^t\|\mathcal K(t,s)\|\,ds
 \le\frac{\mu_0^2}{\alpha_0}.
\label{eq:v3v9-memory-L1}
\end{equation}
\end{corollary}

\begin{proof}
Apply the operator norm to \eqref{eq:v3v9-memory}.  Under the uniform
assumptions, integrate
\(\mu_0^2e^{-\alpha_0(t-s)}\) from \(s=0\) to \(s=t\).
\end{proof}

\begin{assessment}[Two blocks, one transfer path]
\label{ass:v3v9-one-path}
The product in \eqref{eq:v3v9-memory} is one round trip:
low-to-high incidence, high propagation, then high-to-low incidence.
It is not two independent copies of low energy.  V8 pays each incidence
block through the spectral gap; V9 composes them before taking the norm, as
required by the companion single-payer lesson.
\end{assessment}

\subsection{Passive feedback from an adjoint coupling}
\label{sec:v3v9-passivity}

Assume for this subsection that the spaces and coefficients are stationary.
Let
\begin{align}
 \dot x&=A_-x-C^*y+f,
\label{eq:v3v9-passive-low}\\
 \dot y&=Cx+A_+y,
\label{eq:v3v9-passive-high}
\end{align}
where
\begin{equation}
 \operatorname{Re}\langle A_+y,y\rangle
 \le-\alpha\|y\|^2
\label{eq:v3v9-high-dissipative}
\end{equation}
for some \(\alpha\ge0\).  The feedback force on the low sector is
\[
 w(t):=-C^*y(t).
\]

\begin{theorem}[Time-domain passivity identity]
\label{thm:v3v9-passivity}
For every regular solution of \eqref{eq:v3v9-passive-high},
\begin{equation}
 \operatorname{Re}\langle w,x\rangle
 =-\frac12\frac{d}{dt}\|y\|^2
  +\operatorname{Re}\langle A_+y,y\rangle.
\label{eq:v3v9-passivity-identity}
\end{equation}
Consequently,
\begin{equation}
 \int_0^T\operatorname{Re}\langle w(t),x(t)\rangle\,dt
 \le
 \frac12\|y(0)\|^2-\frac12\|y(T)\|^2
 -\alpha\int_0^T\|y(t)\|^2dt.
\label{eq:v3v9-passivity-integrated}
\end{equation}
If \(y(0)=0\), the eliminated sector supplies no positive net energy to the
low sector.
\end{theorem}

\begin{proof}
Taking the real inner product of \eqref{eq:v3v9-passive-high} with \(y\)
gives
\[
 \frac12\frac{d}{dt}\|y\|^2
 =\operatorname{Re}\langle Cx,y\rangle
  +\operatorname{Re}\langle A_+y,y\rangle.
\]
Since
\(\operatorname{Re}\langle-C^*y,x\rangle
=-\operatorname{Re}\langle Cx,y\rangle\),
rearrangement proves \eqref{eq:v3v9-passivity-identity}.
Integrate and use \eqref{eq:v3v9-high-dissipative}.
\end{proof}

\begin{corollary}[Total energy estimate without a small-coupling condition]
\label{cor:v3v9-total-energy}
If also
\[
 \operatorname{Re}\langle A_-x,x\rangle\le0,
\]
then solutions of \eqref{eq:v3v9-passive-low}--\eqref{eq:v3v9-passive-high}
satisfy
\begin{equation}
 \frac12\frac{d}{dt}
 \bigl(\|x\|^2+\|y\|^2\bigr)
 \le\operatorname{Re}\langle f,x\rangle-\alpha\|y\|^2.
\label{eq:v3v9-total-energy}
\end{equation}
No smallness assumption on \(C\) is needed for this combined estimate.
\end{corollary}

\begin{proof}
Take inner products of the two equations with \(x\) and \(y\).  The two
coupling powers are
\(-\operatorname{Re}\langle C^*y,x\rangle\) and
\(\operatorname{Re}\langle Cx,y\rangle\), so they cancel.
\end{proof}

\subsection{Positive-real frequency response}
\label{sec:v3v9-positive-real}

Suppose \(A_+\) generates a contraction semigroup with
\eqref{eq:v3v9-high-dissipative}.  For \(\operatorname{Re}z>0\), define
\begin{equation}
 \widehat{\mathcal M}(z)
 :=C^*(zI-A_+)^{-1}C.
\label{eq:v3v9-transfer}
\end{equation}
The eliminated feedback in the low equation is
\(-\widehat{\mathcal M}(z)\widehat x(z)\).

\begin{theorem}[Positive-real memory transfer]
\label{thm:v3v9-positive-real}
For every \(\operatorname{Re}z>0\),
\begin{equation}
 \operatorname{Re}
 \langle\widehat{\mathcal M}(z)\xi,\xi\rangle
 \ge0.
\label{eq:v3v9-positive-real}
\end{equation}
More precisely, if
\(y=(zI-A_+)^{-1}C\xi\), then
\begin{equation}
 \operatorname{Re}
 \langle\widehat{\mathcal M}(z)\xi,\xi\rangle
 \ge(\operatorname{Re}z+\alpha)\|y\|^2.
\label{eq:v3v9-positive-real-margin}
\end{equation}
\end{theorem}

\begin{proof}
Since \(C\xi=(zI-A_+)y\),
\begin{align*}
 \operatorname{Re}
 \langle\widehat{\mathcal M}(z)\xi,\xi\rangle
 &=\operatorname{Re}\langle y,C\xi\rangle\\
 &=\operatorname{Re}z\,\|y\|^2
   -\operatorname{Re}\langle y,A_+y\rangle\\
 &\ge(\operatorname{Re}z+\alpha)\|y\|^2.
\end{align*}
\end{proof}

\begin{firewall}[Passivity depends on the adjoint sign]
\label{fw:v3v9-adjoint-sign}
The conclusion uses the pair \(C\) and \(-C^*\).  Arbitrary off-diagonal
blocks need not be passive.  Their memory kernel is still given by
\eqref{eq:v3v9-memory}, but it may inject energy.  The adjoint sign must be
derived from the common action or symplectic structure, not imposed after
estimating the two equations separately.
\end{firewall}

\subsection{Static Feshbach--Schur reduction}
\label{sec:v3v9-static-schur}

Let \(I\) be the finite inflow space and \(Y\) a Hilbert space of additional
boundary modes.  Consider the quadratic form
\begin{equation}
 \mathcal E(i,y)
 :=\frac12\langle Gi,i\rangle
   +\operatorname{Re}\langle Di,y\rangle
   +\frac12\langle Hy,y\rangle,
\label{eq:v3v9-static-energy}
\end{equation}
where \(G>0\), \(H\ge h_0I>0\), and \(D:I\to Y\) is bounded.  Define
\begin{equation}
 G_{\rm eff}:=G-D^*H^{-1}D.
\label{eq:v3v9-Geff}
\end{equation}

\begin{theorem}[Exact static elimination]
\label{thm:v3v9-static-schur}
For fixed \(i\), the unique minimizing high mode is
\begin{equation}
 y_*(i)=-H^{-1}Di,
\label{eq:v3v9-ystar}
\end{equation}
and
\begin{equation}
 \mathcal E(i,y)
 =\frac12\langle G_{\rm eff}i,i\rangle
 +\frac12
  \|H^{1/2}(y+H^{-1}Di)\|^2.
\label{eq:v3v9-square-completion}
\end{equation}
In particular,
\begin{equation}
 \inf_y\mathcal E(i,y)
 =\frac12\langle G_{\rm eff}i,i\rangle.
\label{eq:v3v9-reduced-energy}
\end{equation}
\end{theorem}

\begin{proof}
Expand the last squared norm in
\eqref{eq:v3v9-square-completion}:
\[
 \frac12\langle Hy,y\rangle
 +\operatorname{Re}\langle Di,y\rangle
 +\frac12\langle D^*H^{-1}Di,i\rangle.
\]
The remaining term is
\(\frac12\langle(G-D^*H^{-1}D)i,i\rangle\), giving
\eqref{eq:v3v9-static-energy}.  Strict positivity of \(H\) makes the square
vanish only at \eqref{eq:v3v9-ystar}.
\end{proof}

\begin{theorem}[Sharp coercivity margin after elimination]
\label{thm:v3v9-coercivity}
Set
\begin{equation}
 \rho:=\|H^{-1/2}DG^{-1/2}\|.
\label{eq:v3v9-rho}
\end{equation}
If \(\rho<1\), then
\begin{equation}
 (1-\rho^2)G\le G_{\rm eff}\le G.
\label{eq:v3v9-Geff-bound}
\end{equation}
Hence the reduced physical metric is positive and uniformly equivalent to
\(G\).  Conversely, in finite dimensions, \(G_{\rm eff}>0\) if and only if
\(\rho<1\).
\end{theorem}

\begin{proof}
For every \(i\),
\[
 \langle D^*H^{-1}Di,i\rangle
 =\|H^{-1/2}Di\|^2
 \le\rho^2\|G^{1/2}i\|^2.
\]
Subtraction gives the lower bound; positivity of the subtracted operator
gives the upper bound.  For the converse, conjugate:
\[
 G^{-1/2}G_{\rm eff}G^{-1/2}
 =I-(H^{-1/2}DG^{-1/2})^*
     (H^{-1/2}DG^{-1/2}).
\]
This operator is positive definite exactly when the largest singular value
of \(H^{-1/2}DG^{-1/2}\) is strictly below one.
\end{proof}

\begin{corollary}[Updated obstruction primitive]
\label{cor:v3v9-updated-primitive}
Let \(L:I\to\mathcal R\) be onto and assume \(\rho<1\).  After static
elimination of \(Y\), the least-energy primitive of \(Li=a\) is
\begin{equation}
 i_{\rm eff}(a)
 =G_{\rm eff}^{-1}L^*
  (LG_{\rm eff}^{-1}L^*)^{-1}a,
\label{eq:v3v9-effective-primitive}
\end{equation}
and its reduced cost is
\begin{equation}
 \frac12\left\langle
  (LG_{\rm eff}^{-1}L^*)^{-1}a,a
 \right\rangle.
\label{eq:v3v9-effective-cost}
\end{equation}
Moreover,
\begin{equation}
 LG^{-1}L^*
 \le LG_{\rm eff}^{-1}L^*
 \le(1-\rho^2)^{-1}LG^{-1}L^*.
\label{eq:v3v9-obstruction-comparison}
\end{equation}
\end{corollary}

\begin{proof}
Apply the V5 least-energy theorem with \(G_{\rm eff}\) in place of \(G\).
Inequality \eqref{eq:v3v9-Geff-bound} reverses under inversion; multiply
the resulting inequality on the left by \(L\) and on the right by \(L^*\).
\end{proof}

\begin{counterexample}[Threshold collapse of the reduced metric]
\label{cex:v3v9-threshold}
Take \(I=Y=\mathbb R\), \(G=H=1\), and \(D=\rho\).  Then
\[
 \mathcal E(i,y)=\frac12i^2+\rho iy+\frac12y^2,
 \qquad
 G_{\rm eff}=1-\rho^2.
\]
At \(\rho=1\), the direction \(y=-i\) has zero energy.  For \(\rho>1\),
the quadratic form is indefinite.  Thus strictness of the margin in
\eqref{eq:v3v9-Geff-bound} is necessary.
\end{counterexample}

\subsection{Finite-dimensional Gaussian check}
\label{sec:v3v9-gaussian}

\begin{proposition}[Bosonic Gaussian elimination]
\label{prop:v3v9-gaussian}
If \(Y=\mathbb R^N\), \(H>0\), and Lebesgue measure is normalized in the
usual way, then
\begin{align}
 &\int_{\mathbb R^N}
 \exp\bigl(-\mathcal E(i,y)\bigr)\,dy
\notag\\
 &\qquad
 =(2\pi)^{N/2}(\det H)^{-1/2}
 \exp\left(-\frac12\langle G_{\rm eff}i,i\rangle\right).
\label{eq:v3v9-gaussian}
\end{align}
\end{proposition}

\begin{proof}
Use the square completion
\eqref{eq:v3v9-square-completion}, translate
\(y\mapsto y-H^{-1}Di\), diagonalize the positive matrix \(H\), and
multiply the one-dimensional Gaussian integrals.
\end{proof}

\begin{firewall}[The determinant factor is not yet renormalized]
\label{fw:v3v9-determinant}
Equation \eqref{eq:v3v9-gaussian} is finite-dimensional and bosonic.  In a
continuum limit, \(\det H\) requires regularization and local counterterms;
fermionic and ghost sectors use graded determinants and may carry phase or
global anomaly data.  Positivity of \(G_{\rm eff}\) does not settle those
questions.
\end{firewall}

\begin{firewall}[Static positivity does not imply dynamic decay]
\label{fw:v3v9-static-dynamic}
The condition \(\rho<1\) makes the static Hessian positive.  Dynamic
passivity additionally requires the adjoint coupling and a dissipative or
conservative high evolution.  A nonnormal unstable \(A_+\) can violate the
memory estimates even when the instantaneous static Schur complement is
positive.
\end{firewall}

\subsection{V9 conclusion and V10 audit mandate}
\label{sec:v3v9-conclusion}

V9 gives the exact causal feedback of the boundary complement.  The memory
kernel contains both incidence maps and the high propagator in their actual
order.  With adjoint signs it is passive, and its positive-real transform
cannot inject hidden energy.  Static elimination preserves the V5 inflow
metric precisely below the sharp Schur threshold
\(\|H^{-1/2}DG^{-1/2}\|<1\).

V10 is the next mandatory companion audit.  It will fingerprint the current
Navier--Stokes and Yang--Mills notes before choosing the next route.  The
leading unresolved alternatives are a regulator-uniform memory/Schur
margin, a rough moving-projector construction requiring fewer metric
derivatives, and the variational Einstein junction law needed to turn the
abstract port identity into a constraint-compatible coupled boundary
system.
\clearpage
\section{Third Volume, Version V10: Companion Audit and the Variational Einstein--Inflow Junction Law}
\label{sec:v3v10}

\subsection{Mandatory ten-version companion audit}
\label{sec:v3v10-audit}

The required read-only audit was performed before the V10 direction was
chosen.  The companion inventory was:
\begin{center}
\begin{tabular}{p{0.48\textwidth}r p{0.33\textwidth}}
file & bytes & SHA--256\\
\hline
\texttt{Renewal\_NS\_Stability\_Research\_Notes\_V31.tex}
&887537&
\texttt{F1121B62238AD5491BB7CF03635EA9934942EDF573ED8FAAA9EE79E1D38A6696}
\\
\texttt{Renewal\_Yang\_Mills\_Mass\_Gap\_Research\_Notes\_V80.tex}
&1211024&
\texttt{44266D80D20D153F2DE7CB70AF3E53610B7AABF0119A515C2C77661D7A83104F}
\\
\texttt{Renewal\_Yang\_Mills\_Mass\_Gap\_Research\_Notes\_V81.tex}
&1243265&
\texttt{2FFA4AD294CEEFCB871F18E17F88780F3465D4D36130C749FF25612F1BAF4ECB}
\end{tabular}
\end{center}
The Navier--Stokes note and its fingerprint are unchanged from the V5
audit.  The two Yang--Mills files are no longer duplicates: V80 contains the
new conditional internal/crossing factorization and V81 appends the paid
internal parallel-carrier ledger.  Neither companion file was modified.

The NS lesson remains that a persistent resource is used once and that
changing the source norm cannot manufacture a missing formation moment.
The new YM results sharpen the common proof order:
\begin{enumerate}
\item condition the complete carrier into its exact internal and crossing
      algebras before estimating it;
\item form the complete product multiplier and its single denominator
      before taking moments;
\item contract an unused parallel carrier instead of charging both copies;
      and
\item attach output ancestry to the already allocated payer rather than
      adding it as new input mass.
\end{enumerate}
For the present programme this says that the bulk Einstein variation,
boundary kinetic variation, and exchange term must first be assembled into
one action.  The boundary stress and port flux are then derivatives of that
single carrier.  Adding a phenomenological junction stress after separately
estimating the boundary field would risk counting the same degree of freedom
twice.

\begin{firewall}[Audit type boundary]
\label{fw:v3v10-audit-types}
The YM conditional covariance is a Wilson-source identity and the NS
formation cost is a dyadic heat estimate.  Neither proves an Einstein
junction equation.  Only their complete-before-norm and one-use proof order
is transferred.  The geometric equations below are derived independently
from a variational principle.
\end{firewall}

\subsection{Geometric conventions and the complete action}
\label{sec:v3v10-action}

Let \((M,g)\) be a smooth Lorentzian manifold with smooth timelike boundary
\(\mathcal N\).  Let \(n\) be the outward spacelike unit normal,
\(\gamma=g-n^\flat\otimes n^\flat\) the induced Lorentzian metric, and
\begin{equation}
 K_{ab}:=\gamma_a{}^\mu\gamma_b{}^\nu\nabla_\mu n_\nu,
 \qquad K:=\gamma^{ab}K_{ab}.
\label{eq:v3v10-K}
\end{equation}
Fix \(\kappa>0\).  The gravitational action with its
Gibbons--Hawking--York term is
\begin{equation}
 S_{\rm grav}[g]
 :=\frac1{2\kappa}\int_M(R_g-2\Lambda)\,dV_g
   +\frac1\kappa\int_{\mathcal N}K\,dV_\gamma .
\label{eq:v3v10-grav-action}
\end{equation}
Let \(S_{\rm bulk}[g,\Psi]\) be the bulk Standard Model matter action at the
classical regularized level and let
\(S_\partial[\gamma,\phi,\Psi|_{\mathcal N}]\) be a boundary/inflow action.
The complete variational carrier is
\begin{equation}
 S_{\rm tot}:=S_{\rm grav}+S_{\rm bulk}+S_\partial.
\label{eq:v3v10-total-action}
\end{equation}
Define the bulk and boundary stress tensors by
\begin{align}
 \delta_gS_{\rm bulk}
 &:=-\frac12\int_M T_{\mu\nu}\,\delta g^{\mu\nu}\,dV_g,
\label{eq:v3v10-bulk-stress}\\
 \delta_\gamma S_\partial
 &:=-\frac12\int_{\mathcal N}
 S_{ab}\,\delta\gamma^{ab}\,dV_\gamma,
\label{eq:v3v10-boundary-stress}
\end{align}
where all boundary contributions, including metric dependence of an
interaction term, are included in \(S_{ab}\).

\subsection{First variation and the junction equation}
\label{sec:v3v10-first-variation}

\begin{lemma}[Einstein--Hilbert plus GHY variation]
\label{lem:v3v10-GHY}
For variations which preserve the location of \(\mathcal N\),
\begin{align}
 \delta S_{\rm grav}
 &=
 \frac1{2\kappa}\int_M
  (G_{\mu\nu}+\Lambda g_{\mu\nu})
  \delta g^{\mu\nu}\,dV_g
\notag\\
 &\quad
 +\frac1{2\kappa}\int_{\mathcal N}
  (K_{ab}-K\gamma_{ab})
  \delta\gamma^{ab}\,dV_\gamma.
\label{eq:v3v10-GHY-variation}
\end{align}
\end{lemma}

\begin{proof}
The Palatini identity gives
\[
 \delta(\sqrt{|g|}R)
 =\sqrt{|g|}G_{\mu\nu}\delta g^{\mu\nu}
  +\sqrt{|g|}\nabla_\mu
   \bigl(\nabla_\nu\delta g^{\mu\nu}
         -g_{\alpha\beta}\nabla^\mu\delta g^{\alpha\beta}\bigr).
\]
The divergence theorem converts the second term to a boundary expression
containing normal derivatives of \(\delta g\).  Direct variation of
\(\sqrt{|\gamma|}K\) gives the negative of those normal-derivative terms,
plus
\(\frac12\sqrt{|\gamma|}(K_{ab}-K\gamma_{ab})
\delta\gamma^{ab}\) and a tangential divergence.  The latter integrates to
zero when the boundary has no corners, or is canceled by the standard
corner term when corners are present.  Variation of the cosmological term
adds \(\Lambda g_{\mu\nu}\delta g^{\mu\nu}\), proving
\eqref{eq:v3v10-GHY-variation}.
\end{proof}

\begin{definition}[Brown--York momentum]
\label{def:v3v10-BY}
With the above orientation and sign convention, set
\begin{equation}
 \tau^{\rm BY}_{ab}
 :=\frac1\kappa(K_{ab}-K\gamma_{ab}).
\label{eq:v3v10-BY}
\end{equation}
\end{definition}

\begin{theorem}[Variational Einstein--inflow junction law]
\label{thm:v3v10-junction}
Suppose the induced metric is allowed to vary freely and all field
Euler--Lagrange equations hold.  Stationarity of the complete action
\eqref{eq:v3v10-total-action} is equivalent to the bulk Einstein equation
\begin{equation}
 G_{\mu\nu}+\Lambda g_{\mu\nu}
 =\kappa T_{\mu\nu}
\label{eq:v3v10-Einstein}
\end{equation}
and the boundary junction equation
\begin{equation}
 \boxed{\tau^{\rm BY}_{ab}=S_{ab}.}
\label{eq:v3v10-junction}
\end{equation}
\end{theorem}

\begin{proof}
Add \eqref{eq:v3v10-GHY-variation},
\eqref{eq:v3v10-bulk-stress}, and
\eqref{eq:v3v10-boundary-stress}.  Arbitrary compactly supported bulk
variations give \eqref{eq:v3v10-Einstein}.  The remaining coefficient of
the arbitrary boundary variation is
\[
 \frac12(\tau^{\rm BY}_{ab}-S_{ab})\delta\gamma^{ab},
\]
which vanishes for all \(\delta\gamma^{ab}\) exactly when
\eqref{eq:v3v10-junction} holds.
\end{proof}

\begin{assessment}[Dirichlet and dynamical boundary metrics]
\label{ass:v3v10-Dirichlet}
If \(\gamma\) is fixed, then \(\delta\gamma=0\) and no junction equation is
obtained.  Equation \eqref{eq:v3v10-junction} belongs to the dynamical
boundary-metric problem.  It must not be appended as an extra condition to
a Dirichlet Einstein problem without checking that the combined boundary
system is neither overdetermined nor characteristic.
\end{assessment}

\subsection{The V6 boundary field stress}
\label{sec:v3v10-scalar-stress}

On \((\mathcal N,\gamma)\), let \(I\) carry the positive metric \(G\) and
consider
\begin{equation}
 S_\phi
 :=-\frac12\int_{\mathcal N}
 \left(
  \gamma^{ab}\langle G D_a\phi,D_b\phi\rangle
  +m^2\langle G\phi,\phi\rangle
 \right)dV_\gamma .
\label{eq:v3v10-scalar-action}
\end{equation}
Define
\begin{equation}
 K_\gamma:=-\Box_\gamma+m^2.
\label{eq:v3v10-Kgamma}
\end{equation}

\begin{proposition}[Stress and its divergence]
\label{prop:v3v10-scalar-stress}
Metric variation of \eqref{eq:v3v10-scalar-action} gives
\begin{equation}
 S^\phi_{ab}
 =\langle GD_a\phi,D_b\phi\rangle
 -\frac12\gamma_{ab}
 \left(
  \langle GD^c\phi,D_c\phi\rangle
  +m^2\langle G\phi,\phi\rangle
 \right).
\label{eq:v3v10-scalar-stress}
\end{equation}
If \(K_\gamma\phi=f\), then
\begin{equation}
 D^aS^\phi_{ab}
 =-\langle Gf,D_b\phi\rangle.
\label{eq:v3v10-scalar-divergence}
\end{equation}
\end{proposition}

\begin{proof}
Use
\(\delta\sqrt{|\gamma|}
=-\frac12\sqrt{|\gamma|}\gamma_{ab}\delta\gamma^{ab}\)
and vary the inverse metric in the gradient contraction.  This gives
\eqref{eq:v3v10-scalar-stress}.  Taking its divergence, the two terms
containing \(D_aD_b\phi\) cancel by symmetry, leaving
\[
 \langle G(\Box_\gamma-m^2)\phi,D_b\phi\rangle
 =-\langle GK_\gamma\phi,D_b\phi\rangle.
\]
Substitute the field equation.
\end{proof}

\subsection{Codazzi compatibility is the boundary constraint}
\label{sec:v3v10-Codazzi}

\begin{lemma}[Contracted Codazzi identity]
\label{lem:v3v10-Codazzi}
With the convention \eqref{eq:v3v10-K},
\begin{equation}
 D^a(K_{ab}-K\gamma_{ab})
 =G_{\mu\nu}n^\mu\gamma^\nu{}_b.
\label{eq:v3v10-Codazzi}
\end{equation}
\end{lemma}

\begin{proof}
The Codazzi equation is
\[
 D_aK_{bc}-D_bK_{ac}
 =\gamma_a{}^\alpha\gamma_b{}^\beta\gamma_c{}^\gamma
  n^\delta R_{\delta\gamma\alpha\beta}.
\]
Contract the first and third tangential indices.  The result is
\(D^aK_{ab}-D_bK=R_{\mu\nu}n^\mu\gamma^\nu{}_b\).
Because \(g_{\mu\nu}n^\mu\gamma^\nu{}_b=0\), the Ricci contraction equals
the corresponding Einstein-tensor contraction.
\end{proof}

\begin{theorem}[Flux compatibility of the junction law]
\label{thm:v3v10-flux-compatibility}
Every smooth solution of the bulk Einstein equation and the junction law
necessarily satisfies
\begin{equation}
 \boxed{
 D^aS_{ab}
 =T_{\mu\nu}n^\mu\gamma^\nu{}_b.}
\label{eq:v3v10-flux-compatibility}
\end{equation}
If \(S_{ab}=S^\phi_{ab}\) and \(K_\gamma\phi=f\), this becomes
\begin{equation}
 T_{\mu\nu}n^\mu\gamma^\nu{}_b
 =-\langle Gf,D_b\phi\rangle.
\label{eq:v3v10-scalar-flux}
\end{equation}
\end{theorem}

\begin{proof}
Take the boundary divergence of
\eqref{eq:v3v10-junction}.  Apply
Lemma~\ref{lem:v3v10-Codazzi} and then
\eqref{eq:v3v10-Einstein}; the cosmological term has zero mixed
normal--tangential contraction.  This proves
\eqref{eq:v3v10-flux-compatibility}.  The scalar specialization follows
from Proposition~\ref{prop:v3v10-scalar-stress}.
\end{proof}

\begin{counterexample}[An arbitrary boundary stress is incompatible]
\label{cex:v3v10-incompatible-stress}
Let a proposed boundary tensor \(\widetilde S_{ab}\) satisfy
\[
 D^a\widetilde S_{ab}
 \ne T_{\mu\nu}n^\mu\gamma^\nu{}_b
\]
on an open boundary set.  Then no smooth metric can satisfy simultaneously
the bulk Einstein equation and
\(\tau^{\rm BY}_{ab}=\widetilde S_{ab}\) there.  This follows immediately
by taking the divergence and using the Codazzi identity.  Thus the abstract
V6 port cannot be converted into a junction stress by norm matching alone.
\end{counterexample}

\subsection{One variational port and its adjoint signs}
\label{sec:v3v10-port}

Let \(z\) denote a bulk boundary variable and let
\(\mathsf B\) be the trace map into its dual port space.  At the linear
level, take the single interaction carrier
\begin{equation}
 S_{\rm int}[\phi,z]
 :=-\int_{\mathcal N}\langle\mathsf B\phi,z\rangle\,dV_\gamma.
\label{eq:v3v10-interaction}
\end{equation}
Any metric variation of this expression is included in the total
\(S_{ab}\).

\begin{proposition}[The adjoint coupling is fixed before norms]
\label{prop:v3v10-adjoint-port}
Variation in \(\phi\) contributes the source
\begin{equation}
 f_{\rm port}=-G^{-1}\mathsf B^\dagger z
\label{eq:v3v10-port-force}
\end{equation}
to \(K_\gamma\phi=f_{\rm port}\).  Variation in \(z\) contributes
\(-\mathsf B\phi\) to the bulk boundary equation.  In a time foliation, the
associated instantaneous exchange powers have equal magnitude and opposite
sign:
\begin{equation}
 -\langle\mathsf B^\dagger z,\partial_t\phi\rangle
 =-\langle z,\mathsf B\partial_t\phi\rangle.
\label{eq:v3v10-port-adjoint}
\end{equation}
\end{proposition}

\begin{proof}
Differentiate the same bilinear functional
\eqref{eq:v3v10-interaction} in its two variables.  The definition of the
Hilbert adjoint gives \eqref{eq:v3v10-port-adjoint}.  No independent
coupling coefficient is introduced on either side.
\end{proof}

\begin{corollary}[Compatibility of the port with Codazzi flux]
\label{cor:v3v10-port-Codazzi}
For the scalar boundary stress, the port force
\eqref{eq:v3v10-port-force} gives
\begin{equation}
 D^aS^\phi_{ab}
 =\langle\mathsf B^\dagger z,D_b\phi\rangle
 =\langle z,\mathsf B D_b\phi\rangle.
\label{eq:v3v10-port-flux}
\end{equation}
Therefore the Einstein junction constraint is compatible precisely when
the mixed bulk stress flux equals the same port pairing:
\begin{equation}
 T_{\mu\nu}n^\mu\gamma^\nu{}_b
 =\langle z,\mathsf B D_b\phi\rangle.
\label{eq:v3v10-bulk-port-flux}
\end{equation}
\end{corollary}

\begin{proof}
Substitute \eqref{eq:v3v10-port-force} into
\eqref{eq:v3v10-scalar-divergence}, then apply
Theorem~\ref{thm:v3v10-flux-compatibility}.
\end{proof}

\begin{assessment}[The V9 passive sign is variational]
\label{ass:v3v10-passive-sign}
The pair \(\mathsf B\) and \(-\mathsf B^\dagger\) used in the V9 passive
memory theorem is not an auxiliary estimate.  It is the pair of first
derivatives of the single interaction
\eqref{eq:v3v10-interaction}.  Forming the action before eliminating the
high sector therefore supplies the adjoint sign required for passivity.
\end{assessment}

\subsection{Integrated boundary energy balance}
\label{sec:v3v10-integrated-balance}

Let \(t^a\) be a future timelike vector field tangent to \(\mathcal N\), and
define the boundary current
\begin{equation}
 J^a:=S^a{}_bt^b.
\label{eq:v3v10-current}
\end{equation}

\begin{proposition}[Geometric boundary work formula]
\label{prop:v3v10-boundary-work}
On a boundary slab \(\mathcal N_{[t_0,t_1]}\),
\begin{align}
 &\int_{\Sigma_{t_1}}J^a\nu_a\,d\mu
 -\int_{\Sigma_{t_0}}J^a\nu_a\,d\mu
\notag\\
 &\quad
 =\int_{\mathcal N_{[t_0,t_1]}}
 \left(
  T_{\mu\nu}n^\mu\gamma^\nu{}_bt^b
  +S_{ab}D^at^b
 \right)dV_\gamma .
\label{eq:v3v10-boundary-work}
\end{align}
The first term on the right is the unique bulk--boundary port power; the
second is geometric deformation work.
\end{proposition}

\begin{proof}
Compute
\[
 D_aJ^a=(D_aS^a{}_b)t^b+S^a{}_bD_at^b
\]
and use \eqref{eq:v3v10-flux-compatibility}.  The divergence theorem on the
boundary slab gives \eqref{eq:v3v10-boundary-work}, with \(\nu\) the future
normal to its spatial cuts and the displayed orientation.
\end{proof}

\begin{corollary}[Recovery of the V6--V7 ledger]
\label{cor:v3v10-recovery}
If \(t^a\) is Killing for the boundary metric, then
\(S_{ab}D^at^b=0\) for symmetric \(S\), and the boundary energy changes
only by the mixed bulk flux.  If \(t^a\) is not Killing, the deformation
term is the covariant form of the V7 coefficient
\(\mathcal D_{h,G}\).  In both cases the exchange flux occurs once.
\end{corollary}

\begin{proof}
For a symmetric stress tensor,
\[
 S_{ab}D^at^b=\frac12S_{ab}(\mathcal L_t\gamma)^{ab}.
\]
This vanishes for a Killing field.  In an ADM splitting, expanding
\(\mathcal L_t\gamma\) differentiates the spatial metric, lapse, and shift,
which are precisely the deformation terms in the V7 energy calculation.
\end{proof}

\subsection{Constraint and quantum firewalls}
\label{sec:v3v10-firewalls}

\begin{firewall}[Variational compatibility is not hyperbolic well-posedness]
\label{fw:v3v10-wellposedness}
The junction equation passes the necessary Codazzi constraint test.  This
does not prove that a gauge-fixed Einstein initial-boundary value problem
with \eqref{eq:v3v10-junction} satisfies the uniform Lopatinski condition or
generates a causal evolution.  That principal-symbol problem is the next
analytic gate.
\end{firewall}

\begin{firewall}[The classical Standard Model action is not the quantum
continuum]
\label{fw:v3v10-quantum}
The notation \(S_{\rm bulk}\) packages the classical regularized
Standard Model fields.  Removing the regulator requires gauge-invariant
renormalization, BRST identities, fermion determinant control, anomaly
cancellation in the relevant integral theory, and a measure on the coupled
metric--matter configuration space.  V10 proves none of those limits.
\end{firewall}

\begin{firewall}[Counterterms alter the boundary stress]
\label{fw:v3v10-counterterms}
Any local boundary counterterm contributes to
\(S_{ab}=-2\,\delta S_\partial/\delta\gamma^{ab}\).  It must be inserted
in the complete action before the junction and Schur reductions are formed.
Adding its stress afterward can break both the Codazzi Ward identity and the
single-payer allocation.
\end{firewall}

\subsection{V10 conclusion and V11 target}
\label{sec:v3v10-conclusion}

V10 completes the scheduled audit and derives the geometric law missing
from the V6--V9 port model.  Free variation of the induced metric equates
the Brown--York momentum to the complete boundary stress.  Contracted
Codazzi then forces the divergence of that stress to equal the mixed bulk
stress flux.  For the local inflow field, the same bilinear interaction
produces the adjoint port force, the boundary flux, and the passive sign
used in V9.  The exchange is therefore one variational carrier rather than
two separately normalized sources.

The next bottleneck is well-posedness.  V11 will analyze a linear
symmetric-hyperbolic bulk system with a dynamic boundary field and the
variational one-port condition.  It will prove the maximal-dissipativity
criterion and exhibit the boundary-symbol failure when a free junction law
is appended to incompatible incoming data.  This supplies a rigorous model
for the Einstein gauge-boundary test before attempting the full nonlinear
system.
\clearpage
\section{Third Volume, Version V11: Maximal Dissipativity of a Dynamic One-Port Boundary System}
\label{sec:v3v11}

\subsection{The well-posedness question left by V10}
\label{sec:v3v11-purpose}

The V10 junction law is variationally and constraint compatible, but those
facts alone do not make an Einstein initial-boundary value problem
well-posed.  A hyperbolic boundary must prescribe precisely the incoming
characteristic data, and a dynamic boundary field must exchange energy with
those data through a passive port.  V11 proves the complete statement for a
linear characteristic model.  The model is simple enough that maximal
dissipativity follows from an explicit resolvent rather than a formal energy
calculation.

Let \(\mathsf U\) and \(Z\) be finite-dimensional Hilbert spaces.  On the
half-line \(x>0\), let
\[
 w_+(t,x),w_-(t,x)\in\mathsf U,
 \qquad z(t)\in Z.
\]
The bulk characteristic equations are
\begin{align}
 \partial_tw_++\partial_xw_+&=0,
\label{eq:v3v11-win}\\
 \partial_tw_- -\partial_xw_-&=0.
\label{eq:v3v11-wout}
\end{align}
The \(+\) field enters the domain from \(x=0\), while the \(-\) field leaves
it.  Let
\[
 A_b:Z\to Z,\quad B:\mathsf U\to Z,\quad
 C:Z\to\mathsf U,\quad D:\mathsf U\to\mathsf U
\]
be bounded operators.  The dynamic boundary node is
\begin{align}
 \dot z&=A_bz+B\,w_-(t,0),
\label{eq:v3v11-node}\\
 w_+(t,0)&=Cz+D\,w_-(t,0).
\label{eq:v3v11-scattering}
\end{align}

\subsection{The exact passivity matrix}
\label{sec:v3v11-passivity}

Let \(G>0\) be the physical boundary-state metric inherited from the V5
inflow block.  The total energy is
\begin{equation}
 \mathcal E(t)
 :=\frac12\int_0^\infty
  \bigl(\|w_+\|^2+\|w_-\|^2\bigr)\,dx
  +\frac12\langle Gz,z\rangle.
\label{eq:v3v11-energy}
\end{equation}

\begin{lemma}[Characteristic boundary power]
\label{lem:v3v11-boundary-power}
For regular solutions decaying at infinity,
\begin{equation}
 \frac{d}{dt}\frac12\int_0^\infty
  \bigl(\|w_+\|^2+\|w_-\|^2\bigr)\,dx
 =\frac12\bigl(
  \|w_+(0)\|^2-\|w_-(0)\|^2
 \bigr).
\label{eq:v3v11-boundary-power}
\end{equation}
\end{lemma}

\begin{proof}
Multiply \eqref{eq:v3v11-win} and
\eqref{eq:v3v11-wout} by \(w_+\) and \(w_-\), respectively, take real
parts, and integrate by parts on the half-line.
\end{proof}

\begin{definition}[Scattering passivity matrix]
\label{def:v3v11-passivity-matrix}
Define the self-adjoint block operator on \(Z\oplus\mathsf U\)
\begin{equation}
 \mathcal P_G:=
 \begin{pmatrix}
  GA_b+A_b^*G+C^*C & GB+C^*D\\
  B^*G+D^*C & D^*D-I
 \end{pmatrix}.
\label{eq:v3v11-passivity-matrix}
\end{equation}
\end{definition}

\begin{theorem}[Energy passivity is a single block inequality]
\label{thm:v3v11-passivity}
For a solution of \eqref{eq:v3v11-win}--\eqref{eq:v3v11-scattering},
with \(y=w_-(t,0)\),
\begin{equation}
 \frac{d\mathcal E}{dt}
 =\frac12
 \left\langle
  \mathcal P_G\binom zy,\binom zy
 \right\rangle.
\label{eq:v3v11-energy-block}
\end{equation}
Consequently, the coupled node is passive for every boundary state and
outgoing trace if and only if
\begin{equation}
 \boxed{\mathcal P_G\le0.}
\label{eq:v3v11-passivity-condition}
\end{equation}
\end{theorem}

\begin{proof}
The state energy contributes
\[
 \frac12\frac{d}{dt}\langle Gz,z\rangle
 =\operatorname{Re}\langle G(A_bz+By),z\rangle.
\]
Lemma~\ref{lem:v3v11-boundary-power} and
\(w_+(0)=Cz+Dy\) contribute
\(\frac12(\|Cz+Dy\|^2-\|y\|^2)\).
Expanding their sum gives exactly
\eqref{eq:v3v11-energy-block}.  A Hermitian quadratic form is nonpositive
on every vector precisely when its representing operator is nonpositive.
\end{proof}

\begin{assessment}[Complete carrier before separate estimates]
\label{ass:v3v11-complete-carrier}
The cross block \(GB+C^*D\) is formed by adding the dynamic-state power to
the characteristic boundary power.  Estimating those two powers separately
would replace cancellation by a smallness condition.  The exact passivity
test \eqref{eq:v3v11-passivity-condition} is the boundary-symbol analogue
of the complete-product rule found in the V10 companion audit.
\end{assessment}

\subsection{The generator}
\label{sec:v3v11-generator}

Set
\begin{equation}
 \mathscr H
 :=L^2(\mathbb R_+;\mathsf U\oplus\mathsf U)\oplus Z
\label{eq:v3v11-H}
\end{equation}
with norm squared
\[
 \|(w_+,w_-,z)\|_{\mathscr H}^2
 =\|w_+\|_2^2+\|w_-\|_2^2+\langle Gz,z\rangle.
\]
Define
\begin{align}
 \mathscr A(w_+,w_-,z)
 &:=
 \bigl(-\partial_xw_+,\partial_xw_-,
       A_bz+B w_-(0)\bigr),
\label{eq:v3v11-generator}\\
 \mathcal D(\mathscr A)
 &:=
 \left\{
 (w_+,w_-,z)\in
 H^1(\mathbb R_+;\mathsf U)^2\oplus Z:
 w_+(0)=Cz+Dw_-(0)
 \right\}.
\label{eq:v3v11-domain}
\end{align}

\begin{lemma}[Density and closedness]
\label{lem:v3v11-closed}
The domain \(\mathcal D(\mathscr A)\) is dense in \(\mathscr H\), and
\(\mathscr A\) is closed.
\end{lemma}

\begin{proof}
For density, approximate the two \(L^2\) fields by smooth functions
supported away from \(x=0\), approximate \(z\) in the finite-dimensional
space, and add to \(w_+\) a boundary layer
\(\chi_n(x)(Cz)\), where \(\chi_n(0)=1\) and
\(\|\chi_n\|_{L^2}\to0\).  Add the analogous term
\(\chi_nD w_-(0)\) after first choosing an \(H^1\) approximation for
\(w_-\).  This enforces the trace relation while converging in
\(\mathscr H\).

For closedness, suppose \(X_n\to X\) and
\(\mathscr AX_n\to Y\) in \(\mathscr H\).  Convergence of each function and
its weak derivative gives convergence in \(H^1\) on every bounded interval,
so the boundary traces converge.  The finite-dimensional state converges,
and the trace relation passes to the limit.  The derivative and state
components of \(Y\) then equal those in \eqref{eq:v3v11-generator}.
\end{proof}

\begin{theorem}[Dissipativity criterion for the generator]
\label{thm:v3v11-generator-dissipative}
The operator \(\mathscr A\) is dissipative on \(\mathscr H\) if and only if
\(\mathcal P_G\le0\).
\end{theorem}

\begin{proof}
Integration by parts gives
\[
 \operatorname{Re}\langle\mathscr AX,X\rangle_{\mathscr H}
 =\frac12
 \left\langle\mathcal P_G\binom zy,\binom zy\right\rangle,
 \qquad y=w_-(0).
\]
Every pair \((z,y)\) is realized by some \(H^1\) traces satisfying
\eqref{eq:v3v11-scattering}.  Thus the quadratic form is nonpositive on the
operator domain exactly when \(\mathcal P_G\le0\).
\end{proof}

\subsection{Explicit resolvent and maximal dissipativity}
\label{sec:v3v11-resolvent}

\begin{theorem}[Surjectivity of the resolvent]
\label{thm:v3v11-resolvent}
Let \(\lambda>0\) be such that
\(\lambda I-A_b\) is invertible.  For every
\((f_+,f_-,g)\in\mathscr H\), the equation
\begin{equation}
 (\lambda I-\mathscr A)(w_+,w_-,z)
 =(f_+,f_-,g)
\label{eq:v3v11-resolvent-equation}
\end{equation}
has a unique solution in \(\mathcal D(\mathscr A)\).  It is given by
\begin{align}
 w_-(x)
 &=\int_x^\infty e^{-\lambda(s-x)}f_-(s)\,ds,
\label{eq:v3v11-wminus-resolvent}\\
 y:=w_-(0)
 &=\int_0^\infty e^{-\lambda s}f_-(s)\,ds,
\label{eq:v3v11-y-resolvent}\\
 z&=(\lambda I-A_b)^{-1}(g+By),
\label{eq:v3v11-z-resolvent}\\
 \xi&:=Cz+Dy,
\label{eq:v3v11-xi-resolvent}\\
 w_+(x)
 &=e^{-\lambda x}\xi
   +\int_0^x e^{-\lambda(x-s)}f_+(s)\,ds.
\label{eq:v3v11-wplus-resolvent}
\end{align}
\end{theorem}

\begin{proof}
The two transport resolvent equations are
\[
 \lambda w_++\partial_xw_+=f_+,
 \qquad
 \lambda w_- -\partial_xw_-=f_-.
\]
The unique \(L^2\) solution of the second equation is
\eqref{eq:v3v11-wminus-resolvent}; its trace is
\eqref{eq:v3v11-y-resolvent}.  The finite-dimensional equation is
\((\lambda I-A_b)z=g+By\), giving
\eqref{eq:v3v11-z-resolvent}.  The boundary condition fixes
\(\xi=w_+(0)\), and forward integration gives
\eqref{eq:v3v11-wplus-resolvent}.  Young's inequality shows that the
integral formulas lie in \(L^2\); the equations then place their weak
derivatives in \(L^2\), so they lie in \(H^1\).  Reversing the calculation
proves the equation, and the same formulas with zero data prove uniqueness.
\end{proof}

\begin{theorem}[Maximal dissipativity and causal evolution]
\label{thm:v3v11-maximal}
If \(\mathcal P_G\le0\), then \(\mathscr A\) is maximally dissipative and
generates a strongly continuous contraction semigroup on \(\mathscr H\).
Thus
\begin{equation}
 \|e^{t\mathscr A}X_0\|_{\mathscr H}
 \le\|X_0\|_{\mathscr H},
 \qquad t\ge0.
\label{eq:v3v11-contraction}
\end{equation}
\end{theorem}

\begin{proof}
Theorem~\ref{thm:v3v11-generator-dissipative} gives dissipativity.
For every sufficiently large \(\lambda>0\),
\(\lambda I-A_b\) is invertible, and
Theorem~\ref{thm:v3v11-resolvent} gives
\(\operatorname{Ran}(\lambda I-\mathscr A)=\mathscr H\).
The operator is densely defined and closed by
Lemma~\ref{lem:v3v11-closed}.  The Lumer--Phillips theorem now gives maximal
dissipativity and the contraction semigroup.
\end{proof}

\subsection{A conservative variational node}
\label{sec:v3v11-conservative-node}

Let \(K:Z\to\mathsf U\) and let \(J:Z\to Z\) satisfy
\begin{equation}
 GJ+J^*G=0.
\label{eq:v3v11-G-skew}
\end{equation}
Choose
\begin{equation}
 D=I,\qquad C=K,\qquad
 B=-G^{-1}K^*,\qquad
 A_b=J-\frac12G^{-1}K^*K.
\label{eq:v3v11-conservative-coefficients}
\end{equation}

\begin{proposition}[Exact lossless colligation]
\label{prop:v3v11-lossless}
For the coefficients \eqref{eq:v3v11-conservative-coefficients},
\begin{equation}
 \mathcal P_G=0.
\label{eq:v3v11-lossless}
\end{equation}
Hence the total energy is conserved, while the state damping
\(-\frac12G^{-1}K^*K\) exactly pays the interference between
\(Kz\) and the direct outgoing-to-incoming channel.
\end{proposition}

\begin{proof}
The lower-right block is \(D^*D-I=0\).  The off-diagonal block is
\[
 GB+C^*D=-K^*+K^*=0.
\]
The upper-left block is
\[
 GJ+J^*G-K^*K+K^*K=0.
\]
Thus every block of \(\mathcal P_G\) vanishes.
\end{proof}

\begin{assessment}[Physical metric cannot be renormalized independently]
\label{ass:v3v11-physical-metric}
The coefficient \(B=-G^{-1}K^*\) and the damping
\(-\frac12G^{-1}K^*K\) are fixed by the same physical state metric that
appears in the V5 primitive.  Rescaling \(G\) while leaving \(B\) and
\(A_b\) unchanged generally destroys \(\mathcal P_G\le0\).  This is the
dynamic version of the V5 collapsing-physical-margin counterexample.
\end{assessment}

\subsection{Boundary-symbol check}
\label{sec:v3v11-symbol}

Take the Laplace transform in time with
\(\operatorname{Re}s>0\) and consider a homogeneous decaying bulk mode.
Equations \eqref{eq:v3v11-win}--\eqref{eq:v3v11-wout} give
\begin{equation}
 w_+(x)=e^{-sx}a,\qquad w_-(x)=0.
\label{eq:v3v11-stable-mode}
\end{equation}
The homogeneous boundary symbol is
\begin{equation}
 (sI-A_b)\zeta=0,\qquad a=C\zeta.
\label{eq:v3v11-boundary-symbol}
\end{equation}

\begin{proposition}[No unstable boundary mode under passivity]
\label{prop:v3v11-symbol}
If \(\mathcal P_G\le0\), then
\(sI-A_b\) is injective for every \(\operatorname{Re}s>0\).
Consequently the decaying homogeneous boundary-symbol problem has only the
zero solution.
\end{proposition}

\begin{proof}
Set \(y=0\) in the passivity inequality.  This gives
\[
 2\operatorname{Re}\langle GA_bz,z\rangle+\|Cz\|^2\le0,
\]
so \(A_b\) is dissipative in the \(G\)-metric.  If
\((sI-A_b)\zeta=0\), then
\[
 \operatorname{Re}s\,\|\zeta\|_G^2
 =\operatorname{Re}\langle GA_b\zeta,\zeta\rangle\le0.
\]
Since \(\operatorname{Re}s>0\), \(\zeta=0\), and then \(a=0\).
\end{proof}

\subsection{Why an extra incoming condition fails}
\label{sec:v3v11-overdetermination}

\begin{counterexample}[A junction law plus Dirichlet incoming data is not
maximal]
\label{cex:v3v11-extra-condition}
Take \(Z=\{0\}\), \(D=I\), so the passive junction law is
\begin{equation}
 w_+(0)=w_-(0).
\label{eq:v3v11-reflection}
\end{equation}
Append the independent condition \(w_+(0)=0\).  The two conditions force
\(w_-(0)=0\).  For any \(\lambda>0\), choose
\(f_-\in L^2(\mathbb R_+;\mathsf U)\) with
\[
 \int_0^\infty e^{-\lambda s}f_-(s)\,ds\ne0.
\]
Equation \eqref{eq:v3v11-wminus-resolvent} then forces a nonzero trace
\(w_-(0)\), contradicting the appended boundary conditions.  Therefore
\(\lambda I-\mathscr A\) is not onto and the restricted dissipative
operator is not maximal.
\end{counterexample}

\begin{assessment}[Meaning for the Einstein junction problem]
\label{ass:v3v11-Einstein-meaning}
A variational junction law must be expressed in the incoming
characteristic variables of a chosen hyperbolic Einstein gauge.  One may
also prescribe gauge and constraint data, but the total boundary symbol
must have the correct rank.  Keeping a full Dirichlet metric condition and
then appending the free Brown--York law has the same structural defect as
the counterexample: it can make the resolvent range smaller than the state
space.
\end{assessment}

\subsection{A sign error can create boundary amplification}
\label{sec:v3v11-bad-sign}

\begin{counterexample}[Reversing the adjoint port sign]
\label{cex:v3v11-bad-sign}
Let \(G=I\), \(D=I\), \(C=K\ne0\), and keep
\(A_b=-\frac12K^*K\), but take \(B=+K^*\) instead of \(-K^*\).
Then
\[
 \mathcal P_I=
 \begin{pmatrix}0&2K^*\\2K&0\end{pmatrix},
\]
which has positive spectrum.  Suitable boundary traces make
\(\frac{d\mathcal E}{dt}>0\).  Thus the sign fixed by the V10 interaction
action is necessary for passivity.
\end{counterexample}

\subsection{Uniform families}
\label{sec:v3v11-uniform}

\begin{corollary}[Regulator-uniform contraction criterion]
\label{cor:v3v11-uniform}
Let \(G_n,A_{b,n},B_n,C_n,D_n\) be a regulator family.  If
\[
 g_-I\le G_n\le g_+I,\qquad
 \mathcal P_{G_n}\le0
\]
with \(g_\pm\) independent of \(n\), then every generated semigroup is a
contraction in its physical norm, and those norms are uniformly equivalent
to one fixed Hilbert norm.  If the coefficients converge in operator norm,
the explicit resolvents \eqref{eq:v3v11-wminus-resolvent}--%
\eqref{eq:v3v11-wplus-resolvent} converge for every sufficiently large
fixed \(\lambda\).
\end{corollary}

\begin{proof}
Theorem~\ref{thm:v3v11-maximal} applies for every \(n\).
The \(G_n\) bounds give uniform norm equivalence.  Matrix inversion is
continuous at \(\lambda I-A_b\), and all other operations in the explicit
resolvent are bounded integral maps independent of the regulator.
\end{proof}

\begin{firewall}[The model is not the Einstein principal symbol]
\label{fw:v3v11-model}
The two transport fields represent normalized incoming and outgoing
characteristics.  A gauge-fixed Einstein--Standard Model system has
multiple speeds, zero-speed constraint fields, gauge modes, and
frequency-dependent boundary symbols.  V11 proves the exact operator
criterion for the one-port model and identifies what must be checked; it
does not assert that a particular Einstein gauge already satisfies it.
\end{firewall}

\subsection{V11 conclusion and V12 target}
\label{sec:v3v11-conclusion}

V11 turns the V10 variational port into a complete linear
initial-boundary evolution.  The single matrix
\(\mathcal P_G\) is equivalent to energy passivity; its nonpositivity plus
the explicit resolvent proves maximal dissipativity and a contraction
semigroup.  The lossless colligation shows how the physical metric fixes
the adjoint port and state damping.  The extra-condition counterexample
proves that a junction law cannot simply be appended to independently fixed
incoming data.

V12 will lift this calculation from unit-speed transport to a general
constant-coefficient symmetric-hyperbolic bulk system.  It will diagonalize
the boundary matrix, include its zero-speed constraint block, and give a
frequency-uniform Kreiss/Lopatinski margin for a dynamic boundary node.
That is the closest finite-symbol precursor to a gauge-fixed linearized
Einstein--Standard Model boundary theorem.
\clearpage
\section{Third Volume, Version V12: General Symmetric-Hyperbolic Boundary Symbols and the Zero-Speed Constraint Gate}
\label{sec:v3v12}

\subsection{Constant-coefficient bulk system}
\label{sec:v3v12-setup}

V11 proved maximal dissipativity for two normalized unit-speed
characteristics.  A gauge-fixed Einstein--Standard Model system has a
matrix-valued principal symbol, tangential frequencies, unequal
characteristic speeds, and often a zero-speed constraint sector.  V12
isolates those features at the frozen-coefficient level.

On the half-space
\(\mathbb R_t\times\mathbb R^{d-1}_y\times\mathbb R_{+,x}\), consider
\begin{equation}
 A^0\partial_tu+A^n\partial_xu
 +\sum_{j=1}^{d-1}A^j\partial_{y_j}u=0,
\label{eq:v3v12-system}
\end{equation}
where \(A^0>0\) and all \(A^n,A^j\) are self-adjoint matrices.  The bulk
energy is
\begin{equation}
 E_{\rm bulk}(t)
 :=\frac12\int_{x>0}
 \langle A^0u,u\rangle\,dx\,dy .
\label{eq:v3v12-bulk-energy}
\end{equation}

\begin{lemma}[Frozen bulk energy flux]
\label{lem:v3v12-flux}
For regular solutions decaying at spatial infinity,
\begin{equation}
 \frac{dE_{\rm bulk}}{dt}
 =\frac12\int_{x=0}\langle A^nu,u\rangle\,dy .
\label{eq:v3v12-flux}
\end{equation}
\end{lemma}

\begin{proof}
Take the real inner product of \eqref{eq:v3v12-system} with \(u\).
Self-adjointness converts every spatial term into a divergence.  Tangential
divergences integrate to zero, while integration on \(x>0\) gives the
displayed boundary sign.
\end{proof}

\subsection{Canonical incoming, outgoing, and zero-speed traces}
\label{sec:v3v12-characteristics}

Set
\begin{equation}
 H_n:=(A^0)^{-1/2}A^n(A^0)^{-1/2}.
\label{eq:v3v12-Hn}
\end{equation}
Choose a unitary \(V\) such that
\begin{equation}
 VH_nV^*
 =
 \begin{pmatrix}
  \Lambda_+&0&0\\
  0&-\Lambda_-&0\\
  0&0&0
 \end{pmatrix},
\label{eq:v3v12-diagonal}
\end{equation}
where \(\Lambda_\pm>0\).  For
\(\widetilde u=V(A^0)^{1/2}u|_{x=0}\), write
\(\widetilde u=(u_+,u_-,r)\) and define normalized port variables
\begin{equation}
 w_+:=\Lambda_+^{1/2}u_+,
 \qquad
 w_-:=\Lambda_-^{1/2}u_-.
\label{eq:v3v12-normalized-ports}
\end{equation}

\begin{proposition}[Flux normal form]
\label{prop:v3v12-flux-normal}
Pointwise on the boundary,
\begin{equation}
 \langle A^nu,u\rangle
 =\|w_+\|^2-\|w_-\|^2.
\label{eq:v3v12-flux-normal}
\end{equation}
The zero-speed trace \(r\) has no sign and no norm in the bulk boundary
flux.
\end{proposition}

\begin{proof}
Substitute \eqref{eq:v3v12-diagonal} and
\eqref{eq:v3v12-normalized-ports} into the quadratic form.
\end{proof}

\subsection{Why an uncontrolled zero-speed trace cannot enter a passive
port}
\label{sec:v3v12-zero-passivity}

Suppose one tries to use the dynamic node
\begin{align}
 \dot z&=A_bz+Bw_-+Fr,
\label{eq:v3v12-node-zero}\\
 w_+&=Cz+Dw_-+Er
\label{eq:v3v12-scatter-zero}
\end{align}
with state metric \(G>0\), while assigning no energy or constraint to
\(r\).

\begin{theorem}[Zero-flux decoupling forced by passivity]
\label{thm:v3v12-zero-decoupling}
If
\begin{align}
 &2\operatorname{Re}
 \langle G(A_bz+By+Fr),z\rangle
\notag\\
 &\qquad
 +\|Cz+Dy+Er\|^2-\|y\|^2\le0
\label{eq:v3v12-zero-passivity}
\end{align}
for every \(z,y,r\), then
\begin{equation}
 E=0,\qquad F=0.
\label{eq:v3v12-zero-decoupling}
\end{equation}
Thus a passive boundary law cannot depend on a freely variable trace which
has zero bulk flux and no boundary energy.
\end{theorem}

\begin{proof}
Set \(z=y=0\) in \eqref{eq:v3v12-zero-passivity}.  This gives
\(\|Er\|^2\le0\) for every \(r\), hence \(E=0\).  With \(E=0\), fix
\(z,r\), set \(y=0\), and replace \(r\) by \(tr\) for arbitrary real
\(t\).  The only term linear in \(t\) is
\(2t\operatorname{Re}\langle GFr,z\rangle\).  A real affine function
bounded above for both signs and all magnitudes of \(t\) has zero linear
coefficient.  Replacing \(z\) by \(iz\) in the complex case also kills the
imaginary part.  Hence \(\langle GFr,z\rangle=0\) for all \(r,z\), so
\(F=0\).
\end{proof}

\begin{assessment}[Constraint or energy, rather than a silent zero block]
\label{ass:v3v12-zero-choice}
A zero-speed trace may couple after it is determined by a propagated
constraint, or after a positive boundary energy is assigned to it.  It
cannot be inserted as an arbitrary input to the V11 passivity matrix.  In
an Einstein gauge this is the point at which gauge and constraint
characteristics must be separated from freely specifiable radiation data.
\end{assessment}

\subsection{Uniform algebraic elimination of the zero block}
\label{sec:v3v12-zero-elimination}

Let the transformed zero row of the frozen Laplace--Fourier system and any
boundary constraint give
\begin{equation}
 R_0(s,\xi)r=S_0(s,\xi)z+T_0(s,\xi)w_-+g_0.
\label{eq:v3v12-zero-constraint}
\end{equation}
Assume
\begin{equation}
 \sigma_{\min}(R_0(s,\xi))\ge\rho_0>0
\label{eq:v3v12-zero-margin}
\end{equation}
on the normalized frequency set under consideration.  Then
\begin{equation}
 r=R_0^{-1}S_0z+R_0^{-1}T_0w_-+R_0^{-1}g_0.
\label{eq:v3v12-zero-solution}
\end{equation}

\begin{proposition}[Effective dynamic node after zero-block elimination]
\label{prop:v3v12-effective-node}
Substituting \eqref{eq:v3v12-zero-solution} into
\eqref{eq:v3v12-node-zero}--\eqref{eq:v3v12-scatter-zero} gives
\begin{align}
 \dot z&=A_b^{\rm eff}z+B^{\rm eff}w_-+f_0^{\rm eff},
\label{eq:v3v12-effective-dynamics}\\
 w_+&=C^{\rm eff}z+D^{\rm eff}w_-+g_0^{\rm eff},
\label{eq:v3v12-effective-scattering}
\end{align}
where
\begin{align}
 A_b^{\rm eff}&=A_b+FR_0^{-1}S_0,
 &B^{\rm eff}&=B+FR_0^{-1}T_0,
\label{eq:v3v12-effective-AB}\\
 C^{\rm eff}&=C+ER_0^{-1}S_0,
 &D^{\rm eff}&=D+ER_0^{-1}T_0,
\label{eq:v3v12-effective-CD}\\
 f_0^{\rm eff}&=FR_0^{-1}g_0,
 &g_0^{\rm eff}&=ER_0^{-1}g_0.
\label{eq:v3v12-effective-data}
\end{align}
Every correction is bounded by the actual inverse constraint margin
\(\rho_0^{-1}\).
\end{proposition}

\begin{proof}
This is direct substitution and collection of the \(z\), \(w_-\), and
\(g_0\) coefficients.  Bound \(\|R_0^{-1}\|\le\rho_0^{-1}\).
\end{proof}

\begin{counterexample}[Collapsing zero-block margin]
\label{cex:v3v12-zero-collapse}
Take scalar \(R_{0,\varepsilon}=\varepsilon\), \(S_0=1\), and
\(T_0=g_0=0\).  Then \(r=\varepsilon^{-1}z\).  If \(E\ne0\) or \(F\ne0\),
the corresponding effective node coefficient diverges as
\(\varepsilon\downarrow0\).  A regulator-uniform boundary theorem therefore
requires the positive margin \eqref{eq:v3v12-zero-margin}; formal
solvability at each cutoff is insufficient.
\end{counterexample}

\subsection{Laplace--Fourier stable subspace}
\label{sec:v3v12-stable-subspace}

Take the Laplace transform in \(t\) and Fourier transform in \(y\).  For
\(\operatorname{Re}s>0\), the normal equation is
\begin{equation}
 A^n\partial_x\widehat u
 =-\left(sA^0+iA(\xi)\right)\widehat u,
\qquad
 A(\xi):=\sum_{j=1}^{d-1}\xi_jA^j.
\label{eq:v3v12-normal-ode}
\end{equation}
Assume after the zero-block reduction that the decaying solutions form a
stable bundle \(\mathbb E^s(s,\xi)\) of constant dimension equal to the
number of incoming variables.  Let
\begin{equation}
 \mathcal T^s(s,\xi)a
 =\bigl(W_+(s,\xi)a,W_-(s,\xi)a\bigr)
\label{eq:v3v12-stable-trace}
\end{equation}
be their normalized boundary traces.

We impose the following explicit regularity gate:
\begin{equation}
 \mathbb E^s,\ W_+,\ W_-\ \hbox{extend continuously to}\quad
 \mathfrak S:=
 \{(s,\xi):\operatorname{Re}s\ge0,\
 |s|^2+|\xi|^2=1\}.
\label{eq:v3v12-stable-continuity}
\end{equation}
This excludes variable-multiplicity and unresolved glancing singularities;
it is a hypothesis to be verified for a chosen Einstein gauge.

\begin{lemma}[Stable-mode flux identity]
\label{lem:v3v12-stable-flux}
For every decaying solution of \eqref{eq:v3v12-normal-ode},
\begin{equation}
 \|W_+a\|^2-\|W_-a\|^2
 =2\operatorname{Re}s
 \int_0^\infty
 \langle A^0\widehat u,\widehat u\rangle\,dx
 \ge0.
\label{eq:v3v12-stable-flux}
\end{equation}
The nonnegative identity persists on \(\mathfrak S\) by the continuous
stable-subspace limit.
\end{lemma}

\begin{proof}
Take the real inner product of
\eqref{eq:v3v12-normal-ode} with \(\widehat u\) and integrate in \(x\).
The tangential term is purely imaginary.  The normal derivative integrates
to minus one half of the boundary quadratic form.  Rearrangement and
Proposition~\ref{prop:v3v12-flux-normal} give
\eqref{eq:v3v12-stable-flux}.  Pass to the boundary of the frequency set by
\eqref{eq:v3v12-stable-continuity}.
\end{proof}

\subsection{Dynamic Lopatinski operator}
\label{sec:v3v12-Lopatinski}

After zero-block elimination, use the effective coefficients and consider
the homogeneous dynamic boundary symbol
\begin{align}
 W_+a&=Cz+DW_-a,
\label{eq:v3v12-symbol-scatter}\\
 (sI-A_b)z&=BW_-a.
\label{eq:v3v12-symbol-node}
\end{align}
Define
\begin{equation}
 \mathfrak L(s,\xi)
 \binom az
 :=
 \binom{
  (W_+-DW_-)a-Cz
 }{
  -BW_-a+(sI-A_b)z
 }.
\label{eq:v3v12-L}
\end{equation}
The domain and target have the same finite dimension under the stable-count
hypothesis.

\begin{definition}[Uniform dynamic Lopatinski margin]
\label{def:v3v12-gamma}
The boundary system has margin \(\gamma_{\rm Lop}>0\) if
\begin{equation}
 \|\mathfrak L(s,\xi)(a,z)\|
 \ge\gamma_{\rm Lop}\|(a,z)\|
\label{eq:v3v12-Lopatinski-margin}
\end{equation}
for every \((s,\xi)\in\mathfrak S\).
\end{definition}

\subsection{Strict passivity yields a uniform symbol margin}
\label{sec:v3v12-strict-passivity}

Let the effective coefficients satisfy the strict physical passivity
condition
\begin{equation}
 \mathcal P_G
 :=
 \begin{pmatrix}
  GA_b+A_b^*G+C^*C&GB+C^*D\\
  B^*G+D^*C&D^*D-I
 \end{pmatrix}
 \le-\varepsilon
 \begin{pmatrix}G&0\\0&I\end{pmatrix}
\label{eq:v3v12-strict-passivity}
\end{equation}
for some \(\varepsilon>0\), uniformly on \(\mathfrak S\) if the effective
coefficients depend on frequency.

\begin{theorem}[No closed-half-plane surface mode]
\label{thm:v3v12-no-surface-mode}
Under \eqref{eq:v3v12-strict-passivity}, the homogeneous equations
\eqref{eq:v3v12-symbol-scatter}--\eqref{eq:v3v12-symbol-node} have only
\(a=0\), \(z=0\) for every \((s,\xi)\in\mathfrak S\).
\end{theorem}

\begin{proof}
Put \(y=W_-a\).  Evaluate the passivity quadratic form at \((z,y)\).
Using the two homogeneous boundary equations gives
\begin{align}
 &2\operatorname{Re}\langle G(A_bz+By),z\rangle
 +\|Cz+Dy\|^2-\|y\|^2
\notag\\
 &\qquad
 =2\operatorname{Re}s\,\|z\|_G^2
  +\|W_+a\|^2-\|W_-a\|^2.
\label{eq:v3v12-symbol-energy}
\end{align}
Lemma~\ref{lem:v3v12-stable-flux} makes the right-hand side nonnegative.
Strict passivity makes the left-hand side at most
\(-\varepsilon(\|z\|_G^2+\|y\|^2)\).  Hence \(z=0\) and \(y=0\).
The scattering equation then gives \(W_+a=0\).  The full nonzero-speed trace
vanishes, and the uniformly solved zero constraint gives \(r=0\).
Uniqueness for the normal ODE with zero boundary trace yields \(a=0\).
\end{proof}

\begin{theorem}[Compactness produces the uniform Kreiss margin]
\label{thm:v3v12-uniform-Lopatinski}
Assume \eqref{eq:v3v12-stable-continuity}, the stable-count condition, the
uniform zero-block margin \(\rho_0>0\), and uniform strict passivity
\eqref{eq:v3v12-strict-passivity}.  Then
\begin{equation}
 \gamma_{\rm Lop}
 :=\min_{(s,\xi)\in\mathfrak S}
 \sigma_{\min}(\mathfrak L(s,\xi))>0.
\label{eq:v3v12-gamma-formula}
\end{equation}
\end{theorem}

\begin{proof}
Theorem~\ref{thm:v3v12-no-surface-mode} gives injectivity of the square
matrix \(\mathfrak L(s,\xi)\), hence invertibility, at every frequency.
The zero-block inverse and stable traces are continuous by hypothesis, so
\(\mathfrak L\) and its least singular value are continuous.
The normalized frequency set \(\mathfrak S\) is compact.  A positive
continuous function on a compact set has a positive minimum.
\end{proof}

\begin{assessment}[Strictness supplies the missing endpoint compactness]
\label{ass:v3v12-strictness}
Ordinary passivity excludes modes with \(\operatorname{Re}s>0\), as in V11,
but it may allow neutral surface modes on \(\operatorname{Re}s=0\).
Strict passivity excludes those endpoint modes and turns pointwise
injectivity into the uniform margin
\eqref{eq:v3v12-gamma-formula}.  This distinction is essential for
variable-coefficient localization and regulator removal.
\end{assessment}

\subsection{Two failures of a merely formal symbol test}
\label{sec:v3v12-failures}

\begin{counterexample}[Passive but nonuniform boundary oscillator]
\label{cex:v3v12-neutral-node}
Let the boundary state decouple from the bulk:
\[
 B=C=0,\qquad D=0,\qquad GA_b+A_b^*G=0.
\]
Then \(\mathcal P_G=\operatorname{diag}(0,-I)\le0\).  If
\(A_bz=i\omega z\) for some nonzero \(z\), the boundary symbol
\((sI-A_b)z\) vanishes at \(s=i\omega\).  There is no unstable
right-half-plane mode, but
\(\gamma_{\rm Lop}=0\).  Passivity without strictness is therefore
insufficient for a uniform Kreiss estimate.
\end{counterexample}

\begin{counterexample}[Pointwise zero-block solvability without a margin]
\label{cex:v3v12-pointwise-zero}
The family \(R_{0,n}=n^{-1}I\) is invertible at every cutoff, yet
\(\|R_{0,n}^{-1}\|=n\).  Even if the unreduced matrices are uniformly
bounded, the effective coefficients in
\eqref{eq:v3v12-effective-AB}--\eqref{eq:v3v12-effective-CD} may diverge,
and the Lopatinski margin may collapse.  Checking rank at each regulator
does not prove a uniform continuum boundary estimate.
\end{counterexample}

\begin{firewall}[The stable-bundle hypothesis is a real Einstein gate]
\label{fw:v3v12-stable-gate}
At glancing frequencies or variable characteristic multiplicity, the
stable subspace may fail to extend continuously to
\(\operatorname{Re}s=0\).  The compactness proof then does not apply.
A full Einstein--Standard Model theorem must analyze the chosen
gauge-reduced principal symbol and cannot inherit
\eqref{eq:v3v12-stable-continuity} from this abstract model.
\end{firewall}

\begin{firewall}[Kreiss stability is classical]
\label{fw:v3v12-quantum}
The uniform Lopatinski margin controls a classical linear boundary symbol.
It does not construct the interacting quantum measure, prove BRST
renormalization, or establish reflection positivity.
\end{firewall}

\subsection{V12 conclusion and V13 target}
\label{sec:v3v12-conclusion}

V12 diagonalizes the general symmetric-hyperbolic normal flux and shows why
zero-speed traces require their own uniformly invertible constraint block.
After that block is eliminated, strict passivity of the physical dynamic
node excludes all closed-half-plane surface modes.  Continuity of the
stable bundle and compactness of normalized frequency space then give a
uniform dynamic Lopatinski margin.  The neutral-node example proves that
ordinary energy passivity alone does not give this margin.

The next upstream obstacle is glancing and variable multiplicity, where the
spectral stable bundle may not be continuous.  V13 will replace the raw
eigenvector bundle by a resolvent/symmetrizer formulation, prove a
frequency-local energy estimate under a Kreiss symmetrizer, and identify
the derivative loss when the zero-speed or glancing margin degenerates.
This is required before any variable-coefficient Einstein localization can
be justified.
\clearpage
\section{Third Volume, Version V13: Resolvent Symmetrizers and Quantified Glancing Loss}
\label{sec:v3v13}

\subsection{Replacing eigenvectors by an energy multiplier}
\label{sec:v3v13-purpose}

V12 obtained a uniform boundary margin by assuming that the stable
eigenbundle extends continuously to the closed normalized frequency
hemisphere.  At glancing or variable multiplicity, individual eigenvectors
may fail to provide such a bundle.  The correct replacement is an
operator-valued symmetrizer: a self-adjoint multiplier whose bulk real part
and boundary quadratic form remain coercive.  V13 proves the resolvent
estimate directly from those two inequalities and records exactly how a
degenerating coercivity constant becomes derivative or trace loss.

After Laplace transformation in time and Fourier transformation in the
tangential variables, write the normal problem as
\begin{align}
 \partial_xU(x)&=\mathcal G(\zeta)U(x)+F(x),
 \qquad x>0,
\label{eq:v3v13-normal}\\
 \mathcal B(\zeta)U(0)&=g,
\label{eq:v3v13-boundary}
\end{align}
where
\[
 \zeta=(s,\xi),\qquad \eta:=\operatorname{Re}s>0.
\]
The vector \(U\) may already include the V12 dynamic boundary state and the
uniformly eliminated zero-speed variables.

\subsection{A resolvent construction without eigenvectors}
\label{sec:v3v13-Lyapunov}

\begin{lemma}[Stable Lyapunov symmetrizer]
\label{lem:v3v13-Lyapunov}
Let \(G\) be a finite-dimensional matrix satisfying
\begin{equation}
 \|e^{xG}\|\le M e^{-\alpha x},
 \qquad x\ge0,
\label{eq:v3v13-semigroup-decay}
\end{equation}
for \(M\ge1\) and \(\alpha>0\).  Define
\begin{equation}
 R:=\int_0^\infty e^{xG^*}e^{xG}\,dx,
 \qquad S:=-R.
\label{eq:v3v13-RS}
\end{equation}
Then \(R>0\),
\begin{equation}
 G^*R+RG=-I,
 \qquad
 G^*S+SG=I,
\label{eq:v3v13-Lyapunov}
\end{equation}
and
\begin{equation}
 \|S\|=\|R\|\le\frac{M^2}{2\alpha}.
\label{eq:v3v13-S-bound}
\end{equation}
\end{lemma}

\begin{proof}
The decay hypothesis makes the integral norm convergent and gives
\eqref{eq:v3v13-S-bound}.  For \(v\ne0\), continuity at \(x=0\) makes
\(\int_0^\infty\|e^{xG}v\|^2dx>0\), so \(R>0\).  Differentiate
\(e^{xG^*}e^{xG}\), integrate from zero to infinity, and use decay:
\[
 G^*R+RG
 =\left[e^{xG^*}e^{xG}\right]_{x=0}^{x=\infty}
 =-I.
\]
The identity for \(S\) follows by changing sign.
\end{proof}

\begin{assessment}[Resolvent separation is the real input]
\label{ass:v3v13-resolvent-input}
Formula \eqref{eq:v3v13-RS} uses the stable semigroup as a whole.  It does
not select or normalize individual eigenvectors.  However, if the stable
decay rate \(\alpha\) collapses at glancing, the bound
\eqref{eq:v3v13-S-bound} degenerates.  A symmetrizer removes an artificial
eigenvector singularity; it cannot remove a genuine spectral collision.
\end{assessment}

\begin{proposition}[Riesz projector form]
\label{prop:v3v13-Riesz}
If the stable spectrum of \(\mathcal G(\zeta)\) is separated from the rest
by a contour \(\Gamma(\zeta)\), then its projector is
\begin{equation}
 P^s(\zeta)
 =\frac1{2\pi i}\int_{\Gamma(\zeta)}
  (zI-\mathcal G(\zeta))^{-1}\,dz.
\label{eq:v3v13-Riesz}
\end{equation}
Consequently a uniform contour length and a uniform resolvent bound give a
bounded projector without choosing eigenvectors.
\end{proposition}

\begin{proof}
This is the finite-dimensional holomorphic functional calculus.  On a
generalized eigenvector, the contour integral is one if its eigenvalue is
inside \(\Gamma\) and zero otherwise; linearity and the Jordan
decomposition prove the formula.  Taking norms under the integral gives the
last assertion.
\end{proof}

\subsection{Abstract Kreiss symmetrizer hypotheses}
\label{sec:v3v13-hypotheses}

Let \(\mathcal S(\zeta)=\mathcal S(\zeta)^*\) be a bounded matrix.  Assume
there are positive constants \(M_S,c_{\rm bulk},c_{\rm tr},C_{\rm B}\)
such that
\begin{align}
 \|\mathcal S(\zeta)\|&\le M_S,
\label{eq:v3v13-S-uniform}\\
 \operatorname{Re}
 \langle\mathcal S\mathcal Gv,v\rangle
 &\ge c_{\rm bulk}\eta\|v\|^2,
\label{eq:v3v13-bulk-coercivity}\\
 \langle\mathcal Sv,v\rangle
 &\ge c_{\rm tr}\|v\|^2
   -C_{\rm B}\|\mathcal Bv\|^2
\label{eq:v3v13-boundary-coercivity}
\end{align}
for every vector in the decaying normal solution space.  The last
inequality is the symmetrizer version of the V12 Lopatinski margin.

\begin{theorem}[Frequency-local Kreiss estimate]
\label{thm:v3v13-Kreiss}
Under \eqref{eq:v3v13-S-uniform}--%
\eqref{eq:v3v13-boundary-coercivity}, every decaying solution of
\eqref{eq:v3v13-normal}--\eqref{eq:v3v13-boundary} satisfies
\begin{equation}
 c_{\rm tr}\|U(0)\|^2
 +c_{\rm bulk}\eta\|U\|_{L^2_x}^2
 \le
 C_{\rm B}\|g\|^2
 \frac{M_S^2}{c_{\rm bulk}\eta}\|F\|_{L^2_x}^2.
\label{eq:v3v13-Kreiss}
\end{equation}
\end{theorem}

\begin{proof}
Differentiate the symmetrized normal energy:
\[
 \partial_x\langle\mathcal SU,U\rangle
 =2\operatorname{Re}\langle\mathcal S\mathcal GU,U\rangle
  +2\operatorname{Re}\langle\mathcal SF,U\rangle.
\]
Integration from zero to infinity and decay of \(U\) give
\begin{align*}
 -\langle\mathcal SU(0),U(0)\rangle
 &=
 2\int_0^\infty
  \operatorname{Re}\langle\mathcal S\mathcal GU,U\rangle dx\\
 &\quad
 +2\int_0^\infty
  \operatorname{Re}\langle\mathcal SF,U\rangle dx.
\end{align*}
Use \eqref{eq:v3v13-bulk-coercivity} and move the boundary term to the
other side.  The boundary inequality and
\(\mathcal BU(0)=g\) give
\[
 c_{\rm tr}\|U(0)\|^2
 +2c_{\rm bulk}\eta\|U\|_2^2
 \le C_{\rm B}\|g\|^2+2M_S\|F\|_2\|U\|_2.
\]
Young's inequality with parameter \(c_{\rm bulk}\eta\) absorbs one bulk
term and yields \eqref{eq:v3v13-Kreiss}.
\end{proof}

\begin{corollary}[Weighted physical-space estimate]
\label{cor:v3v13-Plancherel}
Suppose the symmetrizer constants are uniform in
\((\tau,\xi)\), where \(s=\eta+i\tau\), for a fixed \(\eta>0\).
Then Plancherel's theorem gives
\begin{align}
 &c_{\rm tr}
 \|e^{-\eta t}U|_{x=0}\|_{L^2_{t,y}}^2
 +c_{\rm bulk}\eta
 \|e^{-\eta t}U\|_{L^2_{t,x,y}}^2
\notag\\
 &\qquad
 \le C_{\rm B}
 \|e^{-\eta t}g\|_{L^2_{t,y}}^2
 \frac{M_S^2}{c_{\rm bulk}\eta}
 \|e^{-\eta t}F\|_{L^2_{t,x,y}}^2.
\label{eq:v3v13-physical}
\end{align}
\end{corollary}

\begin{proof}
Integrate \eqref{eq:v3v13-Kreiss} over \((\tau,\xi)\) and apply Plancherel
to the Laplace--Fourier transform.
\end{proof}

\subsection{Degenerate symmetrizers and derivative loss}
\label{sec:v3v13-loss}

Let
\[
 \langle\zeta\rangle
 :=(1+|s|^2+|\xi|^2)^{1/2}.
\]
Suppose instead of \eqref{eq:v3v13-boundary-coercivity} one has
\begin{equation}
 \langle\mathcal Sv,v\rangle
 \ge c_{\rm tr}\langle\zeta\rangle^{-2r}\|v\|^2
 -C_{\rm B}\|\mathcal Bv\|^2
\label{eq:v3v13-degenerate-boundary}
\end{equation}
for \(r\ge0\).

\begin{theorem}[Exact loss encoded by the boundary margin]
\label{thm:v3v13-derivative-loss}
Under \eqref{eq:v3v13-degenerate-boundary} and the uniform bulk hypotheses,
\begin{equation}
 \langle\zeta\rangle^{-2r}\|U(0)\|^2
 +\eta\|U\|_{L^2_x}^2
 \le C\left(
  \|g\|^2+\eta^{-1}\|F\|_{L^2_x}^2
 \right).
\label{eq:v3v13-loss-estimate}
\end{equation}
Equivalently, controlling the unweighted trace \(U(0)\) requires
\(r\) additional tangential derivatives of the boundary and source data.
\end{theorem}

\begin{proof}
Repeat the proof of Theorem~\ref{thm:v3v13-Kreiss} with
\eqref{eq:v3v13-degenerate-boundary}.  The first term on the left is the
square of the Fourier multiplier
\(\langle D_{t,y}\rangle^{-r}U(0)\).  Multiplying the transformed equation
by \(\langle\zeta\rangle^r\) shows that an unweighted trace estimate asks
for the \(H^r_{t,y}\) norm of the data.
\end{proof}

\begin{assessment}[A loss must be carried into the nonlinear ledger]
\label{ass:v3v13-loss-ledger}
If \(r>0\), the boundary estimate does not close at the same Sobolev order
as the bulk energy.  A nonlinear iteration must supply the missing
derivatives through additional regularity, a good-unknown cancellation, or
a Nash--Moser scheme.  Declaring the glancing margin to be a harmless
constant would create an unearned regularity gain.
\end{assessment}

\subsection{Explicit glancing regression}
\label{sec:v3v13-wave-regression}

Consider the scalar wave equation
\begin{equation}
 \partial_t^2u-\partial_x^2u-\Delta_yu=0,
 \qquad x>0.
\label{eq:v3v13-wave}
\end{equation}
For a Laplace--Fourier mode, the decaying solution is
\begin{equation}
 \widehat u(x)=a e^{-\kappa x},
 \qquad
 \kappa=(s^2+|\xi|^2)^{1/2},
 \qquad \operatorname{Re}\kappa>0.
\label{eq:v3v13-wave-mode}
\end{equation}

\begin{proposition}[Neumann trace degeneracy at glancing]
\label{prop:v3v13-Neumann-glancing}
For Neumann data
\(-\partial_x\widehat u(0)=g\),
\begin{equation}
 a=\kappa^{-1}g.
\label{eq:v3v13-Neumann-symbol}
\end{equation}
At a normalized glancing point
\(s=\eta+i\tau\), \(|\xi|=|\tau|\ne0\), one has
\begin{equation}
 |\kappa|
 =|\,\eta^2+2i\eta\tau\,|^{1/2}
 \asymp (|\tau|\eta)^{1/2}
\label{eq:v3v13-kappa-glancing}
\end{equation}
as \(\eta\downarrow0\).  Hence the boundary value satisfies
\begin{equation}
 |a|\asymp(|\tau|\eta)^{-1/2}|g|.
\label{eq:v3v13-half-loss}
\end{equation}
\end{proposition}

\begin{proof}
Differentiate \eqref{eq:v3v13-wave-mode} at \(x=0\) to obtain
\eqref{eq:v3v13-Neumann-symbol}.  When \(|\xi|=|\tau|\),
\[
 s^2+|\xi|^2
 =(\eta+i\tau)^2+\tau^2
 =\eta^2+2i\eta\tau.
\]
Taking the square root of its modulus proves
\eqref{eq:v3v13-kappa-glancing}, then
\eqref{eq:v3v13-half-loss}.
\end{proof}

\begin{firewall}[The wave equation remains well posed]
\label{fw:v3v13-Neumann}
The Neumann wave problem has a standard energy estimate.  The regression
shows that its boundary value is not uniformly controlled by Neumann data
at glancing; it does not claim bulk ill-posedness.  The quantity which loses
half an order must be stated explicitly.
\end{firewall}

\subsection{Block symmetrizers across a multiple root}
\label{sec:v3v13-blocks}

Suppose the normal matrix decomposes by resolvent projectors into separated
clusters
\[
 \mathcal H=\bigoplus_{\alpha=1}^N\mathcal H_\alpha
\]
which remain continuous even though eigenvalues inside one cluster cross.
Let \(\mathcal G_\alpha\) be the restriction to a cluster.

\begin{proposition}[Clusterwise construction]
\label{prop:v3v13-cluster}
If every stable cluster satisfies a semigroup estimate
\[
 \|e^{x\mathcal G_\alpha}\|
 \le M_\alpha e^{-\alpha_\alpha x}
\]
and the cluster projectors have uniformly bounded resolvents, then the
direct sum of the Lyapunov symmetrizers
\[
 \mathcal S_\alpha
 =-\int_0^\infty
 e^{x\mathcal G_\alpha^*}e^{x\mathcal G_\alpha}\,dx
\]
is bounded by
\begin{equation}
 \|\mathcal S\|
 \le
 C_{\rm proj}
 \max_\alpha\frac{M_\alpha^2}{2\alpha_\alpha}.
\label{eq:v3v13-cluster-bound}
\end{equation}
Eigenvalue crossings inside a separated stable cluster do not by themselves
destroy the construction.
\end{proposition}

\begin{proof}
Apply Lemma~\ref{lem:v3v13-Lyapunov} on each invariant cluster.  Conjugating
the block direct sum back to the original coordinates costs the projector
condition number \(C_{\rm proj}\).  No diagonalization inside a cluster is
used.
\end{proof}

\begin{counterexample}[A spectral collision still destroys the bound]
\label{cex:v3v13-collision}
For \(\mathcal G_\varepsilon=-\varepsilon I\), the stable semigroup has
\(\alpha=\varepsilon\), and
\[
 R_\varepsilon=\int_0^\infty e^{-2\varepsilon x}dx\,I
 =\frac1{2\varepsilon}I.
\]
Thus the Lyapunov symmetrizer norm diverges as the stable root reaches the
imaginary axis.  Grouping eigenvectors into a cluster cannot repair an
actual loss of normal decay.
\end{counterexample}

\subsection{Uniform continuum gate}
\label{sec:v3v13-uniform-gate}

\begin{theorem}[Regulator-uniform symmetrizer criterion]
\label{thm:v3v13-uniform}
Let a regulator family have symmetrizers \(\mathcal S_n\) satisfying
\eqref{eq:v3v13-S-uniform}--%
\eqref{eq:v3v13-boundary-coercivity} with the same positive constants.
Assume also the V12 zero-block margins satisfy
\(\rho_{0,n}\ge\rho_*>0\).  Then the transformed estimates
\eqref{eq:v3v13-Kreiss} and physical estimates
\eqref{eq:v3v13-physical} are uniform in \(n\).

If any one of
\[
 M_{S,n}^{-1},\quad c_{{\rm bulk},n},\quad
 c_{{\rm tr},n},\quad \rho_{0,n}
\]
tends to zero, the argument supplies no uniform continuum estimate.
\end{theorem}

\begin{proof}
The constants in Theorem~\ref{thm:v3v13-Kreiss} depend only on the four
displayed margins and \(C_{\rm B}\).  Uniformity therefore passes directly
through the proof and through Plancherel.  If a margin collapses, the same
formula displays the divergent inverse factor.
\end{proof}

\begin{firewall}[Existence of a symmetrizer remains gauge dependent]
\label{fw:v3v13-Einstein}
V13 proves what follows from a Kreiss symmetrizer and how to build one from
a uniformly stable cluster.  It does not construct such a symmetrizer for a
specific gauge-fixed Einstein--Standard Model principal symbol.  The gauge,
constraint additions, and junction variables must first be written
explicitly.
\end{firewall}

\subsection{V13 conclusion and V14 target}
\label{sec:v3v13-conclusion}

V13 replaces fragile eigenvector tracking by resolvent projectors and
Lyapunov/Kreiss multipliers.  Uniform bulk and boundary coercivity give the
frequency-local estimate and, by Plancherel, the weighted physical-space
estimate.  A boundary margin of order
\(\langle\zeta\rangle^{-2r}\) produces exactly \(r\) derivatives of trace
loss.  The Neumann wave regression exhibits the square-root glancing
degeneracy explicitly.

V14 will address variable coefficients.  It will quantize a smooth
frequency-dependent symmetrizer, compute the commutators with the evolving
Einstein coefficients, and give the Sobolev regularity threshold at which
the localized boundary estimate closes.  If the V13 loss \(r>0\) persists,
the note will state the resulting tame estimate instead of claiming an
ordinary Picard iteration.
\clearpage
\section{Third Volume, Version V14: Quantized Symmetrizers and a Tame Variable-Coefficient Boundary Estimate}
\label{sec:v3v14}

\subsection{Scope of the localization step}
\label{sec:v3v14-purpose}

V13 was a frozen-frequency theorem.  In a curved Einstein background, the
normal operator and boundary node depend on position and on the evolving
fields.  Quantizing a frequency-dependent symmetrizer produces commutators
and lower-order remainders.  V14 makes those costs explicit.  The conclusion
is conditional on a smooth uniform symmetrizer symbol for the chosen
gauge-reduced system; it does not assume that such a symbol has already been
constructed for Einstein--Standard Model variables.

Let \(Y\) denote the tangential spacetime variables after inserting a
Laplace weight \(e^{-\eta t}\), and let
\[
 \zeta=(\tau,\xi),\qquad
 \langle\zeta\rangle=(1+|\tau|^2+|\xi|^2)^{1/2}.
\]
The normal system on \(x>0\) is written
\begin{align}
 \partial_xU
 &=\operatorname{Op}(\mathcal G)U+F,
\label{eq:v3v14-normal}\\
 \operatorname{Op}(\mathcal B)U|_{x=0}
 &=g.
\label{eq:v3v14-boundary}
\end{align}
The vector \(U\) includes any dynamic boundary variables after the V12
zero-speed constraint has been eliminated with uniform margin.

\subsection{Symbol classes and three calculus estimates}
\label{sec:v3v14-calculus}

For a matrix symbol \(a(Y,\zeta)\), write \(a\in S^m\) when
\begin{equation}
 \|\partial_Y^\alpha\partial_\zeta^\beta a(Y,\zeta)\|
 \le C_{\alpha\beta}\langle\zeta\rangle^{m-|\beta|}
\label{eq:v3v14-symbol-class}
\end{equation}
for the finitely many derivatives used below.  Quantization is the
tangential Weyl quantization.

\begin{lemma}[Bounded order-zero quantization]
\label{lem:v3v14-CV}
If \(a\in S^0\) and its derivatives through an order
\(N>N_Y/2+1\) are bounded, then
\begin{equation}
 \|\operatorname{Op}(a)v\|_{L^2_Y}
 \le C_N
 \max_{|\alpha|+|\beta|\le N}C_{\alpha\beta}
 \|v\|_{L^2_Y}.
\label{eq:v3v14-CV}
\end{equation}
\end{lemma}

\begin{proof}
Decompose frequency space into unit cubes and insert a smooth partition of
unity.  On each cube, integrate the oscillatory kernel by parts
\(N\) times in \(\zeta\); the resulting kernel is bounded by
\(C(1+|Y-Y'|)^{-N}\).  The same bound holds after summing the finitely
overlapping cubes because the symbol has order zero.  Since
\((1+|Y|)^{-N}\) is integrable for \(N>N_Y\), Schur's kernel test proves
\eqref{eq:v3v14-CV}.  The stated larger-than-half-dimension threshold can
be obtained by the standard almost-orthogonality refinement; taking
\(N>N_Y+1\) also suffices for every application below.
\end{proof}

\begin{lemma}[One derivative gained by a coefficient commutator]
\label{lem:v3v14-commutator}
Let \(a\in S^0\) and let \(q=q(Y)\) be Lipschitz with the higher derivatives
needed by \eqref{eq:v3v14-symbol-class}.  Then
\begin{equation}
 [\operatorname{Op}(a),q]\in\operatorname{Op}(S^{-1})
\label{eq:v3v14-commutator-order}
\end{equation}
and
\begin{equation}
 \|[\operatorname{Op}(a),q]v\|_{L^2}
 \le C\|\nabla_Yq\|_{L^\infty}\|v\|_{H^{-1}}
 +C_{\rm high}(q)\|v\|_{H^{-2}}.
\label{eq:v3v14-commutator-bound}
\end{equation}
\end{lemma}

\begin{proof}
The commutator kernel contains
\[
 \bigl(q(Y')-q(Y)\bigr)
 e^{i(Y-Y')\cdot\zeta}
 a\left(\frac{Y+Y'}2,\zeta\right).
\]
Write
\[
 q(Y')-q(Y)
 =(Y'-Y)\cdot
 \int_0^1\nabla q(Y+\theta(Y'-Y))\,d\theta.
\]
Move the factor \(Y'-Y\) onto the exponential as
\(-i\partial_\zeta\).  One \(\zeta\)-derivative lowers the symbol order by
one.  Repeating the kernel estimate of
Lemma~\ref{lem:v3v14-CV} proves the leading \(H^{-1}\) bound; derivatives
falling on \(\nabla q\) produce the lower-order
\(C_{\rm high}(q)H^{-2}\) remainder.
\end{proof}

\begin{lemma}[Quantized positivity with a lower-order remainder]
\label{lem:v3v14-Garding}
If \(a=a^*\in S^m\) and
\begin{equation}
 a(Y,\zeta)\ge c\langle\zeta\rangle^m I,
\label{eq:v3v14-symbol-positive}
\end{equation}
then
\begin{equation}
 \operatorname{Re}\langle\operatorname{Op}(a)v,v\rangle
 \ge c\|v\|_{H^{m/2}}^2
 -C_a\|v\|_{H^{(m-1)/2}}^2.
\label{eq:v3v14-Garding}
\end{equation}
\end{lemma}

\begin{proof}
Write
\(a=c\langle\zeta\rangle^mI+b^*b\) at principal-symbol level using the
positive square root.  Symbol composition gives
\[
 \operatorname{Op}(b)^*\operatorname{Op}(b)
 =\operatorname{Op}(b^*b)+\operatorname{Op}(r),
\qquad r\in S^{m-1}.
\]
The square is nonnegative, and
Lemma~\ref{lem:v3v14-CV} applied after conjugation by
\(\langle D_Y\rangle^{(m-1)/2}\) bounds the remainder by the last term in
\eqref{eq:v3v14-Garding}.  A smooth regularization of the square root
handles semidefinite \(b^*b\), followed by a limit.
\end{proof}

\subsection{Smooth symmetrizer hypotheses}
\label{sec:v3v14-symmetrizer}

Assume a self-adjoint symbol
\(\mathcal S(Y,x,\zeta)\in S^0\) satisfies uniformly:
\begin{align}
 \|\mathcal S\|_{S^0_N}&\le M_N,
\label{eq:v3v14-S-seminorm}\\
 \operatorname{Re}
 \bigl(\mathcal S\mathcal G\bigr)
 &\ge c_{\rm bulk}\eta I
\quad\hbox{on the decaying normal bundle},
\label{eq:v3v14-symbol-bulk}\\
 \langle\mathcal Sv,v\rangle
 &\ge c_{\rm tr}\|v\|^2
 -C_{\rm B}\|\mathcal Bv\|^2
\quad\hbox{at }x=0.
\label{eq:v3v14-symbol-boundary}
\end{align}
Also assume
\begin{equation}
 \|\partial_x\mathcal S\|_{S^0_N}
 +\|\mathcal G\|_{S^1_N}
 +\|\mathcal B\|_{S^0_N}
 \le K_N.
\label{eq:v3v14-coefficient-size}
\end{equation}

\begin{proposition}[Quantized bulk coercivity]
\label{prop:v3v14-quantized-bulk}
There is \(C_0=C(M_N,K_N)\) such that
\begin{align}
 &2\operatorname{Re}
 \langle\operatorname{Op}(\mathcal S)
       \operatorname{Op}(\mathcal G)U,U\rangle
 +\langle\operatorname{Op}(\partial_x\mathcal S)U,U\rangle
\notag\\
 &\qquad
 \ge (2c_{\rm bulk}\eta-C_0)\|U\|_{L^2_Y}^2
 -C_0\|U\|_{H^{-1/2}_Y}^2.
\label{eq:v3v14-quantized-bulk}
\end{align}
\end{proposition}

\begin{proof}
The self-adjoint part of the product has principal symbol
\(\operatorname{Re}(\mathcal S\mathcal G)\).  Weyl composition differs
from that symbol by an order-zero term involving one \(Y\)- or
\(\zeta\)-derivative of the factors.  Apply
Lemma~\ref{lem:v3v14-Garding} to the positive principal part and
Lemma~\ref{lem:v3v14-CV} to the order-zero composition and
\(\partial_x\mathcal S\) terms.  The harmless
\(H^{-1/2}\) term records the sharp-Garding remainder; enlarging \(C_0\)
also bounds every order-zero remainder.
\end{proof}

\begin{proposition}[Quantized boundary coercivity]
\label{prop:v3v14-quantized-boundary}
There is \(C_1=C(M_N,K_N)\) such that
\begin{align}
 \langle\operatorname{Op}(\mathcal S) v,v\rangle
 &\ge \frac{c_{\rm tr}}2\|v\|_{L^2_Y}^2
 -C_{\rm B}'\|\operatorname{Op}(\mathcal B)v\|_{L^2_Y}^2
\notag\\
 &\quad
 -C_1\|v\|_{H^{-1/2}_Y}^2.
\label{eq:v3v14-quantized-boundary}
\end{align}
\end{proposition}

\begin{proof}
Quantize \eqref{eq:v3v14-symbol-boundary}.  The positive part is controlled
by Lemma~\ref{lem:v3v14-Garding}.  Symbol composition replaces
\(\operatorname{Op}(\mathcal B^*\mathcal B)\) by
\(\operatorname{Op}(\mathcal B)^*\operatorname{Op}(\mathcal B)\) plus an
operator of order \(-1\).  Bound that remainder by
\(\|v\|_{H^{-1/2}}^2\), and absorb a fixed fraction of the trace coefficient
into \(c_{\rm tr}/2\).
\end{proof}

\subsection{High-frequency absorption and the low-frequency gate}
\label{sec:v3v14-high-low}

Let \(\chi_R(D_Y)\) project to
\(\langle\zeta\rangle\ge R\).  Then
\begin{equation}
 \|\chi_Rv\|_{H^{-1/2}}
 \le R^{-1/2}\|\chi_Rv\|_{L^2}.
\label{eq:v3v14-high-absorption}
\end{equation}
Choose \(R\) so large that the last term in
\eqref{eq:v3v14-quantized-boundary} is at most
\((c_{\rm tr}/4)\|\chi_Rv\|^2\).  The complementary frequency ball is
compact.  Its required hypothesis is the finite-frequency resolvent margin
\begin{equation}
 \|\mathcal L_{\rm low}(\zeta)q\|
 \ge\gamma_{\rm low}\|q\|,
 \qquad
 \langle\zeta\rangle\le2R,
\label{eq:v3v14-low-margin}
\end{equation}
where \(\mathcal L_{\rm low}\) is the complete bulk--dynamic-boundary
symbol, including the V12 zero block.

\begin{firewall}[High-frequency symmetrization does not settle low modes]
\label{fw:v3v14-low}
The Gårding remainder is absorbable at high frequency.  It is not small on
a fixed compact frequency ball.  The separate margin
\(\gamma_{\rm low}>0\) must be checked; omitting it could hide a neutral
boundary state of the kind exhibited in V12.
\end{firewall}

\subsection{Variable-coefficient resolvent estimate}
\label{sec:v3v14-resolvent}

\begin{theorem}[Quantized Kreiss estimate]
\label{thm:v3v14-quantized-Kreiss}
Assume \eqref{eq:v3v14-S-seminorm}--%
\eqref{eq:v3v14-coefficient-size}, the uniform V12 zero-block margin, and
the low-frequency margin \eqref{eq:v3v14-low-margin}.  There is
\(\eta_0=\eta_0(M_N,K_N,c_{\rm bulk})\) such that for
\(\eta\ge\eta_0\), every decaying solution of
\eqref{eq:v3v14-normal}--\eqref{eq:v3v14-boundary} satisfies
\begin{equation}
 \|U(0)\|_{L^2_Y}^2
 +\eta\|U\|_{L^2_{x,Y}}^2
 \le C\left(
  \|g\|_{L^2_Y}^2
  +\eta^{-1}\|F\|_{L^2_{x,Y}}^2
 \right),
\label{eq:v3v14-quantized-Kreiss}
\end{equation}
where \(C\) depends only on the displayed uniform margins and symbol
seminorms.
\end{theorem}

\begin{proof}
Differentiate
\(\langle\operatorname{Op}(\mathcal S)U,U\rangle_{L^2_Y}\) in \(x\) and
use \eqref{eq:v3v14-normal}.  Proposition
\ref{prop:v3v14-quantized-bulk} supplies
\((2c_{\rm bulk}\eta-C_0)\|U\|^2\).  Choose
\(\eta_0\ge C_0/c_{\rm bulk}+1\), then absorb the order-zero bulk
remainder.  Cauchy--Schwarz and Young absorb the source term with cost
\(C\eta^{-1}\|F\|^2\).

At the boundary, Proposition~\ref{prop:v3v14-quantized-boundary},
\eqref{eq:v3v14-high-absorption}, and the choice of \(R\) control the high
frequencies.  The compact low-frequency block is controlled by
\eqref{eq:v3v14-low-margin}.  Commutators of the high/low cutoffs with the
coefficients are of order \(-1\) by
Lemma~\ref{lem:v3v14-commutator}; enlarge \(\eta_0\) once to absorb them.
Adding both frequency pieces and using
\(\operatorname{Op}(\mathcal B)U(0)=g\) proves
\eqref{eq:v3v14-quantized-Kreiss}.
\end{proof}

\subsection{High-order tame estimate}
\label{sec:v3v14-tame}

Let \(s_0>N_Y/2+1\).  For \(s\ge s_0\), let \(K_s(t)\) denote the
\(H^s\) size of the coefficients and symmetrizer symbols, and let
\(K_{s_0}(t)\) denote their Lipschitz-level size.

\begin{lemma}[Tame commutator form]
\label{lem:v3v14-tame-commutator}
For a first-order tangential operator \(P(a)\) with coefficient \(a\),
\begin{equation}
 \|[\Lambda^s,P(a)]u\|_{L^2}
 \le C\left(
  \|a\|_{W^{1,\infty}}\|u\|_{H^s}
  +\|a\|_{H^{s+1}}\|u\|_{H^{s_0}}
 \right).
\label{eq:v3v14-tame-commutator}
\end{equation}
\end{lemma}

\begin{proof}
Write the Fourier transform of the commutator.  Its multiplier contains
\(\langle\xi\rangle^s-\langle\xi-\eta\rangle^s\), which by the mean-value
theorem is bounded by
\[
 C|\eta|
 \bigl(\langle\xi-\eta\rangle^{s-1}
       +\langle\eta\rangle^{s-1}\bigr).
\]
Convolution, Sobolev embedding at \(s_0>N_Y/2+1\), and distribution of the
high frequency to either \(u\) or \(a\) give
\eqref{eq:v3v14-tame-commutator}.
\end{proof}

\begin{theorem}[Tame localized boundary estimate]
\label{thm:v3v14-tame}
Under the hypotheses of Theorem~\ref{thm:v3v14-quantized-Kreiss}, suppose
the symbol seminorms are finite through order \(s+1\).  Then the
time-dependent physical solution satisfies an estimate of the form
\begin{align}
 E_s(t)^{1/2}
 &\le
 C e^{C\int_0^t(1+K_{s_0}(\tau))d\tau}
 \Bigl[
  E_s(0)^{1/2}
  \|f\|_{L^1_tH^s}
  \|g\|_{H^s_{t,Y}}
\notag\\
 &\hspace{38mm}
  +\int_0^tK_s(\tau)E_{s_0}(\tau)^{1/2}d\tau
 \Bigr].
\label{eq:v3v14-tame}
\end{align}
The high coefficient norm multiplies only the low solution norm.
\end{theorem}

\begin{proof}
Apply \(\Lambda^s\) to the localized system and use
Theorem~\ref{thm:v3v14-quantized-Kreiss} on each time slab.  Principal
terms reproduce the \(s=0\) estimate.  Lemma
\ref{lem:v3v14-tame-commutator} bounds every differentiated-coefficient
remainder by
\[
 C K_{s_0}E_s^{1/2}+C K_sE_{s_0}^{1/2}.
\]
The first term enters the Gronwall coefficient; the second is retained as
the tame source.  Sum a finite partition of unity, including its
order-zero commutators, and apply Gronwall's inequality.  Trace terms are
controlled by the same quantized boundary estimate, giving
\eqref{eq:v3v14-tame}.
\end{proof}

\subsection{What happens when V13 has loss}
\label{sec:v3v14-loss}

\begin{corollary}[Tame estimate with \(r\)-derivative boundary loss]
\label{cor:v3v14-loss}
If the symbolic boundary coercivity is only
\[
 \langle\mathcal Sv,v\rangle
 \ge c\langle\zeta\rangle^{-2r}\|v\|^2
 -C\|\mathcal Bv\|^2,
\]
then the same argument estimates \(E_s\) using boundary data at order
\(s+r\) and coefficient commutators through order \(s+r+1\).
\end{corollary}

\begin{proof}
Conjugate the boundary equation by
\(\Lambda^{s+r}\) and use the V13 derivative-loss theorem.  Every
commutator is then evaluated at the shifted order.  The tame distribution
of high derivatives is unchanged, but the top boundary norm is
\(H^{s+r}\).
\end{proof}

\begin{assessment}[Picard versus Nash--Moser]
\label{ass:v3v14-Nash-Moser}
When \(r=0\), estimate \eqref{eq:v3v14-tame} has the form required by an
ordinary quasilinear energy iteration, subject to compatibility conditions.
When \(r>0\), the linear solution map loses derivatives of the boundary
data.  This estimate alone cannot make a contraction on one fixed Sobolev
space; a structural cancellation or a smoothing iteration is required.
\end{assessment}

\subsection{A coefficient regression}
\label{sec:v3v14-coefficient-regression}

\begin{counterexample}[Small amplitude does not make the commutator small]
\label{cex:v3v14-small-amplitude}
On \(S^1\), let
\[
 q_N(y):=N^{-1}\sin(Ny),\qquad u(y):=1,
\]
and let \(\Lambda=|D_y|\).  Then
\begin{equation}
 \|q_N\|_{L^\infty}=N^{-1}\longrightarrow0,
 \qquad
 [\Lambda,q_N]u=\Lambda q_N=\sin(Ny)
\label{eq:v3v14-commutator-counterexample}
\end{equation}
up to the fixed Fourier normalization.  Hence the commutator remains order
one.  Uniform \(L^\infty\) closeness of the coefficients to a frozen value
does not replace control of their derivatives.
\end{counterexample}

\begin{firewall}[No pseudodifferential import into the quantum measure]
\label{fw:v3v14-quantum}
The symmetrizer is a classical energy multiplier.  It is unrelated to the
Faddeev--Popov determinant or to a quantum effective action unless an
additional construction proves such a relation.
\end{firewall}

\begin{firewall}[Boundary compatibility conditions remain]
\label{fw:v3v14-compatibility}
Estimate \eqref{eq:v3v14-tame} assumes data in the domain of the dynamic
boundary operator.  At high regularity, the initial data must satisfy the
boundary equations and their time derivatives to the required order.
The estimate does not manufacture those corner compatibility conditions.
\end{firewall}

\subsection{V14 conclusion and V15 audit mandate}
\label{sec:v3v14-conclusion}

V14 quantizes the V13 symmetrizer.  Symbol composition and metric motion
create lower-order commutators, absorbed at large Laplace weight and high
frequency; a separate compact low-frequency margin excludes neutral modes.
The high-order estimate is tame: high coefficient norms multiply only the
low solution norm.  A surviving glancing loss shifts the required boundary
regularity by exactly \(r\), and the oscillatory coefficient regression
proves why derivative control cannot be replaced by amplitude smallness.

V15 is the next mandatory companion audit.  It will fingerprint the newest
Yang--Mills and Navier--Stokes notes before choosing between two concrete
upstream tasks: writing an explicit gauge-reduced linearized Einstein
boundary symbol to test the V12--V14 gates, or returning to the BRST
boundary complex to make the dynamic junction node compatible with the
quantum Ward differential.  No result from the companion programmes will
be imported without a type-correct map.
\clearpage
\section{Third Volume, Version V15: Companion Audit and the Explicit Linearized Harmonic--Einstein Junction Block}
\label{sec:v3v15}

\subsection{Mandatory fifteen-version companion audit}
\label{sec:v3v15-audit}

The scheduled audit was completed before the V15 construction.  The
read-only companion inventory was:
\begin{center}
\begin{tabular}{p{0.48\textwidth}r p{0.33\textwidth}}
file & bytes & SHA--256\\
\hline
\texttt{Renewal\_NS\_Stability\_Research\_Notes\_V31.tex}
&887537&
\texttt{F1121B62238AD5491BB7CF03635EA9934942EDF573ED8FAAA9EE79E1D38A6696}
\\
\texttt{Renewal\_Yang\_Mills\_Mass\_Gap\_Research\_Notes\_V81.tex}
&1243239&
\texttt{1610676A7CABAC1C83CE05CF1DEFAD8A01FC463AD0F11E6F735AEE6D6A4D9C34}
\\
\texttt{Renewal\_Yang\_Mills\_Mass\_Gap\_Research\_Notes\_V82.tex}
&1273518&
\texttt{69729185D276766616F80A8FBB336D49BCCF7C2C1B30BDEF1F3E2D124BEDC2A6}
\end{tabular}
\end{center}
The Navier--Stokes state is unchanged.  Yang--Mills V82 proves the
previously targeted two-endpoint bounds when both gains are backed by one
actual common incidence-two coordinate subbank.  The selected subbank is
extracted by an exact covariance identity before a positive envelope is
taken; the remainder is then a contraction or the unique payer.  The
result is deliberately not assigned to an arbitrary scalar row-pair card.
Higher-incidence, mismatched-partner, overlap, and continuum sectors remain
open.  No companion file was modified.

\begin{assessment}[V15 transfer decision]
\label{ass:v3v15-audit-decision}
For the Einstein boundary symbol, the corresponding actual physical block
is the six-dimensional tangential symmetric tensor on which the
Brown--York momentum acts.  The four harmonic constraints supply the
remaining normal-index equations.  V15 forms those two geometric blocks
before taking characteristic norms.  It does not distribute ten abstract
scalar boundary conditions among metric components by a counting
argument.
\end{assessment}

\begin{firewall}[Strict audit type boundary]
\label{fw:v3v15-audit-types}
The Yang--Mills subbank theorem is a finite-cutoff Wilson covariance
estimate.  The calculation below is a flat-background differential-symbol
calculation.  Their only shared rule is to identify and extract the genuine
operator block before applying a norm.
\end{firewall}

\subsection{Linearized harmonic Einstein equations}
\label{sec:v3v15-harmonic}

Work in four-dimensional Minkowski half-space
\begin{equation}
 M=\{(t,x,y^1,y^2):x>0\},
 \qquad
 \eta=\operatorname{diag}(-1,1,1,1).
\label{eq:v3v15-halfspace}
\end{equation}
Indices \(a,b\) are tangent to the boundary
\(\mathcal N=\{x=0\}\), and Greek indices are spacetime indices.  Let
\(h_{\mu\nu}\) be a metric perturbation and define its trace reverse
\begin{equation}
 \bar h_{\mu\nu}
 :=h_{\mu\nu}-\frac12\eta_{\mu\nu}h,
 \qquad
 h:=\eta^{\alpha\beta}h_{\alpha\beta}.
\label{eq:v3v15-trace-reverse}
\end{equation}
The linear harmonic constraint is
\begin{equation}
 C_\mu:=\partial^\nu\bar h_{\mu\nu}.
\label{eq:v3v15-C}
\end{equation}

\begin{proposition}[Reduced field equation]
\label{prop:v3v15-reduced-Einstein}
The linearized Einstein tensor is
\begin{equation}
 G^{(1)}_{\mu\nu}
 =-\frac12\Box\bar h_{\mu\nu}
 +\partial_{(\mu}C_{\nu)}
 -\frac12\eta_{\mu\nu}\partial^\alpha C_\alpha.
\label{eq:v3v15-linear-Einstein}
\end{equation}
In harmonic gauge \(C_\mu=0\), the sourced Einstein equation becomes
\begin{equation}
 \Box\bar h_{\mu\nu}=-2\kappa T_{\mu\nu}^{(1)}.
\label{eq:v3v15-reduced-wave}
\end{equation}
\end{proposition}

\begin{proof}
The variations of the Ricci tensor and scalar about flat space are
\begin{align*}
 R^{(1)}_{\mu\nu}
 &=\frac12\left(
  \partial^\alpha\partial_\mu h_{\alpha\nu}
 +\partial^\alpha\partial_\nu h_{\alpha\mu}
 -\Box h_{\mu\nu}
 -\partial_\mu\partial_\nu h
 \right),\\
 R^{(1)}
 &=\partial^\mu\partial^\nu h_{\mu\nu}-\Box h.
\end{align*}
Substitute \(h_{\mu\nu}=\bar h_{\mu\nu}
-\frac12\eta_{\mu\nu}\bar h\), collect the divergence
\(C_\mu\), and subtract
\(\frac12\eta_{\mu\nu}R^{(1)}\).  This gives
\eqref{eq:v3v15-linear-Einstein}.  Setting \(C=0\) yields
\eqref{eq:v3v15-reduced-wave}.
\end{proof}

\begin{theorem}[Constraint propagation]
\label{thm:v3v15-constraint-propagation}
If the source is conserved,
\begin{equation}
 \partial^\nu T_{\mu\nu}^{(1)}=0,
\label{eq:v3v15-source-conservation}
\end{equation}
then every solution of \eqref{eq:v3v15-reduced-wave} satisfies
\begin{equation}
 \Box C_\mu=0.
\label{eq:v3v15-C-wave}
\end{equation}
Consequently zero initial constraint data and a homogeneous
constraint-preserving incoming boundary condition imply \(C_\mu=0\).
\end{theorem}

\begin{proof}
Apply \(\partial^\nu\) to \eqref{eq:v3v15-reduced-wave} and commute the
constant-coefficient derivatives:
\[
 \Box(\partial^\nu\bar h_{\mu\nu})
 =-2\kappa\partial^\nu T_{\mu\nu}^{(1)}=0.
\]
The uniqueness conclusion follows from the scalar wave energy estimate for
each component of \(C\).
\end{proof}

\subsection{Linearized Brown--York block}
\label{sec:v3v15-BY}

The background boundary is the plane \(x=0\).  For this local symbol
calculation choose \(n=+\partial_x\).  The geometric outward normal for the
half-space \(x>0\) is the opposite choice and reverses every Brown--York
sign without changing any rank or margin below.  The background extrinsic
curvature vanishes.  Tangential derivatives are denoted by \(\partial_a\).

\begin{lemma}[Variation of the extrinsic curvature]
\label{lem:v3v15-dK}
At the flat boundary,
\begin{equation}
 K^{(1)}_{ab}
 =\frac12\left(
  \partial_xh_{ab}
 -\partial_ah_{xb}
 -\partial_bh_{xa}
 \right),
\label{eq:v3v15-dK}
\end{equation}
and
\begin{equation}
 K^{(1)}
 =\frac12\left(
  \partial_xh^c{}_c-2\partial^ch_{xc}
 \right).
\label{eq:v3v15-dKtrace}
\end{equation}
\end{lemma}

\begin{proof}
The background normal covector is \(dx\).  To first order, tangential
projection removes the normal component, and
\[
 K^{(1)}_{ab}=-\Gamma^{x(1)}_{ab}.
\]
The linearized Christoffel symbol is
\[
 \Gamma^{x(1)}_{ab}
 =\frac12\left(
  \partial_ah^x{}_b+\partial_bh^x{}_a-\partial^xh_{ab}
 \right).
\]
Since the \(x\) direction is spacelike, this gives
\eqref{eq:v3v15-dK}.  Contract with the flat induced metric to obtain
\eqref{eq:v3v15-dKtrace}.
\end{proof}

\begin{proposition}[Explicit linearized junction operator]
\label{prop:v3v15-linear-junction}
The linearized Brown--York tensor is
\begin{align}
 \tau_{ab}^{(1)}
 &=\frac1{2\kappa}
 \Bigl[
  \partial_xh_{ab}
 -\partial_ah_{xb}
 -\partial_bh_{xa}
\notag\\
 &\hspace{18mm}
 -\gamma_{ab}\bigl(
  \partial_xh^c{}_c-2\partial^ch_{xc}
 \bigr)
 \Bigr].
\label{eq:v3v15-linear-BY}
\end{align}
The variational boundary equation is
\begin{equation}
 \tau_{ab}^{(1)}=S_{ab}^{(1)}.
\label{eq:v3v15-linear-junction}
\end{equation}
\end{proposition}

\begin{proof}
Linearize
\(\tau_{ab}^{\rm BY}=\kappa^{-1}(K_{ab}-K\gamma_{ab})\).
The background \(K\) vanishes, so the variation of \(\gamma_{ab}\) in the
second term contributes nothing at first order.  Insert
\eqref{eq:v3v15-dK}--\eqref{eq:v3v15-dKtrace}.
\end{proof}

\subsection{The tangential six-block is invertible}
\label{sec:v3v15-six-block}

Let \(\operatorname{Sym}^2(T^*\mathcal N)\) have dimension six.  Define
\begin{equation}
 (\mathcal Qk)_{ab}
 :=k_{ab}-\gamma_{ab}\operatorname{tr}_\gamma k.
\label{eq:v3v15-Q}
\end{equation}

\begin{lemma}[Inverse of the Brown--York trace map]
\label{lem:v3v15-Qinverse}
In boundary dimension three, \(\mathcal Q\) is invertible and
\begin{equation}
 (\mathcal Q^{-1}q)_{ab}
 =q_{ab}-\frac12\gamma_{ab}\operatorname{tr}_\gamma q.
\label{eq:v3v15-Qinverse}
\end{equation}
\end{lemma}

\begin{proof}
Taking the boundary trace of \(q=\mathcal Qk\) gives
\[
 \operatorname{tr}_\gamma q
 =(1-3)\operatorname{tr}_\gamma k
 =-2\operatorname{tr}_\gamma k.
\]
Thus
\[
 k=q+\gamma\operatorname{tr}_\gamma k
 =q-\frac12\gamma\operatorname{tr}_\gamma q.
\]
This is \eqref{eq:v3v15-Qinverse}.
\end{proof}

\begin{corollary}[Junction equation solves six normal derivatives]
\label{cor:v3v15-six-normal}
For prescribed \(S_{ab}^{(1)}\) and tangential derivatives, the six junction
equations uniquely determine \(\partial_xh_{ab}\).  Explicitly, set
\begin{equation}
 q_{ab}:=
 2\kappa S_{ab}^{(1)}
 +\partial_ah_{xb}+\partial_bh_{xa}
 -2\gamma_{ab}\partial^ch_{xc}.
\label{eq:v3v15-qdata}
\end{equation}
Then
\begin{equation}
 \partial_xh_{ab}
 =q_{ab}-\frac12\gamma_{ab}\operatorname{tr}_\gamma q.
\label{eq:v3v15-six-solution}
\end{equation}
\end{corollary}

\begin{proof}
Rearrange \eqref{eq:v3v15-linear-BY} as
\(\mathcal Q(\partial_xh)_{ab}=q_{ab}\), then apply
Lemma~\ref{lem:v3v15-Qinverse}.
\end{proof}

\subsection{The four harmonic constraints complete the normal block}
\label{sec:v3v15-four-block}

At the boundary,
\begin{equation}
 C_\mu
 =-\partial_t\bar h_{\mu t}
  +\partial_x\bar h_{\mu x}
  +\partial_A\bar h_{\mu A},
\qquad A=1,2.
\label{eq:v3v15-C-boundary}
\end{equation}

\begin{theorem}[Exact six-plus-four normal-derivative resolution]
\label{thm:v3v15-ten-block}
The six equations
\(\tau_{ab}^{(1)}=S_{ab}^{(1)}\) and the four equations
\(C_\mu=c_\mu\) uniquely determine all ten normal derivatives
\begin{equation}
 \{\partial_xh_{\mu\nu}:\mu\le\nu\}
\label{eq:v3v15-ten-derivatives}
\end{equation}
from tangential derivatives and the data \((S^{(1)},c)\).
\end{theorem}

\begin{proof}
Corollary~\ref{cor:v3v15-six-normal} first determines the six quantities
\(\partial_xh_{ab}\).  For a tangential index \(a\),
\[
 \partial_x\bar h_{ax}=\partial_xh_{ax},
\]
so the three equations \(C_a=c_a\) determine
\(\partial_xh_{ax}\).

For \(\mu=x\),
\[
 \partial_x\bar h_{xx}
 =\partial_xh_{xx}
 -\frac12\partial_x
  (h_{xx}+h^a{}_a)
 =\frac12\left(
  \partial_xh_{xx}-\partial_xh^a{}_a
 \right).
\]
The tangential trace derivative is already known, so \(C_x=c_x\)
determines \(\partial_xh_{xx}\).  The solution is triangular and unique.
\end{proof}

\begin{corollary}[Uniform frozen algebraic margin]
\label{cor:v3v15-algebraic-margin}
On a compact set of background boundary frames with uniformly
nondegenerate induced metric, the inverse of the six-plus-four normal map
is uniformly bounded.
\end{corollary}

\begin{proof}
In an orthonormal frame the inverse consists only of
\eqref{eq:v3v15-Qinverse}, three identity solves, and the scalar factor two
in the last equation.  Change of frame on a compact uniformly
nondegenerate set has uniformly bounded condition number.
\end{proof}

\begin{assessment}[This is the V12 zero-block acquisition]
\label{ass:v3v15-zero-acquisition}
Theorem~\ref{thm:v3v15-ten-block} supplies an explicit algebraic normal
resolution for the flat harmonic Einstein symbol.  It is the concrete
counterpart of the abstract V12 zero-speed inverse \(R_0^{-1}\).  The
calculation fixes the actual six-dimensional Brown--York block before the
four constraint equations are used.
\end{assessment}

\subsection{Characteristic form of the constraint equations}
\label{sec:v3v15-characteristics}

For each tensor component define
\begin{equation}
 W_{\mu\nu}^{\rm in}
 :=\partial_th_{\mu\nu}-\partial_xh_{\mu\nu},
 \qquad
 W_{\mu\nu}^{\rm out}
 :=\partial_th_{\mu\nu}+\partial_xh_{\mu\nu}.
\label{eq:v3v15-characteristics}
\end{equation}
Then
\[
 \partial_t=\frac12(W^{\rm in}+W^{\rm out}),
 \qquad
 \partial_x=\frac12(W^{\rm out}-W^{\rm in})
\]
componentwise.

\begin{proposition}[Four incoming constraint combinations]
\label{prop:v3v15-incoming-constraints}
The boundary equations \(C_\mu=c_\mu\) determine four independent linear
combinations of the incoming trace-reversed characteristics.  Define
\(\bar W_{\mu\nu}^{\rm in/out}\) from \(\bar h\) as in
\eqref{eq:v3v15-characteristics}.  Then
\begin{equation}
 \bar W_{\mu t}^{\rm in}+\bar W_{\mu x}^{\rm in}
 =
 -\bar W_{\mu t}^{\rm out}+\bar W_{\mu x}^{\rm out}
 +2\partial_A\bar h_{\mu A}-2c_\mu.
\label{eq:v3v15-incoming-constraint}
\end{equation}
\end{proposition}

\begin{proof}
Substitute the characteristic identities into
\eqref{eq:v3v15-C-boundary} and multiply by two.  The map
\[
 X_{\mu\nu}\longmapsto X_{\mu t}+X_{\mu x}
\]
from symmetric tensors to covectors has rank four: if its value is
prescribed, choose only the components \(X_{\mu t}\), resolving the
symmetry at the \(t,x\) intersection in any fixed complementary basis.
Thus the four displayed combinations are independent.
\end{proof}

\begin{firewall}[Rank is not yet the Lopatinski margin]
\label{fw:v3v15-rank-not-Lopatinski}
Theorem~\ref{thm:v3v15-ten-block} proves invertibility with respect to
normal derivatives at a frozen boundary point.  Tangential
Laplace--Fourier coupling can still create glancing or neutral surface
modes.  The V12--V14 Kreiss gates must be tested on the complete symbol.
\end{firewall}

\subsection{The boundary stress must come from one Hessian}
\label{sec:v3v15-Hessian}

Let \(q_{ab}=h_{ab}|_{\mathcal N}\) and let \(z\) be a linear boundary
inflow field.  A quadratic boundary action has the form
\begin{equation}
 S_\partial^{(2)}[q,z]
 =\frac12\langle q,H_{qq}q\rangle
 +\operatorname{Re}\langle Lq,z\rangle
 +\frac12\langle z,H_{zz}z\rangle,
\label{eq:v3v15-boundary-Hessian}
\end{equation}
where the pairings include tangential integration.

\begin{proposition}[Adjoint metric--inflow blocks]
\label{prop:v3v15-Hessian-adjoint}
The linearized boundary stress contribution and the boundary field equation
contain the cross operators
\begin{equation}
 \frac{\delta S_\partial^{(2)}}{\delta q}
 \supset L^*z,
 \qquad
 \frac{\delta S_\partial^{(2)}}{\delta z}
 \supset Lq.
\label{eq:v3v15-Hessian-adjoint}
\end{equation}
Thus the metric junction source and inflow forcing are adjoints before any
positive envelope or energy norm is chosen.
\end{proposition}

\begin{proof}
Differentiate the single bilinear term
\(\operatorname{Re}\langle Lq,z\rangle\) in each argument.
\end{proof}

\begin{assessment}[Minimal scalar matter around the zero field]
\label{ass:v3v15-scalar-background}
The minimally coupled scalar stress is quadratic in \(\phi\).  Linearizing
about \(\phi=0\) gives \(S_{ab}^{\phi(1)}=0\), so it does not alter the
linear boundary symbol.  A nontrivial linear dynamic junction block
requires a nonzero background or a quadratic cross interaction such as
\eqref{eq:v3v15-boundary-Hessian}.  V15 keeps this distinction explicit.
\end{assessment}

\subsection{Linearized Codazzi compatibility}
\label{sec:v3v15-Codazzi}

\begin{proposition}[Divergence of the linearized junction]
\label{prop:v3v15-linear-Codazzi}
The flat-boundary Brown--York tensor satisfies
\begin{equation}
 \partial^a\tau_{ab}^{(1)}
 =\kappa^{-1}G_{xb}^{(1)}.
\label{eq:v3v15-linear-Codazzi}
\end{equation}
Consequently a linear junction source must obey
\begin{equation}
 \partial^aS_{ab}^{(1)}=T_{xb}^{(1)}
\label{eq:v3v15-linear-flux}
\end{equation}
on solutions of the linearized Einstein equation.
\end{proposition}

\begin{proof}
Linearize the exact contracted Codazzi identity of V10 about the flat plane,
where all background curvature and extrinsic curvature terms vanish.  Then
use \(G_{xb}^{(1)}=\kappa T_{xb}^{(1)}\).
\end{proof}

\begin{counterexample}[Six arbitrary junction functions are not free data]
\label{cex:v3v15-arbitrary-junction}
Choose \(S_{ab}^{(1)}\) with
\(\partial^aS_{ab}^{(1)}\ne T_{xb}^{(1)}\).  Although
Corollary~\ref{cor:v3v15-six-normal} algebraically solves six normal
derivatives at a time slice, no solution of the full linearized Einstein
equations can realize those data.  The four divergence relations are the
boundary constraints on the six junction components.
\end{counterexample}

\subsection{A surviving glancing regression}
\label{sec:v3v15-glancing}

\begin{proposition}[The tangential trace sector contains a Neumann-type
subsymbol]
\label{prop:v3v15-Neumann-subsymbol}
Set the normal-index perturbations and tangential derivatives to zero and
choose a tangential trace-free component \(h_{ab}^{\rm TF}\).  With
\(S_{ab}^{(1)}=0\), the linearized junction equation reduces to
\begin{equation}
 \partial_xh_{ab}^{\rm TF}=0.
\label{eq:v3v15-Neumann-subsymbol}
\end{equation}
That component obeys the scalar wave equation in harmonic reduction and
therefore has the Neumann glancing trace degeneracy computed in V13.
\end{proposition}

\begin{proof}
In \eqref{eq:v3v15-linear-BY}, the terms containing \(h_{xa}\) vanish by
the restriction, and the trace term vanishes on the tangential trace-free
component.  The remaining equation is
\eqref{eq:v3v15-Neumann-subsymbol}.  Equation
\eqref{eq:v3v15-reduced-wave} acts componentwise on the flat background.
The V13 Laplace--Fourier calculation applies.
\end{proof}

\begin{assessment}[Algebraic closure does not remove glancing]
\label{ass:v3v15-glancing}
The explicit ten-block resolves normal derivatives but does not provide a
uniform trace estimate for the free Brown--York problem.  A dissipative
boundary term, a boundary kinetic state, or a weaker trace norm may be
needed for the tangential trace-free sector.
\end{assessment}

\subsection{Claim boundary and V16 target}
\label{sec:v3v15-conclusion}

V15 completes the companion audit and writes the actual flat harmonic
Einstein boundary block.  The six Brown--York equations invert the
tangential trace map, and the four harmonic constraints then determine the
three mixed and one double-normal derivatives.  The result supplies a
uniform frozen algebraic margin and preserves the Codazzi restriction on
the junction source.  A quadratic boundary action fixes the metric--inflow
cross blocks as adjoints.

The tangential trace-free regression also shows that the free junction
contains a Neumann glancing subsymbol.  V16 will compute the complete
Laplace--Fourier boundary energy for the six-plus-four system and add the
smallest variational boundary kinetic/damping term that removes the neutral
trace degeneracy without violating Codazzi or the BRST cochain structure.
The added term will be tested through the strict-passivity matrix from V12,
not by componentwise damping.
\clearpage
\section{Third Volume, Version V16: Induced Boundary Gravity, Passive Reservoirs, and the Glancing No-Go}
\label{sec:v3v16}

\subsection{What V15 leaves unresolved}
\label{sec:v3v16-purpose}

V15 identified the genuine six-dimensional Brown--York block and proved
that the four harmonic constraints complete its normal-derivative solve.
Its tangential trace-free sector nevertheless contains a Neumann glancing
symbol.  V16 tests the smallest local covariant kinetic correction and
finds that its principal symbol vanishes on the same boundary null cone.
It then realizes a passive impedance by adding, rather than hiding, a
conservative reservoir.  The reservoir construction keeps the energy
carrier single and makes clear which repair is covariant and which repair
depends on a chosen hyperbolic time flow.

\subsection{Intrinsic induced-gravity action}
\label{sec:v3v16-induced-action}

Let \((\mathcal N,\gamma)\) be the timelike boundary and let
\begin{equation}
 S_{\rm ind}[\gamma]
 :=\frac{\ell}{2\kappa}
 \int_{\mathcal N}(R_\gamma-2\Lambda_\partial)\,dV_\gamma,
 \qquad \ell>0.
\label{eq:v3v16-induced-action}
\end{equation}
Assume either that \(\mathcal N\) has no boundary or that the corresponding
boundary-of-boundary term is included.

\begin{proposition}[Variation of intrinsic boundary gravity]
\label{prop:v3v16-induced-variation}
The metric variation of \eqref{eq:v3v16-induced-action} is
\begin{equation}
 \delta S_{\rm ind}
 =\frac{\ell}{2\kappa}\int_{\mathcal N}
 \bigl(G^\partial_{ab}+\Lambda_\partial\gamma_{ab}\bigr)
 \delta\gamma^{ab}\,dV_\gamma.
\label{eq:v3v16-induced-variation}
\end{equation}
The complete free-metric junction equation becomes
\begin{equation}
 \tau^{\rm BY}_{ab}
 +\frac{\ell}{\kappa}
  \bigl(G^\partial_{ab}+\Lambda_\partial\gamma_{ab}\bigr)
 =S^{\rm matter}_{ab}.
\label{eq:v3v16-induced-junction}
\end{equation}
\end{proposition}

\begin{proof}
Apply the Einstein--Hilbert first-variation formula intrinsically on
\(\mathcal N\).  The boundary-of-boundary divergence is absent by
hypothesis.  Add the result to the V10 GHY variation and the matter
variation.  The coefficient of arbitrary \(\delta\gamma^{ab}\) is zero
exactly when \eqref{eq:v3v16-induced-junction} holds.
\end{proof}

\begin{theorem}[Codazzi compatibility is unchanged]
\label{thm:v3v16-Codazzi}
The induced-gravity correction does not change the mixed-flux constraint:
\begin{equation}
 D^aS^{\rm matter}_{ab}
 =T_{\mu\nu}n^\mu\gamma^\nu{}_b.
\label{eq:v3v16-Codazzi}
\end{equation}
\end{theorem}

\begin{proof}
Take the divergence of \eqref{eq:v3v16-induced-junction}.  The contracted
boundary Bianchi identity and metric compatibility give
\[
 D^aG^\partial_{ab}=0,
 \qquad
 D^a(\Lambda_\partial\gamma_{ab})=0.
\]
The divergence of the Brown--York term is the mixed bulk Einstein tensor by
the contracted Codazzi identity.  Use the bulk Einstein equation.
\end{proof}

\subsection{Linear tangential trace-free symbol}
\label{sec:v3v16-TT}

Return to the flat half-space and use the geometric outward normal
\(n=-\partial_x\).  Let \(q_{ab}\) be a tangential trace-free perturbation
which also satisfies the intrinsic harmonic constraint.  Put
\begin{equation}
 K_\partial:=\partial_t^2-\Delta_y.
\label{eq:v3v16-Kpartial}
\end{equation}

\begin{lemma}[Induced Einstein tensor on the TT block]
\label{lem:v3v16-induced-TT}
On the tangential trace-free, divergence-free block,
\begin{equation}
 (G^\partial_{ab})^{(1)}
 =\frac12K_\partial q_{ab}.
\label{eq:v3v16-induced-TT}
\end{equation}
\end{lemma}

\begin{proof}
The intrinsic analogue of the V15 trace-reversed formula is
\[
 (G^\partial_{ab})^{(1)}
 =-\frac12\Box_\gamma\bar q_{ab}
 +D_{(a}C^\partial_{b)}
 -\frac12\gamma_{ab}D^cC^\partial_c.
\]
On the trace-free block \(\bar q=q\), and the intrinsic harmonic constraint
makes the last two terms zero.  Since
\(-\Box_\gamma=\partial_t^2-\Delta_y=K_\partial\), the formula follows.
\end{proof}

\begin{proposition}[Wentzell form of the free TT junction]
\label{prop:v3v16-Wentzell}
With zero matter stress and the above orientation, the TT part of
\eqref{eq:v3v16-induced-junction} has the scalar form
\begin{equation}
 \ell K_\partial q-\partial_xq=0
\label{eq:v3v16-Wentzell}
\end{equation}
componentwise, up to the common positive factor \(2\kappa\).
\end{proposition}

\begin{proof}
For the outward normal \(-\partial_x\), the TT Brown--York variation is
\(-\frac1{2\kappa}\partial_xq\).  Lemma
\ref{lem:v3v16-induced-TT} contributes
\(\frac{\ell}{2\kappa}K_\partial q\).  Their sum vanishes precisely in
\eqref{eq:v3v16-Wentzell}.
\end{proof}

\subsection{Energy and the covariant kinetic no-go}
\label{sec:v3v16-no-go}

Consider one TT component as a scalar bulk wave
\begin{equation}
 \partial_t^2q-\partial_x^2q-\Delta_yq=0,\qquad x>0.
\label{eq:v3v16-bulk-wave}
\end{equation}
Define
\begin{align}
 E_{\rm b}(t)
 &:=\frac12\int_{x>0}
  \bigl(|q_t|^2+|q_x|^2+|\nabla_yq|^2\bigr)\,dx\,dy,
\label{eq:v3v16-bulk-energy}\\
 E_\partial(t)
 &:=\frac{\ell}{2}\int_{x=0}
  \bigl(|q_t|^2+|\nabla_yq|^2\bigr)\,dy.
\label{eq:v3v16-boundary-energy}
\end{align}

\begin{theorem}[Conservation for the induced kinetic junction]
\label{thm:v3v16-induced-energy}
Under \eqref{eq:v3v16-Wentzell},
\begin{equation}
 \frac{d}{dt}(E_{\rm b}+E_\partial)=0.
\label{eq:v3v16-induced-energy}
\end{equation}
\end{theorem}

\begin{proof}
Integration by parts in the bulk gives
\[
 \frac{dE_{\rm b}}{dt}
 =-\int_{x=0}q_tq_x\,dy.
\]
Tangential integration by parts gives
\[
 \frac{dE_\partial}{dt}
 =\ell\int_{x=0}q_tK_\partial q\,dy.
\]
Equation \eqref{eq:v3v16-Wentzell} makes the two terms cancel.
\end{proof}

\begin{theorem}[Local covariant kinetic terms do not remove principal
glancing]
\label{thm:v3v16-kinetic-no-go}
Let \(P_\partial(D)\) be a local, boundary-covariant, second-order kinetic
operator whose principal symbol is a nonzero scalar multiple of
\(\gamma^{ab}\zeta_a\zeta_b\) on the TT block.  The junction
\begin{equation}
 -\partial_xq+P_\partial(D)q=0
\label{eq:v3v16-general-kinetic}
\end{equation}
has zero principal boundary symbol at every common glancing covector
\begin{equation}
 \gamma^{ab}\zeta_a\zeta_b=0,
 \qquad \kappa(\zeta)=0,
\label{eq:v3v16-common-glancing}
\end{equation}
where \(\kappa\) is the decaying bulk normal root.  Therefore such a kinetic
term alone cannot give a homogeneous zero-loss Lopatinski margin.
\end{theorem}

\begin{proof}
For a decaying mode \(q=a e^{st+i\xi\cdot y-\kappa x}\),
the normal derivative contributes \(\kappa a\).  The order-two tangential
term contributes its principal symbol
\(c\,\gamma^{ab}\zeta_a\zeta_ba\).  Both vanish under
\eqref{eq:v3v16-common-glancing}.  Lower-order terms do not change this
homogeneous principal-symbol statement.
\end{proof}

\begin{corollary}[Induced boundary gravity retains the V15 glancing point]
\label{cor:v3v16-induced-glancing}
For \eqref{eq:v3v16-Wentzell}, the Laplace--Fourier symbol is
\begin{equation}
 b_0(s,\xi)
 =\kappa+\ell(s^2+|\xi|^2),
 \qquad
 \kappa=(s^2+|\xi|^2)^{1/2}.
\label{eq:v3v16-b0}
\end{equation}
It vanishes at the closed-frequency glancing set
\(s=i\tau\), \(|\xi|=|\tau|\).
\end{corollary}

\begin{proof}
Substitute the decaying mode in
\eqref{eq:v3v16-Wentzell}.  At glancing, both displayed factors vanish.
\end{proof}

\subsection{A conservative reservoir realization}
\label{sec:v3v16-reservoir}

Add an auxiliary half-line \(r>0\) over every boundary point and a field
\(\psi(t,y,r)\) with mass \(\mu\ge0\):
\begin{equation}
 \psi_{tt}-\psi_{rr}-\Delta_y\psi+\mu^2\psi=0.
\label{eq:v3v16-reservoir-wave}
\end{equation}
Its energy is
\begin{equation}
 E_{\rm R}
 :=\frac12\int_{r>0}
 \bigl(
  |\psi_t|^2+|\psi_r|^2+|\nabla_y\psi|^2+\mu^2|\psi|^2
 \bigr)\,dr\,dy.
\label{eq:v3v16-reservoir-energy-def}
\end{equation}
Couple the fields by the trace constraint
\begin{equation}
 \psi(t,y,0)=\alpha q(t,y,0),
\qquad \alpha\ge0,
\label{eq:v3v16-reservoir-trace}
\end{equation}
and the natural junction equation
\begin{equation}
 \ell K_\partial q-\partial_xq-\alpha\partial_r\psi|_{r=0}=0.
\label{eq:v3v16-reservoir-junction}
\end{equation}

\begin{theorem}[The enlarged variational carrier is conservative]
\label{thm:v3v16-reservoir-energy}
Solutions of \eqref{eq:v3v16-bulk-wave},
\eqref{eq:v3v16-reservoir-wave}, and
\eqref{eq:v3v16-reservoir-trace}--%
\eqref{eq:v3v16-reservoir-junction} satisfy
\begin{equation}
 \frac{d}{dt}(E_{\rm b}+E_\partial+E_{\rm R})=0.
\label{eq:v3v16-reservoir-energy}
\end{equation}
\end{theorem}

\begin{proof}
As above,
\[
 E_{\rm b}'=-\int q_tq_x\,dy,
 \qquad
 E_\partial'=\ell\int q_tK_\partial q\,dy.
\]
For the reservoir half-line,
\[
 E_{\rm R}'=-\int\psi_t(0)\psi_r(0)\,dy
 =-\alpha\int q_t\psi_r(0)\,dy.
\]
The sum is the integral of \(q_t\) times the left-hand side of
\eqref{eq:v3v16-reservoir-junction}, and hence vanishes.
\end{proof}

\begin{assessment}[Dissipation is exported energy]
\label{ass:v3v16-exported-energy}
The reservoir is not a second estimate for the same boundary state.  It is
an additional local degree of freedom with its own positive energy.  What
appears dissipative after the reservoir is eliminated is exactly the energy
transported to \(r=\infty\).
\end{assessment}

\subsection{Retarded elimination and positive-real impedance}
\label{sec:v3v16-impedance}

For \(\operatorname{Re}s>0\), the decaying reservoir solution with trace
\eqref{eq:v3v16-reservoir-trace} is
\begin{equation}
 \widehat\psi(r)
 =\alpha e^{-\rho r}\widehat q(0),
 \qquad
 \rho:=(s^2+|\xi|^2+\mu^2)^{1/2},
 \quad \operatorname{Re}\rho>0.
\label{eq:v3v16-reservoir-mode}
\end{equation}

\begin{proposition}[Effective boundary symbol]
\label{prop:v3v16-effective-symbol}
Eliminating the no-incoming reservoir gives
\begin{equation}
 b_{\rm eff}(s,\xi)
 =\kappa+\ell(s^2+|\xi|^2)+\alpha^2\rho.
\label{eq:v3v16-effective-symbol}
\end{equation}
The reservoir self-energy \(\alpha^2\rho\) has nonnegative real part for
\(\operatorname{Re}s>0\).
\end{proposition}

\begin{proof}
Equation \eqref{eq:v3v16-reservoir-mode} gives
\(\partial_r\widehat\psi(0)=-\alpha\rho\widehat q(0)\).
Substitute it and
\(\partial_x\widehat q(0)=-\kappa\widehat q(0)\) into
\eqref{eq:v3v16-reservoir-junction}.  The branch choice gives
\(\operatorname{Re}\rho>0\).
\end{proof}

\begin{theorem}[Massive reservoir removes the exact closed-frequency zero]
\label{thm:v3v16-massive-no-zero}
Assume \(\alpha>0\), \(\mu>0\), and \(\ell\ge0\).  Then
\begin{equation}
 b_{\rm eff}(s,\xi)\ne0
\label{eq:v3v16-no-zero}
\end{equation}
for every nonzero closed frequency with
\(\operatorname{Re}s\ge0\), using limiting outgoing square-root branches.
\end{theorem}

\begin{proof}
For \(\operatorname{Re}s>0\), a zero would produce a nonzero exponentially
growing normal mode of the conservative positive-energy system.  Multiply
the transformed equations by the conjugate time velocity and use the
energy identity to obtain
\[
 \operatorname{Re}s\,
 (E_{\rm b}+E_\partial+E_{\rm R})=0,
\]
which is impossible.

It remains to inspect \(s=i\tau\).  Put
\(a=|\xi|^2-\tau^2\).  If \(a\ge0\), then
\(\kappa=\sqrt a\) and
\(\rho=\sqrt{a+\mu^2}\) are nonnegative real, and the strictly positive
term \(\alpha^2\rho\) prevents a zero.  If
\(-\mu^2<a<0\), then \(\kappa\) is nonzero imaginary while \(\rho\) is
positive real, so the real and imaginary parts cannot both vanish.  If
\(a\le-\mu^2\), both square roots are imaginary with the same outgoing
sign, while the real part \(\ell a\) is nonzero when \(\ell>0\); if
\(\ell=0\), their nonzero imaginary parts have the same sign and cannot
cancel.  This proves \eqref{eq:v3v16-no-zero}.
\end{proof}

\begin{firewall}[A lower-order gap is not a homogeneous principal margin]
\label{fw:v3v16-lower-order}
At exact bulk glancing,
\[
 b_{\rm eff}(i\tau,\xi)=\alpha^2\mu
\quad\hbox{when }|\xi|=|\tau|.
\]
This removes the exact finite-frequency zero.  Under high-frequency
homogeneous rescaling, however, the fixed mass is lower order and
\(\alpha^2\mu/\langle\zeta\rangle\to0\).  The covariant massive reservoir
therefore does not by itself prove the zero-loss principal Kreiss margin
required by V14; it may yield only a weaker trace estimate.
\end{firewall}

\subsection{Order-one impedance from a transmission-line bath}
\label{sec:v3v16-transmission}

If the reservoir has no tangential spatial propagation,
\begin{equation}
 \psi_{tt}-\psi_{rr}=0,
\label{eq:v3v16-line-wave}
\end{equation}
the retarded solution is
\(\psi(t,r)=\alpha q(t-r)\), and hence
\begin{equation}
 \partial_r\psi(0)=-\alpha\partial_tq.
\label{eq:v3v16-line-DtN}
\end{equation}
The effective boundary condition becomes
\begin{equation}
 \ell K_\partial q-\partial_xq+\gamma\partial_tq=0,
\qquad \gamma:=\alpha^2>0.
\label{eq:v3v16-damped-Wentzell}
\end{equation}

\begin{theorem}[Strict glancing symbol for the transmission line]
\label{thm:v3v16-line-margin}
The effective symbol is
\begin{equation}
 b_{\rm line}(s,\xi)
 =\kappa+\ell(s^2+|\xi|^2)+\gamma s.
\label{eq:v3v16-line-symbol}
\end{equation}
On the normalized closed-frequency hemisphere, it has no zero.  In
particular, at nonzero glancing frequency,
\begin{equation}
 b_{\rm line}(i\tau,\xi)=i\gamma\tau\ne0.
\label{eq:v3v16-line-glancing}
\end{equation}
The reduced primary energy satisfies
\begin{equation}
 \frac{d}{dt}(E_{\rm b}+E_\partial)
 =-\gamma\int_{x=0}|q_t|^2\,dy.
\label{eq:v3v16-line-dissipation}
\end{equation}
\end{theorem}

\begin{proof}
Equation \eqref{eq:v3v16-line-DtN} gives
\eqref{eq:v3v16-damped-Wentzell} and
\eqref{eq:v3v16-line-symbol}.  The energy proof of
Theorem~\ref{thm:v3v16-reservoir-energy} shows that the missing primary
energy equals the positive flux carried down the line, giving
\eqref{eq:v3v16-line-dissipation}.

For a closed imaginary frequency \(s=i\tau\), put
\(a=|\xi|^2-\tau^2\).  If \(a>0\), the real part
\(\sqrt a+\ell a\) is positive.  If \(a=0\), use
\eqref{eq:v3v16-line-glancing}.  If \(a<0\), the outgoing root
\(\kappa\) and \(\gamma i\tau\) have the same imaginary sign, while
\(\ell a\le0\) is real; they cannot sum to zero.  The positive-real energy
identity excludes zeros for \(\operatorname{Re}s>0\).  Compactness gives a
positive pointwise minimum on the normalized hemisphere.
\end{proof}

\begin{firewall}[The line bath chooses a time foliation]
\label{fw:v3v16-line-covariance}
Equation \eqref{eq:v3v16-line-wave} propagates independently at each
spatial boundary point and singles out the chosen boundary time.  It is
local and variational on the enlarged gauge-fixed system, but it is not an
intrinsic Lorentz-covariant boundary matter action.  Its use in gravity
requires incorporation into the gauge-fixing/BRST sector or replacement by
a covariant reservoir with a proved principal margin.
\end{firewall}

\subsection{Cochain compatibility of the enlarged reservoir}
\label{sec:v3v16-cochain}

Let \((E^\bullet,q)\) be the finite boundary BRST complex.  Suppose all
bulk, boundary, and reservoir wave operators act as scalar differential
operators on the \(E^\bullet\) factor, and let the trace identification be
\(\psi(0)=\alpha q_{\rm fld}\), where \(\alpha\) commutes with the cochain
differential.

\begin{proposition}[The reservoir boundary system is a subcomplex]
\label{prop:v3v16-subcomplex}
The equations
\begin{align}
 K_{\rm bulk}q_{\rm fld}&=0,\qquad
 K_{\rm R}\psi=0,
\label{eq:v3v16-cochain-waves}\\
 \psi(0)&=\alpha q_{\rm fld},\qquad
 \ell K_\partial q_{\rm fld}
 -\partial_xq_{\rm fld}-\alpha\partial_r\psi(0)=0
\label{eq:v3v16-cochain-boundary}
\end{align}
are preserved by \(q\).  Hence BRST-closed data produce a BRST-closed
solution whenever the hyperbolic problem is unique.
\end{proposition}

\begin{proof}
Every differential operator in
\eqref{eq:v3v16-cochain-waves}--%
\eqref{eq:v3v16-cochain-boundary} commutes with \(q\), and so does
\(\alpha\).  Applying \(q\) to a solution gives a solution of the same
homogeneous system with cochain-differentiated data.  If the original data
are closed, uniqueness makes the differentiated solution zero.
\end{proof}

\begin{firewall}[Classical cochain commutation is not the quantum Ward
identity]
\label{fw:v3v16-quantum-BRST}
Proposition~\ref{prop:v3v16-subcomplex} is a linear classical statement.
The reservoir measure, statistics, gauge-fixing fermion, counterterms, and
determinant phase must also satisfy the Slavnov--Taylor identity before the
construction can be called a quantum BRST realization.
\end{firewall}

\subsection{V16 conclusion and V17 target}
\label{sec:v3v16-conclusion}

V16 adds the unique elementary covariant kinetic correction supplied by an
intrinsic boundary Einstein--Hilbert term and proves that the exact Codazzi
constraint survives.  Its TT block has a positive conserved Wentzell
energy, but its principal symbol vanishes at the same null glancing set as
the bulk Neumann root.  This is a sharp no-go for kinetic correction alone.

An enlarged conservative reservoir gives a positive-real boundary
self-energy without duplicate accounting.  A covariant massive reservoir
removes exact finite-frequency zeros but leaves a lower-order homogeneous
margin.  A transmission-line reservoir gives the order-one impedance and
strict glancing symbol, while selecting a boundary time foliation.  The
linear cochain system commutes with BRST.

V17 will attack the remaining covariance gap.  It will formulate the
reservoir in ADM boundary variables, introduce its lapse and shift
dependence through a BRST gauge-fixing fermion, and prove at the classical
linear level that the impedance acts only on the physical tangential
trace-free quotient while the four harmonic constraints retain their
homogeneous propagation.  If that projection is necessarily nonlocal, the
note will state the obstruction rather than hiding it in a componentwise
damping term.
\clearpage
\section{Third Volume, Version V17: ADM--BRST Boundary Variables and the Nonlocal Physical-Projector Obstruction}
\label{sec:v3v17}

\subsection{The remaining covariance question}
\label{sec:v3v17-purpose}

V16 obtained an order-one passive impedance from a local transmission-line
reservoir, but that construction used a chosen boundary time.  Applying the
impedance only to selected tensor components would generally break linear
diffeomorphism covariance and the harmonic constraints.  V17 compares two
routes:
\begin{enumerate}
\item project exactly onto the physical boundary quotient and damp only
      that quotient; or
\item couple the complete gauge-fixed BRST complex to a chain-equivariant
      reservoir and let the impedance descend to cohomology.
\end{enumerate}
The first route is necessarily spatially nonlocal.  The second is local at
the linear level and preserves harmonic-constraint propagation.

\subsection{Linear ADM variables on the timelike boundary}
\label{sec:v3v17-ADM}

Write the boundary metric in ADM form
\begin{equation}
 \gamma
 =-N^2dt^2
 +\sigma_{AB}(dy^A+\beta^A dt)(dy^B+\beta^Bdt).
\label{eq:v3v17-ADM}
\end{equation}
Linearize about
\[
 N=1,\qquad \beta^A=0,\qquad \sigma_{AB}=\delta_{AB}.
\]
Let
\begin{equation}
 N=1+\nu,\qquad
 \beta_A=b_A,\qquad
 \sigma_{AB}=\delta_{AB}+k_{AB}.
\label{eq:v3v17-ADM-perturbations}
\end{equation}
Then the tangential metric perturbation is
\begin{equation}
 q_{tt}=-2\nu,\qquad q_{tA}=b_A,\qquad q_{AB}=k_{AB}.
\label{eq:v3v17-q-ADM}
\end{equation}

\subsection{Linear boundary BRST complex}
\label{sec:v3v17-BRST}

Let \(c^a\) be the boundary diffeomorphism ghost,
\(\bar c_a\) the antighost, and \(B_a\) the Nakanishi--Lautrup field.
Define the odd differential \(s\) by
\begin{align}
 sq_{ab}&=\partial_ac_b+\partial_bc_a,
\label{eq:v3v17-sq}\\
 sc_a&=0,
\label{eq:v3v17-sc}\\
 s\bar c_a&=B_a,
\label{eq:v3v17-sbarc}\\
 sB_a&=0.
\label{eq:v3v17-sB}
\end{align}

\begin{lemma}[Nilpotence and ADM component rules]
\label{lem:v3v17-nilpotence}
One has \(s^2=0\).  In the ADM variables
\eqref{eq:v3v17-ADM-perturbations},
\begin{align}
 s\nu&=\partial_tc^t,
\label{eq:v3v17-snu}\\
 sb_A&=\partial_tc_A-\partial_Ac^t,
\label{eq:v3v17-sb}\\
 sk_{AB}&=\partial_Ac_B+\partial_Bc_A.
\label{eq:v3v17-sk}
\end{align}
\end{lemma}

\begin{proof}
Equations \eqref{eq:v3v17-sc} and \eqref{eq:v3v17-sB} make a second
application of \(s\) vanish on every generator.  For the component rules,
use \(c_t=-c^t\) on the flat background and compare
\(sq_{tt}=2\partial_tc_t\),
\(sq_{tA}=\partial_tc_A+\partial_Ac_t\), and
\(sq_{AB}=2\partial_{(A}c_{B)}\) with
\eqref{eq:v3v17-q-ADM}.
\end{proof}

\subsection{Harmonic gauge from a gauge-fixing fermion}
\label{sec:v3v17-gauge-fermion}

Define the boundary de Donder functional
\begin{equation}
 F_a(q)
 :=\partial^bq_{ab}-\frac12\partial_aq^b{}_b
 =\partial^b\bar q_{ab}.
\label{eq:v3v17-F}
\end{equation}
For a gauge parameter \(\alpha_{\rm gf}>0\), take
\begin{equation}
 \Psi_{\rm gf}
 :=\int_{\mathcal N}
 \bar c^a\left(F_a(q)+\frac{\alpha_{\rm gf}}2B_a\right)dV.
\label{eq:v3v17-gauge-fermion}
\end{equation}

\begin{proposition}[BRST-exact harmonic gauge sector]
\label{prop:v3v17-gauge-fixed}
The BRST variation is
\begin{equation}
 s\Psi_{\rm gf}
 =\int_{\mathcal N}
 \left(
  B^aF_a+\frac{\alpha_{\rm gf}}2B^aB_a
  -\bar c^a\Box_\partial c_a
 \right)dV,
\label{eq:v3v17-gauge-fixed}
\end{equation}
up to the conventional overall sign for odd integration.  Eliminating
\(B_a\) gives the harmonic gauge-fixing term
\(-\frac1{2\alpha_{\rm gf}}F^aF_a\), while the ghost principal operator is
\(\Box_\partial\).
\end{proposition}

\begin{proof}
The linear gauge variation of \eqref{eq:v3v17-F} is
\begin{align*}
 sF_a
 &=\partial^b(\partial_ac_b+\partial_bc_a)
   -\partial_a\partial^bc_b\\
 &=\Box_\partial c_a.
\end{align*}
Apply the graded Leibniz rule to
\eqref{eq:v3v17-gauge-fermion}.  The algebraic equation for \(B\) is
\(B_a=-\alpha_{\rm gf}^{-1}F_a\); completing the square gives the stated
gauge-fixing term.
\end{proof}

\begin{assessment}[The lapse and shift are gauge-sector fields]
\label{ass:v3v17-lapse-shift}
Equations \eqref{eq:v3v17-snu}--\eqref{eq:v3v17-sb} show that damping
\(\nu\) or \(b_A\) independently is not BRST invariant.  A local impedance
must either be written on a gauge-invariant observable or extended to the
ghost and auxiliary partners through the complete gauge-fixed complex.
\end{assessment}

\subsection{The exact physical projector is nonlocal}
\label{sec:v3v17-projector}

At a nonzero tangential Fourier covector \(\xi\), define the linear Killing
symbol
\begin{equation}
 d_\xi v:=i(\xi\otimes v+v\otimes\xi)
\label{eq:v3v17-Killing-symbol}
\end{equation}
from covectors to symmetric two-tensors.  With fixed positive Euclidean
fibre metrics, the orthogonal projector onto the complement of gauge
directions is
\begin{equation}
 P_{\rm phys}(\xi)
 :=I-d_\xi(d_\xi^*d_\xi)^{-1}d_\xi^*.
\label{eq:v3v17-physical-projector}
\end{equation}

\begin{lemma}[Inverse derivative in the quotient projector]
\label{lem:v3v17-projector-inverse}
For \(\xi\ne0\), the operator \(d_\xi^*d_\xi\) is positive and its inverse
is homogeneous of degree \(-2\).  Consequently
\(P_{\rm phys}(\xi)\) is homogeneous of degree zero and depends
nontrivially on the direction \(\xi/|\xi|\).
\end{lemma}

\begin{proof}
One computes
\[
 \|d_\xi v\|^2
 =2|\xi|^2|v|^2+2|\xi\cdot v|^2,
\]
which is positive for \(v\ne0\).  The operator is quadratic in \(\xi\), so
its inverse has degree \(-2\).  The range of \(d_\xi\) rotates with the
direction of \(\xi\), so its orthogonal complement and projector are not
constant on the unit sphere.
\end{proof}

\begin{theorem}[No finite-order local physical projector]
\label{thm:v3v17-no-local-projector}
There is no constant-coefficient finite-order differential operator whose
Fourier symbol equals \(P_{\rm phys}(\xi)\) for every \(\xi\ne0\).
More generally, no order-zero differential principal symbol can equal the
direction-dependent projector \eqref{eq:v3v17-physical-projector}.
\end{theorem}

\begin{proof}
The symbol of a finite-order constant-coefficient differential operator is
a matrix polynomial in \(\xi\).  If it defines an \(L^2\)-bounded
order-zero projection, that polynomial is bounded as
\(|\xi|\to\infty\), hence is constant.  But
Lemma~\ref{lem:v3v17-projector-inverse} shows that
\(P_{\rm phys}\) is direction dependent.  For variable coefficients, the
principal symbol of an order-zero differential operator is likewise
independent of \(\xi\); the same pointwise contradiction applies.
\end{proof}

\begin{corollary}[Physical-only damping is pseudodifferential]
\label{cor:v3v17-nonlocal-damping}
An exact impedance of the form
\begin{equation}
 \gamma P_{\rm phys}\partial_tq
\label{eq:v3v17-physical-damping}
\end{equation}
is spatially pseudodifferential.  Its kernel contains the elliptic inverse
\((d^*d)^{-1}\) and is not a local componentwise boundary term.
\end{corollary}

\begin{firewall}[A Fourier projector is not a local covariant action]
\label{fw:v3v17-Fourier-projector}
Formula \eqref{eq:v3v17-physical-projector} can be used in a microlocal
estimate, but inserting it into a fundamental boundary action changes the
locality problem.  The instruction to damp only the physical modes cannot
hide this projector.  One must accept the nonlocal effective action or
introduce local auxiliary fields which realize the inverse.
\end{firewall}

\subsection{Local alternative: a reservoir for the complete complex}
\label{sec:v3v17-complete-reservoir}

Let \(\mathcal E^\bullet\) denote the linear boundary BRST complex
containing
\[
 (c_a;\ q_{ab};\ \bar c_a,B_a).
\]
Introduce a reservoir copy
\[
 (\chi_a;\ \psi_{ab};\ \bar\chi_a,\mathcal B_a)
\]
on \(r>0\), with
\begin{align}
 s\psi_{ab}&=\partial_a\chi_b+\partial_b\chi_a,
\label{eq:v3v17-spsi}\\
 s\chi_a&=0,\qquad
 s\bar\chi_a=\mathcal B_a,\qquad
 s\mathcal B_a=0.
\label{eq:v3v17-reservoir-quartet}
\end{align}
Impose the chainwise trace conditions
\begin{equation}
 \psi(0)=\alpha q,\quad
 \chi(0)=\alpha c,\quad
 \bar\chi(0)=\alpha\bar c,\quad
 \mathcal B(0)=\alpha B.
\label{eq:v3v17-chain-trace}
\end{equation}

\begin{proposition}[The trace coupling is a BRST subcomplex]
\label{prop:v3v17-trace-subcomplex}
The four relations \eqref{eq:v3v17-chain-trace} are preserved by \(s\).
\end{proposition}

\begin{proof}
Apply \(s\) to each relation.  For the tensor relation,
\[
 s(\psi(0)-\alpha q)
 =d(\chi(0)-\alpha c)=0.
\]
The ghost relation is closed, and the antighost relation maps to the
auxiliary-field relation, which is closed.
\end{proof}

Let the gauge-fixed reservoir kinetic operator act degreewise as
\begin{equation}
 K_{\rm R}:=\partial_t^2-\partial_r^2+\Delta_{\rm gf},
\label{eq:v3v17-KR}
\end{equation}
where \(\Delta_{\rm gf}\) is the boundary Hodge/BRST Laplacian induced by
the quadratic invariant action plus
\(s\Psi_{\rm gf}\).  At the flat principal level,
\begin{equation}
 [K_{\rm R},s]=0.
\label{eq:v3v17-KR-commutes}
\end{equation}

\begin{theorem}[Local impedance descends to BRST cohomology]
\label{thm:v3v17-descends}
Assume the complete bulk, boundary, and reservoir equations commute with
\(s\), and use the trace subcomplex
\eqref{eq:v3v17-chain-trace}.  Retarded elimination of the transmission
line induces the same scalar impedance multiplier \(\gamma\partial_t\) on
each degree.  This multiplier:
\begin{enumerate}
\item maps closed fields to closed fields;
\item maps exact fields to exact fields; and
\item therefore induces a well-defined passive operator on
      \(H^\bullet(\mathcal E,s)\).
\end{enumerate}
\end{theorem}

\begin{proof}
The no-incoming solution of the line equation gives
\(\partial_r\Psi(0)=-\alpha\partial_t\Phi\) in every degree, where
\(\Phi\) denotes the corresponding boundary field.  Because
\([s,\partial_t]=0\), the impedance commutes with \(s\).  If
\(s\Phi=0\), then
\(s(\gamma\partial_t\Phi)=\gamma\partial_ts\Phi=0\).
If \(\Phi=s\Lambda\), then
\(\gamma\partial_t\Phi=s(\gamma\partial_t\Lambda)\).
Thus the operator is well defined on the quotient
\(\ker s/\operatorname{Ran}s\).  Its passivity is the V16 reservoir energy
identity applied degreewise to the bosonic physical sector.
\end{proof}

\begin{assessment}[Locality is recovered before quotienting]
\label{ass:v3v17-local-before-quotient}
The full reservoir contains gauge, ghost, and auxiliary partners.  Its
equations and trace maps are local.  Only after solving the local complex is
the quotient by exact states taken.  This reverses the nonlocal order in
\eqref{eq:v3v17-physical-damping}: quotienting first requires the elliptic
projector.
\end{assessment}

\subsection{Harmonic-constraint propagation with impedance}
\label{sec:v3v17-constraint}

Let \(C_\mu=\partial^\nu\bar h_{\mu\nu}\) be the V15 bulk harmonic
constraint.  At the flat principal level define the componentwise
Wentzell--impedance operator
\begin{equation}
 \mathcal W
 :=\ell(\partial_t^2-\Delta_y)
   -\partial_x+\gamma\partial_t.
\label{eq:v3v17-W}
\end{equation}

\begin{lemma}[Constraint commutator]
\label{lem:v3v17-constraint-commutator}
For constant coefficients,
\begin{equation}
 \partial^\nu\overline{\mathcal Wh}_{\mu\nu}
 =\mathcal W C_\mu.
\label{eq:v3v17-constraint-commutator}
\end{equation}
\end{lemma}

\begin{proof}
Trace reversal is algebraic and every derivative in \(\mathcal W\) commutes
with every constant-coefficient derivative in \(C_\mu\).
\end{proof}

\begin{theorem}[Homogeneous constraint propagation at the damped junction]
\label{thm:v3v17-constraint-propagation}
Suppose
\[
 \Box C_\mu=0
\]
in the bulk, the differentiated boundary source is co-closed so that
\begin{equation}
 \mathcal WC_\mu=0
\label{eq:v3v17-C-boundary}
\end{equation}
at \(x=0\), and the initial position and velocity constraints vanish.
Then \(C_\mu=0\) for all times for which the linear problem exists.
\end{theorem}

\begin{proof}
For each component, use the V16 bulk plus Wentzell boundary energy and the
transmission-line reservoir energy.  With zero source the total constraint
energy is conserved, while the reduced primary energy obeys
\[
 E_C'(t)=-\gamma\|\partial_tC_\mu|_{x=0}\|_{L^2_y}^2\le0.
\]
The initial energy is zero, so every positive term remains zero and
\(C_\mu=0\).  Equivalently, uniqueness for the homogeneous dissipative
boundary problem gives the same conclusion.
\end{proof}

\begin{corollary}[Linear Codazzi compatibility]
\label{cor:v3v17-Codazzi}
If the complete reservoir stress is obtained by varying the same
gauge-fixed action and its linear source satisfies the BRST/Noether
co-closure used above, the V15 relation
\[
 \partial^aS_{ab}^{(1)}=T_{xb}^{(1)}
\]
is preserved.  The impedance does not create an independent divergence
source.
\end{corollary}

\begin{proof}
The linear Noether identity for the complete action identifies the
divergence of its boundary metric variation with the boundary field
equations.  Those equations include the reservoir flux with the opposite
sign, as in the V16 energy identity.  On shell, only the mixed bulk flux
remains, giving the linear Codazzi relation.
\end{proof}

\begin{firewall}[Gauge-fixed Noether is not full nonlinear covariance]
\label{fw:v3v17-nonlinear-covariance}
The proof uses the flat linear BRST differential and constant-coefficient
commutators.  With evolving lapse, shift, and boundary metric,
\([s,D_t]\) contains ghost and connection terms.  A nonlinear
Slavnov--Taylor identity and covariant reservoir stress have not yet been
constructed.
\end{firewall}

\subsection{Regulator and determinant firewalls}
\label{sec:v3v17-firewalls}

\begin{firewall}[Ghost energy is not positive]
\label{fw:v3v17-ghost-energy}
The positive reservoir energy proves passivity on bosonic physical
components.  Grassmann ghost fields do not carry a positive Hilbert energy
of the same kind.  Their role is algebraic cancellation in the BRST
complex, to be checked by graded determinants rather than by the bosonic
energy inequality.
\end{firewall}

\begin{firewall}[Acyclic partners may still leave a phase]
\label{fw:v3v17-phase}
Adding a local contractible quartet can cancel nonzero determinant
magnitudes formally, but eta phases, zero modes, and global anomaly
holonomy may survive.  Cohomological descent of the impedance does not
trivialize the determinant line.
\end{firewall}

\begin{counterexample}[Projecting after an approximate inverse]
\label{cex:v3v17-approximate-projector}
Let \(P_n=I-d(d^*d+n^{-1})^{-1}d^*\).  Each \(P_n\) is bounded and local
only after introducing an elliptic auxiliary solve, but
\begin{equation}
 P_n^2-P_n
 =-n^{-1}d(d^*d+n^{-1})^{-2}d^*
\label{eq:v3v17-projector-defect}
\end{equation}
is generally nonzero.  Thus finite regulator mass turns the exact quotient
projector into a non-idempotent filter.  BRST exactness cannot be inferred
from norm convergence alone; the chain relation must be checked at each
regulator.
\end{counterexample}

\subsection{V17 conclusion and V18 target}
\label{sec:v3v17-conclusion}

V17 writes the boundary lapse, shift, and spatial metric in the linear BRST
complex and derives harmonic gauge from a gauge-fixing fermion.  The exact
orthogonal projector away from diffeomorphism directions contains
\((d^*d)^{-1}\), so a physical-only impedance is necessarily
pseudodifferential.  A local alternative is to attach the reservoir to the
complete gauge-fixed complex.  Its chainwise trace condition is BRST
stable, the retarded impedance descends to cohomology, and the harmonic
constraints retain a homogeneous dissipative boundary problem.

V18 will test the quantum-linearized part of this proposal.  At finite
spectral cutoff it will compute the graded Gaussian determinant of one
reservoir BRST quartet, identify the exact cancellation of nonzero
longitudinal modes, and isolate zero modes and determinant phases as the
remaining data.  That calculation is required before the local complete-
complex route can be claimed to preserve the quantum boundary measure.
\clearpage
\section{Third Volume, Version V18: Finite-Cutoff Quartet Determinants and the Surviving Quantum Boundary Data}
\label{sec:v3v18}

\subsection{Purpose of the determinant test}
\label{sec:v3v18-purpose}

V17 attached a local reservoir to the complete gauge-fixed BRST complex.
Classical commutation with the BRST differential proves that the impedance
descends to cohomology, but it does not prove that the unphysical reservoir
variables cancel in the quantum measure.  V18 performs the exact
finite-dimensional Gaussian calculation.  The cancellation holds
mode-by-mode only when the bosonic and fermionic operators, cutoffs, and
boundary domains are the same chainwise.  Physical cohomology, zero modes,
and determinant-line phases remain.

\subsection{A balanced complex quartet}
\label{sec:v3v18-quartet}

Let \(V\cong\mathbb C^N\) and let \(A:V\to V\) be positive Hermitian.
Introduce complex bosonic variables \(\phi,\bar\phi\) and Grassmann
variables \(c,\bar c\).  Define
\begin{equation}
 s\phi=c,\qquad sc=0,\qquad
 s\bar c=\bar\phi,\qquad s\bar\phi=0.
\label{eq:v3v18-quartet-differential}
\end{equation}

\begin{lemma}[Nilpotence and acyclicity]
\label{lem:v3v18-quartet-acyclic}
The differential in \eqref{eq:v3v18-quartet-differential} satisfies
\(s^2=0\), and its polynomial cohomology is represented by constants.
\end{lemma}

\begin{proof}
Nilpotence holds on each generator.  Define the counting operator
\[
 N_{\rm q}
 =\phi\partial_\phi+c\partial_c
  +\bar c\partial_{\bar c}+\bar\phi\partial_{\bar\phi}
\]
and the homotopy
\[
 h=\phi\partial_c+\bar c\partial_{\bar\phi}.
\]
A direct graded calculation gives
\[
 sh+hs=N_{\rm q}.
\]
If a closed homogeneous polynomial has positive quartet degree \(m\), then
it equals \(m^{-1}s(hP)\) and is exact.  Only degree-zero constants remain.
\end{proof}

The quadratic action is
\begin{equation}
 S_{\rm q}
 :=\bar\phi A\phi-\bar c A c.
\label{eq:v3v18-quartet-action}
\end{equation}
With the sign convention for the Berezin exponent used below, this is
\(s(\bar c A\phi)\).

\begin{proposition}[BRST exactness of the quadratic action]
\label{prop:v3v18-action-exact}
If \(A\) commutes with \(s\), then
\begin{equation}
 s(\bar c A\phi)
 =\bar\phi A\phi-\bar c A c=S_{\rm q}.
\label{eq:v3v18-action-exact}
\end{equation}
\end{proposition}

\begin{proof}
Use the graded Leibniz rule:
\[
 s(\bar c A\phi)
 =(s\bar c)A\phi-\bar c A(s\phi)
 =\bar\phi A\phi-\bar cAc.
\]
\end{proof}

\subsection{Exact finite-dimensional Gaussian cancellation}
\label{sec:v3v18-Gaussian}

Normalize the complex bosonic measure by
\[
 d\mu_{\rm B}
 :=\prod_{j=1}^N\frac{d\operatorname{Re}\phi_j\,
 d\operatorname{Im}\phi_j}{\pi},
\]
and fix the Berezin orientation by
\[
 \int\prod_{j=1}^N d\bar c_j\,dc_j\,
 \exp(\bar cAc)=\det A.
\]

\begin{theorem}[One quartet has unit nonzero-mode determinant]
\label{thm:v3v18-quartet-cancellation}
For \(A>0\),
\begin{align}
 Z_{\rm q}(A)
 &:=
 \int d\mu_{\rm B}\,d\bar c\,dc\,
 \exp\bigl(-\bar\phi A\phi+\bar cAc\bigr)
\notag\\
 &= (\det A)^{-1}\det A
 =1.
\label{eq:v3v18-quartet-cancellation}
\end{align}
The equality is exact for every finite cutoff \(N\).
\end{theorem}

\begin{proof}
Diagonalize \(A=U^*\operatorname{diag}(\lambda_1,\ldots,\lambda_N)U\).
The unitary change of bosonic variables preserves \(d\mu_{\rm B}\), and
the corresponding inverse changes in \(c,\bar c\) preserve the Berezin
measure.  For each eigenvalue,
\[
 \int_{\mathbb C}\frac{d^2z}{\pi}e^{-\lambda|z|^2}
 =\lambda^{-1},
 \qquad
 \int d\bar c\,dc\,e^{\lambda\bar cc}
 =\lambda.
\]
Multiplication over \(j\) gives \eqref{eq:v3v18-quartet-cancellation}.
\end{proof}

\begin{corollary}[A common reservoir impedance cancels on the quartet]
\label{cor:v3v18-common-impedance}
Let \(A_\gamma=A+\gamma R\) be positive, where \(R\) is the same regulated
reservoir self-energy in all four quartet variables and commutes with
\(s\).  Then
\begin{equation}
 Z_{\rm q}(A_\gamma)=1.
\label{eq:v3v18-common-impedance}
\end{equation}
\end{corollary}

\begin{proof}
Apply Theorem~\ref{thm:v3v18-quartet-cancellation} to the positive matrix
\(A_\gamma\).  The proof uses equality of the complete operators, not a
large-eigenvalue approximation.
\end{proof}

\begin{assessment}[Two real bosons and one complex ghost pair]
\label{ass:v3v18-degree-count}
A complex boson contains two real Gaussian variables and contributes
\((\det A)^{-1}\).  The complex Grassmann pair contributes
\(\det A\).  Omitting one real bosonic partner would leave
\((\det A)^{1/2}\); cohomological acyclicity by itself does not correct a
mismatched functional measure.
\end{assessment}

\subsection{Chainwise spectral cutoff}
\label{sec:v3v18-cutoff}

Let \(K\ge0\) be a self-adjoint elliptic cutoff operator which commutes with
\(s\).  Its spectral projection is
\begin{equation}
 \Pi_\Lambda:=\mathbf 1_{[0,\Lambda]}(K).
\label{eq:v3v18-spectral-cutoff}
\end{equation}

\begin{lemma}[The cutoff is a chain map]
\label{lem:v3v18-cutoff-chain}
If \([K,s]=0\), then
\begin{equation}
 [\Pi_\Lambda,s]=0.
\label{eq:v3v18-cutoff-chain}
\end{equation}
\end{lemma}

\begin{proof}
In finite volume, \(K\) has a spectral decomposition.  Commutation with
\(s\) makes every eigenspace invariant under \(s\), so the sum of
eigenspaces with eigenvalue at most \(\Lambda\) is invariant.  Equivalently,
use the Riesz formula for \(\Pi_\Lambda\) and commute \(s\) through the
resolvent.
\end{proof}

\begin{theorem}[Exact cancellation at a common spectral cutoff]
\label{thm:v3v18-cutoff-cancellation}
Suppose the quartet kinetic operator \(A\) and the cutoff \(K\) commute with
\(s\), and restrict to the nonzero eigenspaces of \(A\) inside
\(\operatorname{Ran}\Pi_\Lambda\).  Then the quartet contribution is
\begin{equation}
 Z_{{\rm q},\Lambda}'=1.
\label{eq:v3v18-cutoff-cancellation}
\end{equation}
No limit \(\Lambda\to\infty\) is needed for this equality.
\end{theorem}

\begin{proof}
Lemma~\ref{lem:v3v18-cutoff-chain} restricts the quartet in complete
BRST multiplets.  The restricted positive operator is a finite matrix.
Theorem~\ref{thm:v3v18-quartet-cancellation} applies to that matrix.
\end{proof}

\begin{counterexample}[Different cutoffs leave an artificial determinant]
\label{cex:v3v18-mismatched-cutoff}
Let the bosonic cutoff retain eigenvalues
\(\lambda_1,\ldots,\lambda_N\), while the ghost cutoff retains only the
first \(N-1\).  Then
\begin{equation}
 Z_{\rm mismatch}
 =\left(\prod_{j=1}^N\lambda_j^{-1}\right)
  \left(\prod_{j=1}^{N-1}\lambda_j\right)
 =\lambda_N^{-1}.
\label{eq:v3v18-mismatched-cutoff}
\end{equation}
The residual factor is a regulator artifact created by cutting a quartet
in half.
\end{counterexample}

\subsection{Boundary domains must also form a complex}
\label{sec:v3v18-domains}

Let \(A_{\rm B}\) and \(A_{\rm F}\) denote the bosonic and ghost kinetic
operators with their boundary domains included.

\begin{theorem}[Domain equality is part of determinant cancellation]
\label{thm:v3v18-domain-equality}
If \(A_{\rm B}=A_{\rm F}=A\) as operators, including their domains, then the
nonzero quartet determinant is one.  If their differential expressions
agree but their domains differ, the ratio is
\begin{equation}
 Z_{\rm q}
 =\frac{\det{}'A_{\rm F}}{\det{}'A_{\rm B}},
\label{eq:v3v18-domain-ratio}
\end{equation}
which need not equal one.
\end{theorem}

\begin{proof}
An operator determinant is the product of eigenvalues satisfying its
domain conditions.  Equal differential expressions and equal domains give
the same spectrum with multiplicity and Theorem
\ref{thm:v3v18-quartet-cancellation}.  Different domains can change the
eigenvalues or their multiplicities; the two Gaussian integrations then
give the displayed ratio.
\end{proof}

\begin{counterexample}[One-mode boundary mismatch]
\label{cex:v3v18-boundary-mismatch}
For a single retained mode, take
\[
 A_{\rm B}=\lambda,\qquad
 A_{\rm F}=\lambda+\delta,
 \qquad \lambda>0,\quad\delta\ne0.
\]
Then
\begin{equation}
 Z_{\rm q}=1+\frac{\delta}{\lambda}\ne1.
\label{eq:v3v18-boundary-mismatch}
\end{equation}
Thus asymptotic agreement as \(\lambda\to\infty\) does not prove exact Ward
cancellation at a finite regulator.
\end{counterexample}

\begin{corollary}[V17 trace conditions are quantum-relevant]
\label{cor:v3v18-trace-domains}
The chainwise reservoir trace relations of V17 are necessary for the
bosonic and ghost reservoir operators to have paired boundary domains.
It is not enough that their interior principal symbols agree.
\end{corollary}

\subsection{Heat supertrace on the acyclic sector}
\label{sec:v3v18-heat}

Let \(\mathcal H_{\rm q}\) be a finite quartet Hilbert superspace, with
bosonic parity on \(\phi,\bar\phi\) and fermionic parity on
\(c,\bar c\).  Let \(A\) be positive, even, and commute with \(s\).

\begin{proposition}[Modewise vanishing of the quartet heat supertrace]
\label{prop:v3v18-heat-supertrace}
On a balanced acyclic quartet,
\begin{equation}
 \operatorname{Str}_{\mathcal H_{\rm q}}e^{-tA}=0,
 \qquad t>0.
\label{eq:v3v18-heat-supertrace}
\end{equation}
Consequently all finite-cutoff heat coefficients cancel between the paired
quartet variables when their operators and domains coincide.
\end{proposition}

\begin{proof}
Each eigenvalue eigenspace is an \(s\)-invariant quartet by
\([A,s]=0\).  It contains equal bosonic and fermionic multiplicities, so
its contribution is
\[
 (m_{\rm B}-m_{\rm F})e^{-t\lambda}=0.
\]
Sum over the finitely many retained eigenvalues.
\end{proof}

\begin{corollary}[Nonzero quartet logarithmic determinants cancel]
\label{cor:v3v18-logdet}
For any proper-time interval
\(0<\varepsilon<T<\infty\),
\begin{equation}
 \int_\varepsilon^T\frac{dt}{t}
 \operatorname{Str}_{\mathcal H_{\rm q}}e^{-tA}=0.
\label{eq:v3v18-logdet}
\end{equation}
Thus no local heat divergence is produced by the exactly paired nonzero
quartet sector.
\end{corollary}

\begin{proof}
Integrate the pointwise identity
\eqref{eq:v3v18-heat-supertrace}.
\end{proof}

\begin{firewall}[Continuum cancellation requires uniform domains]
\label{fw:v3v18-continuum-domains}
At finite cutoff the cancellation is exact.  Passing to a continuum still
requires a common self-adjoint extension, uniform heat bounds, and a
regulator which preserves the chain domain.  Boundary heat coefficients do
not cancel if the bosonic and ghost Robin or impedance parameters differ.
\end{firewall}

\subsection{Physical cohomology does not cancel}
\label{sec:v3v18-physical}

Let the regulated gauge-fixed space split as
\begin{equation}
 \mathcal H_\Lambda
 =\mathcal H_{{\rm phys},\Lambda}
 \oplus\mathcal H_{{\rm q},\Lambda}
 \oplus\mathcal H_{0,\Lambda},
\label{eq:v3v18-splitting}
\end{equation}
where the second summand is a direct sum of acyclic quartets and the third
contains zero modes.

\begin{theorem}[Finite-cutoff factorization]
\label{thm:v3v18-factorization}
If the quadratic action is block diagonal with respect to
\eqref{eq:v3v18-splitting} and the quartet hypotheses hold, then
\begin{equation}
 Z_\Lambda
 =Z_{{\rm phys},\Lambda}\,Z_{0,\Lambda}.
\label{eq:v3v18-factorization}
\end{equation}
The reservoir determinant and Schur self-energy on physical cohomology
remain in \(Z_{{\rm phys},\Lambda}\).
\end{theorem}

\begin{proof}
Gaussian integration factorizes over the three blocks.  Every acyclic
quartet factor equals one by
Theorem~\ref{thm:v3v18-cutoff-cancellation}.  No cancellation theorem
applies to physical cohomology or to zero modes.
\end{proof}

\begin{assessment}[The physical reservoir is not disposable]
\label{ass:v3v18-physical-reservoir}
The local complete-complex construction adds reservoir variables in every
degree, but only the unphysical balanced quartets cancel.  The reservoir
coupled to a physical BRST class produces the V9 Schur complement and its
bosonic determinant.  Removing that factor would erase the passive
impedance itself.
\end{assessment}

\subsection{Zero modes}
\label{sec:v3v18-zero-modes}

Let \(A\ge0\) have kernel \(Z_0\) of dimension \(k>0\).

\begin{proposition}[Nonzero determinants do not define the zero-mode
measure]
\label{prop:v3v18-zero-modes}
Writing \(V=Z_0\oplus Z_0^\perp\), the primed quartet determinant on
\(Z_0^\perp\) equals one, but the integral over \(Z_0\) is not determined
by that cancellation.  Bosonic zero modes give volume factors or collective
coordinates, while unsaturated Grassmann zero modes make the Berezin
integral vanish.
\end{proposition}

\begin{proof}
Diagonalize \(A\).  Every positive eigenvalue cancels as before.  At
\(\lambda=0\), the bosonic factor
\(\int_{\mathbb C}d^2z\) diverges unless a quotient or compact collective
coordinate is specified.  The Grassmann integral
\(\int d\bar c\,dc\,1=0\) unless insertions saturate both variables.
Neither expression is the limit of the primed determinant alone.
\end{proof}

\begin{counterexample}[A common small mass can hide zero-mode data]
\label{cex:v3v18-small-mass}
Replacing \(A\) by \(A+\varepsilon I\) makes the quartet determinant one
even on the former kernel for every \(\varepsilon>0\).  Nevertheless, the
bosonic zero-mode variance is \(\varepsilon^{-1}\) and diverges as
\(\varepsilon\downarrow0\), while Grassmann correlation functions require
zero-mode saturation.  The number one obtained by multiplying determinants
does not specify the limiting measure.
\end{counterexample}

\subsection{Determinant lines and phases}
\label{sec:v3v18-phases}

\begin{firewall}[Positive Euclidean matrices have no phase information]
\label{fw:v3v18-phase}
Theorem~\ref{thm:v3v18-quartet-cancellation} uses a positive Hermitian
matrix and a fixed Berezin orientation.  Lorentzian first-order operators,
chiral fermions, and families of boundary Dirac operators can carry eta
phases, spectral flow, and determinant-line holonomy.  Cancellation of
positive determinant magnitudes does not determine those phases.
\end{firewall}

\begin{firewall}[Local pairing need not be global]
\label{fw:v3v18-global-pairing}
Even if spectra pair in every coordinate chart or parameter neighborhood,
the identifications can return with holonomy around a parameter loop.  A
global quantum cancellation requires an orientation and trivialization of
the determinant line compatible with the BRST pairing.
\end{firewall}

\begin{firewall}[Retarded elimination and Euclidean integration differ]
\label{fw:v3v18-retarded-Euclidean}
The V16 impedance used a retarded, no-incoming reservoir condition.  The
Gaussian calculation in V18 is Euclidean.  Relating them requires an
analytic continuation and a spectral condition selecting the same
boundary domain.  Equality of their formal symbols is not a proof of
reflection positivity or causal reconstruction.
\end{firewall}

\subsection{Regulator-stable cancellation criterion}
\label{sec:v3v18-regulator}

\begin{theorem}[Exact finite-regulator Ward package]
\label{thm:v3v18-regulator-package}
For each regulator \(\Lambda\), assume:
\begin{enumerate}
\item \(s_\Lambda^2=0\);
\item the kinetic operator and spectral cutoff commute with \(s_\Lambda\);
\item bosonic and fermionic quartet operators have identical nonzero
      spectra, multiplicities, and boundary domains;
\item the functional measures use compatible complex and Berezin
      orientations; and
\item zero modes are split off before determinants are formed.
\end{enumerate}
Then every nonzero acyclic quartet contributes exactly one to the regulated
partition function and zero to the regulated heat supertrace.
\end{theorem}

\begin{proof}
Items one and two split the cutoff space into invariant finite-dimensional
complexes.  Item three identifies the matrices in the bosonic and
fermionic Gaussian integrals.  Item four fixes their reciprocal
normalizations.  Apply Theorem~\ref{thm:v3v18-quartet-cancellation} on each
nonzero block and Proposition~\ref{prop:v3v18-heat-supertrace}.  Item five
prevents the undefined zero-mode integrals from being included in the
determinant identity.
\end{proof}

\begin{assessment}[What remains for cutoff removal]
\label{ass:v3v18-cutoff-removal}
Theorem~\ref{thm:v3v18-regulator-package} is exact at every cutoff.  A
continuum theorem still needs uniform physical-sector heat estimates,
local counterterms, convergence of the V9 Schur kernel, a zero-mode
measure, and determinant-line phase control.  These requirements are
separate from the already closed acyclic nonzero quartet cancellation.
\end{assessment}

\subsection{V18 conclusion and V19 target}
\label{sec:v3v18-conclusion}

V18 proves the finite-cutoff quantum cancellation available from the V17
complete-complex reservoir.  One complex boson and one complex ghost pair
with the identical positive operator give reciprocal determinants, and a
chainwise spectral cutoff preserves the equality mode-by-mode.  Boundary
domains are part of the operator: mismatched impedance conditions leave an
explicit determinant ratio.  The heat supertrace of the nonzero acyclic
quartet vanishes, while physical reservoir modes and zero modes remain.

V19 will move from algebraic cutoff cancellation to analytic cutoff
removal.  It will formulate a relative heat-kernel comparison for the
physical boundary/reservoir Schur pair, prove which heat coefficients are
local and subtractable, and identify a trace-class criterion for the
finite relative determinant.  Zero modes and determinant-line phases will
remain explicit gates rather than being absorbed into a primed
determinant.
\clearpage
\section{Third Volume, Version V19: Relative Heat Renormalization and the Schatten Gate for the Physical Reservoir Determinant}
\label{sec:v3v19}

\subsection{From exact quartet cancellation to the physical determinant}
\label{sec:v3v19-purpose}

V18 removes the nonzero acyclic quartet at every chainwise cutoff.  The
physical reservoir does not cancel: integrating it out changes the
physical boundary Hessian by the V9 Schur complement and produces a
determinant.  V19 separates two analytic situations.  A trace-class
relative correction has an ordinary Fredholm determinant.  The
order-minus-one correction expected from a boundary impedance in three
Euclidean dimensions is generally not trace class; it requires local heat
subtractions and a modified determinant.

\subsection{Finite-cutoff Schur determinant identity}
\label{sec:v3v19-finite-Schur}

Let \(X_\Lambda\) be the regulated physical boundary space and
\(Y_\Lambda\) the regulated reservoir space.  Consider the positive block
Hessian
\begin{equation}
 \mathbb A_\Lambda
 =
 \begin{pmatrix}
  G_\Lambda&D_\Lambda^*\\
  D_\Lambda&H_\Lambda
 \end{pmatrix},
\qquad H_\Lambda>0.
\label{eq:v3v19-block-Hessian}
\end{equation}
Define
\begin{equation}
 G_{{\rm eff},\Lambda}
 :=G_\Lambda-D_\Lambda^*H_\Lambda^{-1}D_\Lambda.
\label{eq:v3v19-Geff}
\end{equation}

\begin{theorem}[Exact block determinant]
\label{thm:v3v19-block-determinant}
If \(G_{{\rm eff},\Lambda}>0\), then
\begin{equation}
 \det\mathbb A_\Lambda
 =(\det H_\Lambda)(\det G_{{\rm eff},\Lambda}).
\label{eq:v3v19-block-determinant}
\end{equation}
The normalized physical correction is
\begin{equation}
 G_\Lambda^{-1/2}G_{{\rm eff},\Lambda}G_\Lambda^{-1/2}
 =I-K_\Lambda,
\label{eq:v3v19-IminusK}
\end{equation}
where
\begin{equation}
 K_\Lambda
 :=G_\Lambda^{-1/2}D_\Lambda^*H_\Lambda^{-1}
   D_\Lambda G_\Lambda^{-1/2}
 =T_\Lambda^*T_\Lambda,
\quad
 T_\Lambda:=H_\Lambda^{-1/2}D_\Lambda G_\Lambda^{-1/2}.
\label{eq:v3v19-K}
\end{equation}
\end{theorem}

\begin{proof}
Multiply \(\mathbb A_\Lambda\) on the left and right by the triangular
matrices
\[
 L=
 \begin{pmatrix}I&-D_\Lambda^*H_\Lambda^{-1}\\0&I\end{pmatrix},
 \qquad
 R=
 \begin{pmatrix}I&0\\-H_\Lambda^{-1}D_\Lambda&I\end{pmatrix}.
\]
A direct multiplication gives
\[
 L\mathbb A_\Lambda R
 =
 \begin{pmatrix}G_{{\rm eff},\Lambda}&0\\0&H_\Lambda\end{pmatrix}.
\]
Both triangular matrices have determinant one, proving
\eqref{eq:v3v19-block-determinant}.  Conjugating
\eqref{eq:v3v19-Geff} by \(G_\Lambda^{-1/2}\) gives
\eqref{eq:v3v19-IminusK}--\eqref{eq:v3v19-K}.
\end{proof}

\begin{corollary}[Finite-cutoff coercivity and determinant]
\label{cor:v3v19-finite-margin}
If \(\|T_\Lambda\|\le\rho<1\), then
\begin{equation}
 (1-\rho^2)I\le I-K_\Lambda\le I
\label{eq:v3v19-finite-margin}
\end{equation}
and
\begin{equation}
 \frac{\det G_{{\rm eff},\Lambda}}{\det G_\Lambda}
 =\det(I-K_\Lambda)>0.
\label{eq:v3v19-finite-relative}
\end{equation}
\end{corollary}

\begin{proof}
\(K_\Lambda=T_\Lambda^*T_\Lambda\) is positive and bounded above by
\(\rho^2I\).  Apply the determinant to
\eqref{eq:v3v19-IminusK}.
\end{proof}

\subsection{Ordinary Fredholm determinant criterion}
\label{sec:v3v19-trace-class}

Let \(X\) now be an infinite-dimensional Hilbert space.  Recall that
\(\mathfrak S_1\) and \(\mathfrak S_2\) denote the trace-class and
Hilbert--Schmidt ideals.

\begin{lemma}[Hilbert--Schmidt incidence gives trace-class feedback]
\label{lem:v3v19-HS}
If
\begin{equation}
 T:=H^{-1/2}DG^{-1/2}\in\mathfrak S_2,
\label{eq:v3v19-T-HS}
\end{equation}
then
\begin{equation}
 K:=T^*T\in\mathfrak S_1,
\qquad
 \|K\|_{\mathfrak S_1}=\|T\|_{\mathfrak S_2}^2.
\label{eq:v3v19-K-trace}
\end{equation}
\end{lemma}

\begin{proof}
Let \(s_j(T)\) be the singular values of \(T\).  The eigenvalues of
\(T^*T\) are \(s_j(T)^2\).  Therefore
\[
 \operatorname{Tr}(T^*T)=\sum_js_j(T)^2
 =\|T\|_{\mathfrak S_2}^2<\infty.
\]
\end{proof}

\begin{theorem}[Trace-class physical relative determinant]
\label{thm:v3v19-Fredholm}
Assume \(K\in\mathfrak S_1\), \(K=K^*\ge0\), and
\(\|K\|\le\rho^2<1\).  Then
\begin{equation}
 \det_{\rm F}(I-K)
 :=\prod_{j=1}^\infty(1-\lambda_j(K))
\label{eq:v3v19-Fredholm}
\end{equation}
converges to a strictly positive number and
\begin{equation}
 \log\det_{\rm F}(I-K)
 =-\sum_{m=1}^\infty\frac1m\operatorname{Tr}(K^m).
\label{eq:v3v19-log-Fredholm}
\end{equation}
\end{theorem}

\begin{proof}
Trace class gives
\(\sum_j\lambda_j(K)<\infty\), while
\(0\le\lambda_j(K)\le\rho^2<1\).  The scalar bound
\[
 |\log(1-\lambda)|\le\frac{\lambda}{1-\rho^2}
\]
makes the product and logarithm absolutely convergent.  Expand
\(\log(1-\lambda_j)=-\sum_{m\ge1}\lambda_j^m/m\).
All summands are nonnegative after changing sign, so Tonelli's theorem
interchanges the two sums and gives
\eqref{eq:v3v19-log-Fredholm}.
\end{proof}

\begin{corollary}[Cutoff convergence in trace norm]
\label{cor:v3v19-trace-convergence}
If \(K_\Lambda\to K\) in \(\mathfrak S_1\) and
\(\sup_\Lambda\|K_\Lambda\|\le\rho^2<1\), then
\begin{equation}
 \det(I-K_\Lambda)\longrightarrow\det_{\rm F}(I-K).
\label{eq:v3v19-trace-convergence}
\end{equation}
\end{corollary}

\begin{proof}
The Fredholm determinant is continuous in trace norm.  In the present
positive setting this follows directly by matching the eigenvalue lists in
\(\ell^1\) and using the Lipschitz bound for \(\log(1-\lambda)\) on
\([0,\rho^2]\).
\end{proof}

\subsection{Why the physical boundary correction is usually not trace
class}
\label{sec:v3v19-order}

Let the Euclidean boundary spacetime have dimension \(n\).  Suppose \(G\)
is elliptic of order two and the reservoir Dirichlet-to-Neumann
self-energy \(\Sigma:=D^*H^{-1}D\) has order one.  Then
\begin{equation}
 K=G^{-1/2}\Sigma G^{-1/2}
\label{eq:v3v19-K-order}
\end{equation}
is a pseudodifferential operator of order \(-1\).

\begin{lemma}[Schatten threshold for a negative-order operator]
\label{lem:v3v19-Schatten}
On a compact \(n\)-dimensional manifold, a classical
pseudodifferential operator of order \(-r<0\) belongs to
\(\mathfrak S_p\) whenever
\begin{equation}
 rp>n.
\label{eq:v3v19-Schatten-threshold}
\end{equation}
An elliptic operator of exact order \(-r\) is generically not in
\(\mathfrak S_p\) when \(rp\le n\).
\end{lemma}

\begin{proof}
The singular values of an elliptic order-\(-r\) operator obey the Weyl
scale
\[
 s_j\asymp j^{-r/n}.
\]
Therefore
\(\sum_js_j^p\) has the same convergence threshold as
\(\sum_jj^{-rp/n}\), which converges exactly when \(rp/n>1\).
For a nonelliptic operator, the upper Weyl bound still proves membership
under the strict inequality.
\end{proof}

\begin{corollary}[Three-dimensional impedance needs a modified determinant]
\label{cor:v3v19-dim-three}
For \(n=3\) and \(K\) of order \(-1\),
\begin{equation}
 K\in\mathfrak S_p\quad\hbox{for every }p>3,
\label{eq:v3v19-Sp}
\end{equation}
but \(K\notin\mathfrak S_1\) generically.  The first integer modified
determinant available without additional smoothing is \(p=4\).
\end{corollary}

\begin{proof}
Apply Lemma~\ref{lem:v3v19-Schatten} with \(r=1,n=3\).
\end{proof}

\begin{counterexample}[The raw determinant product diverges]
\label{cex:v3v19-order-minus-one}
Let \(K\) have eigenvalues
\[
 \lambda_j=cj^{-1/3},\qquad 0<c<1.
\]
Then \(K\in\mathfrak S_p\) for every \(p>3\), but
\(\sum_j\lambda_j=\infty\).  Hence
\[
 \sum_j\log(1-\lambda_j)=-\infty
\]
and the unmodified product \(\prod_j(1-\lambda_j)\) vanishes.  This is not
a finite physical determinant; it is the unsubtracted ultraviolet
divergence.
\end{counterexample}

\subsection{Modified determinants}
\label{sec:v3v19-modified}

For an integer \(p\ge2\) and \(K\in\mathfrak S_p\), define
\begin{equation}
 \det{}_p(I-K)
 :=
 \det_{\rm F}\left[
  (I-K)\exp\left(
   K+\frac{K^2}{2}+\cdots+\frac{K^{p-1}}{p-1}
  \right)
 \right].
\label{eq:v3v19-detp}
\end{equation}

\begin{theorem}[Convergent logarithm of the modified determinant]
\label{thm:v3v19-detp}
If \(K\in\mathfrak S_p\) and \(\|K\|<1\), then
\begin{equation}
 \log\det{}_p(I-K)
 =-\sum_{m=p}^\infty\frac1m\operatorname{Tr}(K^m),
\label{eq:v3v19-log-detp}
\end{equation}
and the series is absolutely convergent.
\end{theorem}

\begin{proof}
Schatten Hölder implies \(K^m\in\mathfrak S_1\) for every \(m\ge p\).
Moreover,
\[
 |\operatorname{Tr}(K^m)|
 \le\|K^p\|_{\mathfrak S_1}\|K\|^{m-p}.
\]
The resulting geometric series converges.  The scalar logarithm of the
operator in \eqref{eq:v3v19-detp} begins at power \(K^p\); termwise trace
therefore gives \eqref{eq:v3v19-log-detp}.
\end{proof}

\begin{corollary}[Continuity under Schatten convergence]
\label{cor:v3v19-detp-continuity}
If \(K_\Lambda\to K\) in \(\mathfrak S_p\) and
\(\sup_\Lambda\|K_\Lambda\|\le\rho<1\), then
\begin{equation}
 \det{}_p(I-K_\Lambda)\longrightarrow\det{}_p(I-K).
\label{eq:v3v19-detp-continuity}
\end{equation}
\end{corollary}

\begin{proof}
Schatten Hölder makes every trace
\(\operatorname{Tr}(K_\Lambda^m)\) continuous for \(m\ge p\).
The uniform geometric majorant from the proof of
Theorem~\ref{thm:v3v19-detp} permits dominated convergence in the series.
\end{proof}

\begin{assessment}[The missing low powers are counterterm data]
\label{ass:v3v19-low-powers}
The modified determinant omits
\(\operatorname{Tr}K,\ldots,\operatorname{Tr}K^{p-1}\), precisely the
powers which need not be trace class.  Their renormalized values are not
zero by definition; they correspond to local ultraviolet counterterms and
must be fixed by a renormalization condition.
\end{assessment}

\subsection{Relative heat expansion}
\label{sec:v3v19-relative-heat}

Let \(A_0,A_1\) be positive self-adjoint elliptic operators of order two on
a compact \(n\)-manifold, with elliptic local boundary conditions.  Assume
their principal symbols and bundle ranks agree.  Define
\begin{equation}
 \Theta(t):=\operatorname{Tr}(e^{-tA_1}-e^{-tA_0}).
\label{eq:v3v19-relative-heat}
\end{equation}

\begin{theorem}[Local short-time relative heat coefficients]
\label{thm:v3v19-heat-expansion}
As \(t\downarrow0\),
\begin{equation}
 \Theta(t)
 \sim\sum_{j=0}^\infty c_jt^{(j-n)/2}.
\label{eq:v3v19-heat-expansion}
\end{equation}
Each \(c_j\) is an integral of a local polynomial in the interior symbols,
curvatures, boundary symbols, second fundamental form, and their finite
jets.  If \(A_0,A_1\) and their boundary domains agree in a sector, every
coefficient from that sector cancels.
\end{theorem}

\begin{proof}
Construct the heat resolvent parametrices in interior coordinate charts
and boundary collar charts.  Homogeneous symbol recursion determines each
coefficient from finitely many jets of the operator and boundary symbol.
The inverse Laplace transform and Gaussian integration in cotangent
variables produce powers \(t^{(j-n)/2}\); collar terms account for the odd
values of \(j\).  Taking the trace and subtracting the two parametrices
gives \eqref{eq:v3v19-heat-expansion}.  Identical local data give identical
recursion coefficients, proving sectorwise cancellation.
\end{proof}

\subsection{Renormalized relative determinant}
\label{sec:v3v19-renormalized}

Choose a smooth cutoff \(\chi(t)\) equal to one near \(t=0\) and zero for
large \(t\).  Subtract through the logarithmically divergent coefficient:
\begin{equation}
 \Theta_{\rm ren}(t)
 :=\Theta(t)
 -\chi(t)\sum_{j=0}^{n}c_jt^{(j-n)/2}.
\label{eq:v3v19-renormalized-heat}
\end{equation}

\begin{theorem}[Finite proper-time relative determinant]
\label{thm:v3v19-renormalized-det}
Assume \(A_0,A_1\ge m_0^2I>0\).  Then
\begin{equation}
 \int_0^\infty\frac{|\Theta_{\rm ren}(t)|}{t}\,dt<\infty.
\label{eq:v3v19-proper-time-finite}
\end{equation}
Thus
\begin{equation}
 \log\det_{\rm ren}(A_1,A_0)
 :=-\int_0^\infty\frac{\Theta_{\rm ren}(t)}{t}\,dt
\label{eq:v3v19-renormalized-det}
\end{equation}
is finite for the chosen subtraction scheme.
\end{theorem}

\begin{proof}
After subtracting through \(j=n\), the short-time expansion starts at
\(O(t^{1/2})\), so
\(t^{-1}\Theta_{\rm ren}(t)=O(t^{-1/2})\), which is integrable at zero.
The positive spectral gap gives
\[
 |\Theta(t)|\le
 \operatorname{Tr}(e^{-tA_0})+\operatorname{Tr}(e^{-tA_1})
 \le Ce^{-m_0^2t/2}
\]
for \(t\ge1\), after absorbing the finite heat trace at time \(1/2\) into
\(C\).  The subtraction vanishes there, so the integral is also finite at
infinity.
\end{proof}

\begin{assessment}[Renormalization ambiguity is local]
\label{ass:v3v19-local-ambiguity}
Changing \(\chi\), the subtraction scale, or finite parts changes
\eqref{eq:v3v19-renormalized-det} by a linear combination of the local
coefficients \(c_0,\ldots,c_n\).  These are precisely local bulk and
boundary counterterms.  The finite modified determinant controls the
nonlocal remainder after those local choices are fixed.
\end{assessment}

\begin{corollary}[Exact quartet cancellation survives heat subtraction]
\label{cor:v3v19-quartet-heat}
For a V18 nonzero quartet with identical operators and chainwise boundary
domains,
\begin{equation}
 \Theta_{\rm q}(t)=0
\label{eq:v3v19-quartet-heat}
\end{equation}
at every finite cutoff and in every heat coefficient.  No quartet
counterterm is required in that exactly paired sector.
\end{corollary}

\begin{proof}
The bosonic and fermionic heat traces agree eigenvalue by eigenvalue with
opposite graded multiplicity.  Apply
Theorem~\ref{thm:v3v19-heat-expansion}.
\end{proof}

\subsection{Zero-mode and phase gates}
\label{sec:v3v19-gates}

\begin{firewall}[The gap hypothesis excludes zero modes]
\label{fw:v3v19-zero}
Theorem~\ref{thm:v3v19-renormalized-det} assumes
\(A_i\ge m_0^2I\).  If zero modes are present, subtract their projections
from the heat trace and specify their collective-coordinate measure before
forming a primed determinant.  The proper-time subtraction does not supply
that measure.
\end{firewall}

\begin{firewall}[Positive relative determinants do not fix chiral phases]
\label{fw:v3v19-phase}
The operators in V19 are positive Euclidean Hessians.  Chiral Standard
Model fermions produce determinant lines and phases governed by spectral
flow and eta data.  The positive Schur determinant and its heat
counterterms do not settle those phases or global anomalies.
\end{firewall}

\begin{firewall}[Schatten membership must be proved for the actual
boundary operator]
\label{fw:v3v19-actual-operator}
The order count leading to
Corollary~\ref{cor:v3v19-dim-three} is a model for an elliptic order-two kinetic
operator plus order-one impedance.  A chosen Einstein--Standard Model
gauge may have mixed orders, constraints, or glancing blocks.  Its actual
normalized Schur correction must be constructed and placed in a Schatten
ideal; notation alone is not evidence.
\end{firewall}

\subsection{V19 conclusion and V20 audit mandate}
\label{sec:v3v19-conclusion}

V19 gives the determinant gate for the physical reservoir.  At finite
cutoff the block determinant factors through the Schur metric
\(G_{\rm eff}\).  If the normalized incidence is Hilbert--Schmidt, the
feedback is trace class and the ordinary Fredholm determinant converges.
For the expected order-minus-one correction on a three-dimensional
boundary, trace class fails generically; \(K\in\mathfrak S_p\) only for
\(p>3\), so a fourth modified determinant and the first three local
counterterms are required.  Relative heat subtraction makes those local
ambiguities explicit.

V20 is the next mandatory companion audit.  After fingerprinting the
current Yang--Mills and Navier--Stokes notes, it will choose between proving
Schatten convergence for a concrete reservoir Dirichlet-to-Neumann map and
building the zero-mode/determinant-line package needed by the chiral
Standard Model sector.  The choice will follow the strongest type-correct
lesson from the companion state.

\section{V20: companion audit and exact cylindrical reservoir determinant}
\label{sec:v3v20}

\subsection{Purpose of this installment}
\label{sec:v3v20-purpose}

V19 identified a determinant gate but deliberately left the physical
boundary operator abstract.  The mandatory fifth-version audit below makes
that abstraction the present bottleneck.  The strongest transferable lesson
from the current Yang--Mills programme is that an order or capacity assigned
before the actual intermediate output is formed can give a formally plausible
but false estimate.  Here the analogous unsafe move would be to infer the
determinant from the nominal order of a reservoir impedance without first
constructing that impedance and separating its harmonic mode.

Accordingly, V20 goes upstream to a reservoir for which the entire
Dirichlet-to-Neumann map can be diagonalized.  This has three purposes.  It
proves, rather than assumes, the Schatten threshold used in V19; it identifies
the three divergent traces which the fourth modified determinant removes; and
it distinguishes the positive Euclidean energy of a trace-constrained
reservoir from the negative Schur correction of an unconstrained bilinear
coupling.  This is a benchmark block, not yet the complete Einstein--Standard
Model boundary operator.

\subsection{Mandatory V20 companion audit}
\label{sec:v3v20-audit}

The audit was performed before the new argument was selected.  The inspected
files and their exact fingerprints were
\begin{center}
\begin{tabular}{|p{0.30\textwidth}|p{0.12\textwidth}|p{0.43\textwidth}|}
\hline
programme and active file & bytes & SHA--256 \\
\hline
Navier--Stokes,
\texttt{Renewal\_NS\_Stability\_Research\_Notes\_V31.tex}
& 887537
& \texttt{F1121B62238AD5491BB7CF03635EA9934942EDF573ED8FAAA9EE79E1D38A6696}
\\
\hline
Yang--Mills,
\texttt{Renewal\_Yang\_Mills\_Mass\_Gap\_Research\_Notes\_V84.tex}
& 1325375
& \texttt{01C79C9BD6CD741EC512C456115CB37DF1A5FAF042EBA1EEF9F35EAF72FC37AB}
\\
\hline
\end{tabular}
\end{center}
The Navier--Stokes file is unchanged from the V15 audit.  The Yang--Mills
file has advanced.  Its V83--V84 analysis shows that one must first form the
complete row-partition state and its actual selected physical bank.  It also
gives an exact \(SU(2)\) counterexample to replacing a low intermediate
representation output by a smaller crossing stock while retaining a nominal
square denominator.  The surviving currency is representation-dimension
weighted output, and repeated mass ascent does not manufacture a fresh bank.

\begin{transferprinciple}[Operator analogue of the V84 currency rule]
\label{tp:v3v20-currency}
For the SM--Einstein boundary programme, first construct the complete
reservoir extension and its Dirichlet-to-Neumann operator.  Only then take
singular values, heat traces, or determinants.  The nonzero high-frequency
tail and the zero boundary harmonic mode are different currencies and must
not be substituted for one another.
\end{transferprinciple}

No Yang--Mills or Navier--Stokes theorem is imported as a theorem about the
present operator.  The audit changes only the order of proof: actual operator,
then spectral currency, then summation.

\subsection{The massive half-cylinder and its exact Poisson extension}
\label{sec:v3v20-cylinder}

Let
\[
 N=\mathbb T^3=(\mathbb R/2\pi\mathbb Z)^3,
 \qquad
 X=N\times[0,\infty),
\]
with the flat product metric.  On \(L^2(N)\) set
\begin{equation}
 A_\mu=-\Delta_N+\mu^2,
 \qquad \mu>0.
\label{eq:v3v20-Amu}
\end{equation}
For \(k\in\mathbb Z^3\), let
\(e_k(x)=(2\pi)^{-3/2}e^{ik\cdot x}\).  Then
\begin{equation}
 A_\mu e_k=a_k e_k,
 \qquad a_k=|k|^2+\mu^2.
\label{eq:v3v20-spectrum-A}
\end{equation}
For a field \(\psi\) on \(X\), define the reservoir energy
\begin{equation}
 \mathcal E_{\rm res}(\psi)
 =\frac12\int_0^\infty
 \left(\|\partial_r\psi(r)\|_{L^2(N)}^2
       +\langle A_\mu\psi(r),\psi(r)\rangle_{L^2(N)}\right)\,dr.
\label{eq:v3v20-res-energy}
\end{equation}

\begin{theorem}[Exact Poisson extension and Dirichlet energy]
\label{thm:v3v20-poisson}
For every \(q\in H^{1/2}(N)\), the unique finite-energy critical point of
\eqref{eq:v3v20-res-energy} with trace \(\psi(0)=q\) and decay at infinity is
\begin{equation}
 \psi_q(r)=e^{-rA_\mu^{1/2}}q.
\label{eq:v3v20-poisson}
\end{equation}
It is the unique minimizer, and
\begin{equation}
 \inf_{\operatorname{Tr}\psi=q}\mathcal E_{\rm res}(\psi)
 =\mathcal E_{\rm res}(\psi_q)
 =\frac12\langle A_\mu^{1/2}q,q\rangle.
\label{eq:v3v20-dirichlet-energy}
\end{equation}
With the outward normal of \(X\) at \(r=0\), the positive
Dirichlet-to-Neumann map is
\begin{equation}
 \Sigma_\mu q=-\partial_r\psi_q(0)=A_\mu^{1/2}q.
\label{eq:v3v20-DtN}
\end{equation}
It maps \(H^s(N)\) continuously to \(H^{s-1}(N)\).
\end{theorem}

\begin{proof}
Write \(q=\sum_k q_ke_k\).  The Euler equation is
\(-\partial_r^2\psi+A_\mu\psi=0\).  In the \(k\)-th mode its decaying
solution with boundary value \(q_k\) is
\(q_ke^{-\sqrt{a_k}r}\).  Thus the mode sum is
\eqref{eq:v3v20-poisson}.  Direct integration gives
\[
 \frac12\int_0^\infty
 \left(a_k|q_k|^2e^{-2\sqrt{a_k}r}
      +a_k|q_k|^2e^{-2\sqrt{a_k}r}\right)dr
 =\frac12\sqrt{a_k}|q_k|^2.
\]
Summing proves \eqref{eq:v3v20-dirichlet-energy}; finiteness is precisely
the spectral characterization of \(H^{1/2}(N)\).

For any other admissible extension write \(\psi=\psi_q+\eta\) with
\(\eta(0)=0\).  Integration by parts, first for smooth compactly supported
fields and then by density, makes the cross term vanish because
\((-\partial_r^2+A_\mu)\psi_q=0\).  Hence
\[
 \mathcal E_{\rm res}(\psi)
 =\mathcal E_{\rm res}(\psi_q)+\mathcal E_{\rm res}(\eta).
\]
Since \(A_\mu\ge\mu^2I\), equality forces \(\eta=0\).  Differentiating
\eqref{eq:v3v20-poisson} at the boundary proves
\eqref{eq:v3v20-DtN}.  Finally, multiplication of the \(k\)-th Fourier
coefficient by \(\sqrt{a_k}\) proves the Sobolev mapping statement.
\end{proof}

\subsection{The physical trace attachment and its sign}
\label{sec:v3v20-physical-sign}

Take the bare boundary quadratic form
\begin{equation}
 \mathcal E_{\rm bd}(q)=\frac12\langle A_\mu q,q\rangle
\label{eq:v3v20-bare-boundary}
\end{equation}
and impose the physical trace relation
\(\psi(0)=\sqrt{\rho}\,q\), where \(\rho>0\).  Eliminating the reservoir by
Theorem~\ref{thm:v3v20-poisson} gives
\begin{equation}
 \mathcal E_{\rm eff}(q)
 =\frac12\langle B_{\mu,\rho}q,q\rangle,
 \qquad
 B_{\mu,\rho}=A_\mu+\rho A_\mu^{1/2}.
\label{eq:v3v20-effective-Hessian}
\end{equation}

\begin{proposition}[Positive impedance and the exact gap]
\label{prop:v3v20-positive-gap}
The trace-constrained reservoir adds a positive impedance.  More precisely,
\begin{equation}
 B_{\mu,\rho}e_k
 =b_ke_k,
 \qquad
 b_k=a_k+\rho\sqrt{a_k},
\label{eq:v3v20-b-spectrum}
\end{equation}
and
\begin{equation}
 B_{\mu,\rho}\ge(\mu^2+\rho\mu)I.
\label{eq:v3v20-gap}
\end{equation}
The constant mode attains equality, so the lower bound is sharp.
\end{proposition}

\begin{proof}
Equation~\eqref{eq:v3v20-b-spectrum} follows from the common Fourier
eigenbasis.  The function \(x\mapsto x+\rho\sqrt{x}\) is strictly increasing
for \(x>0\), and \(\min_k a_k=\mu^2\) at \(k=0\).  This proves the claim.
\end{proof}

\begin{proposition}[Why the V19 negative Schur sign is a different model]
\label{prop:v3v20-sign-distinction}
Let \(G,H>0\) and consider an unconstrained block form
\begin{equation}
 \frac12\langle Gq,q\rangle
 +\operatorname{Re}\langle Dq,z\rangle
 +\frac12\langle Hz,z\rangle.
\label{eq:v3v20-unconstrained-block}
\end{equation}
Minimization over the independent variable \(z\) gives the Schur Hessian
\begin{equation}
 G-D^*H^{-1}D.
\label{eq:v3v20-negative-schur}
\end{equation}
By contrast, minimization of a positive reservoir energy subject to a trace
constraint gives \eqref{eq:v3v20-effective-Hessian}, with a plus sign.
Consequently a negative Schur threshold from V9 or V19 cannot be assigned to
the trace-constrained physical reservoir without deriving an independent
bilinear coupling of the form \eqref{eq:v3v20-unconstrained-block}.
\end{proposition}

\begin{proof}
Complete the square:
\[
 \frac12\|H^{1/2}(z+H^{-1}Dq)\|^2
 +\frac12\langle(G-D^*H^{-1}D)q,q\rangle.
\]
This proves \eqref{eq:v3v20-negative-schur}.  Equation
\eqref{eq:v3v20-dirichlet-energy} proves the constrained statement because
the minimized positive reservoir energy is added to the bare energy.
\end{proof}

\begin{firewall}[Coupling signs belong to actions, not analogies]
\label{fw:v3v20-sign}
The plus attachment and the negative Schur complement are both correct in
their respective variational problems.  They are not interchangeable
``impedance conventions.''  Every later SM--Einstein block must display its
quadratic action or boundary constraint before using either determinant.
\end{firewall}

\subsection{Exact normalized spectrum and Schatten threshold}
\label{sec:v3v20-schatten}

The determinant ratio suggested by
\eqref{eq:v3v20-effective-Hessian} is normalized by the bare Hessian:
\begin{equation}
 A_\mu^{-1/2}B_{\mu,\rho}A_\mu^{-1/2}
 =I+K_{\mu,\rho},
 \qquad
 K_{\mu,\rho}=\rho A_\mu^{-1/2}.
\label{eq:v3v20-normalized-K}
\end{equation}
The complete operator was formed before taking its spectral currency.

\begin{lemma}[Dyadic lattice count]
\label{lem:v3v20-lattice}
For \(j\ge1\), let
\[
 \mathcal A_j
 =\{k\in\mathbb Z^3:2^j\le |k|_\infty<2^{j+1}\}.
\]
There are constants \(0<c<C<\infty\), independent of \(j\), such that
\begin{equation}
 c,2^{3j}\le\#\mathcal A_j\le C,2^{3j}.
\label{eq:v3v20-lattice-count}
\end{equation}
Moreover \(2^j\le |k|\le\sqrt3,2^{j+1}\) on \(\mathcal A_j\).
\end{lemma}

\begin{proof}
The number of lattice points in the cube \(|k|_\infty<R\) is
\((2\lceil R\rceil-1)^3\).  Subtract the counts for radii \(2^{j+1}\) and
\(2^j\) and absorb the endpoint convention into fixed constants.  The norm
comparison is elementary.
\end{proof}

\begin{theorem}[Sharp ideal membership of the physical correction]
\label{thm:v3v20-sharp-schatten}
The normalized reservoir correction is positive and compact, with
\begin{equation}
 K_{\mu,\rho}e_k=\kappa_ke_k,
 \qquad
 \kappa_k=\frac{\rho}{\sqrt{|k|^2+\mu^2}}.
\label{eq:v3v20-kappa}
\end{equation}
For every \(p>0\),
\begin{equation}
 K_{\mu,\rho}\in\mathfrak S_p
 \quad\Longleftrightarrow\quad p>3.
\label{eq:v3v20-iff-sp}
\end{equation}
In particular it is neither trace class nor Hilbert--Schmidt nor in
\(\mathfrak S_3\), but it lies in \(\mathfrak S_4\).
\end{theorem}

\begin{proof}
Equation~\eqref{eq:v3v20-kappa} follows from
\eqref{eq:v3v20-spectrum-A}.  Its eigenvalues tend to zero, proving
compactness.  For all sufficiently large \(j\), Lemma
\ref{lem:v3v20-lattice} gives two constants such that
\[
 c_p2^{(3-p)j}
 \le \sum_{k\in\mathcal A_j}\kappa_k^p
 \le C_p2^{(3-p)j}.
\]
The sum over dyadic shells is a geometric series precisely when \(p>3\).
At \(p=3\) every large shell contributes an amount bounded below by a
positive constant, and for \(p<3\) its contribution grows.  This proves the
equivalence.
\end{proof}

This theorem supplies the actual-operator proof required by Firewall
\ref{fw:v3v19-actual-operator}.  The answer agrees with the order-minus-one
prediction of V19, but agreement is now a conclusion rather than an input.

\subsection{Cutoff determinant and the three unpaid traces}
\label{sec:v3v20-det4}

Let \(P_R\) project onto the Fourier modes with \(|k|_\infty\le R\), and put
\(K_R=P_RK_{\mu,\rho}P_R\).  The exact finite-cutoff determinant ratio is
\begin{equation}
 \frac{\det(P_RB_{\mu,\rho}P_R)}
      {\det(P_RA_\mu P_R)}
 =\det(I+K_R)
 =\prod_{|k|_\infty\le R}(1+\kappa_k).
\label{eq:v3v20-cutoff-product}
\end{equation}

\begin{proposition}[Divergence degrees of the first three traces]
\label{prop:v3v20-three-traces}
As \(R\to\infty\),
\begin{align}
 \operatorname{Tr}K_R&\asymp R^2,
 &\operatorname{Tr}K_R^2&\asymp R,
 &\operatorname{Tr}K_R^3&\asymp\log R.
\label{eq:v3v20-trace-growth}
\end{align}
For every integer \(m\ge4\), \(\operatorname{Tr}K_{\mu,\rho}^m<\infty\).
Here \(f(R)\asymp g(R)\) means two-sided bounds by positive constants for
all sufficiently large \(R\); no unproved lattice asymptotic coefficient is
being asserted.
\end{proposition}

\begin{proof}
Choose \(J\) with \(2^J\le R<2^{J+1}\).  By the estimates in the proof of
Theorem~\ref{thm:v3v20-sharp-schatten}, the contribution of the \(j\)-th
complete dyadic shell to \(\operatorname{Tr}K_R^m\) is comparable to
\(2^{(3-m)j}\).  Summing from a fixed shell through \(J\) gives respectively
\(2^{2J}\), \(2^J\), and \(J\) for \(m=1,2,3\).  An interior cutoff
\(2^{J-1}\) supplies the lower bound and the enclosing cutoff
\(2^{J+1}\) the upper bound.  For \(m\ge4\), the infinite geometric series
converges.
\end{proof}

\begin{theorem}[Convergence of the fourth modified determinant]
\label{thm:v3v20-det4}
The limit
\begin{equation}
 \log\det{}_4(I+K_{\mu,\rho})
 :=\sum_{k\in\mathbb Z^3}
 \left[
  \log(1+\kappa_k)-\kappa_k
  +\frac12\kappa_k^2-\frac13\kappa_k^3
 \right]
\label{eq:v3v20-det4}
\end{equation}
is absolutely convergent.  Equivalently,
\begin{equation}
 \begin{split}
 \log\det(I+K_R)
 ={}&\operatorname{Tr}K_R-\frac12\operatorname{Tr}K_R^2
     +\frac13\operatorname{Tr}K_R^3\\
 &+\log\det{}_4(I+K_{\mu,\rho})+o(1).
 \end{split}
\label{eq:v3v20-cutoff-renormalization}
\end{equation}
Thus precisely the first three high-frequency trace currencies must be
subtracted from the raw cutoff product.
\end{theorem}

\begin{proof}
Let
\(F(x)=\log(1+x)-x+x^2/2-x^3/3\).  Taylor's theorem gives
\(|F(x)|\le Cx^4\) for \(0\le x\le1/2\).  Only finitely many
\(\kappa_k\) exceed \(1/2\), while Theorem
\ref{thm:v3v20-sharp-schatten} gives \(\sum_k\kappa_k^4<\infty\).
This proves absolute convergence of \eqref{eq:v3v20-det4}.  Applying the
same identity to the finite product \eqref{eq:v3v20-cutoff-product} and
letting \(R\to\infty\) proves
\eqref{eq:v3v20-cutoff-renormalization}.
\end{proof}

\begin{assessment}[What is local in this benchmark]
\label{ass:v3v20-local}
The terms of degrees one, two, and three in
\eqref{eq:v3v20-cutoff-renormalization} are the flat-torus manifestations of
the local counterterms predicted in V19.  The fourth modified determinant is
the cutoff-independent spectral remainder for this specified Fourier
cutoff.  A comparison of cutoff schemes on a curved boundary still requires
the local symbol and heat-coefficient calculation planned for V21.
\end{assessment}

\subsection{The harmonic mode is a separate gate}
\label{sec:v3v20-zero-mode}

The high-frequency ideal and the infrared coercive gap are controlled by
different ends of the same exact spectrum:
\begin{equation}
 \|K_{\mu,\rho}\|=\kappa_0=\frac{\rho}{\mu},
 \qquad
 \inf\sigma(B_{\mu,\rho})=\mu^2+\rho\mu.
\label{eq:v3v20-two-ends}
\end{equation}

\begin{proposition}[Massless degeneration despite ultraviolet
renormalizability]
\label{prop:v3v20-massless}
For every fixed \(\mu>0\), the correction belongs to \(\mathfrak S_4\) and
\(B_{\mu,\rho}>0\).  As \(\mu\downarrow0\), however,
\begin{equation}
 \|K_{\mu,\rho}\|\longrightarrow\infty,
 \qquad
 \inf\sigma(B_{\mu,\rho})\longrightarrow0.
\label{eq:v3v20-massless-collapse}
\end{equation}
At \(\mu=0\), the constant mode lies in the kernel of both
\(-\Delta_N\) and its square root; the normalized operator
\((-\Delta_N)^{-1/2}\) is not defined on that mode.
\end{proposition}

\begin{proof}
The first sentence is Theorem~\ref{thm:v3v20-sharp-schatten} and
Proposition~\ref{prop:v3v20-positive-gap}.  The limits follow from
\eqref{eq:v3v20-two-ends}.  The final assertion follows directly from the
Fourier eigenvalue at \(k=0\).
\end{proof}

\begin{firewall}[A UV determinant does not pay for a zero mode]
\label{fw:v3v20-uv-ir}
The convergence of \(\det_4\) is a statement about the large-
\(|k|\) tail.  It neither removes the massless constant mode nor supplies
its collective-coordinate measure.  A primed determinant may omit that
mode only after the physical quotient, gauge volume, or modulus measure has
been specified.
\end{firewall}

\begin{counterexample}[Nominal order misses the gap]
\label{cex:v3v20-order-gap}
Fix \(\rho>0\) and take \(\mu_n\downarrow0\).  Every
\(K_{\mu_n,\rho}\) has the same order minus one and the same sharp
Schatten membership \(p>3\).  Nevertheless the normalized zero-mode
eigenvalue tends to infinity and the absolute Hessian gap tends to zero.
Thus the order class alone cannot give a cutoff-uniform inverse or a
massless determinant.
\end{counterexample}

\subsection{V20 conclusion and V21 target}
\label{sec:v3v20-conclusion}

V20 has constructed a complete physical reservoir block before taking any
spectral trace.  On the flat massive half-cylinder the exact
Dirichlet-to-Neumann map is \(A_\mu^{1/2}\).  Its trace-constrained Euclidean
attachment adds \(\rho A_\mu^{1/2}\) to the boundary Hessian.  The normalized
correction \(\rho A_\mu^{-1/2}\) lies in \(\mathfrak S_p\) exactly for
\(p>3\), and the finite cutoff determinant separates into three divergent
traces and an absolutely convergent fourth modified determinant.  The
constant mode independently controls the coercive gap and becomes singular
in the massless limit.

This resolves the concrete-operator alternative set in V19 and records an
essential sign correction: the negative Schur complement belongs to an
unconstrained bilinear auxiliary field, whereas the positive
Dirichlet-to-Neumann energy belongs to the trace-constrained reservoir used
here.  V21 should transfer this benchmark to a compact curved
three-boundary.  Its proof obligations are to construct the curved
Dirichlet-to-Neumann map, identify its principal and subprincipal symbols,
prove the same Schatten threshold for the actual normalized operator, and
show that changes of local cutoff alter only the first three local
counterterms.  Only after that geometric step should the programme return to
matrix-valued SM polarizations and the chiral determinant-line phase.

\section{V21: curved product collars, Weyl currency, and local regulator
dependence}
\label{sec:v3v21}

\subsection{Scope and geometric setup}
\label{sec:v3v21-setup}

V20 gave an exact Fourier calculation on \(\mathbb T^3\).  The present
installment removes flatness while retaining a product collar, so that the
reservoir operator can still be constructed exactly.  This separates the
spectral issue from the nonproduct Riccati correction that will be treated in
V22.

Let \((N,h)\) be a closed smooth Riemannian three-manifold.  Let \(E\to N\)
be a Hermitian vector bundle of finite rank \(r\), and let
\begin{equation}
 A=\nabla^*\nabla+\mathcal V
\label{eq:v3v21-A}
\end{equation}
be a positive self-adjoint Laplace-type operator on \(L^2(N;E)\), with
smooth self-adjoint endomorphism \(\mathcal V\) and
\begin{equation}
 A\ge a_0 I,\qquad a_0>0.
\label{eq:v3v21-gap}
\end{equation}
The half-cylinder is \(X=N\times[0,\infty)\) with product metric and pulled
back bundle.  The reservoir quadratic form is
\begin{equation}
 \mathcal E_X(\psi)
 =\frac12\int_0^\infty
 \bigl(\|\partial_r\psi\|^2+\langle A\psi,\psi\rangle\bigr)\,dr.
\label{eq:v3v21-energy}
\end{equation}

\subsection{Exact curved product Dirichlet-to-Neumann map}
\label{sec:v3v21-DtN}

\begin{theorem}[Spectral Poisson extension on a curved product collar]
\label{thm:v3v21-poisson}
For \(q\in H^{1/2}(N;E)\), the unique finite-energy minimizer of
\eqref{eq:v3v21-energy} with trace \(q\) is
\begin{equation}
 \psi_q(r)=e^{-rA^{1/2}}q.
\label{eq:v3v21-poisson}
\end{equation}
Its energy and positive Dirichlet-to-Neumann map are
\begin{equation}
 \mathcal E_X(\psi_q)=\frac12\langle A^{1/2}q,q\rangle,
 \qquad
 \Sigma=A^{1/2}.
\label{eq:v3v21-DtN}
\end{equation}
\end{theorem}

\begin{proof}
Compactness and ellipticity give an orthonormal eigenbasis
\(\{u_j\}_{j\ge1}\) with \(Au_j=\lambda_ju_j\) and
\(\lambda_j\ge a_0\).  Expanding \(q=\sum_jq_ju_j\), the decaying solution
of
\(-\partial_r^2\psi+A\psi=0\) in the \(j\)-th eigenspace is
\(q_je^{-\sqrt{\lambda_j}r}\).  The same one-dimensional calculation as in
Theorem~\ref{thm:v3v20-poisson} yields
\[
 \mathcal E_X(\psi_q)
 =\frac12\sum_j\sqrt{\lambda_j}|q_j|^2.
\]
This sum is finite exactly on the form domain of \(A^{1/2}\), namely
\(H^{1/2}(N;E)\) with equivalent norm.  If
\(\psi=\psi_q+\eta\) and \(\eta(0)=0\), integration by parts gives
\(\mathcal E_X(\psi)=\mathcal E_X(\psi_q)+\mathcal E_X(\eta)\).
The gap \eqref{eq:v3v21-gap} makes the second term strictly positive unless
\(\eta=0\).  Finally
\(-\partial_r\psi_q(0)=A^{1/2}q\).
\end{proof}

Attach the reservoir by the trace condition
\(\psi(0)=\sqrt{\rho}\,q\), \(\rho>0\), to a boundary kinetic operator \(A\).
The exact effective Hessian and normalized correction are
\begin{equation}
 B_\rho=A+\rho A^{1/2},
 \qquad
 K_\rho=A^{-1/2}(B_\rho-A)A^{-1/2}
 =\rho A^{-1/2}.
\label{eq:v3v21-K}
\end{equation}
Thus \(B_\rho\ge A\ge a_0I\).  The sign distinction proved in
Proposition~\ref{prop:v3v20-sign-distinction} survives curvature without
change.

\subsection{Principal symbol and the absence of an order-zero product
correction}
\label{sec:v3v21-symbol}

To state subprincipal symbols invariantly, conjugate scalar operators to
half-densities and use the geometric Weyl symbol.  For a connection
Laplace-type operator, the degree-two symbol is
\(|\xi|_h^2 I_E\), and its covariant degree-one symbol vanishes.

\begin{proposition}[First two symbol degrees of the product DtN map]
\label{prop:v3v21-symbol}
The operator \(\Sigma=A^{1/2}\) is a classical elliptic
pseudodifferential operator of order one.  In covariant half-density Weyl
calculus,
\begin{equation}
 \sigma_1(\Sigma)(x,\xi)=|\xi|_h I_E,
 \qquad
 \operatorname{sub}(\Sigma)=0.
\label{eq:v3v21-symbol}
\end{equation}
Consequently \(K_\rho\) has order minus one and principal symbol
\begin{equation}
 \sigma_{-1}(K_\rho)(x,\xi)
 =\rho|\xi|_h^{-1}I_E.
\label{eq:v3v21-K-symbol}
\end{equation}
The potential, connection curvature, and Riemann curvature enter only in
lower homogeneous degrees.
\end{proposition}

\begin{proof}
The parameter-dependent resolvent construction for \(A\), followed by the
Cauchy integral formula for the positive square root, gives a classical
operator \(Q=A^{1/2}\) satisfying \(Q^2=A\).  This can also be obtained
recursively by solving the symbol equation \(q\# q=a\).
At degree two,
\(q_1^2=|\xi|_h^2I_E\); positivity selects
\(q_1=|\xi|_hI_E\).  At degree one, the covariant Weyl product gives
\[
 2|\xi|_h q_0+\frac{1}{2i}\{|\xi|_h,|\xi|_h\}I_E=a_1.
\]
The Poisson bracket of a function with itself is zero and the covariant
subprincipal term \(a_1\) of a connection Laplacian on half-densities is
zero.  Hence \(q_0=0\).  Inverting \(Q\) in the classical calculus gives
principal symbol \(|\xi|_h^{-1}I_E\), proving
\eqref{eq:v3v21-K-symbol}.
\end{proof}

\begin{firewall}[The vanishing subprincipal term uses the product collar]
\label{fw:v3v21-product}
Equation~\eqref{eq:v3v21-symbol} is not a claim about a general moving
collar.  Mean curvature and the first normal variation of the tangential
metric enter the factorization at order zero when \(h=h(r)\).  That term
must be derived from the nonproduct operator, not appended to the product
answer by analogy.
\end{firewall}

\subsection{Weyl law and the sharp curved Schatten result}
\label{sec:v3v21-weyl}

\begin{lemma}[Weyl currency for a rank-\(r\) Laplace-type operator]
\label{lem:v3v21-weyl}
Let
\[
 N_A(L)=\#\{j:\lambda_j\le L\},
\]
counting multiplicity.  Then
\begin{equation}
 N_A(L)
 =C_{N,E}L^{3/2}+O(L),
 \qquad
 C_{N,E}=\frac{r\,\operatorname{Vol}_h(N)}{6\pi^2}.
\label{eq:v3v21-weyl}
\end{equation}
In particular,
\begin{equation}
 \lambda_j\sim (j/C_{N,E})^{2/3}.
\label{eq:v3v21-eigen-asymptotic}
\end{equation}
\end{lemma}

\begin{proof}
In a coordinate chart and a local trivialization, freeze the principal
symbol \(|\xi|_h^2I_E\) on cubes of side \(\delta\).  Dirichlet--Neumann
bracketing compares \(A\), up to \(O(\delta)\) changes of its quadratic
form and lower-order bounded terms, with constant-coefficient operators on
those cubes.  Their lattice counts equal
\((2\pi)^{-3}\) times the phase-space volume
\[
 r\int_N\int_{|\xi|_h^2\le L}d\xi\,d\operatorname{vol}_h
 =r\,\operatorname{Vol}_h(N)\frac{4\pi}{3}L^{3/2}.
\]
The standard boundary-cube estimate is \(O(L)\) in dimension three.
First let \(L\to\infty\) at fixed \(\delta\), and then let
\(\delta\downarrow0\), to obtain
\eqref{eq:v3v21-weyl}.  Monotonic inversion of the counting function gives
\eqref{eq:v3v21-eigen-asymptotic}.
\end{proof}

\begin{theorem}[Curved product Schatten threshold]
\label{thm:v3v21-schatten}
For \(p>0\),
\begin{equation}
 K_\rho\in\mathfrak S_p
 \quad\Longleftrightarrow\quad p>3.
\label{eq:v3v21-schatten}
\end{equation}
The fourth modified determinant
\begin{equation}
 \log\det{}_4(I+K_\rho)
 =\sum_j\left[
 \log(1+\rho\lambda_j^{-1/2})
 -\rho\lambda_j^{-1/2}
 +\frac{\rho^2}{2}\lambda_j^{-1}
 -\frac{\rho^3}{3}\lambda_j^{-3/2}
 \right]
\label{eq:v3v21-det4}
\end{equation}
is absolutely convergent.
\end{theorem}

\begin{proof}
The singular values of \(K_\rho\) are
\(\rho\lambda_j^{-1/2}\).  By
\eqref{eq:v3v21-eigen-asymptotic}, their \(p\)-th powers are comparable,
for large \(j\), to \(j^{-p/3}\).  The resulting series converges exactly
for \(p>3\).  Taylor's theorem bounds the bracket in
\eqref{eq:v3v21-det4} by \(C\rho^4\lambda_j^{-2}\) once
\(\rho\lambda_j^{-1/2}\le1/2\).  The exceptional set is finite and
\(\sum_j\lambda_j^{-2}<\infty\) by the already proved criterion with
\(p=4\).
\end{proof}

\subsection{Sharp spectral cutoffs and universal leading currencies}
\label{sec:v3v21-sharp-cutoff}

Let \(P_\Lambda=\mathbf 1_{[0,\Lambda^2]}(A)\), and define
\begin{equation}
 T_m(\Lambda)=\operatorname{Tr}(P_\Lambda A^{-m/2}).
\label{eq:v3v21-Tm}
\end{equation}

\begin{proposition}[Leading curved cutoff divergences]
\label{prop:v3v21-leading}
As \(\Lambda\to\infty\),
\begin{align}
 T_1(\Lambda)
 &=\frac{r\,\operatorname{Vol}_h(N)}{4\pi^2}\Lambda^2+O(\Lambda),
\label{eq:v3v21-T1}\\
 T_2(\Lambda)
 &=\frac{r\,\operatorname{Vol}_h(N)}{2\pi^2}\Lambda+O(\log\Lambda),
\label{eq:v3v21-T2}\\
 T_3(\Lambda)
 &=\frac{r\,\operatorname{Vol}_h(N)}{2\pi^2}\log\Lambda+O(1).
\label{eq:v3v21-T3}
\end{align}
The exact cutoff determinant therefore satisfies
\begin{equation}
 \begin{split}
 \log\det\!\left(I+P_\Lambda K_\rho P_\Lambda\right)
 ={}&\rho T_1(\Lambda)-\frac{\rho^2}{2}T_2(\Lambda)
 +\frac{\rho^3}{3}T_3(\Lambda)\\
 &+\log\det{}_4(I+K_\rho)+o(1).
 \end{split}
\label{eq:v3v21-sharp-det}
\end{equation}
\end{proposition}

\begin{proof}
Stieltjes summation gives
\[
 T_m(\Lambda)
 =\int_{a_0}^{\Lambda^2}L^{-m/2}\,dN_A(L).
\]
Insert \eqref{eq:v3v21-weyl} and integrate the principal term.  Its
coefficient is
\[
 \frac32C_{N,E}\int^{\Lambda^2}L^{(1-m)/2}\,dL.
\]
For \(m=1,2,3\), this yields respectively the displayed
\(\Lambda^2\), \(\Lambda\), and \(\log\Lambda\) coefficients.
Integration by parts against the \(O(L)\) remainder gives
\(O(\Lambda)\), \(O(\log\Lambda)\), and \(O(1)\), respectively.
For \eqref{eq:v3v21-sharp-det}, apply the scalar Taylor identity to each
eigenvalue below the cutoff.  The sum of fourth-order remainders converges
by Theorem~\ref{thm:v3v21-schatten}, so its cutoff tail is \(o(1)\).
\end{proof}

\begin{assessment}[What the sharp Weyl statement does and does not give]
\label{ass:v3v21-sharp}
The leading coefficients in
\eqref{eq:v3v21-T1}--\eqref{eq:v3v21-T3} are universal volume terms.
The \(O(\Lambda)\) and \(O(\log\Lambda)\) remainders of a sharp spectral
cutoff can contain spectral oscillations.  They must not be advertised as a
local heat expansion.  Locality is obtained with a smooth
pseudodifferential or proper-time regulator, as follows.
\end{assessment}

\subsection{Why regulator dependence is confined to three local terms}
\label{sec:v3v21-locality}

Call \(\chi\in C^\infty([0,\infty))\) admissible if
\(\chi=1\) near zero and \(\chi\) is rapidly decreasing with all
derivatives.  Define the smoothly regulated relative logarithm
\begin{equation}
 \mathcal L_{\chi,\Lambda}
 =\operatorname{Tr}\!\left[
 \chi(A/\Lambda^2)\log(I+\rho A^{-1/2})
 \right].
\label{eq:v3v21-smooth-log}
\end{equation}

\begin{theorem}[Finite local regulator dependence]
\label{thm:v3v21-local-regulator}
For every admissible \(\chi\),
\begin{equation}
 \begin{split}
 \mathcal L_{\chi,\Lambda}
 ={}&
 \rho\,\operatorname{Tr}\!\left[\chi(A/\Lambda^2)A^{-1/2}\right]
 -\frac{\rho^2}{2}\operatorname{Tr}\!\left[\chi(A/\Lambda^2)A^{-1}\right]\\
 &+\frac{\rho^3}{3}\operatorname{Tr}\!\left[\chi(A/\Lambda^2)A^{-3/2}\right]
 +\log\det{}_4(I+K_\rho)+o(1).
 \end{split}
\label{eq:v3v21-local-split}
\end{equation}
Each divergent coefficient in the first line and the first term on the
second line is an integral over \(N\) of a universal polynomial in a finite
jet of the complete symbol of \(A\), multiplied by a moment of \(\chi\).
Consequently changing \(\chi\) changes only these three local
counterterm families; the \(\det_4\) remainder is unchanged.
\end{theorem}

\begin{proof}
The scalar identity
\[
 \log(1+x)=x-\frac{x^2}{2}+\frac{x^3}{3}+F(x),
 \qquad F(x)=O(x^4),
\]
applied by functional calculus gives
\[
 \log(I+\rho A^{-1/2})
 =\rho A^{-1/2}-\frac{\rho^2}{2}A^{-1}
 +\frac{\rho^3}{3}A^{-3/2}+F(\rho A^{-1/2}).
\]
The last operator is trace class because its eigenvalues are
\(O(\lambda_j^{-2})\).  Since
\(\chi(A/\Lambda^2)\to I\) strongly and is uniformly bounded, dominated
convergence in trace norm sends its trace to
\(\log\det_4(I+K_\rho)\).

It remains to locate the other terms.  For a classical operator \(P\) of
order \(-m\), the parameter-dependent symbolic composition of
\(\chi(A/\Lambda^2)P\), followed by the diagonal-kernel trace, has the form
\[
 (2\pi)^{-3}\int_N\int_{T_x^*N}
 \operatorname{tr}_E\!
 \left[\chi(|\xi|_h^2/\Lambda^2)
       \sum_{\ell=0}^{M}p_{-m-\ell}(x,\xi)\right]d\xi\,dx
 +o(1)
\]
after enough homogeneous degrees are retained.  Rescaling
\(\xi=\Lambda\eta\) shows that only degrees with
\(-m-\ell\ge-3\) can diverge.  Their coefficients depend on finitely many
symbol jets and on moments of \(\chi\), hence are local.  For
\(P=A^{-1/2},A^{-1},A^{-3/2}\), these are exactly the three families in
\eqref{eq:v3v21-local-split}.  Terms below degree \(-3\) are integrable in
\(\xi\), while the degree \(-3\) term gives the possible logarithm.  A
partition of unity completes the argument globally.
\end{proof}

\begin{firewall}[Locality requires the complete symbol]
\label{fw:v3v21-complete-symbol}
The leading symbol \(\rho|\xi|^{-1}\) proves the Schatten threshold but
does not determine every power and logarithmic counterterm.  Curvature,
bundle curvature, the potential, and later the second fundamental form
enter lower symbol degrees.  They must be retained before the local
counterterm is named.
\end{firewall}

\subsection{V21 conclusion and V22 target}
\label{sec:v3v21-conclusion}

V21 extends the exact physical reservoir calculation to every compact
curved product three-boundary and every finite-rank Laplace-type
polarization bundle with a positive gap.  The Poisson extension gives
\(\Sigma=A^{1/2}\) exactly.  Weyl currency proves
\(K_\rho=\rho A^{-1/2}\in\mathfrak S_p\) exactly for \(p>3\), and hence
gives a canonical fourth modified determinant.  The first three terms in
the logarithmic expansion contain all regulator dependence; with a smooth
regulator their divergent coefficients are local complete-symbol
integrals.  The cutoff-independent spectral remainder is
\(\det_4(I+K_\rho)\).

The remaining geometric bottleneck is the nonproduct collar actually met by
a moving Einstein boundary.  V22 should factor the scalar Laplace-type
operator when \(h=h(r)\), derive the operator Riccati equation for its
Dirichlet-to-Neumann family, and calculate the order-zero mean-curvature
term and the first lower symbol correction.  That calculation will decide
which new local boundary invariants enter before the construction is
promoted to matrix-valued gauge, Higgs, and fermion complexes.

\section{V22: nonproduct collars, Riccati factorization, and extrinsic
symbol currency}
\label{sec:v3v22}

\subsection{The nonproduct geometric bottleneck}
\label{sec:v3v22-setup}

The product hypothesis in V21 suppresses the first normal jet of the
boundary metric.  An Einstein boundary generally has no such suppression.
This installment computes the missing term from the actual scalar
Laplace-type equation.  The result is a direction-dependent order-zero
symbol involving the second fundamental form; its spherical average is a
mean-curvature term.  That distinction matters because taking the trace
before forming the complete symbol would erase anisotropic extrinsic
information.

Let
\[
 X=N\times[0,\infty),\qquad
 g=dr^2+h(r),
\]
where \(N\) is a closed three-manifold and \(h(r)\) is a smooth family of
Riemannian metrics.  The coordinate \(r\) points into the reservoir, so the
outward normal at \(N=N\times\{0\}\) is \(-\partial_r\).  Define
\begin{equation}
 \mathbb I_{ij}(r)=\frac12\partial_rh_{ij}(r),
 \qquad
 H(r)=\operatorname{tr}_{h(r)}\mathbb I(r)
 =\partial_r\log\sqrt{\det h(r)}.
\label{eq:v3v22-II-H}
\end{equation}
These signs are tied to the inward vector \(+\partial_r\).

For clarity the main calculation is scalar.  Let
\begin{equation}
 L=-\Delta_g+\mu^2
 =-\partial_r^2-H(r)\partial_r+A(r),
 \qquad
 A(r)=-\Delta_{h(r)}+\mu^2,
\label{eq:v3v22-L}
\end{equation}
with \(\mu>0\).  Assume that the completed reservoir end has bounded
geometry and that the Dirichlet problem for \(L\) has a unique decaying
finite-energy solution.  These global assumptions select the
Dirichlet-to-Neumann operator; the symbol calculation below is local at the
boundary.

\subsection{Exact Riccati equation}
\label{sec:v3v22-riccati}

For each slice \(N_r=N\times\{r\}\), let \(\Sigma(r)\) be the inward-tail
Dirichlet-to-Neumann operator: if \(u\) is the decaying solution on
\(N\times[r,\infty)\), then
\begin{equation}
 -\partial_ru(r)=\Sigma(r)u(r).
\label{eq:v3v22-DtN-def}
\end{equation}

\begin{theorem}[Operator Riccati equation for the tail DtN family]
\label{thm:v3v22-riccati}
On smooth boundary data, the family \(\Sigma(r)\) satisfies
\begin{equation}
 \partial_r\Sigma
 =\Sigma^2-H\Sigma-A.
\label{eq:v3v22-riccati}
\end{equation}
At \(r=0\), \(\Sigma(0)\) is positive and self-adjoint with respect to
\(L^2(N,d\operatorname{vol}_{h(0)})\), and
\begin{equation}
 \frac12\langle\Sigma(0)q,q\rangle
 =\inf_{\operatorname{Tr}u=q}
 \frac12\int_X\left(|du|_g^2+\mu^2|u|^2\right)
 d\operatorname{vol}_g.
\label{eq:v3v22-energy-DtN}
\end{equation}
\end{theorem}

\begin{proof}
Along a decaying solution, differentiate
\(-u'=\Sigma u\):
\[
 -u''=\Sigma'u+\Sigma u'
 =\Sigma'u-\Sigma^2u.
\]
Substitute this identity and \(u'=-\Sigma u\) directly into
\(Lu=-u''-Hu'+Au=0\):
\[
 0=(\Sigma'-\Sigma^2+H\Sigma+A)u.
\]
Since the boundary value \(u(r)\) is arbitrary in the domain of the tail
Poisson operator, this is equivalent to
\(\Sigma'=\Sigma^2-H\Sigma-A\), which is
\eqref{eq:v3v22-riccati}.

Green's identity and the outward normal \(-\partial_r\) give
\[
 \int_X(|du|^2+\mu^2|u|^2)d\operatorname{vol}_g
 =\int_Nq(-\partial_ru)|_{r=0}\,d\operatorname{vol}_{h(0)}
 =\langle q,\Sigma(0)q\rangle.
\]
The usual zero-trace orthogonal decomposition of the energy gives the
minimum property.  Positivity and symmetry follow from the same identity;
the mass and uniqueness assumptions exclude a kernel.
\end{proof}

\begin{assessment}[Riccati sign check]
\label{ass:v3v22-sign-check}
If \(h\), \(H\), and \(A\) are constant and a mode has
\(Aq=aq\), a decaying solution \(e^{-sr}q\) obeys
\(s^2-Hs-a=0\).  Thus
\[
 s=\frac{H+\sqrt{H^2+4a}}2
 =\sqrt a+\frac H2+O(a^{-1/2}),
\]
which agrees with \eqref{eq:v3v22-riccati} and fixes the sign of the
order-zero mean-curvature contribution.
\end{assessment}

\subsection{Half-density conjugation and the boundary relation}
\label{sec:v3v22-half-density}

The first normal derivative in \eqref{eq:v3v22-L} can be removed without
guesswork.  Locally let
\(J(r,x)=\sqrt{\det h(r,x)}/\sqrt{\det h(0,x)}\), so
\(\partial_r\log J=H\), and set
\begin{equation}
 v=J^{1/2}u.
\label{eq:v3v22-conjugation}
\end{equation}
This identifies the varying slice \(L^2\) spaces with the fixed boundary
half-density space.

\begin{proposition}[Conjugated normal operator]
\label{prop:v3v22-conjugated}
Under \eqref{eq:v3v22-conjugation},
\begin{equation}
 J^{1/2}LJ^{-1/2}
 =-\partial_r^2+\widetilde A(r),
\label{eq:v3v22-conjugated-L}
\end{equation}
where
\begin{equation}
 \widetilde A(r)
 =J^{1/2}A(r)J^{-1/2}
 +\frac14H(r)^2+\frac12\partial_rH(r).
\label{eq:v3v22-Atilde}
\end{equation}
If \(\widetilde\Sigma(r)v=-\partial_rv\) on decaying solutions, then
\begin{equation}
 \Sigma(r)
 =J^{-1/2}\widetilde\Sigma(r)J^{1/2}+\frac12H(r).
\label{eq:v3v22-Sigma-relation}
\end{equation}
\end{proposition}

\begin{proof}
From \(u=J^{-1/2}v\),
\[
 u'=J^{-1/2}(v'-Hv/2)
\]
and
\[
 u''=J^{-1/2}
 \left(v''-Hv'+(H^2/4-H'/2)v\right).
\]
Substitution into \(-u''-Hu'+Au\) cancels the \(v'\) terms and gives
\eqref{eq:v3v22-conjugated-L}--\eqref{eq:v3v22-Atilde}.
Moreover,
\[
 -u'=J^{-1/2}(-v'+Hv/2)
 =J^{-1/2}(\widetilde\Sigma+H/2)v,
\]
which is \eqref{eq:v3v22-Sigma-relation}.
\end{proof}

\subsection{The complete order-zero DtN symbol}
\label{sec:v3v22-symbol}

All symbols below are covariant Weyl symbols on half-densities.  Put
\begin{equation}
 p(r,x,\xi)=|\xi|_{h(r)}
 =\sqrt{h^{ij}(r,x)\xi_i\xi_j}.
\label{eq:v3v22-p}
\end{equation}

\begin{theorem}[Extrinsic order-zero symbol]
\label{thm:v3v22-order-zero}
The boundary operator \(\Sigma=\Sigma(0)\) is a classical elliptic
pseudodifferential operator of order one.  Its degree-one and degree-zero
symbols are
\begin{align}
 \sigma_1(\Sigma)(x,\xi)
 &=p(0,x,\xi),
\label{eq:v3v22-sigma1}\\
 \sigma_0(\Sigma)(x,\xi)
 &=\frac12H
 -\frac12
 \frac{\mathbb I^{ij}\xi_i\xi_j}{|\xi|_h^2},
\qquad \xi\ne0.
\label{eq:v3v22-sigma0}
\end{align}
All tensors on the right of \eqref{eq:v3v22-sigma0} are evaluated at
\(r=0\).  Thus the order-zero symbol is generally anisotropic and cannot be
replaced pointwise by a scalar multiple of \(H\).
\end{theorem}

\begin{proof}
The elliptic factorization of
\eqref{eq:v3v22-conjugated-L} selects a positive classical symbol
\(\widetilde s\sim s_1+s_0+s_{-1}+\cdots\) for
\(\widetilde\Sigma\).  Its Riccati equation is
\[
 \partial_r\widetilde\Sigma
 =\widetilde\Sigma^2-\widetilde A.
\]
The principal symbol of \(\widetilde A\) is \(p^2\), and its covariant
degree-one symbol vanishes.  At degree two the Riccati equation gives
\(s_1^2=p^2\), so positivity selects \(s_1=p\).  At degree one it gives
\[
 \partial_rp=2ps_0.
\]
There is no Poisson-bracket term because \(\{p,p\}=0\).  Hence
\[
 s_0=\frac{\partial_rp}{2p}.
\]
Differentiating \(h^{ij}h_{jk}=\delta^i_k\) and using
\(\partial_rh_{ij}=2\mathbb I_{ij}\) gives
\[
 \partial_rp
 =-\frac{\mathbb I^{ij}\xi_i\xi_j}{p}.
\]
Finally \eqref{eq:v3v22-Sigma-relation} adds \(H/2\), proving
\eqref{eq:v3v22-sigma0}.  Standard elliptic recursion supplies all lower
degrees and proves classicality.
\end{proof}

\begin{corollary}[Mean curvature is the angular average, not the full
symbol]
\label{cor:v3v22-angular-average}
On the unit cotangent sphere \(S_x^*N\),
\begin{equation}
 \frac1{4\pi}\int_{S_x^*N}\sigma_0(\Sigma)(x,\omega)\,dS(\omega)
 =\frac13H(x).
\label{eq:v3v22-angular-average}
\end{equation}
\end{corollary}

\begin{proof}
Rotational invariance gives
\[
 \frac1{4\pi}\int_{S_x^*N}
 \mathbb I^{ij}\omega_i\omega_j\,dS(\omega)
 =\frac13\operatorname{tr}_h\mathbb I=\frac13H.
\]
Insert this in \eqref{eq:v3v22-sigma0}.
\end{proof}

\begin{counterexample}[Mean curvature alone loses directional data]
\label{cex:v3v22-tracefree-II}
At a point choose principal curvatures
\((\kappa,-\kappa,0)\).  Then \(H=0\), but for a covector in the first
principal direction
\(\sigma_0(\Sigma)=-\kappa/2\), while in the second it is
\(\kappa/2\).  Any substitution \(\sigma_0=cH\) gives zero in both
directions and is false before angular integration.
\end{counterexample}

\subsection{Effect on the normalized determinant operator}
\label{sec:v3v22-normalized}

Let \(A_0=-\Delta_{h(0)}+\mu^2\) be the bare boundary kinetic operator and
attach the reservoir with strength \(\rho>0\).  Define
\begin{equation}
 B=A_0+\rho\Sigma,
 \qquad
 K=\rho A_0^{-1/2}\Sigma A_0^{-1/2}.
\label{eq:v3v22-K}
\end{equation}
The positive-energy identity \eqref{eq:v3v22-energy-DtN} gives
\(B\ge A_0>0\).

\begin{proposition}[First two degrees of the normalized correction]
\label{prop:v3v22-K-symbol}
The operator \(K\) is positive, compact, and classical of order minus one.
Its first two homogeneous symbol degrees are
\begin{align}
 \sigma_{-1}(K)
 &=\rho|\xi|_h^{-1},
\label{eq:v3v22-Kminus1}\\
 \sigma_{-2}(K)
 &=\rho|\xi|_h^{-2}
 \left(
 \frac12H
 -\frac12\frac{\mathbb I^{ij}\xi_i\xi_j}{|\xi|_h^2}
 \right).
\label{eq:v3v22-Kminus2}
\end{align}
Consequently
\begin{equation}
 K\in\mathfrak S_p
 \quad\Longleftrightarrow\quad p>3,
\label{eq:v3v22-schatten}
\end{equation}
and \(\det_4(I+K)\) is defined.
\end{proposition}

\begin{proof}
In covariant half-density Weyl calculus, \(A_0^{-1/2}\) has leading symbol
\(|\xi|^{-1}\) and vanishing degree-minus-two term.  At total degree
minus one, composition with \(\sigma_1(\Sigma)=|\xi|\) gives
\(\rho|\xi|^{-1}\).  At degree minus two, the only nonzero contribution is
\(|\xi|^{-1}\sigma_0(\Sigma)|\xi|^{-1}\).  The Weyl brackets among the
scalar reciprocal leading symbols cancel because
\(\{|\xi|^{-1},|\xi|\}=0\).  This proves
\eqref{eq:v3v22-Kminus1}--\eqref{eq:v3v22-Kminus2}.

Positivity follows from
\(\langle Kf,f\rangle
=\rho\langle\Sigma A_0^{-1/2}f,A_0^{-1/2}f\rangle\ge0\).
The Birman--Solomyak singular-value law for a classical order-minus-one
operator on a compact three-manifold, or equivalently the Weyl bracketing
argument applied to its positive elliptic principal symbol, gives
\(s_j(K)\asymp j^{-1/3}\).  Hence
\(\sum_js_j(K)^p\) converges exactly for \(p>3\).
\end{proof}

\begin{firewall}[Positivity does not make the order-zero symbol scalar]
\label{fw:v3v22-positivity}
The quadratic form of \(\Sigma\) is positive, but individual lower
homogeneous symbols need not be pointwise positive.  The tracefree part of
\(\mathbb I\) changes sign with the cotangent direction.  Positivity may be
used for the full operator and its eigenvalues, not to discard those local
symbol components.
\end{firewall}

\subsection{New local counterterm channels}
\label{sec:v3v22-counterterms}

Theorem~\ref{thm:v3v21-local-regulator} applies to the new \(K\) because
\(K\in\mathfrak S_4\).  The symbol
\eqref{eq:v3v22-Kminus2} identifies a local channel absent from the product
calculation.

\begin{proposition}[Where extrinsic curvature first enters]
\label{prop:v3v22-extrinsic-counterterm}
For a smooth radial high-frequency regulator, the order-minus-two part of
\(\operatorname{Tr}K\) produces a power-divergent local coefficient
proportional to
\begin{equation}
 \rho\int_NH\,d\operatorname{vol}_h.
\label{eq:v3v22-H-counterterm}
\end{equation}
Before angular integration its density is the complete anisotropic symbol
\eqref{eq:v3v22-Kminus2}.  Further logarithmic terms receive contributions
from the order-minus-three symbol of \(K\), from the order-minus-three
symbol of \(K^2\), and from the principal symbol of \(K^3\).
\end{proposition}

\begin{proof}
In dimension three, integrating a homogeneous degree-minus-two symbol over
\(|\xi|\lesssim\Lambda\) gives a term of order \(\Lambda\).
By Corollary~\ref{cor:v3v22-angular-average}, the sphere integral of the
parenthesis in \eqref{eq:v3v22-Kminus2} is
\(4\pi H/3\).  Radial integration and the regulator moment supply a
scheme-dependent scalar coefficient, leaving precisely
\eqref{eq:v3v22-H-counterterm}.  Homogeneous degree minus three gives a
radial integral \(\int^\Lambda r^{-1}dr\), hence a logarithm.  The listed
sources are exactly the total degree-minus-three pieces in the first three
terms of
\(\log(I+K)=K-K^2/2+K^3/3+O(K^4)\).
\end{proof}

\begin{assessment}[Relation to the gravitational boundary action]
\label{ass:v3v22-GHY}
The appearance of \(\int_NH\) is structurally compatible with the
Gibbons--Hawking--York channel derived in V10, after translating the normal
orientation.  V22 does not identify their coefficients: the reservoir
coefficient depends on \(\rho\), the regulated field content, and statistics.
Those contributions must be summed over the actual SM--Einstein complex
before any renormalization of Newton's constant or boundary tension is
claimed.
\end{assessment}

\subsection{V22 conclusion and V23 target}
\label{sec:v3v22-conclusion}

V22 removes the product-collar assumption at the first nontrivial symbol
degree.  The exact tail Dirichlet-to-Neumann family obeys
\(\Sigma'=\Sigma^2-H\Sigma-A\).  Its principal symbol is
\(|\xi|_h\), while its order-zero symbol is
\[
 \frac H2-\frac12
 \frac{\mathbb I^{ij}\xi_i\xi_j}{|\xi|_h^2}.
\]
The angular mean is \(H/3\), but the tracefree second fundamental form
survives direction by direction.  After normalization, this becomes the
degree-minus-two symbol of the physical determinant operator.  The sharp
\(p>3\) Schatten threshold and the fourth modified determinant remain
unchanged, while new local extrinsic counterterms appear.

V23 should now promote the scalar factorization to a finite elliptic
BRST complex with matrix-valued lower symbols.  The immediate proof gate is
to show that the chain differential intertwines the Poisson semigroups and
Dirichlet-to-Neumann maps when the bulk Laplacians and boundary domains
really form a Hilbert complex.  That will determine whether the nonzero
quartet cancellation of V18 survives the curved reservoir symbol by symbol,
and will expose the precise curvature or domain defects when it does not.

\section{V23: Hilbert-complex Poisson transport and exact nonzero
superdeterminant cancellation}
\label{sec:v3v23}

\subsection{Why the complex must precede the determinant}
\label{sec:v3v23-purpose}

V18 proved finite-cutoff quartet cancellation under an identical-operator
and identical-domain hypothesis.  V20--V22 have now constructed the actual
reservoir Dirichlet-to-Neumann operator in scalar sectors.  The present
question is whether those operators form a chain when the boundary fields
form an elliptic BRST complex.  The answer is exact on a product collar with
one common Hodge Laplacian construction.  It fails as soon as masses,
couplings, cutoffs, or domains cease to be chainwise compatible.

Let
\begin{equation}
 0\longrightarrow\mathcal H^0
 \xrightarrow{d_0}\mathcal H^1
 \xrightarrow{d_1}\cdots
 \xrightarrow{d_{\ell-1}}\mathcal H^\ell
 \longrightarrow0
\label{eq:v3v23-complex}
\end{equation}
be a closed Hilbert complex arising from an elliptic differential complex
on a closed compact three-manifold \(N\).  Thus
\(d_{q+1}d_q=0\), each \(d_q\) is densely defined and closed, and the Hodge
Laplacian
\begin{equation}
 \Delta_q=d_{q-1}d_{q-1}^*+d_q^*d_q
\label{eq:v3v23-hodge}
\end{equation}
is self-adjoint with compact resolvent.  Put
\begin{equation}
 A_q=\Delta_q+\mu^2I,\qquad \mu>0.
\label{eq:v3v23-Aq}
\end{equation}
The degree-\(q\) product reservoir is
\(-\partial_r^2+A_q\) on \(N\times[0,\infty)\).

\subsection{Intertwining of functional calculus and Poisson extensions}
\label{sec:v3v23-intertwining}

\begin{lemma}[Hodge Laplacians intertwine the differential]
\label{lem:v3v23-hodge-intertwine}
On the natural common core of smooth sections,
\begin{equation}
 d_q\Delta_q=\Delta_{q+1}d_q,
 \qquad
 d_qA_q=A_{q+1}d_q.
\label{eq:v3v23-hodge-intertwine}
\end{equation}
The identities extend to the corresponding graph domains.
\end{lemma}

\begin{proof}
Using \(d_{q+1}d_q=0\),
\[
 d_q\Delta_q
 =d_qd_{q-1}d_{q-1}^*+d_qd_q^*d_q
 =d_qd_q^*d_q,
\]
whereas
\[
 \Delta_{q+1}d_q
 =d_qd_q^*d_q+d_{q+1}^*d_{q+1}d_q
 =d_qd_q^*d_q.
\]
Adding the same scalar mass proves the second identity.  Elliptic
regularity makes smooth sections a core, so closure gives the graph-domain
extension.
\end{proof}

\begin{theorem}[Chainwise Poisson transport]
\label{thm:v3v23-poisson-chain}
For every \(r\ge0\),
\begin{align}
 d_qe^{-rA_q^{1/2}}
 &=e^{-rA_{q+1}^{1/2}}d_q,
\label{eq:v3v23-semigroup-intertwine}\\
 d_qA_q^{1/2}
 &=A_{q+1}^{1/2}d_q.
\label{eq:v3v23-sqrt-intertwine}
\end{align}
Consequently the product-collar Poisson operators
\(\mathcal P_q(r)=e^{-rA_q^{1/2}}\) and the exact
Dirichlet-to-Neumann maps
\(\Sigma_q=A_q^{1/2}\) form chain maps:
\begin{equation}
 d_q\mathcal P_q(r)=\mathcal P_{q+1}(r)d_q,
 \qquad
 d_q\Sigma_q=\Sigma_{q+1}d_q.
\label{eq:v3v23-chain-maps}
\end{equation}
\end{theorem}

\begin{proof}
The resolvent form of Lemma~\ref{lem:v3v23-hodge-intertwine} is
\[
 d_q(A_q+z)^{-1}=(A_{q+1}+z)^{-1}d_q
\]
on the graph domain for every \(z>0\).  Approximate the bounded functions
\(e^{-r\sqrt{x}}\) and \(\sqrt{x}(1+x)^{-1}\) uniformly on
\([\,\mu^2,\infty)\) by linear combinations of resolvent functions, or use
the spectral integral directly.  Functional calculus then gives
\eqref{eq:v3v23-semigroup-intertwine} and
\[
 d_qA_q^{1/2}(I+A_q)^{-1}
 =A_{q+1}^{1/2}(I+A_{q+1})^{-1}d_q.
\]
Removing the resolvent on the natural graph domain yields
\eqref{eq:v3v23-sqrt-intertwine}.  The final identities use the exact
product Poisson construction from Theorem~\ref{thm:v3v21-poisson}.
\end{proof}

\begin{corollary}[Chainwise physical attachment]
\label{cor:v3v23-chain-attachment}
If the same coupling \(\rho>0\) is used in every degree, then
\begin{equation}
 B_q=A_q+\rho A_q^{1/2},
 \qquad
 K_q=\rho A_q^{-1/2}
\label{eq:v3v23-BK}
\end{equation}
obey
\begin{equation}
 d_qB_q=B_{q+1}d_q,
 \qquad
 d_qK_q=K_{q+1}d_q
\label{eq:v3v23-BK-chain}
\end{equation}
on their natural domains.
\end{corollary}

\begin{proof}
Apply Lemma~\ref{lem:v3v23-hodge-intertwine},
Theorem~\ref{thm:v3v23-poisson-chain}, and the same functional-calculus
argument to \(x^{-1/2}\), which is bounded on
\([\,\mu^2,\infty)\).
\end{proof}

\subsection{Exactness of every positive eigensubcomplex}
\label{sec:v3v23-positive-exact}

For \(\lambda\ge0\), let
\begin{equation}
 E_\lambda^q=\ker(\Delta_q-\lambda I).
\label{eq:v3v23-Elambda}
\end{equation}
Ellipticity makes these spaces finite dimensional, and
Lemma~\ref{lem:v3v23-hodge-intertwine} makes
\((E_\lambda^\bullet,d)\) a finite complex.

\begin{lemma}[Positive Hodge eigenspaces are exact]
\label{lem:v3v23-positive-exact}
If \(\lambda>0\), then the complex
\((E_\lambda^\bullet,d)\) is exact.  Hence
\begin{equation}
 \sum_{q=0}^{\ell}(-1)^q\dim E_\lambda^q=0.
\label{eq:v3v23-positive-euler}
\end{equation}
For \(\lambda=0\),
\begin{equation}
 E_0^q=\ker\Delta_q\cong H^q(\mathcal H^\bullet,d),
\label{eq:v3v23-harmonic-cohomology}
\end{equation}
and its dimension is the Betti number \(b_q\).
\end{lemma}

\begin{proof}
Let \(u\in E_\lambda^q\) with \(d_qu=0\).  The Hodge identity gives
\[
 \lambda u
 =d_{q-1}d_{q-1}^*u+d_q^*d_qu
 =d_{q-1}d_{q-1}^*u.
\]
Thus
\(u=d_{q-1}(\lambda^{-1}d_{q-1}^*u)\), and
\(d_{q-1}^*u\in E_\lambda^{q-1}\) because the adjoint also intertwines the
Hodge Laplacians.  Therefore every closed vector is exact.  A finite exact
complex has zero alternating dimension, proving
\eqref{eq:v3v23-positive-euler}.  The \(\lambda=0\) statement is the Hodge
theorem for the elliptic Hilbert complex.
\end{proof}

\subsection{The exact spectral supertrace identity}
\label{sec:v3v23-supertrace}

\begin{theorem}[Only cohomology survives an Euler supertrace]
\label{thm:v3v23-supertrace}
Let \(F:[0,\infty)\to\mathbb C\) be such that
\(F(\Delta_q)\) is trace class in every degree.  Then
\begin{equation}
 \sum_{q=0}^{\ell}(-1)^q
 \operatorname{Tr}F(\Delta_q)
 =\chi(\mathcal H^\bullet,d)\,F(0),
\label{eq:v3v23-supertrace}
\end{equation}
where
\begin{equation}
 \chi(\mathcal H^\bullet,d)
 =\sum_q(-1)^qb_q.
\label{eq:v3v23-euler}
\end{equation}
The same identity holds at a common finite spectral cutoff
\(\mathbf1_{[0,\Lambda^2]}(\Delta_q)\) without a trace-class hypothesis.
\end{theorem}

\begin{proof}
Decompose each trace into the finite-dimensional eigenspaces
\(E_\lambda^q\).  At every \(\lambda>0\), the coefficient of \(F(\lambda)\)
is zero by \eqref{eq:v3v23-positive-euler}.  At \(\lambda=0\), its
coefficient is \(\sum_q(-1)^qb_q=\chi\).  Absolute convergence justifies
the infinite regrouping.  With a common finite cutoff, the sum is finite
and no convergence argument is needed.
\end{proof}

\subsection{Reservoir superdeterminant and harmonic remainder}
\label{sec:v3v23-superdet}

Let \(P_{q,\Lambda}=\mathbf1_{[0,\Lambda^2]}(\Delta_q)\).  Define the
Euler-weighted finite determinant ratio
\begin{equation}
 \mathcal D_\Lambda
 =\prod_{q=0}^{\ell}
 \det\!\left(
 I+P_{q,\Lambda}K_qP_{q,\Lambda}
 \right)^{(-1)^q}.
\label{eq:v3v23-DLambda}
\end{equation}

\begin{theorem}[Exact cancellation of every nonzero reservoir mode]
\label{thm:v3v23-superdet-cancel}
For every \(\Lambda\ge0\),
\begin{equation}
 \mathcal D_\Lambda
 =\left(1+\frac{\rho}{\mu}\right)^{
 \chi(\mathcal H^\bullet,d)}.
\label{eq:v3v23-superdet-cancel}
\end{equation}
In particular, the right side is independent of the ultraviolet cutoff.
Every nonzero eigenvalue cancels before a limit or counterterm is taken.
\end{theorem}

\begin{proof}
On \(E_\lambda^q\), the eigenvalue of \(K_q\) is
\(\rho(\lambda+\mu^2)^{-1/2}\).  Taking the logarithm of
\eqref{eq:v3v23-DLambda} gives the finite supertrace in
Theorem~\ref{thm:v3v23-supertrace} with
\[
 F_\rho(\lambda)
 =\log\left(1+\rho(\lambda+\mu^2)^{-1/2}\right).
\]
Only \(\lambda=0\) remains, and
\(F_\rho(0)=\log(1+\rho/\mu)\).  Exponentiation proves
\eqref{eq:v3v23-superdet-cancel}.
\end{proof}

\begin{corollary}[The first three counterterm supertraces also cancel]
\label{cor:v3v23-counterterm-cancel}
For \(m=1,2,3\), at every common cutoff,
\begin{equation}
 \sum_q(-1)^q
 \operatorname{Tr}\!\left[
 P_{q,\Lambda}K_q^mP_{q,\Lambda}
 \right]
 =\chi(\mathcal H^\bullet,d)
 \left(\frac{\rho}{\mu}\right)^m.
\label{eq:v3v23-power-cancel}
\end{equation}
Thus the nonzero-sector power and logarithmic divergences cancel exactly.
The fourth modified superdeterminant has only the finite harmonic
contribution
\begin{equation}
 \sum_q(-1)^q\log\det{}_4(I+K_q)
 =\chi\left[
 \log\left(1+\frac{\rho}{\mu}\right)
 -\frac{\rho}{\mu}
 +\frac{\rho^2}{2\mu^2}
 -\frac{\rho^3}{3\mu^3}
 \right].
\label{eq:v3v23-det4-harmonic}
\end{equation}
\end{corollary}

\begin{proof}
Use Theorem~\ref{thm:v3v23-supertrace} with
\(F(\lambda)=\rho^m(\lambda+\mu^2)^{-m/2}\) at the common finite cutoff.
For the last formula, apply it to the trace-class fourth-order Taylor
remainder used in \eqref{eq:v3v21-det4}.
\end{proof}

\begin{firewall}[Euler superdeterminant versus the physical one-loop
weight]
\label{fw:v3v23-physical-weight}
Equations~\eqref{eq:v3v23-superdet-cancel} and
\eqref{eq:v3v23-det4-harmonic} use the algebraic parity weight
\((-1)^q\).  They apply to a physical BRST quartet only after its bosonic,
fermionic, real, complex, and square-root multiplicities reduce to that
same weight.  V18 supplies the finite-dimensional bookkeeping gate.
V23 does not replace it.
\end{firewall}

\begin{firewall}[Cohomology is not a cancelled quartet]
\label{fw:v3v23-cohomology}
The harmonic factor in \eqref{eq:v3v23-superdet-cancel} is finite but
physical.  It must be combined with gauge-volume normalization,
collective-coordinate measures, and determinant-line orientation.  It may
not be deleted by extending the positive-eigenvalue cancellation to
\(\lambda=0\).
\end{firewall}

\subsection{Quantitative defect when the chain identity fails}
\label{sec:v3v23-defect}

The exact result is useful only if its failure has a measurable source.
Consider positive self-adjoint operators
\(A_0,A_1\ge a_0I>0\), a chain map \(d\), and the defect
\begin{equation}
 R=A_1d-dA_0.
\label{eq:v3v23-R}
\end{equation}
Assume for this estimate that \(R\) extends to a bounded operator between
the relevant Hilbert spaces; the same formula applies between graph norms
for differential \(d\).

\begin{proposition}[Square-root and Poisson chain-defect bounds]
\label{prop:v3v23-defect}
Put
\begin{equation}
 C=A_1^{1/2}d-dA_0^{1/2}.
\label{eq:v3v23-C}
\end{equation}
Then
\begin{equation}
 C=\frac1\pi\int_0^\infty
 t^{1/2}(A_1+t)^{-1}R(A_0+t)^{-1}\,dt
\label{eq:v3v23-C-integral}
\end{equation}
and
\begin{equation}
 \|C\|\le\frac{\|R\|}{2\sqrt{a_0}}.
\label{eq:v3v23-C-bound}
\end{equation}
Moreover,
\begin{equation}
 \left\|
 e^{-rA_1^{1/2}}d-de^{-rA_0^{1/2}}
 \right\|
 \le
 r e^{-r\sqrt{a_0}}\frac{\|R\|}{2\sqrt{a_0}}.
\label{eq:v3v23-poisson-defect}
\end{equation}
\end{proposition}

\begin{proof}
For \(x>0\),
\[
 \sqrt{x}=\frac1\pi\int_0^\infty
 t^{-1/2}\frac{x}{x+t}\,dt.
\]
The resolvent identity gives
\[
 (A_1+t)^{-1}d-d(A_0+t)^{-1}
 =-(A_1+t)^{-1}R(A_0+t)^{-1}.
\]
Since \(A(A+t)^{-1}=I-t(A+t)^{-1}\), substitution in the square-root
integral yields \eqref{eq:v3v23-C-integral}.  Hence
\[
 \|C\|
 \le\frac{\|R\|}{\pi}
 \int_0^\infty\frac{t^{1/2}}{(a_0+t)^2}\,dt
 =\frac{\|R\|}{2\sqrt{a_0}},
\]
because the rescaled beta integral is
\(\int_0^\infty u^{1/2}(1+u)^{-2}du=\pi/2\).

Duhamel's formula gives
\[
 e^{-rA_1^{1/2}}d-de^{-rA_0^{1/2}}
 =-\int_0^r e^{-(r-s)A_1^{1/2}}
 C e^{-sA_0^{1/2}}\,ds.
\]
Both semigroups have norm at most the corresponding factor
\(e^{-r\sqrt{a_0}}\), which proves
\eqref{eq:v3v23-poisson-defect}.
\end{proof}

\begin{counterexample}[Unequal attachment data destroys modewise
cancellation]
\label{cex:v3v23-unequal}
Take a two-term positive eigensubcomplex in which
\(d:E_\lambda^0\to E_\lambda^1\) is an isomorphism.  Attach parameters
\((\rho_0,\mu_0)\) and \((\rho_1,\mu_1)\).  Its logarithmic contribution is
\[
 \log\left(1+\frac{\rho_0}{\sqrt{\lambda+\mu_0^2}}\right)
 -\log\left(1+\frac{\rho_1}{\sqrt{\lambda+\mu_1^2}}\right),
\]
which is nonzero unless the spectral attachment functions agree.  At high
\(\lambda\) its leading term is
\((\rho_0-\rho_1)\lambda^{-1/2}\).  In dimension three that unpaid
order-minus-one tail has a quadratic cutoff divergence.  Exactness of the
underlying complex alone therefore does not cancel mismatched reservoir
data.
\end{counterexample}

\subsection{Nonproduct warning}
\label{sec:v3v23-nonproduct-warning}

If one applies the scalar Riccati equation of V22 degree by degree and sets
\(C_q=\Sigma_{q+1}d_q-d_q\Sigma_q\), a formal differentiation produces,
even when \(A_{q+1}d_q=d_qA_q\),
\begin{equation}
 \partial_rC_q
 =\Sigma_{q+1}C_q+C_q\Sigma_q-HC_q
 +[d_q,H]\Sigma_q
\label{eq:v3v23-moving-defect}
\end{equation}
when \(d_q\) is independent of \(r\).  The commutator
\([d_q,H]\) need not vanish for a differential operator \(d_q\).

\begin{firewall}[The full normal--tangential complex is required]
\label{fw:v3v23-full-complex}
Equation~\eqref{eq:v3v23-moving-defect} is not an anomaly of the de Rham or
BRST complex.  It signals that a degreewise scalar normal form has omitted
the normal components and connection terms that participate in the full
bulk chain identity.  On a moving collar the programme must factor the
complete block complex before claiming symbolwise quartet cancellation.
\end{firewall}

\subsection{V23 conclusion and V24 target}
\label{sec:v3v23-conclusion}

V23 proves the product-collar chain gate.  A common massive Hodge
construction makes the Poisson semigroups, Dirichlet-to-Neumann maps,
effective Hessians, and normalized determinant operators into chain maps.
Every positive Hodge eigensubcomplex is exact, so all nonzero reservoir
modes cancel in the Euler superdeterminant at every common cutoff.  Only the
finite cohomological factor
\((1+\rho/\mu)^\chi\) remains.  A resolvent formula quantifies the failure
when the Laplacians cease to intertwine.

The next bottleneck is no longer the product determinant.  V24 should build
the full normal--tangential de Rham prototype on a moving collar, including
the second fundamental form in the block differential and adjoint.  It
should prove the bulk chain identity before factorization, derive the
matrix Riccati system, and identify which apparent
\([d,H]\) defects cancel between tangential and normal components.  This is
the smallest rigorous model for the corresponding BRST boundary complex.

\section{V24: the full moving-collar de Rham block and its induced
boundary complex}
\label{sec:v3v24}

\subsection{Why the scalar degreewise split was incomplete}
\label{sec:v3v24-purpose}

V23 isolated the formal defect \([d,H]\Sigma\) that appears if one applies
the scalar Riccati equation separately to tangential form degrees.  The
source is now made explicit.  A bulk \(q\)-form has both a tangential
\(q\)-form and a normal-wedged tangential \((q-1)\)-form.  The moving
Hodge inner product couples their normal adjoints through the full second
fundamental form.  Once these components are retained, the bulk Hodge
Laplacians form an exact chain by construction, and the decaying Poisson
solutions induce a nilpotent boundary differential.

This de Rham calculation is a prototype for the linear BRST complex.  It
does not yet include the nonabelian algebraic ghost terms or the Einstein
trace-reversal block.

\subsection{Tangential--normal splitting and the shape action}
\label{sec:v3v24-splitting}

Use the moving collar and sign convention of V22:
\[
 g=dr^2+h(r),\qquad
 \mathbb I=\frac12\dot h,\qquad H=\operatorname{tr}_h\mathbb I.
\]
Every smooth \(q\)-form has a unique decomposition
\begin{equation}
 \omega=\alpha+dr\wedge\beta,
 \qquad
 \alpha(r)\in\Omega^q(N),\quad
 \beta(r)\in\Omega^{q-1}(N).
\label{eq:v3v24-split}
\end{equation}
Let \(S=h^{-1}\mathbb I\) be the shape endomorphism.  Its induced action on
tangential \(q\)-forms is
\begin{equation}
 \mathcal S_q
 =\sum_{i=1}^3\theta^i\wedge\iota_{S e_i},
\label{eq:v3v24-shape-action}
\end{equation}
in any local \(h(r)\)-orthonormal frame.  Define the self-adjoint
zero-order endomorphism
\begin{equation}
 M_q=H I_{\Lambda^q}-2\mathcal S_q.
\label{eq:v3v24-Mq}
\end{equation}

\begin{lemma}[Variation of the form inner product]
\label{lem:v3v24-inner-variation}
For tangential \(q\)-forms \(a(r),b(r)\),
\begin{equation}
 \partial_r\langle a,b\rangle_{L^2(h(r))}
 =\langle\partial_ra,b\rangle
 +\langle a,\partial_rb\rangle
 +\langle M_qa,b\rangle.
\label{eq:v3v24-inner-variation}
\end{equation}
Consequently the formal adjoint of \(\partial_r\), in the collar integral,
is
\begin{equation}
 (\partial_r)_q^*=-\partial_r-M_q
\label{eq:v3v24-normal-adjoint}
\end{equation}
away from the boundary term.
\end{lemma}

\begin{proof}
The volume density varies by
\(\partial_rd\operatorname{vol}_{h}=H\,d\operatorname{vol}_{h}\).
The inverse metric varies by
\(\partial_rh^{-1}=-2S h^{-1}\); differentiating every contracted covector
index contributes
\(-2\langle\mathcal S_qa,b\rangle\).  Adding the volume variation gives
\eqref{eq:v3v24-inner-variation}.  Integrating that identity in \(r\)
gives \eqref{eq:v3v24-normal-adjoint}.
\end{proof}

\subsection{The exact bulk differential and its adjoint}
\label{sec:v3v24-D}

Let \(d_q=d_N:\Omega^q(N)\to\Omega^{q+1}(N)\), independent of the metric,
and let \(\delta_q(r):\Omega^q(N)\to\Omega^{q-1}(N)\) be its
\(h(r)\)-adjoint.  From \eqref{eq:v3v24-split},
\begin{equation}
 d_X(\alpha+dr\wedge\beta)
 =d_q\alpha+dr\wedge(\partial_r\alpha-d_{q-1}\beta).
\label{eq:v3v24-d-split}
\end{equation}
Thus on column vectors \(U_q=(\alpha,\beta)^{\mathsf T}\),
\begin{equation}
 D_q=
 \begin{pmatrix}
  d_q&0\\
  \partial_r&-d_{q-1}
 \end{pmatrix},
 \qquad
 D_qU_q=\operatorname{split}(d_X\omega).
\label{eq:v3v24-Dq}
\end{equation}

\begin{proposition}[Nilpotence and the full formal adjoint]
\label{prop:v3v24-D-adjoint}
The block differential obeys
\begin{equation}
 D_{q+1}D_q=0.
\label{eq:v3v24-D-square}
\end{equation}
Its collar formal adjoint is
\begin{equation}
 D_q^*=
 \begin{pmatrix}
  \delta_{q+1}&-\partial_r-M_q\\
  0&-\delta_q
 \end{pmatrix}.
\label{eq:v3v24-Dstar}
\end{equation}
\end{proposition}

\begin{proof}
The upper component of \(D_{q+1}D_qU\) is \(d_{q+1}d_q\alpha=0\).
The lower component is
\[
 \partial_rd_q\alpha-d_q(\partial_r\alpha-d_{q-1}\beta)=0,
\]
because \(d_N\) is independent of \(r\) and \(d_qd_{q-1}=0\).
This proves \eqref{eq:v3v24-D-square}.  For the adjoint, the tangential
blocks use \(d_q^*=\delta_{q+1}\), the normal block uses
Lemma~\ref{lem:v3v24-inner-variation}, and the minus sign in
\(-d_{q-1}\) is preserved.  The result is
\eqref{eq:v3v24-Dstar}.
\end{proof}

\begin{lemma}[Variation of the tangential codifferential]
\label{lem:v3v24-delta-dot}
The metric derivative of \(\delta_q\) is
\begin{equation}
 \dot\delta_q=\delta_qM_q-M_{q-1}\delta_q.
\label{eq:v3v24-delta-dot}
\end{equation}
\end{lemma}

\begin{proof}
Differentiate
\(\langle d_{q-1}a,b\rangle_q
=\langle a,\delta_qb\rangle_{q-1}\)
for \(r\)-independent test forms \(a,b\), and use
\eqref{eq:v3v24-inner-variation} on both sides.  Moving the
\(M_q\) term on the left through the adjoint relation gives
\[
 \langle a,\delta_qM_qb\rangle
 =\langle a,M_{q-1}\delta_qb+\dot\delta_qb\rangle.
\]
Arbitrariness of \(a,b\) proves the identity.
\end{proof}

\subsection{The complete moving Hodge-Laplacian block}
\label{sec:v3v24-laplacian}

Let
\begin{equation}
 \Delta_{X,q}=D_{q-1}D_{q-1}^*+D_q^*D_q
\label{eq:v3v24-full-laplacian}
\end{equation}
and
\(\Delta_{N,q}=d_{q-1}\delta_q+\delta_{q+1}d_q\).

\begin{theorem}[Explicit normal--tangential Hodge block]
\label{thm:v3v24-hodge-block}
In the splitting \eqref{eq:v3v24-split},
\begin{equation}
 \Delta_{X,q}=
 \begin{pmatrix}
 -\partial_r^2-M_q\partial_r+\Delta_{N,q}
 &
 M_qd_{q-1}-d_{q-1}M_{q-1}
 \\
 \dot\delta_q
 &
 -\partial_r^2-M_{q-1}\partial_r-\dot M_{q-1}
 +\Delta_{N,q-1}
 \end{pmatrix}.
\label{eq:v3v24-hodge-block}
\end{equation}
The two off-diagonal entries are slice adjoints:
\begin{equation}
 \dot\delta_q
 =(M_qd_{q-1}-d_{q-1}M_{q-1})^*.
\label{eq:v3v24-offdiag-adjoint}
\end{equation}
\end{theorem}

\begin{proof}
Multiply the matrices in
\eqref{eq:v3v24-full-laplacian}.  The upper-left entry is
\[
 d_{q-1}\delta_q+\delta_{q+1}d_q
 +(-\partial_r-M_q)\partial_r
 =\Delta_{N,q}-\partial_r^2-M_q\partial_r.
\]
The upper-right entry is
\[
 d_{q-1}(-\partial_r-M_{q-1})
 -(-\partial_r-M_q)d_{q-1}
 =M_qd_{q-1}-d_{q-1}M_{q-1},
\]
because \([\partial_r,d_N]=0\).  The lower-left entry is
\([\partial_r,\delta_q]=\dot\delta_q\).  The lower-right entry is
\[
 \partial_r(-\partial_r-M_{q-1})+\Delta_{N,q-1}
 =-\partial_r^2-M_{q-1}\partial_r-\dot M_{q-1}
 +\Delta_{N,q-1}.
\]
Finally \eqref{eq:v3v24-offdiag-adjoint} is exactly
Lemma~\ref{lem:v3v24-delta-dot}.
\end{proof}

\begin{corollary}[The omitted scalar terms that repair the formal defect]
\label{cor:v3v24-repair}
If \(\mathbb I\ne0\), the off-diagonal blocks in
\eqref{eq:v3v24-hodge-block} generally do not vanish.  They contain both
\([d,H]\) and the tracefree shape action.  Therefore the scalar
degreewise evolution used only as a warning in
\eqref{eq:v3v23-moving-defect} is not the full Hodge evolution.
\end{corollary}

\begin{proof}
Insert \(M_q=HI-2\mathcal S_q\):
\[
 M_qd-dM_{q-1}
 =- [d,H]-2(\mathcal S_qd-d\mathcal S_{q-1}).
\]
These are precisely the missing first-order tangential couplings.
\end{proof}

\subsection{Bulk chain identity before factorization}
\label{sec:v3v24-bulk-chain}

Add a common mass and set
\begin{equation}
 L_q=\Delta_{X,q}+\mu^2I,\qquad\mu>0.
\label{eq:v3v24-Lq}
\end{equation}

\begin{theorem}[Exact massive bulk chain]
\label{thm:v3v24-bulk-chain}
On the common smooth core,
\begin{equation}
 D_qL_q=L_{q+1}D_q.
\label{eq:v3v24-bulk-chain}
\end{equation}
This identity includes all normal, mean-curvature, and tracefree-shape
terms in \eqref{eq:v3v24-hodge-block}.
\end{theorem}

\begin{proof}
Use \(D_{q+1}D_q=0\):
\[
 D_q\Delta_{X,q}
 =D_qD_{q-1}D_{q-1}^*+D_qD_q^*D_q
 =D_qD_q^*D_q,
\]
whereas
\[
 \Delta_{X,q+1}D_q
 =D_qD_q^*D_q+D_{q+1}^*D_{q+1}D_q
 =D_qD_q^*D_q.
\]
The common scalar mass commutes with \(D_q\).
\end{proof}

\begin{firewall}[Do not verify the chain by selected block terms]
\label{fw:v3v24-selected-blocks}
Individual entries of \eqref{eq:v3v24-hodge-block} contain derivatives of
\(H\), shape actions, and derivatives of \(\delta\).  Cancelling only the
\([d,H]\) term is neither necessary nor sufficient.  The algebraic identity
\eqref{eq:v3v24-bulk-chain} applies to the complete block and is the
type-correct source of the cancellation.
\end{firewall}

\subsection{Matrix Riccati system}
\label{sec:v3v24-matrix-riccati}

Write \eqref{eq:v3v24-Lq} as
\begin{equation}
 L_q=-\partial_r^2-\mathcal M_q\partial_r+\mathcal A_q,
\qquad
 \mathcal M_q=
 \begin{pmatrix}M_q&0\\0&M_{q-1}\end{pmatrix},
\label{eq:v3v24-normal-form}
\end{equation}
where
\begin{equation}
 \mathcal A_q=
 \begin{pmatrix}
 \Delta_{N,q}+\mu^2
 &
 M_qd_{q-1}-d_{q-1}M_{q-1}
 \\
 \dot\delta_q
 &
 \Delta_{N,q-1}+\mu^2-\dot M_{q-1}
 \end{pmatrix}.
\label{eq:v3v24-A-matrix}
\end{equation}
Let \(\boldsymbol\Sigma_q(r)\) be the logarithmic derivative of the
decaying matrix Poisson solution:
\begin{equation}
 -\partial_rU_q(r)=\boldsymbol\Sigma_q(r)U_q(r).
\label{eq:v3v24-matrix-DtN}
\end{equation}

\begin{proposition}[Full matrix Riccati equation]
\label{prop:v3v24-matrix-riccati}
The complete Dirichlet-to-Neumann matrix satisfies
\begin{equation}
 \partial_r\boldsymbol\Sigma_q
 =\boldsymbol\Sigma_q^2
 -\mathcal M_q\boldsymbol\Sigma_q-\mathcal A_q.
\label{eq:v3v24-matrix-riccati}
\end{equation}
Its principal symbol is
\begin{equation}
 \sigma_1(\boldsymbol\Sigma_q)
 =|\xi|_{h(r)}
 \begin{pmatrix}I_{\Lambda^q}&0\\0&I_{\Lambda^{q-1}}\end{pmatrix}.
\label{eq:v3v24-matrix-principal}
\end{equation}
If \(a_1(\mathcal A_q)\) denotes the degree-one covariant symbol of
\(\mathcal A_q\), its degree-zero symbol is
\begin{equation}
 \sigma_0(\boldsymbol\Sigma_q)
 =\frac12\mathcal M_q
 +\frac{\partial_r|\xi|_h}{2|\xi|_h}I
 +\frac{a_1(\mathcal A_q)}{2|\xi|_h}.
\label{eq:v3v24-matrix-sub}
\end{equation}
\end{proposition}

\begin{proof}
Differentiate \eqref{eq:v3v24-matrix-DtN} and substitute in
\eqref{eq:v3v24-normal-form}, exactly as in
Theorem~\ref{thm:v3v22-riccati}; matrix multiplication preserves the
displayed order and gives \eqref{eq:v3v24-matrix-riccati}.
The degree-two symbol equation is
\(s_1^2=|\xi|^2I\), and decay selects
\(s_1=|\xi|I\).  At degree one, the covariant Weyl product has no
\(\{|\xi|,|\xi|\}\) term and gives
\[
 \partial_r|\xi|\,I
 =2|\xi|s_0-\mathcal M_q|\xi|-a_1(\mathcal A_q).
\]
Solving proves \eqref{eq:v3v24-matrix-sub}.
\end{proof}

\begin{assessment}[The scalar V22 symbol is a diagonal projection]
\label{ass:v3v24-scalar-projection}
The first two terms in \eqref{eq:v3v24-matrix-sub} reduce on a scalar
tangential component to the V22 expression
\(H/2-\mathbb I(\xi^\sharp,\xi^\sharp)/(2|\xi|^2)\).
The last term contains the off-diagonal shape-differential symbols required
by the full complex.  Dropping it recreates the false degreewise defect.
\end{assessment}

\subsection{The boundary differential induced by decaying bulk fields}
\label{sec:v3v24-boundary-complex}

Let \(\mathcal P_q\) send a two-component boundary value
\(f=(\alpha_0,\beta_0)\) to the unique decaying solution of
\(L_qU=0\) with \(U(0)=f\).  Define
\begin{equation}
 \mathcal Q_q
 =\operatorname{Tr}\,D_q\mathcal P_q.
\label{eq:v3v24-Q-def}
\end{equation}
Because \(-\partial_r\mathcal P_qf|_{r=0}
=\boldsymbol\Sigma_qf\), this is the explicit boundary operator
\begin{equation}
 \mathcal Q_q
 \binom{\alpha}{\beta}
 =
 \binom{
 d_q\alpha
 }{
 -\pi_{\mathrm t}\boldsymbol\Sigma_q
       \binom{\alpha}{\beta}
 -d_{q-1}\beta
 },
\label{eq:v3v24-Q-explicit}
\end{equation}
where \(\pi_{\mathrm t}\) selects the tangential \(q\)-form component of
the matrix DtN output.

\begin{theorem}[Exact induced boundary complex]
\label{thm:v3v24-boundary-complex}
The bulk differential maps decaying Poisson solutions to decaying Poisson
solutions and
\begin{align}
 D_q\mathcal P_q
 &=\mathcal P_{q+1}\mathcal Q_q,
\label{eq:v3v24-P-chain}\\
 \mathcal Q_{q+1}\mathcal Q_q&=0.
\label{eq:v3v24-Q-square}
\end{align}
\end{theorem}

\begin{proof}
By Theorem~\ref{thm:v3v24-bulk-chain},
\[
 L_{q+1}D_q\mathcal P_qf
 =D_qL_q\mathcal P_qf=0.
\]
The field on the left decays and its boundary value is, by definition,
\(\mathcal Q_qf\).  Uniqueness of the decaying Dirichlet problem therefore
gives \eqref{eq:v3v24-P-chain}.  Apply
\(\operatorname{Tr}D_{q+1}\) and use
\(D_{q+1}D_q=0\):
\[
 \mathcal Q_{q+1}\mathcal Q_qf
 =\operatorname{Tr}D_{q+1}\mathcal P_{q+1}\mathcal Q_qf
 =\operatorname{Tr}D_{q+1}D_q\mathcal P_qf=0.
\]
\end{proof}

\begin{corollary}[Principal boundary symbol is already a complex]
\label{cor:v3v24-Q-principal}
Let \(\varepsilon_\xi=i\xi\wedge\).  The order-one principal symbol of
\(\mathcal Q_q\) is
\begin{equation}
 q_q^{(1)}(\xi)
 \binom{\alpha}{\beta}
 =
 \binom{
 \varepsilon_\xi\alpha
 }{
 -|\xi|\alpha-\varepsilon_\xi\beta
 }.
\label{eq:v3v24-Q-principal}
\end{equation}
It obeys
\begin{equation}
 q_{q+1}^{(1)}(\xi)q_q^{(1)}(\xi)=0
\qquad(\xi\ne0).
\label{eq:v3v24-Q-principal-square}
\end{equation}
\end{corollary}

\begin{proof}
Use \eqref{eq:v3v24-Q-explicit} and
\eqref{eq:v3v24-matrix-principal}.  Since
\(\varepsilon_\xi^2=0\),
\[
 -|\xi|\varepsilon_\xi\alpha
 -\varepsilon_\xi(-|\xi|\alpha-\varepsilon_\xi\beta)=0,
\]
while the upper component also squares to zero.
\end{proof}

\begin{firewall}[The boundary BRST operator depends on the Calderon data]
\label{fw:v3v24-Calderon}
On a moving collar, the correct induced boundary differential is
\(\mathcal Q=\operatorname{Tr}D\mathcal P\).  Its lower component contains
the full matrix Dirichlet-to-Neumann map.  The tangential exterior
differential by itself is not the boundary complex of decaying bulk
solutions.
\end{firewall}

\subsection{V24 conclusion and V25 audit mandate}
\label{sec:v3v24-conclusion}

V24 constructs the complete moving-collar de Rham prototype.  The shape
variation of the form inner product is
\(M_q=HI-2\mathcal S_q\).  It generates off-diagonal terms
\(M_qd-dM_{q-1}\) and \(\dot\delta_q\) in the Hodge Laplacian; these terms
repair the apparent scalar \([d,H]\) defect.  The full massive bulk
Laplacian intertwines the exterior differential exactly.  Its matrix
Dirichlet-to-Neumann map satisfies a matrix Riccati equation, and decaying
bulk solutions induce the exact nilpotent boundary differential
\(\mathcal Q=\operatorname{Tr}D\mathcal P\).

V25 is the next mandatory companion audit.  It will fingerprint the active
Yang--Mills and Navier--Stokes notes before selecting the next step.  The
present candidates are: proving ellipticity and Hodge decomposition for
the induced boundary complex \((\mathcal Q_q)\), or transferring the full
block construction to the linear Yang--Mills BRST differential.  The
choice must follow the strongest current companion lesson and must retain
the complete matrix Calderon data.

\section{V25: companion audit, ellipticity of the full Calderon boundary
complex, and massive contraction}
\label{sec:v3v25}

\subsection{Mandatory companion audit}
\label{sec:v3v25-audit}

The fifth-version audit was performed after V24 and before the next proof
route was selected.  The active companion files were
\begin{center}
\begin{tabular}{|p{0.30\textwidth}|p{0.12\textwidth}|p{0.43\textwidth}|}
\hline
programme and active file & bytes & SHA--256 \\
\hline
Navier--Stokes,
\texttt{Renewal\_NS\_Stability\_Research\_Notes\_V31.tex}
& 887537
& \texttt{F1121B62238AD5491BB7CF03635EA9934942EDF573ED8FAAA9EE79E1D38A6696}
\\
\hline
Yang--Mills,
\texttt{Renewal\_Yang\_Mills\_Mass\_Gap\_Research\_Notes\_V85.tex}
& 1357382
& \texttt{2D60E74FE8F6A16E26A7527F1F4F86DBD09FDAEEF3C00365836010CBC7D32F94}
\\
\hline
\end{tabular}
\end{center}
The Navier--Stokes snapshot is unchanged from the V20 audit.  The
Yang--Mills programme has advanced from V84 to V85.  V85 constructs
complete-bank output carriers without a multiplicity-copy sum and proves
that exact channel signs cancel the first-order component.  It also proves
that this local signed covariance is still too large for the proposed
crossing square tail.  The necessary next cancellation must occur before
the local joining coordinate is isolated: the binary prefixes, joining
partition signs, incidence-three coordinate, and suffix must remain in one
complete source packet.  Its overlap branch similarly reduces to one
identified repeated bank, but that bank must be controlled by one joint
two-occurrence carrier rather than two separately multiplied first-moment
bounds.

\begin{transferprinciple}[Complete boundary packet before a positive norm]
\label{tp:v3v25-complete-packet}
The SM analogue of the V85 obstruction is to retain the full
normal--tangential Calderon boundary differential
\(\mathcal Q=\operatorname{Tr}D\mathcal P\) before forming
\(\mathcal Q^*\mathcal Q\), singular values, or determinants.  Repeated
field components in different blocks must be paid for by the joint graph
norm of the complete operator, not by multiplying unrelated componentwise
bounds.
\end{transferprinciple}

No Yang--Mills estimate is imported.  The audit selects the proof order:
first establish the complete symbol complex, then take its positive Hodge
operator.

\subsection{The full principal symbol complex}
\label{sec:v3v25-symbol-complex}

Retain the notation and hypotheses of V24.  The degree-\(q\) boundary
bundle is
\begin{equation}
 \mathcal F^q
 =\Lambda^qT^*N\oplus\Lambda^{q-1}T^*N.
\label{eq:v3v25-Fq}
\end{equation}
The induced differential
\(\mathcal Q_q:C^\infty(\mathcal F^q)\to
C^\infty(\mathcal F^{q+1})\) is the order-one classical
pseudodifferential operator in
\eqref{eq:v3v24-Q-explicit}.  For \(\xi\ne0\), put
\[
 p=|\xi|_h,\qquad \varepsilon_\xi=i\xi\wedge.
\]
Its principal symbol is
\begin{equation}
 q_q(\xi)
 \binom{\alpha}{\beta}
 =
 \binom{\varepsilon_\xi\alpha}{
 -p\alpha-\varepsilon_\xi\beta}.
\label{eq:v3v25-q-symbol}
\end{equation}

\begin{theorem}[Explicit contraction of the principal symbol complex]
\label{thm:v3v25-symbol-contraction}
For every \(\xi\ne0\), define
\begin{equation}
 h_q(\xi):
 \mathcal F_x^q\longrightarrow\mathcal F_x^{q-1},
 \qquad
 h_q(\xi)\binom{\alpha}{\beta}
 =\binom{-p^{-1}\beta}{0}.
\label{eq:v3v25-h-symbol}
\end{equation}
Then
\begin{equation}
 q_{q-1}(\xi)h_q(\xi)
 +h_{q+1}(\xi)q_q(\xi)
 =I_{\mathcal F_x^q}.
\label{eq:v3v25-symbol-homotopy}
\end{equation}
Consequently the symbol sequence
\begin{equation}
 0\longrightarrow\mathcal F_x^0
 \xrightarrow{q_0(\xi)}\mathcal F_x^1
 \xrightarrow{q_1(\xi)}\cdots
 \xrightarrow{q_3(\xi)}\mathcal F_x^4
 \longrightarrow0
\label{eq:v3v25-symbol-sequence}
\end{equation}
is exact for every nonzero cotangent vector.
\end{theorem}

\begin{proof}
Apply the two compositions to \(f=(\alpha,\beta)^{\mathsf T}\).
First,
\[
 q_{q-1}h_qf
 =q_{q-1}\binom{-p^{-1}\beta}{0}
 =\binom{-p^{-1}\varepsilon_\xi\beta}{\beta}.
\]
Second,
\[
 h_{q+1}q_qf
 =h_{q+1}\binom{\varepsilon_\xi\alpha}{
 -p\alpha-\varepsilon_\xi\beta}
 =\binom{\alpha+p^{-1}\varepsilon_\xi\beta}{0}.
\]
Their sum is \(f\), proving \eqref{eq:v3v25-symbol-homotopy}.
Exactness follows because a chain with a contracting homotopy has zero
cohomology.
\end{proof}

\begin{corollary}[The induced boundary chain is an elliptic complex]
\label{cor:v3v25-elliptic-complex}
The sequence \((C^\infty(\mathcal F^\bullet),\mathcal Q)\) is a classical
elliptic pseudodifferential complex.
\end{corollary}

\begin{proof}
Nilpotence is Theorem~\ref{thm:v3v24-boundary-complex}; exactness of its
principal symbol away from the zero section is
Theorem~\ref{thm:v3v25-symbol-contraction}.  These are the two defining
properties.
\end{proof}

\begin{assessment}[Four-dimensional origin of the symbol]
\label{ass:v3v25-four-symbol}
Under the identification
\((\alpha,\beta)\leftrightarrow\alpha+dr\wedge\beta\),
\eqref{eq:v3v25-q-symbol} is exterior multiplication by the complex
characteristic covector
\[
 \zeta=i\xi-p\,dr.
\]
The simple homotopy \eqref{eq:v3v25-h-symbol} contracts with
\(-p^{-1}\partial_r\).  Thus the exactness belongs to the complete
four-dimensional characteristic packet; neither tangential wedge
multiplication nor the normal scalar factor is exact by itself.
\end{assessment}

\subsection{The positive boundary Hodge operator}
\label{sec:v3v25-boundary-hodge}

Equip \(\mathcal F^q\) with the direct-sum boundary \(L^2\) metric and let
\(\mathcal Q_q^*\) be the Hilbert adjoint.  Define
\begin{equation}
 \Box_{\mathcal Q,q}
 =\mathcal Q_{q-1}\mathcal Q_{q-1}^*
 +\mathcal Q_q^*\mathcal Q_q.
\label{eq:v3v25-box}
\end{equation}

\begin{theorem}[Strong ellipticity of the complete positive packet]
\label{thm:v3v25-box-elliptic}
The boundary Hodge operator is a positive self-adjoint elliptic
pseudodifferential operator of order two, with scalar principal symbol
\begin{equation}
 \sigma_2(\Box_{\mathcal Q,q})(x,\xi)
 =2|\xi|_h^2I_{\mathcal F_x^q}.
\label{eq:v3v25-box-symbol}
\end{equation}
In particular it has compact resolvent and finite-dimensional kernel.
\end{theorem}

\begin{proof}
With respect to the Hermitian form metric,
\(\varepsilon_\xi^*=-i\iota_{\xi^\sharp}\), and
\[
 \varepsilon_\xi^*\varepsilon_\xi
 +\varepsilon_\xi\varepsilon_\xi^*
 =|\xi|^2I.
\]
Writing \eqref{eq:v3v25-q-symbol} as the matrix
\[
 q_q(\xi)=
 \begin{pmatrix}
 \varepsilon_\xi&0\\
 -p&-\varepsilon_\xi
 \end{pmatrix},
\]
a direct multiplication gives
\[
 q_{q-1}q_{q-1}^*+q_q^*q_q
 =
 \begin{pmatrix}
 |\xi|^2+p^2&0\\0&|\xi|^2+p^2
 \end{pmatrix}
 =2p^2I.
\]
This is \eqref{eq:v3v25-box-symbol}.  Positivity follows from
\[
 \langle\Box_{\mathcal Q,q}f,f\rangle
 =\|\mathcal Q_{q-1}^*f\|^2+\|\mathcal Q_qf\|^2.
\]
Standard elliptic parametrix theory on compact \(N\) gives a self-adjoint
Fredholm realization with compact resolvent and smooth finite-dimensional
kernel.
\end{proof}

\begin{firewall}[The positive operator is formed after exact cancellation]
\label{fw:v3v25-positive-order}
The off-diagonal entries of
\(q_{q-1}q_{q-1}^*\) and \(q_q^*q_q\) cancel only in their sum.
Estimating either positive block first loses that cancellation and can
double-charge the shared normal component.  Equation
\eqref{eq:v3v25-box-symbol} is therefore a complete-packet result in the
sense selected by the V25 audit.
\end{firewall}

\subsection{A full massive contracting homotopy}
\label{sec:v3v25-massive-homotopy}

Ellipticity alone gives only finite-dimensional boundary cohomology.  The
common massive bulk Hodge construction makes the induced complex actually
contractible.

\begin{theorem}[Exact contraction of the full induced boundary complex]
\label{thm:v3v25-full-contraction}
Let \(\mathcal P_q\) be the decaying Poisson operators of V24 for
\[
 L_q=\Delta_{X,q}+\mu^2I,\qquad\mu>0,
\]
and suppose the adjoint bulk differential preserves the selected decaying
class.  Define
\begin{equation}
 \mathcal H_q
 =-\mu^{-2}\operatorname{Tr}
 D_{q-1}^*\mathcal P_q:
 C^\infty(\mathcal F^q)\longrightarrow
 C^\infty(\mathcal F^{q-1}).
\label{eq:v3v25-Hq}
\end{equation}
Then
\begin{equation}
 \mathcal Q_{q-1}\mathcal H_q
 +\mathcal H_{q+1}\mathcal Q_q
 =I.
\label{eq:v3v25-full-homotopy}
\end{equation}
Hence the smooth cohomology of the full induced boundary complex vanishes
in every degree.
\end{theorem}

\begin{proof}
The Hodge identities imply
\[
 D_{q-1}^*L_q=L_{q-1}D_{q-1}^*.
\]
Thus \(D_{q-1}^*\mathcal P_qf\) is a decaying solution of
\(L_{q-1}=0\), and uniqueness of the Dirichlet problem gives
\begin{equation}
 \mathcal P_{q-1}\mathcal H_qf
 =-\mu^{-2}D_{q-1}^*\mathcal P_qf.
\label{eq:v3v25-PH}
\end{equation}
Similarly, Theorem~\ref{thm:v3v24-boundary-complex} gives
\(\mathcal P_{q+1}\mathcal Q_q=D_q\mathcal P_q\).  Therefore
\[
 \begin{split}
 \mathcal Q_{q-1}\mathcal H_qf
 +\mathcal H_{q+1}\mathcal Q_qf
 &=-\mu^{-2}\operatorname{Tr}
 \left(D_{q-1}D_{q-1}^*+D_q^*D_q\right)\mathcal P_qf\\
 &=-\mu^{-2}\operatorname{Tr}\Delta_{X,q}\mathcal P_qf.
 \end{split}
\]
Since \(L_q\mathcal P_qf=0\),
\(\Delta_{X,q}\mathcal P_qf=-\mu^2\mathcal P_qf\).  Its boundary trace is
\(f\), proving \eqref{eq:v3v25-full-homotopy}.
\end{proof}

\begin{assessment}[Why V23 retained cohomology while V25 does not]
\label{ass:v3v25-v23}
V23 used only the tangential Hilbert complex
\((\Omega^\bullet(N),d_N)\), whose harmonic subspace survives and gives
\((1+\rho/\mu)^\chi\).  V25 uses the larger
\(\Lambda^q\oplus\Lambda^{q-1}\) boundary complex induced from the full
bulk differential.  Its normal component supplies the contracting
homotopy.  The two determinants belong to different complexes and their
cohomological conclusions must not be identified.
\end{assessment}

\subsection{Hodge decomposition and a fixed-geometry boundary gap}
\label{sec:v3v25-Hodge}

\begin{corollary}[Invertible boundary Hodge operator]
\label{cor:v3v25-box-invertible}
Under the hypotheses of
Theorem~\ref{thm:v3v25-full-contraction},
\begin{equation}
 \ker\Box_{\mathcal Q,q}=\{0\}.
\label{eq:v3v25-kernel-zero}
\end{equation}
Consequently there is a constant \(c_q>0\), depending on the completed
geometry, \(\mu\), and the boundary realization, such that
\begin{equation}
 \langle\Box_{\mathcal Q,q}f,f\rangle
 \ge c_q\|f\|_{L^2}^2
\label{eq:v3v25-boundary-gap}
\end{equation}
on the operator domain.
\end{corollary}

\begin{proof}
If \(f\in\ker\Box_{\mathcal Q,q}\), positivity in the proof of
Theorem~\ref{thm:v3v25-box-elliptic} gives
\(\mathcal Q_qf=0\) and \(\mathcal Q_{q-1}^*f=0\).
The homotopy identity gives \(f=\mathcal Q_{q-1}\mathcal H_qf\).
Hence
\[
 \|f\|^2
 =\langle\mathcal H_qf,\mathcal Q_{q-1}^*f\rangle=0.
\]
The positive elliptic operator has discrete spectrum; absence of zero in
that spectrum gives a positive first eigenvalue \(c_q\).
\end{proof}

\begin{corollary}[Exact boundary Hodge decomposition]
\label{cor:v3v25-Hodge}
Let \(G_q=\Box_{\mathcal Q,q}^{-1}\).  Then
\begin{equation}
 I
 =\mathcal Q_{q-1}\mathcal Q_{q-1}^*G_q
 +\mathcal Q_q^*\mathcal Q_qG_q,
\label{eq:v3v25-Hodge-identity}
\end{equation}
and
\begin{equation}
 L^2(\mathcal F^q)
 =\operatorname{ran}\mathcal Q_{q-1}
 \mathbin{\widehat\oplus}
 \operatorname{ran}\mathcal Q_q^*.
\label{eq:v3v25-Hodge-decomposition}
\end{equation}
Both ranges are closed and the sum is orthogonal.
\end{corollary}

\begin{proof}
Equation~\eqref{eq:v3v25-Hodge-identity} is
\(\Box_{\mathcal Q,q}G_q=I\).  Nilpotence makes the two displayed ranges
orthogonal.  Elliptic regularity and boundedness of the corresponding
order-zero Hodge projections give closedness and the completed direct sum.
\end{proof}

\begin{firewall}[The gap is not yet uniform]
\label{fw:v3v25-nonuniform-gap}
The constant \(c_q\) in \eqref{eq:v3v25-boundary-gap} is positive for one
fixed massive completed reservoir.  V25 does not bound it uniformly under
collar degeneration, variation of the remote end, \(\mu\downarrow0\), or
nonabelian background fields.  Such a uniform bound requires quantitative
control of the Poisson and homotopy operators.
\end{firewall}

\begin{firewall}[The massless limit reopens cohomology]
\label{fw:v3v25-massless}
The exact homotopy \(\mathcal H_q\) contains \(\mu^{-2}\).
It gives no estimate at \(\mu=0\), where topological and gauge zero modes
may return.  Principal-symbol ellipticity survives at nonzero \(\xi\), but
it does not determine the global massless cohomology or its measure.
\end{firewall}

\subsection{V25 conclusion and V26 target}
\label{sec:v3v25-conclusion}

The mandatory audit has led to a complete-packet boundary theorem.  The
principal symbol of the full Calderon-induced complex has the explicit
contraction \(h(\alpha,\beta)=(-|\xi|^{-1}\beta,0)\).  Its positive Hodge
operator must be formed from both neighboring degrees; only then do its
off-diagonal symbols cancel, leaving the strongly elliptic scalar symbol
\(2|\xi|^2I\).  For a common positive bulk mass, the full boundary complex
has the exact operator homotopy
\(-\mu^{-2}\operatorname{Tr}D^*\mathcal P\), so its cohomology vanishes and
its Hodge operator has a positive fixed-geometry spectral gap.

V26 should transfer this construction to the linear Abelian
Yang--Mills BRST prototype, where the differential contains both the de
Rham gauge map and the algebraic ghost/antighost multiplier pair.  The
proof must form the full statistics-weighted complex before taking
determinants, identify its massive contractible quartets, and leave the
massless physical transverse and harmonic sectors explicit.  This is the
smallest setting in which the Calderon boundary complex, V18 determinant
bookkeeping, and V17 physical-projector firewall meet.

\section{V26: Abelian BRST spectral bookkeeping and the surviving physical
reservoir determinant}
\label{sec:v3v26}

\subsection{Purpose and claim boundary}
\label{sec:v3v26-purpose}

V25 proves that the full massive Calderon boundary de Rham complex is
elliptic and contractible.  That statement concerns the chain.  A quantum
gauge determinant also contains statistics and Gaussian powers.  The
present installment computes those powers in the linear Abelian prototype
and shows exactly what cancels.  The result is deliberately more restrictive
than saying \emph{BRST quartets cancel}: the exact one-form sector is
isospectral to the nonconstant scalar ghost sector, but their conventional
one-loop weights leave a square-root scalar factor unless an independently
derived balanced auxiliary determinant pays it.

Let \(N\) be a closed connected Riemannian three-manifold.  Write
\[
 \Delta_0=\delta d\quad\hbox{on functions},\qquad
 \Delta_1=d\delta+\delta d\quad\hbox{on one-forms}.
\]
All determinants in V26 are finite spectral-cutoff determinants until the
modified determinant is explicitly introduced.

\subsection{Linear Abelian BRST differential}
\label{sec:v3v26-BRST}

The minimal and nonminimal Abelian fields are
\[
 A\in\Omega^1,\qquad
 c\in\Omega^0[1],\qquad
 \bar c\in\Omega^0[-1],\qquad
 b\in\Omega^0,
\]
where the brackets record ghost number.  The linear BRST differential is
\begin{equation}
 sA=dc,\qquad sc=0,\qquad s\bar c=b,\qquad sb=0.
\label{eq:v3v26-BRST}
\end{equation}

\begin{proposition}[Nilpotence and the two contractible directions]
\label{prop:v3v26-nilpotence}
The differential in \eqref{eq:v3v26-BRST} satisfies \(s^2=0\).
The pair \((\bar c,b)\) is algebraically contractible.  The pair
\((c,d c)\) is contractible on the orthogonal complement of
\(\ker\Delta_0\), but constant ghosts are stabilizer zero modes and do not
have an exact one-form partner.
\end{proposition}

\begin{proof}
Nilpotence follows from \(d^2=0\) and the last three definitions in
\eqref{eq:v3v26-BRST}.  The map \(\bar c\mapsto b\mapsto0\) has the
algebraic homotopy \(b\mapsto\bar c\), \(\bar c\mapsto0\).  If
\(\Delta_0u=\lambda u\) with \(\lambda>0\), then \(du\ne0\) and
\(\lambda^{-1}\delta(du)=u\), giving a contraction of
\(u\mapsto du\).  If \(u\) is constant, \(du=0\), so this argument is
unavailable.
\end{proof}

\begin{assessment}[Moving-collar realization]
\label{ass:v3v26-moving-realization}
On the moving collar, the gauge part of
\eqref{eq:v3v26-BRST} is represented by the full boundary differential
\(\mathcal Q_0=\operatorname{Tr}D_0\mathcal P_0\) from V24, not by
tangential \(d_N\) alone.  Theorem~\ref{thm:v3v24-boundary-complex} proves
its nilpotence before a physical projection is taken.  V26 uses the
tangential closed-manifold spectrum only to make the Gaussian
statistics calculation transparent.
\end{assessment}

\subsection{Exact--coexact Hodge splitting}
\label{sec:v3v26-Hodge}

Let \(\mathcal H^q=\ker\Delta_q\), let
\[
 \Omega^0_\perp=(\mathcal H^0)^\perp,
\]
and take closures in \(L^2\).  Hodge theory gives
\begin{equation}
 L^2\Omega^1
 =\overline{d\Omega^0_\perp}
 \mathbin{\widehat\oplus}
 \overline{\delta\Omega^2}
 \mathbin{\widehat\oplus}
 \mathcal H^1.
\label{eq:v3v26-Hodge}
\end{equation}
Call the first two summands
\(\mathcal E^1\) and \(\mathcal C^1\), respectively.

\begin{theorem}[Unitary scalar--exact spectral equivalence]
\label{thm:v3v26-unitary}
On the positive spectral subspace of \(\Delta_0\), the operator
\begin{equation}
 U=d\Delta_0^{-1/2}:
 \Omega^0_\perp\longrightarrow\mathcal E^1
\label{eq:v3v26-U}
\end{equation}
is unitary and
\begin{equation}
 U\Delta_0=\Delta_1U.
\label{eq:v3v26-U-intertwine}
\end{equation}
Thus, with multiplicity, the positive spectrum of \(\Delta_0\) equals
the spectrum of \(\Delta_1\) restricted to exact one-forms.
\end{theorem}

\begin{proof}
For \(f\in\Omega^0_\perp\),
\[
 \|Uf\|^2
 =\langle d\Delta_0^{-1/2}f,
          d\Delta_0^{-1/2}f\rangle
 =\langle f,f\rangle.
\]
Its range is dense in and closed onto \(\mathcal E^1\), so it is unitary.
The Hodge identity
\(\Delta_1d=d\Delta_0\) gives
\eqref{eq:v3v26-U-intertwine} on smooth eigenvectors and hence on the
spectral domain.
\end{proof}

\begin{corollary}[Functional calculus matches only the exact sector]
\label{cor:v3v26-functional}
For every bounded Borel function \(F\),
\begin{equation}
 F(\Delta_1)|_{\mathcal E^1}
 =UF(\Delta_0)|_{\Omega^0_\perp}U^*.
\label{eq:v3v26-functional}
\end{equation}
No such identity identifies the coexact or harmonic one-form sectors with
scalar ghosts.
\end{corollary}

\subsection{The reservoir factors and a common spectral cutoff}
\label{sec:v3v26-reservoir}

Fix \(\mu,\rho>0\), and define
\begin{equation}
 R_q
 =I+K_q,\qquad
 K_q=\rho(\Delta_q+\mu^2)^{-1/2}.
\label{eq:v3v26-Rq}
\end{equation}
Let
\[
 P_{0,\Lambda}
 =\mathbf1_{(0,\Lambda^2]}(\Delta_0)
\]
and let \(P_{1,\Lambda}^{\rm ex}\) and
\(P_{1,\Lambda}^{\rm coex}\) be the corresponding positive spectral
cutoffs on \(\mathcal E^1\) and \(\mathcal C^1\).

\begin{proposition}[Exact reservoir matching in the gauge direction]
\label{prop:v3v26-reservoir-match}
At every cutoff,
\begin{equation}
 \det\!\left(
 P_{1,\Lambda}^{\rm ex}R_1P_{1,\Lambda}^{\rm ex}
 \right)
 =
 \det\!\left(
 P_{0,\Lambda}R_0P_{0,\Lambda}
 \right).
\label{eq:v3v26-det-match}
\end{equation}
The equality holds mode by mode, with common eigenvalue
\begin{equation}
 r_\lambda
 =1+\frac{\rho}{\sqrt{\lambda+\mu^2}},
 \qquad 0<\lambda\le\Lambda^2.
\label{eq:v3v26-r-lambda}
\end{equation}
\end{proposition}

\begin{proof}
Apply Corollary~\ref{cor:v3v26-functional} to
\(F(x)=1+\rho(x+\mu^2)^{-1/2}\).  The unitary \(U\) also intertwines the
common spectral cutoffs, so the two finite matrices are unitarily
equivalent.
\end{proof}

\begin{firewall}[A cutoff by field label is not a common cutoff]
\label{fw:v3v26-common-cutoff}
The equality \eqref{eq:v3v26-det-match} requires paired spectral
subspaces.  Independent mode counts, different masses, different reservoir
couplings, or different boundary domains destroy the unitary equivalence
even if the two factors have the same nominal differential order.
\end{firewall}

\subsection{The conventional Abelian one-loop power}
\label{sec:v3v26-one-loop}

For a real gauge one-form and one complex Faddeev--Popov ghost pair, the
relative Gaussian reservoir factor at finite cutoff has the conventional
power
\begin{equation}
 Z_{\rm FP,\Lambda}^{\rm rel}
 =
 \det{}'(R_1)_{\Lambda}^{-1/2}
 \det{}'(R_0)_{\Lambda}.
\label{eq:v3v26-ZFP}
\end{equation}
The prime excludes \(\mathcal H^0\) and \(\mathcal H^1\), whose measures
are treated separately below.

\begin{theorem}[Exact factor left after gauge--ghost matching]
\label{thm:v3v26-FP-factor}
The finite-cutoff factor \eqref{eq:v3v26-ZFP} is
\begin{equation}
 Z_{\rm FP,\Lambda}^{\rm rel}
 =
 \det(R_1|_{\mathcal C^1})_{\Lambda}^{-1/2}
 \det(R_0|_{\Omega^0_\perp})_{\Lambda}^{1/2}.
\label{eq:v3v26-FP-reduced}
\end{equation}
Thus the ghost cancels the exact gauge determinant only to one half of its
power.  The surviving scalar square root is part of the conventional
Faddeev--Popov weight and cannot be deleted by the isospectrality theorem.
\end{theorem}

\begin{proof}
The Hodge decomposition is invariant under \(R_1\), so
\[
 \det{}'(R_1)_\Lambda
 =\det(R_1|_{\mathcal E^1})_\Lambda
  \det(R_1|_{\mathcal C^1})_\Lambda.
\]
Use \eqref{eq:v3v26-det-match} in
\eqref{eq:v3v26-ZFP}:
\[
 \det(R_1|_{\mathcal E^1})_\Lambda^{-1/2}
 \det(R_0)_\Lambda
 =\det(R_0)_\Lambda^{1/2}.
\]
The coexact factor remains unchanged, proving
\eqref{eq:v3v26-FP-reduced}.
\end{proof}

\begin{counterexample}[Isospectral does not mean fully cancelled]
\label{cex:v3v26-half-power}
For one paired eigenvalue, the exact real gauge mode contributes
\(r_\lambda^{-1/2}\), while the complex ghost pair contributes
\(r_\lambda\).  Their product is \(r_\lambda^{1/2}\), not one.  A claim of
unit quartet cancellation that lists only these two fields is false even
at a one-dimensional cutoff.
\end{counterexample}

\begin{proposition}[Conditional balanced quartet]
\label{prop:v3v26-balanced}
Suppose an independently derived real commuting auxiliary mode has the
same factor \(r_\lambda\) and Gaussian power \(-1/2\).  Then the three
nonphysical factors obey
\begin{equation}
 r_\lambda^{-1/2}\,r_\lambda\,r_\lambda^{-1/2}=1
\label{eq:v3v26-balanced-mode}
\end{equation}
at every positive scalar eigenvalue.  Hence a chainwise common cutoff gives
exact nonzero-mode cancellation as in V18.
\end{proposition}

\begin{proof}
This is the displayed scalar identity applied on every paired eigenspace.
\end{proof}

\begin{firewall}[The auxiliary payer must come from the action and measure]
\label{fw:v3v26-auxiliary}
Proposition~\ref{prop:v3v26-balanced} is a determinant gate, not a license
to insert a field.  The Nakanishi--Lautrup field is often ultralocal before
elimination, and contour choices can change its phase.  It pays the last
half-power only if the gauge-fixing action, boundary domain, reservoir
coupling, cutoff, and integration contour derive precisely that common
factor.
\end{firewall}

\subsection{Fourth modified determinant in the surviving sectors}
\label{sec:v3v26-det4}

Both \(K_0\) and \(K_1\) are order-minus-one operators on a compact
three-manifold and belong to \(\mathfrak S_p\) exactly for \(p>3\).
Their fourth modified determinants are therefore defined.  With primes
and restrictions as above, set
\begin{equation}
 \log Z_{\rm FP,4}^{\rm rel}
 =-\frac12\log\det{}_4(R_1)'
 +\log\det{}_4(R_0)'.
\label{eq:v3v26-Z4}
\end{equation}

\begin{corollary}[Renormalized exact-sector reduction]
\label{cor:v3v26-Z4-reduced}
The nonlocal modified-determinant remainder satisfies
\begin{equation}
 \log Z_{\rm FP,4}^{\rm rel}
 =-\frac12\log\det{}_4(R_1|_{\mathcal C^1})
 +\frac12\log\det{}_4(R_0|_{\Omega^0_\perp}).
\label{eq:v3v26-Z4-reduced}
\end{equation}
The first three local counterterm traces obey the same exact-sector
half-power reduction when one common smooth regulator is used.
\end{corollary}

\begin{proof}
The logarithm defining \(\det_4\) is a trace-class functional calculus of
\(\Delta_q\).  Apply the unitary equivalence
\eqref{eq:v3v26-functional} to its exact one-form restriction.  The same
argument applies separately to each of the first three regulated powers
\(K,K^2,K^3\).
\end{proof}

\begin{proposition}[A nonzero ultraviolet polarization currency remains]
\label{prop:v3v26-uv-currency}
At principal-symbol level, the coexact one-form projector has rank two and
the nonconstant scalar sector has rank one.  Therefore the weighted
principal trace in \eqref{eq:v3v26-Z4-reduced} is
\begin{equation}
 -\frac12(2)\frac{\rho}{|\xi|}
 +\frac12(1)\frac{\rho}{|\xi|}
 =-\frac{\rho}{2|\xi|}.
\label{eq:v3v26-uv-currency}
\end{equation}
The Abelian gauge--ghost system does not cancel all three local reservoir
counterterms.
\end{proposition}

\begin{proof}
For \(\xi\ne0\), exact one-forms span the line generated by \(\xi\);
its orthogonal complement in \(T_x^*N\) has dimension two.  The principal
symbol of each normalized reservoir correction is
\(\rho|\xi|^{-1}\) times the identity on its polarization space.  Insert
the powers from \eqref{eq:v3v26-Z4-reduced}.
\end{proof}

\begin{assessment}[Physical interpretation is postponed]
\label{ass:v3v26-uv-interpretation}
Equation~\eqref{eq:v3v26-uv-currency} is a boundary spectral count for this
specified tangential model.  It is not yet the four-dimensional photon
degree-of-freedom count.  Normal components, the full Calderon complex,
auxiliary normalization, and the Einstein boundary conditions must be
assembled before assigning a physical vacuum-energy coefficient.
\end{assessment}

\subsection{Zero modes and the nonlocal physical projector}
\label{sec:v3v26-zero-projector}

\begin{proposition}[The two Abelian zero-mode measures]
\label{prop:v3v26-zero-modes}
For connected \(N\),
\(\mathcal H^0\) consists of constants and is the Lie algebra of the global
\(U(1)\) stabilizer.  The space \(\mathcal H^1\) parametrizes infinitesimal
flat Abelian connections.  Neither space belongs to the primed determinant:
the first is treated by the stabilizer volume, and the second by the
flat-connection moduli measure, including its integral lattice for compact
\(U(1)\).
\end{proposition}

\begin{proof}
The first statement is the maximum principle for harmonic functions on a
connected closed manifold.  A harmonic one-form is closed and coclosed and
therefore gives an infinitesimal flat connection orthogonal to gauge-exact
directions.  Large \(U(1)\) gauge transformations shift its cohomology
class by the integral lattice.  Gaussian inversion is absent in both
zero-eigenvalue sectors, so a determinant cannot supply either quotient
measure.
\end{proof}

\begin{firewall}[The reservoir mass does not erase gauge topology]
\label{fw:v3v26-reservoir-mass}
The parameter \(\mu>0\) makes
\((\Delta_q+\mu^2)^{-1/2}\) bounded, but it does not turn a constant gauge
transformation into a nontrivial exact one-form or turn a flat-connection
modulus into a Gaussian fluctuation.  Zero-mode classification is made
with \(\Delta_q\) before the reservoir function is applied.
\end{firewall}

Let \(G_0=(\Delta_0|_{\Omega^0_\perp})^{-1}\).  The exact one-form
projector is
\begin{equation}
 P_{\rm ex}=dG_0\delta,
\qquad
 P_{\rm phys}=I-P_{\rm ex}-P_{\mathcal H^1}.
\label{eq:v3v26-projectors}
\end{equation}

\begin{proposition}[Physical projection is elliptic and nonlocal]
\label{prop:v3v26-projector-nonlocal}
The principal symbol of \(P_{\rm ex}\) is
\begin{equation}
 \sigma_0(P_{\rm ex})(x,\xi)
 =\frac{\xi\otimes\xi}{|\xi|^2}.
\label{eq:v3v26-Pex-symbol}
\end{equation}
The exact operator contains the Green operator \(G_0\) and is nonlocal.
Consequently a local boundary action cannot be obtained by simply replacing
the full gauge field with \(P_{\rm phys}A\).
\end{proposition}

\begin{proof}
The leading symbols of \(d\), \(G_0\), and \(\delta\) are respectively
\(i\xi\wedge\), \(|\xi|^{-2}\), and
\(-i\iota_{\xi^\sharp}\), whose product on one-forms is
\(\xi\otimes\xi/|\xi|^2\).  The Schwartz kernel of the inverse elliptic
operator \(G_0\) is not supported on the diagonal, proving nonlocality.
\end{proof}

\begin{assessment}[Compatibility with V17 and V25]
\label{ass:v3v26-V17-V25}
Proposition~\ref{prop:v3v26-projector-nonlocal} reproduces the V17
projector firewall in the explicit Abelian spectrum.  The local route is
the full BRST/Calderon complex of V24--V25; the physical projector is used
after solving or in the final spectral decomposition, not inserted as a
local field-level replacement.
\end{assessment}

\subsection{V26 conclusion and V27 target}
\label{sec:v3v26-conclusion}

V26 supplies exact Abelian statistics bookkeeping.  Positive scalar
eigenmodes are unitarily equivalent to exact one-form modes, and their
reservoir factors match at every common cutoff.  With the conventional
real-vector and complex-ghost powers, this matching leaves a scalar
square-root factor and the coexact one-form determinant.  A fully balanced
unit cancellation requires an additional real commuting half-power derived
from the actual gauge-fixing action and measure.  The surviving sectors
have well-defined fourth modified determinants but retain local
counterterms.  Constant ghosts and harmonic one-forms remain quotient and
moduli measures, and the physical projector is necessarily nonlocal.

V27 should derive, rather than assume, the nonminimal auxiliary factor from
a BRST-exact Euclidean gauge-fixing action with a specified contour and
boundary domain.  It must diagonalize the longitudinal, ghost,
antighost, and multiplier block at finite cutoff, calculate its determinant
including phases, and compare it with the conditional balanced identity
\eqref{eq:v3v26-balanced-mode}.  If the common reservoir factor cannot be
generated by a local gauge-fixing fermion, the programme must retain the
half-power in \eqref{eq:v3v26-FP-reduced} instead of declaring a quartet.

\section{V27: BRST-exact nonminimal gauge fixing, contour-fixed Gaussian
blocks, and exact reservoir cancellation}
\label{sec:v3v27}

\subsection{The missing determinant in V26}
\label{sec:v3v27-purpose}

V26 kept the Nakanishi--Lautrup payer conditional because its ordinary
form is ultralocal.  The correct question is not whether the multiplier has
the same post hoc kinetic operator as the ghost.  It is whether one common
spectral factor occurs in the entire BRST-exact gauge-fixing fermion.  If it
does, the longitudinal gauge amplitude and multiplier form a coupled
two-by-two Gaussian block whose determinant contains two powers of that
factor.  The complex ghost supplies the inverse statistics power, and the
relative reservoir contribution cancels exactly.

V27 proves this at every finite spectral cutoff and fixes the Euclidean
contour by a convergent regularization.  It also proves that the desired
inverse-square-root factor cannot be produced by a finite-order local
boundary gauge-fixing operator.  A local realization therefore requires an
attached reservoir complex, rather than a nonlocal boundary instruction.

\subsection{A spectrally weighted gauge-fixing fermion}
\label{sec:v3v27-gauge-fermion}

Use the Abelian fields of V26, but choose the Euclidean convention
\begin{equation}
 sA=dc,\qquad sc=0,\qquad
 s\bar c=i b,\qquad sb=0,
\label{eq:v3v27-BRST}
\end{equation}
where \(A,b\) are integrated as real bosonic fields and
\((c,\bar c)\) are independent complex Grassmann fields.  Let
\(F=F(\Delta_0)>0\) be a bounded positive self-adjoint spectral multiplier
which commutes with \(\Delta_0\).  For gauge parameter \(\xi>0\), define
\begin{equation}
 \Psi_F
 =\langle\bar c,F\delta A\rangle
 -\frac{i\xi}{2}\langle\bar c,Fb\rangle.
\label{eq:v3v27-Psi}
\end{equation}

\begin{proposition}[The complete BRST-exact nonminimal action]
\label{prop:v3v27-sPsi}
The gauge-fixing and ghost action is
\begin{equation}
 S_{\rm nm,F}=s\Psi_F
 =i\langle b,F\delta A\rangle
 +\frac{\xi}{2}\langle b,Fb\rangle
 -\langle\bar c,F\Delta_0c\rangle.
\label{eq:v3v27-Snm}
\end{equation}
It is BRST closed at a common finite spectral cutoff invariant under
\(\Delta_0\).
\end{proposition}

\begin{proof}
The graded Leibniz rule gives
\[
 s\langle\bar c,F\delta A\rangle
 =i\langle b,F\delta A\rangle
 -\langle\bar c,F\delta dc\rangle.
\]
Since \(\delta d=\Delta_0\) on functions, this is the first and third
term of \eqref{eq:v3v27-Snm}.  The second part of
\eqref{eq:v3v27-Psi} contributes
\[
 -\frac{i\xi}{2}\langle ib,Fb\rangle
 =\frac{\xi}{2}\langle b,Fb\rangle.
\]
Nilpotence gives \(sS_{\rm nm,F}=s^2\Psi_F=0\).  A spectral cutoff commutes
with \(F,\Delta_0\), and the scalar--exact unitary of V26, so it preserves
the finite-dimensional BRST block.
\end{proof}

\begin{firewall}[One factor must multiply the whole gauge fermion]
\label{fw:v3v27-one-F}
Placing \(F\) only in the ghost term, only in the multiplier quadratic
term, or only in a post-elimination longitudinal Hessian is not
\eqref{eq:v3v27-Snm}.  BRST exactness and the determinant cancellation
below use the same \(F\) in all three terms.
\end{firewall}

\subsection{Exact positive-eigenvalue block}
\label{sec:v3v27-mode}

Let \(u\) be a normalized scalar eigenfunction with
\(\Delta_0u=\lambda u\), \(\lambda>0\), and put
\begin{equation}
 e_\lambda=\lambda^{-1/2}du.
\label{eq:v3v27-e}
\end{equation}
Then \(e_\lambda\) is a normalized exact one-form and
\(\delta e_\lambda=\sqrt\lambda\,u\).  Write
\[
 A_{\rm ex}=a e_\lambda,\qquad
 b=b_\lambda u,\qquad
 c=c_\lambda u,\qquad
 \bar c=\bar c_\lambda u,
\]
and let \(f_\lambda=F(\lambda)>0\).  The mode action is
\begin{equation}
 S_{\lambda,F}
 =\frac{\xi f_\lambda}{2}b_\lambda^2
 +if_\lambda\sqrt\lambda\,b_\lambda a
 -f_\lambda\lambda\,\bar c_\lambda c_\lambda.
\label{eq:v3v27-mode-action}
\end{equation}

\begin{theorem}[Contour-fixed nonminimal mode integral]
\label{thm:v3v27-mode-integral}
Define the bosonic integral by adding
\(\varepsilon a^2/2\), \(\varepsilon>0\), integrating
\((a,b_\lambda)\in\mathbb R^2\), and taking
\(\varepsilon\downarrow0\) after the Gaussian is evaluated.  Then
\begin{align}
 \lim_{\varepsilon\downarrow0}
 \int_{\mathbb R^2}
 e^{-S_{\lambda,F}-\varepsilon a^2/2}\,da\,db_\lambda
 &=\frac{2\pi}{f_\lambda\sqrt\lambda},
\label{eq:v3v27-boson-integral}\\
 \int d\bar c_\lambda\,dc_\lambda\,
 e^{\,f_\lambda\lambda\bar c_\lambda c_\lambda}
 &=f_\lambda\lambda,
\label{eq:v3v27-ghost-integral}
\end{align}
up to the fixed orientation sign of the Grassmann measure.  With that
orientation fixed once, the complete mode integral is
\begin{equation}
 Z_{\lambda,F}=2\pi\sqrt\lambda,
\label{eq:v3v27-mode-cancel}
\end{equation}
independent of \(F\) and \(\xi\).
\end{theorem}

\begin{proof}
For \(\varepsilon>0\), integrate \(b_\lambda\) first:
\[
 \int_{\mathbb R}
 \exp\left(
 -\frac{\xi f_\lambda}{2}b^2
 -if_\lambda\sqrt\lambda\,ba
 \right)db
 =
 \sqrt{\frac{2\pi}{\xi f_\lambda}}\,
 \exp\left(-\frac{f_\lambda\lambda}{2\xi}a^2\right).
\]
The remaining \(a\)-integral, including its regulator, is
\[
 \sqrt{\frac{2\pi}{
 \varepsilon+f_\lambda\lambda/\xi}}.
\]
Their product tends to
\(2\pi/(f_\lambda\sqrt\lambda)\), proving
\eqref{eq:v3v27-boson-integral}.  Berezin integration extracts the
coefficient \(f_\lambda\lambda\), proving
\eqref{eq:v3v27-ghost-integral}.  Multiplication gives
\eqref{eq:v3v27-mode-cancel}; the factors of \(\xi\) have already
cancelled inside the bosonic block.
\end{proof}

\begin{assessment}[Equivalent steepest-descent contour]
\label{ass:v3v27-contour}
Completing the square gives
\[
 \frac{\xi f_\lambda}{2}
 \left(b_\lambda+\frac{i\sqrt\lambda}{\xi}a\right)^2
 +\frac{f_\lambda\lambda}{2\xi}a^2.
\]
The \(\varepsilon\)-prescription is equivalent to translating the
\(b\)-contour through this entire Gaussian and then using the positive
longitudinal action.  It fixes the square-root branch in
\eqref{eq:v3v27-boson-integral}; no unrecorded phase is left.
\end{assessment}

\begin{corollary}[Derived balanced reservoir identity]
\label{cor:v3v27-balanced}
Relative to \(F=I\), every positive scalar eigenmode contributes
\begin{equation}
 \frac{Z_{\lambda,F}}{Z_{\lambda,I}}=1.
\label{eq:v3v27-relative-one}
\end{equation}
Equivalently, the coupled real \((a,b)\) block contributes
\(f_\lambda^{-1}\) and the complex ghost block contributes
\(f_\lambda\).  This derives the balanced cancellation anticipated in
Proposition~\ref{prop:v3v26-balanced}.
\end{corollary}

\begin{proof}
Divide \eqref{eq:v3v27-mode-cancel} by its value at \(f_\lambda=1\).
\end{proof}

\subsection{Finite cutoff determinant of the coupled block}
\label{sec:v3v27-finite-cutoff}

The complex symmetric bosonic Hessian in the ordered variables
\((a,b_\lambda)\) is
\begin{equation}
 \mathbb B_{\lambda,F}
 =
 \begin{pmatrix}
 0&i f_\lambda\sqrt\lambda\\
 i f_\lambda\sqrt\lambda&\xi f_\lambda
 \end{pmatrix},
 \qquad
 \det\mathbb B_{\lambda,F}=f_\lambda^2\lambda.
\label{eq:v3v27-Bmatrix}
\end{equation}
The regulator in Theorem~\ref{thm:v3v27-mode-integral}, rather than a
formal square root of this complex matrix, fixes its Gaussian branch.

\begin{theorem}[Exact common-cutoff cancellation]
\label{thm:v3v27-cutoff-cancel}
Let \(P_\Lambda=\mathbf1_{(0,\Lambda^2]}(\Delta_0)\), and apply
\eqref{eq:v3v27-Snm} to the paired exact gauge, multiplier, and complex
ghost modes in \(P_\Lambda\).  Then
\begin{equation}
 \frac{Z_{{\rm nm},F,\Lambda}}
      {Z_{{\rm nm},I,\Lambda}}
 =1
\label{eq:v3v27-cutoff-cancel}
\end{equation}
for every finite \(\Lambda\).  The identity holds before a heat expansion,
modified determinant, or continuum limit.
\end{theorem}

\begin{proof}
The cutoff space is an orthogonal sum of the mode blocks above, including
multiplicity.  Equation~\eqref{eq:v3v27-relative-one} equals one on every
summand, so their finite product is one.
\end{proof}

\begin{corollary}[No nonminimal reservoir counterterm]
\label{cor:v3v27-no-counterterm}
If one common smooth regulator is used and \(F\) is inserted through the
BRST-exact action \eqref{eq:v3v27-Snm}, the exact-gauge,
multiplier, and ghost sectors have zero relative heat coefficients and
zero relative modified determinant.  They generate no reservoir
counterterm.
\end{corollary}

\begin{proof}
The regulated logarithm of
\eqref{eq:v3v27-cutoff-cancel} is identically zero at every cutoff.
Every coefficient and the cutoff-independent remainder therefore vanish.
\end{proof}

\begin{assessment}[Correction of the V26 conditional branch]
\label{ass:v3v27-v26-correction}
Equation~\eqref{eq:v3v26-FP-reduced} assigned a reservoir factor to the
gauge and ghost Hessians while leaving the multiplier normalization outside
the relative determinant.  V27 shows that this is not the determinant of
the common-\(F\) BRST-exact gauge fermion.  Once the multiplier is retained
inside the coupled block, its additional factor completes the cancellation.
The V26 half-power remains correct for the determinant assignment stated
there, but that assignment is not the balanced action
\eqref{eq:v3v27-Snm}.
\end{assessment}

\subsection{The physical coexact factor remains}
\label{sec:v3v27-physical}

The Maxwell action is nonzero on coexact one-forms and the gauge-fixing
fermion does not pair them with scalar ghosts.  Suppose their physical
reservoir ratio is
\begin{equation}
 R_1^{\rm phys}
 =I+\rho(\Delta_1+\mu^2)^{-1/2}
 \quad\hbox{on }\mathcal C^1.
\label{eq:v3v27-Rphys}
\end{equation}

\begin{corollary}[Abelian relative determinant after the exact
nonminimal cancellation]
\label{cor:v3v27-physical-det}
At common finite cutoff, the nonzero Abelian relative determinant reduces
to
\begin{equation}
 Z_{\rm rel,\Lambda}
 =\det(R_1^{\rm phys}|_{\mathcal C^1})_\Lambda^{-1/2}.
\label{eq:v3v27-physical-det}
\end{equation}
After subtracting its first three local terms, the nonlocal remainder is
\begin{equation}
 \log Z_{\rm rel,4}
 =-\frac12\log\det{}_4
 (R_1^{\rm phys}|_{\mathcal C^1}).
\label{eq:v3v27-physical-det4}
\end{equation}
\end{corollary}

\begin{proof}
Theorem~\ref{thm:v3v27-cutoff-cancel} removes the relative contribution of
the exact gauge, multiplier, and ghost block.  The orthogonal Hodge
splitting leaves the coexact real Gaussian, whose determinant power is
\(-1/2\).  The order-minus-one Schatten result of V21 defines the fourth
modified determinant.
\end{proof}

\begin{firewall}[Harmonic and stabilizer modes remain outside]
\label{fw:v3v27-zero}
Theorem~\ref{thm:v3v27-mode-integral} assumes \(\lambda>0\).  At
\(\lambda=0\), there is no normalized exact gauge vector
\(\lambda^{-1/2}du\), the ghost quadratic term vanishes, and the Berezin
integral is unsaturated.  Constant gauge transformations and harmonic
one-forms retain the quotient and moduli measures identified in V26.
\end{firewall}

\subsection{Locality obstruction for the desired spectral factor}
\label{sec:v3v27-locality}

The physical reservoir factor from V20--V21 is
\begin{equation}
 F_{\mu,\rho}
 =I+\rho(\Delta_0+\mu^2)^{-1/2}.
\label{eq:v3v27-Fres}
\end{equation}

\begin{theorem}[No finite-order local boundary operator equals the
reservoir factor]
\label{thm:v3v27-nonlocal-F}
If \(\rho\ne0\), the operator \(F_{\mu,\rho}\) in
\eqref{eq:v3v27-Fres} is not a finite-order differential operator on
\(N\).  Its classical symbol begins
\begin{equation}
 \sigma(F_{\mu,\rho})(x,\xi)
 \sim I+\rho|\xi|^{-1}I+O(|\xi|^{-2}).
\label{eq:v3v27-F-symbol}
\end{equation}
Therefore the gauge fermion \(\Psi_{F_{\mu,\rho}}\), viewed as a
boundary-only functional, is pseudodifferential and nonlocal.
\end{theorem}

\begin{proof}
Functional calculus for the positive Laplacian gives the classical
order-minus-one operator
\((\Delta_0+\mu^2)^{-1/2}\) with principal symbol \(|\xi|^{-1}\).
A finite-order differential operator has a polynomial symbol in \(\xi\).
If its total order is zero, it is multiplication by a smooth endomorphism
and has no negative homogeneous tail.  If its order is positive, it has a
nonzero positive-degree principal symbol.  Neither possibility matches
\eqref{eq:v3v27-F-symbol} when \(\rho\ne0\).
\end{proof}

\begin{firewall}[Nonlocal gauge fixing must be localized before the full
continuum claim]
\label{fw:v3v27-localize}
The finite-cutoff cancellation is mathematically exact for
\(F_{\mu,\rho}\), but writing that operator directly in a boundary gauge
fermion does not give a local four-dimensional SM--Einstein action.  A
local continuum construction must realize the same factor by extra
half-cylinder fields arranged in BRST doublets, with their own local
action, trace relation, domains, and measures.
\end{firewall}

\subsection{V27 conclusion and V28 target}
\label{sec:v3v27-conclusion}

V27 derives the missing payer from a complete BRST-exact nonminimal
action.  On every positive scalar eigenmode, the coupled longitudinal
gauge and multiplier Gaussian contributes \(f_\lambda^{-1}\), while the
complex ghost contributes \(f_\lambda\).  A convergent
\(\varepsilon\)-prescription fixes the Euclidean contour and phase.  The
relative reservoir factor therefore cancels exactly at every common cutoff,
leaving only the physical coexact determinant and the separate zero-mode
measures.

The desired factor
\(I+\rho(\Delta_0+\mu^2)^{-1/2}\) is nonlocal as a boundary operator.
V28 should construct a local half-cylinder BRST doublet whose elimination
produces the common \(F_{\mu,\rho}\) simultaneously in the multiplier,
longitudinal, and ghost terms.  It must prove locality and nilpotence before
elimination, reproduce \eqref{eq:v3v27-cutoff-cancel} after elimination,
and test whether the trace constraints overdetermine any member of the
doublet.  If one local reservoir cannot generate the inverse-square-root
factor in all three terms, that failure must be recorded instead of using
the boundary pseudodifferential gauge fermion as a continuum action.

\section{V28: local Neumann reservoir quartet realizing the nonlocal
gauge-fixing multiplier}
\label{sec:v3v28}

\subsection{Going upstream from the inverse square root}
\label{sec:v3v28-purpose}

V27 proved that the common multiplier
\[
 F_{\mu,\rho}=I+\rho(\Delta_0+\mu^2)^{-1/2}
\]
cancels from the nonminimal Abelian determinant, but also proved that
writing \(F_{\mu,\rho}\) directly in a boundary gauge fermion is nonlocal.
The inverse square root is not a Dirichlet-to-Neumann map.  It is the
Neumann-to-Dirichlet map of the same massive half-cylinder.  V28 uses that
fact to construct a local bulk quartet whose elimination produces the
entire common multiplier in one step.

The construction is first made at a common finite tangential and radial
cutoff, where every integral is finite dimensional.  The operator
identities then pass to the energy spaces by spectral convergence.  This
keeps determinant normalization visible.

\subsection{Exact Neumann-to-Dirichlet map}
\label{sec:v3v28-NtD}

Let \(N\) be the closed boundary and set
\begin{equation}
 A=\Delta_0+\mu^2,\qquad\mu>0.
\label{eq:v3v28-A}
\end{equation}
On \(X=N\times[0,\infty)\), define the local symmetric energy form
\begin{equation}
 \mathfrak a(u,v)
 =\int_0^\infty
 \left(
 \langle\partial_ru,\partial_rv\rangle_N
 +\langle A^{1/2}u,A^{1/2}v\rangle_N
 \right)\,dr.
\label{eq:v3v28-a}
\end{equation}
Let \(T u=u(0)\) be the trace.  For a boundary source \(f\), define
\(\mathcal Rf\) weakly by
\begin{equation}
 \mathfrak a(v,\mathcal Rf)=\langle Tv,f\rangle_N
 \qquad\hbox{for every energy test field }v.
\label{eq:v3v28-Rweak}
\end{equation}

\begin{theorem}[The reservoir response is the inverse square root]
\label{thm:v3v28-NtD}
For every smooth \(f\),
\begin{equation}
 (\mathcal Rf)(r)
 =e^{-rA^{1/2}}A^{-1/2}f,
\label{eq:v3v28-R}
\end{equation}
and
\begin{equation}
 \mathcal N:=T\mathcal R=A^{-1/2}.
\label{eq:v3v28-N}
\end{equation}
Equivalently, \(\mathcal Rf\) solves
\begin{equation}
 (-\partial_r^2+A)\mathcal Rf=0,\qquad
 -\partial_r\mathcal Rf(0)=f,
\label{eq:v3v28-Neumann}
\end{equation}
and decays at infinity.  The response is unique in the energy space.
\end{theorem}

\begin{proof}
Expand \(f=\sum_jf_ju_j\) in an eigenbasis
\(Au_j=a_ju_j\), \(a_j\ge\mu^2\).  The decaying solution of the
\(j\)-th equation with outward Neumann datum \(f_j\) is
\[
 f_ja_j^{-1/2}e^{-r\sqrt{a_j}}u_j.
\]
Its boundary trace is \(a_j^{-1/2}f_ju_j\), proving
\eqref{eq:v3v28-R}--\eqref{eq:v3v28-N}.  Integration by parts gives
\[
 \mathfrak a(v,\mathcal Rf)
 =\langle v(0),-\partial_r\mathcal Rf(0)\rangle
 =\langle Tv,f\rangle,
\]
so it satisfies \eqref{eq:v3v28-Rweak}.  If the source is zero, testing
with the response itself gives zero energy; the mass gap forces the field
to vanish.
\end{proof}

\begin{corollary}[Positive source energy]
\label{cor:v3v28-source-energy}
The on-shell response obeys
\begin{equation}
 \mathfrak a(\mathcal Rf,\mathcal Rf)
 =\langle f,A^{-1/2}f\rangle.
\label{eq:v3v28-source-energy}
\end{equation}
Thus \(\mathcal N\) is positive and self-adjoint.
\end{corollary}

\begin{proof}
Set \(v=\mathcal Rf\) in \eqref{eq:v3v28-Rweak} and use
\eqref{eq:v3v28-N}.  Polarization gives self-adjointness.
\end{proof}

\subsection{The local reservoir quartet}
\label{sec:v3v28-quartet}

Retain the boundary BRST fields and convention of V27.  Put
\begin{equation}
 Y=\delta A_{\rm gauge}-\frac{i\xi}{2}b,
\qquad
 sY=\Delta_0c.
\label{eq:v3v28-Y}
\end{equation}
The subscript on \(A_{\rm gauge}\) distinguishes the gauge field from the
positive scalar operator \(A\) in \eqref{eq:v3v28-A}.

Introduce four scalar reservoir fields on \(X\):
\[
 \chi,\beta\ \hbox{bosonic},\qquad
 C,\bar\chi\ \hbox{fermionic},
\]
with ghost numbers \(0,0,1,-1\), respectively, and
\begin{equation}
 s\chi=C,\qquad sC=0,\qquad
 s\bar\chi=i\beta,\qquad s\beta=0.
\label{eq:v3v28-res-BRST}
\end{equation}
This is a pair of algebraic BRST doublets.

Define the local extended gauge fermion
\begin{equation}
 \begin{split}
 \Psi_{\rm loc}
 ={}&\langle\bar c,Y\rangle_N\\
 &+\rho\left[
 -\mathfrak a(\bar\chi,\chi)
 +\langle\bar c,T\chi\rangle_N
 +\langle T\bar\chi,Y\rangle_N
 \right].
 \end{split}
\label{eq:v3v28-Psi-local}
\end{equation}
The mixed energy pairing is interpreted by complex bilinear extension for
the Grassmann field \(\bar\chi\).

\begin{proposition}[Locality, ghost number, and nilpotence]
\label{prop:v3v28-local}
Every term of \(\Psi_{\rm loc}\) is a local bulk or boundary integral of
finite-order differential expressions.  The gauge fermion has ghost number
minus one and odd parity.  The enlarged BRST differential is nilpotent, so
\begin{equation}
 S_{\rm loc}=s\Psi_{\rm loc}
\label{eq:v3v28-Sloc}
\end{equation}
is BRST closed.
\end{proposition}

\begin{proof}
The form \(\mathfrak a\) contains one normal or tangential derivative on
each factor and the local mass term; \(T\) is restriction to the boundary.
The stated ghost numbers make every bracket in
\eqref{eq:v3v28-Psi-local} have ghost number minus one and odd parity.
Equation~\eqref{eq:v3v28-res-BRST} squares to zero, as does the boundary
differential in \eqref{eq:v3v27-BRST}.  The two act on disjoint generators,
so their sum is nilpotent and \(sS_{\rm loc}=s^2\Psi_{\rm loc}=0\).
\end{proof}

\subsection{Elimination of the local quartet}
\label{sec:v3v28-elimination}

\begin{theorem}[The local action imposes the two Neumann response
constraints]
\label{thm:v3v28-constraints}
The reservoir-dependent part of \(S_{\rm loc}\) is
\begin{equation}
 \begin{split}
 S_{\rm res}
 =\rho\bigl[
 &-i\mathfrak a(\beta,\chi)
 +\mathfrak a(\bar\chi,C)
 +i\langle b,T\chi\rangle
 -\langle\bar c,TC\rangle\\
 &+i\langle T\beta,Y\rangle
 -\langle T\bar\chi,\Delta_0c\rangle
 \bigr].
 \end{split}
\label{eq:v3v28-Sres}
\end{equation}
With the real multiplier contour for \(\beta\) and the standard Berezin
orientation for \(\bar\chi\), integration imposes, in weak form,
\begin{align}
 \mathfrak a(v,\chi)&=\langle Tv,Y\rangle
 &&\hbox{for all bosonic }v,
\label{eq:v3v28-chi-constraint}\\
 \mathfrak a(\bar v,C)&=\langle T\bar v,\Delta_0c\rangle
 &&\hbox{for all fermionic }\bar v.
\label{eq:v3v28-C-constraint}
\end{align}
Hence
\begin{equation}
 \chi=\mathcal RY,\qquad
 C=\mathcal R\Delta_0c.
\label{eq:v3v28-responses}
\end{equation}
\end{theorem}

\begin{proof}
Apply the graded Leibniz rule to
\eqref{eq:v3v28-Psi-local}.  Since
\(\bar\chi\) is odd,
\[
 s[-\mathfrak a(\bar\chi,\chi)]
 =-i\mathfrak a(\beta,\chi)+\mathfrak a(\bar\chi,C).
\]
The two boundary pairings give the remaining four terms in
\eqref{eq:v3v28-Sres}.  The field \(\beta\) occurs linearly as
\[
 -i\rho\left[
 \mathfrak a(\beta,\chi)-\langle T\beta,Y\rangle
 \right],
\]
so its Fourier integral is the delta functional
\eqref{eq:v3v28-chi-constraint}.  The Grassmann field \(\bar\chi\) occurs
linearly in
\[
 \rho\left[
 \mathfrak a(\bar\chi,C)
 -\langle T\bar\chi,\Delta_0c\rangle
 \right],
\]
whose Berezin integral imposes the fermionic linear constraint
\eqref{eq:v3v28-C-constraint}.  Theorem~\ref{thm:v3v28-NtD} gives
\eqref{eq:v3v28-responses}.
\end{proof}

\begin{theorem}[Exact effective boundary gauge fermion and action]
\label{thm:v3v28-effective}
After the local reservoir quartet is eliminated with the preceding
contours,
\begin{equation}
 \Psi_{\rm eff}
 =\left\langle
 \bar c,
 \left(I+\rho A^{-1/2}\right)Y
 \right\rangle_N.
\label{eq:v3v28-Psi-effective}
\end{equation}
Equivalently, the effective nonminimal action is
\begin{equation}
 \begin{split}
 S_{\rm eff}
 ={}&
 i\langle b,F_{\mu,\rho}\delta A_{\rm gauge}\rangle
 +\frac{\xi}{2}\langle b,F_{\mu,\rho}b\rangle\\
 &-\langle\bar c,F_{\mu,\rho}\Delta_0c\rangle,
 \qquad
 F_{\mu,\rho}=I+\rho(\Delta_0+\mu^2)^{-1/2}.
 \end{split}
\label{eq:v3v28-S-effective}
\end{equation}
\end{theorem}

\begin{proof}
It is useful first to eliminate the reservoir fields in the gauge fermion
itself.  Stationarity in \(\bar\chi\) and \(\chi\) gives
\(\chi=\mathcal RY\) and
\(\bar\chi=\mathcal R\bar c\), with the second identity understood by
duality.  At that stationary point,
\[
 -\mathfrak a(\bar\chi,\chi)
 +\langle\bar c,T\chi\rangle
 +\langle T\bar\chi,Y\rangle
 =\langle\bar c,\mathcal NY\rangle.
\]
This uses the weak identity
\(\mathfrak a(\mathcal R\bar c,\mathcal RY)
=\langle\bar c,\mathcal NY\rangle\).
By Theorem~\ref{thm:v3v28-NtD},
\(\mathcal N=A^{-1/2}\), proving
\eqref{eq:v3v28-Psi-effective}.

The direct action elimination in
Theorem~\ref{thm:v3v28-constraints} gives the same result:
the remaining terms in \eqref{eq:v3v28-Sres} are
\[
 i\rho\langle b,\mathcal NY\rangle
 -\rho\langle\bar c,\mathcal N\Delta_0c\rangle.
\]
Add \(s\langle\bar c,Y\rangle
=i\langle b,Y\rangle-\langle\bar c,\Delta_0c\rangle\), insert
\(Y=\delta A_{\rm gauge}-i\xi b/2\), and obtain
\eqref{eq:v3v28-S-effective}.
\end{proof}

\begin{corollary}[The V27 cancellation now has a local parent]
\label{cor:v3v28-local-parent}
At every common finite cutoff, eliminating the local quartet and then
integrating the boundary nonminimal block gives
\begin{equation}
 \frac{Z_{{\rm nm},\rho,\Lambda}}
      {Z_{{\rm nm},0,\Lambda}}=1.
\label{eq:v3v28-local-cancel}
\end{equation}
Thus no nonphysical reservoir determinant or local counterterm remains.
\end{corollary}

\begin{proof}
Theorem~\ref{thm:v3v28-effective} produces exactly the common spectral
multiplier used in Theorem~\ref{thm:v3v27-cutoff-cancel}.  Apply that
theorem mode by mode.
\end{proof}

\subsection{Quartet Jacobians at finite cutoff}
\label{sec:v3v28-jacobians}

\begin{proposition}[Bosonic and fermionic reservoir Jacobians cancel]
\label{prop:v3v28-jacobians}
Choose identical finite-dimensional Galerkin spaces and the same energy
matrix \(H_\Lambda>0\) for
\((\chi,\beta)\) and \((C,\bar\chi)\).  Integrating the constraints in
\eqref{eq:v3v28-Sres} produces the factors
\begin{equation}
 (\det H_\Lambda)^{-1}
 \quad\hbox{and}\quad
 \det H_\Lambda,
\label{eq:v3v28-jacobians}
\end{equation}
respectively.  Their product is one, including all multiplicities.
\end{proposition}

\begin{proof}
The finite bosonic Fourier delta
\(\delta(H_\Lambda\chi-j)\) changes variables with Jacobian
\(|\det H_\Lambda|^{-1}\).  The fermionic linear integral
\(\int d\bar\chi\,dC\,
\exp[-\bar\chi(H_\Lambda C-\eta)]\) yields
\(\det H_\Lambda\).  Positivity fixes the bosonic sign, and the common
Berezin orientation fixes the fermionic sign.  Multiplication proves the
claim.
\end{proof}

\begin{firewall}[The same domains and discretization are part of the
quartet]
\label{fw:v3v28-domains}
The cancellation in \eqref{eq:v3v28-jacobians} fails if the bosonic and
fermionic reservoirs use different radial boundary conditions, tangential
cutoffs, masses, or remote-end conditions.  Equality of formal differential
expressions is insufficient; the energy matrices and their domains must be
identical.
\end{firewall}

\subsection{No overdetermination}
\label{sec:v3v28-boundary-count}

\begin{proposition}[One source condition per reservoir field]
\label{prop:v3v28-not-overdetermined}
The weak constraints
\eqref{eq:v3v28-chi-constraint}--\eqref{eq:v3v28-C-constraint} impose the
outward Neumann data
\begin{equation}
 -\partial_r\chi(0)=Y,\qquad
 -\partial_rC(0)=\Delta_0c,
\label{eq:v3v28-Neumann-data}
\end{equation}
together with decay at infinity.  They impose no independent Dirichlet
trace.  The positive mass gives a unique solution for every source.
\end{proposition}

\begin{proof}
Integration by parts in the weak equations gives the interior homogeneous
equation and exactly the boundary data
\eqref{eq:v3v28-Neumann-data}.  Theorem~\ref{thm:v3v28-NtD} proves
existence and uniqueness.
\end{proof}

\begin{counterexample}[Simultaneous trace and source data are generically
inconsistent]
\label{cex:v3v28-overdetermined}
For a tangential eigenmode \(Au=au\), \(a>0\), the decaying solution with
Neumann source \(f u\) is
\[
 \chi(r)=f a^{-1/2}e^{-\sqrt a r}u.
\]
Its trace is forced to be \(f a^{-1/2}u\).  Adding the independent
condition \(\chi(0)=0\) makes the problem inconsistent whenever \(f\ne0\).
Thus the local realization must use the Neumann source relation alone.
\end{counterexample}

\subsection{Sign comparison with the physical reservoir}
\label{sec:v3v28-sign}

\begin{assessment}[Two maps, two roles]
\label{ass:v3v28-two-maps}
The physical trace-constrained energy in V20 produces the positive
Dirichlet-to-Neumann map \(A^{1/2}\) and adds
\(\rho A^{1/2}\) to a boundary Hessian.  The gauge-fixing localization in
V28 uses a Neumann source and produces \(A^{-1/2}\) inside a fermionic gauge
fermion.  There is no sign contradiction: the first is an on-shell positive
energy, while the second is a response operator enforced by a BRST quartet
whose bosonic and fermionic Jacobians cancel.
\end{assessment}

\begin{firewall}[Do not use the Neumann quartet as a physical mass
reservoir]
\label{fw:v3v28-not-physical}
The fields \((\chi,C,\bar\chi,\beta)\) are contractible gauge-fixing
variables.  Their transfer function localizes a nonminimal multiplier but
does not create a physical pole or a positive coexact gauge-field gap.  The
physical determinant in \eqref{eq:v3v27-physical-det4} remains a separate
sector.
\end{firewall}

\subsection{V28 conclusion and V29 target}
\label{sec:v3v28-conclusion}

V28 supplies the missing local parent of the V27 cancellation.  The
Neumann-to-Dirichlet map of a massive half-cylinder is exactly
\((\Delta_0+\mu^2)^{-1/2}\).  Two local reservoir BRST doublets, coupled
through one energy form and local boundary sources, generate
\[
 F_{\mu,\rho}=I+\rho(\Delta_0+\mu^2)^{-1/2}
\]
in the complete gauge-fixing fermion.  Their bosonic and fermionic
Jacobians cancel at every common cutoff, their source problems are not
overdetermined, and subsequent integration of the boundary nonminimal
block gives unit relative determinant.  Locality and nilpotence therefore
hold before elimination.

V29 should extend the construction from \(d\) to a background covariant
differential \(d_{\mathcal A}\).  It must distinguish a flat background,
where \(d_{\mathcal A}^2=0\) and the entire proof transfers, from a curved
Yang--Mills background, where
\(d_{\mathcal A}^2=[F_{\mathcal A},\,\cdot\,]\) obstructs the naive
two-term complex.  The next proof should calculate that curvature defect,
identify the additional BV or Hessian terms that restore the linearized
gauge identity on shell, and retain the physical coexact/nonzero-curvature
sector rather than forcing an Abelian Hodge cancellation.

\section{V29: background-covariant Yang--Mills Jacobi complex and
on-shell reservoir cancellation}
\label{sec:v3v29}

\subsection{Two different curvature questions}
\label{sec:v3v29-purpose}

Replacing \(d\) by a covariant differential \(d_{\mathcal A}\) creates two
issues that must not be conflated.  First, the covariant de Rham sequence is
not a complex unless the background connection is flat.  Second, gauge
invariance still gives a two-term linearized gauge identity for the actual
Yang--Mills Hessian at every classical background, including backgrounds
with nonzero curvature.  It is the second identity, not a false assertion
\(d_{\mathcal A}^2=0\), that pairs ghost and longitudinal fluctuation
spectra.

V29 derives this identity from gauge covariance, transfers the V28 local
Neumann quartet to the adjoint ghost bundle, and quantifies the failure of
reservoir matching away from the Yang--Mills equation.

\subsection{Background notation and the failed covariant de Rham complex}
\label{sec:v3v29-background}

Let \(P\to N\) be a principal bundle with compact structure group \(G\) and
invariant inner product on its Lie algebra.  Let
\(\mathfrak g_P=\operatorname{Ad}P\), let \(\mathcal A\) be a smooth
background connection, and write
\[
 d_{\mathcal A}:\Omega^q(N;\mathfrak g_P)
 \longrightarrow\Omega^{q+1}(N;\mathfrak g_P).
\]

\begin{proposition}[Exact curvature defect of the covariant de Rham
sequence]
\label{prop:v3v29-d-square}
For every adjoint-valued form \(\omega\),
\begin{equation}
 d_{\mathcal A}^2\omega
 =[F_{\mathcal A}\wedge\omega].
\label{eq:v3v29-d-square}
\end{equation}
Thus \((\Omega^\bullet,d_{\mathcal A})\) is a differential complex for all
adjoint-valued forms if \(F_{\mathcal A}=0\).  For a generic noncentral
curvature it is not a complex.
\end{proposition}

\begin{proof}
In a local trivialization,
\(d_{\mathcal A}=d+[\mathcal A\wedge\,\cdot\,]\).  Expanding the square,
using the graded Leibniz rule and the Jacobi identity, gives
\[
 d_{\mathcal A}^2\omega
 =[d\mathcal A+\tfrac12[\mathcal A\wedge\mathcal A],\omega]
 =[F_{\mathcal A}\wedge\omega].
\]
If \(F_{\mathcal A}=0\), the square vanishes.  Conversely, any curvature
component whose adjoint action is nonzero supplies a form on which the
right side is nonzero.
\end{proof}

\begin{firewall}[Linear BRST nilpotence does not assert
\(d_{\mathcal A}^2=0\)]
\label{fw:v3v29-linear-BRST}
The linearized gauge transformation is
\(s_0a=d_{\mathcal A}c,\ s_0c=0\), and hence \(s_0^2a=0\) because the
background is fixed.  This two-term nilpotence does not extend
\(d_{\mathcal A}\) to a covariant de Rham complex through higher form
degrees.  Equation~\eqref{eq:v3v29-d-square} remains present.
\end{firewall}

\subsection{Gauge covariance of the equation and its Hessian identity}
\label{sec:v3v29-Jacobi}

On a compact Riemannian four-region or on a collar with compatible domains,
write the Yang--Mills equation as
\begin{equation}
 \mathcal E(\mathcal A)
 =d_{\mathcal A}^*F_{\mathcal A}.
\label{eq:v3v29-E}
\end{equation}
Let
\begin{equation}
 \mathcal M_{\mathcal A}
 =D\mathcal E_{\mathcal A}:
 \Omega^1(\mathfrak g_P)\longrightarrow\Omega^1(\mathfrak g_P)
\label{eq:v3v29-M}
\end{equation}
be the Jacobi operator, with domains chosen so that it is the symmetric
second-variation operator.

\begin{theorem}[Off-shell gauge-orbit defect]
\label{thm:v3v29-offshell}
With the infinitesimal convention
\(\delta_\epsilon\mathcal A=d_{\mathcal A}\epsilon\),
\begin{equation}
 \mathcal M_{\mathcal A}d_{\mathcal A}\epsilon
 =[\mathcal E(\mathcal A),\epsilon].
\label{eq:v3v29-offshell}
\end{equation}
In particular, at a Yang--Mills background
\(\mathcal E(\mathcal A)=0\),
\begin{equation}
 \mathcal M_{\mathcal A}d_{\mathcal A}=0.
\label{eq:v3v29-Jacobi-complex}
\end{equation}
\end{theorem}

\begin{proof}
The Yang--Mills equation is gauge covariant:
\[
 \mathcal E(\mathcal A^g)
 =g^{-1}\mathcal E(\mathcal A)g.
\]
Differentiate a one-parameter gauge transformation
\(g(t)\) at \(t=0\) using the stated infinitesimal convention.  The
derivative of the left side is
\(\mathcal M_{\mathcal A}d_{\mathcal A}\epsilon\); the derivative of the
right side is
\([\mathcal E(\mathcal A),\epsilon]\), with the same convention.
This proves \eqref{eq:v3v29-offshell}; the on-shell identity is immediate.
\end{proof}

\begin{corollary}[Adjoint Noether identity on shell]
\label{cor:v3v29-Noether}
If \(\mathcal A\) is Yang--Mills and
\(\mathcal M_{\mathcal A}\) is self-adjoint, then
\begin{equation}
 d_{\mathcal A}^*\mathcal M_{\mathcal A}=0.
\label{eq:v3v29-Noether}
\end{equation}
\end{corollary}

\begin{proof}
For smooth test fields \(a,\epsilon\),
\[
 \langle d_{\mathcal A}^*\mathcal M_{\mathcal A}a,\epsilon\rangle
 =\langle a,\mathcal M_{\mathcal A}d_{\mathcal A}\epsilon\rangle=0
\]
by \eqref{eq:v3v29-Jacobi-complex}.
\end{proof}

\subsection{Gauge-fixed Laplacians and exact spectral pairing}
\label{sec:v3v29-gauge-fixed}

Define the ghost and Feynman-gauge fluctuation operators
\begin{equation}
 L_0=d_{\mathcal A}^*d_{\mathcal A},
 \qquad
 L_1=\mathcal M_{\mathcal A}
 +d_{\mathcal A}d_{\mathcal A}^*.
\label{eq:v3v29-L01}
\end{equation}

\begin{theorem}[On-shell gauge--ghost intertwining]
\label{thm:v3v29-intertwining}
For every background,
\begin{equation}
 L_1d_{\mathcal A}-d_{\mathcal A}L_0
 =\operatorname{ad}_{\mathcal E(\mathcal A)},
\qquad
 \operatorname{ad}_{\mathcal E}(\epsilon)=[\mathcal E,\epsilon].
\label{eq:v3v29-L-defect}
\end{equation}
At a Yang--Mills background,
\begin{align}
 L_1d_{\mathcal A}&=d_{\mathcal A}L_0,
\label{eq:v3v29-L-intertwine}\\
 d_{\mathcal A}^*L_1&=L_0d_{\mathcal A}^*.
\label{eq:v3v29-L-adjoint-intertwine}
\end{align}
\end{theorem}

\begin{proof}
By definition,
\[
 L_1d_{\mathcal A}
 =\mathcal M_{\mathcal A}d_{\mathcal A}
 +d_{\mathcal A}d_{\mathcal A}^*d_{\mathcal A}
 =\operatorname{ad}_{\mathcal E(\mathcal A)}
 +d_{\mathcal A}L_0,
\]
using Theorem~\ref{thm:v3v29-offshell}.  This is
\eqref{eq:v3v29-L-defect}.  On shell the defect vanishes.  Taking adjoints
gives \eqref{eq:v3v29-L-adjoint-intertwine}.
\end{proof}

Assume from now on that \(\mathcal A\) is Yang--Mills and that, after adding
a common regulator mass,
\begin{equation}
 A_0=L_0+\mu^2I,\qquad
 A_1=L_1+\mu^2I
\label{eq:v3v29-A01}
\end{equation}
are self-adjoint and obey
\begin{equation}
 A_0,A_1\ge a_0I>0.
\label{eq:v3v29-positive}
\end{equation}
For \(L_0u=\lambda u\), \(\lambda>0\), define
\begin{equation}
 Uu=\lambda^{-1/2}d_{\mathcal A}u.
\label{eq:v3v29-U}
\end{equation}

\begin{proposition}[Non-Abelian longitudinal unitary]
\label{prop:v3v29-longitudinal-unitary}
The map \(U\) is unitary from
\((\ker L_0)^\perp\) onto the closure of
\(\operatorname{ran}d_{\mathcal A}\), and
\begin{equation}
 UL_0=L_1U,\qquad
 UA_0=A_1U.
\label{eq:v3v29-U-intertwine}
\end{equation}
\end{proposition}

\begin{proof}
For \(u\perp\ker L_0\),
\[
 \|d_{\mathcal A}L_0^{-1/2}u\|^2
 =\langle u,L_0^{-1/2}
 d_{\mathcal A}^*d_{\mathcal A}L_0^{-1/2}u\rangle
 =\|u\|^2.
\]
Its range is the closed exact subspace after completion.  Equation
\eqref{eq:v3v29-L-intertwine} gives the first intertwining identity, and
the common scalar mass gives the second.
\end{proof}

\begin{corollary}[On-shell matching of reservoir functions]
\label{cor:v3v29-reservoir-match}
For every bounded Borel function \(f\),
\begin{equation}
 f(A_1)|_{\overline{\operatorname{ran}d_{\mathcal A}}}
 =Uf(A_0)U^*.
\label{eq:v3v29-functional}
\end{equation}
In particular, the positive longitudinal fluctuation and ghost sectors
have identical Dirichlet-to-Neumann, Neumann-to-Dirichlet, and normalized
reservoir spectra at every common spectral cutoff.
\end{corollary}

\begin{proof}
Apply spectral functional calculus to
\eqref{eq:v3v29-U-intertwine}.
\end{proof}

\subsection{Local adjoint-valued Neumann quartet}
\label{sec:v3v29-local-quartet}

Replace the scalar bundle in V28 by the adjoint bundle and its ghost
operator \(A_0=L_0+\mu^2\).  On the product half-cylinder define
\begin{equation}
 \mathfrak a_{\mathcal A}(u,v)
 =\int_0^\infty
 \left(
 \langle\partial_ru,\partial_rv\rangle
 +\langle A_0^{1/2}u,A_0^{1/2}v\rangle
 \right)\,dr.
\label{eq:v3v29-energy}
\end{equation}
Although written spectrally in the second term,
\(\langle A_0^{1/2}u,A_0^{1/2}v\rangle
=\langle d_{\mathcal A}u,d_{\mathcal A}v\rangle
+\mu^2\langle u,v\rangle\), so the energy is local.

\begin{proposition}[Adjoint Neumann response]
\label{prop:v3v29-adjoint-N}
The weak Neumann response defined by
\[
 \mathfrak a_{\mathcal A}(v,\mathcal R_{\mathcal A}f)
 =\langle Tv,f\rangle
\]
has boundary transfer
\begin{equation}
 T\mathcal R_{\mathcal A}
 =(L_0+\mu^2)^{-1/2}.
\label{eq:v3v29-adjoint-N}
\end{equation}
\end{proposition}

\begin{proof}
The proof of Theorem~\ref{thm:v3v28-NtD} uses only positivity and the
spectral theorem.  Apply it to the adjoint-bundle operator \(A_0\).
\end{proof}

Use adjoint-valued reservoir fields
\((\chi,C,\bar\chi,\beta)\) with the BRST doublets
\eqref{eq:v3v28-res-BRST}, replace
\(\mathfrak a\) by \(\mathfrak a_{\mathcal A}\), and put
\begin{equation}
 Y=d_{\mathcal A}^*a-\frac{i\xi}{2}b.
\label{eq:v3v29-Y}
\end{equation}
At the linearized level
\begin{equation}
 s_0Y=L_0c.
\label{eq:v3v29-sY}
\end{equation}

\begin{theorem}[Local non-Abelian linearized gauge-fixing multiplier]
\label{thm:v3v29-local-multiplier}
At an on-shell background, the local quartet construction of V28
eliminates to the BRST-exact quadratic action with common multiplier
\begin{equation}
 F_{\mathcal A,\mu,\rho}
 =I+\rho(L_0+\mu^2)^{-1/2}.
\label{eq:v3v29-F}
\end{equation}
At every common finite cutoff, its longitudinal, multiplier, ghost, and
reservoir-quartet relative determinant is one.
\end{theorem}

\begin{proof}
Equation~\eqref{eq:v3v29-sY} has the same form as
\eqref{eq:v3v28-Y}; Proposition~\ref{prop:v3v29-adjoint-N} supplies the
same Neumann response.  Therefore the local gauge fermion
\eqref{eq:v3v28-Psi-local}, with all pairings in the adjoint bundle,
eliminates to \(F_{\mathcal A,\mu,\rho}\).  The bosonic and fermionic
reservoir Jacobians cancel by
Proposition~\ref{prop:v3v28-jacobians}.  Finally,
Corollary~\ref{cor:v3v29-reservoir-match} pairs the ghost and longitudinal
spectral factors, and the finite-mode calculation of V27 cancels the
multiplier normalization.
\end{proof}

\begin{firewall}[The background equation is a determinant hypothesis]
\label{fw:v3v29-onshell}
Theorem~\ref{thm:v3v29-local-multiplier} uses
\(\mathcal E(\mathcal A)=0\).  Gauge invariance of the nonlinear action
does not make the off-shell Hessian annihilate gauge directions.  Expanding
the continuum measure about an arbitrary background requires the explicit
defect \(\operatorname{ad}_{\mathcal E(\mathcal A)}\) or a full BV source
term.
\end{firewall}

\subsection{Quantitative off-shell mismatch}
\label{sec:v3v29-offshell-bound}

Let
\begin{equation}
 R=A_1d_{\mathcal A}-d_{\mathcal A}A_0
 =\operatorname{ad}_{\mathcal E(\mathcal A)}
\label{eq:v3v29-R}
\end{equation}
off shell, and suppose \(A_0,A_1\ge a_0I>0\).  If
\(\mathcal E(\mathcal A)\in L^\infty\), the adjoint bracket defines a
bounded zero-order operator with
\begin{equation}
 \|R\|\le C_{\mathfrak g}
 \|\mathcal E(\mathcal A)\|_{L^\infty}.
\label{eq:v3v29-R-bound}
\end{equation}

\begin{proposition}[Off-shell DtN chain defect]
\label{prop:v3v29-DtN-defect}
The massive square-root maps obey
\begin{equation}
 \left\|
 A_1^{1/2}d_{\mathcal A}
 -d_{\mathcal A}A_0^{1/2}
 \right\|
 \le
 \frac{C_{\mathfrak g}}{2\sqrt{a_0}}
 \|\mathcal E(\mathcal A)\|_{L^\infty},
\label{eq:v3v29-DtN-defect}
\end{equation}
between the corresponding graph spaces.
\end{proposition}

\begin{proof}
Use Proposition~\ref{prop:v3v23-defect} with the defect
\eqref{eq:v3v29-R}, followed by
\eqref{eq:v3v29-R-bound}.
\end{proof}

\begin{counterexample}[Small curvature is not the same as small
equation defect]
\label{cex:v3v29-curvature-vs-E}
A background can have \(F_{\mathcal A}\ne0\) and still satisfy
\(d_{\mathcal A}^*F_{\mathcal A}=0\); instantons are the standard
geometric class of examples.  For such a background,
\(d_{\mathcal A}^2\ne0\) on generic adjoint forms, but the Hessian
intertwining defect in \eqref{eq:v3v29-L-defect} is zero.  Conversely,
an arbitrarily small-looking connection in one gauge need not give a
useful bound unless the gauge-invariant equation defect
\(\mathcal E(\mathcal A)\) is controlled.
\end{counterexample}

\subsection{Physical slice and moduli}
\label{sec:v3v29-physical}

At an on-shell background, define the background gauge slice
\begin{equation}
 \mathcal T_{\mathcal A}
 =\ker d_{\mathcal A}^*
\label{eq:v3v29-slice}
\end{equation}
and the stabilizer algebra
\begin{equation}
 \mathfrak h_{\mathcal A}=\ker L_0.
\label{eq:v3v29-stabilizer}
\end{equation}

\begin{proposition}[Invariant physical slice]
\label{prop:v3v29-slice-invariant}
The gauge-fixed Hessian \(L_1\) preserves
\(\mathcal T_{\mathcal A}\), and the orthogonal decomposition
\begin{equation}
 L^2\Omega^1(\mathfrak g_P)
 =\overline{\operatorname{ran}d_{\mathcal A}}
 \mathbin{\widehat\oplus}\mathcal T_{\mathcal A}
\label{eq:v3v29-splitting}
\end{equation}
reduces \(L_1\).  Zero modes of
\(\mathcal M_{\mathcal A}|_{\mathcal T_{\mathcal A}}\) are infinitesimal
Yang--Mills moduli and remain outside the primed determinant.
\end{proposition}

\begin{proof}
If \(a\in\ker d_{\mathcal A}^*\), then
\[
 d_{\mathcal A}^*L_1a
 =L_0d_{\mathcal A}^*a=0
\]
by \eqref{eq:v3v29-L-adjoint-intertwine}.  Self-adjointness gives
invariance of its orthogonal complement
\(\overline{\operatorname{ran}d_{\mathcal A}}\).
The final statement is the definition of the kernel of the Jacobi operator
modulo gauge directions.
\end{proof}

\begin{assessment}[What remains after the quartet]
\label{ass:v3v29-remaining}
The local reservoir quartet removes only the nonminimal and longitudinal
relative determinant.  The determinant of the physical slice
\(\mathcal T_{\mathcal A}\), the stabilizer volume
\(\mathfrak h_{\mathcal A}\), negative modes of a nonminimal critical
background, and the moduli kernel remain.  Adding \(\mu^2\) makes the
reservoir response bounded but does not change those Yang--Mills
classifications.
\end{assessment}

\subsection{V29 conclusion and V30 audit mandate}
\label{sec:v3v29-conclusion}

V29 replaces the false covariant de Rham complex by the correct on-shell
Jacobi complex.  Although
\(d_{\mathcal A}^2=[F_{\mathcal A},\,\cdot\,]\), gauge covariance gives
\[
 \mathcal M_{\mathcal A}d_{\mathcal A}
 =[\mathcal E(\mathcal A),\,\cdot\,].
\]
At a Yang--Mills background the gauge-fixed fluctuation operator
intertwines the ghost Laplacian exactly.  The adjoint-valued local Neumann
quartet of V28 then produces the common inverse-square-root multiplier and
cancels the longitudinal, multiplier, ghost, and reservoir determinants at
every common cutoff.  Off shell, the square-root mismatch is quantitatively
controlled by the Yang--Mills equation defect.

V30 is the next mandatory companion audit.  After fingerprinting the
active Yang--Mills and Navier--Stokes notes, it should choose between
constructing the physical-slice determinant uniformly over a compact
on-shell background family and transferring the Jacobi-complex argument
to the Einstein diffeomorphism Hessian.  The latter is the direct route
toward the full SM--Einstein complex, but it must derive the off-shell
Bianchi defect and boundary Calderon domains rather than copying the
Yang--Mills formula.

\section{V30: companion audit and the complete Einstein--matter
diffeomorphism Jacobi identity}
\label{sec:v3v30}

\subsection{Mandatory companion audit}
\label{sec:v3v30-audit}

The V30 audit was performed after V29 and before choosing between a
physical Yang--Mills determinant family and the Einstein diffeomorphism
Hessian.  The active companion files were
\begin{center}
\begin{tabular}{|p{0.30\textwidth}|p{0.12\textwidth}|p{0.43\textwidth}|}
\hline
programme and active file & bytes & SHA--256 \\
\hline
Navier--Stokes,
\texttt{Renewal\_NS\_Stability\_Research\_Notes\_V31.tex}
& 887537
& \texttt{F1121B62238AD5491BB7CF03635EA9934942EDF573ED8FAAA9EE79E1D38A6696}
\\
\hline
Yang--Mills,
\texttt{Renewal\_Yang\_Mills\_Mass\_Gap\_Research\_Notes\_V86.tex}
& 1384322
& \texttt{1183C7C7670E4FC8EC05AF3DC1DB48AB749E0FEE1E59907DDDF2C6CA81050115}
\\
\hline
\end{tabular}
\end{center}
The Navier--Stokes note remains unchanged.  Yang--Mills V86 restores the
two mandatory proper-prefix factors to the exact minimal \(2+2\) packet.
That complete packet neutralizes the V84--V85 adversarial sequence at the
sharp target scale.  The repair is not yet uniform: arbitrary row-cell data
still require a hot-prefix or two-cold-prefix dichotomy, and an overlap
operator may use the adjacent joint card only after algebra proves that the
two occurrences are actually adjacent.

\begin{transferprinciple}[Complete Euler--Lagrange packet]
\label{tp:v3v30-complete-EL}
For Einstein coupled to matter, form the diffeomorphism identity using the
metric and every matter Euler--Lagrange expression before taking a Hessian
block, norm, or determinant.  The matter equation is a mandatory prefix in
the metric Bianchi identity.  A metric-only cancellation obtained by
discarding that term is no more reliable than the isolated joining
coordinate refuted and repaired in Yang--Mills V84--V86.
\end{transferprinciple}

No companion theorem is transferred as an Einstein theorem.  The audit
selects the complete coupled action as the next source object.

\subsection{Abstract natural field theory}
\label{sec:v3v30-abstract}

Let \(M\) be a smooth four-manifold and let
\begin{equation}
 \Phi=(g,\varphi)
\label{eq:v3v30-Phi}
\end{equation}
denote the metric together with all matter fields.  The symbol \(\varphi\)
may include gauge connections, Higgs fields, spinors, and auxiliary fields.
Let \(S[\Phi]\) be a diffeomorphism-invariant local action, including the
boundary terms required by its chosen variational problem.  For a vector
field \(X\), the infinitesimal diffeomorphism generator is
\begin{equation}
 \mathcal R_\Phi X
 =\mathcal L_X\Phi
 =(\mathcal L_Xg,\mathcal L_X\varphi).
\label{eq:v3v30-R}
\end{equation}
For compactly supported variations, write
\begin{equation}
 DS_\Phi(V)=\langle\mathcal E(\Phi),V\rangle,
\qquad
 \mathcal E=(\mathcal E_g,\mathcal E_\varphi).
\label{eq:v3v30-E}
\end{equation}
The pairings include the volume density and the natural bundle dualities.

\begin{theorem}[Complete diffeomorphism Noether identity]
\label{thm:v3v30-Noether}
For compactly supported \(X\), or for \(X\) satisfying a boundary
condition that removes the diffeomorphism flux,
\begin{equation}
 \mathcal R_\Phi^*\mathcal E(\Phi)=0.
\label{eq:v3v30-Noether}
\end{equation}
The identity contains the Euler--Lagrange expressions of every field on
which the Lie derivative acts.
\end{theorem}

\begin{proof}
Diffeomorphism invariance gives
\[
 0=\left.\frac{d}{dt}\right|_{t=0}
 S[(\exp tX)^*\Phi]
 =DS_\Phi(\mathcal R_\Phi X)
 =\langle\mathcal E(\Phi),\mathcal R_\Phi X\rangle.
\]
Move the differential generator to its formal adjoint.  Since \(X\) is
arbitrary in the admitted class,
\(\mathcal R_\Phi^*\mathcal E(\Phi)=0\).
\end{proof}

\begin{firewall}[A boundary flux is not zero by notation]
\label{fw:v3v30-boundary-flux}
If \(X\) is nonzero at a physical boundary, integration by parts in
Theorem~\ref{thm:v3v30-Noether} produces Brown--York, matter-flux, and
possibly corner terms.  The bulk identity alone applies either to
boundary-vanishing gauge transformations or after those boundary
currencies and edge variables have been included.  V10 supplies the
Einstein--GHY stress channel; it must not be silently discarded.
\end{firewall}

\subsection{Differentiated Noether identity and the Jacobi packet}
\label{sec:v3v30-Jacobi}

Let
\begin{equation}
 \mathbb H_\Phi=D\mathcal E_\Phi
\label{eq:v3v30-H}
\end{equation}
be the full coupled Jacobi operator.  At a critical point it is the
operator representing the symmetric second variation, after the field
metrics and boundary domains are fixed.

\begin{theorem}[Off-shell differentiated Noether identity]
\label{thm:v3v30-differentiated-Noether}
For every variation \(V\),
\begin{equation}
 \mathcal R_\Phi^*\mathbb H_\Phi V
 =-\left(D\mathcal R_\Phi^*[V]\right)\mathcal E(\Phi).
\label{eq:v3v30-differentiated-Noether}
\end{equation}
At a full Einstein--matter solution
\(\mathcal E(\Phi)=0\),
\begin{equation}
 \mathcal R_\Phi^*\mathbb H_\Phi=0.
\label{eq:v3v30-left-Jacobi}
\end{equation}
If \(\mathbb H_\Phi\) is self-adjoint on the selected domain, then also
\begin{equation}
 \mathbb H_\Phi\mathcal R_\Phi=0.
\label{eq:v3v30-right-Jacobi}
\end{equation}
\end{theorem}

\begin{proof}
Differentiate the identity
\(\mathcal R_\Phi^*\mathcal E(\Phi)=0\) in direction \(V\).  The product
rule gives
\[
 (D\mathcal R_\Phi^*[V])\mathcal E(\Phi)
 +\mathcal R_\Phi^*\mathbb H_\Phi V=0,
\]
which is \eqref{eq:v3v30-differentiated-Noether}.  At a critical point the
first term vanishes.  Taking the adjoint of
\eqref{eq:v3v30-left-Jacobi} gives
\eqref{eq:v3v30-right-Jacobi}.
\end{proof}

\begin{corollary}[Block identities require the matter cross-Hessian]
\label{cor:v3v30-block}
Write
\begin{equation}
 \mathbb H_\Phi=
 \begin{pmatrix}
 H_{gg}&H_{g\varphi}\\
 H_{\varphi g}&H_{\varphi\varphi}
 \end{pmatrix},
 \qquad
 \mathcal R_\Phi X=
 \binom{\mathcal L_Xg}{\mathcal L_X\varphi}.
\label{eq:v3v30-H-block}
\end{equation}
At a full solution,
\begin{align}
 H_{gg}\mathcal L_Xg
 +H_{g\varphi}\mathcal L_X\varphi&=0,
\label{eq:v3v30-metric-row}\\
 H_{\varphi g}\mathcal L_Xg
 +H_{\varphi\varphi}\mathcal L_X\varphi&=0.
\label{eq:v3v30-matter-row}
\end{align}
Neither diagonal term is required to vanish separately.
\end{corollary}

\begin{proof}
Expand \(\mathbb H_\Phi\mathcal R_\Phi X=0\) by block rows.
\end{proof}

\begin{counterexample}[The metric block is not the coupled Jacobi
complex]
\label{cex:v3v30-metric-only}
Suppose
\(H_{g\varphi}\mathcal L_X\varphi\ne0\) for some background matter field
and vector field \(X\).  Then
\eqref{eq:v3v30-metric-row} forces
\[
 H_{gg}\mathcal L_Xg
 =-H_{g\varphi}\mathcal L_X\varphi\ne0.
\]
Thus replacing the full gauge vector
\((\mathcal L_Xg,\mathcal L_X\varphi)\) by its metric component produces a
false kernel statement even on shell.
\end{counterexample}

\subsection{Concrete Einstein--scalar identity}
\label{sec:v3v30-scalar}

The complete-packet point can be seen without abstract notation.  Consider
the Euclidean action
\begin{equation}
 S[g,\phi]
 =\int_M\sqrt g\left[
 \frac1{2\kappa}(R-2\Lambda)
 +\frac12|d\phi|_g^2+V(\phi)
 \right]dx
 +S_{\rm GHY},
\label{eq:v3v30-scalar-action}
\end{equation}
with compactly supported variations in the calculation below.  Define
Euler--Lagrange tensors by
\begin{equation}
 \delta S
 =\int_M\sqrt g\left(
 E_g^{\mu\nu}\delta g_{\mu\nu}
 +E_\phi\delta\phi
 \right)dx.
\label{eq:v3v30-scalar-E}
\end{equation}

\begin{proposition}[Matter-completed contracted Bianchi identity]
\label{prop:v3v30-scalar-Bianchi}
The Euler--Lagrange expressions in
\eqref{eq:v3v30-scalar-E} obey
\begin{equation}
 2\nabla_\mu E_g^{\mu\nu}
 =E_\phi\nabla^\nu\phi.
\label{eq:v3v30-scalar-Bianchi}
\end{equation}
\end{proposition}

\begin{proof}
Under an infinitesimal diffeomorphism,
\[
 \delta_Xg_{\mu\nu}=2\nabla_{(\mu}X_{\nu)},
 \qquad
 \delta_X\phi=X^\nu\nabla_\nu\phi.
\]
Insert these in \eqref{eq:v3v30-scalar-E} and integrate the metric term by
parts:
\[
 0=\delta_XS
 =\int_M\sqrt g\,X_\nu
 \left(-2\nabla_\mu E_g^{\mu\nu}
 +E_\phi\nabla^\nu\phi\right)dx.
\]
Arbitrariness of compactly supported \(X\) proves the identity.
\end{proof}

\begin{corollary}[Why the matter equation is mandatory]
\label{cor:v3v30-matter-mandatory}
If \(E_\phi\ne0\) and \(d\phi\ne0\), then the metric Euler tensor is not
divergence free by itself.  At a coupled solution,
\(E_g=E_\phi=0\), and differentiating
\eqref{eq:v3v30-scalar-Bianchi} yields a linearized identity containing
both the metric and scalar Hessian blocks.
\end{corollary}

\begin{proof}
The first statement is immediate from
\eqref{eq:v3v30-scalar-Bianchi}.  Differentiating its two sides uses the
variations of \(E_g,E_\phi,\nabla\), and \(d\phi\), which are precisely the
coupled Hessian and lower background terms.  At a solution all terms
proportional to the undifferentiated Euler tensors vanish, leaving the
linearized coupled identity.
\end{proof}

\begin{assessment}[Standard Model extension]
\label{ass:v3v30-SM-extension}
For the Standard Model, the right side of the metric divergence identity
is the sum of each matter Euler expression contracted with its
diffeomorphism generator.  Gauge connections also carry their internal
gauge Noether identity, and spinors require the spin lift of the Lie
derivative.  The abstract identity
\eqref{eq:v3v30-Noether} includes these terms, but V30 has not yet written
their complete matrix Hessian.
\end{assessment}

\subsection{Gauge fixing and the Einstein--matter ghost intertwining}
\label{sec:v3v30-gauge-fixing}

Choose field-space metrics and let
\(\mathcal R_\Phi^*\) be the corresponding gauge-fixing operator.  Define
the ghost and gauge-fixed coupled Jacobi operators
\begin{equation}
 \mathbb L_0=\mathcal R_\Phi^*\mathcal R_\Phi,
 \qquad
 \mathbb L_1=\mathbb H_\Phi
 +\mathcal R_\Phi\mathcal R_\Phi^*.
\label{eq:v3v30-L01}
\end{equation}

\begin{theorem}[On-shell diffeomorphism intertwining]
\label{thm:v3v30-intertwining}
At a full Einstein--matter solution,
\begin{align}
 \mathbb L_1\mathcal R_\Phi
 &=\mathcal R_\Phi\mathbb L_0,
\label{eq:v3v30-L-intertwine}\\
 \mathcal R_\Phi^*\mathbb L_1
 &=\mathbb L_0\mathcal R_\Phi^*.
\label{eq:v3v30-L-adjoint}
\end{align}
Off shell,
\begin{equation}
 \mathbb L_1\mathcal R_\Phi
 -\mathcal R_\Phi\mathbb L_0
 =\mathbb H_\Phi\mathcal R_\Phi,
\label{eq:v3v30-offshell-defect}
\end{equation}
and the right side is linear in the complete Euler packet and its natural
diffeomorphism action.
\end{theorem}

\begin{proof}
Expand
\[
 \mathbb L_1\mathcal R_\Phi
 =\mathbb H_\Phi\mathcal R_\Phi
 +\mathcal R_\Phi\mathcal R_\Phi^*\mathcal R_\Phi.
\]
The second term is \(\mathcal R_\Phi\mathbb L_0\), proving
\eqref{eq:v3v30-offshell-defect}.  At a solution,
\eqref{eq:v3v30-right-Jacobi} removes the first term and proves
\eqref{eq:v3v30-L-intertwine}.  Taking adjoints gives
\eqref{eq:v3v30-L-adjoint}.  Naturality of the Euler expression, or the
differentiated Noether identity, shows that the off-shell term is
proportional to \(\mathcal E(\Phi)\).
\end{proof}

\begin{corollary}[Conditional transfer of the local Neumann quartet]
\label{cor:v3v30-conditional-quartet}
Suppose the gauge-fixed Euclidean operators admit self-adjoint domains,
the ghost operator is positive after removal of stabilizers, and a common
positive shift makes the required reservoir response bounded.  Then the
V28 Neumann quartet, with
\[
 A_0=\mathbb L_0+\mu^2,
\]
locally realizes
\[
 I+\rho(\mathbb L_0+\mu^2)^{-1/2}
\]
in the complete diffeomorphism gauge fermion.  At a full on-shell
background and a common cutoff, the longitudinal metric--matter,
multiplier, ghost, and reservoir-quartet relative determinant cancels
mode by mode.
\end{corollary}

\begin{proof}
The construction of V28 uses only a positive ghost energy form, nilpotent
nonminimal doublets, and the intertwining between gauge and ghost
operators.  These are supplied by the hypotheses and
Theorem~\ref{thm:v3v30-intertwining}.  The finite-dimensional determinant
calculation is Theorem~\ref{thm:v3v27-cutoff-cancel}.
\end{proof}

\begin{firewall}[Euclidean gravity is not yet a positive Laplace problem]
\label{fw:v3v30-conformal}
Corollary~\ref{cor:v3v30-conditional-quartet} is conditional because the
Euclidean Einstein Hessian has the conformal-factor problem, and mixed
matter blocks or boundary conditions may obstruct a positive
self-adjoint realization.  Adding a scalar \(\mu^2\) does not repair a
wrong-sign principal direction.  A contour or physical reduction must be
derived before applying positive functional calculus to
\(\mathbb L_1\).
\end{firewall}

\subsection{Boundary gauge group and edge sector}
\label{sec:v3v30-boundary-gauge}

\begin{proposition}[Two inequivalent boundary diffeomorphism choices]
\label{prop:v3v30-boundary-choices}
Let \(\partial M=N\).  The following define different quantum problems.
\begin{enumerate}
\item If \(X|_N=0\), diffeomorphisms generated by \(X\) are pure gauge at
the boundary; the ghost receives a boundary-vanishing domain and no
boundary charge occurs in the Noether identity.
\item If \(X|_N\ne0\) is admitted, its action changes boundary data or
generates a boundary symmetry.  The variational identity contains the
Brown--York and matter boundary currents, and the associated modes require
edge variables, charges, or a quotient by an explicitly specified
boundary group.
\end{enumerate}
The two determinant domains are not interchangeable.
\end{proposition}

\begin{proof}
Apply the diffeomorphism variation without discarding the integration-by-
parts boundary term.  It vanishes in the first case.  In the second, the
metric part is the boundary stress paired with the induced metric
variation, while matter supplies its boundary symplectic flux.  These
terms are exactly the Hamiltonian generators of the admitted boundary
transformation.  Hence the domain and quotient data differ.
\end{proof}

\begin{firewall}[The Calderon boundary complex needs a declared gauge
group]
\label{fw:v3v30-declared-group}
The full Calderon complex of V24--V25 describes boundary traces of decaying
bulk solutions.  Before identifying its exact directions with
diffeomorphism ghosts, the programme must state which boundary vector
fields are quotiented and which are physical symmetries.  Otherwise a
Brown--York charge can be incorrectly cancelled as a ghost mode.
\end{firewall}

\subsection{V30 conclusion and V31 target}
\label{sec:v3v30-conclusion}

The V30 audit directs the programme to the complete Einstein--matter
Euler packet.  Diffeomorphism invariance gives
\(\mathcal R_\Phi^*\mathcal E(\Phi)=0\).  Differentiation proves that the
full coupled Hessian annihilates the full diffeomorphism generator at a
solution.  The metric Hessian need not annihilate
\(\mathcal L_Xg\) separately; the matter cross-Hessian is a mandatory part
of the identity.  Adding the gauge-fixing square yields exact on-shell
intertwining between the coupled fluctuation operator and the
diffeomorphism ghost operator, in direct analogy with the correctly
derived Yang--Mills Jacobi complex.

The next bottleneck is the conditional word in
Corollary~\ref{cor:v3v30-conditional-quartet}.  V31 should compute the
principal symbol of the Euclidean Einstein--matter Hessian in de Donder
gauge, isolate its trace and tracefree directions, and state an explicit
contour that makes the finite-cutoff Gaussian well defined.  It must also
keep the boundary gauge-group choice visible.  Only after that calculation
may the positive Neumann reservoir and modified determinant machinery be
applied to the gravitational longitudinal sector.

\section{V31: de Donder gravitational principal form and the
finite-cutoff conformal contour}
\label{sec:v3v31}

\subsection{Scope}
\label{sec:v3v31-scope}

V30 reduced the diffeomorphism quartet to a positivity and domain problem
for the gauge-fixed Einstein--matter Hessian.  V31 treats the bulk
gravitational principal form.  It proves that de Donder gauge makes the
metric Hessian Laplace type but leaves one negative conformal direction.
Rotating that trace variable gives a positive principal Gaussian and a
definite finite-cutoff phase.  This is a contour prescription, not a proof
of strong ellipticity for a chosen gravitational boundary condition and
not a construction of the chiral matter determinant.

Use the Euclidean Einstein action convention
\begin{equation}
 S_E[g]=-\frac1{2\kappa}\int_M(R-2\Lambda)\,d\operatorname{vol}_g
 +S_{\rm GHY},
\qquad \kappa>0.
\label{eq:v3v31-SE}
\end{equation}
Let \(h_{\mu\nu}\) be a real symmetric metric fluctuation and
\(h=g^{\mu\nu}h_{\mu\nu}\).  The de Donder gauge function is
\begin{equation}
 C_\nu(h)=\nabla^\mu h_{\mu\nu}-\frac12\nabla_\nu h.
\label{eq:v3v31-C}
\end{equation}
The gauge-fixing coefficient is chosen to cancel the nonminimal
second-derivative terms in the Einstein Hessian.

\subsection{Ghost principal operator}
\label{sec:v3v31-ghost}

For a vector field \(X\), the metric gauge generator is
\begin{equation}
 (\mathcal R_gX)_{\mu\nu}
 =\nabla_\mu X_\nu+\nabla_\nu X_\mu.
\label{eq:v3v31-Rg}
\end{equation}

\begin{proposition}[de Donder Faddeev--Popov operator]
\label{prop:v3v31-ghost}
The variation of the gauge function along a diffeomorphism is
\begin{equation}
 C_\nu(\mathcal R_gX)
 =\nabla^\mu\nabla_\mu X_\nu+R_{\nu\mu}X^\mu.
\label{eq:v3v31-C-R}
\end{equation}
Thus, with positive Euclidean principal sign, the ghost operator may be
taken as
\begin{equation}
 (L_{\rm gh}X)_\nu
 =-\nabla^\mu\nabla_\mu X_\nu-R_{\nu\mu}X^\mu,
\label{eq:v3v31-Lghost}
\end{equation}
whose principal symbol is
\begin{equation}
 \sigma_2(L_{\rm gh})(x,\xi)=|\xi|_g^2I_{T_x^*M}.
\label{eq:v3v31-ghost-symbol}
\end{equation}
\end{proposition}

\begin{proof}
Insert \eqref{eq:v3v31-Rg} in \eqref{eq:v3v31-C}:
\[
 C_\nu(\mathcal R_gX)
 =\nabla^\mu\nabla_\mu X_\nu
 +\nabla^\mu\nabla_\nu X_\mu-\nabla_\nu\nabla^\mu X_\mu.
\]
The last two terms differ by the Ricci commutator
\(R_{\nu\mu}X^\mu\), proving \eqref{eq:v3v31-C-R}.  Negating gives the
positive Laplace principal convention in \eqref{eq:v3v31-Lghost}; curvature
is zero order and does not change \eqref{eq:v3v31-ghost-symbol}.
\end{proof}

\begin{firewall}[Ghost positivity is spectral, not principal]
\label{fw:v3v31-ghost-positive}
Equation~\eqref{eq:v3v31-ghost-symbol} proves ellipticity but not
nonnegativity of \(L_{\rm gh}\).  Ricci curvature and the boundary domain
can create a finite negative spectrum, while Killing fields create a
kernel.  The Neumann response of V28 requires removal of stabilizers and a
positive shift larger than the negative spectral bottom.
\end{firewall}

\subsection{Gauge-fixed metric principal form}
\label{sec:v3v31-metric-principal}

Introduce the DeWitt fibre form in four dimensions,
\begin{equation}
 \mathcal G(h,k)
 =h_{\mu\nu}k^{\mu\nu}-\frac12hk.
\label{eq:v3v31-DeWitt}
\end{equation}

\begin{theorem}[Laplace-type de Donder principal Hessian]
\label{thm:v3v31-principal-Hessian}
After adding the de Donder gauge-fixing square with the matching
coefficient, the second-order bulk part of the quadratic metric action is,
up to the overall positive factor \(1/(4\kappa)\),
\begin{equation}
 Q_{\rm prin}(h)
 =\int_M
 \left(
 \nabla_\alpha h_{\mu\nu}\nabla^\alpha h^{\mu\nu}
 -\frac12\nabla_\alpha h\nabla^\alpha h
 \right)d\operatorname{vol}_g.
\label{eq:v3v31-Q-prin}
\end{equation}
Equivalently, the gauge-fixed Hessian has scalar Laplace principal symbol
\begin{equation}
 \sigma_2(L_{\rm grav}^{\rm gf})(x,\xi)
 =|\xi|_g^2\mathcal G
\label{eq:v3v31-grav-symbol}
\end{equation}
on symmetric tensors.  Curvature, cosmological, and on-shell
Einstein--matter mixing terms are lower order for minimally coupled
second-order bosonic matter.
\end{theorem}

\begin{proof}
The principal variation of the Ricci tensor is
\[
 (D\operatorname{Ric}\,h)_{\mu\nu}
 =\frac12\left(
 -\nabla^2h_{\mu\nu}
 \nabla_\mu\nabla^\alpha h_{\alpha\nu}
 \nabla_\nu\nabla^\alpha h_{\alpha\mu}
 -\nabla_\mu\nabla_\nu h
 \right)
 +\text{lower order}.
\]
The scalar-curvature variation supplies the trace reversal of this
expression.  The quadratic gauge term built from
\eqref{eq:v3v31-C} cancels the mixed
\(\nabla_\mu\nabla^\alpha h_{\alpha\nu}\) and
\(\nabla_\mu\nabla_\nu h\) principal terms.  Integrating the remaining
rough Laplacian by parts yields
\eqref{eq:v3v31-Q-prin}.  The coefficient of the trace is exactly the
DeWitt form \eqref{eq:v3v31-DeWitt}, proving
\eqref{eq:v3v31-grav-symbol}.
\end{proof}

\begin{firewall}[Nonminimal curvature couplings require a new matrix
symbol]
\label{fw:v3v31-nonminimal-matter}
A term such as \(\zeta R|\phi|^2\) has metric--scalar second-derivative
mixing and is not covered by the last sentence of
Theorem~\ref{thm:v3v31-principal-Hessian}.  Such counterterms are generally
allowed by renormalization.  Their coupled principal matrix must be
diagonalized before the conformal contour is extended to the full
Einstein--Higgs block.
\end{firewall}

\subsection{Trace--tracefree inertia}
\label{sec:v3v31-trace}

Decompose
\begin{equation}
 h_{\mu\nu}=h_{\mu\nu}^{\rm TF}+\frac14g_{\mu\nu}\tau,
\qquad
 g^{\mu\nu}h_{\mu\nu}^{\rm TF}=0,
\qquad
 \tau=h.
\label{eq:v3v31-decomposition}
\end{equation}

\begin{proposition}[One conformal negative fibre direction]
\label{prop:v3v31-conformal-sign}
The DeWitt form and principal quadratic form split as
\begin{align}
 \mathcal G(h,h)
 &=|h^{\rm TF}|^2-\frac14\tau^2,
\label{eq:v3v31-G-split}\\
 Q_{\rm prin}(h)
 &=\int_M\left(
 |\nabla h^{\rm TF}|^2-\frac14|\nabla\tau|^2
 \right)d\operatorname{vol}_g.
\label{eq:v3v31-Q-split}
\end{align}
Thus the tracefree principal directions are positive and the scalar trace
direction is negative.
\end{proposition}

\begin{proof}
Orthogonality and \(g_{\mu\nu}g^{\mu\nu}=4\) give
\[
 |h|^2=|h^{\rm TF}|^2+\frac14\tau^2.
\]
Insert this and \(h=\tau\) into
\eqref{eq:v3v31-DeWitt}; this gives
\eqref{eq:v3v31-G-split}.  Metric compatibility makes the same
orthogonal decomposition after one covariant derivative, proving
\eqref{eq:v3v31-Q-split}.
\end{proof}

\begin{counterexample}[A positive mass shift does not repair the
conformal principal sign]
\label{cex:v3v31-mass-shift}
For rapidly oscillating trace fluctuations
\(\tau_k=\chi e^{ik\cdot x}\) supported in one coordinate chart, the
negative term in \eqref{eq:v3v31-Q-split} is of order
\(-|k|^2\|\tau_k\|^2\).  Any fixed positive zero-order mass contributes
only \(O(\|\tau_k\|^2)\).  Hence the quadratic form remains unbounded below
as \(|k|\to\infty\).
\end{counterexample}

\subsection{Finite-cutoff conformal rotation}
\label{sec:v3v31-contour}

Let a common elliptic cutoff reduce the real trace sector to
\(\mathbb R^{N_\Lambda}\).  After separating lower-order mixing by a real
linear transformation with recorded determinant, write its trace quadratic
form as
\begin{equation}
 -\frac12\langle\tau,K_{\rm tr}\tau\rangle,
\label{eq:v3v31-negative-trace}
\end{equation}
where \(K_{\rm tr}\) has positive principal symbol.  At fixed cutoff add,
if necessary, a positive scalar shift so that
\(K_{\rm tr}>0\).

\begin{theorem}[Contour-rotated trace Gaussian]
\label{thm:v3v31-conformal-contour}
Orient each trace-mode contour from \(-i\infty\) to \(+i\infty\), set
\begin{equation}
 \tau=i\sigma,\qquad \sigma\in\mathbb R^{N_\Lambda},
\label{eq:v3v31-rotation}
\end{equation}
and use the product orientation.  Then
\begin{equation}
 \int_{(i\mathbb R)^{N_\Lambda}}
 \exp\left(
 +\frac12\langle\tau,K_{\rm tr}\tau\rangle
 \right)d\tau
 =
 i^{N_\Lambda}(2\pi)^{N_\Lambda/2}
 (\det K_{\rm tr})^{-1/2},
\label{eq:v3v31-trace-Gaussian}
\end{equation}
where the determinant square root is the positive one.  The rotated
principal quadratic form is positive.
\end{theorem}

\begin{proof}
Under \(\tau=i\sigma\),
\[
 \frac12\langle\tau,K_{\rm tr}\tau\rangle
 =-\frac12\langle\sigma,K_{\rm tr}\sigma\rangle
\]
and the product measure is \(d\tau=i^{N_\Lambda}d\sigma\).
The standard positive real Gaussian gives
\((2\pi)^{N_\Lambda/2}(\det K_{\rm tr})^{-1/2}\), proving the formula.
\end{proof}

\begin{corollary}[The conformal phase is cutoff data]
\label{cor:v3v31-conformal-phase}
The trace rotation contributes
\begin{equation}
 \exp\left(\frac{i\pi}{2}N_\Lambda\right).
\label{eq:v3v31-phase}
\end{equation}
It cannot be removed from a finite-cutoff comparison unless the same trace
cutoff and contour occur in numerator and denominator.
\end{corollary}

\begin{proof}
This is \(i^{N_\Lambda}\) in
\eqref{eq:v3v31-trace-Gaussian}.
\end{proof}

\begin{firewall}[Contour rotation is not a positivity theorem for the
unrotated action]
\label{fw:v3v31-contour-status}
Theorem~\ref{thm:v3v31-conformal-contour} defines a finite-dimensional
integration cycle on which the principal trace Gaussian converges.  It does
not show that the original real Euclidean Einstein action is bounded below,
and it does not choose among globally inequivalent Picard--Lefschetz cycles
when lower-order saddles are present.
\end{firewall}

\subsection{Compatibility with the diffeomorphism quartet}
\label{sec:v3v31-quartet}

\begin{proposition}[The contour does not alter the on-shell Jacobi
identity]
\label{prop:v3v31-Jacobi-contour}
The algebraic identities
\eqref{eq:v3v30-left-Jacobi}--\eqref{eq:v3v30-L-intertwine} hold before a
Gaussian contour is selected.  At finite cutoff, a complex-linear extension
of the coupled Hessian and gauge generator preserves those identities
after the trace variable is rotated.
\end{proposition}

\begin{proof}
The identities follow from differentiating diffeomorphism invariance and
are polynomial linear-operator equalities.  Complexification preserves
such equalities.  Restricting the complexified trace coordinate to
\(i\mathbb R^{N_\Lambda}\) changes the integration cycle but not the
operator compositions.
\end{proof}

\begin{assessment}[What can now use positive functional calculus]
\label{ass:v3v31-positive-calculus}
After stabilizers are removed, the ghost operator can be shifted to a
positive Laplace-type operator.  After trace rotation, the metric principal
Gaussian has positive tracefree and trace symbols; any remaining finite
negative spectrum can likewise be isolated or shifted at a fixed
background.  Under a self-adjoint strongly elliptic boundary realization,
these are the hypotheses needed for the local Neumann quartet of V28 and
the fourth modified determinant of V21.  The boundary realization remains
the unresolved qualifier.
\end{assessment}

\subsection{Boundary gate}
\label{sec:v3v31-boundary}

\begin{firewall}[Bulk Laplace type does not imply strong boundary
ellipticity]
\label{fw:v3v31-boundary-ellipticity}
The principal symbols
\eqref{eq:v3v31-ghost-symbol} and
\eqref{eq:v3v31-grav-symbol} are bulk ellipticity statements.  A local
metric boundary condition combined with de Donder gauge can still fail the
complementing condition at glancing covectors, as the harmonic boundary
analysis in V15 already warns.  No gravitational determinant is defined
until the ghost and metric domains form the same on-shell Jacobi complex
and pass the boundary symbol test.
\end{firewall}

\begin{proposition}[Boundary-vanishing diffeomorphisms preserve fixed
induced metric data]
\label{prop:v3v31-boundary-vanishing}
If the induced boundary metric is fixed and the gauge vector field
satisfies \(X|_N=0\), then the tangential--tangential part of
\(\mathcal L_Xg\) vanishes at \(N\).  Normal derivatives of \(X\) still
enter mixed and normal metric components and must be paired with the
gauge-fixing boundary conditions.
\end{proposition}

\begin{proof}
For tangential vectors \(Y,Z\) on \(N\),
\[
 (\mathcal L_Xg)(Y,Z)
 =g(\nabla_YX,Z)+g(Y,\nabla_ZX).
\]
The tangential derivative of the boundary restriction \(X|_N=0\) is zero,
so the expression vanishes.  Components containing the normal derivative
of \(X\) are not fixed by its trace, proving the second statement.
\end{proof}

\subsection{V31 conclusion and V32 target}
\label{sec:v3v31-conclusion}

V31 makes the Euclidean gravitational sign gate explicit.  De Donder gauge
produces a Laplace-type bulk Hessian with DeWitt principal form
\(|h|^2-h^2/2\).  In four dimensions this is positive on tracefree metric
fluctuations and negative on the scalar trace.  A positive mass shift
cannot change that high-frequency inertia.  At finite cutoff the contour
\(\tau=i\sigma\) gives a convergent positive trace Gaussian and the
recorded phase \(i^{N_\Lambda}\).  The diffeomorphism Jacobi identity
survives complexification.

The remaining obstruction is the boundary domain.  V32 should assemble
the principal half-space boundary symbol for the de Donder metric operator,
the vector ghost, fixed induced metric data, and the additional gauge
conditions on mixed and normal components.  It must calculate the stable
normal roots and test the complementing determinant, then compare any
glancing kernel with the V15 flat harmonic block.  If the local conditions
fail, the programme should attach the V11--V16 dynamic boundary reservoir
to the complete metric--ghost complex rather than modifying one component
at a time.

\section{V32: flat de Donder boundary symbol, its scalar kernel, and the
order-one repair gate}
\label{sec:v3v32}

\subsection{The boundary problem being tested}
\label{sec:v3v32-setup}

V31 proved bulk Laplace type after de Donder gauge but left the
complementing condition open.  V32 freezes the coefficients at a flat
boundary point and tests the most direct local boundary system:
\begin{enumerate}
\item fix the six components of the induced metric fluctuation,
      \(h_{ab}|_N=0\);
\item impose the four de Donder gauge functions,
      \(C_\mu(h)|_N=0\);
\item impose Dirichlet data on the vector ghost.
\end{enumerate}
Here \(a,b=1,2,3\) are tangential indices and \(r\) is the inward normal
coordinate.  The ten metric conditions have the correct numerical count,
but the calculation below shows that their principal symbols are not
independent on the stable root.

Work on the Euclidean half-space
\[
 X=\{r\ge0\}\subset\mathbb R^4,\qquad
 g=dr^2+\delta_{ab}dx^adx^b.
\]
For the trace reversal
\begin{equation}
 \bar h_{\mu\nu}
 =h_{\mu\nu}-\frac12\delta_{\mu\nu}h,
\label{eq:v3v32-trace-reverse}
\end{equation}
the frozen gauge function is
\begin{equation}
 C_\mu(h)=\partial^\nu\bar h_{\mu\nu}.
\label{eq:v3v32-C}
\end{equation}
The gauge-fixed metric and ghost principal operators are copies of the
positive scalar Laplacian.

\subsection{Stable normal root}
\label{sec:v3v32-root}

Let \(\xi\in\mathbb R^3\setminus\{0\}\) be a tangential covector and let
\(z\) be a complex spectral parameter.  A stable mode has the form
\begin{equation}
 h_{\mu\nu}(r,x)
 =H_{\mu\nu}e^{i\xi\cdot x-pr},
\qquad
 p=p(\xi,z)=\sqrt{|\xi|^2-z},
\qquad
 \operatorname{Re}p>0.
\label{eq:v3v32-mode}
\end{equation}
It solves
\((-\partial_r^2-\Delta_x-z)h=0\).

\begin{lemma}[Multiplicity of the stable root]
\label{lem:v3v32-root}
For every \((\xi,z)\) for which the branch in
\eqref{eq:v3v32-mode} has positive real part, the stable space of metric
amplitudes has dimension ten and the stable ghost space has dimension
four.  At \(z=0\),
\begin{equation}
 p=|\xi|.
\label{eq:v3v32-zero-root}
\end{equation}
\end{lemma}

\begin{proof}
The gauge-fixed principal symbols are scalar multiples of the identity on
the ten-dimensional symmetric-tensor fibre and the four-dimensional
covector fibre.  Each component has the unique decaying root with
\(\operatorname{Re}p>0\).  The last formula follows from the chosen
positive branch.
\end{proof}

\subsection{Reduction by the six induced-metric conditions}
\label{sec:v3v32-reduction}

The conditions \(H_{ab}=0\) leave the four amplitudes
\begin{equation}
 v=(H_{r1},H_{r2},H_{r3},H_{rr})^{\mathsf T}.
\label{eq:v3v32-v}
\end{equation}
On this subspace,
\begin{align}
 h&=H_{rr},
\label{eq:v3v32-trace}\\
 \bar H_{rr}&=\frac12H_{rr},
\qquad
 \bar H_{ra}=H_{ra},
\qquad
 \bar H_{ab}=-\frac12\delta_{ab}H_{rr}.
\label{eq:v3v32-bar-components}
\end{align}

\begin{proposition}[The reduced de Donder boundary matrix]
\label{prop:v3v32-boundary-matrix}
On the stable mode \eqref{eq:v3v32-mode}, the four gauge conditions are
\begin{align}
 C_r&=-\frac p2H_{rr}+i\xi^aH_{ra},
\label{eq:v3v32-Cr}\\
 C_a&=-pH_{ra}-\frac i2\xi_aH_{rr}.
\label{eq:v3v32-Ca}
\end{align}
In the ordering \eqref{eq:v3v32-v}, their matrix is
\begin{equation}
 \mathcal B(\xi,z)=
 \begin{pmatrix}
 -pI_3&-\frac i2\xi\\
 i\xi^{\mathsf T}&-\frac p2
 \end{pmatrix},
\label{eq:v3v32-B}
\end{equation}
where the first three rows are \(C_a\) and the last row is \(C_r\).
\end{proposition}

\begin{proof}
On the mode,
\(\partial_r=-p\) and \(\partial_a=i\xi_a\).  Use
\eqref{eq:v3v32-bar-components} in
\(C_r=\partial_r\bar h_{rr}+\partial^a\bar h_{ra}\) and
\(C_a=\partial_r\bar h_{ar}+\partial^b\bar h_{ab}\).
This gives \eqref{eq:v3v32-Cr}--\eqref{eq:v3v32-Ca}, and their matrix is
\eqref{eq:v3v32-B}.
\end{proof}

\subsection{Exact complementing determinant}
\label{sec:v3v32-determinant}

\begin{theorem}[The boundary determinant is proportional to the spectral
parameter]
\label{thm:v3v32-determinant}
The reduced boundary matrix obeys
\begin{equation}
 \det\mathcal B(\xi,z)
 =-\frac12zp^2.
\label{eq:v3v32-det}
\end{equation}
Consequently the ordinary zero-parameter complementing condition fails
for every \(\xi\ne0\).
\end{theorem}

\begin{proof}
The upper-left block \(-pI_3\) is invertible on the stable branch.  Its
Schur complement is
\[
 -\frac p2
 -i\xi^{\mathsf T}(-pI_3)^{-1}
 \left(-\frac i2\xi\right)
 =-\frac p2+\frac{|\xi|^2}{2p}
 =\frac{|\xi|^2-p^2}{2p}
 =\frac z{2p}.
\]
Multiplying by \(\det(-pI_3)=-p^3\) gives
\eqref{eq:v3v32-det}.  At \(z=0\), the determinant vanishes although
\(\xi\ne0\), which is exactly failure of the complementing condition.
\end{proof}

\begin{corollary}[Explicit one-dimensional stable kernel]
\label{cor:v3v32-kernel}
At \(z=0\), the kernel consists of
\begin{equation}
 H_{rr}=t,\qquad
 H_{ra}=-\frac{i\xi_a}{2|\xi|}t,\qquad
 H_{ab}=0,
\qquad t\in\mathbb C.
\label{eq:v3v32-kernel}
\end{equation}
It is one dimensional.
\end{corollary}

\begin{proof}
The three equations \eqref{eq:v3v32-Ca} give the displayed \(H_{ra}\).
Substitution in \eqref{eq:v3v32-Cr} gives
\[
 -\frac{|\xi|}{2}t
 +\frac{|\xi|^2}{2|\xi|}t=0.
\]
Conversely every kernel vector has this form because the
\(C_a\) block is invertible in \(H_{ra}\).
\end{proof}

\begin{corollary}[Collapse of the parameter margin]
\label{cor:v3v32-margin}
For \(z\ne0\), the four reduced conditions are independent, but their
Schur factor is \(z/(2p)\).  Hence the inverse boundary matrix cannot have
a uniform bound as \(z\to0\) with fixed \(\xi\ne0\).
\end{corollary}

\begin{proof}
This is the Schur-complement calculation in the proof of
Theorem~\ref{thm:v3v32-determinant}.
\end{proof}

\subsection{Geometric identity of the kernel}
\label{sec:v3v32-geometric-kernel}

\begin{proposition}[The kernel is a normal boundary displacement]
\label{prop:v3v32-normal-displacement}
Let
\begin{equation}
 X_r=-\frac{t}{2|\xi|}
 e^{i\xi\cdot x-|\xi|r},
\qquad
 X_a=0.
\label{eq:v3v32-X}
\end{equation}
Then
\begin{equation}
 h_{\mu\nu}=\partial_\mu X_\nu+\partial_\nu X_\mu
\label{eq:v3v32-pure-diffeo}
\end{equation}
has exactly the amplitudes in \eqref{eq:v3v32-kernel}.  Its tangential
metric trace vanishes, but \(X_r|_N\ne0\) when \(t\ne0\).
\end{proposition}

\begin{proof}
Since \(X_a=0\),
\[
 h_{ab}=0,\qquad
 h_{rr}=2\partial_rX_r=t e^{i\xi\cdot x-|\xi|r},
\]
and
\[
 h_{ra}=\partial_aX_r
 =-\frac{i\xi_a}{2|\xi|}t
 e^{i\xi\cdot x-|\xi|r}.
\]
These are the stated amplitudes.  The boundary value of \(X_r\) is
\(-t/(2|\xi|)\).
\end{proof}

\begin{assessment}[Why ghost Dirichlet does not remove this metric mode]
\label{ass:v3v32-ghost-Dirichlet}
The vector ghost Dirichlet symbol is complementing: a decaying ghost mode
\(Xe^{i\xi\cdot x-pr}\) with \(X|_N=0\) has zero amplitude.  The metric
kernel in Proposition~\ref{prop:v3v32-normal-displacement}, however, is
generated by a vector with nonzero normal boundary trace.  It moves the
location of the boundary while preserving its tangential metric to
principal order.  Such a transformation is excluded from the
boundary-vanishing gauge group but its metric image is still admitted by
the ten local metric conditions.  This mismatch is the source of the
kernel.
\end{assessment}

\begin{firewall}[A boundary symmetry cannot be both excluded and
cancelled]
\label{fw:v3v32-edge}
If normal boundary displacements are excluded from the ghost domain, the
mode \eqref{eq:v3v32-kernel} must be fixed by an additional physical
boundary equation or represented by an edge variable.  It cannot be
called a ghost quartet.  If such displacements are instead quotiented,
their boundary generator and Brown--York charge must be included as in
V30.
\end{firewall}

\subsection{Distinction from the V15 regression}
\label{sec:v3v32-v15}

\begin{proposition}[V15 and V32 expose different blocks]
\label{prop:v3v32-v15-distinction}
The V15 flat harmonic junction calculation used six Brown--York-type
normal-derivative equations plus four harmonic constraints and found a
tangential tracefree Neumann glancing subsymbol.  V32 instead fixes the
six tangential metric values and finds a scalar normal-displacement kernel
in the remaining four de Donder conditions.  Neither result implies the
other.
\end{proposition}

\begin{proof}
The V15 kernel test lies in the tangential tracefree sector and its
boundary operator is of Neumann/junction type.  In V32,
\(H_{ab}=0\) identically and the kernel has only \(H_{rr},H_{ra}\).
The two subspaces and the six boundary operators are therefore different.
\end{proof}

\begin{assessment}[Common structural lesson]
\label{ass:v3v32-common-lesson}
Both calculations show that algebraically solving ten normal data at one
frozen point is weaker than a uniform boundary-symbol estimate.  V15 loses
the tracefree Neumann factor at hyperbolic glancing; V32 loses a scalar
Schur factor at the Euclidean zero spectral parameter.
\end{assessment}

\subsection{Minimal order required for a repair}
\label{sec:v3v32-repair}

Consider modifying only the last reduced condition by a scalar boundary
operator \(Z\) acting on \(h_{rr}\):
\begin{equation}
 C_r+Zh_{rr}=0.
\label{eq:v3v32-Z-condition}
\end{equation}

\begin{proposition}[Zero-order additions do not give a homogeneous
high-frequency margin]
\label{prop:v3v32-zero-order}
If \(Z\) has order zero with bounded principal amplitude \(z_0(x,\xi)\),
then at \(z=0\) the repaired Schur factor is \(z_0\), while the other rows
have size \(|\xi|\).  After homogeneous order-one normalization its margin
is \(O(|\xi|^{-1})\) and tends to zero.  An order-one principal contribution
is necessary for a uniform order-one complementing estimate.
\end{proposition}

\begin{proof}
At \(z=0\), the unmodified Schur factor is zero.  The addition contributes
the symbol \(z_0\).  The natural boundary norm for the four first-order
gauge conditions divides the final scalar row by \(|\xi|\); hence the
normalized contribution is \(z_0/|\xi|\).  Boundedness of an order-zero
symbol makes this tend to zero.
\end{proof}

\begin{proposition}[A positive order-one impedance removes the
zero-parameter kernel]
\label{prop:v3v32-order-one}
Let \(Z\) be a positive elliptic order-one boundary operator with principal
symbol
\begin{equation}
 \sigma_1(Z)(x,\xi)=\zeta|\xi|,
\qquad \zeta>0.
\label{eq:v3v32-Z-symbol}
\end{equation}
Then the zero-parameter Schur factor of
\eqref{eq:v3v32-Z-condition} is \(\zeta|\xi|\), and the reduced boundary
matrix is invertible for every \(\xi\ne0\), with a scale-invariant
principal margin depending only on \(\zeta\).
\end{proposition}

\begin{proof}
The original Schur factor is zero at \(z=0\); adding
\eqref{eq:v3v32-Z-symbol} gives \(\zeta|\xi|\).  The upper block
\(-|\xi|I_3\) is invertible, so the Schur criterion proves invertibility.
After dividing every first-order row by \(|\xi|\), the matrix depends only
on \(\widehat\xi=\xi/|\xi|\) and \(\zeta\).  Compactness of the unit sphere
and the nonzero Schur factor give a positive minimum singular value.
\end{proof}

\begin{firewall}[The symbol \(|D_N|\) is not yet a local boundary action]
\label{fw:v3v32-nonlocal-repair}
Proposition~\ref{prop:v3v32-order-one} identifies the necessary principal
order but uses a pseudodifferential impedance.  As in V27--V28, locality
requires a bulk or boundary reservoir whose elimination produces that
symbol.  The reservoir must be attached to the complete metric--ghost
boundary complex; adding \(|D_N|h_{rr}\) by hand does not prove BRST
compatibility or a local SM--Einstein continuum.
\end{firewall}

\subsection{Domain-level BRST gate}
\label{sec:v3v32-domain}

\begin{assessment}[On-shell symbol compatibility versus operator-domain
compatibility]
\label{ass:v3v32-domain}
For a ghost stable mode solving
\((L_{\rm gh}-z)X=0\), Dirichlet data imply \(X=0\), so its metric image is
trivial and the principal boundary chain closes.  For a general ghost in
the Dirichlet operator domain, the trace of
\(C(\mathcal R_gX)= -L_{\rm gh}X\) is not determined by \(X|_N=0\) alone.
A closed BRST operator domain may require a graph boundary condition or an
extended boundary multiplier.  V32 proves the stable-symbol failure
independently of that additional domain issue.
\end{assessment}

\subsection{V32 conclusion and V33 target}
\label{sec:v3v32-conclusion}

V32 computes the complete flat half-space boundary determinant for fixed
induced metric plus de Donder gauge.  After the six Dirichlet conditions,
the four gauge rows have determinant
\[
 \det\mathcal B(\xi,z)=-\frac12zp^2.
\]
At \(z=0\) there is one stable mode,
\(H_{ra}=-i\xi_aH_{rr}/(2|\xi|)\).  It is the metric image of a normal
boundary displacement whose ghost parameter has nonzero boundary trace.
Ghost Dirichlet therefore excludes the generator without excluding its
metric image.  This scalar kernel is distinct from V15's tangential
tracefree Neumann glancing sector.

A bounded zero-order boundary term cannot give a homogeneous margin.
A positive order-one impedance removes the frozen kernel, but writing
\(|D_N|\) directly is nonlocal.  V33 should construct a local dynamic edge
reservoir whose eliminated symbol supplies this order-one scalar
impedance, include its BRST partner for normal diffeomorphisms, and compute
the full parameter-dependent complementing determinant.  The proof must
test both the \(z\to0\) scalar kernel found here and the V15 tracefree
glancing limit before claiming a strongly elliptic domain.

\section{V33: local scalar edge reservoir, parameter ellipticity, and a
principal BRST edge completion}
\label{sec:v3v33}

\subsection{Purpose and limits}
\label{sec:v3v33-purpose}

V32 proved that fixed induced metric plus de Donder gauge loses one scalar
normal-displacement mode and that a homogeneous repair must have order one.
V33 constructs that order from a local half-cylinder energy and computes
the full resolvent-parameter determinant.  The sign is not cosmetic: one
choice moves every new zero to the excluded positive spectral ray, while
the other creates a negative-real zero and fails strong ellipticity.

The second part gives a nilpotent principal-symbol edge completion for the
normal diffeomorphism.  It proves compatibility at the frozen leading
order.  Curved lower symbols, the exact gravitational boundary action, and
the separate V15 tracefree junction glancing problem remain open.

\subsection{A local edge Dirichlet-to-Neumann operator}
\label{sec:v3v33-edge-DtN}

At a frozen boundary point let the edge reservoir occupy
\[
 Y=N\times[0,\infty)_y
\]
with a second inward coordinate \(y\).  For a scalar \(\psi\), set
\begin{equation}
 \mathcal E_{\rm e}(\psi)
 =\frac12\int_0^\infty
 \left(
 \|\partial_y\psi\|_{L^2(N)}^2
 +\langle(-\Delta_N+m_{\rm e}^2)\psi,\psi\rangle
 \right)dy,
\qquad m_{\rm e}\ge0.
\label{eq:v3v33-edge-energy}
\end{equation}
At \(m_{\rm e}=0\), the constant boundary mode is separated from the
homogeneous principal calculation.

\begin{theorem}[Exact local parent of the order-one edge impedance]
\label{thm:v3v33-edge-DtN}
For boundary trace \(u\), the decaying minimizer of
\eqref{eq:v3v33-edge-energy} is
\begin{equation}
 \psi_u(y)
 =e^{-yQ_{\rm e}}u,
\qquad
 Q_{\rm e}=\sqrt{-\Delta_N+m_{\rm e}^2}.
\label{eq:v3v33-edge-Poisson}
\end{equation}
Its energy and outward Dirichlet-to-Neumann map are
\begin{equation}
 \mathcal E_{\rm e}(\psi_u)
 =\frac12\langle Q_{\rm e}u,u\rangle,
\qquad
 -\partial_y\psi_u(0)=Q_{\rm e}u.
\label{eq:v3v33-edge-on-shell}
\end{equation}
The principal symbol of \(Q_{\rm e}\) is \(|\xi|\).
\end{theorem}

\begin{proof}
For a tangential Fourier mode of frequency \(\xi\), put
\[
 q=q(\xi)=\sqrt{|\xi|^2+m_{\rm e}^2}.
\]
The unique decaying solution with trace \(u_\xi\) is
\(e^{-qy}u_\xi\).  Direct integration gives
\[
 \frac12\int_0^\infty
 \left(q^2e^{-2qy}+q^2e^{-2qy}\right)|u_\xi|^2dy
 =\frac q2|u_\xi|^2.
\]
Summing the spectral modes proves
\eqref{eq:v3v33-edge-Poisson}--\eqref{eq:v3v33-edge-on-shell}.
The high-frequency expansion
\(q=|\xi|+O(|\xi|^{-1})\) proves the principal-symbol statement.
\end{proof}

\begin{corollary}[Positive local edge energy]
\label{cor:v3v33-positive-edge}
Imposing the trace relation
\(\psi(0)=\sqrt{\zeta}\,u\), \(\zeta>0\), adds the positive on-shell
quadratic form
\begin{equation}
 \frac{\zeta}{2}\langle Q_{\rm e}u,u\rangle.
\label{eq:v3v33-positive-form}
\end{equation}
\end{corollary}

\begin{proof}
Apply \eqref{eq:v3v33-edge-on-shell} to
\(\sqrt{\zeta}\,u\).
\end{proof}

\subsection{The two possible signs in the gravitational boundary row}
\label{sec:v3v33-signs}

Retain the stable metric mode of V32:
\[
 p=\sqrt{|\xi|^2-z},\qquad\operatorname{Re}p>0.
\]
After the three \(C_a\) rows are eliminated, the unmodified \(C_r\) Schur
factor is \(z/(2p)\).  Consider
\begin{equation}
 C_r\pm\zeta Q_{\rm e}h_{rr}=0.
\label{eq:v3v33-two-signs}
\end{equation}

\begin{proposition}[Exact augmented boundary determinants]
\label{prop:v3v33-augmented-det}
The two reduced boundary matrices have determinants
\begin{equation}
 \det\mathcal B_\pm(\xi,z)
 =(-p)^3\left(
 \frac{z}{2p}\pm\zeta q
 \right),
\qquad
 q=\sqrt{|\xi|^2+m_{\rm e}^2}.
\label{eq:v3v33-det-pm}
\end{equation}
\end{proposition}

\begin{proof}
The upper \(3\times3\) block and its elimination are unchanged from
Theorem~\ref{thm:v3v32-determinant}.  The added scalar row changes only its
Schur complement, from \(z/(2p)\) to
\(z/(2p)\pm\zeta q\).  Multiply by \(\det(-pI_3)=(-p)^3\).
\end{proof}

For a positive self-adjoint bulk operator, use the resolvent cone
\begin{equation}
 \mathfrak C=\mathbb C\setminus[0,\infty).
\label{eq:v3v33-cone}
\end{equation}

\begin{theorem}[The minus row is parameter complementing]
\label{thm:v3v33-good-sign}
For \(\zeta>0\), the boundary matrix
\(\mathcal B_-\) has no zero for
\begin{equation}
 z\in\mathfrak C,\qquad
 (\xi,z)\ne(0,0),\qquad
 \operatorname{Re}p>0.
\label{eq:v3v33-parameter-region}
\end{equation}
At \(z=0\) and \(\xi\ne0\), its scalar Schur factor is
\begin{equation}
 -\zeta\sqrt{|\xi|^2+m_{\rm e}^2},
\label{eq:v3v33-zero-margin}
\end{equation}
so the V32 kernel is removed.
\end{theorem}

\begin{proof}
Suppose the Schur factor of \(\mathcal B_-\) vanishes.  Then
\[
 z=2\zeta q p.
\]
Since \(p^2=|\xi|^2-z\), the stable root obeys
\begin{equation}
 p^2+2\zeta q p-|\xi|^2=0.
\label{eq:v3v33-good-quadratic}
\end{equation}
This quadratic has real roots
\[
 p=-\zeta q\pm
 \sqrt{\zeta^2q^2+|\xi|^2}.
\]
The root with positive real part is real and nonnegative; it is strictly
positive when \(\xi\ne0\).  Hence
\(z=2\zeta q p\) is positive real, contrary to
\(z\in\mathfrak C\).  If \(\xi=0\), the positive-real-part condition and
\eqref{eq:v3v33-good-quadratic} are incompatible unless \(p=0\), which is
excluded by \((\xi,z)\ne(0,0)\).  Thus no zero occurs.  Setting \(z=0\)
in the Schur factor gives \eqref{eq:v3v33-zero-margin}.
\end{proof}

\begin{corollary}[Uniformity on closed parameter subsectors]
\label{cor:v3v33-uniform}
In the homogeneous case \(m_{\rm e}=0\), normalize
\(|\xi|^2+|z|=1\).  On every closed conic subsector of
\(\mathfrak C\) separated from the positive real ray, the smallest singular
value of the row-normalized matrix \(\mathcal B_-\) has a positive lower
bound depending on the subsector and \(\zeta\).
\end{corollary}

\begin{proof}
After homogeneous row normalization, the matrix is continuous on the
compact normalized parameter set.  Theorem~\ref{thm:v3v33-good-sign}
shows that it is invertible there.  Its smallest singular value therefore
has a positive minimum.
\end{proof}

\begin{theorem}[The plus row creates a forbidden negative-real zero]
\label{thm:v3v33-bad-sign}
For every \(\xi\ne0\), \(\zeta>0\), and \(q>0\), the plus boundary matrix
\(\mathcal B_+\) has a zero at
\begin{align}
 p_*&=\zeta q+\sqrt{\zeta^2q^2+|\xi|^2}>0,
\label{eq:v3v33-bad-p}\\
 z_*&=-2\zeta q p_*<0.
\label{eq:v3v33-bad-z}
\end{align}
Thus the plus row fails parameter ellipticity inside the resolvent cone.
\end{theorem}

\begin{proof}
The plus Schur factor vanishes when \(z=-2\zeta q p\).  Combining this with
\(p^2=|\xi|^2-z\) gives
\[
 p^2-2\zeta q p-|\xi|^2=0.
\]
Its positive root is \eqref{eq:v3v33-bad-p}; the corresponding
\eqref{eq:v3v33-bad-z} is negative and hence belongs to
\(\mathfrak C\).  Substitution in
\eqref{eq:v3v33-det-pm} proves the zero.
\end{proof}

\begin{firewall}[The action must produce the good canonical orientation]
\label{fw:v3v33-action-sign}
The positive reservoir form \eqref{eq:v3v33-positive-form} contributes
\(+\zeta Q_{\rm e}h_{rr}\) to its canonical boundary momentum equation.
The parameter-elliptic row is therefore obtained when the unperturbed
canonical momentum is \(-C_r\):
\[
 -C_r+\zeta Q_{\rm e}h_{rr}=0,
\]
which is equivalent to the minus row in
\eqref{eq:v3v33-two-signs}.  V33 proves the required orientation.  The
complete gauge-fixed Einstein--GHY boundary variation must still verify
that its scalar canonical row is indeed normalized this way; a sign chosen
after elimination is not evidence.
\end{firewall}

\subsection{Principal normal-diffeomorphism edge pair}
\label{sec:v3v33-BRST-edge}

The V32 kernel is generated by a normal boundary displacement.  A repair
which simply fixes its metric image would not respect a boundary gauge
extension.  Introduce on the edge half-cylinder a bosonic Stückelberg field
\(U(y,x)\) and a fermionic ghost extension \(C(y,x)\), with
\begin{equation}
 sU=C,\qquad sC=0.
\label{eq:v3v33-edge-doublet}
\end{equation}
Define the boundary scalar
\begin{equation}
 \widehat h_{rr}
 =h_{rr}-2\partial_yU|_{y=0}
\label{eq:v3v33-hhat}
\end{equation}
and extend the principal BRST rule by
\begin{equation}
 sh_{rr}=2\partial_yC|_{y=0}.
\label{eq:v3v33-shrr}
\end{equation}

\begin{proposition}[Exact nilpotent edge invariant]
\label{prop:v3v33-edge-invariant}
The rules \eqref{eq:v3v33-edge-doublet}--\eqref{eq:v3v33-shrr} satisfy
\begin{equation}
 s^2U=s^2C=s^2h_{rr}=0,
\qquad
 s\widehat h_{rr}=0.
\label{eq:v3v33-edge-nilpotent}
\end{equation}
If \(C\) is the decaying massless edge extension with trace
\(c_\perp\), then at principal symbol
\begin{equation}
 sh_{rr}=-2|\xi|c_\perp.
\label{eq:v3v33-normal-gauge-symbol}
\end{equation}
This is the normal-diffeomorphism symbol of the V32 kernel.
\end{proposition}

\begin{proof}
Nilpotence follows immediately from \(sC=0\).  Moreover,
\[
 s\widehat h_{rr}
 =2\partial_yC-2\partial_y(sU)=0.
\]
For the decaying massless extension,
\(C(y)=e^{-|\xi|y}c_\perp\), so
\(2\partial_yC(0)=-2|\xi|c_\perp\).  In
Proposition~\ref{prop:v3v32-normal-displacement}, a stable normal vector
with boundary amplitude \(c_\perp\) gives
\(h_{rr}=-2|\xi|c_\perp\), proving the match.
\end{proof}

\begin{corollary}[BRST-invariant local trace attachment]
\label{cor:v3v33-invariant-trace}
Let a physical edge reservoir \(\psi\) be BRST invariant and impose the
local trace relation
\begin{equation}
 \psi(0)=\sqrt{\zeta}\,\widehat h_{rr}.
\label{eq:v3v33-invariant-trace}
\end{equation}
Then the constraint and the local energy
\eqref{eq:v3v33-edge-energy} are BRST invariant under the extended
principal rules.  In the edge gauge \(U|_N=0\), their eliminated scalar
form is \(\zeta\langle h_{rr},Q_{\rm e}h_{rr}\rangle/2\).
\end{corollary}

\begin{proof}
Both sides of \eqref{eq:v3v33-invariant-trace} are BRST invariant by
Proposition~\ref{prop:v3v33-edge-invariant}.  The energy depends only on
\(\psi\).  In the stated gauge,
\(\widehat h_{rr}=h_{rr}\), and
Corollary~\ref{cor:v3v33-positive-edge} gives the eliminated form.
\end{proof}

\begin{firewall}[Principal edge closure is not the curved BRST domain]
\label{fw:v3v33-principal-only}
The rule \eqref{eq:v3v33-shrr} matches the normal diffeomorphism only at
the frozen massless principal symbol, where the metric and edge stable
roots both equal \(|\xi|\).  On a curved collar, the vector ghost
Dirichlet-to-Neumann map has matrix lower symbols involving curvature and
the second fundamental form.  Those terms must replace
\(\partial_yC\) through an actual local coupled extension before an exact
curved BRST operator domain is claimed.
\end{firewall}

\subsection{The V15 tracefree sector is not paid by this scalar edge}
\label{sec:v3v33-v15}

\begin{proposition}[Scalar repair leaves tangential tracefree data
unchanged]
\label{prop:v3v33-not-V15}
The added form \eqref{eq:v3v33-positive-form} and the condition
\eqref{eq:v3v33-two-signs} depend only on \(h_{rr}\).  They vanish
identically on the V15 tangential tracefree test sector
\[
 h_{rr}=h_{ra}=0,\qquad
 h_{ab}\ne0,\qquad \delta^{ab}h_{ab}=0.
\]
Consequently V33 does not repair the V15 Neumann glancing subsymbol.
\end{proposition}

\begin{proof}
Substitution of \(h_{rr}=0\) annihilates the scalar trace attachment and
its boundary variation.  The remaining V15 boundary rows are therefore
unchanged.
\end{proof}

\begin{firewall}[One reservoir cannot be charged to two disjoint
obstructions]
\label{fw:v3v33-two-obstructions}
The V32 scalar displacement and the V15 tracefree junction mode lie in
different invariant fibre sectors.  The scalar edge energy pays only the
first.  A full Einstein boundary construction needs a separate tracefree
dynamic channel or a single matrix reservoir whose positive form contains
both sectors explicitly.
\end{firewall}

\subsection{V33 conclusion and V34 target}
\label{sec:v3v33-conclusion}

V33 gives a local parent for the missing scalar order-one symbol.  A
half-cylinder Dirichlet energy produces
\(Q_{\rm e}=\sqrt{-\Delta_N+m_{\rm e}^2}\).  Appending it with the canonical
row \(-C_r+\zeta Q_{\rm e}h_{rr}=0\) yields the reduced determinant
\[
 (-p)^3\left(\frac{z}{2p}-\zeta q\right),
\]
which has no zero in the resolvent cone
\(\mathbb C\setminus[0,\infty)\).  The opposite sign has an explicit
negative-real zero.  A local edge Stückelberg--ghost pair makes the
trace attachment BRST invariant at the frozen principal level.

The proof leaves two genuine gates.  The curved vector-ghost Calderon map
must replace the frozen edge derivative to close the BRST domain beyond
principal order, and the V15 tangential tracefree glancing block needs its
own dynamic channel.  V34 should construct the matrix edge reservoir on
the direct sum of the scalar normal-displacement line and the tangential
tracefree bundle, derive its positive energy and BRST transformation from
the complete metric--ghost Calderon complex, and test the two diagonal
limits before estimating mixed curvature symbols.

\section{V34: positive matrix edge impedance, mixing margins, and the
Euclidean--Lorentzian signature obstruction}
\label{sec:v3v34}

\subsection{Why a direct-sum construction is necessary}
\label{sec:v3v34-purpose}

V33 proves that the scalar normal-displacement line needs a positive
order-one edge impedance.  V15--V16 show that the tangential tracefree
junction sector also needs an order-one contribution at Lorentzian
glancing.  These facts do not yet imply that one scalar reservoir pays both
obstructions.  V34 forms a genuine matrix reservoir on the direct sum of
the scalar and tracefree fibres, proves its positivity and stability under
mixing, and then tests its Lorentzian principal symbol.

The test exposes a signature obstruction.  The covariant Euclidean
Dirichlet-to-Neumann symbol is \(|\xi|\), but the covariant Lorentzian
symbol is the outgoing square root of the boundary wave symbol and
vanishes at null glancing.  A mass removes an exact finite-frequency zero
but is lower order under homogeneous scaling.  Thus the matrix reservoir
solves the Euclidean V32 problem but does not, by itself, solve the V15
zero-loss hyperbolic problem.

\subsection{The scalar--tracefree boundary bundle}
\label{sec:v3v34-bundle}

On a Riemannian three-boundary \((N,h)\), define
\begin{equation}
 E_{\rm e}
 =E_{\rm s}\oplus E_{\rm tf},
\qquad
 E_{\rm s}=N\times\mathbb R,
\qquad
 E_{\rm tf}=S_0^2T^*N.
\label{eq:v3v34-E}
\end{equation}
The tracefree bundle is kept in its full local rank-five form.  The
frequency-dependent transverse--tracefree projection is not inserted into
the action.

Let \(\mathbb A\) be a positive self-adjoint Laplace-type operator on
\(L^2(N;E_{\rm e})\), with
\begin{equation}
 \mathbb A\ge a_0I,\qquad a_0>0,
\qquad
 \sigma_2(\mathbb A)(x,\xi)=|\xi|_h^2I_{E_{\rm e}}.
\label{eq:v3v34-A}
\end{equation}
On \(N\times[0,\infty)_y\), use the local matrix energy
\begin{equation}
 \mathcal E_{\mathbb A}(\Psi)
 =\frac12\int_0^\infty
 \left(
 \|\partial_y\Psi\|^2
 +\langle\mathbb A\Psi,\Psi\rangle
 \right)dy.
\label{eq:v3v34-energy}
\end{equation}

\begin{theorem}[Exact matrix Poisson extension]
\label{thm:v3v34-matrix-Poisson}
For \(u\in H^{1/2}(N;E_{\rm e})\), the unique decaying energy minimizer
with trace \(u\) is
\begin{equation}
 \Psi_u(y)=e^{-y\mathbb A^{1/2}}u.
\label{eq:v3v34-Poisson}
\end{equation}
Its energy and positive Dirichlet-to-Neumann map are
\begin{equation}
 \mathcal E_{\mathbb A}(\Psi_u)
 =\frac12\langle\mathbb A^{1/2}u,u\rangle,
\qquad
 \boldsymbol\Sigma_{\rm e}=\mathbb A^{1/2}.
\label{eq:v3v34-DtN}
\end{equation}
Moreover,
\begin{equation}
 \sigma_1(\boldsymbol\Sigma_{\rm e})(x,\xi)
 =|\xi|_hI_{E_{\rm e}}.
\label{eq:v3v34-DtN-symbol}
\end{equation}
\end{theorem}

\begin{proof}
The spectral theorem gives an orthonormal eigenbasis
\(\mathbb Au_j=a_ju_j\), \(a_j\ge a_0\).  The decaying extension of the
\(j\)-th trace coefficient is \(e^{-\sqrt{a_j}y}u_j\), and its energy is
\(\sqrt{a_j}|u_j|^2/2\).  Summation proves
\eqref{eq:v3v34-Poisson}--\eqref{eq:v3v34-DtN}.  Positivity gives
uniqueness.  The classical square-root calculus applied to the scalar
principal symbol in \eqref{eq:v3v34-A} gives
\eqref{eq:v3v34-DtN-symbol}.
\end{proof}

\subsection{Positive coupling and the no-cross-subsidy rule}
\label{sec:v3v34-coupling}

Let \(\Gamma\) be a smooth nonnegative self-adjoint bundle endomorphism and
impose the local trace relation
\begin{equation}
 \Psi(0)=\Gamma^{1/2}u,
\qquad
 u=(u_{\rm s},u_{\rm tf}).
\label{eq:v3v34-trace}
\end{equation}

\begin{proposition}[Matrix impedance from a positive local energy]
\label{prop:v3v34-impedance}
Eliminating \(\Psi\) adds the positive boundary form
\begin{equation}
 \frac12\langle\mathbb Zu,u\rangle,
\qquad
 \mathbb Z=\Gamma^{1/2}\mathbb A^{1/2}\Gamma^{1/2}.
\label{eq:v3v34-Z}
\end{equation}
The impedance is classical of order one and has principal symbol
\begin{equation}
 \sigma_1(\mathbb Z)(x,\xi)=|\xi|_h\Gamma(x).
\label{eq:v3v34-Z-symbol}
\end{equation}
\end{proposition}

\begin{proof}
Apply Theorem~\ref{thm:v3v34-matrix-Poisson} to the trace
\(\Gamma^{1/2}u\).  Composition of principal symbols gives
\eqref{eq:v3v34-Z-symbol}.
\end{proof}

\begin{theorem}[No principal cross subsidy]
\label{thm:v3v34-no-cross-subsidy}
Suppose \(\Gamma\) vanishes on a nonzero fibre subspace \(V\subset E_{\rm
e}\).  Then the order-one symbol of \(\mathbb Z\) vanishes on \(V\), even
if lower-order terms of \(\mathbb A\) mix \(V\) with its orthogonal
complement.  In particular, a coupling only to \(E_{\rm s}\) cannot supply
an order-one tracefree margin, and a coupling only to \(E_{\rm tf}\) cannot
repair the V32 scalar line.
\end{theorem}

\begin{proof}
For \(v\in V\), \(\Gamma v=0\).  Equation
\eqref{eq:v3v34-Z-symbol} gives
\(\sigma_1(\mathbb Z)v=|\xi|\Gamma v=0\).  Lower-order mixing in
\(\mathbb A\) changes only degree zero and below in its square root and
cannot create the missing degree-one restriction.
\end{proof}

\begin{firewall}[A shared bulk field is not a shared payment]
\label{fw:v3v34-shared-field}
Writing one matrix field \(\Psi\) does not pay every boundary sector.
The positive trace incidence \(\Gamma^{1/2}\) must have a nonzero principal
restriction on each sector requiring an order-one margin.  This is the
operator version of retaining each actual payer incidence.
\end{firewall}

\subsection{Positivity under scalar--tracefree mixing}
\label{sec:v3v34-mixing}

Write the tangential operator in blocks,
\begin{equation}
 \mathbb A=
 \begin{pmatrix}
 A_{\rm s}&W\\
 W^*&A_{\rm tf}
 \end{pmatrix},
\qquad
 A_{\rm s},A_{\rm tf}>0.
\label{eq:v3v34-A-block}
\end{equation}
Assume the normalized incidence
\begin{equation}
 C=A_{\rm s}^{-1/2}WA_{\rm tf}^{-1/2}
\label{eq:v3v34-C}
\end{equation}
extends boundedly.

\begin{theorem}[Exact block positivity criterion]
\label{thm:v3v34-block-positive}
If
\begin{equation}
 \|C\|=\theta<1,
\label{eq:v3v34-theta}
\end{equation}
then
\begin{equation}
 \langle\mathbb A(u_{\rm s},u_{\rm tf}),
                 (u_{\rm s},u_{\rm tf})\rangle
 \ge
 (1-\theta)\left(
 \|A_{\rm s}^{1/2}u_{\rm s}\|^2
 +\|A_{\rm tf}^{1/2}u_{\rm tf}\|^2
 \right).
\label{eq:v3v34-block-gap}
\end{equation}
Conversely, if \(C\) has a unit singular pair with singular value greater
than one, the corresponding two-dimensional block of \(\mathbb A\) has a
negative direction.
\end{theorem}

\begin{proof}
Put \(x=A_{\rm s}^{1/2}u_{\rm s}\) and
\(y=A_{\rm tf}^{1/2}u_{\rm tf}\).  The quadratic form is
\[
 \|x\|^2+\|y\|^2+2\operatorname{Re}\langle x,Cy\rangle.
\]
The bound
\(2|\langle x,Cy\rangle|
\le\theta(\|x\|^2+\|y\|^2)\)
proves \eqref{eq:v3v34-block-gap}.  If
\(Cy=\sigma x\), \(C^*x=\sigma y\) with unit \(x,y\) and
\(\sigma>1\), choose the relative sign so the cross term is
\(-2\sigma\); the value \(2-2\sigma\) is negative.
\end{proof}

\begin{counterexample}[Diagonal gaps do not control arbitrary mixing]
\label{cex:v3v34-mixing}
In a two-dimensional fibre, take
\[
 A_{\rm s}=A_{\rm tf}=1,\qquad W=w.
\]
The eigenvalues of \(\mathbb A\) are \(1\pm|w|\).  Both diagonal blocks
remain positive while the full energy becomes indefinite when
\(|w|>1\).  Thus separate scalar and tracefree energy estimates cannot
replace the joint criterion \eqref{eq:v3v34-theta}.
\end{counterexample}

\subsection{Boundary-symbol margins under bounded mixing}
\label{sec:v3v34-margin}

\begin{lemma}[Joint singular-value margin]
\label{lem:v3v34-margin}
Let \(B_0(\omega)\) be a continuous family of square matrices on a compact
normalized frequency set \(\Omega\), with
\begin{equation}
 \inf_{\omega\in\Omega}s_{\min}(B_0(\omega))=c_0>0.
\label{eq:v3v34-c0}
\end{equation}
If \(E(\omega)\) is continuous and
\begin{equation}
 \sup_{\omega\in\Omega}\|E(\omega)\|\le\varepsilon<c_0,
\label{eq:v3v34-epsilon}
\end{equation}
then
\begin{equation}
 s_{\min}(B_0(\omega)+E(\omega))
 \ge c_0-\varepsilon>0.
\label{eq:v3v34-mixed-margin}
\end{equation}
\end{lemma}

\begin{proof}
For every unit vector \(v\),
\[
 \|(B_0+E)v\|
 \ge\|B_0v\|-\|Ev\|
 \ge c_0-\varepsilon.
\]
Take the infimum over \(v\) and \(\omega\).
\end{proof}

\begin{assessment}[What this lemma permits]
\label{ass:v3v34-mixing-use}
If the scalar V33 row and a separately proved tracefree row form a
block-diagonal normalized symbol with margins \(c_{\rm s},c_{\rm tf}>0\),
then curvature and second-fundamental-form mixing is harmless whenever its
complete normalized matrix norm is smaller than
\(\min(c_{\rm s},c_{\rm tf})\).  The lemma does not prove that smallness;
it states the exact joint gate and forbids estimating repeated mixed blocks
independently.
\end{assessment}

\subsection{Lorentzian continuation of the covariant matrix reservoir}
\label{sec:v3v34-Lorentzian}

On a timelike boundary with time \(t\) and spatial covector \(\xi\), the
local covariant reservoir wave equation has outgoing Laplace--Fourier root
\begin{equation}
 \rho(s,\xi)
 =\sqrt{s^2+|\xi|^2+m_{\rm e}^2},
\qquad
 \operatorname{Re}\rho>0\quad(\operatorname{Re}s>0).
\label{eq:v3v34-rho}
\end{equation}
Its homogeneous principal part is
\begin{equation}
 \rho_1(s,\xi)=\sqrt{s^2+|\xi|^2}.
\label{eq:v3v34-rho-principal}
\end{equation}

\begin{theorem}[Covariant matrix impedance vanishes at Lorentzian
glancing]
\label{thm:v3v34-signature-obstruction}
At every nonzero closed-frequency glancing point
\begin{equation}
 s=i\tau,\qquad |\xi|=|\tau|,
\label{eq:v3v34-glancing}
\end{equation}
the homogeneous symbol \(\rho_1\) vanishes.  Therefore the principal
symbol of every covariant matrix impedance obtained from
\eqref{eq:v3v34-energy} by Lorentzian continuation also vanishes on the
V15 tracefree glancing fibre.  A fixed mass gives
\(\rho=m_{\rm e}\) there, but this contribution is lower order under
homogeneous high-frequency scaling.
\end{theorem}

\begin{proof}
At \eqref{eq:v3v34-glancing},
\[
 s^2+|\xi|^2=-\tau^2+|\xi|^2=0,
\]
so \(\rho_1=0\).  The matrix coupling multiplies this scalar leading root
by \(\Gamma\) and cannot change its zero.  With mass,
\(\rho=m_{\rm e}\), while scaling
\((s,\xi)\mapsto R(s,\xi)\) divides the row-normalized mass contribution
by \(R\); it tends to zero as \(R\to\infty\).
\end{proof}

\begin{corollary}[The V33 Euclidean theorem does not imply the V15
hyperbolic margin]
\label{cor:v3v34-no-Wick-inference}
The Euclidean symbol \(|\xi|_h\Gamma\) in
\eqref{eq:v3v34-Z-symbol} is elliptic for \(\xi\ne0\), but its Lorentzian
covariant continuation is characteristic on the null cone.  Strong
Euclidean parameter ellipticity of the scalar row cannot be Wick-rotated
into the zero-loss hyperbolic tracefree estimate without an additional
argument.
\end{corollary}

\begin{proof}
Compare \eqref{eq:v3v34-Z-symbol} with
\eqref{eq:v3v34-rho-principal} and apply
Theorem~\ref{thm:v3v34-signature-obstruction}.
\end{proof}

\subsection{The transmission-line alternative and its cost}
\label{sec:v3v34-line}

V16's local transmission-line reservoir has effective order-one symbol
\(\gamma s\), \(\gamma>0\), on the tracefree block.

\begin{proposition}[The line symbol pays glancing but selects a
foliation]
\label{prop:v3v34-line}
At \eqref{eq:v3v34-glancing} with \(\tau\ne0\),
\begin{equation}
 \gamma s=i\gamma\tau\ne0.
\label{eq:v3v34-line-glancing}
\end{equation}
It therefore supplies the missing homogeneous tracefree currency.
However, it is not the symbol of an intrinsic Lorentz-covariant scalar
function of the boundary metric, because it singles out the chosen time
covector \(dt\).
\end{proposition}

\begin{proof}
The first statement is direct substitution.  An intrinsic scalar
first-order symbol built only from the Lorentz metric depends on the
quadratic invariant \(-\tau^2+|\xi|^2\) through a square-root branch and
vanishes on its null cone.  The nonzero linear factor \(i\tau\) requires a
chosen time orientation and foliation.
\end{proof}

\begin{firewall}[A hybrid matrix is not yet the SM--Einstein boundary
action]
\label{fw:v3v34-hybrid}
One may place the V33 Euclidean scalar edge channel and the V16
transmission-line tracefree channel on a direct sum and then apply
Lemma~\ref{lem:v3v34-margin}.  That hybrid is local and can have positive
energy in a chosen foliation, but it is not an intrinsic covariant boundary
matter action.  Its BRST transformation, Wick rotation, and
Osterwalder--Schrader interpretation must be constructed before it can
belong to the full continuum theory.
\end{firewall}

\subsection{V34 conclusion and V35 audit mandate}
\label{sec:v3v34-conclusion}

V34 proves the positive matrix reservoir theorem on
\(E_{\rm s}\oplus E_{\rm tf}\).  Its eliminated impedance is
\(\Gamma^{1/2}\mathbb A^{1/2}\Gamma^{1/2}\), with principal symbol
\(|\xi|\Gamma\).  Positivity under scalar--tracefree mixing is controlled
by the single joint incidence norm
\(\|A_{\rm s}^{-1/2}WA_{\rm tf}^{-1/2}\|<1\), and a uniform boundary
margin survives a complete normalized perturbation smaller than the
diagonal margin.  No sector can borrow an order-one payment from a coupling
which vanishes on that sector.

The Euclidean matrix construction repairs the V32 scalar kernel, but its
covariant Lorentzian principal root vanishes at the V15 null glancing set.
The transmission-line factor pays that set only by selecting a boundary
time foliation.  V35 is the next mandatory companion audit.  It should use
the newest Yang--Mills and Navier--Stokes results to choose between
constructing a reflection-positive Euclidean origin for the line bath and
returning upstream to the boundary gauge/edge phase space where a preferred
time flow may arise from physical boundary data rather than being inserted
by hand.

\section{V35: companion audit and a reflection-positive Euclidean origin
for subluminal edge baths}
\label{sec:v3v35}

\subsection{Mandatory V35 companion audit}
\label{sec:v3v35-audit}

The audit was performed after V34 and before choosing between a
reflection-positive origin for the transmission line and a direct
Hamiltonian edge-phase-space construction.  The newest companion snapshots
were
\begin{center}
\begin{tabular}{|p{0.30\textwidth}|p{0.12\textwidth}|p{0.43\textwidth}|}
\hline
programme and newest file & bytes & SHA--256 \\
\hline
Navier--Stokes,
\texttt{Renewal\_NS\_Stability\_Research\_Notes\_V31.tex}
& 887537
& \texttt{F1121B62238AD5491BB7CF03635EA9934942EDF573ED8FAAA9EE79E1D38A6696}
\\
\hline
Yang--Mills,
\texttt{Renewal\_Yang\_Mills\_Mass\_Gap\_Research\_Notes\_V87.tex}
& 1413054
& \texttt{EBAFD9961174D65073AE8DB02229454B6667109EB6E7475A119CEA9DBE3A2606}
\\
\hline
\end{tabular}
\end{center}
The Navier--Stokes snapshot is unchanged.  The Yang--Mills folder also
contains V86; it was left untouched.  V87 is newer and is the file used for
this audit.

Yang--Mills V87 proves unequal-spin heat cards and a certified all-payer
sector, but it also constructs a mesoscopic family which is input-Casimir
hot while its binary covariance remains unsuppressed.  Thus the V86
hot-prefix/cold-prefix dichotomy was incomplete.  V87 replaces it with
coordinate-resolved cold, mesoscopic, and superhot regimes and retains one
paired carrier through all payer positions.

\begin{transferprinciple}[Do not force a two-regime impedance choice]
\label{tp:v3v35-three-regimes}
The V34 alternatives \emph{covariant bath} and \emph{transmission line}
are endpoints of a larger local family.  Introduce the actual bath
propagation speed and compute its Euclidean reflection property and
Lorentzian glancing symbol before classifying it.  A middle speed regime
may retain both locality and the required margin even though it is neither
Lorentz-covariant nor spatially ultralocal.
\end{transferprinciple}

No Yang--Mills bound is imported.  The audit changes the parameter space
of the edge-bath construction and forces the intermediate regime to be
tested explicitly.

\subsection{A local Euclidean bath with adjustable boundary speed}
\label{sec:v3v35-Euclidean-bath}

Let the Euclidean timelike boundary have coordinates
\((\tau,x)\in\mathbb R\times\Sigma\), with \(\Sigma\) a compact
two-manifold for the global statement and \(\mathbb R^2\) for the frozen
symbol.  Add a bath coordinate \(y\ge0\).  For
\(0\le c_{\rm e}\le1\) and \(m_{\rm e}\ge0\), define
\begin{equation}
 S_{\rm bath}^{E}(\psi)
 =\frac12\int_{\mathbb R_\tau}\int_\Sigma\int_0^\infty
 \left(
 |\partial_\tau\psi|^2+|\partial_y\psi|^2
 +c_{\rm e}^2|\nabla_x\psi|^2+m_{\rm e}^2|\psi|^2
 \right)dy\,dx\,d\tau.
\label{eq:v3v35-E-action}
\end{equation}
Couple a real boundary component \(q(\tau,x)\) by
\begin{equation}
 \psi(\tau,x,0)=\alpha q(\tau,x),
\qquad \alpha>0.
\label{eq:v3v35-trace}
\end{equation}
The case \(c_{\rm e}=0\) is the transmission-line bath of V16 after Wick
rotation.  The case \(c_{\rm e}=1\) is the boundary-Lorentz-covariant
reservoir at principal level.

\begin{theorem}[Exact Euclidean bath kernel]
\label{thm:v3v35-E-kernel}
For Euclidean frequency \(\omega\) and spatial covector \(\xi\), the
decaying solution of the bath equation with trace
\eqref{eq:v3v35-trace} is
\begin{equation}
 \widehat\psi(y)
 =\alpha e^{-q_Ey}\widehat q,
\qquad
 q_E(\omega,\xi)
 =\sqrt{\omega^2+c_{\rm e}^2|\xi|^2+m_{\rm e}^2}.
\label{eq:v3v35-E-Poisson}
\end{equation}
Eliminating the bath adds
\begin{equation}
 S_{\rm eff}^{E}(q)
 =\frac{\alpha^2}{2}
 \int q_E(\omega,\xi)|\widehat q(\omega,\xi)|^2
 \,d\omega\,d\xi
\label{eq:v3v35-E-effective}
\end{equation}
with the usual spectral measure.  The kernel is nonnegative.
\end{theorem}

\begin{proof}
The transformed Euler equation is
\[
 -\partial_y^2\widehat\psi
 +(\omega^2+c_{\rm e}^2|\xi|^2+m_{\rm e}^2)
 \widehat\psi=0.
\]
The positive root gives \eqref{eq:v3v35-E-Poisson}.  The energy calculation
of Theorem~\ref{thm:v3v33-edge-DtN} gives
\(\alpha^2q_E|\widehat q|^2/2\) in each mode, proving
\eqref{eq:v3v35-E-effective}.
\end{proof}

\begin{firewall}[The massless joint zero mode]
\label{fw:v3v35-zero-mode}
If \(m_{\rm e}=0\) and \((\omega,\xi)=(0,0)\), the decaying root vanishes
and a constant bath profile is not selected by finite energy.  One must
remove the joint mean, work at finite bath length, or keep
\(m_{\rm e}>0\) until the infrared measure is specified.  The nonzero
homogeneous glancing calculation below does not pay this zero mode.
\end{firewall}

\subsection{Reflection positivity of the local parent}
\label{sec:v3v35-RP}

Let \(\Theta\) be Euclidean time reflection,
\begin{equation}
 (\Theta\psi)(\tau,x,y)=\psi(-\tau,x,y),
\qquad
 (\Theta q)(\tau,x)=q(-\tau,x).
\label{eq:v3v35-reflection}
\end{equation}
First keep a finite spatial and bath cutoff and, if needed, a positive
infrared mass.  The field configuration at fixed \(\tau\) lies in a real
Hilbert space \(\mathcal K\) incorporating \(q\), \(\psi\), and the
time-local trace constraint.  Its quadratic action has the form
\begin{equation}
 \frac12\int_{\mathbb R}
 \left(
 \|\partial_\tau\Phi\|_{\mathcal K}^2
 +\langle K\Phi,\Phi\rangle_{\mathcal K}
 \right)d\tau,
\qquad K\ge0.
\label{eq:v3v35-abstract-action}
\end{equation}

\begin{lemma}[Positive reflected covariance]
\label{lem:v3v35-reflected-covariance}
If \(K>0\), the covariance of
\(-\partial_\tau^2+K\) is
\begin{equation}
 C(\tau,\tau')
 =\frac1{2K^{1/2}}
 e^{-K^{1/2}|\tau-\tau'|}.
\label{eq:v3v35-covariance}
\end{equation}
For every compactly supported
\(f:(0,\infty)\to\mathcal K\),
\begin{equation}
 \int_0^\infty\!\!\int_0^\infty
 \langle f(\tau),C(-\tau,\tau')f(\tau')\rangle
 d\tau\,d\tau'
 =
 \frac12\left\|
 \int_0^\infty
 K^{-1/4}e^{-\tau K^{1/2}}f(\tau)d\tau
 \right\|^2
 \ge0.
\label{eq:v3v35-RP-kernel}
\end{equation}
\end{lemma}

\begin{proof}
Spectral calculus reduces the inverse of
\(-\partial_\tau^2+K\) to the scalar Green function
\(e^{-\sqrt\lambda|\tau-\tau'|}/(2\sqrt\lambda)\), proving
\eqref{eq:v3v35-covariance}.  For positive \(\tau,\tau'\),
\[
 C(-\tau,\tau')
 =\frac1{2K^{1/2}}
 e^{-(\tau+\tau')K^{1/2}}.
\]
Factor the two semigroups and move \(K^{-1/4}\) to each side.  This gives
the squared norm in \eqref{eq:v3v35-RP-kernel}.
\end{proof}

\begin{theorem}[Reflection positivity survives the trace constraint and
bath elimination]
\label{thm:v3v35-RP}
At every finite cutoff, the joint Gaussian measure defined by
\eqref{eq:v3v35-E-action} together with any reflection-positive primary
quadratic action and the time-local constraint
\eqref{eq:v3v35-trace} is reflection positive.  Its marginal measure on
the primary boundary field \(q\), whose action contains
\eqref{eq:v3v35-E-effective}, is also reflection positive.
\end{theorem}

\begin{proof}
The constraint at time \(\tau\) is mapped to the identical constraint at
\(-\tau\), so the constrained configuration space is a
\(\Theta\)-invariant closed subspace of \(\mathcal K\).
Lemma~\ref{lem:v3v35-reflected-covariance} applies to the restriction of
\(K\) to that subspace.  Gaussian reflection positivity for polynomial
functionals follows from the positive reflected two-point kernel by Wick's
rule and completion.

Now let \(F(q)\) depend only on positive-time primary fields.  Integrating
the joint expectation
\(\langle\Theta F\,F\rangle\ge0\) over the bath variables is exactly the
same expectation in the marginal \(q\)-measure.  Hence marginalization
preserves the inequality.
\end{proof}

\begin{assessment}[A nonlocal effective action can have a local
reflection-positive parent]
\label{ass:v3v35-nonlocal-parent}
The square root in \eqref{eq:v3v35-E-effective} is nonlocal when written
only on the boundary.  Reflection positivity is proved from the enlarged
local action before elimination.  It would not follow merely from the
pointwise inequality \(q_E\ge0\).
\end{assessment}

\subsection{Lorentzian retarded response}
\label{sec:v3v35-Lorentzian}

Wick rotate the local parent rather than the absolute-value function in
isolation.  The Lorentzian bath equation is
\begin{equation}
 \partial_t^2\psi-\partial_y^2\psi
 -c_{\rm e}^2\Delta_x\psi+m_{\rm e}^2\psi=0.
\label{eq:v3v35-L-wave}
\end{equation}
For a Laplace mode \(e^{st+i\xi\cdot x}\), \(\operatorname{Re}s>0\), its
no-incoming root is
\begin{equation}
 q_L(s,\xi)
 =\sqrt{s^2+c_{\rm e}^2|\xi|^2+m_{\rm e}^2},
\qquad
 \operatorname{Re}q_L>0.
\label{eq:v3v35-L-root}
\end{equation}

\begin{proposition}[Retarded impedance and passivity]
\label{prop:v3v35-passive}
Eliminating the Lorentzian bath adds the retarded boundary impedance
\begin{equation}
 Z_L(s,\xi)=\alpha^2q_L(s,\xi).
\label{eq:v3v35-ZL}
\end{equation}
For \(\operatorname{Re}s>0\), the outgoing branch has
\(\operatorname{Re}Z_L>0\).  The enlarged primary-plus-bath energy is
conserved; after bath elimination, the lost primary energy is the
nonnegative flux transported toward \(y=\infty\).
\end{proposition}

\begin{proof}
The decaying Laplace solution is
\(\widehat\psi(y)=\alpha e^{-q_Ly}\widehat q\), so
\(-\partial_y\widehat\psi(0)=\alpha q_L\widehat q\).  The trace coupling
supplies the second factor of \(\alpha\), giving
\eqref{eq:v3v35-ZL}.  The chosen square root maps the right half-plane to
the right half-plane.  Multiplying
\eqref{eq:v3v35-L-wave} by \(\partial_t\psi\), integrating, and combining
the boundary flux with the primary junction gives the energy statement,
as in Theorem~\ref{thm:v3v16-reservoir-energy}.
\end{proof}

\begin{corollary}[The transmission line has Euclidean kernel
\(\alpha^2|\omega|\)]
\label{cor:v3v35-line}
For \(c_{\rm e}=m_{\rm e}=0\),
\begin{equation}
 q_E=|\omega|,
\qquad
 q_L=s\quad(\operatorname{Re}s>0).
\label{eq:v3v35-line-pair}
\end{equation}
Thus the V16 line impedance \(\alpha^2s\) is the retarded response of the
local reflection-positive Euclidean half-plane bath.
\end{corollary}

\begin{proof}
Substitute \(c_{\rm e}=m_{\rm e}=0\) in
\eqref{eq:v3v35-E-Poisson} and \eqref{eq:v3v35-L-root}; the outgoing
right-half-plane square root of \(s^2\) is \(s\).
\end{proof}

\subsection{The intermediate-speed glancing calculation}
\label{sec:v3v35-speed}

The primary relativistic boundary glancing set is
\begin{equation}
 s=i\tau,\qquad |\xi|=|\tau|\ne0.
\label{eq:v3v35-primary-glancing}
\end{equation}

\begin{theorem}[Every uniformly subluminal bath pays the primary null
cone]
\label{thm:v3v35-subluminal}
Let \(0\le c_{\rm e}<1\) and \(m_{\rm e}=0\).  On
\eqref{eq:v3v35-primary-glancing}, the limiting outgoing homogeneous root
is
\begin{equation}
 q_L(i\tau,\xi)
 =i\,\operatorname{sgn}(\tau)
 \sqrt{1-c_{\rm e}^2}\,|\tau|,
\label{eq:v3v35-subluminal-root}
\end{equation}
up to the fixed outgoing convention.  In particular,
\begin{equation}
 |Z_L(i\tau,\xi)|
 =\alpha^2\sqrt{1-c_{\rm e}^2}\,|\tau|>0.
\label{eq:v3v35-subluminal-margin}
\end{equation}
If
\begin{equation}
 1-c_{\rm e}^2\ge\delta>0,
\label{eq:v3v35-speed-gap}
\end{equation}
the normalized glancing margin is at least
\(\alpha^2\sqrt\delta\).
\end{theorem}

\begin{proof}
At \eqref{eq:v3v35-primary-glancing},
\[
 s^2+c_{\rm e}^2|\xi|^2
 =-(1-c_{\rm e}^2)\tau^2.
\]
The limiting outgoing square root is the imaginary number in
\eqref{eq:v3v35-subluminal-root}.  Its modulus gives
\eqref{eq:v3v35-subluminal-margin}; division by \(|\tau|\) and
\eqref{eq:v3v35-speed-gap} give the uniform bound.
\end{proof}

\begin{corollary}[The covariant endpoint is exactly the margin-collapse
point]
\label{cor:v3v35-covariant-endpoint}
At \(c_{\rm e}=1\), the homogeneous root in
\eqref{eq:v3v35-subluminal-root} is zero.  As
\(c_{\rm e}\uparrow1\), the normalized margin collapses like
\(\sqrt{1-c_{\rm e}^2}\).
\end{corollary}

\begin{proof}
Set \(c_{\rm e}=1\) and take the indicated limit in
\eqref{eq:v3v35-subluminal-margin}.
\end{proof}

\begin{counterexample}[A nonzero but unseparated speed does not give a
uniform family]
\label{cex:v3v35-speed-collapse}
Let \(c_n\uparrow1\) while \(\alpha\) is fixed.  Every individual bath has
a nonzero glancing response, yet its normalized margin is
\(\alpha^2\sqrt{1-c_n^2}\to0\).  Pointwise nonvanishing therefore does not
give a uniform continuum estimate over an unseparated speed family.
\end{counterexample}

\begin{assessment}[Cold, intermediate, and covariant regimes]
\label{ass:v3v35-regimes}
The line bath \(c_{\rm e}=0\), the intermediate baths
\(0<c_{\rm e}<1\), and the covariant endpoint \(c_{\rm e}=1\) are
analytically distinct.  The first two have homogeneous glancing currency;
the last does not.  This is the boundary-impedance analogue of the V87
warning that endpoint labels do not control an unresolved middle regime.
\end{assessment}

\subsection{Euclidean symmetry retained and lost}
\label{sec:v3v35-symmetry}

\begin{proposition}[Locality and reflection symmetry]
\label{prop:v3v35-symmetry}
For every \(0\le c_{\rm e}\le1\), the enlarged action
\eqref{eq:v3v35-E-action} is local, invariant under Euclidean time
translations and reflection, and invariant under isometries of \(\Sigma\).
For \(c_{\rm e}=1\) its principal derivatives are rotationally symmetric
among the boundary coordinates.  For \(c_{\rm e}\ne1\), the coefficient of
\(|\nabla_x\psi|^2\) distinguishes the reflection-time direction.
\end{proposition}

\begin{proof}
All claims follow directly from the derivative contractions in
\eqref{eq:v3v35-E-action}.  Equality of the \(\tau\) and \(x\) derivative
coefficients occurs exactly at \(c_{\rm e}=1\).
\end{proof}

\begin{firewall}[Reflection positivity is not full Euclidean covariance]
\label{fw:v3v35-not-covariant}
Theorem~\ref{thm:v3v35-RP} supplies the positivity axiom associated with a
chosen Euclidean time reflection.  For \(c_{\rm e}<1\), the action lacks
full boundary Euclidean rotational invariance, so the standard
Osterwalder--Schrader route to a Lorentz-covariant boundary theory is not
yet complete.  The preferred time must arise from physical boundary data,
a gauge choice with proved independence of observables, or an enlarged
covariant medium field.
\end{firewall}

\subsection{Matrix extension}
\label{sec:v3v35-matrix}

Let the bath take values in the scalar--tracefree bundle \(E_{\rm e}\) of
V34, and let \(C_{\rm e}\) be a positive self-adjoint endomorphism with
spectrum in \([0,1-\delta]\).  Replace
\(c_{\rm e}^2|\nabla_x\psi|^2\) by the corresponding positive bundle
quadratic form.

\begin{proposition}[Matrix speed-gap margin]
\label{prop:v3v35-matrix-speed}
If the speed operator commutes with the frozen principal polarization
splitting, the outgoing root on a speed eigenspace with eigenvalue
\(\nu=c_{\rm e}^2\) has normalized primary-glancing modulus
\(\sqrt{1-\nu}\).  Hence
\begin{equation}
 \sigma(C_{\rm e})\subset[0,1-\delta]
\quad\Longrightarrow\quad
 |q_L|/|\tau|\ge\sqrt\delta
\label{eq:v3v35-matrix-gap}
\end{equation}
on every coupled scalar and tracefree principal polarization.
\end{proposition}

\begin{proof}
Diagonalize the finite-dimensional self-adjoint principal speed
endomorphism and apply Theorem~\ref{thm:v3v35-subluminal} on each
eigenspace.  The spectral upper bound gives
\eqref{eq:v3v35-matrix-gap}.
\end{proof}

\begin{firewall}[Noncommuting lower symbols require the V34 joint norm]
\label{fw:v3v35-matrix-mixing}
Curvature, shape, and BRST lower symbols need not commute with the speed
endomorphism.  Their stability is governed by the complete normalized
matrix perturbation in Lemma~\ref{lem:v3v34-margin}, not by applying
\eqref{eq:v3v35-matrix-gap} separately to diagonal entries.
\end{firewall}

\subsection{V35 conclusion and V36 target}
\label{sec:v3v35-conclusion}

The mandatory audit has replaced a false two-endpoint choice by an exact
speed family.  The transmission-line impedance \(\alpha^2s\) has a local
Euclidean parent with kernel \(\alpha^2|\omega|\), and the enlarged
Gaussian measure is reflection positive before and after bath elimination.
More generally, every bath with speed
\(0\le c_{\rm e}<1\) is local and reflection positive and supplies the
homogeneous glancing margin
\(\alpha^2\sqrt{1-c_{\rm e}^2}\).  The Lorentz-covariant endpoint
\(c_{\rm e}=1\) is exactly where that margin collapses.

This resolves the reflection-positivity objection to the V16 line bath but
not the covariance objection.  V36 should go upstream to the boundary
clock/normal data of the Einstein variational problem.  It should determine
whether a unit timelike boundary flow can be constructed covariantly from
admitted boundary fields, derive the speed operator and energy from that
flow, and prove that BRST-invariant observables are independent of its
gauge representative.  If no such clock field is present, the subluminal
bath must remain a gauge-fixed regulator rather than a constituent of the
full covariant continuum.

\section{V36: covariant boundary clocks, subluminal bath geometry, and
the physical-versus-gauge dichotomy}
\label{sec:v3v36}

\subsection{The source of a preferred boundary time}
\label{sec:v3v36-purpose}

V35 proves reflection positivity and a homogeneous glancing margin for a
bath whose propagation speed is strictly below the primary boundary light
speed.  Its coefficients distinguish a time direction.  V36 asks whether
that direction can be constructed covariantly rather than inserted as a
coordinate label.

A scalar clock with timelike gradient gives an exact covariant answer on
the open set where its gradient stays uniformly timelike.  The resulting
bath action is diffeomorphism invariant and reduces to V35 in clock-adapted
coordinates.  This does not by itself solve the Standard Model--Einstein
programme: a physical clock is an additional field or a restriction on
matter backgrounds, whereas a gauge-fixing clock must occur in a
BRST-exact regulator if physical observables are to be independent of it.

\subsection{Clock geometry on a timelike boundary}
\label{sec:v3v36-clock}

Let \((N,\gamma)\) be a time-oriented Lorentzian three-manifold with
signature \((-++)\).  Let \(T:N\to\mathbb R\) be a smooth scalar and define
\begin{equation}
 X_T=-\gamma^{ab}\nabla_aT\nabla_bT.
\label{eq:v3v36-X}
\end{equation}
Assume
\begin{equation}
 X_T\ge\chi^2>0.
\label{eq:v3v36-clock-gap}
\end{equation}
Choose the time orientation so that
\begin{equation}
 u_a=-\frac{\nabla_aT}{\sqrt{X_T}}
\label{eq:v3v36-u}
\end{equation}
is future directed, and define the spatial projector
\begin{equation}
 q_{ab}=\gamma_{ab}+u_au_b,
\qquad
 q_a{}^b=\delta_a{}^b+u_au^b.
\label{eq:v3v36-q}
\end{equation}

\begin{proposition}[Exact clock identities]
\label{prop:v3v36-clock-identities}
The fields in \eqref{eq:v3v36-u}--\eqref{eq:v3v36-q} obey
\begin{align}
 \gamma^{ab}u_au_b&=-1,
\label{eq:v3v36-unit}\\
 q_{ab}u^b&=0,
\label{eq:v3v36-orthogonal}\\
 q_a{}^cq_c{}^b&=q_a{}^b.
\label{eq:v3v36-projector}
\end{align}
The distribution orthogonal to \(u\) is integrable, and \(u\) is unchanged
under every increasing reparametrization \(T\mapsto f(T)\).
\end{proposition}

\begin{proof}
The first three identities follow by substituting
\(\gamma^{ab}\nabla_aT\nabla_bT=-X_T\).  Since \(u\) is proportional to
the exact one-form \(dT\),
\[
 u\wedge du=0,
\]
so Frobenius' theorem gives integrability.  If \(f'>0\), both
\(dT\) and \(\sqrt{X_T}\) acquire the same factor \(f'(T)\), leaving
\eqref{eq:v3v36-u} unchanged.
\end{proof}

\begin{proposition}[Diffeomorphism covariance]
\label{prop:v3v36-covariance}
Under a diffeomorphism \(\phi\),
\begin{equation}
 (\gamma,T)\longmapsto(\phi^*\gamma,\phi^*T)
\label{eq:v3v36-diffeo}
\end{equation}
implies
\begin{equation}
 u[\phi^*\gamma,\phi^*T]=\phi^*u[\gamma,T],
\qquad
 q[\phi^*\gamma,\phi^*T]=\phi^*q[\gamma,T].
\label{eq:v3v36-natural}
\end{equation}
\end{proposition}

\begin{proof}
Exterior differentiation, inverse-metric contraction, square roots of
positive scalars, and tensor products all commute with pullback.  Apply
these operations to \eqref{eq:v3v36-X}--\eqref{eq:v3v36-q}.
\end{proof}

\subsection{Covariant subluminal bath action}
\label{sec:v3v36-bath}

Let \(y\ge0\) be the reservoir coordinate.  For a real scalar reservoir
\(\psi\), define
\begin{equation}
 \begin{split}
 S_{\rm bath}^{L}[\psi;\gamma,T]
 =\frac12\int_N\int_0^\infty\sqrt{-\gamma}\,
 \bigl[
 &(u^a\nabla_a\psi)^2-(\partial_y\psi)^2\\
 &-c_{\rm e}^2q^{ab}\nabla_a\psi\nabla_b\psi
 -m_{\rm e}^2\psi^2
 \bigr]\,dy\,d^3x,
 \end{split}
\label{eq:v3v36-bath-action}
\end{equation}
where \(0\le c_{\rm e}<1\).

\begin{theorem}[Local covariance and frozen reduction]
\label{thm:v3v36-covariant-bath}
The action \eqref{eq:v3v36-bath-action} is a local diffeomorphism scalar
when \(\psi\) and \(T\) transform as scalars and \(\gamma\) as a metric.
At a frozen point in clock-adapted orthonormal coordinates,
\begin{equation}
 u=\partial_t,\qquad
 q^{ab}\nabla_a\psi\nabla_b\psi=|\nabla_x\psi|^2,
\label{eq:v3v36-adapted}
\end{equation}
and its Euler equation is
\begin{equation}
 \partial_t^2\psi-\partial_y^2\psi
 -c_{\rm e}^2\Delta_x\psi+m_{\rm e}^2\psi=0,
\label{eq:v3v36-wave}
\end{equation}
the V35 bath equation.
\end{theorem}

\begin{proof}
Every index in \eqref{eq:v3v36-bath-action} is contracted using the natural
tensors \(u,q,\gamma\), and the measure is invariant.  This proves
diffeomorphism covariance.  A unit timelike vector and its orthogonal
projector admit the local normal form
\eqref{eq:v3v36-adapted}.  Varying the frozen quadratic action and
integrating by parts gives \eqref{eq:v3v36-wave}.
\end{proof}

\begin{corollary}[Covariant origin of the speed-gap symbol]
\label{cor:v3v36-speed-symbol}
At the frozen point, the retarded outgoing normal root is
\begin{equation}
 q_L(s,\xi)
 =\sqrt{s^2+c_{\rm e}^2|\xi|^2+m_{\rm e}^2}.
\label{eq:v3v36-root}
\end{equation}
For \(m_{\rm e}=0\), the primary null cone
\(s=i\tau,\ |\xi|=|\tau|\ne0\) has normalized modulus
\begin{equation}
 |q_L|/|\tau|=\sqrt{1-c_{\rm e}^2}.
\label{eq:v3v36-gap}
\end{equation}
\end{corollary}

\begin{proof}
Fourier--Laplace transform \eqref{eq:v3v36-wave} and use the decaying
branch.  The second statement is
Theorem~\ref{thm:v3v35-subluminal}.
\end{proof}

\begin{assessment}[Covariant does not mean Lorentz symmetric at a clock
background]
\label{ass:v3v36-spontaneous}
The formula \eqref{eq:v3v36-bath-action} is diffeomorphism covariant
because \(u\) is a field-derived tensor.  A background with
\(dT\ne0\) nevertheless selects a timelike direction.  When
\(c_{\rm e}\ne1\), local boost symmetry of that background is broken by
the bath coefficients.  Covariance of the equations and invariance of a
particular state are different assertions.
\end{assessment}

\subsection{Positive energy in the clock frame}
\label{sec:v3v36-energy}

At a frozen clock background, define
\begin{equation}
 E_{\rm bath}(t)
 =\frac12\int_\Sigma\int_0^\infty
 \left(
 |\partial_t\psi|^2+|\partial_y\psi|^2
 +c_{\rm e}^2|\nabla_x\psi|^2+m_{\rm e}^2|\psi|^2
 \right)dy\,dx.
\label{eq:v3v36-energy}
\end{equation}

\begin{proposition}[Positive conserved bath energy]
\label{prop:v3v36-energy}
For solutions of \eqref{eq:v3v36-wave},
\begin{equation}
 \frac{dE_{\rm bath}}{dt}
 =-\int_\Sigma
 \partial_t\psi(t,x,0)\,\partial_y\psi(t,x,0)\,dx.
\label{eq:v3v36-flux}
\end{equation}
Thus the bath energy is positive, and its only change is the boundary flux
exchanged with the primary system.
\end{proposition}

\begin{proof}
Differentiate \eqref{eq:v3v36-energy}, integrate the spatial and
\(y\)-derivatives by parts, and use \eqref{eq:v3v36-wave}.  The bulk terms
cancel and the endpoint at \(y=\infty\) vanishes, leaving
\eqref{eq:v3v36-flux}.
\end{proof}

\begin{firewall}[Clock acceleration and expansion are lower-order energy
terms]
\label{fw:v3v36-variable-clock}
For variable \((\gamma,T)\), differentiation of the energy associated with
\(u\) produces acceleration, expansion, and deformation terms
\(\nabla u\).  V7 gives the corresponding moving-metric mechanism.
Positivity of the instantaneous density survives, but a uniform energy
estimate requires those deformation terms to be bounded relative to the
principal bath energy.  V36 proves the frozen principal statement, not
that global estimate.
\end{firewall}

\subsection{Stability and collapse of the clock construction}
\label{sec:v3v36-stability}

\begin{proposition}[Linear variation of the normalized clock]
\label{prop:v3v36-clock-variation}
For variations \((k_{\mu\nu},\vartheta)=(\delta\gamma_{\mu\nu},\delta T)\),
\begin{align}
 \delta X_T
 &=
 k^{ab}\nabla_aT\nabla_bT
 -2\gamma^{ab}\nabla_aT\nabla_b\vartheta,
\label{eq:v3v36-dX}\\
 \delta u_a
 &=
 -X_T^{-1/2}\nabla_a\vartheta
 +\frac12X_T^{-3/2}\nabla_aT\,\delta X_T.
\label{eq:v3v36-du}
\end{align}
On a set satisfying \eqref{eq:v3v36-clock-gap}, the map
\((\gamma,T)\mapsto u\) is smooth in Sobolev topologies above the
multiplication threshold, with constants depending on inverse powers of
\(\chi\).
\end{proposition}

\begin{proof}
Differentiate
\(X_T=-\gamma^{ab}T_aT_b\), using
\(\delta\gamma^{ab}=-k^{ab}\), to obtain
\eqref{eq:v3v36-dX}.  Differentiate
\(u_a=-T_aX_T^{-1/2}\) to obtain
\eqref{eq:v3v36-du}.  If \(X_T\ge\chi^2\), the scalar functions
\(X_T^{-1/2}\) and \(X_T^{-3/2}\) and their derivatives are bounded by
the smooth composition theorem, proving the Sobolev statement.
\end{proof}

\begin{counterexample}[Clock degeneration]
\label{cex:v3v36-clock-collapse}
Let \(T_n=\varepsilon_nt\) on a fixed flat boundary with
\(\varepsilon_n\downarrow0\).  Then
\(X_{T_n}=\varepsilon_n^2\), while the formulas defining \(u\) contain
\(X_{T_n}^{-1/2}\).  Although the normalized vector remains
\(\partial_t\) for this exact family, perturbations \(\vartheta\) have
linear response of size \(O(\varepsilon_n^{-1})\).  No uniform clock map
exists as the timelike-gradient gap closes.
\end{counterexample}

\begin{firewall}[The Standard Model vacuum does not supply this clock]
\label{fw:v3v36-SM-clock}
A gauge-invariant Standard Model scalar such as \(H^\dagger H\) is constant
in the homogeneous electroweak vacuum, so its gradient cannot satisfy
\eqref{eq:v3v36-clock-gap}.  V36 therefore does not derive a global clock
from the ordinary SM vacuum.  Using an added khronon, dust, or reference
scalar changes the field content unless it is proved to be a contractible
gauge-fixing sector.
\end{firewall}

\subsection{BRST transformation of the physical clock construction}
\label{sec:v3v36-BRST}

Let \(c^a\) be the diffeomorphism ghost.  On covariant fields, take the
classical BRST action
\begin{equation}
 s\gamma=\mathcal L_c\gamma,\qquad
 sT=\mathcal L_cT,\qquad
 s\psi=\mathcal L_c\psi,
\qquad
 sc=\frac12[c,c]
\label{eq:v3v36-BRST}
\end{equation}
with the sign convention making \(s^2=0\).

\begin{proposition}[BRST closure by naturality]
\label{prop:v3v36-BRST-closure}
The derived tensors satisfy
\begin{equation}
 su=\mathcal L_cu,\qquad sq=\mathcal L_cq,
\label{eq:v3v36-BRST-uq}
\end{equation}
and the covariant bath action obeys
\begin{equation}
 sS_{\rm bath}^{L}=0
\label{eq:v3v36-BRST-action}
\end{equation}
up to the physical boundary flux already displayed in
\eqref{eq:v3v36-flux}.
\end{proposition}

\begin{proof}
Differentiate the naturality identities
\eqref{eq:v3v36-natural} along the ghost-generated diffeomorphism.  This
gives \eqref{eq:v3v36-BRST-uq}.  The integrand of
\eqref{eq:v3v36-bath-action} is a scalar density, so its Lie derivative is
a divergence.  Integration leaves only the admitted boundary flux.
\end{proof}

\begin{assessment}[BRST closed is not BRST exact]
\label{ass:v3v36-closed-not-exact}
If \(T\) and \(\psi\) are physical fields, the action can be BRST closed
without being a gauge-fixing variation.  Observables may then depend on
their couplings and states.  Diffeomorphism covariance does not make the
clock or bath cohomologically invisible.
\end{assessment}

\subsection{Gauge-fixing clocks and the Ward identity}
\label{sec:v3v36-Ward}

\begin{theorem}[Finite-cutoff gauge-fermion independence]
\label{thm:v3v36-Ward}
Let a finite-dimensional BRST cutoff have nilpotent differential \(s\), an
\(s\)-invariant integration measure, and a convergent contour without
boundary contributions in field space.  Let
\begin{equation}
 S_\lambda=S_{\rm inv}+s\Psi_\lambda
\label{eq:v3v36-Slambda}
\end{equation}
and let \(O\) be \(s\)-closed.  Then
\begin{equation}
 \frac{d}{d\lambda}\langle O\rangle_\lambda=0.
\label{eq:v3v36-Ward}
\end{equation}
\end{theorem}

\begin{proof}
Differentiate the normalized expectation.  Its numerator gives
\[
 -\langle O\,s(\partial_\lambda\Psi_\lambda)\rangle
 +\langle O\rangle
  \langle s(\partial_\lambda\Psi_\lambda)\rangle.
\]
The second expectation vanishes by integration of an \(s\)-exact
derivative.  Since \(sO=0\),
\[
 O\,s(\partial_\lambda\Psi_\lambda)
 =s(O\,\partial_\lambda\Psi_\lambda)
\]
up to the fixed parity sign, and its expectation also vanishes.  This
proves \eqref{eq:v3v36-Ward}.
\end{proof}

\begin{corollary}[The precise condition for clock-gauge independence]
\label{cor:v3v36-clock-independence}
If the preferred clock and speed enter only through a gauge fermion
\(\Psi_{\lambda}\), and the hypotheses of
Theorem~\ref{thm:v3v36-Ward} hold with a common cutoff and boundary domain,
then BRST-closed observables are independent of that clock representative
and speed parameter.  If the bath action is the physical closed action
\eqref{eq:v3v36-bath-action}, the theorem does not apply.
\end{corollary}

\begin{proof}
The first statement is Theorem~\ref{thm:v3v36-Ward}; the second follows
because varying a coupling in \(S_{\rm bath}^{L}\) has not been shown to be
an \(s\)-exact variation.
\end{proof}

\begin{firewall}[Continuum Ward identities require anomaly and boundary
control]
\label{fw:v3v36-Ward-continuum}
Theorem~\ref{thm:v3v36-Ward} is exact at a regulator preserving the
measure, contour, and BRST domain.  Removing the cutoff requires proof that
there is no diffeomorphism/BRST anomaly, no determinant-line phase defect,
and no uncancelled boundary charge.  V30 and V19 identify those remaining
gates.
\end{firewall}

\subsection{A covariant source from existing boundary stress}
\label{sec:v3v36-stress-clock}

One might try to avoid a new scalar by extracting \(u\) from the
Brown--York plus matter boundary stress endomorphism
\(\mathsf T^a{}_b\).

\begin{proposition}[Local stress-frame construction]
\label{prop:v3v36-stress-frame}
Suppose \(\mathsf T^a{}_b\) has a simple timelike eigenvector with eigenvalue
\(\lambda_0\), separated algebraically from the other two eigenvalues.  Then,
after choosing time orientation, a normalized timelike eigenvector
\(u(\mathsf T)\) exists smoothly in a neighborhood of that stress tensor
and transforms covariantly under boundary diffeomorphisms.
\end{proposition}

\begin{proof}
The simple-root implicit-function theorem applied to
\((\mathsf T-\lambda I)u=0\) together with
\(\gamma(u,u)=-1\) gives local smooth dependence on the tensor entries.
The equations are tensorial, so uniqueness with fixed time orientation
implies covariance.
\end{proof}

\begin{counterexample}[Vacuum stress degeneracy]
\label{cex:v3v36-vacuum-stress}
If \(\mathsf T^a{}_b=\lambda\delta^a_b\), every vector is an eigenvector
and no preferred timelike line is determined.  This includes cosmological
or tension-like boundary stress.  Therefore the stress-frame construction
cannot provide a uniform clock on a configuration space containing such
vacua.
\end{counterexample}

\subsection{V36 conclusion and V37 target}
\label{sec:v3v36-conclusion}

V36 gives a covariant geometric source for the V35 speed family on the
open sector where a scalar clock has a uniformly timelike gradient.  The
unit flow \(u=-dT/\sqrt{-|dT|^2}\) and its spatial projector produce a
local diffeomorphism-invariant bath action whose frozen response has the
strict subluminal glancing margin.  The clock construction is stable only
while its timelike-gradient gap remains positive.

The result exposes a physical dichotomy.  A physical clock or stress frame
makes the bath covariant but changes the state or field content and can
degenerate on the Standard Model vacuum.  A gauge-fixing clock can be
unobservable only if the entire speed-dependent bath is placed in an
\(s\)-exact sector and the regulated measure has an exact Ward identity.
V37 should therefore construct a subluminal bath as a BRST quartet driven
by a clock-dependent gauge fermion, prove its finite-cutoff energy and
determinant cancellations, and identify whether its tracefree impedance
survives in the classical metric boundary equation after the quartet is
removed.  If it cancels there as well, a BRST-exact bath cannot repair the
physical V15 glancing block and the programme must instead retain a
physical boundary medium or change the junction condition upstream.
\clearpage
\section{V37: finite-cutoff BRST quartet elimination and the tracefree no-payment theorem}
\label{sec:v3v37}

\subsection{Purpose and the distinction to be proved}
\label{sec:v3v37-purpose}

V33--V36 produced two mathematically different ways to attach an
order-one edge response.  A physical half-cylinder field has a positive
Dirichlet energy and leaves its Dirichlet-to-Neumann operator in the
metric boundary equation.  A gauge-fixing field may instead belong to a
BRST doublet and enter through a gauge fermion.  The latter can improve a
gauge-fixed boundary matrix before all auxiliary variables are removed,
but V36 did not determine whether that improvement survives on physical
metric polarisations.

This note resolves that question at a finite BRST cutoff.  The calculation
is deliberately finite dimensional, so every determinant, translation,
and contour displacement can be checked directly.  It proves that a
complete BRST quartet removes every quadratic response lying in the gauge
range.  A residual Schur term can occur only in a cokernel not paired with
the gauge ghost.  The V32 scalar displacement is in the normal
diffeomorphism range, whereas the V15 tangential transverse-tracefree
polarisation is orthogonal to the full tangential diffeomorphism symbol.
Consequently an exact quartet can regularize the former as a gauge choice
but cannot pay the latter as a physical junction impedance.

\subsection{A complete finite-dimensional quartet}
\label{sec:v3v37-quartet}

Let \(U\) and \(Q\) be finite-dimensional real Hilbert spaces.  The
coordinate \(u\in U\) parametrizes a gauge orbit and \(q\in Q\) is held
fixed as a BRST-closed physical coordinate.  Let \(c\) and \(\bar c\) be
Grassmann-odd copies of \(U\), let \(b\in U\) be bosonic, and define
\begin{equation}
 su=c,\qquad sc=0,\qquad
 s\bar c=ib,\qquad sb=0,\qquad sq=0.
\label{eq:v3v37-BRST}
\end{equation}
Thus \(s^2=0\).  Let \(M:U\to U\) be positive self-adjoint, let
\(L:U\to U\) be invertible, let \(J:Q\to U\), and let \(a>0\).  With the
ordinary complex-bilinear extension of the real inner product, set
\begin{equation}
 \Psi_q
 =
 \left\langle\bar c,\,
 M\left(Lu+Jq-\frac{ia}{2}b\right)\right\rangle .
\label{eq:v3v37-gauge-fermion}
\end{equation}

\begin{lemma}[Exact gauge-fixed action]
\label{lem:v3v37-exact-action}
The BRST variation of \eqref{eq:v3v37-gauge-fermion} is
\begin{equation}
 S_q=s\Psi_q
 =
 i\langle b,M(Lu+Jq)\rangle
 +\frac a2\langle b,Mb\rangle
 -\langle\bar c,MLc\rangle .
\label{eq:v3v37-action}
\end{equation}
In particular, the dependence on the BRST-closed variable \(q\) is itself
part of an \(s\)-exact action.
\end{lemma}

\begin{proof}
The graded Leibniz rule gives
\[
 s\langle\bar c,MF\rangle
 =
 \langle s\bar c,MF\rangle
 -
 \langle\bar c,M\,sF\rangle
\]
when \(F\) is even.  Apply this with
\(F=Lu+Jq-ia b/2\), use \eqref{eq:v3v37-BRST}, and note that
\(sF=Lc\).  The multiplier term is
\(\langle ib,M(-ia b/2)\rangle=a\langle b,Mb\rangle/2\).
This is \eqref{eq:v3v37-action}.
\end{proof}

Fix compatible Lebesgue and Berezin orientations and define the cutoff
integral
\begin{equation}
 Z(q)=
 \int_{U\times U}
 \int
 e^{-S_q}\,
 \mathrm du\,\mathrm db\,
 \mathrm D\bar c\,\mathrm Dc .
\label{eq:v3v37-Z}
\end{equation}
The \(b\)-contour is the real subspace, displaced in the proof below by
Cauchy's theorem for the entire Gaussian.

\begin{theorem}[Quartet cancellation]
\label{thm:v3v37-quartet-cancellation}
Under the preceding hypotheses, \(Z(q)\) is independent of \(q\), \(J\),
\(M\), and \(a\), up to the fixed overall normalization and the orientation
sign of \(L\).  More precisely, integrating all four members of the
quartet gives
\begin{equation}
 Z(q)=C_{\dim U}\,\operatorname{sgn}\det L,
\label{eq:v3v37-Z-constant}
\end{equation}
where \(C_{\dim U}\) depends only on the convention for the elementary
Gaussian and Berezin measures.  Hence no nonconstant effective quadratic
form in \(q\) survives a complete quartet.
\end{theorem}

\begin{proof}
Write \(F=Lu+Jq\).  Completing the square gives
\begin{equation}
 \frac a2\langle b,Mb\rangle+i\langle b,MF\rangle
 =
 \frac a2
 \left\langle b+\frac{i}{a}F,\,
 M\left(b+\frac{i}{a}F\right)\right\rangle
 +\frac1{2a}\langle F,MF\rangle .
\label{eq:v3v37-complete-square}
\end{equation}
The allowed contour displacement therefore makes the \(b\)-integral a
constant times
\[
 a^{-n/2}(\det M)^{-1/2}
 \exp\left[-\frac1{2a}\langle Lu+Jq,M(Lu+Jq)\rangle\right],
 \qquad n=\dim U.
\]
Now translate \(v=Lu+Jq\).  This removes \(q\) and \(J\), while
\(\mathrm du=|\det L|^{-1}\mathrm dv\).  The remaining \(v\)-Gaussian
contributes a constant times
\(a^{n/2}(\det M)^{-1/2}\).  The bosonic factor is consequently
\[
 C_n(\det M)^{-1}|\det L|^{-1}.
\]
The Berezin integral is \(\det(ML)=\det M\,\det L\), with a convention
dependent common sign.  Multiplication proves
\eqref{eq:v3v37-Z-constant}.
\end{proof}

\begin{corollary}[An order-one gauge kernel still cancels]
\label{cor:v3v37-order-one-cancels}
Theorem~\ref{thm:v3v37-quartet-cancellation} remains valid at every finite
spectral or lattice cutoff when \(M\) is the positive cutoff matrix of an
order-one edge operator, including the half-cylinder
Dirichlet-to-Neumann matrices constructed in V33--V35.  Its eigenvalues
may improve the gauge-fixed Hessian, but they do not leave a physical
\(q\)-dependent determinant or impedance after the quartet is integrated.
\end{corollary}

\begin{proof}
At a finite cutoff, positivity makes \(M\) an ordinary invertible positive
matrix.  The proof of Theorem~\ref{thm:v3v37-quartet-cancellation} uses no
bound on its eigenvalues and cancels \(\det M\) algebraically.
\end{proof}

\subsection{What an unmatched row leaves behind}
\label{sec:v3v37-cokernel}

The preceding theorem requires a complete square gauge system: after
gauge stabilizers are removed, the linearized gauge condition is an
isomorphism from gauge coordinates to antighost rows.  A useful regression
is obtained by dropping that requirement at the bosonic level.

Let \(U,Y,Q\) be finite-dimensional Hilbert spaces, let \(M:Y\to Y\) be
positive self-adjoint, and let \(L:U\to Y\), \(J:Q\to Y\).  Equip \(Y\)
with
\begin{equation}
 \langle y_1,y_2\rangle_M=\langle y_1,My_2\rangle,
 \qquad
 \|y\|_M^2=\langle y,y\rangle_M,
\label{eq:v3v37-M-inner}
\end{equation}
and let \(P_M\) be the \(M\)-orthogonal projection onto
\((\operatorname{Ran}L)^{\perp_M}\).

\begin{proposition}[Cokernel Schur remainder]
\label{prop:v3v37-cokernel}
After quotienting \(\ker L\), integration of the nonzero gauge coordinates
in
\begin{equation}
 \exp\left[-\frac1{2a}\|Lu+Jq\|_M^2\right]
\label{eq:v3v37-bosonic-square}
\end{equation}
leaves, up to a \(q\)-independent factor,
\begin{equation}
 \exp\left[-S_{\rm rem}(q)\right],
 \qquad
 S_{\rm rem}(q)=\frac1{2a}\|P_MJq\|_M^2.
\label{eq:v3v37-remainder}
\end{equation}
The remainder vanishes for every \(q\) if and only if
\begin{equation}
 \operatorname{Ran}J\subseteq\operatorname{Ran}L.
\label{eq:v3v37-range-condition}
\end{equation}
\end{proposition}

\begin{proof}
Decompose
\[
 Jq=(I-P_M)Jq+P_MJq
\]
into \(\operatorname{Ran}L\) and its \(M\)-orthogonal complement.  On
\(U/\ker L\), the map \(L\) is an isomorphism onto
\(\operatorname{Ran}L\), so translation of the range coordinate removes
\((I-P_M)Jq\).  Orthogonality gives
\[
 \|Lu+Jq\|_M^2
 =
 \|Lu+(I-P_M)Jq\|_M^2+\|P_MJq\|_M^2.
\]
The first term integrates to a constant and the second proves
\eqref{eq:v3v37-remainder}.  Finally \(P_MJ=0\) is equivalent to
\eqref{eq:v3v37-range-condition}.
\end{proof}

\begin{firewall}[A cokernel remainder is not a complete quartet]
\label{fw:v3v37-unmatched}
If \(P_MJ\ne0\), the antighost-row space contains directions not paired
with the gauge derivative \(L\).  The rectangular ghost form has
unsaturated Grassmann directions, or an equal-dimensional singular form
has ghost zero modes.  A nonzero expression
\eqref{eq:v3v37-remainder} can become a valid action only after those
directions receive additional physical variables, additional gauge
generators with their ghosts, or a separate measure prescription.
Calling the remainder BRST exact without completing that complex would
misapply Theorem~\ref{thm:v3v36-Ward}.
\end{firewall}

\subsection{Application to the V32 scalar displacement}
\label{sec:v3v37-v32}

At a nonzero tangential covector \(\xi\), let \(u\) be the boundary
amplitude of a normal infinitesimal diffeomorphism.  The flat stable
extension used in Proposition~\ref{prop:v3v32-normal-displacement} has
principal metric image
\begin{equation}
 H_{rr}=-2|\xi|u,\qquad
 H_{ra}=i\xi_a u,\qquad
 H_{ab}=0.
\label{eq:v3v37-normal-image}
\end{equation}
The sign agrees with \eqref{eq:v3v32-X} after identifying \(u=X_r|_N\).

\begin{proposition}[The scalar line is an invertible gauge image]
\label{prop:v3v37-scalar-gauge}
For every \(\xi\ne0\), the scalar symbol
\begin{equation}
 L_\perp(\xi):u\longmapsto H_{rr}=-2|\xi|u
\label{eq:v3v37-Lperp}
\end{equation}
is an isomorphism of one-dimensional complex amplitude spaces.  Therefore
any cutoff quadratic gauge row of the form
\begin{equation}
 \frac1{2a}
 \left\|L_\perp(\xi)u+J_\perp(\xi)q\right\|_{M(\xi)}^2
\label{eq:v3v37-scalar-square}
\end{equation}
loses all \(q\)-dependence after the normal gauge coordinate and its
complete quartet are integrated.
\end{proposition}

\begin{proof}
The inverse of \(L_\perp(\xi)\) is multiplication by
\(-1/(2|\xi|)\).  Apply
Theorem~\ref{thm:v3v37-quartet-cancellation} mode by mode at the finite
cutoff.  Equivalently,
\(\operatorname{Ran}L_\perp\) is the entire scalar row, so
\(P_M=0\) in Proposition~\ref{prop:v3v37-cokernel}.
\end{proof}

\begin{assessment}[The precise status of the V33 Stückelberg pair]
\label{ass:v3v37-V33-status}
If normal boundary displacements are admitted as gauge transformations,
the V33 variables \(U,C\) can complete their principal scalar doublet and
an \(s\)-exact edge kernel can select a complementing gauge.  Its positive
square is then removed by Proposition~\ref{prop:v3v37-scalar-gauge} from
normalized BRST-closed amplitudes.  If normal boundary displacements are
excluded because the boundary location and Brown--York charge are
physical, Proposition~\ref{prop:v3v32-normal-displacement} is not a gauge
orbit in the allowed domain; the V33 impedance must then be derived from
a physical boundary equation and cannot be justified as a disposable
quartet.  These are different boundary theories.
\end{assessment}

\subsection{The tangential transverse-tracefree cokernel}
\label{sec:v3v37-TT}

Let \(E_\xi=\operatorname{Sym}^2(\mathbb C^3)\) with the Euclidean fibre
inner product.  Define the physical transverse-tracefree symbol space
\begin{equation}
 \mathcal T_\xi
 =
 \left\{
 A_{ab}\in E_\xi:
 \delta^{ab}A_{ab}=0,\quad
 \xi^aA_{ab}=0
 \right\}.
\label{eq:v3v37-TT-space}
\end{equation}
For a tangential diffeomorphism amplitude \(v_b\), its tangential metric
symbol is
\begin{equation}
 (R_\xi v)_{ab}=i(\xi_av_b+\xi_bv_a).
\label{eq:v3v37-Rxi}
\end{equation}
A normal diffeomorphism has no \(H_{ab}\) contribution at the flat
principal boundary symbol.

\begin{lemma}[Orthogonality of physical TT data and diffeomorphisms]
\label{lem:v3v37-TT-orthogonal}
For \(\xi\ne0\),
\begin{equation}
 \mathcal T_\xi\subseteq
 \bigl(\operatorname{Ran}R_\xi\bigr)^\perp .
\label{eq:v3v37-TT-orthogonal}
\end{equation}
In fact, for every \(A\in\mathcal T_\xi\) and every \(v\),
\begin{equation}
 \langle A,R_\xi v\rangle
 =2i\,\overline{A}^{ab}\xi_a v_b=0.
\label{eq:v3v37-TT-pairing}
\end{equation}
\end{lemma}

\begin{proof}
Use the symmetry of \(A\) to combine the two terms in
\eqref{eq:v3v37-Rxi}.  The resulting contraction vanishes by
\(\xi_a\overline A^{ab}=0\).  The trace condition identifies the
tracefree physical fibre but is not needed for the orthogonality itself.
\end{proof}

\begin{theorem}[BRST-exact no-payment theorem for the V15 block]
\label{thm:v3v37-no-payment}
Fix \(\xi\ne0\) and a finite cutoff.  Suppose an auxiliary edge sector is
a complete BRST quartet for the boundary diffeomorphism symbol, and its
clock- and speed-dependent quadratic kernel enters only through a gauge
fermion of the form \eqref{eq:v3v37-gauge-fermion}.  Then complete
elimination of that quartet induces no nonzero quadratic impedance on
\(\mathcal T_\xi\).  In particular, it cannot supply the positive
order-one tracefree glancing response required by the V15--V16 junction
block.

A nonzero positive form on \(\mathcal T_\xi\) can survive only if at least
one of the following changes is made:
\begin{enumerate}
 \item physical boundary or edge degrees of freedom are retained;
 \item the physical junction action is changed directly; or
 \item the gauge symmetry is enlarged so that the former TT direction is
       no longer BRST cohomology.
\end{enumerate}
The third alternative changes the physical complex and requires its own
constraint and charge analysis.
\end{theorem}

\begin{proof}
Choose local linear coordinates consisting of gauge-orbit variables \(u\)
and representatives \(q\) of the quotient.  Proper gauge fixing makes the
derivative \(L\) from \(u\) to the antighost rows invertible after
stabilizers are removed.  Any dependence of the gauge fermion on the TT
coordinate is then the map \(Jq\) in
\eqref{eq:v3v37-gauge-fermion}.  Theorem~\ref{thm:v3v37-quartet-cancellation}
shows that translation of \(Lu\) removes \(Jq\), while the multiplier and
ghost determinants cancel.  Thus the reduced normalized integral has no
TT quadratic kernel from that quartet.

The same conclusion is visible geometrically.  By
Lemma~\ref{lem:v3v37-TT-orthogonal}, no boundary diffeomorphism symbol
reaches \(\mathcal T_\xi\).  If a proposed row does leave the remainder
\eqref{eq:v3v37-remainder} on \(\mathcal T_\xi\), it lies in the cokernel
of the gauge derivative.  Firewall~\ref{fw:v3v37-unmatched} shows that it
is not supplied by a complete quartet.  Making it well defined therefore
requires retained physical data or new gauge generators.  The positive
order-one response demanded at V15 glancing is precisely such a
nonvanishing form, proving the claim.
\end{proof}

\begin{corollary}[A physical V35 bath does not cancel]
\label{cor:v3v37-physical-bath}
If the subluminal V35 half-space field is retained as a physical
BRST-closed field and coupled to the tracefree metric fibre, its
Dirichlet-to-Neumann form is outside the hypotheses of
Theorem~\ref{thm:v3v37-no-payment}.  Its order-one response can therefore
survive, together with its energy, state space, stress tensor, and
reflection-positive Euclidean covariance.
\end{corollary}

\begin{proof}
The physical bath action is BRST closed by covariance but was not written
as \(s\Psi\); this is the distinction established in
Assessment~\ref{ass:v3v36-closed-not-exact}.  There is no gauge-coordinate
translation and ghost determinant that cancels its physical Gaussian
determinant.
\end{proof}

\subsection{Why elimination order can be misleading}
\label{sec:v3v37-order}

\begin{proposition}[The apparent positive square is an intermediate
quantity]
\label{prop:v3v37-elimination-order}
For the complete quartet \eqref{eq:v3v37-action}, integrating \(b\) first
produces the positive form
\begin{equation}
 \frac1{2a}\|Lu+Jq\|_M^2.
\label{eq:v3v37-apparent-square}
\end{equation}
Integrating \(u\) next removes its entire component in
\(\operatorname{Ran}L\), and the ghost determinant cancels the associated
Jacobian and \(M\)-determinant.  Hence positivity of
\eqref{eq:v3v37-apparent-square} alone is not evidence for a physical
spectral gap.
\end{proposition}

\begin{proof}
The first statement is \eqref{eq:v3v37-complete-square}.  The translation
and determinant cancellations are the remaining steps in the proof of
Theorem~\ref{thm:v3v37-quartet-cancellation}.  Proposition
\ref{prop:v3v37-cokernel} gives the corresponding range/cokernel
decomposition when the row is not complete.
\end{proof}

\begin{firewall}[Complementing margin and physical gap are distinct]
\label{fw:v3v37-two-gaps}
A gauge-fixed boundary operator may acquire a uniform complementing margin
from an \(s\)-exact square, and that improvement can be essential for
constructing propagators.  The quartet calculation says that the same
square does not become a new positive term in the physical effective
action.  These two uses of the word \emph{gap} must remain separate.
Strong ellipticity of a chosen gauge and coercivity of the reduced
physical junction dynamics require separate proofs.
\end{firewall}

\subsection{V37 conclusion and V38 target}
\label{sec:v3v37-conclusion}

V37 completes the finite-cutoff determinant calculation requested at the
end of V36.  A proper quartet gives an exact cancellation: the multiplier
creates a positive square, the gauge coordinate translates away every
physical argument in the gauge range, and the ghosts cancel the
Jacobian.  An unmatched positive remainder is the \(M\)-orthogonal
cokernel term \(\|P_MJq\|_M^2/(2a)\); it is not a complete quartet.

The V32 scalar boundary displacement has invertible symbol
\(-2|\xi|\) along the normal gauge line.  It may be fixed by a BRST-exact
edge gauge if that boundary motion is admitted as gauge, but the resulting
impedance is not a physical determinant.  The V15 transverse-tracefree
fibre is orthogonal to the diffeomorphism image and remains physical.
Therefore a disposable clock-dependent quartet cannot pay its glancing
loss.

V38 must now classify the first honest physical alternative.  It should
construct the most economical local, reflection-positive boundary medium
whose eliminated tracefree response is passive and order one, compute its
stress contribution to the Einstein junction equation, and determine
whether Standard Model boundary fields already provide that response.
If they do not, the programme must record the added medium as physical
field content or return upstream and replace the V15 junction condition.
\clearpage
\section{V38: a physical material tensor bath and its tracefree junction response}
\label{sec:v3v38}

\subsection{Why a material reference is necessary}
\label{sec:v3v38-purpose}

V37 proves that a complete BRST quartet cannot leave a physical
order-one form on the V15 transverse-tracefree quotient.  A physical bath
can leave such a form, but a minimally coupled bath with zero background
has stress quadratic in the bath amplitude and therefore has no linear
junction response.  To obtain a nonzero metric--bath Hessian from a
diffeomorphism-covariant action, the boundary medium must contain a
physical reference against which metric strain is measured.

This note supplies the smallest local reference at the frozen principal
level: the clock \(T\) of V36 and two scalar material coordinates
\(X^I\), \(I=1,2\).  Their gradients form a spatial material metric.
The difference between that metric and the clock-spatial metric has a
covariant tracefree part.  Coupling a physical half-line tensor field to
that strain produces the V35 subluminal impedance.  A direct null-symbol
calculation proves that the spatial tracefree incidence sees every
tangential TT polarization at primary glancing.

\subsection{Clock-spatial geometry and material strain}
\label{sec:v3v38-strain}

Let \((\mathcal N,\gamma)\) be a time-oriented three-dimensional timelike
boundary.  Assume the V36 clock gap
\begin{equation}
 X_T=-\gamma^{ab}\nabla_aT\nabla_bT\ge\chi^2>0
\label{eq:v3v38-clock-gap}
\end{equation}
and use
\begin{equation}
 u_a=-\frac{\nabla_aT}{\sqrt{X_T}},
 \qquad
 q_{ab}=\gamma_{ab}+u_au_b.
\label{eq:v3v38-uq}
\end{equation}
Let \((\mathcal B,G)\) be a two-dimensional Riemannian material manifold
and let \(X:\mathcal N\to\mathcal B\) have rank two on
\(\ker u^\flat\).  Impose the comoving condition
\begin{equation}
 u^a\nabla_aX^I=0.
\label{eq:v3v38-comoving}
\end{equation}
The pulled-back material metric is
\begin{equation}
 k_{ab}
 =
 G_{IJ}(X)\nabla_aX^I\nabla_bX^J.
\label{eq:v3v38-k}
\end{equation}
It is spatial by \eqref{eq:v3v38-comoving}.

Define the spatial symmetric-tracefree projector
\begin{equation}
 \mathbb P_{ab}{}^{cd}
 =
 q_{(a}{}^cq_{b)}{}^d-\frac12q_{ab}q^{cd}
\label{eq:v3v38-P}
\end{equation}
and the material strain
\begin{equation}
 e_{ab}
 =
 \mathbb P_{ab}{}^{cd}(k_{cd}-q_{cd})
 =
 \mathbb P_{ab}{}^{cd}k_{cd}.
\label{eq:v3v38-e}
\end{equation}

\begin{proposition}[Covariance and positive fibre]
\label{prop:v3v38-covariant-strain}
The tensor \(e\) is boundary-diffeomorphism covariant, satisfies
\begin{equation}
 u^ae_{ab}=0,\qquad q^{ab}e_{ab}=0,
\label{eq:v3v38-STF}
\end{equation}
and lies in a rank-two vector bundle
\(\mathcal S_u\to\mathcal N\) with positive fibre norm
\begin{equation}
 |e|_q^2=q^{ac}q^{bd}e_{ab}e_{cd}.
\label{eq:v3v38-positive-norm}
\end{equation}
It is invariant under changes of material coordinates when \(G\) is
transformed as a material tensor.
\end{proposition}

\begin{proof}
All ingredients of \eqref{eq:v3v38-e} are obtained by tensorial pullback,
contraction, and projection, so naturality proves diffeomorphism
covariance.  Contracting \eqref{eq:v3v38-P} with \(u^a\) gives zero because
\(q_a{}^cu^a=0\).  Contracting with \(q^{ab}\) gives
\(q^{cd}-\frac12(2)q^{cd}=0\), proving
\eqref{eq:v3v38-STF}.  A symmetric tensor on a two-dimensional Euclidean
space has three components and its tracefree subspace has dimension two;
the restriction of the induced \(q\)-inner product is positive.  The last
claim is the coordinate invariance of the pullback metric
\eqref{eq:v3v38-k}.
\end{proof}

\begin{assessment}[The medium is physical]
\label{ass:v3v38-physical-medium}
The scalars \(T,X^1,X^2\) may be used as coordinates where their joint
Jacobian is nonzero, but this does not make their strain a disposable
quartet.  Their values specify a material state.  In the corresponding
unitary coordinate gauge their degrees of freedom are encoded in the
metric relative to that material state, just as phonons may be represented
by strain.  Removing the scalars removes the covariant reference needed
by \eqref{eq:v3v38-e}.
\end{assessment}

\subsection{Local tensor bath}
\label{sec:v3v38-bath}

Let \(y\ge0\) be an auxiliary half-line.  A bath field
\(\Psi_{ab}(y,\cdot)\) takes values in \(\mathcal S_u\).  Let
\(\mathcal D\) denote the connection obtained by projecting the
Levi--Civita connection to \(\mathcal S_u\), and let
\(\mathcal D_u\) be the projected derivative along \(u\).  For constants
\begin{equation}
 \alpha>0,\qquad 0\le c_{\rm e}^2\le1-\delta<1,
 \qquad m_{\rm e}\ge0,
\label{eq:v3v38-parameters}
\end{equation}
take the Lorentzian quadratic bath action
\begin{align}
 S_{\rm bath}^L
 =\frac12\int_{\mathcal N}\!\sqrt{-\gamma}\,d^3x
 \int_0^\infty\!dy\,
 \bigl(
 &|\mathcal D_u\Psi|_q^2
 -|\partial_y\Psi|_q^2
 -c_{\rm e}^2q^{cd}
   \langle\mathcal D_c\Psi,\mathcal D_d\Psi\rangle_q
 \notag\\
 &-m_{\rm e}^2|\Psi|_q^2
 \bigr)
\label{eq:v3v38-L-action}
\end{align}
with the local trace incidence
\begin{equation}
 \Psi_{ab}(0,x)=\alpha e_{ab}(x).
\label{eq:v3v38-trace}
\end{equation}
The trace relation may equivalently be imposed by a local boundary
multiplier.  The Euclidean action is obtained in an adapted clock chart
by changing all four displayed quadratic terms to positive signs.

\begin{proposition}[Locality, covariance, and positive frozen energy]
\label{prop:v3v38-local-energy}
Equations \eqref{eq:v3v38-L-action}--\eqref{eq:v3v38-trace} define a local
boundary-plus-half-line system and are invariant under boundary
diffeomorphisms acting on \(\gamma,T,X,\Psi\).  At a frozen point in an
orthonormal clock frame, the bath energy is
\begin{equation}
 E_{\rm bath}
 =
 \frac12\int_{\Sigma}\!\int_0^\infty
 \left(
 |\partial_t\Psi|^2+|\partial_y\Psi|^2
 +c_{\rm e}^2|\nabla_x\Psi|^2+m_{\rm e}^2|\Psi|^2
 \right)dy\,d^2x,
\label{eq:v3v38-energy}
\end{equation}
which is nonnegative.
\end{proposition}

\begin{proof}
The action contains only first derivatives and tensorial contractions of
fields at the same point of \(\mathcal N\times[0,\infty)\); the trace
condition is also pointwise.  Naturality of \(u,q,k,\mathbb P\) and of the
projected connection proves covariance.  In a frozen clock frame,
\(\mathcal S_u\) has the positive norm
\eqref{eq:v3v38-positive-norm}, and the standard multiplier
\(\partial_t\Psi\) gives \eqref{eq:v3v38-energy}.
\end{proof}

\begin{theorem}[Reflection-positive Euclidean parent]
\label{thm:v3v38-RP}
At a finite spatial and half-line cutoff, around a static flat material
background invariant under Euclidean time reflection, the joint Gaussian
defined by the Euclidean version of
\eqref{eq:v3v38-L-action} and the time-local trace constraint
\eqref{eq:v3v38-trace} is reflection positive.  Its marginal on the
linearized strain is reflection positive.
\end{theorem}

\begin{proof}
In the flat background the two components of \(\mathcal S_u\) are copies
of the V35 Hilbert-space field, with positive spatial operator
\(-c_{\rm e}^2\Delta_x-\partial_y^2+m_{\rm e}^2\).  The reflected
covariance factorization of
Lemma~\ref{lem:v3v35-reflected-covariance} applies componentwise.  The
trace constraint is local in Euclidean time and invariant under
reflection.  The restriction and marginalization argument in
Theorem~\ref{thm:v3v35-RP} then proves both claims.
\end{proof}

\begin{firewall}[Scope of reflection positivity]
\label{fw:v3v38-RP-scope}
Theorem~\ref{thm:v3v38-RP} is a quadratic theorem around a static material
background.  Integrating the nonlinear clock and material-coordinate
fields requires a positive Euclidean action for those fields, a
reflection-invariant admissible set preserving
\eqref{eq:v3v38-clock-gap}, and control of the material rank condition.
Those continuum measure gates are not supplied by the frozen Gaussian
proof.
\end{firewall}

\subsection{Exact flat Hessian and junction source}
\label{sec:v3v38-Hessian}

Use a flat clock frame with boundary coordinates
\((t,x^1,x^2)\),
\begin{equation}
 \gamma^{(0)}=-dt^2+\delta_{ij}dx^idx^j,\qquad
 T=t,\qquad X^I=x^I,\qquad G_{IJ}=\delta_{IJ}.
\label{eq:v3v38-background}
\end{equation}
Let \(h_{ij}=\delta\gamma_{ij}\) and
\(X^I=x^I+\pi^I\).  Write
\begin{equation}
 (\mathbb P_0A)_{ij}
 =
 \frac12(A_{ij}+A_{ji})-\frac12\delta_{ij}\delta^{kl}A_{kl}.
\label{eq:v3v38-P0}
\end{equation}

\begin{lemma}[Linear material strain]
\label{lem:v3v38-linear-strain}
At the background \eqref{eq:v3v38-background},
\begin{equation}
 e_{ij}^{(1)}
 =
 \mathbb P_0
 \left(
 \partial_i\pi_j+\partial_j\pi_i-h_{ij}
 \right).
\label{eq:v3v38-linear-strain}
\end{equation}
\end{lemma}

\begin{proof}
The material metric varies by
\[
 \delta k_{ij}=\partial_i\pi_j+\partial_j\pi_i,
\]
while \(\delta q_{ij}=h_{ij}\) when the clock is held fixed in the
adapted frame.  Since \(k^{(0)}-q^{(0)}=0\), the variation of the projector
multiplies zero.  Thus only \(\mathbb P_0(\delta k-\delta q)\) remains.
\end{proof}

For Euclidean frequency \(\omega\) and spatial covector \(\xi\in
\mathbb R^2\), define
\begin{equation}
 Q_E(\omega,\xi)
 =
 \sqrt{\omega^2+c_{\rm e}^2|\xi|^2+m_{\rm e}^2}.
\label{eq:v3v38-QE}
\end{equation}

\begin{theorem}[Eliminated physical strain form]
\label{thm:v3v38-effective}
The decaying bath solution with trace
\eqref{eq:v3v38-trace} gives the Euclidean quadratic action
\begin{equation}
 S_{\rm eff}^{E,(2)}
 =
 \frac{\alpha^2}{2}
 \int
 Q_E(\omega,\xi)
 \left|\widehat e^{(1)}(\omega,\xi)\right|^2
 \,d\omega\,d^2\xi.
\label{eq:v3v38-effective}
\end{equation}
With the canonical linearized junction source defined by
\begin{equation}
 \delta S_{\rm eff}^{(2)}
 =
 \int
 \left\langle\mathcal J^{ij},\delta h_{ij}^{\rm TF}\right\rangle,
\label{eq:v3v38-J-definition}
\end{equation}
one has
\begin{equation}
 \widehat{\mathcal J}^{\,\rm TF}
 =
 -\alpha^2Q_E\,\widehat e^{(1)}.
\label{eq:v3v38-J}
\end{equation}
In material unitary gauge \(\pi=0\),
\begin{equation}
 \widehat e^{(1)}=-\widehat h^{\rm STF},
 \qquad
 \widehat{\mathcal J}^{\,\rm TF}
 =\alpha^2Q_E\,\widehat h^{\rm STF}.
\label{eq:v3v38-unitary-J}
\end{equation}
\end{theorem}

\begin{proof}
The two fibre components satisfy the scalar Poisson-extension calculation
of Theorem~\ref{thm:v3v35-E-kernel}; summing their energies gives
\eqref{eq:v3v38-effective}.  Self-adjointness of \(Q_E\) yields
\[
 \delta S_{\rm eff}^{E,(2)}
 =
 \alpha^2\langle Q_Ee^{(1)},\delta e^{(1)}\rangle.
\]
At fixed \(\pi\), Lemma~\ref{lem:v3v38-linear-strain} gives
\(\delta e^{(1)}=-\delta h^{\rm TF}\), proving
\eqref{eq:v3v38-J}.  Setting \(\pi=0\) proves
\eqref{eq:v3v38-unitary-J}.
\end{proof}

\begin{assessment}[One Hessian supplies both reciprocal couplings]
\label{ass:v3v38-one-Hessian}
The metric junction source \eqref{eq:v3v38-J} and the material displacement
equation are the two variations of the single quadratic form
\eqref{eq:v3v38-effective}.  Their cross blocks are adjoints, as required
by Proposition~\ref{prop:v3v15-Hessian-adjoint}.  The bath is therefore an
actual energy carrier rather than a damping term appended only to the
metric equation.
\end{assessment}

\subsection{Null TT incidence}
\label{sec:v3v38-null-incidence}

Let \(\zeta_a=(-\tau,\xi_1,\xi_2)\) be a nonzero null boundary covector,
\begin{equation}
 |\xi|=|\tau|\ne0,
\label{eq:v3v38-null}
\end{equation}
and let
\begin{equation}
 \mathcal T_\zeta
 =
 \left\{
 A_{ab}=A_{ba}:
 \gamma^{ab}A_{ab}=0,\quad
 \zeta^aA_{ab}=0
 \right\}.
\label{eq:v3v38-null-TT}
\end{equation}
Let \(B_\zeta A=\mathbb P_0(A_{ij})\) be the spatial tracefree
restriction.

\begin{lemma}[Spatial strain sees every null TT polarization]
\label{lem:v3v38-null-injective}
For every \(\zeta\) satisfying \eqref{eq:v3v38-null}, the map
\begin{equation}
 B_\zeta:\mathcal T_\zeta\longrightarrow
 \operatorname{Sym}_0^2(\mathbb C^2)
\label{eq:v3v38-B}
\end{equation}
is an isomorphism.  In a rotated and rescaled frame
\(\tau=1,\xi=(1,0)\), every \(A\in\mathcal T_\zeta\) has the form
\begin{equation}
 (A_{ab})=
 \begin{pmatrix}
 a&-a&b\\
 -a&a&-b\\
 b&-b&0
 \end{pmatrix},
\qquad
 B_\zeta A=
 \begin{pmatrix}
 a/2&-b\\
 -b&-a/2
 \end{pmatrix}.
\label{eq:v3v38-null-matrix}
\end{equation}
Consequently, with the amplitude norm
\(\|A\|_*^2=|a|^2+|b|^2\),
\begin{equation}
 \|B_\zeta A\|_{\rm HS}^2
 =\frac12|a|^2+2|b|^2
 \ge\frac12\|A\|_*^2.
\label{eq:v3v38-B-margin}
\end{equation}
The constant is uniform on the normalized null circle.
\end{lemma}

\begin{proof}
Rotations and nonzero rescaling reduce to the stated covector.  Since
\(\zeta^a=(1,1,0)\), transversality gives
\[
 A_{01}=-A_{00},\qquad
 A_{11}=A_{00},\qquad
 A_{12}=-A_{02}.
\]
The Lorentzian trace equation
\(-A_{00}+A_{11}+A_{22}=0\) then gives \(A_{22}=0\).
Put \(a=A_{00}\), \(b=A_{02}\); this proves the first matrix in
\eqref{eq:v3v38-null-matrix}.  Taking the spatial tracefree part gives the
second.  It vanishes only when \(a=b=0\), so the map is injective; both
spaces have dimension two, hence it is an isomorphism.  Direct
Hilbert--Schmidt contraction gives \eqref{eq:v3v38-B-margin}.  Rotation
invariance makes the constant uniform.
\end{proof}

Define the retarded bath root
\begin{equation}
 Q_L(s,\xi)
 =
 \sqrt{s^2+c_{\rm e}^2|\xi|^2+m_{\rm e}^2},
 \qquad \operatorname{Re}Q_L>0
\label{eq:v3v38-QL}
\end{equation}
for \(\operatorname{Re}s>0\).

\begin{theorem}[Physical order-one TT glancing margin]
\label{thm:v3v38-TT-margin}
Assume \(m_{\rm e}=0\) and
\(c_{\rm e}^2\le1-\delta\).  In material unitary gauge, the bath
contribution to the retarded TT junction symbol at
\(s=i\tau\), \(|\xi|=|\tau|\ne0\), is
\begin{equation}
 \mathbb Z_{\rm TT}(i\tau,\xi)
 =
 \alpha^2 B_\zeta^*Q_L(i\tau,\xi)B_\zeta.
\label{eq:v3v38-ZTT}
\end{equation}
It obeys the homogeneous lower bound
\begin{equation}
 \|\mathbb Z_{\rm TT}A\|_*
 \ge
 \frac{\alpha^2\sqrt\delta}{2}\,|\tau|\,\|A\|_*
\qquad(A\in\mathcal T_\zeta),
\label{eq:v3v38-TT-margin}
\end{equation}
after using the fixed norms induced by
\eqref{eq:v3v38-null-matrix}.  Thus no null TT polarization is left
uncoupled.
\end{theorem}

\begin{proof}
Theorem~\ref{thm:v3v35-subluminal} gives
\[
 |Q_L(i\tau,\xi)|
 =
 \sqrt{1-c_{\rm e}^2}\,|\tau|
 \ge\sqrt\delta\,|\tau|.
\]
The scalar root acts identically on the two spatial STF components.
Lemma~\ref{lem:v3v38-null-injective} gives
\(B_\zeta^*B_\zeta\ge\frac12I\) in the displayed amplitude norm.
Taking moduli in \eqref{eq:v3v38-ZTT} proves
\eqref{eq:v3v38-TT-margin}.
\end{proof}

\begin{firewall}[The canonical junction sign still belongs to gravity]
\label{fw:v3v38-junction-sign}
Theorem~\ref{thm:v3v38-TT-margin} proves a nonzero physical incidence and
the magnitude of the passive response.  The sign with which
\(\mathcal J^{\rm TF}\) enters the outward Brown--York equation depends on
the GHY orientation and on whether the material action is placed on the
chosen side of the junction.  The complete Einstein--medium variation
must reproduce the energy-conserving sign of
\eqref{eq:v3v16-reservoir-junction}; it may not be selected after the
symbol is eliminated.
\end{firewall}

\subsection{Diffeomorphism Ward identity and Codazzi}
\label{sec:v3v38-Ward}

Let \(\mathcal E_T,\mathcal E_I,\mathcal E_\Psi\) denote the Euler
derivatives of the complete local medium action with respect to
\(T,X^I,\Psi\), including a multiplier for
\eqref{eq:v3v38-trace}.  Define its metric stress by
\begin{equation}
 \delta_\gamma S_{\rm med}
 =
 \frac12\int_{\mathcal N}\sqrt{-\gamma}\,
 \mathcal S^{ab}\delta\gamma_{ab}.
\label{eq:v3v38-stress-definition}
\end{equation}

\begin{proposition}[On-shell boundary stress conservation]
\label{prop:v3v38-Ward}
Diffeomorphism covariance of the local parent implies the Noether identity
\begin{equation}
 \nabla_a\mathcal S^a{}_b
 =
 \mathcal E_T\nabla_bT
 +\mathcal E_I\nabla_bX^I
 +\mathcal W_b(\mathcal E_\Psi,\Psi),
\label{eq:v3v38-Ward}
\end{equation}
where \(\mathcal W_b\) is linear in the bath Euler derivative and its
covariant derivatives.  Hence
\begin{equation}
 \nabla_a\mathcal S^a{}_b=0
\label{eq:v3v38-conserved}
\end{equation}
when all medium equations hold.  The resulting junction source is
compatible with the homogeneous linearized Codazzi condition
\eqref{eq:v3v15-linear-flux}.
\end{proposition}

\begin{proof}
Vary the covariant action under a compactly supported infinitesimal
boundary diffeomorphism generated by \(v^b\).  Substitute
\(\delta\gamma_{ab}=2\nabla_{(a}v_{b)}\),
\(\delta T=v^a\nabla_aT\),
\(\delta X^I=v^a\nabla_aX^I\), and the tensor Lie derivative of \(\Psi\).
Integrating derivatives of \(v\) by parts and using arbitrariness of
\(v\) gives \eqref{eq:v3v38-Ward}; the precise
\(\mathcal W_b\) depends only on the standard integrations in the tensor
Lie derivative and vanishes with \(\mathcal E_\Psi\).
Equation \eqref{eq:v3v38-conserved} follows on shell.  With zero normal
bulk matter flux, it is exactly
\(\partial^aS_{ab}^{(1)}=0\) in
\eqref{eq:v3v15-linear-flux}.
\end{proof}

\subsection{Classical Standard Model vacuum test}
\label{sec:v3v38-SM-test}

Consider the ordinary minimally coupled classical Standard Model at a
Lorentz-invariant vacuum,
\begin{equation}
 F_{\mu\nu}=0,\qquad D_\mu H=0,\qquad
 \psi_{\rm fermion}=0,\qquad V'(H)=0.
\label{eq:v3v38-SM-vacuum}
\end{equation}

\begin{theorem}[No tree-level order-one SM vacuum impedance]
\label{thm:v3v38-SM-no-impedance}
At the background \eqref{eq:v3v38-SM-vacuum}, the first variation of the
classical Standard Model stress with respect to matter perturbations
vanishes.  Its pure metric linearization is algebraic and proportional to
the vacuum potential.  Consequently the classical SM vacuum contributes
no degree-one tangential symbol to the linearized TT junction equation and
cannot realize \eqref{eq:v3v38-ZTT}.
\end{theorem}

\begin{proof}
The Yang--Mills stress is quadratic in \(F\), the Higgs kinetic stress is
quadratic in \(DH\), and the fermion stress is bilinear in the fermion.
Their first variations vanish when the corresponding background fields
in \eqref{eq:v3v38-SM-vacuum} vanish.  The linear term from the Higgs
potential is zero by \(V'(H)=0\).  Varying the remaining constant vacuum
value \(V(H_0)\) with respect to the metric produces only a
cosmological/tension term proportional to the metric perturbation, with no
tangential derivative.  Its symbol has degree zero.  An order-one
Dirichlet-to-Neumann response such as \eqref{eq:v3v38-ZTT} is therefore
absent.
\end{proof}

\begin{corollary}[What must be added or demonstrated]
\label{cor:v3v38-SM-alternatives}
Within the classical vacuum sector, the material fields \(T,X^I\) and the
tensor bath of V38 are additional physical content.  Avoiding that
addition requires either a nonvacuum SM background with a proved
nondegenerate material frame and the same positive symbol, a quantum SM
stress response with a proved reflection-positive local parent and
order-one TT spectral weight, or a change of the gravitational junction
condition.
\end{corollary}

\begin{proof}
Theorem~\ref{thm:v3v38-SM-no-impedance} excludes the standard tree-level
vacuum source.  The listed alternatives are the remaining sources allowed
by Theorem~\ref{thm:v3v37-no-payment}.
\end{proof}

\begin{firewall}[No quantum conclusion from the classical test]
\label{fw:v3v38-quantum-open}
Vacuum polarization and boundary states can generate nonlocal stress
kernels after quantum fields are integrated out.  V38 does not calculate
their spectral measure and therefore does not exclude an order-one
component from a specified quantum boundary state.  Such a claim requires
the actual stress two-point function, its Ward identities, its
renormalized local counterterms, and a uniform high-frequency estimate.
\end{firewall}

\subsection{V38 conclusion and V39 target}
\label{sec:v3v38-conclusion}

V38 constructs a physical local parent for the missing TT impedance.
A timelike clock and two material-coordinate scalars define a covariant
rank-two positive strain bundle.  A subluminal half-line tensor field
coupled to that strain has a reflection-positive flat Euclidean Gaussian,
a conserved enlarged Lorentzian energy, and the eliminated response
\(\alpha^2Q_L\).  The spatial tracefree incidence is an isomorphism on the
null TT fibre and supplies the explicit normalized margin
\(\alpha^2\sqrt\delta/2\).

The construction also identifies its cost.  The clock, material
coordinates, and bath are physical state and field content; their
response does not cancel as a quartet.  The ordinary Lorentz-invariant
classical SM vacuum has no linear derivative stress and does not supply
this content.

V39 should compute the complete coupled flat boundary symbol including
the material displacements rather than fixing unitary gauge at the start.
It must prove that the new phonon rows introduce no unstable surface mode,
recover the TT margin on the diffeomorphism quotient, and track the exact
energy-conserving Brown--York sign.  Failure of that determinant test
would force another upstream change before the material construction can
be admitted into the continuum programme.
\clearpage
\section{V39: the coupled TT--material--bath boundary symbol}
\label{sec:v3v39}

\subsection{The screening issue left by V38}
\label{sec:v3v39-purpose}

V38 computed the tracefree response in material unitary gauge.  A complete
symbol test must retain the material-coordinate perturbation
\(\pi^I\).  This is essential because, at a nonzero spatial covector, the
two-dimensional tracefree symmetrized-gradient map from \(\pi^I\) to
spatial strain has the same rank as the spatial STF fibre.  If the
material coordinates could relax algebraically at every time frequency,
they could set the strain to zero and screen the bath.

The covariant construction already contains the required distinction:
the labels are comoving,
\(u^a\nabla_aX^I=0\).  Its multiplier equation prevents a time-dependent
material label from cancelling a time-dependent gravitational wave.  This
note retains that constraint and its conjugate equation, fixes the
energy-conserving Brown--York sign, and proves absence of unstable or
closed-frequency surface roots in the complete TT--material--bath block
away from the static origin.

\subsection{The spatial strain symbol is onto}
\label{sec:v3v39-screening}

In the flat frame \eqref{eq:v3v38-background}, define
\begin{equation}
 (D_\xi\pi)_{ij}
 =
 \mathbb P_0\bigl(i\xi_i\pi_j+i\xi_j\pi_i\bigr),
 \qquad
 D_\xi:\mathbb C^2\longrightarrow
 \operatorname{Sym}_0^2(\mathbb C^2).
\label{eq:v3v39-D}
\end{equation}

\begin{lemma}[Invertibility of the two-dimensional strain gradient]
\label{lem:v3v39-D-invertible}
For every \(\xi\ne0\), \(D_\xi\) is an isomorphism.  In a frame
\(\xi=(k,0)\), \(k>0\),
\begin{equation}
 D_\xi
 \begin{pmatrix}\pi_1\\ \pi_2\end{pmatrix}
 =
 \begin{pmatrix}
 ik\pi_1&ik\pi_2\\
 ik\pi_2&-ik\pi_1
 \end{pmatrix}.
\label{eq:v3v39-D-matrix}
\end{equation}
Moreover,
\begin{equation}
 \|D_\xi\pi\|_{\rm HS}^2
 =2|\xi|^2|\pi|^2.
\label{eq:v3v39-D-norm}
\end{equation}
\end{lemma}

\begin{proof}
Equation \eqref{eq:v3v39-D-matrix} follows by taking the tracefree part of
\(i\xi_i\pi_j+i\xi_j\pi_i\).  Its two independent entries vanish only if
\(\pi_1=\pi_2=0\), and the domain and target both have dimension two.
Direct contraction gives \eqref{eq:v3v39-D-norm}; rotation invariance
extends it to every \(\xi\).
\end{proof}

\begin{counterexample}[Unconstrained quasistatic screening]
\label{cex:v3v39-screening}
If \(\pi\) is an unconstrained variable with no independent time equation,
then for every \(\xi\ne0\) and every spatial STF metric amplitude \(H\)
there is a unique
\begin{equation}
 \pi=D_\xi^{-1}H
\label{eq:v3v39-screening-pi}
\end{equation}
such that the linear strain
\(e^{(1)}=D_\xi\pi-H\) vanishes.  The eliminated bath form
\(\alpha^2\langle e,Qe\rangle/2\) then gives no metric impedance after
\(\pi\) is minimized.
\end{counterexample}

\begin{proof}
Use Lemma~\ref{lem:v3v39-D-invertible} and substitute
\eqref{eq:v3v39-screening-pi}.
\end{proof}

\begin{assessment}[Why a reference must have an evolution law]
\label{ass:v3v39-reference-law}
A material metric is not an anchor merely because it is written with
scalar fields.  Its label evolution, inertia, elasticity, or prescribed
state must prevent arbitrary pointwise relaxation.  V38 uses the
covariant comoving law.  Other physical media may replace it by a
hyperbolic phonon equation, but their dispersion must then be included in
the boundary determinant.
\end{assessment}

\subsection{Comoving multiplier block}
\label{sec:v3v39-comoving}

Impose \eqref{eq:v3v38-comoving} by a boundary multiplier
\(\lambda_I\):
\begin{equation}
 S_{\rm com}
 =
 \int_{\mathcal N}\sqrt{-\gamma}\,
 \lambda_Iu^a\nabla_aX^I\,d^3x.
\label{eq:v3v39-Scom}
\end{equation}
At the flat background and in the TT block, variations of the clock flow
do not enter the material-label equation.  For a Laplace--Fourier mode
\(e^{st+i\xi\cdot x}\),
\begin{equation}
 s\widehat\pi=0.
\label{eq:v3v39-comoving-symbol}
\end{equation}

Let \(\Lambda^{ij}\) denote the boundary force conjugate to the trace
constraint
\(\Psi(0)=\alpha(D_\xi\pi-B_\zeta A)\).  The principal material equation
has the form
\begin{equation}
 -s\widehat\lambda+D_\xi^*\widehat\Lambda=0,
\label{eq:v3v39-material-symbol}
\end{equation}
up to an inessential fixed sign in the definition of \(\lambda\).

\begin{proposition}[No nonstatic material screening]
\label{prop:v3v39-no-screening}
For \(s\ne0\), the multiplier rows
\eqref{eq:v3v39-comoving-symbol}--%
\eqref{eq:v3v39-material-symbol} uniquely give
\begin{equation}
 \widehat\pi=0,\qquad
 \widehat\lambda=s^{-1}D_\xi^*\widehat\Lambda.
\label{eq:v3v39-pi-lambda}
\end{equation}
They therefore introduce no additional compatibility condition on the TT
metric or bath amplitudes and cannot create a root with
\(\operatorname{Re}s>0\).  Their determinant factor is a nonzero power of
\(s\).
\end{proposition}

\begin{proof}
Equation \eqref{eq:v3v39-comoving-symbol} gives \(\widehat\pi=0\).
The remaining equation then solves uniquely for \(\widehat\lambda\).
In an ordering with the comoving row first and the material Euler row
second, the \((\pi,\lambda)\) principal block is triangular after the
constraint force is regarded as an already determined boundary amplitude;
both diagonal maps are scalar multiples of \(sI_2\).
\end{proof}

\begin{firewall}[Static material relabellings remain]
\label{fw:v3v39-static}
At \(s=0\), the comoving row does not determine \(\pi\), and
Counterexample~\ref{cex:v3v39-screening} applies at nonzero \(\xi\) unless
a static elastic energy or a prescribed material configuration is added.
V39 proves the nonzero-frequency TT glancing estimate and excludes
unstable modes; it does not prove an elliptic zero-frequency estimate for
the material labels.  That static sector must be quotiented by an
identified material relabelling symmetry or stabilized by a physical
elastic modulus.
\end{firewall}

\subsection{Local equations and the energy-selected sign}
\label{sec:v3v39-energy-sign}

Consider one finite-dimensional TT amplitude space with bulk field
\(A(t,x,z)\), \(x>0\), boundary spatial coordinate \(z\in\mathbb R^2\),
and bath \(\Psi(t,z,y)\), \(y>0\).  Let \(B(D)\) be the spatial STF
incidence of V38.  The equations on the nonstatic comoving sector are
\begin{align}
 A_{tt}-A_{xx}-\Delta_zA&=0,
\label{eq:v3v39-bulk-wave}\\
 \Psi_{tt}-\Psi_{yy}-c_{\rm e}^2\Delta_z\Psi
 +m_{\rm e}^2\Psi&=0,
\label{eq:v3v39-bath-wave}\\
 \Psi(0)&=-\alpha B(D)A(0),
\label{eq:v3v39-trace}\\
 \ell(\partial_t^2-\Delta_z)A
 -\partial_xA+\alpha B(D)^*\partial_y\Psi(0)&=0,
 \qquad \ell\ge0.
\label{eq:v3v39-junction}
\end{align}
The minus sign in \eqref{eq:v3v39-trace} is the unitary-gauge identity
\(e^{(1)}=-B(D)A\).

Define
\begin{align}
 E_{\rm g}
 &=
 \frac12\int_{x>0}
 \left(|A_t|^2+|A_x|^2+|\nabla_zA|^2\right)dx\,d^2z,
\label{eq:v3v39-Eg}\\
 E_\partial
 &=
 \frac\ell2\int_{x=0}
 \left(|A_t|^2+|\nabla_zA|^2\right)d^2z,
\label{eq:v3v39-Epartial}\\
 E_{\rm R}
 &=
 \frac12\int_{y>0}
 \left(
 |\Psi_t|^2+|\Psi_y|^2
 +c_{\rm e}^2|\nabla_z\Psi|^2+m_{\rm e}^2|\Psi|^2
 \right)dy\,d^2z.
\label{eq:v3v39-ER}
\end{align}

\begin{theorem}[The Brown--York sign is fixed by energy]
\label{thm:v3v39-energy}
Smooth finite-energy solutions of
\eqref{eq:v3v39-bulk-wave}--%
\eqref{eq:v3v39-junction} satisfy
\begin{equation}
 \frac d{dt}(E_{\rm g}+E_\partial+E_{\rm R})=0.
\label{eq:v3v39-energy}
\end{equation}
Reversing only the sign of the last term in
\eqref{eq:v3v39-junction} fails this cancellation.
\end{theorem}

\begin{proof}
Integration by parts gives
\begin{align*}
 E_{\rm g}'&=-\langle A_t,\partial_xA\rangle_{x=0},\\
 E_\partial'&=
 \ell\langle A_t,(\partial_t^2-\Delta_z)A\rangle_{x=0},\\
 E_{\rm R}'&=-\langle\Psi_t,\partial_y\Psi\rangle_{y=0}.
\end{align*}
Differentiate \eqref{eq:v3v39-trace} in time and use adjointness of \(B\):
\[
 E_{\rm R}'
 =
 \alpha\langle A_t,B^*\partial_y\Psi\rangle_{x=0}.
\]
The sum is the pairing of \(A_t\) with the left side of
\eqref{eq:v3v39-junction}, proving
\eqref{eq:v3v39-energy}.  Changing only the indicated sign makes the two
bath flux terms add instead of cancel.
\end{proof}

\subsection{The eliminated coupled symbol}
\label{sec:v3v39-symbol}

For \(\operatorname{Re}s>0\), define the outgoing roots
\begin{equation}
 \kappa(s,\xi)=\sqrt{s^2+|\xi|^2},
 \qquad
 \rho(s,\xi)=
 \sqrt{s^2+c_{\rm e}^2|\xi|^2+m_{\rm e}^2},
 \qquad
 \operatorname{Re}\kappa,\operatorname{Re}\rho>0.
\label{eq:v3v39-roots}
\end{equation}
The stable solutions are
\begin{equation}
 \widehat A(x)=e^{-\kappa x}A_0,
 \qquad
 \widehat\Psi(y)=
 -\alpha e^{-\rho y}B_\zeta A_0.
\label{eq:v3v39-stable}
\end{equation}

\begin{proposition}[Exact TT boundary matrix]
\label{prop:v3v39-boundary-matrix}
For \(s\ne0\), elimination of \(\pi,\lambda,\Psi\) gives
\begin{equation}
 \mathcal B_{\rm TT}(s,\xi)
 =
 \left[
 \kappa+\ell(s^2+|\xi|^2)
 \right]I
 +\alpha^2\rho\,G_\zeta,
 \qquad
 G_\zeta=B_\zeta^*B_\zeta\ge0.
\label{eq:v3v39-BTT}
\end{equation}
At primary glancing,
\begin{equation}
 s=i\tau,\quad|\xi|=|\tau|\ne0,
\qquad
 G_\zeta\ge\frac12I
\label{eq:v3v39-G-glancing}
\end{equation}
in the amplitude norm of Lemma~\ref{lem:v3v38-null-injective}.
\end{proposition}

\begin{proof}
The trace derivative in \eqref{eq:v3v39-stable} is
\[
 \partial_y\widehat\Psi(0)
 =\alpha\rho B_\zeta A_0.
\]
Substitute this and
\(\partial_x\widehat A(0)=-\kappa A_0\) into
\eqref{eq:v3v39-junction}.  This gives
\eqref{eq:v3v39-BTT}.  Positivity of \(G_\zeta\) is immediate, and
\eqref{eq:v3v39-G-glancing} is
\eqref{eq:v3v38-B-margin}.
\end{proof}

\begin{theorem}[No unstable TT--material--bath surface mode]
\label{thm:v3v39-no-unstable}
Assume
\begin{equation}
 \alpha>0,\qquad
 \ell\ge0,\qquad
 0\le c_{\rm e}^2<1,\qquad
 m_{\rm e}\ge0.
\label{eq:v3v39-assumptions}
\end{equation}
There is no nonzero finite-energy stable normal mode of
\eqref{eq:v3v39-bulk-wave}--%
\eqref{eq:v3v39-junction} with \(\operatorname{Re}s>0\).
Equivalently,
\begin{equation}
 \ker\mathcal B_{\rm TT}(s,\xi)=\{0\}
 \qquad(\operatorname{Re}s>0).
\label{eq:v3v39-no-unstable}
\end{equation}
The comoving material rows add no such root.
\end{theorem}

\begin{proof}
A stable normal mode with \(\operatorname{Re}s>0\) has finite spatial
energy.  Apply the polarized form of
Theorem~\ref{thm:v3v39-energy} to the real and imaginary parts, or
multiply the Laplace-transformed equations by the conjugate time
velocities.  The result is
\[
 \operatorname{Re}s\,
 (E_{\rm g}+E_\partial+E_{\rm R})=0.
\]
All three energies are nonnegative.  If their sum vanishes, the stable
mode and its derivatives vanish; the trace relation then also vanishes.
This contradicts a nonzero kernel vector.  Proposition
\ref{prop:v3v39-no-screening} already shows that the material
\((\pi,\lambda)\) block is invertible for such \(s\).
\end{proof}

\begin{theorem}[No nonzero closed-frequency root]
\label{thm:v3v39-no-closed-root}
Let \(m_{\rm e}=0\) and take limiting outgoing roots on
\(s=i\tau\), with \((\tau,\xi)\ne(0,0)\).  On the TT quotient,
\begin{equation}
 \ker\mathcal B_{\rm TT}(i\tau,\xi)=\{0\}.
\label{eq:v3v39-no-closed-root}
\end{equation}
At glancing this conclusion uses \(c_{\rm e}<1\) and the injectivity of
\(B_\zeta\).
\end{theorem}

\begin{proof}
Put \(a=|\xi|^2-\tau^2\).  If \(a>0\), then
\(\kappa=\sqrt a>0\).  When
\(c_{\rm e}^2|\xi|^2-\tau^2\ge0\), every term in
\eqref{eq:v3v39-BTT} is nonnegative self-adjoint and the \(\kappa I\)
term is strictly positive.  When that bath radicand is negative, the bath
term is purely imaginary while the real part
\((\sqrt a+\ell a)I\) is strictly positive.  No kernel exists.

If \(a=0\), then \(\kappa=0\) and the induced-gravity term vanishes.
Because \(c_{\rm e}<1\),
\[
 \rho=i\,\operatorname{sgn}(\tau)
 \sqrt{1-c_{\rm e}^2}\,|\tau|\ne0.
\]
Equation \eqref{eq:v3v39-G-glancing} makes
\(\alpha^2\rho G_\zeta\) invertible.

If \(a<0\), both \(\kappa\) and \(\rho\) are purely imaginary with the
same outgoing sign.  If \(\ell>0\), the nonzero real part
\(\ell aI\) prevents a kernel.  If \(\ell=0\), the imaginary part is the
same-sign sum of the positive definite \(|\kappa|I\) and the nonnegative
operator \(\alpha^2|\rho|G_\zeta\), so it is positive definite.  These
cases prove \eqref{eq:v3v39-no-closed-root}.  Since glancing has
\(\tau\ne0\), Proposition~\ref{prop:v3v39-no-screening} also removes the
material variables there.
\end{proof}

\begin{corollary}[Uniform glancing margin]
\label{cor:v3v39-glancing-margin}
If \(c_{\rm e}^2\le1-\delta\), then at normalized primary glancing
\begin{equation}
 \sigma_{\min}
 \bigl(\mathcal B_{\rm TT}(i\tau,\xi)\bigr)
 \ge
 \frac{\alpha^2\sqrt\delta}{2}\,|\tau|.
\label{eq:v3v39-glancing-margin}
\end{equation}
\end{corollary}

\begin{proof}
At glancing only the last term in
\eqref{eq:v3v39-BTT} remains.  Apply
\eqref{eq:v3v39-G-glancing} and the value of \(|\rho|\) in the proof of
Theorem~\ref{thm:v3v39-no-closed-root}.
\end{proof}

\subsection{Diffeomorphism quotient}
\label{sec:v3v39-quotient}

At the background \eqref{eq:v3v38-background}, an infinitesimal boundary
diffeomorphism \(v^a\) acts on the three scalar reference fields by
\begin{equation}
 \delta_vT=v^0,\qquad
 \delta_vX^I=v^I.
\label{eq:v3v39-reference-FP}
\end{equation}

\begin{proposition}[Unitary reference gauge has an invertible FP symbol]
\label{prop:v3v39-unitary-gauge}
The map
\begin{equation}
 v^a\longmapsto(\delta_vT,\delta_vX^1,\delta_vX^2)
\label{eq:v3v39-FP-map}
\end{equation}
is the identity in the adapted frame.  Thus the gauge
\(T=t,\ X^I=x^I\) fixes the boundary diffeomorphism symbol with a
frequency-independent nonzero Faddeev--Popov determinant.  The strain
\eqref{eq:v3v38-e} is covariant before this gauge is imposed, and the
physical bath determinant in \eqref{eq:v3v39-BTT} is not cancelled by the
gauge ghosts.
\end{proposition}

\begin{proof}
Equation \eqref{eq:v3v39-reference-FP} is immediate from the scalar Lie
derivative and the background gradients.  Its matrix is \(I_3\).
Diffeomorphism covariance of the strain was proved in
Proposition~\ref{prop:v3v38-covariant-strain}.  The bath and material
fields are physical variables in the invariant action, whereas the FP
ghosts cancel only the Jacobian of the coordinate gauge, as distinguished
by Theorem~\ref{thm:v3v37-quartet-cancellation}.
\end{proof}

\begin{firewall}[The theorem is the complete TT block, not the ten-field
Einstein determinant]
\label{fw:v3v39-scope}
Theorems~\ref{thm:v3v39-no-unstable} and
\ref{thm:v3v39-no-closed-root} include every TT, comoving-material,
multiplier, and bath amplitude in the frozen tracefree quotient.  They do
not recompute the scalar and vector harmonic boundary rows of V32--V33,
nor the lower-order mixing with lapse, shift, clock, and curved
connections.  A full Einstein complementing theorem must assemble those
blocks and control their off-diagonal perturbations.
\end{firewall}

\subsection{V39 conclusion and V40 target}
\label{sec:v3v39-conclusion}

V39 closes the first coupled regression of the physical medium.  The
spatial symmetrized-gradient map is onto, so unconstrained material labels
would screen the metric response.  The covariant comoving multiplier
removes that screening at every nonzero time frequency.  The exact
energy identity fixes the junction sign, and the eliminated TT matrix is
\[
 \mathcal B_{\rm TT}
 =
 [\kappa+\ell(s^2+|\xi|^2)]I
 +\alpha^2\rho B_\zeta^*B_\zeta.
\]
It has no growing root, no nonzero closed-frequency root, and a uniform
subluminal glancing margin.  Static label modes remain an honest separate
gate.

V40 is a mandatory companion audit.  After incorporating any relevant
Yang--Mills or Navier--Stokes lesson, it should assemble the V33 scalar
row, the V39 physical TT block, the remaining harmonic gauge rows, and
the reference-field FP block into one normalized principal matrix.  The
proof must quantify the allowable off-diagonal lower-symbol mixing and
must keep the unresolved \(s=0\) material sector explicit.
\clearpage
\section{V40: companion audit and the boundary-domain assembly obstruction}
\label{sec:v3v40}

\subsection{Mandatory read-only companion audit}
\label{sec:v3v40-audit}

Before choosing the V40 direction, the shared notes folder was inspected
without changing any companion file.  The newest active companion notes
were
\begin{align}
 &\texttt{Renewal\_NS\_Stability\_Research\_Notes\_V31.tex},
 \notag\\
 &\qquad 887537\ {\rm bytes},\qquad
 \texttt{F1121B62238AD5491BB7CF03635EA9934942EDF573ED8FAAA9EE79E1D38A6696},
\label{eq:v3v40-NS-fingerprint}\\
 &\texttt{Renewal\_Yang\_Mills\_Mass\_Gap\_Research\_Notes\_V89.tex},
 \notag\\
 &\qquad 1453371\ {\rm bytes},\qquad
 \texttt{0631BCF94AFCD677ABC65DA5B907D546A4DEE6991902896C270ED3A5699DE703}.
\label{eq:v3v40-YM-fingerprint}
\end{align}
The hashes are SHA--256 digests.  The frozen Yang--Mills volume
\texttt{V135\_f1} and the two frozen SM volumes were also present and were
left untouched.

The relevant Yang--Mills V89 result is methodological but exact: its
joining operator must be split into a signed first-order term and a
quadratic remainder before absolute product-state capacity is taken.
The positive envelope loses a half power, whereas the quadratic remainder
is already paid and only a signed rank-one vector moment remains.  The
relevant Navier--Stokes V31 results are similarly restrictive: removing a
flat dyadic mark restores exactly the missing first moment, and the signed
helicity hinge has a nontrivial nullspace while the total critical coarea
remains positive.

\begin{assessment}[Imported audit rules]
\label{ass:v3v40-audit-rules}
For the Einstein boundary symbol these results impose two useful
regressions.
\begin{enumerate}
 \item Principal off-diagonal rows must be retained in the signed coupled
       matrix.  Replacing them by separate positive norms can lose the
       cancellation or create a false margin.
 \item A factor which vanishes on a frequency face cannot be removed by
       normalization and then declared paid.  Likewise, a TT margin cannot
       be charged to a scalar or vector boundary sector on which its
       incidence vanishes.
\end{enumerate}
No Yang--Mills or Navier--Stokes estimate is imported as an Einstein
theorem; only these proved structural tests are used.
\end{assessment}

\subsection{Two boundary problems had been kept separate}
\label{sec:v3v40-two-domains}

At a flat boundary, let \(H_{\mu\nu}\) denote the ten-dimensional stable
metric amplitude and let \(C_\mu(H)\) be the four de Donder boundary rows.
The V32--V33 Dirichlet problem uses
\begin{equation}
 \mathcal B_DH
 =
 \bigl(
 H_{ab},\,
 C_a(H),\,
 C_x(H)-\zeta Q_{\rm e}H_{xx}
 \bigr),
\label{eq:v3v40-BD}
\end{equation}
where the six conditions \(H_{ab}=0\) are imposed before the last four
rows are reduced.  Its scalar Schur factor is
\begin{equation}
 \frac z{2p}-\zeta q
\label{eq:v3v40-V33-Schur}
\end{equation}
by Proposition~\ref{prop:v3v33-augmented-det}.

The V15--V16 and V38--V39 junction problem instead uses
\begin{equation}
 \mathcal B_JH
 =
 \bigl(
 \tau_{ab}^{(1)}(H)-\mathcal S_{ab}^{(1)}(H,\Psi,X,T),\,
 C_\mu(H)
 \bigr).
\label{eq:v3v40-BJ}
\end{equation}
Here the six tangential metric amplitudes are unknown boundary traces.
The physical bath changes the TT part of the first six rows.

\begin{theorem}[The V33 and V39 margins are not blocks of one boundary
matrix]
\label{thm:v3v40-domain-mismatch}
The scalar determinant of V33 and the TT determinant of V39 cannot be
combined by taking their block direct sum.  They are Schur complements of
different boundary domains:
\begin{equation}
 \mathcal D_D:\ H_{ab}=0,
 \qquad
 \mathcal D_J:\ \tau_{ab}^{(1)}=\mathcal S_{ab}^{(1)}.
\label{eq:v3v40-domains}
\end{equation}
In particular:
\begin{enumerate}
 \item on \(\mathcal D_D\), the TT trace is identically zero, so the V39
       junction block is absent;
 \item on \(\mathcal D_J\), \(H_{ab}\) is not zero, and the three rows
       \(C_a\) contain its tangential-divergence symbol, so the V33
       \(3+1\) reduction and the Schur factor
       \eqref{eq:v3v40-V33-Schur} do not apply; and
 \item imposing both sets of six rows gives sixteen metric boundary rows
       for ten stable metric amplitudes before a different mixed-domain
       analysis is supplied.
\end{enumerate}
Therefore separate nonzero margins in V33 and V39 do not prove
complementarity of a single Einstein boundary problem.
\end{theorem}

\begin{proof}
The first statement is immediate from the six leading entries of
\eqref{eq:v3v40-BD}.  For the second, inspect
\eqref{eq:v3v32-Ca}: its reduced form
\(-pH_{xa}-i\xi_aH_{xx}/2\) was obtained only after every
\(H_{ab}\) was set to zero.  Without that restriction,
\(\partial^b\bar h_{ab}\) contributes a nonzero term involving
\(H_{ab}\).  Thus the upper block whose inverse produced
\(z/(2p)\) has changed.  The last statement is the boundary-row count:
both systems separately have ten metric rows, while their simultaneous
six-component tangential conditions are independent prescriptions unless
a new constrained system proves otherwise.
\end{proof}

\begin{counterexample}[A direct sum can certify a nonexistent problem]
\label{cex:v3v40-false-direct-sum}
Let \(b_D\ne0\) be the scalar Schur factor of a Dirichlet reduction and
let \(b_J\ne0\) be a TT junction factor on a free trace.  The abstract
matrix \(\operatorname{diag}(b_D,b_J)\) is invertible.  If the first
domain imposes the TT trace to be zero and the second domain requires that
trace as its input, however, there is no original amplitude space on which
those two coordinates are independent.  Invertibility of the displayed
matrix says nothing about either original boundary map.
\end{counterexample}

\begin{proof}
The direct sum has been formed after two incompatible restrictions of the
same trace coordinate.  Since it is not the restriction or Schur
complement of either original map, its determinant is not a boundary
determinant.
\end{proof}

\begin{firewall}[This is the V40 upstream correction]
\label{fw:v3v40-upstream}
The conclusion requested at the end of V39 must be corrected before more
perturbation estimates are added.  V33 remains a valid theorem for the
fixed-induced-metric Euclidean domain, and V39 remains a valid theorem for
the physical TT junction quotient.  Neither is revoked.  What fails is
their proposed use as diagonal blocks of one full SM--Einstein domain.
The scalar and vector sectors of the junction problem must be recomputed
with the Brown--York rows present.
\end{firewall}

\subsection{The correct flat ten-row junction symbol}
\label{sec:v3v40-correct-symbol}

The starting point for the recomputation can be written explicitly.  Use
the flat Lorentzian half-space \(x>0\), boundary metric
\(\gamma_{ab}\), and a stable mode
\begin{equation}
 h_{\mu\nu}
 =
 H_{\mu\nu}e^{st+i\eta\cdot y-\kappa x},
 \qquad
 \kappa=\sqrt{s^2+|\eta|^2},
 \qquad
 \operatorname{Re}\kappa>0.
\label{eq:v3v40-stable-mode}
\end{equation}
Put
\begin{equation}
 \zeta_a=(s,i\eta_A),\qquad
 V_a=H_{xa},\qquad
 w=H_{xx},\qquad
 \theta=\gamma^{ab}H_{ab}.
\label{eq:v3v40-amplitudes}
\end{equation}

\begin{proposition}[Exact stable de Donder rows before restriction]
\label{prop:v3v40-full-C}
On \eqref{eq:v3v40-stable-mode},
\begin{align}
 C_a(H)
 &=
 \zeta^bH_{ab}
 -\frac12\zeta_a(w+\theta)
 -\kappa V_a,
\label{eq:v3v40-Ca}\\
 C_x(H)
 &=
 \zeta^aV_a-\frac\kappa2(w-\theta).
\label{eq:v3v40-Cx}
\end{align}
\end{proposition}

\begin{proof}
The trace is \(h=w+\theta\), so
\[
 \bar H_{ab}=H_{ab}-\frac12\gamma_{ab}(w+\theta),
 \quad
 \bar H_{xa}=V_a,
 \quad
 \bar H_{xx}=\frac12(w-\theta).
\]
Use
\(C_\mu=\partial^\nu\bar h_{\mu\nu}\),
\(\partial_a=\zeta_a\), and \(\partial_x=-\kappa\).
\end{proof}

For the normal \(+\partial_x\) used in the V15 algebraic calculation,
define the unscaled Brown--York row
\begin{equation}
 \mathcal J_{ab}^{+}(H)
 :=
 -\kappa H_{ab}
 -\zeta_aV_b-\zeta_bV_a
 +\gamma_{ab}\bigl(\kappa\theta+2\zeta^cV_c\bigr).
\label{eq:v3v40-Jplus}
\end{equation}

\begin{proposition}[Exact stable Brown--York rows]
\label{prop:v3v40-full-J}
The linearized Brown--York equation for the \(+\partial_x\) convention is
\begin{equation}
 \mathcal J_{ab}^{+}(H)
 =
 2\kappa_{\rm g}\mathcal S_{ab}^{(1)}.
\label{eq:v3v40-J-source}
\end{equation}
Changing to the geometric outward normal of \(x>0\) multiplies the
Brown--York side by \(-1\); the medium incidence must be changed
consistently so that the V39 energy identity is retained.
\end{proposition}

\begin{proof}
Substitute
\(\partial_x=-\kappa\) and \(\partial_a=\zeta_a\) into
\eqref{eq:v3v15-linear-BY} and multiply by
\(2\kappa_{\rm g}\).  This is \eqref{eq:v3v40-Jplus}.  Reversing the
normal reverses \(K_{ab}-K\gamma_{ab}\), hence the Brown--York tensor.
\end{proof}

\begin{definition}[The actual frozen metric junction map]
\label{def:v3v40-frozen-map}
Before auxiliary bath elimination, the square metric part of the correct
boundary map is
\begin{equation}
 \mathfrak B_{\rm met}(s,\eta):
 (H_{ab},V_a,w)
 \longmapsto
 \left(
 \mathcal J_{ab}^{+}(H)-2\kappa_{\rm g}\mathcal S_{ab}^{(1)},
 C_a(H),C_x(H)
 \right).
\label{eq:v3v40-Bmet}
\end{equation}
It has six junction rows and four de Donder rows acting on ten stable
metric amplitudes.  The bath traces, comoving multiplier rows, and their
stable normal equations must be appended before their Schur complement is
taken.
\end{definition}

\begin{assessment}[What remains valid inside the correct map]
\label{ass:v3v40-valid-subresults}
Restriction of \eqref{eq:v3v40-Bmet} to the divergence-free TT quotient
recovers the V39 block with the energy-selected orientation.  At a fixed
boundary point, the coefficient of the ten normal derivatives remains
invertible by Theorem~\ref{thm:v3v15-ten-block}.  These facts give the
zero-speed algebraic solve and the TT glancing solve.  They do not fill
the intervening scalar and vector Laplace--Fourier sectors.
\end{assessment}

\subsection{Reference fields and row counting}
\label{sec:v3v40-reference-count}

The unitary reference gauge
\(T=t,\ X^I=x^I\) has the invertible FP map
\eqref{eq:v3v39-FP-map}.  It acts in the boundary-diffeomorphism ghost
sector.  The physical comoving equations act on \(\pi^I,\lambda_I\), and
the half-line bath has its own stable amplitudes.

\begin{proposition}[Reference gauge does not add metric boundary rows]
\label{prop:v3v40-reference-count}
At the flat principal symbol:
\begin{enumerate}
 \item the three unitary reference conditions replace three boundary
       diffeomorphism gauge choices and their FP determinant is nonzero;
 \item the four comoving material rows form the nonstatic
       \((\pi,\lambda)\) block of
       Proposition~\ref{prop:v3v39-no-screening}; and
 \item neither set may be counted as additional equations for the ten
       metric stable amplitudes in
       \eqref{eq:v3v40-Bmet}.
\end{enumerate}
Thus the enlarged system is square only after variables and equations are
counted in their own metric, ghost, material, multiplier, and bath
spaces.
\end{proposition}

\begin{proof}
Item one is Proposition~\ref{prop:v3v39-unitary-gauge}.  Item two is
Proposition~\ref{prop:v3v39-no-screening}.  Their equations vary with
respect to reference and ghost variables, whereas
\eqref{eq:v3v40-Bmet} maps metric stable amplitudes to the ten gravitational
rows.  Counting either as an extra metric row would confuse a block
extension with an overdetermined restriction.
\end{proof}

\subsection{Normalization and the honest perturbation gate}
\label{sec:v3v40-normalization}

Let
\begin{equation}
 \lambda=(|s|^2+|\eta|^2)^{1/2}
\label{eq:v3v40-lambda}
\end{equation}
and restrict first to \(\lambda=1\) with \(s\ne0\).  Brown--York,
de Donder, and massless bath Dirichlet-to-Neumann rows all have degree one,
so divide these rows by \(\lambda\).  The reference FP block has degree
zero and is left unscaled.  The comoving block contains the factor
\(s/\lambda\), which is retained.

\begin{theorem}[Quantified lower-symbol stability]
\label{thm:v3v40-lower-stability}
Suppose the correctly assembled normalized frozen principal matrix
\(\widehat{\mathfrak B}_0(\omega)\) is square on a compact normalized
frequency set \(\Omega\) and satisfies
\begin{equation}
 \inf_{\omega\in\Omega}
 s_{\min}\bigl(\widehat{\mathfrak B}_0(\omega)\bigr)=c_0>0.
\label{eq:v3v40-c0}
\end{equation}
Let curvature, second fundamental form, variable clock, connection, and
mass terms produce a normalized lower-symbol perturbation \(E_\lambda\)
with
\begin{equation}
 \sup_{\omega\in\Omega}\|E_\lambda(\omega)\|
 \le\frac K\lambda.
\label{eq:v3v40-Elower}
\end{equation}
Then
\begin{equation}
 \lambda\ge\frac{2K}{c_0}
 \quad\Longrightarrow\quad
 s_{\min}
 \bigl(\widehat{\mathfrak B}_0+E_\lambda\bigr)
 \ge\frac{c_0}{2}.
\label{eq:v3v40-half-margin}
\end{equation}
A degree-one off-diagonal principal block is not covered by
\eqref{eq:v3v40-Elower}; it must be retained exactly or separately proved
to have norm smaller than \(c_0\).
\end{theorem}

\begin{proof}
The reverse triangle inequality for singular values gives
\[
 s_{\min}(\widehat{\mathfrak B}_0+E_\lambda)
 \ge c_0-\|E_\lambda\|.
\]
Insert \eqref{eq:v3v40-Elower} and the displayed lower bound on
\(\lambda\).  A degree-one block remains order one after normalization, so
it does not decay as \(\lambda^{-1}\) and the argument does not apply to
it.
\end{proof}

\begin{firewall}[The static comoving face cannot be normalized away]
\label{fw:v3v40-static-face}
The material block contains \(s/\lambda\).  On the normalized face
\(s=0,\eta\ne0\), its smallest singular value is zero.  Dividing the row
by \(s\) changes the graph norm and discards precisely the static
screening mode of Firewall~\ref{fw:v3v39-static}.  This is the SM analogue
of the NS V31 unmarking rule: the missing factor must be paid by a static
elastic term, a proved material relabelling quotient, or a weaker
frequency-dependent domain.  It cannot be removed by notation.
\end{firewall}

\begin{firewall}[TT positivity cannot pay a missing scalar row]
\label{fw:v3v40-sector-nullspace}
The V38 incidence \(B_\zeta^*B_\zeta\) is positive on null TT data, but it
has no proved positive restriction on the scalar and vector pieces of
\eqref{eq:v3v40-Bmet}.  The YM signed-row lesson and the NS hinge-nullspace
test both forbid replacing that missing calculation by the norm of the TT
block.  The full signed matrix must be decomposed before absolute operator
bounds are taken.
\end{firewall}

\subsection{V40 conclusion and V41 target}
\label{sec:v3v40-conclusion}

The mandatory audit prevents a false assembly.  V33 and V39 solve
different boundary domains: fixed induced metric plus de Donder gauge,
and Brown--York junction plus de Donder gauge.  Their scalar and TT
determinants cannot be put on the diagonal of one operator.  The correct
junction calculation starts from the ten explicit rows
\eqref{eq:v3v40-Ca}--\eqref{eq:v3v40-J-source}, with the physical bath and
material equations appended before elimination.

The TT restriction, energy sign, and algebraic normal-derivative solve
remain valid.  The scalar/vector junction symbol and the static material
face remain open.  Lower-order curvature mixing will be controlled by
Theorem~\ref{thm:v3v40-lower-stability} only after a genuine frozen margin
is proved.

V41 should now decompose
\(\mathfrak B_{\rm met}\) under rotations about a nonzero spatial
covector into tensor, vector, and scalar sectors.  It should retain the
signed Brown--York/de Donder couplings, eliminate only invertible stable
blocks, compute the scalar and vector determinants on the normalized
Laplace hemisphere, and test the result at glancing and at
\(\eta=0\).  The physical bath may be inserted only through its actual
incidence in each sector.  If a new kernel appears, the programme should
change the junction or material action upstream rather than importing the
V33 Dirichlet Schur factor.
\clearpage
\section{V41: free junction kernel classification and the completed comoving incidence}
\label{sec:v3v41}

\subsection{Purpose}
\label{sec:v3v41-purpose}

V40 replaced an invalid direct-sum assembly by the correct ten-row
Brown--York/de Donder map.  V41 begins its exact reduction.  It first
classifies the kernel of the free metric map.  Away from glancing the
kernel consists exactly of three tangential boundary diffeomorphisms.  At
glancing there is one additional gauge-quotient scalar mode.  The
apparently TT null tensors at a null boundary covector are contained in
the tangential gauge image.

The second calculation corrects the material block used in V39.  The
linearized comoving equation contains the mixed metric component
\(H_{0I}\), and its multiplier contributes to the mixed junction stress.
After the material variables are eliminated at \(s\ne0\), the correct
bath form is \(E_s^*\rho E_s\), where \(E_s\) contains both spatial strain
and the mixed metric components.  The V39 matrix is its restriction to
\(H_{0I}=0\).

\subsection{Rotational vector block}
\label{sec:v3v41-vector}

Choose the nonzero spatial covector along \(y^1\):
\begin{equation}
 \eta=(k,0),\qquad k\ge0,
 \qquad
 \zeta_a=(s,ik,0),
 \qquad
 \kappa^2=s^2+k^2.
\label{eq:v3v41-frame}
\end{equation}
The odd sector under reflection \(y^2\mapsto-y^2\) has variables
\begin{equation}
 (H_{02},H_{12},V_2)
\label{eq:v3v41-vector-vars}
\end{equation}
and rows
\((\mathcal J^+_{02},\mathcal J^+_{12},C_2)\).

\begin{proposition}[Exact free vector matrix]
\label{prop:v3v41-vector-matrix}
The free vector matrix is
\begin{equation}
 \mathcal B_{\rm vec}^{0}
 =
 \begin{pmatrix}
 -\kappa&0&-s\\
 0&-\kappa&-ik\\
 -s&ik&-\kappa
 \end{pmatrix},
\label{eq:v3v41-vector-matrix}
\end{equation}
and
\begin{equation}
 \det\mathcal B_{\rm vec}^{0}
 =
 -\kappa(\kappa^2-s^2-k^2).
\label{eq:v3v41-vector-det}
\end{equation}
Thus its determinant vanishes identically on the stable bulk
characteristic relation.
\end{proposition}

\begin{proof}
Insert the variables \eqref{eq:v3v41-vector-vars} in
\eqref{eq:v3v40-Ca} and \eqref{eq:v3v40-Jplus}.  Expansion of the
three-by-three determinant gives
\[
 -\kappa(\kappa^2-k^2)+\kappa s^2
 =
 -\kappa(\kappa^2-s^2-k^2).
\]
\end{proof}

\begin{corollary}[The vector zero is a tangential diffeomorphism]
\label{cor:v3v41-vector-gauge}
For \(\kappa\ne0\), the vector kernel is one dimensional and is generated
by
\begin{equation}
 H_{02}=s\chi_2,\qquad
 H_{12}=ik\chi_2,\qquad
 V_2=-\kappa\chi_2.
\label{eq:v3v41-vector-gauge}
\end{equation}
This is the stable metric image of the tangential vector field with
boundary amplitude \(\chi_2\).
\end{corollary}

\begin{proof}
The first two rows of \eqref{eq:v3v41-vector-matrix} give the displayed
relations after setting \(V_2=-\kappa\chi_2\).  The third holds because
\(\kappa^2=s^2+k^2\).  These amplitudes are precisely
\(H_{ab}=\zeta_a\chi_b+\zeta_b\chi_a\) and
\(V_a=-\kappa\chi_a\) with only \(\chi_2\ne0\).
\end{proof}

\subsection{The complete free kernel away from glancing}
\label{sec:v3v41-free-kernel}

For a tangential covector amplitude \(\chi_a\), define
\begin{equation}
 (\mathcal R_\zeta\chi)_{ab}
 =\zeta_a\chi_b+\zeta_b\chi_a,
 \qquad
 (\mathcal R_\zeta\chi)_{xa}=-\kappa\chi_a,
 \qquad
 (\mathcal R_\zeta\chi)_{xx}=0.
\label{eq:v3v41-R}
\end{equation}

\begin{theorem}[Exact free kernel for \(\kappa\ne0\)]
\label{thm:v3v41-free-kernel}
On the stable characteristic relation and for \(\kappa\ne0\),
\begin{equation}
 \ker\mathfrak B_{\rm met}^{0}
 =
 \operatorname{Ran}\mathcal R_\zeta,
 \qquad
 \dim\ker\mathfrak B_{\rm met}^{0}=3,
\label{eq:v3v41-free-kernel}
\end{equation}
where the superscript zero means zero junction source.  Consequently the
ten-by-ten free map has rank seven.
\end{theorem}

\begin{proof}
A tangential diffeomorphism of the flat boundary preserves the zero
background Brown--York tensor.  Direct substitution of
\eqref{eq:v3v41-R} in
\eqref{eq:v3v40-Ca}--\eqref{eq:v3v40-Jplus} also gives zero, using
\(\kappa^2=s^2+k^2\).  Hence the displayed range lies in the kernel.

Conversely, let \((H_{ab},V_a,w)\) lie in the kernel.  Subtract the gauge
image \eqref{eq:v3v41-R} with
\begin{equation}
 \chi_a=-\kappa^{-1}V_a.
\label{eq:v3v41-gauge-choice}
\end{equation}
The new representative has \(V_a=0\).  Its six Brown--York rows reduce to
\[
 -\kappa H_{ab}+\kappa\gamma_{ab}\theta=0.
\]
Since \(\kappa\ne0\), \(H_{ab}=\gamma_{ab}\theta\).  Taking the
three-dimensional boundary trace gives
\(\theta=3\theta\), so \(\theta=0\) and \(H_{ab}=0\).  The de Donder rows
then give
\[
 -\frac12\zeta_aw=0,\qquad -\frac\kappa2w=0,
\]
and hence \(w=0\).  Every kernel element was therefore the subtracted
gauge image.  Injectivity of \(\chi\mapsto\mathcal R_\zeta\chi\) for
\(\kappa\ne0\) gives dimension three and rank seven.
\end{proof}

\begin{assessment}[Codazzi and gauge remove equal dimensions]
\label{ass:v3v41-rank-seven}
The rank-seven result is consistent with the three tangential
diffeomorphism directions in the domain and the three boundary
diffeomorphism Ward/Codazzi identities in the row space.  A complementing
determinant must therefore be formed on the quotient by the former and
the compatible quotient of the codomain by the latter, or by a bordered
gauge/identity system.  The determinant of the unquotiented ten-by-ten
matrix is identically zero.
\end{assessment}

\subsection{Exact glancing kernel}
\label{sec:v3v41-glancing}

Take a normalized null frequency
\begin{equation}
 s=i,\qquad k=1,\qquad\kappa=0.
\label{eq:v3v41-normalized-glancing}
\end{equation}
Other nonzero glancing frequencies follow by rotation and homogeneous
rescaling.

\begin{theorem}[One additional quotient mode at glancing]
\label{thm:v3v41-glancing-kernel}
At \eqref{eq:v3v41-normalized-glancing}, the free kernel has dimension
four.  The three-dimensional tangential gauge subspace has
\begin{equation}
 H_{ab}=\zeta_a\chi_b+\zeta_b\chi_a,
 \qquad V_a=0,\qquad w=0.
\label{eq:v3v41-glancing-gauge}
\end{equation}
Modulo that subspace, the remaining kernel is one dimensional and has the
representative
\begin{equation}
 H_{22}=a,\qquad w=-a,
 \qquad
 H_{00}=H_{01}=H_{02}=H_{11}=H_{12}=V_a=0.
\label{eq:v3v41-glancing-representative}
\end{equation}
Thus the free quotient loses exactly one margin at glancing.
\end{theorem}

\begin{proof}
At \(\kappa=0\), the \(C_x\) row gives
\(V_1=V_0\).  The Brown--York rows are
\[
 -\zeta_aV_b-\zeta_bV_a
 +2\gamma_{ab}\zeta^cV_c=0.
\]
Together with \(\zeta^cV_c=0\), already supplied by \(C_x\), these imply
\(V=0\).  The remaining three de Donder equations are
\begin{align*}
 -H_{00}+H_{01}-\frac12(w+\theta)&=0,\\
 -H_{01}+H_{11}-\frac12(w+\theta)&=0,\\
 -H_{02}+H_{12}&=0.
\end{align*}
They impose three independent conditions on the seven amplitudes
\((H_{ab},w)\), so the kernel has dimension four.

The map \(\chi\mapsto\zeta_{(a}\chi_{b)}\) is injective and gives three
solutions with \(w=0\).  Subtracting them can set
\(H_{00},H_{01},H_{02},H_{11},H_{12}\) to zero subject to the three
displayed equations.  Those equations then reduce to
\(w+H_{22}=0\), which is exactly
\eqref{eq:v3v41-glancing-representative}.  Its nonzero \(w\) prevents it
from belonging to the tangential gauge subspace.
\end{proof}

\begin{proposition}[Null boundary TT tensors are gauge-contaminated]
\label{prop:v3v41-null-TT-gauge}
In three-dimensional boundary Minkowski space, every tensor satisfying
\begin{equation}
 \gamma^{ab}A_{ab}=0,\qquad \zeta^aA_{ab}=0,
 \qquad \gamma^{ab}\zeta_a\zeta_b=0
\label{eq:v3v41-null-TT}
\end{equation}
can be written
\begin{equation}
 A_{ab}=\zeta_a\chi_b+\zeta_b\chi_a,
 \qquad
 \zeta^a\chi_a=0.
\label{eq:v3v41-null-TT-gauge}
\end{equation}
Hence the two-dimensional null TT space of V38 is contained in the
limiting tangential gauge image.  The additional physical quotient kernel
is the scalar-normal representative
\eqref{eq:v3v41-glancing-representative}.
\end{proposition}

\begin{proof}
In the normalized frame, the general null TT tensor is the two-parameter
matrix \eqref{eq:v3v38-null-matrix}.  Solving
\(A=\zeta\otimes\chi+\chi\otimes\zeta\) gives a two-parameter family
\(\chi\in\zeta^\perp\), and direct substitution reproduces that matrix.
Equivalently, the right side with \(\zeta\cdot\chi=0\) is tracefree and
transverse because \(\zeta^2=0\); both spaces have dimension two and the
map is injective.
\end{proof}

\begin{firewall}[Meaning of the earlier TT regression]
\label{fw:v3v41-TT-meaning}
V15--V16 correctly identify loss of the componentwise Neumann trace and
V38 correctly proves that spatial STF incidence is injective on the null
TT fibre.  V41 shows that the null TT fibre is not yet the physical
quotient of the complete Brown--York/de Donder map: at exact boundary
glancing it is contained in the tangential gauge image.  The physical
quotient still has a one-dimensional glancing loss, but its representative
mixes \(H_{22}\) with the double-normal amplitude \(w\).  Future
complementing claims must use this quotient representative.
\end{firewall}

\subsection{Linearizing the comoving equation correctly}
\label{sec:v3v41-comoving}

Allow a clock perturbation \(T=t+\vartheta\) and material perturbations
\(X^I=x^I+\pi^I\).  At the background
\eqref{eq:v3v38-background}, the spatial variation of the clock flow is
\begin{equation}
 \delta u^I=-\partial_I\vartheta-h_{0I}.
\label{eq:v3v41-du}
\end{equation}

\begin{lemma}[Completed comoving row]
\label{lem:v3v41-comoving-row}
The linearization of \(u^a\nabla_aX^I=0\) is
\begin{equation}
 \partial_t\pi^I-\partial_I\vartheta-h_{0I}=0.
\label{eq:v3v41-comoving-linear}
\end{equation}
In Laplace--Fourier variables,
\begin{equation}
 s\widehat\pi_I-i\xi_I\widehat\vartheta-H_{0I}=0.
\label{eq:v3v41-comoving-symbol}
\end{equation}
Thus the V39 row \(s\pi=0\) is valid only on the restricted subspace
\(\vartheta=H_{0I}=0\).
\end{lemma}

\begin{proof}
Vary
\(u^a\partial_aX^I\).  The background has \(u=\partial_t\),
\(\partial_tX^I=0\), and
\(\partial_jX^I=\delta_j^I\).  The variation is therefore
\[
 \partial_t\pi^I+\delta u^I.
\]
Equation \eqref{eq:v3v41-du} gives
\eqref{eq:v3v41-comoving-linear}, and Fourier transformation gives the
last formula.
\end{proof}

\begin{proof}[Proof of \eqref{eq:v3v41-du}]
From \(u_a=-\partial_aT/\sqrt{X_T}\), the spatial covector variation is
\(\delta u_I=-\partial_I\vartheta\).  Raising the index also varies the
inverse metric:
\[
 \delta u^I
 =
 \delta\gamma^{I0}u_0+\delta^{IJ}\delta u_J
 =
 -h_{0I}-\partial_I\vartheta.
\]
\end{proof}

\subsection{Eliminating the completed material block}
\label{sec:v3v41-completed-incidence}

For \(s\ne0\), solve \eqref{eq:v3v41-comoving-symbol}:
\begin{equation}
 \widehat\pi
 =
 s^{-1}\bigl(H_{0\cdot}+i\xi\widehat\vartheta\bigr).
\label{eq:v3v41-pi-solve}
\end{equation}
Using Lemma~\ref{lem:v3v38-linear-strain}, define the full strain
incidence
\begin{equation}
 E_s(H,\vartheta)
 :=
 D_\xi s^{-1}
 \bigl(H_{0\cdot}+i\xi\vartheta\bigr)
 -\mathbb P_0(H_{ij}).
\label{eq:v3v41-Es}
\end{equation}

\begin{theorem}[Correct eliminated physical bath form]
\label{thm:v3v41-correct-bath}
At a finite cutoff and for \(s\ne0\), eliminating the comoving material
coordinate and the half-line bath gives the quadratic response
\begin{equation}
 S_{\rm bath,eff}^{(2)}
 =
 \frac{\alpha^2}{2}
 \left\langle
 E_s(H,\vartheta),\,
 \rho(s,\xi)E_s(H,\vartheta)
 \right\rangle
\label{eq:v3v41-effective-form}
\end{equation}
with the retarded or Euclidean branch appropriate to the problem.  Its
complete metric--clock Hessian is
\begin{equation}
 \mathbb Z_{\rm full}
 =
 \alpha^2E_s^*\rho E_s.
\label{eq:v3v41-Zfull}
\end{equation}
The mixed components in \(E_s^*\) are the stress contribution of the
comoving multiplier.  On
\(H_{0I}=\vartheta=0\), this reduces to the V39 spatial form
\(\alpha^2B^*\rho B\).
\end{theorem}

\begin{proof}
Equation \eqref{eq:v3v41-pi-solve} inserted in
\eqref{eq:v3v38-linear-strain} gives exactly
\eqref{eq:v3v41-Es}.  The Poisson extension calculation of
Theorem~\ref{thm:v3v38-effective} gives
\eqref{eq:v3v41-effective-form}.  Varying that single quadratic form in
all metric and clock arguments produces
\eqref{eq:v3v41-Zfull}.  Equivalently, varying
\[
 S_{\rm com}^{(2)}
 =
 \int\lambda_I
 \bigl(\partial_t\pi^I-\partial_I\vartheta-h_{0I}\bigr)
\]
shows directly that \(-\lambda_I\) contributes to the mixed metric stress;
the material Euler equation determines \(\lambda\) from the adjoint strain
force.  Setting \(H_{0I}=\vartheta=0\) leaves
\(E_s=-\mathbb P_0H_{ij}=-B H\), proving the restriction.
\end{proof}

\begin{proposition}[Gauge directions lie in the completed incidence
kernel]
\label{prop:v3v41-Es-gauge}
Let a tangential diffeomorphism have amplitudes
\(\chi^0,\chi^I\), so that
\begin{equation}
 \vartheta=\chi^0,\qquad
 \pi^I=\chi^I,\qquad
 H_{0I}=s\chi_I-i\xi_I\chi^0,\qquad
 H_{ij}=i\xi_i\chi_j+i\xi_j\chi_i.
\label{eq:v3v41-gauge-amplitudes}
\end{equation}
Then
\begin{equation}
 E_s(H,\vartheta)=0.
\label{eq:v3v41-Es-gauge}
\end{equation}
Thus the physical medium does not assign energy to a pure
diffeomorphism before reference gauge fixing.
\end{proposition}

\begin{proof}
Equation \eqref{eq:v3v41-pi-solve} gives \(\pi=\chi\).
Consequently \(D_\xi\pi=\mathbb P_0H_{ij}\), and the two terms in
\eqref{eq:v3v41-Es} cancel.
\end{proof}

\begin{proposition}[Nonzero incidence on the physical glancing quotient]
\label{prop:v3v41-Es-glancing}
For the quotient representative
\eqref{eq:v3v41-glancing-representative} in clock gauge,
\begin{equation}
 E_i=
 E_s(H,0)
 =
 \begin{pmatrix}
 a/2&0\\
 0&-a/2
 \end{pmatrix}
\label{eq:v3v41-Es-glancing}
\end{equation}
up to the choice of which spatial axis is transverse.  Hence
\begin{equation}
 \|E_i\|_{\rm HS}^2=\frac12|a|^2.
\label{eq:v3v41-Es-glancing-norm}
\end{equation}
For \(c_{\rm e}^2\le1-\delta\), the bath form is therefore nonzero of
homogeneous size at least
\(\alpha^2\sqrt\delta\,|\tau||a|^2/4\) on that quotient line.
\end{proposition}

\begin{proof}
The representative has \(H_{0I}=0\) and spatial matrix
\(\operatorname{diag}(0,a)\).  Its negative tracefree part is the matrix
in \eqref{eq:v3v41-Es-glancing}, up to overall sign and axis order.
The norm follows by contraction.  At primary glancing,
\(|\rho|\ge\sqrt\delta|\tau|\); insert this and the norm into
\eqref{eq:v3v41-effective-form}.
\end{proof}

\begin{firewall}[Nonzero quotient form is not yet the bordered
determinant]
\label{fw:v3v41-not-determinant}
Proposition~\ref{prop:v3v41-Es-glancing} proves that the corrected physical
medium sees the actual one-dimensional free quotient kernel.  It does not
by itself prove that the complete bordered boundary matrix is invertible.
The multiplier stress fills mixed junction rows, and the three Ward
identities must be removed or bordered together with the three reference
gauge conditions.  Their signed Schur complement could still change the
orientation or introduce a root away from exact glancing.
\end{firewall}

\subsection{V41 conclusion and V42 target}
\label{sec:v3v41-conclusion}

V41 gives the first exact decomposition of the correct free junction map.
Its vector determinant vanishes identically because of a tangential
diffeomorphism.  For every \(\kappa\ne0\), the full kernel consists
exactly of the three tangential gauge images and the map has rank seven.
At glancing its kernel gains one physical quotient line represented by
\(H_{22}=a,\ H_{xx}=-a\).  Null boundary TT tensors themselves lie in the
limiting tangential gauge image.

The comoving material equation also changes from the V39 restricted row
to
\(s\pi-i\xi\vartheta-H_{0\cdot}=0\).  Complete elimination produces the
covariant incidence \(\alpha^2E_s^*\rho E_s\).  It annihilates gauge
directions and is strictly nonzero on the actual glancing quotient
kernel.  This preserves the physical-bath route but changes the matrix
which must be tested.

V42 should construct an explicit seven-dimensional quotient, or an
equivalent bordered square system with three reference gauge rows and
three Ward rows removed.  It must insert
\(\mathbb Z_{\rm full}\) before taking any determinant, factor the
resulting vector and scalar polynomials, and verify the energy-selected
orientation throughout the closed Laplace hemisphere.  The calculation
must keep \(s=0\) outside the comoving inversion until the static elastic
or relabelling gate is resolved.
\clearpage
\section{V42: an explicit quotient determinant and the intermediate-speed surface root}
\label{sec:v3v42}

\subsection{Purpose and corrected claim boundary}
\label{sec:v3v42-purpose}

V41 identified the three tangential gauge directions, the three
Ward/Codazzi row relations, and the additional one-dimensional glancing
kernel.  It also showed that the physical material response is the
completed Hessian \(E_s^*\rho E_s\), rather than the spatial restriction
used in V39.  V42 constructs an explicit seven-dimensional quotient chart
for \(s\ne0\) and nonzero spatial frequency, derives every multiplier
stress coefficient from the covariant quadratic action, and computes its
determinant.

The result changes the speed-family conclusion.  The completed physical
bath still removes the exact glancing zero.  Nevertheless, for every
strictly positive subluminal speed \(0<c_{\rm e}<1\), the quotient
determinant has a zero at an imaginary frequency strictly inside both the
bulk and bath evanescent regions.  It is a conservative trapped surface
mode.  The V39 no-closed-root theorem therefore applies only to the
restricted equations written there, not to the completed covariant
material system.

\subsection{A local domain quotient chart}
\label{sec:v3v42-chart}

Use the frame \eqref{eq:v3v41-frame} and assume
\begin{equation}
 s\ne0,\qquad k>0.
\label{eq:v3v42-chart-region}
\end{equation}
The tangential gauge image satisfies
\begin{align}
 H_{00}&=2s\chi_0,
&
 H_{01}&=s\chi_1+ik\chi_0,
&
 H_{02}&=s\chi_2.
\label{eq:v3v42-gauge-triplet}
\end{align}

\begin{lemma}[Uniform local gauge slice]
\label{lem:v3v42-gauge-slice}
On \eqref{eq:v3v42-chart-region}, every tangential gauge orbit has a
unique representative satisfying
\begin{equation}
 H_{00}=H_{01}=H_{02}=0.
\label{eq:v3v42-gauge-slice}
\end{equation}
The gauge-to-slice matrix in
\eqref{eq:v3v42-gauge-triplet} has determinant \(2s^3\).
\end{lemma}

\begin{proof}
In the ordering \((\chi_0,\chi_1,\chi_2)\), the matrix is triangular after
the first column is eliminated:
\[
 \begin{pmatrix}
 2s&0&0\\
 ik&s&0\\
 0&0&s
 \end{pmatrix}.
\]
Its determinant is \(2s^3\), which is nonzero in the stated chart.
\end{proof}

On the slice, use the seven variables
\begin{equation}
 X_{\rm q}
 =
 (H_{11},H_{12},H_{22},V_0,V_1,V_2,w)^{\mathsf T}
\label{eq:v3v42-variables}
\end{equation}
and retain the seven rows
\begin{equation}
 (C_0,C_1,C_2,C_x,\mathcal J^+_{01},
 \mathcal J^+_{02},\mathcal J^+_{22}).
\label{eq:v3v42-rows}
\end{equation}
The other three junction rows are recovered from the covariant
Ward/Codazzi relations when the material and multiplier equations hold.

\begin{assessment}[Chart scope]
\label{ass:v3v42-chart-scope}
The selected rows form a local quotient chart, not a new boundary
condition.  Lemma~\ref{lem:v3v42-gauge-slice} removes the domain gauge
directions, while the three omitted rows are the dependent codomain
directions.  The factors \(s\) and \(k\) in the determinant below record
where this chart degenerates.  They may not be interpreted as physical
zeros.
\end{assessment}

\subsection{Relative normalization of the completed material stress}
\label{sec:v3v42-stress}

Work in clock gauge \(\vartheta=0\).  The linear comoving and strain
relations are
\begin{equation}
 s\pi_I-H_{0I}=0,
 \qquad
 e=D_\xi\pi-\mathbb P_0H_{\rm sp}.
\label{eq:v3v42-comoving-strain}
\end{equation}
Let
\begin{equation}
 R=\alpha^2\rho e
\label{eq:v3v42-R}
\end{equation}
be the force conjugate to \(e\).  The quadratic multiplier term is
\begin{equation}
 S_{\rm com}^{(2)}
 =
 \int\lambda_I
 (\partial_t\pi^I-h_{0I}).
\label{eq:v3v42-Scom}
\end{equation}
Variation in \(\pi\) gives
\begin{equation}
 -s\lambda+D_\xi^\dagger R=0.
\label{eq:v3v42-lambda-equation}
\end{equation}

On the slice \eqref{eq:v3v42-gauge-slice}, \(\pi=0\) and
\begin{equation}
 e=-\mathbb P_0H_{\rm sp}.
\label{eq:v3v42-e-slice}
\end{equation}
Put
\begin{equation}
 d=H_{11}-H_{22}.
\label{eq:v3v42-d}
\end{equation}
For \(\xi=(k,0)\),
\begin{equation}
 \mathbb P_0H_{\rm sp}
 =
 \begin{pmatrix}
 d/2&H_{12}\\
 H_{12}&-d/2
 \end{pmatrix},
\qquad
 D_\xi^\dagger R
 =
 \alpha^2\rho
 \begin{pmatrix}
 ikd\\
 2ikH_{12}
 \end{pmatrix}.
\label{eq:v3v42-DstarR}
\end{equation}

\begin{lemma}[Completed mixed and spatial junction rows]
\label{lem:v3v42-completed-rows}
Let
\begin{equation}
 \mu=4\kappa_{\rm g}\alpha^2\rho.
\label{eq:v3v42-mu}
\end{equation}
With the \(+\partial_x\) Brown--York convention of V40 and the
energy-selected sign obtained by reversing the full outward equation, the
material contributions to the retained junction rows are
\begin{align}
 \Delta\mathcal J^+_{01}
 &=-\frac{ik\mu}{2s}\,d,
\label{eq:v3v42-dJ01}\\
 \Delta\mathcal J^+_{02}
 &=-\frac{ik\mu}{s}\,H_{12},
\label{eq:v3v42-dJ02}\\
 \Delta\mathcal J^+_{22}
 &=\frac\mu2\,d.
\label{eq:v3v42-dJ22}
\end{align}
These relative coefficients are fixed by the single covariant action and
cannot be varied independently.
\end{lemma}

\begin{proof}
The tensor variation of the strain form gives
\[
 \mathcal S_{ij}^{\rm strain}=-2R_{ij}.
\]
Since \(R=-\alpha^2\rho\,\mathbb P_0H_{\rm sp}\) on the slice, the row
\(\mathcal J^+-2\kappa_{\rm g}\mathcal S=0\) receives
\(-\mu\mathbb P_0H_{\rm sp}\).  Its \(22\) component is
\(\mu d/2\), proving \eqref{eq:v3v42-dJ22}.

From \eqref{eq:v3v42-lambda-equation} and
\eqref{eq:v3v42-DstarR},
\[
 \lambda_1=\frac{ik\alpha^2\rho}{s}d,
 \qquad
 \lambda_2=\frac{2ik\alpha^2\rho}{s}H_{12}.
\]
Variation of \eqref{eq:v3v42-Scom} in \(h_{0I}\) gives
\(\mathcal S^{0I}=-\lambda_I\), hence
\(\mathcal S_{0I}=\lambda_I\) at the flat Lorentzian background.
The row \(\mathcal J^+-2\kappa_{\rm g}\mathcal S\) therefore receives
\(-2\kappa_{\rm g}\lambda_I\).  Substitution and
\eqref{eq:v3v42-mu} give
\eqref{eq:v3v42-dJ01}--\eqref{eq:v3v42-dJ02}.
\end{proof}

\begin{firewall}[Why the V39 restricted determinant cannot be reused]
\label{fw:v3v42-V39-restricted}
Even though \(H_{0I}=0\) on the quotient slice, variation with respect to
\(H_{0I}\) is not zero there.  It produces the multiplier stress
\eqref{eq:v3v42-dJ01}--\eqref{eq:v3v42-dJ02}.  Dropping those rows after
setting their arguments to zero does not compute the Hessian of the
covariant action.  The full signed row coupling must be retained before
the determinant is taken.
\end{firewall}

\subsection{The seven-by-seven matrix}
\label{sec:v3v42-matrix}

Multiply the two mixed junction rows
\(\mathcal J^+_{01}\) and \(\mathcal J^+_{02}\) by \(s\).  This clears the
comoving denominators and does not change their zero set on
\eqref{eq:v3v42-chart-region}.  In the variable and row order
\eqref{eq:v3v42-variables}--\eqref{eq:v3v42-rows}, the resulting matrix is
\begin{equation}
 \widetilde{\mathcal B}_{\rm q}
 =
 \begin{pmatrix}
 -s/2&0&-s/2&-\kappa&0&0&-s/2\\
 ik/2&0&-ik/2&0&-\kappa&0&-ik/2\\
 0&ik&0&0&0&-\kappa&0\\
 \kappa/2&0&\kappa/2&-s&ik&0&-\kappa/2\\
 -ik\mu/2&0&ik\mu/2&-iks&-s^2&0&0\\
 0&-ik\mu&0&0&0&-s^2&0\\
 \kappa+\mu/2&0&-\mu/2&-2s&2ik&0&0
 \end{pmatrix}.
\label{eq:v3v42-matrix}
\end{equation}

\begin{theorem}[Exact quotient-chart determinant]
\label{thm:v3v42-determinant}
On the stable bulk relation \(\kappa^2=s^2+k^2\),
\begin{align}
 s^2\det\mathcal B_{\rm q}
 &=
 \det\widetilde{\mathcal B}_{\rm q}
\notag\\
 &=
 -\frac{k^2}{2}
 \Bigl[
 \kappa\mu^2(k^2+2s^2)^2
 +s^2\mu(k^4+8s^2k^2+8s^4)
 +4s^6\kappa
 \Bigr].
\label{eq:v3v42-determinant}
\end{align}
For \(\mu=0\), this reduces to
\begin{equation}
 \det\mathcal B_{\rm q}^{0}
 =-2s^4k^2\kappa,
\label{eq:v3v42-free-determinant}
\end{equation}
which agrees with Theorem~\ref{thm:v3v41-free-kernel}: the physical
quotient is invertible when \(\kappa sk\ne0\).
\end{theorem}

\begin{proof}
Insert \eqref{eq:v3v40-Ca}--\eqref{eq:v3v40-Jplus} and
Lemma~\ref{lem:v3v42-completed-rows} on the slice
\eqref{eq:v3v42-gauge-slice}; this gives
\eqref{eq:v3v42-matrix}.  Direct determinant expansion before the stable
relation gives a polynomial in \(s,k,\kappa,\mu\).  Replace each occurrence
of \(\kappa^2\) by \(s^2+k^2\) and collect powers of \(\mu\).  The
quadratic coefficient is
\[
 -\frac{k^2}{2}\kappa(k^2+2s^2)^2,
\]
the linear coefficient is
\[
 -\frac{k^2}{2}s^2(k^4+8s^2k^2+8s^4),
\]
and the constant coefficient is
\(-2k^2s^6\kappa\).  This is
\eqref{eq:v3v42-determinant}.  Divide by \(s^2\) and set \(\mu=0\) to
obtain \eqref{eq:v3v42-free-determinant}.
\end{proof}

\begin{corollary}[Exact glancing zero is lifted]
\label{cor:v3v42-glancing-lift}
At \(s=i\tau\), \(k=|\tau|\ne0\), \(\kappa=0\),
\begin{equation}
 \det\widetilde{\mathcal B}_{\rm q}
 =\frac12\mu k^8\ne0
\label{eq:v3v42-glancing-lift}
\end{equation}
whenever \(\mu\ne0\).
\end{corollary}

\begin{proof}
Use \(s^2=-k^2\) in
\eqref{eq:v3v42-determinant}.  The polynomial multiplying the linear
term becomes \(k^4\), and the other two terms contain \(\kappa\).
\end{proof}

\subsection{The intermediate-speed surface mode}
\label{sec:v3v42-surface-mode}

Take a massless bath with
\begin{equation}
 0<c_{\rm e}<1,\qquad
 \rho(s,k)=\sqrt{s^2+c_{\rm e}^2k^2},
 \qquad
 \mu=\beta\rho,\qquad \beta>0.
\label{eq:v3v42-speed-bath}
\end{equation}
For \(s=ixk\) with \(0<x<c_{\rm e}\), both normal roots are positive real:
\begin{equation}
 \kappa=k\sqrt{1-x^2},
 \qquad
 \rho=k\sqrt{c_{\rm e}^2-x^2}.
\label{eq:v3v42-evanescent-roots}
\end{equation}
After division by the nonzero common power of \(k\), the bracket in
\eqref{eq:v3v42-determinant} is
\begin{align}
 F_{c,\beta}(x)
 =&\,
 \beta^2\sqrt{1-x^2}(c^2-x^2)(1-2x^2)^2
\notag\\
 &-\beta x^2\sqrt{c^2-x^2}
       (1-8x^2+8x^4)
 -4x^6\sqrt{1-x^2},
\qquad c=c_{\rm e}.
\label{eq:v3v42-F}
\end{align}

\begin{theorem}[Every positive subluminal speed creates a neutral surface
root]
\label{thm:v3v42-surface-root}
For every
\begin{equation}
 0<c_{\rm e}<1,\qquad\beta>0,
\label{eq:v3v42-root-assumptions}
\end{equation}
there exists
\begin{equation}
 x_*\in(0,c_{\rm e})
\label{eq:v3v42-xstar}
\end{equation}
such that
\begin{equation}
 F_{c_{\rm e},\beta}(x_*)=0.
\label{eq:v3v42-Fzero}
\end{equation}
Consequently
\[
 s_*=ix_*k,\qquad k>0,
\]
is a nonzero closed-frequency zero of the completed quotient determinant.
Both the bulk and bath profiles decay exponentially in their respective
normal half-lines.  The zero is therefore a trapped neutral surface mode,
not a coordinate-chart zero.
\end{theorem}

\begin{proof}
The function in \eqref{eq:v3v42-F} is continuous on
\([0,c)\) and
\begin{equation}
 F_{c,\beta}(0)=\beta^2c^2>0.
\label{eq:v3v42-F0}
\end{equation}
As \(x\uparrow c\), both terms containing
\(\sqrt{c^2-x^2}\) vanish, while
\begin{equation}
 F_{c,\beta}(x)\longrightarrow
 -4c^6\sqrt{1-c^2}<0.
\label{eq:v3v42-Fc}
\end{equation}
The intermediate value theorem gives \(x_*\in(0,c)\).
At that point \(s_*,k,\kappa,\rho\) and the chart factors are nonzero, so
Theorem~\ref{thm:v3v42-determinant} makes the quotient determinant zero.
Equation \eqref{eq:v3v42-evanescent-roots} gives strict normal decay in
both half-lines.
\end{proof}

\begin{assessment}[Why positive energy permits this root]
\label{ass:v3v42-neutral-energy}
The enlarged metric--material--bath system is conservative.  Positive
energy excludes roots with \(\operatorname{Re}s>0\), but it does not
exclude a purely imaginary eigenfrequency carrying finite energy.  The
mode in Theorem~\ref{thm:v3v42-surface-root} is the boundary analogue of a
Rayleigh wave.  Passivity and exact glancing nonvanishing are therefore
insufficient for a uniform closed-frequency Lopatinski estimate.
\end{assessment}

\begin{corollary}[The strict speed interval is eliminated]
\label{cor:v3v42-speed-eliminated}
No fixed speed
\(0<c_{\rm e}<1\) in the V35 family can supply a zero-free principal
boundary symbol for the completed physical material system of V38.
A uniform family containing any such speed also fails.
\end{corollary}

\begin{proof}
Apply Theorem~\ref{thm:v3v42-surface-root} to the chosen speed and any
positive incidence \(\beta\).
\end{proof}

\begin{firewall}[The cold endpoint requires a separate calculation]
\label{fw:v3v42-cold-open}
The proof of Theorem~\ref{thm:v3v42-surface-root} uses the interval
\(0<x<c_{\rm e}\), on which the bath root is positive real.  That interval
collapses at \(c_{\rm e}=0\).  The cold line response
\(\rho=s\) is not covered by the counterexample and remains the only
member of the V35 speed family not excluded by V42.  Its quotient
determinant, the \(k=0\) chart, and the static \(s=0\) material sector must
be tested directly.
\end{firewall}

\begin{firewall}[Superseded scope of V39]
\label{fw:v3v42-V39-superseded}
Theorem~\ref{thm:v3v39-no-closed-root} is correct for the restricted
system \eqref{eq:v3v39-bulk-wave}--\eqref{eq:v3v39-junction}, in which the
multiplier-induced mixed stress is absent.  It is not a theorem about the
completed covariant material Hessian.  The completed determinant is
\eqref{eq:v3v42-determinant}, and
Theorem~\ref{thm:v3v42-surface-root} is the controlling result for every
positive subluminal speed.
\end{firewall}

\subsection{V42 conclusion and V43 target}
\label{sec:v3v42-conclusion}

V42 constructs a genuine local quotient chart and retains the complete
signed stress produced by the comoving multiplier.  Its exact determinant
is \eqref{eq:v3v42-determinant}.  The physical bath lifts the
one-dimensional exact glancing kernel, but that local success does not
extend across the closed hemisphere.  Every speed
\(0<c_{\rm e}<1\) creates a trapped neutral surface wave at a frequency
strictly below the bath cone.

The result forces another upstream restriction: among the V35 speed
family, only the cold line endpoint \(c_{\rm e}=0\) remains viable for
this material action.  V43 should substitute \(\rho=s\) into the exact
determinant, factor its closed-frequency polynomial, and test for all
\(\tau/k\), including the \(k=0\) axis in a second quotient chart.  It
must then add either a covariant static elastic modulus or a proved
material-relabeling quotient at \(s=0\).  If the cold determinant also has
a surface root, the material bath action must be changed rather than
retuning its speed.
\clearpage
\section{V43: the cold-line quotient is zero-free at every nonstatic closed frequency}
\label{sec:v3v43}

\subsection{Purpose}
\label{sec:v3v43-purpose}

V42 excludes every positive speed in the physical material bath by an
explicit neutral surface mode.  The argument does not include the cold
endpoint \(c_{\rm e}=0\), whose retarded root is
\(\rho=s\).  V43 substitutes that root in the completed quotient
determinant and proves that its phase prevents a closed-frequency zero.
Because the V42 quotient chart degenerates at zero spatial frequency, a
second independent seven-row chart is constructed on that axis.

The result covers every nonzero imaginary frequency, including exact
bulk glancing and the spatially homogeneous axis.  It does not divide by
the comoving factor at \(s=0\); the static multiplier sector remains a
separate gate.

\subsection{Cold determinant for nonzero spatial frequency}
\label{sec:v3v43-cold-kpositive}

Take
\begin{equation}
 c_{\rm e}=0,\qquad m_{\rm e}=0,\qquad
 \rho(s,k)=s\quad(\operatorname{Re}s>0),
\label{eq:v3v43-cold-root}
\end{equation}
with its limiting retarded value on the imaginary axis.  Absorb the
positive coupling constants in
\begin{equation}
 \mu=\beta s,\qquad \beta>0.
\label{eq:v3v43-mu}
\end{equation}

\begin{proposition}[Cold quotient factor]
\label{prop:v3v43-cold-factor}
On the V42 chart \(s\ne0,\ k>0\), substitution of
\eqref{eq:v3v43-mu} in
\eqref{eq:v3v42-determinant} gives
\begin{equation}
 \det\mathcal B_{\rm q}^{\rm cold}
 =
 -\frac{k^2}{2}\mathcal P_\beta(s,k),
\label{eq:v3v43-cold-det}
\end{equation}
where
\begin{align}
 \mathcal P_\beta(s,k)
 ={}&
 \kappa\beta^2(k^2+2s^2)^2
 +\beta s(k^4+8s^2k^2+8s^4)
 +4s^4\kappa,
\notag\\
 &\hspace{35mm}
 \kappa=\sqrt{s^2+k^2}.
\label{eq:v3v43-P}
\end{align}
The apparent powers of \(s\) which were introduced by clearing the two
comoving denominators cancel exactly.
\end{proposition}

\begin{proof}
Insert \(\mu=\beta s\) in the bracket of
\eqref{eq:v3v42-determinant}.  Every term contains \(s^2\):
\[
 \kappa\beta^2s^2(k^2+2s^2)^2
 +\beta s^3(k^4+8s^2k^2+8s^4)
 +4s^6\kappa
 =
 s^2\mathcal P_\beta.
\]
Cancel this factor against the left side
\(s^2\det\mathcal B_{\rm q}\).  This proves
\eqref{eq:v3v43-cold-det}.
\end{proof}

\begin{theorem}[No cold-line closed-frequency root for \(k>0\)]
\label{thm:v3v43-no-closed-kpositive}
For every \(\beta>0\),
\begin{equation}
 \mathcal P_\beta(i\tau,k)\ne0
\label{eq:v3v43-no-closed-kpositive}
\end{equation}
whenever \(k>0\) and \(\tau\ne0\), using the limiting outgoing branch of
\(\kappa\).
\end{theorem}

\begin{proof}
By conjugation symmetry it is enough to take \(\tau>0\).  Put
\(x=\tau/k\).

If \(0<x<1\), then
\(\kappa=k\sqrt{1-x^2}>0\).  Dividing
\eqref{eq:v3v43-P} by \(k^5\), its real part is
\begin{equation}
 \sqrt{1-x^2}
 \left[
 \beta^2(1-2x^2)^2+4x^4
 \right]>0.
\label{eq:v3v43-subglancing-real}
\end{equation}
The middle term is purely imaginary, so no zero is possible.

If \(x=1\), then \(\kappa=0\) and
\begin{equation}
 \frac{\mathcal P_\beta(ik,k)}{k^5}
 =
 i\beta(1-8+8)=i\beta\ne0.
\label{eq:v3v43-glancing-value}
\end{equation}

If \(x>1\), then
\(\kappa=ik\sqrt{x^2-1}\).  All of
\(\mathcal P_\beta/k^5\) is imaginary and its coefficient is
\begin{align}
 &\sqrt{x^2-1}
 \left[
 \beta^2(1-2x^2)^2+4x^4
 \right]
 \notag\\
 &\hspace{18mm}
 +\beta x(1-8x^2+8x^4).
\label{eq:v3v43-superglancing-imag}
\end{align}
The first line is strictly positive.  The quadratic
\(8y^2-8y+1\) has roots
\((2\pm\sqrt2)/4<1\), so its value is positive for
\(y=x^2>1\).  The second line is therefore positive as well.
Thus \eqref{eq:v3v43-superglancing-imag} cannot vanish.
\end{proof}

\begin{corollary}[Cold line removes the completed quotient glancing zero]
\label{cor:v3v43-cold-glancing}
At \(s=i\tau,\ k=|\tau|\ne0\), the completed quotient determinant has
the nonzero value obtained from
\eqref{eq:v3v43-glancing-value}.  This is the exact covariant-multiplier
version of the cold-line glancing response.
\end{corollary}

\begin{proof}
Combine Proposition~\ref{prop:v3v43-cold-factor} with
\eqref{eq:v3v43-glancing-value}.
\end{proof}

\subsection{No growing cold-line root}
\label{sec:v3v43-growing}

\begin{theorem}[Positive energy excludes the open right half-plane]
\label{thm:v3v43-no-growing}
The completed cold-line quotient has no nonzero stable normal mode with
\begin{equation}
 \operatorname{Re}s>0.
\label{eq:v3v43-RHP}
\end{equation}
\end{theorem}

\begin{proof}
Before elimination, the cold half-line field obeys the local wave
equation
\[
 \Psi_{tt}-\Psi_{yy}=0
\]
at each boundary spatial point.  The gravitational bulk energy, any
nonnegative induced boundary energy, and the half-line energy are
nonnegative, and the relative junction signs in
Lemma~\ref{lem:v3v42-completed-rows} were derived from their single
variational flux cancellation.  The comoving multiplier imposes a
constraint and contributes no independent negative quadratic energy.

A stable mode with \(\operatorname{Re}s>0\) and nonzero metric or bath
amplitude would make this conserved positive energy grow exponentially.
Equivalently, the Laplace-transformed energy identity gives
\[
 \operatorname{Re}s\,
 (E_{\rm g}+E_\partial+E_{\rm R})=0.
\]
The energies must vanish, forcing the metric stable mode and bath trace
to vanish.  The quotient gauge slice then removes the remaining
tangential diffeomorphism amplitudes.  Hence no nonzero quotient mode
exists.
\end{proof}

\begin{assessment}[Energy and determinant have different jobs]
\label{ass:v3v43-energy-vs-det}
The energy identity excludes exponential growth but permits conservative
imaginary-frequency surface waves.  Theorem
\ref{thm:v3v43-no-closed-kpositive} is therefore essential.  Conversely,
the closed-axis polynomial does not replace the open-half-plane energy
argument.
\end{assessment}

\subsection{A second quotient chart at \(k=0\)}
\label{sec:v3v43-axis}

When \(k=0\) and \(s\ne0\), the gauge slice
\eqref{eq:v3v42-gauge-slice} remains valid because its FP determinant is
\(2s^3\).  The V42 row choice degenerates, so retain instead
\begin{equation}
 (\mathcal J^+_{00},\mathcal J^+_{01},
 \mathcal J^+_{02},\mathcal J^+_{11},
 \mathcal J^+_{12},\mathcal J^+_{22},C_0).
\label{eq:v3v43-axis-rows}
\end{equation}
At \(k=0\), the material gradient \(D_\xi\) and the multiplier-induced
mixed stress vanish.  The spatial strain stress remains.

\begin{proposition}[Exact axis determinant]
\label{prop:v3v43-axis-det}
Let \(\mu\) denote the completed spatial response coefficient.  In the
variables \eqref{eq:v3v42-variables}, the determinant of the row chart
\eqref{eq:v3v43-axis-rows} is
\begin{equation}
 \det\mathcal B_{\rm axis}
 =
 -2s^5(s+\mu)^2.
\label{eq:v3v43-axis-det}
\end{equation}
For the cold line, \(\mu=\beta s\), and hence
\begin{equation}
 \det\mathcal B_{\rm axis}^{\rm cold}
 =
 -2(1+\beta)^2s^7\ne0
 \qquad(s\ne0).
\label{eq:v3v43-cold-axis}
\end{equation}
\end{proposition}

\begin{proof}
On the stable branch at \(k=0\), \(\kappa=s\).  The free retained rows
reduce to
\begin{align*}
 \mathcal J^+_{00}&=-s(H_{11}+H_{22}),&
 \mathcal J^+_{01}&=-sV_1,&
 \mathcal J^+_{02}&=-sV_2,\\
 \mathcal J^+_{11}&=s(H_{22}-2V_0),&
 \mathcal J^+_{12}&=-sH_{12},&
 \mathcal J^+_{22}&=s(H_{11}-2V_0),\\
 C_0&=-\frac s2(H_{11}+H_{22}+w)-sV_0.
\end{align*}
The material form adds
\(-\mu\mathbb P_0H_{\rm sp}\) to the three spatial junction rows.
Expansion in the order
\((V_1,V_2,H_{12})\) followed by
\((H_{11},H_{22},V_0,w)\) gives
\(-2s^5(s+\mu)^2\).  Substitution of
\(\mu=\beta s\) proves \eqref{eq:v3v43-cold-axis}.
\end{proof}

\begin{theorem}[All nonstatic closed frequencies are covered]
\label{thm:v3v43-all-nonstatic}
For the completed massless cold-line material system with \(\beta>0\),
the physical quotient boundary symbol has no zero at any
\begin{equation}
 s=i\tau,\qquad \tau\ne0,\qquad k\ge0.
\label{eq:v3v43-all-frequency}
\end{equation}
\end{theorem}

\begin{proof}
For \(k>0\), use
Theorem~\ref{thm:v3v43-no-closed-kpositive}.  For \(k=0\), use
\eqref{eq:v3v43-cold-axis}.
\end{proof}

\subsection{The static face}
\label{sec:v3v43-static}

At \(s=0,\ k>0\), the retarded cold root also vanishes.  The local
half-line equation becomes
\begin{equation}
 -\partial_y^2\Psi=0.
\label{eq:v3v43-static-line}
\end{equation}
A decaying finite-energy solution has zero trace.  Hence the strain trace
must satisfy \(e=0\).  By
Lemma~\ref{lem:v3v39-D-invertible}, this determines
\(\pi=D_\xi^{-1}\mathbb P_0H_{\rm sp}\) for a given metric trace.
However, the material equation
\(-s\lambda+D_\xi^\dagger R=0\) loses its \(\lambda\) term.

\begin{firewall}[Static multiplier degeneracy]
\label{fw:v3v43-static-multiplier}
At \(s=0\), the strain force \(R\) vanishes on the finite-energy cold-line
trace, and the material Euler equation does not determine the multiplier
\(\lambda\).  That multiplier enters the mixed metric junction stress.
Therefore neither the \(1/s\) quotient used in V42 nor the nonstatic axis
chart proves invertibility of the enlarged static system.  A Dirac
constraint reduction, a material-relabeling quotient, or a positive
inertial/elastic replacement for the hard comoving multiplier is required.
The joint origin \((s,k)=(0,0)\) also retains the infrared half-line zero
mode already isolated in V35.
\end{firewall}

\begin{assessment}[Current speed classification]
\label{ass:v3v43-speed-classification}
For the completed physical material action:
\begin{enumerate}
 \item every \(0<c_{\rm e}<1\) is excluded by the V42 neutral surface
       mode;
 \item \(c_{\rm e}=1\) loses the homogeneous glancing response by V34--V35;
 \item \(c_{\rm e}=0\) has no growing mode and no nonstatic
       closed-frequency quotient zero by V43; and
 \item the static multiplier face of the cold theory remains unresolved.
\end{enumerate}
Thus the speed family has collapsed to a single analytically viable
endpoint, subject to the static constraint reduction.
\end{assessment}

\subsection{V43 conclusion and V44 target}
\label{sec:v3v43-conclusion}

V43 proves that the massless cold line survives the completed material
stress calculation.  Its \(k>0\) quotient polynomial has no imaginary-axis
zero because the bulk-normal and line responses occupy incompatible
phases below glancing and have the same noncancelling phase above
glancing.  A second quotient chart proves
\(-2(1+\beta)^2s^7\ne0\) on the zero-spatial-frequency axis.
Positive energy excludes the open right half-plane.

The remaining obstruction is no longer the nonstatic boundary symbol.
It is the static degeneracy of the hard comoving multiplier.  V44 should
perform the finite-dimensional Dirac reduction of
\((\pi,\lambda)\) at \(s=0\), determine whether \(\lambda\) is a pure
constraint reaction removed with a Ward row or a genuine undetermined
junction source, and compare that result with a positive quadratic
velocity penalty replacing
\(\lambda_Iu^a\nabla_aX^I\).  Any replacement must be retested for neutral
surface modes before the cold-line boundary domain is accepted.
\clearpage
\section{V44: static cold-line kernel and the elastic-strain repair}
\label{sec:v3v44}

\subsection{Why the zero-frequency line domain matters}
\label{sec:v3v44-purpose}

V43 proved zero-freeness for every nonstatic closed frequency and left the
static face to a constraint reduction.  At \(s=0\), the cold half-line
equation is \(-\partial_y^2\Psi=0\).  Requiring literal decay at infinity
forces zero trace, but the natural cold-line energy contains only
\(|\partial_y\Psi|^2\) at zero frequency.  A constant profile has zero
energy density and is the limit of the retarded solution as
\(s\to0\).  Thus finite energy alone does not select the decaying
zero-trace domain.

V44 keeps that constant zero-energy profile and computes the complete
static unitary-reference kernel.  The multiplier itself is zero on the
kernel.  The missing term is a static strain modulus.  A local positive
elastic energy removes the finite-frequency static kernel through the
material Euler equation.  Uniform high-frequency control still requires
the mixed-order material norm and is not claimed here.

\subsection{Static unitary-reference system}
\label{sec:v3v44-static-system}

Fix
\begin{equation}
 s=0,\qquad k>0,\qquad \kappa=k.
\label{eq:v3v44-static-frequency}
\end{equation}
Use the physical reference fields as boundary coordinates:
\begin{equation}
 \vartheta=0,\qquad \pi^I=0.
\label{eq:v3v44-unitary}
\end{equation}
The comoving rows \eqref{eq:v3v41-comoving-symbol} then impose
\begin{equation}
 H_{01}=H_{02}=0.
\label{eq:v3v44-comoving-static}
\end{equation}
Let \(\lambda_1,\lambda_2\) be retained as independent constraint
reactions.  When the static cold-line force is zero, the variables are
\begin{equation}
 Y=
 (H_{00},H_{11},H_{12},H_{22},
 V_0,V_1,V_2,w,\lambda_1,\lambda_2)^{\mathsf T}.
\label{eq:v3v44-Y}
\end{equation}
Use all six Brown--York rows and all four de Donder rows.

\begin{lemma}[Exact static rows]
\label{lem:v3v44-static-rows}
After division by the common nonzero factor \(k\) and after absorbing the
nonzero junction normalization into \(\lambda\), the ten rows are
\begin{align}
 \mathcal J_{00}&=-H_{11}-H_{22}-2iV_1,
&
 \mathcal J_{01}&=-iV_0-\lambda_1,
\label{eq:v3v44-J00-J01}\\
 \mathcal J_{02}&=-\lambda_2,
&
 \mathcal J_{11}&=-H_{00}+H_{22},
\label{eq:v3v44-J02-J11}\\
 \mathcal J_{12}&=-H_{12}-iV_2,
&
 \mathcal J_{22}&=-H_{00}+H_{11}+2iV_1,
\label{eq:v3v44-J12-J22}\\
 C_0&=-V_0,
&
 C_1&=iH_{11}-\frac i2(w+\theta)-V_1,
\label{eq:v3v44-C0-C1}\\
 C_2&=iH_{12}-V_2,
&
 C_x&=iV_1-\frac12(w-\theta),
\label{eq:v3v44-C2-Cx}\\
 \theta&=-H_{00}+H_{11}+H_{22}.
\notag
\end{align}
\end{lemma}

\begin{proof}
Set \(s=0\), \(\zeta=(0,ik,0)\), and \(\kappa=k\) in
\eqref{eq:v3v40-Ca}--\eqref{eq:v3v40-Jplus}.  The comoving multiplier
stress contributes only to the \(01\) and \(02\) junction rows.  Divide by
\(k\) and rescale the two nonzero multiplier coefficients.  This gives
the displayed system.
\end{proof}

\begin{theorem}[Two-dimensional static cold-line kernel]
\label{thm:v3v44-static-kernel}
The matrix in Lemma~\ref{lem:v3v44-static-rows} has rank eight.  Its kernel
is
\begin{align}
 H_{11}&=a,&
 V_1&=\frac{i}{2}a,
\label{eq:v3v44-kernel-long}\\
 H_{12}&=b,&
 V_2&=ib,
\label{eq:v3v44-kernel-trans}
\end{align}
with all other components in \eqref{eq:v3v44-Y} equal to zero and
\(a,b\in\mathbb C\) arbitrary.  In particular,
\begin{equation}
 \lambda_1=\lambda_2=0
\label{eq:v3v44-lambda-zero}
\end{equation}
on the kernel.
\end{theorem}

\begin{proof}
The \(C_0,\mathcal J_{01},\mathcal J_{02}\) rows give
\(V_0=\lambda_1=\lambda_2=0\).  The
\((C_2,\mathcal J_{12})\) rows give
\(V_2=iH_{12}\), leaving \(H_{12}=b\) free.  The
\(\mathcal J_{11}\) row gives \(H_{22}=H_{00}\).
The remaining scalar rows yield
\[
 H_{00}=H_{22}=w=0,\qquad
 V_1=\frac i2H_{11},
\]
and leave \(H_{11}=a\) free.  Direct substitution proves that both
displayed directions lie in the kernel.  All other variables have been
determined, so the kernel has dimension two and the rank is eight.
\end{proof}

\begin{proposition}[Geometric identity of the static modes]
\label{prop:v3v44-static-gauge-images}
The two kernel directions are the metric images of the spatial
tangential diffeomorphisms
\begin{equation}
 \chi_1=\frac{a}{2ik},\qquad
 \chi_2=\frac{b}{ik},
\label{eq:v3v44-chi}
\end{equation}
while the physical material coordinates are held fixed by
\eqref{eq:v3v44-unitary}.  Their strains are nonzero:
\begin{equation}
 e=-D_\xi\chi\ne0
\label{eq:v3v44-nonzero-strain}
\end{equation}
unless \(a=b=0\).
\end{proposition}

\begin{proof}
The stable tangential gauge symbol is
\[
 H_{ij}=i\xi_i\chi_j+i\xi_j\chi_i,
 \qquad
 V_i=-k\chi_i.
\]
With \(\xi=(k,0)\), the choices
\eqref{eq:v3v44-chi} reproduce
\eqref{eq:v3v44-kernel-long}--%
\eqref{eq:v3v44-kernel-trans}.  In the fully covariant theory the same
diffeomorphism also changes \(\pi\) by \(\chi\) and leaves the strain
zero.  Unitary reference gauge holds \(\pi=0\); consequently
\(e=-\mathbb P_0H_{\rm sp}=-D_\xi\chi\), which is nonzero by
Lemma~\ref{lem:v3v39-D-invertible}.
\end{proof}

\begin{assessment}[The obstruction is zero static stiffness]
\label{ass:v3v44-zero-stiffness}
The multiplier reactions do not cause the kernel:
\eqref{eq:v3v44-lambda-zero} shows that they vanish on it.  The cold line
has retarded impedance proportional to \(s\), so its force vanishes at
\(s=0\) even on nonzero strain.  The two former gauge directions become
physical relative-strain directions after the material reference is
fixed, and no static term pays them.
\end{assessment}

\begin{firewall}[Decay at exactly zero frequency is a different domain]
\label{fw:v3v44-zero-domain}
Imposing \(\Psi(y)\to0\) at \(s=0\) forces the constant cold-line trace to
zero and removes \eqref{eq:v3v44-kernel-long}--\eqref{eq:v3v44-kernel-trans}.  However, the retarded
solutions for \(s>0\) are
\(e^{-sy}\Psi(0)\) and converge locally to a nonzero constant as
\(s\downarrow0\).  The zero-trace decay domain is therefore not the closed
graph limit of the retarded domains in the energy seminorm.  A uniform
continuum theorem cannot switch to it at one frequency without specifying
a stronger \(L^2_y\) infrared domain or a regulator and proving uniform
removal.
\end{firewall}

\subsection{A covariant local elastic term}
\label{sec:v3v44-elastic}

Add to the physical boundary action
\begin{equation}
 S_{\rm el}
 =
 -\frac{G}{2}\int_{\mathcal N}\sqrt{-\gamma}\,
 |e|_q^2\,d^3x,
 \qquad G>0,
\label{eq:v3v44-Sel-L}
\end{equation}
where the Lorentzian sign is potential-energy sign.  Its Euclidean
continuation is
\begin{equation}
 S_{\rm el}^{E}
 =
 \frac{G}{2}\int_{\mathcal N_E}\sqrt{\gamma_E}\,
 |e|_q^2\,d^3x\ge0.
\label{eq:v3v44-Sel-E}
\end{equation}

\begin{proposition}[Covariance and reciprocal elastic rows]
\label{prop:v3v44-elastic-covariant}
The elastic action is boundary-diffeomorphism covariant and
reparametrization invariant in the same sense as
Proposition~\ref{prop:v3v38-covariant-strain}.  At the flat quadratic
level its strain force is
\begin{equation}
 R_{\rm el}=G e.
\label{eq:v3v44-Rel}
\end{equation}
The material and metric equations contain respectively
\begin{equation}
 D_\xi^\dagger R_{\rm el},
 \qquad
 -\mathbb P_0R_{\rm el},
\label{eq:v3v44-reciprocal}
\end{equation}
which are adjoints arising from one Hessian.
\end{proposition}

\begin{proof}
The scalar \(|e|_q^2\) is built from the covariant tensor
\eqref{eq:v3v38-e} and the positive clock-spatial metric.  This proves the
invariance claims.  Varying the quadratic form gives
\(\langle Ge,\delta e\rangle\).  Since
\(\delta e=D_\xi\delta\pi-\mathbb P_0\delta H_{\rm sp}\), the two Euler
rows are exactly \eqref{eq:v3v44-reciprocal}.
\end{proof}

\begin{theorem}[Elastic term removes the finite-frequency static kernel]
\label{thm:v3v44-elastic-removes}
In unitary reference gauge at \(s=0,\ k>0\), append the material Euler rows
from \(S_{\rm el}\).  Then the homogeneous static system has no nonzero
solution:
\begin{equation}
 H=V=w=\lambda=0.
\label{eq:v3v44-static-zero}
\end{equation}
\end{theorem}

\begin{proof}
On the static kernel of the gravitational and comoving rows,
Proposition~\ref{prop:v3v44-static-gauge-images} gives
\(e=-D_\xi\chi\).  The material Euler equation is
\[
 D_\xi^\dagger G D_\xi\chi=0.
\]
Pair it with \(\chi\).  Lemma~\ref{lem:v3v39-D-invertible} gives
\[
 G\|D_\xi\chi\|_{\rm HS}^2
 =2Gk^2|\chi|^2=0,
\]
so \(\chi=0\), hence \(a=b=0\).  The proof of
Theorem~\ref{thm:v3v44-static-kernel} then forces every remaining metric
and multiplier amplitude to vanish.
\end{proof}

\begin{corollary}[Positive static strain gap]
\label{cor:v3v44-static-strain-gap}
On the two-dimensional kernel subspace of the unmodified static system,
the elastic material row has the exact coercivity
\begin{equation}
 \langle\chi,D_\xi^\dagger GD_\xi\chi\rangle
 =2Gk^2|\chi|^2.
\label{eq:v3v44-static-strain-gap}
\end{equation}
\end{corollary}

\begin{proof}
This is \eqref{eq:v3v39-D-norm} multiplied by \(G\).
\end{proof}

\subsection{What changes dynamically}
\label{sec:v3v44-dynamic}

For \(s\ne0\), the elastic force and cold-line force add before the
material and metric variations:
\begin{equation}
 R_{\rm tot}
 =
 (G+\alpha^2s)e.
\label{eq:v3v44-total-force}
\end{equation}
Thus the completed quotient formula of V42 has
\begin{equation}
 \mu(s)=g+\beta s,
 \qquad g=4\kappa_{\rm g}G>0,\quad\beta>0,
\label{eq:v3v44-mu-dynamic}
\end{equation}
with the same relative mixed and spatial row coefficients.

\begin{firewall}[Static repair must be retested dynamically]
\label{fw:v3v44-retest}
Theorem~\ref{thm:v3v44-elastic-removes} is a static no-kernel theorem.
It does not prove that substituting
\(\mu=g+\beta s\) in
\eqref{eq:v3v42-determinant} leaves the whole closed-frequency hemisphere
zero-free.  A conservative elastic term can support surface waves, while
the cold channel can move or damp them.  The exact determinant must be
tested before the repair is admitted.
\end{firewall}

\begin{firewall}[Mixed-order uniformity]
\label{fw:v3v44-mixed-order}
The metric junction contribution of a fixed \(G\) is order zero in the
metric trace, whereas the material Euler row
\(D^\dagger GD\pi\) is order two in the material displacement and order
one in the metric strain.  Uniform control therefore belongs to a
Douglis--Nirenberg or graph norm with different weights for metric and
material variables.  Dividing only the metric rows by \(k\) makes the
static elastic coefficient appear to vanish at high frequency and does
not test the complete material row.  V44 proves finite-frequency
injectivity, not the required uniform mixed-order estimate.
\end{firewall}

\subsection{V44 conclusion and V45 target}
\label{sec:v3v44-conclusion}

V44 resolves the nature of the cold static face.  When constant
zero-energy line profiles are retained, the unitary-reference boundary
matrix has a two-dimensional kernel.  Its elements are spatial
diffeomorphism metric images with the physical material coordinates held
fixed, hence nonzero relative strains.  The multiplier reactions vanish
on the kernel.  The missing currency is static stiffness.

The covariant positive elastic form \(G\|e\|^2/2\) removes that kernel
through the exact coercivity \(2Gk^2|\chi|^2\).  It changes the nonstatic
response to \(\mu=g+\beta s\), which must still pass the complete
determinant, and its uniform estimate requires mixed metric/material
weights.

V45 is a mandatory companion audit.  It should then substitute
\(\mu=g+\beta s\) into the V42 quotient polynomial, determine whether any
closed-frequency elastic surface root survives, and formulate the
Douglis--Nirenberg weights of the local metric--material--line system.
If the local elastic repair creates an unavoidable neutral root, the next
upstream candidate is a two-channel passive bath whose cold component
supplies damping and whose spatial component supplies homogeneous static
stiffness.
\clearpage
\section{V45: companion audit and a zero-free two-channel physical bath}
\label{sec:v3v45}

\subsection{Mandatory read-only companion audit}
\label{sec:v3v45-audit}

The V45 audit inspected the newest companion notes without changing them.
Their identities were
\begin{align}
 &\texttt{Renewal\_NS\_Stability\_Research\_Notes\_V31.tex},
 \notag\\
 &\qquad 887537\ {\rm bytes},\qquad
 \texttt{F1121B62238AD5491BB7CF03635EA9934942EDF573ED8FAAA9EE79E1D38A6696},
\label{eq:v3v45-NS-fingerprint}\\
 &\texttt{Renewal\_Yang\_Mills\_Mass\_Gap\_Research\_Notes\_V91.tex},
 \notag\\
 &\qquad 1477832\ {\rm bytes},\qquad
 \texttt{307FCCDF99747D923860D4129A0299317A6B4B30934083DBD66E28A068D75174}.
\label{eq:v3v45-YM-fingerprint}
\end{align}
The hashes are SHA--256 digests.  All frozen companion volumes and SM
volumes were left untouched.

Yang--Mills V90 constructs a tilted coherent counterexample to a proposed
local signed card.  V91 then proves that a singleton crossing bank and
mesoscopic prefixes transmit the obstruction: moving one level upstream
does not help unless an exact coupled cancellation is present.
Navier--Stokes V31 proves that removal of a probability mark restores the
exact missing first-moment scale, and that a signed hinge with a nullspace
cannot pay the positive total sector.

\begin{assessment}[Consequences for the static repair]
\label{ass:v3v45-audit-consequence}
The fixed elastic modulus \(G\) in V44 removes the finite-frequency static
kernel, but in metric junction rows it is order zero.  After homogeneous
row normalization its coefficient is \(G/\lambda\).  The NS scale lesson
forbids treating that vanishing factor as an order-one static payment.
The YM counterexample lesson requires the V42 determinant to be tested
after all physical channels are combined; a claim that damping will
presumably move a surface wave is insufficient.
\end{assessment}

\subsection{Why the fixed elastic modulus is not the principal repair}
\label{sec:v3v45-fixed-G}

Use the frequency scale
\begin{equation}
 \lambda=(|s|^2+k^2)^{1/2}.
\label{eq:v3v45-lambda}
\end{equation}
Rescale material and multiplier variables by
\begin{equation}
 \widehat\pi=\lambda\pi,\qquad
 \widehat\lambda=\lambda^{-1}\lambda_{\rm mult}.
\label{eq:v3v45-rescaling}
\end{equation}
Then strain has order zero:
\begin{equation}
 e=
 \frac{D_\xi}{\lambda}\widehat\pi
 -\mathbb P_0H_{\rm sp}.
\label{eq:v3v45-normalized-strain}
\end{equation}
Junction rows are divided by \(\lambda\), material Euler rows by
\(\lambda^2\), and comoving rows are left at order zero.

\begin{proposition}[A fixed local modulus is lower order in the complete
graph norm]
\label{prop:v3v45-G-lower}
Under \eqref{eq:v3v45-rescaling}, the fixed elastic contributions satisfy
\begin{align}
 \lambda^{-1}G e&=O(G/\lambda)
 &&\text{in the metric junction rows},
\label{eq:v3v45-G-junction}\\
 \lambda^{-2}D_\xi^\dagger Ge&=O(G/\lambda)
 &&\text{in the material Euler rows}.
\label{eq:v3v45-G-material}
\end{align}
Thus \(G\) supplies finite-frequency static injectivity but no homogeneous
principal margin as \(\lambda\to\infty\) on \(s=0\).
\end{proposition}

\begin{proof}
The normalized strain in
\eqref{eq:v3v45-normalized-strain} has order zero.
The junction force \(Ge\) has no derivative, so division by the
order-one row scale gives \(G/\lambda\).  The material force has one
spatial derivative \(D_\xi^\dagger\), so it has size \(G\lambda\);
division by the order-two row scale again gives \(G/\lambda\).
\end{proof}

\begin{firewall}[Finite injectivity and uniform complementarity]
\label{fw:v3v45-finite-vs-uniform}
Theorem~\ref{thm:v3v44-elastic-removes} remains a valid
finite-frequency no-kernel result.  Proposition
\ref{prop:v3v45-G-lower} shows that its smallest normalized singular value
can collapse along high spatial frequencies.  The full continuum
programme requires the latter uniform estimate, so the fixed modulus
cannot be the final payer.
\end{firewall}

\subsection{Two local physical channels}
\label{sec:v3v45-two-channel}

Attach two independent physical tensor baths to the same covariant strain
\(e\):
\begin{enumerate}
 \item a cold line with speed zero and positive incidence \(\beta>0\);
 \item a spatially propagating bath with speed
       \(0<c<1\) and positive incidence \(\gamma>0\).
\end{enumerate}
Both are local half-line systems of the V35 type, have positive enlarged
energy, and share the clock/material covariance of V38.  Their retarded
roots are
\begin{equation}
 \rho_0=s,\qquad
 \rho_c=\sqrt{s^2+c^2k^2},
\qquad
 \operatorname{Re}\rho_0,\operatorname{Re}\rho_c>0
\label{eq:v3v45-roots}
\end{equation}
in the open right half-plane.  After absorbing positive gravitational
normalizations, the total completed response in the V42 determinant is
\begin{equation}
 \mu(s,k)=\beta s+\gamma\rho_c(s,k).
\label{eq:v3v45-mu}
\end{equation}

\begin{proposition}[Locality, energy, and reflection positivity]
\label{prop:v3v45-two-channel-parent}
The direct sum of the two bath actions is local and has nonnegative energy
equal to the sum of the two channel energies.  At a finite Euclidean
cutoff its Gaussian measure is reflection positive, and eliminating the
two baths adds the sum of their positive Euclidean
Dirichlet-to-Neumann forms.
\end{proposition}

\begin{proof}
Each summand has the properties proved in
Theorems~\ref{thm:v3v35-RP} and
\ref{thm:v3v38-RP}.  A direct sum of positive quadratic actions has
product Gaussian measure and additive energy.  Reflection positivity of
the product follows by tensoring the two reflected covariance
factorizations; marginalization over both baths preserves it.
\end{proof}

\begin{proposition}[Homogeneous static stiffness]
\label{prop:v3v45-static-stiffness}
At \(s=0,\ k>0\),
\begin{equation}
 \mu(0,k)=\gamma c k>0.
\label{eq:v3v45-static-mu}
\end{equation}
On the static strain kernel parameterized by \(\chi\), the material Euler
row has coercivity
\begin{equation}
 \left\langle
 \chi,D_\xi^\dagger(\gamma ck)D_\xi\chi
 \right\rangle
 =
 2\gamma c k^3|\chi|^2.
\label{eq:v3v45-static-coercivity}
\end{equation}
After the mixed-order normalization
\eqref{eq:v3v45-rescaling}, this is a nonzero homogeneous principal
coefficient.
\end{proposition}

\begin{proof}
Set \(s=0\) in \eqref{eq:v3v45-roots}--\eqref{eq:v3v45-mu}.
Then use \eqref{eq:v3v39-D-norm}.  The material row has order two in the
rescaled variables; division by \(k^2\) leaves the positive coefficient
\(2\gamma c\) after measuring \(\widehat\pi=k\pi\).
\end{proof}

\subsection{Exact closed-frequency determinant}
\label{sec:v3v45-closed}

Use the quotient determinant
\eqref{eq:v3v42-determinant}.  Let
\begin{equation}
 s=ixk,\qquad x>0,\qquad k>0,
\label{eq:v3v45-x}
\end{equation}
and define
\begin{equation}
 Q(x)=(1-2x^2)^2,\qquad
 P(x)=1-8x^2+8x^4.
\label{eq:v3v45-QP}
\end{equation}
After division by the nonzero common power of \(k\), the determinant
bracket is
\begin{equation}
 \mathcal F(x)
 =
 \widehat\kappa\,\widehat\mu^{\,2}Q(x)
 -x^2\widehat\mu\,P(x)
 -4x^6\widehat\kappa,
\label{eq:v3v45-F}
\end{equation}
where
\begin{equation}
 \widehat\kappa=\kappa/k,\qquad
 \widehat\mu=\mu/k.
\label{eq:v3v45-hat}
\end{equation}

\begin{theorem}[The two-channel determinant has no nonstatic closed root]
\label{thm:v3v45-no-closed-root}
For every
\begin{equation}
 \beta>0,\qquad\gamma>0,\qquad0<c<1,
\label{eq:v3v45-positive-parameters}
\end{equation}
one has
\begin{equation}
 \mathcal F(x)\ne0
 \qquad(x>0).
\label{eq:v3v45-no-closed-root}
\end{equation}
Consequently the completed quotient boundary symbol has no
imaginary-frequency zero for \(s\ne0,\ k>0\).
\end{theorem}

\begin{proof}
There are four frequency regions.

If \(0<x<c\), put
\[
 r=\sqrt{1-x^2},\qquad
 h=\gamma\sqrt{c^2-x^2}>0,\qquad
 m=\beta x>0.
\]
Then
\(\widehat\kappa=r\) and
\(\widehat\mu=h+im\).  Suppose
\(\mathcal F=0\).  Its imaginary part is
\begin{equation}
 m\bigl(2rhQ-x^2P\bigr)=0.
\label{eq:v3v45-imag-subcone}
\end{equation}
If \(Q=0\), then \(x^2=1/2\) and \(P=-1\), so
\eqref{eq:v3v45-imag-subcone} is nonzero.  If \(Q>0\),
the imaginary equation gives
\[
 h=\frac{x^2P}{2rQ}.
\]
Positivity of \(h\) forces \(P>0\).  Substitute
\(x^2P=2rhQ\) in the real part:
\begin{align*}
 \operatorname{Re}\mathcal F
 &=
 r(h^2-m^2)Q-x^2hP-4x^6r\\
 &=
 -r\bigl(h^2Q+m^2Q+4x^6\bigr)<0,
\end{align*}
a contradiction.

If \(c\le x<1\), write
\[
 \widehat\mu=iM,\qquad
 M=\beta x+\gamma\sqrt{x^2-c^2}>0,
 \qquad
 \widehat\kappa=\sqrt{1-x^2}>0.
\]
Then
\[
 \operatorname{Re}\mathcal F
 =
 -\widehat\kappa\bigl(M^2Q+4x^6\bigr)<0.
\]
At \(x=1\), \(\widehat\kappa=0\), \(P(1)=1\), and
\[
 \mathcal F(1)=-\widehat\mu(1)\ne0.
\]

Finally, if \(x>1\), put
\[
 \widehat\kappa=i\sqrt{x^2-1},\qquad
 \widehat\mu=iM,\qquad
 M=\beta x+\gamma\sqrt{x^2-c^2}>0.
\]
Then
\[
 \mathcal F
 =
 -i\left[
 \sqrt{x^2-1}\bigl(M^2Q+4x^6\bigr)
 +x^2MP
 \right].
\]
The first term in brackets is positive.  The polynomial
\(P=8x^4-8x^2+1\) is positive for \(x^2>1\), so the second term is
positive as well.  Hence \(\mathcal F\ne0\) in every case.
\end{proof}

\begin{assessment}[The cold channel removes the V42 surface wave]
\label{ass:v3v45-cold-removes}
Below the positive-speed bath cone, that channel contributes a positive
real root while the cold channel contributes \(i\beta x\).  The imaginary
equation then forces a relation which makes the real part strictly
negative.  This is an exact coupled cancellation obstruction to a zero,
not an appeal to damping after separate estimates.  Setting \(\beta=0\)
removes the decisive imaginary component and restores the V42
intermediate-speed surface root.
\end{assessment}

\subsection{Axes, static face, and the open half-plane}
\label{sec:v3v45-remaining-faces}

\begin{proposition}[Zero-spatial-frequency axis]
\label{prop:v3v45-axis}
At \(k=0\), both bath roots equal \(s\), so
\begin{equation}
 \mu(s,0)=(\beta+\gamma)s.
\label{eq:v3v45-axis-mu}
\end{equation}
The V43 axis determinant is
\begin{equation}
 \det\mathcal B_{\rm axis}
 =
 -2(1+\beta+\gamma)^2s^7\ne0
 \qquad(s\ne0).
\label{eq:v3v45-axis-det}
\end{equation}
\end{proposition}

\begin{proof}
Use \eqref{eq:v3v45-roots} in
\eqref{eq:v3v45-mu}, then apply
Proposition~\ref{prop:v3v43-axis-det} with the combined line
coefficient \(\beta+\gamma\).
\end{proof}

\begin{proposition}[Static face]
\label{prop:v3v45-static-face}
At \(s=0,\ k>0\), the unitary-reference static kernel of
Theorem~\ref{thm:v3v44-static-kernel} is removed by the positive-speed
channel.  The homogeneous material-row margin on that kernel is
\eqref{eq:v3v45-static-coercivity}.
\end{proposition}

\begin{proof}
Replace \(G\) by the positive order-one coefficient
\(\gamma ck\) in the proof of
Theorem~\ref{thm:v3v44-elastic-removes}.  The pairing with \(\chi\) is
\eqref{eq:v3v45-static-coercivity}, so \(\chi=0\); all remaining static
amplitudes then vanish as in V44.
\end{proof}

\begin{theorem}[No open-half-plane root]
\label{thm:v3v45-no-growing}
The local two-channel system has no finite-energy stable normal mode with
\(\operatorname{Re}s>0\).
\end{theorem}

\begin{proof}
The gravitational, cold-line, and positive-speed bath energies are
nonnegative, and the trace and junction signs come from one variational
flux identity.  The material comoving multiplier imposes the constraint
without adding a negative quadratic energy.  A mode with
\(\operatorname{Re}s>0\) would make the conserved positive energy grow.
The Laplace energy identity therefore forces every metric and bath
amplitude to vanish, after which the reference gauge removes the
tangential diffeomorphism amplitudes.
\end{proof}

\begin{theorem}[Uniform frozen principal margin]
\label{thm:v3v45-uniform-margin}
Equip the metric, material, and multiplier variables with the graph
weights \eqref{eq:v3v45-rescaling}, quotient the three tangential
diffeomorphism directions, and remove the three Ward/Codazzi row
relations.  On the normalized closed Laplace hemisphere excluding the
joint origin, the frozen two-channel principal boundary symbol has a
strictly positive minimum singular value.
\end{theorem}

\begin{proof}
The hemisphere is covered by three finite-dimensional quotient charts:
\begin{enumerate}
 \item \(s\ne0,\ k>0\), where
       Theorem~\ref{thm:v3v45-no-closed-root} gives injectivity;
 \item \(s\ne0,\ k=0\), where
       Proposition~\ref{prop:v3v45-axis} gives injectivity; and
 \item \(s=0,\ k>0\), where
       Proposition~\ref{prop:v3v45-static-face} gives injectivity in the
       mixed-order graph norm.
\end{enumerate}
The local augmented matrices and their quotient transition maps are
continuous on slightly smaller closed neighborhoods of these faces.  The
two-channel response is homogeneous of degree one, so no coefficient
vanishes under normalization.  Finite dimensionality makes injectivity
equivalent to a positive smallest singular value in each chart.
Compactness and a finite subcover give a positive global minimum.
\end{proof}

\begin{firewall}[Frozen principal theorem only]
\label{fw:v3v45-frozen-only}
Theorem~\ref{thm:v3v45-uniform-margin} is a flat, frozen, linear principal
boundary theorem.  Curved lower symbols, the exact clock/material
nonlinear equations, anomaly-free BRST domains, renormalized quantum
stress, and removal of all cutoffs remain open.  The theorem permits the
lower-symbol perturbation gate
\eqref{eq:v3v40-half-margin}; it does not prove its hypotheses on a global
curved configuration space.
\end{firewall}

\subsection{V45 conclusion and V46 target}
\label{sec:v3v45-conclusion}

The mandatory audit prevents promotion of the fixed elastic modulus to a
homogeneous payer.  V45 replaces it by a direct sum of two physical local
baths.  The cold line supplies a nonzero imaginary response at every
nonstatic frequency, while the positive-speed channel supplies
order-one static stiffness.  Their exact completed determinant has no
closed-frequency zero; the axis and static charts also close, and positive
energy excludes the open right half-plane.  With mixed metric/material
weights, the frozen principal quotient has a uniform margin.

V46 should now return from the flat symbol to a variable-coefficient
collar.  It must write the two-channel local action with projected bundle
connections, compute the complete first lower symbol generated by
curvature, second fundamental form, clock acceleration, material strain,
and noncommuting speed endomorphisms, and bound that perturbation relative
to the V45 margin.  The proof must state a concrete open set of geometric
backgrounds on which the clock gap, material-rank gap, coupling gap, and
boundary-symbol gap remain simultaneous.
\clearpage
\section{V46: variable-coefficient two-channel collars and the auxiliary glancing gate}
\label{sec:v3v46}

\subsection{Purpose and correction to the next inference}
\label{sec:v3v46-purpose}

V45 proves that the completed flat quotient determinant for a cold bath
plus a positive-speed bath has no root on the normalized closed Laplace
hemisphere.  Passing to a variable-coefficient collar requires more than
algebraic nonvanishing.  The positive-speed bath has its own characteristic
cone
\[
 s^2+c^2|\xi|_q^2=0,
\]
where its normal root vanishes.  The total determinant remains nonzero
there because the cold channel is active, but the usual
Dirichlet-to-Neumann lower-symbol expansion divides by that vanishing
root.

V46 writes the covariant local parent, computes the complete structure of
its first lower symbol, proves a quantitative perturbation theorem away
from the auxiliary cone, and gives an explicit regression showing why
that theorem cannot be extended across the cone by symbol continuity.
A local augmented glancing symmetrizer is the next required step.

\subsection{Covariant two-channel local action}
\label{sec:v3v46-action}

Retain the clock-spatial geometry
\((u,q)\), material strain \(e\), and positive fibre
\(\mathcal S_u\) of V38.  Let
\(\mathcal D\) be the connection obtained by orthogonally projecting the
boundary Levi--Civita connection to \(\mathcal S_u\).  For
\[
 A\in\{0,c\},\qquad c_0=0,\qquad0<c_c=c<1,
\]
let \(\Psi_A(y,x)\) be two independent fields on
\([0,\infty)_y\times\mathcal N\).  Take
\begin{align}
 S_A^L
 =
 \frac12\int_{\mathcal N}\sqrt{-\gamma}\,d^3x
 \int_0^\infty dy\,
 \Bigl(
 &|\mathcal D_u\Psi_A|_q^2
 -|\partial_y\Psi_A|_q^2
 \notag\\
 &-c_A^2q^{ab}
 \langle\mathcal D_a\Psi_A,\mathcal D_b\Psi_A\rangle_q
 \Bigr)
\label{eq:v3v46-action}
\end{align}
with trace conditions
\begin{equation}
 \Psi_0(0)=\alpha_0e,\qquad
 \Psi_c(0)=\alpha_ce,
\qquad
 \alpha_0,\alpha_c>0.
\label{eq:v3v46-traces}
\end{equation}
Mass terms are omitted from the homogeneous calculation; they are order
zero.

\begin{proposition}[Local covariance and energy]
\label{prop:v3v46-action-properties}
The sum
\[
 S_{\rm 2ch}^L=S_0^L+S_c^L
\]
is local and invariant under boundary diffeomorphisms acting on all
fields.  In a frozen clock frame its energy is the sum
\begin{align}
 E_{\rm 2ch}
 =
 \frac12\sum_{A\in\{0,c\}}
 \int_{\Sigma}\!\int_0^\infty
 \left(
 |\partial_t\Psi_A|^2+|\partial_y\Psi_A|^2
 +c_A^2|\nabla_x\Psi_A|^2
 \right)dy\,d^2x,
\label{eq:v3v46-energy}
\end{align}
and is nonnegative.
\end{proposition}

\begin{proof}
Every derivative and contraction in
\eqref{eq:v3v46-action} is natural, and the trace relations use the
covariant tensor \(e\).  This proves locality and covariance.  The
clock-spatial fibre metric is positive on \(\mathcal S_u\); the standard
time multiplier gives \eqref{eq:v3v46-energy} channel by channel.
\end{proof}

\subsection{Frozen principal and first lower symbols}
\label{sec:v3v46-symbols}

Fix a boundary point and choose an orthonormal clock frame there.  Let
\(\mathcal A_a\) denote the matrix of the projected bundle connection:
\begin{equation}
 \mathcal D_a=\partial_a+\mathcal A_a.
\label{eq:v3v46-connection}
\end{equation}
After the Euler equation is written in divergence form, its tangential
part has the local expression
\begin{equation}
 P_A
 =
 -\partial_y^2
 +a_A^{ab}\partial_a\partial_b
 +B_A^a\partial_a+C_A,
\label{eq:v3v46-local-operator}
\end{equation}
where
\begin{equation}
 a_A^{ab}=u^au^b+c_A^2q^{ab}
\label{eq:v3v46-aA}
\end{equation}
after Laplace sign convention, and \(B_A^a,C_A\) are bundle
endomorphisms.  The first-order coefficient contains both the scalar
divergence terms and the connection terms:
\begin{equation}
 B_A^a
 =
 b_A^aI
 +2u^au^b\mathcal A_b
 +2c_A^2q^{ab}\mathcal A_b.
\label{eq:v3v46-BA}
\end{equation}
Here \(b_A^a\) is obtained by differentiating
\(\sqrt{-\gamma}u^au^b\) and
\(\sqrt{-\gamma}c_A^2q^{ab}\); derivatives of the STF projector are
included in \(\mathcal A\).

\begin{proposition}[Complete first lower-symbol structure]
\label{prop:v3v46-first-lower}
For a Laplace covector \(s\,u^\flat+i\xi\), with
\(\xi(u)=0\), the principal tangential symbols are
\begin{equation}
 p_{A,2}(s,\xi)
 =
 s^2+c_A^2|\xi|_q^2,
\label{eq:v3v46-p2}
\end{equation}
and the complete degree-one symbols have the form
\begin{align}
 p_{A,1}(s,\xi)
 ={}&
 b_{A,u}s+i\,b_{A,\perp}^a\xi_a
 +2\mathcal A_us
 +2ic_A^2q^{ab}\mathcal A_a\xi_b.
\label{eq:v3v46-p1}
\end{align}
Every contribution of clock expansion, clock acceleration, derivatives of
the spatial projector, material-frame rotation, and the projected STF
connection occurs in one of the four terms in
\eqref{eq:v3v46-p1}.  Curvature, mass, and derivatives of these
coefficients occur in \(C_A\) and have degree zero.
\end{proposition}

\begin{proof}
Expand each divergence-form second derivative in
\eqref{eq:v3v46-action}.  Differentiating its scalar density and leading
tensor produces \(b_A^a\partial_a\).  Replacing one derivative in
\((\partial+\mathcal A)^2\) by the connection produces twice the two
connection contractions displayed in \eqref{eq:v3v46-BA}.  All terms with
two derivatives give \eqref{eq:v3v46-p2}; all terms with one derivative
give \eqref{eq:v3v46-p1}.  Commutators of projected derivatives and
undifferentiated coefficients have degree zero and belong to \(C_A\).
\end{proof}

Let
\begin{equation}
 \rho_{A,1}=\sqrt{p_{A,2}}
\label{eq:v3v46-rho1}
\end{equation}
with the outgoing branch.

\begin{lemma}[First Dirichlet-to-Neumann correction away from glancing]
\label{lem:v3v46-DtN-lower}
On a conic set satisfying
\begin{equation}
 |p_{A,2}(s,\xi)|\ge\varepsilon\lambda^2,
 \qquad
 \lambda=(|s|^2+|\xi|_q^2)^{1/2},
\label{eq:v3v46-separated}
\end{equation}
the order-zero correction \(r_{A,0}\) to the scalar-principal
Dirichlet-to-Neumann root solves
\begin{equation}
 \rho_{A,1}r_{A,0}+r_{A,0}\rho_{A,1}=p_{A,1}.
\label{eq:v3v46-Sylvester}
\end{equation}
Since \(\rho_{A,1}\) is scalar on the STF fibre,
\begin{equation}
 r_{A,0}=\frac{p_{A,1}}{2\rho_{A,1}}.
\label{eq:v3v46-r0}
\end{equation}
If the degree-one coefficient norm is bounded by
\(K_1\lambda\), then
\begin{equation}
 \|r_{A,0}\|
 \le\frac{K_1}{2\sqrt\varepsilon}.
\label{eq:v3v46-r0-bound}
\end{equation}
\end{lemma}

\begin{proof}
Factor the normal symbol to one lower order:
\[
 (\partial_y+\rho_{A,1}+r_{A,0})
 (-\partial_y+\rho_{A,1}+r_{A,0}).
\]
The degree-one cross term is the left side of
\eqref{eq:v3v46-Sylvester}.  Scalarity of the principal root gives
\eqref{eq:v3v46-r0}.  The separation condition implies
\(|\rho_{A,1}|\ge\sqrt\varepsilon\lambda\); combine it with the
hypothesis on \(p_{A,1}\).
\end{proof}

\subsection{Geometric perturbation size}
\label{sec:v3v46-geometric-size}

On a collar patch define
\begin{align}
 \mathfrak K_1
 :=\sup\bigl(
 &|K_{\mathcal N}|+|\nabla u|+|\mathcal A|
 +|\nabla\log\alpha_0|+|\nabla\log\alpha_c|
 +|\nabla c|
 \bigr),
\label{eq:v3v46-K1}\\
 \mathfrak K_2
 :=\sup\bigl(
 &|\operatorname{Rm}_\gamma|+|\nabla K_{\mathcal N}|
 +|\nabla\mathcal A|+|\mathcal A|^2
 +m_0^2+m_c^2
 \bigr).
\label{eq:v3v46-K2}
\end{align}
The material rank and inverse-frame bounds are kept separately below.

\begin{proposition}[Regular-cone normalized perturbation bound]
\label{prop:v3v46-regular-bound}
Fix \(\varepsilon>0\) and suppose both channel roots used in an eliminated
description satisfy \eqref{eq:v3v46-separated}.  Let \(E_\lambda\) be the
difference between the row-normalized variable-coefficient boundary
symbol and its exactly frozen principal symbol, in the mixed graph weights
of V45.  Then
\begin{equation}
 \|E_\lambda\|
 \le
 C_{\varepsilon,\mathcal P}
 \left(
 \frac{\mathfrak K_1}{\lambda}
 +\frac{\mathfrak K_2}{\lambda^2}
 +\|e_0\|\frac{\mathfrak K_1}{\lambda}
 \right),
\label{eq:v3v46-regular-bound}
\end{equation}
where \(\mathcal P\) denotes uniform bounds on the clock gap, material
frame condition number, speeds, and incidences.
\end{proposition}

\begin{proof}
The Brown--York and de Donder rows have order one.  Their connection and
second-fundamental-form corrections have order zero bounded by
\(C\mathfrak K_1\), while curvature and differentiated coefficients have
order minus one after the stable normal solve and are bounded by
\(C\mathfrak K_2/\lambda\).  Row normalization gives the first two ratios.

Lemma~\ref{lem:v3v46-DtN-lower} gives an order-zero bath correction
bounded by \(C_\varepsilon\mathfrak K_1\); division by the order-one
junction scale gives the same first ratio.  Linearizing the strain
projector at a nonzero background strain \(e_0\) produces only
undifferentiated metric/frame mixing, bounded by the last term.  The
material Euler row has one additional tangential derivative and is
divided by \(\lambda^2\), producing identical normalized ratios.
\end{proof}

\subsection{A simultaneous open background set}
\label{sec:v3v46-open-set}

For positive constants
\[
 \chi,r_-,r_+,c_-,c_+,\beta_-,\beta_+,\gamma_-,\gamma_+
\]
with \(0<c_-<c_+<1\), let \(\mathscr U\) be the set of backgrounds
satisfying:
\begin{enumerate}
 \item \(X_T\ge\chi^2\);
 \item the two singular values of the material coframe on \(\ker u^\flat\)
       lie in \([r_-,r_+]\);
 \item \(c\in[c_-,c_+]\),
       \(\beta\in[\beta_-,\beta_+]\), and
       \(\gamma\in[\gamma_-,\gamma_+]\);
 \item the exactly frozen normalized quotient symbol is within
       \(\delta_*\) in operator norm of a V45 flat symbol; and
 \item \(\mathfrak K_1,\mathfrak K_2\) are finite.
\end{enumerate}

\begin{theorem}[High-frequency margin on regular cones]
\label{thm:v3v46-regular-margin}
Let \(c_{45}>0\) be the minimum flat algebraic margin on the compact V45
parameter set.  Choose
\begin{equation}
 0<\delta_*<\frac{c_{45}}4.
\label{eq:v3v46-delta}
\end{equation}
On every regular conic subset satisfying
\eqref{eq:v3v46-separated} for the eliminated channels, there is
\(\Lambda_*\) depending only on the displayed background bounds such that
\begin{equation}
 \lambda\ge\Lambda_*
 \quad\Longrightarrow\quad
 s_{\min}(\widehat{\mathfrak B})
 \ge\frac{c_{45}}2.
\label{eq:v3v46-regular-margin}
\end{equation}
One may take \(\Lambda_*\) so that the right side of
\eqref{eq:v3v46-regular-bound} is at most \(c_{45}/4\).
\end{theorem}

\begin{proof}
The frozen-coefficient displacement costs at most \(\delta_*\), leaving
margin at least \(3c_{45}/4\) by
Lemma~\ref{lem:v3v34-margin}.  Choose \(\Lambda_*\) so the normalized
lower-symbol perturbation costs at most \(c_{45}/4\).  A second
application of the same lemma gives
\eqref{eq:v3v46-regular-margin}.
\end{proof}

\subsection{Failure of the regular expansion at the warm cone}
\label{sec:v3v46-warm-glancing}

The positive-speed channel has the closed-frequency glancing set
\begin{equation}
 s=i\tau,\qquad |\tau|=c|\xi|_q\ne0,
\qquad
 \rho_{c,1}=0.
\label{eq:v3v46-warm-cone}
\end{equation}
At this set, the cold root is
\begin{equation}
 \rho_{0,1}=s=i\tau\ne0,
\label{eq:v3v46-cold-at-warm}
\end{equation}
and V45 proves that the total quotient determinant is nonzero.

\begin{proposition}[The DtN lower symbol is singular at auxiliary
glancing]
\label{prop:v3v46-DtN-singular}
Suppose \(p_{c,1}\) has a nonzero limit at a point of
\eqref{eq:v3v46-warm-cone}.  Approach it through
\begin{equation}
 s_\epsilon=\epsilon+i\tau,\qquad\epsilon\downarrow0.
\label{eq:v3v46-approach}
\end{equation}
Then
\begin{equation}
 \rho_{c,1}(s_\epsilon,\xi)
 =
 \sqrt{2i\tau\epsilon+O(\epsilon^2)}
 =
 O(\epsilon^{1/2}),
\label{eq:v3v46-root-collapse}
\end{equation}
while
\begin{equation}
 r_{c,0}
 =
 \frac{p_{c,1}}{2\rho_{c,1}}
 =
 O(\epsilon^{-1/2}).
\label{eq:v3v46-subprincipal-blowup}
\end{equation}
Consequently the positive-speed lower symbol is not a uniformly bounded
order-zero perturbation across its glancing cone.
\end{proposition}

\begin{proof}
Use \(\tau^2=c^2|\xi|^2\):
\[
 s_\epsilon^2+c^2|\xi|^2
 =
 2i\tau\epsilon+\epsilon^2.
\]
Taking the outgoing square root gives
\eqref{eq:v3v46-root-collapse}.  Formula
\eqref{eq:v3v46-r0} then gives
\eqref{eq:v3v46-subprincipal-blowup}.
\end{proof}

\begin{counterexample}[Algebraic determinant margin does not bound the
eliminated lower symbol]
\label{cex:v3v46-algebraic-not-lower}
At \eqref{eq:v3v46-warm-cone}, the total principal response contains the
nonzero cold term \(\beta s\), so the V45 determinant has a nonzero
limit.  Nevertheless a fixed nonzero connection coefficient in
\(p_{c,1}\) makes the eliminated warm-channel correction diverge as
\(\epsilon^{-1/2}\) by
Proposition~\ref{prop:v3v46-DtN-singular}.  Thus no estimate of the form
\[
 \|E_\lambda\|\le C\mathfrak K_1/\lambda
\]
can hold uniformly in angular distance to the warm cone after that
channel is eliminated.
\end{counterexample}

\begin{firewall}[Correct scope of the V45 uniform statement]
\label{fw:v3v46-V45-scope}
Theorem~\ref{thm:v3v45-uniform-margin} proves a uniform flat algebraic
margin for the limiting principal quotient matrices in their graph norms.
It does not supply the differentiable stable bundle or the
variable-coefficient Kreiss symmetrizer at
\eqref{eq:v3v46-warm-cone}.  V46 proves the high-frequency perturbation
theorem only on separated regular cones.  Across auxiliary glancing, the
warm bath must remain as a local PDE state and be treated by an augmented
glancing estimate.
\end{firewall}

\subsection{V46 conclusion and V47 target}
\label{sec:v3v46-conclusion}

V46 writes the local covariant two-channel parent and identifies every
degree-one lower-symbol contribution through
\eqref{eq:v3v46-p1}.  On conic regions separated from both channel
characteristic sets, the normalized curvature, clock, material, and
connection perturbation is
\(O(\mathfrak K_1/\lambda+\mathfrak K_2/\lambda^2)\), so the V45 flat
margin persists on a concrete open set of backgrounds at high frequency.

The positive-speed bath cone is a genuine remaining gate.  The cold
channel keeps the total algebraic determinant nonzero there, but the
eliminated warm Dirichlet-to-Neumann subprincipal symbol blows up like the
inverse square root of angular distance.  No ordinary small-perturbation
argument crosses that set.

V47 should retain the warm bath normal state, reduce its normal equation
near \(\rho_c=0\) to a first-order doubled system, and construct a
Kreiss/glancing symmetrizer using the cold boundary dissipation as the
strict boundary term.  The proof must quantify the trace norm at the warm
cone and determine whether it is zero-loss or carries the V13
square-root loss.  If the loss survives, the nonlinear iteration must be
reweighted or the warm channel replaced by a static local carrier without
a propagating characteristic cone.
\clearpage
\section{V47: replace the warm bath by a static elliptic carrier}
\label{sec:v3v47}

\subsection{Why the upstream replacement is preferable}
\label{sec:v3v47-purpose}

V46 shows that the positive-speed bath solves the frozen algebraic
problem but creates an auxiliary glancing cone at which its eliminated
lower symbol is singular.  One could attempt a full glancing symmetrizer,
but the positive-speed channel was introduced only to supply homogeneous
static stiffness.  That function can instead be performed by a local
spatial elliptic carrier on each clock slice.  Such a carrier has no
propagating characteristic cone.

V47 constructs this replacement and tests the complete V42 quotient
determinant again.  The resulting response
\[
 \mu(s,\xi)=\beta s+\gamma|\xi|_q
\]
is zero-free on every nonzero closed frequency, supplies order-one static
coercivity, and has a local positive parent.  It selects the physical
clock foliation and is an instantaneous constraint sector; those features
are stated rather than hidden.

\subsection{Local static half-line carrier}
\label{sec:v3v47-static-carrier}

Let \(\Phi(y,x)\) take values in the same positive spatial STF bundle
\(\mathcal S_u\) as the material strain.  On every clock slice, take the
Lorentzian potential action
\begin{equation}
 S_{\rm stat}^L
 =
 -\frac12\int_{\mathcal N}\sqrt{-\gamma}\,d^3x
 \int_0^\infty dy\,
 \left(
 |\partial_y\Phi|_q^2
 +c_*^2q^{ab}
 \langle\mathcal D_a\Phi,\mathcal D_b\Phi\rangle_q
 \right),
 \qquad c_*>0,
\label{eq:v3v47-static-action}
\end{equation}
with no \(\mathcal D_u\Phi\) term, and impose
\begin{equation}
 \Phi(0,x)=\alpha_*e(x),\qquad\alpha_*>0.
\label{eq:v3v47-static-trace}
\end{equation}
Its stored energy is the positive functional
\begin{equation}
 E_{\rm stat}
 =
 \frac12\int_{\Sigma}\!\int_0^\infty
 \left(
 |\partial_y\Phi|_q^2
 +c_*^2q^{ab}
 \langle\mathcal D_a\Phi,\mathcal D_b\Phi\rangle_q
 \right)dy\,d^2x.
\label{eq:v3v47-static-energy}
\end{equation}

\begin{proposition}[Local covariance and constraint character]
\label{prop:v3v47-static-properties}
The action \eqref{eq:v3v47-static-action} is local and covariant under
boundary diffeomorphisms acting on \((\gamma,T,X,\Phi)\).  Its Euler
equation is elliptic in \((y,\Sigma_T)\) on each clock slice and contains
no second clock-time derivative.  Thus \(\Phi\) is a constrained physical
carrier with stored energy, not an additional propagating bath mode.
\end{proposition}

\begin{proof}
All contractions use the covariant clock-spatial projector and projected
bundle connection, so naturality proves covariance.  Varying \(\Phi\)
gives
\[
 -\partial_y^2\Phi+c_*^2\mathcal D_\Sigma^*
 \mathcal D_\Sigma\Phi=0.
\]
Its principal symbol in the spatial variables is
\(\eta_y^2+c_*^2|\xi|_q^2>0\) away from the joint spatial zero.  The
absence of \(\mathcal D_u\Phi\) proves the constraint statement, and
\eqref{eq:v3v47-static-energy} is manifestly nonnegative.
\end{proof}

\begin{theorem}[Static Dirichlet-to-Neumann response]
\label{thm:v3v47-static-DtN}
At a frozen flat point and for \(\xi\ne0\), the decaying solution with
trace \eqref{eq:v3v47-static-trace} is
\begin{equation}
 \widehat\Phi(y)
 =
 \alpha_*e^{-c_*|\xi|y}\widehat e.
\label{eq:v3v47-static-Poisson}
\end{equation}
Eliminating \(\Phi\) adds
\begin{equation}
 \frac{\alpha_*^2c_*}{2}
 \int|\xi|\,|\widehat e|^2\,d\xi
\label{eq:v3v47-static-effective}
\end{equation}
and the positive order-one response
\begin{equation}
 \mu_{\rm stat}(\xi)=\gamma|\xi|,
 \qquad
 \gamma>0
\label{eq:v3v47-static-mu}
\end{equation}
after fixed gravitational normalization.
\end{theorem}

\begin{proof}
Fourier transformation in the two spatial directions gives
\[
 -\partial_y^2\widehat\Phi+c_*^2|\xi|^2\widehat\Phi=0.
\]
The unique decaying solution is
\eqref{eq:v3v47-static-Poisson}.  Direct integration of
\eqref{eq:v3v47-static-energy} gives
\(\alpha_*^2c_*|\xi||\widehat e|^2/2\).
Its variation supplies the response
\eqref{eq:v3v47-static-mu}.
\end{proof}

\begin{firewall}[Clock-local constraint, not Lorentz-covariant matter]
\label{fw:v3v47-clock-constraint}
The static carrier propagates no independent signal, but after it is
eliminated its response is instantaneous and spatially nonlocal on each
clock slice.  The enlarged parent is local and diffeomorphism covariant
because the physical clock is retained; it is not boundary
Lorentz-invariant matter.  Its acceptability therefore depends on the
physical clock/material sector already admitted in V38.
\end{firewall}

\subsection{Euclidean reflection positivity}
\label{sec:v3v47-RP}

The Euclidean action is the positive version of
\eqref{eq:v3v47-static-action}.  At a finite spatial and half-line cutoff,
let \(A_{\rm stat}>0\) be its spatial quadratic matrix.  There is no
coupling between distinct Euclidean time slices.

\begin{proposition}[Reflection positivity of the static carrier]
\label{prop:v3v47-RP}
At finite cutoff, the static-carrier Gaussian is reflection positive.
Its covariance has the form
\begin{equation}
 C_{\rm stat}(\tau,\tau')
 =
 \delta(\tau-\tau')A_{\rm stat}^{-1}.
\label{eq:v3v47-static-covariance}
\end{equation}
For test functions supported at strictly positive Euclidean time,
\begin{equation}
 \int_{\tau,\tau'>0}
 \langle f(\tau),
 C_{\rm stat}(-\tau,\tau')f(\tau')\rangle
 d\tau d\tau'=0\ge0.
\label{eq:v3v47-static-RP}
\end{equation}
The time-zero contribution is nonnegative.
\end{proposition}

\begin{proof}
On a time lattice, the positive Gaussian factorizes into independent
copies at every time slice, giving the diagonal covariance
\eqref{eq:v3v47-static-covariance}.  Reflection pairs distinct positive
and negative slices, whose covariance is zero.  A slice fixed by
reflection has covariance \(A_{\rm stat}^{-1}\ge0\).  The continuum
formula is the distributional limit and proves
\eqref{eq:v3v47-static-RP}.
\end{proof}

\begin{assessment}[Degenerate time reconstruction]
\label{ass:v3v47-degenerate-time}
Reflection positivity holds, but the static carrier has no independent
time transfer semigroup.  It reconstructs as a constraint variable rather
than a propagating particle.  The cold line remains the channel carrying
exported energy and retarded time response.
\end{assessment}

\subsection{Combined cold-plus-static response}
\label{sec:v3v47-combined}

Combine the static carrier with the V43 cold line.  The completed material
force enters the V42 determinant through
\begin{equation}
 \mu(s,k)=\beta s+\gamma k,
 \qquad
 \beta,\gamma>0,\quad k=|\xi|.
\label{eq:v3v47-mu}
\end{equation}
This symbol is homogeneous of degree one and has no square-root
characteristic cone.

For a closed frequency \(s=ixk\), \(x>0\), define \(Q,P\) as in
\eqref{eq:v3v45-QP}.  The normalized determinant bracket remains
\begin{equation}
 \mathcal F_{\rm cs}(x)
 =
 \widehat\kappa\,\widehat\mu^{\,2}Q
 -x^2\widehat\mu P
 -4x^6\widehat\kappa,
\qquad
 \widehat\mu=\gamma+i\beta x.
\label{eq:v3v47-F}
\end{equation}

\begin{theorem}[Cold-plus-static closed-frequency theorem]
\label{thm:v3v47-no-closed}
For every \(\beta,\gamma>0\),
\begin{equation}
 \mathcal F_{\rm cs}(x)\ne0
 \qquad(x>0).
\label{eq:v3v47-no-closed}
\end{equation}
Hence the completed quotient determinant has no nonzero closed-frequency
root for \(k>0\).
\end{theorem}

\begin{proof}
First suppose \(0<x<1\).  Put
\[
 r=\sqrt{1-x^2},\qquad h=\gamma>0,\qquad m=\beta x>0.
\]
Then \(\widehat\kappa=r\) and
\(\widehat\mu=h+im\).  If
\(\mathcal F_{\rm cs}=0\), its imaginary part gives
\[
 m(2rhQ-x^2P)=0.
\]
When \(Q=0\), one has \(x^2=1/2\) and \(P=-1\), so this is impossible.
When \(Q>0\), it gives
\(x^2P=2rhQ\), which also forces \(P>0\).  The real part then becomes
\[
 r(h^2-m^2)Q-x^2hP-4x^6r
 =
 -r(h^2Q+m^2Q+4x^6)<0,
\]
again impossible.

At \(x=1\), \(\widehat\kappa=0\), \(P=1\), and
\[
 \mathcal F_{\rm cs}(1)=-(\gamma+i\beta)\ne0.
\]

Finally let \(x>1\), so
\(\widehat\kappa=i\sqrt{x^2-1}\).  The real part of
\eqref{eq:v3v47-F} is
\begin{equation}
 -2\sqrt{x^2-1}\,\gamma\beta x\,Q
 -x^2\gamma P.
\label{eq:v3v47-super-real}
\end{equation}
For \(x>1\), \(P=8x^4-8x^2+1>0\).  Both terms in
\eqref{eq:v3v47-super-real} are nonpositive and the second is strictly
negative.  Therefore the determinant cannot vanish.
\end{proof}

\begin{proposition}[Frequency axes]
\label{prop:v3v47-axes}
At \(s=0,\ k>0\),
\begin{equation}
 \mu=\gamma k>0
\label{eq:v3v47-static-axis}
\end{equation}
and the static material-row coercivity is
\(2\gamma k^3|\chi|^2\).  At \(k=0,\ s\ne0\),
\begin{equation}
 \mu=\beta s
\label{eq:v3v47-time-axis}
\end{equation}
and the axis determinant is
\begin{equation}
 -2(1+\beta)^2s^7\ne0.
\label{eq:v3v47-time-axis-det}
\end{equation}
\end{proposition}

\begin{proof}
The first statement follows from
Theorem~\ref{thm:v3v47-static-DtN} and
\eqref{eq:v3v39-D-norm}.  The second is the cold-line axis calculation
\eqref{eq:v3v43-cold-axis}.
\end{proof}

\begin{theorem}[Frozen quotient margin without auxiliary glancing]
\label{thm:v3v47-frozen-margin}
In the mixed metric/material graph weights of V45, the flat
cold-plus-static principal quotient is injective on every nonzero closed
Laplace frequency and has a positive normalized minimum singular value.
It has no auxiliary normal-root glancing set away from \(k=0\).
\end{theorem}

\begin{proof}
Theorem~\ref{thm:v3v47-no-closed} covers
\(s\ne0,\ k>0\), while
Proposition~\ref{prop:v3v47-axes} covers the two axes.  The response is
homogeneous of degree one.  The local augmented quotient matrices and
their transition maps are continuous in the graph weights; compactness of
the normalized hemisphere gives a positive minimum.  The static normal
root is \(c_*k\), which vanishes only at \(k=0\), where the cold axis
matrix is invertible.
\end{proof}

\begin{theorem}[No growing mode]
\label{thm:v3v47-no-growing}
The local cold-plus-static system has no finite-energy quotient mode with
\(\operatorname{Re}s>0\).
\end{theorem}

\begin{proof}
The cold line carries nonnegative wave energy, the static carrier stores
the nonnegative potential energy \eqref{eq:v3v47-static-energy}, and the
metric/material junction signs are variations of their sum.  The total
energy identity rules out exponential growth exactly as in
Theorem~\ref{thm:v3v45-no-growing}.
\end{proof}

\subsection{Variable-coefficient advantage}
\label{sec:v3v47-variable}

The static carrier has tangential principal symbol
\(c_*^2|\xi|_q^2\).  On a cone \(k\ge\varepsilon\lambda\), its
Dirichlet-to-Neumann root \(c_*|\xi|_q\) is elliptic and its lower symbols
obey the regular square-root expansion of
Lemma~\ref{lem:v3v46-DtN-lower}.  On the complementary cone
\(k<\varepsilon\lambda\), the cold response \(\beta s\) is uniformly
order one.

\begin{proposition}[Two-cone lower-symbol cover]
\label{prop:v3v47-two-cone}
Fix \(0<\varepsilon<1/2\).  The normalized frequency hemisphere is
covered by
\begin{align}
 \Omega_{\rm sp}&=\{k\ge\varepsilon\lambda\},
\label{eq:v3v47-Omega-sp}\\
 \Omega_{\rm tm}&=\{|s|\ge\sqrt{1-\varepsilon^2}\,\lambda\}.
\label{eq:v3v47-Omega-tm}
\end{align}
On \(\Omega_{\rm sp}\), the static root is separated by
\(c_*\varepsilon\lambda\).  On \(\Omega_{\rm tm}\), the cold root is
separated by \(\sqrt{1-\varepsilon^2}\lambda\).  Therefore the
variable-coefficient lower-symbol perturbation obeys a regular bound of
the form \eqref{eq:v3v46-regular-bound} on a finite conic cover, with no
warm-bath glancing exception.
\end{proposition}

\begin{proof}
If \(k<\varepsilon\lambda\), then
\(|s|^2=\lambda^2-k^2>(1-\varepsilon^2)\lambda^2\), proving the cover.
The two lower bounds follow directly.  Apply
Lemma~\ref{lem:v3v46-DtN-lower} only to the carrier whose root is separated
on the relevant cone; keep the other local channel uneliminated if
necessary.
\end{proof}

\begin{firewall}[Remaining nonlinear and physical gates]
\label{fw:v3v47-remaining}
V47 removes the auxiliary glancing obstruction at the frozen and first
lower-symbol levels.  It does not yet prove the nonlinear Einstein
initial-boundary value problem, the global persistence of the clock and
material-rank gaps, the curved BRST domain, the renormalized quantum
measure, or equivalence of the added constraint carrier to ordinary
Standard Model matter.  The static carrier is additional physical
boundary structure.
\end{firewall}

\subsection{V47 conclusion and V48 target}
\label{sec:v3v47-conclusion}

V47 replaces the propagating positive-speed bath by a local spatial
elliptic constraint carrier.  Together with the cold line it has response
\(\beta s+\gamma|\xi|\), positive stored energy, a reflection-positive
Euclidean parent, homogeneous static stiffness, and no auxiliary
characteristic cone.  The exact completed V42 determinant is zero-free on
the entire nonzero closed hemisphere, and the two frequency axes are
controlled separately.

V48 should carry this cold-plus-static parent through the full
variable-coefficient boundary operator.  It should construct the conic
parametrices on \(\Omega_{\rm sp}\) and \(\Omega_{\rm tm}\), patch their
symmetrizers, and prove a high-frequency estimate with constants depending
only on the explicit clock, material-rank, coupling, and curvature bounds
of V46.  The resulting estimate must then be inserted into the harmonic
Einstein constraint-propagation system rather than treated as an isolated
TT/material block.
\clearpage
\section{V48: a time-domain glancing estimate and harmonic constraint closure}
\label{sec:v3v48}

\subsection{Why the time-domain estimate is the correct next step}
\label{sec:v3v48-purpose}

V47 removes the auxiliary warm-bath cone, but the Einstein bulk normal
root
\(\kappa=\sqrt{s^2+|\xi|^2}\) still vanishes at primary glancing.
Differentiating the stable root across that set recreates the V13
square-root singularity.  The cold-plus-static boundary law has a stronger
structure which does not require such differentiation: the cold line
exports positive energy and the static carrier stores the positive
boundary energy \(\langle q,|D_\Sigma|q\rangle/2\).

V48 proves the resulting time-domain estimate first for the exact scalar
regression, then for the physical quotient bundle with variable
coefficients.  It also inserts the estimate into the harmonic constraint
system.  The proof gives zero-loss control in the natural energy and
boundary half-derivative norm without assuming smoothness of the stable
root at glancing.

\subsection{Exact scalar regression}
\label{sec:v3v48-scalar}

Let \(x>0\) be the gravitational normal coordinate and
\(z\in\mathbb R^2\) the boundary spatial variable.  Put
\begin{equation}
 \Lambda=(-\Delta_z)^{1/2}.
\label{eq:v3v48-Lambda}
\end{equation}
Consider
\begin{align}
 q_{tt}-q_{xx}-\Delta_zq&=f,
 &&x>0,
\label{eq:v3v48-wave}\\
 -q_x+\beta q_t+\gamma\Lambda q&=g,
 &&x=0,
\qquad \beta,\gamma>0.
\label{eq:v3v48-boundary}
\end{align}

Define
\begin{align}
 E_{\rm bulk}(t)
 &=
 \frac12\int_{x>0}
 \left(
 |q_t|^2+|q_x|^2+|\nabla_zq|^2
 \right)dx\,d^2z,
\label{eq:v3v48-Ebulk}\\
 E_{\rm stat}(t)
 &=
 \frac\gamma2
 \langle q(t,0),\Lambda q(t,0)\rangle_{L^2_z}.
\label{eq:v3v48-Estat}
\end{align}

\begin{theorem}[Exact cold-plus-static energy identity]
\label{thm:v3v48-scalar-energy}
Every smooth decaying solution of
\eqref{eq:v3v48-wave}--\eqref{eq:v3v48-boundary} satisfies
\begin{equation}
 \frac d{dt}(E_{\rm bulk}+E_{\rm stat})
 +\beta\|q_t(t,0)\|_{L^2_z}^2
 =
 \langle f,q_t\rangle_{x>0}
 +\langle g,q_t(0)\rangle_z.
\label{eq:v3v48-energy-identity}
\end{equation}
In the homogeneous problem, the total reduced energy is nonincreasing.
When the cold line is retained rather than eliminated, the lost term is
exactly the positive energy flux down that line and the enlarged energy is
conserved.
\end{theorem}

\begin{proof}
Integration by parts in the half-space gives
\begin{equation}
 E_{\rm bulk}'
 =
 -\langle q_t,q_x\rangle_{x=0}
 +\langle f,q_t\rangle.
\label{eq:v3v48-Ebulk-prime}
\end{equation}
Equation \eqref{eq:v3v48-boundary} gives
\[
 q_x=\beta q_t+\gamma\Lambda q-g.
\]
Substitute this in \eqref{eq:v3v48-Ebulk-prime}.  Self-adjointness of
\(\Lambda\) gives
\[
 \gamma\langle q_t,\Lambda q\rangle
 =
 \frac d{dt}E_{\rm stat}.
\]
Rearrangement proves \eqref{eq:v3v48-energy-identity}.  The last statement
is the cold-line flux identity of V16 and V39.
\end{proof}

\begin{corollary}[Finite-time energy estimate]
\label{cor:v3v48-scalar-estimate}
For \(0\le t\le T\),
\begin{align}
 &E_{\rm bulk}(t)+E_{\rm stat}(t)
 +\frac\beta2\int_0^t\|q_t(\sigma,0)\|_2^2\,d\sigma
\notag\\
 &\qquad\le
 E_{\rm bulk}(0)+E_{\rm stat}(0)
 +\int_0^t\|f(\sigma)\|_2
          \|q_t(\sigma)\|_2\,d\sigma
 +\frac1{2\beta}\int_0^t\|g(\sigma)\|_2^2\,d\sigma.
\label{eq:v3v48-scalar-estimate}
\end{align}
\end{corollary}

\begin{proof}
Integrate \eqref{eq:v3v48-energy-identity} and use
\[
 |\langle g,q_t(0)\rangle|
 \le\frac1{2\beta}\|g\|^2+\frac\beta2\|q_t(0)\|^2.
\]
\end{proof}

\begin{proposition}[The estimate crosses primary glancing]
\label{prop:v3v48-cross-glancing}
The constants in
\eqref{eq:v3v48-energy-identity}--%
\eqref{eq:v3v48-scalar-estimate} are independent of
\(\kappa=\sqrt{s^2+|\xi|^2}\).  In particular, they do not deteriorate as
\(\kappa\to0\).  The controlled boundary quantities are
\begin{equation}
 \Lambda^{1/2}q|_{x=0}\in L^\infty_tL^2_z,
 \qquad
 q_t|_{x=0}\in L^2_{t,z}.
\label{eq:v3v48-boundary-norms}
\end{equation}
\end{proposition}

\begin{proof}
The proof uses integration by parts and positivity of
\(\beta,\gamma,\Lambda\); no stable-root factor is divided.  The two
boundary norms occur explicitly in
\(E_{\rm stat}\) and the dissipative integral.
\end{proof}

\begin{assessment}[Meaning of zero loss]
\label{ass:v3v48-zero-loss}
The bulk wave energy controls one derivative of \(q\), and the boundary
carrier controls the natural half derivative
\(\Lambda^{1/2}q(0)\).  This is the trace order dictated by the positive
order-one boundary form.  V48 calls the estimate zero-loss because no
additional negative power of the glancing root is required beyond these
energy spaces.
\end{assessment}

\subsection{Physical quotient bundle}
\label{sec:v3v48-bundle}

Let \(\mathcal H_{\rm phys}\) denote the seven-dimensional frozen
metric/material quotient constructed in V41--V45, equipped with a smooth
positive symmetrizer \(M(t,z)\).  Let
\(\mathsf B(t,z)\) be the cold incidence and
\(\mathsf A(t)\) the positive self-adjoint order-one operator obtained by
eliminating the static carrier.  Assume
\begin{align}
 m_-I\le M\le m_+I,
\label{eq:v3v48-M-bounds}\\
 \mathsf B^*M\mathsf B\ge\beta_0I,
 \qquad \beta_0>0,
\label{eq:v3v48-B-positive}\\
 \langle U,\mathsf A U\rangle
 \ge\gamma_0\|\Lambda^{1/2}U\|^2-C_0\|U\|^2,
 \qquad \gamma_0>0.
\label{eq:v3v48-A-positive}
\end{align}
The possible finite-dimensional spatial zero mode is placed in the
\(C_0\|U\|^2\) term and treated by the infrared sector.

For a normally hyperbolic physical quotient operator
\(\mathcal L_{\rm g}\), write the boundary law schematically as
\begin{equation}
 -\nabla_nU+\mathsf B^*M\mathsf B\,\nabla_uU
 +\mathsf A U
 =
 G+\mathsf L_0U,
\label{eq:v3v48-bundle-boundary}
\end{equation}
where \(\mathsf L_0\) contains order-zero curvature and frame terms.

\begin{theorem}[Variable-coefficient quotient energy estimate]
\label{thm:v3v48-variable-energy}
Suppose the clock gap, material-rank gap, coefficient bounds
\eqref{eq:v3v48-M-bounds}--\eqref{eq:v3v48-A-positive}, and the geometric
bounds \eqref{eq:v3v46-K1}--\eqref{eq:v3v46-K2} hold on
\(0\le t\le T\).  Assume also
\begin{equation}
 \left|
 \langle U,\dot{\mathsf A}U\rangle
 \right|
 \le
 K_A\left(
 \langle U,\mathsf A U\rangle+\|U\|^2
 \right).
\label{eq:v3v48-Adot}
\end{equation}
Then smooth solutions of the linear physical quotient system obey
\begin{align}
 \mathcal E(t)
 &+
 \frac{\beta_0}{2}
 \int_0^t\|\nabla_uU(\sigma)|_{\mathcal N}\|^2\,d\sigma
\notag\\
 &\le
 C_T\left[
 \mathcal E(0)
 +\int_0^t\|F(\sigma)\|\,\|\nabla_uU(\sigma)\|\,d\sigma
 +\int_0^t\|G(\sigma)\|^2\,d\sigma
 \right],
\label{eq:v3v48-variable-estimate}
\end{align}
where
\begin{equation}
 \mathcal E
 \simeq
 \|\nabla_uU\|^2+\|\nabla_nU\|^2
 +\|\nabla_\Sigma U\|^2
 +\|\mathsf A^{1/2}U|_{\mathcal N}\|^2
 +\|U\|^2
\label{eq:v3v48-variable-energy}
\end{equation}
and \(C_T\) depends only on \(T\) and the displayed uniform coefficient
bounds.  It has no inverse power of the primary glancing root.
\end{theorem}

\begin{proof}
Pair the normally hyperbolic equation with
\(M\nabla_uU\).  The principal bulk terms give the first three positive
terms in \eqref{eq:v3v48-variable-energy}; coefficient derivatives,
clock acceleration, and commutators contribute at most
\(C\mathfrak K_1\mathcal E\).  The boundary flux is the pairing of
\(\nabla_uU\) with \(\nabla_nU\).  Substitute
\eqref{eq:v3v48-bundle-boundary}.

The cold term gives at least
\(\beta_0\|\nabla_uU\|^2\).
Self-adjointness of \(\mathsf A\) gives
\[
 \langle\nabla_uU,\mathsf AU\rangle
 =
 \frac12\frac d{dt}\langle U,\mathsf AU\rangle
 -\frac12\langle U,\dot{\mathsf A}U\rangle
\]
up to the bounded change between \(u\) and the chosen time coordinate.
Use \eqref{eq:v3v48-Adot} and
\eqref{eq:v3v48-A-positive}.  The order-zero boundary term is controlled
by the boundary energy and a small part of the cold dissipation.
Cauchy--Schwarz and Young inequalities control \(G\); the bulk source
pairs with \(\nabla_uU\).  Gronwall's inequality gives
\eqref{eq:v3v48-variable-estimate}.
\end{proof}

\begin{corollary}[Uniqueness and continuous dependence]
\label{cor:v3v48-uniqueness}
On every finite interval on which the background gaps remain positive,
the linear physical quotient problem has at most one solution in the
energy class, and its solution map is continuous in the norms of
\eqref{eq:v3v48-variable-estimate}.
\end{corollary}

\begin{proof}
Apply Theorem~\ref{thm:v3v48-variable-energy} to the difference of two
solutions.  With zero data and sources, Gronwall gives
\(\mathcal E=0\).
\end{proof}

\subsection{Local parent versus the nonlocal boundary form}
\label{sec:v3v48-parent}

The operator \(\mathsf A\) is nonlocal after the static carrier is
eliminated.  Retaining \(\Phi\) makes the system local and replaces
\(\langle U,\mathsf AU\rangle/2\) by the stored energy
\eqref{eq:v3v47-static-energy}.

\begin{proposition}[Equivalent enlarged energy identity]
\label{prop:v3v48-parent-energy}
For smooth data satisfying the trace relation, the time derivative of the
static-carrier energy equals the boundary pairing
\begin{equation}
 \frac d{dt}E_{\rm stat}
 =
 \langle\nabla_uU,\mathsf AU\rangle
\label{eq:v3v48-parent-energy}
\end{equation}
up to the same variable-coefficient commutators bounded in
Theorem~\ref{thm:v3v48-variable-energy}.  Hence the reduced and enlarged
energy estimates are equivalent on the constraint surface.
\end{proposition}

\begin{proof}
Differentiate \eqref{eq:v3v47-static-energy}, integrate the spatial and
\(y\)-derivatives by parts, and use the static carrier Euler equation.
Only the \(y=0\) flux remains.  Differentiate its trace
\(\Phi(0)=\alpha_*e(U)\) in time.  The resulting normal flux is the
Dirichlet-to-Neumann force \(\mathsf AU\), giving
\eqref{eq:v3v48-parent-energy}.  Derivatives of the frame, projector, and
volume density give the bounded commutator terms.
\end{proof}

\subsection{Harmonic constraint propagation}
\label{sec:v3v48-constraints}

Let
\begin{equation}
 C_\mu=\nabla^\nu\bar h_{\mu\nu}
\label{eq:v3v48-C}
\end{equation}
be the linearized harmonic constraint on a background satisfying the
Einstein equations.  The reduced Einstein equation and the linearized
Bianchi identity give a normally hyperbolic constraint equation
\begin{equation}
 \Box_g C_\mu+\mathcal R_\mu{}^\nu C_\nu=0
\label{eq:v3v48-C-wave}
\end{equation}
when the complete bulk source is conserved.

The medium Ward identity
\eqref{eq:v3v38-Ward} and its cold-plus-static analogue imply, on the
material and carrier equations,
\begin{equation}
 D^a\mathcal S_{ab}=T_{nb}^{\rm bulk}
\label{eq:v3v48-boundary-Ward}
\end{equation}
with the appropriate normal exchange term.  This is exactly the Codazzi
compatibility required by
\eqref{eq:v3v15-linear-flux}.

\begin{theorem}[Constraint closure of the energy-controlled boundary
system]
\label{thm:v3v48-constraint-closure}
Assume:
\begin{enumerate}
 \item the bulk Standard Model source is covariantly conserved;
 \item the material, cold-line, and static-carrier Euler equations hold;
 \item the junction source satisfies
       \eqref{eq:v3v48-boundary-Ward};
 \item the four boundary de Donder rows impose a homogeneous
       constraint-preserving condition on the incoming constraint
       characteristics; and
 \item \(C_\mu\) and its normal time derivative vanish initially.
\end{enumerate}
Then
\begin{equation}
 C_\mu=0
\label{eq:v3v48-C-zero}
\end{equation}
throughout the interval of existence.  The physical quotient estimate
\eqref{eq:v3v48-variable-estimate} therefore evolves a solution of the
unreduced linearized Einstein equations, not merely the harmonic reduced
system.
\end{theorem}

\begin{proof}
Bulk conservation and the Bianchi identity give
\eqref{eq:v3v48-C-wave}.  Boundary Ward/Codazzi compatibility removes the
inhomogeneous constraint source in the normal Einstein equation.  The
four de Donder rows give the homogeneous incoming boundary condition
described in Proposition~\ref{prop:v3v15-incoming-constraints}.  Apply the
standard energy estimate for the vector wave equation
\eqref{eq:v3v48-C-wave}.  Zero initial and incoming data imply
\(C=0\) by uniqueness.  With \(C=0\), the harmonic reduction terms vanish
in the metric equation, proving the final statement.
\end{proof}

\begin{firewall}[What is assumed in the constraint theorem]
\label{fw:v3v48-constraint-assumptions}
Theorem~\ref{thm:v3v48-constraint-closure} uses an exact curved Ward
identity and a closed constraint-preserving boundary domain.  V38 proves
the classical covariant Ward identity for the local medium action, while
V32 warns that a complete curved BRST operator domain needs the actual
ghost Calderon map.  V48 does not yet prove anomaly-free quantum Ward
identities or the nonlinear compatibility conditions at corners.
\end{firewall}

\subsection{V48 conclusion and V49 target}
\label{sec:v3v48-conclusion}

V48 bypasses the primary square-root stable-root singularity.  The cold
boundary channel gives a positive \(L^2\) time-derivative trace and the
static carrier gives the positive half-derivative trace energy.  Their
time-domain identity is uniform through primary glancing.  The
variable-coefficient quotient estimate persists on any finite interval
where the clock, material-rank, coupling, and geometric coefficient bounds
remain controlled.

The medium Ward identity also supplies the Codazzi source compatibility.
Together with the four de Donder rows, it propagates zero harmonic
constraints and converts the quotient energy estimate into a linear
unreduced Einstein estimate.

V49 should lift this linear estimate to the quasilinear
Einstein--Standard Model--material system by deriving the differentiated
energy hierarchy.  It must track derivative loss in the order-one static
trace operator, prove tame commutator estimates for the moving clock
projector and material coframe, and identify the regularity at which the
boundary and bulk nonlinearities close.  A local existence claim may be
made only after compatibility conditions and persistence of the clock and
rank gaps are proved.
\clearpage
\section{V49: differentiated quasilinear energies and persistence of the open geometric set}
\label{sec:v3v49}

\subsection{Purpose and regularity convention}
\label{sec:v3v49-purpose}

V48 proves a zero-loss linear energy estimate through primary glancing.
The next obstruction is nonlinear: tangential differentiation changes the
clock projector, the material coframe, the cold incidence, and the static
elliptic operator.  An estimate for the eliminated square-root operator
would require a separate parameter-dependent pseudodifferential calculus.
There is no need to introduce that calculus.  The half-line field
\(\Phi\) in \eqref{eq:v3v47-static-action} is local, and ordinary elliptic
regularity supplies exactly the boundary half derivative occurring in the
V48 energy.

This note proves a tame differentiated \emph{a priori} estimate for a
smooth solution of the reduced classical system.  It also proves that the
clock, material-rank, and positivity gaps persist for a quantitative time.
The result is the energy part of a local existence theorem.  Construction
of solutions and convergence of an iteration remain separate gates.

Let \(S_t\) be a three-dimensional spacelike slice meeting the timelike
boundary in the compact two-dimensional surface \(\Sigma_t\).  Local
statements on noncompact slices are obtained with a uniformly locally
finite partition of unity.  Fix an integer
\begin{equation}
 m\ge4.
\label{eq:v3v49-m}
\end{equation}
Let \(\mathcal Z\) be a smooth family of derivatives tangent to the clock
foliation and tangent to the boundary when restricted to it.  Multi-index
notation \(\mathcal Z^I\) includes the transported clock derivative
\(\nabla_u\).  The local static carrier is never eliminated in this note.

Write \(W\) for the collection of hyperbolic unknowns: the harmonic
metric quotient, the clock and material fields, the cold carrier, and the
gauge-fixed classical Standard Model fields.  Dirac components are placed
in first-order symmetric-hyperbolic form; wave, Higgs, and gauge components
are placed in their standard first-order energy variables.  Let
\(\mathbf c(W,\nabla W)\) denote the finite coefficient array consisting
of the inverse metric, clock projector, material coframe, incidence maps,
connections, and zeroth-order couplings which occur in the reduced
equations.

\subsection{The two tame estimates}
\label{sec:v3v49-tame}

The following elementary estimates are the only high-order product input.

\begin{lemma}[Interior and boundary tame products]
\label{lem:v3v49-tame-products}
Let \(d\le3\), \(m\ge4\), and let \(D\) be a smooth covariant derivative
on a vector bundle over a compact \(d\)-manifold.  For \(|I|\le m\),
\begin{align}
 \|D^I(av)-aD^Iv\|_{L^2}
 &\le C\left(
   \|Da\|_{L^\infty}\|v\|_{H^{m-1}}
   +\|a\|_{H^m}\|v\|_{L^\infty}
 \right),
\label{eq:v3v49-commutator}\
 \|av\|_{H^m}
 &\le C\left(
   \|a\|_{L^\infty}\|v\|_{H^m}
   +\|a\|_{H^m}\|v\|_{L^\infty}
 \right).
\label{eq:v3v49-product}
\end{align}
The same estimates hold on \(\Sigma_t\), and the trace map satisfies
\begin{equation}
 H^r(S_t)\longrightarrow H^{r-1/2}(\Sigma_t),
 \qquad r>1/2.
\label{eq:v3v49-trace}
\end{equation}
All constants are uniform when the coordinate and bundle metrics range in
a bounded uniformly elliptic family.
\end{lemma}

\begin{proof}
Expand the commutator by the covariant Leibniz rule.  In a term
\(D^Ja\,D^{I-J}v\), put the factor with at most one derivative below the
Sobolev threshold in \(L^\infty\), and the other in \(L^2\).  The remaining
intermediate cases follow from Holder and Sobolev interpolation.  Since
\(m\ge4>d/2+2\), \(H^{m-1}\) embeds into \(W^{1,\infty}\).  Summing the
terms proves \eqref{eq:v3v49-commutator}; the same expansion proves
\eqref{eq:v3v49-product}.  Boundary dimension is smaller, so the same
argument applies there.  Equation \eqref{eq:v3v49-trace} is the standard
local half-space trace theorem, patched by the finite atlas.
\end{proof}

Let \(L_t\) be the static carrier operator on
\(Y_t=[0,\infty)_y\times\Sigma_t\):
\begin{equation}
 L_t\Phi
 =-\partial_y(a_y\partial_y\Phi)
   +\mathcal D_A^*(a^{AB}\mathcal D_B\Phi)+a_0\Phi,
 \qquad
 \Phi|_{y=0}=\alpha_*e(W).
\label{eq:v3v49-static-operator}
\end{equation}
Assume \(a_y\ge a_->0\), \(a^{AB}\zeta_A\zeta_B\ge
a_-|\zeta|^2\), and impose decay, or a fixed coercive infrared condition,
as \(y\to\infty\).

\begin{lemma}[Tame local-carrier estimate]
\label{lem:v3v49-carrier-tame}
Suppose the coefficients of \(L_t\) are bounded in \(H^m\), are uniformly
elliptic, and their first derivatives are bounded in \(L^\infty\).  For
the solution of \eqref{eq:v3v49-static-operator},
\begin{align}
 \|\nabla_{y,\Sigma}\Phi\|_{H^{m-1}(Y_t)}
 &\le C_{\rm ell}\left(
   \|e(W)\|_{H^{m-1/2}(\Sigma_t)}
   +\|\Phi\|_{L^2(Y_t)}
 \right),
\label{eq:v3v49-ell-upper}\
 \|e(W)\|_{H^{m-1/2}(\Sigma_t)}^2
 &\le C_{\rm tr}\left(
   \|\nabla_{y,\Sigma}\Phi\|_{H^{m-1}(Y_t)}^2
   +\|\Phi\|_{L^2(Y_t)}^2
 \right).
\label{eq:v3v49-ell-trace}
\end{align}
After applying \(\mathcal Z^I\), \(|I|\le m-1\), every commutator is
bounded by
\begin{equation}
 C\left(
  \|D\mathbf c\|_{L^\infty}
       \|\nabla\Phi\|_{H^{m-2}}
  +\|\mathbf c\|_{H^m}
       \|\nabla\Phi\|_{L^\infty}
 \right).
\label{eq:v3v49-static-comm}
\end{equation}
Thus the carrier costs no derivative beyond the boundary
\(H^{m-1/2}\) norm already dictated by the V48 energy.
\end{lemma}

\begin{proof}
Choose an extension \(E e\in H^m(Y_t)\) of the Dirichlet datum, with
\(\|Ee\|_{H^m}\le C\|e\|_{H^{m-1/2}}\), and set
\(\Phi_0=\Phi-\alpha_*Ee\).  The field \(\Phi_0\) has homogeneous
Dirichlet data.  Difference quotients in \(y\) and in boundary charts,
followed by coercivity of the bilinear form of \(L_t\), give the
\(H^m\) elliptic estimate \eqref{eq:v3v49-ell-upper}.  Lower-order or
spatial zero modes are retained in the displayed \(L^2\) term.  The trace
theorem applied to \(\Phi\) gives \eqref{eq:v3v49-ell-trace}.  Commuting
the local divergence-form equation produces only ordinary products of a
coefficient derivative with a derivative of \(\Phi\); apply
Lemma~\ref{lem:v3v49-tame-products}.  No spectral square root has been
differentiated, which proves the last statement.
\end{proof}

\subsection{Differentiated energy}
\label{sec:v3v49-energy}

For \(|I|\le m-1\), apply the V48 energy to
\(W_I=\mathcal Z^IW\), and include the corresponding differentiated
static energy.  Denote the sum by
\begin{align}
 \mathcal E_m(t)
 \simeq{}&
 \|W(t)\|_{H^m(S_t)}^2
 +\|\nabla_uW(t)\|_{H^{m-1}(S_t)}^2
\notag\\
 &+\|e(W)(t)\|_{H^{m-1/2}(\Sigma_t)}^2
 +\|\nabla_{y,\Sigma}\Phi(t)\|_{H^{m-1}(Y_t)}^2
 +\mathcal E_{\rm IR}(t),
\label{eq:v3v49-Em}
\end{align}
where \(\mathcal E_{\rm IR}\) controls the finite-dimensional spatial
zero sector.  Let
\begin{equation}
 \mathcal D_m(t)
 =\|\nabla_u e(W)|_{\Sigma_t}\|_{H^{m-1}}^2
  +\mathcal D_{{\rm SM},m}(t)
\label{eq:v3v49-Dm}
\end{equation}
be the cold-line boundary dissipation plus any nonnegative maximally
dissipative Standard Model boundary flux.

We state explicitly the coefficient hypothesis needed by the estimate:
\begin{equation}
 \mathcal K_m(t)
 =1+\|D\mathbf c(t)\|_{L^\infty}
   +\|\partial_t\mathbf c(t)\|_{L^\infty}
   +\|\mathbf c(t)\|_{H^m}
   +\|\Phi(t)\|_{W^{1,\infty}(Y_t)}.
\label{eq:v3v49-Km}
\end{equation}

\begin{theorem}[Tame differentiated a priori estimate]
\label{thm:v3v49-high-energy}
Let \(m\ge4\).  Let a smooth solution of the reduced classical
Einstein--Standard Model--clock/material system with the cold and static
carriers exist on \([0,T]\).  Suppose:
\begin{enumerate}
 \item the V48 clock, rank, ellipticity, incidence, and coupling gaps hold
       uniformly;
 \item the harmonic, Standard Model gauge, material, cold, and static
       boundary domains are constraint compatible and maximally
       dissipative in their hyperbolic sectors;
 \item the coefficients in \eqref{eq:v3v49-Km} are finite; and
 \item the data satisfy the corner compatibility conditions through order
       \(m-1\).
\end{enumerate}
Then
\begin{align}
 \frac d{dt}\mathcal E_m(t)+\frac{\beta_0}{4}\mathcal D_m(t)
 \le{}& C_{\rm low}(t)\mathcal E_m(t)
 +C\mathcal E_m(t)^{1/2}\mathcal E_{m_0}(t)^{1/2}
       \mathcal E_m(t)
\notag\\
 &+C\|\mathcal F(t)\|_{H^{m-1}}^2,
\label{eq:v3v49-high-diff}
\end{align}
where \(m_0=3\), \(\mathcal F\) denotes prescribed forcing, and
\(C_{\rm low}\) depends on the reciprocal gaps and the
\(W^{1,\infty}\) coefficient norms, but not on a glancing root.  In
particular, on every interval on which
\(\mathcal E_{m_0}^{1/2}\le R\),
\begin{align}
 \mathcal E_m(t)
 +\frac{\beta_0}{4}\int_0^t\mathcal D_m(\sigma)d\sigma
 \le
 \exp\!\left(\int_0^t C_R(\sigma)d\sigma\right)
 \left(
  \mathcal E_m(0)+C\int_0^t\|\mathcal F\|_{H^{m-1}}^2d\sigma
 \right).
\label{eq:v3v49-high-Gronwall}
\end{align}
\end{theorem}

\begin{proof}
Commute \(\mathcal Z^I\), \(|I|\le m-1\), through every hyperbolic
equation.  Its principal commutator is a sum of terms
\[
 (\mathcal Z^J\mathbf c)\,D^2\mathcal Z^KW,
 \qquad |J|\ge1,\quad |J|+|K|\le m-1.
\]
Integrate once by parts when the two derivatives fall on the top-order
factor.  Lemma~\ref{lem:v3v49-tame-products} then bounds the result by a
low coefficient norm times \(\mathcal E_m\), plus a product with one
top-order coefficient factor and only low derivatives of \(W\).  Since
\(H^{m_0}\hookrightarrow W^{1,\infty}\) for \(m_0=3\) in three spatial
dimensions, the latter contribution has the tame form displayed in
\eqref{eq:v3v49-high-diff}.  First-order symmetric-hyperbolic Standard
Model and Dirac blocks give the same or lower-order products.

At the boundary, the differentiated cold incidence has principal part
\(\mathsf B^*M\mathsf B\nabla_u\mathcal Z^Ie\).  Its symmetric part
retains at least \(\beta_0\mathcal D_m/2\).  Incidence commutators are
bounded by the boundary version of
Lemma~\ref{lem:v3v49-tame-products}; Young's inequality absorbs one half
of the available dissipation.  The remaining positive amount is
\(\beta_0\mathcal D_m/4\).

For the static term, commute the local equation
\eqref{eq:v3v49-static-operator} rather than the reduced
Dirichlet-to-Neumann map.  Lemma~\ref{lem:v3v49-carrier-tame} bounds all
carrier commutators in the norms of \eqref{eq:v3v49-Em}.  Differentiating
its local stored energy produces the same coefficient-time-derivative
terms as \eqref{eq:v3v48-Adot}, controlled by \(C_{\rm low}\mathcal E_m\).
The trace source has precisely order \(H^{m-1/2}\), so there is no
uncontrolled top derivative.

Sum over \(I\).  Source pairings are bounded by Cauchy--Schwarz and Young.
This proves \eqref{eq:v3v49-high-diff}; Gronwall gives
\eqref{eq:v3v49-high-Gronwall}.  No step divides by
\(\sqrt{s^2+|\xi|^2}\), so the constants remain uniform through primary
glancing.
\end{proof}

\begin{lemma}[Recovery of normal derivatives]
\label{lem:v3v49-normal-recovery}
Under the noncharacteristic collar and ellipticity hypotheses, the
tangential energy \eqref{eq:v3v49-Em}, the equations, and the sources
control every bulk normal derivative of total order at most \(m\), and
every carrier \(y\)-derivative of total order at most \(m\), with a tame
constant depending on the same gaps.
\end{lemma}

\begin{proof}
The normally hyperbolic equation has a nonzero coefficient multiplying
\(\nabla_n^2W\).  Solve it for that derivative and proceed inductively on
the number of normal derivatives; commutators are controlled by
Lemma~\ref{lem:v3v49-tame-products}.  First-order blocks solve for their
incoming normal derivatives on the noncharacteristic subspace, while
their boundary conditions determine the incoming traces.  For \(\Phi\),
solve \eqref{eq:v3v49-static-operator} for \(\partial_y^2\Phi\) and use
elliptic induction.  The reciprocal principal coefficients are bounded by
the stated noncharacteristic and ellipticity gaps.
\end{proof}

\subsection{Corner compatibility and constraint differentiation}
\label{sec:v3v49-corners}

Let \(\mathcal B(W,\Phi)=0\) denote all boundary equations, including the
trace \(\Phi(0)=\alpha_*e(W)\).  Define the formal time jets
\(W_j=\partial_t^jW|_{t=0}\) recursively from the reduced field equations
and initial data.

\begin{definition}[Compatibility through order \(m-1\)]
\label{def:v3v49-compatibility}
Initial data are compatible through order \(m-1\) when
\begin{equation}
 \left.
 \partial_t^j\mathcal B(W,\Phi)
 \right|_{t=0}=0,
 \qquad 0\le j\le m-1,
\label{eq:v3v49-compatibility}
\end{equation}
after every time derivative has been replaced by its recursively computed
jet, and when the Einstein, Standard Model gauge, and material constraints
hold for the corresponding initial jets.
\end{definition}

\begin{proposition}[No new differentiated constraint source]
\label{prop:v3v49-diff-constraints}
Under Definition~\ref{def:v3v49-compatibility}, the differentiated
harmonic constraints satisfy homogeneous initial and incoming boundary
data through order \(m-1\).  Consequently the conclusion
\(C_\mu=0\) of Theorem~\ref{thm:v3v48-constraint-closure} holds in the
\(H^{m-1}\) energy class.
\end{proposition}

\begin{proof}
Differentiate the covariant source Ward identity and the boundary
Ward--Codazzi identity together with the material and carrier equations.
Because these identities are tensor identities on solutions, each
commutator is a homogeneous linear combination of lower differentiated
constraints.  The compatibility recursion makes their initial and
incoming boundary jets zero.  Induction in the derivative order and the
vector-wave energy estimate in \eqref{eq:v3v48-C-wave} give zero at every
order up to \(m-1\).
\end{proof}

\subsection{Persistence of the clock, rank, and positivity gaps}
\label{sec:v3v49-gaps}

Let \(\vartheta=X_T\) be the clock scalar and let
\(F:T\Sigma_t\to\mathbb R^2\) be the material coframe.  Denote its
smallest singular value by \(\sigma_-(F)\).  Let \(P\) range over the
finite collection of positive coefficient endomorphisms in the V48
energy, including \(M\), \(\mathsf B^*M\mathsf B\), and the spatial
elliptic tensor of the static carrier.

\begin{theorem}[Quantitative persistence of the open set]
\label{thm:v3v49-gap-persistence}
Assume initially
\begin{equation}
 \vartheta(0)\ge4\chi^2,
 \qquad
 \sigma_-(F(0))\ge2r_-,
 \qquad
 \lambda_{\min}(P(0))\ge2p_-
\label{eq:v3v49-initial-gaps}
\end{equation}
pointwise.  If
\begin{align}
 \int_0^t\|\nabla_u\vartheta\|_{L^\infty}d\sigma
 &\le2\chi^2,
\label{eq:v3v49-clock-small}\
 \int_0^t\|\nabla_uF\|_{L^\infty_{\rm op}}d\sigma
 &\le r_-,
\label{eq:v3v49-rank-small}\
 \int_0^t\|\nabla_uP\|_{L^\infty_{\rm op}}d\sigma
 &\le p_-
\label{eq:v3v49-positive-small}
\end{align}
for every such \(P\), then throughout \([0,t]\)
\begin{equation}
 \vartheta\ge2\chi^2,
 \qquad
 \sigma_-(F)\ge r_-,
 \qquad
 \lambda_{\min}(P)\ge p_-.
\label{eq:v3v49-persisted-gaps}
\end{equation}
For \(m\ge4\), the integrands are controlled by a continuous function of
\(\mathcal E_m^{1/2}\) and the reciprocal metric lapse gap.  Hence every
smooth initial configuration satisfying strict gaps has a nonzero time
on which all gaps persist.
\end{theorem}

\begin{proof}
The fundamental theorem of calculus along the clock flow gives
\[
 \vartheta(t)\ge\vartheta(0)
 -\int_0^t|\nabla_u\vartheta|d\sigma,
\]
which proves the first bound.  The singular-value perturbation inequality
\[
 |\sigma_-(A)-\sigma_-(B)|\le\|A-B\|_{\rm op}
\]
and the analogous Weyl inequality for the least eigenvalue of a
self-adjoint endomorphism give the other two bounds after integration
along the same flow.  Sobolev embedding from
Lemma~\ref{lem:v3v49-tame-products} controls all displayed
\(L^\infty\) derivatives by \(\mathcal E_m^{1/2}\), with the conversion
between coordinate time and clock time bounded by the lapse gap.
Continuity then makes the three integrals small on a sufficiently short
interval.
\end{proof}

\begin{corollary}[A priori continuation alternative]
\label{cor:v3v49-continuation}
Let a smooth solution exist on \([0,T_*)\).  If the reciprocal clock,
rank, positivity, ellipticity, and noncharacteristic gaps stay bounded,
and
\begin{equation}
 \int_0^{T_*}\left(
 1+\|D\mathbf c\|_{L^\infty}
  +\|\partial_t\mathbf c\|_{L^\infty}
  +\mathcal E_{m_0}^{1/2}
 \right)dt<\infty,
 \qquad
 \int_0^{T_*}\|\mathcal F\|_{H^{m-1}}^2dt<\infty,
\label{eq:v3v49-continuation-condition}
\end{equation}
then \(\sup_{t<T_*}\mathcal E_m(t)<\infty\).  Thus failure of the
differentiated estimate at a finite time requires collapse of one of the
open-set gaps or divergence of a norm in
\eqref{eq:v3v49-continuation-condition}.
\end{corollary}

\begin{proof}
Apply \eqref{eq:v3v49-high-Gronwall} on \([0,t]\) and let
\(t\uparrow T_*\).  The assumed integral bounds make its right-hand side
uniform.  This is an a priori alternative; extending the solution past
\(T_*\) additionally requires the construction theorem not yet proved.
\end{proof}

\begin{firewall}[What V49 does and does not construct]
\label{fw:v3v49-scope}
V49 proves the high-order estimate for an existing smooth reduced
classical solution and proves quantitative persistence of its open
geometric set.  It does not construct that solution.  In particular, it
does not yet prove uniform solvability of the linearized iteration,
contraction of successive iterates, existence of compatible nonlinear
corner data, curved BRST closure, or removal of ultraviolet and infrared
regulators.
\end{firewall}

\subsection{V49 conclusion and V50 target}
\label{sec:v3v49-conclusion}

The local static parent closes the differentiated boundary energy without
loss: bulk order \(H^m\), boundary order \(H^{m-1/2}\), and carrier order
\(H^m\) form one tame hierarchy.  Cold-line positivity absorbs incidence
commutators, normal derivatives are recovered from the hyperbolic and
elliptic equations, and differentiated Ward identities preserve the
harmonic constraints.  Strict clock, coframe-rank, and coefficient
positivity gaps persist for a quantitative time controlled by the same
energy.

V50 is a mandatory cross-program audit.  It should inspect the newest
Navier--Stokes and Yang--Mills notes, record their exact fingerprints, and
test whether their newest obstruction or construction changes the SM
iteration.  It should then formulate the linearized Picard step on a
fixed moving-domain reference geometry, including the static elliptic
solve, corner compatibility, and a difference estimate strong enough to
decide whether an ordinary contraction suffices or a Nash--Moser scheme
is required.
\clearpage
\section{V50: compulsory companion audit and a fixed-cylinder Picard map}
\label{sec:v3v50}

\subsection{Compulsory NS/YM audit}
\label{sec:v3v50-audit}

The audit was performed immediately after V49.  The active companion
files at that time were:
\begin{center}
\begin{tabular}{p{0.54\textwidth}r p{0.28\textwidth}}
\hline
file & bytes & SHA256\\
\hline
\texttt{Renewal\_NS\_Stability\_Research\_Notes\_V31.tex}
&887537&\texttt{F1121B62238AD5491BB7CF03635EA9934942EDF573ED8FAAA9EE79E1D38A6696}\\
\texttt{Renewal\_Yang\_Mills\_Mass\_Gap\_Research\_Notes\_V91.tex}
&1496919&\texttt{354CB2844CEE45E572E9FDC2806B07E130BEC0B3252D70BB489E45D33F686CAE}\\
\hline
\end{tabular}
\end{center}
The NS fingerprint is unchanged from the V45 audit.  The YM file retains
the name V91 but has been extended in place since V45; the fingerprint
above is the audited state used here.

The NS V31 source-to-tail theorem makes the missing scale explicit.  Heat
formation gives the critical first moment only at the
\(A^{-1/4}\mathcal B_{\rm NS}\) source level.  Rewriting it with the weak
energy-paid source returns the same missing powers in dyadic weights, and
the high-parent identity returns the continuation quantity
\(\int\|\Lambda u\|_2^4dt\).  The lesson for V50 is exact: if the material
coframe enters a principal coefficient through one derivative of \(X\),
the Picard difference norm must pay that derivative.  It may not be
recovered by attaching a smoothing or trace label to a weaker norm.

The updated YM V91 constructs a mesoscopic-prefix/coherent-join packet on
which a crossing-only square-tail normalization fails even after the full
connected carrier is restored.  The same packet is paid by a legal
marked-tree path which retains the actual prefix incidences.  The lesson
for V50 is to keep the cold line, static carrier, strain incidence, and
junction flux assembled until their common energy identity is used.
Separate absolute estimates for these rows can discard the incidence at
which the positive derivative is actually paid.

\begin{audit}[Direction selected by the V50 companion audit]
\label{audit:v3v50-direction}
The SM iteration will therefore use:
\begin{enumerate}
 \item a weak contraction graph norm containing the material coframe
       difference at its true derivative order; and
 \item one assembled cold--static--junction bilinear form, estimated by
       polarization of its positive energy before any rowwise bound.
\end{enumerate}
Neither modification changes the V49 high-order norm.  They determine the
correct one-order-weaker distance for the iteration.
\end{audit}

\subsection{One reference cylinder for every iterate}
\label{sec:v3v50-cylinder}

Let \((S,\Sigma)\) be the initial Cauchy slice and its boundary.  Choose
once, from the initial data, a smooth timelike reference field
\(u_{\rm ref}\) transverse to \(S\) and tangent to the timelike boundary.
Its flow defines
\begin{equation}
 \mathscr C_T=[0,T]\times S,
 \qquad
 \mathscr N_T=[0,T]\times\Sigma,
 \qquad
 \mathscr Y_T=[0,T]\times[0,\infty)_y\times\Sigma.
\label{eq:v3v50-reference-cylinder}
\end{equation}
All iterates are tensor fields on these three fixed manifolds.  The
physical clock \(T(x)\), lapse, shift, and projector are coefficients and
unknowns on \(\mathscr C_T\); no iterate-dependent change of variables is
used.

\begin{lemma}[Uniform equivalence on the fixed cylinder]
\label{lem:v3v50-cylinder-equivalence}
Let \(\mathcal U_R\) be a ball of coefficient states satisfying the
persisted gaps \eqref{eq:v3v49-persisted-gaps} and an \(H^m\) bound \(R\).
For sufficiently small \(T\), the physical energy and Sobolev norms of
every state in \(\mathcal U_R\) are uniformly equivalent to fixed
reference norms on \eqref{eq:v3v50-reference-cylinder}.  The equivalence
constants depend only on \(R\), the reference atlas, and the reciprocal
gaps.
\end{lemma}

\begin{proof}
The inverse metric, lapse, induced spatial metric, material fibre metric,
and static elliptic tensor range in a compact subset of the corresponding
positive matrix cones.  Pointwise quadratic forms are therefore uniformly
equivalent.  Their derivatives through order \(m\) are bounded by the
\(H^m\) ball and the tame composition theorem.  A fixed finite partition
of unity then gives uniform equivalence of the Sobolev norms.  Theorem
\ref{thm:v3v49-gap-persistence} supplies the required small time on which
the positive cones are not left.
\end{proof}

This fixed-cylinder choice avoids a composition loss in differences.
Pulling iterate \(r+1\) back by the flow of iterate \(r\) would make the
difference of two pullbacks contain a top derivative of the transported
field.  Here the flow is fixed once, so differences act only on
coefficients.

\subsection{The assembled linear step}
\label{sec:v3v50-linear-step}

Let \(\overline W\in\mathcal U_R\) be a coefficient state.  Freeze every
principal coefficient at \(\overline W\), but leave the unknowns
\((W,\Phi,\Psi_0)\) coupled.  The Picard step is
\begin{align}
 \mathcal L(\overline W)W
 &=\mathcal F(\overline W),
 &&\text{on }\mathscr C_T,
\label{eq:v3v50-linear-bulk}\
 \mathcal L_{\rm stat}(\overline W)\Phi
 &=0,
 &&\text{on }\mathscr Y_T,
\label{eq:v3v50-linear-static}\
 \Phi|_{y=0}
 &=\alpha_*(\overline W)e_{\overline W}(W),
\label{eq:v3v50-linear-trace}\
 \mathcal B_{\rm ass}(\overline W)(W,\Phi,\Psi_0)
 &=\mathcal G(\overline W),
 &&\text{on }\mathscr N_T.
\label{eq:v3v50-linear-boundary}
\end{align}
Here \(e_{\overline W}(W)\) is the strain linearization used in the V42
quotient, and \(\mathcal B_{\rm ass}\) is the complete junction,
de Donder, cold-line, static-flux, and Standard Model boundary operator.
The static equation is solved on each reference-time slice, but its flux
is not inserted as an independently estimated nonlocal row.

Let \(a_{\rm stat}^{\overline W}\) be the positive bilinear form of
\eqref{eq:v3v49-static-operator} and let
\(d_{\rm cold}^{\overline W}\) be the positive cold-line flux form.  The
boundary part of the linear energy identity is the assembled form
\begin{align}
 \mathfrak a_{\overline W}[(W,\Phi),(V,\Xi)]
 :={}&
 a_{\rm stat}^{\overline W}(\Phi,\Xi)
 +d_{\rm cold}^{\overline W}(\nabla_u e(W),\nabla_u e(V))
\notag\\
 &+\mathfrak j_{\overline W}[(W,\Phi),(V,\Xi)],
\label{eq:v3v50-assembled-form}
\end{align}
where \(\mathfrak j\) is the polarized junction flux.  On the trace
constraint surface, variation of the first and third terms produces the
static normal force, while their time derivative is the static stored
energy identity of Proposition~\ref{prop:v3v48-parent-energy}.

\begin{proposition}[Coercivity is an assembled statement]
\label{prop:v3v50-assembled-coercivity}
Uniformly for \(\overline W\in\mathcal U_R\), the diagonal energy flux of
\eqref{eq:v3v50-assembled-form} satisfies
\begin{align}
 \mathfrak a_{\overline W}[(W,\Phi),(W,\Phi)]
 \ge{}&
 c_R\|e(W)\|_{H^{1/2}(\Sigma)}^2
 +\beta_R\|\nabla_ue(W)\|_{L^2(\Sigma)}^2
\notag\\
 &-C_R\|W\|_{L^2(\Sigma)}^2,
\label{eq:v3v50-assembled-coercivity}
\end{align}
with \(c_R,\beta_R>0\).  The same estimate holds after tangential
differentiation through order \(m-1\), up to the tame commutators already
bounded in V49.
\end{proposition}

\begin{proof}
Retain \(\Phi\) and integrate its positive bulk form over \(y\).  The
trace relation and Lemma~\ref{lem:v3v49-carrier-tame} give the positive
\(H^{1/2}\) trace form, modulo its finite-dimensional infrared sector.
The cold flux gives the second positive term by the incidence gap.  The
junction pairing is the common boundary variation of these energies, so
it is included in the time derivative rather than estimated in absolute
value.  Order-zero curvature and frame terms give the final negative
\(L^2\) term.  Commutation and summation use the proof of
Theorem~\ref{thm:v3v49-high-energy}.
\end{proof}

\begin{assessment}[The precise linear-solvability gate]
\label{ass:v3v50-linear-gate}
V48--V50 give uniqueness and uniform a priori estimates for
\eqref{eq:v3v50-linear-bulk}--\eqref{eq:v3v50-linear-boundary}.  A Picard
map is defined only after existence for this linear mixed
hyperbolic--elliptic problem is proved on the fixed cylinder.  V50 does
not turn an energy estimate into an existence assertion.  A compatible
Galerkin or regularized construction is the next analytic gate.
\end{assessment}

\subsection{The correct weak graph distance}
\label{sec:v3v50-distance}

Because principal material coefficients depend on the coframe \(DX\), a
weak norm on \(X\) alone is insufficient.  Define
\begin{align}
 \mathfrak d_{m-1}(W_1,W_2)^2
 :={}&
 \|W_1-W_2\|_{H^{m-1}(S)}^2
 +\|D(X_1-X_2)\|_{H^{m-2}(S)}^2
\notag\\
 &+\|e(W_1)-e(W_2)\|_{H^{m-3/2}(\Sigma)}^2
 +\|\nabla(\Phi_1-\Phi_2)\|_{H^{m-2}(Y)}^2
\notag\\
 &+\mathfrak d_{{\rm cold},m-1}^2
 +\mathfrak d_{{\rm gauge},m-1}^2.
\label{eq:v3v50-weak-distance}
\end{align}
The apparently repeated coframe term records the graph component needed
when the notation \(W\) counts \(X\) at wave order.  It may be suppressed
only after the chosen first-order reduction proves norm equivalence.

\begin{lemma}[Lipschitz coefficient map in the weak graph distance]
\label{lem:v3v50-coefficient-Lipschitz}
For \(W_1,W_2\in\mathcal U_R\),
\begin{equation}
 \|\mathbf c(W_1,DW_1)-\mathbf c(W_2,DW_2)\|_{H^{m-2}}
 \le C_R\mathfrak d_{m-1}(W_1,W_2).
\label{eq:v3v50-coefficient-Lipschitz}
\end{equation}
The same estimate holds for reciprocal lapse, inverse coframe, normalized
incidence, and fibre projectors.
\end{lemma}

\begin{proof}
Join the two states by the straight segment in a coordinate chart of the
open coefficient set.  The fundamental theorem of calculus writes the
difference as the derivative of the smooth coefficient function times
\((W_1-W_2,DW_1-DW_2)\).  Since \(m-2\ge2>3/2\), \(H^{m-2}\) is an
algebra in three dimensions.  The Moser product estimate therefore proves
\eqref{eq:v3v50-coefficient-Lipschitz}.  Matrix inversion and normalized
projection are smooth on a compact set separated from zero by the V49
gaps, so the same proof applies to them.
\end{proof}

\subsection{Difference estimate and the scheme selection}
\label{sec:v3v50-difference}

Assume temporarily that the linear problem in
\eqref{eq:v3v50-linear-bulk}--\eqref{eq:v3v50-linear-boundary} has a
solution operator \(\mathcal S_{\overline W}\) satisfying the V49
estimate uniformly on \(\mathcal U_R\).  Define
\begin{equation}
 \mathcal P(\overline W)
 :=\mathcal S_{\overline W}
   (\mathcal F(\overline W),\mathcal G(\overline W)).
\label{eq:v3v50-Picard-map}
\end{equation}

\begin{lemma}[Static-solve difference without loss]
\label{lem:v3v50-static-difference}
Let \(\Phi_i\) solve the static equation with coefficient state
\(W_i\in\mathcal U_R\) and trace \(e(W_i)\).  Then
\begin{align}
 \|\nabla(\Phi_1-\Phi_2)\|_{H^{m-2}(Y)}
 \le C_R\bigl(
 &\|e(W_1)-e(W_2)\|_{H^{m-3/2}(\Sigma)}
\notag\\
 &+\mathfrak d_{m-1}(W_1,W_2)\bigr).
\label{eq:v3v50-static-difference}
\end{align}
\end{lemma}

\begin{proof}
Subtract the two elliptic equations.  The difference solves the first
operator with boundary datum \(e(W_1)-e(W_2)\) and interior source
\((L_{W_2}-L_{W_1})\Phi_2\).  Apply the order \(m-1\) version of
Lemma~\ref{lem:v3v49-carrier-tame}.  Estimate the source with
Lemma~\ref{lem:v3v50-coefficient-Lipschitz}, the \(H^m\) bound on
\(\Phi_2\), and the algebra property of \(H^{m-2}\).  This gives
\eqref{eq:v3v50-static-difference} without a derivative of the eliminated
square-root map.
\end{proof}

\begin{theorem}[One-order-weaker Picard difference estimate]
\label{thm:v3v50-Picard-difference}
Let \(\overline W_1,\overline W_2\in\mathcal U_R\) have identical initial
data and compatible jets.  Under the temporary linear-solvability
hypothesis above, their Picard images satisfy
\begin{equation}
 \sup_{0\le t\le T}
 \mathfrak d_{m-1}
 \bigl(\mathcal P(\overline W_1),
       \mathcal P(\overline W_2)\bigr)^2
 \le
 C_R T e^{C_RT}
 \sup_{0\le t\le T}
 \mathfrak d_{m-1}(\overline W_1,\overline W_2)^2.
\label{eq:v3v50-Picard-difference}
\end{equation}
Therefore an ordinary Picard contraction, rather than a Nash--Moser
scheme, is analytically sufficient if the linear existence theorem and
the invariant-ball estimate are established with the V49 bounds.
\end{theorem}

\begin{proof}
Let \(Z=\mathcal P(\overline W_1)-
\mathcal P(\overline W_2)\).  Subtraction gives
\begin{align*}
 \mathcal L(\overline W_1)Z
 ={}&
 [\mathcal L(\overline W_2)-\mathcal L(\overline W_1)]
       \mathcal P(\overline W_2)
 +\mathcal F(\overline W_1)-\mathcal F(\overline W_2),
\end{align*}
with the analogous assembled boundary equation.  Estimate at order
\(m-1\), so the source is required only in \(H^{m-2}\).  Lemma
\ref{lem:v3v50-coefficient-Lipschitz}, the uniform high norm \(R\), and
the tame product estimate give
\begin{equation}
 \|\text{bulk difference source}\|_{H^{m-2}}
 \le C_R\mathfrak d_{m-1}(\overline W_1,\overline W_2).
\label{eq:v3v50-bulk-difference-source}
\end{equation}

For the boundary source, first subtract the complete bilinear forms
\(\mathfrak a_{\overline W_1}-
\mathfrak a_{\overline W_2}\).  Polarization, Lemma
\ref{lem:v3v50-coefficient-Lipschitz}, and Lemma
\ref{lem:v3v50-static-difference} give the same right side in the
\(H^{m-3/2}\) trace dual norm.  The positive diagonal part remains on the
left by Proposition~\ref{prop:v3v50-assembled-coercivity}.  This is where
estimating the junction, cold, and static rows separately would lose the
common payer.

The difference has zero initial data because the jets agree.  Apply the
V49 estimate at order \(m-1\), integrate the squared source bound over an
interval of length \(T\), and use Gronwall.  This proves
\eqref{eq:v3v50-Picard-difference}.  Choosing
\(C_RT e^{C_RT}<1\) gives a strict contraction.  Since no estimate loses a
derivative of the input difference, smoothing iterations are not required
at this stage.
\end{proof}

\subsection{Compatibility-preserving affine iterates}
\label{sec:v3v50-compatible-iterates}

Suppose compatible nonlinear initial data through order \(m-1\) have been
given in the sense of Definition~\ref{def:v3v49-compatibility}.  Let
\begin{equation}
 J_{m-1}(t)=\sum_{j=0}^{m-1}\frac{t^j}{j!}W_j
\label{eq:v3v50-jet-polynomial}
\end{equation}
in a collar of \(t=0\), multiplied by a fixed time cutoff.  Write every
iterate as
\begin{equation}
 W^{(r)}=J_{m-1}+t^mV^{(r)}.
\label{eq:v3v50-affine-iterates}
\end{equation}

\begin{proposition}[Common jets remove difference corner sources]
\label{prop:v3v50-common-jets}
If the linear step is written for \(V^{(r+1)}\) after substitution of
\eqref{eq:v3v50-affine-iterates}, every iterate has the prescribed initial
jets through order \(m-1\), and the difference of two consecutive iterates
has homogeneous corner data to that order.
\end{proposition}

\begin{proof}
For \(0\le j\le m-1\),
\(\partial_t^j(t^mV)|_{t=0}=0\).  Hence all iterates have the jets of
\(J_{m-1}\).  Subtracting two representations cancels the polynomial and
leaves \(t^m(V^{(r+1)}-V^{(r)})\), whose first \(m-1\) jets vanish.  The
boundary residual created by the fixed polynomial is included in the
known inhomogeneity of the linear step.
\end{proof}

\begin{firewall}[Conditional conclusion about Picard versus Nash--Moser]
\label{fw:v3v50-Picard-scope}
Theorem~\ref{thm:v3v50-Picard-difference} decides the derivative-counting
question: the local-parent formulation admits an ordinary one-order-weaker
contraction.  It is conditional on a uniform linear solution operator and
an invariant high-order ball.  V50 has not yet constructed that operator,
proved the ball invariant, or produced compatible nonlinear initial data.
It therefore does not yet prove local existence.
\end{firewall}

\subsection{V50 conclusion and V51 target}
\label{sec:v3v50-conclusion}

The compulsory audit has changed the proof architecture rather than merely
being recorded.  The weak distance explicitly pays the coframe derivative
identified by the NS missing-scale lesson.  The boundary difference keeps
the full cold--static--junction carrier assembled until coercivity is used,
as required by the YM mark-routing lesson.  On a single fixed reference
cylinder these choices yield a Lipschitz Picard map in one lower norm and
select ordinary contraction over Nash--Moser, conditional on linear
solvability.

V51 should construct the missing linear solution operator.  A viable route
is a coercive variational/Galerkin approximation for the wave and static
blocks, coupled to the symmetric-hyperbolic Standard Model blocks, with
the normal junction law appearing as a natural boundary term.  The proof
must obtain estimates uniform in the Galerkin cutoff, pass the static trace
constraint to the limit, recover the strong boundary equations, and retain
the differentiated compatibility jets.  Only after that construction may
the invariant-ball and Banach fixed-point arguments be closed.
\clearpage
\section{V51: construction of the mixed linear solution operator}
\label{sec:v3v51}

\subsection{Purpose and abstract fixed-cylinder setting}
\label{sec:v3v51-purpose}

V50 obtained the correct Picard difference estimate but left its linear
solution operator unconstructed.  This note fills that gate for the
classical reduced system under the structural hypotheses already isolated
in V48--V50.  The proof is variational.  It uses the local static carrier
as a positive elliptic form and therefore remains uniform through primary
glancing.

All spaces below live on the fixed cylinder
\eqref{eq:v3v50-reference-cylinder}.  After the trace of the cold field is
removed by a fixed bounded lifting, let
\begin{equation}
 V_{\rm h}\hookrightarrow H_{\rm h}\cong H_{\rm h}^*
 \hookrightarrow V_{\rm h}^*
\label{eq:v3v51-hyperbolic-triple}
\end{equation}
be the Gelfand triple for the second-order hyperbolic variables, including
the cold half-line variable with homogeneous lifted trace.  The norm of
\(V_{\rm h}\) contains one bulk spatial derivative and the infrared
sector.  Let
\begin{equation}
 V_{\rm s}=H^1_0([0,\infty)_y\times\Sigma;E_{\rm STF})
\label{eq:v3v51-static-space}
\end{equation}
be the zero-trace static-carrier space, with the chosen coercive infrared
condition.  Gauge-fixed Dirac and first-order Standard Model variables
belong to a symmetric-hyperbolic energy space \(H_{\rm f}\).

For a frozen coefficient state \(\overline W\in\mathcal U_R\), suppress
it in the notation and assume the following time-dependent forms:
\begin{align}
 m_t(u,v)&=\overline{m_t(v,u)},
 &m_-\|u\|_{H_{\rm h}}^2
 &\le m_t(u,u)\le m_+\|u\|_{H_{\rm h}}^2,
\label{eq:v3v51-mass-form}\
 a_{{\rm h},t}(u,v)&=\overline{a_{{\rm h},t}(v,u)},
 &a_{{\rm h},t}(u,u)+c_0\|u\|_{H_{\rm h}}^2
 &\ge a_0\|u\|_{V_{\rm h}}^2,
\label{eq:v3v51-hyperbolic-form}\
 d_t(v,v)&\ge\beta_0\|\Gamma_uv\|_{L^2(\Sigma)}^2,
\label{eq:v3v51-damping-form}\
 a_{{\rm s},t}(\phi,\phi)&\ge s_0\|\phi\|_{V_{\rm s}}^2.
\label{eq:v3v51-static-form}
\end{align}
All forms are bounded on their displayed spaces, measurable in time, and
their coefficient-time derivatives satisfy the form bounds used in
\eqref{eq:v3v48-Adot}.  The first-order block has a positive symmetrizer
\(S_t\) and a maximally dissipative boundary space.  Cross couplings are
bounded from the energy space to its dual and their principal pieces obey
the variational Ward adjointness which produces the common energy identity.

\subsection{The static minimizer and its Schur form}
\label{sec:v3v51-static-minimizer}

Let \(C_t\Gamma u\in H^{1/2}(\Sigma)\) be the linearized strain trace in
\eqref{eq:v3v50-linear-trace}.  Choose once a bounded extension
\begin{equation}
 E:H^{1/2}(\Sigma)\longrightarrow H^1([0,\infty)\times\Sigma)
\label{eq:v3v51-extension}
\end{equation}
with compact support in \(y\).  Write
\begin{equation}
 \Phi=EC_t\Gamma u+\phi,
 \qquad \phi\in V_{\rm s}.
\label{eq:v3v51-static-lift}
\end{equation}

\begin{lemma}[Static minimizer]
\label{lem:v3v51-static-minimizer}
For every \(t\) and \(u\in V_{\rm h}\), there is a unique
\(\phi_t[u]\in V_{\rm s}\) such that
\begin{equation}
 a_{{\rm s},t}
 \bigl(EC_t\Gamma u+\phi_t[u],\eta\bigr)=0
 \quad\text{for every }\eta\in V_{\rm s}.
\label{eq:v3v51-static-weak}
\end{equation}
The map
\(K_tu:=EC_t\Gamma u+\phi_t[u]\) is linear and bounded
\(V_{\rm h}\to H^1(Y)\).  The Schur form
\begin{equation}
 q_t(u,v):=a_{{\rm s},t}(K_tu,K_tv)
\label{eq:v3v51-Schur-form}
\end{equation}
is bounded, symmetric, and nonnegative on \(V_{\rm h}\).
\end{lemma}

\begin{proof}
Equation \eqref{eq:v3v51-static-weak} is equivalent to
\[
 a_{{\rm s},t}(\phi_t[u],\eta)
 =-a_{{\rm s},t}(EC_t\Gamma u,\eta).
\]
Coercivity \eqref{eq:v3v51-static-form}, boundedness of the trace and
extension maps, and Lax--Milgram give existence, uniqueness, linearity,
and the asserted norm bound.  Symmetry of \(a_{{\rm s},t}\) gives symmetry
of \(q_t\), and its diagonal value is the nonnegative static stored energy.
\end{proof}

\begin{lemma}[Time regularity of the Schur form]
\label{lem:v3v51-Schur-time}
If the coefficients of \(a_{{\rm s},t}\) and \(C_t\) are Lipschitz in
time in their form norms, then \(K_t\) and \(q_t\) are Lipschitz in time.
For almost every \(t\),
\begin{equation}
 |\dot q_t(u,u)|
 \le C_R\bigl(q_t(u,u)+\|u\|_{V_{\rm h}}^2\bigr).
\label{eq:v3v51-qdot}
\end{equation}
\end{lemma}

\begin{proof}
Subtract \eqref{eq:v3v51-static-weak} at times \(t\) and \(r\), test by
\(\phi_t[u]-\phi_r[u]\), and use coercivity.  The result is
\(\|K_t-K_r\|\le C_R|t-r|\).  Polarize
\eqref{eq:v3v51-Schur-form}, divide its time difference by \(t-r\), and
use the form derivative bounds.  Terms containing \(\dot K_tu\) vanish
against zero-trace variations by the Euler equation
\eqref{eq:v3v51-static-weak}; the remaining terms give
\eqref{eq:v3v51-qdot}.
\end{proof}

Define the total second-order spatial form
\begin{equation}
 A_t(u,v)=a_{{\rm h},t}(u,v)+q_t(u,v)+a_{{\rm low},t}(u,v),
\label{eq:v3v51-total-form}
\end{equation}
where the lower form contains bounded curvature and junction terms.  By
increasing the infrared constant if needed,
\begin{equation}
 A_t(u,u)+C_R\|u\|_{H_{\rm h}}^2
 \ge c_R\|u\|_{V_{\rm h}}^2.
\label{eq:v3v51-total-Garding}
\end{equation}

\subsection{Galerkin construction}
\label{sec:v3v51-Galerkin}

Choose a countable set \(\{v_j\}\subset V_{\rm h}\) dense in
\(V_{\rm h}\) and \(H_{\rm h}\), and finite-dimensional maximally
dissipative approximating spaces for the first-order block.  Let
\(V_N=\operatorname{span}\{v_1,\ldots,v_N\}\).  Seek
\begin{equation}
 u_N(t)=\sum_{j=1}^Nc_j(t)v_j,
 \qquad
 \Phi_N(t)=K_tu_N(t).
\label{eq:v3v51-Galerkin-fields}
\end{equation}
The projected second-order equation is
\begin{align}
 m_t(\ddot u_N,v_i)
 +d_t(\dot u_N,v_i)
 +A_t(u_N,v_i)
 +b_t(z_N,v_i)
 =F_t(v_i),
\label{eq:v3v51-Galerkin-wave}
\end{align}
and the projected first-order equation is
\begin{equation}
 (S_t\dot z_N,\zeta_i)+h_t(z_N,\zeta_i)
 +c_t(u_N,\dot u_N;\zeta_i)=G_t(\zeta_i).
\label{eq:v3v51-Galerkin-first}
\end{equation}

\begin{lemma}[Finite-dimensional solvability]
\label{lem:v3v51-finite-solvability}
For projected initial data, equations
\eqref{eq:v3v51-Galerkin-wave}--\eqref{eq:v3v51-Galerkin-first} have a
unique absolutely continuous solution on \([0,T]\).
\end{lemma}

\begin{proof}
The matrices of \(m_t\) and \(S_t\) are uniformly positive definite.
The static coefficient vector is not an additional differential unknown:
Lemma~\ref{lem:v3v51-static-minimizer} expresses it linearly in
\(c(t)\).  Thus the projected system is a finite-dimensional first-order
ODE for \((c,\dot c,z)\) with bounded measurable, in fact Lipschitz,
coefficients.  The Caratheodory linear ODE theorem gives the unique
solution on the whole finite interval.
\end{proof}

\begin{proposition}[Uniform Galerkin energy estimate]
\label{prop:v3v51-Galerkin-energy}
There is a constant independent of \(N\) such that
\begin{align}
 &\sup_{0\le t\le T}
 \left(
  \|\dot u_N(t)\|_{H_{\rm h}}^2
  +\|u_N(t)\|_{V_{\rm h}}^2
  +\|z_N(t)\|_{H_{\rm f}}^2
  +q_t(u_N(t),u_N(t))
 \right)
\notag\\
 &\qquad
 +\int_0^T\|\Gamma_u\dot u_N\|_{L^2(\Sigma)}^2dt
 +\int_0^T\mathcal D_{\rm f}(z_N)dt
\notag\\
 &\le C_{R,T}\left(
  \|u_0\|_{V_{\rm h}}^2+\|u_1\|_{H_{\rm h}}^2
  +\|z_0\|_{H_{\rm f}}^2
  +\|F\|_{L^2(0,T;V_{\rm h}^*)}^2
  +\|G\|_{L^2(0,T;H_{\rm f})}^2
 \right).
\label{eq:v3v51-uniform-energy}
\end{align}
\end{proposition}

\begin{proof}
Multiply \eqref{eq:v3v51-Galerkin-wave} by \(\dot c_i\) and sum.  Pair
\eqref{eq:v3v51-Galerkin-first} with \(z_N\).  The principal cross terms
cancel by the Ward adjointness hypothesis.  The derivatives of \(m_t\),
\(S_t\), and \(q_t\) are bounded by their energy forms; for \(q_t\), use
Lemma~\ref{lem:v3v51-Schur-time}.  The cold and first-order boundary forms
give the two nonnegative integrals.  Gårding's inequality
\eqref{eq:v3v51-total-Garding}, Cauchy--Schwarz, Young, and Gronwall give
\eqref{eq:v3v51-uniform-energy}.  This is the finite-dimensional version
of the assembled V48 identity, so its constants contain no inverse
glancing root.
\end{proof}

\begin{theorem}[Energy-class linear existence and uniqueness]
\label{thm:v3v51-linear-existence}
Under \eqref{eq:v3v51-mass-form}--\eqref{eq:v3v51-static-form} and the
stated symmetric-hyperbolic boundary hypotheses, for
\begin{align*}
 &(u_0,u_1,z_0)\in V_{\rm h}\times H_{\rm h}\times H_{\rm f},\\
 &F\in L^2(0,T;V_{\rm h}^*),
 \qquad G\in L^2(0,T;H_{\rm f}),
\end{align*}
there is a unique weak solution with
\begin{align}
 u&\in L^\infty(0,T;V_{\rm h}),
 &\dot u&\in L^\infty(0,T;H_{\rm h}),
\label{eq:v3v51-u-regularity}\
 z&\in L^\infty(0,T;H_{\rm f}),
 &\Phi(t)&=K_tu(t)\in L^\infty(0,T;H^1(Y)).
\label{eq:v3v51-zPhi-regularity}
\end{align}
The dissipative traces belong to the \(L^2\) spaces in
\eqref{eq:v3v51-uniform-energy}, the variational equations hold for every
test vector in the energy domains, and the same energy estimate holds for
the limit.
\end{theorem}

\begin{proof}
Proposition~\ref{prop:v3v51-Galerkin-energy} gives bounded sequences in
the indicated weak-star spaces.  Banach--Alaoglu supplies a subsequence
converging weak-star to \((u,\dot u,z)\); the dissipative traces converge
weakly in \(L^2\).  Boundedness of \(K_t\) gives weak-star compactness of
\(\Phi_N\) in \(L^\infty H^1\).  For a finite linear combination of
Galerkin test vectors, integrate the projected equations against a smooth
compactly supported time test function.  Every term is linear in the
weakly convergent unknown, so it passes to the limit.  Density extends the
identity to the full energy domains.  Equation
\eqref{eq:v3v51-static-weak} also passes to the limit and uniqueness of its
solution identifies \(\Phi=K_tu\).

Weak lower semicontinuity gives \eqref{eq:v3v51-uniform-energy}.  Standard
Gelfand-triple time regularity recovers the initial data.  Apply the same
estimate to the difference of two solutions to obtain uniqueness.
\end{proof}

\subsection{Strong boundary equations and high regularity}
\label{sec:v3v51-strong}

\begin{theorem}[Compatible smooth data produce the V49 solution operator]
\label{thm:v3v51-high-linear}
Let \(m\ge4\).  Suppose the coefficients have the V49 regularity, the
data and sources have the corresponding \(H^m\) regularity, and the corner
conditions \eqref{eq:v3v49-compatibility} hold through order \(m-1\).
Then the weak solution of Theorem~\ref{thm:v3v51-linear-existence} has the
regularity measured by \(\mathcal E_m\), satisfies the boundary equations
in their strong trace spaces, and obeys
\eqref{eq:v3v49-high-Gronwall}.  In particular it defines the uniform
linear solution operator assumed in V50.
\end{theorem}

\begin{proof}
Smooth the coefficients and compatible data while preserving their common
finite jets, and perform the Galerkin construction for the smoothed
problem.  Commute tangential derivatives through order \(m-1\).  Theorem
\ref{thm:v3v49-high-energy} gives a cutoff-independent estimate; its proof
uses only form bounds which are uniform under the smoothing.  Weak
compactness at every differentiated level gives a limit with finite
\(\mathcal E_m\).

Lemma~\ref{lem:v3v49-normal-recovery} recovers the missing hyperbolic
normal and carrier \(y\)-derivatives.  Green's identity is now legitimate
in the required trace spaces.  The variational natural boundary term is
therefore exactly
\eqref{eq:v3v50-linear-boundary}; the static weak equation and its trace
give \eqref{eq:v3v50-linear-static}--\eqref{eq:v3v50-linear-trace}
strongly.  Lower semicontinuity preserves the V49 estimate.  Uniqueness
shows that the limit is independent of the smoothing sequence.
\end{proof}

\begin{corollary}[Constraint-preserving linear solve]
\label{cor:v3v51-constraint-solve}
If the initial and incoming constraint data vanish and the frozen source
obeys the linearized bulk and boundary Ward identities, the solution in
Theorem~\ref{thm:v3v51-high-linear} satisfies \(C_\mu=0\) and the analogous
Standard Model gauge constraints.
\end{corollary}

\begin{proof}
The limit satisfies the strong boundary domain, so Theorem
\ref{thm:v3v48-constraint-closure} applies.  The same homogeneous
symmetric-hyperbolic energy argument applies to the gauge constraints.
\end{proof}

\begin{firewall}[Structural hypotheses retained]
\label{fw:v3v51-structural}
V51 constructs the linear solver for coefficient states whose forms obey
the V48--V50 positivity, Ward-adjointness, and maximally dissipative domain
hypotheses.  It does not independently derive a nonlinear Standard Model
boundary domain for every fermion representation, nor prove that all
nonlinear compatible initial data exist.  These hypotheses must be
verified for the concrete classical field content before the final local
theorem is promoted to the full SM system.
\end{firewall}

\subsection{V51 conclusion and V52 target}
\label{sec:v3v51-conclusion}

The static constraint can be solved at every Galerkin level by
Lax--Milgram.  Its positive Schur form is time regular, the assembled
Galerkin energy is uniform, weak compactness produces an energy solution,
and V49's differentiated estimate upgrades compatible data to the strong
solution operator required in V50.  No stable normal root is differentiated
or inverted.

V52 should combine this operator with the V50 difference estimate.  It
must choose a complete metric space, prove that the Picard map preserves a
strict high-order ball and all open-set gaps, choose one time satisfying
both invariance and contraction inequalities, and pass the nonlinear Ward
identities to the fixed point.  It must keep the remaining concrete
Standard Model boundary-domain verification as an explicit hypothesis
rather than conceal it in the abstract form theorem.
\clearpage
\section{V52: invariant Picard ball and a conditional classical local theorem}
\label{sec:v3v52}

\subsection{The iteration space}
\label{sec:v3v52-space}

V49 supplies the high-order estimate, V50 supplies the one-order-weaker
difference estimate, and V51 constructs the corresponding linear solution
operator.  They can now be combined without assuming the conclusion of the
iteration.

Fix compatible initial jets through order \(m-1\), \(m\ge4\), and let
\(J_{m-1}\) be \eqref{eq:v3v50-jet-polynomial}.  Let the initial clock,
rank, and form eigenvalue gaps be four times the target gaps used in the
iteration.  For \(R>0\), define \(\mathbb X_{T,R}\) to be the set of
states \((W,\Phi)\) on the fixed cylinder such that:
\begin{enumerate}
 \item \(W-J_{m-1}\) has zero time jets through order \(m-1\) at
       \(t=0\);
 \item \(\sup_{0\le t\le T}\mathcal E_m[W,\Phi](t)\le R^2\);
 \item the static trace and weak elliptic equations hold on every time
       slice;
 \item the clock, rank, hyperbolicity, ellipticity, incidence, and form
       positivity gaps are at least their target values; and
 \item the gauge and boundary variables lie in the fixed structural
       domains assumed in V51.
\end{enumerate}
Equip this set with
\begin{equation}
 d_T(Y_1,Y_2)
 =\sup_{0\le t\le T}
   \mathfrak d_{m-1}(Y_1(t),Y_2(t)),
\label{eq:v3v52-metric}
\end{equation}
where \(\mathfrak d_{m-1}\) is
\eqref{eq:v3v50-weak-distance}.

\begin{lemma}[Completeness of the closed high ball in the weak metric]
\label{lem:v3v52-complete}
The metric space \((\mathbb X_{T,R},d_T)\) is complete after the target
gap inequalities are taken non-strictly.
\end{lemma}

\begin{proof}
Let \(Y_n\) be Cauchy in \(d_T\).  Completeness of the lower Sobolev
spaces gives a strong limit \(Y\) in the metric norm.  The uniform
\(\mathcal E_m\) bound gives a weak-star convergent subsequence in every
top-order component.  The strong lower-order limit identifies the weak
limit uniquely, and weak lower semicontinuity gives
\(\sup_t\mathcal E_m[Y](t)\le R^2\).

Since \(m-1\ge3\), Sobolev embedding makes the coefficient state converge
uniformly with the derivatives needed to define the clock and coframe.
Least singular values and least eigenvalues are continuous, so the closed
gap inequalities pass to the limit.  The common jets pass in the trace
spaces.  Finally subtract the static weak equations and use Lemma
\ref{lem:v3v50-static-difference}; their limit obeys the same trace and
elliptic relations.  Thus \(Y\in\mathbb X_{T,R}\).
\end{proof}

\subsection{Uniform nonlinear bounds}
\label{sec:v3v52-nonlinear-bounds}

Let \(\mathcal P\) be the Picard map
\eqref{eq:v3v50-Picard-map}, now defined by
Theorem~\ref{thm:v3v51-high-linear}.  Smoothness of the local classical
Lagrangian on the open coefficient set gives nondecreasing functions
\(P_R,Q_R<\infty\) such that, for every \(Y\in\mathbb X_{T,R}\),
\begin{align}
 \|\mathcal F(Y)\|_{H^{m-1}}^2
 +\|\mathcal G(Y)\|_{H^{m-3/2}}^2
 &\le P_R,
\label{eq:v3v52-source-bound}\
 \|\mathcal F(Y_1)-\mathcal F(Y_2)\|_{H^{m-2}}
 +\|\mathcal G(Y_1)-\mathcal G(Y_2)\|_{H^{m-5/2}}
 &\le Q_R\mathfrak d_{m-1}(Y_1,Y_2).
\label{eq:v3v52-source-Lipschitz}
\end{align}
Here boundary norms are interpreted in the dual trace space when the
operator is written variationally.

\begin{lemma}[The bounds are tame consequences, not extra smoothing]
\label{lem:v3v52-Moser}
For the local polynomial and smooth rational nonlinearities of the reduced
Einstein, gauge-fixed Standard Model, material, cold, and static equations,
\eqref{eq:v3v52-source-bound}--\eqref{eq:v3v52-source-Lipschitz} follow
from the \(H^m\) ball, the weak graph distance, and the reciprocal gaps.
No derivative above \(m\) occurs.
\end{lemma}

\begin{proof}
Every coefficient is a smooth function of the finite coefficient array
\(\mathbf c(W,DW)\).  Inverses occur only for matrices whose least
singular values are separated from zero.  Apply the Moser composition and
product estimates on the compact range of that array.  The top-order
estimate puts one factor in \(H^m\) and all remaining factors in
\(W^{1,\infty}\).  For a difference, use the integral of the derivative
of the nonlinear map along the segment joining the two coefficient
states, exactly as in Lemma
\ref{lem:v3v50-coefficient-Lipschitz}.  The coframe component of
\eqref{eq:v3v50-weak-distance} pays every occurrence of \(D(X_1-X_2)\).
Boundary orders follow from the trace theorem.  The carrier contribution
is controlled by Lemma~\ref{lem:v3v50-static-difference}.
\end{proof}

\subsection{One time for invariance, gaps, and contraction}
\label{sec:v3v52-small-time}

Let \(C_I(R,T)\) denote the high-energy constant of
Theorem~\ref{thm:v3v51-high-linear}, and let \(C_D(R)\) be the constant in
the difference estimate.  Choose
\begin{equation}
 R^2=4C_I(R,0)\mathcal E_m(0)+1,
\label{eq:v3v52-R-choice}
\end{equation}
by first fixing any radius larger than twice the initial energy constant;
equivalently one may replace the notation on the right by its value on a
preassigned coefficient ball.  Continuity of the estimate at \(T=0\)
then permits a positive \(T\) satisfying simultaneously
\begin{align}
 C_I(R,T)\bigl(\mathcal E_m(0)+TP_R\bigr)
 &\le R^2,
\label{eq:v3v52-invariant-inequality}\
 C_D(R)T e^{C_D(R)T}&\le\frac12,
\label{eq:v3v52-contraction-inequality}\
 T\,G_R&\le
 \min\{2\chi^2,r_-,p_-,\delta_{\rm hyp},\delta_{\rm ell}\},
\label{eq:v3v52-gap-time}
\end{align}
where \(G_R\) bounds the gap derivatives by Sobolev embedding.

\begin{remark}[Noncircular radius selection]
\label{rem:v3v52-radius}
The symbolic occurrence of \(R\) in \(C_I\) does not define an implicit
equation.  Choose a closed coefficient neighborhood of the initial state,
compute its finite form constants, then choose \(R\) larger than the
initial energy in those reference norms.  Shrinking \(T\) keeps every
iterate inside that coefficient neighborhood by
\eqref{eq:v3v52-gap-time}.  Only after these two choices is
\(C_I(R,T)\) evaluated.
\end{remark}

\begin{proposition}[Invariant ball]
\label{prop:v3v52-invariant}
For \(T\) satisfying
\eqref{eq:v3v52-invariant-inequality} and
\eqref{eq:v3v52-gap-time},
\begin{equation}
 \mathcal P(\mathbb X_{T,R})\subseteq\mathbb X_{T,R}.
\label{eq:v3v52-invariant}
\end{equation}
\end{proposition}

\begin{proof}
The high-order linear estimate and
\eqref{eq:v3v52-source-bound} give
\[
 \sup_t\mathcal E_m[\mathcal P(Y)](t)
 \le C_I(R,T)(\mathcal E_m(0)+TP_R)\le R^2.
\]
The affine-jet construction of Proposition
\ref{prop:v3v50-common-jets} preserves the initial jets.  The linear
solver imposes the static relation and structural boundary domains.
Finally \eqref{eq:v3v52-gap-time} and Theorem
\ref{thm:v3v49-gap-persistence} preserve every open-set gap.  These are
exactly the defining conditions of \(\mathbb X_{T,R}\).
\end{proof}

\begin{proposition}[Strict contraction]
\label{prop:v3v52-contraction}
For \(T\) satisfying \eqref{eq:v3v52-contraction-inequality},
\begin{equation}
 d_T(\mathcal P(Y_1),\mathcal P(Y_2))
 \le2^{-1/2}d_T(Y_1,Y_2)
\label{eq:v3v52-contraction}
\end{equation}
for all \(Y_1,Y_2\in\mathbb X_{T,R}\).
\end{proposition}

\begin{proof}
Theorem~\ref{thm:v3v50-Picard-difference} and the Lipschitz source bound
\eqref{eq:v3v52-source-Lipschitz} give the square of
\eqref{eq:v3v52-contraction}, with coefficient
\(C_D(R)Te^{C_D(R)T}\).  Apply
\eqref{eq:v3v52-contraction-inequality} and take square roots.
\end{proof}

\subsection{The fixed point and the physical constraints}
\label{sec:v3v52-fixed-point}

\begin{theorem}[Local well-posedness for the structural classical system]
\label{thm:v3v52-local-structural}
Assume:
\begin{enumerate}
 \item the nonlinear reduced Einstein--Standard Model--clock/material
       equations arise from the local cold-plus-static classical action;
 \item their concrete gauge and boundary domains satisfy the V51
       symmetric-hyperbolic, maximally dissipative, form-coercive, and Ward
       hypotheses;
 \item the initial data belong to the V49 \(H^m\) graph class, satisfy
       compatibility through order \(m-1\), and have strict clock, rank,
       hyperbolicity, ellipticity, incidence, and positivity gaps.
\end{enumerate}
Then there is \(T>0\) and a unique reduced solution in
\(\mathbb X_{T,R}\).  It depends continuously on the initial data in the
one-order-weaker graph topology.  If the nonlinear bulk and boundary Ward
identities and zero initial gauge constraints hold, the fixed point solves
the unreduced classical Einstein equations and preserves the Standard
Model gauge constraints on \([0,T]\).
\end{theorem}

\begin{proof}
Lemmas and Propositions
\ref{lem:v3v52-complete},
\ref{prop:v3v52-invariant}, and
\ref{prop:v3v52-contraction} put a strict contraction on a complete metric
space.  Banach's fixed-point theorem gives a unique \(Y\) with
\(\mathcal P(Y)=Y\).  Strong convergence in the weak graph norm, the
uniform high-order bound, and the tame estimates permit passage to every
nonlinear term; thus the fixed point satisfies the reduced equations and
the local static carrier equation.  The same difference estimate with
nonzero initial difference gives continuous dependence and uniqueness in
the iteration class.

At the fixed point, the matter equations imply the nonlinear covariant
bulk Ward identity, and the complete material/carrier boundary variation
implies the boundary Ward--Codazzi identity.  The generalized harmonic
constraint satisfies its homogeneous normally hyperbolic propagation
equation with the homogeneous incoming boundary domain.  Zero initial
constraint data and uniqueness imply \(C_\mu=0\).  The analogous
symmetric-hyperbolic argument applies to the Standard Model gauge
constraints.
\end{proof}

\begin{corollary}[A genuine local continuation criterion]
\label{cor:v3v52-local-continuation}
Within the structural class of Theorem
\ref{thm:v3v52-local-structural}, a solution extends past a finite time
\(T_*\) whenever its V49 continuation quantities remain finite and its
open-set gaps remain uniformly positive.
\end{corollary}

\begin{proof}
Corollary~\ref{cor:v3v49-continuation} bounds the \(H^m\) energy up to
\(T_*\).  The equations and time regularity give limiting data at
\(T_*\) with the same graph regularity.  Uniform positive gaps place those
data in another admissible coefficient neighborhood.  Apply Theorem
\ref{thm:v3v52-local-structural} with initial time \(T_*\), using the
limiting compatibility jets inherited from the solution, and concatenate
by uniqueness.
\end{proof}

\begin{firewall}[Why this is conditional for the full Standard Model]
\label{fw:v3v52-conditional-SM}
Theorem~\ref{thm:v3v52-local-structural} is a complete local fixed-point
theorem for the structural PDE class stated in its hypotheses.  Promotion
to the entire concrete Standard Model requires a representation-by-
representation construction of the fermion, gauge, Higgs, and ghost
boundary domains, proof of their common nonlinear Ward identity at the
timelike junction, and construction of compatible Einstein--SM initial
data.  V52 does not declare those concrete gates solved.  It is also a
classical local theorem, not a continuum quantum construction.
\end{firewall}

\subsection{V52 conclusion and V53 target}
\label{sec:v3v52-conclusion}

The high-order ball is complete in the weak graph distance, the V51
solver maps it into itself, and the V50 estimate makes that map a strict
contraction on one time interval which also preserves every open
geometric gap.  This proves local well-posedness and constraint propagation
for the structural classical system and upgrades the V49 a priori
alternative to an actual continuation criterion inside that class.

The most upstream remaining classical gate is now the initial-boundary
constraint manifold.  V53 should construct compatible initial jets rather
than assume them.  It should split the Einstein constraints into conformal
elliptic data, add the Standard Model Gauss constraints and material
boundary Ward conditions, compute the corner map through the first
nontrivial orders, and prove an implicit-function theorem near a simple
background.  If the linearized corner map has a cokernel, that cokernel
must be identified before any further local theorem is claimed for the
full field content.
\clearpage
\section{V53: the bulk constraint manifold and the reduced corner obstruction}
\label{sec:v3v53}

\subsection{Why the bulk and corner maps must be separated}
\label{sec:v3v53-purpose}

V52 proves local well-posedness when compatible constrained data are
given.  The word ``compatible'' contains two analytically different
requirements.  The Einstein and Standard Model Gauss constraints are
elliptic equations on the initial slice.  The time derivatives of the
dynamic junction law are corner equations on its boundary.  Solving the
bulk constraints does not automatically solve the corner equations,
because an elliptic correction with zero boundary value generally changes
its normal derivative.

V53 constructs the bulk constraint manifold near a simple background and
then restricts the full corner map to that manifold.  The resulting Schur
complement identifies the exact remaining obstruction.  This prevents an
implicit-function theorem for the bulk constraints from being misquoted as
a construction of initial-boundary data.

Let \(S\) be a compact three-manifold with boundary \(\Sigma\), and take
\(s>5/2\).  Initial geometric data are \((g,K)\), and the classical
Standard Model data are denoted schematically by
\begin{equation}
 \mathfrak m=(A,E,H,\Pi_H,\psi;T,X,\text{carrier data}).
\label{eq:v3v53-matter-data}
\end{equation}
The constraint map is
\begin{align}
 \mathcal H(g,K,\mathfrak m)
 &=R(g)-|K|_g^2+(\operatorname{tr}_gK)^2
   -2\Lambda-2\kappa_g\rho(\mathfrak m),
\label{eq:v3v53-Hamiltonian}\
 \mathcal M_i(g,K,\mathfrak m)
 &=\nabla^j(K_{ij}-g_{ij}\operatorname{tr}_gK)
   -\kappa_gJ_i(\mathfrak m),
\label{eq:v3v53-momentum}\
 \mathcal G^a(A,E,H,\Pi_H,\psi)
 &=D_iE^{ia}-\rho_{\rm gauge}^a(H,\Pi_H,\psi).
\label{eq:v3v53-Gauss}
\end{align}
There is one Gauss equation for each generator of
\(SU(3)\times SU(2)\times U(1)\).

\subsection{A simple background and the correction variables}
\label{sec:v3v53-background}

Choose a smooth background
\begin{equation}
 (\bar g,\bar K,\bar A,\bar E,\bar H,\bar\Pi_H,\bar\psi)
 =(\bar g,0,0,0,0,0,0)
\label{eq:v3v53-background}
\end{equation}
with \(R(\bar g)=2\Lambda\).  The clock and material background is allowed
to be nonzero but satisfies the strict V49 gaps and has zero initial source
in \eqref{eq:v3v53-Hamiltonian}--\eqref{eq:v3v53-Gauss}.  Assume:
\begin{enumerate}
 \item the scalar Dirichlet operator
       \(L_0=-8\Delta_{\bar g}+V_0\) is an isomorphism
       \(H^{s+2}\cap H^1_0\to H^s\);
 \item the vector York operator
       \(\Delta_L=\operatorname{div}_{\bar g}\mathcal L_{\bar g}\), with
       Dirichlet data, has no kernel; and
 \item the componentwise gauge Dirichlet Laplacian has no kernel.
\end{enumerate}
For the zero-source time-symmetric background, \(V_0\) is the derivative
of the cosmological and prescribed conformal source terms.  The first
assumption includes its sign or nonresonance condition.

Let the free seed be
\begin{equation}
 p=(\widetilde g,\sigma,\tau,
       E_{\rm T},A,H,\Pi_H,\psi,T,X,\ldots),
\label{eq:v3v53-free-seed}
\end{equation}
where \(\sigma\) is trace free and \(E_{\rm T}\) is a freely chosen gauge
electric field.  Introduce correction variables
\begin{equation}
 y=(\varphi-1,W,\chi),
 \qquad
 y|_\Sigma=0,
\label{eq:v3v53-corrections}
\end{equation}
and set, in three spatial dimensions,
\begin{align}
 g&=\varphi^4\widetilde g,
\label{eq:v3v53-conformal-g}\
 K&=\varphi^{-2}(\sigma+\mathcal L_{\widetilde g}W)
     +\frac{\tau}{3}g,
\label{eq:v3v53-conformal-K}\
 E&=E_{\rm T}+D_{\widetilde A}\chi
\label{eq:v3v53-electric-correction}
\end{align}
with the standard conformal weights inserted in the matter sources.  Only
their smooth dependence is needed below; at the zero matter background
their energy and charge contributions have no linear term.

\begin{theorem}[Local Einstein--SM bulk constraint manifold]
\label{thm:v3v53-constraint-manifold}
Under the three background isomorphism assumptions, there are
neighborhoods \(\mathcal P\) of the background seed and \(\mathcal Y\) of
zero, and a unique smooth map
\begin{equation}
 p\longmapsto y(p)\in\mathcal Y
\label{eq:v3v53-correction-map}
\end{equation}
such that the data
\eqref{eq:v3v53-conformal-g}--\eqref{eq:v3v53-electric-correction}
satisfy
\begin{equation}
 \mathcal H=0,
 \qquad \mathcal M=0,
 \qquad \mathcal G=0.
\label{eq:v3v53-all-constraints}
\end{equation}
The solution set is a smooth Banach submanifold parameterized by the free
seed \(p\).  Strict clock, rank, and positivity gaps persist after the
neighborhood is reduced.
\end{theorem}

\begin{proof}
Let \(\mathfrak C(p,y)=(\mathcal H,\mathcal M,\mathcal G)\) after the
substitutions above.  Sobolev multiplication for \(s>5/2\), smoothness of
matrix inversion on the open metric cone, and the polynomial Standard
Model source formulas make
\(\mathfrak C:\mathcal P\times\mathcal Y\to
H^s\times H^s(T^*S)\times H^s(\mathfrak g)\) a smooth map.

At the background, its derivative in the correction variables is block
triangular.  The diagonal blocks are
\begin{equation}
 D_y\mathfrak C
 =operatorname{diag}(L_0,\Delta_L,\Delta_{\bar g}^{\mathfrak g}).
\label{eq:v3v53-constraint-linearization}
\end{equation}
Indeed, matter energy, current, and Higgs/fermion charge are at least
quadratic at the zero matter state; the conformal scalar changes the
Hamiltonian constraint by \(L_0\), the longitudinal correction changes
the momentum constraint by \(\Delta_L\), and the electric potential
changes Gauss by the gauge Laplacian.  The assumed Dirichlet
isomorphisms make \eqref{eq:v3v53-constraint-linearization} invertible.
The Banach implicit-function theorem gives the unique smooth map
\eqref{eq:v3v53-correction-map}.  Continuity in \(C^1\), supplied by
Sobolev embedding, preserves all strict pointwise gaps on a smaller
neighborhood.
\end{proof}

\begin{lemma}[Why the York Dirichlet hypothesis is natural]
\label{lem:v3v53-York-coercive}
If no nonzero conformal Killing field of \((S,\bar g)\) vanishes on
\(\Sigma\), then \(\Delta_L:H^{s+2}\cap H^1_0\to H^s\) is an
isomorphism.
\end{lemma}

\begin{proof}
Integration by parts gives
\begin{equation}
 -\langle W,\Delta_LW\rangle
 =\frac12\|\mathcal LW\|_{L^2}^2
\label{eq:v3v53-York-energy}
\end{equation}
for Dirichlet \(W\).  Korn's inequality for the conformal Killing
operator and the kernel assumption give coercivity on \(H^1_0\).
Lax--Milgram gives weak invertibility, and strong elliptic regularity gives
the displayed Sobolev isomorphism.
\end{proof}

\subsection{The corner map restricted to the constraint manifold}
\label{sec:v3v53-corner-map}

Let \(\mathcal K_j(p,y)\) be the order-\(j\) residual of the complete
boundary domain after time derivatives are replaced recursively using the
reduced equations:
\begin{equation}
 \mathcal K_j(p,y)
 =\left.\partial_t^j
   \mathcal B(W,\Phi,\Psi_0)\right|_{t=0},
 \qquad0\le j\le m-1.
\label{eq:v3v53-corner-jets}
\end{equation}
Define
\begin{equation}
 \mathfrak K(p,y)=(\mathcal K_0,\ldots,\mathcal K_{m-1}),
 \qquad
 \mathfrak K_{\rm red}(p)=\mathfrak K(p,y(p)).
\label{eq:v3v53-reduced-corner-map}
\end{equation}

\begin{proposition}[Exact Schur derivative of the reduced corner map]
\label{prop:v3v53-corner-Schur}
At the background,
\begin{equation}
 D\mathfrak K_{\rm red}
 =\mathfrak K_p
  -\mathfrak K_y
   (\mathfrak C_y)^{-1}\mathfrak C_p.
\label{eq:v3v53-corner-Schur}
\end{equation}
Consequently the bulk-constrained data satisfying all corner conditions
form a Banach submanifold near the background if and only if the operator
in \eqref{eq:v3v53-corner-Schur} is split surjective.
\end{proposition}

\begin{proof}
Differentiate
\(\mathfrak C(p,y(p))=0\):
\[
 Dy=-(\mathfrak C_y)^{-1}\mathfrak C_p.
\]
Apply the chain rule to
\(\mathfrak K_{\rm red}(p)=\mathfrak K(p,y(p))\) to obtain
\eqref{eq:v3v53-corner-Schur}.  The submanifold assertion is the implicit-
function theorem applied to \(\mathfrak K_{\rm red}\).  Split
surjectivity is necessary in Banach spaces and supplies the bounded right
inverse needed by that theorem.
\end{proof}

\begin{assessment}[The normal-derivative term cannot be omitted]
\label{ass:v3v53-normal-Schur}
The second term in \eqref{eq:v3v53-corner-Schur} is generally nonzero.
Although the corrections \(y(p)\) have zero Dirichlet value, their normal
derivatives enter the incoming Einstein, gauge, and carrier
characteristics.  Dropping this term would repeat the error of treating a
bulk constraint solve as automatic corner compatibility.
\end{assessment}

\subsection{Adjoint cokernel and global charge balances}
\label{sec:v3v53-cokernel}

Suppose the joint constraint--corner linearization is elliptic in the
Agmon--Douglis--Nirenberg sense after the incoming characteristic traces
are used as boundary variables.  Then
\(D\mathfrak K_{\rm red}\) is Fredholm.  Define the finite-dimensional
corner obstruction space
\begin{equation}
 \mathscr O_{\rm corner}
 :=\ker(D\mathfrak K_{\rm red})^*.
\label{eq:v3v53-obstruction-space}
\end{equation}

\begin{theorem}[Solvability criterion for constrained corner data]
\label{thm:v3v53-corner-Fredholm}
Under the Fredholm hypothesis, a linear corner residual \(r\) can be
removed by a constrained seed variation if and only if
\begin{equation}
 \langle r,\zeta\rangle=0
 \quad\text{for every }\zeta\in\mathscr O_{\rm corner}.
\label{eq:v3v53-Fredholm-orthogonality}
\end{equation}
If \(\mathscr O_{\rm corner}=0\), the nonlinear constrained compatible
data form a local Banach submanifold by Proposition
\ref{prop:v3v53-corner-Schur}.
\end{theorem}

\begin{proof}
For a Fredholm operator, the closure of the range is the orthogonal
complement of the adjoint kernel, and the range is closed.  This proves
\eqref{eq:v3v53-Fredholm-orthogonality}.  A zero cokernel gives
surjectivity; finite-dimensional kernel and closed complement give a
bounded right inverse.  Apply the preceding proposition.
\end{proof}

Two unavoidable members of an obstruction space, when the boundary domain
freezes their corresponding fluxes, have direct integral descriptions.

\begin{proposition}[Gauge-charge corner balance]
\label{prop:v3v53-charge-balance}
For every Lie-algebra element \(q\) which is covariantly constant on the
background and preserves the boundary gauge domain,
\begin{equation}
 \int_S\langle q,\rho_{\rm gauge}\rangle,dV
 =\int_\Sigma\langle q,n_iE^i\rangle,dS.
\label{eq:v3v53-charge-balance}
\end{equation}
Thus a boundary condition imposing \(n_iE^i=0\) requires vanishing total
\(q\)-charge.  The corresponding constant adjoint gauge parameter lies in
the cokernel unless the normal electric flux is a free corner control.
\end{proposition}

\begin{proof}
Pair \eqref{eq:v3v53-Gauss} with \(q\), integrate over \(S\), and use
covariant integration by parts.  Since \(Dq=0\), only the boundary flux
remains, proving \eqref{eq:v3v53-charge-balance}.
\end{proof}

\begin{proposition}[Momentum/Killing corner balance]
\label{prop:v3v53-momentum-balance}
Let \(Y\) be a Killing field of the initial metric which preserves the
corner domain.  Every constrained datum obeys
\begin{equation}
 \kappa_g\int_SJ_iY^i,dV
 =\int_\Sigma
  (K_{ij}-g_{ij}\operatorname{tr}K)n^jY^i,dS.
\label{eq:v3v53-momentum-balance}
\end{equation}
If the boundary domain fixes the right side, this is a necessary global
momentum balance and its adjoint Killing datum is a corner obstruction.
\end{proposition}

\begin{proof}
Pair \eqref{eq:v3v53-momentum} with \(Y\) and integrate by parts.  The
interior term is the contraction of the symmetric momentum tensor with
\(\nabla_{(i}Y_{j)}\), which vanishes for a Killing field.  The remaining
boundary term is exactly \eqref{eq:v3v53-momentum-balance}.
\end{proof}

\subsection{Higher corner jets}
\label{sec:v3v53-higher-jets}

In incoming/outgoing variables, the principal boundary equation has the
form
\begin{equation}
 w_+=D(W)w_-+C(W)z+g(W).
\label{eq:v3v53-incoming-corner}
\end{equation}
After \(j\) time derivatives, the new highest incoming jet occurs with
the identity coefficient; all other terms depend on lower jets.

\begin{lemma}[Triangular corner recursion after the zeroth constrained
solve]
\label{lem:v3v53-corner-recursion}
Assume the constrained zeroth-order corner operator is split surjective
and that the boundary domain remains maximally dissipative.  Then the
linearized higher-jet map
\((\mathcal K_0,\ldots,\mathcal K_{m-1})\) is upper triangular in the
incoming jets with invertible diagonal.  Subject to the differentiated
bulk constraints, no new infinite-dimensional cokernel appears at higher
orders; only time derivatives of the finite-dimensional obstruction
conditions in \eqref{eq:v3v53-Fredholm-orthogonality} remain.
\end{lemma}

\begin{proof}
Differentiate \eqref{eq:v3v53-incoming-corner} \(j\) times.  Its highest
term is \(\partial_t^jw_+\), while differentiation of coefficients and
the outgoing variables uses jets of order at most \(j\).  The field
equations replace time derivatives by spatial jets and the differentiated
constraints remove the constrained combinations.  Ordering variables by
jet number gives an upper triangular operator with the zeroth incoming
map on each diagonal block.  A split right inverse can therefore be
constructed recursively.  Pairing at each step with the adjoint kernel
shows that only the differentiated finite-dimensional compatibility
conditions survive.
\end{proof}

\begin{firewall}[The remaining uncomputed operator]
\label{fw:v3v53-corner-uncomputed}
V53 constructs the bulk Einstein--SM constraint manifold near a simple
background and gives the exact reduced corner derivative and its Fredholm
solvability criterion.  It does not yet verify the complementing condition
or compute \(\mathscr O_{\rm corner}\) for the complete cold--static
junction with every Standard Model representation.  The gauge-charge and
Killing balances above exhibit possible cokernel elements; they are not a
proof that no others exist.
\end{firewall}

\subsection{V53 conclusion and V54 target}
\label{sec:v3v53-conclusion}

Near a time-symmetric simple background, conformal, York, and gauge
potential corrections solve the bulk Einstein and Standard Model
constraints by a genuine Banach implicit-function theorem.  Restricting
the dynamic boundary jets to this constraint manifold produces the Schur
operator \eqref{eq:v3v53-corner-Schur}.  Its adjoint kernel is the exact
corner obstruction space; fixed electric and gravitational fluxes already
force the global charge and momentum balances
\eqref{eq:v3v53-charge-balance} and
\eqref{eq:v3v53-momentum-balance}.

V54 should calculate the frozen half-space normal family of the joint
constraint--corner operator.  It must include the normal derivatives
created by the conformal/York/Gauss corrections, prove or disprove the
complementing condition, and count the adjoint kernel after global charge
and Killing modes are removed.  That calculation, rather than another
bulk implicit-function theorem, decides whether V52 supplies a local
theorem for an open set of concrete constrained SM data.
\clearpage
\section{V54: frozen constraint--corner normal families and the Schur gap}
\label{sec:v3v54}

\subsection{Flat half-space convention}
\label{sec:v3v54-convention}

V53 reduces construction of compatible data to the complementing
condition for the joint constraint--corner operator.  This note computes
the constraint-correction part of that normal family.  Work on
\(x>0\), put the boundary at \(x=0\), and take the geometric outward
normal to be
\begin{equation}
 n=-\partial_x.
\label{eq:v3v54-outward-normal}
\end{equation}
For a nonzero tangential covector, rotate coordinates so that
\(\xi=(k,0)\), \(k=|\xi|>0\).  A decaying scalar harmonic mode is
\(b e^{-kx+i k y}\), and hence
\begin{equation}
 \partial_n(be^{-kx+i k y})|_{x=0}=kb e^{iky}.
\label{eq:v3v54-scalar-DtN}
\end{equation}

The correction operators at the time-symmetric background of V53 are
\begin{equation}
 -8\Delta\quad\text{for the conformal scalar},
 \qquad
 \Delta_LW=\Delta W+\frac13\nabla\operatorname{div}W,
 \qquad
 -\Delta\quad\text{for each Gauss potential}.
\label{eq:v3v54-correction-operators}
\end{equation}
Zeroth-order potentials do not enter the homogeneous normal family.

\subsection{Scalar and gauge blocks}
\label{sec:v3v54-scalar-gauge}

\begin{proposition}[Positive scalar and Gauss normal families]
\label{prop:v3v54-scalar-gauge-DtN}
For \(k>0\), the natural outward flux maps of the decaying correction
modes are
\begin{equation}
 \mathcal T_H(k)=8k
 \quad\text{and}\quad
 \mathcal T_G(k)=kI_{\mathfrak g}.
\label{eq:v3v54-scalar-gauge-DtN}
\end{equation}
They are positive and have no kernel.
\end{proposition}

\begin{proof}
Equation \eqref{eq:v3v54-scalar-DtN} gives the outward derivative.
The weak form of \(-8\Delta\) has boundary flux
\(8\partial_n\); the componentwise Gauss form has flux
\(\partial_n\).  This proves \eqref{eq:v3v54-scalar-gauge-DtN}.
\end{proof}

\subsection{Exact York Dirichlet-to-Neumann matrix}
\label{sec:v3v54-York}

Write
\begin{equation}
 W=(X,Y,Z)e^{iky},
\label{eq:v3v54-York-components}
\end{equation}
where \(X\) is normal, \(Y\) is parallel to \(\xi\), and \(Z\) is the
remaining tangential component.  The decaying solution with boundary
value \((A,C,D)\) is
\begin{align}
 X(x)&=(A+Bx)e^{-kx},
\label{eq:v3v54-X-solution}\
 Y(x)&=(C-iBx)e^{-kx},
\label{eq:v3v54-Y-solution}\
 Z(x)&=De^{-kx},
\label{eq:v3v54-Z-solution}\
 B&=\frac{k}{7}(A-iC).
\label{eq:v3v54-B}
\end{align}
Let
\begin{equation}
 (\mathcal LW)_{ij}
 =\partial_iW_j+\partial_jW_i
  -\frac23\delta_{ij}\partial^aW_a
\label{eq:v3v54-conformal-Killing}
\end{equation}
and define its natural outward traction by
\begin{equation}
 (\mathcal T_LW)_j=n^i(\mathcal LW)_{ij}.
\label{eq:v3v54-York-traction}
\end{equation}

\begin{theorem}[Exact positive York normal family]
\label{thm:v3v54-York-DtN}
The Dirichlet-to-traction map of \(\Delta_L\) is
\begin{equation}
 \mathcal T_L(k)
 =k\begin{pmatrix}
  8/7&6i/7&0\\
  -6i/7&8/7&0\\
  0&0&1
 \end{pmatrix}
\label{eq:v3v54-York-DtN}
\end{equation}
in the \((A,C,D)\) basis.  Its eigenvalues are
\begin{equation}
 \frac{2k}{7},\qquad k,\qquad2k.
\label{eq:v3v54-York-eigenvalues}
\end{equation}
In particular,
\begin{equation}
 \langle b,\mathcal T_L(k)b\rangle
 \ge\frac{2k}{7}|b|^2.
\label{eq:v3v54-York-positive}
\end{equation}
\end{theorem}

\begin{proof}
The equations \(\Delta_LW=0\) reduce to
\begin{align*}
 4X''-3k^2X+ikY'&=0,\\
 3Y''-4k^2Y+ikX'&=0,\\
 Z''-k^2Z&=0.
\end{align*}
Substitution of \eqref{eq:v3v54-X-solution}--%
\eqref{eq:v3v54-Z-solution} makes the coefficients of \(x\) vanish when
the coefficient of \(x\) in \(Y\) is \(-iB\).  The constant terms both
give \eqref{eq:v3v54-B}; hence these are the unique decaying modes with the
prescribed boundary values.

Since \(n=-\partial_x\), direct substitution in
\eqref{eq:v3v54-York-traction} gives
\begin{align*}
 (\mathcal T_LW)_x
 &=\frac{k}{7}(8A+6iC),\\
 (\mathcal T_LW)_y
 &=\frac{k}{7}(-6iA+8C),\\
 (\mathcal T_LW)_z&=kD,
\end{align*}
which is \eqref{eq:v3v54-York-DtN}.  The longitudinal \(2\times2\)
Hermitian block has eigenvalues \((8-6)k/7\) and \((8+6)k/7\); the
transverse eigenvalue is \(k\).  This proves
\eqref{eq:v3v54-York-eigenvalues}--\eqref{eq:v3v54-York-positive}.
\end{proof}

\begin{corollary}[Constraint-correction coercivity]
\label{cor:v3v54-correction-coercivity}
Let \(b=(b_H,b_L,b_G)\) collect the conformal, York, and Gauss Dirichlet
amplitudes.  Their direct-sum flux family
\begin{equation}
 \mathcal T_C(k)
 =8k\oplus\mathcal T_L(k)\oplus kI_{\mathfrak g}
\label{eq:v3v54-total-correction-DtN}
\end{equation}
satisfies
\begin{equation}
 \mathcal T_C(k)\ge\frac{2k}{7}I.
\label{eq:v3v54-total-correction-positive}
\end{equation}
\end{corollary}

\begin{proof}
Combine Proposition~\ref{prop:v3v54-scalar-gauge-DtN} with
Theorem~\ref{thm:v3v54-York-DtN}.
\end{proof}

\subsection{Energy-selected Robin completion}
\label{sec:v3v54-Robin}

Let \(R_C(k)\) be a Hermitian nonnegative order-one boundary form acting
on the correction traces.  It includes any positive static or infrared
corner stiffness allocated to these variables.

\begin{proposition}[Nonzero-frequency complementing condition]
\label{prop:v3v54-Robin-complementing}
For every \(k>0\),
\begin{equation}
 \mathcal N_C(k)=\mathcal T_C(k)+R_C(k)
\label{eq:v3v54-correction-normal-family}
\end{equation}
is invertible and
\begin{equation}
 \|\mathcal N_C(k)^{-1}\|\le\frac{7}{2k}.
\label{eq:v3v54-correction-inverse}
\end{equation}
\end{proposition}

\begin{proof}
For every amplitude \(b\),
\[
 \langle b,\mathcal N_C(k)b\rangle
 \ge\langle b,\mathcal T_C(k)b\rangle
 \ge\frac{2k}{7}|b|^2.
\]
The least singular value is therefore at least \(2k/7\), proving the
claim.
\end{proof}

\begin{counterexample}[The opposite flux sign fails the complementing
condition]
\label{cex:v3v54-wrong-sign}
For the scalar block, the boundary equation
\begin{equation}
 8\partial_n\varphi-8|D_\Sigma|\varphi=0
\label{eq:v3v54-wrong-sign-boundary}
\end{equation}
admits the nonzero decaying mode
\(e^{-kx+iky}\) for every \(k>0\).  Hence the normal orientation and the
variational sign cannot be changed independently.
\end{counterexample}

\begin{proof}
Use \eqref{eq:v3v54-scalar-DtN}: both terms in
\eqref{eq:v3v54-wrong-sign-boundary} equal \(8k\) on the displayed mode.
\end{proof}

\subsection{Coupling to the physical incoming block}
\label{sec:v3v54-coupled}

Theorem~\ref{thm:v3v15-ten-block} and Proposition
\ref{prop:v3v15-incoming-constraints} give an invertible change from
metric normal derivatives to six junction and four incoming-constraint
coordinates.  After this change, let \(D_P(k)\) denote the positive
principal corner form on the physical metric/material/cold/static
quotient, and let \(K(k)\) contain the off-diagonal coupling to the
constraint corrections.  Variational Hessian symmetry forces the full
principal form to be
\begin{equation}
 \mathcal N_{CP}(k)
 =\begin{pmatrix}
   \mathcal N_C(k)&K(k)^*\\
   K(k)&D_P(k)
  \end{pmatrix}.
\label{eq:v3v54-coupled-normal-family}
\end{equation}

\begin{theorem}[Quantitative Schur complementing criterion]
\label{thm:v3v54-Schur-gap}
Suppose, uniformly for \(k>0\),
\begin{equation}
 \mathcal N_C(k)\ge c_CkI,
 \qquad
 D_P(k)\ge c_PkI,
\label{eq:v3v54-diagonal-gaps}
\end{equation}
and
\begin{equation}
 \left\|
 D_P(k)^{-1/2}K(k)\mathcal N_C(k)^{-1/2}
 \right\|\le1-\varepsilon,
 \qquad \varepsilon>0.
\label{eq:v3v54-coupling-gap}
\end{equation}
Then \(\mathcal N_{CP}(k)\) is invertible and
\begin{equation}
 \mathcal N_{CP}(k)
 \ge c_\varepsilon k I
\label{eq:v3v54-full-positive}
\end{equation}
after a uniformly bounded block rescaling.  Thus the joint
constraint--physical corner problem satisfies the nonzero-frequency
complementing condition.
\end{theorem}

\begin{proof}
Complete the square:
\begin{align*}
 \langle(b,p),\mathcal N_{CP}(b,p)\rangle
 ={}&
 \|\mathcal N_C^{1/2}b
   +\mathcal N_C^{-1/2}K^*p\|^2\\
 &+\langle p,
 (D_P-K\mathcal N_C^{-1}K^*)p\rangle.
\end{align*}
Condition \eqref{eq:v3v54-coupling-gap} gives
\[
 K\mathcal N_C^{-1}K^*
 \le(1-\varepsilon)^2D_P.
\]
The Schur complement is therefore at least
\((2\varepsilon-\varepsilon^2)D_P\).  A triangular block change with
uniformly bounded inverse converts the completed square to
\eqref{eq:v3v54-full-positive}.  The bounds
\eqref{eq:v3v54-diagonal-gaps} supply the factor \(k\).
\end{proof}

\begin{assessment}[Relation to the V47 physical margin]
\label{ass:v3v54-V47-relation}
The V47 determinant and V48 time-domain energy provide the candidate
physical gap \(c_P\).  They do not by themselves imply
\eqref{eq:v3v54-coupling-gap}, because V53's conformal/York/Gauss
corrections change normal derivatives in the junction rows.  The Schur
gap is the precise additional corner condition.  At a decoupled simple
background \(K=0\), it holds with \(\varepsilon=1\), and hence persists on
an open neighborhood by continuity.
\end{assessment}

\subsection{Zero tangential frequency and the obstruction projection}
\label{sec:v3v54-zero}

At \(k=0\), all homogeneous order-one Dirichlet-to-Neumann maps in
\eqref{eq:v3v54-total-correction-DtN} vanish.  This does not create a
high-frequency complementing failure; it exposes the global kernel already
identified in V53.

Let \(\Pi_{\mathscr O}\) be the projection onto constant admissible gauge
parameters and gravitational Killing initial data which preserve the
corner domain.  Impose the balance conditions
\begin{equation}
 \Pi_{\mathscr O}
 \left(
  \int_S\rho_{\rm gauge}-\int_\Sigma n\cdot E,
  \quad
  \kappa_g\int_SJ-\int_\Sigma(K-\operatorname{tr}K,g)(n,\cdot)
 \right)=0.
\label{eq:v3v54-zero-balances}
\end{equation}
The finite-dimensional infrared energy fixes the remaining spatial
constant carrier modes.

\begin{theorem}[Corner manifold near a decoupled simple background]
\label{thm:v3v54-compatible-manifold}
Assume the V53 bulk constraint isomorphisms, the concrete maximally
dissipative incoming domains, the zero-mode balances
\eqref{eq:v3v54-zero-balances}, and the Schur gap
\eqref{eq:v3v54-coupling-gap}.  Then the reduced corner derivative
\eqref{eq:v3v53-corner-Schur}, restricted to the orthogonal complement of
\(\mathscr O_{\rm corner}\), is split surjective.  Consequently the
Einstein--SM data satisfying the bulk constraints and the corner equations
through order \(m-1\) form a local Banach submanifold near the decoupled
simple background.
\end{theorem}

\begin{proof}
For \(k>0\), Theorem~\ref{thm:v3v54-Schur-gap} gives the uniform normal
family inverse of order \(-1\).  The incoming constraint coordinate change
is uniformly invertible by Corollary
\ref{cor:v3v15-algebraic-margin}.  Standard elliptic boundary estimates
then make the joint constraint--corner operator Fredholm.  At \(k=0\),
project away the adjoint charge/Killing kernel and use the fixed infrared
inverse on the remaining finite-dimensional sector.  The resulting
operator has zero cokernel and closed range, hence a bounded right inverse.
Proposition~\ref{prop:v3v53-corner-Schur} and Lemma
\ref{lem:v3v53-corner-recursion} give the nonlinear submanifold and all
higher compatible jets.
\end{proof}

\begin{corollary}[Removal of the compatible-data hypothesis near the
simple background]
\label{cor:v3v54-local-data}
For concrete classical field boundary domains satisfying the hypotheses
of Theorem~\ref{thm:v3v54-compatible-manifold}, Theorem
\ref{thm:v3v52-local-structural} applies to an open Banach manifold of
balanced constrained initial data, rather than to isolated assumed data.
\end{corollary}

\begin{firewall}[What remains representation dependent]
\label{fw:v3v54-representation}
The scalar, York, and Gauss normal families are computed exactly in V54.
The full theorem still assumes that the concrete Standard Model boundary
domains reduce to the maximally dissipative blocks used in V51 and that
their off-diagonal corner incidence satisfies
\eqref{eq:v3v54-coupling-gap}.  Chiral fermion boundary conditions and
nonabelian ghost/BRST domains must be checked representation by
representation.  The argument is classical and says nothing yet about a
renormalized continuum measure.
\end{firewall}

\subsection{V54 conclusion and V55 target}
\label{sec:v3v54-conclusion}

The correction Schur term left open in V53 is not sign-indefinite at a
decoupled simple background.  Its scalar and Gauss blocks are \(8k\) and
\(k\), while the exact York traction has eigenvalues
\(2k/7,k,2k\).  The energy-selected Robin sign therefore satisfies the
nonzero-frequency complementing condition.  Coupling to the physical
boundary sector is controlled by the explicit normalized Schur gap
\eqref{eq:v3v54-coupling-gap}; zero modes are precisely the global
charge/Killing balances plus the chosen infrared sector.

V55 is a compulsory companion audit.  After that audit it should begin
the representation-by-representation classical Standard Model boundary
domain: first the gauge and Higgs bosonic sector, where covariant electric
and magnetic characteristic variables make maximal dissipativity and the
Gauss flux explicit.  Fermion and ghost domains should be postponed until
the bosonic Ward and corner forms have been computed exactly.
\clearpage
\section{V55: compulsory audit and the concrete bosonic Standard Model boundary domain}
\label{sec:v3v55}

\subsection{Compulsory companion audit}
\label{sec:v3v55-audit}

The audit immediately preceding this note found:
\begin{center}
\begin{tabular}{p{0.54\textwidth}r p{0.28\textwidth}}
\hline
file & bytes & SHA256\\
\hline
\texttt{Renewal\_NS\_Stability\_Research\_Notes\_V31.tex}
&887537&\texttt{F1121B62238AD5491BB7CF03635EA9934942EDF573ED8FAAA9EE79E1D38A6696}\\
\texttt{Renewal\_Yang\_Mills\_Mass\_Gap\_Research\_Notes\_V92.tex}
&1529019&\texttt{FAA1A8F09443B7AD80D607B2F551BF011B447333C738681A63A731D8572ABB29}\\
\hline
\end{tabular}
\end{center}

NS V31 is unchanged.  Its balanced-helicity plane wave has zero signed
hinge variables but positive total critical coarea.  The SM translation is
that a Gauss or polarization constraint may vanish while the full bosonic
energy remains large.  The iteration must retain the complete electric,
magnetic, and Higgs energy; it may not use the constraint sector as its
positive norm.

YM V92 proves an exact expansion of assembled edge-bank products into
legal marks, and its overlap theorem keeps a repeated bank inside a
noncommutative sandwich until a one-occurrence capacity is available.  The
SM translation is that gauge and Higgs boundary contributions sharing the
same covariant derivative must first be assembled in the Hessian of one
action.  Estimating the Yang--Mills flux, Higgs current, and Gauss row as
independent absolute terms can charge the same normal covariant derivative
twice.

\begin{audit}[Direction selected at V55]
\label{audit:v3v55-direction}
V55 therefore starts with the complete bosonic action and chooses one
variational boundary domain.  Positivity is proved from its full stress
flux.  Gauss propagation is checked afterward and is not used as an energy
payer.
\end{audit}

\subsection{Bosonic action and boundary geometry}
\label{sec:v3v55-action}

Let
\begin{equation}
 G_{\rm SM}=SU(3)\times SU(2)\times U(1)
\label{eq:v3v55-gauge-group}
\end{equation}
with a positive invariant inner product on its Lie algebra, including the
three positive gauge couplings.  Let \(A\) be a connection, \(F\) its
curvature, and \(H\) the electroweak Higgs doublet in its unitary
representation.  On a spacetime region with timelike boundary
\(\mathcal N\), take
\begin{align}
 S_{\rm bos}
 =\int_M\sqrt{-g}\left(
 -\frac14\langle F_{\mu\nu},F^{\mu\nu}\rangle
 -\langle D_\mu H,D^\mu H\rangle
 -V(H)
 \right)d^4x.
\label{eq:v3v55-bosonic-action}
\end{align}
Let \(n\) be the outward spacelike unit normal, \(u\) the boundary clock,
and \(e_A\), \(A=1,2\), an orthonormal spatial frame tangent to
\(\mathcal N\).

The concrete absolute/Neumann boundary domain is
\begin{align}
 A_n&=0,
\label{eq:v3v55-An}\
 n^\mu F_{\mu a}&=0
 \quad\text{for every }a\in T\mathcal N,
\label{eq:v3v55-absolute-F}\
 D_nH&=0.
\label{eq:v3v55-Higgs-Neumann}
\end{align}
The first condition fixes the normal gauge component.  The residual gauge
transformations obey the covariant Neumann condition
\begin{equation}
 D_n\epsilon=0.
\label{eq:v3v55-gauge-parameter-Neumann}
\end{equation}

\begin{proposition}[Variational boundary domain]
\label{prop:v3v55-variational}
The first variation of \eqref{eq:v3v55-bosonic-action} has no boundary
term on variations preserving
\eqref{eq:v3v55-An}--\eqref{eq:v3v55-Higgs-Neumann}.  The domain is
preserved by infinitesimal gauge transformations satisfying
\eqref{eq:v3v55-gauge-parameter-Neumann} and the corresponding tangential
compatibility equation.
\end{proposition}

\begin{proof}
Integration by parts gives the boundary variation
\begin{equation}
 -\int_{\mathcal N}\sqrt{-\gamma}left(
  \langle n_\mu F^{\mu a},\delta A_a\rangle
  +2\operatorname{Re}\langle D_nH,\delta H\rangle
 \right)d^3x.
\label{eq:v3v55-boundary-variation}
\end{equation}
Condition \eqref{eq:v3v55-An} restricts \(\delta A_n=0\); the remaining
terms vanish by \eqref{eq:v3v55-absolute-F} and
\eqref{eq:v3v55-Higgs-Neumann}.  Under
\(\delta_\epsilon A_n=D_n\epsilon\), equation
\eqref{eq:v3v55-gauge-parameter-Neumann} preserves \(A_n=0\).
Gauge covariance gives
\(\delta_\epsilon(n\lrcorner F)=[n\lrcorner F,\epsilon]\) and
\(\delta_\epsilon(D_nH)=\epsilon\cdot D_nH\), so their zero sets are
preserved.  In a curved collar, preservation of the chosen gauge-fixing
representative adds the standard tangential compatibility equation; it is
lower order and does not change the covariant zero sets.
\end{proof}

\subsection{Exact energy and stress flux}
\label{sec:v3v55-flux}

Define
\begin{align}
 E_n&=F_{un},& E_A&=F_{uA},
\label{eq:v3v55-electric}\
 B_A&=F_{nA},& B_n&=F_{12},
\label{eq:v3v55-magnetic}\
 \Pi_H&=D_uH,& N_H&=D_nH.
\label{eq:v3v55-Higgs-characteristics}
\end{align}
Up to the positive coupling weights, the bosonic energy density is
\begin{equation}
 \mathcal E_{\rm bos}
 =\frac12\left(
  |E_n|^2+|E_A|^2+|B_n|^2+|B_A|^2
 \right)
 +|\Pi_H|^2+|D_AH|^2+|N_H|^2+V(H).
\label{eq:v3v55-bosonic-energy}
\end{equation}

\begin{theorem}[Zero normal energy and momentum exchange]
\label{thm:v3v55-zero-flux}
On \eqref{eq:v3v55-absolute-F}--\eqref{eq:v3v55-Higgs-Neumann}, the
bosonic stress satisfies
\begin{equation}
 T^{\rm bos}_{na}=0
 \quad\text{for every }a\in T\mathcal N.
\label{eq:v3v55-zero-stress-flux}
\end{equation}
In particular the energy flux \(T^{\rm bos}_{nu}\) vanishes exactly.
\end{theorem}

\begin{proof}
For Yang--Mills,
\[
 T^{\rm YM}_{na}
 =\langle F_{n\lambda},F_a{}^\lambda\rangle
\]
because \(g_{na}=0\).  Every nonzero component \(F_{n\lambda}\) has
\(\lambda\) tangent to the boundary, and is zero by
\eqref{eq:v3v55-absolute-F}.  For the Higgs field,
\[
 T^H_{na}=2\operatorname{Re}\langle D_nH,D_aH\rangle,
\]
which vanishes by \eqref{eq:v3v55-Higgs-Neumann}.  Summing the two proves
\eqref{eq:v3v55-zero-stress-flux}.
\end{proof}

Introduce physical tangential characteristics
\begin{equation}
 w_A^\pm=E_A\pm B_A,
 \qquad
 z^\pm=\Pi_H\pm N_H.
\label{eq:v3v55-bosonic-characteristics}
\end{equation}

\begin{proposition}[Maximal conservative physical domain]
\label{prop:v3v55-maximal-domain}
The physical part of
\eqref{eq:v3v55-absolute-F}--\eqref{eq:v3v55-Higgs-Neumann} is
\begin{equation}
 w_A^+=w_A^- ,
 \qquad z^+=z^-.
\label{eq:v3v55-reflection-law}
\end{equation}
It is a maximal isotropic, hence maximally dissipative, subspace for the
bosonic boundary flux form.  More generally
\begin{equation}
 w^+=R_gw^- ,
 \qquad z^+=R_Hz^- ,
 \qquad \|R_g\|,\|R_H\|\le1
\label{eq:v3v55-contractive-law}
\end{equation}
is maximally dissipative, while the variational domain above is the
unitary-reflection case \(R_g=R_H=I\).
\end{proposition}

\begin{proof}
The gauge flux is a fixed positive multiple of
\(|w^+|^2-|w^-|^2\), and the Higgs flux is a fixed positive multiple of
\(|z^+|^2-|z^-|^2\), with sign determined by the outward convention.
Equation \eqref{eq:v3v55-reflection-law} makes both differences zero and
has dimension equal to the number of incoming characteristics.  It is
therefore maximal isotropic.  For \eqref{eq:v3v55-contractive-law}, the
corresponding difference has the dissipative sign because each reflection
is a contraction; maximality follows from the same dimension count.
\end{proof}

\begin{assessment}[Gauss is not the energy norm]
\label{ass:v3v55-Gauss-not-energy}
The Gauss constraint controls the covariant divergence of \(E\).  It does
not control divergence-free electric modes, magnetic modes, or Higgs
energy.  The positive norm used in V49--V52 is therefore
\eqref{eq:v3v55-bosonic-energy}, not a norm of \(\mathcal G\).  This is the
direct consequence of the NS balanced-sector nullspace in the V55 audit.
\end{assessment}

\subsection{Gauss and Lorenz constraint propagation}
\label{sec:v3v55-constraints}

The Yang--Mills--Higgs equations are
\begin{align}
 D_\mu F^{\mu\nu}&=j_H^\nu,
\label{eq:v3v55-YM-equation}\
 D_\mu D^\mu H-V'(H)&=0.
\label{eq:v3v55-Higgs-equation}
\end{align}
Gauge invariance gives
\begin{equation}
 D_\nu j_H^\nu=0
\label{eq:v3v55-Higgs-current-conservation}
\end{equation}
on the Higgs equation.

\begin{proposition}[No normal gauge-charge leakage]
\label{prop:v3v55-no-charge-leak}
The Higgs gauge current obeys
\begin{equation}
 j_H^n=0
\label{eq:v3v55-normal-current-zero}
\end{equation}
on \eqref{eq:v3v55-Higgs-Neumann}.  Together with
\(E_n=F_{un}=0\), the integrated Gauss balance is preserved and requires
zero total gauge charge in every covariantly constant admissible charge
direction.
\end{proposition}

\begin{proof}
The normal Higgs current is the anti-Hermitian part of
\(\langle H,D_nH\rangle\), so it vanishes when \(D_nH=0\).  The time
derivative of integrated Gauss charge is the boundary integral of
\(j_H^n\) plus the normal electric flux contribution.  Both vanish.  The
remaining total-charge condition is exactly
Proposition~\ref{prop:v3v53-charge-balance} with \(n\cdot E=0\).
\end{proof}

For the gauge-fixed potential equation, let \(C_A\) be the Lorenz/background
gauge constraint.  The gauge-fixed Euler equation and the current Ward
identity give a homogeneous covariant wave equation
\begin{equation}
 D^\mu D_\mu C_A+\mathcal R_A(C_A)=0,
\label{eq:v3v55-gauge-constraint-wave}
\end{equation}
where \(\mathcal R_A\) is zeroth order in \(C_A\).

\begin{theorem}[Bosonic gauge-constraint closure]
\label{thm:v3v55-gauge-constraint-closure}
Use the curvature-corrected absolute Hodge boundary domain whose flat
principal part is \eqref{eq:v3v55-An}--\eqref{eq:v3v55-absolute-F}.  It
induces the homogeneous Neumann constraint condition
\begin{equation}
 D_nC_A+L_0C_A=0
\label{eq:v3v55-constraint-Neumann}
\end{equation}
with \(L_0\) of order zero.  If \(C_A\) and its initial time derivative
vanish, then \(C_A=0\) on the interval of existence.
\end{theorem}

\begin{proof}
In a flat collar, \(A_n=0\) and
\(F_{na}=0\) imply \(\partial_nA_a=0\).  Differentiating
\(C_A=\partial^\mu A_\mu\) normally and using the normal component of the
gauge-fixed wave equation gives \(\partial_nC_A=0\).  In a curved collar,
commutation with the normal and differentiation of the frame add only
second-fundamental-form and connection multiples of \(C_A\), giving
\eqref{eq:v3v55-constraint-Neumann}.  Pair
\eqref{eq:v3v55-gauge-constraint-wave} with \(D_uC_A\).  The Robin term is
lower order and Gronwall controls it.  Zero initial data imply zero energy
and hence \(C_A=0\).
\end{proof}

\subsection{Frozen boundary symbol and the V54 corner form}
\label{sec:v3v55-symbol}

At the trivial bosonic background, a stable gauge-potential mode has
normal root \(\kappa=\sqrt{s^2+|\xi|^2}\).  Conditions
\eqref{eq:v3v55-An}--\eqref{eq:v3v55-absolute-F} give
\begin{equation}
 A_n=0,
 \qquad -\kappa A_a=0
\label{eq:v3v55-flat-gauge-symbol}
\end{equation}
after \(A_n=0\) removes the tangential derivative term.  Higgs Neumann
gives \(-\kappa H=0\).  Thus the stable kernel is zero for
\(\kappa\ne0\); at primary glancing the time-domain energy argument, not
division by \(\kappa\), supplies uniqueness.

\begin{proposition}[Bosonic contribution to the corner Schur gap]
\label{prop:v3v55-bosonic-corner-gap}
At the decoupled zero-field background, the bosonic boundary Hessian is
block diagonal from the gravitational constraint corrections and is
nonnegative.  Hence it does not reduce the V54 Schur gap.  For sufficiently
small background \((F,DH,H)\) in \(C^1\), the mixed stress and current
blocks are a continuous perturbation and preserve
\eqref{eq:v3v54-coupling-gap} on a smaller open neighborhood.
\end{proposition}

\begin{proof}
At zero bosonic field, the first variation of the stress and gauge current
with respect to metric constraint corrections vanishes because both are
quadratic.  The diagonal quadratic form is the positive energy
\eqref{eq:v3v55-bosonic-energy}.  At a nonzero small background, every
mixed block is linear in \(F,DH,H\) and depends continuously on the
coefficients.  Singular-value continuity preserves the strict normalized
Schur gap.
\end{proof}

\begin{firewall}[Scope of the bosonic domain]
\label{fw:v3v55-bosonic-scope}
V55 verifies a concrete classical boundary domain for the gauge and Higgs
bosons.  It is reflecting and forces the global charge balance associated
with zero normal electric flux.  It does not yet include chiral fermions,
Yukawa boundary terms, Faddeev--Popov ghosts, BRST cohomology, or a
renormalized quantum measure.  Other physical boundary interactions may
be possible, but their stress and current must be recomputed from one
action.
\end{firewall}

\subsection{V55 conclusion and V56 target}
\label{sec:v3v55-conclusion}

The absolute gauge domain \(A_n=0\), \(n\lrcorner F=0\), together with
covariant Higgs Neumann data, is variational, maximally conservative, has
zero normal stress and gauge-current flux, propagates the bosonic gauge
constraint, and contributes a nonnegative block to the V54 corner form.
It therefore verifies the bosonic part of the structural boundary-domain
hypothesis in V52 near the simple background.

V56 should treat the chiral fermions.  It must choose a local boundary
projector for each Standard Model representation, compute the Dirac normal
current and stress flux, test compatibility with the Higgs Yukawa maps,
and determine whether a single gauge-invariant maximal isotropic domain
exists without adding mirror boundary degrees of freedom.  If chirality
and gauge invariance obstruct such a projector, the obstruction must be
stated before ghosts or quantum determinants are introduced.
\clearpage
\section{V56: chiral fermion projectors from the material boundary frame}
\label{sec:v3v56}

\subsection{The apparent chiral obstruction and the available structure}
\label{sec:v3v56-purpose}

A Lorentz-invariant MIT bag projector pairs opposite four-dimensional
chiralities.  The left- and right-handed Standard Model multiplets are not
isomorphic gauge representations in the electroweak symmetric phase, so
that construction is not a gauge-equivariant boundary condition for the
minimal Standard Model.

The renewal system contains more boundary structure than an ordinary bag:
the clock \(u\) and the rank-two material coframe supply an oriented
orthonormal spatial frame \((e_1,e_2)\) tangent to \(\mathcal N\).  Spatial
Clifford multiplication by \(e_1\) acts within each Weyl bundle and
anticommutes with the normal characteristic matrix.  V56 uses that fact to
construct a gauge-equivariant maximal isotropic projector for every chiral
multiplet separately.  The construction selects the physical material
frame and is therefore not boundary Lorentz invariant.

Let \(S_\chi\), \(\chi=L,R\), be a Weyl spin bundle and let \(R\) be any
unitary Standard Model representation.  The one-particle boundary fibre is
\begin{equation}
 \mathcal E_{\chi,R}=S_\chi|_{\mathcal N}\otimes R.
\label{eq:v3v56-boundary-fibre}
\end{equation}
With respect to the clock inner product, define the Hermitian spatial
Clifford matrices
\begin{equation}
 J_\chi=\alpha_\chi(n),
 \qquad
 B_\chi=\alpha_\chi(e_1).
\label{eq:v3v56-JB}
\end{equation}
They obey
\begin{equation}
 J_\chi^*=J_\chi,
 \quad B_\chi^*=B_\chi,
 \quad J_\chi^2=B_\chi^2=I,
 \quad B_\chi J_\chi+J_\chi B_\chi=0.
\label{eq:v3v56-Clifford-relations}
\end{equation}

For \(\varepsilon_\chi\in\{+1,-1\}\), impose
\begin{equation}
 P_{\chi,-}\psi_\chi=0,
 \qquad
 P_{\chi,-}=\frac12(I-\varepsilon_\chi B_\chi),
\label{eq:v3v56-Weyl-projector}
\end{equation}
or equivalently
\(B_\chi\psi_\chi=\varepsilon_\chi\psi_\chi\).

\subsection{Green form and maximal isotropy}
\label{sec:v3v56-Green}

The symmetric Weyl action has Green boundary form
\begin{equation}
 \mathfrak b_\chi(\psi,\phi)
 =\int_\Sigma
  \langle\psi,J_\chi\phi\rangle,dS.
\label{eq:v3v56-Green-form}
\end{equation}
This is also the normal probability-current form and, after multiplication
by the frequency, the principal energy-flux form.

\begin{theorem}[Gauge-equivariant maximal isotropic Weyl domain]
\label{thm:v3v56-maximal-isotropic}
For every unitary gauge representation \(R\), the subbundle
\begin{equation}
 \mathcal L_{\chi,R}^{\varepsilon}
 =\ker(B_\chi-\varepsilon I)\otimes R
\label{eq:v3v56-Lagrangian-subbundle}
\end{equation}
is maximal isotropic for \(\mathfrak b_\chi\).  It is invariant under the
full Standard Model gauge group.  Hence
\eqref{eq:v3v56-Weyl-projector} is a local gauge-equivariant conservative,
and therefore maximally dissipative, boundary condition for a single
chiral multiplet.
\end{theorem}

\begin{proof}
If \(B_\chi\psi=\varepsilon\psi\) and
\(B_\chi\phi=\varepsilon\phi\), then
\begin{align*}
 \langle\psi,J_\chi\phi\rangle
 &=\langle B_\chi\psi,B_\chi J_\chi\phi\rangle\\
 &=-\langle\psi,J_\chi\phi\rangle
\end{align*}
by \eqref{eq:v3v56-Clifford-relations}.  Thus the Green form vanishes.
Anticommutation with \(J_\chi\) gives an isomorphism between the two
eigenspaces of \(B_\chi\), so each has half the spinor dimension.  An
isotropic subspace of a nondegenerate balanced Hermitian form has dimension
at most one half, proving maximality.  The projector acts only on the
spinor factor and therefore commutes with the unitary action on \(R\).
\end{proof}

\begin{corollary}[Zero normal gauge and fermion-number currents]
\label{cor:v3v56-zero-current}
For every gauge generator \(T^a\) and every fermion multiplet satisfying
\eqref{eq:v3v56-Weyl-projector},
\begin{equation}
 j_n=\langle\psi,J_\chi\psi\rangle=0,
 \qquad
 j_n^a=\langle\psi,J_\chi T^a\psi\rangle=0.
\label{eq:v3v56-zero-currents}
\end{equation}
\end{corollary}

\begin{proof}
The first identity is the diagonal case of Theorem
\ref{thm:v3v56-maximal-isotropic}.  Since \(T^a\) acts on \(R\), it
commutes with \(B_\chi\), so \(T^a\psi\) belongs to the same maximal
isotropic subspace.  Apply the off-diagonal identity with
\(\phi=T^a\psi\).
\end{proof}

\subsection{The complete Standard Model fermion list}
\label{sec:v3v56-SM-list}

Per generation, the minimal Standard Model representations are
\begin{align}
 Q_L&:(\mathbf3,\mathbf2)_{1/6},
 &L_L&:(\mathbf1,\mathbf2)_{-1/2},
\label{eq:v3v56-left-reps}\
 u_R&:(\mathbf3,\mathbf1)_{2/3},
 &d_R&:(\mathbf3,\mathbf1)_{-1/3},
 &e_R&:(\mathbf1,\mathbf1)_{-1}.
\label{eq:v3v56-right-reps}
\end{align}
Apply \eqref{eq:v3v56-Weyl-projector} to each flavour copy and each
representation.  No gauge intertwiner between different representations
is required.

The Yukawa endomorphism acts on internal and chirality indices but contains
no normal derivative.  In Hamiltonian form its left-to-right block
\(Y(H)\) obeys the Clifford relation
\begin{equation}
 B_RY(H)=-Y(H)B_L
\label{eq:v3v56-Yukawa-Clifford}
\end{equation}
because the restrictions of \(\alpha(e_1)\) to opposite Weyl chiralities
have opposite Pauli signs, while \(Y(H)\) is spin scalar.

\begin{proposition}[Yukawa-compatible sign choice]
\label{prop:v3v56-Yukawa-sign}
Choose one sign \(\varepsilon\) for all left-handed multiplets and the
opposite sign for all right-handed multiplets:
\begin{equation}
 \varepsilon_L=\varepsilon,
 \qquad
 \varepsilon_R=-\varepsilon.
\label{eq:v3v56-chiral-signs}
\end{equation}
Then every quark and charged-lepton Yukawa block maps the left boundary
subspace into the right boundary subspace and its adjoint maps back.  The
unpaired minimal-SM neutrino remains in its own left Weyl domain and causes
no rank defect in the boundary form.
\end{proposition}

\begin{proof}
If \(B_L\psi_L=\varepsilon\psi_L\), then
\[
 B_RY(H)\psi_L
 =-Y(H)B_L\psi_L
 =-\varepsilon Y(H)\psi_L.
\]
Thus \(Y(H)\psi_L\) lies in the right eigenspace prescribed by
\eqref{eq:v3v56-chiral-signs}.  Taking adjoints gives the reverse map.
The neutrino has no right Yukawa image in the minimal model, but Theorem
\ref{thm:v3v56-maximal-isotropic} already gives a complete half-dimensional
domain for its left Weyl operator.
\end{proof}

\begin{assessment}[Why the material frame changes the usual obstruction]
\label{ass:v3v56-material-frame}
Without a distinguished tangent vector, there is no natural local map
between the two normal-characteristic lines of one Weyl spinor; a
boundary-Lorentz-invariant bag condition normally pairs opposite
chiralities.  The rank-two material coframe supplies \(e_1\) and hence
\(B_\chi\).  The construction is diffeomorphism covariant when the material
fields transform, but it intentionally selects their physical frame.
It is defined only while the V49 rank gap persists.
\end{assessment}

\subsection{Variational and energy identities}
\label{sec:v3v56-energy}

For the symmetrized fermion action
\begin{equation}
 S_F=\frac i2\int_M\sqrt{-g}left(
 \overline\psi\gamma^\mu D_\mu\psi
 -(D_\mu\overline\psi)\gamma^\mu\psi
 \right)d^4x+S_{\rm Yuk},
\label{eq:v3v56-fermion-action}
\end{equation}
the Yukawa term is algebraic and produces no boundary variation.

\begin{proposition}[Stationarity of the chiral domain]
\label{prop:v3v56-stationary}
If \(\psi\) and its allowed variation \(\delta\psi\) satisfy the same
projector \eqref{eq:v3v56-Weyl-projector}, the boundary term in
\(\delta S_F\) vanishes.  Thus the direct sum of all domains with signs
\eqref{eq:v3v56-chiral-signs} is a variational Standard Model fermion
domain.
\end{proposition}

\begin{proof}
The boundary variation is the skew polarization of the Green form
\eqref{eq:v3v56-Green-form}.  Both arguments lie in the maximal isotropic
subbundle of Theorem~\ref{thm:v3v56-maximal-isotropic}, so each polarized
term vanishes.  The internal Yukawa action has no derivative and hence no
boundary term.
\end{proof}

Assume first that the material spin projector is transported along the
clock:
\begin{equation}
 \nabla_uB_\chi=0.
\label{eq:v3v56-B-transported}
\end{equation}

\begin{theorem}[Conservative fermion energy boundary]
\label{thm:v3v56-fermion-energy}
Under \eqref{eq:v3v56-B-transported}, the normal fermion energy flux
vanishes.  The direct sum of the Standard Model Weyl Hamiltonians with
\eqref{eq:v3v56-Weyl-projector} is symmetric on the boundary domain, and
the Yukawa blocks with \eqref{eq:v3v56-chiral-signs} preserve the common
energy pairing.
\end{theorem}

\begin{proof}
Differentiate \(B_\chi\psi=\varepsilon_\chi\psi\) along \(u\).  Condition
\eqref{eq:v3v56-B-transported} gives
\(B_\chi D_u\psi=\varepsilon_\chi D_u\psi\).  Therefore the polarized
normal current between \(\psi\) and \(D_u\psi\) vanishes by maximal
isotropy.  That polarization is the normal energy flux of the first-order
Dirac system.  Symmetry follows from the Green identity.  The Yukawa
Hamiltonian is Hermitian only after a left-to-right block and its adjoint
are added; Proposition~\ref{prop:v3v56-Yukawa-sign} puts both in the paired
boundary domain, so their bulk energy contributions cancel in the usual
way.
\end{proof}

If \(B_\chi\) is obtained directly from a moving material frame, then
\(\nabla_uB_\chi\) need not vanish.

\begin{proposition}[Moving-projector error is lower order]
\label{prop:v3v56-moving-projector}
For a time-dependent material projector,
\begin{equation}
 |\mathcal F_{n,u}^{F}|
 \le C|\nabla_uB_\chi|\,|\psi|^2
\label{eq:v3v56-moving-flux}
\end{equation}
pointwise, and the corresponding boundary term is controlled by the V49
low coefficient norm times the fermion energy.  It introduces no
derivative loss.
\end{proposition}

\begin{proof}
Differentiation of the projector equation gives
\[
 (B_\chi-\varepsilon_\chi)D_u\psi
 =-(\nabla_uB_\chi)\psi.
\]
The component of \(D_u\psi\) outside the isotropic subspace is therefore
bounded by \(|\nabla_uB_\chi||\psi|\).  Only this component contributes to
the polarized normal flux, yielding
\eqref{eq:v3v56-moving-flux}.  Trace and Sobolev estimates place it among
the lower-order coefficient terms in Theorem
\ref{thm:v3v49-high-energy}.
\end{proof}

\subsection{Gauge constraint and corner contribution}
\label{sec:v3v56-gauge-corner}

Corollary~\ref{cor:v3v56-zero-current} adds no fermionic normal charge
source to Proposition~\ref{prop:v3v55-no-charge-leak}.  At the zero-fermion
background, the fermion stress, gauge current, and Yukawa source are
quadratic, so their first derivatives with respect to the V54 constraint
corrections vanish.

\begin{proposition}[Fermion block preserves the simple-background Schur
gap]
\label{prop:v3v56-Schur-gap}
At \(\psi=0\), the fermion projector supplies a maximal conservative
first-order boundary block and has zero off-diagonal entry in
\eqref{eq:v3v54-coupled-normal-family}.  For sufficiently small fermion
data and bounded material-frame derivatives, it is a continuous lower-
order perturbation of the V54--V55 corner form.
\end{proposition}

\begin{proof}
Maximal conservativity is Theorem
\ref{thm:v3v56-maximal-isotropic}.  All classical fermion currents and
stress components are quadratic in \(\psi\), so their linearization at
zero is zero.  The principal Dirac normal matrix and projector depend
smoothly on the metric and material frame while the clock and rank gaps
hold.  Continuity and the strict Schur margin prove the perturbative
statement.
\end{proof}

\subsection{The remaining tangential stress gate}
\label{sec:v3v56-stress-gate}

The projector kills normal gauge current and, when transported, normal
energy flux.  It does not force every tangential momentum flux
\(T^F_{nA}\) to vanish.  A reflecting fermion can exert pressure and shear
on the material frame which defines its reflector.

\begin{firewall}[A symmetric Dirac domain is not yet the boundary Ward
identity]
\label{fw:v3v56-stress-Ward}
V56 proves the variational Green-domain, gauge-current condition, energy
estimate, Yukawa sign compatibility, and corner-margin perturbation for all
minimal Standard Model chiral multiplets.  It does not yet construct the
boundary stress whose divergence equals the generally nonzero fermion
momentum flux \(T^F_{nA}\).  Because the projector depends on the material
frame, varying that frame produces a boundary reaction force.  This force
must be derived from a local multiplier or edge action before the complete
Codazzi Ward identity assumed in V52 is verified.
\end{firewall}

\subsection{V56 conclusion and V57 target}
\label{sec:v3v56-conclusion}

The rank-two material coframe removes the local chiral projector
obstruction: \(B_\chi=\alpha_\chi(e_1)\) anticommutes with the normal
current matrix and defines a gauge-equivariant maximal isotropic line in
every Weyl spinor.  Opposite projector signs for left- and right-handed
multiplets respect every Yukawa block, including the unpaired minimal-SM
neutrino.  The price is dependence on the physical material frame.

V57 should derive that price variationally.  It should introduce the
minimal local boundary multiplier imposing
\eqref{eq:v3v56-Weyl-projector}, vary the metric and material coframe,
compute the resulting edge stress and material force, and prove the exact
identity
\(D^aS^{F}_{ab}=T^F_{nb}\) on the fermion and multiplier equations.  A
failure of this identity would invalidate the use of the projector in the
Einstein junction, even though its Dirac energy estimate is sound.
\clearpage
\section{V57: the fermion projector multiplier and its boundary Ward identity}
\label{sec:v3v57}

\subsection{A local multiplier action}
\label{sec:v3v57-multiplier}

V56 gives a sound Dirac energy domain but leaves the tangential momentum
reaction of the material-frame reflector unaccounted for.  The projector
must therefore be imposed before the stress is varied.  Let
\begin{equation}
 P_-=\frac12(I-\varepsilon B),
 \qquad B=\alpha(e_1),
\label{eq:v3v57-projector}
\end{equation}
for one Weyl multiplet, and let \(\lambda\) be a boundary spinor in
\(\operatorname{Ran}P_-\) transforming in the same gauge representation
as \(\psi\).  With the real Hermitian pairing
\((v,w)_{\mathbb R}=\operatorname{Re}\langle v,w\rangle\), define
\begin{equation}
 S_{\rm con}[\psi,\lambda;\gamma,T,X]
 =\int_{\mathcal N}\sqrt{-\gamma}\,
  (\lambda,P_-\psi)_{\mathbb R},d^3x.
\label{eq:v3v57-constraint-action}
\end{equation}
Sum \eqref{eq:v3v57-constraint-action} over every Standard Model Weyl
multiplet with the signs \eqref{eq:v3v56-chiral-signs}.

\begin{proposition}[The multiplier imposes exactly the V56 domain]
\label{prop:v3v57-exact-domain}
Variation with respect to \(\lambda\) gives
\begin{equation}
 P_-\psi=0.
\label{eq:v3v57-projector-equation}
\end{equation}
Variation with respect to the unconstrained boundary trace of \(\psi\)
determines \(\lambda\) algebraically and imposes no second condition on
the allowed component \(P_+\psi\).
\end{proposition}

\begin{proof}
The first assertion follows immediately from
\eqref{eq:v3v57-constraint-action}.  Let \(\Theta_D(\psi,\delta\psi)\) be
the boundary term in the first variation of the symmetrized Weyl action.
Its polarization is the Green form
\eqref{eq:v3v56-Green-form}.  On
\(P_-\psi=P_-\delta\psi=0\), maximal isotropy makes \(\Theta_D=0\).
For an arbitrary variation, the component \(P_-\delta\psi\) is paired
with the normal Green covector of \(P_+\psi\).  The multiplier variation
contributes \((\lambda,P_-\delta\psi)_{\mathbb R}\), so its equation sets
\(\lambda\) equal to the negative of that covector.  The
\(P_+\delta\psi\) coefficient vanishes by isotropy.  Thus one half of the
trace is constrained and the multiplier is determined; there is no
overdetermination.
\end{proof}

\begin{corollary}[No new propagating boundary degree of freedom]
\label{cor:v3v57-no-propagating}
The field \(\lambda\) has no time or spatial derivative in
\eqref{eq:v3v57-constraint-action}.  Eliminating it recovers precisely the
V56 maximal isotropic domain and adds no boundary characteristic.
\end{corollary}

\begin{proof}
Its Euler equation is algebraic, as is the equation determining it in
Proposition~\ref{prop:v3v57-exact-domain}.  Algebraic elimination therefore
leaves the same characteristic count.
\end{proof}

\subsection{Frame reaction force}
\label{sec:v3v57-reaction}

The multiplier action vanishes on solutions because \(P_-\psi=0\), but
its variation with respect to the projector does not.  For an arbitrary
variation of the material tangent vector,
\begin{equation}
 \delta B=\delta\alpha(e_1)
 =\alpha(\delta e_1)+\delta_{\gamma,u}\alpha(e_1),
\label{eq:v3v57-delta-B}
\end{equation}
where the second term is the spin-metric variation needed to preserve the
Clifford relations.

\begin{proposition}[Exact on-shell projector variation]
\label{prop:v3v57-projector-variation}
On \eqref{eq:v3v57-projector-equation},
\begin{equation}
 \left.\delta S_{\rm con}\right|_{\psi,\lambda}
 =-\frac{\varepsilon}{2}
  \int_{\mathcal N}\sqrt{-\gamma}\,
  (\lambda,\delta B\,\psi)_{\mathbb R},d^3x.
\label{eq:v3v57-on-shell-variation}
\end{equation}
The volume-density variation contributes zero on shell.  Hence the
material-frame force is a bilinear expression in the reflected fermion
and its algebraically determined reaction multiplier.
\end{proposition}

\begin{proof}
Differentiate \(P_-=(I-\varepsilon B)/2\).  Terms varying
\(\sqrt{-\gamma}\) are proportional to
\((\lambda,P_-\psi)_{\mathbb R}=0\).  Holding the Euler fields fixed leaves
only \((\lambda,\delta P_-\psi)_{\mathbb R}\), which is
\eqref{eq:v3v57-on-shell-variation}.
\end{proof}

Define the fermion reaction stress and clock/material forces by the exact
variational formulas
\begin{align}
 S^{F,\partial}_{ab}
 &:=-\frac{2}{\sqrt{-\gamma}}
   \frac{\delta S_{\rm con}}{\delta\gamma^{ab}},
\label{eq:v3v57-reaction-stress}\
 \mathcal F_T
 &:=\frac{1}{\sqrt{-\gamma}}
   \frac{\delta S_{\rm con}}{\delta T},
 &
 \mathcal F_I
 &:=\frac{1}{\sqrt{-\gamma}}
   \frac{\delta S_{\rm con}}{\delta X^I}.
\label{eq:v3v57-material-forces}
\end{align}
Equations \eqref{eq:v3v57-delta-B}--%
\eqref{eq:v3v57-on-shell-variation} are an explicit local formula for all
three: differentiate the Gram--Schmidt map from \((\gamma,T,X)\) to
\((u,e_1,e_2)\), then insert the Clifford variation.

\begin{lemma}[Tame order of the reaction terms]
\label{lem:v3v57-reaction-tame}
While the clock and material-rank gaps remain positive, the maps in
\eqref{eq:v3v57-reaction-stress}--%
\eqref{eq:v3v57-material-forces} are smooth zeroth-order functions of
\((\gamma,DT,DX,\psi,\lambda)\).  For \(m\ge4\), their differentiated
\(H^{m-1/2}\) boundary norms obey the tame product estimates of V49 and
introduce no derivative loss.
\end{lemma}

\begin{proof}
Gram--Schmidt orthonormalization is a smooth rational matrix function on
the open set where the clock norm and the two coframe singular values are
bounded away from zero.  Clifford multiplication and the spin-metric
variation are smooth algebraic functions of the orthonormal frame.  The
remaining dependence is bilinear in \((\psi,\lambda)\).  Apply the boundary
Moser estimate in Lemma~\ref{lem:v3v49-tame-products}; the only derivative
of a material field is the coframe \(DX\), already present in the V50
graph norm.
\end{proof}

\subsection{Tangential diffeomorphism Ward identity}
\label{sec:v3v57-Ward}

Let \(\mathcal E_\psi,\mathcal E_\lambda\) denote the boundary Euler
expressions obtained from the bulk fermion action plus
\(S_{\rm con}\).  For a vector field \(\zeta\) tangent to
\(\mathcal N\), apply the simultaneous infinitesimal diffeomorphism to all
fields, including the material frame from which \(B\) is built.

\begin{theorem}[Exact fermion projector Ward identity]
\label{thm:v3v57-projector-Ward}
The covariant local action satisfies the off-shell identity
\begin{align}
 D^aS^{F,\partial}_{ab}-T^F_{nb}
 ={}&
 \mathcal F_TD_bT+\mathcal F_ID_bX^I
 +\mathcal W_b(\mathcal E_\psi,\psi)
 +\mathcal W_b(\mathcal E_\lambda,\lambda),
\label{eq:v3v57-offshell-Ward}
\end{align}
where \(\mathcal W_b\) is the standard spinorial Euler pairing generated
by the covariant Lie derivative.  On the fermion and multiplier equations,
\begin{equation}
 D^aS^{F,\partial}_{ab}-T^F_{nb}
 =\mathcal F_TD_bT+\mathcal F_ID_bX^I.
\label{eq:v3v57-onshell-Ward}
\end{equation}
\end{theorem}

\begin{proof}
Diffeomorphism invariance gives zero total variation.  The on-shell bulk
fermion variation contributes the normal flux
\(T^F_{nb}\zeta^b\).  The boundary metric variation is
\(-\frac12\int\sqrt{-\gamma}S^{F,\partial}_{ab}
\mathcal L_\zeta\gamma^{ab}\); integration by parts gives
\(\int\sqrt{-\gamma}(D^aS^{F,\partial}_{ab})\zeta^b\).
The clock and material variations give
\(\mathcal F_T\zeta^bD_bT+
\mathcal F_I\zeta^bD_bX^I\).  Spinorial Lie derivatives give the two
Euler pairings.  Since \(\zeta\) is arbitrary, their coefficient is
\eqref{eq:v3v57-offshell-Ward}.  Setting the Euler expressions to zero
gives \eqref{eq:v3v57-onshell-Ward}.
\end{proof}

The material action has the complementary Noether identity.  Once
\(S_{\rm con}\) is included in the total boundary action, its clock and
material equations contain \(-\mathcal F_T\) and \(-\mathcal F_I\).

\begin{corollary}[Closure of the fermion--material Codazzi exchange]
\label{cor:v3v57-total-Ward}
Let \(S^{\rm mat}_{ab}\) be the stress of the clock/material/cold/static
boundary system and include \(S_{\rm con}\) before varying its equations.
On all boundary Euler equations,
\begin{equation}
 D^a\left(S^{\rm mat}_{ab}+S^{F,\partial}_{ab}\right)
 =T^{\rm other}_{nb}+T^F_{nb},
\label{eq:v3v57-total-Ward}
\end{equation}
where \(T^{\rm other}_{nb}\) denotes the nonfermionic bulk port flux.
Thus the tangential fermion momentum which did not vanish in V56 is paid
once by the reaction stress and material force of the same projector
action.
\end{corollary}

\begin{proof}
The material Ward identity contains the negative of the right side of
\eqref{eq:v3v57-onshell-Ward}, because the forces enter its Euler
equations from the same total boundary action.  Add the two identities.
The force terms cancel, leaving \eqref{eq:v3v57-total-Ward}.
\end{proof}

\subsection{Gauge Ward identity and current balance}
\label{sec:v3v57-gauge-Ward}

\begin{proposition}[The multiplier creates no anomalous classical gauge
flux]
\label{prop:v3v57-gauge-Ward}
If \(\lambda\) transforms in the same unitary representation as \(\psi\),
then \(S_{\rm con}\) is exactly gauge invariant.  Its gauge Noether
identity is
\begin{equation}
 (\mathcal E_\lambda,T^a\lambda)_{\mathbb R}
 +(\mathcal E_\psi,T^a\psi)_{\mathbb R}=0.
\label{eq:v3v57-gauge-Ward}
\end{equation}
On the boundary equations it supplies no normal gauge source, consistently
with \eqref{eq:v3v56-zero-currents}.
\end{proposition}

\begin{proof}
The spin projector commutes with every gauge generator and the Hermitian
pairing is invariant.  Vary \eqref{eq:v3v57-constraint-action} under an
arbitrary infinitesimal gauge parameter.  The coefficient of the parameter
is \eqref{eq:v3v57-gauge-Ward}.  On shell it vanishes, and the bulk normal
current is already zero by Corollary~\ref{cor:v3v56-zero-current}.
\end{proof}

\subsection{Energy and corner estimates with the reaction block}
\label{sec:v3v57-estimates}

\begin{proposition}[Reaction terms preserve the V49--V54 estimates]
\label{prop:v3v57-estimates}
At the zero-fermion simple background,
\(\lambda=0\) and the reaction Hessian has no linear off-diagonal block in
the constraint--corner normal family.  In a sufficiently small fermion
and frame neighborhood, its contribution is a tame order-zero perturbation
and preserves both the differentiated energy estimate and the strict V54
Schur gap.
\end{proposition}

\begin{proof}
Proposition~\ref{prop:v3v57-exact-domain} makes \(\lambda\) linear in
\(\psi\), so the reaction stress and forces are quadratic at
\(\psi=0\).  Their first variation there vanishes.  Lemma
\ref{lem:v3v57-reaction-tame} controls the differentiated terms by the
low coefficient norm times the high energy.  The strict finite-dimensional
Schur margin is open under such coefficient perturbations.
\end{proof}

\begin{firewall}[Classical Ward closure is not BRST closure]
\label{fw:v3v57-not-BRST}
V57 closes the classical diffeomorphism and gauge Ward identities for the
material-frame Weyl projector.  The multiplier is an algebraic reaction
field, not a Faddeev--Popov ghost, and its functional determinant has not
been defined.  Quantum gauge independence, anomaly cancellation with a
boundary, eta phases, and reflection positivity remain open.
\end{firewall}

\subsection{V57 conclusion and V58 target}
\label{sec:v3v57-conclusion}

Imposing the V56 projector through the local multiplier action makes its
reaction observable.  The multiplier is algebraic, introduces no new
characteristic, and its frame/metric variation supplies the exact boundary
stress and material force.  Tangential diffeomorphism invariance proves
that these terms convert the reflected fermion momentum flux into the
combined material boundary Ward identity.  Classical gauge invariance and
zero normal fermion current preserve the Gauss balance.

V58 should construct the ghost and gauge-fixing boundary complex for the
absolute bosonic domain of V55.  It must specify ghost, antighost, and
Nakanishi--Lautrup boundary conditions, prove that the BRST differential
maps the domain into itself, compute the flat Calderon/normal complex, and
test whether the fermion projector multiplier changes the classical BRST
charge.  Determinant or anomaly claims must wait until this domain-level
complex is exact.
\clearpage
\section{V58: the absolute BRST boundary complex and its flat normal family}
\label{sec:v3v58}

\subsection{Gauge fixing and BRST differential}
\label{sec:v3v58-BRST}

V55 proves the absolute classical gauge/Higgs domain, and V57 proves that
the material-frame fermion projector is classically gauge invariant.  The
next question is whether one gauge-fixed field domain is preserved by the
BRST differential.  Let \(c,\bar c,b\) be the ghost, antighost, and
Nakanishi--Lautrup fields.  With a compact gauge algebra and anti-Hermitian
action on matter, use
\begin{align}
 \mathsf sA_\mu&=D_\mu c,
 &\mathsf sc&=-\frac12[c,c],
\label{eq:v3v58-BRST-gauge}\
 \mathsf s\bar c&=b,
 &\mathsf sb&=0,
\label{eq:v3v58-BRST-doublet}\
 \mathsf sH&=-cH,
 &\mathsf s\psi&=-c\psi,
 &\mathsf s\lambda&=-c\lambda.
\label{eq:v3v58-BRST-matter}
\end{align}
The metric, clock, and material fields are gauge singlets.  The gauge-fixing
fermion is
\begin{equation}
 \Psi_{\rm gf}
 =\int_M\sqrt{-g}\,
  \left\langle\bar c,
  \nabla^\mu A_\mu+\frac\alpha2b\right\rangle d^4x,
 \qquad \alpha>0,
\label{eq:v3v58-gauge-fixing-fermion}
\end{equation}
and
\begin{equation}
 S_{\rm gf+gh}=\mathsf s\Psi_{\rm gf}.
\label{eq:v3v58-exact-action}
\end{equation}

The complete absolute BRST domain is
\begin{align}
 A_n&=0,
 &n^\mu F_{\mu a}&=0,
\label{eq:v3v58-absolute-A}\
 D_nc&=0,
 &D_n\bar c&=0,
 &D_nb&=0,
\label{eq:v3v58-Neumann-triple}\
 D_nH&=0,
 &P_-\psi&=0,
 &P_-\lambda&=\lambda.
\label{eq:v3v58-matter-domain}
\end{align}
The last multiplier convention is equivalent to the V57 choice
\(\lambda\in\operatorname{Ran}P_-\).

\begin{theorem}[The absolute domain is a BRST subcomplex]
\label{thm:v3v58-domain-invariant}
The differential \(\mathsf s\) maps
\eqref{eq:v3v58-absolute-A}--\eqref{eq:v3v58-matter-domain} into itself and
is nilpotent on that domain.
\end{theorem}

\begin{proof}
First,
\(\mathsf sA_n=D_nc=0\).  Curvature covariance gives
\[
 \mathsf s(n\lrcorner F)=[n\lrcorner F,c]=0.
\]
The graded Leibniz rule and the Jacobi identity give
\begin{equation}
 \mathsf s(D_nc)=-[D_nc,c],
\label{eq:v3v58-sDnc}
\end{equation}
up to the equivalent sign fixed by the anti-Hermitian convention; it
vanishes when \(D_nc=0\).  Equations
\eqref{eq:v3v58-BRST-doublet} give
\(\mathsf s(D_n\bar c)=D_nb\) and
\(\mathsf s(D_nb)=0\) on \(A_n=0\), so the antighost--auxiliary domain is
preserved.

For the Higgs field,
\[
 \mathsf s(D_nH)=-cD_nH-(D_nc)H=0.
\]
The spin projector is a gauge singlet and commutes with the action on every
Standard Model representation.  Hence
\(\mathsf s(P_-\psi)=-cP_-\psi=0\), and the same calculation preserves
\(P_-\lambda=\lambda\).  Bulk nilpotency follows from the Jacobi identity;
the calculations above show that it descends to the boundary domain.
\end{proof}

\begin{proposition}[Variational closure of the gauge-fixed action]
\label{prop:v3v58-variational}
The Yang--Mills--Higgs action plus
\eqref{eq:v3v58-exact-action} has no uncancelled integration-by-parts term
on the domain
\eqref{eq:v3v58-absolute-A}--\eqref{eq:v3v58-Neumann-triple}.  The
auxiliary equation is algebraic in the bulk; its Neumann condition is the
boundary compatibility condition required by
\(\mathsf s(D_n\bar c)=D_nb\).
\end{proposition}

\begin{proof}
The invariant action boundary term vanishes by Proposition
\ref{prop:v3v55-variational}.  In the gauge-fixing term, variation of
\(A\) produces a normal contribution proportional to
\(b\,\delta A_n\), and \(A_n=\delta A_n=0\).  Write the ghost quadratic
term in its symmetric first-derivative form.  Its boundary variations are
proportional to \(D_nc\) and \(D_n\bar c\), hence vanish.  The field \(b\)
has no differentiated bulk equation; imposing \(D_nb=0\) only selects the
BRST-compatible boundary trace of its algebraic solution.
\end{proof}

\subsection{Linear boundary complex}
\label{sec:v3v58-linear-complex}

At the trivial connection and matter background, the linear BRST complex
is
\begin{equation}
 c\xrightarrow{\ d\ }A
 \xrightarrow{\ d\ }F,
 \qquad
 \bar c\xrightarrow{\ \mathsf s\ }b\longrightarrow0.
\label{eq:v3v58-linear-complex}
\end{equation}
Let
\begin{equation}
 p=(s,i\xi,-\kappa),
 \qquad
 \kappa=\sqrt{s^2+|\xi|^2},
 \qquad \operatorname{Re}\kappa>0
\label{eq:v3v58-Lorentz-covector}
\end{equation}
be a decaying Lorentzian stable covector.  The symbol maps are
\begin{equation}
 d_0C=pC,
 \qquad
 d_1a=p\wedge a.
\label{eq:v3v58-symbol-complex}
\end{equation}

\begin{lemma}[Algebraic exactness away from the zero covector]
\label{lem:v3v58-symbol-exact}
For every nonzero covector \(p\),
\begin{equation}
 \ker d_1=\operatorname{ran}d_0.
\label{eq:v3v58-symbol-exact}
\end{equation}
\end{lemma}

\begin{proof}
Choose a vector \(v\) with \(p(v)=1\).  If \(p\wedge a=0\), contract with
\(v\):
\[
 a-p\,a(v)=0.
\]
Thus \(a=pC\) with \(C=a(v)\).  Conversely
\(p\wedge pC=0\).
\end{proof}

The absolute boundary operator on a pure gauge mode is
\begin{equation}
 A_n=-\kappa C,
 \qquad
 F_{na}=0.
\label{eq:v3v58-pure-gauge-boundary}
\end{equation}
The ghost Neumann operator is the same scalar factor
\begin{equation}
 D_nc=\kappa C
\label{eq:v3v58-ghost-normal}
\end{equation}
for the outward convention of V54.

\begin{proposition}[Lorentzian stable-domain exactness off glancing]
\label{prop:v3v58-Lorentz-exact}
If \(\kappa\ne0\), the intersection of the pure-gauge stable Calderon
space with the absolute gauge domain is zero, exactly when the stable ghost
mode satisfying Neumann data is zero.  If \(\kappa=0\), the surviving
pure tangential gauge mode is the BRST image of the surviving residual
Neumann ghost mode.
\end{proposition}

\begin{proof}
Equations \eqref{eq:v3v58-pure-gauge-boundary} and
\eqref{eq:v3v58-ghost-normal} both reduce to \(\kappa C=0\).  This forces
\(C=0\) off glancing.  At glancing neither equation removes \(C\), and
\eqref{eq:v3v58-symbol-complex} identifies the gauge amplitude with
\(d_0C\).  Thus the residual mode occurs in the complex rather than as an
unpaired gauge amplitude.
\end{proof}

\subsection{Euclidean normal family}
\label{sec:v3v58-Euclidean}

For determinant and elliptic-domain questions, Wick rotate the tangential
time frequency.  Let \(p_E=(i\omega,i\xi,-\rho)\), where
\begin{equation}
 \rho=(\omega^2+|\xi|^2)^{1/2}>0.
\label{eq:v3v58-Euclidean-root}
\end{equation}
A decaying one-form amplitude satisfies
\begin{equation}
 A_n=0,
 \qquad
 F_{na}=-\rho A_a-i p_aA_n=0,
\label{eq:v3v58-Euclidean-absolute}
\end{equation}
and a ghost amplitude satisfies
\begin{equation}
 \partial_nc=\rho C=0.
\label{eq:v3v58-Euclidean-ghost}
\end{equation}

\begin{theorem}[Strong ellipticity of the nonzero Euclidean normal family]
\label{thm:v3v58-Euclidean-elliptic}
For every \((\omega,\xi)\ne0\), the absolute one-form and Neumann ghost
normal families have zero decaying kernel.  The linear BRST symbol complex
is exact on their Calderon spaces.  The only unremoved scalar ghost sector
is the joint zero mode \((\omega,\xi)=(0,0)\).
\end{theorem}

\begin{proof}
Equation \eqref{eq:v3v58-Euclidean-absolute} first sets \(A_n=0\), then
sets every tangential component to zero because \(\rho>0\).  Equation
\eqref{eq:v3v58-Euclidean-ghost} sets \(C=0\) for the same reason.  Exactness
of the interior symbol is Lemma~\ref{lem:v3v58-symbol-exact}; the boundary
calculation shows that restriction introduces no nonzero-frequency
cohomology.  At joint zero frequency, \(\rho=0\), so a covariantly constant
ghost remains.
\end{proof}

\begin{corollary}[Removal of the global gauge zero mode]
\label{cor:v3v58-based-gauge}
Choose either a based gauge group or the orthogonality condition
\begin{equation}
 \int_M c\,dV=0
\label{eq:v3v58-ghost-mean-zero}
\end{equation}
in every covariantly constant unbroken gauge direction.  Then the
Euclidean Neumann ghost Laplacian is invertible on the resulting space,
and no zero eigenvalue is silently included in its determinant.
\end{corollary}

\begin{proof}
On a connected compact regulated region, the kernel of the scalar Neumann
Laplacian consists of covariantly constant sections.  A based condition or
orthogonal projection removes exactly that finite-dimensional kernel.
\end{proof}

\subsection{Constraint propagation and the BRST charge}
\label{sec:v3v58-charge}

At the linearized level, the gauge-fixing constraint transforms as
\begin{equation}
 \mathsf s(\nabla^\mu A_\mu)=\Box c.
\label{eq:v3v58-gauge-slice-transform}
\end{equation}
Thus the ghost equation makes the gauge slice BRST invariant on shell.
The nonlinear formula adds the standard adjoint lower terms and remains
proportional to the ghost Euler equation.

\begin{theorem}[No BRST-charge leakage through the absolute boundary]
\label{thm:v3v58-charge-conserved}
Let \(J_{\rm B}^\mu\) be the classical Noether current for
\eqref{eq:v3v58-BRST-gauge}--\eqref{eq:v3v58-exact-action}, including the
matter and projector-multiplier sectors.  On the complete boundary domain,
\begin{equation}
 n_\mu J_{\rm B}^\mu=0.
\label{eq:v3v58-no-BRST-flux}
\end{equation}
Consequently the regulated classical BRST charge
\begin{equation}
 Q_{\rm B}(t)=\int_{S_t}u_\mu J_{\rm B}^\mu,dV
\label{eq:v3v58-BRST-charge}
\end{equation}
is time independent.
\end{theorem}

\begin{proof}
Promote the constant odd BRST parameter to a test function.  The normal
Noether flux is the sum of canonical normal momenta paired with the BRST
variations.  The Yang--Mills term contains
\(n\lrcorner F\) and vanishes by
\eqref{eq:v3v58-absolute-A}.  The gauge-fixing normal term contains
\(\mathsf sA_n=D_nc\), while the ghost terms contain \(D_nc\) or
\(D_n\bar c\); all vanish by
\eqref{eq:v3v58-Neumann-triple}.  The Higgs contribution contains
\(D_nH\), and the fermion contribution is the normal gauge current of
Corollary~\ref{cor:v3v56-zero-current}.  The multiplier action has no
normal derivative and is exactly gauge, hence BRST, invariant by
Proposition~\ref{prop:v3v57-gauge-Ward}; it contributes no normal current.
This proves \eqref{eq:v3v58-no-BRST-flux}.  Integrate current conservation
over the spacetime slab to obtain constancy of \(Q_{\rm B}\).
\end{proof}

\begin{proposition}[The antighost--auxiliary sector is contractible]
\label{prop:v3v58-doublet-contractible}
On the common Neumann domain, the pair \((\bar c,b)\) is a BRST doublet and
has zero cohomology away from constants fixed by
\eqref{eq:v3v58-ghost-mean-zero}.
\end{proposition}

\begin{proof}
Let \(N\) count total polynomial degree in \((\bar c,b)\), and define the
homotopy \(h=\bar c\,\partial/\partial b\), with graded left derivatives.
On positive doublet degree,
\(\mathsf sh+h\mathsf s=N\).  Hence a closed element of degree
\(N>0\) is exact after division by \(N\).  The Neumann domain is preserved
by both maps because \(D_n\bar c=D_nb=0\).
\end{proof}

\begin{firewall}[Domain exactness is not a determinant theorem]
\label{fw:v3v58-no-determinant}
V58 proves classical BRST invariance of the complete boundary domain,
nonzero-frequency exactness of its flat normal complex, strong ellipticity
of the Euclidean absolute/Neumann pair, removal of the residual global
gauge mode, and conservation of the regulated classical BRST charge.  It
does not define the interacting ghost determinant, prove cancellation of
its phase, establish a renormalized Slavnov--Taylor identity, or exclude a
boundary gauge anomaly.
\end{firewall}

\subsection{V58 conclusion and V59 target}
\label{sec:v3v58-conclusion}

The absolute gauge/Higgs domain extends to an exact classical BRST domain
by covariant Neumann conditions on \(c,\bar c,b\).  The fermion projector
and its reaction multiplier are stable under the same differential.  The
flat Lorentzian complex pairs its glancing residual gauge modes with
Neumann ghosts, while the Euclidean normal family is strongly elliptic at
every nonzero boundary frequency.  Only the finite-dimensional global
gauge algebra remains, and it is removed explicitly by a based or
mean-zero condition.

V59 should attack the first quantum gate at finite regulator.  It should
define the gauge-fixed Gaussian Hessian on the absolute/Neumann domain,
separate its zero modes, and prove the one-loop BRST cancellation for
nonzero modes by an isospectral exact-complex argument.  Chiral fermion
phases and local boundary anomaly polynomials must be kept separate from
that modulus calculation; failure of their cancellation would block the
full Standard Model continuum even if the Gaussian determinants pair.
\clearpage
\section{V59: finite-cutoff one-loop BRST pairing and the chiral phase firewall}
\label{sec:v3v59}

\subsection{Regulated Euclidean Hessians}
\label{sec:v3v59-Hessians}

V58 establishes the domain-level BRST complex.  This note asks only for
the Gaussian integral at a finite, complex-compatible spectral cutoff.
Let \(M_E\) be a compact Euclidean regulator with boundary and use the
absolute one-form domain
\begin{equation}
 \iota_nA=0,
 \qquad \iota_ndA=0,
\label{eq:v3v59-absolute-domain}
\end{equation}
and the scalar Neumann domain
\begin{equation}
 \partial_nf=0.
\label{eq:v3v59-Neumann-domain}
\end{equation}
Remove the covariantly constant scalar modes as in
Corollary~\ref{cor:v3v58-based-gauge}.  At the trivial flat connection and
zero Higgs background, put
\begin{align}
 \Delta_0&=\delta d,
\label{eq:v3v59-Delta0}\
 \Delta_1&=d\delta+\delta d,
\label{eq:v3v59-Delta1}\
 P_\alpha&=\delta d+\alpha^{-1}d\delta,
 \qquad \alpha>0.
\label{eq:v3v59-Palpha}
\end{align}
The operators act componentwise in the Standard Model gauge algebra.

\subsection{Exact and coexact spectra}
\label{sec:v3v59-Hodge}

Let \(f_j\) be an orthonormal nonzero Neumann eigenfunction,
\begin{equation}
 \Delta_0f_j=\lambda_jf_j,
 \qquad \lambda_j>0,
\label{eq:v3v59-scalar-eigen}
\end{equation}
and define
\begin{equation}
 e_j=\lambda_j^{-1/2}df_j.
\label{eq:v3v59-exact-one-form}
\end{equation}

\begin{lemma}[Isospectral exact-mode pairing]
\label{lem:v3v59-isospectral}
The form \(e_j\) has unit norm, belongs to the absolute domain, and obeys
\begin{equation}
 \Delta_1e_j=\lambda_je_j,
 \qquad
 P_\alpha e_j=\alpha^{-1}\lambda_je_j.
\label{eq:v3v59-exact-eigenvalues}
\end{equation}
Conversely, every nonzero exact one-form in the absolute domain is obtained
from a nonzero Neumann scalar in this way.
\end{lemma}

\begin{proof}
Neumann data give \(\iota_ndf_j=0\), while
\(d(df_j)=0\), so \(e_j\) satisfies
\eqref{eq:v3v59-absolute-domain}.  Integration by parts gives
\[
 \|df_j\|^2=\langle f_j,\Delta_0f_j\rangle=\lambda_j,
\]
and hence \(\|e_j\|=1\).  Since \(d\) commutes with the Hodge Laplacian,
\(\Delta_1df_j=d\Delta_0f_j\).  Also
\(\delta d(df_j)=0\) and
\(d\delta(df_j)=d\Delta_0f_j\), proving
\eqref{eq:v3v59-exact-eigenvalues}.  If \(A=df\) is absolute, its normal
component is \(\partial_nf=0\); removing the constant part of \(f\) gives
the claimed converse by scalar spectral expansion.
\end{proof}

Let \(\mathcal H^1_{\rm abs}\) be the finite-dimensional space of absolute
harmonic one-forms, and let
\begin{equation}
 \mathscr T
 =\ker\delta\cap(\mathcal H^1_{\rm abs})^\perp
\label{eq:v3v59-transverse-space}
\end{equation}
be the nonharmonic co-closed space.

\begin{proposition}[Absolute Hodge splitting]
\label{prop:v3v59-Hodge-splitting}
There is an orthogonal decomposition
\begin{equation}
 L^2\Omega^1_{\rm abs}
 =d\bigl(H^1_N/\ker\Delta_0\bigr)
  \oplus\mathscr T
  \oplus\mathcal H^1_{\rm abs}.
\label{eq:v3v59-Hodge-splitting}
\end{equation}
Both \(\Delta_1\) and \(P_\alpha\) preserve the three summands;
\(P_\alpha=\Delta_1\) on \(\mathscr T\).
\end{proposition}

\begin{proof}
This is the Hodge decomposition with absolute boundary conditions.  It can
be proved directly by minimizing \(\|A-df\|^2\) over Neumann scalars,
which gives \(\delta(A-df)=0\), and then separating the finite harmonic
kernel.  The commutation relations in Lemma
\ref{lem:v3v59-isospectral} show invariance.  On \(\ker\delta\), the term
\(d\delta\) vanishes, so \(P_\alpha=\delta d=\Delta_1\).
\end{proof}

\subsection{A BRST-compatible cutoff}
\label{sec:v3v59-cutoff}

For \(\Lambda<\infty\), retain:
\begin{enumerate}
 \item every nonzero scalar Neumann eigenfunction with
       \(\lambda_j\le\Lambda^2\), simultaneously in the
       \(c,\bar c,b\) sectors;
 \item its exact partner \(e_j\) in the gauge sector; and
 \item every transverse absolute one-form eigenmode with eigenvalue at most
       \(\Lambda^2\).
\end{enumerate}
Harmonic one-forms are recorded separately as a finite moduli factor.  Let
\(N_0(\Lambda)\) be the number of retained nonzero scalar modes, including
the gauge-algebra multiplicity.

\begin{proposition}[The cutoff is a finite BRST complex]
\label{prop:v3v59-cutoff-complex}
The linearized transformations
\begin{equation}
 \mathsf sA=dc,
 \qquad \mathsf sc=0,
 \qquad \mathsf s\bar c=b,
 \qquad \mathsf sb=0
\label{eq:v3v59-linear-BRST}
\end{equation}
preserve the cutoff space and the boundary domain.
\end{proposition}

\begin{proof}
Lemma~\ref{lem:v3v59-isospectral} pairs a scalar eigenmode with an exact
one-form of the same \(\Delta\)-eigenvalue, so the first map stays inside
the cutoff.  The antighost and auxiliary field use the identical scalar
basis.  Domain preservation is Theorem
\ref{thm:v3v58-domain-invariant}.
\end{proof}

\subsection{Exact finite determinant identity}
\label{sec:v3v59-determinant}

Let \(\Delta_{0,\Lambda}\) be the positive scalar matrix and let
\(\Delta_{T,\Lambda}\) be the restriction of \(\Delta_1\) to retained
transverse modes.  Excluding the separately recorded harmonic sector,
Lemma~\ref{lem:v3v59-isospectral} gives
\begin{equation}
 \det P_{\alpha,\Lambda}
 =\alpha^{-N_0(\Lambda)}
  \det\Delta_{0,\Lambda}\,
  \det\Delta_{T,\Lambda}.
\label{eq:v3v59-det-factorization}
\end{equation}

Normalize the finite auxiliary-field Gaussian so that its pure
gauge-parameter factor is
\begin{equation}
 \mathcal N_{b,\Lambda}(\alpha)=\alpha^{-N_0(\Lambda)/2}.
\label{eq:v3v59-b-normalization}
\end{equation}
This is the normalization obtained by using the same Lebesgue measure
before and after completing the square in \(b\).

\begin{theorem}[Finite-cutoff gauge-parameter cancellation]
\label{thm:v3v59-one-loop-modulus}
The nonzero-mode gauge/ghost/auxiliary Gaussian factor is
\begin{align}
 Z_{\rm ggh}^{(\Lambda)}(\alpha)
 &:={}
 \mathcal N_{b,\Lambda}(\alpha)
 \det\Delta_{0,\Lambda}
 \bigl(\det P_{\alpha,\Lambda}\bigr)^{-1/2}
\notag\\
 &=
 \bigl(\det\Delta_{0,\Lambda}\bigr)^{1/2}
 \bigl(\det\Delta_{T,\Lambda}\bigr)^{-1/2}.
\label{eq:v3v59-one-loop-modulus}
\end{align}
It is exactly independent of \(\alpha\) at every finite cutoff.
\end{theorem}

\begin{proof}
Substitute \eqref{eq:v3v59-det-factorization} and
\eqref{eq:v3v59-b-normalization}.  The powers of \(\alpha\) cancel, and
the remaining scalar and transverse determinants give
\eqref{eq:v3v59-one-loop-modulus}.
\end{proof}

\begin{assessment}[What cancels and what remains]
\label{ass:v3v59-cancellation-scope}
The ghost spectrum is exactly paired with the exact one-form spectrum, but
the final scalar determinant occurs with power \(+1/2\).  It must not be
deleted by saying that all unphysical modes cancel.  Together with the
three Euclidean co-closed vector components, this ratio has the two
physical gauge-polarization degree count.  Equation
\eqref{eq:v3v59-one-loop-modulus} is the correct Gaussian modulus.
\end{assessment}

\subsection{Finite-dimensional Slavnov identity}
\label{sec:v3v59-Slavnov}

Let \(d\mu_\Lambda\) be Lebesgue--Berezin measure on the finite complex of
Proposition~\ref{prop:v3v59-cutoff-complex}, with the positive auxiliary
contour fixed by the Gaussian completion.

\begin{theorem}[One-loop BRST-exact insertions vanish]
\label{thm:v3v59-finite-Ward}
For every polynomial \(K\) on the finite cutoff space,
\begin{equation}
 \int d\mu_\Lambda\,
  \mathsf s\left(K e^{-S_\Lambda^{(2)}}\right)=0.
\label{eq:v3v59-finite-Ward}
\end{equation}
Consequently expectation values of \(\mathsf sK\) vanish and the Gaussian
partition function is invariant under a finite variation of the gauge-
fixing fermion which preserves the cutoff complex and convergence contour.
\end{theorem}

\begin{proof}
The transformation \eqref{eq:v3v59-linear-BRST} is triangular and
nilpotent, so its finite-dimensional Berezinian is one.  The quadratic
action is BRST invariant because it is the sum of the invariant Hessian
and a BRST-exact gauge-fixing Hessian.  Berezin integration of a total odd
derivative and ordinary integration of a convergent total derivative both
vanish.  Expanding \eqref{eq:v3v59-finite-Ward} and using
\(\mathsf sS_\Lambda^{(2)}=0\) gives
\(\langle\mathsf sK\rangle_\Lambda=0\).  Replacing \(K\) by the variation
of the gauge-fixing fermion gives gauge-fixing independence.
\end{proof}

\begin{firewall}[Why this theorem is only one loop]
\label{fw:v3v59-one-loop-only}
The nonlinear transformation \(\mathsf sc=-[c,c]/2\) multiplies cutoff
modes and generally leaves a sharp spectral subspace.  V59 proves the
linearized Gaussian Ward identity because the Hodge-compatible cutoff is
closed under \(d\), not an interacting finite-cutoff Slavnov identity.  A
nonlinear regulator must preserve the full master equation or carry
explicit symmetry-restoring counterterms.
\end{firewall}

\subsection{Bulk Standard Model anomaly arithmetic}
\label{sec:v3v59-anomaly-arithmetic}

Write all fermions as left-handed Weyl fields.  Per generation the list is
\begin{equation}
 Q_L:(\mathbf3,\mathbf2)_{1/6},\quad
 u_R^c:(\overline{\mathbf3},\mathbf1)_{-2/3},\quad
 d_R^c:(\overline{\mathbf3},\mathbf1)_{1/3},\quad
 L_L:(\mathbf1,\mathbf2)_{-1/2},\quad
 e_R^c:(\mathbf1,\mathbf1)_1.
\label{eq:v3v59-left-list}
\end{equation}

\begin{theorem}[Cancellation of the perturbative bulk gauge anomaly]
\label{thm:v3v59-bulk-anomaly}
For the representations \eqref{eq:v3v59-left-list}, the coefficients of
\(SU(3)^3\), \(SU(3)^2U(1)\), \(SU(2)^2U(1)\), \(U(1)^3\), and the mixed
gravitational--\(U(1)\) anomaly vanish generation by generation.  The
number of \(SU(2)\) doublets is four per generation, so the mod-two Witten
anomaly also vanishes.
\end{theorem}

\begin{proof}
For \(SU(3)^3\), the two fundamental copies in \(Q_L\) cancel the two
antifundamentals \(u_R^c,d_R^c\).  With
\(T(\mathbf3)=T(\mathbf2)=1/2\),
\begin{align*}
 SU(3)^2U(1):&\quad
 2\frac16\frac12-\frac23\frac12+\frac13\frac12=0,\\
 SU(2)^2U(1):&\quad
 3\frac16\frac12-\frac12\frac12=0.
\end{align*}
The cubic hypercharge sum is
\begin{equation}
 6\left(\frac16\right)^3
 +3\left(-\frac23\right)^3
 +3\left(\frac13\right)^3
 +2\left(-\frac12\right)^3+1=0,
\label{eq:v3v59-Y3-sum}
\end{equation}
and the linear sum is
\begin{equation}
 6\frac16+3\left(-\frac23\right)
 +3\frac13+2\left(-\frac12\right)+1=0.
\label{eq:v3v59-Y-sum}
\end{equation}
There are three color copies of \(Q_L\) and one \(L_L\), hence four
\(SU(2)\) doublets.
\end{proof}

\subsection{The Weyl determinant phase remains separate}
\label{sec:v3v59-phase}

At finite cutoff, integrating the V57 multiplier imposes the chiral
projector and leaves the Weyl operator on that domain.  Its modulus is
controlled by the positive squared operator, but its phase depends on
spectral asymmetry of the boundary problem.

\begin{firewall}[Bulk anomaly cancellation does not determine the boundary
eta phase]
\label{fw:v3v59-eta}
Theorem~\ref{thm:v3v59-bulk-anomaly} removes the local four-dimensional
perturbative anomaly polynomial and the \(SU(2)\) mod-two obstruction.  It
does not prove that the material-frame local projector has zero eta phase,
that all boundary Chern--Simons variations cancel, or that a gauge- and
diffeomorphism-invariant phase convention exists over the full background
configuration space.  Those are global properties of the determinant
line, not consequences of the arithmetic sums alone.
\end{firewall}

\begin{firewall}[No cutoff removal]
\label{fw:v3v59-no-limit}
Every determinant and Ward identity in V59 is finite dimensional.  No
uniform \(\Lambda\to\infty\) estimate, counterterm construction,
renormalized measure, or Osterwalder--Schrader limit has been proved.
Gauge-field harmonic modes and moduli are recorded separately rather than
included in a primed determinant by notation alone.
\end{firewall}

\subsection{V59 conclusion and V60 target}
\label{sec:v3v59-conclusion}

The Hodge-compatible regulator makes the nonzero ghost and exact one-form
spectra identical.  It gives the exact finite determinant factorization
\eqref{eq:v3v59-det-factorization}, cancels the gauge parameter, and proves
the Gaussian Slavnov identity by finite-dimensional BRST integration.  The
remaining modulus has the correct transverse-vector/scalar ratio.  The
Standard Model bulk anomaly coefficients and the Witten anomaly cancel,
but the boundary determinant-line phase remains open.

V60 is a compulsory companion audit.  After recording the current NS and
YM states, it should construct the diffeomorphism ghost boundary complex
for the harmonic Einstein junction and combine it with V58's internal
gauge complex.  The target is an exact finite-regulator product complex
whose zero modes are the genuine joint isometries and gauge stabilizers.
Only after that construction should the boundary eta line be attacked.
\clearpage
\section{V60: compulsory audit and the joint diffeomorphism--gauge ghost complex}
\label{sec:v3v60}

\subsection{Compulsory companion audit}
\label{sec:v3v60-audit}

The V60 audit found the following active companion states:
\begin{center}
\begin{tabular}{p{0.54\textwidth}r p{0.28\textwidth}}
\hline
file & bytes & SHA256\\
\hline
\texttt{Renewal\_NS\_Stability\_Research\_Notes\_V31.tex}
&887537&\texttt{F1121B62238AD5491BB7CF03635EA9934942EDF573ED8FAAA9EE79E1D38A6696}\\
\texttt{Renewal\_Yang\_Mills\_Mass\_Gap\_Research\_Notes\_V94.tex}
&1574286&\texttt{D4C894AD796376A88E1D87FAA764D73D74D3D7A5FB4C9A5AE0DEF4481C2015B1}\\
\hline
\end{tabular}
\end{center}

NS remains at V31.  Its exact balanced-sector nullspace again forbids
using a signed or constraint quantity as a norm for the full field.  In
V60 this means that exactness of the ghost complex cannot control the
transverse-traceless graviton Hessian; that physical operator keeps its own
energy and Euclidean-domain tests.

YM V94 replaces a false fine-pinching route by an exact unsplit dual
carrier.  Its geometric-mean estimate works only while the two bank
occurrences and their separator remain assembled.  The type-correct lesson
for V60 is to retain the semidirect diffeomorphism--internal-gauge
Faddeev--Popov operator as one triangular block.  On backgrounds with
stabilizers, its kernel is a joint kernel and may not be obtained by
subtracting two independent zero-mode counts.

\begin{audit}[Direction selected at V60]
\label{audit:v3v60-direction}
V60 constructs the full classical ghost domain and its assembled normal
operator.  It proves the triangular finite-cutoff determinant identity only
after the joint zero space is projected.  It makes no inference about the
physical graviton determinant from that ghost identity.
\end{audit}

\subsection{Boundary-preserving diffeomorphism ghosts}
\label{sec:v3v60-diff-ghost}

Let \(\xi^\mu,\bar\xi_\mu,B_\mu\) be the diffeomorphism ghost, antighost,
and auxiliary vector.  The gravitational BRST differential is
\begin{align}
 \mathsf sg_{\mu\nu}&=\mathcal L_\xi g_{\mu\nu},
 &
 \mathsf s\xi&=\frac12[\xi,\xi],
\label{eq:v3v60-sg}\
 \mathsf s\bar\xi&=B,
 &
 \mathsf sB&=0.
\label{eq:v3v60-diff-doublet}
\end{align}
On internal fields, \(\mathsf s\) also contains their covariant Lie
derivative.  A fixed timelike boundary requires
\begin{equation}
 \iota_n\xi^\flat=0.
\label{eq:v3v60-xi-normal}
\end{equation}
Choose the absolute vector boundary domain
\begin{equation}
 \iota_n\xi^\flat=0,
 \qquad
 \iota_nd\xi^\flat=0,
\label{eq:v3v60-xi-absolute}
\end{equation}
and impose the same two conditions on \(\bar\xi\) and \(B\).  In a collar,
the tangential part of the second condition is
\begin{equation}
 \nabla_n\xi_a+K_a{}^b\xi_b=0
\label{eq:v3v60-xi-Robin}
\end{equation}
for the second-fundamental-form convention induced by
\eqref{eq:v3v60-xi-absolute}.

Let
\begin{equation}
 C_\mu(g)=g^{\alpha\beta}
 \left(\Gamma_{\alpha\beta\mu}
       -\widehat\Gamma_{\alpha\beta\mu}\right)
\label{eq:v3v60-harmonic-gauge}
\end{equation}
be generalized harmonic gauge.  Its linear gauge variation is
\begin{equation}
 \delta_\xi C_\mu
 =\nabla^\alpha\nabla_\alpha\xi_\mu
  +R_\mu{}^\nu\xi_\nu
 =:\mathcal M_{\rm d}\xi_\mu.
\label{eq:v3v60-diff-FP}
\end{equation}
Use the gauge-fixing fermion
\begin{equation}
 \Psi_{\rm d}
 =\int_M\sqrt{-g}\,
  \bar\xi^\mu\left(C_\mu+\frac{\alpha_{\rm d}}2B_\mu\right)d^4x.
\label{eq:v3v60-diff-gf}
\end{equation}

\begin{proposition}[The geometric boundary domain is BRST invariant]
\label{prop:v3v60-geometric-domain}
Suppose every boundary field, including \((T,X,\Phi,\Psi_0)\), is
transformed by the same tangential diffeomorphism ghost.  Then the zero set
of the complete Brown--York/material junction and the fermion projector
constraint is preserved by \(\mathsf s\) whenever
\eqref{eq:v3v60-xi-normal} holds.  On the ghost equation
\(\mathcal M_{\rm d}\xi=0\), the harmonic boundary rows are preserved as
well.
\end{proposition}

\begin{proof}
The complete junction is an equality of natural tangential tensors.  Its
BRST variation on a boundary-preserving ghost is its tangential Lie
derivative, so the variation vanishes on its zero set.  The static and cold
trace relations are natural bundle equations and transform identically.
The fermion projector is built from \((g,T,X)\); hence
\(\mathsf s(P_-\psi)=\mathcal L_\xi(P_-\psi)\) plus the already closed
internal-gauge term.  Finally \eqref{eq:v3v60-diff-FP} shows that the
variation of \(C_\mu=0\) vanishes on the diffeomorphism-ghost equation.
\end{proof}

\begin{proposition}[Absolute diffeomorphism-ghost subcomplex]
\label{prop:v3v60-absolute-subcomplex}
The domain \eqref{eq:v3v60-xi-absolute}, together with the same absolute
domain for \((\bar\xi,B)\), is preserved by the linearized BRST differential
at a background satisfying the geometric boundary equations.  The
nonlinear corrections are sums of Lie brackets with fields satisfying the
same tangency condition and hence remain tangent; preservation of the
second condition is the covariant Lie derivative of
\(\iota_nd\xi^\flat=0\).
\end{proposition}

\begin{proof}
At the linearized level \(\mathsf s\xi=0\) and the doublet map sends
\(\bar\xi\) to \(B\) in the identical domain.  Nonlinearly, vector fields
tangent to a fixed hypersurface form a Lie algebra, so
\([\xi,\xi]^n=0\).  Naturality of exterior differentiation and pullback to
the boundary gives
\(\mathsf s(\iota_nd\xi^\flat)\) as a sum of Lie derivatives and brackets
of quantities which vanish on the domain.  The same argument applies to
the doublet.
\end{proof}

\subsection{Flat diffeomorphism-ghost normal family}
\label{sec:v3v60-flat-normal}

At a flat boundary, Wick rotate the tangential time frequency and put
\begin{equation}
 \rho=(\omega^2+|\eta|^2)^{1/2}>0.
\label{eq:v3v60-rho}
\end{equation}
A decaying vector ghost is
\(\xi_\mu=v_\mu e^{-\rho x+i\omega\tau+i\eta\cdot y}\).  Its absolute
conditions reduce to
\begin{equation}
 v_n=0,
 \qquad -\rho v_a-i p_av_n=0.
\label{eq:v3v60-vector-normal-family}
\end{equation}

\begin{theorem}[Strong ellipticity of the diffeomorphism-ghost normal
family]
\label{thm:v3v60-vector-elliptic}
For every nonzero Euclidean tangential covector, the decaying kernel of
\eqref{eq:v3v60-vector-normal-family} is zero.  The estimate
\begin{equation}
 |v|\le\rho^{-1}|\mathcal B_{\rm abs}v|
\label{eq:v3v60-vector-normal-bound}
\end{equation}
holds after the Dirichlet normal component and three Robin tangential
components are placed in homogeneous order-one graph weights.
\end{theorem}

\begin{proof}
The first row gives \(v_n=0\).  The remaining three rows give
\(-\rho v_a=0\), hence \(v_a=0\).  Solving the same triangular equations
with inhomogeneous rows yields \eqref{eq:v3v60-vector-normal-bound} in the
standard mixed boundary weights.
\end{proof}

For a ghost amplitude \(v\), the pure metric-gauge symbol is
\begin{equation}
 (\mathcal Kp(v))_{\mu\nu}
 =p_\mu v_\nu+p_\nu v_\mu.
\label{eq:v3v60-Killing-symbol}
\end{equation}

\begin{lemma}[No nonzero symbol stabilizer]
\label{lem:v3v60-Killing-injective}
For \(p\ne0\), the map \(\mathcal Kp\) in
\eqref{eq:v3v60-Killing-symbol} is injective.
\end{lemma}

\begin{proof}
Choose an index \(r\) with \(p_r\ne0\).  The \((r,r)\) component gives
\(2p_rv_r=0\).  The \((r,\mu)\) component then gives
\(p_rv_\mu=0\) for every \(\mu\).
\end{proof}

At joint zero frequency, the kernel of \(\mathcal M_{\rm d}\) with
absolute data contains boundary-preserving Killing fields.  Assume the
standard no-KID condition that the energy identity for
\(\mathcal M_{\rm d}\) has no further kernel.  Denote this finite space by
\begin{equation}
 \mathfrak{isom}_\partial(g)
 =\{\xi:\mathcal L_\xi g=0,\ \xi^n|_\partial=0\}.
\label{eq:v3v60-boundary-isometries}
\end{equation}
It is removed by a based diffeomorphism group or orthogonal projection.

\subsection{The semidirect product complex}
\label{sec:v3v60-semidirect}

Let \(c\) be the internal ghost of V58.  On a gauge connection, the
combined infinitesimal action is
\begin{equation}
 \mathcal D_0(\xi,c)
 =\left(
  \mathcal L_\xi g,
  \mathcal L_\xi A+D_Ac
 \right).
\label{eq:v3v60-combined-action}
\end{equation}
Because internal gauge transformations do not change the metric gauge
function, while diffeomorphisms change the Lorenz gauge function when
\(A\ne0\), the combined Faddeev--Popov operator has the exact triangular
form
\begin{equation}
 \mathcal M_{\rm joint}
 =\begin{pmatrix}
   \mathcal M_{\rm d}&0\\
   \mathcal K_A&\mathcal M_{\rm g}
  \end{pmatrix}.
\label{eq:v3v60-joint-FP}
\end{equation}
Here \(\mathcal M_{\rm g}=\nabla^\mu D_\mu\) and
\(\mathcal K_A\xi\) is the variation of the internal gauge-fixing function
under \(\mathcal L_\xi A\).

\begin{theorem}[Joint boundary BRST complex]
\label{thm:v3v60-joint-complex}
Equip \(\xi,\bar\xi,B\) with the absolute vector domain and
\(c,\bar c,b\) with the Neumann scalar domain of V58.  With the semidirect
BRST rules
\begin{align}
 \mathsf sc&=\mathcal L_\xi c-\frac12[c,c],
 &
 \mathsf sA&=\mathcal L_\xi A+D_Ac,
\label{eq:v3v60-semidirect-BRST}
\end{align}
and the corresponding covariant rules for matter, the direct-sum boundary
domain is a nilpotent BRST subcomplex.  The complete junction and all
V55--V57 matter domains lie in that subcomplex.
\end{theorem}

\begin{proof}
Internal-gauge closure is Theorem
\ref{thm:v3v58-domain-invariant}; diffeomorphism closure is Propositions
\ref{prop:v3v60-geometric-domain}--%
\ref{prop:v3v60-absolute-subcomplex}.  The cross term
\(\mathcal L_\xi c\) preserves Neumann data because \(\xi\) is tangent and
its absolute normal jet preserves the collar normal.  Nilpotency is the
Jacobi identity for the semidirect Lie algebra of boundary-preserving
diffeomorphisms and gauge transformations.  Every material and fermion
projector is built naturally from transforming fields, so its zero set is
preserved.
\end{proof}

\subsection{Finite-cutoff determinant and joint zero modes}
\label{sec:v3v60-determinant}

Choose diagonal spectral cutoff spaces for
\(\mathcal M_{\rm d}\) and \(\mathcal M_{\rm g}\) on their respective
absolute and Neumann domains, and project \(\mathcal K_A\) between them.
The resulting finite matrix retains the triangular form
\eqref{eq:v3v60-joint-FP}.

\begin{proposition}[Exact triangular ghost determinant]
\label{prop:v3v60-triangular-determinant}
If both diagonal blocks are invertible on the projected spaces, then
\begin{equation}
 \det\mathcal M_{{\rm joint},\Lambda}
 =\det\mathcal M_{{\rm d},\Lambda}
  \det\mathcal M_{{\rm g},\Lambda},
\label{eq:v3v60-triangular-determinant}
\end{equation}
independently of \(\mathcal K_{A,\Lambda}\), and
\begin{equation}
 \mathcal M_{{\rm joint},\Lambda}^{-1}
 =\begin{pmatrix}
  \mathcal M_{\rm d}^{-1}&0\\
  -\mathcal M_{\rm g}^{-1}\mathcal K_A
   \mathcal M_{\rm d}^{-1}&\mathcal M_{\rm g}^{-1}
 \end{pmatrix}.
\label{eq:v3v60-triangular-inverse}
\end{equation}
\end{proposition}

\begin{proof}
These are the determinant and inverse formulas for a finite block lower-
triangular matrix.  Projection does not alter triangularity because the
upper-right block is identically zero before projection.
\end{proof}

When either diagonal has a kernel, the joint kernel is
\begin{equation}
 \ker\mathcal M_{\rm joint}
 =\left\{(\xi,c):
  \mathcal M_{\rm d}\xi=0,\quad
  \mathcal M_{\rm g}c=-\mathcal K_A\xi
 \right\}.
\label{eq:v3v60-joint-kernel}
\end{equation}

\begin{assessment}[Zero modes must be assembled]
\label{ass:v3v60-joint-zero}
Equation \eqref{eq:v3v60-joint-kernel} is not generally the direct sum of
the two diagonal kernels.  A boundary isometry may require a compensating
internal gauge transformation to stabilize \(A\), and an internal
stabilizer can overlap that compensation.  The correct finite-dimensional
zero group is the stabilizer of the pair \((g,A)\) under the semidirect
action \eqref{eq:v3v60-combined-action}.  This is the precise application
of the V94 assembled-carrier lesson.
\end{assessment}

\begin{corollary}[Simple-background zero sector]
\label{cor:v3v60-simple-zero}
At the V53 simple background \(A=0\), one has \(\mathcal K_A=0\), and the
joint zero algebra is
\begin{equation}
 \mathfrak{isom}_\partial(\bar g)
 \oplus\mathfrak g_{\rm constant}.
\label{eq:v3v60-simple-zero}
\end{equation}
A based semidirect gauge group removes it.  On its orthogonal complement,
the nonzero Euclidean ghost normal family is strongly elliptic.
\end{corollary}

\begin{proof}
At \(A=0\), the variation of the Lorenz gauge under a diffeomorphism of
\(A\) is zero, so \(\mathcal K_A=0\).  Combine
Theorem~\ref{thm:v3v60-vector-elliptic} with
Theorem~\ref{thm:v3v58-Euclidean-elliptic} and remove the two finite
stabilizer spaces.
\end{proof}

\subsection{The physical graviton firewall}
\label{sec:v3v60-graviton-firewall}

\begin{firewall}[Ghost ellipticity does not prove graviton ellipticity]
\label{fw:v3v60-graviton-ellipticity}
V60 proves the Euclidean complementing condition for the diffeomorphism
and internal ghost operators.  The Euclidean metric Hessian with the full
Brown--York/material junction and harmonic rows is a different tensor
boundary problem.  Its transverse-traceless and physical quotient sectors
are invisible to the ghost normal family, just as the NS V31 hinge
nullspace is invisible to total energy.  Strong ellipticity and a lower
spectral bound for that tensor Hessian remain unproved.
\end{firewall}

\begin{firewall}[No interacting joint determinant]
\label{fw:v3v60-no-interacting-det}
The determinant identity
\eqref{eq:v3v60-triangular-determinant} is finite dimensional and
one-loop.  A sharp cutoff is not closed under nonlinear semidirect BRST
products.  V60 does not define the renormalized joint ghost determinant,
its determinant line over backgrounds, or its coupling to the chiral eta
phase.
\end{firewall}

\subsection{V60 conclusion and V61 target}
\label{sec:v3v60-conclusion}

The harmonic Einstein ghost has the absolute vector domain
\(\iota_n\xi=\iota_nd\xi=0\), whose nonzero Euclidean normal family is
injective.  Together with V58's Neumann internal ghosts it forms a
nilpotent semidirect BRST boundary complex preserving the complete
junction and all bosonic and fermionic matter domains.  Its finite
Faddeev--Popov matrix is triangular, but its zero modes are the assembled
stabilizer of \((g,A)\), not an automatic sum of separate kernels.

V61 should attack the physical Euclidean metric Hessian left outside the
ghost calculation.  It should Wick rotate the exact V42 cold-plus-static
junction, specify the Euclidean continuation of the cold line, compute the
full tensor/material normal family, and test strong ellipticity on every
nonzero covector.  If the dissipative cold term has no reflection-positive
Euclidean continuation, that obstruction must be isolated before any
gravitational determinant or eta-line argument continues.
\clearpage
\section{V61: Euclidean cold-plus-static ellipticity and the conformal-factor obstruction}
\label{sec:v3v61}

\subsection{Continue the local parent, not the reduced impedance}
\label{sec:v3v61-parent}

The Lorentzian cold response is \(\beta s\) in the open Laplace half-plane.
Substituting an imaginary frequency into that reduced expression is not the
Euclidean construction.  Wick rotation must be performed before the
half-line field is eliminated.

For the cold carrier, the Euclidean action is
\begin{equation}
 S_0^E
 =\frac12\int_{\mathcal N_E}d\tau,d^2z
  \int_0^\infty dy\,
  \left(
   |\partial_\tau\Psi_0|^2+|\partial_y\Psi_0|^2
  \right),
\label{eq:v3v61-cold-Euclidean}
\end{equation}
with trace \(\Psi_0(0)=\alpha_0e\).  The static carrier remains
\begin{equation}
 S_*^E
 =\frac12\int d\tau,d^2z\int_0^\infty dy\,
  \left(
   |\partial_y\Phi|^2+c_*^2|D_\Sigma\Phi|^2
  \right),
\label{eq:v3v61-static-Euclidean}
\end{equation}
with trace \(\Phi(0)=\alpha_*e\).

\begin{proposition}[Exact Euclidean impedances]
\label{prop:v3v61-Euclidean-impedances}
At Fourier frequency \((\omega,\xi)\), the decaying solutions are
\begin{align}
 \widehat\Psi_0(y)&=alpha_0e^{-|\omega|y}\widehat e,
\label{eq:v3v61-cold-Poisson}\
 \widehat\Phi(y)&=alpha_*e^{-c_*|\xi|y}\widehat e.
\label{eq:v3v61-static-Poisson}
\end{align}
Elimination produces the positive boundary form
\begin{equation}
 \frac12\left(
  \beta|\omega|+\gamma|\xi|
 \right)|\widehat e|^2,
 \qquad \beta,\gamma>0.
\label{eq:v3v61-Euclidean-response}
\end{equation}
\end{proposition}

\begin{proof}
Fourier transformation reduces the cold Euler equation to
\(-\partial_y^2\widehat\Psi_0+\omega^2\widehat\Psi_0=0\) and the static
equation to
\(-\partial_y^2\widehat\Phi+c_*^2|\xi|^2\widehat\Phi=0\).  Decay and the
traces give \eqref{eq:v3v61-cold-Poisson}--%
\eqref{eq:v3v61-static-Poisson}.  Direct integration of the two positive
energies gives coefficients proportional to \(|\omega|\) and \(|\xi|\);
absorbing the fixed incidences into \(\beta,\gamma\) gives
\eqref{eq:v3v61-Euclidean-response}.
\end{proof}

\begin{assessment}[Retarded and Euclidean responses share one parent]
\label{ass:v3v61-parent-continuation}
The Lorentzian response \(\beta s\) and Euclidean response
\(\beta|\omega|\) are two boundary values of the same local half-line
problem with different time signatures.  They need not be obtained by a
pointwise analytic substitution in an already eliminated kernel.  The
local parent fixes the sign and removes that ambiguity.
\end{assessment}

\subsection{The augmented physical normal problem}
\label{sec:v3v61-physical-normal}

Let \(U\) denote the metric/material physical quotient of V48, with
positive fibre metric \(M\).  At a flat Euclidean point put
\begin{equation}
 \rho=(\omega^2+|\xi|^2)^{1/2}>0.
\label{eq:v3v61-rho}
\end{equation}
Before eliminating either carrier, the principal quadratic form per
tangential Fourier mode is
\begin{align}
 \mathcal Q_{\omega,\xi}[U,\Psi_0,\Phi]
 :={}&\frac12\int_0^\infty
  \left(
   |\partial_xU|_M^2+\rho^2|U|_M^2
  \right)dx
\notag\\
 &+\frac12\int_0^\infty
  \left(
   |\partial_y\Psi_0|^2+\omega^2|\Psi_0|^2
  \right)dy
\notag\\
 &+\frac12\int_0^\infty
  \left(
   |\partial_y\Phi|^2+c_*^2|\xi|^2|\Phi|^2
  \right)dy,
\label{eq:v3v61-augmented-form}
\end{align}
subject to the incidence traces
\begin{equation}
 \Psi_0(0)=C_0U(0),
 \qquad
 \Phi(0)=C_*U(0).
\label{eq:v3v61-incidence-traces}
\end{equation}
The natural boundary equation for \(U\) is the Euclidean continuation of
the assembled Brown--York/material junction.

\begin{lemma}[One-dimensional trace coercivity]
\label{lem:v3v61-trace-coercivity}
For \(r>0\) and every decaying \(v\in H^1(0,\infty)\),
\begin{equation}
 \int_0^\infty\left(|v'|^2+r^2|v|^2\right)dx
 \ge r|v(0)|^2.
\label{eq:v3v61-trace-coercivity}
\end{equation}
\end{lemma}

\begin{proof}
Expand the nonnegative integral
\[
 \int_0^\infty|v'+rv|^2dx
 =\int_0^\infty(|v'|^2+r^2|v|^2)dx-r|v(0)|^2.
\]
\end{proof}

\begin{theorem}[Strong ellipticity of the augmented physical quotient]
\label{thm:v3v61-physical-elliptic}
For every \((\omega,\xi)\ne0\), the homogeneous decaying Euler--boundary
system of \eqref{eq:v3v61-augmented-form} has only the zero solution.
Moreover
\begin{equation}
 \mathcal Q_{\omega,\xi}
 \ge c\rho|U(0)|^2
  +\frac{\beta}{2}|\omega|,|C_0U(0)|^2
  +\frac{\gamma}{2}|\xi|,|C_*U(0)|^2
\label{eq:v3v61-physical-coercivity}
\end{equation}
for a uniform \(c>0\) on the compact positive parameter set.  Hence the
augmented physical boundary problem satisfies the Euclidean complementing
condition with an order-one inverse bound.
\end{theorem}

\begin{proof}
Apply Lemma~\ref{lem:v3v61-trace-coercivity} to \(U\) with \(r=\rho\), to
\(\Psi_0\) with \(r=|\omega|\) when \(\omega\ne0\), and to \(\Phi\) with
\(r=c_*|\xi|\) when \(\xi\ne0\).  The trace incidences give
\eqref{eq:v3v61-physical-coercivity}; the bulk term alone controls every
component of \(U(0)\) because \(M>0\) and \(\rho>0\).

If a homogeneous solution existed, test its three Euler equations by the
solution and integrate by parts.  The natural junction terms cancel the
three boundary variations because they arise from the common quadratic
form.  The result is
\(2\mathcal Q_{\omega,\xi}=0\).  Positivity forces
\(U=\Psi_0=\Phi=0\).  The quantitative lower bound and standard finite-
dimensional normal-family compactness give the order-one inverse estimate.
\end{proof}

\begin{corollary}[Ellipticity after carrier elimination]
\label{cor:v3v61-reduced-elliptic}
Eliminating \((\Psi_0,\Phi)\) gives the reduced positive response
\(\beta|\omega|+\gamma|\xi|\) and preserves the zero kernel and
complementing estimate of Theorem
\ref{thm:v3v61-physical-elliptic}.
\end{corollary}

\begin{proof}
For each trace, the fields in Proposition
\ref{prop:v3v61-Euclidean-impedances} uniquely minimize their positive
half-line energies.  Substitution is therefore a positive Schur
complement.  A zero of the reduced problem would lift to a zero of the
augmented problem, which the theorem excludes.
\end{proof}

\subsection{Finite-regulator reflection positivity of the carriers}
\label{sec:v3v61-reflection}

Choose a reflection-symmetric time lattice and positive nearest-neighbor
discretization of \eqref{eq:v3v61-cold-Euclidean}.  Keep the static action
diagonal in Euclidean time, as in V47.

\begin{theorem}[Carrier reflection positivity]
\label{thm:v3v61-carrier-RP}
At every finite spatial, half-line, and reflection-symmetric time cutoff,
the Gaussian measure of the direct-sum cold and static carriers is
reflection positive.  Imposing the linear trace constraints
\eqref{eq:v3v61-incidence-traces} preserves reflection positivity when the
boundary strain is reflected tensorially.
\end{theorem}

\begin{proof}
For the cold carrier, slice the local positive lattice action at the
reflection plane.  Its transfer matrix is the heat kernel of a positive
quadratic Hamiltonian.  For every polynomial \(F\) supported at positive
Euclidean times, the Markov composition law gives
\[
 \langle\Theta F\,F\rangle
 =\|e^{-\tau_0H}F\Omega\|^2\ge0.
\]
The static carrier has no coupling between distinct time slices, so its
reflection form is the square of the integral over positive-time variables,
as proved in the V47 finite-cutoff argument.  The direct sum preserves
positivity.  The trace constraints identify fields at equal Euclidean time
through a real linear map commuting with reflection; restricting a positive
quadratic Gaussian form to such an invariant linear subspace preserves the
same square representation.
\end{proof}

\begin{firewall}[A spectral time cutoff need not preserve reflection]
\label{fw:v3v61-cutoff-choice}
Theorem~\ref{thm:v3v61-carrier-RP} uses a reflection-symmetric local time
regulator.  A sharp Fourier cutoff in \(\omega\) is convenient for the V59
determinant identity but is not automatically reflection positive.  The two
finite regulators may be compared only after a common reflection-positive,
BRST-compatible scheme is constructed.
\end{firewall}

\subsection{The Euclidean Einstein conformal factor}
\label{sec:v3v61-conformal}

The preceding positivity is a theorem on the physical quotient and its
carriers.  The full gauge-fixed metric integral still contains the
conformal trace.  At a flat background, write
\begin{equation}
 h_{\mu\nu}=h^{\rm TL}_{\mu\nu}+\frac14g_{\mu\nu}\sigma.
\label{eq:v3v61-trace-split}
\end{equation}
With the conventional positive sign for the transverse-traceless Euclidean
Einstein quadratic action, the trace contribution has the opposite sign:
\begin{equation}
 Q_{\rm conf}[\sigma]
 =-c_g\int_{M_E}|d\sigma|^2dV
  +\text{lower-order and boundary terms},
 \qquad c_g>0.
\label{eq:v3v61-conformal-form}
\end{equation}

\begin{theorem}[Boundary carriers cannot cure the conformal-factor
instability]
\label{thm:v3v61-conformal-obstruction}
The real Euclidean quadratic form of the complete metric/carrier system is
unbounded below, even though the physical quotient and both carriers are
positive.  More precisely, there is a sequence of smooth trace
perturbations supported in a fixed interior coordinate ball, vanishing near
the boundary, for which
\begin{equation}
 Q_E[h^{(N)},0,0]\longrightarrow-\infty.
\label{eq:v3v61-conformal-minus-infinity}
\end{equation}
\end{theorem}

\begin{proof}
Choose a nonzero compactly supported smooth function \(\chi\) in an
interior coordinate ball and set
\(\sigma_N=N^{-1/2}\chi(x)\sin(Nx^1)\).  Then
\(\|d\sigma_N\|_2^2\asymp N\), while its lower-order quadratic norms are
\(O(N^{-1})\) or \(O(1)\).  Since \(\sigma_N\) vanishes in a collar of the
boundary, every Brown--York, cold, static, and boundary material term is
zero.  Equation \eqref{eq:v3v61-conformal-form} therefore gives
\(Q_E\le-cN+O(1)\), proving
\eqref{eq:v3v61-conformal-minus-infinity}.
\end{proof}

\begin{corollary}[No real Gaussian gravitational measure]
\label{cor:v3v61-no-real-Gaussian}
No finite-dimensional regulator which contains arbitrarily high interior
conformal modes has a convergent real Gaussian integral with weight
\(e^{-Q_E}\).  Boundary damping, static stiffness, and ghost determinants
do not alter this conclusion.
\end{corollary}

\begin{proof}
Along every negative conformal eigenvector, the real Gaussian factor grows
exponentially.  The carriers vanish on the sequence used in the theorem,
and Grassmann determinants multiply rather than change convergence of a
bosonic real contour.
\end{proof}

The formal repair is a contour rotation
\begin{equation}
 \sigma=i\varphi,
\label{eq:v3v61-conformal-rotation}
\end{equation}
which makes the principal trace form positive.  This is not yet a theorem
about the nonlinear configuration space: the trace mixes with boundary
conditions, zero modes, and the determinant line.

\begin{firewall}[The full Euclidean Hessian remains open]
\label{fw:v3v61-full-Hessian}
V61 proves strong ellipticity and finite-regulator reflection positivity
for the augmented physical cold-plus-static quotient.  It disproves
convergence of the full real Euclidean metric Gaussian.  It does not prove
that the contour \eqref{eq:v3v61-conformal-rotation} is globally defined,
BRST invariant, compatible with the complete junction domain, or
reflection positive.  Therefore no gravitational determinant or coupled
Euclidean measure is claimed.
\end{firewall}

\subsection{V61 conclusion and V62 target}
\label{sec:v3v61-conclusion}

Wick rotation of the local carriers produces the positive response
\(\beta|\omega|+\gamma|\xi|\).  The augmented physical quotient is
coercive at every nonzero Euclidean tangential covector and its local
finite regulator is reflection positive.  These results remove the cold
dissipation as the Euclidean obstruction.  The actual obstruction is the
interior Einstein conformal factor, which no boundary term can control.

V62 should test the conformal rotation on the complete linearized boundary
domain.  It must split the ten metric amplitudes into trace, pure-gauge,
and physical quotient variables, compute how the Brown--York and de Donder
rows transform under \(\sigma=i\varphi\), and determine whether a constant
complex contour is preserved by the diffeomorphism BRST differential.  If
the boundary Hessian becomes non-sectorial or the contour is not BRST
closed, the programme must move upstream to a different Euclidean
definition rather than assigning a formal determinant.
\clearpage
\section{V62: conformal rotation of the longitudinal gravitational BRST quartet}
\label{sec:v3v62}

\subsection{Trace action of the diffeomorphism differential}
\label{sec:v3v62-trace-BRST}

V61 shows that rotating the metric trace is necessary for convergence of
the Euclidean Gaussian.  Rotating that variable alone is not compatible
with the linear BRST differential.  In four dimensions write
\begin{equation}
 h_{\mu\nu}=h^{\rm TL}_{\mu\nu}
  +\frac14g_{\mu\nu}\sigma.
\label{eq:v3v62-trace-split}
\end{equation}
At a fixed background,
\begin{align}
 \mathsf sh_{\mu\nu}&=\nabla_\mu\xi_\nu+\nabla_\nu\xi_\mu,
\label{eq:v3v62-sh}\
 \mathsf s\sigma&=2\nabla_\mu\xi^\mu,
\label{eq:v3v62-ssigma}\
 \mathsf sh^{\rm TL}_{\mu\nu}
 &=\nabla_\mu\xi_\nu+\nabla_\nu\xi_\mu
   -\frac12g_{\mu\nu}\nabla\cdot\xi.
\label{eq:v3v62-sTL}
\end{align}

On the Euclidean regulator, use the absolute Hodge decomposition of the
diffeomorphism ghost,
\begin{equation}
 \xi=\xi_{\rm T}+d\upsilon+\xi_{\rm harm},
 \qquad \delta\xi_{\rm T}=0,
\label{eq:v3v62-ghost-Hodge}
\end{equation}
where \(\upsilon\) has Neumann data and zero mean.  Remove the harmonic
stabilizer as in V60.  Then only \(d\upsilon\) contributes to
\eqref{eq:v3v62-ssigma}.

\begin{proposition}[Trace and exact ghost are one BRST block]
\label{prop:v3v62-trace-exact-block}
On a scalar eigenmode
\(\Delta_0\upsilon=\lambda\upsilon\), \(\lambda>0\), the linear BRST map
has the form
\begin{equation}
 \mathsf s\sigma=-2\lambda\upsilon,
 \qquad \mathsf s\upsilon=0,
\label{eq:v3v62-mode-BRST}
\end{equation}
up to the fixed sign convention for \(\delta d=\Delta_0\).  Therefore the
real trace contour is not an invariant linear cycle after
\(\sigma=i\varphi\) unless the exact ghost contour is rotated at the same
time.
\end{proposition}

\begin{proof}
The co-closed and harmonic terms in
\eqref{eq:v3v62-ghost-Hodge} have zero divergence.  Since
\(\nabla\cdot d\upsilon=-\Delta_0\upsilon\) in the stated convention,
\eqref{eq:v3v62-ssigma} gives
\eqref{eq:v3v62-mode-BRST}.  If \(\upsilon\) remains on its original real
Grassmann real form, \(i\mathsf s\varphi=-2\lambda\upsilon\), so
\(\mathsf s\varphi\) does not have the rotated reality convention.
\end{proof}

\subsection{The quartet rotation}
\label{sec:v3v62-quartet}

Let \((\upsilon,\bar\upsilon,b_\upsilon)\) be the scalar potentials of the
exact diffeomorphism ghost, antighost, and auxiliary field.  On each
nonzero scalar mode perform the constant complex change
\begin{equation}
 \sigma=i\varphi,
 \qquad
 \upsilon=i\eta,
 \qquad
 \bar\upsilon=i\bar\eta,
 \qquad
 b_\upsilon=i\beta_\eta.
\label{eq:v3v62-quartet-rotation}
\end{equation}
The transverse ghost and traceless physical variables are not rotated.

\begin{theorem}[BRST-tangent finite-dimensional conformal cycle]
\label{thm:v3v62-BRST-cycle}
At every Hodge-compatible finite cutoff, the supercycle
\eqref{eq:v3v62-quartet-rotation} is preserved by the linear gravitational
BRST differential.  In rotated variables,
\begin{equation}
 \mathsf s\varphi=-2\lambda\eta,
 \qquad
 \mathsf s\eta=0,
 \qquad
 \mathsf s\bar\eta=\beta_\eta,
 \qquad
 \mathsf s\beta_\eta=0.
\label{eq:v3v62-rotated-BRST}
\end{equation}
The combined bosonic--Berezin Jacobian of the quartet rotation is one per
mode.
\end{theorem}

\begin{proof}
Insert \eqref{eq:v3v62-quartet-rotation} into
\eqref{eq:v3v62-mode-BRST} and into
\(\mathsf s\bar\upsilon=b_\upsilon\).  Cancelling the common factor \(i\)
gives \eqref{eq:v3v62-rotated-BRST}; thus the differential is tangent to
the rotated real form of the supervector space.

The two bosonic changes \((\sigma,b_\upsilon)=
i(\varphi,\beta_\eta)\) have ordinary Jacobian \(i^2=-1\).  Berezin
measure transforms inversely, so the two fermionic changes
\((\upsilon,\bar\upsilon)=i(\eta,\bar\eta)\) contribute
\(i^{-2}=-1\).  Their product is one.
\end{proof}

\begin{corollary}[Finite-cutoff BRST Ward identity on the rotated cycle]
\label{cor:v3v62-rotated-Ward}
If the rotated Gaussian has positive real part, then for every polynomial
\(K\) in the finite cutoff variables,
\begin{equation}
 \int_{\Gamma_{\rm rot}}d\mu_\Lambda\,
 \mathsf s\left(Ke^{-S_\Lambda^{(2)}}\right)=0.
\label{eq:v3v62-rotated-Ward}
\end{equation}
\end{corollary}

\begin{proof}
Theorem~\ref{thm:v3v62-BRST-cycle} makes \(\mathsf s\) tangent to the
integration supercycle and gives unit measure Jacobian.  Positive real part
gives decay at its bosonic ends.  Finite-dimensional integration of the
total BRST derivative then vanishes, exactly as in Theorem
\ref{thm:v3v59-finite-Ward}.
\end{proof}

\subsection{Exact trace contribution to the boundary rows}
\label{sec:v3v62-boundary-rows}

At a flat Euclidean boundary, take a decaying conformal amplitude
\begin{equation}
 H_{\mu\nu}^{\rm conf}
 =\frac14\delta_{\mu\nu}\sigma e^{-\rho x+i p\cdot y},
 \qquad \rho=|p|>0.
\label{eq:v3v62-conformal-amplitude}
\end{equation}
In the notation of V40,
\begin{equation}
 V_a=0,
 \qquad w=\frac14\sigma,
 \qquad \theta=\frac34\sigma.
\label{eq:v3v62-conformal-components}
\end{equation}

\begin{proposition}[Conformal Brown--York and de Donder rows]
\label{prop:v3v62-conformal-rows}
The conformal contribution to the Euclideanized principal boundary rows is
\begin{align}
 C_a^{\rm conf}&=-\frac14(ip_a)\sigma,
 &C_x^{\rm conf}&=\frac{\rho}{4}\sigma,
\label{eq:v3v62-conformal-C}\
 \mathcal J_{ab}^{\rm conf}&=\frac{\rho}{2}\delta_{ab}\sigma
\label{eq:v3v62-conformal-J}
\end{align}
for the V40 normal convention.  Under \(\sigma=i\varphi\), every displayed
trace column is multiplied by \(i\), and no other boundary column changes.
\end{proposition}

\begin{proof}
Insert \eqref{eq:v3v62-conformal-components} in the stable de Donder rows
\eqref{eq:v3v40-Ca}--\eqref{eq:v3v40-Cx}, replacing the Euclidean normal
root by \(\rho\).  This gives
\[
 C_a=\frac14ip_a\sigma-\frac12ip_a\sigma,
 \qquad
 C_x=-\frac\rho2\left(\frac14-\frac34\right)\sigma.
\]
For Brown--York, insert \(H_{ab}=\delta_{ab}\sigma/4\),
\(V_a=0\), and \(\theta=3\sigma/4\) into
\eqref{eq:v3v40-Jplus}.  The result is
\((-1/4+3/4)\rho\delta_{ab}\sigma\).  The rotation statement is
linearity.
\end{proof}

\begin{corollary}[The conformal column has no nonzero normal kernel]
\label{cor:v3v62-conformal-column}
For \(\rho>0\), a pure conformal decaying mode satisfying the homogeneous
de Donder rows is zero.  The same is true after rotation.
\end{corollary}

\begin{proof}
Equation \eqref{eq:v3v62-conformal-C} contains
\(C_x=\rho\sigma/4\).  Thus \(C_x=0\) implies \(\sigma=0\).  Multiplication
of the column by \(i\) does not change its kernel.
\end{proof}

\subsection{Normal-family invariance under contour rotation}
\label{sec:v3v62-normal-invariance}

Let \(R\) be the complex linear map which is the identity on traceless,
physical, material, and carrier amplitudes and multiplication by \(i\) on
the conformal amplitude.  If \(\mathcal B_E(p)\) is the unrotated full
principal boundary matrix, the rotated matrix is
\begin{equation}
 \mathcal B_E^{\rm rot}(p)=\mathcal B_E(p)R.
\label{eq:v3v62-rotated-boundary-matrix}
\end{equation}

\begin{lemma}[Complementing kernels are unchanged]
\label{lem:v3v62-kernel-invariance}
For every nonzero tangential covector,
\begin{equation}
 \ker\mathcal B_E^{\rm rot}(p)
 =R^{-1}\ker\mathcal B_E(p).
\label{eq:v3v62-kernel-invariance}
\end{equation}
In particular, a constant invertible contour rotation cannot create or
remove a complementing root.
\end{lemma}

\begin{proof}
The equation \(\mathcal B_E Rv=0\) is equivalent to
\(Rv\in\ker\mathcal B_E\), and \(R\) is invertible.
\end{proof}

At the decoupled simple background, use the direct decomposition into the
physical quotient of V61, pure diffeomorphism directions of V60, and the
conformal column of Proposition~\ref{prop:v3v62-conformal-rows}.

\begin{theorem}[Strong ellipticity after the finite conformal rotation]
\label{thm:v3v62-rotated-elliptic}
Assume the V54 Schur gap and the simple-background block decomposition.
Then the complete rotated metric/material/carrier normal family has zero
decaying kernel at every nonzero Euclidean tangential covector.  The result
persists under sufficiently small coefficient perturbations preserving the
normalized Schur gap.
\end{theorem}

\begin{proof}
The physical block has zero kernel by Theorem
\ref{thm:v3v61-physical-elliptic}.  Pure gauge amplitudes are controlled by
the V60 absolute ghost/gauge coordinate block, and the remaining conformal
line is killed by Corollary~\ref{cor:v3v62-conformal-column}.  The V54
Schur gap controls the off-diagonal constraint--physical coupling.  Thus
the unrotated principal matrix is injective.  Lemma
\ref{lem:v3v62-kernel-invariance} transfers injectivity to the rotated
matrix.  Compactness of the normalized covector sphere and continuity of
the least singular value give stability under small perturbations.
\end{proof}

\subsection{Sectoriality of the rotated quadratic form}
\label{sec:v3v62-sectorial}

At the real decoupled background, trace and traceless bulk modes are
orthogonal for the DeWitt Hessian.  Rotation changes the negative trace
term \eqref{eq:v3v61-conformal-form} to
\begin{equation}
 +c_g\int|d\varphi|^2.
\label{eq:v3v62-positive-trace}
\end{equation}
Real symmetric trace--physical boundary cross terms acquire one factor of
\(i\).

\begin{theorem}[Positive real part near the simple background]
\label{thm:v3v62-sectorial}
At the decoupled simple background, the rotated finite-cutoff quadratic
form satisfies
\begin{equation}
 \operatorname{Re}Q_E^{\rm rot}
 \ge c\left(
  \|h^{\rm TL}_{\rm phys}\|_{H^1}^2
  +\|\varphi\|_{H^1}^2
  +\mathcal Q_{\rm carriers}
 \right)-C\|h\|_{L^2}^2
\label{eq:v3v62-sectorial-bound}
\end{equation}
after the gauge and finite zero modes are removed.  The estimate persists
on a sufficiently small complex coefficient neighborhood.
\end{theorem}

\begin{proof}
V61 gives positivity of the physical and carrier blocks, while
\eqref{eq:v3v62-positive-trace} gives positivity of the conformal block.
At the real decoupled point, a cross term containing exactly one rotated
trace variable is purely imaginary and drops out of the real part.  Gauge-
fixing supplies the positive nonzero pure-gauge block; harmonic stabilizers
were removed in V60.  Lower-order curvature terms give the displayed
\(L^2\) remainder.  The strict principal and Schur gaps absorb sufficiently
small perturbations by the form version of the relative-bound argument.
\end{proof}

\begin{corollary}[Convergent finite Gaussian on the rotated supercycle]
\label{cor:v3v62-finite-Gaussian}
After the finite infrared modes are treated and the positive mass shift
used to control the Gårding remainder is removed by ordinary finite-
dimensional continuation, the regulated quadratic integral on
\(\Gamma_{\rm rot}\) converges and obeys the BRST Ward identity
\eqref{eq:v3v62-rotated-Ward}.
\end{corollary}

\begin{proof}
Theorem~\ref{thm:v3v62-sectorial} makes the real part positive definite on
the finite nonzero-mode space after an admissible shift.  The resulting
Gaussian converges.  If no eigenvalue crosses zero while the shift is
removed, analytic continuation defines the unshifted integral; any
crossing is instead recorded as a finite zero mode.  Apply Corollary
\ref{cor:v3v62-rotated-Ward}.
\end{proof}

\subsection{Global and reflection-positivity firewalls}
\label{sec:v3v62-firewalls}

\begin{firewall}[The Hodge rotation is background dependent]
\label{fw:v3v62-background-dependent}
The decomposition \eqref{eq:v3v62-ghost-Hodge} depends on the Euclidean
metric and its boundary.  Along a general configuration-space loop,
eigenvalues can cross zero and exact, harmonic, and coexact bundles can
have nontrivial determinant-line holonomy.  V62 proves a finite-cutoff
cycle in one simple-background neighborhood.  It does not construct a
global nonlinear integration cycle.
\end{firewall}

\begin{firewall}[Physical reflection positivity is not yet reconstructed]
\label{fw:v3v62-RP}
The rotated trace is imaginary in the original metric variable.  Although
it belongs to a BRST quartet and the physical cold/static quotient is
reflection positive by V61, V62 does not prove Osterwalder--Schrader
positivity for the full gauge-fixed measure or for its BRST cohomology.
Such a result requires a common reflection-positive regulator and a
positive physical reconstruction map.
\end{firewall}

\begin{firewall}[No determinant-line phase]
\label{fw:v3v62-phase}
The unit super-Jacobian in Theorem~\ref{thm:v3v62-BRST-cycle} removes a
local mode-counting phase of the quartet rotation.  It does not determine
the global phase of the graviton, ghost, or chiral Weyl determinant lines.
\end{firewall}

\subsection{V62 conclusion and V63 target}
\label{sec:v3v62-conclusion}

The conformal trace, exact diffeomorphism ghost, antighost, and auxiliary
field form one longitudinal BRST quartet.  Rotating all four variables by
the same factor \(i\) gives a BRST-tangent finite-dimensional supercycle
with unit super-Jacobian.  The exact Brown--York and de Donder trace columns
are merely multiplied by \(i\), so the complementing kernel is unchanged;
the rotated quadratic form has positive real part near the simple
background.

V63 should replace the incompatible pair of sharp spectral and local time
regulators by one finite regulator which preserves the linear BRST complex,
the boundary domains, and reflection positivity on physical observables.
A cell or finite-element de Rham complex with an exact discrete gradient is
a plausible route.  The construction must reproduce the V59 determinant
pairing and V61 transfer-matrix positivity in the same finite model before
any cutoff limit is attempted.
\clearpage
\section{V63: correction of the conformal-quartet claim and the York scalar contour}
\label{sec:v3v63}

\subsection{The missing scalar-longitudinal metric component}
\label{sec:v3v63-correction}

Theorem~\ref{thm:v3v62-BRST-cycle} rotated the trace and the exact
diffeomorphism ghost but left the traceless metric variables on their real
contour.  That cycle is not preserved by the full BRST differential:
an exact ghost also generates a scalar-longitudinal traceless metric
variation.  This note corrects that statement before the regulator
programme continues.

Fix a nonzero Euclidean covector \(p\) and define the traceless scalar
tensor
\begin{equation}
 S_{\mu\nu}(p)
 =p_\mu p_\nu-\frac14\delta_{\mu\nu}p^2.
\label{eq:v3v63-S-tensor}
\end{equation}
Write the scalar part of a metric amplitude as
\begin{equation}
 h^{\rm sc}_{\mu\nu}
 =S_{\mu\nu}(p)\psi
  +\frac14\delta_{\mu\nu}\sigma.
\label{eq:v3v63-scalar-decomposition}
\end{equation}
For an exact ghost \(\xi_\mu=p_\mu\epsilon\),
\begin{equation}
 \mathsf sh_{\mu\nu}=2p_\mu p_\nu\epsilon.
\label{eq:v3v63-exact-gauge-image}
\end{equation}

\begin{proposition}[Full scalar BRST transformation]
\label{prop:v3v63-scalar-BRST}
The two scalar coordinates in
\eqref{eq:v3v63-scalar-decomposition} transform as
\begin{equation}
 \mathsf s\psi=2\epsilon,
 \qquad
 \mathsf s\sigma=2p^2\epsilon.
\label{eq:v3v63-scalar-BRST}
\end{equation}
Consequently
\begin{equation}
 \Phi:=\sigma-p^2\psi
\label{eq:v3v63-Phi}
\end{equation}
is BRST invariant:
\begin{equation}
 \mathsf s\Phi=0.
\label{eq:v3v63-Phi-closed}
\end{equation}
\end{proposition}

\begin{proof}
Substitute \eqref{eq:v3v63-S-tensor} into
\eqref{eq:v3v63-scalar-decomposition}.  The variations in
\eqref{eq:v3v63-scalar-BRST} give
\begin{align*}
 \mathsf sh^{\rm sc}_{\mu\nu}
 &=2\left(p_\mu p_\nu-\frac14\delta_{\mu\nu}p^2\right)\epsilon
   +\frac12\delta_{\mu\nu}p^2\epsilon\\
 &=2p_\mu p_\nu\epsilon,
\end{align*}
which is \eqref{eq:v3v63-exact-gauge-image}.  Subtraction gives
\eqref{eq:v3v63-Phi-closed}.
\end{proof}

\begin{counterexample}[The V62 trace-only quartet cycle is not BRST
tangent]
\label{cex:v3v63-V62-cycle}
Take a nonzero exact ghost mode and impose the V62 rotations
\(\sigma=i\varphi\), \(\epsilon=i\eta\), while keeping \(\psi\) real.
Then
\begin{equation}
 \mathsf s\psi=2i\eta,
\label{eq:v3v63-off-real-psi}
\end{equation}
which is not tangent to the real \(\psi\)-contour.  Thus
Theorem~\ref{thm:v3v62-BRST-cycle} and its unit-super-Jacobian corollary do
not describe the full scalar metric complex.  They may not be used in
later determinant or Ward arguments.
\end{counterexample}

\begin{proof}
Equation \eqref{eq:v3v63-off-real-psi} is the first formula in
\eqref{eq:v3v63-scalar-BRST} after \(\epsilon=i\eta\).  A nonzero imaginary
variation leaves the real integration line.
\end{proof}

\begin{assessment}[Parts of V62 which remain valid]
\label{ass:v3v63-V62-valid}
The explicit pure-trace Brown--York/de Donder columns in
Proposition~\ref{prop:v3v62-conformal-rows} remain correct.  The elementary
kernel identity in Lemma~\ref{lem:v3v62-kernel-invariance} also remains
correct for any specified invertible complex change of variables.  What is
withdrawn is the assertion that the particular trace/ghost quartet
rotation \eqref{eq:v3v62-quartet-rotation} is a BRST-tangent cycle of the
complete metric complex, and every conclusion depending on that assertion.
\end{assessment}

\subsection{The correct scalar coordinates}
\label{sec:v3v63-York}

Equation \eqref{eq:v3v63-scalar-decomposition} can be rewritten as
\begin{equation}
 h^{\rm sc}_{\mu\nu}
 =p_\mu p_\nu\psi
  +\frac14\delta_{\mu\nu}\Phi.
\label{eq:v3v63-York-form}
\end{equation}
The first term is a pure diffeomorphism coordinate, while \(\Phi\) is the
scalar coordinate on the quotient by exact diffeomorphisms.

\begin{theorem}[Scalar Einstein Hessian factors through \(\Phi\)]
\label{thm:v3v63-scalar-Hessian}
At a flat on-shell background, the gauge-invariant Euclidean Einstein
quadratic form restricted to the scalar sector depends only on \(\Phi\).
With the sign convention of \eqref{eq:v3v61-conformal-form},
\begin{equation}
 Q_{E,\rm inv}^{\rm sc}
 =-c_g p^2|\Phi|^2,
 \qquad c_g>0,
\label{eq:v3v63-negative-Phi}
\end{equation}
up to the fixed normalization used to define \(\Phi\).  Gauge fixing adds
a positive quadratic form on the independent orbit coordinate \(\psi\)
and does not change the negative sign of the quotient scalar.
\end{theorem}

\begin{proof}
Linearized diffeomorphism invariance makes the invariant Hessian vanish
whenever either argument is a pure gauge tensor.  By
\eqref{eq:v3v63-York-form}, its scalar restriction therefore factors
through the one-dimensional quotient coordinate \(\Phi\).  Evaluate that
factor on \(\psi=0\), where \(h_{\mu\nu}=\delta_{\mu\nu}\Phi/4\).  The
conformal calculation \eqref{eq:v3v61-conformal-form} gives a negative
multiple of \(p^2|\Phi|^2\), proving
\eqref{eq:v3v63-negative-Phi}.  A harmonic gauge-fixing square is positive
and nondegenerate along the pure-gauge orbit, but BRST exactness prevents it
from altering the invariant quotient form.
\end{proof}

\begin{assessment}[The negative scalar is not a ghost quartet]
\label{ass:v3v63-not-quartet}
The conformal instability belongs to the BRST-closed quotient coordinate
\(\Phi\), not solely to the exact diffeomorphism orbit.  Ghost determinants
cannot make its real Gaussian converge.  This is stronger than the V61
interior-support counterexample and corrects the interpretation in V62.
\end{assessment}

\subsection{The corrected local contour}
\label{sec:v3v63-contour}

Keep \((\psi,\epsilon,\bar\epsilon,b_\epsilon)\) on the original
gauge-fixing supercycle and rotate only the BRST-invariant quotient scalar:
\begin{equation}
 \Phi=i\phi,
 \qquad \phi\in\mathbb R.
\label{eq:v3v63-correct-rotation}
\end{equation}

\begin{theorem}[BRST tangency of the York scalar contour]
\label{thm:v3v63-York-cycle}
At every finite nonzero-mode cutoff in a fixed flat-background York
decomposition, the cycle \eqref{eq:v3v63-correct-rotation} is preserved by
the linear BRST differential.  The scalar invariant form becomes
\begin{equation}
 Q_{E,\rm inv}^{\rm sc}[i\phi]
 =+c_gp^2|\phi|^2.
\label{eq:v3v63-positive-Phi}
\end{equation}
The orbit gauge-fixing and ghost blocks remain on their original domains.
\end{theorem}

\begin{proof}
Equation \eqref{eq:v3v63-Phi-closed} gives
\(\mathsf s\phi=0\), so the differential is tangent to the rotated line.
The orbit coordinate and its ghost doublets were not rotated and retain
their ordinary BRST transformations.  Substitution of
\(\Phi=i\phi\) in \eqref{eq:v3v63-negative-Phi} gives
\eqref{eq:v3v63-positive-Phi}.
\end{proof}

\begin{proposition}[The rotation has an uncancelled local phase]
\label{prop:v3v63-contour-phase}
For \(N_\Phi\) retained real quotient-scalar modes, the change
\eqref{eq:v3v63-correct-rotation} contributes the ordinary contour factor
\begin{equation}
 i^{N_\Phi}.
\label{eq:v3v63-contour-phase}
\end{equation}
There is no corresponding ghost rotation which cancels this phase, because
\(\Phi\) is BRST invariant.  A regulator-independent meaning for
\eqref{eq:v3v63-contour-phase} belongs to the gravitational determinant
line and is not fixed by a local mode count.
\end{proposition}

\begin{proof}
Each one-dimensional bosonic contour change \(d\Phi=i,d\phi\) contributes
one factor of \(i\).  The ghost variables are unchanged in the correct
cycle, so there is no Berezin Jacobian.  Multiplication over the retained
modes gives \eqref{eq:v3v63-contour-phase}.
\end{proof}

\subsection{Boundary normal family in York coordinates}
\label{sec:v3v63-boundary}

Let \(R_Y(p)\) be the invertible scalar change
\((\psi,\Phi)\mapsto(\psi,\sigma=\Phi+p^2\psi)\) for \(p\ne0\), followed by
\(\Phi=i\phi\).  The complete boundary matrix in the new coordinates is
\begin{equation}
 \mathcal B_E^Y(p)=\mathcal B_E(p)R_Y(p).
\label{eq:v3v63-York-boundary}
\end{equation}

\begin{proposition}[Complementing kernels survive the corrected rotation]
\label{prop:v3v63-York-kernel}
For \(p\ne0\),
\begin{equation}
 \ker\mathcal B_E^Y(p)
 =R_Y(p)^{-1}\ker\mathcal B_E(p).
\label{eq:v3v63-York-kernel}
\end{equation}
Hence the corrected rotation preserves every nonzero-frequency
complementing result proved independently in V54, V60, and V61.
\end{proposition}

\begin{proof}
The coordinate map is invertible at \(p\ne0\).  Apply the same elementary
composition argument as Lemma~\ref{lem:v3v62-kernel-invariance}.
\end{proof}

The change \(R_Y(p)\) is homogeneous but depends on tangential frequency
through the York projector.  In position space it is pseudodifferential
and depends on the background metric.  Thus normal-family injectivity does
not yet give a global contour.

\subsection{Constraint meaning of the quotient scalar}
\label{sec:v3v63-constraint}

In Lorentzian vacuum gravity, the scalar quotient mode is removed from the
propagating phase space by the Hamiltonian and momentum constraints.  With
matter it is determined elliptically by the energy and longitudinal
momentum sources.  It is therefore absent from the two-polarization
physical energy of V61 even though it appears in the off-shell Euclidean
metric integral.

\begin{proposition}[BRST cohomology alone is too large off shell]
\label{prop:v3v63-BRST-too-large}
The scalar \(\Phi\) is BRST closed by
\eqref{eq:v3v63-Phi-closed}, but it is not a propagating gravitational
observable.  A physical Euclidean reconstruction must impose or resolve
the Einstein constraint complex in addition to taking BRST cohomology.
\end{proposition}

\begin{proof}
BRST closure follows from Proposition
\ref{prop:v3v63-scalar-BRST}.  The ADM Hamiltonian constraint is the Euler
equation of the lapse and fixes the scalar curvature combination; the
momentum constraint fixes the longitudinal conjugate momentum.  Linearized
around the simple background, these are elliptic equations for the scalar
and longitudinal York data, as shown by the isomorphisms in V53.  Thus
\(\Phi\) is constrained despite being invariant under infinitesimal
diffeomorphisms.
\end{proof}

\begin{firewall}[Correction to V62]
\label{fw:v3v63-correction}
Later notes must not cite Theorem~\ref{thm:v3v62-BRST-cycle}, Corollary
\ref{cor:v3v62-rotated-Ward}, or the unit-super-Jacobian conclusion as
statements about the complete scalar metric sector.  The valid replacement
is Theorem~\ref{thm:v3v63-York-cycle}: rotate the BRST-invariant York scalar
\(\Phi\), retain the gauge orbit complex, and record the uncancelled
determinant-line phase.
\end{firewall}

\begin{firewall}[No reflection-positive metric measure yet]
\label{fw:v3v63-no-RP}
The corrected contour makes the finite scalar Gaussian convergent and
preserves the nonzero normal kernel.  It does not prove reflection
positivity.  The variable \(\Phi\) is imaginary in the original metric,
and its removal from physical observables uses the Einstein constraints,
not BRST exactness alone.
\end{firewall}

\subsection{V63 conclusion and V64 target}
\label{sec:v3v63-conclusion}

The negative conformal direction is the gauge-invariant York scalar
\(\Phi=\sigma-p^2\psi\).  The trace-only quartet rotation proposed in V62
missed the scalar-longitudinal traceless metric variation and is not BRST
tangent.  Rotating \(\Phi\) alone is BRST compatible and makes its finite
Gaussian positive, but produces an uncancelled contour phase and requires
the Einstein constraint complex to identify physical observables.

V64 should build the common finite regulator using both complexes: a
discrete de Rham complex for internal gauge symmetry and a discrete
Killing--constraint complex for linearized gravity.  Its incidence maps
must commute exactly, its boundary parities must reproduce the absolute
domains, and its physical observable space must impose the discrete
Einstein constraints before reflection positivity is tested.  The V59
determinant pairing and V61 carrier transfer matrix must then be recovered
inside that one corrected regulator.
\clearpage
\section{V64: one finite local regulator for the gauge, constraint, and carrier complexes}
\label{sec:v3v64}

V59 treated the internal Faddeev--Popov cancellation in a finite spectral
cutoff, whereas V61 obtained reflection positivity from a local time
carrier.  Those two regulators cannot simply be identified: a sharp
spectral projection is nonlocal in Euclidean time.  V63 also showed that
BRST cohomology does not by itself remove the gauge-invariant conformal
scalar.  This note constructs a single finite, nearest-neighbour regulator
for the linear theory.  The Einstein constraints are imposed before
reflection positivity is asserted.

The construction is deliberately restricted to a flat background, a flat
boundary face, and the free quadratic theory.  Within that scope every
complex and every determinant below is finite dimensional.

\subsection{Reflection double and the discrete de Rham complex}
\label{sec:v3v64-derham}

Let \(\widetilde X_h\) be a finite cubical cellulation invariant under
reflection \(r\) in a flat spatial face, and let \(X_h\) be one half of the
double.  Write \(C^j=C^j(\widetilde X_h;\mathfrak g)\) for Lie-algebra-valued
cellular cochains, with incidence maps
\begin{equation}
 D_0:C^0\longrightarrow C^1,
 \qquad D_1:C^1\longrightarrow C^2,
 \qquad D_1D_0=0.
\label{eq:v3v64-cochain}
\end{equation}
Choose reflection-invariant positive cell weights and denote adjoints by a
star.  The absolute subspaces on \(X_h\) are obtained from the following
parities on the double:
\begin{equation}
 c\circ r=c,
 \qquad A_{\mathrm{tan}}\circ r=A_{\mathrm{tan}},
 \qquad A_{\mathrm{nor}}\circ r=-A_{\mathrm{nor}}.
\label{eq:v3v64-absolute-parity}
\end{equation}
Thus scalar ghosts have discrete Neumann parity, while normal gauge links
vanish on the fixed face.  Because the cellular differential commutes with
pullback, \(D_0\) maps the even scalar subspace into the absolute one-form
subspace and \(D_1D_0=0\) remains exact at the boundary.

Remove \(\ker D_0\) by the based, equivalently mean-zero on each connected
component, convention already used in V58.  Put
\begin{equation}
 \Delta_0=D_0^*D_0,
 \qquad E^1=\operatorname{ran}D_0,
 \qquad T^1=\ker D_0^*.
\label{eq:v3v64-hodge-data}
\end{equation}
Finite-dimensional orthogonality gives
\(C^1=E^1\mathbin{\oplus}T^1\).

\begin{lemma}[Exact lattice isospectrality]
\label{lem:v3v64-isospectrality}
The positive spectra, with multiplicity, of \(\Delta_0\) on
\((\ker D_0)^\perp\) and of \(D_0D_0^*\) on \(E^1\) coincide.
\end{lemma}

\begin{proof}
Take a singular-value decomposition
\(D_0=U\Sigma V^*\) after restricting the domain to
\((\ker D_0)^\perp\) and the range to \(E^1\).  Both restricted spaces have
dimension \(r=\operatorname{rank}D_0\), and all diagonal entries
\(s_1,\ldots,s_r\) of \(\Sigma\) are positive.  The two operators are
unitarily equivalent to \(\Sigma^2\), so their eigenvalues are
\(s_j^2\) with identical multiplicities.
\end{proof}

For \(\alpha>0\), the regulated gauge operator is
\begin{equation}
 P_{\alpha,h}=D_1^*D_1+\alpha^{-1}D_0D_0^*.
\label{eq:v3v64-gauge-operator}
\end{equation}
The two summands preserve \(T^1\) and \(E^1\), respectively; the other
summand vanishes on each of those spaces.  Let
\(\Delta_{T,h}=D_1^*D_1|_{T^1}\), with any genuine harmonic zero modes
removed and recorded separately.

\begin{proposition}[V59 determinant inside the local regulator]
\label{prop:v3v64-determinant}
Let \(N_0=\dim(\ker D_0)^\perp\).  The product of the bosonic gauge
Gaussian, the complex scalar ghost integral, and the normalized auxiliary
Gaussian is
\begin{equation}
 \begin{split}
 &\det(P_{\alpha,h})^{-1/2}\,
   \det(\Delta_0)\,\alpha^{-N_0/2}\\
 &\hspace{18mm}=
   \det(\Delta_0)^{1/2}\det(\Delta_{T,h})^{-1/2},
 \end{split}
\label{eq:v3v64-determinant-product}
\end{equation}
and is independent of \(\alpha\).
\end{proposition}

\begin{proof}
On \(E^1\), Lemma~\ref{lem:v3v64-isospectrality} gives determinant
\(\alpha^{-N_0}\det\Delta_0\).  On \(T^1\), the determinant is
\(\det\Delta_{T,h}\).  Substitute this block factorization into the left
side of \eqref{eq:v3v64-determinant-product}; the powers of \(\alpha\) and
one square root of \(\det\Delta_0\) cancel exactly.
\end{proof}

\subsection{The finite ADM constraint--gauge complex}
\label{sec:v3v64-ADM}

We next give a concrete complex rather than assuming that a discretized
Einstein complex is exact.  Use a periodic spatial double for the symbol
calculation.  A forward difference has Fourier multiplier
\(h^{-1}(e^{ip_jh}-1)\).  Conjugating each cell space by harmless unitary
half-cell phases replaces it by \(i k_j\), where
\begin{equation}
 k_j=\frac{2}{h}\sin\frac{p_jh}{2}.
\label{eq:v3v64-lattice-momentum}
\end{equation}
Hence \(k=0\) only on the constant mode in the first Brillouin zone.  Fix
\(k\ne0\), and let
\begin{equation}
 \mathcal X_k=\operatorname{Sym}^2(\mathbb C^3)_q
 \oplus\operatorname{Sym}^2(\mathbb C^3)_\pi,
 \qquad \mathcal E_k=\mathbb C^3_\zeta\oplus\mathbb C_\eta .
\label{eq:v3v64-phase-space}
\end{equation}
Define the linearized spatial-diffeomorphism and lapse gauge map by
\begin{equation}
 K_k(\zeta,\eta)=
 \left(i(k_i\zeta_j+k_j\zeta_i),
 (k_i k_j-\delta_{ij}|k|^2)\eta\right),
\label{eq:v3v64-K}
\end{equation}
and define the Hamiltonian and momentum constraint map by
\begin{equation}
 C_k(q,\pi)=
 \left(k_i k_jq_{ij}-|k|^2\operatorname{tr}q,
       k_j\pi_{ij}\right)\in\mathbb C\oplus\mathbb C^3.
\label{eq:v3v64-C}
\end{equation}
Repeated spatial indices are summed.

\begin{theorem}[Exact nonzero-mode constraint quotient]
\label{thm:v3v64-ADM-complex}
For every \(k\ne0\),
\begin{equation}
 C_kK_k=0,
 \qquad \operatorname{rank}K_k=4,
 \qquad \operatorname{rank}C_k=4.
\label{eq:v3v64-ranks}
\end{equation}
Moreover
\begin{equation}
 \mathcal X_k=
 \operatorname{ran}K_k\ \mathbin{\oplus}^{\perp}
 \mathcal P_k\ \mathbin{\oplus}^{\perp}
 \operatorname{ran}C_k^*,
 \qquad
 \mathcal P_k=\ker C_k\cap\ker K_k^*,
\label{eq:v3v64-Hodge}
\end{equation}
and \(\dim_{\mathbb C}\mathcal P_k=4\).  The quotient
\(\ker C_k/\operatorname{ran}K_k\) is canonically represented by
\(\mathcal P_k\); it consists of two transverse-traceless components of
\(q\) and their two transverse-traceless momenta.
\end{theorem}

\begin{proof}
For the first component of \(C_kK_k\),
\[
 k_i k_j i(k_i\zeta_j+k_j\zeta_i)
 -|k|^2 2i k_j\zeta_j=0.
\]
For its momentum component,
\[
 k_j(k_i k_j-\delta_{ij}|k|^2)\eta=0.
\]
Rotate coordinates so that \(k=(\kappa,0,0)\), \(\kappa>0\).  The three
parameters \(\zeta_1,\zeta_2,\zeta_3\) independently produce
\(q_{11},q_{12},q_{13}\), while \(\eta\) produces the transverse trace
\(\pi_{22}+\pi_{33}\); hence \(K_k\) has rank four.  In the same frame,
\begin{equation}
 C_k(q,\pi)=
 \left(-\kappa^2(q_{22}+q_{33}),
       \kappa\pi_{11},\kappa\pi_{21},\kappa\pi_{31}\right),
\label{eq:v3v64-frame-C}
\end{equation}
so \(C_k\) has rank four.

Since \(C_kK_k=0\), \(\operatorname{ran}K_k\) is orthogonal to
\(\operatorname{ran}C_k^*\).  The orthogonal complement of their sum is
exactly \(\ker K_k^*\cap\ker C_k=\mathcal P_k\), proving
\eqref{eq:v3v64-Hodge}.  Its dimension is \(12-4-4=4\).  In the chosen
frame, \(K_k^*q=0\) removes \(q_{11},q_{12},q_{13}\), and the Hamiltonian
constraint removes \(q_{22}+q_{33}\).  Thus \(q_{22}-q_{33}\) and
\(q_{23}\) remain.  Similarly, the momentum constraint removes the first
row of \(\pi\), and orthogonality to the lapse direction removes
\(\pi_{22}+\pi_{33}\), leaving \(\pi_{22}-\pi_{33}\) and \(\pi_{23}\).
\end{proof}

The constant mode \(k=0\) is not covered by the theorem.  It contains the
finite-dimensional Killing and moduli sector and must be fixed or
integrated separately.  This is the discrete counterpart of the joint
stabilizer qualification in V60.

\subsection{Where the conformal scalar goes}
\label{sec:v3v64-conformal}

Equation \eqref{eq:v3v64-frame-C} identifies the transverse trace
\(q_{22}+q_{33}\) as the scalar normal to the Hamiltonian-constraint
surface.  It is invariant under the spatial gauge directions in
\(\operatorname{ran}K_k\), but it is not in \(\ker C_k\).  This is the
finite counterpart of the York scalar \(\Phi\) in V63.

\begin{corollary}[Constraint removal of the bad scalar]
\label{cor:v3v64-conformal-removal}
On a nonzero mode, the restriction of the gravitational Gaussian to
\(\ker C_k\), followed by quotienting \(\operatorname{ran}K_k\), contains
no conformal scalar.  Its configuration part is precisely the
two-dimensional transverse-traceless space.
\end{corollary}

\begin{proof}
Theorem~\ref{thm:v3v64-ADM-complex} identifies the quotient with
\(\mathcal P_k\).  The displayed frame calculation shows directly that
the transverse trace is absent from \(\mathcal P_k\).
\end{proof}

If one keeps the off-constraint Gaussian, its component in
\(\operatorname{ran}C_k^*\) is the finite York scalar sector.  Rotating
that component gives the ordinary phase described in V63.  Imposing the
constraint by restriction to \(\ker C_k\) removes that integration
direction; it does not cancel the phase by a ghost determinant.  These are
two distinct definitions of a regulated functional and must not be mixed.

\subsection{A common local Euclidean time lattice}
\label{sec:v3v64-time}

Let the Euclidean times be \(-M+1,\ldots,M\), reflected through the link
between \(0\) and \(1\) by \(\vartheta n=1-n\).  For a finite-dimensional
real slice space \(V\), a positive matrix \(G\), and a nonnegative matrix
\(L\), define the one-link kernel
\begin{equation}
 K_L(x,y)=\exp\left[-\frac12\langle x-y,G(x-y)\rangle
 -\frac14\langle x,Lx\rangle-\frac14\langle y,Ly\rangle\right].
\label{eq:v3v64-link-kernel}
\end{equation}
Products of this kernel over adjacent time slices define the finite
nearest-neighbour measure.  The slice space may be the physical TT space
or the cold carrier space, including the latter's finite auxiliary normal
chain.  The static carrier has no Euclidean-time derivative: it is a
positive spatial Gaussian independently on each time slice.  Its measure
therefore factorizes across the reflection plane.  Both carrier boundary
couplings are the linear trace couplings used in V61.

\begin{lemma}[Positive transfer operator]
\label{lem:v3v64-positive-transfer}
The integral operator with kernel \(K_L\) is positive and self-adjoint on
\(L^2(V,dx)\).  Consequently the finite chain satisfies
Osterwalder--Schrader reflection positivity for functions supported at
positive times.
\end{lemma}

\begin{proof}
Write \(M_L(x)=\exp[-\langle x,Lx\rangle/4]\).  The operator is
\(M_LTM_L\), where \(T\) is convolution with
\(\exp[-\langle x,Gx\rangle/2]\).  The Fourier transform of the convolution
kernel is a positive Gaussian, so \(T\ge0\), and hence \(M_LTM_L\ge0\).
It is self-adjoint because its kernel is real and symmetric.

Cut the chain at the reflection link.  Integrating all links strictly in
the positive and negative halves turns the reflection form
\(\int\overline{F(\vartheta x)}F(x)d\mu\) into
\(\langle f,K_L f\rangle\) after absorbing the remaining half-chain
kernels into \(f\).  This is nonnegative by positivity of the transfer
operator.  Iterating the same argument covers observables on arbitrarily
many positive slices.
\end{proof}

For each nonzero spatial mode, restrict the gravitational slice variable
to the transverse-traceless configuration space supplied by
Corollary~\ref{cor:v3v64-conformal-removal}; integrating its conjugate
momentum produces a kernel of the form \eqref{eq:v3v64-link-kernel} with
positive \(G\) and \(L=|k|_h^2\).  The cold carrier also has a positive
time-transfer kernel.  The static carrier contributes a factorized
positive Gaussian on the two time halves.

\begin{theorem}[Finite physical reflection positivity]
\label{thm:v3v64-physical-RP}
After removing the constant stabilizer modes and imposing \(C_k=0\) on
every spatial mode \(k\ne0\), the free TT metric modes together with the
local cold and static carriers have a reflection-positive finite
Euclidean functional.  Eliminating the carrier variables by Schur
complement gives the same finite-lattice boundary response; in the
continuum normal-family limit it converges modewise to
\begin{equation}
 \beta|\omega|+\gamma|\xi|.
\label{eq:v3v64-response}
\end{equation}
\end{theorem}

\begin{proof}
Each TT and cold mode has a positive transfer operator by
Lemma~\ref{lem:v3v64-positive-transfer}.  The static carrier measure
factorizes between positive and negative times, so its reflection form is
an absolute square after the negative variables are reflected.  Tensor
products preserve positivity, which proves the first assertion.
Gaussian elimination of a positive block gives its positive Schur
complement.  For a half-line nearest-neighbour chain in an auxiliary
extension coordinate, with lattice spacing \(a\), the decaying root
\(r_a\) satisfies
\(r_a+r_a^{-1}-2=a^2m^2\); hence
\begin{equation}
 \frac{1-r_a}{a}\longrightarrow m
 \quad(a\downarrow0).
\label{eq:v3v64-DN-limit}
\end{equation}
Taking \(m=|\omega|\) for the cold extension and \(m=|\xi|\) for the
static spatial extension gives \eqref{eq:v3v64-response}, with the
positive weights \(\beta\) and \(\gamma\).
\end{proof}

\subsection{Finite BRST Ward identity on the same grid}
\label{sec:v3v64-Ward}

On the free grid define
\begin{equation}
 sA=D_0c,
 \quad sc=0,
 \quad s\bar c=b,
 \quad sb=0,
 \qquad sh=K\epsilon,
 \quad s\epsilon=0,
 \quad s\bar\epsilon=B,
 \quad sB=0.
\label{eq:v3v64-free-BRST}
\end{equation}
Here \(K\) is the direct sum of the maps \(K_k\), with the boundary
parities induced by reflecting vectors and tensors on the double.  In
particular the normal vector ghost is odd and its tangential components
are even, as in V60.  Nilpotence is immediate, and
\(CK=0\) makes the constraint surface BRST invariant.

\begin{proposition}[Constrained finite Ward identity]
\label{prop:v3v64-Ward}
Let \(S_h=S_{\mathrm{inv}}+s\Psi\) be a quadratic action on these finite
spaces, with reflection-parity-compatible gauge fermion \(\Psi\).  Let
\(d\mu_h\) be the product of Lebesgue and Berezin measures with all
stabilizer zero modes removed.  For every polynomial \(F\),
\begin{equation}
 \int_{\ker C} sF\,e^{-S_h}\,d\mu_h=0.
\label{eq:v3v64-Ward}
\end{equation}
\end{proposition}

\begin{proof}
The finite BRST change of variables has unit Berezinian: every displayed
transformation is triangular with identity diagonal.  The action is
invariant because \(sS_{\mathrm{inv}}=0\) and \(s^2\Psi=0\).  The domain
\(\ker C\) is invariant since \(Cs h=CK\epsilon=0\).  Berezin integration
and ordinary integration by parts therefore give the integral of the
total BRST derivative \(s(Fe^{-S_h})\) as zero.  Since \(sS_h=0\), this is
exactly \eqref{eq:v3v64-Ward}.
\end{proof}

The proposition and Theorem~\ref{thm:v3v64-physical-RP} concern different
structures on the same regulator.  The Ward identity uses the full
gauge-fixed finite superintegral.  Ordinary reflection positivity is
asserted only after passage to the constrained TT observable sector and
the positive carrier variables; no positivity is assigned to Berezin
ghosts or to the rotated off-constraint York direction.

\begin{firewall}[Scope of the common regulator]
\label{fw:v3v64-scope}
This construction proves compatibility at finite cutoff for the free
simple-background theory.  It does not construct a nonlinear
diffeomorphism-invariant lattice measure, prove convergence as
\(h\downarrow0\), control the chiral determinant phase, or establish an
interacting Osterwalder--Schrader reconstruction.  Pointwise products and
the nonlinear BRST algebra do not preserve a finite Fourier subspace, and
the flat-wall reflection parity does not by itself cover curved junctions.
\end{firewall}

\subsection{V64 conclusion and V65 audit target}
\label{sec:v3v64-conclusion}

The internal gauge complex, the ADM constraint quotient, and the
cold/static carriers now coexist on one finite local time regulator.  The
singular-value pairing reproduces the V59 gauge determinant.  The exact
finite ADM calculation isolates two gravitational polarizations and
removes the conformal scalar by the Hamiltonian constraint.  A positive
transfer kernel then proves reflection positivity for those physical free
variables and yields the V61 response in the continuum normal-family
limit.

V65 must perform the scheduled NS/YM audit.  It should then attack the
first genuinely interacting obstruction: construct a regulator-stable
Slavnov defect equation and determine which local boundary counterterms
can absorb it without changing the positive physical principal symbol.
\clearpage
\section{V65: companion audit and the regulator-stable Slavnov defect}
\label{sec:v3v65}

This is the compulsory five-version audit following V60.  It is also the
first step beyond the free common regulator of V64.  The central question
is whether a finite-cutoff Slavnov defect can be removed by counterterms
that remain local as the cutoff is released and that preserve the positive
physical boundary symbol.

\subsection{Read-only NS/YM checkpoint}
\label{sec:v3v65-audit}

The companion files read for this checkpoint were
\begin{enumerate}
\item
\texttt{Renewal\_NS\_Stability\_Research\_Notes\_V31.tex}, with
\(887537\) bytes and SHA--256
\begin{center}\scriptsize\ttfamily
F1121B62238AD5491BB7CF03635EA9934942EDF573ED8FAAA9EE79E1D38A6696;
\end{center}
\item
\texttt{Renewal\_Yang\_Mills\_Mass\_Gap\_Research\_Notes\_V95.tex},
with \(1598351\) bytes and SHA--256
\begin{center}\scriptsize\ttfamily
576EA390F65781807289220449F8D9874C43841AB0A36AE832173794BB3DC7F7.
\end{center}
\end{enumerate}
Both reads were passive.  The two frozen SM files and all companion files
were left unchanged.

NS V31 proves that a flat dyadic mark can remove stopping-time
multiplicity, but exact unmarking restores the missing first moment.  A
weaker source norm merely hides the same critical weight.  The relevant
lesson here is that finite-mode solvability cannot replace an estimate in
the topology required by the continuum limit.

YM V95 proves its repeated-bank estimate by multiplying an independent
uniform norm with a capacity carrying the actual one-use stock.  It also
keeps a different-row geometric mean open when no common stock has been
proved.  The corresponding SM rule is to require one counterterm family
that is uniformly local across regulators; separate modewise primitives
are not automatically a common local primitive.

\begin{assessment}[Nonduplication decision at V65]
\label{ass:v3v65-nonduplication}
V59 already proved the finite internal determinant pairing, V60 assembled
the joint ghost kernel, and V64 put the free Ward identity and physical
transfer matrix on one grid.  V65 does not repeat those calculations.  It
studies the new implication needed for interactions: when a regulated
Slavnov defect has a local counterterm primitive, and which such
counterterms leave the V61/V64 physical principal symbol intact.
\end{assessment}

\subsection{Loopwise consistency of the defect}
\label{sec:v3v65-consistency}

Let \(\mathcal F_h\) be the finite BV algebra of fields, ghosts,
antifields, and boundary multipliers on the V64 grid.  Its odd bracket is
denoted by \((\ ,\ )_h\).  Let \(S_h\) have ghost number zero and satisfy
the finite classical master equation
\begin{equation}
 \frac12(S_h,S_h)_h=0.
\label{eq:v3v65-CME}
\end{equation}
The linearized Slavnov operator is
\begin{equation}
 \mathcal B_h X=(S_h,X)_h,
 \qquad \mathcal B_h^2=0.
\label{eq:v3v65-B}
\end{equation}
The second identity follows from the graded Jacobi identity and
\eqref{eq:v3v65-CME}.

Write a regulated effective action as
\(\Gamma_h=S_h+\sum_{j\ge1}\hbar^j\Gamma_{h,j}\).  Suppose counterterms
have restored the Slavnov identity through order \(n-1\), and write
\begin{equation}
 \frac12(\Gamma_h,\Gamma_h)_h
 =\hbar^n\mathcal A_{h,n}+O(\hbar^{n+1}).
\label{eq:v3v65-defect}
\end{equation}

\begin{theorem}[Finite Wess--Zumino consistency and restoration]
\label{thm:v3v65-WZ}
The leading defect in \eqref{eq:v3v65-defect} has ghost number one and
obeys
\begin{equation}
 \mathcal B_h\mathcal A_{h,n}=0.
\label{eq:v3v65-WZ}
\end{equation}
If \(\mathcal A_{h,n}=\mathcal B_hX_{h,n}\) for a ghost-number-zero
functional \(X_{h,n}\), then replacing
\(\Gamma_h\) by \(\Gamma_h-\hbar^nX_{h,n}\) restores the identity through
order \(n\).
\end{theorem}

\begin{proof}
For any even \(\Gamma_h\), the graded Jacobi identity gives
\[
 (\Gamma_h,(\Gamma_h,\Gamma_h)_h)_h=0.
\]
Insert \eqref{eq:v3v65-defect}.  At its first nonzero order,
\(\Gamma_h\) may be replaced by \(S_h\), which yields
\eqref{eq:v3v65-WZ}.  The counterterm changes the coefficient of
\(\hbar^n\) in one half of the master bracket by
\(- (S_h,X_{h,n})_h=-\mathcal B_hX_{h,n}\), cancelling the defect.
\end{proof}

This theorem is purely algebraic.  It does not say that \(X_{h,n}\) is
local, uniformly bounded, compatible with the boundary domain, or
convergent as \(h\downarrow0\).

\subsection{What the finite doublet argument does prove}
\label{sec:v3v65-doublet}

The standard doublet lemma is exact in the finite algebra.  We record it
because it separates harmless gauge-fixing variables from the derivative
ghost complex.

\begin{lemma}[Algebraic doublet contraction]
\label{lem:v3v65-doublet}
Suppose \(su_i=v_i\), \(sv_i=0\), and \(s\) on all other generators is
independent of \(u_i,v_i\).  Let \(N\) count total polynomial degree in
these variables.  Every \(s\)-closed polynomial of positive \(N\)-degree
is \(s\)-exact modulo its component independent of all \(u_i,v_i\).
\end{lemma}

\begin{proof}
Set \(\rho=\sum_i u_i\partial/\partial v_i\).  Direct calculation on
generators gives the graded anticommutator
\begin{equation}
 s\rho+\rho s=N.
\label{eq:v3v65-doublet-homotopy}
\end{equation}
If \(sa=0\) and \(Na=ma\) with \(m>0\), then
\(a=m^{-1}s(\rho a)\).  Decompose a general polynomial into its finitely
many \(N\)-homogeneous pieces.
\end{proof}

Thus the antighost--auxiliary pairs
\((\bar c,b)\) and \((\bar\epsilon,B)\) do not contribute to the finite
cohomology.  For nonzero modes one can also manufacture orbit doublets.  In
the internal sector set
\begin{equation}
 u_g=(D_0^*D_0)^{-1}D_0^*A,
 \qquad s u_g=c,
\label{eq:v3v65-gauge-orbit-homotopy}
\end{equation}
and in the gravitational sector set
\begin{equation}
 u_d=(K^*K)^{-1}K^*h,
 \qquad s u_d=\epsilon,
\label{eq:v3v65-diffeo-orbit-homotopy}
\end{equation}
after restricting away from the stabilizer kernels.  Equations
\eqref{eq:v3v65-gauge-orbit-homotopy} and
\eqref{eq:v3v65-diffeo-orbit-homotopy} prove modewise algebraic
contractibility, but their inverse elliptic operators are nonlocal.

\subsection{A sharp obstruction to turning modewise exactness into locality}
\label{sec:v3v65-locality-obstruction}

Call a translation-invariant lattice operator uniformly finite range if
its convolution kernel is supported in \(|j|\le R\), with \(R\) independent
of the number of sites.  Its Fourier symbol is then a Laurent polynomial
of degree at most \(R\).

\begin{proposition}[No uniformly local inverse of a difference]
\label{prop:v3v65-no-local-inverse}
Let \(D_N\) be the forward difference on the \(N\)-cycle,
\((D_Nf)_j=f_{j+1}-f_j\), restricted to mean-zero fields.  There is no
family of translation-invariant right inverses \(H_N\) satisfying
\begin{equation}
 H_ND_N=I-P_{0,N}
\label{eq:v3v65-right-inverse}
\end{equation}
whose range is bounded by one fixed \(R\) for arbitrarily large \(N\).
\end{proposition}

\begin{proof}
If such a family existed, the symbol of \(H_N\) would be a Laurent
polynomial \(q_N(z)\) of degree at most \(R\), and
\begin{equation}
 q_N(z)(z-1)=1
\label{eq:v3v65-root-identity}
\end{equation}
at every nontrivial \(N\)-th root of unity.  The coefficients may depend on
\(N\), so first choose \(2R+2\) fixed arcs separated from \(z=1\) and, for
each sufficiently large \(N\), choose one root in each arc.  Interpolation
of Laurent polynomials of degree \(R\) from \(2R+1\) of those points shows
that their coefficients remain bounded, because the limiting Vandermonde
matrix is invertible and the values \((z-1)^{-1}\) are bounded on the
chosen arcs.  Pass to a coefficientwise convergent subsequence
\(q_N\to q\).

Every point of the unit circle other than \(1\) is a limit of nontrivial
roots.  Equation \eqref{eq:v3v65-root-identity} therefore gives
\(q(z)(z-1)=1\) on infinitely many points, hence as a Laurent-polynomial
identity.  Evaluating after multiplication by a suitable power of \(z\))
at \(z=1\) gives \(0=1\), a contradiction.
\end{proof}

\begin{corollary}[Finite BRST exactness is not a local restoration theorem]
\label{cor:v3v65-finite-not-local}
The primitives constructed from
\eqref{eq:v3v65-gauge-orbit-homotopy} or
\eqref{eq:v3v65-diffeo-orbit-homotopy} cannot, solely by that
construction, yield a support-uniform local counterterm family.  A
continuum restoration argument must compute local BRST cohomology or give
a separate uniform locality estimate.
\end{corollary}

\begin{proof}
On any lattice line the orbit homotopy contains a right inverse of the
first difference after projection off constants.  Proposition
\ref{prop:v3v65-no-local-inverse} excludes a uniformly finite-range family.
The same obstruction is present in the vector complex after restriction
to a longitudinal line.
\end{proof}

This is the exact analogue of the two companion warnings.  The inverse
difference is the missing first-moment cost in NS language; a separately
chosen primitive on every mode is the different-row estimate in YM
language.  Neither supplies the one uniformly local object required here.

\subsection{Counterterms that preserve the physical boundary symbol}
\label{sec:v3v65-symbol-protection}

The locality problem is separate from principal-symbol stability.  The
latter admits a clean sufficient criterion.  At a nonzero boundary
covector \(p\), let \(P(p)\) be the orthogonal projector onto the V64
physical space
\begin{equation}
 \mathcal P(p)=\ker C(p)\cap\ker K(p)^*.
\label{eq:v3v65-physical-projector}
\end{equation}
Let \(Q_1(p)\) be the Hermitian physical boundary symbol including the
cold/static response, and suppose
\begin{equation}
 P Q_1 P\ge c|p|P
 \qquad(c>0).
\label{eq:v3v65-original-gap}
\end{equation}

\begin{theorem}[Protected counterterm class]
\label{thm:v3v65-protected-class}
Suppose a quadratic boundary counterterm has symbol
\begin{equation}
 R(p)=K(p)M(p)K(p)^*+C(p)^*N(p)C(p)+R_0(p)+A(p),
\label{eq:v3v65-safe-decomposition}
\end{equation}
where
\begin{equation}
 \|P R_0(p)P\|\le M_0,
 \qquad \operatorname{Re}\bigl(PA(p)P\bigr)=0.
\label{eq:v3v65-safe-bounds}
\end{equation}
Then on the physical quotient
\begin{equation}
 \operatorname{Re}P(Q_1+R)P
 \ge(c|p|-M_0)P.
\label{eq:v3v65-protected-gap}
\end{equation}
In particular the order-one complementing gap is unchanged at high
frequency, and is at least \(c|p|/2\) for \(|p|\ge2M_0/c\).
\end{theorem}

\begin{proof}
By definition of \(P\), \(K^*P=0\) and \(CP=0\).  The first two terms in
\eqref{eq:v3v65-safe-decomposition} therefore vanish after compression by
\(P\).  The last term has zero Hermitian real part by
\eqref{eq:v3v65-safe-bounds}.  For \(v=Pv\),
\[
 \operatorname{Re}\langle v,(Q_1+R)v\rangle
 \ge c|p|\|v\|^2-M_0\|v\|^2,
\]
which proves both assertions.
\end{proof}

\begin{corollary}[Relative order-one version]
\label{cor:v3v65-relative-order-one}
If instead
\(\|\operatorname{Re}PR_1(p)P\|\le\varepsilon|p|\) with
\(0\le\varepsilon<c\), then the new physical gap is
\((c-\varepsilon)|p|\).
\end{corollary}

\begin{proof}
Repeat the last quadratic-form estimate with
\(M_0\) replaced by \(\varepsilon|p|\).
\end{proof}

Orbit gauge-fixing counterterms and squared-constraint counterterms occupy
the first two summands of \eqref{eq:v3v65-safe-decomposition}.  Bounded
mass/potential renormalizations occupy \(R_0\).  A purely skew contribution
to the physical Euclidean symbol occupies \(A\).  Low frequencies still
require finitely many renormalization conditions; the theorem does not
make their sign automatic.

\subsection{The remaining anomaly and locality gate}
\label{sec:v3v65-gate}

The V59 Standard Model charge arithmetic removes the perturbative bulk
gauge and mixed gauge--gravitational anomaly polynomial generation by
generation, and its even number of \(SU(2)\) doublets removes the familiar
mod-two obstruction.  Those calculations do not yet orient the chiral
determinant line for the local material-frame boundary projector.  Nor do
they prove that every boundary cocycle satisfying
\eqref{eq:v3v65-WZ} has a local primitive.

\begin{firewall}[No inference from algebraic to local cohomology]
\label{fw:v3v65-no-local-inference}
The finite grid has no stabilizer-free orbit cohomology after use of the
inverse homotopies, but this does not prove vanishing of local BRST
cohomology.  Proposition~\ref{prop:v3v65-no-local-inverse} gives an explicit
reason.  No all-loop Slavnov restoration, chiral phase, or cutoff-removal
claim follows from Theorem~\ref{thm:v3v65-WZ}.
\end{firewall}

\subsection{V65 conclusion and V66 target}
\label{sec:v3v65-conclusion}

The first interacting problem is now split into two precise tests.  A
leading loop defect is a closed class of the linearized Slavnov operator,
and it is removable exactly when it has a counterterm primitive.  A
modewise primitive is insufficient: inverse differences have no uniformly
finite-range realization.  Independently, Theorem
\ref{thm:v3v65-protected-class} identifies counterterms that preserve the
positive physical principal symbol.

V66 should compute the local, infinitesimal boundary cohomology relevant to
the chiral fermion phase by descent on a collar.  It must distinguish
small-gauge Slavnov anomalies from the global determinant-line holonomy,
insert the actual Standard Model representations, and retain only
counterterms belonging to the protected class above.  If a local primitive
requires an inverse tangential derivative, the argument must move upstream
to a local bulk--collar inflow construction rather than use the nonlocal
homotopy.
\clearpage
\section{V66: local anomaly descent and the flat determinant-line remainder}
\label{sec:v3v66}

V65 showed that modewise BRST homotopies are nonlocal.  The correct
upstream question is therefore whether the chiral Standard Model defect
cancels at the level of its local descent polynomial, before any inverse
differential is introduced.  This note performs that calculation and then
separates it from the global phase of the boundary determinant line.

\subsection{One-generation anomaly arithmetic in left-handed variables}
\label{sec:v3v66-arithmetic}

Use the left-handed fermion list
\begin{equation}
 \begin{array}{c|ccccc}
  &Q&u^c&d^c&L&e^c\\ \hline
 SU(3)&\mathbf3&\overline{\mathbf3}&\overline{\mathbf3}&\mathbf1&\mathbf1\\
 SU(2)&\mathbf2&\mathbf1&\mathbf1&\mathbf2&\mathbf1\\
 Y&1/6&-2/3&1/3&-1/2&1.
 \end{array}
\label{eq:v3v66-SM-list}
\end{equation}
A sterile \(\nu^c\), if present, has \(Y=0\) and changes none of the
following coefficients.  Normalize
\(T(\mathbf3)=T(\mathbf2)=1/2\) and the cubic \(SU(3)\) index by
\(A(\mathbf3)=1\), \(A(\overline{\mathbf3})=-1\).

\begin{proposition}[Coefficientwise cancellation]
\label{prop:v3v66-coefficient-cancellation}
For one Standard Model generation, the perturbative coefficients satisfy
\begin{align}
 \mathcal A_{SU(3)^3}&=2-1-1=0,
\label{eq:v3v66-SU3-cubic}\\
 \mathcal A_{SU(3)^2Y}
 &=2\left(\frac16\right)\frac12
   +\left(-\frac23\right)\frac12
   +\left(\frac13\right)\frac12=0,
\label{eq:v3v66-SU3Y}\\
 \mathcal A_{SU(2)^2Y}
 &=3\left(\frac16\right)\frac12
   +\left(-\frac12\right)\frac12=0,
\label{eq:v3v66-SU2Y}\\
 \mathcal A_{Y^3}
 &=6\left(\frac16\right)^3
 +3\left(-\frac23\right)^3
 +3\left(\frac13\right)^3
 +2\left(-\frac12\right)^3+1=0,
\label{eq:v3v66-Y3}\\
 \mathcal A_{Y\text{-}\mathrm{grav}^2}
 &=6\left(\frac16\right)
 +3\left(-\frac23\right)
 +3\left(\frac13\right)
 +2\left(-\frac12\right)+1=0.
\label{eq:v3v66-Ygrav}
\end{align}
There is no perturbative \(SU(2)^3\) coefficient, and the Witten
\(SU(2)\) sign is trivial because there are
\begin{equation}
 3\cdot Q+L=4
\label{eq:v3v66-four-doublets}
\end{equation}
left-handed doublets per generation.
\end{proposition}

\begin{proof}
Equations \eqref{eq:v3v66-SU3-cubic}--\eqref{eq:v3v66-Ygrav} are direct
sums over the multiplicities in \eqref{eq:v3v66-SM-list}.  For example,
with denominator \(36\), the numerator in \eqref{eq:v3v66-Y3} is
\(1-32+4-9+36=0\).  The fundamental of \(SU(2)\) is pseudoreal, so its
symmetric cubic trace vanishes.  A large transformation in the nontrivial
class of \(\pi_4(SU(2))\cong\mathbb Z_2\) contributes one sign per doublet;
four copies give \((-1)^4=1\).
\end{proof}

The same calculation holds for every generation separately.  Yukawa
couplings and masses do not alter these representation indices.

\subsection{Cancellation before taking a BRST primitive}
\label{sec:v3v66-descent}

For a left-handed representation \(R\), the degree-six part of
\(\widehat A(TM)\operatorname{ch}_R(F)\) has the schematic normalized form
\begin{equation}
 I_6(R)=a\,\operatorname{tr}_R(F^3)
       +b\,p_1(TM)\operatorname{tr}_R(F),
\label{eq:v3v66-I6}
\end{equation}
where the nonzero universal constants \(a,b\) depend only on curvature
conventions.  Mixed product-group terms are obtained by expanding the
trace in \eqref{eq:v3v66-I6}.  Locally one has the descent equations
\begin{equation}
 I_6=dQ_5,
 \qquad sQ_5+dQ_4^{,1}=0,
 \qquad sQ_4^{,1}+dQ_3^{,2}=0.
\label{eq:v3v66-descent}
\end{equation}

\begin{theorem}[Vanishing of the perturbative bulk descent class]
\label{thm:v3v66-bulk-descent-zero}
For the direct sum of one Standard Model generation, the invariant
polynomial \(I_6^{\mathrm{SM}}\) vanishes coefficientwise.  Consequently a
common local descent scheme may be chosen with
\begin{equation}
 Q_{5}^{\mathrm{SM}}=Q_{4}^{1,\mathrm{SM}}
 =Q_{3}^{2,\mathrm{SM}}=0.
\label{eq:v3v66-zero-descent}
\end{equation}
No inverse tangential derivative is used.
\end{theorem}

\begin{proof}
The independent invariant monomials in the expansion of
\eqref{eq:v3v66-I6} are precisely the cubic nonabelian term, the two
nonabelian-square--hypercharge terms, the hypercharge cube, and the
mixed hypercharge--Pontryagin term.  Their coefficients vanish by
Proposition~\ref{prop:v3v66-coefficient-cancellation}; the perturbative
\(SU(2)^3\) polynomial vanishes identically.  Hence the sum of the
polynomials is the zero differential form for arbitrary gauge and metric
backgrounds.  Descent is linear in the invariant polynomial, so applying
one common transgression prescription to the sum gives
\eqref{eq:v3v66-zero-descent}.
\end{proof}

\begin{corollary}[The V65 inverse-difference obstruction is bypassed for
the bulk anomaly]
\label{cor:v3v66-local-bypass}
The perturbative bulk chiral defect does not require the nonlocal orbit
homotopies \eqref{eq:v3v65-gauge-orbit-homotopy} or
\eqref{eq:v3v65-diffeo-orbit-homotopy}.  It cancels in the local
representation sum before a Slavnov primitive is sought.
\end{corollary}

\begin{proof}
The candidate ghost-number-one local density is the descent form
\(Q_4^{1,\mathrm{SM}}\), which is zero in the common scheme by
Theorem~\ref{thm:v3v66-bulk-descent-zero}.
\end{proof}

This removes the local bulk anomaly class.  A boundary condition can still
carry a relative eta phase, so it remains necessary to understand the
determinant line rather than declare the full phase trivial.

\subsection{Local eta variation versus global spectral flow}
\label{sec:v3v66-eta}

Let \(D_z\) be a smooth family of self-adjoint elliptic boundary operators
obtained from the collar reduction of the Euclidean chiral operator and the
material-frame projector.  On a parameter chart where \(D_z\) is
invertible, define the eta function by
\begin{equation}
 \eta(D_z;s)=\frac1{\Gamma((s+1)/2)}
 \int_0^\infty t^{(s-1)/2}
 \operatorname{Tr}(D_z e^{-tD_z^2})\,dt
\label{eq:v3v66-eta-integral}
\end{equation}
for large \(\operatorname{Re}s\), followed by meromorphic continuation to
\(s=0\).

\begin{proposition}[The local/global split of the eta phase]
\label{prop:v3v66-eta-split}
Along an invertible path \(z(t)\), the derivative of
\(\eta(D_{z(t)};0)\) is the integral of a local heat-kernel coefficient.
For a closed path, the continuation of
\begin{equation}
 \exp\left(-\frac{i\pi}{2}\eta(D_z;0)\right)
\label{eq:v3v66-eta-phase}
\end{equation}
may additionally acquire the spectral-flow sign
\((-1)^{\operatorname{SF}(D_{z(t)})}\).  Thus a local counterterm controls
the infinitesimal variation on a contractible invertible chart, while
spectral flow is a global datum.
\end{proposition}

\begin{proof}
Differentiate \eqref{eq:v3v66-eta-integral}.  Duhamel's formula for the
derivative of \(e^{-tD^2}\), cyclicity of the trace, and one integration by
parts express the derivative at \(s=0\) as the constant term in the
small-\(t\) expansion of a trace containing \(\dot D\).  Heat-kernel
coefficients are integrals of local polynomials in the jets of the operator
and boundary projector.  Invertibility removes a large-\(t\) kernel term.

When an eigenvalue crosses zero, the eta invariant jumps by twice the
signed crossing multiplicity.  The phase
\eqref{eq:v3v66-eta-phase} therefore changes by \(-1\) for each unit of
spectral flow.  Integrating the local derivative accounts for the smooth
part; the product of crossing signs gives the asserted factor.
\end{proof}

The proposition explains why perturbative anomaly cancellation does not
by itself settle the boundary phase: it sets the curvature/local variation
class to zero, but a flat line may still have nontrivial holonomy.

\subsection{Flat determinant-line reduction}
\label{sec:v3v66-flat-line}

Let \(\mathscr P\) be the parameter space of admissible gauge, metric, and
material-frame backgrounds modulo based transformations, restricted to a
component on which the boundary problem is Fredholm.  Denote the total
fermion determinant line by \(\mathscr L_{\mathrm{SM}}\to\mathscr P\).

\begin{theorem}[Local cancellation reduces the phase problem to holonomy]
\label{thm:v3v66-flat-line}
Assume the collar family index formula applies to the V56 local projector,
so that the curvature of the determinant-line connection is represented
by the sum of the bulk polynomial \(I_6\) and its local boundary
transgression.  After subtracting the local boundary transgression on an
invertible chart, Proposition
\ref{prop:v3v66-coefficient-cancellation} makes the Standard Model
determinant connection flat.  It admits a parallel nonvanishing section on
every simply connected chart.  A global section exists if and only if the
holonomy representation
\begin{equation}
 \operatorname{Hol}:\pi_1(\mathscr P)\longrightarrow U(1)
\label{eq:v3v66-holonomy}
\end{equation}
is trivial.
\end{theorem}

\begin{proof}
Under the stated family-index hypothesis, Theorem
\ref{thm:v3v66-bulk-descent-zero} removes the bulk curvature class and the
subtracted transgression removes the exact boundary contribution.  Hence
the connection is flat.  Parallel transport from a base point is
path-independent on a simply connected chart and supplies a nonzero
parallel section.  On the full component, path independence is equivalent
to unit holonomy around every loop, which is exactly triviality of
\eqref{eq:v3v66-holonomy}.
\end{proof}

The theorem has an explicit hypothesis because the complementing analysis
of V56--V61 does not yet prove the family index formula for this
field-dependent material projector.  Establishing that formula is part of
the remaining analytic work.

\begin{proposition}[Traditional \(SU(2)\) global loop cancels]
\label{prop:v3v66-Witten-loop}
For the loop generated by the nontrivial class of
\(\pi_4(SU(2))\), the Standard Model fermion phase has unit holonomy in
every generation.
\end{proposition}

\begin{proof}
The mod-two index, equivalently spectral-flow parity, is additive over
Weyl doublets.  Equation \eqref{eq:v3v66-four-doublets} gives four
doublets, so the total parity is zero and the sign is \(+1\).
\end{proof}

This proposition covers one distinguished loop.  It does not classify
loops involving the material frame, the boundary projector, spacetime
diffeomorphisms, or possible quotients of the global Standard Model gauge
group.

\subsection{Compatibility with the protected physical symbol}
\label{sec:v3v66-protection}

Local counterterms that change only gauge fixing or impose the constraints
have the first two forms in
\eqref{eq:v3v65-safe-decomposition}.  A parity-odd local transgression has
purely imaginary Euclidean quadratic form on real fields and hence zero
Hermitian real part.  Therefore, whenever its differential order does not
raise the boundary problem above the V61 order-one class, it belongs to the
\(A\)-sector of Theorem~\ref{thm:v3v65-protected-class}.  The positive real
physical gap is then unchanged.

Higher-derivative parity-odd terms, including a gravitational
Chern--Simons representative, can change the differential order even when
their real quadratic part vanishes.  They must be fixed by a
renormalization condition or incorporated into a new complementing
analysis; V65's real-part estimate alone is not enough.

\begin{firewall}[What has and has not been cancelled]
\label{fw:v3v66-boundary}
V66 proves coefficientwise cancellation of the local four-dimensional
Standard Model anomaly polynomial and the traditional \(SU(2)\) mod-two
sign.  It does not prove the collar family-index hypothesis for the
material projector, trivialize the full determinant line, compute all
mapping-torus holonomies, or construct the continuum chiral measure.
Flatness is not the same as a chosen global orientation.
\end{firewall}

\subsection{V66 conclusion and V67 target}
\label{sec:v3v66-conclusion}

The local perturbative anomaly is removed at its source: all independent
degree-six coefficients vanish in the Standard Model representation sum,
so no nonlocal BRST primitive is required.  The remaining phase problem is
global.  On a collar-compatible determinant family it is the holonomy of a
flat line, detected by spectral flow.  The traditional \(SU(2)\) loop is
already trivial because there are four doublets per generation.

V67 should construct the precise parameter-space loop invariant for the
V56 material-frame projector.  It must express its holonomy as a
five-dimensional mapping-torus index plus an explicit projector spectral
flow, prove homotopy invariance through complementing families, and test
whether the frame sector introduces a new mod-two class.  If the projector
family fails to remain Fredholm along the required homotopy, the analysis
must return upstream to the boundary-domain construction rather than
assign a formal eta phase.
\clearpage
\section{V67: the material-projector loop and its spectral-flow invariant}
\label{sec:v3v67}

The material vector in V56 is a unit vector in the oriented spatial
two-plane tangent to the timelike boundary.  Its pointwise parameter space
is therefore a circle.  This makes the frame part of the determinant-line
problem concrete: a closed material rotation has an integer winding, and
the eigenspace of the Weyl projector has a computable Berry sign.  The sign
must still be combined with the spectral flow of the full boundary Dirac
operator.

\subsection{Clifford lift of a material-frame rotation}
\label{sec:v3v67-frame-loop}

Fix an oriented orthonormal pair \((e_1^0,e_2^0)\) in the material spatial
plane and put
\begin{equation}
 B_1=\alpha_\chi(e_1^0),
 \qquad B_2=\alpha_\chi(e_2^0).
\label{eq:v3v67-B12}
\end{equation}
Then \(B_i^*=B_i\), \(B_i^2=I\), and
\(B_1B_2+B_2B_1=0\).  For a rotation angle \(\theta\), define
\begin{align}
 e_1(\theta)&=\cos\theta\,e_1^0+\sin\theta\,e_2^0,
\label{eq:v3v67-e-rotation}\\
 B(\theta)&=\cos\theta\,B_1+\sin\theta\,B_2,
\qquad
 U(\theta)=\exp\left(\frac\theta2B_2B_1\right).
\label{eq:v3v67-BU}
\end{align}

\begin{lemma}[Closed projector and antiperiodic spin lift]
\label{lem:v3v67-spin-lift}
For every \(\theta\),
\begin{equation}
 B(\theta)=U(\theta)B_1U(\theta)^{-1},
 \qquad U(2\pi)=-I.
\label{eq:v3v67-lift-identities}
\end{equation}
Consequently
\begin{equation}
 P_\varepsilon(\theta)=\frac12(I+\varepsilon B(\theta))
 =U(\theta)P_\varepsilon(0)U(\theta)^{-1}
\label{eq:v3v67-projector-loop}
\end{equation}
is a closed loop of local maximal-isotropic projectors, although its
canonical spin lift closes with sign \(-1\).
\end{lemma}

\begin{proof}
Set \(A=B_2B_1\).  The Clifford relations give \(A^*=-A\), \(A^2=-I\),
and \([A,B_1]=2B_2\).  Differentiating
\(U(\theta)B_1U(\theta)^{-1}\) yields the rotation equation with initial
value \(B_1\), proving the first identity.  Since \(A^2=-I\),
\(U(2\pi)=\cos\pi+ A\sin\pi=-I\).  Conjugation removes this central sign,
so the projector is periodic.
\end{proof}

For a loop \(e_1(t,x)\) based at \(e_1^0(x)\), define on each connected
boundary component
\begin{equation}
 w(e_1)=\frac1{2\pi}\int_0^1
 \langle e_2(t,x),\partial_t e_1(t,x)\rangle\,dt\in\mathbb Z.
\label{eq:v3v67-winding}
\end{equation}
The integer is independent of \(x\): it varies continuously with \(x\)
and is integer valued.

\begin{proposition}[Pointwise projector determinant holonomy]
\label{prop:v3v67-projector-Berry}
Let \(L_\varepsilon(\theta)=\operatorname{ran}P_\varepsilon(\theta)\) in
one Weyl spin fibre.  The projected Berry connection on this line has
holonomy
\begin{equation}
 \operatorname{Hol}(L_\varepsilon)=(-1)^{w(e_1)},
\label{eq:v3v67-line-holonomy}
\end{equation}
independently of \(\varepsilon\).  After tensoring with an internal
representation of complex dimension \(d_R\), the determinant of the
projector subspace has holonomy \((-1)^{d_Rw(e_1)}\).
\end{proposition}

\begin{proof}
Choose a unit vector \(v_0\in\operatorname{ran}P_\varepsilon(0)\) and set
\(v(\theta)=U(\theta)v_0\).  Its projected connection coefficient is
\begin{equation}
 \langle v,\partial_\theta v\rangle
 =\frac12\langle v_0,B_2B_1v_0\rangle=0.
\label{eq:v3v67-horizontal}
\end{equation}
Indeed \(B_2B_1\) maps the \(\varepsilon\)-eigenspace of \(B_1\) to its
orthogonal complement.  Thus \(v\) is horizontal.  Lemma
\ref{lem:v3v67-spin-lift} gives \(v(2\pi)=-v(0)\), which proves the sign
for one winding and hence for every winding.  The determinant of
\(d_R\) identical line transports is the \(d_R\)-th power.
\end{proof}

\begin{corollary}[Minimal-Standard-Model frame parity]
\label{cor:v3v67-SM-frame-parity}
The sum of internal dimensions in one minimal generation is
\begin{equation}
 6+3+3+2+1=15.
\label{eq:v3v67-fifteen}
\end{equation}
Thus the pointwise determinant of all projector subspaces has sign
\((-1)^{15N_gw}=(-1)^{N_gw}\) for \(N_g\) generations.  With three minimal
generations this candidate sign is \((-1)^w\).  Adding one sterile
right-handed neutrino per generation changes \(15\) to \(16\) and makes
this pointwise sign trivial.
\end{corollary}

\begin{proof}
Sum the dimensions in \eqref{eq:v3v66-SM-list} and apply Proposition
\ref{prop:v3v67-projector-Berry}.  The sign choice
\(\varepsilon_L=-\varepsilon_R\) does not change the result because
\eqref{eq:v3v67-line-holonomy} is independent of \(\varepsilon\).
\end{proof}

The corollary computes a finite-rank boundary-domain factor.  It is not yet
the holonomy of the infinite-dimensional fermion determinant.  Bulk
spectral flow and the regularized nonzero boundary modes can contribute a
second factor.

\subsection{Fixed-domain reduction and mapping-torus index}
\label{sec:v3v67-mapping-torus}

Let \(D_t\) be a smooth loop of self-adjoint Fredholm Dirac realizations
whose domains are given by the projectors \(P_t\).  Suppose the rank and
complementing gaps of V49 and V56 remain positive.  A smooth unitary
\(U_t\) with \(P_t=U_tP_0U_t^{-1}\) identifies the domains, giving
\begin{equation}
 \widetilde D_t=U_t^{-1}D_tU_t
\label{eq:v3v67-fixed-domain}
\end{equation}
on the fixed domain of \(P_0\).  The endpoint is glued by the unitary
\(G=U_1^{-1}U_0\), including the central sign found in
Lemma~\ref{lem:v3v67-spin-lift}.  Let
\begin{equation}
 \mathscr D_\gamma=\partial_t+\widetilde D_t
\label{eq:v3v67-suspension}
\end{equation}
on the resulting mapping torus.

\begin{theorem}[Spectral flow equals the suspension index]
\label{thm:v3v67-SF-index}
If the endpoints are invertible, then
\begin{equation}
 \operatorname{SF}(D_t)
 =\operatorname{ind}(\mathscr D_\gamma).
\label{eq:v3v67-SF-index}
\end{equation}
The right side is independent of the chosen domain trivialization and is
constant under homotopies through Fredholm complementing families with
invertible endpoints.
\end{theorem}

\begin{proof}
After a generic arbitrarily small perturbation, all zero crossings of
\(\widetilde D_t\) are simple and transverse.  Near a crossing, split off
the crossing eigenline.  On it, \(\partial_t+\widetilde D_t\) is homotopic
to \(\partial_t+\lambda(t)\), whose index is
\(\operatorname{sign}\dot\lambda\).  On the uniformly invertible
orthogonal complement it has index zero.  Summing the local indices gives
the signed crossing number, which is spectral flow.  General paths follow
by homotopy invariance of both sides.

Changing \(U_t\) conjugates the suspension by a periodic unitary after its
endpoint gluing is included, so the index is unchanged.  A Fredholm
homotopy cannot change the index; invertible endpoints prevent crossings
from escaping through the ends.
\end{proof}

When the family has a boundary, the suspension inherits a boundary on the
mapping torus of \(\mathcal N\).  The statement assumes the suspended
projector is itself complementing.  That is the family-domain hypothesis
left explicit in V66.

\subsection{Projector spectral flow as a Maslov index}
\label{sec:v3v67-Maslov}

Let \(\mathcal C_t\) be the Calderon space of boundary Cauchy data of the
interior Dirac operator, and let
\(\mathcal L_t=\operatorname{ran}P_t\) be the allowed maximal-isotropic
boundary space.  At an isolated time when
\(\mathcal C_t\cap\mathcal L_t\ne0\), the boundary problem has a zero mode.
The crossing form is the derivative of the Green pairing on this finite
intersection.

\begin{proposition}[Boundary crossing formula]
\label{prop:v3v67-Maslov}
For a path with only regular crossings,
\begin{equation}
 \operatorname{SF}(D_t)
 =\sum_{t_*}\operatorname{sign}\mathfrak q_{t_*}
 =:\operatorname{Mas}(\mathcal C_t,\mathcal L_t).
\label{eq:v3v67-Maslov}
\end{equation}
The equality extends to all Fredholm paths by homotopy.
\end{proposition}

\begin{proof}
At a regular crossing, Green's identity identifies the derivative of the
small eigenvalue of \(D_t\) with the boundary crossing form on
\(\mathcal C_t\cap\mathcal L_t\).  Its signature is therefore the local
spectral-flow contribution.  Summing gives the first equality.  Both sides
are invariant under endpoint-fixed Fredholm homotopy, which removes the
regular-crossing assumption.
\end{proof}

This formula shows why the pointwise Berry sign in Corollary
\ref{cor:v3v67-SM-frame-parity} is not by itself decisive.  The determinant
phase depends on the relative motion of the material projector and the
Calderon space, not only on the former.

\subsection{The loop invariant and its domain of definition}
\label{sec:v3v67-invariant}

For the sign part of the phase, record the pair
\begin{equation}
 \mathfrak I(\gamma)=
 \left(
  \operatorname{ind}\mathscr D_\gamma\bmod2,
  \sum_R d_Rw_R\bmod2
 \right)\in\mathbb Z_2\oplus\mathbb Z_2.
\label{eq:v3v67-invariant-pair}
\end{equation}
The first component is the regularized operator spectral flow; the second
is the explicitly computed boundary-domain transport.  A determinant-line
gluing theorem fixes how these two components combine in a chosen phase
convention.  Recording them separately prevents double counting.

\begin{theorem}[Homotopy invariance within the admissible domain]
\label{thm:v3v67-homotopy}
The pair \(\mathfrak I(\gamma)\) is constant under a homotopy of loops for
which
\begin{enumerate}
\item the clock and material rank gaps stay positive;
\item the local projectors remain maximal isotropic and complementing;
\item the suspended operators stay Fredholm; and
\item the endpoint operators remain invertible after removal of the fixed
stabilizer sector.
\end{enumerate}
\end{theorem}

\begin{proof}
The first component is constant by Theorem~\ref{thm:v3v67-SF-index}.  The
second is a sum of winding parities.  Winding is invariant under a homotopy
through nondegenerate oriented frames, which is ensured by the material
rank gap.  The other assumptions guarantee that both operator and domain
invariants remain defined throughout.
\end{proof}

\begin{firewall}[Candidate frame parity is not yet an anomaly]
\label{fw:v3v67-candidate-only}
For three minimal generations the projector subspace has an odd Berry sign
under one material-frame winding.  V67 does not identify this sign with the
full determinant-line holonomy.  It may be cancelled or modified by the
Calderon-relative spectral flow in
\eqref{eq:v3v67-Maslov}, by the mapping-torus eta invariant, or by an
allowed material Wess--Zumino term.  None of those factors may be omitted.
\end{firewall}

\begin{firewall}[Gap failure ends the invariant]
\label{fw:v3v67-gap}
If the material coframe loses rank, or if the projector ceases to be
complementing, the winding and Fredholm homotopy arguments do not extend
through that point.  Assigning a phase by continuity across such a failure
would assume the very global domain theorem still missing.
\end{firewall}

\subsection{V67 conclusion and V68 target}
\label{sec:v3v67-conclusion}

The material projector has a nontrivial loop already at the level of one
Weyl fibre: its spin lift is antiperiodic and its eigenspace Berry holonomy
is \((-1)^w\).  The minimal three-generation representation has an odd
pointwise domain sign, whereas a sterile neutrino per generation would
make that finite-rank sign even.  The full operator contribution is the
index of the mapping-torus suspension, equivalently the Maslov spectral
flow relative to the Calderon space.

V68 should compute this relative spectral flow in the frozen flat collar
model.  It must Fourier decompose the tangential Weyl operator, rotate the
material projector through one winding, identify every zero crossing and
its sign, and test whether the regularized Calderon contribution cancels
the finite-rank projector sign.  Only after that explicit model is known
should a curved-family gluing theorem be attempted.
\clearpage
\section{V68: Euclidean normal-family obstruction for the material Weyl projector}
\label{sec:v3v68}

V67 proposed computing the material-loop spectral flow in the frozen
collar.  The normal family fails before that computation can be made: the
V56 material projector intersects the decaying Euclidean Weyl Calderon
line at a nonzero tangential covector.  This note gives the explicit
intersection and the associated local single-Weyl no-go theorem.  The
Lorentzian energy results of V56--V58 are unaffected.

\subsection{Flat Euclidean Weyl normal family}
\label{sec:v3v68-normal-family}

Work on the half-space \(x_3\ge0\), with Euclidean tangential coordinates
\((x_4,x_1,x_2)\).  Up to multiplication by an invertible constant matrix,
take the positive-chirality Weyl operator to be
\begin{equation}
 D_+=\partial_4+i\sigma_1\partial_1
              +i\sigma_2\partial_2+i\sigma_3\partial_3.
\label{eq:v3v68-Weyl}
\end{equation}
For a tangential covector \(p=(p_4,p_1,p_2)\ne0\), seek a decaying normal
mode
\begin{equation}
 \psi=e^{i(p_4x_4+p_1x_1+p_2x_2)}e^{-\rho x_3}v,
 \qquad \rho>0.
\label{eq:v3v68-decaying-ansatz}
\end{equation}
The normal equation is
\begin{equation}
 \left(ip_4I-p_1\sigma_1-p_2\sigma_2-i\rho\sigma_3\right)v=0.
\label{eq:v3v68-normal-equation}
\end{equation}
Its determinant is \(\rho^2-|p|^2\), so the unique decaying root is
\(\rho=|p|\), and its kernel is one complex line
\(\mathcal C_+(p)\subset\mathbb C^2\).

Set \(p_4=0\), write
\((p_1,p_2)=k(\cos\phi,\sin\phi)\) with \(k>0\), and take
\(\rho=k\).  Then
\begin{equation}
 v_\phi=\frac1{\sqrt2}
 \begin{pmatrix}1\\-i e^{i\phi}\end{pmatrix}
\label{eq:v3v68-Calderon-vector}
\end{equation}
spans \(\mathcal C_+(0,k\cos\phi,k\sin\phi)\).

The material projector at angle \(\theta\) is
\begin{equation}
 B(\theta)=\cos\theta\,\sigma_1+\sin\theta\,\sigma_2,
 \qquad B(\theta)v=\varepsilon v.
\label{eq:v3v68-material-projector}
\end{equation}
Its allowed line is spanned by
\begin{equation}
 w_{\varepsilon,\theta}=\frac1{\sqrt2}
 \begin{pmatrix}1\\\varepsilon e^{i\theta}\end{pmatrix}.
\label{eq:v3v68-material-vector}
\end{equation}

\begin{theorem}[Explicit complementing failure]
\label{thm:v3v68-explicit-failure}
For every material angle \(\theta\) and either projector sign
\(\varepsilon\), choose
\begin{equation}
 \phi=\theta+\varepsilon\frac\pi2\pmod{2\pi}.
\label{eq:v3v68-bad-direction}
\end{equation}
Then \(v_\phi=w_{\varepsilon,\theta}\) up to a nonzero scalar.  Hence the
homogeneous Euclidean normal family has a nonzero decaying solution
satisfying the V56 boundary condition.  The Lopatinski--Shapiro
complementing condition fails at every boundary point.
\end{theorem}

\begin{proof}
Equation \eqref{eq:v3v68-bad-direction} gives
\[
 -i e^{i\phi}
 =-i e^{i\theta}e^{i\varepsilon\pi/2}
 =\varepsilon e^{i\theta}.
\]
The two vectors \eqref{eq:v3v68-Calderon-vector} and
\eqref{eq:v3v68-material-vector} therefore coincide.  The ansatz
\eqref{eq:v3v68-decaying-ansatz} with \(\rho=k>0\) is a nonzero decaying
solution at the nonzero tangential covector
\((0,k\cos\phi,k\sin\phi)\).  Complementing requires the intersection of
the decaying Calderon line with the allowed boundary kernel to be zero for
every nonzero tangential covector, so it fails.
\end{proof}

The bad direction is spatial and perpendicular, with orientation fixed by
\(\varepsilon\), to the selected material vector.  It is therefore not a
zero-frequency stabilizer that can be removed by a based convention.

\subsection{The symbol-level single-Weyl no-go theorem}
\label{sec:v3v68-no-go}

The preceding intersection is not peculiar to the equatorial material
family.  The full Calderon-line map is the Hopf line over the tangential
cosphere.

\begin{lemma}[Surjectivity of the Weyl Calderon map]
\label{lem:v3v68-Calderon-surjective}
The map
\begin{equation}
 S^2\ni \widehat p\longmapsto
 \mathcal C_+(\widehat p)\in\mathbb{CP}^1
\label{eq:v3v68-Calderon-map}
\end{equation}
defined by \eqref{eq:v3v68-normal-equation} is a bijection, up to the
orientation convention for \(S^2\).  Its line bundle has first Chern
number \(\pm1\).
\end{lemma}

\begin{proof}
On the chart \(p_1-ip_2\ne0\), the first row of
\eqref{eq:v3v68-normal-equation} gives the projective coordinate
\begin{equation}
 \frac{v_2}{v_1}
 =\frac{i(p_4-|p|)}{p_1-ip_2}.
\label{eq:v3v68-stereographic}
\end{equation}
After division by \(|p|\), this is stereographic projection from the unit
two-sphere to the Riemann sphere.  The omitted pole is obtained on the
second projective chart.  Thus the map is bijective.  The transition
function between the two normalized eigensections winds once around the
equator, so its first Chern number is \(1\) or \(-1\), depending on the
normal and chirality conventions.
\end{proof}

\begin{theorem}[No covector-independent local rank-one condition for one
Euclidean Weyl field]
\label{thm:v3v68-local-Weyl-no-go}
At a fixed boundary point, let \(L\subset\mathbb C^2\) be any complex line
chosen independently of the tangential covector.  There exists
\(p\ne0\) such that
\begin{equation}
 L=\mathcal C_+(p).
\label{eq:v3v68-line-hit}
\end{equation}
Consequently no local zeroth-order rank-one projector on a single
Euclidean Weyl field is complementing for all nonzero tangential
covectors.
\end{theorem}

\begin{proof}
The line \(L\) is one point of \(\mathbb{CP}^1\).  Lemma
\ref{lem:v3v68-Calderon-surjective} supplies a tangential direction whose
Calderon line is that point.  At this covector the allowed line intersects
the decaying line nontrivially, contradicting the complementing condition.
\end{proof}

This elementary normal-family statement is the flat representative of the
Atiyah--Bott obstruction.  A covector-dependent spectral projector can
avoid the intersection, and a full Dirac pair can use opposite chiral
symbols, but a fixed local line for one chirality cannot.

\subsection{Gauge-equivariant collection obstruction}
\label{sec:v3v68-representation}

At high frequency, Yukawa and mass matrices are lower order.  A local
elliptic completion must therefore cancel the Weyl symbol obstruction in
each gauge representation before symmetry breaking is used.

\begin{proposition}[Representationwise necessity of a mirror symbol]
\label{prop:v3v68-mirror-necessity}
Suppose a local zeroth-order boundary condition is required to commute with
the unbroken Standard Model gauge group.  For every irreducible gauge
representation \(R\), a positive-chirality Weyl principal symbol in \(R\)
must be paired with an opposite-chirality symbol in an isomorphic copy of
\(R\), or else its nontrivial Calderon class remains in that isotypic
component.  Lower-order Higgs/Yukawa maps cannot remove this obstruction.
\end{proposition}

\begin{proof}
Gauge equivariance makes the boundary symbol block diagonal on isotypic
components and permits principal mixing only between equivalent
representations.  By Lemma~\ref{lem:v3v68-Calderon-surjective}, one Weyl
block contributes a line bundle of Chern number \(\pm1\) over the
tangential cosphere; the opposite chirality contributes the inverse class.
A covector-independent local complement can exist only after these classes
cancel inside the same gauge isotypic component.  Algebraic Yukawa and
mass terms do not enter the principal normal symbol, so they cannot change
its Chern class.
\end{proof}

The chiral Standard Model does not contain such an opposite-chirality copy
in every unbroken representation.  The Lorentzian material projector
avoids the need for a gauge intertwiner because maximal dissipativity is a
different condition from Euclidean ellipticity; Theorem
\ref{thm:v3v68-explicit-failure} shows that the distinction is essential.

\subsection{Correction of the quantum route}
\label{sec:v3v68-correction}

\begin{firewall}[The V56 domain is Lorentzian, not Euclidean elliptic]
\label{fw:v3v68-Lorentzian-only}
Theorems in V56--V58 establishing maximal isotropy, zero normal current,
the energy estimate, the material reaction force, and classical BRST
domain invariance remain valid for the Lorentzian initial-boundary problem.
They must not be cited as an elliptic Euclidean boundary theorem.
\end{firewall}

\begin{firewall}[Status of V59 and V66--V67]
\label{fw:v3v68-quantum-correction}
The gauge/ghost determinants of V59 and the bosonic Euclidean results of
V60--V65 are unaffected.  A Euclidean Weyl determinant on the V56 local
domain has not been defined.  The determinant-line conclusions of V66 and
the mapping-torus/spectral-flow construction of V67 were explicitly
conditional on a complementing collar family; Theorem
\ref{thm:v3v68-explicit-failure} shows that this hypothesis fails for the
frozen material projector.  Their abstract anomaly and loop calculations
may be reused only after an elliptic fermion completion is supplied.
\end{firewall}

Three mathematically distinct repairs remain:
\begin{enumerate}
\item use the nonlocal Atiyah--Patodi--Singer tangential spectral
projector;
\item add opposite-chirality mirror or collar regulator fields in each
gauge representation, producing a local full-Dirac principal symbol; or
\item avoid a Euclidean fermion determinant and construct the quantum
theory directly from the Lorentzian Hamiltonian domain.
\end{enumerate}
The first threatens Euclidean time locality and reflection positivity; the
second requires a decoupling and anomaly-inflow proof; the third requires a
nonperturbative Lorentzian construction not presently available.

\subsection{V68 conclusion and V69 target}
\label{sec:v3v68-conclusion}

The explicit decaying mode
\eqref{eq:v3v68-Calderon-vector} proves that the material Weyl projector is
not Euclidean complementing.  More generally, the Weyl Calderon line maps
the tangential cosphere bijectively to \(\mathbb{CP}^1\), so every fixed
local rank-one line is hit.  Gauge equivariance forces cancellation of
this symbol obstruction representation by representation, which the
chiral Standard Model cannot provide without an auxiliary completion.

V69 should compare the two Euclidean repairs at the principal-symbol
level.  It should prove that the APS projector is elliptic but nonlocal in
Euclidean time, then construct the smallest representationwise mirror
collar block with a local complementing bag projector.  The decisive next
test is whether its mirror determinant can be divided out or driven to
infinite mass while leaving the V66 anomaly cancellation, V65 protected
physical symbol, and V64 reflection-positive bosonic sector intact.
\clearpage
\section{V69: APS nonlocality and the local mirror-collar alternative}
\label{sec:v3v69}

V68 leaves two Euclidean routes worth testing.  The APS route chooses the
Calderon complement directly and is elliptic, but is pseudodifferential in
the tangential variables.  The mirror route cancels the Weyl symbol class
with an opposite chirality in the same gauge representation and permits a
local bag projector.  This note proves both principal-symbol statements
and identifies the obstruction to decoupling a finite four-dimensional
mirror by an ordinary mass.

\subsection{The APS complement is elliptic and time-nonlocal}
\label{sec:v3v69-APS}

Let \(A(\widehat p)\) be the Hermitian Calderon involution of the flat Weyl
normal family, normalized so that the decaying data are
\begin{equation}
 \mathcal C_+(p)=\ker(A(\widehat p)-I),
 \qquad \widehat p=p/|p|.
\label{eq:v3v69-Calderon-involution}
\end{equation}
The principal APS boundary projector kills this line, or equivalently
allows
\begin{equation}
 \mathcal L_{\mathrm{APS}}(p)=\ker(A(\widehat p)+I).
\label{eq:v3v69-APS-line}
\end{equation}

\begin{proposition}[APS principal complementing condition]
\label{prop:v3v69-APS-complementing}
For every \(p\ne0\),
\begin{equation}
 \mathcal C_+(p)\cap\mathcal L_{\mathrm{APS}}(p)=\{0\}.
\label{eq:v3v69-APS-intersection}
\end{equation}
Thus the APS principal boundary symbol is complementing for the Euclidean
Weyl operator.
\end{proposition}

\begin{proof}
An involution has disjoint \(+1\) and \(-1\) eigenspaces.  Equations
\eqref{eq:v3v69-Calderon-involution} and
\eqref{eq:v3v69-APS-line} identify the two spaces with those eigenspaces.
\end{proof}

\begin{proposition}[APS is nonlocal in Euclidean time]
\label{prop:v3v69-APS-time-nonlocal}
For a boundary that is timelike before Wick rotation, the tangential
covector \(p\) contains the Euclidean frequency \(p_4\).  The symbol
\begin{equation}
 P_{\mathrm{APS}}(p)=\frac12(I-A(p/|p|))
\label{eq:v3v69-APS-symbol}
\end{equation}
is not the symbol of a finite-order differential operator and is not
finite range on the V64 time lattice.  In particular it does not define a
one-link transfer kernel of the form \eqref{eq:v3v64-link-kernel}.
\end{proposition}

\begin{proof}
The entries of \(A\) contain ratios \(p_a/|p|\).  A finite-order
differential operator has a polynomial symbol, whereas a uniformly
finite-range lattice operator has a trigonometric-polynomial symbol.
Neither can equal \eqref{eq:v3v69-APS-symbol} on all nonzero covectors:
restriction to \((p_4,p_1,p_2)=(s,1,0)\) contains
\((1+s^2)^{-1/2}\), which is neither polynomial nor a fixed
trigonometric polynomial.  Dependence on \(p_4/|p|\) therefore couples
arbitrarily separated Euclidean time slices.  The link-factorization
hypothesis used in Lemma~\ref{lem:v3v64-positive-transfer} is absent.
\end{proof}

The proposition does not prove that every APS functional violates
reflection positivity.  It proves that the local transfer-matrix proof of
V64 cannot be applied.

\subsection{A local full-Dirac completion}
\label{sec:v3v69-bag}

Add to a physical Weyl field in representation \(R\) an auxiliary Weyl
field of opposite chirality in the same representation.  Their direct sum
is a Euclidean Dirac spinor \(\Psi\in(S_+\oplus S_-)\otimes R\).  Let
\(\gamma_n\) be Clifford multiplication by the boundary normal and let
\(\gamma_5\) be chirality.  Define two Hermitian involutions
\begin{equation}
 A(\widehat p)=i\gamma_n\gamma(\widehat p),
 \qquad B_{\mathrm{bag}}=i\gamma_n\gamma_5.
\label{eq:v3v69-A-B}
\end{equation}
Here \(A\) selects the decaying normal data of the full Dirac operator.

\begin{lemma}[Bag--Calderon anticommutation]
\label{lem:v3v69-anticommutation}
For every unit tangential covector \(\widehat p\),
\begin{equation}
 A^2=B_{\mathrm{bag}}^2=I,
 \qquad A^*=A,
 \qquad B_{\mathrm{bag}}^*=B_{\mathrm{bag}},
 \qquad AB_{\mathrm{bag}}+B_{\mathrm{bag}}A=0.
\label{eq:v3v69-AB-relations}
\end{equation}
\end{lemma}

\begin{proof}
Euclidean Clifford multiplication by orthogonal unit vectors anticommutes
and is Hermitian.  Thus \(\gamma_n\gamma(\widehat p)\) and
\(\gamma_n\gamma_5\) are anti-Hermitian and square to \(-I\), proving the
involution and adjoint identities after multiplication by \(i\).  Since
\(\gamma_5\) anticommutes with both \(\gamma_n\) and
\(\gamma(\widehat p)\), direct reordering gives the last identity.
\end{proof}

\begin{theorem}[Local complementing bag block]
\label{thm:v3v69-local-bag}
Impose the local zeroth-order condition
\begin{equation}
 B_{\mathrm{bag}}\Psi=\varepsilon\Psi,
 \qquad \varepsilon\in\{+1,-1\}.
\label{eq:v3v69-bag-condition}
\end{equation}
Then the allowed boundary space has zero intersection with the decaying
Calderon space at every \(p\ne0\).  The condition is gauge equivariant
because the two chiralities carry the same representation \(R\).
\end{theorem}

\begin{proof}
Suppose \(A v=v\) and \(B_{\mathrm{bag}}v=\varepsilon v\).  By
anticommutation,
\[
 \varepsilon v=AB_{mathrm{bag}}v
 =-B_{mathrm{bag}}Av=-\varepsilon v,
\]
so \(v=0\).  This is the complementing condition.  The Clifford
involution acts on spin indices and the identity acts on \(R\), hence the
projector commutes with gauge transformations.
\end{proof}

Unlike the V56 material projector, \(B_{\mathrm{bag}}\) exchanges the two
four-dimensional chiralities.  Proposition
\ref{prop:v3v68-mirror-necessity} explains why that exchange is necessary.

\subsection{Why an ordinary mirror mass removes the physical field too}
\label{sec:v3v69-mass}

Consider \(n_+\) positive-chirality and \(n_-\) negative-chirality Weyl
fields in one fixed complex gauge representation.  A gauge-invariant
quadratic Dirac mass is a matrix
\begin{equation}
 M:\mathbb C^{n_-}\longrightarrow\mathbb C^{n_+}
\label{eq:v3v69-mass-map}
\end{equation}
and its adjoint.  Singular-value decomposition reduces the massive sector
to independent opposite-chirality pairs.

\begin{proposition}[Quadratic mirror-decoupling obstruction]
\label{prop:v3v69-mass-obstruction}
For one physical Weyl field and one opposite-chirality mirror in the same
representation, either \(M=0\), in which case both remain massless, or
\(M\ne0\), in which case both form one massive Dirac field.  There is no
gauge-invariant quadratic mass that makes only the mirror heavy while
leaving the original Weyl field massless.
\end{proposition}

\begin{proof}
Here \(n_+=n_-=1\), so \(M\) is one complex number.  If it vanishes, the
mass matrix has two chiral zero modes.  If it is nonzero, its sole singular
value is positive and pairs the two fields into one massive Dirac spinor.
No same-chirality complex mass exists without an additional invariant
antilinear pairing, which the general complex Standard Model
representations do not possess.
\end{proof}

\begin{proposition}[Mass-matrix index]
\label{prop:v3v69-mass-index}
For the general map \eqref{eq:v3v69-mass-map}, the numbers of unpaired
massless chiral modes obey
\begin{equation}
 \dim\ker M^*-\dim\ker M=n_+-n_-.
\label{eq:v3v69-mass-index}
\end{equation}
In particular, a representationwise local bag completion with
\(n_+=n_-\) can be fully gapped, but a purely quadratic mass cannot gap
only one member of each completing pair.
\end{proposition}

\begin{proof}
Rank--nullity gives
\(\dim\ker M=n_--\operatorname{rank}M\) and
\(\dim\ker M^*=n_+-\operatorname{rank}M\).  Subtraction proves
\eqref{eq:v3v69-mass-index}.  The final assertion is the case
\(n_+=n_-\).
\end{proof}

Thus the smallest local elliptic completion is vectorlike, while the
desired light spectrum is chiral.  A four-dimensional hard mass does not
interpolate between those two facts.

\subsection{Eliminating a local mirror produces a nonlocal Schur complement}
\label{sec:v3v69-Schur}

Let a local quadratic completion have block operator
\begin{equation}
 \mathcal D_M=
 \begin{pmatrix}
  D_+&Y_M\\ Z_M&D_-
 \end{pmatrix},
\label{eq:v3v69-block-operator}
\end{equation}
with a complementing bag domain.  If the mirror block \(D_-\) is
invertible, formal Gaussian elimination gives the physical Schur
complement
\begin{equation}
 D_{\mathrm{eff}}=D_+-Y_MD_-^{-1}Z_M.
\label{eq:v3v69-Schur}
\end{equation}

\begin{lemma}[Local completion does not imply a local reduced Weyl
operator]
\label{lem:v3v69-Schur-nonlocal}
Unless the off-diagonal factors cancel the inverse by a special exact
factorization, the second term in \eqref{eq:v3v69-Schur} is
pseudodifferential and has noncompact kernel.  Dividing out the mirror
determinant therefore reintroduces the locality problem.
\end{lemma}

\begin{proof}
The inverse of a first-order elliptic operator has symbol of order \(-1\)
and a Green kernel supported away from the diagonal.  Multiplication by
local differential operators \(Y_M,Z_M\) does not make that kernel
diagonal unless their product contains the exact operator factor needed to
cancel \(D_-^{-1}\).  No such factor is present in an algebraic mass
coupling.  Hence the generic Schur complement is nonlocal.
\end{proof}

This is the fermionic counterpart of the V65 inverse-difference
obstruction.  It also explains why a determinant ratio must be analyzed as
a single regulated object rather than by separately declaring the mirror
harmless.

\subsection{Consequences for the continuum route}
\label{sec:v3v69-consequences}

\begin{theorem}[Principal-symbol decision after V69]
\label{thm:v3v69-decision}
At the presently established level:
\begin{enumerate}
\item APS supplies an elliptic chiral boundary condition but is nonlocal in
Euclidean time, so the V64 transfer argument does not apply;
\item a representationwise mirror supplies a local complementing full
Dirac bag condition; and
\item an ordinary four-dimensional quadratic mass cannot decouple that
mirror while retaining the target chiral field, while Gaussian elimination
generically yields a nonlocal reduced operator.
\end{enumerate}
\end{theorem}

\begin{proof}
The three items are Propositions
\ref{prop:v3v69-APS-complementing}--%
\ref{prop:v3v69-APS-time-nonlocal}, Theorem
\ref{thm:v3v69-local-bag}, and Propositions
\ref{prop:v3v69-mass-obstruction}--%
\ref{prop:v3v69-mass-index} together with Lemma
\ref{lem:v3v69-Schur-nonlocal}.
\end{proof}

The remaining local option must carry chirality as an interface index, not
as an unpaired field in a finite four-dimensional mass matrix.  A
domain-wall or overlap-type collar is the natural next construction: its
bulk is vectorlike and local, while a mass sign change supports a chiral
interface mode.

\begin{firewall}[No mirror decoupling yet]
\label{fw:v3v69-no-decoupling}
V69 proves local complementing ellipticity of the doubled principal block.
It does not construct a reflection-positive mirror action, prove
decoupling, define a determinant ratio, or show that a domain-wall limit
reproduces the V56 Lorentzian boundary theory.  The finite mirror may not
be omitted from degree-of-freedom or anomaly counts.
\end{firewall}

\subsection{V69 conclusion and V70 audit target}
\label{sec:v3v69-conclusion}

The APS and mirror repairs exchange two difficulties.  APS retains the
chiral field but introduces a frequency-dependent nonlocal projector.  A
mirror restores a local elliptic bag block, but a conventional mass either
leaves both chiralities light or gaps both, and eliminating the mirror is
nonlocal.  The route with a chance to preserve locality is therefore an
interface-index construction in an auxiliary collar dimension.

V70 is the next compulsory NS/YM audit.  After that audit it should build
the one-dimensional domain-wall normal operator in each Standard Model
representation, prove its chiral zero-mode count and bulk gap exactly, and
state the boundary conditions required to combine it with the V69 bag
block.  Anomaly inflow and reflection positivity must remain separate
tests.
\clearpage
\section{V70: companion audit and the exact domain-wall normal index}
\label{sec:v3v70}

This is the compulsory five-version audit following V65.  It is followed
by the first local interface-index construction suggested by V69.  The
one-dimensional calculation is exact and identifies both the chiral zero
mode and the obstruction to realizing a single wall on a compact periodic
regulator.

\subsection{Read-only NS/YM checkpoint}
\label{sec:v3v70-audit}

The active companion snapshots read at this checkpoint were
\begin{enumerate}
\item
\texttt{Renewal\_NS\_Stability\_Research\_Notes\_V31.tex}, with
\(887537\) bytes and SHA--256
\begin{center}\scriptsize\ttfamily
F1121B62238AD5491BB7CF03635EA9934942EDF573ED8FAAA9EE79E1D38A6696;
\end{center}
\item
\texttt{Renewal\_Yang\_Mills\_Mass\_Gap\_Research\_Notes\_V97.tex},
with \(1643309\) bytes and SHA--256
\begin{center}\scriptsize\ttfamily
1BB58202555A58DE468EAB7B5FE8EB1597CA0ED5FB4B478A0348A262CAC90368.
\end{center}
\end{enumerate}
The reads were passive; no companion or frozen file was changed.

NS V31 remains at the sharp conclusion used in V65: a modewise or marked
estimate can hide the exact critical first moment that reappears on
unmarking.  For the fermion problem this forbids proving decoupling only
after projecting onto the light mode; the norm of the reconstruction from
the complete higher-dimensional field must also be controlled.

YM V97 keeps a repeated occurrence, its distinct adjacent response, and
the payer in one five-factor Fourier sandwich.  It closes only the
distinct-adjacent branch and leaves the self-adjacent three-occurrence bank
for a separate higher-moment estimate.  The corresponding rule here is to
keep the domain-wall zero mode, the heavy bulk modes, and the determinant
phase in one regulated operator.  Multiplying a separately normalized
zero-mode determinant by an independently bounded heavy determinant would
repeat or erase the common spectral stock.

\begin{assessment}[V70 nonduplication and audit transfer]
\label{ass:v3v70-audit-transfer}
V68--V69 already prove the Euclidean Weyl obstruction, APS nonlocality,
and local full-Dirac bag completion.  V70 does not revisit those symbols.
It moves to the missing interface index and keeps all modes in the single
supersymmetric normal operator below, as required by the two companion
sharpness lessons.
\end{assessment}

\subsection{The supersymmetric wall operator}
\label{sec:v3v70-wall-operator}

Let \(m:\mathbb R\to\mathbb R\) be locally absolutely continuous and
satisfy
\begin{equation}
 \lim_{y\to-\infty}m(y)=-m_0,
 \qquad
 \lim_{y\to+\infty}m(y)=m_0,
 \qquad m_0>0.
\label{eq:v3v70-mass-asymptotics}
\end{equation}
Define closed operators on \(L^2(\mathbb R)\)
\begin{equation}
 Q=\partial_y+m(y),
 \qquad Q^*=-\partial_y+m(y),
\label{eq:v3v70-Q}
\end{equation}
with Sobolev domain \(H^1(\mathbb R)\), and the self-adjoint normal Dirac
operator
\begin{equation}
 \mathcal H_m=
 \begin{pmatrix}0&Q^*\\Q&0\end{pmatrix}.
\label{eq:v3v70-H}
\end{equation}

\begin{theorem}[Exact one-wall index]
\label{thm:v3v70-wall-index}
Under \eqref{eq:v3v70-mass-asymptotics},
\begin{equation}
 \dim\ker Q=1,
 \qquad \dim\ker Q^*=0,
 \qquad \operatorname{ind}Q=1.
\label{eq:v3v70-index}
\end{equation}
The normalized zero mode is a scalar multiple of
\begin{equation}
 f_0(y)=\exp\left[-\int_0^y m(s)\,ds\right].
\label{eq:v3v70-zero-mode}
\end{equation}
Reversing both asymptotic signs gives index \(-1\) and the opposite
chirality.
\end{theorem}

\begin{proof}
The first-order equation \(Qf=0\) has exactly the one-dimensional solution
space spanned by \eqref{eq:v3v70-zero-mode}.  Its exponent tends to
\(-m_0|y|+O(1)\) at both ends, so the solution lies in \(H^1\).  The
solutions of \(Q^*g=0\) are multiples of
\(\exp[\int_0^y m]\), which grow like \(e^{m_0|y|}\) and are not in
\(L^2\).  This proves \eqref{eq:v3v70-index}.  Reversing \(m\) interchanges
\(Q\) and \(-Q^*\).
\end{proof}

The two components in \eqref{eq:v3v70-H} are the two four-dimensional
chiralities.  Thus the sign change, rather than a rank defect of a finite
mass matrix, selects one light Weyl chirality.

\subsection{Exact sharp-wall spectrum}
\label{sec:v3v70-spectrum}

For the step profile
\begin{equation}
 m(y)=m_0\operatorname{sign}y,
\label{eq:v3v70-step-mass}
\end{equation}
the squared operator is
\begin{equation}
 \mathcal H_m^2=
 \begin{pmatrix}
  Q^*Q&0\\0&QQ^*
 \end{pmatrix},
\quad
 Q^*Q=-\partial_y^2+m_0^2-2m_0\delta_0,
\quad
 QQ^*=-\partial_y^2+m_0^2+2m_0\delta_0.
\label{eq:v3v70-squared-step}
\end{equation}

\begin{proposition}[One zero mode and a bulk gap]
\label{prop:v3v70-step-spectrum}
For \eqref{eq:v3v70-step-mass}, \(Q^*Q\) has the single bound state
\begin{equation}
 f_0(y)=\sqrt{m_0}\,e^{-m_0|y|}
\label{eq:v3v70-step-zero}
\end{equation}
at eigenvalue zero and otherwise has spectrum \([m_0^2,\infty)\).
The operator \(QQ^*\) has spectrum \([m_0^2,\infty)\) and no bound state.
Consequently \(\mathcal H_m\) has one zero mode and all nonzero normal
spectrum satisfies \(|\lambda|\ge m_0\).
\end{proposition}

\begin{proof}
Away from zero, a bound state below \(m_0^2\) has the form
\(ce^{-\kappa|y|}\), \(\kappa>0\).  The attractive delta jump condition
for \(Q^*Q\) is
\(f'(0^+)-f'(0^-)=-2m_0f(0)\), which gives
\(-2\kappa=-2m_0\) and hence the unique value \(\kappa=m_0\), with
eigenvalue \(m_0^2-\kappa^2=0\).  Normalization gives
\eqref{eq:v3v70-step-zero}.  The repulsive delta jump for \(QQ^*\) would
require \(-2\kappa=2m_0\), impossible for \(\kappa>0\).  The essential
spectrum of both constant-asymptotic Schr\"odinger operators is
\([m_0^2,\infty)\).  Taking square roots proves the last assertion.
\end{proof}

\subsection{Four-dimensional Weyl mode from the complete block}
\label{sec:v3v70-effective-Weyl}

Let \(D_4\) be the gauge-covariant four-dimensional Dirac operator in one
Standard Model representation \(R\), extended independently of \(y\) in a
neighbourhood of the wall.  The local five-dimensional operator has the
schematic form
\begin{equation}
 \mathcal D_5=D_4+\gamma_5\partial_y+m(y),
\label{eq:v3v70-D5}
\end{equation}
with the same wall factorization as \(\mathcal H_m\).

\begin{proposition}[Chiral interface reduction]
\label{prop:v3v70-interface-compression}
Let \(\psi_\chi(x)\) have the chirality selected by
Theorem~\ref{thm:v3v70-wall-index}.  At the principal level, the separated
field \(\Psi(x,y)=f_0(y)\psi_\chi(x)\) obeys
\begin{equation}
 \mathcal D_5\Psi=f_0(y)D_4\psi_\chi.
\label{eq:v3v70-separated-D5}
\end{equation}
Thus its tangential equation is the target four-dimensional Weyl equation
in \(R\).  The orthogonal normal sector has mass gap \(m_0\) for the sharp
wall.
\end{proposition}

\begin{proof}
The normal part annihilates \(f_0\) in the selected chirality, while the
tangential operator acts only on \(x\), proving
\eqref{eq:v3v70-separated-D5}.  The normalized integral
\(\int_{\mathbb R}|f_0|^2dy=1\) gives the same Weyl kinetic coefficient in
the reduced bilinear with its independent Grassmann conjugate.
Proposition~\ref{prop:v3v70-step-spectrum} gives the gap on the orthogonal
normal complement.
\end{proof}

For the full Standard Model, one takes a separate wall block in every
representation in \eqref{eq:v3v66-SM-list}, reversing the wall orientation
for the desired four-dimensional chirality.  The complete direct sum is
local and vectorlike in five dimensions, while its interface index is the
chiral representation list.

\subsection{Why a compact periodic fifth direction restores mirrors}
\label{sec:v3v70-compact}

\begin{proposition}[Zero net wall index on a circle]
\label{prop:v3v70-circle-zero}
Let \(m:S^1\to\mathbb R\) be continuous with only simple zeros.  The
numbers of sign changes from negative to positive and from positive to
negative are equal.  Hence the sum of the local wall indices is zero.
Every periodic local mass profile carrying a positive-chirality wall also
carries an opposite-chirality wall.
\end{proposition}

\begin{proof}
Traversing a circle returns to the initial sign.  Sign changes therefore
alternate, so their two orientations occur equally often.  Theorem
\ref{thm:v3v70-wall-index} assigns \(+1\) to one orientation and \(-1\) to
the other, and the sum vanishes.
\end{proof}

\begin{corollary}[A single wall needs an index sink]
\label{cor:v3v70-index-sink}
A compact local regulator can realize one unpaired interface chirality only
if a physical boundary condition, a nonlocal projection, or another
topological sector carries the opposite index.  Simply making the fifth
direction finite and periodic restores a mirror.
\end{corollary}

This is the auxiliary-dimensional form of the V69 mass-matrix index.  The
infinite line realizes the missing index at its two asymptotic phases; a
compact circle cannot.

\subsection{One-block requirement and present claim boundary}
\label{sec:v3v70-one-block}

The orthogonal decomposition
\begin{equation}
 I=P_0+(I-P_0)
\label{eq:v3v70-decomposition}
\end{equation}
is useful for estimates, but it is not a factorization of the regulated
determinant into independently normalized physical and heavy measures.
Gauge-field dependence mixes the two sectors away from the frozen product
background, and the parity-odd heavy determinant supplies anomaly inflow.

\begin{firewall}[No premature determinant split]
\label{fw:v3v70-no-split}
V70 proves the normal index and gap.  It does not prove that the heavy
determinant is a field-independent constant, nor may that determinant be
divided out before its phase and mixing with \(P_0\) are controlled.  Such
a split would violate the common-block audit lesson and could discard the
inflow that makes the interface theory gauge consistent.
\end{firewall}

\begin{firewall}[No finite local chiral regulator yet]
\label{fw:v3v70-no-finite-regulator}
The exact single-wall construction uses a noncompact auxiliary line.
Proposition~\ref{prop:v3v70-circle-zero} shows why a finite periodic
replacement contains a mirror.  V70 therefore does not yet supply the
finite local regulator required for cutoff removal or reflection
positivity.
\end{firewall}

\subsection{V70 conclusion and V71 target}
\label{sec:v3v70-conclusion}

A sign-changing mass produces exactly one normalizable chiral mode, and
the sharp wall has a strict \(m_0\) gap to every other normal excitation.
This solves the local principal-symbol completion on the noncompact
auxiliary line.  Compact periodic localization has zero net index and
restores an opposite wall, so decoupling must be proved as a limit of one
complete two-wall operator.

V71 should solve the finite two-wall normal problem.  It must compute the
exponential light-mode mixing as a function of wall separation, give a
uniform resolvent bound on the heavy complement, and identify an order of
limits in which the unwanted wall decouples without splitting the
determinant prematurely.  Reflection positivity and anomaly inflow must be
tested on that same finite block.
\clearpage
\section{V71: exact kink--antikink mixing and the heavy resolvent}
\label{sec:v3v71}

V70 proves a single-wall index on a noncompact auxiliary line and shows
that compact periodic localization restores an opposite wall.  The first
controlled approximation is therefore a kink and antikink separated by a
finite distance \(L\), kept inside one normal Dirac operator.  For the
sharp profile, its light splitting and heavy resolvent can be computed
exactly.

\subsection{The two-wall operator}
\label{sec:v3v71-operator}

Translate the physical wall to \(y=0\) and the mirror wall to \(y=L\).
Take
\begin{equation}
 m_L(y)=
 \begin{cases}
  -m_0,&y<0,\\
  +m_0,&0<y<L,\\
  -m_0,&y>L,
 \end{cases}
 \qquad m_0>0,
\label{eq:v3v71-two-wall-mass}
\end{equation}
and define \(Q_L=\partial_y+m_L\) and
\begin{equation}
 \mathcal H_L=
 \begin{pmatrix}0&Q_L^*\\Q_L&0\end{pmatrix}.
\label{eq:v3v71-HL}
\end{equation}
Its square has the partner blocks
\begin{align}
 Q_L^*Q_L
 &=-\partial_y^2+m_0^2-2m_0\delta_0+2m_0\delta_L,
\label{eq:v3v71-QstarQ}\\
 Q_LQ_L^*
 &=-\partial_y^2+m_0^2+2m_0\delta_0-2m_0\delta_L.
\label{eq:v3v71-QQstar}
\end{align}
The attractive delta in the first line is localized at the physical wall;
the attractive delta in the second is localized at the mirror wall.

\subsection{Exact light singular value}
\label{sec:v3v71-light}

Let \(E^2<m_0^2\) be a bound-state eigenvalue of \(Q_L^*Q_L\) and put
\begin{equation}
 \kappa=(m_0^2-E^2)^{1/2}>0.
\label{eq:v3v71-kappa}
\end{equation}
A bound-state eigenfunction has the form
\begin{equation}
 f(y)=
 \begin{cases}
  A e^{\kappa y},&y<0,\\
  B e^{-\kappa y}+C e^{\kappa y},&0<y<L,\\
  D e^{-\kappa(y-L)},&y>L.
 \end{cases}
\label{eq:v3v71-piecewise-f}
\end{equation}

\begin{theorem}[Exact two-wall splitting equation]
\label{thm:v3v71-splitting}
For every \(L>0\), there is exactly one bound eigenvalue
\(E_L^2\in(0,m_0^2)\) of \(Q_L^*Q_L\).  It is characterized by
\begin{equation}
 E_L^2=m_0^2e^{-2\kappa_LL},
 \qquad
 \kappa_L^2=m_0^2-E_L^2,
\label{eq:v3v71-splitting-equation}
\end{equation}
or equivalently
\begin{equation}
 \kappa_L^2=m_0^2\bigl(1-e^{-2\kappa_LL}\bigr).
\label{eq:v3v71-kappa-equation}
\end{equation}
The partner \(Q_LQ_L^*\) has the same positive bound eigenvalue.  Hence
\(\mathcal H_L\) has precisely the two light eigenvalues
\(\pm E_L\).
\end{theorem}

\begin{proof}
Continuity at zero gives \(A=B+C\).  The attractive delta jump condition
there is
\[
 f'(0^+)-f'(0^-)=-2m_0f(0),
\]
which reduces to
\begin{equation}
 \frac CB=\frac{\kappa-m_0}{m_0}.
\label{eq:v3v71-first-ratio}
\end{equation}
Continuity at \(L\) and the repulsive jump condition
\(f'(L^+)-f'(L^-)=2m_0f(L)\) give
\begin{equation}
 \frac CB=-\frac{m_0}{\kappa+m_0}e^{-2\kappa L}.
\label{eq:v3v71-second-ratio}
\end{equation}
Equating the two ratios yields
\(m_0^2-\kappa^2=m_0^2e^{-2\kappa L}\), which is
\eqref{eq:v3v71-splitting-equation}.

Let
\(F(\kappa)=\kappa^2-m_0^2+m_0^2e^{-2\kappa L}\).  Then \(F(0)=0\),
\(F'(0)=-2m_0^2L<0\), \(F(m_0)>0\), and
\(F''(\kappa)=2+4m_0^2L^2e^{-2\kappa L}>0\).  Strict convexity gives at
most one positive zero; the sign change gives exactly one.  The zero at
\(\kappa=0\) represents the continuum threshold rather than an
\(L^2\) state, whereas the positive root gives the strict bound state.

For \(E_L>0\), the map \(f\mapsto E_L^{-1}Q_Lf\) is an isometry from this
eigenspace of \(Q_L^*Q_L\) to the corresponding eigenspace of
\(Q_LQ_L^*\).  The two partner components give eigenvectors of
\(\mathcal H_L\) with eigenvalues \(\pm E_L\).  The same convexity argument
excludes any other bound singular value.
\end{proof}


\begin{corollary}[Exponential mixing bounds]
\label{cor:v3v71-exponential}
If \(m_0L\ge1\), then \(\kappa_L\ge m_0/2\) and
\begin{equation}
 m_0e^{-m_0L}\le E_L\le m_0e^{-m_0L/2}.
\label{eq:v3v71-rough-exponential}
\end{equation}
Moreover
\begin{equation}
 E_L=m_0e^{-m_0L}
 \exp\bigl((m_0-\kappa_L)L\bigr),
 \qquad
 m_0-\kappa_L=\frac{E_L^2}{m_0+\kappa_L}.
\label{eq:v3v71-sharp-exponential}
\end{equation}
Thus \(E_L/(m_0e^{-m_0L})\to1\) as \(L\to\infty\).
\end{corollary}

\begin{proof}
At \(\kappa=m_0/2\), the right side of
\eqref{eq:v3v71-kappa-equation} is at least
\(m_0^2(1-e^{-1})>m_0^2/4\), so the positive crossing occurs above
\(m_0/2\).  Since \(m_0/2\le\kappa_L<m_0\), the first equation in
\eqref{eq:v3v71-splitting-equation} gives
\eqref{eq:v3v71-rough-exponential}.  Rearrangement gives
\eqref{eq:v3v71-sharp-exponential}.  The rough upper bound makes
\((m_0-\kappa_L)L\le 2m_0Le^{-m_0L}\to0\), proving the limit.
\end{proof}

\subsection{Localized chiral basis and one-use mixing}
\label{sec:v3v71-localized-basis}

Let \(f_L\) be the normalized bound state of \(Q_L^*Q_L\) and set
\begin{equation}
 g_L=E_L^{-1}Q_Lf_L.
\label{eq:v3v71-partner-mode}
\end{equation}
Then \(f_L\) is localized at \(y=0\) and \(g_L\) at \(y=L\).  In the
orthonormal two-dimensional space spanned by
\((f_L,0)\) and \((0,g_L)\),
\begin{equation}
 \mathcal H_L^{\mathrm{light}}=
 \begin{pmatrix}0&E_L\\E_L&0\end{pmatrix}.
\label{eq:v3v71-light-matrix}
\end{equation}

\begin{proposition}[The residual mass is one common-wall matrix element]
\label{prop:v3v71-one-use}
The off-diagonal coefficient in \eqref{eq:v3v71-light-matrix} is both
\(\langle g_L,Q_Lf_L\rangle\) and
\(\langle f_L,Q_L^*g_L\rangle\), and equals \(E_L\).  These are adjoint
descriptions of one mixing event, not two independent exponentially small
factors.
\end{proposition}

\begin{proof}
Equation \eqref{eq:v3v71-partner-mode} and normalization give the first
matrix element as \(E_L\).  Taking the adjoint gives the second.  Both are
the two entries of the same self-adjoint block
\eqref{eq:v3v71-light-matrix}.
\end{proof}

This is the exact implementation of the V70/YM audit rule: multiplying the
two adjoint entries as if they arose from distinct walls would charge the
same singular value twice.

\subsection{Heavy spectral and resolvent gap}
\label{sec:v3v71-heavy}

\begin{theorem}[Uniform heavy resolvent]
\label{thm:v3v71-heavy-resolvent}
Let \(P_L\) be the spectral projection of \(\mathcal H_L\) onto
\(\{ -E_L,E_L\}\) and \(P_L^\perp=I-P_L\).  Then
\begin{equation}
 \sigma(\mathcal H_L|_{P_L^\perp})
 \subset(-\infty,-m_0]\cup[m_0,\infty).
\label{eq:v3v71-heavy-spectrum}
\end{equation}
For every \(|z|<m_0\),
\begin{equation}
 \left\|P_L^\perp(\mathcal H_L-z)^{-1}P_L^\perp\right\|
 \le\frac1{m_0-|z|}.
\label{eq:v3v71-heavy-resolvent}
\end{equation}
The bound is independent of \(L\).
\end{theorem}

\begin{proof}
Theorem~\ref{thm:v3v71-splitting} gives the only bound singular value.
Both partner Schr\"odinger operators have constant asymptotic potential
\(m_0^2\), so the rest of their spectrum is
\([m_0^2,\infty)\).  Taking the signed square root gives
\eqref{eq:v3v71-heavy-spectrum}.  The spectral theorem then bounds the
resolvent by the reciprocal distance from \(z\) to that set, which is at
least \(m_0-|z|\).
\end{proof}

\subsection{Stable Feshbach reduction under local interactions}
\label{sec:v3v71-Feshbach}

Let \(V\) be a bounded perturbation containing frozen tangential momentum,
gauge, Higgs, and weak curvature terms, and set
\(\mathcal H_{L,V}=\mathcal H_L+V\).  Denote \(P=P_L\) and
\(Q=I-P\).

\begin{proposition}[One-block effective operator]
\label{prop:v3v71-Feshbach}
If \(\|V\|\le m_0/4\) and \(|z|\le m_0/4\), then
\(Q(\mathcal H_{L,V}-z)Q\) is invertible and
\begin{equation}
 \left\|[Q(\mathcal H_{L,V}-z)Q]^{-1}\right\|
 \le\frac2{m_0}.
\label{eq:v3v71-perturbed-heavy}
\end{equation}
The light resolvent is governed exactly by the Feshbach operator
\begin{equation}
 \mathcal F_L(z)=
 P(\mathcal H_{L,V}-z)P
 -PVQ[Q(\mathcal H_{L,V}-z)Q]^{-1}QVP,
\label{eq:v3v71-Feshbach}
\end{equation}
and the heavy correction satisfies
\begin{equation}
 \left\|PVQ[Q(\mathcal H_{L,V}-z)Q]^{-1}QVP\right\|
 \le\frac{2\|V\|^2}{m_0}.
\label{eq:v3v71-Feshbach-bound}
\end{equation}
\end{proposition}

\begin{proof}
On \(Q\), Theorem~\ref{thm:v3v71-heavy-resolvent} gives
\(\|(Q(\mathcal H_L-z)Q)^{-1}\|\le4/(3m_0)\).  The Neumann series for
the addition of \(QVQ\) converges because its norm relative to the inverse
is at most \(1/3\), and the resulting bound is at most \(2/m_0\).
Block Gaussian elimination gives \eqref{eq:v3v71-Feshbach}; applying the
operator norm proves \eqref{eq:v3v71-Feshbach-bound}.
\end{proof}

Unlike the formal Schur complement in V69, this reduction comes with a
uniform heavy inverse.  It remains one identity for the complete operator;
it does not independently normalize the two determinant factors.

\subsection{Order of limits and present boundary}
\label{sec:v3v71-limits}

At the free normal level, the controlled order is
\begin{equation}
 m_0L\longrightarrow\infty
 \quad\hbox{with fixed }m_0>0,
 \qquad E_L\longrightarrow0,
\label{eq:v3v71-first-limit}
\end{equation}
while the heavy resolvent remains bounded by \(m_0^{-1}\).  Only after
correlators localized near \(y=0\), the determinant phase, and the
Feshbach correction are controlled uniformly may a regulator limit in the
four physical directions or a large-\(m_0\) limit be taken.

\begin{firewall}[A small residual mass is not mirror decoupling]
\label{fw:v3v71-not-decoupling}
The estimate \(E_L\to0\) separates the two light wall modes but does not
remove the mirror mode from the Hilbert space or determinant.  Local
observables near one wall require an off-diagonal propagation estimate,
and gauge fields that extend to both walls can still couple to both.  No
chiral determinant or anomaly-inflow limit has yet been proved.
\end{firewall}

\begin{firewall}[Reflection positivity remains a same-block test]
\label{fw:v3v71-no-RP}
The mass profile is real and independent of Euclidean time, but that fact
alone is not a fermionic Osterwalder--Schrader proof.  A time-local lattice
Dirac discretization, its reflection map, and positivity of the complete
kink--antikink transfer operator must be established before the bosonic
V64 result can be extended to this sector.
\end{firewall}

\subsection{V71 conclusion and V72 target}
\label{sec:v3v71-conclusion}

The sharp two-wall operator has exactly one light singular value satisfying
\(E_L=m_0e^{-\kappa_LL}\), two Dirac eigenvalues \(\pm E_L\), and an
\(L\)-independent gap \(m_0\) on the heavy complement.  The Feshbach map
controls weak local interactions without splitting the regulated block.
These estimates make the separation limit quantitative, while showing
that the mirror remains a light state at the remote wall.

V72 should derive a Combes--Thomas or direct Green-kernel estimate from the
physical wall to the mirror wall, including bounded tangential gauge/Higgs
perturbations.  It should then formulate the precise observable algebra
localized near the physical wall and decide whether its connected
correlators converge as \(L\to\infty\) while the complete determinant
phase is retained.  The phase should be matched to the V66 descent only
after this locality estimate is proved.
\clearpage
\section{V72: exponential wall locality and convergence of local fermion correlators}
\label{sec:v3v72}

V71 controls the light singular value and the heavy complement, but a
small eigenvalue alone does not prove that the remote mirror is invisible
to observables near the physical wall.  This note proves the required
off-diagonal resolvent estimate directly.  It then compares the two-wall
and single-wall covariances on a fixed physical-wall neighbourhood.

\subsection{A first-order Combes--Thomas estimate}
\label{sec:v3v72-CT}

Let \(H\) be a self-adjoint one-dimensional Dirac operator on
\(L^2(\mathbb R;\mathbb C^N)\) of the form
\begin{equation}
 H=J\partial_y+W(y),
 \qquad J^*=-J,
 \qquad \|J\|=1,
 \qquad W(y)^*=W(y),
\label{eq:v3v72-general-Dirac}
\end{equation}
with bounded measurable \(W\).  This includes \(\mathcal H_L\) after a
constant unitary change of basis, and it also permits bounded frozen
tangential gauge/Higgs perturbations.

\begin{theorem}[Off-diagonal Dirac resolvent]
\label{thm:v3v72-CT}
Let \(z\in\mathbb C\setminus\mathbb R\), put
\(\delta=|\operatorname{Im}z|\), and let \(X,Y\subset\mathbb R\) be
measurable sets at distance \(d\).  For every \(0<a<\delta\),
\begin{equation}
 \left\|\mathbf1_X(H-z)^{-1}\mathbf1_Y\right\|
 \le\frac{e^{-ad}}{\delta-a}.
\label{eq:v3v72-CT-bound}
\end{equation}
The constant is independent of the wall separation and of
\(\|W\|_\infty\).
\end{theorem}

\begin{proof}
Choose a real Lipschitz function \(\varphi\) with
\(|\varphi'|\le1\), \(\varphi=0\) on \(Y\), and
\(\varphi\ge d\) on \(X\); truncation of the distance to \(Y\) gives such
a function.  Conjugation on compactly supported smooth fields gives
\begin{equation}
 e^{a\varphi}He^{-a\varphi}=H-aJ\varphi'.
\label{eq:v3v72-conjugated-H}
\end{equation}
The perturbation has norm at most \(a\).  Since \(H\) is self-adjoint,
\(\|(H-z)^{-1}\|\le\delta^{-1}\).  The resolvent Neumann series for
\(H-aJ\varphi'-z\) therefore gives
\begin{equation}
 \left\|e^{a\varphi}(H-z)^{-1}e^{-a\varphi}\right\|
 \le(\delta-a)^{-1}.
\label{eq:v3v72-weighted-resolvent}
\end{equation}
Multiplication on the left by \(\mathbf1_Xe^{-a\varphi}\) supplies
\(e^{-ad}\), while \(e^{a\varphi}\mathbf1_Y=\mathbf1_Y\).  Approximation
by bounded truncations of \(\varphi\) justifies the conjugation without an
unbounded-weight assumption.
\end{proof}

Taking \(a=\delta/2\) gives the convenient estimate
\begin{equation}
 \left\|\mathbf1_X(H-z)^{-1}\mathbf1_Y\right\|
 \le\frac2\delta e^{-\delta d/2}.
\label{eq:v3v72-CT-half}
\end{equation}
This bound includes the light pair \(\pm E_L\); it does not rely on a
separate normalization of that pair.

\subsection{Comparison with the single wall}
\label{sec:v3v72-comparison}

Let \(H_\infty\) be the single-kink operator with mass
\begin{equation}
 m_\infty(y)=
 \begin{cases}-m_0,&y<0,\\+m_0,&y>0,\end{cases}
\label{eq:v3v72-single-mass}
\end{equation}
and let \(H_L\) be the kink--antikink operator of V71.  They agree on
\(( -\infty,L)\), and their bounded difference
\begin{equation}
 V_L=H_L-H_\infty
\label{eq:v3v72-remote-difference}
\end{equation}
is supported in \([L,\infty)\) with \(\|V_L\|\le2m_0\).

Fix the physical-wall neighbourhood
\begin{equation}
 X_R=[-R,R],
 \qquad 0<R<L.
\label{eq:v3v72-XR}
\end{equation}

\begin{theorem}[Local resolvent convergence]
\label{thm:v3v72-local-resolvent}
For \(z\notin\mathbb R\), \(\delta=|\operatorname{Im}z|\),
\begin{equation}
 \left\|\mathbf1_{X_R}
 \bigl[(H_L-z)^{-1}-(H_\infty-z)^{-1}\bigr]
 \mathbf1_{X_R}\right\|
 \le
 \frac{8m_0}{\delta^2}e^{-\delta(L-R)}.
\label{eq:v3v72-local-resolvent-convergence}
\end{equation}
Thus the two-wall covariance converges exponentially in operator norm on
every fixed physical-wall neighbourhood.
\end{theorem}

\begin{proof}
The resolvent identity gives
\begin{equation}
 (H_L-z)^{-1}-(H_\infty-z)^{-1}
 =-(H_L-z)^{-1}V_L(H_\infty-z)^{-1}.
\label{eq:v3v72-resolvent-identity}
\end{equation}
Insert \(\mathbf1_{[L,\infty)}\) on both sides of \(V_L\).  The distance
from \(X_R\) to that half-line is \(L-R\).  Apply
\eqref{eq:v3v72-CT-half} to each resolvent and use
\(\|V_L\|\le2m_0\).  The product is
\((2/\delta)^2(2m_0)e^{-\delta(L-R)}\), which is the displayed bound.
\end{proof}

The same proof applies when both operators contain the same bounded local
gauge/Higgs potential on \(( -\infty,L)\) and differ only in the remote
region.  Self-adjointness is the only spectral input in Theorem
\ref{thm:v3v72-CT}.

\subsection{Separated wall-to-wall propagation}
\label{sec:v3v72-cross-wall}

Let \(Y_{L,R}=[L-R,L+R]\) be a fixed-width mirror-wall neighbourhood.

\begin{corollary}[Physical-to-mirror resolvent suppression]
\label{cor:v3v72-cross-wall}
For \(L>2R\),
\begin{equation}
 \left\|\mathbf1_{X_R}(H_L-z)^{-1}mathbf1_{Y_{L,R}}\right\|
 \le\frac2\delta e^{-\delta(L-2R)/2}.
\label{eq:v3v72-cross-wall}
\end{equation}
\end{corollary}

\begin{proof}
The two sets are at distance \(L-2R\).  Apply
\eqref{eq:v3v72-CT-half}.
\end{proof}

At a fixed nonzero Euclidean spectral distance, the remote mirror is
therefore invisible exponentially fast.  The decay rate in
\eqref{eq:v3v72-cross-wall} is governed by \(\delta\), not by the sharper
normal mass \(m_0\); this general form remains valid in the presence of
arbitrary bounded self-adjoint local perturbations.

\subsection{The localized Gaussian observable algebra}
\label{sec:v3v72-observables}

Fix \(R\) and \(\delta_0>0\).  Let \(\mathfrak A_{R,\delta_0}\) be the
finite polynomial Grassmann algebra generated by smeared fields whose
auxiliary-coordinate test functions are supported in \(X_R\), and whose
spectral transforms are supported where
\begin{equation}
 |\operatorname{Im}z|\ge\delta_0.
\label{eq:v3v72-IR-window}
\end{equation}
At finite tangential regulator, Gaussian expectations are finite sums of
determinants of covariance matrix elements.

\begin{theorem}[Convergence of localized free correlators]
\label{thm:v3v72-correlator-convergence}
For every fixed polynomial \(F\in\mathfrak A_{R,\delta_0}\), the normalized
two-wall free expectation converges to the single-wall expectation as
\(L\to\infty\).  For a monomial of degree \(2q\), the difference is bounded
by
\begin{equation}
 C_{F,q,m_0,\delta_0}
 e^{-\delta_0(L-R)},
\label{eq:v3v72-correlator-rate}
\end{equation}
where the constant is independent of \(L\).
\end{theorem}

\begin{proof}
Wick's rule writes the expectation of a degree-\(2q\) monomial as the
determinant of a \(q\times q\) matrix of smeared covariance entries.  By
Theorem~\ref{thm:v3v72-local-resolvent}, every entry differs from its
single-wall value by at most a constant times
\(e^{-\delta_0(L-R)}\), while all entries are uniformly bounded by
\(\delta_0^{-1}\) times the test-function norms.  Expanding the difference
of the two determinants one column at a time gives at most \(q\) terms,
each with one covariance difference and \(q-1\) uniformly bounded
columns.  Hadamard's inequality gives
\eqref{eq:v3v72-correlator-rate}.  Linearity treats every polynomial.
\end{proof}

This is a genuine local decoupling theorem, but it contains the infrared
restriction \eqref{eq:v3v72-IR-window}.  The constants diverge as
\(\delta_0\downarrow0\), so the massless zero-frequency sector has not
been controlled.

\subsection{Stability under a common bounded interaction}
\label{sec:v3v72-interaction}

Let \(W_L\) and \(W_\infty\) be self-adjoint bounded local potentials with
\begin{equation}
 W_L=W_\infty\quad\hbox{on }(-\infty,L),
 \qquad \sup_L(\|W_L\|+\|W_\infty\|)<\infty.
\label{eq:v3v72-common-potential}
\end{equation}
They may contain a finite-cutoff tangential gauge field, Higgs Yukawa
matrix, and weak frozen geometric coefficients.  Put
\(\widehat H_L=H_L+W_L\) and
\(\widehat H_\infty=H_\infty+W_\infty\).

\begin{proposition}[Interacting frozen-coefficient locality]
\label{prop:v3v72-interacting-locality}
The bounds
\eqref{eq:v3v72-local-resolvent-convergence} and
\eqref{eq:v3v72-cross-wall} remain valid with a constant proportional to
\(\|\widehat H_L-\widehat H_\infty\|\), provided their difference remains
supported in \([L,\infty)\).  No smallness assumption on the common
self-adjoint bounded potential is required.
\end{proposition}

\begin{proof}
The commutator in \eqref{eq:v3v72-conjugated-H} depends only on the
first-order derivative term; every multiplication potential commutes with
the weight.  Thus Theorem~\ref{thm:v3v72-CT} has the same constants.  Use
the resolvent identity with
\(\widehat H_L-\widehat H_\infty\) in place of \(V_L\).
\end{proof}

This proposition is still a Gaussian/frozen-background statement.  A
dynamical gauge field is integrated rather than held fixed, and uniform
control of the potential norm under that integration is a later problem.

\subsection{What local convergence does not control}
\label{sec:v3v72-phase}

The determinant phase is a trace over the complete auxiliary line.  The
operator-norm estimate on \(X_R\) does not bound that trace, and the remote
region has volume independent of the local observable support.

\begin{firewall}[Local correlators do not determine the determinant line]
\label{fw:v3v72-determinant}
Theorem~\ref{thm:v3v72-correlator-convergence} proves convergence of a
specified local free observable algebra at nonzero spectral distance.  It
does not prove convergence of the partition function, eta invariant,
determinant phase, anomaly inflow, or the zero-frequency chiral correlator.
Those quantities may retain a topological contribution from the remote
wall even when every local resolvent matrix element converges.
\end{firewall}

\begin{firewall}[No four-dimensional cutoff removal]
\label{fw:v3v72-no-cutoff}
The tangential regulator remains finite in the Wick-determinant argument.
Uniformity in its dimension, renormalized composite fields, and the
simultaneous limit \(\delta_0\downarrow0\) are not established.
\end{firewall}

\subsection{V72 conclusion and V73 target}
\label{sec:v3v72-conclusion}

At any fixed nonzero Euclidean spectral distance, the complete two-wall
resolvent is exponentially local in the auxiliary direction.  Its
restriction to a fixed physical-wall neighbourhood converges in norm to
the single-wall resolvent, and every finite Gaussian polynomial correlator
in that localized algebra converges exponentially.  This remains true for
common bounded self-adjoint gauge/Higgs backgrounds.

V73 should retain the complete two-wall determinant and differentiate its
phase along a compactly supported gauge/metric path.  A relative
heat-kernel formula should separate the local wall contribution from the
remote wall and show how their descent terms combine.  The Standard Model
sum from V66 must be inserted before the limit \(L\to\infty\); any residual
term should be identified as a mapping-torus holonomy rather than estimated
by the local resolvent norm.
\clearpage
\section{V73: the complete determinant phase and local anomaly inflow}
\label{sec:v3v73}

V72 proves locality of resolvents and Gaussian correlators but deliberately
does not split the determinant.  This note differentiates the phase of the
complete two-wall determinant at finite regulator.  Gauge covariance then
gives an exact Ward identity for the whole block, while the V72 estimate
controls the effect of the remote wall on compactly supported phase
variations.

\subsection{Finite-regulator phase connection}
\label{sec:v3v73-phase-connection}

Fix the tangential and auxiliary local lattices, remove the stabilizer
sector, and let \(D_L(t)\) be the finite matrix representing the complete
kink--antikink fermion operator along a \(C^1\) path of backgrounds.  Add
an infrared spectral displacement \(i\mu\), \(\mu>0\), and define
\begin{equation}
 \Theta_L(t)=\operatorname{Im}\log\det(D_L(t)+i\mu)
\label{eq:v3v73-Theta}
\end{equation}
on a continuous logarithm branch along the path.

\begin{lemma}[Exact determinant differential]
\label{lem:v3v73-det-differential}
For every \(t\),
\begin{equation}
 \dot\Theta_L(t)=
 \operatorname{Im}\operatorname{Tr}\left[
 (D_L(t)+i\mu)^{-1}\dot D_L(t)
 \right].
\label{eq:v3v73-phase-derivative}
\end{equation}
\end{lemma}

\begin{proof}
The shift makes the matrix invertible because \(D_L(t)\) is self-adjoint.
Jacobi's determinant formula gives
\(d\log\det A=\operatorname{Tr}(A^{-1}dA)\).  Taking imaginary parts
proves the identity.
\end{proof}

No light/heavy determinant factorization occurs in
\eqref{eq:v3v73-phase-derivative}.

\subsection{Remote-wall independence of a localized phase variation}
\label{sec:v3v73-local-phase}

Let \(D_\infty(t)\) be the single-wall operator with the same background
on \(X_R=[-R,R]\), and suppose
\begin{equation}
 \dot D_L(t)=\dot D_\infty(t)=\dot D(t),
 \qquad \dot D(t)=\mathbf1_{X_R}\dot D(t)\mathbf1_{X_R}.
\label{eq:v3v73-local-variation}
\end{equation}
Let \(r_h\) be the rank of the finite regulator subspace supported in
\(X_R\).

\begin{theorem}[Exponential convergence of the phase connection]
\label{thm:v3v73-phase-convergence}
Under \eqref{eq:v3v73-local-variation},
\begin{equation}
 |\dot\Theta_L(t)-\dot\Theta_\infty(t)|
 \le
 \frac{8m_0r_h}{\mu^2}
 \|\dot D(t)\|e^{-\mu(L-R)}.
\label{eq:v3v73-phase-convergence}
\end{equation}
For every fixed local regulator and \(\mu>0\), the complete two-wall phase
connection therefore converges exponentially to the single-wall
connection on compactly supported background variations.
\end{theorem}

\begin{proof}
Subtract the two identities
\eqref{eq:v3v73-phase-derivative}.  The absolute value of the imaginary
part is bounded by
\begin{equation}
 \left|\operatorname{Tr}\left[
 \mathbf1_{X_R}\bigl((D_L+i\mu)^{-1}-(D_\infty+i\mu)^{-1}\bigr)
 \mathbf1_{X_R}\dot D
 \right]\right|.
\label{eq:v3v73-trace-bound-start}
\end{equation}
On an \(r_h\)-dimensional space,
\(|\operatorname{Tr}(AB)|\le r_h\|A\|\|B\|\).  Apply Theorem
\ref{thm:v3v72-local-resolvent} with \(\delta=\mu\).
\end{proof}

If the path length is measured by
\(\int_0^1\|\dot D(t)\|dt\), integration of
\eqref{eq:v3v73-phase-convergence} gives the same exponential comparison
for the total phase change along the path.

\subsection{Exact one-block gauge Ward identity}
\label{sec:v3v73-gauge-Ward}

Let a finite based gauge transformation act on the complete wall field by
a unitary matrix \(G\), extended in the auxiliary direction.  A
gauge-covariant local regulator obeys
\begin{equation}
 D_L(A^G,H^G,g)=G^{-1}D_L(A,H,g)G.
\label{eq:v3v73-gauge-covariance}
\end{equation}

\begin{theorem}[Finite complete-block gauge invariance]
\label{thm:v3v73-finite-gauge-invariance}
At every finite gauge-covariant regulator and for every \(\mu>0\),
\begin{equation}
 \det(D_L(A^G)+i\mu)=\det(D_L(A)+i\mu).
\label{eq:v3v73-det-gauge-invariance}
\end{equation}
Infinitesimally, if \(\delta_cD_L=[D_L,c]\), then
\begin{equation}
 \operatorname{Tr}\bigl[(D_L+i\mu)^{-1}[D_L,c]\bigr]=0.
\label{eq:v3v73-trace-Ward}
\end{equation}
\end{theorem}

\begin{proof}
Equation \eqref{eq:v3v73-gauge-covariance} is a similarity transformation,
and the scalar shift commutes with \(G\), proving the determinant identity.
For the infinitesimal formula, write
\([
 D_L,c]=[D_L+i\mu,c]\).  Cyclicity of the finite trace gives
\begin{align*}
 \operatorname{Tr}[(D_L+i\mu)^{-1}(D_L+i\mu)c]
 -\operatorname{Tr}[(D_L+i\mu)^{-1}c(D_L+i\mu)]
 =\operatorname{Tr}c-\operatorname{Tr}c=0.
\end{align*}
\end{proof}

The physical-wall chiral mode alone need not satisfy
\eqref{eq:v3v73-trace-Ward}; its variation and the parity-odd heavy-wall
phase are two pieces of this one zero.

\subsection{Anomaly-inflow ledger}
\label{sec:v3v73-inflow}

At fixed regulator, let \(P_L\) be the light projection of V71 and
\(Q_L=I-P_L\).  Inserting \(P_L+Q_L\) inside the trace is an exact additive
identity, but neither term is separately a determinant connection once
background variations mix the subspaces.  Define only the trace ledger
\begin{align}
 \mathcal A_{\mathrm{light},L}(c)
 &=\operatorname{Tr}\bigl[P_L(D_L+i\mu)^{-1}\delta_cD_L\bigr],
\label{eq:v3v73-light-ledger}\\
 \mathcal A_{\mathrm{heavy},L}(c)
 &=\operatorname{Tr}\bigl[Q_L(D_L+i\mu)^{-1}\delta_cD_L\bigr].
\label{eq:v3v73-heavy-ledger}
\end{align}

\begin{proposition}[Exact inflow balance]
\label{prop:v3v73-inflow-balance}
For every finite regulator,
\begin{equation}
 \mathcal A_{\mathrm{light},L}(c)
 +\mathcal A_{\mathrm{heavy},L}(c)=0.
\label{eq:v3v73-inflow-balance}
\end{equation}
If a local continuum expansion exists, the first term has the chiral wall
descent density and the second has its negative, together with
BRST-exact scheme terms.
\end{proposition}

\begin{proof}
The sum is the trace in \eqref{eq:v3v73-trace-Ward}, so it vanishes
exactly.  The statement about a continuum expansion is the only local
ghost-number-one possibility allowed by the Wess--Zumino consistency
condition in Theorem~\ref{thm:v3v65-WZ}; its cohomologically nontrivial
coefficient is the representation polynomial \eqref{eq:v3v66-I6}.
Scheme changes add local Slavnov-exact terms.
\end{proof}

The second sentence is a classification conditional on existence of the
local expansion, not a cutoff-removal proof.

\begin{corollary}[Standard Model local inflow cancels in the representation
sum]
\label{cor:v3v73-SM-inflow}
In a common regulator for all fields in one Standard Model generation,
the sum of the cohomologically nontrivial local wall terms vanishes before
\(L\to\infty\).  The corresponding sum of nontrivial local inflow terms
also vanishes.
\end{corollary}

\begin{proof}
The wall coefficient is the sum of the anomaly polynomials.  It is zero by
Theorem~\ref{thm:v3v66-bulk-descent-zero}.  Equation
\eqref{eq:v3v73-inflow-balance} then makes the opposite heavy coefficient
zero as well.  A common regulator is essential because the cancellation is
a sum of traces before limits or separate phase choices.
\end{proof}

\subsection{A diagonal wall-separation condition}
\label{sec:v3v73-diagonal}

The rank \(r_h\) in \eqref{eq:v3v73-phase-convergence} grows when the
four-dimensional regulator is removed.  The finite-cutoff result is
uniform along a diagonal family if
\begin{equation}
 r_h\,e^{-\mu(L(h)-R)}\longrightarrow0.
\label{eq:v3v73-diagonal-condition}
\end{equation}

\begin{proposition}[Sufficient separation schedule]
\label{prop:v3v73-separation-schedule}
For any positive sequence \(\varepsilon_h\downarrow0\), the choice
\begin{equation}
 L(h)\ge R+\frac1\mu
 \left(\log r_h+\log\frac1{\varepsilon_h}\right)
\label{eq:v3v73-separation-schedule}
\end{equation}
gives
\(r_he^{-\mu(L(h)-R)}\le\varepsilon_h\).
\end{proposition}

\begin{proof}
Exponentiate the negative of
\eqref{eq:v3v73-separation-schedule} and multiply by \(r_h\).
\end{proof}

This schedule is only a sufficient finite-rank bookkeeping result.  It
becomes useless as \(\mu\downarrow0\), and it does not replace a
regulator-uniform trace-ideal estimate.

\subsection{Global phase and claim boundary}
\label{sec:v3v73-global}

For \(\mu>0\), the logarithm can be continued along any fixed path, but a
closed loop may accumulate a phase when \(\mu\downarrow0\) through
spectral flow.  The local derivative estimate does not determine that
holonomy.

\begin{firewall}[Inflow balance is not determinant factorization]
\label{fw:v3v73-no-factorization}
Equations \eqref{eq:v3v73-light-ledger}--%
\eqref{eq:v3v73-heavy-ledger} split one trace by an identity resolution.
They do not define separate light and heavy determinants.  Background
mixing, basis transport, and eta phases remain in the complete determinant.
\end{firewall}

\begin{firewall}[Local cancellation leaves flat holonomy]
\label{fw:v3v73-flat-holonomy}
The Standard Model sum cancels the nontrivial local descent coefficient.
It does not prove triviality of the mapping-torus holonomy isolated in V66
and V67, nor does it control the simultaneous \(h\downarrow0\),
\(\mu\downarrow0\), and \(L\to\infty\) limit.
\end{firewall}

\subsection{V73 conclusion and V74 target}
\label{sec:v3v73-conclusion}

The complete finite two-wall determinant is exactly gauge invariant, and
its light and heavy gauge variations cancel in one trace.  Compactly
supported phase variations converge exponentially to the single-wall
connection.  In the Standard Model representation sum, the nontrivial
local descent and inflow coefficients vanish before wall separation.  The
present bound carries the local regulator rank \(r_h\).

V74 should replace that rank factor by a trace-ideal estimate uniform in
the tangential cutoff.  It should use heat-kernel or Schatten bounds for
the product of two separated resolvents and a local variation, determine
the dimension-dependent short-time subtraction, and show that the
renormalized phase connection obeys a wall-separation estimate in a fixed
Sobolev class.  If the estimate restores a power of the cutoff, the
counterterm must be identified before a diagonal limit is claimed.
\clearpage
\section{V74: ultraviolet-summable wall separation in the flat product model}
\label{sec:v3v74}

The finite-rank estimate of V73 used
\(|\operatorname{Tr}AB|\le r_h\|A\|\|B\|\) and therefore lost the
tangential ultraviolet decay.  In the flat product model, tangential
Fourier momentum increases the decay rate between the two walls.  This
note proves that the remote-wall contribution to a localized phase
variation is trace class uniformly in the tangential cutoff.

\subsection{Tangential fibre decomposition}
\label{sec:v3v74-fibres}

Put the four physical Euclidean coordinates on a flat torus of side
\(\ell\), with momenta
\begin{equation}
 q\in\mathcal Q_\ell=\frac{2\pi}{\ell}\mathbb Z^4.
\label{eq:v3v74-momenta}
\end{equation}
Assume the background path is translation invariant in these directions
for the present model calculation.  The complete domain-wall operator
decomposes as
\begin{equation}
 \mathcal D_L=\bigoplus_{q\in\mathcal Q_\ell}\mathcal D_L(q),
\qquad
 \mathcal D_L(q)^2=\mathcal H_L^2+|q|^2,
\label{eq:v3v74-fibre-square}
\end{equation}
where the tangential Clifford term anticommutes with the normal operator
\(\mathcal H_L\).  Set
\begin{equation}
 \kappa(q)=\sqrt{|q|^2+\mu^2}.
\label{eq:v3v74-kappa-q}
\end{equation}

\subsection{Fibre trace-ideal decay}
\label{sec:v3v74-fibre-decay}

Let \(X_R=[-R,R]\), let \(Y_L=[L,\infty)\), and put
\(d=L-R>0\).  Denote by \(R_L(q)=(\mathcal D_L(q)+i\mu)^{-1}\) and
\(R_\infty(q)=(\mathcal D_\infty(q)+i\mu)^{-1}\) the two-wall and
single-wall fibre resolvents.

\begin{lemma}[Momentum-enhanced separated resolvent]
\label{lem:v3v74-fibre-HS}
For the sharp step profiles and fixed \(R,m_0,\mu>0\), there are constants
\(C,c>0\), independent of \(L\) and \(q\), such that
\begin{align}
 \|\mathbf1_{X_R}R_L(q)\mathbf1_{Y_L}\|_2
 &\le C(1+\kappa(q))\kappa(q)^{-1/2}
 e^{-c\kappa(q)d},
\label{eq:v3v74-left-HS}\\
 \|\mathbf1_{Y_L}R_\infty(q)\mathbf1_{X_R}\|_2
 &\le C(1+\kappa(q))\kappa(q)^{-1/2}
 e^{-c\kappa(q)d}.
\label{eq:v3v74-right-HS}
\end{align}
Here \(\|\cdot\|_2\) is the Hilbert--Schmidt norm in the auxiliary
coordinate and spin indices.
\end{lemma}

\begin{proof}
Use
\begin{equation}
 R_L(q)=\bigl(\mathcal D_L(q)-i\mu\bigr)
 \bigl(\mathcal H_L^2+\kappa(q)^2\bigr)^{-1}.
\label{eq:v3v74-resolvent-square}
\end{equation}
Each diagonal block of \(\mathcal H_L^2\) is the one-dimensional
Schr\"odinger operator in
\eqref{eq:v3v71-QstarQ} or \eqref{eq:v3v71-QQstar}.  Away from the two
delta interfaces its resolvent kernel is a linear combination of
\(e^{-\omega|y-y'|}/(2\omega)\), where
\begin{equation}
 \omega=\sqrt{m_0^2+\kappa(q)^2}\ge\kappa(q).
\label{eq:v3v74-omega}
\end{equation}
The delta matching conditions form a two-by-two linear system.  At the
spectral point \(-\kappa(q)^2<0\), positivity of
\(\mathcal H_L^2\) bounds its inverse by \(\kappa(q)^{-2}\); hence the
matching coefficients are bounded by a fixed polynomial in
\(1+\kappa(q)\) and \(\kappa(q)^{-1}\), uniformly in \(L\).  Differentiating
the scalar kernel once in \eqref{eq:v3v74-resolvent-square} supplies at
most one factor \(\omega\).

For \(y\in X_R\) and \(y'\in Y_L\), every direct or reflected kernel path
has length at least \(d\).  Squaring the resulting bound and integrating
over the finite interval \(X_R\) and the half-line \(Y_L\) gives a factor
\(\omega^{-1}e^{-2\omega d}\).  Absorbing the fixed rational powers into
the displayed polynomial and using \(\omega\ge\kappa(q)\) proves
\eqref{eq:v3v74-left-HS}.  The single-wall proof is identical and gives
\eqref{eq:v3v74-right-HS}.
\end{proof}

The precise polynomial in this lemma is not optimal.  Exponential decay in
\(|q|d\) is the decisive feature.

\subsection{Uniform trace bound for the remote-wall difference}
\label{sec:v3v74-trace}

Let \(V_L=\mathcal D_L-\mathcal D_\infty\), supported in \(Y_L\), and let
\(\dot D\) be a localized tangential differential variation of order
\(r\ge0\), supported in \(X_R\), whose fibre symbol satisfies
\begin{equation}
 \|\dot D(q)\|\le C_{\dot D}(1+|q|)^r.
\label{eq:v3v74-variation-order}
\end{equation}

\begin{theorem}[Cutoff-uniform phase-difference trace]
\label{thm:v3v74-uniform-trace}
There are constants \(C,c>0\), independent of the tangential Fourier
cutoff \(\Lambda\) and wall separation \(L>R\), such that
\begin{align}
 &\left|\sum_{\substack{q\in\mathcal Q_\ell\\|q|\le\Lambda}}
 \operatorname{Tr}_y\left[
 \mathbf1_{X_R}\bigl(R_L(q)-R_\infty(q)\bigr)
 \mathbf1_{X_R}\dot D(q)
 \right]\right|
\notag\\
 &\hspace{20mm}\le
 C C_{\dot D}\|V_L\|
 e^{-c\mu d}
 \left(1+d^{-(6+r)}\right).
\label{eq:v3v74-uniform-trace}
\end{align}
In particular the sum converges absolutely as \(\Lambda\to\infty\), and
the remote-wall contribution tends to zero as \(L\to\infty\) uniformly in
that cutoff.
\end{theorem}

\begin{proof}
The fibre resolvent identity is
\begin{equation}
 R_L(q)-R_\infty(q)=-R_L(q)V_LR_\infty(q).
\label{eq:v3v74-fibre-resolvent-identity}
\end{equation}
Insert \(\mathbf1_{Y_L}\) around \(V_L\).  The Schatten inequality
\(|\operatorname{Tr}(AB)|\le\|A\|_2\|B\|_2\), Lemma
\ref{lem:v3v74-fibre-HS}, and
\eqref{eq:v3v74-variation-order} bound the absolute value of the
\(q\)-summand by
\begin{equation}
 C C_{\dot D}\|V_L\|
 (1+|q|)^{r+2}\kappa(q)^{-1}
 e^{-2c\kappa(q)d}.
\label{eq:v3v74-q-summand}
\end{equation}
Use
\begin{equation}
 \kappa(q)\ge\frac{\mu+|q|}{\sqrt2}.
\label{eq:v3v74-kappa-split}
\end{equation}
The number of lattice points with \(n\le|q|<n+1\) is at most
\(C_\ell(1+n)^3\).  Summing
\eqref{eq:v3v74-q-summand} by shells is therefore bounded by a constant
times
\begin{equation}
 e^{-c\mu d}
 \sum_{n\ge0}(1+n)^{r+5}e^{-cnd}
 \le C_{\ell,r}e^{-c\mu d}(1+d^{-(6+r)}).
\label{eq:v3v74-shell-sum}
\end{equation}
The majorant is independent of \(\Lambda\), so absolute convergence and
the uniform \(L\)-limit follow.
\end{proof}

The power \(6+r\) is deliberately nonsharp and arises from the polynomial
loss in Lemma~\ref{lem:v3v74-fibre-HS}.  For the wall-separation limit
\(d\to\infty\), it is harmless.

\subsection{Removal of the V73 rank factor}
\label{sec:v3v74-rank-removal}

\begin{corollary}[Commutation of the tangential cutoff and wall limit]
\label{cor:v3v74-commuting-limits}
For the flat translation-invariant product model, fixed \(\mu>0\), and a
localized variation satisfying \eqref{eq:v3v74-variation-order},
\begin{equation}
 \lim_{L\to\infty}\lim_{\Lambda\to\infty}
 (\dot\Theta_{L,\Lambda}-\dot\Theta_{\infty,\Lambda})
 =
 \lim_{\Lambda\to\infty}\lim_{L\to\infty}
 (\dot\Theta_{L,\Lambda}-\dot\Theta_{\infty,\Lambda})=0.
\label{eq:v3v74-commuting-limits}
\end{equation}
No logarithmic separation schedule involving the cutoff rank is needed.
\end{corollary}

\begin{proof}
Theorem~\ref{thm:v3v74-uniform-trace} supplies one absolutely summable
majorant for every \(\Lambda\), and its right side tends to zero with
\(L\).  Dominated convergence gives both iterated limits and their common
value.
\end{proof}

This result concerns the difference between the two-wall and single-wall
phase connections.  The single-wall connection still has local
short-distance divergences.  Their nontrivial ghost-number-one part is the
V66 anomaly polynomial, which cancels only after the common Standard Model
representation sum; BRST-exact and gauge-invariant local counterterms must
still be renormalized.

\subsection{A relative heat representation}
\label{sec:v3v74-heat}

For positive \(A\),
\begin{equation}
 (A+\kappa^2)^{-1}=\int_0^\infty e^{-t\kappa^2}e^{-tA}\,dt.
\label{eq:v3v74-heat-resolvent}
\end{equation}
Applied to \(A=\mathcal H_L^2\), the off-diagonal heat kernel carries the
Gaussian factor \(e^{-d^2/(ct)}\).  The relative trace therefore has a
small-time majorant of the form
\begin{equation}
 t^{-N}e^{-d^2/(ct)},
\label{eq:v3v74-relative-small-time}
\end{equation}
which is integrable at \(t=0\) for every \(N\) when \(d>0\).

\begin{proposition}[No new ultraviolet counterterm from wall separation]
\label{prop:v3v74-no-new-counterterm}
In the flat product model, the difference between the two-wall and
single-wall localized phase connections has no coincident-point
short-time divergence.  All ultraviolet counterterms are local to one
wall and are already present in the single-wall problem.
\end{proposition}

\begin{proof}
The resolvent identity forces every relative kernel path to travel from
\(X_R\) to \(Y_L\) and back.  The heat representation then contains the
off-diagonal Gaussian factor in
\eqref{eq:v3v74-relative-small-time}.  Exponential suppression beats every
power as \(t\downarrow0\), so the relative trace has no local heat
coefficient at the diagonal.  The remaining small-time coefficients come
from the separate single-wall connection.
\end{proof}

\subsection{Scope and next obstruction}
\label{sec:v3v74-scope}

\begin{firewall}[Flat-product theorem]
\label{fw:v3v74-flat-only}
The Fourier proof assumes a flat product collar and tangentially
translation-invariant coefficients.  It does not establish uniform
Schatten estimates for variable gauge, Higgs, metric, or material fields.
Those require covariant heat-kernel bounds with boundary/interface
conditions on one common local lattice or continuum domain.
\end{firewall}

\begin{firewall}[The infrared limit is still singular]
\label{fw:v3v74-IR}
The constants depend on \(\mu>0\).  At \(q=0\), the decay supplied by
\(\kappa(q)\) degenerates as \(\mu\downarrow0\), reflecting the genuine
massless chiral mode.  V74 does not interchange the zero-mass/zero-frequency
limit with wall separation.
\end{firewall}

\begin{firewall}[Relative finiteness is not a global phase]
\label{fw:v3v74-no-holonomy}
Ultraviolet finiteness of the relative connection does not orient the
determinant line or compute its mapping-torus holonomy.  The local Standard
Model anomaly cancellation from V66 and the global loop problem from V67
remain logically separate.
\end{firewall}

\subsection{V74 conclusion and V75 audit target}
\label{sec:v3v74-conclusion}

High tangential momentum strengthens decay across the auxiliary wall
separation.  The resulting Schatten sum removes V73's regulator-rank
factor and proves that the tangential cutoff and wall-separation limit
commute for the relative phase connection at fixed nonzero infrared
parameter in the flat product model.  The relative trace has no new local
ultraviolet counterterm.

V75 is the next compulsory NS/YM audit.  It should then replace Fourier
decomposition by covariant off-diagonal heat-kernel estimates for a weakly
curved product collar with bounded gauge/Higgs coefficients.  The target
is a trace-class relative phase connection uniform in the local regulator,
with every required single-wall counterterm placed in the protected V65
class before the Standard Model representation sum is taken.
\clearpage
\section{V75: companion audit and covariant relative heat-kernel locality}
\label{sec:v3v75}

This is the compulsory audit following V70.  It is followed by a
variable-coefficient replacement for the flat Fourier estimate of V74.
The new theorem concerns the relative two-wall/single-wall phase
connection on a bounded-geometry product collar.  It does not estimate the
unsubtracted single-wall determinant.

\subsection{Read-only NS/YM checkpoint}
\label{sec:v3v75-audit}

The active companion snapshots read at this checkpoint were
\begin{enumerate}
\item
\texttt{Renewal\_NS\_Stability\_Research\_Notes\_V31.tex}, with
\(887537\) bytes and SHA--256
\begin{center}\scriptsize\ttfamily
F1121B62238AD5491BB7CF03635EA9934942EDF573ED8FAAA9EE79E1D38A6696;
\end{center}
\item
\texttt{Renewal\_Yang\_Mills\_Mass\_Gap\_Research\_Notes\_V100.tex},
with \(1706880\) bytes and SHA--256
\begin{center}\scriptsize\ttfamily
EA766F10712335DCDFF8EACF2C8CB190C7968361EEC93CB37F7D740161BB51C3.
\end{center}
\end{enumerate}
The YM folder also contained its immediate V99 predecessor while the
parallel V100 freeze was in progress.  Neither it nor any other companion
or frozen file was modified.

NS V31 remains at its exact separation between total critical coarea and
the signed polarization hinge: the latter has a balanced-sector nullspace
and cannot control the former.  The analogue here is that ultraviolet
finiteness of a relative wall trace cannot control the unsubtracted
single-wall determinant.  These are different sectors, not two estimates
of one quantity.

YM V100 proves that a disconnected matching word contributes only after a
first component-breaking coordinate is retained, and it keeps the matching
responses and breaker in one complete packet.  The resolvent identity in
V74 has exactly that structure: the coefficient difference supported at
the remote wall is the breaker, and the two separated resolvents connect it
to the local variation.  Estimating either resolvent before this product is
assembled would lose the off-diagonal ultraviolet gain.

\begin{assessment}[V75 audit transfer]
\label{ass:v3v75-audit-transfer}
V75 keeps
\[
 \mathbf1_K R_L\mathbf1_Y\,V_L\,
 \mathbf1_YR_\infty\mathbf1_K\dot D
\]
as one trace-ideal product.  The NS audit forbids using the resulting
relative estimate as a single-wall renormalization theorem, and the YM
audit fixes the order in which the remote support and both propagators are
assembled.
\end{assessment}

\subsection{Bounded-geometry collar hypotheses}
\label{sec:v3v75-hypotheses}

Let \(Z=M^4\times\mathbb R_y\) carry a complete metric of bounded geometry,
uniformly equivalent to a product metric.  Let \(D_L\) and \(D_\infty\) be
self-adjoint Dirac-type operators with common principal Clifford symbol.
Assume:
\begin{enumerate}
\item their connection, curvature, gauge, Higgs, and endomorphism
coefficients have uniform bounds through the finite derivative order used
below;
\item \(D_L=D_\infty\) on \(y<L\);
\item \(V_L=D_L-D_\infty\) is a bounded endomorphism supported in
\(Y_L=\{y\ge L\}\), with \(\sup_L\|V_L\|<\infty\); and
\item the local background variation \(\dot D\) is a differential operator
of order \(r\) with compact support
\(K\subset M^4\times[-R,R]\), where \(L>R\).
\end{enumerate}
Put \(d=\operatorname{dist}(K,Y_L)\), so uniform product equivalence gives
\(d\ge c_0(L-R)\).

The squares \(A_j=D_j^2\) are nonnegative Laplace-type operators.
Their local heat-kernel constants are uniform under the stated
bounded-geometry assumptions.

\subsection{Davies--Gaffney and local Hilbert--Schmidt bounds}
\label{sec:v3v75-DG}

\begin{lemma}[Covariant off-diagonal heat bound]
\label{lem:v3v75-DG}
There are uniform constants \(C,c>0\) such that for measurable
\(E,F\subset Z\),
\begin{equation}
 \|\mathbf1_Ee^{-tA_j}\mathbf1_F\|
 \le C\exp\left(Ct-\frac{\operatorname{dist}(E,F)^2}{ct}\right),
 \qquad t>0.
\label{eq:v3v75-DG}
\end{equation}
For compact \(K\) and every differential operator \(B\) of order \(s\)
with bounded coefficients supported in \(K\),
\begin{equation}
 \|B e^{-tA_j}\|_2
 \le C_{K,B}t^{-(5+2s)/4}e^{Ct},
 \qquad0<t\le1.
\label{eq:v3v75-local-HS}
\end{equation}
\end{lemma}

\begin{proof}
For a real Lipschitz \(\varphi\) with \(|d\varphi|\le1\), conjugating the
quadratic form of \(A_j\) by \(e^{a\varphi}\) changes it by terms bounded
by \(Ca^2\|u\|^2\).  The energy inequality for the conjugated heat
equation gives
\begin{equation}
 \|e^{a\varphi}e^{-tA_j}e^{-a\varphi}\|
 \le Ce^{Ca^2t+Ct}.
\label{eq:v3v75-weighted-heat}
\end{equation}
Choose \(\varphi=0\) on \(F\),
\(\varphi\ge\operatorname{dist}(E,F)\) on \(E\), and optimize in \(a\)
to obtain \eqref{eq:v3v75-DG}.

On a five-dimensional bounded-geometry manifold the on-diagonal heat
kernel satisfies \(|K_t(x,x)|\le Ct^{-5/2}e^{Ct}\).  Covariant derivatives
of total order \(s\) cost at most \(Ct^{-s/2}\).  Integrating the squared
kernel over the compact support of \(B\) gives
\eqref{eq:v3v75-local-HS}.  Uniform bounded geometry makes all constants
independent of \(L\).
\end{proof}

\begin{corollary}[Separated heat kernel is smoothing to every order]
\label{cor:v3v75-separated-smoothing}
For every \(N,s\ge0\),
\begin{equation}
 \|B e^{-tA_j}\mathbf1_{Y_L}\|_2
 \le C_{N,K,B}t^{-N}
 \exp\left(Ct-\frac{d^2}{ct}\right),
 \qquad0<t\le1.
\label{eq:v3v75-separated-HS}
\end{equation}
\end{corollary}

\begin{proof}
Split the heat semigroup at \(t/2\).  Apply the local Hilbert--Schmidt
bound to the factor adjacent to \(B\) and the Davies--Gaffney bound to the
factor crossing from a fixed enlargement of \(K\) to \(Y_L\).  Polynomial
powers may be enlarged to any stated \(N\) because the off-diagonal
exponential dominates them as \(t\downarrow0\).
\end{proof}

\subsection{Covariant trace-class relative resolvent}
\label{sec:v3v75-relative-resolvent}

For \(\mu>0\), use
\begin{equation}
 (D_j+i\mu)^{-1}
 =(D_j-i\mu)
 \int_0^\infty e^{-t(D_j^2+\mu^2)}\,dt.
\label{eq:v3v75-resolvent-heat}
\end{equation}
\begin{theorem}[Variable-coefficient relative phase trace]
\label{thm:v3v75-relative-trace}
Under the hypotheses above,
\begin{equation}
 T_L:=mathbf1_K
 \bigl[(D_L+i\mu)^{-1}-(D_\infty+i\mu)^{-1}\bigr]
 \mathbf1_K\dot D
\label{eq:v3v75-TL}
\end{equation}
is trace class.  There are constants \(C,c>0\), depending on the bounded
geometry, \(K\), finitely many coefficient bounds, \(r\), and \(\mu\), but
not on \(L\), such that
\begin{equation}
 \|T_L\|_1
 \le C\|V_L\|
 e^{-c\mu d}\left(1+d^{-N_r}\right)
\label{eq:v3v75-relative-trace}
\end{equation}
for some finite \(N_r\).  In particular \(\|T_L\|_1\to0\) as
\(L\to\infty\).
\end{theorem}

\begin{proof}
The resolvent identity gives
\begin{equation}
 T_L=-\mathbf1_KR_L\mathbf1_{Y_L}V_L
 \mathbf1_{Y_L}R_\infty\mathbf1_K\dot D,
\label{eq:v3v75-breaker-product}
\end{equation}
where \(R_j=(D_j+i\mu)^{-1}\).  This is the complete breaker product
identified in Assessment~\ref{ass:v3v75-audit-transfer}.

Apply \eqref{eq:v3v75-resolvent-heat} to both separated resolvents.
Corollary~\ref{cor:v3v75-separated-smoothing}, with one additional
derivative for each numerator \(D_j-i\mu\) and \(r\) derivatives for
\(\dot D\), makes both factors adjacent to \(V_L\) Hilbert--Schmidt.  The
Schatten inequality
\begin{equation}
 \|AB\|_1\le\|A\|_2\|B\|_2
\label{eq:v3v75-Schatten}
\end{equation}
then proves trace class.

For the norm, integrate the heat bounds.  At small time the exponent is
\(-\mu^2t-d^2/(ct)\); its minimum over \(t>0\) is a positive multiple of
\(\mu d\).  Polynomial heat powers produce only a finite factor
\(1+d^{-N_r}\).  At large time, split each heat factor at time one; the local
Hilbert--Schmidt factor is then uniformly bounded by its time-one value,
the remaining semigroup is a contraction, and \(e^{-\mu^2t}\) is
integrable for every \(\mu>0\).  Multiplication by \(V_L\) contributes
\(\|V_L\|\), yielding
\eqref{eq:v3v75-relative-trace}.
\end{proof}

\begin{corollary}[Regulator-independent wall separation]
\label{cor:v3v75-regulator-independent}
Any sequence of local lattice or spectral approximants whose heat kernels
obey the uniform versions of
\eqref{eq:v3v75-DG}--\eqref{eq:v3v75-local-HS} has the same bound
\eqref{eq:v3v75-relative-trace}.  The wall limit may therefore be taken
without the rank schedule \eqref{eq:v3v73-separation-schedule}.
\end{corollary}

\begin{proof}
The proof of Theorem~\ref{thm:v3v75-relative-trace} uses only the stated
uniform heat bounds, the resolvent identity, and Schatten inequalities.
It passes unchanged to every approximant.
\end{proof}

The corollary is a criterion.  V75 does not yet construct a chiral lattice
operator satisfying those estimates and the exact gauge covariance of
V73 simultaneously.

\subsection{Renormalization and the protected symbol}
\label{sec:v3v75-renormalization}

Equation \eqref{eq:v3v75-breaker-product} forces every relative heat path
to traverse the positive distance \(d\) twice.  It therefore has no
coincident small-time coefficient.  All ultraviolet subtractions belong
to the single-wall determinant.

The local single-wall counterterms split into:
\begin{enumerate}
\item Slavnov-exact gauge-fixing terms and squared-constraint terms, which
compress to zero on the V64 physical quotient and lie in the first two
classes of \eqref{eq:v3v65-safe-decomposition};
\item parity-odd phase representatives, whose Euclidean quadratic real
part vanishes but whose differential order must still be checked as in
V66; and
\item gauge-invariant kinetic and potential renormalizations, whose
physical coefficients require explicit renormalization conditions and
cannot be declared positive by cohomology.
\end{enumerate}

\begin{proposition}[Protected counterterms remain separated from the
relative estimate]
\label{prop:v3v75-protected-separation}
The relative trace in Theorem~\ref{thm:v3v75-relative-trace} requires no
new local counterterm.  Counterterms of the first class above preserve the
physical principal symbol exactly; bounded order-zero terms preserve its
high-frequency gap by Theorem~\ref{thm:v3v65-protected-class}.
\end{proposition}

\begin{proof}
Off-diagonal smoothing proves the first assertion.  The remaining two are
the compression identities and coercive estimate in Theorem
\ref{thm:v3v65-protected-class}.
\end{proof}

\subsection{Claim boundary}
\label{sec:v3v75-boundary}

\begin{firewall}[Relative trace versus single-wall determinant]
\label{fw:v3v75-relative-only}
V75 proves a trace-class estimate for the two-wall/single-wall difference.
It does not define or renormalize the single-wall determinant, remove its
local counterterms, or control its mapping-torus holonomy.  This is the
sector distinction required by the NS audit.
\end{firewall}

\begin{firewall}[Heat bounds are analytic hypotheses for future
regulators]
\label{fw:v3v75-regulator-hypothesis}
The continuum bounded-geometry theorem is proved under standard uniform
heat-kernel estimates.  A finite gauge-covariant, reflection-positive
fermion regulator realizing the domain wall and satisfying the same bounds
has not been constructed.  Corollary
\ref{cor:v3v75-regulator-independent} must not be cited as its existence
proof.
\end{firewall}

\begin{firewall}[Infrared and global phase remain]
\label{fw:v3v75-IR-global}
The constants depend on \(\mu>0\).  The massless zero mode, global
determinant-line orientation, and simultaneous zero-infrared limit remain
open even though wall separation is ultraviolet uniform.
\end{firewall}

\subsection{V75 conclusion and V76 target}
\label{sec:v3v75-conclusion}

The flat Fourier mechanism of V74 survives bounded curvature and variable
gauge/Higgs coefficients at the continuum heat-kernel level.  The
relative resolvent identity supplies one remote-wall breaker, and covariant
Davies--Gaffney plus Hilbert--Schmidt estimates make the complete product
trace class with exponential wall-separation decay.  It introduces no new
ultraviolet counterterm.

V76 should construct a concrete finite, time-local domain-wall lattice
operator with exact internal gauge covariance and prove the uniform heat
bounds required by Corollary~\ref{cor:v3v75-regulator-independent}.  It
must then test fermionic reflection positivity on the complete two-wall
block.  The Wilson term needed to remove lattice doublers must be tracked
through the wall index and anomaly inflow; it cannot be treated as a
harmless lower-order perturbation.
\clearpage
\section{V76: a gauge-covariant Wilson domain-wall lattice and its finite reflection form}
\label{sec:v3v76}

V75 reduced the regulator question to a concrete existence problem.  This
note constructs a five-dimensional nearest-neighbour Wilson operator,
tracks the Wilson term through the free interface index, proves
volume-uniform off-diagonal heat estimates under an admissibility bound,
and establishes reflection positivity for the complete two-wall
Grassmann block.  The last statement includes the mirror wall; no chiral
determinant quotient is taken.

\subsection{Finite covariant lattice operator}
\label{sec:v3v76-operator}

Let
\begin{equation}
 \Lambda_{a,N}\times S_{a_5,N_5}
\label{eq:v3v76-lattice}
\end{equation}
be a finite four-dimensional reflection-symmetric lattice and a periodic
fifth lattice.  Write \(z=(x,s)\), and let \(A=1,\ldots,5\) label all
directions, with spacings \(a_A=a\) for \(A\le4\) and \(a_5\) for the
last direction.  Gauge links are unitary matrices \(U_A(z)\); for the
minimal domain-wall construction one may take \(U_5=I\) and extend the
four-dimensional links independently of \(s\).

Define unitary covariant shifts
\begin{align}
 (T_A^+\psi)(z)&=U_A(z)\psi(z+\widehat A),
\label{eq:v3v76-forward-shift}\\
 (T_A^-\psi)(z)&=U_A(z-\widehat A)^*\psi(z-\widehat A),
\label{eq:v3v76-backward-shift}
\end{align}
so \((T_A^+)^*=T_A^-\).  Put
\begin{equation}
 \nabla_A=\frac{T_A^+-T_A^-}{2a_A},
 \qquad
 \Delta_A=\frac{2-T_A^+-T_A^-}{a_A^2}.
\label{eq:v3v76-differences}
\end{equation}
Then \(\nabla_A^*=-\nabla_A\) and \(\Delta_A\ge0\).  With Hermitian
Euclidean Clifford matrices \(\gamma_A\), Wilson parameter \(r>0\), and a
real two-wall mass \(m(s)\), set
\begin{equation}
 D_{\mathrm{DW}}(U,m)=
 \sum_{A=1}^5\gamma_A\nabla_A
 +\sum_{A=1}^5\frac{r a_A}{2}\Delta_A+m(s).
\label{eq:v3v76-DW}
\end{equation}
Every term couples either one lattice link or one site.

\begin{proposition}[Exact finite gauge covariance]
\label{prop:v3v76-gauge-covariance}
For a sitewise unitary gauge transformation \(g(z)\), let
\begin{equation}
 U_A^g(z)=g(z)U_A(z)g(z+\widehat A)^{-1},
 \qquad \psi^g(z)=g(z)\psi(z).
\label{eq:v3v76-gauge-action}
\end{equation}
Then
\begin{equation}
 D_{\mathrm{DW}}(U^g,m)=gD_{\mathrm{DW}}(U,m)g^{-1}.
\label{eq:v3v76-exact-covariance}
\end{equation}
Consequently its finite determinant and every polynomial in
\(D_{\mathrm{DW}}^*D_{\mathrm{DW}}\) are exactly invariant under based
gauge transformations.
\end{proposition}

\begin{proof}
Substitution in \eqref{eq:v3v76-forward-shift} gives
\(T_A^+(U^g)=gT_A^+(U)g^{-1}\); the backward identity follows by taking
adjoints.  Equations \eqref{eq:v3v76-differences} and
\eqref{eq:v3v76-DW} inherit the same conjugation, while the real scalar
mass commutes with \(g\).  Determinants and functional calculus are
invariant under unitary conjugation.
\end{proof}

\subsection{The Wilson one-species window}
\label{sec:v3v76-index}

The fifth-direction wall interpolates between two bulk Wilson phases.  To
count their difference, consider the free four-dimensional Wilson symbol
\begin{equation}
 D_W(p;M)=
 \frac{i}{a}\sum_{\mu=1}^4\gamma_\mu\sin(ap_\mu)
 +\frac r a\sum_{\mu=1}^4(1-\cos(ap_\mu))-M.
\label{eq:v3v76-Wilson-symbol}
\end{equation}
Away from its zeros, normalize the five-vector
\begin{equation}
 h_M(p)=\left(sin(ap_1),\ldots,sin(ap_4),
 r\sum_{\mu}(1-\cos(ap_\mu))-aM\right)
\label{eq:v3v76-h-map}
\end{equation}
to obtain \(n_M:T^4\to S^4\).

\begin{theorem}[Free Wilson phase degree]
\label{thm:v3v76-Wilson-degree}
If \(aM\notin\{0,2r,4r,6r,8r\}\), the degree of
\(n_M\), with the orientation chosen so the first window is positive, is
\begin{equation}
 \nu(M)=-\frac12\sum_{k=0}^4
 (-1)^k{4\choose k}\operatorname{sgn}(2rk-aM).
\label{eq:v3v76-degree-formula}
\end{equation}
In particular
\begin{equation}
 \nu(M)=1\quad\hbox{for}\quad0<aM<2r,
 \qquad
 \nu(M)=0\quad\hbox{for}\quad aM<0.
\label{eq:v3v76-one-species-window}
\end{equation}
Thus a wall between these phases carries one net chiral species; the
fifteen nonzero Brillouin-zone corners do not create light wall doublers.
\end{theorem}

\begin{proof}
The first four components of \(h_M\) vanish only at corners with each
\(ap_\mu\in\{0,\pi\}\).  A corner containing \(k\) entries equal to
\(\pi\) has fifth component \(2rk-aM\), multiplicity \({4\choose k}\),
and local orientation \((-1)^k\), because differentiating
\(\sin(ap_\mu)\) at \(\pi\) reverses one orientation.  Summing the local
degrees at the two poles gives \eqref{eq:v3v76-degree-formula}.

If \(aM<0\), every sign is positive and the binomial sum is
\((1-1)^4=0\).  If \(0<aM<2r\), the \(k=0\) sign is negative and every
other sign is positive.  Hence the sum in
\eqref{eq:v3v76-degree-formula} is \(-2\), giving \(\nu=1\).  Only the
zero-momentum corner changes phase across this wall; all other corners
retain a nonzero Wilson mass.
\end{proof}

\begin{corollary}[Finite periodic fifth direction has two interfaces]
\label{cor:v3v76-two-interfaces}
On the periodic fifth lattice, a mass profile which enters the
\(\nu=1\) phase must later leave it.  The two jumps have opposite index,
so the complete finite operator is vectorlike.  If the bulk symbols on
both plateaux have a gap \(g_5>0\), the free transfer roots give a
light-wall mixing bound
\begin{equation}
 |E_{N_5}|\le Cq^{N_{\rm sep}},
 \qquad0<q<1,
\label{eq:v3v76-lattice-mixing}
\end{equation}
where \(N_{\rm sep}\) is the number of fifth-direction links between the
walls and \(q\) depends on the bulk gap but not on the total volume.
\end{corollary}

\begin{proof}
Periodicity makes the signed sum of phase jumps zero, as in Proposition
\ref{prop:v3v70-circle-zero}.  On either constant plateau, the zero-energy
fifth recurrence has a finite transfer matrix with no unit-circle
eigenvalue because of the gap \(g_5\).  Compactness of the Brillouin torus
gives a uniform maximum modulus \(q<1\) for all stable roots.  Propagation
from one interface to the other therefore costs at most
\(Cq^{N_{\rm sep}}\).  The off-diagonal two-mode matrix element has the
same bound.
\end{proof}

The corollary is an upper bound.  Its continuum sharp-wall analogue is the
exact equation in V71.

\subsection{Uniform discrete heat estimates}
\label{sec:v3v76-heat}

Assume comparable spacings \(c_a a\le a_5\le C_a a\), and let
\(\rho=\max(a,a_5)\).  Put
\begin{equation}
 A_{a,L}=D_{\mathrm{DW}}^*D_{\mathrm{DW}}.
\label{eq:v3v76-positive-square}
\end{equation}
Assume the physical links satisfy the admissibility and bounded-coefficient
conditions
\begin{equation}
 \|I-U_{AB}(z)\|\le C_Fa_Aa_B,
 \qquad
 \|m\|_\infty+\rho\|\nabla_5m\|_\infty\le C_m,
\label{eq:v3v76-admissibility}
\end{equation}
with constants independent of volume and lattice spacing.  The first
condition bounds the spin-curvature terms in the lattice Weitzenbock
formula.

\begin{lemma}[Uniform lattice G\aa rding comparison]
\label{lem:v3v76-Garding}
In a closed subwindow of \(0<aM<2r\), and for an admissibility constant
small enough relative to its distance from the Wilson thresholds, there
are constants \(c,C>0\), independent of volume and spacing, such that
\begin{equation}
 c\sum_A\|\nabla_Au\|_2^2-C\|u\|_2^2
 \le\langle u,A_{a,L}u\rangle
 \le C\left(\sum_A\|\nabla_Au\|_2^2+\|u\|_2^2
 +\sum_Aa_A^2\|\Delta_Au\|_2^2\right).
\label{eq:v3v76-Garding}
\end{equation}
\end{lemma}

\begin{proof}
For trivial links, Fourier transformation reduces the claim to the sum of
squares of the sine symbol and the Wilson mass.  Near the zero corner the
sine term controls the lattice gradient up to a bounded mass term; near
every other corner the Wilson mass is bounded below by the distance from
the thresholds.  Compactness of the Brillouin torus supplies uniform
constants.  Also \(a_A^2\Delta_A^2\le C\Delta_A\) on the lattice spectrum.
For admissible links, commuting covariant shifts produces plaquette
curvatures of size \(O(a_Aa_B)\); the mass commutator is bounded by the
second condition in \eqref{eq:v3v76-admissibility}.  Cauchy's inequality
absorbs these terms into a small part of the free gradient form plus
\(C\|u\|_2^2\).  This proves the comparison.
\end{proof}

\begin{theorem}[Volume-uniform lattice heat bounds]
\label{thm:v3v76-lattice-heat}
There are \(C,c>0\), independent of the finite volume, such that for
lattice sets \(E,F\), physical distance \(d(E,F)\), and \(t\ge c\rho^2\),
\begin{equation}
 \|\mathbf1_Ee^{-tA_{a,L}}\mathbf1_F\|
 \le C\exp\left(Ct-\frac{d(E,F)^2}{ct}\right).
\label{eq:v3v76-discrete-DG}
\end{equation}
For \(0<t<c\rho^2\), finite hopping range gives
\begin{equation}
 \|\mathbf1_Ee^{-tA_{a,L}}\mathbf1_F\|
 \le C\exp\left[-c\frac d\rho
 \log\left(1+\frac{d\rho}{t}\right)
ight].
\label{eq:v3v76-short-hop}
\end{equation}
With lattice measure \(a^4a_5\), the on-diagonal bound is
\begin{equation}
 |e^{-tA_{a,L}}(z,z)|
 \le C(t+\rho^2)^{-5/2}e^{Ct}.
\label{eq:v3v76-on-diagonal}
\end{equation}
The same constants apply to wall separations and volumes for which
\eqref{eq:v3v76-admissibility} is uniform.
\end{theorem}

\begin{proof}
The quadratic form of \(A_{a,L}\) is the norm square of the finite-range
operator \eqref{eq:v3v76-DW}.  Conjugating by a lattice weight
\(e^{\alpha\varphi}\), with \(\varphi\) one-Lipschitz in physical distance
and \(\alpha\rho\le1\), changes the form by at most
\(C(\alpha^2+1)\|u\|^2\).  Unitary links do not change this estimate;
\eqref{eq:v3v76-admissibility} bounds the commutator and potential terms.
The weighted semigroup argument used in Lemma~\ref{lem:v3v75-DG}, optimized
in \(\alpha\), gives \eqref{eq:v3v76-discrete-DG} once diffusion has
crossed one lattice spacing.

For shorter time, \(A_{a,L}\) has hopping range at most \(2\rho\).  In the
series
\(e^{-tA}=\sum_{n\ge0}(-tA)^n/n!\), a matrix element between sets at
distance \(d\) vanishes for \(n<d/(2\rho)\).  Stirling's inequality and
\(\|A\|\le C\rho^{-2}\) yield \eqref{eq:v3v76-short-hop} after adjusting
constants.

Lemma~\ref{lem:v3v76-Garding} and the scalar lattice Nash inequality
give
\begin{equation}
 \|u\|_2^{2+4/5}
 \le C\langle u,(A_{a,L}+C)u\rangle\|u\|_1^{4/5},
\label{eq:v3v76-Nash}
\end{equation}
uniformly under \eqref{eq:v3v76-admissibility}.  The standard Nash
differential inequality for
\(\|e^{-tA}u\|_2^2\) gives the ultracontractive factor
\((t+\rho^2)^{-5/4}\) from \(L^1\) to \(L^2\).  Squaring by the semigroup
property proves \eqref{eq:v3v76-on-diagonal}.
\end{proof}

\begin{corollary}[V75 relative trace criterion is realized at fixed
admissibility]
\label{cor:v3v76-V75-realized}
The lattice family \eqref{eq:v3v76-DW} satisfies uniform versions of the
off-diagonal and local Hilbert--Schmidt estimates used in V75.  At fixed
\(\mu>0\) and positive physical wall separation, the relative
two-wall/single-wall phase connection is trace class with a bound
independent of volume and of a refining admissible lattice.
\end{corollary}

\begin{proof}
Combine \eqref{eq:v3v76-discrete-DG}--%
\eqref{eq:v3v76-on-diagonal} exactly as in Corollary
\ref{cor:v3v75-separated-smoothing} and Theorem
\ref{thm:v3v75-relative-trace}.  The short-time hopping estimate replaces
the continuum off-diagonal Gaussian below the lattice scale and is even
more strongly summable at fixed positive separation.
\end{proof}

\subsection{Finite fermionic reflection form}
\label{sec:v3v76-reflection}

Use an even open Euclidean-time lattice and reflect through the link
between times \(0\) and \(1\).  Let \(\vartheta\) be the standard
antilinear Grassmann reflection, reversing product order and satisfying
\begin{equation}
 \vartheta\psi(x,s)=\bar\psi(\vartheta x,s)\gamma_4,
 \qquad
 \vartheta\bar\psi(x,s)=\gamma_4\psi(\vartheta x,s).
\label{eq:v3v76-Grassmann-reflection}
\end{equation}
Choose temporal gauge on the links crossing the reflection plane and put
\(P_\pm=(I\pm\gamma_4)/2\).

\begin{lemma}[Positive Grassmann cone]
\label{lem:v3v76-Grassmann-cone}
If the cross-plane exponential is a finite sum
\begin{equation}
 e^{-S_\times}=\sum_j c_j\,\vartheta(F_j)F_j,
 \qquad c_j\ge0,
\label{eq:v3v76-positive-cone}
\end{equation}
then, for every polynomial \(F\) supported at positive times,
\begin{equation}
 \int \vartheta(F)F e^{-S}\,d\bar\psi d\psi\ge0
\label{eq:v3v76-fermion-RP}
\end{equation}
whenever \(S=S_++\vartheta S_++S_\times\) and the Berezin orientation is
chosen by the reflected pairs.
\end{lemma}

\begin{proof}
Insert \eqref{eq:v3v76-positive-cone}.  The integral factorizes into
negative and positive Grassmann variables.  Reflection and reversed order
make the negative-half integral the complex conjugate of the positive-half
integral.  Each summand is therefore
\(c_j|\int_+F F_j e^{-S_+}|^2\ge0\).
\end{proof}

\begin{theorem}[Reflection positivity of the complete two-wall Wilson
block]
\label{thm:v3v76-complete-RP}
For nonnegative temporal hopping coefficient, the Wilson action associated
with \eqref{eq:v3v76-DW} has a cross-plane exponential of the form
\eqref{eq:v3v76-positive-cone}.  Hence the complete finite two-wall
Grassmann functional is reflection positive.  The fifth-direction Wilson
term and the sign-changing mass do not alter the proof because they are
contained within individual time halves.
\end{theorem}

\begin{proof}
The only terms crossing the time-reflection link are the Wilson temporal
hops.  In temporal gauge they are sums over spatial and fifth-direction
sites of
\begin{equation}
 -2\kappa_4\left(
  \bar\psi_0P_-\psi_1+\bar\psi_1P_+\psi_0
 \right),
 \qquad \kappa_4\ge0.
\label{eq:v3v76-cross-hopping}
\end{equation}
Choose orthonormal bases of the ranges of \(P_+\) and \(P_-\).  With the
reflection \eqref{eq:v3v76-Grassmann-reflection}, every component bilinear
in \(-S_\times\) is \(2\kappa_4\vartheta(C_j)C_j\) for a positive-time
linear Grassmann generator \(C_j\), after the fixed reflected-pair
orientation sign is included.  Expanding the finite Grassmann exponential
produces \eqref{eq:v3v76-positive-cone} with coefficients that are products
of \(2\kappa_4\), hence nonnegative.  Apply Lemma
\ref{lem:v3v76-Grassmann-cone}.  All spatial, fifth-direction, mass, Yukawa,
and on-slice gauge terms belong to \(S_++\vartheta S_+\) when their
coefficients are reflection covariant.
\end{proof}

For dynamical gauge links, the usual reflection-positive plaquette action
can be multiplied into the same cone.  This statement remains at fixed
background metric and does not yet include the Euclidean Einstein contour.

\subsection{Claim boundary and next target}
\label{sec:v3v76-boundary}

\begin{firewall}[Complete block, not a chiral quotient]
\label{fw:v3v76-complete-only}
Theorem~\ref{thm:v3v76-complete-RP} applies to the finite vectorlike
kink--antikink Grassmann system.  Dividing by a Pauli--Villars determinant,
taking an overlap limit, or discarding the remote wall changes the
functional.  Reflection positivity of any such chiral determinant ratio
has not been proved.
\end{firewall}

\begin{firewall}[Admissible gauge sector]
\label{fw:v3v76-admissible-only}
The regulator-uniform heat bounds use the plaquette admissibility condition
\eqref{eq:v3v76-admissibility}.  A proof that the full interacting gauge
measure concentrates in, or can be restricted to, this sector without
changing the continuum theory remains open.
\end{firewall}

\begin{firewall}[Index degree is a free principal result]
\label{fw:v3v76-free-index}
Theorem~\ref{thm:v3v76-Wilson-degree} proves the free Wilson phase and
doubler count.  Stability under admissible links requires a uniform gap of
the Hermitian Wilson kernel.  Exact gauge covariance alone does not prove
that gap.
\end{firewall}

\subsection{V76 conclusion and V77 target}
\label{sec:v3v76-conclusion}

The programme now has one concrete finite, nearest-neighbour domain-wall
lattice with exact internal gauge covariance.  The Wilson degree formula
selects one interface species in \(0<aM<2r\), the admissible positive square
has the uniform heat bounds required by V75, and the complete two-wall
Grassmann block is reflection positive.  These statements retain both
walls and their heavy spectrum in one regulator.

V77 should derive the finite-wall Schur complement seen by the physical
boundary and take the infinite-fifth-direction limit to the overlap
operator.  It must prove the Ginsparg--Wilson relation algebraically,
identify the exact lattice chiral Jacobian, and compare its Standard Model
representation sum with V66.  The Pauli--Villars or determinant-ratio step
must be kept explicit, because its reflection positivity is the next
unresolved gate.
\clearpage
\section{V77: the overlap Schur complement, Ginsparg--Wilson algebra, and measure obstruction}
\label{sec:v3v77}

V76 proves reflection positivity only for the complete two-wall Grassmann
block.  To identify the physical-wall operator, one must eliminate the
fifth-direction interior and normalize the heavy determinant.  This note
derives the finite transfer approximation to the sign function, proves its
overlap limit and Ginsparg--Wilson relation, and writes the chiral measure
curvature explicitly.  Reflection positivity of the determinant ratio is
not assumed.

\subsection{Transfer approximation to the sign function}
\label{sec:v3v77-transfer}

Let \(H_W\) be a finite Hermitian Wilson kernel in the one-species window
of V76.  Assume for this note that it has the uniform gap and normalization
\begin{equation}
 \sigma(H_W)\subset[-1,-\delta]\cup[\delta,1],
 \qquad0<\delta<1.
\label{eq:v3v77-gap}
\end{equation}
The normalization is obtained by multiplying the kernel by a positive
constant.  Define
\begin{equation}
 \epsilon_N(H_W)=
 \frac{(I+H_W)^N-(I-H_W)^N}
      {(I+H_W)^N+(I-H_W)^N}.
\label{eq:v3v77-rational-sign}
\end{equation}
The denominator is positive and invertible by spectral calculus.

\begin{lemma}[Uniform sign approximation]
\label{lem:v3v77-sign-error}
Put \(q=(1-\delta)/(1+\delta)<1\).  Then
\begin{equation}
 \|\epsilon_N(H_W)-\operatorname{sgn}(H_W)\|
 \le2q^N.
\label{eq:v3v77-sign-error}
\end{equation}
\end{lemma}

\begin{proof}
For an eigenvalue \(\lambda>0\), divide numerator and denominator in
\eqref{eq:v3v77-rational-sign} by \((1+\lambda)^N\).  With
\(t=(1-\lambda)/(1+\lambda)\),
\begin{equation}
 1-\epsilon_N(\lambda)=\frac{2t^N}{1+t^N}\le2q^N.
\label{eq:v3v77-positive-error}
\end{equation}
For \(\lambda<0\), oddness gives the same estimate.  The finite spectral
theorem turns the scalar supremum into the operator norm.
\end{proof}

\subsection{Exact finite-wall Schur identity}
\label{sec:v3v77-Schur}

Use a Shamir form of the V76 nearest-neighbour fifth operator, with boundary
fields \(q,\bar q\) obtained from the two endpoint chiral components.  Let
\(D_{\mathrm{DW},N}(m_f)\) denote its block tridiagonal matrix at wall mass
\(m_f\), and use \(m_f=1\) for the Pauli--Villars heavy normalization.
Define
\begin{equation}
 D_N(m_f)=
 \frac{1+m_f}{2}I+\frac{1-m_f}{2}
 \gamma_5\epsilon_N(H_W).
\label{eq:v3v77-effective-N}
\end{equation}

\begin{theorem}[Finite fifth-direction determinant identity]
\label{thm:v3v77-Schur-identity}
With consistent finite-dimensional determinant orientations,
\begin{equation}
 \frac{\det D_{\mathrm{DW},N}(m_f)}
      {\det D_{\mathrm{DW},N}(1)}
 =\det D_N(m_f).
\label{eq:v3v77-determinant-ratio}
\end{equation}
The boundary-to-boundary propagator is the inverse of the same Schur
complement \(D_N(m_f)\), up to the fixed boundary-field normalization.
\end{theorem}

\begin{proof}
Order the fifth slices from one wall to the other and perform block Gaussian
elimination without taking determinants of the boundary block.  The bulk
recurrence transports one chirality by
\begin{equation}
 T=(I-H_W)(I+H_W)^{-1}.
\label{eq:v3v77-transfer-matrix}
\end{equation}
After \(N\) slices, the boundary relation contains
\((I-T^N)(I+T^N)^{-1}=\epsilon_N(H_W)\), which is
\eqref{eq:v3v77-rational-sign}.  The determinant is the product of the
bulk pivots and the boundary Schur determinant.  The bulk pivots are
independent of \(m_f\) and equal those at \(m_f=1\), so they cancel in the
ratio, leaving \(\det D_N(m_f)\).  The block inverse formula identifies the
boundary block of \(D_{\mathrm{DW},N}^{-1}\) with the inverse Schur
complement.
\end{proof}

This theorem is an identity for the complete finite ratio.  It does not
declare the denominator reflection positive.

\subsection{Overlap limit and exact chiral algebra}
\label{sec:v3v77-GW}

At \(m_f=0\), choose the conventional dimensionful normalization
\begin{equation}
 D_{\mathrm{ov}}=\frac1a\left(I+V\right),
 \qquad
 V=\gamma_5\epsilon,
 \qquad
 \epsilon=\operatorname{sgn}(H_W).
\label{eq:v3v77-overlap}
\end{equation}

\begin{proposition}[Norm convergence of the wall Schur complement]
\label{prop:v3v77-overlap-convergence}
Under \eqref{eq:v3v77-gap},
\begin{equation}
 \left\|a^{-1}(I+\gamma_5\epsilon_N(H_W))-D_{\mathrm{ov}}\right\|
 \le\frac{2}{a}q^N.
\label{eq:v3v77-overlap-convergence}
\end{equation}
\end{proposition}

\begin{proof}
Multiply the estimate in Lemma~\ref{lem:v3v77-sign-error} by the unitary
\(\gamma_5\) and by \(a^{-1}\).
\end{proof}

\begin{theorem}[Ginsparg--Wilson relation]
\label{thm:v3v77-GW}
The overlap operator satisfies
\begin{equation}
 \gamma_5D_{\mathrm{ov}}+D_{\mathrm{ov}}\gamma_5
 =aD_{\mathrm{ov}}\gamma_5D_{\mathrm{ov}}.
\label{eq:v3v77-GW}
\end{equation}
Moreover
\begin{equation}
 V^*V=I,
 \qquad V^*=\gamma_5V\gamma_5.
\label{eq:v3v77-V-properties}
\end{equation}
\end{theorem}

\begin{proof}
Since \(\epsilon^*=\epsilon\) and \(\epsilon^2=I\),
\(V^*V=\epsilon\gamma_5\gamma_5\epsilon=I\) and
\(V^*=\epsilon\gamma_5=\gamma_5V\gamma_5\).  Also
\begin{equation}
 V\gamma_5V=\gamma_5\epsilon^2=\gamma_5.
\label{eq:v3v77-VgammaV}
\end{equation}
Expanding the right side of \eqref{eq:v3v77-GW} gives
\begin{align*}
 aD_{\mathrm{ov}}\gamma_5D_{\mathrm{ov}}
 &=a^{-1}(I+V)\gamma_5(I+V)\\
 &=a^{-1}(2\gamma_5+V\gamma_5+\gamma_5V),
\end{align*}
which is the expanded left side.
\end{proof}

Define the modified chirality
\begin{equation}
 \widehat\gamma_5=\gamma_5(I-aD_{\mathrm{ov}})=-\epsilon.
\label{eq:v3v77-hat-gamma}
\end{equation}
Then \(\widehat\gamma_5^*=widehat\gamma_5\) and
\(\widehat\gamma_5^2=I\).

\begin{corollary}[Exact finite lattice chiral variation]
\label{cor:v3v77-Luscher}
The massless action \(\bar\psi D_{\mathrm{ov}}\psi\) is invariant under
\begin{equation}
 \delta\psi=i\alpha\widehat\gamma_5\psi,
 \qquad
 \delta\bar\psi=i\alpha\bar\psi\gamma_5.
\label{eq:v3v77-chiral-transformation}
\end{equation}
\end{corollary}

\begin{proof}
The variation is proportional to
\(D_{\mathrm{ov}}\widehat\gamma_5+\gamma_5D_{\mathrm{ov}}\).  Substitute
\eqref{eq:v3v77-hat-gamma}; this expression is exactly
\eqref{eq:v3v77-GW} rearranged.
\end{proof}

\subsection{Index and exact measure Jacobian}
\label{sec:v3v77-index}

Define
\begin{equation}
 \operatorname{ind}D_{\mathrm{ov}}
 =\frac12\operatorname{Tr}(\gamma_5+\widehat\gamma_5)
 =-\frac12\operatorname{Tr}\epsilon,
\label{eq:v3v77-index}
\end{equation}
where \(\operatorname{Tr}\gamma_5=0\) on the full finite lattice spin
space.

\begin{proposition}[Integer index and chiral Jacobian]
\label{prop:v3v77-index-Jacobian}
The number in \eqref{eq:v3v77-index} is an integer and equals the chiral
zero-mode difference of \(D_{\mathrm{ov}}\).  Under
\eqref{eq:v3v77-chiral-transformation}, the Grassmann measure has
infinitesimal Jacobian
\begin{equation}
 \delta\log J=-2i\alpha\operatorname{ind}D_{\mathrm{ov}}
 =i\alpha\operatorname{Tr}\epsilon.
\label{eq:v3v77-Jacobian}
\end{equation}
\end{proposition}

\begin{proof}
The Ginsparg--Wilson relation pairs every nonzero nonexceptional eigenmode
with the opposite modified chirality.  Only zero modes and the cutoff modes
at eigenvalue \(2/a\) contribute to the two traces; the latter cancel the
finite-lattice trace of \(\gamma_5\), leaving the difference of positive
and negative chiral zero modes.  Hence \eqref{eq:v3v77-index} is integral.

For Grassmann variables, a linear field change contributes the inverse
determinant.  To first order,
\begin{equation}
 \delta\log J=-i\alpha\operatorname{Tr}
 (\widehat\gamma_5+\gamma_5)
 =-2i\alpha\operatorname{ind}D_{\mathrm{ov}}.
\end{equation}
Use \(\widehat\gamma_5=-\epsilon\) for the second equality in
\eqref{eq:v3v77-Jacobian}.
\end{proof}

\subsection{Gauge-covariant Weyl subspace and its curvature}
\label{sec:v3v77-measure}

Exact gauge covariance of V76 and functional calculus give
\begin{equation}
 \epsilon(U^g)=g\epsilon(U)g^{-1},
 \qquad
 \widehat P_-(U)=\frac12(I-\widehat\gamma_5(U)).
\label{eq:v3v77-gauge-projector}
\end{equation}
The lattice Weyl field lies in \(\operatorname{ran}\widehat P_-\), while
its independent conjugate uses the fixed complementary continuum-chirality
projector.  Choose a local orthonormal basis \(v_j(U)\) of
\(\operatorname{ran}\widehat P_-(U)\).  Its measure connection and
curvature are
\begin{align}
 \mathcal A&=i\sum_j\langle v_j,dv_j\rangle,
\label{eq:v3v77-measure-connection}\\
 \mathcal F&=i\operatorname{Tr}
 \left(\widehat P_-,d\widehat P_-\wedge d\widehat P_-\right).
\label{eq:v3v77-measure-curvature}
\end{align}

\begin{theorem}[Finite integrability criterion for a chiral measure]
\label{thm:v3v77-integrability}
A globally single-valued gauge-covariant Weyl measure on a connected
admissible gauge-field component exists only if:
\begin{enumerate}
\item the curvature \(\mathcal F\) differs from a gauge-invariant local
connection curvature by an exact local two-form; and
\item the resulting flat connection has unit holonomy around every loop.
\end{enumerate}
On a contractible chart, the first condition is sufficient for a local
choice of measure current.  It does not determine global holonomy.
\end{theorem}

\begin{proof}
A change of basis \(v\mapsto vU\) changes \(\mathcal A\) by the differential
of \(i\log\det U\) and leaves \(\mathcal F=d\mathcal A\) invariant.  A
local counterterm shifts the connection by an exact local one-form.  Thus
the curvature class is the local obstruction.  If it vanishes after the
allowed shift, parallel transport defines a basis-independent measure on a
simply connected chart.  On the full component, path independence is
equivalent to unit holonomy on every loop, as in Theorem
\ref{thm:v3v66-flat-line}.
\end{proof}

In the smooth-field continuum expansion, the nontrivial local part of
\eqref{eq:v3v77-measure-curvature} is the descent of the representation
polynomial \eqref{eq:v3v66-I6}.  The Standard Model coefficients cancel by
Theorem~\ref{thm:v3v66-bulk-descent-zero}, and the traditional
\(SU(2)\) mod-two loop cancels by Proposition
\ref{prop:v3v66-Witten-loop}.  These facts remove the known continuum
classes; they do not prove that the finite-lattice curvature is exactly a
local coboundary on every admissible component.

\subsection{Reflection-positivity boundary}
\label{sec:v3v77-RP}

Equation \eqref{eq:v3v77-determinant-ratio} divides the reflection-positive
Grassmann determinant of V76 by a heavy determinant.  A common local
representation introduces commuting Pauli--Villars fields, but their
Gaussian kernel and reflection cone differ from the Grassmann one.

\begin{firewall}[Ginsparg--Wilson symmetry is not reflection positivity]
\label{fw:v3v77-GW-not-RP}
The norm limit, Ginsparg--Wilson relation, integer index, and exact chiral
Jacobian do not prove Osterwalder--Schrader positivity of the determinant
ratio or of the non-ultralocal overlap operator.  V76's positivity theorem
continues to apply only before the heavy normalization is divided out.
\end{firewall}

\begin{firewall}[Gap and locality hypothesis]
\label{fw:v3v77-gap}
Every sign-function and exponential fifth-length estimate in V77 assumes
the Hermitian Wilson gap \eqref{eq:v3v77-gap}.  V76's plaquette
admissibility suggests a stable sector but does not yet prove a uniform gap
on all topological components used by the interacting measure.
\end{firewall}

\begin{firewall}[Euclidean repair versus the Lorentzian projector]
\label{fw:v3v77-Lorentzian-match}
The overlap wall supplies an elliptic Euclidean chiral operator without the
noncomplementing V56 material projector.  V77 does not prove that its
Osterwalder--Schrader reconstruction has the V56 Lorentzian
maximal-isotropic boundary domain or its material reaction stress.
\end{firewall}

\subsection{V77 conclusion and V78 target}
\label{sec:v3v77-conclusion}

The finite fifth-direction Schur complement converges exponentially to the
overlap operator under a Wilson-kernel gap.  The limit obeys the exact
Ginsparg--Wilson relation, has an integer index and exact lattice chiral
Jacobian, and defines a gauge-covariant Weyl subspace with explicit measure
curvature.  Standard Model anomaly arithmetic cancels the known continuum
curvature class, while global finite-lattice integrability and the heavy
determinant reflection form remain open.

V78 should prove a Wilson-kernel gap on a quantitatively admissible gauge
sector and construct the local measure current solving the finite-lattice
integrability equation for the complete anomaly-free Standard Model
representation.  It must test curvature periods and the global gauge-group
quotient explicitly.  In parallel, the Pauli--Villars determinant ratio
must be written as one time-local transfer system before any
reflection-positivity conclusion is attempted.
\clearpage
\section{V78: a quantitative Wilson-gap patch and a local overlap measure connection}
\label{sec:v3v78}

V77 assumes a Hermitian Wilson-kernel gap.  Plaquette admissibility alone
does not make that assumption immediate on every topological component.
This note proves the gap on a concrete small-link component, obtains
exponential locality of the sign function, and constructs a local measure
connection on that contractible patch.  The result is intentionally not
extended across topological-sector boundaries.

\subsection{Small-link component and gap stability}
\label{sec:v3v78-gap}

Work at fixed lattice spacing and write the Hermitian kernel as
\begin{equation}
 H_W(U)=\gamma_5(D_W(U)-M).
\label{eq:v3v78-HW}
\end{equation}
Let \(H_0=H_W(I)\), and suppose the mass lies in a closed subinterval of
the one-species window.  Its free gap is
\begin{equation}
 g_0:=\min_{p\in T^4}\min|\sigma(H_0(p))|>0.
\label{eq:v3v78-free-gap}
\end{equation}
Choose principal link logarithms
\begin{equation}
 U_\mu(x)=e^{A_\mu(x)},
 \qquad A_\mu(x)^*=-A_\mu(x),
 \qquad \|A\|_\infty\le\eta<\pi.
\label{eq:v3v78-small-links}
\end{equation}

\begin{lemma}[Wilson operator perturbation]
\label{lem:v3v78-Wilson-perturbation}
There is a dimensionless constant \(C_W\), depending only on the Wilson
parameter and dimension, such that
\begin{equation}
 \|H_W(U)-H_0\|\le\frac{C_W}{a}(e^\eta-1).
\label{eq:v3v78-operator-perturbation}
\end{equation}
The bound is independent of lattice volume.
\end{lemma}

\begin{proof}
Every covariant shift differs from the free shift by multiplication by
\(U_\mu-I\), followed by a unitary translation.  Hence its norm difference
is at most \(e^\eta-1\).  The Wilson operator is a sum of a fixed finite
number of forward and backward shifts with coefficients bounded by
\(C/a\).  Sum these operator norms and use that multiplication by
\(\gamma_5\) is unitary.
\end{proof}

\begin{theorem}[Quantitative small-link Wilson gap]
\label{thm:v3v78-small-link-gap}
If
\begin{equation}
 \frac{C_W}{a}(e^\eta-1)<g_0,
\label{eq:v3v78-gap-condition}
\end{equation}
then
\begin{equation}
 \sigma(H_W(U))\cap(-g,g)=\varnothing,
 \qquad
 g=g_0-\frac{C_W}{a}(e^\eta-1)>0.
\label{eq:v3v78-interacting-gap}
\end{equation}
The straight logarithmic path \(U(t)=e^{tA}\), \(0\le t\le1\), remains in
the same gapped component.
\end{theorem}

\begin{proof}
For self-adjoint matrices, the Hausdorff distance between spectra is at
most the operator-norm perturbation.  Lemma
\ref{lem:v3v78-Wilson-perturbation} therefore leaves the interval in
\eqref{eq:v3v78-interacting-gap} empty.  Replacing \(\eta\) by \(t\eta\)
only decreases the perturbation bound, so the entire radial path is gapped.
\end{proof}

This component contains the trivial connection and is contractible in the
chosen logarithmic coordinates.  It need not contain gauge fields of
nonzero lattice topological charge.

\subsection{Locality of the sign and overlap operator}
\label{sec:v3v78-locality}

For an invertible Hermitian matrix,
\begin{equation}
 \operatorname{sgn}H=\frac2\pi\int_0^\infty
 H(H^2+t^2)^{-1}\,dt.
\label{eq:v3v78-sign-integral}
\end{equation}

\begin{theorem}[Exponential overlap locality on the gap patch]
\label{thm:v3v78-overlap-locality}
Under Theorem~\ref{thm:v3v78-small-link-gap}, there are \(C,c>0\),
independent of volume, such that
\begin{equation}
 \|\operatorname{sgn}(H_W)(x,y)\|
 \le C e^{-cg\,d(x,y)}.
\label{eq:v3v78-sign-locality}
\end{equation}
The same type of bound holds for a fixed finite number of derivatives with
respect to link logarithms.  Hence \(D_{\mathrm{ov}}\),
\(\widehat P_-\), and the curvature
\eqref{eq:v3v77-measure-curvature} are exponentially local.
\end{theorem}

\begin{proof}
The kernel \(H_W\) has finite range.  Conjugation by
\(e^{\alpha d(\cdot,y)}\) changes it by an operator of norm at most
\(C\alpha\) for \(\alpha a\le1\).  The spectral distance of
\(H_W^2+t^2\) from zero is at least \(g^2+t^2\); the resolvent Neumann
argument therefore gives a Combes--Thomas kernel bound with exponent
\(c\sqrt{g^2+t^2}\,d(x,y)\).  Insert it in
\eqref{eq:v3v78-sign-integral}; the integral is finite and retains
\(e^{-cgd(x,y)}\).

Differentiating \eqref{eq:v3v78-sign-integral} inserts finitely many local
link derivatives separated by the same resolvents.  Repeated convolution
of exponential kernels remains exponential, with a possibly smaller
constant \(c\).  The overlap and projector statements follow from their
algebraic definitions in V77, and the curvature contains two such
derivatives and one projector.
\end{proof}

\subsection{Poincar\'e construction of the measure connection}
\label{sec:v3v78-Poincare}

Let \(\mathscr U_\eta\) be the star-shaped logarithmic patch in
\eqref{eq:v3v78-small-links}, and write the total Standard Model Weyl
curvature as the sum of \eqref{eq:v3v77-measure-curvature} over all
representations and generations:
\begin{equation}
 \mathcal F_{\rm SM}=\frac12
 \mathcal F_{IJ}(A)\,dA^I\wedge dA^J.
\label{eq:v3v78-total-curvature}
\end{equation}
Finite-dimensional Chern--Weil algebra gives \(d\mathcal F_{\rm SM}=0\).
Define
\begin{equation}
 \mathcal A_I^{\rm rad}(A)
 =\int_0^1 t A^J\mathcal F_{JI}(tA)\,dt.
\label{eq:v3v78-radial-connection}
\end{equation}

\begin{theorem}[Local curvature primitive on the gap patch]
\label{thm:v3v78-curvature-primitive}
The one-form \(\mathcal A^{\rm rad}\) is smooth and exponentially local and
satisfies
\begin{equation}
 d\mathcal A^{\rm rad}=\mathcal F_{\rm SM}
 \quad\hbox{on }\mathscr U_\eta.
\label{eq:v3v78-curvature-primitive}
\end{equation}
It therefore defines a single-valued Weyl measure connection on this
contractible patch, once its phase is fixed at \(A=0\).
\end{theorem}

\begin{proof}
Smoothness follows from the gap along the radial path.  Exponential
locality follows from Theorem~\ref{thm:v3v78-overlap-locality}, uniformly
for \(0\le t\le1\), and integration over the finite parameter interval.
For the differential identity, let \(R=A^I\partial_I\) be the radial vector
field and \(h_t(A)=tA\).  Formula \eqref{eq:v3v78-radial-connection} is the
standard homotopy operator
\(K\mathcal F=\int_0^1t^{-1}h_t^*(\iota_R\mathcal F)dt\).  Cartan's formula
gives
\begin{equation}
 dK\mathcal F+K d\mathcal F
 =\mathcal F-h_0^*\mathcal F.
\end{equation}
The pullback of a two-form to the point \(A=0\) is zero and
\(d\mathcal F=0\), proving \eqref{eq:v3v78-curvature-primitive}.  A
connection with this curvature determines parallel measure transport from
the fixed base phase.
\end{proof}

The construction is a primitive in field space.  It is exponentially
local in spacetime because the curvature is, but it is tied to the chosen
small-link logarithmic chart.

\subsection{Gauge directions and the Standard Model sum}
\label{sec:v3v78-gauge}

On a based infinitesimal gauge orbit, the contraction of the measure
curvature is the consistency derivative of the lattice gauge anomaly.  In
the smooth-field expansion its cohomologically nontrivial limit is the V66
descent polynomial.  Proposition
\ref{prop:v3v66-coefficient-cancellation} cancels that polynomial in the
complete Standard Model sum.

\begin{proposition}[Patchwise gauge measure criterion]
\label{prop:v3v78-patch-gauge}
Suppose the remaining finite-lattice gauge contraction of
\(\mathcal A^{\rm rad}\) differs from the desired covariant current by the
gauge variation of an exponentially local functional \(X_\eta(A)\).
Then
\begin{equation}
 \mathcal A^{\rm inv}=\mathcal A^{\rm rad}-dX_\eta
\label{eq:v3v78-invariant-connection}
\end{equation}
defines a gauge-invariant Weyl measure on the based small-link patch.
\end{proposition}

\begin{proof}
Subtracting an exact one-form leaves the curvature unchanged.  By the
hypothesis, the contraction of \(\mathcal A^{\rm inv}\) with every based
gauge vector equals the connection required to cancel the fermion-action
variation.  Integrating along the contractible based gauge orbit gives an
invariant measure phase, with no residual loop because the orbit patch is
simply connected.
\end{proof}

The hypothesis is the finite-lattice local cohomology equation.  V66 proves
its continuum nontrivial coefficient vanishes; it does not construct
\(X_\eta\) at finite spacing.  Thus Theorem
\ref{thm:v3v78-curvature-primitive} solves geometric integrability on the
patch, while exact gauge invariance still requires this local Ward
primitive.

\subsection{Claim boundary and next target}
\label{sec:v3v78-boundary}

\begin{firewall}[Small-link patch only]
\label{fw:v3v78-small-patch}
The norm perturbation proof excludes large pure-gauge representatives and
nontrivial admissible topological sectors unless they can be brought into
the chosen link-logarithm ball.  It is not a global admissibility-gap
theorem.
\end{firewall}

\begin{firewall}[Curvature primitive is not the finite anomaly solution]
\label{fw:v3v78-curvature-not-Ward}
Every closed two-form on a star-shaped finite-dimensional patch has the
primitive \eqref{eq:v3v78-radial-connection}.  Gauge invariance additionally
requires the local functional \(X_\eta\) in Proposition
\ref{prop:v3v78-patch-gauge}.  Standard Model continuum arithmetic alone
does not prove its finite-lattice existence.
\end{firewall}

\begin{firewall}[No new reflection-positivity claim]
\label{fw:v3v78-no-RP}
Exponential locality of the overlap operator is weaker than ultralocality
and does not supply the positive cross-plane cone of V76.  Reflection
positivity of the Pauli--Villars-normalized overlap determinant remains
open.
\end{firewall}

\subsection{V78 conclusion and V79 target}
\label{sec:v3v78-conclusion}

The Wilson gap, overlap locality, and a smooth exponentially local measure
connection are now proved on a quantitative contractible small-link
component.  This removes the gap assumption there and makes the geometric
part of finite Weyl-measure integrability explicit.  The unresolved local
equation is the exact finite-spacing gauge Ward primitive, while global
topological components and reflection positivity remain separate.

V79 should compute the gauge contraction of
\eqref{eq:v3v78-radial-connection} and solve its local cohomology equation
order by order in the small link logarithm for the complete Standard Model
representation.  It must give cutoff-uniform locality bounds on the
counterterm coefficients and identify the first order at which global
group-quotient data, rather than perturbative anomaly coefficients, can
enter.
\clearpage
\section{V79: the overlap moment map and perturbative Standard Model measure defect}
\label{sec:v3v79}

V78 constructs a local curvature primitive but leaves its contraction with
gauge directions unresolved.  The finite overlap projector lives in a
Grassmannian whose gauge action has an exact moment map.  This identity
computes the contraction without using the nonlocal inverse-difference
homotopy excluded in V65.  It also separates perturbative Lie-algebra data
from global gauge-group holonomy.

\subsection{Exact Grassmannian moment map}
\label{sec:v3v79-moment-map}

Let \(P=\widehat P_-(U)\) and
\begin{equation}
 \mathcal F=i\operatorname{Tr}(P,dP\wedge dP)
\label{eq:v3v79-curvature}
\end{equation}
be the curvature of one finite Weyl subspace.  For an anti-Hermitian
infinitesimal gauge generator \(c\), the fundamental vector field on the
projector manifold is
\begin{equation}
 X_cP=[c,P].
\label{eq:v3v79-gauge-vector}
\end{equation}
Define
\begin{equation}
 \mu_c(P)=-i\operatorname{Tr}(cP).
\label{eq:v3v79-moment-map}
\end{equation}

\begin{theorem}[Finite projector moment-map identity]
\label{thm:v3v79-moment-map}
For a field-independent generator \(c\),
\begin{equation}
 \iota_{X_c}\mathcal F=d\mu_c.
\label{eq:v3v79-moment-map-identity}
\end{equation}
The identity is exact at finite lattice spacing and finite volume.
\end{theorem}

\begin{proof}
Differentiate \(P^2=P\) to obtain
\(P(dP)P=(I-P)(dP)(I-P)=0\).  Relative to
\(\operatorname{ran}P\oplus\ker P\), write
\begin{equation}
 dP=\begin{pmatrix}0&b\\b^*&0\end{pmatrix},
 \qquad
 c=\begin{pmatrix}c_{11}&c_{12}\\c_{21}&c_{22}\end{pmatrix}.
\label{eq:v3v79-blocks}
\end{equation}
Then
\([c,P]=\left(\begin{smallmatrix}0&-c_{12}\\c_{21}&0\end{smallmatrix}\right)\).
For an arbitrary tangent \(\dot P\), direct block multiplication gives
\begin{align}
 \mathcal F(X_c,\dot P)
 &=i\operatorname{Tr}P[[c,P],\dot P]\\
 &=-i\operatorname{Tr}(c\dot P)
 =d\mu_c(\dot P),
\end{align}
which proves the formula.
\end{proof}

For a field-dependent BRST ghost, the same calculation is used in the
Cartan equivariant differential; the additional \(dc\) term is precisely
the standard consistency term.

\subsection{Gauge defect of a chosen measure connection}
\label{sec:v3v79-defect}

Let \(\mathcal A\) be any local connection with curvature \(\mathcal F\).
Its infinitesimal equivariant defect is
\begin{equation}
 \mathscr G(c;A)=\iota_{X_c}\mathcal A+\mu_c.
\label{eq:v3v79-gauge-defect}
\end{equation}
The sign is fixed so that \(\mathscr G=0\) is horizontal equivariance of
the determinant section.

\begin{proposition}[Consistency of the finite measure defect]
\label{prop:v3v79-consistency}
For field-independent generators \(c_1,c_2\),
\begin{equation}
 X_{c_1}\mathscr G(c_2)-X_{c_2}\mathscr G(c_1)
 -\mathscr G([c_1,c_2])=0
\label{eq:v3v79-consistency}
\end{equation}
whenever \(\mathcal A\) is invariant under the simultaneous finite gauge
action.  For a general local trivialization, the left side is the gauge
variation of a local scalar and can be shifted by an exact change of
connection.
\end{proposition}

\begin{proof}
Cartan's formula and Theorem~\ref{thm:v3v79-moment-map} give
\(d\mathscr G(c)=\mathcal L_{X_c}\mathcal A\).  For an invariant
connection this vanishes, so the Hamiltonian commutator identity for the
moment maps yields \eqref{eq:v3v79-consistency}.  Replacing
\(\mathcal A\) by \(\mathcal A+dX\) replaces \(\mathscr G\) by
\(\mathscr G+X_cX\), which is the stated local scalar shift.
\end{proof}

The radial connection of V78 is not invariant automatically, but the last
sentence gives the exact equation which its local correction must solve.

\subsection{Convergent small-link expansion}
\label{sec:v3v79-expansion}

On the gap ball of V78, the sign function and projector are analytic in
the link logarithms.  Write
\begin{equation}
 P(A)=\sum_{n\ge0}P^{(n)}(A),
 \qquad
 \mathscr G(c;A)=\sum_{n\ge0}\mathscr G^{(n)}(c;A),
\label{eq:v3v79-series}
\end{equation}
where the superscript is homogeneous degree in \(A\).

\begin{lemma}[Uniform analytic locality bounds]
\label{lem:v3v79-analytic-locality}
On every smaller ball \(\|A\|_\infty\le\eta'<\eta\), there are
\(C,R,c>0\), independent of volume, such that
\begin{equation}
 \|P^{(n)}(x,y)\|
 \le C R^{-n}e^{-cgd(x,y)}\|A\|_infty^n.
\label{eq:v3v79-coefficient-bound}
\end{equation}
The same estimate, with a fixed polynomial prefactor, holds for the
coefficients of \(\mathcal F\), \(\mathcal A^{\rm rad}\), and
\(\mathscr G\).
\end{lemma}

\begin{proof}
The gap persists in a complex operator-norm neighbourhood of every smaller
real ball.  The contour formula
\begin{equation}
 P(A)=\frac1{2\pi i}\int_\Gamma(z-H_W(A))^{-1},dz
\label{eq:v3v79-Riesz}
\end{equation}
uses a contour a fixed distance from both spectral components.  Cauchy's
estimate gives the factor \(R^{-n}\), and the finite-range
Combes--Thomas estimate used in V78 gives the exponential kernel bound.
Products and the finite radial integral preserve exponential locality and
only change the constants.
\end{proof}

Thus the local Ward equation can be addressed coefficient by coefficient
without a merely formal series.

\subsection{Zeroth order and the first nontrivial local class}
\label{sec:v3v79-leading}

For the complete Standard Model representation, sum the moment maps and
defects over all left-handed variables in \eqref{eq:v3v66-SM-list}.  At
\(A=0\), translation and spin symmetry make the diagonal of \(P^{(0)}\)
proportional to the identity in each internal representation.

\begin{proposition}[Zeroth-order gauge trace vanishes]
\label{prop:v3v79-zero-order}
For nonabelian generators, \(\mathscr G^{(0)}\) vanishes by tracelessness.
For hypercharge, its representation factor is
\begin{equation}
 6\left(\frac16\right)+3\left(-\frac23\right)
 +3\left(\frac13\right)+2\left(-\frac12\right)+1=0.
\label{eq:v3v79-linear-hypercharge}
\end{equation}
Hence the complete one-generation measure has no field-independent gauge
character on the based small-link patch.
\end{proposition}

\begin{proof}
The trace factorizes into the common spacetime-spin trace and the internal
trace of the generator.  The latter vanishes for \(SU(3)\) and \(SU(2)\).
For \(U(1)_Y\), insert the dimensions and charges from
\eqref{eq:v3v66-SM-list}; the displayed sum is zero.
\end{proof}

The first cohomologically nontrivial four-dimensional term is the
consistent descent density
\begin{equation}
 \mathscr G_R^{\rm an}(c,A)=
 C_{\rm an}\int
 \operatorname{tr}_R\left[
 c\,d\left(A\,dA+\frac12A^3\right)
 \right],
\label{eq:v3v79-consistent-anomaly}
\end{equation}
with the mixed gravitational term proportional to
\(p_1(T)\operatorname{tr}_R c\).

\begin{theorem}[Cancellation of the first nontrivial local measure class]
\label{thm:v3v79-leading-cancellation}
In the complete one-generation Standard Model sum, every representation
coefficient multiplying \eqref{eq:v3v79-consistent-anomaly} and its mixed
gravitational partner is zero.  Therefore the quadratic and cubic
small-field terms contain no nontrivial local BRST cohomology class; any
remaining regulator-dependent terms at those orders must be local
Slavnov coboundaries or irrelevant lattice artifacts.
\end{theorem}

\begin{proof}
The independent coefficients are
\(SU(3)^3\), \(SU(3)^2Y\), \(SU(2)^2Y\), \(Y^3\), and
\(Y\text{-}\mathrm{grav}^2\).  They vanish separately by Proposition
\ref{prop:v3v66-coefficient-cancellation}.  The perturbative
\(SU(2)^3\) trace is zero because the doublet is pseudoreal.  The local
cohomology classification of a consistent four-dimensional gauge defect
has these descent polynomials as its nontrivial ghost-number-one
representatives.  Removing their coefficients leaves only coboundaries
and terms carrying positive powers of the lattice spacing.
\end{proof}

The theorem classifies the remaining terms; it does not yet display their
finite-spacing primitive or a cutoff-uniform bound for that primitive.

\subsection{Recursive local restoration statement}
\label{sec:v3v79-recursion}

Let \(s=s_0+s_1+\cdots\) be the BRST differential filtered by degree in
\(A\), with \(s_0A=dc\).  Suppose counterterms have cancelled
\(\mathscr G^{(j)}\) for \(j<n\).  Consistency gives
\begin{equation}
 s_0\mathscr G^{(n)}
 =-\sum_{k=1}^n s_k\mathscr G^{(n-k)}=0.
\label{eq:v3v79-recursive-closed}
\end{equation}

\begin{proposition}[Order-by-order restoration criterion]
\label{prop:v3v79-recursion}
If an exponentially local homogeneous functional \(X^{(n)}\) satisfies
\begin{equation}
 \mathscr G^{(n)}=s_0X^{(n)},
\label{eq:v3v79-recursive-primitive}
\end{equation}
then changing the measure connection by \(-dX^{(n)}\) cancels the degree
\(n\) defect without changing lower degrees or its curvature.  If the
bounds on \(X^{(n)}\) have a common convergence radius, the corrections
sum to an exponentially local finite-spacing measure current.
\end{proposition}

\begin{proof}
The exact connection shift changes the gauge contraction by
\(-sX^{(n)}\).  Its degree-\(n\) part is
\(-s_0X^{(n)}\), while every other term has higher degree.  Exact shifts do
not change \(d\mathcal A=\mathcal F\).  A common geometric bound on the
coefficients gives absolute convergence by the Weierstrass test and
preserves exponential locality under summation.
\end{proof}

V79 verifies the absence of the nontrivial class through the leading
anomalous orders.  Constructing \(X^{(n)}\) with a common bound for all
\(n\) is the remaining finite-spacing problem.

\subsection{Global quotient data are nonperturbative}
\label{sec:v3v79-global}

\begin{theorem}[No finite Taylor order detects the global gauge-group
quotient]
\label{thm:v3v79-global-not-perturbative}
Two compact gauge groups with the same Lie algebra and the same local
representation differentials have identical small-link coefficients
\(P^{(n)}\), \(\mathcal F^{(n)}\), and \(\mathscr G^{(n)}\) at every finite
order.  Differences caused by a quotient of
\(SU(3)\times SU(2)\times U(1)\) can enter only through the allowed global
components and loop holonomies.
\end{theorem}

\begin{proof}
Every Taylor coefficient at the identity is obtained by differentiating
link exponentials and representation matrices finitely many times.  These
derivatives depend only on the Lie algebra representation.  A discrete
central quotient changes which finite transformations and bundles exist,
but not any derivative at the identity.  Hence it is invisible on the
contractible logarithmic patch and can appear only through global topology.
\end{proof}

\begin{firewall}[Leading cancellation is not an all-order lattice measure]
\label{fw:v3v79-not-all-order}
Theorem~\ref{thm:v3v79-leading-cancellation} removes the nontrivial
continuum local class at its first possible orders.  It does not construct
the complete sequence \(X^{(n)}\), prove its convergence, or establish
exact gauge invariance at finite lattice spacing.
\end{firewall}

\begin{firewall}[Patch analysis cannot decide global holonomy]
\label{fw:v3v79-no-global}
No perturbative expansion in link logarithms can test the global Standard
Model gauge-group quotient, the material-frame loop of V67, or a
mapping-torus determinant phase.  Those require loops outside the
contractible patch.
\end{firewall}

\subsection{V79 conclusion and V80 audit target}
\label{sec:v3v79-conclusion}

The overlap curvature has an exact finite moment map.  Its small-link
series is convergent and exponentially local, its zeroth-order gauge
character vanishes in the Standard Model, and all leading consistent
anomaly coefficients cancel.  The recursive equation now identifies the
precise missing object: uniformly bounded local primitives for the
finite-spacing coboundary terms at every order.

V80 is the next compulsory NS/YM audit.  It should then build the lattice
algebraic Poincar\'e descent needed to solve
\eqref{eq:v3v79-recursive-primitive} with support and norm bounds uniform
in order, or exhibit the first coefficient where those bounds fail.  The
global quotient and Pauli--Villars reflection problems must remain outside
that local recursion.

\section{V80: audited lattice descent, its first-moment loss, and the cellwise criterion}
\label{sec:v3v80}

This note performs the compulsory fifth-version cross-program audit and then
addresses the precise primitive problem left by V79.  The issue is not the
vanishing of the Standard Model anomaly polynomial, which was already proved
in V66 and used in V79.  The issue is whether the remaining finite-spacing
BRST coboundaries admit primitives with locality constants that can be summed
over all orders.

\subsection{Compulsory NS/YM audit and nonduplication ledger}
\label{sec:v3v80-audit}

The two companion programmes were read passively before this note was written.
The files actually inspected were
\begin{align}
 &\texttt{Renewal\_NS\_Stability\_Research\_Notes\_V31.tex},
 &&887537\text{ bytes},
 &&\operatorname{SHA256}=
 \texttt{F1121B62238AD5491BB7CF03635EA9934942EDF573ED8FAAA9EE79E1D38A6696},
 \label{eq:v3v80-ns-fingerprint}\\
 &\texttt{Renewal\_Yang\_Mills\_Mass\_Gap\_Research\_Notes\_V2.tex},
 &&51823\text{ bytes},
 &&\operatorname{SHA256}=
 \texttt{C8513BAD49F03FDD20B970323950E5F5E78F30A3B946A54996093677AFF60157}.
 \label{eq:v3v80-ym-fingerprint}
\end{align}
The YM programme had frozen its preceding hundred-note volume and restarted
its active sequence.  No companion file was modified.

The useful YM lesson is algebraic: changes of local roles must first be
assembled as a finite signed sum of complete cover covariances; estimating
the pieces before this assembly destroys the cancellation.  The useful NS
lesson is analytic: a marked or anchored primitive can restore a missing
moment, but the cost must be measured in the topology in which the final
summation is required.  In the present problem these become, respectively,
the requirement to assemble a complete local descent identity before taking
norms and the requirement to account for the first spatial moment introduced
by a lattice homotopy.

V78--V79 already establish a gapped analytic overlap projector, exponential
locality of its Taylor coefficients, the exact Grassmannian moment map, the
zeroth-order character cancellation, and cancellation of the leading
nontrivial anomaly classes.  None of those results is repeated below.  This
note isolates the extra hypothesis needed to turn that cohomological
information into uniformly summable primitives.

\subsection{The exact tree primitive}
\label{sec:v3v80-tree}

Let \(G=(V,E)\) be a finite connected graph with a chosen orientation.  For an
edge current \(J\in C^1(G;\mathbb R)\), use the convention
\begin{equation}
 (\operatorname{div}J)(x)
 =\sum_{e:\,e^+=x}J(e)-\sum_{e:\,e^-=x}J(e).
\label{eq:v3v80-divergence}
\end{equation}
This is the adjoint, with the present sign convention, of
\((d_hc)(e)=c(e^+)-c(e^-)\):
\begin{equation}
 \sum_{e\in E}J(e)(d_hc)(e)
 =\sum_{x\in V}c(x)(\operatorname{div}J)(x).
\label{eq:v3v80-summation-by-parts}
\end{equation}

\begin{theorem}[Constructive discrete Poincar\'e primitive]
\label{thm:v3v80-tree-primitive}
Choose a root \(o\) and a spanning tree \(T\), orienting every tree edge away
from \(o\).  If \(f\in C^0(G;\mathbb R)\) satisfies
\(\sum_{x\in V}f(x)=0\), then there is a current \(K_Tf\), supported on the
minimal rooted subtree joining \(o\) to \(\operatorname{supp}f\), such that
\begin{equation}
 \operatorname{div}K_Tf=f,
 \qquad
 \|K_Tf\|_{\ell^1(E)}
 \le \sum_{x\in V}d_T(o,x)|f(x)|.
\label{eq:v3v80-tree-bound}
\end{equation}
More explicitly, if \(e=(p(v),v)\) and \(T_v\) is the descendant subtree
rooted at \(v\), then
\begin{equation}
 (K_Tf)(e)=\sum_{x\in T_v}f(x),
\label{eq:v3v80-tree-formula}
\end{equation}
and the current is zero off \(T\).
\end{theorem}

\begin{proof}
At a nonroot vertex \(v\), the incoming contribution in
\eqref{eq:v3v80-divergence} is the sum of \(f\) over \(T_v\); subtracting
the corresponding sums over the child subtrees leaves exactly \(f(v)\).
At the root, the absence of an incoming edge gives
\(-\sum_{v:\,p(v)=o}\sum_{x\in T_v}f(x)=f(o)\), by the zero-sum hypothesis.
This proves the divergence identity.

An edge \((p(v),v)\) can carry current only if \(T_v\) meets the support of
\(f\), which gives the support statement after deleting branches whose total
charge is zero.  Finally,
\begin{align}
 \sum_{e\in T}|(K_Tf)(e)|
 &\le \sum_{v\ne o}\sum_{x\in T_v}|f(x)|
 =\sum_{x\in V}d_T(o,x)|f(x)|.
\end{align}
Each \(x\) occurs once for every edge on its root path, proving
\eqref{eq:v3v80-tree-bound}.
\end{proof}

The formula is the degree-zero part of the usual cellular cone homotopy
\(d_hK+Kd_h=I-\pi\), where \(\pi\) projects onto constants.  It is enough
for the first filtered BRST equation, because \(s_0A=d_hc\).

\begin{corollary}[Exact linearized BRST primitive]
\label{cor:v3v80-brst-primitive}
Let
\(\mathscr G_f(c)=\sum_x\langle c(x),f(x)\rangle\), with
\(\sum_x f(x)=0\) in each Lie-algebra component.  Put
\begin{equation}
 X_f(A)=\sum_e\langle (K_Tf)(e),A(e)\rangle.
\label{eq:v3v80-brst-counterterm}
\end{equation}
Then \(s_0X_f=\mathscr G_f\).  Its \(\ell^1\) coefficient norm obeys the
first-moment estimate \eqref{eq:v3v80-tree-bound}.
\end{corollary}

\begin{proof}
Insert \(s_0A=d_hc\) in \eqref{eq:v3v80-brst-counterterm}, apply
\eqref{eq:v3v80-summation-by-parts}, and use
\(\operatorname{div}K_Tf=f\).
\end{proof}

\subsection{A sharp obstruction to support-uniform inversion}
\label{sec:v3v80-obstruction}

The moment in \eqref{eq:v3v80-tree-bound} is necessary, rather than an
artifact of the selected tree.

\begin{theorem}[Separated dipole lower bound]
\label{thm:v3v80-dipole}
On the path with vertices \(0,1,\ldots,L\), oriented from left to right, let
\begin{equation}
 f_L=\delta_L-\delta_0.
\label{eq:v3v80-dipole-source}
\end{equation}
Every current satisfying \(\operatorname{div}J=f_L\) has
\begin{equation}
 J(k,k+1)=1\quad(0\le k<L),
 \qquad
 \|J\|_{\ell^1}=L.
\label{eq:v3v80-dipole-current}
\end{equation}
Consequently there is no right inverse of divergence which, on all finite
boxes, has both a volume-independent \(\ell^1\) operator bound and a fixed
finite propagation radius.
\end{theorem}

\begin{proof}
The divergence equation at vertex zero gives \(-J(0,1)=-1\), hence
\(J(0,1)=1\).  At every interior vertex it gives
\(J(k-1,k)-J(k,k+1)=0\).  Induction proves that all \(L\) edge values equal
one; the endpoint equation is then automatic.  Thus the norm and support
statements in \eqref{eq:v3v80-dipole-current} are forced.  Since
\(\|f_L\|_{\ell^1}=2\) while \(L\) is arbitrary, a uniform operator bound
is impossible, and a current of fixed propagation cannot join the two
charges.
\end{proof}

This obstruction occurs already in the abelian, linearized equation.  It
therefore precedes all complications due to nonabelian commutators or high
Taylor order.  Cohomological triviality alone cannot provide the estimates
assumed in Proposition~\ref{prop:v3v79-recursion}.

The lower bound does not destroy exponential locality when the defect is
already decomposed into neutral local clusters.  Indeed, for
\(0<\mu'<\mu\),
\begin{equation}
 L e^{-\mu L}
 \le \frac{1}{e(\mu-\mu')}e^{-\mu'L}.
\label{eq:v3v80-exponent-trade}
\end{equation}
Thus one spatial moment can be paid by a controlled loss of exponential
rate.  What fails is the attempt to infer such a neutral cluster
decomposition from the vanishing of the total cohomology class after global
summation.

\subsection{A sufficient cellwise descent theorem}
\label{sec:v3v80-cellwise}

Let a cluster \(Z\) be a finite connected set of lattice cells and write
\(\operatorname{diam}Z\) for its graph diameter.  A homogeneous local
functional is presented in \emph{assembled cluster form} if
\begin{equation}
 \mathscr G^{(n)}=\sum_Z\mathscr G_Z^{(n)},
 \qquad
 \mathscr G_Z^{(n)}=s_0X_Z^{(n)}+d_hY_Z^{(n)},
\label{eq:v3v80-cluster-descent}
\end{equation}
where the second term sums to zero on the periodic lattice.  The identity
must hold for the complete Standard Model representation at each \(Z\),
before any absolute value is taken.

\begin{theorem}[Convergent cellwise restoration criterion]
\label{thm:v3v80-cellwise-criterion}
Suppose there are constants \(C,B,\mu>0\), independent of \(n\) and the
lattice volume, such that \eqref{eq:v3v80-cluster-descent} holds and
\begin{equation}
 \sup_x\sum_{Z\ni x}e^{\mu\operatorname{diam}Z}
 \|X_Z^{(n)}\|\le CB^n.
\label{eq:v3v80-cluster-bound}
\end{equation}
Then, for every \(\mu'<\mu\),
\(X^{(n)}=\sum_ZX_Z^{(n)}\) is exponentially local with a
volume-independent bound \(CB^n\) at rate \(\mu'\), and
\begin{equation}
 X(A)=\sum_{n\ge0}X^{(n)}(A)
\label{eq:v3v80-counterterm-series}
\end{equation}
converges absolutely whenever \(B\|A\|_\infty<1\).  On that ball the exact
connection shift \(-dX\) cancels the assembled defect, provided the
filtered recursion has already removed lower orders as in V79.

If the local construction of \(X_Z^{(n)}\) uses the tree homotopy of
Theorem~\ref{thm:v3v80-tree-primitive}, the same conclusion holds when the
right-hand side before homotopy is bounded at rate \(\mu\): the factor
\(\operatorname{diam}Z\) is absorbed by
\eqref{eq:v3v80-exponent-trade}.
\end{theorem}

\begin{proof}
For any marked cell \(x\), absolute summation of
\eqref{eq:v3v80-cluster-bound} bounds the contribution of all clusters
meeting \(x\), uniformly in the ambient volume.  Dropping the excess weight
\(e^{(\mu-\mu')\operatorname{diam}Z}\) gives exponential locality at rate
\(\mu'\).  Homogeneity supplies the additional factor
\(\|A\|_\infty^n\); hence the Weierstrass test proves convergence of
\eqref{eq:v3v80-counterterm-series} for
\(B\|A\|_\infty<1\).  Absolute convergence permits \(s_0\), the cluster
sum, and the Taylor sum to be interchanged.  The periodic sum of
\(d_hY_Z^{(n)}\) vanishes, leaving
\(s_0X^{(n)}=\mathscr G^{(n)}\).  Proposition
\ref{prop:v3v79-recursion} then cancels successive orders.

For a tree-produced primitive, Theorem~\ref{thm:v3v80-tree-primitive}
adds at most one diameter factor.  Inequality
\eqref{eq:v3v80-exponent-trade} absorbs it after reducing the exponential
rate, with a constant independent of \(n\) and volume.
\end{proof}

\begin{corollary}[Where the Standard Model cancellation is already local]
\label{cor:v3v80-sm-leading-local}
At the leading nontrivial continuum descent order, the complete one-generation
Standard Model sum has zero coefficient for each independent invariant
polynomial.  At that order the anomalous cluster density cancels before a
Poincar\'e inverse is applied, so the dipole obstruction of
Theorem~\ref{thm:v3v80-dipole} is absent.  The unresolved terms are the
finite-spacing irrelevant overlap artifacts beyond those invariant
polynomials.
\end{corollary}

\begin{proof}
The independent representation coefficients and their separate vanishing
were proved in Proposition~\ref{prop:v3v66-coefficient-cancellation} and
Theorem~\ref{thm:v3v79-leading-cancellation}.  Since each coefficient
multiplies a local invariant polynomial density, the complete
representation sum vanishes at the density level, before spatial
summation.  The conclusion does not apply to higher-order lattice artifacts
whose local factorization has not yet been computed.
\end{proof}

\subsection{What has been decided and the next upstream move}
\label{sec:v3v80-decision}

V80 gives a constructive primitive and a sharp diagnosis.  A box-anchored
homotopy always solves the neutral linearized equation, but necessarily
pays one spatial moment and can spread across the cluster diameter.  A
convergent exponentially local counterterm follows if the overlap defect is
first decomposed into neutral cellwise descents with a common geometric
Taylor bound.  Vanishing of the globally summed anomaly class is too weak
to imply that decomposition.

\begin{firewall}[No all-order measure has been constructed]
\label{fw:v3v80-no-all-order}
Theorem~\ref{thm:v3v80-cellwise-criterion} is conditional.  V80 neither
constructs the cluster identities \eqref{eq:v3v80-cluster-descent} for all
overlap coefficients nor proves their common bound.  It therefore does not
establish exact finite-spacing Standard Model gauge invariance.
\end{firewall}

\begin{firewall}[The global and reflection problems remain separate]
\label{fw:v3v80-global-reflection}
The local small-link argument does not determine holonomy under large gauge
loops or the global Standard Model quotient, and it does not prove
reflection positivity for the determinant ratio or its Pauli--Villars
completion.
\end{firewall}

The next step goes upstream to the source of the cellwise identity.  V81
should compare the overlap-projector curvature with a continuum local
projector along a gapped interpolation and derive an exact transgression
formula.  If the transgression is a uniformly local cluster coboundary, the
continuum anomaly cancellation will lift to the lattice without applying a
global inverse.  If it is not, the first failing irrelevant coefficient
must be displayed explicitly.


\section{V81: equivariant projector transgression before norming}
\label{sec:v3v81}

V80 showed that applying a Poincar\'e inverse to a globally cancelled defect
is the wrong operation: even a separated abelian dipole forces a first
spatial moment.  This note goes upstream.  For a path of covariant chiral
projectors, curvature and gauge moment map satisfy one exact equivariant
transgression identity.  It is assembled at the operator level before any
kernel norm is taken, precisely as required by the V80 audit.

\subsection{Finite-dimensional projector calculus}
\label{sec:v3v81-calculus}

Let \(\mathcal U\) be the finite-lattice small-link patch, let \(d\) denote
the exterior derivative on \(\mathcal U\), and let
\(P_t(A)=P_t(A)^*=P_t(A)^2\), \(0\le t\le1\), be a smooth constant-rank
family.  Put a dot for \(\partial_t\).  Differentiating \(P_t^2=P_t\)
gives
\begin{equation}
 P_t\dot P_tP_t=0,
 \qquad
 P_t(dP_t)P_t=0;
\label{eq:v3v81-offdiagonal}
\end{equation}
thus both \(\dot P_t\) and \(dP_t\) are off diagonal relative to
\(\operatorname{ran}P_t\oplus\ker P_t\).

Define the Grassmannian curvature and the transgression one-form by
\begin{align}
 \mathcal F_t
 &= i\operatorname{Tr}\bigl(P_t,dP_t\wedge dP_t\bigr),
 \label{eq:v3v81-curvature}\\
 \mathcal B_t
 &=i\operatorname{Tr}\bigl(P_t[\dot P_t,dP_t]\bigr).
 \label{eq:v3v81-transgression-form}
\end{align}
The commutator in \eqref{eq:v3v81-transgression-form} is the ordinary
matrix commutator; \(\dot P_t\) has form degree zero and \(dP_t\) has form
degree one.

\begin{theorem}[Exact projector transgression]
\label{thm:v3v81-projector-transgression}
For every smooth constant-rank projector path,
\begin{equation}
 \partial_t\mathcal F_t=d\mathcal B_t.
\label{eq:v3v81-transgression}
\end{equation}
Consequently
\begin{equation}
 \mathcal F_1-\mathcal F_0=d\mathcal B,
 \qquad
 \mathcal B=\int_0^1\mathcal B_t\,dt.
\label{eq:v3v81-integrated-transgression}
\end{equation}
\end{theorem}

\begin{proof}
Differentiate \eqref{eq:v3v81-curvature}.  The term
\(\operatorname{Tr}(\dot P_t,dP_t\wedge dP_t)\) vanishes: the product of
the two one-form factors is block diagonal by
\eqref{eq:v3v81-offdiagonal}, whereas \(\dot P_t\) is off diagonal.  Hence
\begin{equation}
 \partial_t\mathcal F_t
 =i\operatorname{Tr}\bigl(
 P_t,d\dot P_t\wedge dP_t+P_t,dP_t\wedge d\dot P_t
 \bigr).
\label{eq:v3v81-differentiated-curvature}
\end{equation}
On the other hand, the graded Leibniz rule gives
\begin{align}
 d\mathcal B_t/i
 =\operatorname{Tr}\bigl(&dP_t\,\dot P_t\wedge dP_t
 +P_t,d\dot P_t\wedge dP_t\\
 &-dP_t\wedge dP_t\,\dot P_t
 +P_t,dP_t\wedge d\dot P_t\bigr).
\end{align}
The first and third traces cancel by graded cyclicity.  The two remaining
terms are exactly \eqref{eq:v3v81-differentiated-curvature}.  Integrating in
\(t\) proves \eqref{eq:v3v81-integrated-transgression}.
\end{proof}

\subsection{The gauge moment map is the degree-zero partner}
\label{sec:v3v81-equivariant}

Let \(R(c)\) be the infinitesimal anti-Hermitian gauge action on the
one-particle space.  Gauge covariance of the path means
\begin{equation}
 \iota_{X_c}dP_t=X_cP_t=[R(c),P_t],
 \qquad
 [\partial_t,R(c)]=0.
\label{eq:v3v81-covariance}
\end{equation}
With the convention compatible with
\eqref{eq:v3v81-curvature}, define
\begin{equation}
 \mu_t(c)=i\operatorname{Tr}\bigl(R(c)P_t\bigr).
\label{eq:v3v81-moment-map}
\end{equation}
For an abelian trace component one may subtract an \(A\)-independent
reference projector; Proposition~\ref{prop:v3v79-zero-order} shows that this
constant disappears in the complete Standard Model sum.

\begin{theorem}[Equivariant transgression identity]
\label{thm:v3v81-equivariant-transgression}
If \eqref{eq:v3v81-covariance} holds, then
\begin{equation}
 \iota_{X_c}\mathcal B_t=\partial_t\mu_t(c).
\label{eq:v3v81-moment-transgression}
\end{equation}
Writing \(d_c=d-\iota_{X_c}\), equations
\eqref{eq:v3v81-transgression} and
\eqref{eq:v3v81-moment-transgression} combine into
\begin{equation}
 \partial_t\bigl(\mathcal F_t-\mu_t(c)\bigr)
 =d_c\mathcal B_t.
\label{eq:v3v81-equivariant-exact}
\end{equation}
This is an identity for the complete finite-lattice trace; no descent
inverse or termwise norm estimate enters its proof.
\end{theorem}

\begin{proof}
At a fixed \((t,A)\), decompose matrices relative to \(P_t\):
\begin{equation}
 P_t=\begin{pmatrix}1&0\\0&0\end{pmatrix},\quad
 \dot P_t=\begin{pmatrix}0&u\\v&0\end{pmatrix},\quad
 R(c)=\begin{pmatrix}a&b\\c_0&d_0\end{pmatrix}.
\label{eq:v3v81-blocks}
\end{equation}
Then
\([R(c),P_t]=\left(\begin{smallmatrix}0&-b\\c_0&0\end{smallmatrix}\right)\).
Substitution in \eqref{eq:v3v81-transgression-form} gives
\begin{equation}
 \iota_{X_c}\mathcal B_t
 =i\operatorname{Tr}(uc_0+bv)
 =i\operatorname{Tr}(R(c)\dot P_t)
 =\partial_t\mu_t(c).
\end{equation}
Together with Theorem~\ref{thm:v3v81-projector-transgression}, this is
\eqref{eq:v3v81-equivariant-exact}.
\end{proof}

\begin{corollary}[Transport of an invariant determinant connection]
\label{cor:v3v81-connection-transport}
Suppose \(\mathcal A_0\) is a gauge-invariant local connection one-form with
\(d\mathcal A_0=\mathcal F_0\).  Then
\begin{equation}
 \mathcal A_1^{\rm tr}=\mathcal A_0+\mathcal B
\label{eq:v3v81-transported-connection}
\end{equation}
has curvature \(\mathcal F_1\).  If the whole path is gauge covariant, then
\(\mathcal B\) and \(\mathcal A_1^{\rm tr}\) are gauge invariant, and
\begin{equation}
 \iota_{X_c}\mathcal A_1^{\rm tr}
 =\iota_{X_c}\mathcal A_0+\mu_1(c)-\mu_0(c).
\label{eq:v3v81-transported-moment}
\end{equation}
\end{corollary}

\begin{proof}
The curvature statement follows from
\eqref{eq:v3v81-integrated-transgression}.  Naturality of the trace under
simultaneous conjugation makes \(\mathcal B_t\) gauge invariant.  Integrate
\eqref{eq:v3v81-moment-transgression} to obtain
\eqref{eq:v3v81-transported-moment}.
\end{proof}

This corollary identifies the useful object more precisely than an
order-by-order scalar counterterm.  A local covariant path transports the
entire determinant connection, including the moment-map term, through a
single one-form \(\mathcal B\).  The two components are never separated
before cancellation.

\subsection{Uniform locality of the transgression}
\label{sec:v3v81-locality}

For kernels on a bounded-degree lattice, write
\begin{equation}
 \|K\|_\gamma=
 \max\left\{
 \sup_x\sum_y e^{\gamma d(x,y)}\|K(x,y)\|,
 \sup_y\sum_x e^{\gamma d(x,y)}\|K(x,y)\|
 \right\}.
\label{eq:v3v81-kernel-norm}
\end{equation}
This weighted Schur norm is submultiplicative after reducing the exponent
by an arbitrarily small amount; for finite-range bounded-degree lattices it
has constants independent of the volume.

\begin{theorem}[Local transgression bound]
\label{thm:v3v81-local-transgression}
Assume that, uniformly in \(t\), volume, and
\(\|A\|_\infty\le\eta'<\eta\), the kernels of
\(P_t\), \(\dot P_t\), and every first field derivative \(D_aP_t\) obey
weighted Schur bounds at rate \(\gamma>0\).  Then the coefficient of
\(dA_a\) in \(\mathcal B_t\) is exponentially local at every rate
\(\gamma'<\gamma\), uniformly in \(t\) and volume.

For the coefficient assertion, assume in addition the standard joint
locality estimate furnished by differentiated resolvent products: for each
\(n\), the \(n\)-linear kernels of \(P_t\), \(\dot P_t\), and \(D_aP_t\)
are bounded by \(CB^n e^{-\gamma\tau}\), where \(\tau\) is the length of a
minimal lattice tree joining the two kernel sites and all field-insertion
cells.  Then the Taylor coefficients of \(\mathcal B\) obey a common
geometric bound
\begin{equation}
 \sup_x\sum_{Z\ni x}
 e^{\gamma'\operatorname{diam}Z}\|\mathcal B_Z^{(n)}\|
 \le C'B'^n.
\label{eq:v3v81-transgression-cluster-bound}
\end{equation}
\end{theorem}

\begin{proof}
The coefficient in direction \(a\) is the trace density of
\begin{equation}
 iP_t\bigl(\dot P_tD_aP_t-D_aP_t\dot P_t\bigr).
\label{eq:v3v81-local-product}
\end{equation}
Expand its diagonal kernel as the sum over the two intermediate lattice
sites.  The triangle inequality for graph distance distributes the weight
among the three factors, and the weighted Schur estimates make both sums
uniform.  A small loss from \(\gamma\) to \(\gamma'\) absorbs the polynomial
number of possible intermediate cells.

For the coefficient statement, insert the assumed multilinear tree bounds
in \eqref{eq:v3v81-local-product}.  The union of the three factor trees and
the two contracted kernel paths contains a connected tree through all
external insertions.  Summing the contracted sites costs only a smaller
decay rate.  The coefficient of total degree \(n\) is a finite convolution
over three nonnegative degrees; its polynomial multiplicity is absorbed by
increasing \(B\) to \(B'\).  Grouping equal connected supports gives
\eqref{eq:v3v81-transgression-cluster-bound}.  Integration over the compact
parameter interval changes only \(C'\).
\end{proof}

\begin{corollary}[No first-moment inverse along a local path]
\label{cor:v3v81-no-inverse}
Under the hypotheses of Theorem~\ref{thm:v3v81-local-transgression}, the
difference of equivariant curvatures of the endpoint projectors has a local
primitive satisfying the cellwise geometric estimate required in V80.  The
construction uses multiplication, trace, and the bounded \(t\)-integral;
it never connects separated charges by a tree current.
\end{corollary}

\subsection{The topological condition on the reference path}
\label{sec:v3v81-topology}

The transgression theorem is exact, but it does not manufacture a suitable
endpoint \(P_0\).  This point is substantive.  A translation-invariant
gapped finite-range Hamiltonian defines a momentum-space projector and
hence a \(K\)-theory class.  That class is constant along every uniformly
gapped local path.

\begin{proposition}[A trivial projector is not an admissible reference]
\label{prop:v3v81-no-trivial-reference}
In the one-species Wilson window of V76, the free overlap projector has the
nonzero degree computed there.  It cannot be joined to an on-site constant
projector by a translation-invariant uniformly gapped local Hamiltonian
path.
\end{proposition}

\begin{proof}
Spectral flattening turns a uniformly gapped Hamiltonian path into a
continuous path of momentum-space projectors.  The Chern character, or
equivalently the degree of the normalized Wilson symbol used in V76, is
integer-valued and homotopy invariant.  The on-site projector has constant
symbol and zero degree, whereas the overlap projector in the one-species
window has nonzero degree.  Equality would contradict homotopy invariance.
\end{proof}

Thus a useful reference must lie in the same local topological phase as the
overlap projector.  The continuum leading symbol has the correct anomaly
class, but it must be regularized on the same lattice Hilbert space and
connected through a path with a volume-independent gap.  Merely matching
the anomaly polynomial does not prove that such a path exists.

\begin{firewall}[Transgression is conditional on a local gapped path]
\label{fw:v3v81-path-conditional}
V81 proves that any covariant uniformly local projector path transports the
complete equivariant curvature with convergent local bounds.  It does not
construct a same-phase reference projector or prove a gap for an
interpolation from that reference to the interacting overlap projector.
\end{firewall}

\begin{firewall}[No conclusion about large gauge loops]
\label{fw:v3v81-no-global}
The identity \eqref{eq:v3v81-equivariant-exact} holds on the contractible
small-link patch.  It does not trivialize determinant-line holonomy around
large gauge transformations, change the global gauge-group quotient, or
settle the reflection-positive Pauli--Villars ratio.
\end{firewall}

\subsection{V81 conclusion and V82 target}
\label{sec:v3v81-conclusion}

The overlap curvature and its gauge moment map are now tied together by an
exact equivariant transgression.  A uniformly local covariant path would
produce the complete cellwise primitive required by V80 with a common
convergence radius, avoiding the sharp nonlocality of an anchored
Poincar\'e inverse.  The bottleneck has moved upstream to the construction
of a same-topological-phase reference and a volume-uniformly gapped path.

V82 should construct the quasi-adiabatic intertwiner for an arbitrary
finite-range covariant Hamiltonian path with a uniform gap, prove its
volume-independent exponential locality, and identify exactly which
endpoint data remain to be supplied for a Standard Model reference.  This
separates the analytic transport theorem from the genuinely topological
existence problem.


\section{V82: local spectral flow and the obstruction to a free covariant reference}
\label{sec:v3v82}

V81 reduced local determinant transport to the existence of a gauge-covariant
uniformly gapped Hamiltonian path.  This note proves the analytic part of
that statement: such a path has an exact, volume-uniformly local spectral
flow which intertwines its chiral projectors.  It then proves that the
tempting radial endpoint \(P(0)\) cannot supply the required covariant
reference.  The reference must itself depend on the gauge field and lie in
the Wilson topological phase.

\subsection{Uniformly gapped local paths}
\label{sec:v3v82-path}

Let \(H_t(A)=H_t(A)^*\), \(0\le t\le1\), act on the finite-lattice
one-particle space.  Assume constants \(M,g,\gamma_0,C>0\), independent of
the volume, such that
\begin{equation}
 \|H_t\|\le M,
 \qquad
 \operatorname{spec}H_t\cap(-g,g)=\varnothing,
 \qquad
 \|H_t\|_{\gamma_0}+\|\dot H_t\|_{\gamma_0}\le C.
\label{eq:v3v82-path-assumptions}
\end{equation}
Here \(\|\cdot\|_{\gamma_0}\) is the weighted Schur norm of
\eqref{eq:v3v81-kernel-norm}.  Put
\begin{equation}
 P_t=\mathbf1_{(-\infty,0)}(H_t).
\label{eq:v3v82-negative-projector}
\end{equation}

\begin{lemma}[Uniform local Riesz derivative]
\label{lem:v3v82-riesz-locality}
Under \eqref{eq:v3v82-path-assumptions}, there is
\(0<\gamma<\gamma_0\) and \(C_1<\infty\), independent of volume and
\(t\), such that
\begin{equation}
 \|P_t\|_\gamma+\|\dot P_t\|_\gamma\le C_1.
\label{eq:v3v82-projector-locality}
\end{equation}
Moreover
\begin{equation}
 \dot P_t=\frac{1}{2\pi i}\int_\Gamma
 (z-H_t)^{-1}\dot H_t(z-H_t)^{-1}\,dz,
\label{eq:v3v82-riesz-derivative}
\end{equation}
where a single contour \(\Gamma\) may be chosen for all \(t\) and all
volumes.
\end{lemma}

\begin{proof}
The spectral and norm bounds place the negative spectrum in
\([-M,-g]\) and the positive spectrum in \([g,M]\).  Choose a fixed
rectangular contour enclosing the first interval and staying a fixed
distance from both intervals.  Conjugating the resolvent by the lattice
weight \(e^{\gamma d(\,cdot\,,x_0)}\) changes \(H_t\) by an operator of
norm \(O(\gamma)\), uniformly in \(x_0,t\), and volume.  For sufficiently
small \(\gamma\), the resolvent identity and a Neumann series therefore
give a uniform weighted resolvent bound on \(\Gamma\).  The Riesz formula
for \(P_t\) proves its bound.  Differentiating that formula gives
\eqref{eq:v3v82-riesz-derivative}; submultiplicativity of the weighted
Schur norm and the bound on \(\dot H_t\) prove the second estimate.
\end{proof}

\subsection{Exact local spectral flow}
\label{sec:v3v82-flow}

Define the anti-Hermitian generator
\begin{equation}
 Q_t=[\dot P_t,P_t],
 \qquad Q_t^*=-Q_t,
\label{eq:v3v82-generator}
\end{equation}
and let \(U_t\) solve
\begin{equation}
 \dot U_t=Q_tU_t,
 \qquad U_0=I.
\label{eq:v3v82-unitary-ode}
\end{equation}

\begin{theorem}[Volume-uniform local projector intertwiner]
\label{thm:v3v82-local-intertwiner}
Under \eqref{eq:v3v82-path-assumptions}, \(U_t\) is unitary and
\begin{equation}
 P_t=U_tP_0U_t^*.
\label{eq:v3v82-intertwining}
\end{equation}
For every sufficiently small \(\gamma>0\),
\begin{equation}
 \sup_{0\le t\le1}\|U_t\|_\gamma
 \le \exp\left(\int_0^1\|Q_s\|_\gamma\,ds\right)
 \le e^{C_2},
\label{eq:v3v82-unitary-locality}
\end{equation}
with \(C_2\) independent of the volume.  Hence conjugation by \(U_t\)
preserves exponential locality, after an arbitrarily small reduction of the
decay rate.
\end{theorem}

\begin{proof}
Anti-Hermiticity of \(Q_t\) and differentiation of \(U_t^*U_t\) prove
unitarity.  From \(P_t^2=P_t\),
\begin{equation}
 P_t\dot P_tP_t=0,
 \qquad
 \dot P_t=P_t\dot P_t+\dot P_tP_t.
\end{equation}
A direct calculation then gives
\begin{equation}
 [[\dot P_t,P_t],P_t]=\dot P_t.
\label{eq:v3v82-lax}
\end{equation}
Thus both \(P_t\) and \(U_tP_0U_t^*\) solve the same Lax equation
\(\dot P=[Q_t,P]\) with the same initial condition, proving
\eqref{eq:v3v82-intertwining}.

Lemma~\ref{lem:v3v82-riesz-locality} and submultiplicativity give a uniform
bound on \(Q_t\) in the weighted norm.  The Dyson series for
\eqref{eq:v3v82-unitary-ode} is absolutely convergent there, and its
\(k\)-th term is bounded by
\(k!^{-1}(\int_0^t\|Q_s\|_\gamma ds)^k\).  Summation proves
\eqref{eq:v3v82-unitary-locality}.  Two applications of the same product
bound control \(U_tKU_t^*\) for every exponentially local \(K\).
\end{proof}

\begin{proposition}[Covariance and analytic cluster bounds]
\label{prop:v3v82-covariant-analytic-flow}
Suppose, in addition, that
\begin{equation}
 H_t(A^g)=G(g)H_t(A)G(g)^{-1}
\label{eq:v3v82-H-covariance}
\end{equation}
and that the differentiated kernels of \(H_t\) and \(\dot H_t\) have the
uniform multilinear tree bounds used in
Theorem~\ref{thm:v3v81-local-transgression}.  Then
\begin{equation}
 U_t(A^g)=G(g)U_t(A)G(g)^{-1},
\label{eq:v3v82-U-covariance}
\end{equation}
and the Taylor kernels of \(P_t,Q_t,U_t\), and the V81 transgression
one-form have volume-independent connected-cluster bounds
\(C'B'^ne^{-\gamma'\tau}\) for every sufficiently small
\(\gamma'<\gamma\).
\end{proposition}

\begin{proof}
Functional calculus carries \eqref{eq:v3v82-H-covariance} to \(P_t\) and
then to \(Q_t\).  The two sides of \eqref{eq:v3v82-U-covariance} solve the
same initial-value problem, so uniqueness proves covariance.

Differentiate the Riesz formula
\eqref{eq:v3v82-riesz-derivative} repeatedly.  Every derivative is a finite
sum of alternating resolvents and local differentiated vertices.  Joining
the resolvent paths gives a tree through all external insertion cells; the
Combes--Thomas factors give \(e^{-\gamma'\tau}\).  The number of ordered
partitions of the derivatives is absorbed into \(B'^n\).  Differentiating
the Dyson series adds time-ordered products of such kernels; the simplex
factor \(1/k!\) makes their sum convergent.  The product formula
\eqref{eq:v3v81-local-product} then gives the same statement for the
transgression one-form.
\end{proof}

The proposition supplies all analytic hypotheses in V81 from a single
concrete input: a covariant local Hamiltonian path with a common gap.

\subsection{Why radial scaling is not the covariant path}
\label{sec:v3v82-radial}

The small-link gap of V78 does give a path
\(H_t(A)=H_W(tA)\) in logarithmic coordinates.  It is useful for the radial
connection, but local gauge transformations act affinely on a connection:
in continuum notation,
\begin{equation}
 (tA)^g=t\,gAg^{-1}-(dg)g^{-1},
 \qquad
 t(A^g)=t\,gAg^{-1}-t\,(dg)g^{-1},
\end{equation}
whereas gauge transforming the zero endpoint produces a pure gauge field,
not the zero field.  More invariantly, the identity-link configuration is
not fixed by independent gauge rotations at adjacent sites.  Therefore the
radial family does not obey \eqref{eq:v3v82-H-covariance} as a two-variable
identity in \((A,g)\).  This is exactly why the radial connection in V79
was not automatically gauge invariant.

The failure cannot be repaired by selecting another field-independent
projector in the same Wilson phase.

\begin{theorem}[Field-independent covariant references are topologically
trivial]
\label{thm:v3v82-reference-obstruction}
Let \(P_{\rm ref}\) be independent of the gauge links and commute with every
sitewise gauge transformation in a nontrivial irreducible charged
representation.  Then
\begin{equation}
 P_{\rm ref}(x,y)=0\quad(x\ne y),
 \qquad
 P_{\rm ref}(x,x)=p_{\rm spin}\otimes I_{\rm int}
\label{eq:v3v82-onsite-reference}
\end{equation}
for a site-independent translation-invariant reference.  Its momentum
symbol is constant and all positive-degree Chern characters vanish.
Consequently it cannot lie in the one-species Wilson phase of V76.
\end{theorem}

\begin{proof}
For arbitrary independent values of \(g(x)\) and \(g(y)\), covariance and
field independence require
\begin{equation}
 \rho(g(x))P_{\rm ref}(x,y)\rho(g(y))^{-1}
 =P_{\rm ref}(x,y).
\label{eq:v3v82-independent-gauge-values}
\end{equation}
When \(x\ne y\), vary \(g(x)\) while fixing \(g(y)\).  A nontrivial
irreducible representation has no invariant vector in the charged internal
factor, so the off-diagonal intertwiner is zero.  At \(x=y\), Schur's lemma
forces the internal factor to be scalar, leaving an arbitrary spin
projector.  Translation invariance makes that projector site independent.
The Fourier symbol of \eqref{eq:v3v82-onsite-reference} is therefore
constant; its derivatives vanish, so every positive-degree Chern form and
the Wilson-symbol degree vanish.  V76 computed a nonzero degree in the
one-species window, and homotopy invariance rules out a uniformly gapped
local path between the two.
\end{proof}

\begin{corollary}[Necessary form of a usable endpoint]
\label{cor:v3v82-reference-must-vary}
A reference capable of transporting the Standard Model determinant data by
V81 must be a gauge-field-dependent covariant lattice projector with the
same free topological degree as the overlap projector.  Its anomaly
cancellation must be known locally, rather than obtained by declaring its
field dependence absent.
\end{corollary}

\subsection{A quantitative conditional transport theorem}
\label{sec:v3v82-conditional}

\begin{theorem}[Reduction to a same-phase covariant reference]
\label{thm:v3v82-reduction}
For each Standard Model Weyl representation \(R\), suppose there is a
finite-range covariant reference Hamiltonian \(H_{{\rm ref},R}(A)\) and a
covariant interpolation \(H_{t,R}(A)\) to the Wilson kernel such that:
\begin{enumerate}
 \item the bounds \eqref{eq:v3v82-path-assumptions} and the joint tree
 estimates hold with constants common to the finite representation list;
 \item the complete reference determinant connection has a locally
 assembled vanishing gauge defect; and
 \item the reference and Wilson projectors have the same ranks and free
 topological degrees.
\end{enumerate}
Then the integrated V81 transgression transports the reference connection
to an exponentially local connection for the overlap projector with a
common convergent small-link expansion.  Its infinitesimal gauge defect
vanishes on that patch.
\end{theorem}

\begin{proof}
Theorem~\ref{thm:v3v82-local-intertwiner} and
Proposition~\ref{prop:v3v82-covariant-analytic-flow} verify the locality and
cluster hypotheses of Theorem~\ref{thm:v3v81-local-transgression}.
Equation \eqref{eq:v3v81-transported-connection} gives the endpoint
connection and \eqref{eq:v3v81-transported-moment} transports its complete
gauge contraction.  Because the interpolation is covariant, the curvature
and moment-map differences occur as the two components of the same
equivariantly exact local one-form.  Adding that form to the assumed
reference connection therefore preserves the assembled Ward identity.
The finite representation list permits the maximum of all constants to be
used, giving a common convergence radius.
\end{proof}

\begin{firewall}[The reference hypothesis is not yet discharged]
\label{fw:v3v82-reference-open}
No gauge-field-dependent same-phase reference satisfying item 2 of
Theorem~\ref{thm:v3v82-reduction} has yet been constructed.  The theorem is
a rigorous reduction, not a proof of the finite-spacing Standard Model
measure.
\end{firewall}

\begin{firewall}[Spectral flow does not settle reflection positivity]
\label{fw:v3v82-rp-open}
Exponential locality and unitary projector transport do not imply
Osterwalder--Schrader positivity of a chiral determinant ratio.  The
two-wall and Pauli--Villars reflection problem remains independent.
\end{firewall}

\subsection{V82 conclusion and V83 target}
\label{sec:v3v82-conclusion}

Every covariant uniformly gapped local Hamiltonian path now has an exact
volume-uniformly local intertwiner and the joint cluster estimates needed by
the V81 transgression.  The remaining obstruction is not analytic: a
field-independent covariant reference is forced to be on-site and hence
topologically trivial.  A usable reference must carry the gauge field and
the same Wilson degree.

V83 should construct a candidate reference by covariant blocking from a
fine Wilson lattice to a coarse lattice, or show why such blocking cannot
preserve a common gap.  The comparison must remain within the same
topological phase and must expose its local determinant defect before any
norm is taken.


\section{V83: a covariantly blocked same-phase reference}
\label{sec:v3v83}

V82 proved that the reference projector must depend covariantly on the gauge
field and must retain the nonzero Wilson degree.  This note constructs such
a reference on the original fine Hilbert space.  A covariant block isometry
carries one physical coarse Wilson band, while the orthogonal fine modes
receive an explicit massive involution.  The resulting Hamiltonian is local,
covariant, uniformly gapped whenever the coarse kernel is gapped, and has the
correct relative topological class.  Its determinant data reduce to the
coarse determinant plus a finite-range wavelet transgression.

\subsection{Gauge-covariant block isometry}
\label{sec:v3v83-isometry}

Partition the four-dimensional fine lattice into disjoint hypercubes
\(C_b\) of fixed side \(q\), with anchor \(b\) in each block.  For every
\(x\in C_b\), fix once and for all a path inside \(C_b\) from \(b\) to
\(x\), and denote its fine-link parallel transport by \(T_U(x,b)\).  Choose
weights \(w_x\), independent of \(b\), with
\begin{equation}
 \sum_{x\in C_b}|w_x|^2=1.
\label{eq:v3v83-weights}
\end{equation}
For a coarse spinor \(\phi\), define
\begin{equation}
 (B_U\phi)(x)=w_xT_U(x,b)\phi(b),
 \qquad x\in C_b.
\label{eq:v3v83-block-map}
\end{equation}
The map acts trivially on spin and flavour indices.

Let \(g_f(x)\) be an arbitrary fine sitewise gauge transformation and let
\(g_c(b)=g_f(b)\) be its restriction to anchors.  Parallel transport obeys
\begin{equation}
 T_{U^{g_f}}(x,b)=g_f(x)T_U(x,b)g_f(b)^{-1}.
\label{eq:v3v83-transport-covariance}
\end{equation}

\begin{lemma}[Isometry and covariance]
\label{lem:v3v83-isometry}
The block map satisfies
\begin{equation}
 B_U^*B_U=I_c,
 \qquad
 B_{U^{g_f}}G_c(g_c)=G_f(g_f)B_U.
\label{eq:v3v83-isometry-covariance}
\end{equation}
Consequently \(\Pi_U=B_UB_U^*\) is an orthogonal finite-range projector and
\begin{equation}
 \Pi_{U^{g_f}}=G_f(g_f)\Pi_UG_f(g_f)^{-1}.
\label{eq:v3v83-Pi-covariance}
\end{equation}
Its kernel vanishes unless its two sites belong to the same block.
\end{lemma}

\begin{proof}
Unitarity of each parallel transport and
\eqref{eq:v3v83-weights} give
\(\langle B_U\phi,B_U\psi\rangle=\langle\phi,\psi\rangle\).
Equation \eqref{eq:v3v83-transport-covariance} inserted in
\eqref{eq:v3v83-block-map} proves the intertwining identity.  The projector
claims follow immediately.  Both factors in its kernel must lie in the
same image block, proving finite range.
\end{proof}

Define a coarse link between neighbouring anchors as parallel transport
along a fixed fine path joining them:
\begin{equation}
 U_c(b,b')=T_U(b,b').
\label{eq:v3v83-coarse-link}
\end{equation}
It transforms as
\(U_c^{g_f}(b,b')=g_c(b)U_c(b,b')g_c(b')^{-1}\).  If the fine links have
logarithms bounded by \(\eta\), a fixed-block Baker--Campbell--Hausdorff
estimate gives \(\|\log U_c\|\le C_q\eta\) as long as the product remains
inside the selected logarithmic chart.  Thus the coarse small-link gap is
available after shrinking the fine patch by the fixed factor \(C_q\).

\subsection{The blocked reference Hamiltonian}
\label{sec:v3v83-reference}

Let \(H_c(U_c)\) be the Hermitian Wilson kernel on the coarse lattice in the
same one-species mass window as V76.  Let \(S=S^*=S^{-1}\) be a fixed spin
involution commuting with \(B_U\) and \(\Pi_U\), and put
\(P_S^-=(I-S)/2\).  For a constant \(m_s>0\), define on the fine Hilbert
space
\begin{equation}
 H_{\rm blk}(U)
 =B_UH_c(U_c)B_U^*+m_sS(I-\Pi_U).
\label{eq:v3v83-block-Hamiltonian}
\end{equation}

\begin{theorem}[Exact block decomposition and gap]
\label{thm:v3v83-block-gap}
The orthogonal decomposition
\begin{equation}
 \mathcal H_f=\operatorname{ran}B_U\oplus\ker B_U^*
\label{eq:v3v83-Hilbert-split}
\end{equation}
reduces \(H_{\rm blk}\).  On its two summands the restrictions are unitarily
equivalent to \(H_c\) and to \(m_sS\), respectively.  Hence
\begin{equation}
 \operatorname{dist}\bigl(0,\operatorname{spec}H_{\rm blk}\bigr)
 \ge \min\{g_c,m_s\}
\label{eq:v3v83-block-gap}
\end{equation}
whenever the coarse Wilson gap is \(g_c>0\), and its negative spectral
projector is
\begin{equation}
 P_{\rm blk}(U)
 =B_UP_c(U_c)B_U^*+(I-\Pi_U)P_S^-.
\label{eq:v3v83-block-projector}
\end{equation}
The Hamiltonian and projector are gauge covariant and exponentially local
with constants independent of the total volume.
\end{theorem}

\begin{proof}
The first term of \eqref{eq:v3v83-block-Hamiltonian} vanishes on
\(\ker B_U^*\), while the second vanishes on \(\operatorname{ran}B_U\).
The identity \(B_U^*B_U=I_c\) identifies the first restriction with
\(H_c\).  Since \(S\) commutes with \(\Pi_U\), its restriction to the
complement remains an involution.  The spectrum is therefore the disjoint
union of the coarse spectrum and the two spectator masses \(\{\pm m_s\}\),
which proves \eqref{eq:v3v83-block-gap} and
\eqref{eq:v3v83-block-projector}.  Lemma~\ref{lem:v3v83-isometry} and
covariance of \(H_c\) prove covariance.  The block maps have fixed finite
range; coarse nearest-neighbour couplings lift to a range bounded in terms
of \(q\), and the Riesz estimate of Lemma~\ref{lem:v3v82-riesz-locality}
gives projector locality.
\end{proof}

Choose the multiplicity of the negative eigenspace of \(S\) so that the
rank in \eqref{eq:v3v83-block-projector} equals the rank of the fine overlap
projector on the small-link component.  For the usual four-spinor Wilson
kernel one may take a balanced spin involution; rank constancy then extends
the free equality throughout the gapped patch.

\subsection{Relative topological class}
\label{sec:v3v83-topology}

At the free field, regard a \(q^4\)-site block as the unit cell.  The map
\(B_1(k)\) is a smooth periodic isometry over the reduced Brillouin torus.
Because it is spin diagonal,
\begin{equation}
 P_{\rm blk}-P_S^-
 =B_1\bigl(P_c-P_S^-\bigr)B_1^*.
\label{eq:v3v83-relative-projector}
\end{equation}

\begin{theorem}[The block reference retains the Wilson class]
\label{thm:v3v83-same-class}
The reduced \(K\)-theory class of \(P_{\rm blk}\), measured relative to the
constant spectator projector \(P_S^-\), equals the reduced class of the
coarse Wilson projector.  In particular, if the coarse kernel is in the
one-species window, the block reference has the same nonzero degree as the
fine Wilson projector viewed with the enlarged unit cell.
\end{theorem}

\begin{proof}
Equation \eqref{eq:v3v83-relative-projector} identifies the virtual bundle
\([P_{\rm blk}]-[P_S^-]\) with the isometric image of
\([P_c]-[P_S^-]\).  An isometric bundle embedding preserves the stable
class; the orthogonal spectator complement adds the same trivial bundle to
both sides.  The Chern character and its top-degree integral are stable
invariants, so they agree.  Enlarging the unit cell only folds the Brillouin
zone and does not change the physical Wilson degree.  V76 computed that
degree to be the one-species value throughout \(0<aM<2r\).
\end{proof}

This removes the topological obstruction of
Theorem~\ref{thm:v3v82-reference-obstruction}: the reference is not
field-independent, and its physical band carries the required class.

\subsection{Determinant-line factorization}
\label{sec:v3v83-determinant}

On each contractible small-link block patch, extend the normalized vector
\((w_xT_U(x,b))_{x\in C_b}\) to an orthonormal frame by applying a fixed
Householder or Gram--Schmidt chart whose chosen pivot stays nonzero after
shrinking the patch.  Block by block this gives a unitary
\(W_U\), of range at most one block, whose first channel is \(B_U\).  It is
analytic and covariant, and it diagonalizes \(\Pi_U\):
\begin{equation}
 W_U^*P_{\rm blk}W_U
 =P_c(U_c)\oplus P_{S,\rm high}^-.
\label{eq:v3v83-wavelet-split}
\end{equation}

\begin{proposition}[Coarse determinant plus a local wavelet term]
\label{prop:v3v83-determinant-factorization}
On the block chart, the determinant line of \(P_{\rm blk}\) is canonically
isomorphic to
\begin{equation}
 \det\operatorname{ran}P_c
 \otimes\det\operatorname{ran}P_{S,\rm high}^-.
\label{eq:v3v83-line-factorization}
\end{equation}
The second factor has a fixed coordinate frame before applying \(W_U\).
Consequently the block-reference connection equals the pulled-back coarse
connection plus the finite-range one-form
\begin{equation}
 \mathcal A_{\rm wav}
 =i\operatorname{Tr}\bigl(
 (P_c\oplus P_{S,\rm high}^-),W_U^*dW_U
 \bigr).
\label{eq:v3v83-wavelet-connection}
\end{equation}
Every Taylor coefficient of \(\mathcal A_{\rm wav}\) has support in a fixed
number of neighbouring blocks and a geometric bound independent of the
total volume.
\end{proposition}

\begin{proof}
Equation \eqref{eq:v3v83-wavelet-split} identifies the negative subspace
with the displayed orthogonal direct sum, and determinants turn direct sums
into tensor products.  Differentiating the unitary identification gives
\eqref{eq:v3v83-wavelet-connection}.  Each entry of \(W_U\) is an analytic
function of finitely many link matrices in one block.  Cauchy's estimate on
a smaller link chart gives a geometric Taylor bound, while the fixed block
diameter gives the support assertion.
\end{proof}

The complete Standard Model sum must be taken in
\eqref{eq:v3v83-wavelet-connection} before estimating it.  Its possible
gauge character at zero field cancels by
Proposition~\ref{prop:v3v79-zero-order}; the remaining wavelet term is an
explicit local connection contribution rather than a globally inverted
defect.

\subsection{What the construction does and does not solve}
\label{sec:v3v83-status}

The block Hamiltonian supplies the missing kind of endpoint in V82: it is
covariant, local, uniformly gapped, and in the correct free topological
class.  Its determinant problem is exactly a coarse copy of the original
physical determinant, accompanied by a controlled finite-range spectator
term.  This creates a possible renormalization recursion rather than an
illusory trivial endpoint.

\begin{firewall}[Same phase does not prove a straight-line gap]
\label{fw:v3v83-interpolation-gap}
Equality of stable topological classes is necessary but not sufficient for
the specific interpolation
\((1-t)H_{\rm blk}+tH_W\) to remain gapped.  V83 has not established the
covariant gapped path required by Theorem~\ref{thm:v3v82-reduction}.
\end{firewall}

\begin{firewall}[The coarse determinant has not disappeared]
\label{fw:v3v83-coarse-remains}
The factorization \eqref{eq:v3v83-line-factorization} transfers the physical
chiral determinant to the coarse lattice.  It does not prove its Ward
identity or global triviality.  Any iteration must show that the accumulated
local transgressions converge and must control the terminal global modes.
\end{firewall}

\subsection{V83 conclusion and V84 target}
\label{sec:v3v83-conclusion}

A concrete covariant same-phase reference now exists, with an exact gap and
determinant-line factorization.  It reduces the physical sector by one fixed
blocking scale and isolates all eliminated modes in an explicitly local
wavelet connection.  The unresolved analytic input is a uniform gapped path
from the fine Wilson kernel to this block Hamiltonian.

V84 should formulate a Feshbach criterion for that path.  It should bound
the high-mode Schur complement uniformly, identify a checkable coarse-symbol
inequality which prevents a gap closing, and determine whether the first
blocking step can be certified on the small-link patch.  If the straight
path fails the criterion, the note should insert a spectral-flattening
detour without assuming the desired locality.


\section{V84: Schur certification and existence of a stabilized local gapped path}
\label{sec:v3v84}

V83 constructed a covariant block Hamiltonian in the same stable Wilson
phase as the fine kernel, but did not connect the two through a gapped path.
This note gives a quantitative Schur-complement test for the straight path.
More decisively, it proves that equality of the full free \(K\)-class yields,
after finitely many on-site spectator pairs, a finite-range free gapped path.
Covariantizing its finitely many hoppings and shrinking the small-link patch
then gives the path required by V81--V82.  This avoids assuming that the
straight interpolation happens to be gapped.

\subsection{A sharp low/high Feshbach test}
\label{sec:v3v84-feshbach}

Fix the V83 block projector \(\Pi=BB^*\) and put \(Q=I-\Pi\).  At a fixed
gauge field set
\begin{equation}
 H_0=H_{\rm blk},
 \qquad H_1=H_W,
 \qquad H_t=H_0+t(H_1-H_0).
\label{eq:v3v84-straight-path}
\end{equation}
Relative to \(\operatorname{ran}\Pi\oplus\operatorname{ran}Q\), write
\begin{equation}
 H_t=\begin{pmatrix}A_t&C_t^*\\ C_t&D_t\end{pmatrix}.
\label{eq:v3v84-block-matrix}
\end{equation}
Define the endpoint comparison numbers
\begin{align}
 v_L&=\|\Pi H_W\Pi-BH_cB^*\|,
 &v_H&=\|QH_WQ-m_sSQ\|,
 &\kappa&=\|QH_W\Pi\|,
\label{eq:v3v84-comparison-numbers}
\end{align}
and
\begin{equation}
 \alpha=g_c-v_L,
 \qquad \delta=m_s-v_H.
\label{eq:v3v84-alpha-delta}
\end{equation}

\begin{theorem}[Straight-path Schur criterion]
\label{thm:v3v84-schur-criterion}
If \(\alpha>0\), \(\delta>0\), and
\begin{equation}
 \kappa^2<\alpha\delta,
\label{eq:v3v84-schur-condition}
\end{equation}
then every operator in \eqref{eq:v3v84-straight-path} has the uniform gap
\begin{equation}
 \operatorname{spec}H_t\cap(-\varepsilon_*,\varepsilon_*)=\varnothing,
 \qquad
 \varepsilon_*=\frac{\alpha+\delta-
 \sqrt{(\alpha-\delta)^2+4\kappa^2}}{2}>0.
\label{eq:v3v84-schur-gap}
\end{equation}
\end{theorem}

\begin{proof}
The low block is a perturbation of \(BH_cB^*\) of norm at most \(v_L\),
and the high block is a perturbation of \(m_sSQ\) of norm at most \(v_H\).
Thus both are invertible and, for real \(|z|<\min(\alpha,\delta)\),
\begin{equation}
 \|(A_t-z)^{-1}\|\le(\alpha-|z|)^{-1},
 \qquad
 \|(D_t-z)^{-1}\|\le(\delta-|z|)^{-1}.
\label{eq:v3v84-block-resolvents}
\end{equation}
The off-diagonal block is \(C_t=tQH_W\Pi\), so
\(\|C_t\|\le\kappa\).  The Schur complement of \(D_t-z\) is
\begin{equation}
 S_t(z)=A_t-z-C_t^*(D_t-z)^{-1}C_t.
\label{eq:v3v84-schur-complement}
\end{equation}
It is invertible whenever
\begin{equation}
 \frac{\kappa^2}{\delta-|z|}<\alpha-|z|.
\label{eq:v3v84-neumann-condition}
\end{equation}
The smaller root of equality is exactly \(\varepsilon_*\).  Hence
\(H_t-z\) is invertible for every real \(|z|<\varepsilon_*\), which proves
the claimed spectral gap.
\end{proof}

The criterion is sufficient, not necessary.  Its virtue is that every
quantity in \eqref{eq:v3v84-comparison-numbers} is the supremum of a fixed
finite Bloch matrix in the free theory.  It can therefore certify a concrete
choice of block weights, spectator mass, and coarse Wilson parameters
without a volume extrapolation.

\begin{proposition}[Free-symbol certification extends to small links]
\label{prop:v3v84-symbol-to-links}
Suppose the enlarged-cell free Bloch path has
\begin{equation}
 g_*:=\inf_{t\in[0,1]}inf_{k\in T^4_{\rm red}}
 \min|\operatorname{spec}H_t(I;k)|>0.
\label{eq:v3v84-free-path-gap}
\end{equation}
If every hopping in the path is a sum of at most \(N_R\) parallel transports
of length at most \(R\), with row-sum coefficient bound \(C_R/a\), then
\begin{equation}
 \sup_t\|H_t(U)-H_t(I)\|
 \le \frac{C_RN_R}{a}(e^{R\eta}-1).
\label{eq:v3v84-path-perturbation}
\end{equation}
Consequently the covariantized path is gapped by at least \(g_*/2\) on the
volume-independent patch where the right-hand side is less than \(g_*/2\).
\end{proposition}

\begin{proof}
A product of at most \(R\) links with
\(\|U_e-I\|\le e^\eta-1\) differs from the identity by at most
\(e^{R\eta}-1\), by telescoping the product and using unitarity.  Sum the
finitely many hopping norms in each row to obtain
\eqref{eq:v3v84-path-perturbation}; the Schur test bounds the operator norm
by the same row estimate.  Hausdorff stability of the spectrum of a
self-adjoint matrix then leaves a gap of at least \(g_*/2\).
\end{proof}

\subsection{A local projector-geodesic alternative}
\label{sec:v3v84-geodesic}

There is a second directly checkable criterion which need not control the
individual low and high blocks.

\begin{theorem}[Norm-close local projector geodesic]
\label{thm:v3v84-projector-geodesic}
Let \(P_0,P_1\) be gauge-covariant exponentially local projectors with
\begin{equation}
 \rho:=\|P_1-P_0\|<1.
\label{eq:v3v84-projector-close}
\end{equation}
Put \(S_j=I-2P_j\),
\begin{equation}
 Z=S_1S_0,
 \qquad K=\frac12\operatorname{Log}Z,
\label{eq:v3v84-direct-rotation}
\end{equation}
where \(\operatorname{Log}\) is the principal logarithm.  Then \(K^*=-K\),
is gauge covariant and exponentially local, and
\begin{equation}
 P_s=e^{sK}P_0e^{-sK},
 \qquad0\le s\le1,
 \qquad P_1=e^KP_0e^{-K}.
\label{eq:v3v84-projector-path}
\end{equation}
Thus \(I-2P_s\) is a covariant exponentially local Hamiltonian path with
exact gap one.
\end{theorem}

\begin{proof}
The two-projection canonical-angle decomposition writes the nontrivial
two-dimensional blocks using angles \(0\le\theta<\pi/2\), with
\(\sin\theta\le\rho\).  On each such block, \(S_1S_0\) is rotation by
\(2\theta\), so its spectrum avoids the negative real axis uniformly and
the principal logarithm is anti-Hermitian.  Its half-angle rotation maps
\(\operatorname{ran}P_0\) to \(\operatorname{ran}P_1\), proving
\eqref{eq:v3v84-projector-path}.  The common and orthogonal summands are
fixed.

The product \(Z\) is exponentially local.  Because its spectrum stays a
distance depending only on \(1-\rho\) from the logarithm cut, holomorphic
functional calculus in a slightly weaker weighted Schur algebra makes
\(K\) exponentially local with a volume-independent bound.  Its exponential
has the same property by the convergent power series.  Conjugation covariance
of holomorphic functional calculus proves gauge covariance.
\end{proof}

This theorem supplies the spectral-flattening detour requested at the end of
V83 whenever \eqref{eq:v3v84-projector-close} can be checked.  It does not
infer norm closeness from equality of topological classes.

\subsection{Stable topological equivalence gives a finite-range free path}
\label{sec:v3v84-stable-path}

The absence of either elementary inequality is not a topological
obstruction.  The correct existence statement uses stabilization, which is
already natural for the massive spectator sector.

\begin{theorem}[Stabilized finite-range free interpolation]
\label{thm:v3v84-stable-interpolation}
Let \(H_0(k)\) and \(H_1(k)\) be finite-range periodic Hermitian Bloch
Hamiltonians on \(T^4_{\rm red}\), with gaps at zero.  Assume their negative
bundles define the same class in \(K^0(T^4_{\rm red})\).  Then, after adding
finitely many constant positive and negative spectator bands to the two
sides, there is a continuous family of self-adjoint trigonometric-polynomial
matrices \(\widetilde H_t(k)\) such that
\begin{equation}
 \widetilde H_0=H_0\oplus H_{0,\rm sp},
 \qquad
 \widetilde H_1=H_1\oplus H_{1,\rm sp},
 \qquad
 \inf_{t,k}\min|\operatorname{spec}\widetilde H_t(k)|=g_*>0.
\label{eq:v3v84-polynomial-path}
\end{equation}
The hopping range and coefficient bound are finite numbers independent of
the eventual lattice volume.
\end{theorem}

\begin{proof}
Equality in \(K^0\) means that, after adjoining finite trivial bundles, the
two negative bundles are isomorphic.  Their orthogonal complements become
isomorphic after adjoining the corresponding positive trivial bundles.
An isomorphism of the stabilized decompositions gives a continuous path of
projectors after stabilization; equivalently, the two classifying maps are
homotopic in a sufficiently large finite Grassmannian.  Smooth approximation
on the compact manifold \([0,1]\times T^4_{\rm red}\), relative to its two
boundary components, gives a smooth self-adjoint gapped Hamiltonian path
with the original stabilized endpoints.

Matrix-valued trigonometric polynomials are uniformly dense in smooth
periodic matrices.  Apply relative approximation to the interior of the
path, preserving the finite-range endpoints exactly, and then take the
self-adjoint part.  The unapproximated smooth path has a positive minimum
gap by compactness.  An approximation in operator norm smaller than half
that minimum retains a positive gap by spectral stability.  A single
trigonometric degree works on the compact parameter interval, hence the
hopping range and row-sum bound are finite and uniform in \(t\).
\end{proof}

The V83 identity \eqref{eq:v3v83-relative-projector} proves equality of the
block reference with the coarse Wilson class.  Viewing the fine Wilson
operator with the same enlarged unit cell folds, but does not change, its
full stable class.  Choosing both Wilson masses in the same one-species
component therefore supplies the hypothesis of
Theorem~\ref{thm:v3v84-stable-interpolation}; equality of only the top Chern
number would not have been sufficient.

\subsection{Gauge-covariantization of the stable path}
\label{sec:v3v84-covariantization}

Write each matrix Fourier coefficient of \(\widetilde H_t\) as a hopping
between two sites of the enlarged cell.  For each nonzero hopping choose a
fixed fine-lattice path between its endpoints and multiply the coefficient
by parallel transport along that path; pair every path with its reverse to
preserve self-adjointness.  At \(t=0\), choose the path decomposition already
used in \(H_{\rm blk}(U)\), and at \(t=1\) use the nearest-neighbour Wilson
paths.  Relative approximation permits these endpoint choices to remain
fixed.

\begin{theorem}[Covariant gapped path on a quantitative small-link patch]
\label{thm:v3v84-covariant-path}
The covariantization above produces finite-range self-adjoint Hamiltonians
\(\widetilde H_t(U)\) satisfying
\begin{equation}
 \widetilde H_t(U^g)=G(g)\widetilde H_t(U)G(g)^{-1}.
\label{eq:v3v84-path-covariance}
\end{equation}
There is \(\eta_*>0\), independent of lattice volume, such that
\begin{equation}
 \|\log U_e\|<\eta_*quad\hbox{for every link}
 \quad\Longrightarrow\quad
 \operatorname{spec}\widetilde H_t(U)
 \cap(-g_*/2,g_*/2)=\varnothing
\label{eq:v3v84-interacting-path-gap}
\end{equation}
for all \(t\).  The path and all of its link derivatives obey the finite-range
joint cluster bounds required in V82.
\end{theorem}

\begin{proof}
Parallel transport transforms at its endpoints, which proves
\eqref{eq:v3v84-path-covariance}; reverse paths make the operator
self-adjoint.  The polynomial path has a finite maximal hopping length and
finite uniform coefficient sum.  Proposition
\ref{prop:v3v84-symbol-to-links} therefore applies.  Choose \(\eta_*\) so
that its right-hand side is below \(g_*/2\), proving
\eqref{eq:v3v84-interacting-path-gap}.  Repeated link differentiation of a
finite path product is a finite sum supported on that path.  The fixed
range and compact \(t\)-interval give joint cluster bounds and a common
analytic radius.
\end{proof}

\begin{corollary}[The V81 transport hypothesis is realized locally]
\label{cor:v3v84-v81-realized}
After adjoining the finite spectator bands of
Theorem~\ref{thm:v3v84-stable-interpolation}, the fine Wilson projector and
the V83 block projector are connected on the quantitative small-link patch
by a covariant uniformly gapped local path.  Hence the equivariant
transgression of V81 is exponentially local with a common convergent cluster
expansion and transports their determinant connections.
\end{corollary}

\begin{proof}
Apply Theorem~\ref{thm:v3v84-covariant-path}, then
Lemma~\ref{lem:v3v82-riesz-locality},
Theorem~\ref{thm:v3v82-local-intertwiner}, and
Theorem~\ref{thm:v3v81-local-transgression}.  Constant spectator bands have
zero field-space curvature.  Their field-independent gauge character is a
representation trace, which cancels in the complete Standard Model sum by
the same arithmetic as Proposition~\ref{prop:v3v79-zero-order}.
\end{proof}

\subsection{Scope and the remaining multiscale problem}
\label{sec:v3v84-scope}

The local path is now an existence theorem rather than an assumed straight
line.  Its range and small-link radius are finite but not optimized; the
Feshbach and projector-distance tests remain useful for obtaining explicit
shorter paths.  More importantly, the endpoint determinant is the coarse
physical determinant from V83.  Repeating the construction produces a
sequence of local transgressions, and their infinite or continuum-scale sum
has not yet been controlled.

\begin{firewall}[Stable spectators are part of the regulator]
\label{fw:v3v84-spectators}
Theorem~\ref{thm:v3v84-stable-interpolation} may add finitely many on-site
positive and negative bands.  They are locally trivial and cancel in the
complete Standard Model character, but removing them from a determinant
ratio must be tracked together with the Pauli--Villars normalization.  The
theorem does not prove reflection positivity of that removal.
\end{firewall}

\begin{firewall}[Small-link patch only]
\label{fw:v3v84-small-link-only}
The interacting gap follows by a norm perturbation around the free path.
It neither covers all admissible topological sectors nor decides large-gauge
holonomy.  Those global questions cannot be inferred from stable free
\(K\)-theory alone.
\end{firewall}

\subsection{V84 conclusion and V85 audit target}
\label{sec:v3v84-conclusion}

The straight block interpolation now has an explicit Schur-complement gap
test.  Even when that conservative test is unavailable, equality of the
full stable Wilson class yields a finite-range free path after harmless
spectator stabilization; covariantization preserves a uniform gap on a
quantitative small-link patch.  This discharges the local path hypothesis in
V81 and converts the fine determinant problem into a coarse determinant plus
a controlled local transgression.

V85 is the next compulsory NS/YM audit.  It should then formulate the
iterated blocking identity, identify the scaling dimension of each
transgression remainder, and prove a summability theorem under an explicit
irrelevant-order gain.  If no such gain follows from the present construction,
the first marginal term must be isolated rather than hidden in a formal
renormalization recursion.


\section{V85: audited multiscale determinant transport and the marginal obstruction}
\label{sec:v3v85}

This compulsory audit is performed immediately after V84's existence theorem
for a stabilized covariant gapped path.  The purpose of the present note is
to decide what that one-step result gives after repeated blocking.  The
answer is exact at every finite depth, but an infinite-depth limit requires
both transport of the analytic hypotheses and a positive irrelevant-order
gain.  A scale-by-scale local bound of order one is only marginal and
accumulates logarithmically.

\subsection{Compulsory NS/YM audit}
\label{sec:v3v85-audit}

The current Navier--Stokes file was read passively:
\begin{equation}
 \begin{split}
 &\texttt{Renewal\_NS\_Stability\_Research\_Notes\_V31.tex},\\
 &887537\text{ bytes},\qquad
 \operatorname{SHA256}=
 \texttt{F1121B62238AD5491BB7CF03635EA9934942EDF573ED8FAAA9EE79E1D38A6696}.
 \end{split}
\label{eq:v3v85-ns-fingerprint}
\end{equation}
It is unchanged from the V80 audit.  Its relevant lesson remains sharp: a
marked estimate can be exact at every scale while removal of the mark costs
one critical first moment; a weaker source norm merely restores that missing
weight elsewhere.

The newest named Yang--Mills file was also read passively:
\begin{equation}
 \begin{split}
 &\texttt{Renewal\_Yang\_Mills\_Mass\_Gap\_Research\_Notes\_V5.tex},\\
 &98637\text{ bytes},\qquad
 \operatorname{SHA256}=
 \texttt{8163303DD7C0167A0E3DCA82A5C7E28ABCA77123F406926A768E428D400C4901}.
 \end{split}
\label{eq:v3v85-ym-fingerprint}
\end{equation}
At the time of inspection it was byte-identical to
\texttt{Renewal\_Yang\_Mills\_Mass\_Gap\_Research\_Notes\_V4.tex}, including
the same byte count and digest, and its last mathematical section was the V4
ledger.  This shared-folder state was not altered.  The useful V4 lesson is
that acyclicity of a scalar ascent is not a packet-transport theorem: the
local breaker, admissible mark, and unique payer must all survive the change
of row.  Here the analogue is that a one-step stable homotopy does not by
itself transport a common link patch, localization exponent, spectator
normalization, and determinant connection through arbitrarily many scales.

V80--V84 already prove the one-moment obstruction, exact equivariant
transgression, local spectral flow, construction of a covariant same-phase
block reference, and existence of a stabilized finite-range gapped path on
a quantitative small-link patch.  The following argument uses those results
without repeating their proofs.

\subsection{Exact finite-depth blocking identity}
\label{sec:v3v85-recursion}

Let the lattice spacing at level \(n\) be
\begin{equation}
 a_n=q^{-n}a_0,
 \qquad q>1,
\label{eq:v3v85-scales}
\end{equation}
so one V83 block maps level \(n\) to level \(n-1\).  Denote the gauge-field
blocking map by \(r_n:\mathcal U_n\to\mathcal U_{n-1}\).  Let
\(\mathcal A_n\) be a determinant connection for the complete stabilized
Standard Model projector at level \(n\).  Combining the V83 wavelet
connection with the integrated V81 transgression gives a local one-form
\(\Theta_n\) on \(\mathcal U_n\).  After a choice of phase chart there may
also be an exact scalar change \(d\chi_n\), and the one-step identity is
\begin{equation}
 \mathcal A_n=r_n^*\mathcal A_{n-1}+\Theta_n+d\chi_n.
\label{eq:v3v85-one-step}
\end{equation}
All representation sums in \(\Theta_n\) are taken before local norms.

For \(m<n\), put
\begin{equation}
 r_{n:m}=r_{m+1}\circ\cdots\circ r_n,
 \qquad r_{n:n}=\operatorname{id}.
\label{eq:v3v85-composed-block}
\end{equation}

\begin{theorem}[Exact telescoping at finite depth]
\label{thm:v3v85-finite-telescope}
For every finite \(N\),
\begin{equation}
 \mathcal A_N
 =r_{N:0}^*\mathcal A_0
 +\sum_{j=1}^N r_{N:j}^*\Theta_j
 +d\left(\sum_{j=1}^Nr_{N:j}^*\chi_j\right).
\label{eq:v3v85-finite-telescope}
\end{equation}
The corresponding curvatures and gauge contractions telescope with the same
pullbacks.  No transgression term is counted twice.
\end{theorem}

\begin{proof}
Apply \eqref{eq:v3v85-one-step} at level \(N\), substitute the same identity
for \(\mathcal A_{N-1}\), and continue.  Pullback commutes with \(d\), so
all exact phase changes combine into the last term.  Differentiation gives
the curvature identity.  Contraction with a gauge vector commutes with
pullback because every \(r_n\) is equivariant.  Each \(\Theta_j\) belongs to
the unique shell eliminated at step \(j\), which proves the one-use
statement.
\end{proof}

This is the determinant analogue of an exact shell ledger.  It is an
identity, not yet an infinite-series estimate.

\subsection{The gauge patch is not automatically transported}
\label{sec:v3v85-patch}

The small-link ball used in V78 and V84 is defined by individual principal
link logarithms.  A coarse link is a product of \(q\) fine links along a
coordinate segment.

\begin{proposition}[Fixed link balls expand under blocking]
\label{prop:v3v85-link-ball-expansion}
For sufficiently small \(\eta\), one blocking step obeys
\begin{equation}
 \max_{e_c}\|\log r_n(U)(e_c)\|
 \le C_q\max_{e_f}\|\log U(e_f)\|,
\label{eq:v3v85-link-expansion}
\end{equation}
with \(C_q=q+O(\eta)\).  Iterating only this estimate gives
\(C_q^N\eta\).  Hence a fixed-radius fine-link ball is not mapped into the
same fixed-radius ball at arbitrarily many blocking levels.
\end{proposition}

\begin{proof}
Write the links on the path as \(e^{A_1},\ldots,e^{A_q}\).  The
Baker--Campbell--Hausdorff formula on a fixed neighbourhood of zero gives
\begin{equation}
 \left\|\log(e^{A_1}\cdots e^{A_q})\right\|
 \le\sum_{j=1}^q\|A_j\|+C\sum_{i<j}\|A_i\|\|A_j\|,
\end{equation}
which is \((q+O(\eta))\eta\).  Repetition proves the final statement.
\end{proof}

For continuum-scaled fields this is the correct dimensional behaviour, not
a failure.  If
\begin{equation}
 U_n(e)=\mathcal P\exp\int_e A,
 \qquad \|A\|_{L^\infty}\le M_A,
\label{eq:v3v85-continuum-links}
\end{equation}
then a level-\(n\) link has logarithm \(O(a_nM_A)\), and its block product
has size \(O(a_{n-1}M_A)\).  Thus a multiscale statement must be made on a
bounded continuum-field class and stopped at a fixed physical coarse scale.
It cannot be claimed uniformly on the full radius-\(\eta\) link ball at
every level.

\subsection{A norm in which scale summation is meaningful}
\label{sec:v3v85-norm}

Let \(\mathfrak L_\gamma\) be a physical-coordinate local one-form norm:
after pulling every lattice functional to the same smooth background field,
its multilinear kernels are weighted by the length of a minimal tree in
physical distance, with exponential factor \(e^{\gamma\tau_{\rm phys}}\),
and summed over all but one marked cell.  Exact details of the spin and
representation matrix norm are immaterial provided that
\begin{equation}
 \|FG\|_{\mathfrak L_{\gamma'}}
 \le C_{\gamma,\gamma'}
 \|F\|_{\mathfrak L_\gamma}\|G\|_{\mathfrak L_\gamma},
 \qquad0<\gamma'<\gamma,
\label{eq:v3v85-local-algebra}
\end{equation}
and gauge contraction and exterior differentiation are continuous on each
fixed Taylor degree.  The V81--V84 cluster estimates provide this structure
at one scale.

Let \(\mathscr P_{N:j}^*\Theta_j\) denote the level-\(j\) one-form pulled to
the common level-\(N\) continuum background.  This notation includes the
field blocking and the canonical rescaling of the variation \(dA\).

\begin{theorem}[Irrelevant-gain summability]
\label{thm:v3v85-irrelevant-sum}
Assume that one construction of the stabilized paths can be chosen at all
levels with common analytic and locality constants, and that for some
\(\omega>0\)
\begin{equation}
 \|\mathscr P_{N:j}^*\Theta_j^{\rm irr}\|_{\mathfrak L_\gamma}
 \le Cq^{-\omega j}
\label{eq:v3v85-irrelevant-gain}
\end{equation}
for every \(N\ge j\), with \(C\) independent of \(N,j\).  Then
\begin{equation}
 \Theta_\infty^{\rm irr}
 =\sum_{j=1}^\infty\mathscr P_{\infty:j}^*\Theta_j^{\rm irr}
\label{eq:v3v85-infinite-transgression}
\end{equation}
converges absolutely in every weaker norm
\(\mathfrak L_{\gamma'}\), \(0<\gamma'<\gamma\), and
\begin{equation}
 \|\Theta_\infty^{\rm irr}\|_{\mathfrak L_{\gamma'}}
 \le \frac{C'}{1-q^{-\omega}}.
\label{eq:v3v85-geometric-sum}
\end{equation}
Exterior differentiation, infinitesimal gauge contraction, and the
convergent small-field Taylor sum may be interchanged with the scale sum.
Thus the finite identities \eqref{eq:v3v85-finite-telescope} pass to the
limit in the irrelevant sector.
\end{theorem}

\begin{proof}
The assumed bound is a geometric majorant.  Completeness of the weaker local
Banach space and the Weierstrass test give
\eqref{eq:v3v85-infinite-transgression}--%
\eqref{eq:v3v85-geometric-sum}.  Continuity of the listed operations and
uniform convergence justify their interchange with the sum.  A common
analytic radius supplies a second geometric majorant in Taylor degree, so
the scale and field expansions may be interchanged by absolute convergence.
\end{proof}

\subsection{The marginal obstruction is sharp}
\label{sec:v3v85-marginal}

\begin{proposition}[Order-one shell terms need not converge]
\label{prop:v3v85-marginal-obstruction}
The estimate
\begin{equation}
 \|\mathscr P_{N:j}^*\Theta_j\|_{\mathfrak L_\gamma}\le C
\label{eq:v3v85-marginal-bound}
\end{equation}
does not imply convergence as \(N\to\infty\).  More sharply, if a continuous
local test variation \(V\) and a number \(c>0\) satisfy
\begin{equation}
 \operatorname{Re}\bigl(\mathscr P_{N:j}^*\Theta_j(V)\bigr)\ge c
\label{eq:v3v85-coherent-marginal}
\end{equation}
for \(1\le j\le N\), then the norm of the partial sum is at least
\(cN/\|V\|\).
\end{proposition}

\begin{proof}
Evaluate the partial sum on \(V\).  The real parts add and give at least
\(cN\).  Duality bounds the absolute value of this evaluation by the norm of
the sum times \(\|V\|\).  This proves the lower bound.  Alternating examples
also show that bounded summands alone provide no Cauchy criterion.
\end{proof}

This is the exact scale analogue of the NS first-moment warning.  A missing
factor cannot be recovered by renaming a uniform shell estimate as a
summable one.

Let \(\mathfrak M\) denote projection onto the finite-dimensional space of
local terms with nonnegative renormalization-group scaling degree allowed by
gauge covariance, hypercubic symmetry, and the chosen discrete
transformations.  Define
\begin{equation}
 \Theta_j^{\rm marg}=\mathfrak M\Theta_j,
 \qquad
 \Theta_j^{\rm irr}=(I-\mathfrak M)\Theta_j.
\label{eq:v3v85-marginal-split}
\end{equation}
The marginal coefficients require normalization conditions and running
counterterms; they cannot be included in the raw geometric series of
Theorem~\ref{thm:v3v85-irrelevant-sum}.

\begin{proposition}[What anomaly cancellation removes]
\label{prop:v3v85-anomaly-marginal}
In the ghost-number-one gauge variation of the complete one-generation
Standard Model connection, the cohomologically nontrivial marginal component
is the continuum consistent gauge anomaly and its mixed gravitational
partner.  Their coefficients vanish representation by representation sum as
proved in V66 and V79.  Therefore any remaining marginal gauge defect is a
local BRST coboundary and may be fixed by a finite list of normalization
counterterms.  This statement does not imply
\(\Theta_j^{\rm marg}=0\) as a connection one-form.
\end{proposition}

\begin{proof}
The local ghost-number-one cohomology classification used in
Theorem~\ref{thm:v3v79-leading-cancellation} lists precisely the descent
classes of the cubic gauge invariants and the mixed gravitational trace.
Proposition~\ref{prop:v3v66-coefficient-cancellation} sets every Standard
Model coefficient to zero.  The remaining closed marginal terms therefore
represent the zero local cohomology class and are BRST variations of local
functionals.  Gauge-invariant or exact connection terms can nevertheless
have marginal scaling, so vanishing of the anomalous quotient class cannot
set the whole one-form to zero.
\end{proof}

\subsection{A conditional renormalized continuum connection}
\label{sec:v3v85-renormalized}

\begin{theorem}[Multiscale completion criterion]
\label{thm:v3v85-multiscale-criterion}
Suppose:
\begin{enumerate}
 \item the covariant stable paths of V84 are chosen scale-covariantly with
 common locality, gap, and analytic constants on the continuum-field class
 \eqref{eq:v3v85-continuum-links};
 \item a fixed finite family of local counterterms removes
 \(\Theta_j^{\rm marg}\) according to prescribed normalization conditions,
 with the nontrivial anomaly coefficients zero as in
 Proposition~\ref{prop:v3v85-anomaly-marginal};
 \item the normalized remainder satisfies
 \eqref{eq:v3v85-irrelevant-gain} for some \(\omega>0\); and
 \item the terminal coarse determinant connection has a controlled limit on
 the fixed physical coarse lattice.
\end{enumerate}
Then the renormalized version of
\eqref{eq:v3v85-finite-telescope} converges to an exponentially local
infinitesimally gauge-invariant Standard Model determinant connection on the
bounded continuum-field patch.
\end{theorem}

\begin{proof}
Subtract the normalized marginal counterterms in the exact finite-depth
identity.  Proposition~\ref{prop:v3v85-anomaly-marginal} ensures that no
nontrivial local gauge cohomology class obstructs this subtraction.
Theorem~\ref{thm:v3v85-irrelevant-sum} gives convergence of the remainder
and permits passage of curvature and Ward identities to the limit.  The
terminal term converges by item 4.  The common locality and analytic bounds
in item 1 make the resulting connection exponentially local with a nonzero
small-field radius.
\end{proof}

\begin{firewall}[The irrelevant gain has not been proved]
\label{fw:v3v85-gain-open}
V84 gives finite-range paths at each fixed scale by stable approximation,
but it does not make those paths scale-covariant or prove any positive
\(q^{-\omega j}\) gain.  Thus item 3 of
Theorem~\ref{thm:v3v85-multiscale-criterion} remains the principal local
continuum bottleneck.
\end{firewall}

\begin{firewall}[Infinitesimal patch invariance is not the full continuum]
\label{fw:v3v85-global-rp}
Even after the criterion is discharged, large-gauge determinant holonomy,
the global Standard Model quotient, removal of stabilized and
Pauli--Villars bands with reflection positivity, and coupling to the
Einstein contour remain separate gates.
\end{firewall}

\subsection{V85 conclusion and V86 target}
\label{sec:v3v85-conclusion}

Repeated blocking now has an exact one-use transgression ledger.  The audit
prevents a false inference: finite-depth locality is not infinite-depth
transport, and an order-one shell bound is marginal rather than summable.
The nontrivial marginal gauge anomaly cancels in the Standard Model, but the
remaining marginal normalization and the positive irrelevant gain must be
established explicitly.

V86 should choose the block path scale-covariantly and derive the first terms
of its Symanzik expansion.  It must project the complete Standard Model
transgression onto the allowed marginal basis before taking norms, fix those
coefficients by local normalization, and identify the first remainder power
\(\omega\).  If the stable-approximation path of V84 supplies no uniform
power, the construction must be replaced by a functional-calculus path tied
directly to the Wilson symbols.


\section{V86: a scale-covariant path template and the first irrelevant power}
\label{sec:v3v86}

V85 identified two missing hypotheses in the multiscale determinant
transport: the V84 stable paths had not been chosen with common constants,
and no positive power of the ultraviolet lattice spacing had been extracted
after marginal normalization.  Both problems can be addressed at the local
level.  With fixed blocking ratio and fixed dimensionless Wilson parameters,
one stabilized free homotopy is a dimensionless template for every scale.
Uniform exponential locality then gives an ordinary moment expansion on
smooth continuum fields.  After subtracting the complete nondecaying local
jet, the first uncontrolled remainder is \(O(a)\), so the exponent in V85
is \(\omega=1\).  This conclusion concerns the local small-field sector;
global holonomy and reflection positivity remain outside it.

\subsection{One dimensionless path for all scales}
\label{sec:v3v86-template}

Fix the block ratio \(q\), block weights, path system, Wilson parameter, and
dimensionless mass \(m_0=aM\) in a closed subinterval of
\((0,2r)\).  Scale the Hermitian kernels by their fine lattice spacing:
\begin{equation}
 \widehat H_{a,W}=aH_{a,W},
 \qquad
 \widehat H_{a,\rm blk}=aH_{a,\rm blk}.
\label{eq:v3v86-scaled-H}
\end{equation}
The coarse physical block inside \(\widehat H_{a,\rm blk}\) carries the
fixed factor \(a/(qa)=q^{-1}\), and choose the spectator mass in units
\(m_s/a\).  At the free field, after writing momenta in the reduced
dimensionless Brillouin variable \(\widehat k=ak\), both endpoint Bloch
matrices are independent of \(a\).

\begin{theorem}[Scale-covariant stabilized path]
\label{thm:v3v86-template}
After one fixed finite spectator stabilization, there is a self-adjoint
trigonometric-polynomial path
\begin{equation}
 \widehat H_t^{\rm tpl}(\widehat k),
 \qquad (t,\widehat k)\in[0,1]\times T^4_{\rm red},
\label{eq:v3v86-template-path}
\end{equation}
whose endpoints are the dimensionless block and fine Wilson matrices and
whose gap
\begin{equation}
 \widehat g_*=inf_{t,\widehat k}
 \min|\operatorname{spec}\widehat H_t^{\rm tpl}(\widehat k)|
\label{eq:v3v86-template-gap}
\end{equation}
is positive.  The same Fourier coefficients, interpreted in lattice units
at every \(a_n=q^{-n}a_0\), give paths with a common hopping radius, common
row-sum bound, and common dimensionless gap \(\widehat g_*\).

After covariantizing those hoppings by the same path shapes, there is a
single \(\eta_*>0\) such that every scale path is gapped whenever its link
logarithms are bounded by \(\eta_*\).  Its projectors, spectral-flow
generators, and V81 transgressions have common dimensionless exponential
locality and analytic constants.
\end{theorem}

\begin{proof}
The endpoint matrices and their full stable \(K\)-classes are independent
of \(a\) in \(\widehat k\)-coordinates.  Apply
Theorem~\ref{thm:v3v84-stable-interpolation} once, rather than separately at
each scale, to obtain \eqref{eq:v3v86-template-path}.  Its finite Fourier
support and compact parameter interval supply fixed hopping and coefficient
bounds.  Reinterpreting an integer displacement at spacing \(a_n\) changes
its physical length but not the dimensionless matrix or gap.

Covariantization replaces each fixed Fourier hopping by a parallel transport
along the corresponding integer path.  The perturbation estimate
\eqref{eq:v3v84-path-perturbation}, applied to the scaled Hamiltonian, has
the same constants at every level.  Choose \(\eta_*\) so its right-hand
side is below \(\widehat g_*/2\).  The Riesz, spectral-flow, and
transgression estimates in V81--V82 depend only on that gap, the fixed
hopping bounds, and a smaller analytic link ball.  Their dimensionless
constants are therefore common to all scales.
\end{proof}

For a continuum connection bounded as in
\eqref{eq:v3v85-continuum-links}, choose the terminal physical spacing
\(a_0\) so that \(C a_0\|A\|_{L^\infty}<\eta_*\).  Then every finer level
also lies in the common path patch.  This is the packet-transport statement
which was absent in V85.

\subsection{The exponential-kernel moment lemma}
\label{sec:v3v86-moments}

The mechanism producing a positive power of \(a\) is elementary and must be
kept explicit.  First consider a translation-covariant linear kernel on
\(\mathbb R^4\) sampled at spacing \(a\):
\begin{equation}
 (T_af)(x)=\sum_{r\in\mathbb Z^4}K(r)f(x+ar).
\label{eq:v3v86-convolution}
\end{equation}

\begin{lemma}[Moment expansion with a uniform remainder]
\label{lem:v3v86-moment-expansion}
If
\begin{equation}
 \sum_r e^{\gamma|r|}\|K(r)\|\le C_K,
\label{eq:v3v86-exp-moment}
\end{equation}
then, for every integer \(p\ge1\) and
\(f\in C_b^p(\mathbb R^4)\),
\begin{equation}
 T_af(x)=
 \sum_{|\alpha|<p}\frac{a^{|\alpha|}}{\alpha!}
 \left(\sum_r r^\alpha K(r)\right)\partial^\alpha f(x)
 +R_{p,a}f(x),
\label{eq:v3v86-moment-series}
\end{equation}
with
\begin{equation}
 \|R_{p,a}f\|_\infty
 \le C_{p,\gamma}C_Ka^p
 \max_{|\alpha|=p}\|\partial^\alpha f\|_\infty.
\label{eq:v3v86-moment-remainder}
\end{equation}
The same statement holds for a multilinear connected kernel when \(|r|\)
is replaced by the length of a tree joining all insertion points.
\end{lemma}

\begin{proof}
Taylor's formula with integral remainder gives
\begin{equation}
 f(x+ar)=\sum_{|\alpha|<p}
 \frac{a^{|\alpha|}r^\alpha}{\alpha!}\partial^\alpha f(x)
 +O\left(a^p|r|^p
 \max_{|\alpha|=p}\|\partial^\alpha f\|_\infty\right).
\end{equation}
All polynomial moments of \(K\) are finite under
\eqref{eq:v3v86-exp-moment}; in particular
\(\sum_r|r|^p\|K(r)\|\le C_{p,\gamma}C_K\).  Summation proves
\eqref{eq:v3v86-moment-series}--%
\eqref{eq:v3v86-moment-remainder}.  For a multilinear kernel, expand every
external field about one marked point.  The product remainder is bounded by
a polynomial in the relative displacements, and the exponential tree weight
summably controls all such moments.
\end{proof}

\begin{corollary}[Vanishing jets give positive lattice powers]
\label{cor:v3v86-vanishing-moments}
If every moment multiplying a local continuum monomial of scaling degree
less than \(p\) vanishes after the complete representation sum and chosen
counterterm subtraction, then the corresponding sampled local functional is
\(O(a^p)\) in the physical local norm on bounded \(C^p\) fields.
\end{corollary}

\subsection{Covariant Symanzik expansion of the transgression}
\label{sec:v3v86-symanzik}

Let \(\Theta_a\) be the sum of the wavelet connection and the V81
transgression built from the template in
Theorem~\ref{thm:v3v86-template}.  Pull it back to a smooth continuum gauge
field using path-ordered link holonomies, and evaluate it on a smooth
variation \(\dot A\).  Analyticity on the common link ball gives an
absolutely convergent expansion in link logarithms.  Exponential projector
locality gives connected kernels satisfying the tree estimate in
Lemma~\ref{lem:v3v86-moment-expansion}.

\begin{theorem}[Uniform local Symanzik jet]
\label{thm:v3v86-symanzik}
For every desired scaling order \(p\), on a bounded \(C^{N(p)}\) continuum-field set for a finite derivative order \(N(p)\), there are
local continuum one-forms \(\Theta_{[\ell]}\), independent of \(a\), such
that
\begin{equation}
 \Theta_a(A;\dot A)
 =\sum_{\ell=\ell_{\min}}^{p-1}a^\ell
 \Theta_{[\ell]}(A;\dot A)+\mathcal R_{p,a}(A;\dot A),
\label{eq:v3v86-symanzik-expansion}
\end{equation}
after the conventional physical-volume and field rescalings, and
\begin{equation}
 \|\mathcal R_{p,a}\|_{\mathfrak L_{\gamma'}}
 \le C_pa^p,
 \qquad0<\gamma'<\gamma.
\label{eq:v3v86-symanzik-remainder}
\end{equation}
Only finitely many terms have \(\ell\le0\).  They are local polynomials
allowed by gauge covariance of the full equivariant identity and by the
lattice symmetries actually imposed on the selected regulator.
\end{theorem}

\begin{proof}
Expand each path-ordered link product in powers of \(aA\).  At fixed field
degree, apply the multilinear part of
Lemma~\ref{lem:v3v86-moment-expansion} to expand all sampled fields about a
marked physical point.  The scale-covariant template makes every
dimensionless kernel moment independent of \(a\).  Dimensional rescaling of
the lattice sum, the fields, and the single variation groups the resulting
terms into integer powers \(a^\ell\).

The common analytic majorant permits summation over field degree after the
finite spatial Taylor expansion.  The tree moment bound gives
\eqref{eq:v3v86-symanzik-remainder}.  At a fixed engineering-dimension
cutoff, only finitely many monomials in \(A,\dot A\), and their derivatives
exist.  Exact covariance of the unexpanded equivariant transgression forces
their gauge variations to assemble into the local BRST descent.  Projection
onto any finite lattice-symmetry group imposed in the construction removes
forbidden monomials without changing the bounds.
\end{proof}

The theorem is stronger than a formal heat-kernel statement: its remainder
is uniform in the physical local norm and in the ultraviolet level.

\subsection{Marginal projection before norming}
\label{sec:v3v86-marginal}

Let \(\mathfrak M_{\le0}\Theta_a\) be the finite jet in
\eqref{eq:v3v86-symanzik-expansion} with \(\ell\le0\), expressed in a fixed
basis of local continuum operators.  It includes all relevant and marginal
connection terms, not only the anomalous quotient class.  Define the
normalized shell connection by
\begin{equation}
 \Theta_a^{\rm ren}
 =\Theta_a-\mathfrak M_{\le0}\Theta_a-dC_a,
\label{eq:v3v86-renormalized-shell}
\end{equation}
where the finite local scalar \(C_a\) imposes the selected BRST-exact
normalization conditions.

\begin{proposition}[Complete Standard Model marginal Ward projection]
\label{prop:v3v86-marginal-projection}
Take the complete one-generation representation sum before applying
\(\mathfrak M_{\le0}\).  In the ghost-number-one variation of its
\(\ell=0\) jet, every nontrivial local BRST class has zero coefficient.
The remaining finite jet is BRST exact or gauge invariant and can be assigned
by local normalization conditions without using a lattice Poincar\'e inverse.
\end{proposition}

\begin{proof}
The equivariant transgression identity
\eqref{eq:v3v81-equivariant-exact} holds before expansion.  Its marginal
ghost-number-one cohomology therefore lies in the standard four-dimensional
consistent-anomaly basis.  The independent \(SU(3)^3\), \(SU(3)^2Y\),
\(SU(2)^2Y\), \(Y^3\), and mixed gravitational coefficients vanish by
Proposition~\ref{prop:v3v66-coefficient-cancellation}; the perturbative
\(SU(2)^3\) class is absent.  The quotient by BRST-exact local terms is thus
zero.  Because the finite jet was assembled before taking norms, its exact
part is removed by the finite counterterm \(C_a\), while gauge-invariant
terms are fixed by the prescribed normalization basis.
\end{proof}

\begin{theorem}[First positive remainder power]
\label{thm:v3v86-omega-one}
With the complete finite jet \(\ell\le0\) removed as in
\eqref{eq:v3v86-renormalized-shell}, the scale-covariant transgression obeys
\begin{equation}
 \|\Theta_{a_n}^{\rm ren}\|_{\mathfrak L_{\gamma'}}
 \le C a_n=C a_0q^{-n}.
\label{eq:v3v86-a-gain}
\end{equation}
Thus the irrelevant-gain hypothesis
\eqref{eq:v3v85-irrelevant-gain} holds locally with \(\omega=1\).  If an
additional exact symmetry eliminates the \(\ell=1\) jet, the same proof
improves the exponent to \(\omega=2\), but that improvement is not assumed.
\end{theorem}

\begin{proof}
Choose the finite Taylor order \(N(1)\) which expands through scaling degree
zero and leaves scaling order \(p=1\), then subtract every term of
nonpositive scaling degree in Theorem~\ref{thm:v3v86-symanzik}.  The first retained integer power is
\(a^1\), and \eqref{eq:v3v86-symanzik-remainder} gives the uniform bound.
Substitute \(a_n=a_0q^{-n}\).  This is precisely
\eqref{eq:v3v85-irrelevant-gain} with \(\omega=1\).
\end{proof}

\begin{corollary}[Local ultraviolet transgression converges]
\label{cor:v3v86-local-uv-limit}
On the bounded smooth small-field class, the normalized ultraviolet shell
sum
\begin{equation}
 \sum_{n=1}^\infty\Theta_{a_n}^{\rm ren}
\label{eq:v3v86-uv-sum}
\end{equation}
converges absolutely to an exponentially local connection one-form, and its
curvature and infinitesimal Ward identities are the limits of the finite
blocking identities.
\end{corollary}

\begin{proof}
Apply Theorem~\ref{thm:v3v85-irrelevant-sum} with \(\omega=1\), using the
common constants of Theorem~\ref{thm:v3v86-template}.
\end{proof}

\subsection{What remains after the local ultraviolet limit}
\label{sec:v3v86-scope}

V86 discharges the local scale-summability item in
Theorem~\ref{thm:v3v85-multiscale-criterion}: the arbitrary per-scale
stable paths have been replaced by one scale-covariant template, and the
full nondecaying Symanzik jet is removed before the first norm.  The
remainder has the rigorous conservative gain \(O(a)\).

\begin{firewall}[Normalization data are genuine input]
\label{fw:v3v86-normalization}
The coefficients in \(\mathfrak M_{\le0}\Theta_a\) are not asserted to
vanish.  Their anomalous cohomology class cancels, but gauge-invariant and
BRST-exact relevant or marginal terms require a finite set of physical
normalization conditions.  V86 proves summability after those conditions
are imposed.
\end{firewall}

\begin{firewall}[Smooth bounded fields, not every lattice configuration]
\label{fw:v3v86-smooth-only}
The moment remainder uses bounded continuum derivatives and the common
small-link patch.  It does not prove that the full dynamical lattice measure
concentrates on this class, nor does it control rough fields in other
topological components.
\end{firewall}

\begin{firewall}[Local Ward convergence is not global determinant
triviality]
\label{fw:v3v86-global-open}
Corollary~\ref{cor:v3v86-local-uv-limit} proves an infinitesimal local
connection limit.  It does not determine determinant-line periods under
large gauge loops, the precise global Standard Model quotient, or the
reflection positivity of the stabilized Pauli--Villars ratio.
\end{firewall}

\subsection{V86 conclusion and V87 target}
\label{sec:v3v86-conclusion}

The local ultraviolet recursion is no longer merely conditional on an
unknown irrelevant exponent.  A single dimensionless path template has
common gap and locality constants at every scale; exponential kernel moments
give a uniform Symanzik expansion; and subtraction of the complete
nondecaying jet leaves an \(O(a_n)\) shell.  The normalized local
transgressions therefore sum geometrically.

V87 should address the terminal coarse determinant rather than repeat the
ultraviolet expansion.  It should construct its determinant line over the
based gauge-orbit patch, separate constant gauge stabilizers from based
transformations, and compute the holonomy formula for a loop by a
five-dimensional mapping-torus index.  The global Standard Model quotient
must enter through the actual allowed loop lattice, not through Lie-algebra
Taylor coefficients.


\section{V87: the based orbit space and vanishing of the spin--\(G_{\rm SM}\) global anomaly group}
\label{sec:v3v87}

V86 constructs the normalized local ultraviolet determinant connection on a
smooth small-field patch.  The terminal question is whether a locally flat
gauge anomaly line can retain nontrivial holonomy around a large gauge loop.
V66 already counted the ordinary weak doublets, V67 already established the
abstract spectral-flow/mapping-torus identity for a complementing family,
and the monograph already reconstructed the global group and its finite
central descent.  This note does not repeat those results.  It identifies
the correct based lattice orbit space, writes its mapping-torus phase, and
computes the relevant five-dimensional spin bordism group for the actual
quotient.  That group vanishes.

\subsection{The global group and the based finite-lattice action}
\label{sec:v3v87-based}

Use the global form already selected by the determinant seed,
\begin{equation}
 G_{\rm SM}=S(U(3)\times U(2))
 \cong\frac{SU(3)\times SU(2)\times U(1)}{\mathbb Z_6}.
\label{eq:v3v87-global-group}
\end{equation}
The monograph's explicit quotient map is
\begin{equation}
 \Phi(g_3,g_2,z)=(z^{-2}g_3,z^3g_2),
\label{eq:v3v87-quotient-map}
\end{equation}
whose kernel consists of
\((z^2I_3,z^{-3}I_2,z)\) with \(z^6=1\).  Every Standard Model multiplet
used in V66 is trivial on this kernel, so it is a bona fide representation
of \(G_{\rm SM}\).

Let \(\Gamma=(V,E)\) be a connected finite lattice graph and
\(\mathcal G=G_{\rm SM}^{V}\) its sitewise gauge group.  Choose one vertex
\(o\) and put
\begin{equation}
 \mathcal G_*=
 \{g\in\mathcal G:g(o)=1\}.
\label{eq:v3v87-based-group}
\end{equation}

\begin{theorem}[The based gauge action on links is free]
\label{thm:v3v87-based-free}
The action
\begin{equation}
 U_{xy}\longmapsto U_{xy}^g=g(x)U_{xy}g(y)^{-1}
\label{eq:v3v87-link-action}
\end{equation}
of \(\mathcal G_*\) on \(G_{\rm SM}^{E}\) is free.  It remains free on
every gauge-invariant gapped or admissible subset.  Hence the based quotient
has no residual stabilizer singularities.
\end{theorem}

\begin{proof}
Suppose \(U^g=U\) and \(g(o)=1\).  Equation
\eqref{eq:v3v87-link-action} on an edge \((x,y)\) gives
\(g(y)=U_{xy}^{-1}g(x)U_{xy}\).  Along any path from \(o\) to \(y\),
induction from \(g(o)=1\) yields \(g(y)=1\).  Connectedness reaches every
vertex, so \(g\) is the identity.
\end{proof}

The full gauge group may have covariantly constant stabilizers, especially
at reducible connections.  They are retained separately as the constant
global action at the base vertex; they must not be divided out a second time.

\subsection{Loop lifting and the mapping torus}
\label{sec:v3v87-loop}

Let \(\mathcal U^{\rm gap}\) be a gauge-invariant component on which the
regulated chiral projector has constant rank.  A loop in
\(\mathcal U^{\rm gap}/\mathcal G_*\) lifts to a path \(U_t\) with
\begin{equation}
 U_1=U_0^{g}
\label{eq:v3v87-loop-gluing}
\end{equation}
for a based transformation \(g\).  In the continuum, glue
\([0,1]\times X^4\) by \((1,x)\sim(0,x)\) and use \(g\) to glue the
principal \(G_{\rm SM}\)-bundle.  The result is a closed spin
five-manifold \(Y_g\) with a classifying map
\begin{equation}
 f_g:Y_g\longrightarrow BG_{\rm SM}.
\label{eq:v3v87-classifying-map}
\end{equation}

For a four-dimensional chiral representation \(R\), denote the reduced eta
invariant of the associated five-dimensional Dirac operator by
\(\xi_R(Y_g)=(\eta_R+h_R)/2\).  Once the local anomaly polynomial has been
cancelled, the residual holonomy is
\begin{equation}
 \operatorname{Hol}_R(g)
 =\exp\bigl(2\pi i\,\xi_R(Y_g)\bigr),
\label{eq:v3v87-DF-holonomy}
\end{equation}
up to the fixed orientation convention for the determinant line.

\begin{theorem}[Bordism character of the residual phase]
\label{thm:v3v87-bordism-character}
If the complete local anomaly polynomial vanishes, the product of
\eqref{eq:v3v87-DF-holonomy} over the chiral multiplets depends only on
\([Y_g,f_g]\in\Omega^{\rm Spin}_5(BG_{\rm SM})\), and defines a homomorphism
\begin{equation}
 \alpha_{\rm SM}:\Omega^{\rm Spin}_5(BG_{\rm SM})\longrightarrow U(1).
\label{eq:v3v87-anomaly-character}
\end{equation}
\end{theorem}

\begin{proof}
For a spin--\(G_{\rm SM}\) bordism \(Z^6\) between two closed
five-manifolds, the APS index theorem expresses the difference of their
reduced eta invariants modulo integers as the integral over \(Z\) of the
degree-six local index density.  The complete Standard Model sum of that
density is the anomaly polynomial cancelled in V66.  Its integral is zero,
so the exponentiated eta phase is bordism invariant.  Disjoint union adds
eta invariants and therefore multiplies phases, proving the homomorphism
property.
\end{proof}

\subsection{The fifth spin bordism group for the \(\mathbb Z_6\) quotient}
\label{sec:v3v87-bordism}

The calculation uses the group extension
\begin{equation}
 \mathbb Z_3\longrightarrow U(2)\times SU(3)
 \longrightarrow G_{\rm SM},
\label{eq:v3v87-group-extension}
\end{equation}
which yields, through the Puppe sequence, a fibration
\begin{equation}
 B(U(2)\times SU(3))\longrightarrow BG_{\rm SM}
 \longrightarrow K(\mathbb Z_3,2).
\label{eq:v3v87-puppe-fibration}
\end{equation}
In degrees at most five the required base homology is
\begin{align}
 H_p(K(\mathbb Z_3,2);\mathbb Z)
 &\cong
 \begin{cases}
 \mathbb Z,&p=0,\\
 \mathbb Z_3,&p=2,4,\\
 0,&p=1,3,5,
 \end{cases}
\label{eq:v3v87-KZ3-integral}\\
 H_p(K(\mathbb Z_3,2);\mathbb Z_2)
 &\cong
 \begin{cases}
 \mathbb Z_2,&p=0,\\
 0,&1\le p\le5.
 \end{cases}
\label{eq:v3v87-KZ3-mod2}
\end{align}
The fibre calculation needed on the total-degree-five diagonal gives
\begin{equation}
 \Omega^{\rm Spin}_5(B(U(2)\times SU(3)))=0,
 \qquad
 \Omega^{\rm Spin}_3(B(U(2)\times SU(3)))=0,
\label{eq:v3v87-fibre-odd}
\end{equation}
while the relevant nonzero coefficient groups in lower degree contain only
free abelian and \(2\)-primary summands.

\begin{theorem}[Vanishing global-anomaly group]
\label{thm:v3v87-bordism-zero}
For the global group \eqref{eq:v3v87-global-group},
\begin{equation}
 \boxed{\Omega^{\rm Spin}_5(BG_{\rm SM})=0.}
\label{eq:v3v87-bordism-zero}
\end{equation}
\end{theorem}

\begin{proof}
Use the Atiyah--Hirzebruch spectral sequence for
\eqref{eq:v3v87-puppe-fibration},
\begin{equation}
 E^2_{p,q}=H_p\left(K(\mathbb Z_3,2);
 \Omega_q^{\rm Spin}(B(U(2)\times SU(3)))\right).
\label{eq:v3v87-AHSS}
\end{equation}
Consider every pair \((p,q)\) with \(p+q=5\).  The \((0,5)\) and
\((2,3)\) entries vanish by \eqref{eq:v3v87-fibre-odd}.  The entries with
\(p=1,3,5\) vanish by the odd integral homology in
\eqref{eq:v3v87-KZ3-integral}.  The remaining \((4,1)\) entry has a
\(2\)-primary coefficient group and vanishes by
\eqref{eq:v3v87-KZ3-mod2}; equivalently, the positive-degree homology of
the \(3\)-primary base with \(2\)-primary coefficients is zero.  Thus every
entry on the total-degree-five diagonal is already zero on the second page.
No later differential or extension can create a group on an empty diagonal,
which proves \eqref{eq:v3v87-bordism-zero}.
\end{proof}

For independent verification, this is equation (4.45) of J.~Davighi,
B.~Gripaios, and N.~Lohitsiri, \emph{Global anomalies in the Standard
Model(s) and Beyond}, \texttt{arXiv:1910.11277v3}.  The proof above records
the low-degree diagonal used here rather than importing a bare conclusion.

\begin{corollary}[No residual spin--\(G_{\rm SM}\) gauge holonomy]
\label{cor:v3v87-global-anomaly-zero}
For the complete Standard Model representation on honest spin
four-manifolds, after the V66 local anomaly cancellation,
\begin{equation}
 \alpha_{\rm SM}=1.
\label{eq:v3v87-trivial-character}
\end{equation}
Every smooth large-gauge mapping-torus phase of the continuum determinant
line is therefore trivial.
\end{corollary}

\begin{proof}
By Theorem~\ref{thm:v3v87-bordism-character}, the phase is a homomorphism
from the group in Theorem~\ref{thm:v3v87-bordism-zero}.  The only
homomorphism from the zero group sends its sole element to one.
\end{proof}

\subsection{Why the ordinary Witten count is not an extra quotient class}
\label{sec:v3v87-witten}

For \(G_{\rm SM}\) in the \(\mathbb Z_6\) global form, a field of weak
isospin \(j\) and integer-normalized hypercharge \(y\) obeys
\begin{equation}
 y\equiv2j\pmod2.
\label{eq:v3v87-spin-charge}
\end{equation}

\begin{proposition}[Local mixed cancellation enforces weak parity]
\label{prop:v3v87-local-global-interplay}
For any collection of bona fide \(G_{\rm SM}\) Weyl representations,
cancellation of the local \(SU(2)^2U(1)\) anomaly implies
\begin{equation}
 \sum_{j\in2\mathbb Z+1/2}N_j\equiv0\pmod2,
\label{eq:v3v87-witten-parity}
\end{equation}
which is the parity condition for the traditional weak global anomaly.
For the actual one-generation spectrum this is also the direct count of
three coloured quark doublets plus one lepton doublet.
\end{proposition}

\begin{proof}
With Dynkin index normalized by
\begin{equation}
 T(j)=\frac23j(j+1)(2j+1),
\label{eq:v3v87-dynkin-index}
\end{equation}
the mixed local condition is \(\sum_jT(j)\sum_{a=1}^{N_j}y_j^{(a)}=0\).
Modulo two, \(T(j)\) is odd exactly for
\(j\in2\mathbb Z+1/2\).  For precisely those representations,
\eqref{eq:v3v87-spin-charge} makes every \(y_j^{(a)}\) odd; all other
summands are even.  Reducing the mixed condition modulo two gives
\eqref{eq:v3v87-witten-parity}.  In the Standard Model fundamental-doublet
case, \(N_{1/2}=3+1=4\).
\end{proof}

Thus the quotient has not hidden an uncancelled Witten sign.  Its allowed
representation lattice and the local mixed anomaly equation already enforce
the relevant parity, consistently with
Theorem~\ref{thm:v3v87-bordism-zero}.

\subsection{Comparison with the finite overlap line}
\label{sec:v3v87-lattice-comparison}

The bordism theorem is a continuum global result.  To transfer it to the
finite overlap regulator on a large loop, the regulator must remain gapped
along a discretization of the entire mapping torus, which need not lie in
the V78 logarithmic ball.

\begin{theorem}[Conditional regulator-to-Dai--Freed comparison]
\label{thm:v3v87-regulator-comparison}
Let \((Y_g,f_g)\) be a smooth spin--\(G_{\rm SM}\) mapping torus.  Suppose
there is a refining family of V76 domain-wall/overlap lattices on \(Y_g\)
such that:
\begin{enumerate}
 \item the Hermitian Wilson kernels have a gap uniform in the refinement and
 the mapping-torus parameter;
 \item the relative heat-kernel estimates of V75--V76 hold uniformly; and
 \item the same local counterterm normalization as V86 is used.
\end{enumerate}
Then the finite overlap determinant holonomies converge to
\(\alpha_{\rm SM}([Y_g,f_g])=1\).
\end{theorem}

\begin{proof}
The finite-wall Schur formula of V77 expresses the regulated phase through a
finite transfer approximation to the sign function.  The uniform gap makes
that approximation converge uniformly and gives exponential locality.
The relative heat estimates in V75--V76 identify the continuum variation of
the phase with the local index density plus the reduced eta invariant; V86's
normalization subtracts the same local density on every refinement.  The
remaining limit is therefore the Dai--Freed character
\eqref{eq:v3v87-anomaly-character}, which is one by
Corollary~\ref{cor:v3v87-global-anomaly-zero}.
\end{proof}

\begin{firewall}[The all-sector Wilson gap remains open]
\label{fw:v3v87-global-gap-open}
The hypotheses of Theorem~\ref{thm:v3v87-regulator-comparison} are proved in
V78 and V84 only on a quantitative component near the trivial connection.
V87 does not prove a uniform overlap gap on representatives of every smooth
mapping-torus class.  The continuum obstruction vanishes, but the current
lattice comparison is conditional outside the local patch.
\end{firewall}

\begin{firewall}[Other tangential structures require new bordism groups]
\label{fw:v3v87-structure-scope}
Equation \eqref{eq:v3v87-bordism-zero} assumes an honest spin structure and
the direct spin--\(G_{\rm SM}\) background used here.  A spin--gauge
quotient, spin\({}^c\) formulation, unorientable background, or the extra
material-frame boundary loop of V67 has a different tangential structure
and is not decided by this computation.
\end{firewall}

\subsection{V87 conclusion and V88 target}
\label{sec:v3v87-conclusion}

The terminal continuum gauge anomaly line now has no global obstruction for
the manuscript's actual \(\mathbb Z_6\) quotient on spin manifolds.  The
based finite-lattice action is free, every residual phase is a five-dimensional
spin-bordism character, and the entire source group of that character
vanishes.  The local mixed anomaly equation also explains why the familiar
weak parity condition is automatic for allowed anomaly-free quotient
representations.

The remaining regulator gate is constructive rather than cohomological.
V88 should replace the global Wilson-gap hypothesis in
Theorem~\ref{thm:v3v87-regulator-comparison} by a spectral-section
construction: cover a smooth mapping torus by finitely many local gap charts,
patch their determinant lines with exact transition functions, and show that
the zero bordism character removes the final Cech cocycle.  If a Wilson zero
is crossed, it must be recorded by spectral flow rather than excluded by a
false uniform gap.


\section{V88: spectral-section patching across Wilson crossings}
\label{sec:v3v88}

V87 proved that the continuum spin--\(G_{\rm SM}\) anomaly character has a
zero source group, but its lattice comparison still assumed a Wilson-kernel
gap around an entire large loop.  That assumption excludes precisely the
spectral crossings which a determinant line is designed to record.  This
note replaces it by local spectral cuts.  The transition between two cuts is
the determinant of a finite eigenvalue band; spectral flow is their Cech
product.  Trivial continuum holonomy then patches the complete Standard
Model tensor-product line to a global section.  Physical Dirac zero modes
remain zeros of that section rather than singularities of the line.

\subsection{Local spectral cuts}
\label{sec:v3v88-cuts}

Let \(B\) be a compact parameter space of smooth gauge and metric
backgrounds, and let \(D_b=D_b^*\) be a norm-resolvent-continuous family of
self-adjoint Dirac-type operators with compact resolvent.  For
\(\lambda\in\mathbb R\), put
\begin{equation}
 B_\lambda=\{b\in B:\lambda\notin\operatorname{spec}D_b\}.
\label{eq:v3v88-cut-chart}
\end{equation}
Compactness of the resolvent makes the spectral subspace in every bounded
interval finite dimensional.

On \(B_\lambda\), the spectral cut defines a local determinant line
\(\mathcal L_\lambda\).  If \(\lambda<\mu\), then on
\(B_\lambda\cap B_\mu\) let
\begin{equation}
 E_{(\lambda,\mu)}(b)
 =\mathbf1_{(\lambda,\mu)}(D_b)\mathcal H.
\label{eq:v3v88-transition-band}
\end{equation}

\begin{proposition}[Finite-band transition line]
\label{prop:v3v88-transition-line}
There is a canonical isomorphism on every overlap,
\begin{equation}
 \mathcal L_\mu
 \cong\mathcal L_\lambda\otimes
 \det E_{(\lambda,\mu)}^{\,\epsilon},
\label{eq:v3v88-transition-isomorphism}
\end{equation}
where \(\epsilon=1\) or \(-1\) according to the fixed determinant-line cut
convention.  For three ordered cuts \(\lambda<\mu<\nu\), these
isomorphisms obey the exact cocycle identity because
\begin{equation}
 E_{(\lambda,\nu)}
 =E_{(\lambda,\mu)}\oplus E_{(\mu,\nu)}.
\label{eq:v3v88-band-sum}
\end{equation}
\end{proposition}

\begin{proof}
Changing a polarization from cut \(\lambda\) to cut \(\mu\) moves only the
finitely many eigenvectors in \((\lambda,\mu)\) between the two polarized
factors.  Taking their top exterior power gives
\eqref{eq:v3v88-transition-isomorphism}.  Equation
\eqref{eq:v3v88-band-sum} is an orthogonal spectral decomposition, and the
canonical identity
\(\det(E\oplus F)=\det E\otimes\det F\) proves the cocycle law, including
the fixed graded ordering sign.
\end{proof}

Thus no individual cut is required to remain outside the spectrum globally.
A finite cover by cut charts exists by compactness of \(B\).

\subsection{Spectral flow is the obstruction to one global cut}
\label{sec:v3v88-spectral-flow}

Consider a loop \(D_t\), \(t\in S^1\), with only regular crossings after an
arbitrarily small perturbation.  At a crossing time \(t_*\), the crossing
form on \(\ker D_{t_*}\) is
\begin{equation}
 q_{t_*}(v)=\langle v,\dot D_{t_*}v\rangle.
\label{eq:v3v88-crossing-form}
\end{equation}
Its signature is the signed number of eigenvalues crossing zero.

\begin{theorem}[Spectral-section criterion on a loop]
\label{thm:v3v88-spectral-section}
A self-adjoint Fredholm loop admits a global spectral section if and only if
its family index in \(K^1(S^1)\cong\mathbb Z\) vanishes.  Equivalently,
\begin{equation}
 \operatorname{SF}(D_t)
 =\sum_{t_*}\operatorname{sign}q_{t_*}=0.
\label{eq:v3v88-SF-zero}
\end{equation}
When the spectral flow is nonzero, local cut lines still patch by
Proposition~\ref{prop:v3v88-transition-line}; their transition product
records that integer rather than making the determinant line undefined.
\end{theorem}

\begin{proof}
For regular crossings, choose disjoint intervals about the crossing times
and a spectral cut on each complementary interval.  Following the negative
spectral subspace around the circle changes its virtual rank by the sum of
the crossing signatures.  A global spectral section has periodic finite-rank
correction, so this net change must be zero.

Conversely, if the signed crossing sum is zero, pair positive and negative
crossings in cyclic order.  On the arc between each pair, add a finite-rank
continuous projection transported from the crossing eigenspaces; at the
second crossing remove it.  The resulting projection agrees with the
negative spectral projection outside a finite window and is periodic.
This is a spectral section.  The signed crossing sum represents the standard
generator pairing of the family index in \(K^1(S^1)\), proving the first
equivalence.  The final statement is the finite-band cocycle of
Proposition~\ref{prop:v3v88-transition-line} multiplied around the cover.
\end{proof}

The theorem explains why V87's proposed all-loop Wilson gap was too strong.
An individual regulator family may have spectral flow even when the complete
anomaly line has trivial phase.

\subsection{Tensor first, then trivialize}
\label{sec:v3v88-tensor}

For every left-handed Standard Model multiplet \(R\), let \(\mathcal L_R\)
be its determinant line and define the complete line
\begin{equation}
 \mathcal L_{\rm SM}=\bigotimes_R\mathcal L_R,
\label{eq:v3v88-SM-line}
\end{equation}
using the multiplicities and conjugate-representation convention of V66.
The tensor product is formed before transition functions are estimated.

\begin{theorem}[Flat line with trivial holonomy has a global parallel
section]
\label{thm:v3v88-flat-trivial}
Let \(L\to X\) be a complex line with a flat connection over a connected
space.  If its holonomy is one around every loop, then for any
\(x_0\in X\) and \(0\ne v_0\in L_{x_0}\) there is a unique global parallel
section \(s\) satisfying \(s(x_0)=v_0\).
\end{theorem}

\begin{proof}
For \(x\in X\), parallel transport \(v_0\) along any path from \(x_0\) to
\(x\).  Two choices differ by transport around the loop obtained by
concatenating one with the reverse of the other.  Trivial holonomy makes the
two values equal.  Local smooth dependence of parallel transport gives a
parallel section.  Its value at \(x_0\) fixes it uniquely.
\end{proof}

\begin{corollary}[Cech cancellation of the complete Standard Model line]
\label{cor:v3v88-Cech-trivial}
On the continuum background component with honest spin structure and gauge
group \(G_{\rm SM}\), impose the V86 local counterterm normalization.  The
finite-band transition functions of the cut charts for
\(\mathcal L_{\rm SM}\) form a Cech coboundary.  Consequently they patch to
a global smooth determinant-line section, unique up to one constant phase on
each connected component.
\end{corollary}

\begin{proof}
V66 cancels the local curvature in gauge directions, and V86 constructs its
renormalized local connection.  Corollary
\ref{cor:v3v87-global-anomaly-zero} says that the remaining holonomy about
every large gauge loop is one.  Apply
Theorem~\ref{thm:v3v88-flat-trivial} to the gauge-anomaly line.  In a cut
cover, write the resulting parallel section as \(h_i\) times the chosen
local cut section.  If \(g_{ij}\) is the finite-band transition function,
equality of the global section gives
\begin{equation}
 g_{ij}=h_i h_j^{-1},
\label{eq:v3v88-Cech-coboundary}
\end{equation}
so the cocycle is exact and the rescaled local sections patch.
\end{proof}

The constant phase is an ordinary normalization choice, not a gauge anomaly.

\subsection{Zero modes are zeros of the section}
\label{sec:v3v88-zero-modes}

\begin{proposition}[Crossing a physical zero does not destroy the line]
\label{prop:v3v88-zero-section}
Let a finite-dimensional low eigenvalue \(\lambda(b)\) of the physical
chiral Dirac operator cross zero transversely.  In a determinant-line chart,
the fermion determinant section factors as
\begin{equation}
 \mathcal Z_F(b)=\lambda(b)\,\mathcal Z_F'(b),
\label{eq:v3v88-zero-factor}
\end{equation}
where \(\mathcal Z_F'\) is smooth and nonzero at the crossing if no other
zero mode is present.  Hence the section vanishes smoothly; only a formula
which divides by \(\lambda\) is singular.
\end{proposition}

\begin{proof}
Use the Riesz projection onto a small disk containing only the crossing
eigenvalue.  It gives a smooth one-dimensional eigenbundle, while the
restriction of the operator to its orthogonal complement remains invertible
and has a smooth regularized determinant.  Multiplicativity of the
determinant under this spectral decomposition gives
\eqref{eq:v3v88-zero-factor}.
\end{proof}

This distinction is necessary for Yukawa backgrounds: physical massless
fermions can make the partition function vanish without creating a global
orientation obstruction.

\subsection{Finite regulator patching without a global sign projector}
\label{sec:v3v88-regulator}

Let \(H_{W,b}\) be a finite-lattice Hermitian Wilson family on a compact
parameter loop.  Choose local mass or spectral-cut charts on which
\(H_{W,b}-\lambda_i\) is invertible.  Each chart has the V77 overlap
projector and determinant line.  On an overlap, the two sign projectors
differ by the finite spectral band between \(\lambda_i\) and \(\lambda_j\),
so Proposition~\ref{prop:v3v88-transition-line} gives an exact
finite-dimensional transition determinant.

\begin{theorem}[Regulator patching criterion]
\label{thm:v3v88-regulator-patching}
Suppose a refining family of such cut covers has:
\begin{enumerate}
 \item volume-uniform gaps from each local cut on its chart;
 \item the V75--V76 relative heat estimates on chart overlaps;
 \item convergence of every finite-band transition determinant to its
 continuum spectral-section transition; and
 \item the common V86 local normalization.
\end{enumerate}
Then the complete Standard Model lattice determinant lines and connections
converge after Cech patching to the global continuum section of
Corollary~\ref{cor:v3v88-Cech-trivial}.  A global Wilson gap is not required.
\end{theorem}

\begin{proof}
Item 1 defines each local overlap line.  Item 2 gives convergence of its
connection and local phase after the shared heat subtraction.  Item 3 makes
the overlap identifications compatible in the limit, while the exact
finite-band direct-sum identity gives the triple-overlap cocycle at every
finite regulator.  Tensor over the full representation before passing to the
limit.  By item 4 and V87, the limiting cocycle has trivial holonomy and is
the coboundary \eqref{eq:v3v88-Cech-coboundary}.  The rescaled local lattice
sections therefore converge to the asserted global section.
\end{proof}

\begin{firewall}[Transition convergence is a remaining analytic estimate]
\label{fw:v3v88-transition-open}
V88 removes the false need for a global Wilson gap, but item 3 of
Theorem~\ref{thm:v3v88-regulator-patching} has not been proved uniformly for
all smooth mapping-torus families.  It requires a crossing-local comparison
between lattice and continuum eigenbundles, including multiplicities and
crossing forms.
\end{firewall}

\begin{firewall}[Gauge triviality does not prove reflection positivity]
\label{fw:v3v88-rp-open}
A global complex determinant section can exist even when its chosen
Pauli--Villars quotient lacks a positive Osterwalder--Schrader
representation.  Spectral-section patching settles neither that positivity
problem nor the material boundary-domain problem corrected in V68.
\end{firewall}

\subsection{V88 conclusion and V89 target}
\label{sec:v3v88-conclusion}

Large-loop Wilson crossings are now treated as data rather than excluded.
Local spectral cuts differ by finite eigenvalue-band determinants, their
cocycle records spectral flow, and the vanishing spin--\(G_{\rm SM}\)
bordism character makes the complete continuum cocycle exact.  The
renormalized Standard Model determinant line therefore has a global smooth
section on each spin background component, up to one constant phase.

V89 should prove the remaining crossing-local regulator comparison.  Around
one regular continuum crossing it should use norm-resolvent convergence and
a Riesz contour to identify the lattice and continuum crossing bundles,
prove equality of crossing signatures, and then extend the result to a
finite cover by homotopy-stable multiplicity clusters.  This will replace
item 3 of Theorem~\ref{thm:v3v88-regulator-patching} by checkable operator
estimates.


\section{V89: crossing-local Riesz comparison and the Wilson/physical spectral distinction}
\label{sec:v3v89}

V88 reduced global regulator patching to convergence of finite spectral-band
transition determinants.  This note proves the required comparison theorem
under explicit generalized norm-resolvent and derivative estimates.  It also
corrects a possible ambiguity in the target: a bounded physical Dirac
crossing is not a zero of the dimensionless massive Wilson kernel.  The
Riesz argument below applies directly to physical and suspended Dirac
families.  Wilson-kernel transitions require convergence of the lattice
index class, a separate statement left visible at the end.

\subsection{Generalized norm-resolvent hypotheses}
\label{sec:v3v89-hypotheses}

Let \(I\subset\mathbb R\) be compact.  For \(t\in I\), let
\(D(t)=D(t)^*\) have compact resolvent on a Hilbert space \(\mathcal H\), and
let \(D_a(t)=D_a(t)^*\) act on a finite-dimensional lattice space
\(\mathcal H_a\).  Let \(J_a:\mathcal H_a\to\mathcal H\) be an isometry.
For a compact contour \(\Gamma\) contained in the common resolvent set,
assume
\begin{equation}
 \varepsilon_a(\Gamma):=
 \sup_{t\in I}\sup_{z\in\Gamma}
 \left\|J_a(D_a(t)-z)^{-1}J_a^*-(D(t)-z)^{-1}\right\|
 \longrightarrow0.
\label{eq:v3v89-NR}
\end{equation}
This formulation includes approximation of the compact continuum resolvent
outside \(\operatorname{ran}J_a\); it is stronger than strong-resolvent
convergence and rules out spectral pollution in the bounded window enclosed
by \(\Gamma\).

\begin{theorem}[Uniform Riesz-bundle convergence]
\label{thm:v3v89-riesz-bundle}
Let \(\Gamma\) enclose a finite spectral cluster of \(D(t)\), separated from
the rest of the spectrum for every \(t\in I\), and put
\begin{align}
 P(t)&=\frac{1}{2\pi i}\int_\Gamma(z-D(t))^{-1}\,dz,
 \label{eq:v3v89-continuum-Riesz}\\
 P_a(t)&=\frac{1}{2\pi i}\int_\Gamma(z-D_a(t))^{-1}\,dz.
 \label{eq:v3v89-lattice-Riesz}
\end{align}
Under \eqref{eq:v3v89-NR}, for all sufficiently small \(a\),
\begin{equation}
 \sup_{t\in I}\|J_aP_a(t)J_a^*-P(t)\|
 \le\frac{\operatorname{length}\Gamma}{2\pi}
 \varepsilon_a(\Gamma)<1.
\label{eq:v3v89-projection-bound}
\end{equation}
In particular, the two cluster bundles have the same finite rank and are
identified by a continuous canonical unitary \(W_a(t)\) from
\(\operatorname{ran}P(t)\) to
\(J_a\operatorname{ran}P_a(t)\), with
\begin{equation}
 \sup_{t\in I}\|W_a(t)-I\|_{\operatorname{ran}P(t)}\longrightarrow0.
\label{eq:v3v89-unitary-close}
\end{equation}
\end{theorem}

\begin{proof}
Subtract the two contour formulae and use
\eqref{eq:v3v89-NR}; this gives
\eqref{eq:v3v89-projection-bound}.  Two finite-rank orthogonal projections
at distance less than one have equal rank.  More explicitly, if
\(Q_a=J_aP_aJ_a^*\), the polar part of
\begin{equation}
 Q_aP+(I-Q_a)(I-P)
\label{eq:v3v89-direct-rotation-map}
\end{equation}
is a unitary intertwining \(P\) and \(Q_a\).  Functional calculus of the
positive factor is continuous while \(\|P-Q_a\|<1\), and tends to the
identity as that norm tends to zero.  Restrict this direct rotation to
\(\operatorname{ran}P\) to obtain \(W_a\) and
\eqref{eq:v3v89-unitary-close}.
\end{proof}

\subsection{Regular crossings and their signatures}
\label{sec:v3v89-crossings}

Let \(t_0\) be an isolated continuum crossing at zero.  Choose
\(\rho>0\) and \(\delta>0\) such that on
\(I_0=[t_0-\delta,t_0+\delta]\) the interval
\((-\rho,\rho)\) contains only the crossing cluster and its endpoints are
outside every spectrum.  Let \(P(t)\) and \(P_a(t)\) be its Riesz
projections.  The continuum crossing form is
\begin{equation}
 Q_0=P(t_0)\dot D(t_0)P(t_0)\big|_{\operatorname{ran}P(t_0)}.
\label{eq:v3v89-crossing-form}
\end{equation}
Assume it is invertible.  In addition to
\eqref{eq:v3v89-NR}, assume the cluster derivative estimate
\begin{equation}
 \left\|W_a(t_0)^*J_aP_a(t_0)\dot D_a(t_0)P_a(t_0)J_a^*W_a(t_0)-Q_0\right\|
 \longrightarrow0.
\label{eq:v3v89-derivative-convergence}
\end{equation}

\begin{theorem}[Crossing-signature stability]
\label{thm:v3v89-signature-stability}
Under these hypotheses, the total lattice spectral flow through zero on
\(I_0\) equals the continuum crossing signature for all sufficiently small
\(a\):
\begin{equation}
 \operatorname{SF}\bigl(D_a(t);I_0\bigr)
 =\operatorname{sign}Q_0
 =\operatorname{SF}\bigl(D(t);I_0\bigr).
\label{eq:v3v89-SF-equality}
\end{equation}
The lattice crossing may split into several regular crossings; their signed
sum is fixed by \eqref{eq:v3v89-SF-equality}.  If the continuum cluster is
simple and \(\dot\lambda(t_0)\ne0\), uniform \(C^1\) convergence of the
cluster eigenvalue gives a unique lattice zero \(t_a\to t_0\) with
\begin{equation}
 \operatorname{sgn}\dot\lambda_a(t_a)
 =\operatorname{sgn}\dot\lambda(t_0).
\label{eq:v3v89-simple-sign}
\end{equation}
\end{theorem}

\begin{proof}
At \(t_0\pm\delta\), zero is separated from both the cluster and the
remaining spectrum.  Riesz convergence implies convergence of the negative
spectral projections in \((-\rho,\rho)\), hence equality of their ranks for
small \(a\).  Spectral flow on an interval with invertible endpoints is the
negative-rank at the left endpoint minus the negative-rank at the right
endpoint, with the chosen positive-crossing convention.  These endpoint
ranks are the same for the lattice and continuum families, proving the first
and third terms in \eqref{eq:v3v89-SF-equality}.  For an isolated regular
continuum crossing, the latter equals \(\operatorname{sign}Q_0\) by
diagonalizing the finite crossing form.

For a simple cluster, Riesz projection identifies a single continuous
eigenvalue.  Uniform convergence gives opposite endpoint signs, while
uniform derivative convergence makes the derivative retain the nonzero sign
of \(\dot\lambda(t_0)\) in a smaller interval.  The intermediate-value
theorem gives a zero, strict monotonicity gives uniqueness, and the same
derivative bound proves \eqref{eq:v3v89-simple-sign}.
\end{proof}

The derivative estimate is needed only for the local direction of individual
crossings.  Equality of the total spectral flow requires norm-resolvent
convergence and endpoint gaps, and is therefore stable under multiplicity
splitting.

\subsection{Convergence of finite-band determinant transitions}
\label{sec:v3v89-transition}

Fix cuts \(\lambda<\mu\) which avoid the continuum and lattice spectra on a
compact overlap \(K\subset B_\lambda\cap B_\mu\).  Let
\(E_{(\lambda,\mu)}\) and \(E_{a,(\lambda,\mu)}\) be the continuum and
lattice bands.

\begin{theorem}[Finite-band transition convergence]
\label{thm:v3v89-transition-convergence}
If generalized norm-resolvent convergence holds uniformly on contours
surrounding \((\lambda,\mu)\), then for small \(a\):
\begin{enumerate}
 \item the two band bundles have the same rank and a canonical unitary
 identification converging uniformly to the identity after embedding;
 \item their top exterior-power transition lines converge in operator norm;
 and
 \item for a \(C^1\) family satisfying
 \eqref{eq:v3v89-derivative-convergence}, the induced determinant-line
 connections and crossing orientations converge uniformly on \(K\).
\end{enumerate}
Consequently the finite-band transition determinant in
Proposition~\ref{prop:v3v88-transition-line} converges to its continuum
spectral-section transition.
\end{theorem}

\begin{proof}
Apply Theorem~\ref{thm:v3v89-riesz-bundle} to a contour enclosing the band.
The canonical direct rotation gives item 1.  The top exterior-power functor
is a polynomial in the matrix entries on a fixed finite rank, so convergence
of the unitary gives item 2.  Differentiate the Riesz formula; the derivative
is a contour integral of
\((z-D)^{-1}\dot D(z-D)^{-1}\).  The derivative hypothesis therefore gives
convergence of \(P_a\,dP_a\) after direct rotation.  The determinant
connection is its finite trace, which proves item 3.  The transition
determinant is precisely the top exterior action of this identified band.
\end{proof}

\subsection{From one crossing to a compact parameter family}
\label{sec:v3v89-finite-cover}

\begin{theorem}[Finite-cover crossing comparison]
\label{thm:v3v89-finite-cover}
Let \(B\) be compact.  Suppose every point has a neighbourhood with a finite
spectral-band contour on which the generalized norm-resolvent estimate is
uniform, and suppose derivative convergence holds near the finitely many
regular crossing clusters of a generic loop.  Then one finite cut cover
works for all sufficiently small \(a\); its lattice Cech transition lines,
triple-overlap products, and loop spectral flows converge to the continuum
ones.  Nonregular crossings may be split by a small endpoint-fixed
perturbation, and the conclusion is independent of that perturbation.
\end{theorem}

\begin{proof}
Choose one contour neighbourhood at every point and extract a finite
subcover.  The minimum of their contour distances from the continuum spectra
is positive.  Uniform norm-resolvent convergence keeps every lattice
spectrum off the same contours for small \(a\).  Apply
Theorem~\ref{thm:v3v89-transition-convergence} on all pairwise overlaps.
The direct-sum identity \eqref{eq:v3v88-band-sum} makes every triple product
exact before the limit, so its limit has the continuum cocycle.  Apply
Theorem~\ref{thm:v3v89-signature-stability} on the crossing intervals and
ordinary gapped convergence elsewhere to obtain loop spectral-flow
convergence.

A generic finite-rank perturbation makes the crossings regular without
moving the loop endpoints through zero.  Spectral flow and the determinant
line cocycle are invariant under such an endpoint-gapped homotopy, so the
answer does not depend on the perturbation.
\end{proof}

\begin{corollary}[Discharge of the V88 transition hypothesis]
\label{cor:v3v89-v88-discharge}
Item 3 of Theorem~\ref{thm:v3v88-regulator-patching} follows for every
physical or suspended Dirac lattice family satisfying the uniform
generalized norm-resolvent and cluster derivative estimates above.  Together
with the remaining V88 hypotheses, its patched determinant line converges to
the global Standard Model section.
\end{corollary}

\subsection{A checkable route to norm-resolvent convergence}
\label{sec:v3v89-consistency}

The abstract hypothesis can be reduced to consistency and compactness.  Let
\(R(z)=(D-z)^{-1}\) and
\(R_a(z)=(D_a-z)^{-1}\).

\begin{proposition}[Consistency plus collective compactness]
\label{prop:v3v89-collective-compactness}
Suppose for one nonreal \(z\):
\begin{enumerate}
 \item \(J_aR_a(z)J_a^*f\to R(z)f\) for every \(f\in\mathcal H\);
 \item the set
 \(\{J_aR_a(z)J_a^*f:\|f\|\le1,\ 0<a<a_0\}\) is relatively compact; and
 \item the same two statements hold for the adjoint resolvents.
\end{enumerate}
Then
\begin{equation}
 \|J_aR_a(z)J_a^*-R(z)\|\longrightarrow0.
\label{eq:v3v89-collective-NR}
\end{equation}
The resolvent identity extends this convergence uniformly to compact subsets
of the common resolvent set.
\end{proposition}

\begin{proof}
If norm convergence failed, there would be unit vectors \(f_a\) and a
positive lower bound for
\(\|(J_aR_aJ_a^*-R)f_a\|\).  Pass to a weakly convergent subsequence of
\(f_a\).  Collective compactness makes the images under the lattice
resolvents converge strongly; compactness of \(R\) does the same for the
continuum images.  Strong convergence on fixed vectors, applied after
splitting off the weak limit, and the corresponding adjoint hypothesis force
both strong limits to agree, contradicting the lower bound.  This is the
collectively compact convergence theorem.  The resolvent identity propagates
norm convergence from one point to every compact set which remains a fixed
distance from the spectra.
\end{proof}

For a consistent finite-difference Dirac family on a smooth compact
background, item 1 is proved first on a smooth core.  A volume-uniform
discrete Garding estimate gives a bounded discrete \(H^1\) norm of the
resolvent image; interpolation followed by Rellich compactness gives item 2.
Applying the same argument to the adjoint gives item 3.  Thus the exact
operator estimates to be verified are a consistency error and a uniform
discrete Sobolev bound, rather than convergence of individual eigenvectors by
choice of phase.

\subsection{The massive Wilson kernel is a different comparison problem}
\label{sec:v3v89-Wilson-distinction}

The overlap regulator uses
\begin{equation}
 H_{W,a}=\gamma_5(D_{W,a}-m_0/a),
 \qquad0<m_0<2r.
\label{eq:v3v89-massive-Wilson}
\end{equation}
Its bounded dimensionless form is \(aH_{W,a}\), whose physical branch tends
at fixed continuum momentum to the constant mass matrix
\(-m_0\gamma_5\), while the Wilson term separates the doubler branches.
It does not converge in norm resolvent to the massless physical Dirac
operator \(D\).  Therefore Theorem~\ref{thm:v3v89-riesz-bundle} cannot be
used to identify a zero of \(H_{W,a}\) with a physical Dirac zero.

\begin{proposition}[Correct division of spectral tasks]
\label{prop:v3v89-three-spectra}
The three spectral comparisons required by the programme are distinct:
\begin{enumerate}
 \item physical Dirac zero modes are treated by
 Proposition~\ref{prop:v3v88-zero-section} and the norm-resolvent theorem
 above;
 \item the five-dimensional suspended Dirac spectral flow gives the
 Dai--Freed mapping-torus phase and is covered by the same theorem when its
 lattice suspension is consistent; and
 \item the massive Wilson-kernel sign defines the overlap topological class
 and must be compared to the continuum index by a lattice index theorem on
 an admissible bundle sector.
\end{enumerate}
No implication between items 1 and 3 follows from norm-resolvent convergence.
\end{proposition}

\begin{proof}
Items 1 and 2 concern bounded spectral windows of consistent discretizations
of their corresponding continuum differential operators.  Equation
\eqref{eq:v3v89-massive-Wilson} contains a regulator mass of order \(a^{-1}\)
and extra ultraviolet branches.  Its scaled low-momentum limit is not either
continuum operator, while its sign carries the degree computed in V76.
Hence its comparison is topological and uses the Wilson index, not the Riesz
limit of the physical spectrum.
\end{proof}

\begin{firewall}[The Wilson index comparison remains]
\label{fw:v3v89-Wilson-index-open}
V89 proves crossing-bundle and transition convergence for physical and
suspended Dirac discretizations under checkable norm-resolvent hypotheses.
It does not prove that the finite overlap index equals the continuum chiral
index on every admissible \(G_{\rm SM}\)-bundle sector.  That is the
remaining regulator-topology gate.
\end{firewall}

\begin{firewall}[Garding verification is background dependent]
\label{fw:v3v89-Garding-open}
V76 supplies the required type of discrete Garding estimate on its controlled
admissible backgrounds.  Extending it uniformly to arbitrary curved
spin--\(G_{\rm SM}\) mapping tori, with transition-function lattice charts,
has not yet been carried out.
\end{firewall}

\subsection{V89 conclusion and V90 audit target}
\label{sec:v3v89-conclusion}

Finite spectral bands now converge without choosing eigenvector phases:
Riesz projections give canonical close bundle identifications, regular
crossing signatures and total spectral flow are stable, and top exterior
powers give the required determinant transition convergence.  This
discharges the V88 transition hypothesis for consistent physical and
suspended Dirac lattices.  It also prevents a false step: the massive Wilson
kernel is not a norm-resolvent approximation of the physical Dirac operator.

V90 is the next compulsory NS/YM audit.  It should then attack the remaining
topological regulator gate by proving an admissible-sector lattice index
comparison for the overlap projector.  The proof must use bundle transition
data or a gauge-covariant interpolation to a geometric lattice Dirac
representative; it must not identify Wilson and physical eigenvalues
mode by mode.

\clearpage
\section{V90: admissible Wilson sectors and the geometric overlap index}
\label{sec:v3v90}

V89 separated three operators which must not be conflated: the physical
Dirac operator, the suspended Dirac operator, and the massive Wilson kernel.
This note treats the third one.  Its purpose is to compare the finite overlap
index with the continuum chiral index without matching their eigenvalues.
The argument has three parts: a volume-uniform gap for almost commuting
covariant shifts, a covariant local-index expansion on geometrically sampled
smooth fields, and integer stabilization.  The result closes the regulator
topology gate on flat smooth bundle sectors.  Curved spin backgrounds and the
family-valued refinement needed for the Einstein continuum are kept as an
explicit next gate.

\subsection{Compulsory V90 NS/YM audit}
\label{sec:v3v90-audit}

The audit was performed before choosing the argument below.  The exact
read-only fingerprints were:
\begin{align}
 &\texttt{Renewal\_NS\_Stability\_Research\_Notes\_V31.tex},
 &&887537\ \text{bytes},
 &&\texttt{F1121B62238AD5491BB7CF03635EA9934942EDF573ED8FAAA9EE79E1D38A6696},
 \label{eq:v3v90-audit-NS}\\
 &\texttt{Renewal\_Yang\_Mills\_Mass\_Gap\_Research\_Notes\_V11.tex},
 &&205286\ \text{bytes},
 &&\texttt{BF48582BB793648986BF4E20483B1ACCAC090B1F064415C0C580B8EFE40BE69B}.
 \label{eq:v3v90-audit-YM}
\end{align}
At the audit time the named YM V11 and V10 files were byte-identical.  The
two frozen YM volumes were also checked:
\begin{align}
 &\texttt{Renewal\_Yang\_Mills\_Mass\_Gap\_Research\_Notes\_V135\_f1.tex},
 &&3083467\ \text{bytes},
 &&\texttt{82F6BBFFEECDAD3C5C63551660B19A0BDC7F3D6DFFE4C68B4E9BBA3A783A2B71},
 \label{eq:v3v90-audit-YM-f1}\\
 &\texttt{Renewal\_Yang\_Mills\_Mass\_Gap\_Research\_Notes\_V100\_f2.tex},
 &&1706928\ \text{bytes},
 &&\texttt{6444A1044E74DFCB1A72ED9C7248BC0F234EFC4A0E4248C736CD7DB212DFCA9A}.
 \label{eq:v3v90-audit-YM-f2}
\end{align}

The NS lesson is that exact marking and unmarking consumes a first moment;
placing an absolute value before the cancellation merely hides the missing
scale weight.  The newest YM lesson is its exact one-use payer deficit: a
stock endpoint produces \(H_{xq}\), whereas the quartic response requires
\(H_{xq}/\Xi_x\), and private mass can make the missing factor arbitrarily
large.  The cure there is an exact post-prefix depleted identity, not a
stronger norm estimate.

For the present problem the analogous nonnegotiable object is the complete
signed diagonal supertrace.  A comparison of the total number of low
physical modes with a continuum index would omit the Wilson branches; a
termwise absolute bound would erase their oriented cancellation.  We must
first assemble the physical corner and every doubler corner in one sign
kernel.  Only after that assembly may we estimate its local remainder.  This
is the direction adopted below, and it is not contained in the audited NS or
YM notes.

\subsection{Geometric lattice sampler and conventions}
\label{sec:v3v90-sampler}

We first work on a flat closed spin four-torus \(X=\mathbb R^4/\Lambda\),
with an arbitrary spin structure.  Let \(E\to X\) be a Hermitian vector
bundle associated to a principal \(G_{\rm SM}\)-bundle through a unitary
representation \(R\), and let \(\nabla^E\) be smooth.  A cubical mesh
\(X_a\) is chosen so that its edge lengths are \(a\).  The lattice fibre at
\(x\) is \(S_x\otimes E_x\).  Parallel transport along the positively
oriented edge from \(x\) to \(x+a e_\mu\) defines the unitary covariant
shift \(T_{\mu,a}\).  This fibrewise definition does not require a global
frame.  A change of local frames conjugates all shifts by one block-diagonal
unitary.

Put
\begin{equation}
 \mathcal D_a=\frac12\sum_{\mu=1}^4
 \gamma_\mu(T_{\mu,a}-T_{\mu,a}^*),
 \qquad
 \mathcal W_a=\frac r2\sum_{\mu=1}^4
 (2-T_{\mu,a}-T_{\mu,a}^*),
\label{eq:v3v90-DW}
\end{equation}
and, for \(m\) in a compact interval
\(I\Subset(0,2r)\),
\begin{equation}
 \widehat H_a(m)=\gamma_5(\mathcal D_a+\mathcal W_a-m)
 =aH_{W,a}.
\label{eq:v3v90-Hhat}
\end{equation}
Thus \(\widehat H_a\) is dimensionless and Hermitian.  Its overlap index and
site density are
\begin{equation}
 \operatorname{ind}_{\rm ov}(a,E)
 =-\frac12\operatorname{Tr}\operatorname{sgn}\widehat H_a,
 \qquad
 q_a(x)=-\frac12\operatorname{tr}_{S\otimes E_x}
 \operatorname{sgn}\widehat H_a(x,x).
\label{eq:v3v90-index-density}
\end{equation}
This agrees with the V77 convention
\eqref{eq:v3v77-index}.  The trace in the first formula uses counting
measure, so \(\operatorname{ind}_{\rm ov}=\sum_{x\in X_a}q_a(x)\).

For a plaquette based at \(x\), write
\begin{equation}
 P_{\mu\nu,a}(x)=
 T_{\mu,a}T_{\nu,a}T_{\mu,a}^*T_{\nu,a}^*
 \quad\hbox{on the fibre at }x,
 \qquad
 \varepsilon_a=\max_{x,\mu,\nu}\|1-P_{\mu\nu,a}(x)\|.
\label{eq:v3v90-plaquette}
\end{equation}

\begin{lemma}[Smooth holonomy is uniformly admissible]
\label{lem:v3v90-smooth-admissible}
For a bounded set \(\mathscr K\) of smooth connections in \(C^2\), there
are constants \(C_{\mathscr K}\) and \(a_{\mathscr K}>0\) such that
\begin{equation}
 \varepsilon_a\le
 a^2\sup_{A\in\mathscr K}\|F_A\|_\infty
 +C_{\mathscr K}a^3,
 \qquad 0<a<a_{\mathscr K}.
\label{eq:v3v90-plaquette-expansion}
\end{equation}
The estimate is independent of local trivialization and of the number of
lattice sites.
\end{lemma}

\begin{proof}
Choose radial gauge at one vertex of a plaquette.  Parallel transport around
its boundary satisfies the nonabelian Stokes expansion
\begin{equation}
 P_{\mu\nu,a}(x)
 =1+a^2F_{\mu\nu}(x)+a^3R_{\mu\nu,a}(x),
\label{eq:v3v90-nonabelian-Stokes}
\end{equation}
where the integral-remainder form of the covariant Taylor formula bounds
\(R\) by a constant times
\(\|\nabla F\|_\infty+\|F\|_\infty^2\).  Those quantities are uniform on a
bounded \(C^2\) set.  Taking norms proves
\eqref{eq:v3v90-plaquette-expansion}.  A frame change conjugates the
holonomy, so neither side depends on the radial gauge used in the proof.
\end{proof}

\subsection{A global admissibility gap without a small-link gauge}
\label{sec:v3v90-gap}

V78 proved a gap on a contractible small-link chart.  The next lemma removes
that chart restriction.  It is qualitative but uniform in lattice volume
and permits topologically nontrivial almost-flat shift systems.

\begin{lemma}[Almost-commuting Wilson gap]
\label{lem:v3v90-almost-commuting-gap}
Fix \(r>0\) and \(I\Subset(0,2r)\).  There are
\(\varepsilon_*>0\) and \(g_*>0\), depending only on \(r\) and \(I\),
such that every finite-dimensional collection of unitary shifts
\(T_1,\ldots,T_4\) satisfying
\begin{equation}
 \max_{\mu,\nu}\|T_\mu T_\nu-T_\nu T_\mu\|<\varepsilon_*
\label{eq:v3v90-almost-commuting}
\end{equation}
obeys
\begin{equation}
 \widehat H(m)^2\ge g_*^2,1
 \qquad\text{for every }m\in I.
\label{eq:v3v90-uniform-gap}
\end{equation}
No assertion that the \(T_\mu\) are close to commuting unitaries is used.
\end{lemma}

\begin{proof}
Suppose the assertion were false.  There would be finite-dimensional
Hilbert spaces \(\mathcal H_n\), unitaries \(T_{\mu,n}\), masses
\(m_n\in I\), and unit vectors \(v_n\) such that all commutators tend to
zero and
\(\|\widehat H_n(m_n)v_n\|\to0\).  Pass to a subsequence with
\(m_n\to m_\infty\in I\).  In the quotient C*-algebra
\begin{equation}
 \mathcal Q=
 \prod_n\mathcal B(\mathcal H_n)\Big/
 \bigoplus_n\mathcal B(\mathcal H_n),
\label{eq:v3v90-product-quotient}
\end{equation}
the classes \(\overline T_\mu\) commute exactly.  They therefore define a
unital representation of \(C(\mathbb T^4)\).  The class of
\(\widehat H_n(m_n)\) is the image of the free Wilson symbol
\begin{equation}
 h_m(p)=\gamma_5\left(
 i\sum_\mu\gamma_\mu\sin p_\mu
 +r\sum_\mu(1-\cos p_\mu)-m\right)
\label{eq:v3v90-free-symbol}
\end{equation}
at \(m=m_\infty\); replacing \(m_n\) by \(m_\infty\) changes its
representative by a norm-null sequence.

Now
\begin{equation}
 h_m(p)^2=
 \sum_\mu\sin^2p_\mu+
 \left(r\sum_\mu(1-\cos p_\mu)-m\right)^2.
\label{eq:v3v90-free-square}
\end{equation}
It can vanish only when every \(p_\mu\) is \(0\) or \(\pi\) and
\(m=2rk\) for some integer \(0\le k\le4\).  Compactness of
\(I\times\mathbb T^4\) thus supplies a positive lower bound.  Hence the
class of \(\widehat H_n(m_n)\) is invertible in \(\mathcal Q\).

Choose a bounded representative \(B_n\) of its inverse.  Then
\(\|B_n\widehat H_n-1\|\to0\).  For large \(n\), this norm is below
\(1/2\), whence
\(\|\widehat H_nv\|\ge(2\sup_n\|B_n\|)^{-1}\|v\|\).  This contradicts
the choice of \(v_n\).  The contradiction simultaneously gives positive
\(\varepsilon_*\) and \(g_*\); otherwise one could choose a sequence with
\(\varepsilon=1/n\) and gap below \(1/n\).
\end{proof}

For covariant shifts,
\(T_\mu T_\nu-T_\nu T_\mu\) is a unitary translate times
\(P_{\mu\nu}-1\).  Lemma~\ref{lem:v3v90-almost-commuting-gap} and
Lemma~\ref{lem:v3v90-smooth-admissible} therefore give the following.

\begin{theorem}[Admissible-sector Wilson gap]
\label{thm:v3v90-admissible-gap}
If \(\varepsilon_a<\varepsilon_*\), then the Wilson kernel
\eqref{eq:v3v90-Hhat} has the volume-uniform gap
\eqref{eq:v3v90-uniform-gap}.  In particular, every geometrically sampled
connection in a bounded smooth family has this gap for all sufficiently
small \(a\), uniformly over the family.
\end{theorem}

\begin{proof}
The commutator identity just noted verifies
\eqref{eq:v3v90-almost-commuting}.  Apply
Lemma~\ref{lem:v3v90-almost-commuting-gap}.  Uniformity for a smooth family
then follows from \eqref{eq:v3v90-plaquette-expansion}.
\end{proof}

\begin{proposition}[Homotopy constancy on a gapped component]
\label{prop:v3v90-homotopy-index}
Along every norm-continuous path of shifts for which
\eqref{eq:v3v90-uniform-gap} holds,
\(\operatorname{ind}_{\rm ov}\) is constant.  It is invariant under all
lattice gauge transformations.
\end{proposition}

\begin{proof}
The negative spectral projection
\(P_-(s)=1_{(-\infty,0)}(\widehat H(s))\) is norm-continuous by its Riesz
formula.  Its finite rank is an integer-valued continuous function and is
therefore constant.  Since
\(-\tfrac12\operatorname{Tr}\operatorname{sgn}\widehat H\) differs from
\(\operatorname{rank}P_-\) by half the fixed dimension of the lattice
spin space, it is constant as well.  Gauge transformations conjugate
\(\widehat H\), so they leave its sign trace unchanged.
\end{proof}

\subsection{The complete signed local-index expansion}
\label{sec:v3v90-local-index}

We now compute the component selected by the gap.  Denote by
\(d\operatorname{vol}_X\,\mathcal I_R(A)\) the degree-four form
\begin{equation}
 [\operatorname{ch}(E_R,\nabla^E)]_4
 =\mathcal I_R(A)\,d\operatorname{vol}_X.
\label{eq:v3v90-continuum-density-flat}
\end{equation}
The flat tangent bundle contributes \(\widehat A(TX)=1\).  The convention
in \eqref{eq:v3v90-continuum-density-flat} fixes all signs and factors of
\(2\pi\); no separate curvature-component convention is required.

\begin{lemma}[Uniform covariant sign-kernel expansion]
\label{lem:v3v90-sign-expansion}
Let \(\mathscr K\) be bounded in \(C^6\), and sample its connections by
edge holonomy.  Uniformly for \(A\in\mathscr K\), lattice sites \(x\), and
\(m\in I\),
\begin{equation}
 q_a(x)=a^4\nu(m,r)\mathcal I_R(A)(x)+a^5\rho_a(A,x),
 \qquad
 \sup_{a,A,x}|\rho_a(A,x)|<\infty,
\label{eq:v3v90-local-index-expansion}
\end{equation}
where
\begin{equation}
 \nu(m,r)=
 -\frac12\sum_{k=0}^4{4\choose k}(-1)^k
 \operatorname{sgn}(2rk-m).
\label{eq:v3v90-degree}
\end{equation}
In the one-species window \(0<m<2r\), \(\nu=1\).
\end{lemma}

\begin{proof}
We give the analytic steps because the order of cancellation is essential.
The uniform gap from Theorem~\ref{thm:v3v90-admissible-gap} permits the
norm-convergent formula
\begin{equation}
 \operatorname{sgn}\widehat H_a=
 \frac2\pi\int_0^\infty
 \widehat H_a(\widehat H_a^2+s^2)^{-1}\,ds.
\label{eq:v3v90-sign-integral}
\end{equation}
For \(s\le1\), every resolvent is bounded by
\((g_*^2+s^2)^{-1}\); for \(s\ge1\), move one factor of
\(\widehat H_a\) through the resolvent and use the bound
\(C(1+s^2)^{-1}\).  These estimates and their first six covariant
commutators are integrable uniformly in \(a\).

In a synchronous frame at \(x\), exact edge transport has the Taylor
formula, on a smooth test section,
\begin{equation}
 T_{\mu,a}=
 \sum_{j=0}^{5}\frac{a^j}{j!}\nabla_\mu^j+a^6R_{\mu,a},
 \qquad
 \|R_{\mu,a}f\|_{L^\infty}
 \le C_{\mathscr K}\|f\|_{C^6}.
\label{eq:v3v90-shift-Taylor}
\end{equation}
Insert this formula into the resolvent identity in
\eqref{eq:v3v90-sign-integral}.  Iterating the identity through order five
gives a finite covariant difference-symbol parametrix.  Its coefficients
are universal polynomials in \(F\) and its covariant derivatives, multiplied
by rational functions of the free symbol
\eqref{eq:v3v90-free-symbol}.  The preceding integrable bounds control the
integral remainder uniformly.  Taking the diagonal amounts to integrating
the free momentum variable over \([-\pi,\pi]^4\).  Thus
\begin{equation}
 q_a(x)=\sum_{j=0}^4a^jQ_j(A,x)+a^5\rho_a(A,x),
\label{eq:v3v90-parametrix-series}
\end{equation}
with the asserted uniform bound.  This derivation is local, so synchronous
frames on overlapping bundle charts give the same scalar coefficients.

It remains to identify the \(Q_j\) before estimating them.  The spin trace
with \(\gamma_5\) vanishes unless four distinct Clifford generators occur.
Consequently \(Q_0=Q_1=Q_2=Q_3=0\).  At order four, covariance and the
Bianchi identity leave, modulo an exact divergence, only the invariant
four-form \([\operatorname{ch}(E_R)]_4\).  An exact divergence cannot
occur here: the diagonal sign trace is gauge invariant at each site, while
a gauge-invariant local three-form primitive for the second Chern character
does not exist.  Hence
\begin{equation}
 Q_4(A,x)=c(m,r)\mathcal I_R(A)(x).
\label{eq:v3v90-Q4}
\end{equation}

The coefficient is found before taking any absolute value.  In the momentum
integral it is the pullback of the normalized volume form on \(S^4\) by
\begin{equation}
 p\longmapsto
 \frac{(\sin p_1,\ldots,\sin p_4,
 r\sum_\mu(1-\cos p_\mu)-m)}
 {\left[\sum_\mu\sin^2p_\mu+
 (r\sum_\mu(1-\cos p_\mu)-m)^2\right]^{1/2}}.
\label{eq:v3v90-Wilson-map}
\end{equation}
Its degree is the oriented corner sum
\eqref{eq:v3v90-degree}: a corner with \(k\) entries equal to \(\pi\)
has orientation \((-1)^k\) and multiplicity \({4\choose k}\).  This is
the V76 degree computation, now inserted into the interacting covariant
parametrix.  Therefore \(c=\nu\), and V76 gives \(\nu=1\) on
\(0<m<2r\).  Substitution into
\eqref{eq:v3v90-parametrix-series} proves
\eqref{eq:v3v90-local-index-expansion}.
\end{proof}

The proof displays the precise audit lesson: the separate corners contribute
the signed binomial sum.  Bounding them separately would retain a sum of
absolute multiplicities and would not recover \(\nu\).

\begin{firewall}[Scope of the parametrix lemma]
\label{fw:v3v90-parametrix-scope}
Lemma~\ref{lem:v3v90-sign-expansion} is proved for exact holonomy sampling
of a \(C^6\)-bounded smooth connection on a flat torus.  It is not a claim
about arbitrary rough admissible links.  Rough links have the gap and an
integer overlap charge by Theorem~\ref{thm:v3v90-admissible-gap}, but
identifying their continuum bundle requires an admissible interpolation or
a separate geometric reconstruction theorem.
\end{firewall}

\subsection{Exact index equality by integer stabilization}
\label{sec:v3v90-index-equality}

\begin{theorem}[Geometric overlap index theorem on the flat four-torus]
\label{thm:v3v90-overlap-index}
Let \(E_R\to X\) and \(\nabla^E\) be as above, and let the lattice links be
their exact edge holonomies.  For every fixed
\(m\in(0,2r)\), there is \(a_0>0\) such that, for every cubical refinement
with \(a<a_0\),
\begin{equation}
 \operatorname{ind}_{\rm ov}(a,E_R)
 =\int_X[\operatorname{ch}(E_R,\nabla^E)]_4
 =\operatorname{ind}D^+_{E_R}.
\label{eq:v3v90-index-theorem}
\end{equation}
The choice of bundle frames and of representatives of the torus spin
transition functions does not affect the equality.
\end{theorem}

\begin{proof}
For small \(a\), Theorem~\ref{thm:v3v90-admissible-gap} defines the left
side.  Sum \eqref{eq:v3v90-local-index-expansion} over the
\(O(a^{-4})\) sites.  The leading term is a Riemann sum and the uniform
remainder is \(O(a)\), so
\begin{equation}
 \operatorname{ind}_{\rm ov}(a,E_R)
 =\int_X[\operatorname{ch}(E_R,\nabla^E)]_4+O(a).
\label{eq:v3v90-index-limit}
\end{equation}
The first quantity is an integer by V77.  The integral is the integer
\(\operatorname{ind}D^+_{E_R}\) by the Atiyah--Singer index theorem on the
flat spin torus.  Two integers whose difference tends to zero are equal once
the difference has modulus below one.  This proves exact equality for all
sufficiently fine meshes.

A frame change conjugates every shift and hence the sign kernel.  Changing a
spin transition representative does the same after the corresponding
unitary identification of lattice spinors.  The trace and both indices are
therefore unchanged.
\end{proof}

\begin{corollary}[Uniform compact-family equality]
\label{cor:v3v90-family-index}
If the connections range over a compact \(C^6\)-bounded parameter family on
a fixed bundle, one value of \(a_0\) works for the whole family.  The
overlap index is constant on each connected family component and equals the
continuum index there.
\end{corollary}

\begin{proof}
The admissibility, gap, and remainder constants in
Lemmas~\ref{lem:v3v90-smooth-admissible} and
\ref{lem:v3v90-sign-expansion} are uniform on the compact family.  Choose
\(a_0\) so that the error in \eqref{eq:v3v90-index-limit} is below one.
The equality then holds at every parameter.  Homotopy constancy also follows
directly from Proposition~\ref{prop:v3v90-homotopy-index}.
\end{proof}

\begin{corollary}[Standard Model representation sum]
\label{cor:v3v90-SM-index}
For every smooth principal
\(G_{\rm SM}=S(U(3)\times U(2))\)-bundle on the flat spin four-torus, the
sum of the sufficiently fine overlap indices of all Standard Model Weyl
representations equals the corresponding sum of continuum chiral indices.
This statement respects the \(\mathbb Z_6\) quotient because each summand is
formed from an honest associated bundle of \(G_{\rm SM}\).
\end{corollary}

\begin{proof}
Apply Theorem~\ref{thm:v3v90-overlap-index} to the finite list of associated
bundles and take the minimum of their mesh thresholds.  Add the equalities.
No lift to \(SU(3)\times SU(2)\times U(1)\) is chosen, so the construction
descends through the physical quotient by definition.
\end{proof}

This corollary is an index comparison, not anomaly cancellation by itself.
The latter uses the degree-six family index density and the V66
representation arithmetic.  What has now been supplied is the missing
topological identification of the finite Wilson sign on every smooth flat
bundle sector reached by geometric refinement.

\subsection{Consequences for V88--V89 patching}
\label{sec:v3v90-patching}

\begin{proposition}[Separation and recombination of the three spectral tasks]
\label{prop:v3v90-three-task-recombination}
On a compact smooth family of flat \(G_{\rm SM}\)-backgrounds, for all
sufficiently fine lattices:
\begin{enumerate}
 \item physical zero modes and their finite-band transition lines are
 compared by V89's generalized norm-resolvent theorems;
 \item the massive Wilson kernel is uniformly gapped and its overlap index
 is the continuum chiral index by Theorem~\ref{thm:v3v90-overlap-index};
 and
 \item neither conclusion requires a mode-by-mode identification between
 the physical operator and the Wilson kernel.
\end{enumerate}
Thus the numerical index mismatch cannot obstruct the V88 determinant-line
patching on such a family.
\end{proposition}

\begin{proof}
Item 1 is Theorem~\ref{thm:v3v89-finite-cover}; item 2 is the uniform form
of Theorem~\ref{thm:v3v90-overlap-index}.  The two proofs use different
spectral regimes, as required by Proposition~\ref{prop:v3v89-three-spectra}.
Their indices agree with the same continuum Fredholm index, so the ranks of
the kernel and cokernel enter the determinant line with the correct virtual
difference.  This removes an index mismatch as an obstruction, while making
no assertion about equality of individual eigenvectors.
\end{proof}

\begin{firewall}[An integer index is not a family index bundle]
\label{fw:v3v90-family-bundle-open}
Corollary~\ref{cor:v3v90-family-index} identifies the integer index on each
fibre.  It does not yet identify the higher Chern character of a family
index bundle, its determinant connection, or its mapping-torus holonomy.
That refinement requires the expansion
\eqref{eq:v3v90-local-index-expansion} with parameter-space differentials
inserted before the supertrace.
\end{firewall}

\begin{firewall}[The gravitational density has not been discretized]
\label{fw:v3v90-gravity-open}
On a curved spin four-manifold the target is
\begin{equation}
 \operatorname{ind}D^+_{E_R}
 =\int_X[\widehat A(TX)\operatorname{ch}(E_R)]_4.
\label{eq:v3v90-curved-target}
\end{equation}
V90 proves only the flat-tangent contribution.  Recovering
\([\widehat A(TX)]_4\operatorname{rank}E_R\) requires a covariant lattice
spin transport and a curved pseudodifference parametrix.  Replacing the
flat shifts by chartwise coordinate differences without that spin transport
would miss precisely the Einstein-sector term.
\end{firewall}

\begin{firewall}[Dynamical rough fields remain outside the comparison]
\label{fw:v3v90-rough-open}
The uniform gap holds for every link field satisfying the stated
plaquette bound, but the local-index identification was proved only for
smooth holonomy samples.  A continuum measure-concentration theorem which
puts the dynamical gauge field into reconstructible admissible sectors has
not been proved.  This is distinct from the smooth-background regulator
comparison settled here.
\end{firewall}

\subsection{V90 conclusion and V91 target}
\label{sec:v3v90-conclusion}

The massive Wilson branch now has its own comparison theorem.  Almost
commuting covariant shifts give a volume-uniform gap without choosing a
small-link gauge.  On exact holonomy samples, the complete signed sign-kernel
has a local expansion whose coefficient is the oriented Wilson degree.
In the one-species window that degree is one, and integer stabilization
turns convergence into exact equality with the continuum chiral index on
all sufficiently fine flat bundle meshes.  This discharges V89's regulator
topology gate in the stated smooth flat sector.

V91 should go upstream to the geometric discretization needed by the full
SM--Einstein continuum.  It should construct chart-independent spin and
gauge edge transports on a bounded-geometry curved mesh, prove a curved
Wilson gap and diagonal parametrix with coefficient
\([\widehat A(TX)\operatorname{ch}(E_R)]_4\), and then add parameter-space
insertions to identify the family determinant connection.  If a local
coordinate stencil cannot reproduce the spin-curvature term, the correct
repair is to change the discrete spin transport before attempting sharper
resolvent estimates.

\clearpage
\section{V91: the flat-background family index and determinant curvature}
\label{sec:v3v91}

V90 identifies the integer overlap index on each smooth flat-background
bundle sector.  A determinant line contains more information than that
integer: its curvature is the degree-two part of a family index, and its
holonomy can retain torsion information even when the curvature vanishes.
This note inserts parameter-space derivatives into V90's complete Wilson
supertrace.  It proves convergence of the overlap family Chern character
and stabilization of the nontorsion first Chern class.  The calculation
also recovers the flat-gauge part of the degree-six polynomial whose
Standard Model sum was cancelled in V66.  Torsion holonomy and curved
vertical geometry are left separate rather than being inferred from a
curvature calculation.

\subsection{The relative overlap index bundle}
\label{sec:v3v91-index-bundle}

Let \(B\) be a compact smooth parameter manifold and let
\(\mathbb E\to X\times B\) be a Hermitian bundle with a smooth unitary
connection \(\mathbb A\), where \(X\) is the flat spin four-torus of V90.
The restriction at \(b\in B\) is denoted \((E_b,A_b)\).  Assume a bounded
set of derivatives of \(\mathbb A\) large enough for all expansions below.
Restricting \(\mathbb E\) to \(X_a\times B\) gives the finite-dimensional
Hermitian bundle
\begin{equation}
 \mathcal H_a=\bigoplus_{x\in X_a}S_x\otimes\mathbb E_{(x,\cdot)}
 \longrightarrow B
\label{eq:v3v91-Hbundle}
\end{equation}
with its induced connection \(\nabla^{\mathcal H_a}\).  Exact edge
holonomy in the \(X\)-directions defines the family
\(\widehat H_a(b)\) of \eqref{eq:v3v90-Hhat}.  V90 gives a common gap for
all sufficiently small \(a\).

Let
\begin{equation}
 P_a=1_{(-\infty,0)}(\widehat H_a),
 \qquad
 P_a^{\rm ref}=\frac{1+\gamma_5}{2}.
\label{eq:v3v91-projectors}
\end{equation}
The second projector is the negative projector of the on-site kernel
\(-m\gamma_5\).  Define the virtual overlap family index
\begin{equation}
 \operatorname{Ind}_{\rm ov,a}
 =[\operatorname{ran}P_a]-[\operatorname{ran}P_a^{\rm ref}]
 \in K^0(B).
\label{eq:v3v91-virtual-index}
\end{equation}
Its rank is the V90 overlap index because
\begin{equation}
 \operatorname{rank}P_a-\frac12\operatorname{rank}\mathcal H_a
 =-\frac12\operatorname{Tr}\operatorname{sgn}\widehat H_a.
\label{eq:v3v91-rank-index}
\end{equation}

For a smooth projection \(P\) in a Hermitian bundle \(\mathcal H\), the
Grassmann connection and curvature are
\begin{equation}
 \nabla^P=P\nabla^{\mathcal H}P,
 \qquad
 \mathcal F_P=P\mathcal F_{\mathcal H}P
 +P(\nabla^{\mathcal H}P)^2P.
\label{eq:v3v91-Grassmann-curvature}
\end{equation}
We use the relative Chern form
\begin{equation}
 \operatorname{Ch}_a=
 \operatorname{Tr}\exp\left(\frac{i\mathcal F_{P_a}}{2\pi}\right)
 -\operatorname{Tr}\exp\left(
 \frac{i\mathcal F_{P_a^{\rm ref}}}{2\pi}\right).
\label{eq:v3v91-relative-Ch}
\end{equation}
Including the ambient-curvature term in
\eqref{eq:v3v91-Grassmann-curvature} is essential when
\(\mathbb E\) is not pulled back from one fixed bundle on \(X\).

\begin{proposition}[Smoothness, covariance, and closedness]
\label{prop:v3v91-basic-family}
For sufficiently small \(a\), \(P_a\) and \(\operatorname{Ch}_a\) are
smooth on \(B\).  The form \(\operatorname{Ch}_a\) is closed, its de Rham
class is \(\operatorname{ch}(\operatorname{Ind}_{\rm ov,a})\), and all
three objects are invariant under changes of frames of \(\mathbb E\).
\end{proposition}

\begin{proof}
The common spectral gap permits one fixed Riesz contour and hence smooth
functional calculus for \(P_a\).  Formula
\eqref{eq:v3v91-Grassmann-curvature} is the curvature of the induced
connection.  The Bianchi identity and cyclicity of the trace show that each
homogeneous part of \eqref{eq:v3v91-relative-Ch} is closed.  Chern--Weil
theory identifies its class with the difference of the two range-bundle
Chern characters.  A frame change conjugates the kernels, projectors, and
connections, leaving the traces invariant.
\end{proof}

\subsection{Resolvent differentiation before the diagonal trace}
\label{sec:v3v91-resolvent}

Write \(\boldsymbol\nabla\) for the connection on operator-valued forms
induced by \(\nabla^{\mathcal H_a}\).  Differentiating the sign function
under its gapped integral gives
\begin{align}
 \boldsymbol\nabla\operatorname{sgn}\widehat H_a
 &=\frac2\pi\int_0^\infty
 \bigl[(\boldsymbol\nabla\widehat H_a)R_s
 -\widehat H_aR_s
 (\boldsymbol\nabla\widehat H_a^2)R_s\bigr],ds,
 \label{eq:v3v91-sign-derivative}\\
 R_s&=(\widehat H_a^2+s^2)^{-1}.
\end{align}
Repeated derivatives are finite sums of local insertions separated by
resolvents.  They must remain in this ordered sum until after the spin trace;
estimating the summands separately would repeat the audit error identified
in V90.

\begin{lemma}[Uniform differentiated locality]
\label{lem:v3v91-differentiated-locality}
For every fixed \(j\), if \(\mathbb A\) has bounded parameter derivatives
through order \(j\), then
\(\boldsymbol\nabla^jP_a\) has an exponentially decaying lattice kernel
with constants uniform in \(a\), after distances are measured in lattice
units.  Products containing a fixed number of such derivatives have a
uniformly summable first spatial moment.
\end{lemma}

\begin{proof}
The finite-range kernel \(\widehat H_a\) has a common dimensionless gap.
Conjugating it by an exponential lattice weight changes its norm by
\(O(\alpha)\), uniformly in \(a\).  The Combes--Thomas resolvent argument
therefore bounds \(R_s(x,y)\) by
\(C(g_*^2+s^2)^{-1}e^{-c\sqrt{g_*^2+s^2}d_a(x,y)}\).
Every parameter derivative of an edge holonomy is supported on the same
edge and is uniformly bounded on the given family.  Formula
\eqref{eq:v3v91-sign-derivative} and its iterates are convolutions of a
fixed number of these local insertions and resolvents.  The convolution of
exponential kernels remains exponential.  Multiplication by
\(1+d_a(x,y)\) is still summable after decreasing the exponent, proving the
moment assertion.
\end{proof}

The first moment is the exact analogue of the NS marking cost.  It is what
allows a parameter insertion to be transported from one endpoint of a
short lattice path to the base point of the local density with one explicit
factor of \(a\).

\subsection{The family local-index expansion}
\label{sec:v3v91-family-expansion}

Decompose differential forms on \(X\times B\) by vertical and horizontal
degree.  Fibre integration selects vertical degree four.  Let
\begin{equation}
 \mathcal I_B=
 \int_X\operatorname{ch}(\mathbb E,\mathbb A)
 \in\Omega^{\rm even}(B).
\label{eq:v3v91-family-target}
\end{equation}
Because \(TX\) is flat, this is the families-index form
\(\int_X\widehat A(TX)\operatorname{ch}(\mathbb E)\).

\begin{theorem}[Flat family overlap local-index theorem]
\label{thm:v3v91-family-local-index}
Fix an even degree \(2\ell\le\dim B\).  For a sufficiently smooth bounded
family and \(m\in I\Subset(0,2r)\),
\begin{equation}
 \bigl[\operatorname{Ch}_a\bigr]_{2\ell}
 =\left[\int_X\operatorname{ch}(\mathbb E,\mathbb A)\right]_{2\ell}
 +d_B\beta_{a,2\ell-1}+\mathcal R_{a,2\ell},
\label{eq:v3v91-family-expansion}
\end{equation}
where \(\beta_{a,2\ell-1}\) is a globally defined Chern--Simons
transgression depending on the chosen connections and
\begin{equation}
 \|\mathcal R_{a,2\ell}\|_{C^0(B)}\le C_\ell a.
\label{eq:v3v91-family-remainder}
\end{equation}
If the Grassmann and continuum index connections are compared in the same
covariant symbol convention, \(\beta_a\) may be absorbed into that
connection choice and the difference itself is \(O(a)\).
\end{theorem}

\begin{proof}
Use \eqref{eq:v3v91-sign-derivative} to express every occurrence of
\(\boldsymbol\nabla P_a\) and
\eqref{eq:v3v91-Grassmann-curvature} to retain the ambient curvature.
Before taking a norm, collect the complete relative trace at each site.
In synchronous product gauge at \((x,b)\), vertical edge transport and
parameter differentiation have covariant Taylor expansions.  Iterated
resolvent identities produce a finite difference-symbol parametrix through
vertical order five.  Lemma~\ref{lem:v3v91-differentiated-locality}
controls its remainders and, in particular, the one displacement created by
moving a horizontal insertion to \((x,b)\).  Consequently the site density
of its degree-\(2\ell\) part has the uniform form
\begin{equation}
 a^4\Omega_{4,2\ell}(\mathbb F)(x,b)
 +a^5\rho_{a,2\ell}(x,b),
 \qquad \sup|\rho_{a,2\ell}|<\infty,
\label{eq:v3v91-site-family-expansion}
\end{equation}
where \(\mathbb F\) is the curvature of \(\mathbb A\).

The \(\gamma_5\) trace again requires four vertical Clifford factors.
Thus every lower vertical order vanishes only after all Wilson corners and
all terms of \eqref{eq:v3v91-sign-derivative} have been summed.  Chern--Weil
invariance identifies the surviving closed polynomial, up to an exact
transgression, with
\([\operatorname{ch}(\mathbb E)]_{(4,2\ell)}\).  Its multiplicative
coefficient is determined at one point by freezing the horizontal
curvatures and evaluating the vertical momentum integral.  That integral
is the same degree of the Wilson map as in
\eqref{eq:v3v90-Wilson-map}; hence it equals
\(\nu(m,r)=1\).  The reference projector subtracts the vertical-degree-zero
ambient contribution and has no Wilson degree.

Summing \eqref{eq:v3v91-site-family-expansion} over
\(O(a^{-4})\) sites gives the fibre integral plus an \(O(a)\) remainder.
Two Chern--Weil representatives of the same invariant polynomial differ by
the exterior derivative of their Chern--Simons transgression; this is the
form \(d_B\beta_a\) in \eqref{eq:v3v91-family-expansion}.  Uniformity on
the compact parameter space proves \eqref{eq:v3v91-family-remainder}.
\end{proof}

For \(\ell=0\), this theorem reduces to V90.  For \(\ell=1\), it is the
determinant-curvature comparison:
\begin{equation}
 \frac{i}{2\pi}\operatorname{Tr}
 (\mathcal F_{P_a}-\mathcal F_{P_a^{\rm ref}})
 \longrightarrow
 \left[\int_X\operatorname{ch}(\mathbb E,\mathbb A)\right]_2
 =\int_X[\operatorname{ch}(\mathbb E,\mathbb A)]_6
\label{eq:v3v91-det-curvature-limit}
\end{equation}
after the displayed exact transgression is included in the connection
comparison.

\subsection{Integral stabilization and its torsion boundary}
\label{sec:v3v91-integral-stabilization}

Let
\begin{equation}
 \mathcal L_a=
 \det(\operatorname{ran}P_a)\otimes
 \det(\operatorname{ran}P_a^{\rm ref})^{-1},
 \qquad
 \mathcal L_{\rm cont}=\det\operatorname{Ind}D^+_{\mathbb E/B}.
\label{eq:v3v91-det-lines}
\end{equation}

\begin{theorem}[Stabilization of determinant Chern numbers]
\label{thm:v3v91-Chern-stabilization}
Assume \(B\) has the homotopy type of a finite CW complex.  For all
sufficiently small \(a\),
\begin{equation}
 c_1(\mathcal L_a)\otimes1
 =c_1(\mathcal L_{\rm cont})\otimes1
 \quad\hbox{in }H^2(B;\mathbb R).
\label{eq:v3v91-real-c1}
\end{equation}
Equivalently, the two integral first Chern classes can differ only by
torsion.  If \(H^2(B;\mathbb Z)\) is torsion-free, the determinant line
bundles are topologically isomorphic for all sufficiently small \(a\).
\end{theorem}

\begin{proof}
Choose integral two-cycles \(\Sigma_1,\ldots,\Sigma_N\) whose classes span
the free part of \(H_2(B;\mathbb Z)\).  The exact transgression in
\eqref{eq:v3v91-family-expansion} integrates to zero.  The uniform
\(O(a)\) remainder shows that the difference of the two curvature periods
on every \(\Sigma_j\) tends to zero.  Each normalized period is an integer,
namely the pairing of a first Chern class with \([\Sigma_j]\).  It is
therefore exactly zero for small \(a\).  There are only finitely many
generators, so one mesh threshold works for all of them.  This proves
\eqref{eq:v3v91-real-c1}.  The kernel of
\(H^2(B;\mathbb Z)\to H^2(B;\mathbb R)\) is its torsion subgroup, proving
the second statement.  If that subgroup is zero, equality of first Chern
classes classifies the complex line bundles.
\end{proof}

\begin{firewall}[Curvature does not detect a flat torsion line]
\label{fw:v3v91-torsion}
No curvature argument can remove the possible torsion difference in
Theorem~\ref{thm:v3v91-Chern-stabilization}.  A flat line can have zero
curvature and nontrivial finite holonomy.  Its comparison requires loop
transport, equivalently the suspended five-dimensional index or eta phase.
V87 proves the continuum Standard Model bordism character is trivial; a
finite-lattice family comparison must still transport that result to the
overlap line.
\end{firewall}

\subsection{The Standard Model degree-six cancellation after discretization}
\label{sec:v3v91-SM-cancellation}

Let \(\mathbb E_{R_f}\) range over all left-handed Standard Model fermion
representations, with right-handed fields written as left-handed conjugates.
Tensor products and direct sums are taken as honest representations of
\(G_{\rm SM}=S(U(3)\times U(2))\).

\begin{corollary}[Vanishing limiting determinant curvature]
\label{cor:v3v91-SM-curvature}
For the complete Standard Model representation, including a gauge-singlet
right-handed neutrino when that field is present in the chosen matter
content,
\begin{equation}
 \sum_f\left[
 \int_X\operatorname{ch}(\mathbb E_{R_f},\mathbb A)
 \right]_2=0.
\label{eq:v3v91-SM-six-zero}
\end{equation}
After the exact Chern--Simons comparison in
Theorem~\ref{thm:v3v91-family-local-index}, the complete finite overlap
determinant curvature is \(O(a)\) uniformly on every compact smooth family.
Its periods vanish exactly on the free part of \(H_2(B;\mathbb Z)\) for all
sufficiently small \(a\).
\end{corollary}

\begin{proof}
On the flat vertical tangent bundle, the degree-six form contains the cubic
nonabelian, mixed nonabelian--abelian, and abelian cubic gauge-anomaly
coefficients.  V66 computes their complete representation sum and finds
each coefficient zero.  The mixed gauge--gravitational term is absent here
because the vertical Pontryagin form is zero; its cancellation was computed
in V66 but its lattice realization awaits the curved extension.  A
gauge-singlet neutrino has trivial gauge Chern character and changes none of
the coefficients present here.  Equation \eqref{eq:v3v91-SM-six-zero}
follows.
Sum \eqref{eq:v3v91-family-expansion} over the finite representation list;
the leading term cancels before a norm is taken, leaving the uniform
\(O(a)\) remainder.  The integral-period statement follows from the same
integer stabilization used in
Theorem~\ref{thm:v3v91-Chern-stabilization}.
\end{proof}

\begin{proposition}[Compatibility with the V81 and V86 transgressions]
\label{prop:v3v91-transgression-compatibility}
On a contractible smooth parameter patch, choose the Chern--Simons
representative \(\beta_{a,1}\) in
\eqref{eq:v3v91-family-expansion}.  Subtracting it from the finite overlap
determinant connection leaves its curvature class unchanged and makes the
complete Standard Model curvature \(O(a)\).  Under the common exponential
locality bounds of Lemma~\ref{lem:v3v91-differentiated-locality}, this
subtraction is a local field-space transgression of the type used in V81,
and its remainder is summable across the V86 refinement tower.
\end{proposition}

\begin{proof}
Changing a connection by the one-form \(-\beta_{a,1}\) changes its curvature
by \(-d_B\beta_{a,1}\) and does not change its line-bundle Chern class.
Theorem~\ref{thm:v3v91-family-local-index} and
Corollary~\ref{cor:v3v91-SM-curvature} then give an \(O(a)\) curvature.
The Chern--Simons formula is an integral of a polynomial in the two local
connections and their curvatures.  Lemma
\ref{lem:v3v91-differentiated-locality} makes each kernel convolution
exponentially local with a summable first moment.  At scale
\(a_n=q^{-n}a_0\), the remainder is bounded by \(Cq^{-n}a_0\), whose sum is
finite.  This is precisely the positive irrelevant power isolated in V86.
\end{proof}

\begin{firewall}[Local asymptotic gauge invariance is not an exact lattice Ward solution]
\label{fw:v3v91-Ward}
The \(O(a)\) result and vanishing periods prove convergence and remove the
nontorsion topological obstruction.  They do not construct an exactly
gauge-invariant finite-spacing Weyl measure on every admissible component.
That requires the finite local Ward primitive left open in V78--V79, or an
equivalent exact equivariant trivialization.
\end{firewall}

\begin{firewall}[Vertical curvature is still flat]
\label{fw:v3v91-vertical-flat}
The vertical tangent bundle of \(X\) remains flat, so the mixed
gauge--gravitational term does not appear in the lattice calculation of this
note.  The degree-four gravitational factor
\([\widehat A(TX)]_4\) in the rank index and the full curved family form
have not been obtained from a lattice spin operator.
\end{firewall}

\subsection{V91 conclusion and V92 target}
\label{sec:v3v91-conclusion}

The overlap comparison now reaches the determinant curvature.  The
relative negative-band bundle has a covariant Chern form, differentiated
resolvents retain uniform locality and the required first moment, and the
complete Wilson supertrace converges to the flat families-index form.  Its
determinant Chern numbers stabilize exactly; only torsion holonomy lies
beyond curvature.  In the complete Standard Model sum, the degree-six term
cancels before estimation and leaves the \(O(a)\) remainder required by the
V86 scale recursion.

V92 should build the missing vertical geometric representative rather than
write a coordinate stencil and assume it glues.  The correct upstream task
is to construct a stabilized Bott--Wilson symbol over \(T^*X\), prove that
spin-frame changes act through gapped degree-one overlap homotopies, and
quantize that global symbol to a uniformly local finite regulator.  Its
K-theory pairing must give
\(\int_X[\widehat A(TX)\operatorname{ch}(E)]_4\); only after this symbol
gluing is explicit should one seek a curved Combes--Thomas and diagonal
parametrix estimate.

\clearpage
\section{V92: the global Bott--Wilson symbol on a curved spin manifold}
\label{sec:v3v92}

The flat family theorem of V91 cannot be transported to a curved manifold
by writing the same hypercubic stencil in every orthonormal frame.  A general
\(SO(4)\) change of frame does not preserve the momentum torus or its nearest
neighbours, so the putative local operators do not satisfy a Cech gluing
law.  This is a structural defect, not a resolvent-estimate defect.

The repair is to identify what the one-species Wilson symbol represents
before discretizing it.  This note proves that it is the spin Bott--Thom
class.  That class is exactly Spin(4)-equivariant and therefore globalizes
over \(T^*X\).  Its pairing with a gauge bundle is the curved chiral index,
including the \(\widehat A\)-term.  A semiclassical pseudodifferential
quantization gives a global uniformly local Fredholm regulator.  Producing
from it a finite lattice matrix with fixed-range locality is retained as a
separate analytic step.

\subsection{The equivariant Bott involution}
\label{sec:v3v92-Bott}

Let \(V\) be an oriented Euclidean four-space, let \(S=S^+\oplus S^-\) be
its complex spin module, and choose Hermitian Clifford matrices
\(\Gamma_1,\ldots,\Gamma_4,\Gamma_5\).  On the one-point compactification
\(V^+\cong S^4\), define
\begin{equation}
 b(v)=\frac{1}{1+|v|^2}
 \left(2\sum_{j=1}^4v_j\Gamma_j+(|v|^2-1)\Gamma_5\right),
 \qquad b(\infty)=\Gamma_5.
\label{eq:v3v92-Bott-involution}
\end{equation}

\begin{lemma}[Bott involution and equivariance]
\label{lem:v3v92-Bott-equivariance}
The map \(b:V^+\to\operatorname{End}(S)\) is a smooth self-adjoint
involution.  If \(s\in\operatorname{Spin}(V)\), then
\begin{equation}
 b(sv)=\rho(s)b(v)\rho(s)^{-1}.
\label{eq:v3v92-Bott-equivariance}
\end{equation}
The negative projector
\begin{equation}
 p_B(v)=\frac{1-b(v)}2
\label{eq:v3v92-Bott-projector}
\end{equation}
equals the constant projector \((1-\Gamma_5)/2\) at infinity and its
reduced class is the Bott generator of \(K_c^0(V)\).
\end{lemma}

\begin{proof}
The Clifford relations make the square of the numerator in
\eqref{eq:v3v92-Bott-involution} equal
\begin{equation}
 4|v|^2+(|v|^2-1)^2=(1+|v|^2)^2,
\end{equation}
so \(b^*=b\) and \(b^2=1\).  Dividing numerator and denominator by
\(|v|^2\) proves smooth extension with value \(\Gamma_5\) at infinity.
The spin representation obeys
\(\rho(s)c(v)\rho(s)^{-1}=c(sv)\) and preserves \(\Gamma_5\), which proves
\eqref{eq:v3v92-Bott-equivariance}.

On the equatorial sphere, the clutching map from the positive to the
negative eigenspace is \(c(v)/|v|:S^3\to U(2)\).  Under
\(\operatorname{Spin}(4)=SU(2)_+\times SU(2)_-\), this is the identity
quaternion map and has winding number one.  The clutching construction
identifies winding number one with the generator of
\(\widetilde K^0(S^4)=K_c^0(V)\).
\end{proof}

\subsection{The Wilson corner is the Bott class}
\label{sec:v3v92-Wilson-Bott}

For \(m\in(0,2r)\), let \(n_W:T^4\to S^4\) be the normalized Wilson map
\eqref{eq:v3v90-Wilson-map}, and let
\(p_W=(1-n_W\cdot\Gamma)/2\).

\begin{theorem}[One-species Wilson--Bott identification]
\label{thm:v3v92-Wilson-Bott}
There is an open ball \(U\Subset T^4\) containing \(p=0\) and a gapped
homotopy of projector symbols which makes \(p_W\) constant on
\(T^4\setminus U\).  After collapsing \(T^4\setminus U\) to a point, its
reduced class is the Bott generator:
\begin{equation}
 [p_W]-[p_W(\infty)]=[p_B]-[p_B(\infty)]
 \quad\hbox{in }K_c^0(\mathbb R^4).
\label{eq:v3v92-Wilson-Bott-class}
\end{equation}
The identification is stable under every change of \(m\) inside a compact
subinterval of \((0,2r)\).
\end{theorem}

\begin{proof}
The function
\(f(p)=r\sum_\mu(1-\cos p_\mu)\) has a unique critical point below its
first positive critical value \(2r\), namely its nondegenerate minimum at
the origin.  Choose
\begin{equation}
 m<c<2r,
 \qquad U=\{p:f(p)<c\}.
\label{eq:v3v92-Wilson-ball}
\end{equation}
Morse theory below the first critical value makes \(U\) an open four-ball.
On its complement the fifth component \(f(p)-m\) is at least \(c-m>0\).
There we may multiply the first four components of the unnormalized Wilson
vector by \(1-t\), \(0\le t\le1\), without meeting zero.  A collar
cutoff makes this contraction constant outside a slightly smaller ball and
fixed near the origin.  Normalization then gives a homotopy through
self-adjoint involutions.

The resulting based map \(U/\partial U\cong S^4\to S^4\) has degree
\(\nu=1\) by Theorem~\ref{thm:v3v76-Wilson-degree}.  Based homotopy classes
of maps \(S^4\to S^4\) are classified by degree.  The map underlying
\eqref{eq:v3v92-Bott-involution} also has degree one.  The two projector
maps are therefore homotopic after the collapse, which proves
\eqref{eq:v3v92-Wilson-Bott-class}.  On a compact mass subinterval the free
symbol has a common gap; the construction may be chosen continuously in
\(m\), proving the last statement.
\end{proof}

\begin{remark}[Why the full class, rather than only an index number, is fixed]
Both projector maps in Theorem~\ref{thm:v3v92-Wilson-Bott} are constant
off a four-ball.  Their reduced classes therefore lie in
\(\widetilde K^0(S^4)\cong\mathbb Z\), where the degree is a complete
invariant.  This is stronger than comparing only the top Chern number of two
arbitrary bundles over \(T^4\), the gap noted in V84.
\end{remark}

\subsection{Globalization by the spin structure}
\label{sec:v3v92-globalization}

Let \((X,g)\) now be an arbitrary closed oriented Riemannian spin
four-manifold, with spinor bundle \(S_X\).  Apply
\eqref{eq:v3v92-Bott-projector} fibrewise to \(T^*X\).  Equation
\eqref{eq:v3v92-Bott-equivariance} makes the local expressions agree under
spin-frame transitions.

\begin{theorem}[Global Bott--Wilson Thom class]
\label{thm:v3v92-global-Thom}
The fibrewise negative projector defines a canonical class
\begin{equation}
 \beta_X\in K_c^0(T^*X),
\label{eq:v3v92-betaX}
\end{equation}
the complex spin Thom class.  For a Hermitian gauge bundle
\(E_R\to X\),
\begin{equation}
 \beta_X(E_R)=\beta_X\smile\pi^*[E_R]
 \in K_c^0(T^*X)
\label{eq:v3v92-twisted-Thom}
\end{equation}
is independent of the orthonormal atlas and of all local Wilson
representatives.  In every sufficiently small cotangent chart, a
one-species hypercubic Wilson projector is stably gapped-homotopic to this
class.
\end{theorem}

\begin{proof}
On an overlap, two oriented orthonormal frames differ by an \(SO(4)\)-valued
map.  The spin structure supplies a lift to \(\operatorname{Spin}(4)\), and
\eqref{eq:v3v92-Bott-equivariance} says that the two local Bott projectors
are conjugate by the spin transition function.  Hence they descend to a
global compactly supported K-class.  The clutching calculation in
Lemma~\ref{lem:v3v92-Bott-equivariance} identifies its restriction to every
cotangent fibre with the Bott generator, which is the defining property of
the spin Thom class.  Tensoring by \(\pi^*E_R\) gives
\eqref{eq:v3v92-twisted-Thom} and plainly commutes with gauge transitions.

Theorem~\ref{thm:v3v92-Wilson-Bott} supplies a stable gapped homotopy from a
local Wilson projector to the Bott projector.  On a chart overlap, perform
the first homotopy, use the exact Spin(4) conjugacy of the Bott endpoint,
and reverse the homotopy in the other frame.  This gives the asserted local
stable equivalence without attempting to rotate the hypercubic momentum
torus itself.
\end{proof}

\begin{proposition}[Homotopy-coherent overlap data]
\label{prop:v3v92-coherent-descent}
After one finite stabilization on a finite good cover of \(X\), the local
Wilson-to-Bott homotopies can be chosen with coherent homotopies on double,
triple, and higher overlaps.  Their descended K-class is uniquely
\(\beta_X(E_R)\).
\end{proposition}

\begin{proof}
Take the global Bott projector as the common endpoint on every chart.  The
space of sufficiently stabilized projector representatives of a fixed
K-class is a classifying space for complex K-theory.  Equality of the two
endpoint classes gives paths on double overlaps; equality of the induced
loops gives fillings after the next stabilization.  Continue over the
finite nerve of the good cover.  Only finitely many skeleta occur, so a
single finite stabilization accommodates all fillings.  Because every
simplex is based at the same global Bott projector, the resulting descent
class restricts fibrewise to the Bott generator and is therefore the Thom
class established in Theorem~\ref{thm:v3v92-global-Thom}.  Different choices
give homotopic classifying maps and the same K-class.
\end{proof}

This proposition is an existence statement in stable topology.  It does not
yet give a fixed hopping range.  V84's trigonometric approximation proves a
fixed-range path only after a common flat momentum torus has been chosen;
the present construction explains how that local result must be organized
on a curved atlas.

\subsection{Index pairing and the Einstein term}
\label{sec:v3v92-index-pairing}

Let \(\pi:T^*X\to X\) be the projection and
\(\pi_!:K_c^0(T^*X)\to K^0(\mathrm{pt})\cong\mathbb Z\) the spin
K-theory pushforward.

\begin{theorem}[Curved Bott--Wilson index]
\label{thm:v3v92-curved-index}
The topological index of the global Bott--Wilson class is
\begin{equation}
 \pi_!\bigl(\beta_X(E_R)\bigr)
 =\operatorname{ind}D^+_{E_R}
 =\int_X[\widehat A(TX,g)\operatorname{ch}(E_R,A)]_4.
\label{eq:v3v92-curved-index}
\end{equation}
For a compact parameter space \(B\) and a smooth family of spin metrics and
gauge bundles, the corresponding family class satisfies
\begin{equation}
 \operatorname{ch}\operatorname{Ind}(D^+_{\mathbb E/B})
 =\int_{X/B}widehat A(T(X/B))\operatorname{ch}(\mathbb E).
\label{eq:v3v92-curved-family-index}
\end{equation}
\end{theorem}

\begin{proof}
The symbol of the chiral spin Dirac operator is Clifford multiplication
\(c(\xi):S^+\to S^-\).  Its clutching class is precisely the fibrewise
class \(\beta_X\) of
Theorem~\ref{thm:v3v92-global-Thom}.  Twisting the operator by \(E_R\)
multiplies the symbol by \(\pi^*[E_R]\), giving
\(\beta_X(E_R)\).  The analytic index of an elliptic operator equals the
K-theory pushforward of its symbol.  The cohomological spin pushforward is
the Atiyah--Singer formula, which gives
\eqref{eq:v3v92-curved-index}.  Applying the same symbol and pushforward
construction fibrewise over \(B\) gives the families index theorem
\eqref{eq:v3v92-curved-family-index}.
\end{proof}

The degree-four part expands invariantly as
\begin{equation}
 [\widehat A(TX)\operatorname{ch}(E_R)]_4
 =[\operatorname{ch}(E_R)]_4
 -\frac{\operatorname{rank}E_R}{24}p_1(TX),
\label{eq:v3v92-Ahat-expansion}
\end{equation}
with the Chern--Weil conventions implicit in \(\widehat A\) and
\(\operatorname{ch}\).  The second term is exactly the vertical
gravitational contribution missing from V90--V91.

\begin{corollary}[Correct curved target for the Standard Model regulator]
\label{cor:v3v92-SM-curved-target}
Any curved finite Wilson regulator which is locally stably equivalent to the
one-species symbol and whose overlap data descend through
Proposition~\ref{prop:v3v92-coherent-descent} has the symbol class
\(\beta_X(E_R)\).  Its Fredholm index, once analytic quantization preserves
that class, is forced to be \eqref{eq:v3v92-curved-index}; no independent
choice of the coefficient of \(p_1(TX)\) remains.
\end{corollary}

\begin{proof}
Stable gapped homotopy preserves the negative-projector K-class.  The local
class is the Bott generator by
Theorem~\ref{thm:v3v92-Wilson-Bott}, and coherent descent makes it the global
Thom class by Theorem~\ref{thm:v3v92-global-Thom}.  The index depends only on
that symbol class, so Theorem~\ref{thm:v3v92-curved-index} determines both
terms in \eqref{eq:v3v92-Ahat-expansion}.
\end{proof}

\subsection{A global semiclassical quantization}
\label{sec:v3v92-quantization}

Choose a smooth fibre-radial representative of the Bott projector which is
exactly equal to its reference projector outside \(|\xi|\le R\).  With a
finite atlas and partition of unity, let
\(Q_a=\operatorname{Op}_a(p_B\otimes1_{E_R})\) be its covariant
semiclassical quantization.  Symmetrize it so that \(Q_a^*=Q_a\).

\begin{theorem}[Global pseudodifferential Bott--Wilson regulator]
\label{thm:v3v92-psido-regulator}
For sufficiently small \(a\), functional calculus at the circle
\(|z-1|=1/3\) produces an exact projection
\begin{equation}
 P_a^{BW}=\frac{1}{2\pi i}
 \int_{|z-1|=1/3}(z-Q_a)^{-1}\,dz.
\label{eq:v3v92-exact-projector}
\end{equation}
It has the following properties:
\begin{enumerate}
 \item it is globally defined, spin- and gauge-covariant, and its principal
 symbol is \(p_B\otimes1_{E_R}\);
 \item the pair \((P_a^{BW},P^{\rm ref})\) is Fredholm and its relative
 index is \(\operatorname{ind}D^+_{E_R}\);
 \item for every \(N\), away from the diagonal its semiclassical kernel
 obeys
 \begin{equation}
  \|P_a^{BW}(x,y)-P^{\rm ref}(x,y)\|
  \le C_Na^{-4}\left(1+\frac{d_g(x,y)}a\right)^{-N};
 \label{eq:v3v92-rapid-locality}
 \end{equation}
 and
 \item these statements are uniform on a compact bounded-geometry family of
 metrics and connections.
\end{enumerate}
\end{theorem}

\begin{proof}
Symbol composition gives
\(Q_a^2-Q_a=aR_a\) with \(R_a\) uniformly bounded in the order-zero
semiclassical calculus.  The semiclassical Calderon--Vaillancourt estimate
gives \(\|Q_a^2-Q_a\|\le Ca\).  Hence every real spectral value \(\lambda\)
satisfies \(|\lambda(\lambda-1)|\le Ca\), so for small \(a\) the spectrum
lies in disjoint neighbourhoods of zero and one and avoids the stated
contour.  The Riesz formula then defines an exact projection.  Holomorphic
functional calculus preserves the principal symbol, proving item 1.

Since the projector symbol equals the reference at fibre infinity, the
operator pair is Fredholm.  The relative index of two pseudodifferential
projections equals the topological index of their difference symbol.
That symbol is \(\beta_X(E_R)\), so item 2 follows from
Theorem~\ref{thm:v3v92-curved-index}.

In each chart the kernel of a compactly supported smooth semiclassical
symbol is an oscillatory integral.  Applying repeatedly the integration-by-
parts operator
\begin{equation}
 L=\frac{1-i\langle\exp_x^{-1}y,\partial_\xi\rangle/a}
 {1+d_g(x,y)^2/a^2}
\label{eq:v3v92-IBP-operator}
\end{equation}
leaves its phase invariant and gains one factor
\((1+d_g(x,y)/a)^{-1}\).  All differentiated amplitudes are uniformly
bounded, which proves \eqref{eq:v3v92-rapid-locality} for \(Q_a\).
The parameter-dependent pseudodifferential functional calculus gives the
same estimate for its Riesz projection.  A finite bounded-geometry atlas,
uniform symbol seminorms, and a common contour make every constant uniform
on the stated family, proving items 3 and 4.
\end{proof}

\begin{firewall}[Rapid locality is not fixed-range lattice locality]
\label{fw:v3v92-fixed-range}
Theorem~\ref{thm:v3v92-psido-regulator} gives decay faster than every power
of the distance in lattice units.  It is an infinite-dimensional Fredholm
pair, not yet a finite matrix on a mesh, and it does not give the fixed
hopping radius used in V76's reflection-positive transfer construction.
Truncating its kernel at distance \(Ra\log(1/a)\) gives a norm-small
finite-propagation approximation, but the range in lattice units grows.
Therefore neither finite-lattice reflection positivity nor a fixed-range
claim follows from this theorem.
\end{firewall}

\begin{firewall}[Stable Cech descent is not an explicit stencil]
\label{fw:v3v92-stencil}
Proposition~\ref{prop:v3v92-coherent-descent} proves that all spin-frame
overlap obstructions vanish after finite stabilization.  It does not supply
numerical hopping coefficients on an arbitrary triangulation.  Those
coefficients must be obtained by a quantization and compression whose error
is controlled in operator norm, so that the gapped K-class cannot change.
\end{firewall}

\subsection{V92 conclusion and V93 target}
\label{sec:v3v92-conclusion}

The curved topology is now fixed without an invalid rotation of a
hypercubic stencil.  The one-species Wilson projector is the Bott generator;
Spin(4)-equivariance globalizes it to the Thom class on \(T^*X\); and its
gauge-twisted pairing gives the full curved index
\(\int\widehat A\operatorname{ch}\).  A global semiclassical regulator with
rapid off-diagonal decay realizes this class analytically.  Thus the
Einstein coefficient is no longer an undetermined lattice counterterm.

V93 should construct a finite-dimensional compression of
\(P_a^{BW}\) on a quasiuniform bounded-geometry mesh.  It must prove an
operator-norm approximation for the projection and its first parameter
derivatives, preserve the relative index through a gap, and quantify the
locality radius.  If a fixed lattice-unit range is incompatible with
spin-frame descent, the note must state the minimal growing range and test
whether V85--V86 summability and V76 reflection positivity survive it,
rather than hiding that geometric cost in a pointwise consistency bound.

\clearpage
\section{V93: finite-mesh compression of the curved Bott--Wilson regulator}
\label{sec:v3v93}

V92 produces a global curved regulator as a Fredholm pair of rapidly local
pseudodifferential projections.  The present note compresses that pair to a
finite-dimensional geometric mesh while preserving its index.  Two scales
must be distinguished.  The symbol scale \(a\) is the regulator length; the
mesh spacing \(h\) resolves its kernel.  A fixed small ratio \(h/a\) is
enough for the exact index.  Convergence of determinant connections requires
\(h/a\to0\) and a growing truncation radius in lattice units.  An explicit
diagonal schedule remains summable in the V85--V86 refinement tower, but it
does not satisfy V76's fixed-time-range reflection-positive hypothesis.

\subsection{Geometric finite-element frames}
\label{sec:v3v93-FE}

Let \((X,g)\) range over a compact bounded-geometry family of closed spin
four-manifolds, and let \(E_R\) have a bounded smooth family of unitary
connections.  Choose quasiuniform shape-regular geodesic triangulations
\(\mathcal T_h\) with mesh diameter \(h\).  For each vertex \(z\), let
\(\phi_z\) be the scalar continuous piecewise-linear hat function.  Radial
spin--gauge transport from \(z\) over the star of \(z\) turns
\(\phi_zv\), \(v\in(S\otimes E_R)_z\), into a local section.  Let
\(W_h\) be the span of these sections and \(\Pi_h\) its \(L^2\)-orthogonal
projection.

The construction is fibrewise under changes of spin and gauge frames.  It
also commutes with chirality:
\begin{equation}
 [\Pi_h,\gamma_5]=0,
 \qquad
 [\Pi_h,P^{\rm ref}]=0.
\label{eq:v3v93-chiral-FE}
\end{equation}

\begin{lemma}[Uniform approximation and inverse bounds]
\label{lem:v3v93-FE-estimates}
For every integer \(1\le s\le s_*\), the spaces above may be enriched to
fixed polynomial degree \(s\) so that
\begin{align}
 \|(1-\Pi_h)u\|_{L^2}
 &\le C_sh^s\|u\|_{H^s},
 \label{eq:v3v93-FE-approx}\\
 \|u_h\|_{H^s}
 &\le C_sh^{-s}\|u_h\|_{L^2},
 \qquad u_h\in W_h.
\label{eq:v3v93-FE-inverse}
\end{align}
The constants are uniform in \(h\), in the bounded-geometry family, and in
local spin--gauge frames.
\end{lemma}

\begin{proof}
On every rescaled reference simplex, polynomial interpolation and the
Bramble--Hilbert lemma give \eqref{eq:v3v93-FE-approx}; finite-dimensional
norm equivalence gives \eqref{eq:v3v93-FE-inverse}.  Shape regularity makes
the affine rescaling constants uniform.  Radial parallel transport has
uniformly bounded derivatives on a bounded-geometry family, so inserting it
changes only the constants.  A finite-overlap sum over vertex stars proves
the global estimates.  Frame changes are unitary and therefore leave all
norms unchanged.
\end{proof}

The nodal frame is not orthonormal.  Its Gram matrix has a uniformly bounded
condition number after the standard \(h^{-4}\) mass scaling and has a fixed
graph bandwidth.  We use the orthonormalized frame obtained from its inverse
square root.

\begin{lemma}[Localized orthonormal frame]
\label{lem:v3v93-local-frame}
There are \(C,c>0\), uniform in \(h\), such that the coefficients of the
inverse square root of the scaled Gram matrix satisfy
\begin{equation}
 |(G_h^{-1/2})_{zw}|\le Ce^{-c d_{\mathcal T}(z,w)}.
\label{eq:v3v93-Gram-decay}
\end{equation}
Consequently the resulting orthonormal basis of \(W_h\) is exponentially
localized in the mesh graph distance.
\end{lemma}

\begin{proof}
After scaling, the spectrum of \(G_h\) lies in one compact interval
\([c_0,C_0]\subset(0,\infty)\), by quasiuniformity and shape regularity.
Approximate \(x^{-1/2}\) uniformly on this interval by degree-\(n\)
Chebyshev polynomials with error \(Cq^n\), \(q<1\).  A degree-\(n\)
polynomial in a fixed-bandwidth matrix has zero \((z,w)\) entry whenever
the graph distance exceeds a constant times \(n\).  Choose \(n\)
proportional to that distance and use the uniform approximation error.  This
proves \eqref{eq:v3v93-Gram-decay}; multiplying the nodal functions by this
matrix gives the localized orthonormal frame.
\end{proof}

\subsection{Compression of the Fredholm involution}
\label{sec:v3v93-compression}

Let \(P_a^{BW}\) be the V92 projection and set
\begin{equation}
 S_a=1-2P_a^{BW},
 \qquad S_0=1-2P^{\rm ref},
 \qquad K_a=S_a-S_0.
\label{eq:v3v93-involutions}
\end{equation}
The operator \(K_a\) is a covariant smoothing operator concentrated at
frequencies \(|\xi|\lesssim a^{-1}\), and
\(S_a^2=S_0^2=1\).

\begin{lemma}[Two-scale compression estimate]
\label{lem:v3v93-compression-estimate}
For every fixed \(s\le s_*\),
\begin{equation}
 \|(1-\Pi_h)K_a\|+
 \|K_a(1-\Pi_h)\|
 \le C_s\left(\frac ha\right)^s,
 \qquad0<h\le c a.
\label{eq:v3v93-compression-error}
\end{equation}
The same estimate holds for a fixed number of derivatives over a compact
parameter family, with a possibly larger constant.
\end{lemma}

\begin{proof}
The semiclassical symbol bounds in V92 imply
\begin{equation}
 \|K_au\|_{H^s}\le C_sa^{-s}\|u\|_{L^2},
 \qquad
 \|K_a^*u\|_{H^s}\le C_sa^{-s}\|u\|_{L^2}.
\label{eq:v3v93-smoothing-bound}
\end{equation}
Apply \eqref{eq:v3v93-FE-approx} to \(K_au\) for the first term.  Apply
the same argument to \(K_a^*=K_a\) and take adjoints for the second.  A
parameter derivative of the symbol has the same semiclassical order and
uniform seminorm bounds, so the differentiated estimates follow
identically.
\end{proof}

On \(W_h\) define the compressed self-adjoint Hamiltonian
\begin{equation}
 A_{a,h}=\Pi_hS_a\Pi_h\big|_{W_h}.
\label{eq:v3v93-compressed-H}
\end{equation}

\begin{theorem}[Finite compression preserves the curved index]
\label{thm:v3v93-compression-index}
There is \(\theta_*>0\) such that, whenever \(h/a<\theta_*\),
\(A_{a,h}\) has a gap at zero and
\begin{equation}
 \operatorname{rank}1_{(-\infty,0)}(A_{a,h})
 -\operatorname{rank}P^{\rm ref}\big|_{W_h}
 =\operatorname{ind}D^+_{E_R}.
\label{eq:v3v93-finite-index}
\end{equation}
The threshold is uniform on compact bounded-geometry families.
\end{theorem}

\begin{proof}
On the full \(L^2\) space form the block-diagonal operator
\begin{equation}
 \widetilde S_{a,h}=
 \Pi_hS_a\Pi_h+(1-\Pi_h)S_0(1-\Pi_h).
\label{eq:v3v93-block-S}
\end{equation}
Because \(\Pi_h\) commutes with \(S_0\), the difference
\(S_a-\widetilde S_{a,h}\) is the sum of the two off-diagonal blocks of
\(K_a\) and its high--high block.  Lemma
\ref{lem:v3v93-compression-estimate} gives
\begin{equation}
 \|S_a-\widetilde S_{a,h}\|
 \le3C_s(h/a)^s.
\label{eq:v3v93-block-close}
\end{equation}
Choose \(\theta_*\) so the right side is below \(1/2\).  The straight path
from the involution \(S_a\) to \(\widetilde S_{a,h}\) then stays gapped by
at least \(1/2\).  Its relative negative-projection index against \(S_0\)
is constant.  The high block of \(\widetilde S_{a,h}\) equals the reference
and contributes zero, while the low block is \(A_{a,h}\).  V92 identifies
the original Fredholm-pair index with \(\operatorname{ind}D^+_{E_R}\),
which proves \eqref{eq:v3v93-finite-index}.  Uniformity follows from the
uniform constant in \eqref{eq:v3v93-compression-error}.
\end{proof}

This argument proves the exact integer equality directly by a gapped
homotopy.  It does not use convergence followed by rounding and therefore
also retains torsion information when applied to a parameter family whose
compressed projectors form a bundle.

\subsection{Finite propagation by controlled truncation}
\label{sec:v3v93-truncation}

Write matrices in the localized orthonormal frame of
Lemma~\ref{lem:v3v93-local-frame}.  For \(L>1\), delete the entries of
\(A_{a,h}-S_0|_{W_h}\) between basis centres whose Riemannian distance
exceeds \(La\), and symmetrize the result.  Denote the finite matrix by
\(A_{a,h}^{(L)}\).

\begin{lemma}[Weighted Schur tail]
\label{lem:v3v93-Schur-tail}
For every \(N>4\),
\begin{equation}
 \|A_{a,h}^{(L)}-A_{a,h}\|
 \le C_NL^{4-N}+C_Ne^{-cLa/h}.
\label{eq:v3v93-tail-bound}
\end{equation}
The same estimate holds for a fixed number of parameter derivatives.
\end{lemma}

\begin{proof}
Equation \eqref{eq:v3v92-rapid-locality}, integrated against two normalized
finite-element basis functions of support \(O(h)\), gives matrix entries
bounded by
\begin{equation}
 C_N(h/a)^4
 \left(1+\frac{d_g(z,w)}a\right)^{-N},
\label{eq:v3v93-matrix-decay}
\end{equation}
up to the exponential tails caused by Gram orthonormalization.  A
quasiuniform four-dimensional mesh has at most
\(C(a/h)^4t^3dt\) centres in a shell of scaled radius \(t=d_g/a\).
The row sum outside \(t=L\) is therefore bounded by
\(C_N\int_L^\infty t^{3-N}dt=C'_NL^{4-N}\).  The Gram tails give the second
term in \eqref{eq:v3v93-tail-bound}.  The same column bound holds; the Schur
test proves the operator estimate.  Parameter derivatives satisfy the same
kernel bounds by V92 and Lemma
\ref{lem:v3v93-compression-estimate}.
\end{proof}

\begin{theorem}[Finite-dimensional finite-propagation curved regulator]
\label{thm:v3v93-finite-regulator}
Choose \(h/a<\theta_*\) with strict margin in
\eqref{eq:v3v93-block-close}, and then choose \(L\) so the right side of
\eqref{eq:v3v93-tail-bound} is smaller than half that margin.  The finite
matrix \(A_{a,h}^{(L)}\) is spin- and gauge-covariant, has propagation at
most \(La+O(h)\), is uniformly gapped, and satisfies
\begin{equation}
 \operatorname{rank}1_{(-\infty,0)}(A_{a,h}^{(L)})
 -\operatorname{rank}P^{\rm ref}|_{W_h}
 =\int_X[\widehat A(TX)\operatorname{ch}(E_R)]_4.
\label{eq:v3v93-finite-curved-index}
\end{equation}
Its sign projector is exponentially local in the mesh graph metric.
\end{theorem}

\begin{proof}
Truncation by geodesic distance is invariant under changes of fibre frames,
so covariance is preserved.  The construction makes propagation explicit.
The tail estimate and spectral stability preserve the gap and give a gapped
straight path back to \(A_{a,h}\); hence their relative indices agree.
Apply Theorem~\ref{thm:v3v93-compression-index} and then
Theorem~\ref{thm:v3v92-curved-index}.  Finally, a finite-propagation
self-adjoint matrix with a uniform gap has an exponentially local sign by
the Combes--Thomas argument of V78.
\end{proof}

For the exact index alone, \(h/a=\theta\) and \(L\) may be fixed once and
for all.  The propagation radius in mesh units is then the fixed number
\begin{equation}
 R_{\rm lat}\le C\frac{La}{h}=C\frac L\theta.
\label{eq:v3v93-fixed-index-range}
\end{equation}
Thus curvature does not by itself force an unbounded hopping range for the
integer regulator.

\subsection{Parameter derivatives and a summable diagonal schedule}
\label{sec:v3v93-diagonal}

To compare determinant connections rather than only index classes, the
compression and tail errors must tend to zero.  Let
\(\mathcal P_{a,h,L}=1_{(-\infty,0)}(A_{a,h}^{(L)})\).

\begin{theorem}[Family projection comparison]
\label{thm:v3v93-family-comparison}
On every compact smooth parameter family and for \(j=0,1,2\), after the
natural finite-element embedding,
\begin{equation}
 \|\boldsymbol\nabla^j
 (\mathcal P_{a,h,L}-P_a^{BW})\|
 \le C_j\left[
 \left(\frac ha\right)^s+L^{4-N}+e^{-cLa/h}
 \right].
\label{eq:v3v93-family-projector-error}
\end{equation}
The finite determinant connection and curvature satisfy the conservative
trace estimate
\begin{equation}
 \|\mathcal A_{a,h,L}-\mathcal A_a^{BW}\|+
 \|\mathcal F_{a,h,L}-\mathcal F_a^{BW}\|
 \le C a^{-4}\left[
 \left(\frac ha\right)^s+L^{4-N}+e^{-cLa/h}
 \right].
\label{eq:v3v93-family-trace-error}
\end{equation}
\end{theorem}

\begin{proof}
The block comparison in \eqref{eq:v3v93-block-close} and the truncation
bound \eqref{eq:v3v93-tail-bound} hold with two parameter derivatives.
Differentiate the Riesz formula for the negative spectral projection.  Each
derivative is a finite contour integral of resolvents separated by one or
two differentiated Hamiltonians.  The common gap uniformly bounds every
resolvent, giving \eqref{eq:v3v93-family-projector-error}.  The determinant connection and curvature are finite traces of respectively
one and two projector derivatives.  The relative kernels are supported in
the semiclassical phase-space region \(|\xi|\lesssim a^{-1}\), whose trace
volume on a fixed four-manifold is \(O(a^{-4})\).  The weighted kernel
estimates and the Schur bound therefore turn the operator error into the
trace estimate \eqref{eq:v3v93-family-trace-error}.  This factor cannot be
dropped merely because the pointwise kernels are local.
\end{proof}

Let \(a_n=a_0\bar a_n\), where \(\bar a_n=q^{-n}\).  A conservative
diagonal choice which pays the trace-volume factor is
\begin{equation}
 \frac{h_n}{a_n}=\bar a_n,
 \qquad L_n=\bar a_n^{-1/2},
 \qquad s\ge6,\quad N\ge16.
\label{eq:v3v93-diagonal-schedule}
\end{equation}
Then
\begin{equation}
 \sum_{n\ge1}\bar a_n^{-4}\left[
 (h_n/a_n)^s+L_n^{4-N}+e^{-cL_na_n/h_n}
 \right]<\infty.
\label{eq:v3v93-summable-error}
\end{equation}
Indeed the first two traced errors are at most fixed multiples of
\(\bar a_n^{s-4}\) and \(\bar a_n^{(N-4)/2-4}\), both bounded by
\(\bar a_n^2\) at the displayed thresholds, while the third is
\(\bar a_n^{-4}e^{-c\bar a_n^{-3/2}}\).  The physical propagation radius
tends to zero,
\begin{equation}
 L_na_n=a_0\bar a_n^{1/2}\longrightarrow0,
\label{eq:v3v93-physical-locality}
\end{equation}
while the mesh-unit radius grows as
\begin{equation}
 R_{{\rm lat},n}\asymp
 \frac{L_na_n}{h_n}=\bar a_n^{-3/2}.
\label{eq:v3v93-growing-range}
\end{equation}

\begin{corollary}[Compatibility with multiscale determinant transport]
\label{cor:v3v93-multiscale}
Along \eqref{eq:v3v93-diagonal-schedule}, replacing the V92 curved
Fredholm regulator by the finite matrix of
Theorem~\ref{thm:v3v93-finite-regulator} changes the total normalized V86
determinant connection by an absolutely convergent series.  It therefore
does not spoil the irrelevant-gain summability established in V86.
\end{corollary}

\begin{proof}
The trace estimate \eqref{eq:v3v93-family-trace-error} bounds each scale's
connection and curvature error by the summand in
\eqref{eq:v3v93-summable-error}.  The Weierstrass test gives absolute and
uniform convergence on the compact parameter family.  Adding V86's existing
\(O(a_n)\) remainder preserves convergence.
\end{proof}

\subsection{The exact locality tradeoff}
\label{sec:v3v93-tradeoff}

Within this compression construction, a uniformly bounded mesh-unit range
cannot make the right side of
\eqref{eq:v3v93-family-projector-error} tend to zero.  Indeed, convergence
requires \(h/a\to0\) and \(L\to\infty\), so
\begin{equation}
 R_{\rm lat}\simeq L(a/h)\longrightarrow\infty.
\label{eq:v3v93-range-necessity}
\end{equation}
This is a limitation of the geometric compression proof, not a universal
no-go theorem for every possible curved stencil.  A specially designed
finite stencil could evade it only by reproducing the global Bott descent
and its first two family derivatives directly, rather than by resolving and
truncating the pseudodifferential kernel.

\begin{firewall}[Reflection positivity is not inherited by compression]
\label{fw:v3v93-RP}
The matrices \(A_{a,h}^{(L)}\) are Hermitian, covariant, gapped, and of
finite propagation at every scale.  Those properties do not imply the
positive cross-reflection cone used in V76.  Along the connection-convergent
schedule their temporal range can grow as
\eqref{eq:v3v93-growing-range}, so V76's fixed one-step transfer
factorization cannot be invoked.  A reflection-positive construction must
either keep the temporal regulator ultralocal while performing the geometric
compression only in spatial directions, or prove a new positive
factorization for the growing temporal band.
\end{firewall}

\begin{firewall}[Finite-dimensional does not mean a canonical lattice]
\label{fw:v3v93-canonical}
The finite matrices depend on a shape-regular triangulation, finite-element
degree, and truncation convention.  Their indices are independent of those
choices within the controlled gap, but their determinant connections differ
by local transgressions.  V86 normalization must be imposed before declaring
two such regulator choices equivalent in the continuum measure.
\end{firewall}

\subsection{V93 conclusion and V94 target}
\label{sec:v3v93-conclusion}

The global curved Bott--Wilson regulator now has finite-mesh representatives.
A two-scale finite-element compression, followed by a controlled geodesic
tail cut, preserves the full \(\widehat A\operatorname{ch}\) index through
an explicit gapped homotopy.  Fixed ratios give a fixed-range exact-index
regulator.  A diagonal schedule makes the first two family derivatives and
the multiscale determinant connection converge with an absolutely summable
error; its physical support shrinks even though its mesh-unit range grows.

V94 should address the remaining torsion holonomy rather than infer it from
the now-controlled curvature.  It should build the finite compressed
operator on a mapping torus, compare its mod-two or exponentiated eta
transport with the continuum suspended Dirac class, and use the V87 result
\(\Omega_5^{\rm Spin}(BG_{\rm SM})=0\) to trivialize the complete Standard
Model flat line.  The proof must distinguish an exact finite-lattice
holonomy from convergence to one; if only convergence follows, it should
state the phase error and show that its scale sum is finite.

\clearpage
\section{V94: exact determinant-line transport and torsion holonomy}
\label{sec:v3v94}

V91 identified determinant curvature only up to a possible flat torsion
line.  V93 contains more information than was used there: the finite curved
matrix and the global Bott--Wilson Fredholm regulator are joined by an
explicit uniform gapped homotopy.  Kato transport along that homotopy gives
an isomorphism of their determinant lines, including torsion.  Comparing the
two endpoint connections produces an exact transgression one-form.  After
subtracting it as part of the regulator normalization, finite and continuum
holonomies agree exactly on the smooth geometric sector.  Without that
subtraction, the phase is only close to the continuum phase, with the
summable error quantified below.

\subsection{Determinant transport along a gapped Fredholm path}
\label{sec:v3v94-Kato}

Let \(B\) be a compact parameter manifold.  Let \(S_s(b)\),
\(0\le s\le1\), be a smooth family of self-adjoint Fredholm operators with
\begin{equation}
 S_s(b)^2\ge g^2,
 \qquad
 P_s(b)=1_{(-\infty,0)}(S_s(b)).
\label{eq:v3v94-gapped-family}
\end{equation}
Assume \(P_s-P^{\rm ref}\) is trace-class in the relative determinant
calculus.  The V93 path is of this form after its finite block is embedded
and the complementary block is set equal to the reference.

Define the Kato generator and transport by
\begin{equation}
 Q_s=[\partial_sP_s,P_s],
 \qquad
 \partial_sU_s=Q_sU_s,
 \qquad U_0=1.
\label{eq:v3v94-Kato-ODE}
\end{equation}

\begin{lemma}[Kato intertwining]
\label{lem:v3v94-Kato}
The solution \(U_s\) is unitary and satisfies
\begin{equation}
 P_sU_s=U_sP_0.
\label{eq:v3v94-Kato-intertwining}
\end{equation}
It therefore induces an isomorphism of relative index bundles and their
determinant lines
\begin{equation}
 T_s:\mathcal L_0\longrightarrow\mathcal L_s.
\label{eq:v3v94-line-isomorphism}
\end{equation}
This isomorphism preserves the complete integral K-class and hence cannot
lose torsion data.
\end{lemma}

\begin{proof}
Differentiating \(P_s^2=P_s\) gives
\(P_s\dot P_sP_s=0\).  Hence \(Q_s^*=-Q_s\), so the solution of
\eqref{eq:v3v94-Kato-ODE} is unitary.  A direct calculation gives
\begin{equation}
 [Q_s,P_s]=\dot P_s.
\end{equation}
Therefore
\(\partial_s(U_s^{-1}P_sU_s)=0\), proving
\eqref{eq:v3v94-Kato-intertwining}.  Exterior powers of the finite relative
subspaces, or equivalently the determinant functor on the Fredholm pair,
give \eqref{eq:v3v94-line-isomorphism}.  It is an actual bundle isomorphism,
not an equality only after applying the Chern character, so all integral and
torsion information is retained.
\end{proof}

Let \(\nabla^{\mathcal L_s}\) be the determinant connection induced by the
Grassmann connection of \(P_s\).  Pull the endpoint connection back by
\(T_1\) and define the one-form
\begin{equation}
 \alpha_{01}=T_1^*\nabla^{\mathcal L_1}
 -\nabla^{\mathcal L_0}\in\Omega^1(B;i\mathbb R).
\label{eq:v3v94-alpha}
\end{equation}

\begin{theorem}[Exact endpoint connection transgression]
\label{thm:v3v94-exact-transgression}
The connection
\begin{equation}
 \nabla^{\mathcal L_1,\rm corr}
 =\nabla^{\mathcal L_1}-(T_1^{-1})^*\alpha_{01}
\label{eq:v3v94-corrected-connection}
\end{equation}
satisfies
\begin{equation}
 T_1^*\nabla^{\mathcal L_1,\rm corr}
 =\nabla^{\mathcal L_0}.
\label{eq:v3v94-exact-connection-match}
\end{equation}
Consequently, for every closed loop \(\gamma\subset B\),
\begin{equation}
 \operatorname{Hol}_\gamma(
 \nabla^{\mathcal L_1,\rm corr})
 =\operatorname{Hol}_\gamma(\nabla^{\mathcal L_0}).
\label{eq:v3v94-exact-holonomy-match}
\end{equation}
Moreover
\begin{equation}
 d_B\alpha_{01}=T_1^*\mathcal F_1-\mathcal F_0.
\label{eq:v3v94-curvature-transgression}
\end{equation}
\end{theorem}

\begin{proof}
The difference of two unitary connections on the same line bundle is a
globally defined imaginary one-form, so
\eqref{eq:v3v94-alpha} is well defined.  Substitution gives
\eqref{eq:v3v94-exact-connection-match}.  Parallel transport is preserved
by a connection-preserving bundle isomorphism, which proves
\eqref{eq:v3v94-exact-holonomy-match}.  Taking curvatures of
\eqref{eq:v3v94-alpha} gives
\eqref{eq:v3v94-curvature-transgression}.
\end{proof}

An intrinsic formula follows by differentiating the Grassmann connection in
Kato gauge.  Up to the exact change caused by a different determinant-frame
choice,
\begin{equation}
 \alpha_{01}=
 \int_0^1\operatorname{Tr}_{\rm rel}
 \left(P_s[\partial_sP_s,\boldsymbol\nabla_B P_s]\right),ds.
\label{eq:v3v94-local-alpha}
\end{equation}
Its exterior derivative is the integrated projector-curvature variation,
in agreement with \eqref{eq:v3v94-curvature-transgression}.

\subsection{Locality and size of the V93 transgression}
\label{sec:v3v94-locality}

Use for \(S_s\) the concatenation of the two gapped straight paths in V93:
first from the global involution \(S_a\) to its block compression, then from
the compressed block to the finite-propagation truncation.  Smooth the
single joining parameter without changing the endpoints or gap.

\begin{proposition}[Covariant local transgression]
\label{prop:v3v94-local-transgression}
For the V93 path, \(\alpha_{01}\) is spin- and gauge-invariant and is a
local field-space one-form.  On a compact smooth family it obeys the
conservative bound
\begin{equation}
 \|\alpha_{01}\|_{C^0(B)}
 \le C a^{-4}\left[
 (h/a)^s+L^{4-N}+e^{-cLa/h}
 \right].
\label{eq:v3v94-alpha-bound}
\end{equation}
The same estimate holds for its first parameter derivative.  Along the
diagonal schedule \eqref{eq:v3v93-diagonal-schedule}, the series of these
norms converges.
\end{proposition}

\begin{proof}
Every operator on the path is covariant, so its Riesz projection, Kato
generator, and relative trace are covariant.  The trace removes conjugation,
proving invariance.  Differentiated Riesz formulas express
\(\partial_sP_s\) and \(\boldsymbol\nabla_BP_s\) as local perturbations
separated by uniformly gapped resolvents.  Their kernels have the weighted
decay used in Lemma~\ref{lem:v3v93-Schur-tail}; the relative trace has
effective phase-space volume \(O(a^{-4})\).  Insert the compression and
truncation differences and telescope the ordered products in
\eqref{eq:v3v94-local-alpha}.  This gives
\eqref{eq:v3v94-alpha-bound}.  One further parameter derivative is treated
by the twice-differentiated Riesz formula.  Summability is exactly
\eqref{eq:v3v93-summable-error}.
\end{proof}

Thus the correction in \eqref{eq:v3v94-corrected-connection} is not an
arbitrary choice of phase.  It is the rapidly local transgression generated
by the same gapped regulator homotopy which proves the finite index theorem.

\subsection{Separating continuum normalization from compression}
\label{sec:v3v94-normalization}

Let \(\mathcal L_a^{BW}\) carry its pseudodifferential Grassmann connection
\(\nabla^{BW,\rm Gr}\).  The symbol identification of V92 gives an
isomorphism \(I_a:\mathcal L_a^{BW}\to\mathcal L_{\rm cont}\), but it
does not assert equality of \(\nabla^{BW,\rm Gr}\) with the Bismut--Freed
connection.  Record their exact difference as
\begin{equation}
 \eta_a=I_a^*\nabla^{\rm BF}-\nabla^{BW,\rm Gr}
 \in\Omega^1(B;i\mathbb R).
\label{eq:v3v94-eta-normalization}
\end{equation}
This is the continuum regulator-normalization form.  The heat-kernel
transgression for two elliptic representatives of the same symbol expresses
\(\eta_a\), modulo an exact phase convention, as a convergent relative
eta-form.  Its short-time jet is local and contains the marginal
normalization data of V86.  Its large-time part depends on the global low
spectrum and is part of the continuum determinant connection; no locality or
smallness of the unnormalized \(\eta_a\) is assumed here.

Let \(T_1:\mathcal L_a^{BW}\to\mathcal L_{a,h,L}\) be the V93 Kato
isomorphism and let \(\nabla^{\rm fin,Gr}\) be the finite Grassmann
connection.  Define
\begin{align}
 \nabla^{\rm fin,norm}
 &=\nabla^{\rm fin,Gr}+(T_1^{-1})^*\eta_a,
 \label{eq:v3v94-finite-normalized}\\
 \alpha_{\rm comp}
 &=T_1^*\nabla^{\rm fin,Gr}-\nabla^{BW,\rm Gr}.
 \label{eq:v3v94-alpha-comp}
\end{align}
Then
\begin{equation}
 T_1^*\nabla^{\rm fin,norm}
 =I_a^*\nabla^{\rm BF}+\alpha_{\rm comp}.
\label{eq:v3v94-normalized-pullback}
\end{equation}
Thus Proposition~\ref{prop:v3v94-local-transgression} estimates precisely
\(\alpha_{\rm comp}\).  Subtracting its pushforward gives the fully
corrected finite connection.  The raw finite Grassmann connection also
contains the generally non-small form \(\eta_a\), so it must not be compared
directly with Bismut--Freed by the compression bound.

\subsection{Mapping-torus holonomy}
\label{sec:v3v94-mapping-torus}

Let \(\gamma\) be a loop in the quotient of smooth spin metrics and
\(G_{\rm SM}\)-connections by based gauge transformations and based spin
diffeomorphisms.  Its endpoints may be related by a large transformation.
Clutching the cylinder \(X\times[0,1]\) by that endpoint transformation
gives a closed spin five-manifold \(Y_\gamma\) with its
\(G_{\rm SM}\)-bundle.

For one chiral representation \(R\), the Dai--Freed formula for the
continuum determinant line is
\begin{equation}
 \operatorname{Hol}_\gamma(\nabla^{\mathcal L_{\rm cont},R})
 =\exp\bigl(2\pi i\,\xi(D_{Y_\gamma,R})\bigr)
 \exp\left(2\pi i\int_{X\times[0,1]}\omega_{5,R}\right),
\label{eq:v3v94-Dai-Freed}
\end{equation}
where \(d\omega_{5,R}\) is the degree-six local index density and the sign
convention is fixed by the chosen determinant connection.

\begin{theorem}[Finite compressed mapping-torus comparison]
\label{thm:v3v94-mapping-torus-comparison}
On every compact smooth loop family for which the V92--V93 bounded-geometry
hypotheses hold, the Kato isomorphism identifies the corrected finite
compressed determinant holonomy with
\eqref{eq:v3v94-Dai-Freed} exactly.  The continuum-normalized but compression-uncorrected holonomies satisfy
\begin{equation}
 \left|
 \operatorname{Hol}_\gamma(\nabla^{\rm fin,norm})
 \operatorname{Hol}_\gamma(\nabla^{\mathcal L_{\rm cont}})^{-1}-1
 \right|
 \le
 e^{\int_\gamma|\alpha_{\rm comp}|}-1.
\label{eq:v3v94-uncorrected-bound}
\end{equation}
Along \eqref{eq:v3v93-diagonal-schedule}, the logarithmic phase errors form
an absolutely summable series.
\end{theorem}

\begin{proof}
The global Bott--Wilson symbol is the symbol of the continuum chiral Dirac
family by V92, so \(I_a\) identifies its determinant line with the
continuum family-index determinant line.  The normalization
\eqref{eq:v3v94-eta-normalization} equips it with the pullback of the
Bismut--Freed connection, whose loop holonomy is
\eqref{eq:v3v94-Dai-Freed}.  Equation
\eqref{eq:v3v94-normalized-pullback} and
Theorem~\ref{thm:v3v94-exact-transgression} then give exact corrected
equality.

For the normalized but compression-uncorrected connection, the ratio of
holonomies is \(\exp(-\int_\gamma\alpha_{\rm comp})\).  The elementary inequality
\(|e^z-1|\le e^{|z|}-1\) proves
\eqref{eq:v3v94-uncorrected-bound}.  Proposition
\ref{prop:v3v94-local-transgression}, with
\(\alpha_{01}=\alpha_{\rm comp}\), and
\eqref{eq:v3v93-summable-error} prove the last assertion.
\end{proof}

\subsection{Exact disappearance of the Standard Model flat line}
\label{sec:v3v94-SM}

Take the tensor product of the determinant lines of the complete Standard
Model Weyl representation, using dual lines for conjugate chirality.  Denote
the resulting connections by \(\nabla^{\mathcal L_{\rm SM}}\).

\begin{theorem}[Exact corrected Standard Model holonomy]
\label{thm:v3v94-SM-holonomy}
On the smooth geometric sector described above, the corrected finite
compressed Standard Model determinant line has
\begin{equation}
 \operatorname{Hol}_\gamma(
 \nabla^{\mathcal L_{a,h,L,\rm SM},\rm corr})=1
\label{eq:v3v94-SM-holonomy-one}
\end{equation}
for every mapping-torus loop \(\gamma\).  This includes torsion loops and
large transformations of the actual quotient group
\(G_{\rm SM}=S(U(3)\times U(2))\).
\end{theorem}

\begin{proof}
The local degree-six terms in \eqref{eq:v3v94-Dai-Freed} cancel in the
complete Standard Model sum by V66, including the mixed
gauge--gravitational coefficient on curved vertical backgrounds.  The
remaining exponentiated eta phase is the bordism character constructed in
V87.  V87 proves
\begin{equation}
 \Omega_5^{\rm Spin}(BG_{\rm SM})=0,
\label{eq:v3v94-bordism-zero}
\end{equation}
so that character is identically one.  The continuum Standard Model
holonomy is therefore one on every \(\gamma\).  Exact Kato transport in
Theorem~\ref{thm:v3v94-mapping-torus-comparison} transfers this holonomy to
the corrected finite line, proving
\eqref{eq:v3v94-SM-holonomy-one}.  Because the argument uses a line-bundle
isomorphism and the full bordism character, it includes torsion which a
Chern-form computation cannot see.
\end{proof}

\begin{corollary}[Global parallel Standard Model section on the smooth quotient]
\label{cor:v3v94-global-section}
On every connected smooth bounded-geometry parameter component, fix the
phase at one base background.  Parallel transport with the corrected
complete Standard Model connection defines a path-independent nowhere-zero
section of its determinant line.
\end{corollary}

\begin{proof}
Local anomaly cancellation and the exact transgression make the corrected
complete connection flat.  Theorem~\ref{thm:v3v94-SM-holonomy} says its
holonomy around every loop is trivial.  Parallel transport from the base
point is therefore independent of path and gives the stated section.
\end{proof}

\begin{firewall}[The correction is geometric-sector data]
\label{fw:v3v94-geometric-sector}
The exact transgression uses the V92 global Bott regulator and the V93
compression homotopy.  It is defined for lattice fields obtained from the
controlled smooth geometric family.  It has not been extended to every
rough plaquette-admissible link field without a smooth bundle
reconstruction.  Thus \eqref{eq:v3v94-SM-holonomy-one} is not yet an exact
finite Ward solution on the entire unconstrained lattice configuration
space.
\end{firewall}

\begin{firewall}[Continuum normalization is not proved to be a local lattice counterterm]
\label{fw:v3v94-eta-locality}
The exact form \(\eta_a\) in \eqref{eq:v3v94-eta-normalization} includes
the global low-spectrum part of the Bismut--Freed connection.  Pulling it
back proves an exact determinant-line comparison on the smooth sector, but
it does not show that \(\eta_a\) is the variation of an intrinsic
finite-spacing local functional of arbitrary links.  V86 controls its local
short-time normalization; separating the physical global determinant data
from a lattice-local counterterm remains necessary before extending the
construction to rough fields.
\end{firewall}

\begin{firewall}[An uncorrected phase converging to one is not one]
\label{fw:v3v94-not-exact-without-correction}
Equation \eqref{eq:v3v94-uncorrected-bound} proves convergence, not exact
triviality, for the continuum-normalized but compression-uncorrected
connection.  The raw Grassmann connection additionally contains
\(\eta_a\) from \eqref{eq:v3v94-eta-normalization}.  Exact equality in
\eqref{eq:v3v94-SM-holonomy-one} uses the explicitly defined Kato
transgression \eqref{eq:v3v94-corrected-connection}.  Omitting that local
normalization would leave a small regulator-dependent phase at every finite
scale.
\end{firewall}

\begin{firewall}[Reflection positivity remains independent]
\label{fw:v3v94-RP}
Trivial determinant holonomy removes the global gauge-phase obstruction.  It
does not prove reflection positivity of the finite compressed action.  The
growing temporal range identified in V93 remains the relevant transfer
obstruction.
\end{firewall}

\subsection{V94 conclusion and V95 audit target}
\label{sec:v3v94-conclusion}

The torsion ambiguity left by V91 is now removed on the smooth geometric
sector.  The gapped compression homotopy supplies an actual determinant-line
isomorphism; its exact local Kato transgression matches endpoint connections
and all loop holonomies.  Combining this with V66 local cancellation and
V87's vanishing fifth spin bordism group gives a global parallel Standard
Model determinant section, including large quotient-group transformations.
The raw finite phase also converges with an explicitly summable error.

V95 is the next compulsory NS/YM audit.  It should decide between the two
remaining upstream gates: extending smooth geometric reconstruction to the
dynamical rough measure, or obtaining reflection positivity with the curved
finite regulator.  The audit should be used to choose which gate has a
credible exact decomposition.  Any reflection argument must exhibit the
cross-plane positive cone; any reconstruction argument must control a
signed local density before taking absolute values.

\clearpage
\section{V95: audited one-occurrence heat splitting of the regulator eta form}
\label{sec:v3v95}

V94 separates the small Kato compression form \(\alpha_{\rm comp}\) from
the Bott-to--Bismut--Freed normalization form \(\eta_a\), but leaves the
latter as an undivided relative eta form.  The present compulsory-audit note
resolves that ambiguity.  A single exact heat-time partition decomposes
\(\eta_a\) into a short-time local ultraviolet term and a complementary
low-spectrum term.  The same heat occurrence is never assigned to both.
After subtracting the complete nondecaying local jet, the remainder is
\(O(a)\), while the complementary term converges to the physical
Bismut--Freed connection, including its spectral-section transitions.

\subsection{Compulsory V95 NS/YM audit}
\label{sec:v3v95-audit}

The read-only files and exact fingerprints at the audit time were
\begin{align}
 &\texttt{Renewal\_NS\_Stability\_Research\_Notes\_V31.tex},
 &&887537\ \text{bytes},
 &&\texttt{F1121B62238AD5491BB7CF03635EA9934942EDF573ED8FAAA9EE79E1D38A6696},
 \label{eq:v3v95-audit-NS}\\
 &\texttt{Renewal\_Yang\_Mills\_Mass\_Gap\_Research\_Notes\_V16.tex},
 &&302586\ \text{bytes},
 &&\texttt{13531C2BA1698FC13C6E6429FAB181FD33654198417EDD50F1377C116255C6B4},
 \label{eq:v3v95-audit-YM16}\\
 &\texttt{Renewal\_Yang\_Mills\_Mass\_Gap\_Research\_Notes\_V15.tex},
 &&284098\ \text{bytes},
 &&\texttt{5A7D398E9D3FC0C2C62E7CF27BFDE1D21527DA8BA3DAB82B0B6518DB16A571D6}.
 \label{eq:v3v95-audit-YM15}
\end{align}
The two frozen YM volumes remained
\begin{align}
 &\texttt{Renewal\_Yang\_Mills\_Mass\_Gap\_Research\_Notes\_V135\_f1.tex},
 &&3083467\ \text{bytes},
 &&\texttt{82F6BBFFEECDAD3C5C63551660B19A0BDC7F3D6DFFE4C68B4E9BBA3A783A2B71},
 \label{eq:v3v95-audit-YM-f1}\\
 &\texttt{Renewal\_Yang\_Mills\_Mass\_Gap\_Research\_Notes\_V100\_f2.tex},
 &&1706928\ \text{bytes},
 &&\texttt{6444A1044E74DFCB1A72ED9C7248BC0F234EFC4A0E4248C736CD7DB212DFCA9A}.
 \label{eq:v3v95-audit-YM-f2}
\end{align}

NS V31 still shows that a flat mark is removed by an exact first moment and
that separate absolute estimates merely restore the missing critical scale.
Its balanced-helicity example also shows that a signed polarization ledger
cannot be promoted to a positive total-tail estimate.

YM V16 adds the directly relevant algebraic lesson.  Its finite telescope
puts exactly one partition-state difference in every summand; coordinate
restriction cannot clone one covariance occurrence into two marks.  It
localizes every genuine multi-mark request to an explicit higher-incidence
collision locus and preserves the unique payer throughout the expansion.

For the SM programme, the one occurrence is the relative heat
transgression defining \(\eta_a\).  Calling its short-time asymptotic a
local counterterm and then using the entire eta form again as physical
spectral data would count that occurrence twice.  The exact cutoff identity
below assigns every heat time to one side.  The complete signed supertrace is
assembled before either part is estimated.  This route is more concrete
than attempting reflection positivity for V93's growing temporal band and
more upstream than extending an unresolved normalization to rough fields.

\subsection{Relative superconnection transgression}
\label{sec:v3v95-superconnection}

Work first on a determinant-line patch on which a spectral section has been
chosen as in V88.  After finite spectator stabilization, choose a smooth
path \(\mathscr D_{a,s}\), \(0\le s\le1\), of elliptic self-adjoint
representatives with common vertical Dirac symbol, joining the curved
Bott--Wilson representative to the continuum Dirac representative.  The
finite-dimensional crossing band is removed by the spectral section; on
its orthogonal complement there is a common gap.  Let
\begin{equation}
 \mathbb A_{a,s,t}=t^{1/2}\mathscr D_{a,s}+\nabla^{\mathscr H}_{a,s}
\label{eq:v3v95-superconnection}
\end{equation}
be the Bismut superconnection on that complement.  Lower-degree curvature
terms required by a nonproduct horizontal distribution are included in
\(\nabla^{\mathscr H}_{a,s}\); they do not change the argument.

\begin{lemma}[Relative heat representation]
\label{lem:v3v95-heat-representation}
There is an operator supertrace density \(\Theta_a(t)\in\Omega^1(B)\),
obtained by integrating the \(s\)-transgression of
\(\exp(-\mathbb A_{a,s,t}^2)\), such that
\begin{equation}
 \eta_a=\operatorname{FP}\int_0^\infty\Theta_a(t)\,dt+d_B\varphi_a.
\label{eq:v3v95-eta-heat}
\end{equation}
Here \(\operatorname{FP}\) is the finite part at \(t=0\), and
\(d_B\varphi_a\) records the determinant-frame convention.  The reduced
kernel defining \(\Theta_a(t)\) decays exponentially as \(t\to\infty\) on
the chosen spectral-section patch.  On overlaps, the change of
\eqref{eq:v3v95-eta-heat} is exactly the finite-band determinant transition
of V88.
\end{lemma}

\begin{proof}
Duhamel's formula gives
\begin{equation}
 \partial_s e^{-\mathbb A_{a,s,t}^2}
 =-\int_0^1e^{-u\mathbb A^2}
 (\partial_s\mathbb A^2)e^{-(1-u)\mathbb A^2}\,du.
\label{eq:v3v95-Duhamel}
\end{equation}
Cyclicity of the supertrace and
\(\partial_s\mathbb A^2=[\mathbb A,\partial_s\mathbb A]_{\rm gr}\)
turn its degree-two component into the exterior derivative of a degree-one
transgression.  Integrate that transgression in \(s\) and in heat time.  The
small-time heat expansion has finitely many nonintegrable powers and
logarithms; taking their Hadamard finite part gives
\eqref{eq:v3v95-eta-heat}, which is the standard relative eta-form formula
for the difference of determinant connections.

The spectral section removes a finite neighbourhood of zero.  The remaining
squared operators have a common positive lower bound, so their heat kernels
and every fixed differentiated insertion decay exponentially at large time.
Changing spectral section restores precisely the removed finite band.  The
top exterior power of that band is the transition line in V88, proving the
last assertion.
\end{proof}

The formula is relative: the common principal symbol cancels the leading
uncut heat singularities.  Nevertheless its finite part can contain local
power and logarithmic terms in the regulator scale, which must be normalized
before the continuum limit.

\subsection{The exact one-occurrence partition}
\label{sec:v3v95-partition}

Choose once and for all \(\chi\in C^\infty([0,\infty);[0,1])\) with
\(\chi(u)=1\) for \(u\le1\) and \(\chi(u)=0\) for \(u\ge2\).  Define
\begin{align}
 \eta_a^{\rm UV}
 &=\operatorname{FP}\int_0^\infty
 \chi(t/a^2)\Theta_a(t)\,dt,
 \label{eq:v3v95-UV}\\
 \eta_a^{\rm IR}
 &=\int_0^\infty
 (1-\chi(t/a^2))\Theta_a(t)\,dt.
 \label{eq:v3v95-IR}
\end{align}

\begin{theorem}[Exact one-occurrence heat split]
\label{thm:v3v95-exact-split}
On every spectral-section patch,
\begin{equation}
 \eta_a=\eta_a^{\rm UV}+\eta_a^{\rm IR}+d_B\varphi_a.
\label{eq:v3v95-exact-split}
\end{equation}
Every Duhamel insertion and every spectral value occurs with total weight
one.  In particular, neither term on the right can be replaced by the full
relative eta form while retaining the other.
\end{theorem}

\begin{proof}
The pointwise identity
\(\chi(t/a^2)+(1-\chi(t/a^2))=1\) partitions the single integrand in
\eqref{eq:v3v95-eta-heat}.  The second integral vanishes near zero and needs
no finite part; the first has compact heat-time support and contains the
entire finite-part prescription.  Adding them gives
\eqref{eq:v3v95-exact-split}.  Since the partition is applied after the
complete \(s\)-integral and supertrace have been assembled, it changes no
Duhamel multiplicity.  This proves the one-occurrence statement.
\end{proof}

\begin{firewall}[A smooth overlap is not a third heat sector]
\label{fw:v3v95-overlap}
The interval \(a^2<t<2a^2\) is shared by the two smooth weights, but their
sum is exactly one there.  It is not an additional intermediate contribution
and may not be charged to both the UV and IR estimates with unit weight.
\end{firewall}

\subsection{Local ultraviolet jet with a positive remainder power}
\label{sec:v3v95-UV-jet}

Let \(\mathfrak M_{\le0}\) extract from a regulator-dependent local form
all terms which do not vanish as \(a\downarrow0\), including logarithms.
The normalization conditions of V86 fix the gauge-invariant coefficients of
this finite jet.

\begin{theorem}[Curved UV transgression expansion]
\label{thm:v3v95-UV-expansion}
On a compact bounded-geometry family with sufficiently many bounded
derivatives,
\begin{equation}
 \eta_a^{\rm UV}
 =\mathfrak M_{\le0}\eta_a^{\rm UV}+r_a^{\rm UV},
 \qquad
 \|r_a^{\rm UV}\|_{C^1(B)}\le Ca.
\label{eq:v3v95-UV-remainder}
\end{equation}
The jet is a finite sum of integrals over \(X\) of covariant polynomials in
the gauge curvature, Riemann curvature, Higgs--Yukawa coefficients, and
their covariant derivatives.  Its gauge-variant cohomology class is the
descent of
\begin{equation}
 [\widehat A(T(X/B))\operatorname{ch}(\mathbb E)]_6.
\label{eq:v3v95-six-form}
\end{equation}
\end{theorem}

\begin{proof}
Rescale heat time by \(t=a^2u\).  On the compact interval
\(0<u<2\), the covariant heat parametrix for
\(\mathbb A_{a,s,a^2u}^2\) is obtained by freezing the metric and
connection in normal coordinates and iterating Duhamel's formula.  Through
any fixed order it is a finite sum of Gaussian moments multiplying covariant
jets of the coefficients.  The uniform bounded-geometry hypotheses control
the integral remainder and its first parameter derivative.

Take the complete Clifford supertrace before estimating.  Odd Gaussian
moments and Clifford degrees below four cancel.  Integrating the remaining
coefficients in \(u\) against \(\chi(u)\) yields finitely many powers of
\(a\) and logarithms.  Taylor expansion one order beyond the last
nondecaying coefficient leaves one explicit factor of \(a\); Gaussian
moments and the finite \(s\)-interval give the uniform bound in
\eqref{eq:v3v95-UV-remainder}.

The local families index theorem identifies the only nontrivial gauge
variation of the determinant transgression with the descent of
\eqref{eq:v3v95-six-form}.  All other nondecaying terms are gauge-invariant
local normalization terms.  This proves the stated structure of the jet.
\end{proof}

\begin{corollary}[Standard Model UV cancellation]
\label{cor:v3v95-SM-UV}
In the complete Standard Model representation, the cohomologically
nontrivial gauge variation of
\(\sum_f\mathfrak M_{\le0}\eta_{a,f}^{\rm UV}\) vanishes, including the
mixed gauge--gravitational part.  After the V86 gauge-invariant
normalization conditions are imposed,
\begin{equation}
 \left\|sum_fr_{a,f}^{\rm UV}\right\|_{C^1(B)}\le C_{\rm SM}a.
\label{eq:v3v95-SM-UV-remainder}
\end{equation}
\end{corollary}

\begin{proof}
V66 cancels every representation coefficient in the descent of
\eqref{eq:v3v95-six-form} for the physical quotient group.  This
cancellation is made in the complete signed sum before a norm is taken.
The remaining jet coefficients are gauge-invariant and are fixed, not
assumed zero, by the V86 normalization data.  Summing the finite collection
of remainder bounds in \eqref{eq:v3v95-UV-remainder} proves
\eqref{eq:v3v95-SM-UV-remainder}.
\end{proof}

\subsection{The infrared term is the physical determinant occurrence}
\label{sec:v3v95-IR}

Let \(\eta_{\rm phys}\) denote the relative connection one-form determined
by the continuum Bismut--Freed determinant line and the chosen spectral
section.  Its changes between patches are the finite-band transitions of
V88.

\begin{theorem}[IR convergence with crossing compatibility]
\label{thm:v3v95-IR-convergence}
Under V89's uniform generalized norm-resolvent and cluster derivative
hypotheses, and the uniform heat estimates of V92, one has
\begin{equation}
 \eta_a^{\rm IR}=\eta_{\rm phys}+r_a^{\rm IR}+d_B\psi_a,
 \qquad
 \|r_a^{\rm IR}\|_{C^1(B)}\le Ca.
\label{eq:v3v95-IR-limit}
\end{equation}
The exact forms depend on determinant frames.  On an overlap of spectral
sections, both sides change by the same finite-band determinant connection,
so \(r_a^{\rm IR}\) is globally compatible.
\end{theorem}

\begin{proof}
Split the reduced spectrum at a fixed large value \(\Lambda\).  On the
finite low band, V89's Riesz direct rotation identifies the lattice and
continuum bundles, their connections, and their regular crossing
orientations.  The heat integrand converges there in \(C^1(B)\), including
through a crossing after the spectral-section band is removed.

On the complement, the common elliptic Garding bound and heat-kernel
estimate give an integrable majorant.  Expand its short end covariantly to
the same order used in Theorem~\ref{thm:v3v95-UV-expansion}.  The terms
already assigned to \(\eta_a^{\rm UV}\) carry the complementary weight
\(1-\chi\) here and cancel in the exact sum
\eqref{eq:v3v95-exact-split}; the first unassigned moment has one factor of
\(a\).  The large-time tail is exponentially small uniformly after
\(\Lambda\) is fixed.  Dominated convergence followed by the uniform
remainder estimate gives \eqref{eq:v3v95-IR-limit}.

Changing spectral section adds the top exterior connection of the removed
finite band.  V89 proves convergence of that transition on both endpoints,
so the same term occurs on the two sides and cancels from the remainder.
\end{proof}

\begin{firewall}[The IR term is not a counterterm]
\label{fw:v3v95-IR-not-local}
The form \(\eta_{\rm phys}\) contains the low spectrum, its crossing
transitions, and mapping-torus holonomy.  It is global physical determinant
data.  Its presence in \eqref{eq:v3v95-IR-limit} is not a failure of
locality and it may not be subtracted as though it were a local ultraviolet
functional.
\end{firewall}

\subsection{Resolved normalization formula}
\label{sec:v3v95-resolved}

Combining the exact split with the two estimates gives
\begin{equation}
 \boxed{
 \eta_a=
 \mathfrak M_{\le0}\eta_a^{\rm UV}
 +\eta_{\rm phys}
 +r_a+d_B\zeta_a,
 \qquad
 \|r_a\|_{C^1(B)}\le Ca.}
\label{eq:v3v95-resolved-eta}
\end{equation}

\begin{theorem}[No-double-counting continuum normalization]
\label{thm:v3v95-normalization}
Subtract only the local jet
\(\mathfrak M_{\le0}\eta_a^{\rm UV}\) from the regulator connection and
identify \(\eta_{\rm phys}\) with the physical continuum determinant
connection.  Then the remaining regulator mismatch is \(O(a)\) in
\(C^1(B)\), is compatible with all V88 transition functions, and is
summable on the geometric refinement \(a_n=q^{-n}a_0\).  Adding V93's
finite-compression error along
\eqref{eq:v3v93-diagonal-schedule} preserves absolute summability.
\end{theorem}

\begin{proof}
Equation \eqref{eq:v3v95-resolved-eta} gives the first two assertions after
an exact determinant-frame change removes \(d_B\zeta_a\).  The geometric
series \(\sum_na_n\) converges.  V93's traced compression errors are
summable by \eqref{eq:v3v93-summable-error}; the sum of two absolutely
summable sequences is absolutely summable.
\end{proof}

For the complete Standard Model line, V66 cancels the anomalous variation of
the local jet, V86 fixes its gauge-invariant coefficients, V94 trivializes
the holonomy of the physical term, and
Theorem~\ref{thm:v3v95-normalization} controls the remaining continuum
error.  These are four distinct operations on the two terms in the exact
partition; none is used twice.

\begin{firewall}[Smooth heat coefficients do not reconstruct rough links]
\label{fw:v3v95-rough}
The heat expansion uses bounded geometry and finitely many coefficient
derivatives.  Plaquette admissibility alone supplies neither those
derivatives nor a canonical smooth interpolation.  V95 therefore resolves
the normalization on the smooth geometric sector but does not extend it to
the full rough lattice measure.
\end{firewall}

\begin{firewall}[No reflection-positive factorization]
\label{fw:v3v95-RP}
The heat-time partition is an exact analytic decomposition of a determinant
connection.  It is not a time-slice factorization of the fermion action and
does not supply the positive cross-plane cone required for reflection
positivity.
\end{firewall}

\subsection{V95 conclusion and V96 target}
\label{sec:v3v95-conclusion}

The V94 normalization form no longer hides a duplicated spectral
contribution.  One exact cutoff partitions its single heat occurrence into
a local UV part and a physical IR part.  The complete signed short-time
supertrace yields a finite nondecaying jet plus an \(O(a)\) remainder; the
complementary heat term converges, with its spectral-section transitions, to
the Bismut--Freed determinant connection.  Standard Model anomaly
cancellation acts only on the local jet, while V94 bordism triviality acts
only on the global holonomy.

V96 should now attack rough-field reconstruction.  It should define a
gauge-covariant smoothing map from quantitatively admissible link fields to
connections with controlled curvature, prove that the original and sampled
links are joined inside the Wilson-gap component, and identify the exact
topological sector.  The map must preserve one plaquette occurrence under
each scale decomposition; a plaquette defect used to control smoothing may
not be charged again as independent curvature stock.  If such a map cannot
be uniform in four dimensions, the obstruction should be isolated before
returning to reflection positivity.

\clearpage
\section{V96: cubical reconstruction of admissible rough link fields}
\label{sec:v3v96}

V95 resolves the determinant normalization on smooth bounded-geometry
families.  A dynamical lattice gauge field is not such a family: a fixed
plaquette admissibility bound permits curvature of size \(a^{-2}\) and gives
no uniform derivative bounds.  This note constructs the exact geometric
object which an admissible link field does determine.  Cellwise tree gauge,
principal plaquette logarithms, and a relative extension over the cubical
skeleton produce a smooth principal bundle and connection whose edge
holonomies are the original links.  The construction is gauge covariant,
uses each plaquette defect only a bounded number of times, and converts the
Wilson plaquette action into an \(L^2\)-curvature bound.  It also proves
overlap-index equality on every admissible component containing a uniformly
smooth representative.  What remains is a measure theorem excluding or
controlling components and concentration regions without such a
representative.

\subsection{Admissible link data}
\label{sec:v3v96-data}

Let \(K_a\) be a finite cubical decomposition of a closed oriented spin
four-manifold at mesh scale \(a\), with uniformly bounded incidence and
shape regularity in geodesic coordinate charts.  Take the faithful matrix
realization
\begin{equation}
 G_{\rm SM}=S(U(3)\times U(2))\subset U(5).
\label{eq:v3v96-GSM}
\end{equation}
A link field assigns \(U_e\in G_{\rm SM}\) to each oriented edge, with
\(U_{\bar e}=U_e^{-1}\).  Its plaquette holonomy is \(U_p\), based at a
chosen vertex.  Put
\begin{equation}
 \epsilon_p=\|1-U_p\|,
 \qquad
 \epsilon(U)=\max_p\epsilon_p.
\label{eq:v3v96-admissibility}
\end{equation}
Changing the base vertex conjugates \(U_p\) and leaves \(\epsilon_p\)
unchanged.

\begin{lemma}[Principal logarithm in the physical group]
\label{lem:v3v96-principal-log}
There is \(\epsilon_{\log}>0\) such that if
\(g\in G_{\rm SM}\) and \(\|1-g\|<\epsilon_{\log}\), then its principal
matrix logarithm \(X=\log g\) belongs to
\(\mathfrak g_{\rm SM}\) and
\begin{equation}
 \|X\|\le C_{\log}\|1-g\|.
\label{eq:v3v96-log-bound}
\end{equation}
The map \(g\mapsto X\) is analytic and conjugation equivariant on this
ball.
\end{lemma}

\begin{proof}
Choose the ball so the spectrum of \(g\) avoids the negative real axis.
Holomorphic functional calculus defines the principal logarithm and gives
the convergent series estimate
\(\|\log g\|\le\sum_{n\ge1}\|g-1\|^n/n\).  Since
\(G_{\rm SM}\) is a closed matrix Lie subgroup, the logarithm of an element
in its exponential neighbourhood lies in its Lie algebra.  Holomorphic
functional calculus commutes with conjugation, proving equivariance.
\end{proof}

\subsection{One-cell tree gauge and one-use curvature stock}
\label{sec:v3v96-cell}

Fix a four-cell \(Q\), a root vertex, and a spanning tree of its one-skeleton.
Gauge successively along the tree so that every tree link is one.

\begin{lemma}[Cell link reduction]
\label{lem:v3v96-cell-link}
Every remaining link in the tree gauge is a word containing at most
\(N_Q\) plaquette holonomies and their conjugates, where \(N_Q\) depends
only on the reference four-cube.  If \(\epsilon(U)<\epsilon_{\log}/N_Q\),
then
\begin{equation}
 U_e=e^{A_e},
 \qquad
 A_e\in\mathfrak g_{\rm SM},
 \qquad
 \|A_e\|\le C_Q\sum_{p\subset Q}\epsilon_p.
\label{eq:v3v96-edge-log}
\end{equation}
Each plaquette occurs in only a bounded number of these words.
\end{lemma}

\begin{proof}
For a non-tree edge, adjoining it to the tree creates one closed edge loop.
The cellular boundary group of a cube is generated by its square faces.
Successively peel faces from a spanning disk for that loop.  This writes its
holonomy as a product of based plaquette holonomies, with intervening tree
transports equal to one.  The reference cube has finitely many edges and
faces, so the word length and incidence are bounded by fixed constants.

For unitary matrices,
\(\|g_1\cdots g_k-1\|\le\sum_j\|g_j-1\|\), by telescoping and unitarity.
Apply Lemma~\ref{lem:v3v96-principal-log} to the resulting link and obtain
\eqref{eq:v3v96-edge-log}.  The face-peeling construction uses each face in
only finitely many reference loops, proving bounded incidence.
\end{proof}

Let \(w_e\) be the cubical Whitney one-form associated with an oriented
edge of \(Q\), scaled to physical size \(a\).  The preliminary local
connection is
\begin{equation}
 A_Q^{(1)}=\sum_{e\subset Q}A_ew_e.
\label{eq:v3v96-Whitney-connection}
\end{equation}
On the edge \(e\), all other Whitney one-forms pull back to zero and
\(w_e\) pulls back to the unit interval form.  Hence the edge holonomy is
exactly \(e^{A_e}=U_e\) in tree gauge.

\begin{proposition}[Cell curvature estimate]
\label{prop:v3v96-cell-curvature}
There is a relative correction supported away from the one-skeleton which
makes \(A_Q^{(1)}\) compatible with the prescribed lower-skeleton data and
obeys, for \(k\ge0\),
\begin{align}
 \|\nabla^kF_{A_Q}\|_{L^\infty(Q)}
 &\le C_k a^{-k-2}
 \sum_{p\subset Q}\epsilon_p,
 \label{eq:v3v96-cell-Ck}\\
 \int_Q|F_{A_Q}|^2
 &\le C\sum_{p\subset Q}\epsilon_p^2.
\label{eq:v3v96-cell-L2}
\end{align}
The correction does not change any edge holonomy.
\end{proposition}

\begin{proof}
On the reference cube the Whitney forms and all their derivatives are
bounded.  Rescaling gives
\(\|\nabla^kw_e\|_\infty\le C_ka^{-k-1}\).  Since
\(F=dA+A\wedge A\), Lemma~\ref{lem:v3v96-cell-link} gives the first bound
for the preliminary connection; the quadratic term is absorbed by choosing
the admissibility constant small.

Compatibility is imposed inductively over faces of dimensions two, three,
and four.  At each step the difference of two already prescribed boundary
connections is a Lie-algebra-valued form which vanishes on the lower
skeleton.  Extend it with a fixed bounded right inverse to the restriction
map on the reference cube and multiply by a cutoff which is zero in a
neighbourhood of the one-skeleton.  Rescaling preserves the same powers of
\(a\).  The finite number of skeleton steps changes only \(C_k\) and leaves
edge holonomies fixed.  Finally square the \(k=0\) bound and multiply by
\(\operatorname{vol}Q\asymp a^4\).  Bounded face incidence and Cauchy's
inequality yield \eqref{eq:v3v96-cell-L2}.
\end{proof}

The estimate is a one-use ledger: a plaquette defect enters a bounded list
of cell coefficients, and the final square is paid by its own
\(\epsilon_p^2\) up to the fixed incidence constant.  No coordinate
restriction produces a second independent curvature occurrence.

\subsection{Global bundle reconstruction}
\label{sec:v3v96-global}

Choose tree gauges independently on all four-cells.  On a common face, the
two gauges agree on the shared vertices up to a uniquely determined
\(G_{\rm SM}\)-valued vertex transformation.  Extend these transformations
over the face in skeleton order, using principal logarithms for the
near-identity relative part and retaining the exact vertex values.

\begin{theorem}[Gauge-covariant cubical reconstruction]
\label{thm:v3v96-reconstruction}
There are universal \(\epsilon_{\rm rec}>0\) and constants \(C_k\) such
that every link field with \(\epsilon(U)<\epsilon_{\rm rec}\) determines:
\begin{enumerate}
 \item a smooth principal \(G_{\rm SM}\)-bundle
 \(P_U\to X\), given by cell transition functions satisfying the exact
 cocycle law;
 \item a smooth connection \(A_U\) on \(P_U\) whose parallel transport on
 every original edge is \(U_e\); and
 \item the estimates
 \begin{align}
  \|\nabla^kF_{A_U}\|_{L^\infty(Q)}
  &\le C_ka^{-k-2}\sum_{p\cap Q\ne\varnothing}\epsilon_p,
  \label{eq:v3v96-global-Ck}\\
  \|F_{A_U}\|_{L^2(X)}^2
  &\le C\sum_p\epsilon_p^2.
  \label{eq:v3v96-global-L2}
 \end{align}
\end{enumerate}
The assignment is continuous on the admissible set and equivariant under
lattice gauge transformations; gauge-related link fields produce
isomorphic bundles with gauge-related connections.
\end{theorem}

\begin{proof}
Proceed inductively over the dual Cech nerve of the cubical cover.  On each
double overlap, extend the vertex gauge transformation while fixing its
lower-dimensional boundary.  On a triple overlap, the product of the three
candidate transitions is one on the lower skeleton because it is obtained
from the same link transports.  Its near-identity logarithm gives a relative
null-homotopy; adjust the last transition by that null-homotopy.  Repeat on
higher overlaps.  The dimension is four and the cover has bounded
incidence, so the induction terminates after finitely many uniform steps and
produces the exact cocycle law.

Use Proposition~\ref{prop:v3v96-cell-curvature} with the same relative
extension operators to make the local connections obey
\begin{equation}
 A_{Q'}=g_{QQ'}^{-1}A_Qg_{QQ'}+g_{QQ'}^{-1}dg_{QQ'}
\label{eq:v3v96-connection-gluing}
\end{equation}
on overlaps.  Relative smoothing away from the one-skeleton makes the data
smooth without altering edge transports.  The local derivative bounds and
bounded overlap prove \eqref{eq:v3v96-global-Ck}; summing
\eqref{eq:v3v96-cell-L2} proves \eqref{eq:v3v96-global-L2}.

Every step uses analytic logarithms and fixed linear extension operators,
so it depends continuously on the links while admissibility holds.  A
lattice gauge transformation merely changes the initial cell tree gauges;
the two constructions are related by the extended transformation.  This
proves equivariance.
\end{proof}

\begin{firewall}[Transition topology is not forced into a principal logarithm]
\label{fw:v3v96-topology-log}
Only the relative near-identity corrections inside contractible overlaps
are logarithmed.  The full transition function, which can carry nontrivial
bundle topology, is retained exactly.  Taking one global logarithm of all
links would incorrectly collapse the construction to the trivial bundle
sector.
\end{firewall}

\subsection{Component invariants and the overlap index}
\label{sec:v3v96-components}

Define the reconstructed topological charge in representation \(R\) by
\begin{equation}
 Q_R^{\rm geom}(U)=
 \int_X[\widehat A(TX)\operatorname{ch}(E_R(P_U))]_4.
\label{eq:v3v96-geometric-charge}
\end{equation}

\begin{proposition}[Local constancy of reconstructed topology]
\label{prop:v3v96-topology-constant}
Along every continuous path inside
\(\{U:\epsilon(U)<\epsilon_{\rm rec}\}\), the isomorphism class of
\(P_U\) and every integer \(Q_R^{\rm geom}(U)\) are constant.
\end{proposition}

\begin{proof}
Theorem~\ref{thm:v3v96-reconstruction} produces a continuous bundle over
\(X\times[0,1]\).  Restriction of a bundle over a product to the two ends
gives isomorphic K-theory classes.  Characteristic numbers are homotopy
invariants, so \eqref{eq:v3v96-geometric-charge} is constant.
\end{proof}

Call an admissible component \(\mathscr C_a\) geometrically tame if it
contains an exact holonomy sample \(U(A_*)\) of a smooth connection on a
fixed bundle for which \(a\) lies below the V90--V93 mesh threshold.  The
other fields in \(\mathscr C_a\) may be arbitrarily rough subject to the
plaquette bound.

\begin{theorem}[Rough-component regulator index theorem]
\label{thm:v3v96-rough-index}
Let \(\mathscr C_a\) be geometrically tame.
\begin{enumerate}
 \item On a flat hypercubic background, decrease
 \(\epsilon_{\rm rec}\) below the Wilson-gap threshold in
 Theorem~\ref{thm:v3v90-admissible-gap}.  The nearest-neighbour overlap
 operator then satisfies, for every \(U\in\mathscr C_a\),
 \begin{equation}
  \operatorname{ind}_{\rm ov}(U;R)
  =Q_R^{\rm geom}(U)
  =\operatorname{ind}D^+_{E_R(P_U)}.
 \label{eq:v3v96-rough-index}
 \end{equation}
 \item On a curved triangulated background, the same equality holds for the
 V93 compressed Bott regulator if the reconstructed path from \(U\) to the
 tame representative admits one common V93 compression gap.  No
 nearest-neighbour curved Wilson stencil is asserted.
\end{enumerate}
\end{theorem}

\begin{proof}
In the flat case, the admissible path from \(U\) to \(U(A_*)\) stays
inside the Wilson gap by Theorem~\ref{thm:v3v90-admissible-gap}.
Proposition~\ref{prop:v3v90-homotopy-index} makes the overlap index constant
along that path, and Theorem~\ref{thm:v3v90-overlap-index} identifies the
endpoint with the continuum index.  In the curved case, the additional
hypothesis gives a gapped path of the finite V93 matrices; their negative
projection rank is constant and Theorem~\ref{thm:v3v93-finite-regulator}
identifies the endpoint index.  In both cases
Proposition~\ref{prop:v3v96-topology-constant} identifies the endpoint
bundle class and characteristic number with those reconstructed from
\(U\).
\end{proof}

\begin{firewall}[Tame-component membership is a real hypothesis]
\label{fw:v3v96-tame}
Theorem~\ref{thm:v3v96-reconstruction} assigns a smooth connection to every
admissible finite link field, but its derivative bounds scale as
\(a^{-k-2}\).  This alone does not show that the field can be joined within
the admissible set to a representative with bounds uniform as
\(a\to0\).  The geometrically tame condition in
Theorem~\ref{thm:v3v96-rough-index} is therefore not automatic and is not
silently removed.
\end{firewall}

\subsection{What the Wilson action controls}
\label{sec:v3v96-action}

For sufficiently small plaquette defects, the Wilson action is comparable
to their square sum:
\begin{equation}
 S_{W,a}(U)\asymp\frac1{g_a^2}\sum_p\epsilon_p^2.
\label{eq:v3v96-Wilson-action}
\end{equation}
The reconstruction therefore gives
\begin{equation}
 \|F_{A_U}\|_{L^2(X)}^2
 \le Cg_a^2S_{W,a}(U).
\label{eq:v3v96-action-curvature}
\end{equation}

\begin{proposition}[Critical compactness supplied by action control]
\label{prop:v3v96-critical-compactness}
On a region where the right side of
\eqref{eq:v3v96-action-curvature} is below the four-dimensional Uhlenbeck
threshold, one can choose gauges with a uniform \(W^{1,2}\) bound for
\(A_U\).  This control is critical and does not imply the \(C^6\) bounds
used in V95.
\end{proposition}

\begin{proof}
The small-curvature gauge theorem gives, on a ball,
\begin{equation}
 d^*A_U=0,
 \qquad
 \|A_U\|_{W^{1,2}}
 \le C\|F_{A_U}\|_{L^2}.
\label{eq:v3v96-Uhlenbeck}
\end{equation}
Equation \eqref{eq:v3v96-action-curvature} supplies the hypothesis and
bound.  In four dimensions \(W^{1,2}\hookrightarrow L^4\) is scale
critical and does not embed in \(C^0\), much less \(C^6\).  Curvature can
concentrate while its global \(L^2\) norm stays bounded.  Hence the last
statement follows by the sharp Sobolev embedding.
\end{proof}

This identifies the dynamical bottleneck precisely.  The plaquette action
pays one \(L^2\)-curvature occurrence.  It cannot simultaneously be treated
as independent pointwise admissibility, six derivative bounds, and a second
counterterm remainder estimate.

\begin{theorem}[Conditional measure extension criterion]
\label{thm:v3v96-measure-criterion}
Let \(\mu_a\) be the normalized interacting lattice measures.  Suppose
there are gauge-invariant sets \(\mathscr G_a\) such that:
\begin{enumerate}
 \item \(\mu_a(\mathscr G_a^c)\to0\);
 \item every field in \(\mathscr G_a\) is admissible and belongs to a
 geometrically tame component;
 \item after removing finitely many concentration balls of total radius
 tending to zero, the reconstructed connections have the bounded-geometry
 and derivative estimates required in V95; and
 \item the determinant contribution of the concentration balls is uniformly
 integrable and converges to the corresponding topological bubble factor.
\end{enumerate}
Then the normalized finite Standard Model determinant sections of
V93--V95 converge in \(\mu_a\)-probability to the smooth-sector global
section, with the reconstructed topological charge on every tame component.
\end{theorem}

\begin{proof}
On \(\mathscr G_a\) away from the concentration balls, apply the uniform
V93 compression estimate and V95 heat split.  Their errors tend to zero and
are summable along the chosen refinement.  Item 4 supplies convergence and
uniform integrability on the removed balls, while
Theorem~\ref{thm:v3v96-rough-index} fixes their integer sector contribution.
Thus convergence holds on \(\mathscr G_a\).  For any \(\delta>0\), the
probability of an error larger than \(\delta\) is bounded by its probability
inside \(\mathscr G_a\), which tends to zero by the analytic estimates, plus
\(\mu_a(\mathscr G_a^c)\), which tends to zero by item 1.
\end{proof}

\begin{firewall}[The measure criterion has not been verified]
\label{fw:v3v96-measure-open}
The Wilson action and Uhlenbeck compactness provide the critical
\(L^2\)-curvature stock, but they do not prove items 1, 3, or 4 of
Theorem~\ref{thm:v3v96-measure-criterion}.  Establishing the required
concentration probability and bubble determinant control is a joint
Yang--Mills--fermion problem.  The full interacting continuum measure is not
constructed in V96.
\end{firewall}

\begin{firewall}[No use of reconstruction as reflection positivity]
\label{fw:v3v96-no-RP}
The connection \(A_U\) is a geometric interpolation of links.  It does not
factor the original lattice action across a time-reflection plane and does
not alter the V93 growing-range obstruction to reflection positivity.
\end{firewall}

\subsection{V96 conclusion and V97 target}
\label{sec:v3v96-conclusion}

Quantitatively admissible links now determine a genuine smooth
\(G_{\rm SM}\)-bundle and connection with exact original edge holonomies.
The reconstruction is gauge covariant, retains nontrivial transition
topology, and converts the plaquette square sum into an \(L^2\)-curvature
bound without duplicating plaquette stock.  On flat cubical backgrounds the nearest-neighbour overlap index equals the
continuum index throughout every tame admissible component.  On curved
backgrounds the corresponding statement holds for V93's compressed Bott
regulator under its explicitly stated common-gap hypothesis.  The sharp obstruction is that four-dimensional
action control gives critical \(W^{1,2}\), not the smooth bounds used by the
heat parametrix.

V97 should analyze concentration balls in the exact determinant heat split.
It should rescale one ball to unit size, separate its integer bubble charge
from the non-topological relative eta remainder, and test whether the
complete anomaly-free Standard Model supertrace gives a uniform integrable
bound.  If the bound requires a Yang--Mills mass-gap estimate not yet proved
in the companion programme, it should be stated as the precise shared
interface rather than replaced by pointwise admissibility.

\clearpage
\section{V97: concentration balls, weak Schatten control, and bubble determinants}
\label{sec:v3v97}

V96 converts plaquette action into the scale-critical bound
\(F\in L^2\) and isolates concentration balls as the obstruction to the
smooth heat expansion.  This note treats the fermion determinant on those
balls without demanding pointwise curvature bounds.  In Uhlenbeck gauge,
the rescaled connection belongs to \(L^4\).  A Dirac resolvent followed by
multiplication by that connection lies in weak Schatten class
\(\mathcal S_{4,\infty}\).  Subtracting the first four Born occurrences
therefore leaves a fifth regularized determinant with a scale-invariant
upper bound.  Each removed Born occurrence still has the exact V95 split
into a local UV coefficient and a complementary physical IR term; neither
is discarded.  Zero modes form a separate finite determinant line carrying
the bubble index.  On a stopping-ball cover the fifth determinant remainder
is exponentially bounded by a small multiple of the Yang--Mills action.

\subsection{Rescaling and the critical gauge norm}
\label{sec:v3v97-rescaling}

Let \(B_\rho(x_0)\subset X\) be a geodesic ball with radius below the
bounded-geometry scale.  In normal coordinates put
\begin{equation}
 A_\rho(y)=\rho A(x_0+\rho y),
 \qquad
 F_\rho(y)=\rho^2F_A(x_0+\rho y),
 \qquad y\in B_1.
\label{eq:v3v97-rescaling}
\end{equation}
Up to uniformly controlled metric errors,
\begin{equation}
 \|F_\rho\|_{L^2(B_1)}=\|F_A\|_{L^2(B_\rho)},
 \qquad
 \|A_\rho\|_{L^4(B_1)}=\|A\|_{L^4(B_\rho)}.
\label{eq:v3v97-critical-scaling}
\end{equation}

\begin{lemma}[Uhlenbeck stopping-ball gauge]
\label{lem:v3v97-Uhlenbeck-ball}
There is \(\varepsilon_U>0\) such that, if
\begin{equation}
 E(B_\rho):=\int_{B_\rho}|F_A|^2\le\varepsilon_U,
\label{eq:v3v97-small-energy}
\end{equation}
then a gauge exists on the concentric ball of radius \(\rho/2\) with
\begin{equation}
 d^*A=0,
 \qquad
 \|A\|_{L^4}+\|\nabla A\|_{L^2}
 \le C E(B_\rho)^{1/2}.
\label{eq:v3v97-Uhlenbeck-bound}
\end{equation}
The estimate is invariant under \eqref{eq:v3v97-rescaling}.
\end{lemma}

\begin{proof}
This is the small-curvature Coulomb-gauge construction: solve the gauge
equation by the inverse function theorem around the trivial connection,
using \(W^{1,2}\hookrightarrow L^4\) to bound the quadratic commutator.
The linear elliptic estimate controls \(\nabla A\) by \(F_A\); the smallness
in \eqref{eq:v3v97-small-energy} absorbs \(A\wedge A\).  Sobolev embedding
then gives the \(L^4\) term.  Both sides are dimensionless in four
dimensions, proving scale invariance.
\end{proof}

\subsection{The critical weak Schatten estimate}
\label{sec:v3v97-Schatten}

Fix on the unit ball a reference Dirac operator \(D_0\) with a local elliptic
boundary condition or a spectral section.  Remove its finite zero band and
write \(R_0\) for the reduced inverse.  For a unitary representation \(R\),
Clifford multiplication by the connection is denoted \(c(A)\), and
\begin{equation}
 K_A=R_0c(A).
\label{eq:v3v97-KA}
\end{equation}

\begin{lemma}[Four-dimensional resolvent ideal]
\label{lem:v3v97-weak-S4}
If \(A\in L^4(B_1)\), then
\begin{equation}
 K_A\in\mathcal S_{4,\infty},
 \qquad
 \sup_{n\ge1}n^{1/4}s_n(K_A)
 \le C_R\|A\|_{L^4},
\label{eq:v3v97-weak-S4}
\end{equation}
where \(s_n\) are singular values.  Consequently, for every \(p>4\),
\begin{equation}
 \|K_A\|_{\mathcal S_p}
 \le C_{p,R}\|A\|_{L^4}.
\label{eq:v3v97-Sp}
\end{equation}
The constants are uniform under the rescaling
\eqref{eq:v3v97-rescaling} and bounded perturbations of the ball metric.
\end{lemma}

\begin{proof}
The reduced first-order elliptic inverse maps \(L^2\) to \(H^1\).  Decompose
the reference Laplacian into dyadic spectral shells.  Weyl's upper bound in
four dimensions gives \(O(2^{4j})\) modes below frequency \(2^j\), while
the inverse Dirac factor contributes \(2^{-j}\).  Multiplication by
\(A\in L^4\), followed by the Sobolev estimate
\(H^1\hookrightarrow L^4\) and Holder's inequality
\(L^4\cdot L^4\subset L^2\), gives approximation numbers bounded by
\(C\|A\|_4n^{-1/4}\).  This is
\eqref{eq:v3v97-weak-S4}.

If \(s_n\le Cn^{-1/4}\), then
\(\sum_ns_n^p\le C^p\sum_nn^{-p/4}<\infty\) for \(p>4\), proving
\eqref{eq:v3v97-Sp}.  The rescaled elliptic and Sobolev constants are
uniform on the bounded-geometry class.
\end{proof}

The exponent four is sharp by the same Weyl count.  Thus a fourth-order
ordinary determinant expansion is marginal; one more occurrence is needed
for an absolutely trace-class remainder.

\subsection{Exact Born subtraction and the fifth determinant}
\label{sec:v3v97-det5}

For \(K\in\mathcal S_5\), define
\begin{equation}
 \det{}_5(1+K)=
 \det\left((1+K)
 \exp\left[-K+\frac{K^2}{2}-\frac{K^3}{3}+\frac{K^4}{4}\right]
 \right).
\label{eq:v3v97-det5}
\end{equation}
The operator inside the ordinary Fredholm determinant differs from one by a
trace-class operator.

\begin{lemma}[Regularized determinant bound]
\label{lem:v3v97-det5-bound}
There is a universal \(C_5\) such that
\begin{equation}
 |\det{}_5(1+K)|
 \le\exp(C_5\|K\|_{\mathcal S_5}^5).
\label{eq:v3v97-det5-bound}
\end{equation}
No inverse of \(1+K\) is required.
\end{lemma}

\begin{proof}
For each eigenvalue \(z\) of \(K\), the scalar factor is
\((1+z)\exp(-z+z^2/2-z^3/3+z^4/4)\).  Its logarithmic modulus is bounded
above by \(C|z|^5\), first on \(|z|\le1/2\) by the Taylor remainder and
then on the complement by enlarging \(C\).  Apply the canonical-product
definition of \(\det{}_5\), sum the bound over eigenvalues, and use Weyl's
inequality
\(\sum|\lambda_j(K)|^5\le\sum s_j(K)^5\).
\end{proof}

For smooth \(A\), the determinant factorization reads schematically but
exactly
\begin{equation}
 \operatorname{Det}(D_0+c(A))=
 \mathcal Z_{\rm zero}(A),
 \exp\left(sum_{j=1}^4\frac{(-1)^{j+1}}j
 \operatorname{TR}K_A^j+\mathcal C_{\rm loc}(A)
ight)
 \det{}_5(1+K_A).
\label{eq:v3v97-exact-factorization}
\end{equation}
Here \(\operatorname{TR}\) is the common finite-part trace, the local term
\(\mathcal C_{\rm loc}\) fixes its gauge-covariant renormalization, and
\(\mathcal Z_{\rm zero}\) is the finite exterior-power section of the
removed spectral band.  For each \(1\le j\le4\), the exact V95 heat cutoff
splits \(\operatorname{TR}K_A^j\) into its local short-time coefficient
and a complementary physical IR Born functional.  Only the former belongs
to \(\mathfrak M_{\le0}\eta^{\rm UV}\); the latter remains part of the
physical determinant.

\begin{theorem}[One-ball determinant remainder]
\label{thm:v3v97-one-ball}
On a ball satisfying \eqref{eq:v3v97-small-energy}, the nonlocal determinant
remainder obeys
\begin{equation}
 |\det{}_5(1+K_A)|
 \le\exp\left(C_R E(B_\rho)^{5/2}\right).
\label{eq:v3v97-one-ball-bound}
\end{equation}
The bound is independent of \(\rho\).  It remains valid when
\(1+K_A\) has a kernel; the corresponding zero-mode data occur only in
\(\mathcal Z_{\rm zero}\).
\end{theorem}

\begin{proof}
Combine Lemma~\ref{lem:v3v97-Uhlenbeck-ball},
\eqref{eq:v3v97-Sp} at \(p=5\), and
Lemma~\ref{lem:v3v97-det5-bound}.  Critical scaling removes every power of
\(\rho\).  The determinant bound is an upper bound for a canonical product
and never divides by \(1+K_A\), so zero modes cause no singularity in
\eqref{eq:v3v97-one-ball-bound}.
\end{proof}

\begin{firewall}[Zero modes are a line, not a bounded scalar inverse]
\label{fw:v3v97-zero-line}
In a nonzero instanton sector the massless chiral determinant can vanish
unless external fermion insertions saturate its zero modes.  V97 does not
replace \(\mathcal Z_{\rm zero}\) by a nonzero scalar or assume a reduced
Dirac gap.  Its index and transition line are treated by V88--V94; the
Schatten bound controls only the complementary determinant occurrence.
\end{firewall}

\subsection{Stopping-ball product and action absorption}
\label{sec:v3v97-stopping-cover}

Choose a Vitali subcover of the concentration set by disjoint inner balls
\(B_i\) whose fixed dilates have bounded overlap and whose energies satisfy
\begin{equation}
 E_i:=\int_{B_i}|F|^2\le\varepsilon,
 \qquad0<\varepsilon\le\varepsilon_U.
\label{eq:v3v97-stopping-energy}
\end{equation}
At finite lattice spacing the reconstructed curvature is smooth, so the
radii can be stopped when the energy first reaches \(\varepsilon\); a final
lower-energy ball covers each residual point.

\begin{theorem}[Multi-bubble determinant absorption]
\label{thm:v3v97-bubble-absorption}
For the product of fifth regularized determinants over the stopping balls,
\begin{equation}
 \prod_i|\det{}_5(1+K_{A,i})|
 \le
 \exp\left(C_R\varepsilon^{3/2}
 \sum_iE_i\right)
 \le
 \exp\left(C'_R\varepsilon^{3/2}|F\|_{L^2(X)}^2
ight).
\label{eq:v3v97-product-bound}
\end{equation}
For the finite list of Standard Model representations the same estimate
holds with \(C_R\) replaced by \(C_{\rm SM}\).  If the Yang--Mills action
has lower bound
\begin{equation}
 S_{\rm YM}(A)\ge c_g\|F\|_2^2
\label{eq:v3v97-action-lower}
\end{equation}
and \(C_{\rm SM}\varepsilon^{3/2}<c_g/2\), the complete nonlocal bubble
remainder is absorbed by half of the bosonic action weight.
\end{theorem}

\begin{proof}
Theorem~\ref{thm:v3v97-one-ball} gives an exponent
\(CE_i^{5/2}\).  Since \(E_i\le\varepsilon\),
\begin{equation}
 E_i^{5/2}\le\varepsilon^{3/2}E_i.
\end{equation}
Sum over the balls.  Bounded overlap gives
\(\sum_iE_i\le C\|F\|_2^2\), proving
\eqref{eq:v3v97-product-bound}.  There are finitely many matter
representations and generations, so their constants add.  The last
inequality gives
\begin{equation}
 e^{-S_{\rm YM}}
 \prod_{i,f}|\det{}_5(1+K_{A,i,f})|
 \le e^{-(c_g/2)\|F\|_2^2},
\end{equation}
which is the asserted absorption.
\end{proof}

This is an upper-integrability result for the fifth determinant only.  It
does not require a fermion mass gap and is compatible with determinant
zeros.  The first four Born terms are absent from
\eqref{eq:v3v97-product-bound}; their local parts are assigned to V95's UV
occurrence, while their complementary IR parts remain to be bounded as
physical lower-order functionals.

\subsection{Bubble charge and the Standard Model phase}
\label{sec:v3v97-bubble-charge}

On a sequence of shrinking balls with bounded energy, Uhlenbeck compactness
gives, after gauge and subsequence, a limiting connection away from finitely
many points.  The lost second Chern number is an integer bubble charge.

\begin{proposition}[Exact separation of bubble topology and Born sectors]
\label{prop:v3v97-bubble-topology}
For a geometrically tame admissible sequence, the determinant data on a
concentration ball split into:
\begin{enumerate}
 \item the finite zero-mode exterior line with index equal to the
 reconstructed integer bubble charge;
 \item the local UV coefficients of the first four Born terms, assigned to
 V95's normalization;
 \item the complementary IR parts of those same four Born terms, retained as
 physical determinant data; and
 \item the fifth determinant bounded by
 \eqref{eq:v3v97-one-ball-bound}.
\end{enumerate}
The four factors arise from the single exact factorization
\eqref{eq:v3v97-exact-factorization} followed by the exact V95 heat
partition; none may be counted in another item.  For the complete Standard
Model representation, the combined phase of items 1--3 has no residual
gauge or mixed gauge--gravitational bordism character after the V66 and V94
normalizations.
\end{proposition}

\begin{proof}
The spectral section splits the finite crossing band from its complement.
Its exterior power is item 1 and V96 identifies its index with the
reconstructed characteristic number.  Expanding the logarithm through
order four and applying \(\chi+(1-\chi)=1\) gives the disjoint items 2 and
3.  The canonical-product remainder is exactly item 4.  This proves the
exhaustive factorization.  V66 cancels the local anomaly of item 2, while
V87--V95 identify and trivialize the complete Standard Model physical
holonomy carried by items 1 and 3, proving the phase statement.
\end{proof}

\begin{theorem}[Partial discharge of the V96 bubble-integrability item]
\label{thm:v3v97-V96-item4}
Assume the admissible reconstruction and stopping cover above, the action
lower bound \eqref{eq:v3v97-action-lower}, and the V95 normalization of the
local UV occurrences.  Then the fifth determinant remainder on all
concentration balls is uniformly integrable under the bosonic weight.  The
remaining parts of item 4 in
Theorem~\ref{thm:v3v96-measure-criterion} are the four lower-order IR Born
functionals, the finite zero-mode exterior sections, and the bosonic bubble
distribution.
\end{theorem}

\begin{proof}
Theorem~\ref{thm:v3v97-bubble-absorption} supplies an integrable dominating
function independent of the number and radii of stopping balls.  Dominated
convergence applies to the fifth determinants wherever the reconstructed
connections converge weakly and then strongly below the critical exponent.
Proposition~\ref{prop:v3v97-bubble-topology} shows that the lower Born
functionals, zero-mode lines, and bosonic distribution are not included in
that fifth determinant, which proves the stated boundary.
\end{proof}

\begin{firewall}[Admissibility probability is still open]
\label{fw:v3v97-admissibility-probability}
The stopping-ball estimate begins after a link field has entered the
reconstruction domain.  It does not prove that the interacting lattice
measure satisfies
\(\mu_a\{\max_p\|1-U_p\|\ge\epsilon_{\rm rec}\}\to0\).  A global action
expectation alone controls an average plaquette defect, not its maximum over
\(O(a^{-4})\) plaquettes.
\end{firewall}

\begin{firewall}[Bosonic continuum and bubble distribution remain]
\label{fw:v3v97-bosonic-open}
Uniform integrability of the fermion remainder does not construct the
four-dimensional Yang--Mills measure, prove a mass gap, or control the joint
Einstein metric measure.  Those bosonic limits are independent major gates
of the full SM--Einstein continuum.
\end{firewall}

\begin{firewall}[No reflection-positive conclusion]
\label{fw:v3v97-RP}
The regularized determinant inequality is an absolute-value bound.  It does
not provide the cross-reflection positive cone for the finite fermion
action.  Taking a determinant modulus before a reflection decomposition
cannot prove Osterwalder--Schrader positivity.
\end{firewall}

\subsection{V97 conclusion and V98 target}
\label{sec:v3v97-conclusion}

The fermionic concentration-ball remainder now has a scale-invariant and
action-absorbable bound.  Critical Uhlenbeck gauge gives \(A\in L^4\), the
Dirac perturbation lies in weak \(\mathcal S_4\), and the fifth regularized
determinant is controlled after exactly four local Born occurrences are
removed.  Zero modes, local UV terms, and the nonlocal determinant remainder
are kept as three disjoint factors.  This discharges the fifth-order
nonzero-mode part of the V96 bubble-integrability criterion and isolates the
remaining lower-order physical terms.

V98 should analyze the quadratic, cubic, and quartic IR Born functionals on
a stopping ball before attacking maximum-plaquette probability.  It should
use gauge covariance and the complete Standard Model signed sum before
taking absolute values, determine which kernels are controlled by
\(L^2\)-curvature and which require additional spectral input, and retain
the linear term's exact Ward cancellation.  Only after those physical
lower-order occurrences are bounded can the admissibility-probability gate
be assessed without hiding determinant mass in the bosonic action.

\clearpage
\section{V98: the lower Born packet and the logarithmic curvature moment}
\label{sec:v3v98}

V97 controls the fifth regularized determinant but correctly leaves the
first four physical infrared Born terms open.  They cannot be estimated as
four unrelated traces.  Ward identities combine their longitudinal pieces,
and the quadratic, cubic, and quartic terms form one gauge-covariant packet.
This note proves a bound for that packet in terms of a logarithmically
weighted curvature moment and lower powers of the critical energy.  A
single transverse plane wave then proves that ordinary \(L^2\)-curvature,
plaquette admissibility, and Wilson action do not control the logarithmic
moment.  The missing estimate is physical vacuum polarization, not an
anomaly coefficient.

\subsection{The exact lower-order IR packet}
\label{sec:v3v98-packet}

Retain the unit-ball spectral-section convention of V97.  Let
\(K_A=R_0c(A)\) and apply V95's complementary heat weight
\(1-\chi(t/a^2)\) to each finite-part trace.  Denote the resulting physical
parts by \(B_{j,a}^{\rm IR}(A)\), \(1\le j\le4\), and put
\begin{equation}
 \mathcal B_{\le4,a}^{\rm IR}(A)
 =\sum_{j=1}^4\frac{(-1)^{j+1}}jB_{j,a}^{\rm IR}(A)
 +\mathcal C_{\le4,a}^{\rm Ward}(A).
\label{eq:v3v98-lower-packet}
\end{equation}
The last term is the finite local completion required to impose the chosen
renormalized Ward identities.  It is fixed by the same convention as V95
and is not a fifth occurrence.

\begin{lemma}[The linear Standard Model occurrence]
\label{lem:v3v98-linear}
For the complete Standard Model representation, the gauge-dependent linear
term in \(\mathcal B_{\le4,a}^{\rm IR}\) vanishes.  More precisely,
nonabelian generators have zero representation trace, while the abelian
coefficient is proportional to the complete chiral charge trace already
cancelled in the mixed gauge--gravitational anomaly arithmetic of V66.
Gauge-singlet Yukawa mass insertions contribute only to the fixed scalar
normalization and not to a gauge tadpole.
\end{lemma}

\begin{proof}
At first order the internal factor is one generator.  Its trace vanishes in
every semisimple representation.  For hypercharge it is
\(\sum_f\dim(R_{3,f})\dim(R_{2,f})Y_f\), with right-handed fields written
as left-handed conjugates.  This is the coefficient of the mixed
gravitational--hypercharge anomaly and is zero by V66.  A gauge-singlet
endomorphism has no gauge generator, proving the last statement.
\end{proof}

The cubic triangle requires a different statement.  V66 cancels its
gauge-variant descent, but a gauge-invariant cubic physical amplitude need
not vanish.  Likewise the quartic box is not an independent positive
quantity.  We therefore keep all three nonlinear orders inside
\eqref{eq:v3v98-lower-packet}.

\subsection{The gauge-covariant logarithmic norm}
\label{sec:v3v98-log-norm}

On the rescaled ball in Coulomb gauge, let \(\Delta_A\) be the nonnegative
covariant Hodge Laplacian on adjoint-valued two-forms with the inherited
elliptic boundary condition.  For a fixed renormalization scale \(\mu>0\),
define
\begin{equation}
 \mathscr L_\mu(F)=
 \left\langle F,
 \log\left(1+\frac{\Delta_A}{\mu^2}\right)F
 \right\rangle_{L^2}.
\label{eq:v3v98-log-moment}
\end{equation}
It is nonnegative and gauge invariant.  It is finite for every smooth
reconstructed connection but is not bounded on an \(L^2\)-curvature ball.

\begin{theorem}[Ward-assembled lower Born estimate]
\label{thm:v3v98-lower-bound}
There is \(\varepsilon_0\le\varepsilon_U\) such that, whenever
\(E=\|F_A\|_2^2\le\varepsilon_0\), the complete renormalized packet in a
unitary representation \(R\) satisfies
\begin{equation}
 |\mathcal B_{\le4,a}^{\rm IR}(A)|
 \le C_R\left[
 \mathscr L_\mu(F_A)+E+E^{3/2}+E^2
 \right]+C_Ra.
\label{eq:v3v98-lower-bound}
\end{equation}
The estimate is uniform under ball rescaling and bounded-geometry
perturbations.  In the complete Standard Model sum the constants add; the
logarithmic term does not cancel.
\end{theorem}

\begin{proof}
Work first with a smooth compactly supported connection in Coulomb gauge.
The renormalized two-point current kernel is transverse by the exact Ward
identity.  Its principal order-zero symbol on two-forms is
\begin{equation}
 c_R\log(1+|\xi|^2/\mu^2),\Pi_{\rm tr}(\xi),
\label{eq:v3v98-polarization-symbol}
\end{equation}
where \(c_R\ge0\) is the quadratic Dynkin index and
\(\Pi_{\rm tr}\) is the transverse projector.  Covariant symbol calculus
therefore writes the quadratic term as
\begin{equation}
 c_R\left\langle F,
 \log(1+\Delta_A/\mu^2)F\right\rangle
 +\langle F,R_{0,A}F\rangle,
\label{eq:v3v98-quadratic-decomposition}
\end{equation}
with \(R_{0,A}\) bounded on \(L^2\) for
\(\|A\|_4\) sufficiently small.  This gives the first two terms on the
right of \eqref{eq:v3v98-lower-bound}.

At orders three and four, apply the Ward identity before estimating.  Every
longitudinal connection insertion either combines with a neighbouring
insertion to make a curvature or belongs to the cancelled local triangle
descent.  The remaining multilinear symbols are order zero after two
curvature factors have been extracted.  Holder's inequality, the
Calderon--Zygmund bound for those symbols, and Uhlenbeck's estimate
\(\|A\|_4\le C E^{1/2}\) give respectively
\begin{equation}
 C\|F\|_2^2\|A\|_4
 \le CE^{3/2},
 \qquad
 C\|F\|_2^2\|A\|_4^2
 \le CE^2.
\label{eq:v3v98-cubic-quartic}
\end{equation}
The anomaly-free representation sum removes the nontrivial local triangle
descent, but it does not remove the transverse two-point symbol because
\(c_R\) is a positive quadratic index.  The V95 heat-parametrix remainder is
\(O(a)\).  Density of smooth connections in the small-energy Coulomb slice
extends the estimate to \(W^{1,2}\), proving the theorem.
\end{proof}

\begin{remark}[Sign information and the claim made here]
In schemes with a positive spectral representation, the leading quadratic
term in the Euclidean effective action is suppressive.  V98 does not use
that sign to discard the entire nonlinear packet: the cubic and quartic
covariant completion and the chiral phase must remain assembled.  The
absolute estimate \eqref{eq:v3v98-lower-bound} is the scheme-independent
claim used below.
\end{remark}

\subsection{Sharp failure of an energy-only estimate}
\label{sec:v3v98-plane-wave}

Consider a flat unit four-torus and choose a fixed abelian generator
\(T\in\mathfrak g_{\rm SM}\).  For an integer \(N\), set
\begin{equation}
 A_N=\frac{\alpha}{N}T\sin(Nx_1)\,dx_2,
 \qquad
 F_N=\alpha T\cos(Nx_1)\,dx_1\wedge dx_2.
\label{eq:v3v98-plane-wave}
\end{equation}

\begin{proposition}[Logarithmic-moment counterexample]
\label{prop:v3v98-log-counterexample}
For fixed nonzero \(\alpha\),
\begin{align}
 \|F_N\|_2^2&=C_T\alpha^2,
 \label{eq:v3v98-fixed-energy}\\
 \mathscr L_\mu(F_N)
 &=C_T\alpha^2\log(1+N^2/\mu^2).
 \label{eq:v3v98-growing-log}
\end{align}
Hence no function of \(\|F\|_2\) alone bounds
\(\mathscr L_\mu(F)\).  If the field is sampled on a lattice with
\(aN\le c<\pi\), then its plaquette defect is \(O(a^2|\alpha|)\), uniformly
in \(N\).  Thus quantitative plaquette admissibility does not repair the
failure.
\end{proposition}

\begin{proof}
The curvature formula is exact because the chosen generator is abelian.
Orthogonality of the Fourier mode gives
\eqref{eq:v3v98-fixed-energy}.  It is an eigenform of the flat Laplacian
with eigenvalue \(N^2\), which gives
\eqref{eq:v3v98-growing-log} by functional calculus.  Edge holonomy around
one plaquette equals
\(\exp(a^2F_{12}+O(a^3N|\alpha|))\).  Under \(aN\le c\), this is bounded by
\(Ca^2|\alpha|\), proving admissibility for small \(a\).
\end{proof}

The example is smooth, has no instanton charge, and lies in a small abelian
sector.  It therefore tests the physical two-point kernel rather than the
topological or rough-bundle machinery.

\begin{corollary}[No action-only absolute bubble bound]
\label{cor:v3v98-no-action-only}
An estimate of the form
\begin{equation}
 |\mathcal B_{\le4,a}^{\rm IR}(A)|
 \le C\bigl(1+\|F_A\|_2^p\bigr)
\label{eq:v3v98-false-energy-bound}
\end{equation}
with constants uniform in the ultraviolet frequency cannot be proved by
bounding the quadratic packet in absolute value.  Any correct upper
integrability argument must use its sign, a running-coupling compensation,
or a measure estimate for \(\mathscr L_\mu(F)\).
\end{corollary}

\begin{proof}
Insert \eqref{eq:v3v98-plane-wave}.  The right side of
\eqref{eq:v3v98-false-energy-bound} is independent of \(N\), whereas the
quadratic principal term in
\eqref{eq:v3v98-quadratic-decomposition} grows as \(\log N\).  The cubic
and quartic terms are \(O(\alpha^3)+O(\alpha^4)\); choose \(\alpha\) fixed
and sufficiently small.  They cannot cancel the unbounded quadratic term.
\end{proof}

\subsection{Running coupling versus the physical logarithm}
\label{sec:v3v98-running}

The ultraviolet part of the two-point function contains the local logarithm
which renormalizes the gauge kinetic coefficient.  Write
\begin{equation}
 \frac1{g_0(a)^2}\|F\|_2^2
 +b_{\rm SM}\log(a\mu)\|F\|_2^2
 =\frac1{g_R(\mu)^2}\|F\|_2^2.
\label{eq:v3v98-running-coupling}
\end{equation}
This identity absorbs the local cutoff logarithm in V95's UV jet.  It does
not remove the spectral logarithm
\(\log(1+\Delta_A/\mu^2)\) in the IR packet.

\begin{proposition}[Exact allocation of the quadratic logarithms]
\label{prop:v3v98-log-allocation}
With one common heat cutoff, the quadratic determinant occurrence decomposes
into:
\begin{enumerate}
 \item the local coefficient in \eqref{eq:v3v98-running-coupling}; and
 \item the gauge-covariant physical operator in
 \eqref{eq:v3v98-log-moment}.
\end{enumerate}
Changing \(\mu\) transfers a local multiple of \(\|F\|_2^2\) between the
two terms and leaves their sum invariant.  It cannot turn item 2 into a
second local action payment.
\end{proposition}

\begin{proof}
The scalar identity
\begin{equation}
 \log\frac{\lambda}{a^{-2}}
 =\log\frac{\lambda}{\mu^2}+\log(a^2\mu^2)
\end{equation}
splits the two-point symbol into a spectral part and a momentum-independent
part.  Functional calculus gives item 1 from the constant and item 2 from
the spectral logarithm.  Replacing \(\mu\) changes the latter by a constant
on two-forms and the former by the opposite constant, proving the claim.
\end{proof}

\subsection{The precise shared bosonic interface}
\label{sec:v3v98-interface}

\begin{theorem}[Conditional completion of lower-Born integrability]
\label{thm:v3v98-conditional-integrability}
Suppose the reconstructed gauge-field measures satisfy, on every stopping
cover and for some \(\delta>0\),
\begin{equation}
 \sup_a\int
 \exp\left(\delta\sum_i\mathscr L_\mu(F_{A,i})\right)
 \exp(-S_{\rm YM}(A))\,dA<\infty.
\label{eq:v3v98-log-integrability}
\end{equation}
Then the complete lower Born packet is uniformly integrable, after choosing
the representation-dependent Holder exponents and combining
\eqref{eq:v3v98-lower-bound} with the small-energy stopping condition.
Together with V97 this verifies the nonzero-mode fermion determinant part of
item 4 in Theorem~\ref{thm:v3v96-measure-criterion}.
\end{theorem}

\begin{proof}
On each ball, \eqref{eq:v3v98-lower-bound} bounds the lower packet by the log
moment plus
\(E_i+E_i^{3/2}+E_i^2\).  Since \(E_i\le\varepsilon_0\), the latter terms
are bounded by \(C E_i\) and are absorbed into a fixed fraction of the
Yang--Mills action as in V97.  Use Holder's inequality with
\eqref{eq:v3v98-log-integrability} for the log moments.  V97 supplies the
same control for the fifth determinant.  Multiplication of the two uniformly
integrable bounds proves the result.
\end{proof}

\begin{firewall}[Anomaly cancellation does not cancel vacuum polarization]
\label{fw:v3v98-not-anomaly}
V66 cancels the gauge variation of the chiral triangle and the mixed
anomaly coefficients.  The quadratic coefficient in
\eqref{eq:v3v98-polarization-symbol} is a sum of positive quadratic indices
and is nonzero in the Standard Model.  Removing it by anomaly arithmetic
would erase physical running of the gauge couplings.
\end{firewall}

\begin{firewall}[The logarithmic moment estimate is not yet proved]
\label{fw:v3v98-log-open}
Neither Wilson action control nor Uhlenbeck \(L^2\)-compactness proves
\eqref{eq:v3v98-log-integrability};
Proposition~\ref{prop:v3v98-log-counterexample} rules out a deterministic
energy-only substitute.  The missing input is a signed or probabilistic
high-frequency estimate for the interacting gauge measure.  It is a precise
interface with the Yang--Mills continuum programme.
\end{firewall}

\begin{firewall}[Reflection positivity remains separate]
\label{fw:v3v98-RP}
The Ward-assembled packet and logarithmic norm are continuum determinant
objects.  Their bounds do not factor the finite V93 action into a positive
cross-reflection cone.
\end{firewall}

\subsection{V98 conclusion and V99 target}
\label{sec:v3v98-conclusion}

The four lower Born orders now have one gauge-covariant estimate rather than
four unrelated trace bounds.  Their only new critical stock is the physical
logarithmic curvature moment.  A transverse abelian plane wave proves that
this stock is invisible to ordinary Yang--Mills energy and to pointwise
plaquette admissibility.  Local running-coupling renormalization and the
physical spectral logarithm are exactly separated, so neither can pay the
other a second time.

V99 should test whether the needed exponential log-moment follows from a
multiscale Gaussian domination or reflection-positive current estimate for
the bosonic measure.  It should first prove the estimate in the free
gauge-field measure with the correct asymptotically free covariance, then
identify the exact interaction comparison needed for Yang--Mills.  Failure
in the free model would falsify the criterion; success only in the free
model must be recorded as a conditional interface, not as construction of
the interacting continuum.

\clearpage
\section{V99: Gaussian falsification of a separate log-moment bound and the signed quadratic repair}
\label{sec:v3v99}

V98 isolates the logarithmic curvature moment required by an absolute bound
on the lower Born packet and proposes exponential integrability as a
conditional interface.  The required test in the free gauge measure is
negative.  There are \(O(\Lambda^4)\) curvature modes below cutoff
\(\Lambda\), and each contributes a logarithmic weight.  Even after
centering, no cutoff-uniform exponential moment survives.  This is not a
failure of the determinant construction.  It shows that an extensive
fermion polarization determinant must be combined with the bosonic
quadratic form and normalized as one Gaussian occurrence.  The exact
combined covariance criterion is proved below.  Its one-loop signs expose
the separate hypercharge ultraviolet obstruction and motivate a punctured,
rather than maximum-plaquette, reconstruction.

\subsection{Free transverse curvature measure}
\label{sec:v3v99-free}

Work on a flat unit four-torus in Coulomb gauge and truncate Fourier modes to
\(0<|k|\le\Lambda\).  After removing gauge zero modes, diagonalize the real
transverse curvature coordinates as \(Z_1,\ldots,Z_{N_\Lambda}\), with
\begin{equation}
 d\mu_{g,\Lambda}(Z)=
 \prod_{j=1}^{N_\Lambda}
 \frac{1}{\sqrt{2\pi g^2}}
 e^{-Z_j^2/(2g^2)},dZ_j,
 \qquad
 N_\Lambda\asymp C\Lambda^4.
\label{eq:v3v99-Gaussian-measure}
\end{equation}
Up to fixed polarization and Lie-algebra multiplicities, the V98 log moment
is
\begin{equation}
 L_{\mu,\Lambda}(Z)=
 \sum_{j=1}^{N_\Lambda}w_jZ_j^2,
 \qquad
 w_j=\log(1+|k_j|^2/\mu^2).
\label{eq:v3v99-discrete-log-moment}
\end{equation}

\begin{lemma}[Exact Gaussian moment formula]
\label{lem:v3v99-Gaussian-mgf}
For \(\delta>0\),
\begin{equation}
 \mathbb E_{g,\Lambda}e^{\delta L_{\mu,\Lambda}}
 =\prod_{j=1}^{N_\Lambda}
 (1-2\delta g^2w_j)^{-1/2}
\label{eq:v3v99-mgf}
\end{equation}
if \(2\delta g^2\max_jw_j<1\), and the expectation is infinite otherwise.
Furthermore
\begin{align}
 \mathbb E L_{\mu,\Lambda}
 &=g^2\sum_jw_j,
 \label{eq:v3v99-mean}\\
 \operatorname{Var}L_{\mu,\Lambda}
 &=2g^4\sum_jw_j^2.
 \label{eq:v3v99-variance}
\end{align}
\end{lemma}

\begin{proof}
For one real centred Gaussian of variance \(g^2\), direct integration gives
\begin{equation}
 \mathbb E e^{\delta wZ^2}=(1-2\delta g^2w)^{-1/2}
\end{equation}
exactly when the remaining quadratic coefficient is positive.  Independence
gives \eqref{eq:v3v99-mgf}.  Differentiating its logarithm once and twice at
\(\delta=0\) gives \eqref{eq:v3v99-mean} and
\eqref{eq:v3v99-variance}.
\end{proof}

Weyl counting gives
\begin{align}
 \sum_jw_j&\asymp C\Lambda^4\log\Lambda,
 \label{eq:v3v99-weight-sum}\\
 \sum_jw_j^2&\asymp C\Lambda^4(\log\Lambda)^2.
 \label{eq:v3v99-weight-square-sum}
\end{align}

\begin{theorem}[Failure of the V98 exponential-moment hypothesis in the free field]
\label{thm:v3v99-free-falsification}
For fixed \(g>0\), every \(\delta>0\) violates the finiteness condition in
Lemma~\ref{lem:v3v99-Gaussian-mgf} for all sufficiently large \(\Lambda\).
More generally, if \(g_\Lambda^2\asymp1/\log\Lambda\), then the uncentred
expectation in \eqref{eq:v3v99-mgf} is not uniformly bounded for any fixed
\(\delta>0\).  The centred variables are not exponentially tight either:
for every fixed sufficiently small \(\delta>0\),
\begin{equation}
 \log\mathbb E\exp\left[
 \delta(L_{\mu,\Lambda}-\mathbb EL_{\mu,\Lambda})
 \right]
 \ge c\delta^2g_\Lambda^4
 \sum_jw_j^2\longrightarrow\infty.
\label{eq:v3v99-centered-divergence}
\end{equation}
Thus the hypothesis \eqref{eq:v3v98-log-integrability}, interpreted as a
cutoff-uniform exponential moment of the separate determinant functional,
fails already in the free asymptotically scaled Gaussian model.
\end{theorem}

\begin{proof}
Since \(\max_jw_j\sim2\log\Lambda\), fixed \(g\) eventually violates
\(2\delta g^2\max w<1\).  If
\(g_\Lambda^2\asymp1/\log\Lambda\) and the finiteness condition happens to
hold for small \(\delta\), use
\(-\log(1-x)\ge x\) in \eqref{eq:v3v99-mgf}.  Equations
\eqref{eq:v3v99-mean} and \eqref{eq:v3v99-weight-sum} show that its logarithm
grows at least like \(c\delta\Lambda^4\).

For the centred expression, cancel the linear term in the logarithm and use
\(-\log(1-x)-x\ge x^2/4\) for small nonnegative \(x\).  This gives the lower
bound in \eqref{eq:v3v99-centered-divergence}.  With
\eqref{eq:v3v99-weight-square-sum} and
\(g_\Lambda^4\asymp(\log\Lambda)^{-2}\), the right side grows like
\(c\delta^2\Lambda^4\).
\end{proof}

\begin{firewall}[Typical free fields do not have bounded classical action]
\label{fw:v3v99-distribution-support}
The mean free curvature action is proportional to \(N_\Lambda\), not to a
cutoff-independent \(L^2\) bound.  In the continuum the Gaussian gauge field
is a distribution.  Therefore V96's bounded classical-action reconstruction
sets cannot be assumed to have probability tending to one under a quantum
field measure; they describe controlled backgrounds and large-deviation
sectors, not typical configurations.
\end{firewall}

\subsection{The determinant and Gaussian action are one quadratic occurrence}
\label{sec:v3v99-combined}

Let \(C_\Lambda>0\) be a finite covariance matrix for transverse gauge
coordinates, and let \(\Pi_\Lambda=\Pi_\Lambda^*\) be the complete
renormalized quadratic polarization operator, including fermions, gauge and
ghost fields, counterterms, and any gravitational contribution which is
claimed at that scale.  Consider the normalized reweighting
\begin{equation}
 d\nu_\Lambda(Z)=
 \frac{
 e^{-\frac12\langle Z,\Pi_\Lambda Z\rangle}
 d\mu_{C_\Lambda}(Z)}{
 \mathbb E_{C_\Lambda}
 e^{-\frac12\langle Z,\Pi_\Lambda Z\rangle}}.
\label{eq:v3v99-combined-measure}
\end{equation}

\begin{theorem}[Exact signed Gaussian covariance criterion]
\label{thm:v3v99-covariance-criterion}
The denominator in \eqref{eq:v3v99-combined-measure} is finite and nonzero
if and only if
\begin{equation}
 C_\Lambda^{-1}+\Pi_\Lambda>0.
\label{eq:v3v99-positive-quadratic-form}
\end{equation}
When this holds,
\begin{align}
 \mathbb E_{C_\Lambda}
 e^{-\frac12\langle Z,\Pi_\Lambda Z\rangle}
 &=\det(1+C_\Lambda^{1/2}Pi_\Lambda
 C_\Lambda^{1/2})^{-1/2},
 \label{eq:v3v99-Gaussian-ratio}\\
 \operatorname{Cov}_{\nu_\Lambda}
 &=(C_\Lambda^{-1}+\Pi_\Lambda)^{-1}.
 \label{eq:v3v99-renormalized-covariance}
\end{align}
All extensive mode products belong to the normalizing determinant in
\eqref{eq:v3v99-Gaussian-ratio}; no separate exponential moment of
\(|\Pi_\Lambda|\) is required.
\end{theorem}

\begin{proof}
Combine the two quadratic exponents.  The finite-dimensional Gaussian
integral converges precisely when the total real quadratic form is positive.
Diagonalizing
\(C_\Lambda^{1/2}\Pi_\Lambda C_\Lambda^{1/2}\) gives the determinant ratio
in \eqref{eq:v3v99-Gaussian-ratio}.  Completing the square with a linear
source and differentiating its generating function gives
\eqref{eq:v3v99-renormalized-covariance}.
\end{proof}

This theorem is the signed repair of V98's absolute criterion.  A positive
polarization direction improves the covariance; a negative direction must
be paid by the bare bosonic inverse covariance at the same momentum.  Taking
absolute values before adding the two forms destroys this distinction.

\begin{corollary}[Modewise logarithmic condition]
\label{cor:v3v99-modewise}
If the ultraviolet transverse symbols have the scalar form
\begin{equation}
 C_\Lambda(k)^{-1}=g_0(\Lambda)^{-2},
 \qquad
 \Pi_\Lambda(k)=b_{\rm tot}\log(|k|/\mu)+O(1),
\label{eq:v3v99-mode-symbol}
\end{equation}
then a necessary and sufficient high-frequency Gaussian condition is
\begin{equation}
 \inf_{0<|k|\le\Lambda}
 \left[g_0(\Lambda)^{-2}
 +b_{\rm tot}\log(|k|/\mu)+O(1)\right]>0.
\label{eq:v3v99-mode-positivity}
\end{equation}
The sign of \(b_{\rm tot}\), not an absolute log-moment bound, decides the
quadratic ultraviolet stability.
\end{corollary}

\subsection{What the perturbative Standard Model signs do and do not prove}
\label{sec:v3v99-SM-signs}

Use the convention
\begin{equation}
 \mu\frac{dg_i}{d\mu}=-\frac{b_i}{16\pi^2}g_i^3+O(g_i^5).
\label{eq:v3v99-beta-convention}
\end{equation}
For the ordinary Standard Model matter content,
\begin{equation}
 b_3=7,
 \qquad
 b_2=\frac{19}{6},
 \qquad
 b_Y=-\frac{41}{6}
\label{eq:v3v99-SM-one-loop}
\end{equation}
in the non-GUT-normalized hypercharge convention.

\begin{proposition}[Perturbative sign diagnosis]
\label{prop:v3v99-sign-diagnosis}
At one loop, the combined gauge, ghost, Higgs, and fermion quadratic packet
is asymptotically suppressive for \(SU(3)\) and \(SU(2)\), but the
hypercharge coefficient has the Landau sign.  Consequently:
\begin{enumerate}
 \item the fermion polarization term must not be tested separately from the
 gauge--ghost term in the nonabelian factors; and
 \item the known Standard Model one-loop flow does not provide a
 cutoff-independent interacting \(U(1)_Y\) ultraviolet limit at fixed
 nonzero renormalized coupling.
\end{enumerate}
These are perturbative diagnoses, not nonperturbative construction or
triviality theorems.
\end{proposition}

\begin{proof}
The standard background-field calculation gives
\begin{equation}
 b=\frac{11}{3}C_2(G)
 -\frac{2}{3}\sum_{\rm Weyl}T(R_f)
 -\frac{1}{3}\sum_{\rm complex\ scalar}T(R_s).
\label{eq:v3v99-one-loop-formula}
\end{equation}
Insert three generations and one Higgs doublet.  This gives
\eqref{eq:v3v99-SM-one-loop}.  Positive \(b\) makes the inverse coupling
grow toward the ultraviolet; negative \(b\) has the opposite sign.  The
claims follow.  Higher loops and nonperturbative effects are outside this
calculation, proving the stated boundary.
\end{proof}

\begin{firewall}[Gravity has not repaired hypercharge]
\label{fw:v3v99-gravity-hypercharge}
Perturbative gravitational contributions to a gauge beta function can
depend on gauge and regularization choices unless embedded in a complete
ultraviolet fixed-point construction.  No result in V1--V99 proves a
scheme-independent nonperturbative gravitational correction which reverses
the hypercharge Landau sign.  Such a correction cannot be inserted into
\(\Pi_\Lambda\) as an assumption and then counted as completion of the
SM--Einstein continuum.
\end{firewall}

\subsection{Maximum plaquettes versus punctured admissibility}
\label{sec:v3v99-plaquettes}

A toy independent-plaquette Gaussian already shows the relevant extreme
scale.  Let \(M_a\asymp Ca^{-4}\) centred scalar plaquette coordinates have
variance \(\sigma_a^2\).  Then
\begin{equation}
 \mathbb P\left(\max_{1\le j\le M_a}|X_j|\le\epsilon\right)
 =\left(1-2\Phi(-\epsilon/\sigma_a)\right)^{M_a}.
\label{eq:v3v99-max-exact}
\end{equation}

\begin{lemma}[Extreme-value admissibility threshold]
\label{lem:v3v99-extreme-threshold}
In the independent model, maximum admissibility tends to one only if
\begin{equation}
 M_a\,\Phi(-\epsilon/\sigma_a)\longrightarrow0.
\label{eq:v3v99-extreme-condition}
\end{equation}
If \(\sigma_a^2\asymp c/\log(1/a)\), the maximum is of order one rather
than tending to zero.  For a sufficiently small fixed admissibility
threshold \(\epsilon\), the probability in
\eqref{eq:v3v99-max-exact} tends to zero.
\end{lemma}

\begin{proof}
Use \((1-p)^M=\exp[-Mp+O(Mp^2)]\) and the Gaussian tail
\(\Phi(-x)\asymp x^{-1}e^{-x^2/2}\).  With
\(M_a\asymp a^{-4}\) and
\(\sigma_a^2\asymp c/\log(1/a)\), the typical maximum is
\(\sigma_a\sqrt{2\log M_a}\asymp\sqrt{8c}\), proving the assertion.
\end{proof}

Actual plaquettes are correlated and constrained by Bianchi identities, so
Lemma~\ref{lem:v3v99-extreme-threshold} is a diagnostic rather than a theorem
about the interacting measure.  It shows why an average action estimate and
one union bound are not a credible proof of global pointwise admissibility.

Define the bad-plaquette set
\begin{equation}
 \mathfrak B_a(\epsilon)=
 \{p:\|1-U_p\|>\epsilon\}.
\label{eq:v3v99-bad-set}
\end{equation}
The appropriate replacement is a cluster estimate: with high probability,
\(\mathfrak B_a(\epsilon)\) should be coverable by stopping balls of vanishing
physical volume and controlled total energy, while its complement admits
the V96 reconstruction.  V97--V98 then assign determinant bubble factors to
the punctures.

\begin{theorem}[Conditional punctured reconstruction criterion]
\label{thm:v3v99-punctured-criterion}
Suppose the full normalized lattice measure, with the quadratic determinant
already combined as in Theorem~\ref{thm:v3v99-covariance-criterion}, obeys:
\begin{enumerate}
 \item uniform positivity of the total transverse quadratic form;
 \item a bad-plaquette cluster bound whose covering balls have total physical
 volume tending to zero and bounded overlap;
 \item the one-use curvature-action estimate of V96 on the good complement;
 and
 \item convergence of the V97--V98 zero-mode and physical Born factors on the
 covering balls.
\end{enumerate}
Then maximum-plaquette admissibility is unnecessary: the determinant section
converges on the punctured good region and extends across the balls by the
finite bubble factors.
\end{theorem}

\begin{proof}
Apply V96 cell reconstruction only to good cells, retaining the boundary
transition functions around every bad cluster.  Bounded overlap prevents a
plaquette defect from being charged to more than a fixed number of balls.
V95 gives convergence on the good region; V97 controls the fifth determinant
on each ball; item 4 supplies the remaining finite-band and lower-Born data.
The vanishing total physical volume removes the artificial boundaries in
the distributional limit.  Positivity in item 1 makes the combined Gaussian
normalization finite at every cutoff.  These pieces prove the stated
conditional extension.
\end{proof}

\begin{firewall}[The free falsification replaces the V98 criterion]
\label{fw:v3v99-replaces-V98}
The hypothesis \eqref{eq:v3v98-log-integrability} is too strong for a quantum
gauge field and must not be used as a completion gate.  The valid replacement
is the signed combined-covariance condition
\eqref{eq:v3v99-positive-quadratic-form}, followed by normalized control of
the nonlinear remainder and bad clusters.
\end{firewall}

\begin{firewall}[The bosonic and hypercharge gates remain open]
\label{fw:v3v99-open}
V99 proves neither the interacting Yang--Mills cluster estimate nor a
nontrivial hypercharge ultraviolet completion.  It also does not construct
the Einstein metric measure or reflection-positive finite action.  These
are genuine remaining gates, not determinant-normalization errors.
\end{firewall}

\subsection{V99 conclusion and V100 audit target}
\label{sec:v3v99-conclusion}

The free theory falsifies a cutoff-uniform exponential moment for the
separate logarithmic curvature functional, even after centering.  The
correct operation is to combine the signed polarization with the bosonic
inverse covariance and normalize their determinant product once.  This
recovers the asymptotically free sign mechanism for the nonabelian Standard
Model factors and exposes the unresolved hypercharge Landau direction.
Extreme-value scaling also replaces implausible maximum admissibility by a
punctured cluster criterion.

V100 is the next compulsory NS/YM audit and the end of the present note
volume.  It should inventory the complete V1--V99 result chain, prove the
strongest honest conditional SM--Einstein continuum theorem supported by
those results, and list the irreducible open hypotheses: the bosonic
Yang--Mills and Einstein measures, hypercharge ultraviolet completion,
punctured-cluster control, and reflection positivity.  After V100 is
validated, this volume must be frozen as \texttt{\_f3} and the next active
note restarted at V1 with context and a major-results summary, in accordance
with the user's versioning instruction.

\clearpage
\section{V100: audited cutoff-weight ancestry and the conditional continuum closure theorem}
\label{sec:v3v100}

This final note of the third volume does two things.  First, it incorporates
the compulsory NS/YM audit into an exact ancestry ledger for the complete
cutoff weight.  The signed quadratic polarization and the bosonic Gaussian
form are combined and normalized once; local UV terms, physical IR terms,
zero-mode lines, and the fifth determinant remain distinct factors in their
actual multiplication order.  Second, it states and proves the strongest
conditional SM--Einstein continuum theorem supported by V1--V99.  The
theorem shows that the fermionic regulator, anomaly, determinant-line,
boundary-domain, and low-energy causal pieces fit together if the remaining
bosonic, positivity, composite-stress, and hypercharge hypotheses hold.  It
does not assert those open hypotheses.

\subsection{Compulsory V100 NS/YM audit}
\label{sec:v3v100-audit}

The exact read-only companion snapshots were
\begin{align}
 &\texttt{Renewal\_NS\_Stability\_Research\_Notes\_V31.tex},
 &&887537\ \text{bytes},
 &&\texttt{F1121B62238AD5491BB7CF03635EA9934942EDF573ED8FAAA9EE79E1D38A6696},
 \label{eq:v3v100-audit-NS}\\
 &\texttt{Renewal\_Yang\_Mills\_Mass\_Gap\_Research\_Notes\_V20.tex},
 &&381046\ \text{bytes},
 &&\texttt{3E374D83D098433AFC244ED49BFDDCE6EF59F42717041464301E0A395A2FA3C8},
 \label{eq:v3v100-audit-YM20}\\
 &\texttt{Renewal\_Yang\_Mills\_Mass\_Gap\_Research\_Notes\_V19.tex},
 &&353984\ \text{bytes},
 &&\texttt{5690E036F5C6A0706AC77F24D24E0B99639987F2E9A66DA5A383FAE3AF57D5CA}.
 \label{eq:v3v100-audit-YM19}
\end{align}
The frozen Yang--Mills archives remained
\begin{align}
 &\texttt{Renewal\_Yang\_Mills\_Mass\_Gap\_Research\_Notes\_V135\_f1.tex},
 &&3083467\ \text{bytes},
 &&\texttt{82F6BBFFEECDAD3C5C63551660B19A0BDC7F3D6DFFE4C68B4E9BBA3A783A2B71},
 \label{eq:v3v100-audit-YM-f1}\\
 &\texttt{Renewal\_Yang\_Mills\_Mass\_Gap\_Research\_Notes\_V100\_f2.tex},
 &&1706928\ \text{bytes},
 &&\texttt{6444A1044E74DFCB1A72ED9C7248BC0F234EFC4A0E4248C736CD7DB212DFCA9A}.
 \label{eq:v3v100-audit-YM-f2}
\end{align}

NS V31 is unchanged.  Its exact first-moment and balanced-helicity tests
continue to forbid replacement of a signed complete packet by separately
normed marks.  YM V20 constructs a unique ancestry tag for every acquired
response, proves that finite support telescopes preserve its parent
occurrence and sole payer, and distinguishes coordinate order from the
noncommutative multiplication slot.  It closes only the boundary-exposed
collision words; a bounded exterior multiplier cannot be deleted unless an
exact one-sided factorization permits it.

The SM transfer is type-correct and narrow.  V99's quadratic determinant is
not an independent weight after it has been absorbed into the bosonic
inverse covariance.  Its Gaussian normalizing determinant is the parent
normalization of that same occurrence.  Likewise V95's heat partition does
not create two eta forms, and V97's Born subtraction does not create a
second copy of the four physical IR terms.  The ledger below fixes these
slots before any convergence or Holder estimate.

\subsection{The complete cutoff-weight ancestry ledger}
\label{sec:v3v100-ledger}

At finite cutoff, in one determinant-line trivialization and after gauge
fixing, write the exact regulated integrand as
\begin{equation}
 \mathcal W_a=
 \mathcal G_a^{\rm comb}
 \mathcal B_a^{\rm nl}
 \mathcal U_a^{\rm loc}
 \mathcal I_a^{\rm phys}
 \mathcal D_{5,a}
 \mathcal Z_{0,a}.
\label{eq:v3v100-weight-factorization}
\end{equation}
The factors are:
\begin{enumerate}
 \item \(\mathcal G_a^{\rm comb}\), the bosonic Gaussian together with the
 complete signed quadratic polarization, normalized by
 Theorem~\ref{thm:v3v99-covariance-criterion};
 \item \(\mathcal B_a^{\rm nl}\), the remaining nonlinear bosonic,
 gauge-fixing, ghost, Higgs, and metric interaction;
 \item \(\mathcal U_a^{\rm loc}\), the exponentiated local V95 UV jet with
 its V86 normalization;
 \item \(\mathcal I_a^{\rm phys}\), the nonquadratic physical IR Born and
 relative eta occurrence, including V88 crossing transitions;
 \item \(\mathcal D_{5,a}\), the V97 fifth regularized determinant; and
 \item \(\mathcal Z_{0,a}\), the finite zero-mode exterior-line section.
\end{enumerate}
This is a ledger identity: a different renormalization convention may move a
local term between adjacent factors, but it must do so by equal and opposite
transgression and may not duplicate it.

\begin{proposition}[No-double-normalization theorem]
\label{prop:v3v100-no-double-normalization}
Let \(Z_a^{\rm quad}\) be the determinant in
\eqref{eq:v3v99-Gaussian-ratio}.  Once
\(\mathcal G_a^{\rm comb}\) is normalized using \(Z_a^{\rm quad}\), neither
the quadratic Born exponential nor \(Z_a^{\rm quad}\) may occur in any of
\(\mathcal U_a^{\rm loc}\), \(\mathcal I_a^{\rm phys}\), or
\(\mathcal D_{5,a}\).  Conversely, deleting a physical IR factor merely
because it has a uniform operator norm changes the normalized measure.
\end{proposition}

\begin{proof}
The Gaussian identity
\begin{equation}
 e^{-\frac12\langle A,C_a^{-1}A\rangle}
 e^{-\frac12\langle A,\Pi_aA\rangle}
 =e^{-\frac12\langle A,(C_a^{-1}+\Pi_a)A\rangle}
\label{eq:v3v100-one-quadratic-occurrence}
\end{equation}
contains one polarization occurrence.  Integrating it gives one determinant
ratio, \(Z_a^{\rm quad}\).  Multiplying either exponential or determinant
ratio back into the residual weight counts the same occurrence twice.

V95 similarly applies \(\chi+(1-\chi)=1\) to one heat integrand, and V97
defines \(\det_5\) only after the four Born terms have been assigned.  These
are exact algebraic partitions, so their descendants have fixed parent
slots.  Finally, a bounded nonconstant factor changes expectations after
normalization unless it cancels with an explicitly displayed reciprocal;
an operator-norm estimate is no such identity.
\end{proof}

\begin{firewall}[Holder exponents do not create new payers]
\label{fw:v3v100-Holder}
Holder's inequality may distribute an already factorized residual weight
among integrability norms.  It does not permit
\(\mathcal G_a^{\rm comb}\) or the Yang--Mills action to pay the same
determinant factor twice.  Every use of action absorption in V97 occurs
after the quadratic base measure has been fixed and applies only to the
fifth determinant remainder.
\end{firewall}

\subsection{Two closure lemmas for normalized limits}
\label{sec:v3v100-closure-lemmas}

The following elementary results isolate what remains to be supplied by the
bosonic construction.

\begin{theorem}[Vitali closure of normalized cutoff weights]
\label{thm:v3v100-Vitali}
Let \(\nu_n\) be probability measures on a Polish distribution space and
suppose \(\nu_n\Rightarrow\nu\).  Let \(W_n\) be the product of the last
five factors in \eqref{eq:v3v100-weight-factorization}, after the combined
quadratic factor has been incorporated into \(\nu_n\).  Suppose:
\begin{enumerate}
 \item \((\Phi_n,W_n)\) converges in distribution to \((\Phi,W)\) for every
 bounded continuous cylinder observable \(\Phi\);
 \item \(\{W_n\}\) is uniformly integrable;
 \item \(\{\Phi_nW_n\}\) is uniformly integrable for every observable in a
 fixed determining algebra; and
 \item \(Z:=\int W\,d\nu\ne0\).
\end{enumerate}
Then
\begin{equation}
 \frac{\int\Phi_nW_n\,d\nu_n}{\int W_n\,d\nu_n}
 \longrightarrow
 \frac{\int\Phi W\,d\nu}{\int W\,d\nu}.
\label{eq:v3v100-normalized-convergence}
\end{equation}
The assertion remains valid for complex determinant sections after a common
trivialization, provided uniform integrability refers to their moduli.
\end{theorem}

\begin{proof}
Skorokhod representation realizes each convergent pair on a common
probability space with almost-sure convergence along a subsequence.  Uniform
integrability and Vitali's theorem give convergence in \(L^1\) of \(W_n\)
and \(\Phi_nW_n\).  Their expectations therefore converge.  The denominator
converges to the nonzero number \(Z\), so division proves
\eqref{eq:v3v100-normalized-convergence}.  Since every subsequence has a
further subsequence with the same limit, the full sequence converges.
Applying the argument to real and imaginary parts proves the complex case.
\end{proof}

\begin{theorem}[Reflection positivity is closed under Schwinger convergence]
\label{thm:v3v100-RP-closure}
Let \(\mathcal S_n\) be normalized finite-cutoff Schwinger functionals on a
common positive-time cylinder algebra, and assume
\(\mathcal S_n(F)\to\mathcal S(F)\) on that algebra.  If
\begin{equation}
 \mathcal S_n((\Theta F)F)\ge0
\label{eq:v3v100-finite-RP}
\end{equation}
for every \(n\) and every positive-time cylinder polynomial \(F\), then
\begin{equation}
 \mathcal S((\Theta F)F)\ge0.
\label{eq:v3v100-limit-RP}
\end{equation}
If Euclidean covariance, symmetry, and clustering also pass to the limit,
the Osterwalder--Schrader reconstruction applies.
\end{theorem}

\begin{proof}
For fixed \(F\), the left side of
\eqref{eq:v3v100-finite-RP} is a finite linear combination of cutoff
Schwinger functions and converges to the left side of
\eqref{eq:v3v100-limit-RP}.  The closed half-line \([0,\infty)\) contains
the limit.  The last statement is the Osterwalder--Schrader reconstruction
theorem once its remaining axioms hold.
\end{proof}

\subsection{Result map of the third volume}
\label{sec:v3v100-result-map}

The proved chain in V1--V99 is organized as follows; each entry abbreviates
the theorem hypotheses in its original version.
\begin{center}
\begin{tabular}{p{0.14\textwidth}p{0.76\textwidth}}
\toprule
versions & proved acquisition\\
\midrule
V1--V4 & BRST-compatible boundary lifts, Hodge--Poisson chain extension,
relative cohomology obstruction, and minimal inflow attachment.\\
V5--V19 & one-payer boundary inflow estimates, moving-boundary Kato and
Feshbach reductions, dynamic one-port maximal dissipativity, glancing
analysis, physical reservoir determinants, and the relative heat/Schatten
gate.\\
V20--V39 & exact reservoir determinant on curved collars, full moving de
Rham block, Calderon ellipticity, Abelian and Yang--Mills BRST complexes,
complete Einstein--matter boundary symbols, reflection-positive Euclidean
origins for local baths, quartet elimination, and the coupled
TT--material--bath junction.\\
V40--V54 & common boundary-domain obstruction, explicit quotient roots and
repairs, zero-free two-channel bath, time-domain glancing estimate,
quasilinear Picard map, constraint manifold, and frozen normal-family Schur
gap.\\
V55--V64 & concrete bosonic SM boundary domain, chiral material projectors,
absolute and joint ghost complexes, conformal-factor correction, and one
finite local regulator for the gauge, constraint, and carrier complexes.\\
V65--V79 & regulator-stable Slavnov defect, full local Standard Model anomaly
cancellation, material spectral flow, domain-wall/mirror alternative,
Wilson degree, exact wall mixing, local correlator convergence, anomaly
inflow, relative heat locality, finite reflection form, overlap/Ginsparg--
Wilson algebra, local sign locality, and the exact overlap moment map.\\
V80--V86 & audited first-moment obstruction, equivariant projector
transgression, local spectral flow, same-phase covariant block reference,
stable gapped path, exact multiscale telescope, and an \(O(a)\) normalized
irrelevant remainder.\\
V87--V94 & vanishing
\(\Omega_5^{\rm Spin}(BG_{\rm SM})\), spectral-section transition
convergence, separation of physical and Wilson spectra, admissible overlap
index theorem, flat family index, curved Bott--Thom globalization,
finite-mesh compression, and exact torsion-preserving determinant-line Kato
transport.\\
V95--V99 & exact UV/IR eta split, gauge-covariant rough-link bundle
reconstruction, weak-Schatten bubble determinant, the physical logarithmic
curvature moment, free falsification of a separate exponential bound, and
the signed combined-covariance criterion.\\
\bottomrule
\end{tabular}
\end{center}

\subsection{Strongest honest conditional SM--Einstein continuum theorem}
\label{sec:v3v100-conditional-theorem}

Let \(a_n\downarrow0\) follow the V93 diagonal schedule.  The hypotheses
below are deliberately stated at the remaining interfaces rather than by
assuming the conclusion.

\begin{definition}[Outstanding completion hypotheses]
\label{def:v3v100-hypotheses}
The completion package consists of:
\begin{enumerate}[label=\textup{(H\arabic*)},leftmargin=4.5em]
 \item \textbf{Signed bosonic positivity.}  The complete gauge--ghost--Higgs--
 metric quadratic form satisfies
 \eqref{eq:v3v99-positive-quadratic-form} uniformly after renormalization,
 including a nontrivial ultraviolet completion of the hypercharge direction.
 \item \textbf{Bosonic tightness and clusters.}  The normalized nonlinear
 bosonic measures are tight on a distribution space and satisfy the
 punctured bad-cluster conditions of
 Theorem~\ref{thm:v3v99-punctured-criterion}.
 \item \textbf{Fermion residual integrability.}  The physical lower Born
 packet, zero-mode exterior sections, and bosonic bubble distribution obey
 the uniform-integrability and convergence conditions left open in
 V97--V99.  The fifth determinant uses the proved V97 bound.
 \item \textbf{Exact finite reflection positivity.}  One regulator in the
 proven stable Bott--Wilson class has a finite-cutoff cross-plane positive
 factorization compatible with the curved metric and with the local
 normalization.  Its Schwinger functions converge on a determining
 cylinder algebra.
 \item \textbf{BRST/domain stability.}  The zero-trace BRST, harmonic, and
 common boundary gaps required in V1--V64 are uniform on the measure-support
 class, and the regulator Slavnov defects of V65 vanish in the normalized
 limit.
 \item \textbf{Composite stress control.}  The renormalized stress tensor
 satisfies the microlocal, uniform-integrability, diagonal-remainder, and
 metric-equicontinuity hypotheses isolated in the frozen volumes, with no
 concentration defect.
 \item \textbf{Einstein response stability.}  The conserved limiting source
 lies in the graph spaces of V48--V54, the low-energy causal inverse obeys
 \(\kappa\|E_RP\|<1\), and the nonlinear geometric continuation estimates
 hold on the intended time interval.
\end{enumerate}
\end{definition}

\begin{theorem}[Conditional SM--Einstein continuum closure]
\label{thm:v3v100-conditional-closure}
Under hypotheses \textup{(H1)--(H7)}, the cutoff theory has a subsequential
continuum with the following properties:
\begin{enumerate}
 \item the Standard Model fermion determinant defines a global section over
 the smooth spin--\(G_{\rm SM}\) quotient, with the correct
 \(\widehat A\operatorname{ch}\) index, local anomaly cancellation, and
 trivial global gauge holonomy;
 \item normalized Euclidean Schwinger functions converge, satisfy the
 limiting Slavnov--Taylor and diffeomorphism Ward identities, and are
 reflection positive;
 \item Osterwalder--Schrader reconstruction yields a physical Hilbert space
 and Lorentzian field algebra on the reconstructed background class;
 \item the renormalized stress tensor exists as a conserved composite field
 with the response and boundary domain established in the earlier volumes;
 and
 \item the causal Einstein equation with that source has the unique
 cutoff-stable solution supplied by the V48--V54 low-energy theorem on the
 interval covered by \textup{(H7)}.
\end{enumerate}
Different regulator and triangulation choices satisfying the same
normalization conditions give the same subsequential observables whenever
their residual local transgressions converge to zero.
\end{theorem}

\begin{proof}
By \textup{(H1)}, incorporate the complete signed quadratic packet into the
base covariance exactly once, as in
Theorem~\ref{thm:v3v99-covariance-criterion}.  Hypothesis \textup{(H2)}
gives a tight subsequence and the punctured V96 reconstruction.  V95 controls
the smooth complement, V97 controls the fifth determinant on concentration
balls, and \textup{(H3)} supplies the remaining physical and zero-mode
factors.  Theorem~\ref{thm:v3v100-Vitali} therefore gives normalized
Schwinger convergence.

V90--V93 identify the finite regulator index with
\(\int\widehat A\operatorname{ch}\).  V66 cancels the complete local anomaly,
V87 proves the relevant fifth spin bordism group vanishes, and V88--V95
construct compatible crossing transitions and determinant-line
transgressions.  This proves item 1.

Hypothesis \textup{(H5)} passes the finite Slavnov and diffeomorphism Ward
identities to the limit.  Hypothesis \textup{(H4)} and
Theorem~\ref{thm:v3v100-RP-closure} give reflection positivity; the other
assumed Euclidean axioms permit Osterwalder--Schrader reconstruction.  This
proves items 2 and 3.

Hypothesis \textup{(H6)} invokes the frozen composite-field compactness
theorem and removes its explicit stress-defect alternatives.  The affine
Ward identity then gives conservation on the common limiting domain, proving
item 4.  Finally \textup{(H7)} supplies the contraction margin, common graph
domain, and quasilinear continuation required by V48--V54, proving item 5.
V81, V86, and V93 show that two allowed regulator paths differ by a sum of
local transgressions plus a summable remainder; the final sentence follows
when the normalized residual tends to zero.
\end{proof}

\begin{firewall}[The conditional theorem is not the completed continuum]
\label{fw:v3v100-not-complete}
Hypotheses \textup{(H1)--(H7)} include major unsolved constructions.  In
particular, this volume has not proved a nonperturbative four-dimensional
Yang--Mills measure and mass gap, a nontrivial hypercharge ultraviolet
completion, an Einstein metric measure, a punctured bad-cluster estimate,
reflection positivity for the curved compressed fermion regulator,
composite-stress uniform integrability, or all-frequency/global nonlinear
Einstein stability.  Theorem~\ref{thm:v3v100-conditional-closure} is a
dependency closure theorem, not a proof of those hypotheses.
\end{firewall}

\subsection{Irreducible next gates and proof order}
\label{sec:v3v100-next-gates}

The V99 free calculation and the V100 audit determine the next order.
\begin{enumerate}
 \item Treat the combined bosonic--polarization measure as the parent
 occurrence.  Prove positivity and tightness for \(SU(3)\) and \(SU(2)\)
 using the same signed normalization; do not estimate the fermion logarithm
 separately.
 \item Isolate the \(U(1)_Y\) direction.  A nontrivial ultraviolet fixed
 point, a rigorously controlled embedding, or a proven gravitational
 mechanism is required; perturbative anomaly cancellation does not address
 the Landau sign.
 \item Replace maximum plaquette admissibility by a bad-cluster stopping
 cover and prove its probability estimate.  Retain the parent plaquette
 energy of every cluster through reconstruction and bubble factorization.
 \item Construct a reflection-positive representative of the curved
 Bott--Wilson class, preferably with ultralocal temporal propagation and
 geometric compression confined to spatial slices.
 \item Only after the measure and positivity gates close should the programme
 spend the resulting probability bounds on composite stress and the
 Einstein response.  Using the same action stock independently in those
 steps would violate the ancestry ledger.
\end{enumerate}

\subsection{V100 conclusion and archive instruction}
\label{sec:v3v100-conclusion}

The third volume ends with an exact factor ancestry and a conditional
closure theorem.  The fermionic topology is now global and curved, the
Standard Model determinant line has no local or global gauge anomaly, the
finite curved regulator has the correct index, and smooth/bubble determinant
normalization has been reduced to explicit signed measure interfaces.  The
free gauge calculation prevents a false absolute-integrability route and
exposes the hypercharge, bosonic cluster, reflection-positive, stress, and
Einstein gates which remain.

After validation this V100 file is to be frozen as
\texttt{Renewal\_SM\_Einstein\_Continuum\_Research\_Notes\_V100\_f3.tex}.
The next active file restarts at V1, summarizes the major results and open
interfaces of all three frozen volumes, and begins with the signed combined
bosonic--polarization parent measure rather than repeating the determinant
work completed here.

\end{document}
